% د ۳۲۱ واحدونو د v0.7.0 نسخې بشپړ يوځای شوے LaTeX متن؛ د ټولې پروژې له سرچينو سره يې يوځای وساتئ.
% For an exact build, copy this file to releases/reader.tex in the source bundle and run three guarded XeLaTeX passes.
\documentclass[11pt,a4paper,oneside]{report}
\usepackage[margin=20mm,headheight=15pt]{geometry}
\usepackage{fontspec}
\setmainfont[Script=Arabic]{Amiri}
\newfontfamily\latinfont{Latin Modern Roman}
\usepackage{amsmath,amssymb,amsfonts,amsthm,mathrsfs,mathtools,mathdots,nicefrac,stmaryrd}
\usepackage{xparse,etoolbox,xpunctuate,mfirstuc}
\usepackage{graphicx,tikz,xcolor,framed,comment,longtable,multirow,array,booktabs,verbatim}
\usetikzlibrary{arrows,automata,positioning,calc,decorations.pathmorphing}
\usepackage{bussproofs}
\input{../upstream/sty/bussproofs-extra.sty}
\usepackage[tableaux]{prooftrees}
\usepackage[authoryear]{natbib}
\usepackage[unicode,colorlinks,linkcolor=blue,urlcolor=blue,citecolor=blue]{hyperref}
\usepackage[capitalize,nameinlink]{cleveref}
\usepackage{fancyhdr}
\usepackage{bidi}

\hypersetup{
  pdftitle={خلاص منطق — پښتو (پاکستان) — له بنسټونو تر ناتکميلۍ او اثبات‌وړتيا},
  pdfauthor={اوپن لاجک پروژه؛ خپلواکه ماشيني پښتو ژباړه},
  pdfsubject={روان کار: له ۷۲۲ سرچينيزو واحدونو څخه ۳۲۱ ژباړل شوي}
}

\newcommand*{\DeclareDocumentMacro}[2]{\def#1{#2}}
\NewDocumentCommand{\settexttoken}{m s m m o o}{}
\makeatletter
\input{../upstream/sty/open-logic-formulas.sty}
\input{../upstream/sty/open-logic-selective.sty}
\input{../upstream/open-logic-config.sty}
\makeatother
\newbox\gnBoxA
\newdimen\gnCornerHgt
\setbox\gnBoxA=\hbox{$\ulcorner$}
\global\gnCornerHgt=\ht\gnBoxA
\newdimen\gnArgHgt
\def\gn #1{%
  \setbox\gnBoxA=\hbox{$#1$}%
  \gnArgHgt=\ht\gnBoxA
  \ifnum\gnArgHgt<\gnCornerHgt \gnArgHgt=0pt
  \else \advance\gnArgHgt by -\gnCornerHgt
  \fi
  \raise\gnArgHgt\hbox{$\ulcorner$}\box\gnBoxA
  \raise\gnArgHgt\hbox{$\urcorner$}}
\renewcommand{\MP}{\LR{\latinfont MP}}
\renewcommand{\QR}{\LR{\latinfont QR}}
\renewcommand{\Hyp}{\LR{\latinfont Hyp}}

\newcommand{\olfileid}[3]{\gdef\pspart{#1}\gdef\pschapter{#2}\gdef\pssection{#3}}
\NewDocumentCommand{\olpart}{o m m}{%
  \gdef\pspart{#2}\gdef\pschapter{}\gdef\pssection{}%
  \IfNoValueTF{#1}{\part{#3}}{\part[#1]{#3}}%
  \label{#2:::part}}
\NewDocumentCommand{\olchapter}{o m m m}{%
  \gdef\pspart{#2}\gdef\pschapter{#3}\gdef\pssection{}%
  \IfNoValueTF{#1}{\chapter{#4}}{\chapter[#1]{#4}}%
  \label{#2:#3::chap}}
\NewDocumentCommand{\olsection}{o m}{%
  \IfNoValueTF{#1}{\section{#2}}{\section[#1]{#2}}%
  \label{\pspart:\pschapter:\pssection:sec}}
\newcommand{\ollabel}[1]{\label{\pspart:\pschapter:\pssection:#1}}
\makeatletter
\newcommand{\olpseganchor}[2]{%
  \protected@write\@auxout{}{%
    \string\newlabel{olpseg:#1:#2}{{}{\thepage}{}{}{}}}}
\makeatother
\NewDocumentCommand{\olref}{o o o m}{\begingroup
  \IfNoValueTF{#1}{\edef\pskey{\pspart:\pschapter:\pssection:#4}}{%
    \IfNoValueTF{#2}{\edef\pskey{\pspart:\pschapter:#1:#4}}{%
      \IfNoValueTF{#3}{\edef\pskey{\pspart:#1:#2:#4}}{\edef\pskey{#1:#2:#3:#4}}}}%
  \ifcsname r@\pskey\endcsname
    \LR{\ref{\pskey}}%
  \else
    \emph{ناحل شوې بهرنۍ حواله}%
  \fi
  \endgroup}
\let\Olref\olref
\NewDocumentCommand{\oliflabeldef}{m m m}{\ifcsname r@#1\endcsname #2\else #3\fi}
\providecommand{\OLEndChapterHook}{}
\providecommand{\OLEndPartHook}{}
\newcommand{\gitissue}[1]{%
  \href{https://github.com/OpenLogicProject/OpenLogic/issues/#1}{issue \##1}}

\newlength{\olphotowidth}
\setlength{\olphotowidth}{.46\textwidth}
\NewDocumentCommand{\olasset}{O{\olphotowidth} m}{%
  \centering\begin{LTR}\centering\resizebox{#1}{!}{\input{#2}}\end{LTR}}

\renewenvironment{prooftree}
  {\begin{center}\begin{LTR}\bottomAlignProof}
  {\DisplayProof\end{LTR}\end{center}}
\newenvironment{oltableau}
  {\center\LTR\tableau{}}
  {\endtableau\endLTR\endcenter}

\newtheorem{thm}{قضيه}[chapter]
\newtheorem{lem}[thm]{لم}
\newtheorem{prop}[thm]{قضيه}
\newtheorem{cor}[thm]{نتيجه}
\theoremstyle{definition}
\newtheorem{defn}[thm]{تعريف}
\newtheorem{ex}[thm]{بېلګه}
\newtheorem{prob}[thm]{تمرين}
\newtheorem{rem}[thm]{يادونه}
\crefname{part}{برخه}{برخې}
\crefname{chapter}{باب}{بابونه}
\crefname{section}{برخه}{برخې}
\crefname{thm}{قضيه}{قضيې}
\crefname{lem}{لم}{لمې}
\crefname{prop}{قضيه}{قضيې}
\crefname{cor}{نتيجه}{نتيجې}
\crefname{defn}{تعريف}{تعريفونه}
\crefname{ex}{بېلګه}{بېلګې}
\crefname{prob}{تمرين}{تمرينونه}
\crefname{rem}{يادونه}{يادونې}
\newcommand{\crefpairconjunction}{ او }
\newcommand{\crefmiddleconjunction}{، }
\newcommand{\creflastconjunction}{ او }
\newcommand{\crefpairgroupconjunction}{ او }
\newcommand{\crefmiddlegroupconjunction}{، }
\newcommand{\creflastgroupconjunction}{ او }
\newcommand{\crefrangeconjunction}{ تر }
\newenvironment{defish}{\begin{framed}}{\end{framed}}
\newenvironment{editorial}{\begin{framed}\textbf{ادارتي يادونه۔}\enspace}{\end{framed}}
\newenvironment{history}{\begin{framed}\textbf{تاريخي يادونه۔}\enspace}{\end{framed}}
\newenvironment{explain}{}{}
\newenvironment{intro}{}{}
\newenvironment{digress}{\par\medskip}{\par\medskip}
\newenvironment{derivation}{%
  \par\smallskip\begin{LTR}\begin{tabular}{@{}rll@{}}}
  {\end{tabular}\end{LTR}\par\smallskip}

\renewcommand{\proofname}{ثبوت}
\renewcommand{\figurename}{شکل}
\renewcommand{\tablename}{جدول}
\renewcommand{\partname}{برخه}
\renewcommand{\chaptername}{باب}
\renewcommand{\contentsname}{فهرست}
\renewcommand{\bibname}{\RL{\normalfont سرچينې}}
\makeatletter
\renewcommand{\@pnumwidth}{2.8em}
\renewcommand{\@tocrmarg}{4.6em}
\makeatother
\setlength{\parindent}{0pt}
\setlength{\parskip}{5pt}
\setlength{\emergencystretch}{8em}
\hfuzz=5pt
\linespread{1.12}
\pagestyle{fancy}
\fancyhf{}
\fancyhead[RE,LO]{\small خلاص منطق — پښتو (پاکستان)}
\fancyhead[LE,RO]{\small له بنسټونو تر ناتکميلۍ}
\fancyfoot[C]{\thepage}

\begin{document}
\tracinglostchars=3
\setRTL
\begin{titlepage}
\thispagestyle{empty}
\begin{center}
\vspace*{22mm}
{\Huge خلاص منطق}\par\bigskip
{\Large پښتو — پاکستان}\par\bigskip
{\LARGE له بنسټونو تر ناتکميلۍ او اثبات‌وړتيا}\par\bigskip
د ماشيني ژباړې د روان کار ټوليزه نسخه
\end{center}
\bigskip
په دې ټوليزه لوستيزه نسخه کښې له OLP-0001 څخه تر OLP-0321 پورې
ټول ژباړل شوي سرچينيز واحدونه شامل دي۔ دا 321 واحدونه د
خلاص منطق د پروژې لومړنۍ مادې، ساده سټ نظريه، بياني او لومړۍ درجې
منطق، د ثبوت نظامونه، موډل نظريه، محاسبويت، د ټيورينګ ماشينونو
پېژندګلو، د رابنسن حساب او د ناتکميلۍ او اثبات‌وړتيا قضيې رانغاړي۔
د ۷۲۲ واحدونو بشپړه ژباړه لا روانه ده۔ تر OLP-0311 پورې متن
OpenAI Codex، GPT-5.6 Sol، Ultra هڅې ژباړلے او کتلے دے؛
د OLP-0312 تر OLP-0321 متن OpenAI Codex، GPT-6 Sol، Ultra
هڅې ژباړلے او کتلے دے۔ د انسان د تصويب ادعا نۀ کېږي۔

\begin{LTR}\latinfont\small
Independent Pashto (Pakistan) machine translation of Open Logic Project
revision \texttt{9620cc73f9c8e0ad003c514a5d3748f29611c4c0}. This cumulative
work-in-progress reader contains 321 source units (OLP-0001
through OLP-0321). It uses the frozen upstream default profile: first-order
logic, all primitive connectives and quantifiers, all four proof systems,
and the default mathematics, computer-science, philosophy, novice, lambda,
and Turing-machine selections. Pakistani usage is primary; Afghan Pashto
sources are labelled regional comparators. Specialist terminology remains
openly reviewable where exact native technical attestation is sparse. No
native-reader review or upstream endorsement is claimed.

Translation and mathematical review through OLP-0311: OpenAI Codex,
GPT-5.6 Sol, Ultra effort. OLP-0312 through OLP-0321: OpenAI Codex,
GPT-6 Sol, Ultra effort.

Original text and adaptation: Creative Commons Attribution 4.0 International,
subject to the component notices preserved with the source snapshot.
\href{https://openlogicproject.org/}{Open Logic Project}.\quad
\href{https://github.com/KokunoYumeto/OpenLogic-translations}{Translation catalogue}.\quad
\href{https://doi.org/10.5281/zenodo.22307197}{Edition lineage}.
\end{LTR}
\end{titlepage}
\tableofcontents
\clearpage


\olpseganchor{OLP-0001-B001}{start}
\chapter*{د خلاص منطق پروژې په اړه}
\addcontentsline{toc}{chapter}{د خلاص منطق پروژې په اړه}
\olpseganchor{OLP-0001-B001}{end}


\olpseganchor{OLP-0001-B002}{start}
\textit{د خلاص منطق متن} د صوري مابعدمنطق او صوري طريقو د خلاصې سرچينې يو ګډ درسي کتاب دے، چې له منځنۍ کچې پيل کوي (يعنې د صوري منطق له يوه لومړني کورس وروسته)۔ که څه هم د رياضي د غيرمتخصصو لوستونکو دپاره دے (په تېره د فلسفې او کمپيوټر ساينس د زده کوونکو دپاره)، خو دقيق دے۔
\olpseganchor{OLP-0001-B002}{end}

\olpseganchor{OLP-0001-B003}{start}
د ځينو هغو موضوعاتو بيان چې اوس پکښې شامل دي، ښائي لا بشپړ نۀ وي، او ډېرې برخې لا هم پراخې بياکتنې ته اړتيا لري۔ مونږ اراده لرو چې په راتلونکي کښې متن پراخ کړو او نور موضوعات پکښې شامل کړو۔ دا اراده هم لرو چې متن ته نورې اسانتياوې ورزياتې کړو، لکه اصطلاحنامه، د نورې لوستنې لست، تاريخي يادښتونه، انځورونه، غوره وضاحتونه، هغه برخې چې د پايلو تړاو له فلسفې، کمپيوټر ساينس او رياضي سره بيانوي، او نور تمرينونه او بېلګې۔ که کومه تېروتنه ومومئ، يا وړانديز لرئ، \href{https://github.com/OpenLogicProject/OpenLogic/wiki/Contributing}{مهرباني وکړئ د پروژې ډلې ته ئې خبر ورکړئ}۔
\olpseganchor{OLP-0001-B003}{end}

\olpseganchor{OLP-0001-B004}{start}
پروژه د خلاصې سرچينې په روحيه کار کوي۔ يوازې متن په ازاده توګه نۀ وړاندې کوو، بلکې د لاټېک سرچينه هم د کريېټيو کامنز د نوم اخيستنې د جواز لاندې ورکوو۔ دا جواز هر چا ته حق ورکوي چې زمونږ کار ښکته کړي، وکاروي، بدل ئې کړي، بيا ئې ترتيب کړي، بڼې ئې واړوي او بيا ئې ووېشي، په دې شرط چې مناسب نسبت ورکړي۔ د نورو معلوماتو دپاره، مهرباني وکړئ په \href{http://openlogicproject.org/}{openlogicproject.org} د خلاص منطق پروژې وېبپاڼه وګورئ۔
\olpseganchor{OLP-0001-B004}{end}

\olpseganchor{OLP-0002-B004}{start}
\clearpage
\olpseganchor{OLP-0002-B004}{end}

\olpseganchor{OLP-0002-B005}{start}
\begin{editorial}
دا دوتنه ټول هغه مطالب رااخلي چې د خلاص منطق په پروژه کښې شامل دي۔ که دا ډول ادارتي يادښتونه ښکاري، نو دا څرګندوي چې دوتنه د دې مطالبو د وړاندې کولو طريقې ته له پام کولو پرته جوړه شوې ده۔ که دا لولئ، نو غالباً له دې پي ډي اېف څخه درس ورکول يا مطالعه کول \emph{مناسب نۀ دي}۔
\olpseganchor{OLP-0002-B005}{end}

\olpseganchor{OLP-0002-B006}{start}
د خلاص منطق پروژه ډېرې داسې طريقې وړاندې کوي چې د هغو په وسيله د درس ورکولو يا ځان زده کړې دپاره مناسب متن جوړېدل ممکن وي۔ د بېلګې په توګه، په اصلي تنظيم کښې متن ټول منطقي عملګر بنسټيز ګڼي، او په ثبوتونو کښې د ټولو عملګرو ټول حالتونه څېړي۔ خو دا ډېر غوره دي چې ځينې حالتونه د تمرينونو په توګه پرېښودل شي۔ د خلاص منطق پروژه پخپله هم روان کار دے۔ د ګډ کار او ښه والي د هڅولو دپاره، مطالب آن هغه وخت هم شاملېږي چې يوازې د مسودې په بڼه وي، تمرينونه نۀ لري، او داسې نور۔ هغه پي ډي اېف چې د يو کورس دپاره جوړېږي، دا برخې به ترې وباسي۔
\olpseganchor{OLP-0002-B006}{end}

\olpseganchor{OLP-0002-B007}{start}
د درس او مطالعې دپاره د لا مناسبو پي ډي اېفونو د موندلو دپاره، \href{http://builds.openlogicproject.org/}{د خلاص منطق پروژې په وېبپاڼه شته نمونوي کورسونه} وګورئ۔ د خپل متن د جوړولو دپاره، له \href{https://github.com/OpenLogicProject/OpenLogic/tree/master/courses/sample}{نمونوي چلوونکې دوتنې} څخه پيل کولے شئ، يا د لا پرمختللو او مفصلو بېلګو دپاره له دې پروژې څخه د جوړو شوو درسي کتابونو سرچينې کتلے شئ۔
\end{editorial}
\olpseganchor{OLP-0002-B007}{end}

\olpseganchor{OLP-0003-B004}{start}
\olpart{sfr}{د سټونو ساده نظريه}
\olpseganchor{OLP-0003-B004}{end}

\olpseganchor{OLP-0003-B005}{start}
\begin{editorial}
  د دې برخې مطالب د سټونو د بنسټيزې ساده نظريې يوه پېژندګلو ده۔ د تېم بټن د سټونو د خلاصې نظريې په شاملېدو سره، دا برخه د عددونو د نظامونو جوړول او د نامتناهي والي بحث هم رانغاړي، چې د خلاص منطق پروژې د منطقي برخو دپاره اړين نۀ دي۔
\end{editorial}
\olpseganchor{OLP-0003-B005}{end}

\olpseganchor{OLP-0004-B004}{start}
\olchapter{sfr}{set}{سټونه}
\olpseganchor{OLP-0004-B004}{end}

\olpseganchor{OLP-0005-B004}{start}
\olfileid{sfr}{set}{bas}
\olsection{د غړو له مخې برابري}
\olpseganchor{OLP-0005-B004}{end}

\olpseganchor{OLP-0005-B005}{start}
\emph{سټ} د څيزونو يوه ټولګه ده چې د يوه څيز په توګه ورته کتل کېږي۔ هغه څيزونه چې سټ ترې جوړ وي، د سټ \emph{عناصر} يا \emph{غړي} بللے شي۔ که \LR{$x$} د يو سټ~\LR{$A$} غړے وي، \LR{$x \in A$} ليکو؛ که نۀ وي، \LR{$x \notin A$} ليکو۔ هغه سټ چې هېڅ غړي نۀ لري، \emph{تش} سټ بللے شي او په ``\LR{$\emptyset$}'' ښودلے کېږي۔
\olpseganchor{OLP-0005-B005}{end}

\olpseganchor{OLP-0005-B006}{start}
\begin{explain}
دا مهمه نۀ ده چې سټ څنګه \emph{ټاکو}، د هغۀ غړي څنګه \emph{ترتيبوو}، يا په رښتيا د هغۀ غړي \emph{څو ځله} شمارو۔ يوازې دا مهم دي چې د هغۀ غړي کوم دي۔ دا خبره په لاندې اصل کښې راټولوو۔
\end{explain}
\olpseganchor{OLP-0005-B006}{end}

\olpseganchor{OLP-0005-B007}{start}
\begin{defn}[د غړو له مخې برابري]
  که \LR{$A$} او \LR{$B$} سټونه وي، نو \LR{$A = B$} هله او يوازې هله وي چې د~\LR{$A$} هر غړے د~\LR{$B$} هم غړے وي، او برعکس۔
\end{defn}
\olpseganchor{OLP-0005-B007}{end}

\olpseganchor{OLP-0005-B008}{start}
د غړو له مخې برابري د ځينو نښو د کارولو جواز راکوي۔ په عمومي توګه، چې کله ځينې څيزونه \LR{$a_{1}$}، \dots، \LR{$a_{n}$} ولرو، نو \LR{$\{a_{1}, \dots, a_{n}\}$} \emph{هم هغه} سټ دے چې غړي ئې \LR{$a_1, \ldots, a_n$} دي۔ په ``\emph{هم هغه}'' ټينګار کوو، ځکه چې د غړو له مخې برابري راته وائي چې داسې سټ يوازې \emph{يو} کېدے شي۔ په حقيقت کښې، د غړو له مخې برابري د لاندې خبرې جواز هم راکوي:
  \[
    \{a, a, b\} = \{a, b\} = \{b,a\}.
  \]
دا هم هغه خبره څرګندوي چې د سټونو په پام کښې نيولو پر وخت، د هغو د غړو ترتيب يا دا چې څو ځله ښودل شوي دي، مهم نۀ دي۔
\olpseganchor{OLP-0005-B008}{end}

\olpseganchor{OLP-0005-B009}{start}

\begin{ex}
چې کله د څيزونو يوه ډله ولرئ، هغه په يو سټ کښې راټولولے شئ۔ د بېلګې په توګه، د رېچرډ د وروڼو او خوېندو سټ يو داسې سټ دے چې يو کس لري، او هغه د \LR{$S=\{\textrm{Ruth}\}$} په بڼه ليکلے شو۔ له \LR{$4$} څخه د کوچنيو مثبتو صحيح عددونو سټ \LR{$\{1, 2, 3\}$} دے، خو د \LR{$\{3, 2, 1\}$} يا آن د \LR{$\{1, 2, 1, 2, 3\}$} په بڼه هم ليکل کېدے شي۔ د غړو له مخې د برابرۍ له مخې، دا ټول هم يو سټ دے۔ ځکه چې د \LR{$\{1, 2, 3\}$} هر غړے د \LR{$\{3, 2, 1\}$} (او د \LR{$\{1, 2, 1, 2, 3\}$}) هم غړے دے، او برعکس۔
\end{ex}

\olpseganchor{OLP-0005-B009}{end}

\olpseganchor{OLP-0005-B010}{start}
ډېر ځله به يو سټ د يو داسې خاصيت په وسيله ټاکو چې د هغۀ غړي ئې په ګډه لري۔ د دې دپاره به دا لنډه نښه کاروو: \LR{$\Setabs{x}{\phi(x)}$}، چې پکښې \LR{$\phi(x)$} د هغه خاصيت استازيتوب کوي چې~\LR{$x$} ئې بايد ولري، څو د سټ په غړو کښې وشمېرل شي۔
\olpseganchor{OLP-0005-B010}{end}

\olpseganchor{OLP-0005-B011}{start}

\begin{ex}
زمونږ په بېلګه کښې، \LR{$S$} په دې بڼه هم ټاکلے شو:
\[
S = \Setabs{x}{x \text{\RL{ د رېچرډ ورور يا خور دے}}}.
\]
\end{ex}

\olpseganchor{OLP-0005-B011}{end}

\olpseganchor{OLP-0005-B012}{start}

\begin{ex}
يو عدد هله او يوازې هله \emph{کامل} بللے شي چې د خپلو حقيقي قاسمونو له مجموعې سره برابر وي (يعنې هغه عددونه چې دا عدد پرې پوره وېشل کېږي، خو پخپله له دې عدد سره برابر نۀ وي)۔ د بېلګې په توګه، \LR{$6$} کامل دے ځکه چې حقيقي قاسمونه ئې \LR{$1$}، \LR{$2$} او~\LR{$3$} دي، او \LR{$6 = 1 + 2 + 3$}۔ په حقيقت کښې، \LR{$6$} له \LR{$10$} څخه يوازينے کوچنے مثبت صحيح عدد دے چې کامل دے۔ نو د غړو له مخې د برابرۍ په کارولو سره وئيلے شو:
\[
	\{6\} = \Setabs{x}{x\text{\RL{ کامل دے او }}0 \leq x \leq 10}
\]
ښي لاس ته نښه داسې لولو: ``د هغو \LR{$x$} ګانو سټ چې \LR{$x$} کامل وي او \LR{$0 \leq x \leq 10$} وي''۔ دلته برابري دا خبره تائيدوي چې د سټونو په پام کښې نيولو پر وخت، د هغو د ټاکلو طريقه مهمه نۀ ده۔ او په لا عمومي توګه، د غړو له مخې برابري ډاډ راکوي چې تل د هغو \LR{$x$} ګانو يوازې يو سټ وي چې \LR{$\phi(x)$} وي۔ نو د غړو له مخې برابري دا جواز راکوي چې \LR{$\Setabs{x}{\phi(x)}$} ته د هغو \LR{$x$} ګانو \emph{هم هغه} سټ ووايو چې~\LR{$\phi(x)$} وي۔
\end{ex}

\olpseganchor{OLP-0005-B012}{end}

\olpseganchor{OLP-0005-B013}{start}
د غړو له مخې برابري راته د دې ښودلو يوه طريقه راکوي چې سټونه برابر دي: د \LR{$A = B$} د ښودلو دپاره، وښايئ چې هر کله \LR{$x \in A$} وي نو \LR{$x \in B$} هم وي، او هر کله \LR{$y \in B$} وي نو \LR{$y \in A$} هم وي۔
\olpseganchor{OLP-0005-B013}{end}

\olpseganchor{OLP-0005-B014}{start}
\begin{prob}
ثابت کړئ چې له يو څخه زيات تش سټونه نۀ شي کېدے، يعنې وښايئ چې که \LR{$A$} او \LR{$B$} داسې سټونه وي چې هېڅ غړي نۀ لري، نو \LR{$A = B$}۔
\end{prob}
\olpseganchor{OLP-0005-B014}{end}

\olpseganchor{OLP-0006-B004}{start}
\olfileid{sfr}{set}{sub}
\olsection{فرعي سټونه او د فرعي سټونو سټونه}
\olpseganchor{OLP-0006-B004}{end}

\olpseganchor{OLP-0006-B005}{start}
\begin{explain}
مونږ به ډېر ځله سټونه سره پرتله کول غواړو۔ د پرتله کولو يو څرګند ډول داسې دے: \emph{هر هغه څيز چې په يو سټ کښې وي، په بل کښې هم وي}۔ دا حالت دومره مهم دے چې د دې دپاره نوې نښې معرفي کړو۔
\end{explain}
\olpseganchor{OLP-0006-B005}{end}

\olpseganchor{OLP-0006-B006}{start}
\begin{defn}[فرعي سټ]
که د يو سټ \LR{$A$} هر غړے د~\LR{$B$} هم غړے وي، نو وايو چې \LR{$A$} د~\LR{$B$} يو \emph{فرعي سټ} دے، او \LR{$A \subseteq B$} ليکو۔ که \LR{$A$} د~\LR{$B$} فرعي سټ نۀ وي، \LR{$A \not\subseteq B$} ليکو۔ که \LR{$A \subseteq B$} وي خو \LR{$A \neq B$} وي، \LR{$A \subsetneq B$} ليکو او وايو چې \LR{$A$} د \LR{$B$} يو \emph{خاص فرعي سټ} دے۔
\end{defn}
\olpseganchor{OLP-0006-B006}{end}

\olpseganchor{OLP-0006-B007}{start}
\begin{ex}
هر سټ د خپل ځان فرعي سټ دے، او \LR{$\emptyset$} د هر سټ فرعي سټ دے۔ د جفتو طبيعي عددونو سټ د طبيعي عددونو د سټ فرعي سټ دے۔ همدارنګه، \LR{$\{ a, b \} \subseteq \{ a, b, c \}$}۔ خو \LR{$\{ a, b, e \}$} د \LR{$\{ a, b, c \}$} فرعي سټ نۀ دے۔
\end{ex}
\olpseganchor{OLP-0006-B007}{end}

\olpseganchor{OLP-0006-B008}{start}
\begin{ex}
عدد \LR{$2$} د صحيح عددونو د سټ يو غړے دے، حال دا چې د جفتو عددونو سټ د صحيح عددونو د سټ فرعي سټ دے۔ خو يو سټ کېدے شي د بل سټ \emph{هم} غړے وي او هم فرعي سټ، د بېلګې په توګه، \LR{$\{0\} \in \{0, \{0\}\}$} او هم \LR{$\{0\} \subseteq \{0, \{0\}\}$}۔
\end{ex}
\olpseganchor{OLP-0006-B008}{end}

\olpseganchor{OLP-0006-B009}{start}
د غړو له مخې برابري د سټونو د برابرۍ معيار راکوي: \LR{$A = B$} هله او يوازې هله وي چې د~\LR{$A$} هر غړے د~\LR{$B$} هم غړے وي، او برعکس۔ د ``فرعي سټ'' تعريف \LR{$A \subseteq B$} بالکل د دې معيار د لومړۍ نيمې په توګه تعريفوي: د~\LR{$A$} هر غړے د~\LR{$B$} هم غړے وي۔ البته، که \LR{$A$} او \LR{$B$} سره بدل کړو، تعريف بيا هم عملي کېږي: يعنې \LR{$B \subseteq A$} هله او يوازې هله وي چې د~\LR{$B$} هر غړے د~\LR{$A$} هم غړے وي۔ او دا بيا بالکل د غړو له مخې د برابرۍ ``برعکس'' برخه ده۔ په نورو ټکو، د غړو له مخې له برابرۍ نه دا نتيجه راوځي چې سټونه هله او يوازې هله برابر وي چې د يو بل فرعي سټونه وي۔
\olpseganchor{OLP-0006-B009}{end}

\olpseganchor{OLP-0006-B010}{start}
\begin{prop}
\LR{$A = B$} هله او يوازې هله وي چې \LR{$A \subseteq B$} او \LR{$B \subseteq A$} دواړه وي۔
\end{prop}
\olpseganchor{OLP-0006-B010}{end}

\olpseganchor{OLP-0006-B011}{start}
اوس د ځينو نورو ګټورو نښو د معرفي کولو دپاره هم ښه موقع ده۔ د دې په تعريف کښې چې \LR{$A$} کله د~\LR{$B$} فرعي سټ وي، مونږ ووئيل چې ``د~\LR{$A$} هر غړے \dots دے''، او د ``\LR{$\dots$}'' ځاے مو په ``د \LR{$B$} غړے'' ډک کړ۔ خو دا د عبارت دومره عامه \emph{بڼه} ده چې د دې دپاره د صوري نښو معرفي کول به ګټور وي۔
\olpseganchor{OLP-0006-B011}{end}

\olpseganchor{OLP-0006-B012}{start}
\begin{defn}\ollabel{forallxina}
\LR{$(\forall x \in A)\phi$} د \LR{$\forall x(x \in A \lif \phi)$} لنډه بڼه ده۔ همدارنګه، \LR{$(\exists x \in A)\phi$} د \LR{$\exists x(x \in A \land \phi)$} لنډه بڼه ده۔
\end{defn}
\olpseganchor{OLP-0006-B012}{end}

\olpseganchor{OLP-0006-B013}{start}
د دې نښو په کارولو سره وئيلے شو چې \LR{$A \subseteq B$} هله او يوازې هله وي چې \LR{$(\forall x \in A)x \in B$} وي۔
\olpseganchor{OLP-0006-B013}{end}

\olpseganchor{OLP-0006-B014}{start}
اوس د يو ځانګړي ډول سټ څېړنې ته ځو: د يو ټاکلي سټ د ټولو فرعي سټونو سټ۔
\olpseganchor{OLP-0006-B014}{end}

\olpseganchor{OLP-0006-B015}{start}
\begin{defn}[د فرعي سټونو سټ]
هغه سټ چې د يو سټ~\LR{$A$} ټول فرعي سټونه پکښې وي، د~\LR{$A$} \emph{د فرعي سټونو سټ} بللے شي، او \LR{$\Pow{A}$} ليکل کېږي۔
  \[
    \Pow{A} = \Setabs{B}{B \subseteq A}
  \]
\end{defn}
\olpseganchor{OLP-0006-B015}{end}

\olpseganchor{OLP-0006-B016}{start}
\begin{ex}
د \LR{$\{ a, b, c \}$} ټول ممکن فرعي سټونه کوم دي؟ دا دي: \LR{$\emptyset$}، \LR{$\{a \}$}، \LR{$\{b\}$}، \LR{$\{c\}$}، \LR{$\{a, b\}$}، \LR{$\{a, c\}$}، \LR{$\{b, c\}$}، \LR{$\{a, b, c\}$}۔ د دې ټولو فرعي سټونو سټ \LR{$\Pow{\{a,b,c\}}$} دے:
\[
\Pow{\{ a, b, c \}} = \{\emptyset, \{a \}, \{b\}, \{c\}, \{a, b\},
\{b, c\}, \{a, c\}, \{a, b, c\}\}
\]
\end{ex}
\olpseganchor{OLP-0006-B016}{end}

\olpseganchor{OLP-0006-B017}{start}
\begin{prob}
د \LR{$\{a, b, c, d\}$} ټول فرعي سټونه وليکئ۔
\end{prob}
\olpseganchor{OLP-0006-B017}{end}

\olpseganchor{OLP-0006-B018}{start}
\begin{prob}
وښايئ چې که \LR{$A$} د \LR{$n$} په شمېر غړي ولري، نو \LR{$\Pow{A}$} د \LR{$2^n$} په شمېر غړي لري۔
\end{prob}
\olpseganchor{OLP-0006-B018}{end}

\olpseganchor{OLP-0007-B004}{start}
\olfileid{sfr}{set}{imp}
\olsection{ځينې مهم سټونه}
\olpseganchor{OLP-0007-B004}{end}

\olpseganchor{OLP-0007-B005}{start}
\begin{ex}
مونږ به زياتره له هغو سټونو سره کار لرو چې غړي ئې رياضيکي څيزونه وي۔ څلور داسې سټونه دومره مهم دي چې ځانګړي نومونه لري:
\begin{multline*}
    \Nat = \{0, 1, 2, 3, \ldots\} \\
    \shoveright{\text{\RL{د طبيعي عددونو سټ}}}\\
    \shoveleft{\Int = \{\ldots, -2, -1, 0, 1, 2, \ldots\}} \\
    \shoveright{\text{\RL{د صحيح عددونو سټ}}}\\
    \shoveleft{\Rat = \Setabs{\nicefrac{m}{n}}{m, n \in \Int\text{\RL{ او }}n \neq 0}}\\
    \shoveright{\text{\RL{د ناطقو عددونو سټ}}}\\
    \shoveleft{\Real = (-\infty, \infty)}\\
    \text{\RL{د حقيقي عددونو سټ (پيوستار)}}
\end{multline*}
دا ټول \emph{نامتناهي} سټونه دي، يعنې د هر يوه غړي په نامتناهي شمېر دي۔
\olpseganchor{OLP-0007-B005}{end}

\olpseganchor{OLP-0007-B006}{start}
چې کله د دې سټونو له يوه نه بل ته ځو، خپلې زېرمې ته \emph{نور} عددونه ورزياتوو۔ په حقيقت کښې، دا بايد څرګنده وي چې \LR{$\Nat \subseteq \Int \subseteq \Rat \subseteq \Real$}: ځکه چې هر طبيعي عدد صحيح عدد دے؛ هر صحيح عدد ناطق عدد دے؛ او هر ناطق عدد حقيقي عدد دے۔ همدارنګه، دا هم بايد څرګنده وي چې \LR{$\Nat \subsetneq \Int \subsetneq \Rat$}، ځکه چې \LR{$-1$} صحيح عدد دے خو طبيعي عدد نۀ دے، او \LR{$\nicefrac{1}{2}$} ناطق عدد دے خو صحيح عدد نۀ دے۔ دا لږه ناڅرګنده خبره ده چې \LR{$\Rat \subsetneq \Real$}، يعنې داسې حقيقي عددونه شته چې ناطق نۀ دي۔
\olpseganchor{OLP-0007-B006}{end}

\olpseganchor{OLP-0007-B007}{start}
کله ناکله به د مثبتو صحيح عددونو سټ \LR{$\PosInt = \{1, 2, 3, \dots\}$} او هغه سټ هم کاروو چې يوازې لومړي دوه طبيعي عددونه لري، يعنې \LR{$\Bin = \{0, 1\}$}۔
\end{ex}
\olpseganchor{OLP-0007-B007}{end}

\olpseganchor{OLP-0007-B008}{start}

\begin{ex}[توريز تسلسلونه]
يوه بله په زړه پورې بېلګه د يوې الفبا \LR{$A$} د \emph{متناهي توريزو تسلسلونو} سټ \LR{$A^{*}$} دے: د~\LR{$A$} د عناصرو هره متناهي لړۍ د \LR{$A$} يو توريز تسلسل دے۔ د هرې الفبا~\LR{$A$} دپاره، \emph{تش توريز تسلسل \LR{$\Lambda$}} هم د~\LR{$A$} په توريزو تسلسلونو کښې شاملوو۔ د بېلګې په توګه،
\begin{multline*}
\Bin^*
=\{\Lambda,0,1,00,01,10,11,\\
000,001,010,011,100,101,110,111,0000,\ldots\}.
\end{multline*}
که \LR{$x=x_{1}\ldots x_{n}\in A^{*}$} يو داسې توريز تسلسل وي چې د \LR{$A$} له \LR{$n$} ``تورو'' جوړ وي، نو وايو چې د دې تسلسل \emph{اوږدوالے}~\LR{$n$} دے او \LR{$\len{x}=n$} ليکو۔
\end{ex}

\olpseganchor{OLP-0007-B008}{end}

\olpseganchor{OLP-0007-B009}{start}
\begin{ex}[نامتناهي لړۍ]
د هر سټ \LR{$A$} دپاره، د~\LR{$A$} د غړو د نامتناهي لړيو سټ~\LR{$A^\omega$} هم په پام کښې نيولے شو۔ يوه نامتناهي لړۍ \LR{$a_1a_2a_3a_4\dots$} د څيزونو له داسې لست نه جوړه وي چې په يو لوري نامتناهي غځېږي، او هر څيز ئې د~\LR{$A$} يو غړے وي۔
\end{ex}
\olpseganchor{OLP-0007-B009}{end}

\olpseganchor{OLP-0008-B004}{start}
\olfileid{sfr}{set}{uni}
\olsection{اتحادونه او تقاطعې}
\olpseganchor{OLP-0008-B004}{end}

\olpseganchor{OLP-0008-B005}{start}
\begin{explain}
په \olref[sfr][set][bas]{sec} کښې مو د تجريد په وسيله د سټونو تعريفونه معرفي کړل، يعنې د \LR{$\Setabs{x}{\phi(x)}$} په بڼه تعريفونه۔ دلته يو خاصيت~\LR{$\phi$} په کار راوړو، او دا خاصيت د هغو سټونو يادونه کولے شي چې مخکې مو تعريف کړي دي۔ نو د بېلګې په توګه، که \LR{$A$} او~\LR{$B$} سټونه وي، د \LR{$\Setabs{x}{x \in A \lor x \in B}$} سټ له هغو ټولو څيزونو جوړ دے چې د \LR{$A$} يا~\LR{$B$} غړي وي، يعنې دا هغه سټ دے چې د \LR{$A$} او~\LR{$B$} غړي سره يوځاے کوي۔ دا لکه په \olref{fig:union} کښې انځورولے شو، چې رنګ شوې ساحه د دواړو سټونو \LR{$A$} او~\LR{$B$} غړي په ګډه ښيي۔
\olpseganchor{OLP-0008-B005}{end}

\olpseganchor{OLP-0008-B006}{start}
\begin{figure}
  \olasset{../upstream/assets/diagrams/union.tikz}
  \caption{د دوو سټونو اتحاد \LR{$A \cup B$} د~\LR{$A$} د غړو او د~\LR{$B$} د غړو ګډ سټ دے۔}
  \ollabel{fig:union}
\end{figure}
\olpseganchor{OLP-0008-B006}{end}

\olpseganchor{OLP-0008-B007}{start}
په سټونو دا عمل، يعنې يوځاے کول ئې، ډېر ګټور او عام دے؛ نو يو صوري نوم او نښه ورکوو۔
\end{explain}
\olpseganchor{OLP-0008-B007}{end}

\olpseganchor{OLP-0008-B008}{start}
\begin{defn}[اتحاد]
د دوو سټونو \LR{$A$} او \LR{$B$} \emph{اتحاد}، چې \LR{$A \cup B$} ليکل کېږي، د ټولو هغو څيزونو سټ دے چې د \LR{$A$}، \LR{$B$}، يا د دواړو غړي وي۔
\[
A \cup B = \Setabs{x}{x \in A \lor x \in B}
\]
\end{defn}
\olpseganchor{OLP-0008-B008}{end}

\olpseganchor{OLP-0008-B009}{start}
\begin{ex}
څرنګه چې د غړو تکرار مهم نۀ دے، د هغو دوو سټونو اتحاد چې ګډ غړے ولري، هغه غړے يوازې يو ځل لري؛ د بېلګې په توګه، \LR{$\{ a, b, c\} \cup \{ a, 0, 1\} = \{a, b, c, 0, 1\}$}۔
\olpseganchor{OLP-0008-B009}{end}

\olpseganchor{OLP-0008-B010}{start}
د يو سټ او د هغۀ د يو فرعي سټ اتحاد هم هغه لوے سټ دے: \LR{$\{a, b, c \} \cup \{a \} = \{a, b, c\}$}۔
\olpseganchor{OLP-0008-B010}{end}

\olpseganchor{OLP-0008-B011}{start}
له تش سټ سره د يو سټ اتحاد له هم هغه سټ سره برابر دے: \LR{$\{a, b, c \} \cup \emptyset = \{a, b, c \}$}۔
\end{ex}
\olpseganchor{OLP-0008-B011}{end}

\olpseganchor{OLP-0008-B012}{start}
\begin{prob}
ثابت کړئ چې که \LR{$A \subseteq B$} وي، نو \LR{$A \cup B = B$}۔
\end{prob}
\olpseganchor{OLP-0008-B012}{end}

\olpseganchor{OLP-0008-B013}{start}
\begin{explain}
د اتحاد يو ``دوه ګونے'' عمل هم په پام کښې نيولے شو۔ دا هغه عمل دے چې د هغو ټولو غړو سټ جوړوي چې د~\LR{$A$} غړي وي او د~\LR{$B$} هم غړي وي۔ دې عمل ته \emph{تقاطع} وئيل کېږي، او لکه په \olref{fig:intersection} کښې ښودل کېدے شي۔
\begin{figure}
  \olasset{../upstream/assets/diagrams/intersection.tikz}
  \caption{د دوو سټونو تقاطع \LR{$A \cap B$} د هغو د ګډو غړو سټ دے۔}
  \ollabel{fig:intersection}
\end{figure}
\end{explain}
\olpseganchor{OLP-0008-B013}{end}

\olpseganchor{OLP-0008-B014}{start}
\begin{defn}[تقاطع]
د دوو سټونو \LR{$A$} او \LR{$B$} \emph{تقاطع}، چې \LR{$A \cap B$} ليکل کېږي، د ټولو هغو څيزونو سټ دے چې د \LR{$A$} او~\LR{$B$} دواړو غړي وي۔
\[
A \cap B = \Setabs{x}{x \in A \land x \in B}
\]
دوه سټونه هله \emph{بېل} بللے شي چې تقاطع ئې تشه وي۔ مطلب دا چې هېڅ ګډ غړي نۀ لري۔
\end{defn}
\olpseganchor{OLP-0008-B014}{end}

\olpseganchor{OLP-0008-B015}{start}
\begin{ex}
که دوه سټونه هېڅ ګډ غړي ونۀ لري، تقاطع ئې تشه ده: \LR{$\{ a, b, c\} \cap \{ 0, 1\} = \emptyset$}۔
\olpseganchor{OLP-0008-B015}{end}

\olpseganchor{OLP-0008-B016}{start}
که دوه سټونه ګډ غړي ولري، تقاطع ئې د هغو ټولو سټ دے: \LR{$\{a, b, c \} \cap \{a, b, d \} = \{a, b\}$}۔
\olpseganchor{OLP-0008-B016}{end}

\olpseganchor{OLP-0008-B017}{start}
د يو سټ او د هغۀ د يو فرعي سټ تقاطع هم هغه کوچنے سټ دے: \LR{$\{a, b, c\} \cap \{a, b\} = \{a, b\}$}۔
\olpseganchor{OLP-0008-B017}{end}

\olpseganchor{OLP-0008-B018}{start}
له تش سټ سره د هر سټ تقاطع تشه ده: \LR{$\{a, b, c \} \cap \emptyset = \emptyset$}۔
\end{ex}
\olpseganchor{OLP-0008-B018}{end}

\olpseganchor{OLP-0008-B019}{start}
\begin{prob}
په دقيقه توګه ثابت کړئ چې که \LR{$A \subseteq B$} وي، نو \LR{$A \cap B = A$}۔
\end{prob}
\olpseganchor{OLP-0008-B019}{end}

\olpseganchor{OLP-0008-B020}{start}
\begin{explain}
له دوو څخه د زياتو سټونو اتحاد يا تقاطع هم جوړولے شو۔ په عمومي توګه د دې د بيانولو يوه ښکلې طريقه دا ده: فرض کړئ چې ټول هغه سټونه چې اتحاد (يا تقاطع) ئې جوړول غواړئ، په يو سټ کښې راټول کړئ۔ بيا د ټولو اصلي سټونو اتحاد د هغو ټولو څيزونو د سټ په توګه تعريفولے شو چې لږ تر لږه د دې سټ د يو غړي غړي وي، او تقاطع د هغو ټولو څيزونو د سټ په توګه چې د دې سټ د هر غړي غړي وي۔
\end{explain}
\olpseganchor{OLP-0008-B020}{end}

\olpseganchor{OLP-0008-B021}{start}
\begin{defn}
که \LR{$A$} د سټونو يو سټ وي، نو \LR{$\bigcup A$} د~\LR{$A$} د غړو د غړو سټ دے:
\begin{align*}
\bigcup A & = \Setabs{x}{x \text{\RL{ د دې سټ د يو غړي غړے دے: }} A},
\text{\RL{ يعنې،}}\\
& = \Setabs{x}{\text{\RL{داسې يو }} B \in A
  \text{\RL{ شته چې }} x \in B}
\end{align*}
\end{defn}
\olpseganchor{OLP-0008-B021}{end}

\olpseganchor{OLP-0008-B022}{start}
\begin{defn}
که \LR{$A$} د سټونو يو سټ وي، نو \LR{$\bigcap A$} د هغو څيزونو سټ دے چې د~\LR{$A$} ټول عناصر ئې په ګډه لري:
\begin{align*}
\bigcap A & = \Setabs{x}{x \text{\RL{ د دې سټ د هر غړي غړے دے: }} A},
\text{\RL{ يعنې،}}\\
 & = \Setabs{x}{\text{\RL{د هر }} B \in A, x \in B}
\end{align*}
\end{defn}
\olpseganchor{OLP-0008-B022}{end}

\olpseganchor{OLP-0008-B023}{start}
\begin{ex}
فرض کړئ چې \LR{$A = \{ \{ a, b \}, \{ a, d, e \}, \{ a, d \} \}$}۔ بيا \LR{$\bigcup A = \{ a, b, d, e \}$} او \LR{$\bigcap A = \{ a \}$}۔
\end{ex}
\begin{prob}
	وښايئ چې که \LR{$A$} يو سټ وي او \LR{$A \in B$} وي، نو \LR{$A \subseteq \bigcup B$}۔
\end{prob}
\olpseganchor{OLP-0008-B023}{end}

\olpseganchor{OLP-0008-B024}{start}
د سټونو د يوې لړۍ \LR{$A_1$}، \LR{$A_2$}، \dots دپاره هم دا کار کولے شو:
\begin{align*}
\bigcup_i A_i & = \Setabs{x}{x \text{\RL{ له دې سټونو څخه د يوه غړے دے: }} A_i}\\
\bigcap_i A_i & = \Setabs{x}{x \text{\RL{ د دې هر سټ غړے دے: }} A_i}.
\end{align*}
\olpseganchor{OLP-0008-B024}{end}

\olpseganchor{OLP-0008-B025}{start}
چې کله د سټونو يو \emph{اندکس} ولرو، يعنې داسې يو سټ \LR{$I$} چې د هر \LR{$i \in I$} دپاره \LR{$A_i$} په پام کښې نيسو، دا لنډې بڼې هم کارولے شو:
\begin{align*}
	\bigcup_{i \in I} A_i & = \bigcup \Setabs{A_i }{i \in I}\\
	\bigcap_{i \in I} A_i & = \bigcap\Setabs{A_i}{i \in I}
\end{align*}
\olpseganchor{OLP-0008-B025}{end}

\olpseganchor{OLP-0008-B026}{start}
په پای کښې، کېدے شي د~\LR{$A$} د هغو ټولو غړو سټ په پام کښې نيول وغواړو چې په~\LR{$B$} کښې نۀ وي۔ دا لکه په \olref{difference} کښې انځورولے شو۔
\olpseganchor{OLP-0008-B026}{end}

\olpseganchor{OLP-0008-B027}{start}
\begin{figure}
  \olasset{../upstream/assets/diagrams/difference.tikz}
  \caption{د دوو سټونو تفاضل \LR{$A \setminus B$} د~\LR{$A$} د هغو غړو سټ دے چې د~\LR{$B$} غړي نۀ وي۔}
  \ollabel{difference}
\end{figure}
\olpseganchor{OLP-0008-B027}{end}

\olpseganchor{OLP-0008-B028}{start}
\begin{defn}[تفاضل]
د سټونو \emph{تفاضل}~\LR{$A \setminus B$} د \LR{$A$} د هغو ټولو غړو سټ دے چې د~\LR{$B$} غړي نۀ وي، يعنې،
\[
A\setminus B = \Setabs{x}{x\in A \text{\RL{ او }} x \notin B}.
\]
\end{defn}
\olpseganchor{OLP-0008-B028}{end}

\olpseganchor{OLP-0008-B029}{start}
\begin{prob}
	ثابت کړئ چې که \LR{$A \subsetneq B$} وي، نو \LR{$B \setminus A \neq \emptyset$}۔
\end{prob}
\olpseganchor{OLP-0008-B029}{end}

\olpseganchor{OLP-0009-B004}{start}
\olfileid{sfr}{set}{pai}
\olsection{جوړې، څوګونې او کارتيزي حاصل ضربونه}
\olpseganchor{OLP-0009-B004}{end}

\olpseganchor{OLP-0009-B005}{start}
\begin{explain}
د غړو له مخې له برابرۍ نه دا نتيجه راوځي چې د سټونو عناصر ترتيب نۀ لري۔ نو که د ترتيب استازيتوب کول غواړو، \emph{مرتبې جوړې} \LR{$\tuple{x, y}$} کاروو۔ په يوه نامرتبه جوړه \LR{$\{x, y\}$} کښې ترتيب مهم نۀ دے: \LR{$\{x, y\} = \{y, x\}$}۔ په مرتبه جوړه کښې ترتيب مهم دے: که \LR{$x \neq y$} وي، نو \LR{$\tuple{x, y} \neq \tuple{y, x}$}۔
\olpseganchor{OLP-0009-B005}{end}

\olpseganchor{OLP-0009-B006}{start}
د سټونو په نظريه کښې مرتبې جوړې څنګه په پام کښې ونيسو؟ بنيادي خبره دا ده چې دا نظر وساتو: مرتبې جوړې هله او يوازې هله برابرې وي چې لومړے عنصر ئې هم يو وي او دويم عنصر ئې هم يو وي، يعنې:
\[
  \tuple{a, b}= \tuple{c, d}\text{\RL{ هله او يوازې هله چې دواړه }}a = c \text{\RL{ او }}b=d.
\]
د وينر--کوراتوسکي تعريف په کارولو سره د سټونو په نظريه کښې مرتبې جوړې تعريفولے شو۔
\end{explain}
\olpseganchor{OLP-0009-B006}{end}

\olpseganchor{OLP-0009-B007}{start}
\begin{defn}[مرتبه جوړه]\ollabel{wienerkuratowski}
	\LR{$\tuple{a, b} = \{\{a\}, \{a, b\}\}$}۔
\end{defn}
\olpseganchor{OLP-0009-B007}{end}

\olpseganchor{OLP-0009-B008}{start}
\begin{prob}
	د \olref[sfr][set][pai]{wienerkuratowski} په کارولو سره ثابت کړئ چې \LR{$\tuple{a, b}= \tuple{c, d}$} هله او يوازې هله وي چې \LR{$a = c$} او \LR{$b=d$} دواړه وي۔
\end{prob}
\olpseganchor{OLP-0009-B008}{end}

\olpseganchor{OLP-0009-B009}{start}
\begin{explain}
چې د مرتبې جوړې تعريف مو وټاکلو، د نورو سټونو د تعريف دپاره ئې کارولے شو۔ د بېلګې په توګه، کله کله له دوو څخه د زياتو څيزونو مرتبې لړۍ هم غواړو، لکه \emph{درې ګونې} \LR{$\tuple{x, y, z}$}، \emph{څلورګونې} \LR{$\tuple{x, y, z, u}$}، او داسې نورې۔ درې ګونې د ځانګړو مرتبو جوړو په توګه ګڼلے شو، چې لومړے عنصر ئې پخپله يوه مرتبه جوړه وي: \LR{$\tuple{x, y, z}$} هم \LR{$\tuple{\tuple{x, y},z}$} دے۔ د څلورګونو دپاره هم همداسې ده: \LR{$\tuple{x,y,z,u}$} هم \LR{$\tuple{\tuple{\tuple{x,y},z},u}$} دے، او داسې نور۔ په عمومي توګه، د \emph{مرتبو \LR{$n$}-ګونو} \LR{$\tuple{x_1, \dots, x_n}$} خبره کوو۔
\olpseganchor{OLP-0009-B009}{end}

\olpseganchor{OLP-0009-B010}{start}
د مرتبو جوړو، يا د نورو مرتبو \LR{$n$}-ګونو، ځينې سټونه به ګټور وي۔
\end{explain}
\olpseganchor{OLP-0009-B010}{end}

\olpseganchor{OLP-0009-B011}{start}
\begin{defn}[کارتيزي حاصل ضرب]
د ورکړل شوو سټونو \LR{$A$} او \LR{$B$} \emph{کارتيزي حاصل ضرب} \LR{$A \times B$} داسې تعريفېږي:
\[
  A \times B = \Setabs{\tuple{x, y}}{x \in A \text{\RL{ او }} y \in B}.
\]
\end{defn}
\olpseganchor{OLP-0009-B011}{end}

\olpseganchor{OLP-0009-B012}{start}
\begin{ex}
که \LR{$A = \{0, 1\}$} او \LR{$B = \{1, a, b\}$} وي، نو حاصل ضرب ئې دا دے:
\[
A \times B = \{ \tuple{0, 1}, \tuple{0, a}, \tuple{0, b},
    \tuple{1, 1}, \tuple{1, a}, \tuple{1, b} \}.
\]
\end{ex}
\olpseganchor{OLP-0009-B012}{end}

\olpseganchor{OLP-0009-B013}{start}
\begin{ex}
که \LR{$A$} يو سټ وي، له خپل ځان سره د \LR{$A$} حاصل ضرب، \LR{$A \times A$}، د~\LR{$A^2$} په بڼه هم ليکل کېږي۔ دا د هغو \emph{ټولو} جوړو \LR{$\tuple{x, y}$} سټ دے چې \LR{$x, y \in A$} وي۔ د ټولو درې ګونو \LR{$\tuple{x, y, z}$} سټ \LR{$A^3$} دے، او داسې نور۔ يو بازګشتي تعريف ورکولے شو:
\begin{align*}
  A^1 & = A\\
  A^{k+1} & = A^k \times A
\end{align*}
\end{ex}
\olpseganchor{OLP-0009-B013}{end}

\olpseganchor{OLP-0009-B014}{start}
\begin{prob}
د \LR{$\{1, 2, 3\}^3$} ټول غړي وليکئ۔
\end{prob}
\olpseganchor{OLP-0009-B014}{end}

\olpseganchor{OLP-0009-B015}{start}
\begin{prop}\ollabel{cardnmprod}
که \LR{$A$} د \LR{$n$} په شمېر غړي او \LR{$B$} د \LR{$m$} په شمېر غړي ولري، نو \LR{$A \times B$} د \LR{$n\cdot m$} په شمېر عناصر لري۔
\end{prop}
\olpseganchor{OLP-0009-B015}{end}

\olpseganchor{OLP-0009-B016}{start}
\begin{proof}
په~\LR{$A$} کښې د هر غړي~\LR{$x$} دپاره، د \LR{$\tuple{x, y} \in A \times B$} په بڼه د \LR{$m$} په شمېر غړي شته۔ فرض کړئ چې \LR{$B_x = \Setabs{\tuple{x, y}}{y \in B}$}۔ څرنګه چې هر کله \LR{$x_1 \neq x_2$} وي، \LR{$\tuple{x_1, y} \neq \tuple{x_2, y}$} وي، نو \LR{$B_{x_1} \cap B_{x_2} = \emptyset$}۔ خو که \LR{$A = \{x_1, \dots, x_n\}$} وي، نو \LR{$A \times B = B_{x_1} \cup \dots \cup B_{x_n}$}، او ځکه د \LR{$n\cdot m$} په شمېر غړي لري۔
\olpseganchor{OLP-0009-B016}{end}

\olpseganchor{OLP-0009-B017}{start}
د دې د انځورولو دپاره، د~\LR{$A \times B$} غړي په يوه جالۍ کښې ترتيب کړئ:
\[
\begin{array}{rcccc}
  B_{x_1} = & \{\tuple{x_1, y_1} & \tuple{x_1, y_2} & \dots & \tuple{x_1, y_m}\}\\
  B_{x_2} = & \{\tuple{x_2, y_1} & \tuple{x_2, y_2} & \dots & \tuple{x_2, y_m}\}\\
  \vdots & & \vdots\\
  B_{x_n} = & \{\tuple{x_n, y_1} & \tuple{x_n, y_2} & \dots & \tuple{x_n, y_m}\}
\end{array}
\]
څرنګه چې ټول \LR{$x_i$} يو له بله بېل دي، او ټول \LR{$y_j$} يو له بله بېل دي، په دې جالۍ کښې هېڅ دوه جوړې يو شان نۀ دي، او شمېر ئې \LR{$n\cdot m$} دے۔
\end{proof}
\olpseganchor{OLP-0009-B017}{end}

\olpseganchor{OLP-0009-B018}{start}
\begin{prob}
په~\LR{$k$} د استقرا په وسيله وښايئ چې د هر \LR{$k \ge 1$} دپاره، که \LR{$A$} د \LR{$n$} په شمېر غړي ولري، نو \LR{$A^k$} د \LR{$n^k$} په شمېر غړي لري۔
\end{prob}
\olpseganchor{OLP-0009-B018}{end}

\olpseganchor{OLP-0009-B019}{start}
\begin{ex}
که \LR{$A$} يو سټ وي، په~\LR{$A$} يوه \emph{کلمه} د~\LR{$A$} د غړو هره لړۍ ده۔ يوه لړۍ د~\LR{$A$} د غړو د يوې \LR{$n$}-ګونې په توګه ګڼل کېدے شي۔ د بېلګې په توګه، که \LR{$A = \{a, b, c\}$} وي، نو لړۍ ``\LR{$bac$}'' د درې ګونې~\LR{$\tuple{b, a, c}$} په توګه ګڼل کېدے شي۔ کلمې، يعنې د نښو لړۍ، په کمپيوټر ساينس کښې بنيادي اهميت لري۔ د منل شوي دود له مخې، د~\LR{$A$} غړي د~\LR{$1$} اوږدوالي لرونکې لړۍ ګڼو، او \LR{$\emptyset$} د~\LR{$0$} اوږدوالي لرونکې لړۍ ګڼو۔ بيا په~\LR{$A$} د \emph{ټولو} کلمو سټ دا دے:
\[
A^* = \{\emptyset\} \cup A \cup A^2 \cup A^3 \cup \dots
\]
\end{ex}
\olpseganchor{OLP-0009-B019}{end}

\olpseganchor{OLP-0010-B004}{start}
\olfileid{sfr}{set}{rus}
\olsection{د رسل پاراډوکس}
\olpseganchor{OLP-0010-B004}{end}

\olpseganchor{OLP-0010-B005}{start}
د غړو له مخې برابري د هغو \LR{$x$} ګانو د \emph{هم هغه} سټ دپاره د \LR{$\Setabs{x}{\phi(x)}$} نښې د کارولو جواز راکوي چې~\LR{$\phi(x)$} وي۔ خو د غړو له مخې برابري \emph{په اصل کښې} يوازې د دې خبرې جواز راکوي: \emph{که} داسې يو سټ وي چې غړي ئې ټول او يوازې هغه څيزونه وي چې \LR{$\phi$} وي، \emph{نو} داسې سټ يوازې يو وي۔ په بل عبارت: د يو~\LR{$\phi$} له ټاکلو وروسته، سټ \LR{$\Setabs{x}{\phi(x)}$} يوازينے دے، \emph{که موجود وي}۔
\olpseganchor{OLP-0010-B005}{end}

\olpseganchor{OLP-0010-B006}{start}
خو دا شرط مهم دے!{} بنيادي خبره دا ده چې له هر خاصيت څخه د \emph{سټ جوړول} ممکن نۀ دي۔ يعنې ځينې خاصيتونه سټونه \emph{نۀ} تعريفوي۔ که ټولو تعريفولے، نو له ښکاره تناقضونو سره به مخ شوو۔ د دې تر ټولو مشهوره بېلګه د رسل پاراډوکس دے۔
\olpseganchor{OLP-0010-B006}{end}

\olpseganchor{OLP-0010-B007}{start}
سټونه د نورو سټونو غړي کېدے شي؛ د بېلګې په توګه، د يو سټ~\LR{$A$} د فرعي سټونو سټ له سټونو جوړ وي۔ نو دا پوښتل يا څېړل معقول دي چې يو سټ د بل سټ غړے دے که نۀ۔ آيا يو سټ د خپل ځان غړے کېدے شي؟ د سټ په تصور کښې ظاهراً هېڅ داسې خبره نشته چې دا ناممکنه کړي۔ د بېلګې په توګه، که \emph{ټول} سټونه د څيزونو يوه ټولګه جوړوي، يو څوک ښائي فکر وکړي چې دا ټول په يو سټ کښې راټولېدے شي، يعنې د ټولو سټونو په سټ کښې۔ او دا چې پخپله يو سټ دے، نو د ټولو سټونو د سټ غړے به وي۔
\olpseganchor{OLP-0010-B007}{end}

\olpseganchor{OLP-0010-B008}{start}
د رسل پاراډوکس هغه وخت راولاړېږي چې دا خاصيت په پام کښې ونيسو چې يو سټ خپل ځان د غړے په توګه ونۀ لري، يعنې \emph{د خپل ځان غړے نۀ وي}۔ که فرض کړو چې د ټولو هغو سټونو يو سټ شته چې خپل ځان د غړے په توګه نۀ لري، نو څه به وشي؟ آيا
\[
R = \Setabs{x}{x \notin x}
\]
شته؟ څرګندېږي چې ثابتولے شو دا نشته۔
\olpseganchor{OLP-0010-B008}{end}

\olpseganchor{OLP-0010-B009}{start}
\begin{thm}[د رسل پاراډوکس]\ollabel{thm:russells-paradox}
	د \LR{$R = \Setabs{x}{x \notin x}$} په بڼه هېڅ سټ نشته۔
\end{thm}
\olpseganchor{OLP-0010-B009}{end}

\olpseganchor{OLP-0010-B010}{start}
\begin{proof}
که \LR{$R = \Setabs{x}{x \notin x}$} موجود وي، نو \LR{$R \in R$} هله او يوازې هله وي چې \LR{$R \notin R$} وي، او دا يو تناقض دے۔
\end{proof}
\olpseganchor{OLP-0010-B010}{end}

\olpseganchor{OLP-0010-B011}{start}

\begin{explain}
راځئ دا ثبوت لږ ورو وګورو۔ که \LR{$R$} موجود وي، دا پوښتنه معقوله ده چې \LR{$R \in R$} دے که نۀ۔ فرض کړئ چې په رښتيا \LR{$R \in R$}۔ اوس، \LR{$R$}~د هغو ټولو سټونو د سټ په توګه تعريف شوے ؤ چې د خپلو ځانونو غړي نۀ دي۔ نو که \LR{$R \in R$} وي، بيا \LR{$R$} پخپله د \LR{$R$} تعريفوونکے خاصيت نۀ لري۔ خو يوازې هغه سټونه په~\LR{$R$} کښې دي چې دا خاصيت لري؛ ځکه \LR{$R$} د~\LR{$R$} غړے نۀ شي کېدے، يعنې \LR{$R \notin R$}۔ خو \LR{$R$} په يو وخت د~\LR{$R$} غړے هم نۀ شي کېدے او ناغړے هم، نو له تناقض سره مخ يو۔
\olpseganchor{OLP-0010-B011}{end}

\olpseganchor{OLP-0010-B012}{start}
څرنګه چې د \LR{$R \in R$} فرض تناقض ته رسوي، نو \LR{$R \notin R$} لرو۔ خو دا هم تناقض ته رسوي!{} ځکه که \LR{$R \notin R$} وي، نو \LR{$R$} پخپله د \LR{$R$} تعريفوونکے خاصيت لري، او ځکه \LR{$R$} به د هغو ټولو نورو سټونو په شان د \LR{$R$} غړے وي چې د خپل ځان غړي نۀ دي۔ او بيا هم، په يو وخت د~\LR{$R$} ناغړے او غړے دواړه نۀ شي کېدے۔
\end{explain}

\olpseganchor{OLP-0010-B012}{end}

\olpseganchor{OLP-0010-B013}{start}
\begin{digress}
د سټونو داسې نظريه څنګه جوړه کړو چې د رسل له پاراډوکس څخه ځان وساتي، يعنې له دې \emph{متناقضې} دعوې څخه ځان وساتي چې \LR{$R = \Setabs{x}{x \notin x}$} موجود دے؟ د دې دپاره داسې بديهي اصلونه ټاکل پکار دي چې دا بيان کړي چې سټونه کله موجود وي (او کله نۀ وي)، او ډېر دقيق شرطونه ورکړي۔
\olpseganchor{OLP-0010-B013}{end}

\olpseganchor{OLP-0010-B014}{start}
د سټونو هغه نظريه چې په دې باب کښې په لنډه توګه بيان شوې، دا کار نۀ کوي۔ دا \emph{په رښتيا ساده نظريه} ده۔ دا يوازې درته وائي چې سټونه د غړو له مخې د برابرۍ پابند دي، او که ځينې سټونه ولرئ، د هغو اتحاد، تقاطع او داسې نور جوړولے شئ۔ د سټونو نظريه له دې نه په زيات دقت هم جوړېدل ممکن دي۔ 
\end{digress}
\olpseganchor{OLP-0010-B014}{end}

\olpseganchor{OLP-0011-B004}{start}
\olchapter{sfr}{rel}{اړيکې}
\olpseganchor{OLP-0011-B004}{end}

\olpseganchor{OLP-0012-B004}{start}
\olfileid{sfr}{rel}{set}
\olsection{اړيکې د سټونو په توګه}
\olpseganchor{OLP-0012-B004}{end}

\olpseganchor{OLP-0012-B005}{start}
\begin{explain}
په \olref[sfr][set][imp]{sec} کښې مو د ځينو مهمو سټونو يادونه کړې وه:
\LR{$\Nat$}، \LR{$\Int$}، \LR{$\Rat$}، \LR{$\Real$}۔ بې شکه به تاسو ته د دې سټونو
د ځينو د غړو ترمنځ په زړه پورې اړيکې يادې وي۔
د بېلګې په توګه، په دې هر سټ باندې يوه پوره معياري \emph{ترتيبي اړيکه}
شته۔ د \emph{هماغه شے کېدو} اړيکه هم شته چې هر څيز ئې له خپل ځان
سره لري او له بل هېڅ څيز سره ئې نۀ لري۔ مونږ به له ډېرو نورو په زړه
پورو اړيکو سره هم مخامخ شو، او ممکنې اړيکې تر دې هم ډېرې دي۔ خو تر
کتنې وړاندې به په دې خبره پيل وکړو چې اړيکې د سټ يو خاص ډول ګڼلے شو۔
\olpseganchor{OLP-0012-B005}{end}

\olpseganchor{OLP-0012-B006}{start}
د دې دپاره له \olref[sfr][set][pai]{sec} څخه دوه خبرې رايادې کړئ۔
لومړے، د \emph{مرتبې جوړې} مفهوم: که \LR{$a$} او \LR{$b$} راکړل شي، نو
\LR{$\tuple{a, b}$} جوړولے شو۔ دلته د غړو ترتيب \emph{خامخا} اهميت لري۔
نو که \LR{$a \neq b$} وي، \LR{$\tuple{a, b} \neq \tuple{b, a}$}۔
(دا له نامرتبو جوړو، يعنې له \LR{$2$}-غړيزو سټونو، سره پرتله کړئ؛ هلته
\LR{$\{a, b\}=\{b, a\}$}۔) دويم، د \emph{کارتيزي حاصل ضرب} مفهوم
راياد کړئ: که \LR{$A$} او \LR{$B$} سټونه وي، نو \LR{$A \times B$} جوړولے شو۔
دا د ټولو هغو جوړو \LR{$\tuple{x, y}$} سټ دے چې \LR{$x \in A$} او \LR{$y \in B$}
وي۔ په تېره بيا، \LR{$A^{2}= A \times A$} د \LR{$A$} د ټولو مرتبو جوړو سټ دے۔
\olpseganchor{OLP-0012-B006}{end}

\olpseganchor{OLP-0012-B007}{start}
اوس به په يو سټ باندې يوه خاصه اړيکه وګورو: د طبيعي عددونو په سټ
\LR{$\Nat$} باندې د \LR{$<$} اړيکه۔ د عددونو د ټولو هغو جوړو \LR{$\tuple{n, m}$}
سټ په پام کښې ونيسئ چې \LR{$n<m$} وي، يعنې:
\[
  R=\Setabs{\tuple{n, m}}{n, m \in \Nat \text{\RL{ او }} n<m}.
\]
د \LR{$n$} له \LR{$m$} څخه د کوچني کېدو او په \LR{$R$} کښې د جوړې \LR{$\tuple{n, m}$}
د غړي کېدو ترمنځ نژدې تړاو شته، يعنې:
\[
      n<m\text{\RL{ هله او يوازې هله چې }}\tuple{n, m} \in R.
\]
په رښتيا، د معلوماتو له هېڅ ضايع کېدو پرته، سټ \LR{$R$} په \LR{$\Nat$} باندې
د \LR{$<$} اړيکه \emph{ګڼلے} شو۔
\olpseganchor{OLP-0012-B007}{end}

\olpseganchor{OLP-0012-B008}{start}
په همدې طريقه، د عددونو ترمنځ د هرې اړيکې دپاره د \LR{$\Nat^{2}$} يو
فرعي سټ جوړولے شو۔ برعکس، که د عددونو د جوړو هر سټ
\LR{$S \subseteq \Nat^{2}$} راکړل شي، نو د عددونو ترمنځ ورسره سمه يوه
اړيکه شته: \LR{$n$} له \LR{$m$} سره دا اړيکه هله او يوازې هله لري چې
\LR{$\tuple{n, m} \in S$} وي۔ دا خبره لاندې تعريف ته جواز ورکوي:
\end{explain}
\olpseganchor{OLP-0012-B008}{end}

\olpseganchor{OLP-0012-B009}{start}
\begin{defn}[دوه ځاييزه اړيکه]
په سټ \LR{$A$} باندې يوه \emph{دوه ځاييزه اړيکه} د \LR{$A^{2}$} يو فرعي سټ
دے۔ که \LR{$R \subseteq A^{2}$} په \LR{$A$} باندې يوه دوه ځاييزه اړيکه وي او
\LR{$x, y \in A$} وي، نو کله کله د \LR{$\tuple{x, y} \in R$} پر ځاے
\LR{$Rxy$} (يا \LR{$xRy$}) ليکو۔
\end{defn}
\olpseganchor{OLP-0012-B009}{end}

\olpseganchor{OLP-0012-B010}{start}
\begin{ex}
  \ollabel{relations}
د طبيعي عددونو د جوړو سټ \LR{$\Nat^{2}$} په يوه دوه بعدي جدول کښې داسې
ليکلے شو:
\[
  \begin{array}{ccccc}
  \mathbf{\tuple{ 0,0 }} & \tuple{ 0,1 } &
    \tuple{ 0,2 } & \tuple{ 0,3 } & \ldots\\
  \tuple{ 1,0 } & \mathbf{\tuple{ 1,1 }} &
    \tuple{ 1,2 } & \tuple{ 1,3 } & \ldots\\
  \tuple{ 2,0 } & \tuple{ 2,1 } &
    \mathbf{\tuple{ 2,2 }} & \tuple{ 2,3 } & \ldots\\
  \tuple{ 3,0 } & \tuple{ 3,1 } & \tuple{ 3,2 } &
    \mathbf{\tuple{ 3,3 }} & \ldots\\
  \vdots & \vdots & \vdots & \vdots & \mathbf{\ddots}
  \end{array}
\]
دلته مو قطر په پنډ ليک ښودلے، ځکه د \LR{$\Nat^2$} هغه فرعي سټ چې پر
قطر پرتې جوړې لري، يعنې:
\[
  \{\tuple{0,0 }, \tuple{ 1,1 }, \tuple{ 2,2 }, \dots\},
  \]
په \LR{$\Nat$} باندې \emph{د عينيت اړيکه} ده۔ (د عينيت اړيکه ډېره کارېږي؛
راځئ د هر سټ \LR{$A$} دپاره
\LR{$\Id{A}=\Setabs{\tuple{ x,x }}{x \in A}$} تعريف کړو۔) د قطر له پاسه
د ټولو پرتو جوړو فرعي سټ، يعنې:
\[
  L = \{\tuple{ 0,1 },\tuple{ 0,2 },\ldots,\tuple{ 1,2 },
  \tuple{ 1,3 }, \dots, \tuple{ 2,3 }, \tuple{ 2,4 },\ldots\},
\]
د \emph{کوچني کېدو} اړيکه ده؛ يعنې \LR{$Lnm$} هله او يوازې هله چې
\LR{$n<m$} وي۔ د قطر لاندې د پرتو جوړو فرعي سټ، يعنې:
\[
  G=\{ \tuple{ 1,0 },\tuple{ 2,0 },\tuple{
    2,1 }, \tuple{ 3,0 },\tuple{ 3,1 },\tuple{ 3,2 }, \dots\},
\]
د \emph{لويېدو} اړيکه ده؛ يعنې \LR{$Gnm$} هله او يوازې هله چې \LR{$n>m$}
وي۔ د \LR{$L$} او \LR{$I$} اتحاد، چې \LR{$K=L\cup I$} ئې بللے شو، د
\emph{کوچني يا برابر کېدو} اړيکه ده: \LR{$Knm$} هله او يوازې هله چې
\LR{$n \le m$} وي۔ همدارنګه، \LR{$H=G \cup I$} د \emph{لوي يا برابر کېدو
اړيکه} ده۔ دا اړيکې، \LR{$L$}، \LR{$G$}، \LR{$K$} او \LR{$H$}، د اړيکو خاص ډولونه دي
چې \emph{ترتيبونه} بللے شي۔ \LR{$L$} او \LR{$G$} دا خاصيت لري چې هېڅ عدد له
خپل ځان سره \LR{$L$} يا \LR{$G$} اړيکه نۀ لري (يعنې، د هر \LR{$n$} دپاره، نۀ \LR{$Lnn$}
شته او نۀ \LR{$Gnn$})۔ د دې خاصيت لرونکې اړيکې \emph{نفي انعکاسي}
بللے شي، او که ترتيبونه هم وي، نو \emph{سخت ترتيبونه} بللے شي۔
\end{ex}
\olpseganchor{OLP-0012-B010}{end}

\olpseganchor{OLP-0012-B011}{start}
\begin{explain}
که څه هم ترتيبونه او عينيت مهمې او طبيعي اړيکې دي، بايد ټينګار وشي
چې زمونږ د تعريف له مخې د \LR{$A^{2}$} \emph{هر} فرعي سټ په \LR{$A$} باندې
يوه اړيکه ده، که هر څومره غيرطبيعي يا په تکلف جوړه شوې ښکاري۔
په تېره بيا، \LR{$\emptyset$} په هر سټ باندې يوه اړيکه ده (\emph{تشه اړيکه}،
چې د غړو هېڅ جوړه ئې نۀ لري)۔ پخپله \LR{$A^{2}$} هم په \LR{$A$} باندې يوه
اړيکه ده (چې هره جوړه ئې لري)، او \emph{ټولشموله اړيکه} بللے شي۔
خو د دې په شان يو څيز هم اړيکه شمېرلے شي:
\LR{$E=\Setabs{\tuple{n, m}}{n>5 \text{\RL{ يا }} m \times n \ge 34}$}۔
\end{explain}
\olpseganchor{OLP-0012-B011}{end}

\olpseganchor{OLP-0012-B012}{start}
\begin{prob}
  په سټ \LR{$\Pow{\{a, b, c\}}$} باندې د \LR{$\subseteq$} اړيکې
  غړي په لړ کښې وليکئ۔
\end{prob}
\olpseganchor{OLP-0012-B012}{end}

\olpseganchor{OLP-0013-B003}{start}
\olfileid{sfr}{rel}{ref}
\olsection{فلسفي غور}
\olpseganchor{OLP-0013-B003}{end}

\olpseganchor{OLP-0013-B004}{start}
په \olref[set]{sec} کښې مو اړيکې د خاصو سټونو په توګه تعريف کړې۔
لږ تم کېږو او يوه لنډه فلسفي پوښتنه کوو: داسې تعريف څۀ \emph{کوي}؟
دا ډېره د شک خبره ده چې مونږ دې ووايو چې د ماوراءالطبيعتي عينيت ځينې
حقيقتونه مو \emph{کشف} کړي دي؛ مثلاً، دا چې په \LR{$\Nat$} باندې ترتيبي
اړيکه \emph{په اصل کښې} هغه سټ
\LR{$R=  \Setabs{\tuple{n,m}}{n, m \in \Nat\text{\RL{ او }} n < m}$}
\emph{وختله} چې په \olref[set]{sec} کښې مو تعريف کړے و۔ د دې خبرې
درې وجهې دا دي۔
\olpseganchor{OLP-0013-B004}{end}

\olpseganchor{OLP-0013-B005}{start}
لومړۍ: په \olref[set][pai]{wienerkuratowski} کښې مو
\LR{$\tuple{a, b} = \{\{a\}, \{a, b\}\}$} تعريف کړے و۔ د دې پر ځاے
دا تعريف وګورئ:
\LR{$\lVert a, b\rVert = \{\{b\}, \{a, b\}\} = \tuple{b,a}$}۔
کله چې \LR{$a \neq b$} وي، نو \LR{$\tuple{a, b} \neq \lVert a,b\rVert$}۔
خو د \LR{$\tuple{a,b}$} پر ځاے مو همداسې \LR{$\lVert a,b\rVert$} هم د مرتبې
جوړې تعريف ګڼلے شو۔ دواړو تعريفونو به يو شان ښۀ کار کړے و۔ نو اوس
د طبيعي عددونو ترتيبي اړيکه ``کېدو'' ته دوه يو شان ښۀ نوماندان لرو:
\begin{align*}
		R &= \Setabs{\tuple{n,m}}{n, m \in \Nat \text{\RL{ او }}n < m}\\
		S &= \Setabs{\lVert n,m\rVert}{n, m \in \Nat \text{\RL{ او }}n < m}.
\end{align*}
ځکه چې د غړو له مخې د برابرۍ د اصل له مخې \LR{$R \neq S$} دي، ښکاره
ده چې \emph{دواړه} په \LR{$\Nat$} باندې له ترتيبي اړيکې سره عينيت نۀ شي
لرلے۔ خو دا به محض خپله خوښه، او له دې امله يو څه د شرم خبره وي، چې
ادعا وکړو چې په واقعيت کښې \LR{$R$}، د \LR{$S$} پر ځاے (يا برعکس)، ترتيبي
اړيکه \emph{ده}۔ (دا د سټونو ته د راکمونې د نظر خلاف د هغه دليل يوه
ډېره ساده بېلګه ده چې بناسراف په \citeyear{Benacerraf1965} کښې مشهور
کړے و۔ مونږ به څو ځله بيا دې ته راګرځو۔)
\olpseganchor{OLP-0013-B005}{end}

\olpseganchor{OLP-0013-B006}{start}
دويمه: که فکر کوو چې \emph{هره} اړيکه بايد له يو سټ سره عينيت
ولري، نو پخپله د سټ د غړيتوب اړيکه، \LR{$\in$}، هم بايد يو خاص سټ وي۔
په رښتيا، هغه به خامخا سټ \LR{$\Setabs{\tuple{x,y}}{x \in y}$} وي۔
خو آيا دا سټ شته؟ د رسل پاراډوکس ته په کتو، د داسې سټ د موجوديت
ادعا ساده خبره نۀ ده۔ په حقيقت کښې،
د سټونو
نظريه په يوه دقيقه طريقه، د بديهي اصولو د يوې نظريې په توګه، جوړول
ممکن دي، او هغه نظريه به په رښتيا د دې سټ موجوديت رد کړي۔
نو که ځينې اړيکې د سټونو په توګه هم ګڼلے شو، د سټ د غړيتوب اړيکه
به يوه خاصه پېښه وي۔
\olpseganchor{OLP-0013-B006}{end}

\olpseganchor{OLP-0013-B007}{start}
درېيمه: کله چې اړيکې له سټونو سره ``يوې'' ګڼو، ومو وئيل چې ځان ته
به د \LR{$\tuple{x,y} \in R$} پر ځاے د \LR{$Rxy$} ليکلو اجازه ورکوو۔ دا سمه
ده، په دې شرط چې د غړيتوب اړيکه، ``\LR{$\in$}''، د يو پريديکات
\emph{په توګه} وکارول شي۔ خو که فکر کوو چې ``\LR{$\in$}'' د سټ د يو
خاص ډول نښه ده، نو عبارت ``\LR{$\tuple{x,y} \in R$}'' يوازې له درېيو
مفردو اصطلاحاتو جوړ دے چې سټونه ښيي: ``\LR{$\tuple{x,y}$}''، ``\LR{$\in$}''
او ``\LR{$R$}''۔ د نومونو دا ډول لړ، لکه بې معنا توريز تسلسل ``پياله
قلمدان مېز''، د يو بیان د څرګندولو وس نۀ لري۔ نو بيا هم، که ځينې
اړيکې د سټونو په توګه ګڼلے شو، د سټ د غړيتوب اړيکه بايد يوه خاصه
پېښه وي۔ (په دې کښې د فرېګه د \emph{آس} د مفهوم د پاراډوکس يوه
ساده بڼه او هغه مشهور اعتراض سره يوځاے شوي چې وېټګنشټاين يو وخت
پر رسل کړے و۔)
\olpseganchor{OLP-0013-B007}{end}

\olpseganchor{OLP-0013-B008}{start}
نو دې خبرې مونږ کوم ځاے ته ورسولو؟ دا وئيل چې په عددونو باندې
اړيکې سټونه دي، هېڅ \emph{غلط} نۀ دي۔ يوازې بايد په هغې موخې پوه
شو چې دا خبره ورسره کوو۔ مونږ د ماوراءالطبيعتي عينيت يو حقيقت نۀ
بيانوو۔ يوازې دا يادونه کوو چې په ځينو سياقونو کښې (ځينې) اړيکې
خاص سټونه \emph{ګڼلے} شو (او ګڼو به ئې)۔
\olpseganchor{OLP-0013-B008}{end}

\olpseganchor{OLP-0014-B004}{start}
\olfileid{sfr}{rel}{prp}
\olsection{د اړيکو خاص خاصيتونه}
\olpseganchor{OLP-0014-B004}{end}

\olpseganchor{OLP-0014-B005}{start}
\begin{intro}
د اړيکو ځينې ډولونه دومره عام دي چې خاص نومونه ورکړل شوي دي۔
د بېلګې په توګه، \LR{$\le$} او \LR{$\subseteq$} دواړه په خپلو خپلو ساحو کښې
(مثلاً، د \LR{$\le$} دپاره په \LR{$\Nat$} او د \LR{$\subseteq$} دپاره په \LR{$\Pow{A}$}
کښې) په يو شان طريقو اړيکه جوړوي۔ د دې اړيکو د ورته والي او توپير
د دقيق معلومولو دپاره، د هغو خاصو خاصيتونو له مخې ئې وېشو چې
اړيکې ئې لرلے شي۔ څرګندېږي چې د دې ځينو خاصو خاصيتونو (ترکيبونه)
په تېره بيا مهم دي: ترتيبونه او د هم ارزښتۍ اړيکې۔
\end{intro}
\olpseganchor{OLP-0014-B005}{end}

\olpseganchor{OLP-0014-B006}{start}
\begin{defn}[انعکاس]
اړيکه \LR{$R \subseteq A^2$} هله او يوازې هله \emph{انعکاسي} ده چې
د هر \LR{$x \in A$} دپاره \LR{$Rxx$} وي۔
\end{defn}
\olpseganchor{OLP-0014-B006}{end}

\olpseganchor{OLP-0014-B007}{start}
\begin{defn}[انتقال]
اړيکه \LR{$R \subseteq A^2$} هله او يوازې هله \emph{انتقالي} ده چې
کله هم \LR{$Rxy$} او \LR{$Ryz$} وي، نو \LR{$Rxz$} هم وي۔
\end{defn}
\olpseganchor{OLP-0014-B007}{end}

\olpseganchor{OLP-0014-B008}{start}
\begin{defn}[تناظر]
اړيکه \LR{$R \subseteq A^2$} هله او يوازې هله \emph{متناظره} ده چې
کله هم \LR{$Rxy$} وي، نو \LR{$Ryx$} هم وي۔
\end{defn}
\olpseganchor{OLP-0014-B008}{end}

\olpseganchor{OLP-0014-B009}{start}
\begin{defn}[ضد تناظر]
اړيکه \LR{$R \subseteq A^2$} هله او يوازې هله \emph{ضد متناظره} ده چې
کله هم \LR{$Rxy$} او \LR{$Ryx$} دواړه وي، نو \LR{$x=y$} وي (په نورو ټکو: که
\LR{$x\neq y$} وي، نو يا \LR{$\lnot Rxy$} وي او يا \LR{$\lnot Ryx$})۔
\end{defn}
\olpseganchor{OLP-0014-B009}{end}

\olpseganchor{OLP-0014-B010}{start}
\begin{explain}
په متناظره اړيکه کښې \LR{$Rxy$} او \LR{$Ryx$} تل يوځاے وي، يا يو هم نۀ وي۔
په ضد متناظره اړيکه کښې \LR{$Rxy$} او \LR{$Ryx$} يوازې هله يوځاے کېدلے شي
چې \LR{$x = y$} وي۔ پام وکړئ چې دا خبره د \LR{$x = y$} په وخت د \LR{$Rxy$} او
\LR{$Ryx$} موجوديت \emph{نۀ لازموي}؛ يوازې دا چې مخه ئې نۀ نيول کېږي۔
نو ضد متناظره اړيکه انعکاسي کېدلے شي، خو هره ضد متناظره اړيکه
انعکاسي نۀ وي۔ دا هم په پام کښې وساتئ چې ضد متناظره کېدل او محض
متناظره نۀ کېدل بېل شرطونه دي۔ په رښتيا، يوه اړيکه په يوه وخت
متناظره او ضد متناظره دواړه کېدلے شي (مثلاً، د عينيت اړيکه داسې ده)۔
\end{explain}
\olpseganchor{OLP-0014-B010}{end}

\olpseganchor{OLP-0014-B011}{start}
\begin{defn}[تړلتيا]
اړيکه \LR{$R \subseteq A^2$} \emph{تړلې} ده که د ټولو \LR{$x,y\in A$}
دپاره، که \LR{$x \neq y$} وي، نو يا \LR{$Rxy$} وي او يا \LR{$Ryx$}۔
\end{defn}
\olpseganchor{OLP-0014-B011}{end}

\olpseganchor{OLP-0014-B012}{start}
\begin{prob}
د داسې اړيکو بېلګې ورکړئ چې (الف) انعکاسي او متناظرې وي، خو انتقالي
نۀ وي؛ (ب) انعکاسي او ضد متناظرې وي؛ (ج) ضد متناظرې او انتقالي وي،
خو انعکاسي نۀ وي؛ او (د) انعکاسي، متناظرې او انتقالي وي۔ پر عددونو
يا سټونو باندې اړيکې مۀ کاروئ۔
\end{prob}
\olpseganchor{OLP-0014-B012}{end}

\olpseganchor{OLP-0014-B013}{start}
\begin{defn}[د انعکاس نفي]
اړيکه \LR{$R \subseteq A^2$} \emph{نفي انعکاسي} بللے شي که د هر
\LR{$x \in A$} دپاره \LR{$Rxx$} نۀ وي۔
\end{defn}
\olpseganchor{OLP-0014-B013}{end}

\olpseganchor{OLP-0014-B014}{start}
\begin{defn}[نامتناظرتيا]
اړيکه \LR{$R \subseteq A^2$} \emph{نامتناظره} بللے شي که د هېڅ جوړې
\LR{$x,y\in A$} دپاره \LR{$Rxy$} او \LR{$Ryx$} دواړه ونۀ لرو۔
\end{defn}
\olpseganchor{OLP-0014-B014}{end}

\olpseganchor{OLP-0014-B015}{start}
پام وکړئ: که \LR{$A \neq \emptyset$} وي، نو په \LR{$A$} باندې هېڅ نفي انعکاسي
اړيکه انعکاسي نۀ ده، او په \LR{$A$} باندې هره نامتناظره اړيکه ضد متناظره
هم ده۔ خو داسې \LR{$R \subseteq A^2$} شته چې نۀ انعکاسي دي او نۀ نفي
انعکاسي، او داسې ضد متناظرې اړيکې شته چې نامتناظرې نۀ دي۔
\olpseganchor{OLP-0014-B015}{end}

\olpseganchor{OLP-0015-B004}{start}
\olfileid{sfr}{rel}{eqv}
\olpseganchor{OLP-0015-B004}{end}

\olpseganchor{OLP-0015-B005}{start}
\olsection{د هم ارزښتۍ اړيکې}
\olpseganchor{OLP-0015-B005}{end}

\olpseganchor{OLP-0015-B006}{start}
په يوه سټ باندې د عينيت اړيکه انعکاسي، متناظره او انتقالي ده۔
هغه اړيکې \LR{$R$} چې دا درې واړه خاصيتونه لري، ډېرې عامې دي۔
\olpseganchor{OLP-0015-B006}{end}

\olpseganchor{OLP-0015-B007}{start}
\begin{defn}[د هم ارزښتۍ اړيکه]
هغه اړيکه \LR{$R \subseteq A^2$} چې انعکاسي، متناظره او انتقالي وي،
\emph{د هم ارزښتۍ اړيکه} بللے شي۔ د \LR{$A$} غړي، \LR{$x$} او
\LR{$y$}، د \LR{$R$} له مخې \emph{هم ارزښته} بللے شي که \LR{$Rxy$} وي۔
\end{defn}
\olpseganchor{OLP-0015-B007}{end}

\olpseganchor{OLP-0015-B008}{start}
د هم ارزښتۍ اړيکې د \emph{هم ارزښتۍ د ټولګي} مفهوم راپيدا کوي۔
د هم ارزښتۍ يوه اړيکه ساحه په بېلو برخو ``ټوټه کوي''۔ په هره برخه
کښې ټول څيزونه يو له بل سره اړيکه لري؛ او د بېلو برخو څيزونه يو
له بل سره اړيکه نۀ لري۔ کله کله د همدې برخو په اړه \emph{نېغ په نېغه}
خبرې کول ګټور وي۔ د دې دپاره يو تعريف راوړو:
\olpseganchor{OLP-0015-B008}{end}

\olpseganchor{OLP-0015-B009}{start}
\begin{defn}\ollabel{def:equivalenceclass}
فرض کړئ \LR{$R \subseteq A^2$} د هم ارزښتۍ اړيکه ده۔ د هر \LR{$x \in A$}
دپاره، په \LR{$A$} کښې د \LR{$x$} \emph{د هم ارزښتۍ ټولګے} دا سټ دے:
\LR{$\equivrep{x}{R} = \Setabs{y \in A}{Rxy}$}۔ د \LR{$R$} له مخې د \LR{$A$}
\emph{خارج قسمت} دا دے:
\LR{$\equivclass{A}{R} = \Setabs{\equivrep{x}{R}}{x \in A}$}، يعنې
د هم ارزښتۍ د دې ټولګيو سټ۔
\end{defn}
\olpseganchor{OLP-0015-B009}{end}

\olpseganchor{OLP-0015-B010}{start}
راتلونکې نتيجه د هم ارزښتۍ د ټولګي د تعريف سموالى ښيي: ثابتوي چې
د هم ارزښتۍ ټولګي په رښتيا د \LR{$A$} وېشونکې برخې دي۔
\olpseganchor{OLP-0015-B010}{end}

\olpseganchor{OLP-0015-B011}{start}
\begin{prop}
که \LR{$R \subseteq A^2$} د هم ارزښتۍ اړيکه وي، نو \LR{$Rxy$} هله او يوازې
هله چې \LR{$\equivrep{x}{R} = \equivrep{y}{R}$} وي۔
\end{prop}
\olpseganchor{OLP-0015-B011}{end}

\olpseganchor{OLP-0015-B012}{start}
\begin{proof}
د کيڼې خوا نه ښۍ خوا ته د لوري دپاره، فرض کړئ \LR{$Rxy$} او
\LR{$z \in \equivrep{x}{R}$} وي۔ نو د تعريف له مخې \LR{$Rxz$}۔ ځکه چې \LR{$R$}
د هم ارزښتۍ اړيکه ده، \LR{$Ryz$}۔ (په تفصيل سره: له \LR{$Rxy$} او د \LR{$R$}
له تناظر څخه \LR{$Ryx$} لرو؛ او له \LR{$Rxz$} او د \LR{$R$} له انتقالي والي څخه
\LR{$Ryz$} لرو۔) نو \LR{$z \in \equivrep{y}{R}$}۔ په عمومي توګه،
\LR{$\equivrep{x}{R} \subseteq \equivrep{y}{R}$}۔ خو بالکل په همدې
ډول، \LR{$\equivrep{y}{R} \subseteq \equivrep{x}{R}$}۔ نو د غړو له
مخې د برابرۍ د اصل له مخې، \LR{$\equivrep{x}{R} = \equivrep{y}{R}$}۔
\olpseganchor{OLP-0015-B012}{end}

\olpseganchor{OLP-0015-B013}{start}
د ښۍ خوا نه کيڼې خوا ته د لوري دپاره، فرض کړئ
\LR{$\equivrep{x}{R} = \equivrep{y}{R}$}۔ ځکه چې \LR{$R$} انعکاسي ده،
\LR{$Ryy$}، نو \LR{$y \in \equivrep{y}{R}$}۔ له دې فرضيې څخه چې
\LR{$\equivrep{x}{R} = \equivrep{y}{R}$}، دا هم لرو چې
\LR{$y \in \equivrep{x}{R}$}۔ نو \LR{$Rxy$}۔
\end{proof}
\olpseganchor{OLP-0015-B013}{end}

\olpseganchor{OLP-0015-B014}{start}
\begin{ex}
د هم ارزښتۍ د اړيکو يوه ښۀ بېلګه د باقي له حسابه راځي۔ د هر \LR{$a$}،
\LR{$b$} او \LR{$n \in \PosInt$} دپاره، وايو چې \LR{$a \equiv_n b$} هله او
يوازې هله چې د \LR{$a$} په \LR{$n$} د وېش باقي د \LR{$b$} په \LR{$n$} د وېش له باقي
سره يو شان وي۔ (په يو څه زياته نښيزه ژبه: \LR{$a \equiv_n b$} هله او
يوازې هله چې د کوم \LR{$k \in \Int$} دپاره \LR{$a - b = kn$} وي۔)
اوس د هر \LR{$n$} دپاره \LR{$\equiv_n$} د هم ارزښتۍ اړيکه ده۔ او د
\LR{$\equiv_n$} له مخې جوړ شوي د هم ارزښتۍ ټولګي په شمېر پوره \LR{$n$}
بېل ټولګي دي؛ يعنې \LR{$\equivclass{\Nat}{\equiv_n}$} په شمېر \LR{$n$}
غړي لري۔ دا ټولګي دي: د هغو عددونو سټ چې بې له باقي پر
\LR{$n$} وېشل کېږي، يعنې \LR{$\equivrep{0}{\equiv_n}$}؛ د هغو عددونو سټ
چې پر \LR{$n$} د وېش باقي ئې \LR{$1$} وي، يعنې \LR{$\equivrep{1}{\equiv_n}$}؛
\ldots؛ او د هغو عددونو سټ چې پر \LR{$n$} د وېش باقي ئې \LR{$n-1$} وي،
يعنې \LR{$\equivrep{n-1}{\equiv_n}$}۔
\end{ex}
\olpseganchor{OLP-0015-B014}{end}

\olpseganchor{OLP-0015-B015}{start}
\begin{prob}
وښيئ چې د هر \LR{$n \in \PosInt$} دپاره \LR{$\equiv_n$} د هم ارزښتۍ
اړيکه ده، او \LR{$\equivclass{\Nat}{\equiv_n}$} په شمېر پوره \LR{$n$}
غړي لري۔
\end{prob}
\olpseganchor{OLP-0015-B015}{end}

\olpseganchor{OLP-0016-B004}{start}
\olfileid{sfr}{rel}{ord}
\olsection{ترتيبونه}
\olpseganchor{OLP-0016-B004}{end}

\olpseganchor{OLP-0016-B005}{start}
\begin{explain}
په ډېرو پرتلو کښې ځينې څيزونه، له يوه پلوه، تر نورو ``کوچني''،
ورسره ``برابر'' يا ترې ``لوي'' بولو۔ په دې کښې د \emph{ترتيب}
اړيکې شاملې دي۔ خو د ترتيبي اړيکو بېل ډولونه شته۔ مثلاً، ځينې
لازموي چې هر دوه څيزونه د پرتلې وړ وي، خو نورې داسې نۀ دي۔ ځينې
عينيت هم رانغاړي (لکه \LR{$\le$})، او ځينې ئې باسي (لکه \LR{$<$})۔ دلته د
ډولونو وېش لرل راسره مرسته کوي۔
\end{explain}
\olpseganchor{OLP-0016-B005}{end}

\olpseganchor{OLP-0016-B006}{start}
\begin{defn}[مقدماتي ترتيب]
هغه اړيکه چې انعکاسي او انتقالي دواړه وي، \emph{مقدماتي ترتيب}
بللے شي۔
\end{defn}
\olpseganchor{OLP-0016-B006}{end}

\olpseganchor{OLP-0016-B007}{start}
\begin{defn}[جزوي ترتيب]
هغه مقدماتي ترتيب چې ضد متناظر هم وي، \emph{جزوي ترتيب} بللے شي۔
\end{defn}
\olpseganchor{OLP-0016-B007}{end}

\olpseganchor{OLP-0016-B008}{start}
\begin{defn}[خطي ترتيب]\ollabel{def:linearorder}
هغه جزوي ترتيب چې تړلے هم وي، \emph{کلي ترتيب} يا \emph{خطي ترتيب}
بللے شي۔
\end{defn}
\olpseganchor{OLP-0016-B008}{end}

\olpseganchor{OLP-0016-B009}{start}
\begin{ex}
هر خطي ترتيب جزوي ترتيب هم دے، او هر جزوي ترتيب مقدماتي ترتيب
هم دے؛ خو برعکس خبرې سمې نۀ دي۔ په \LR{$A$} باندې ټولشموله اړيکه
مقدماتي ترتيب دے، ځکه انعکاسي او انتقالي ده۔ خو که \LR{$A$} له يوه
غړے څخه زيات ولري، نو ټولشموله اړيکه ضد متناظره نۀ ده،
او له دې امله جزوي ترتيب هم نۀ دے۔
\end{ex}
\olpseganchor{OLP-0016-B009}{end}

\olpseganchor{OLP-0016-B010}{start}
\begin{ex}
په \LR{$\Bin^*$} باندې د \emph{تر بل نۀ اوږدېدو} اړيکه
\LR{$\preccurlyeq$} وګورئ: \LR{$x \preccurlyeq y$} هله او يوازې هله چې
\LR{$\len{x} \le \len{y}$} وي۔ دا مقدماتي ترتيب دے (انعکاسي او انتقالي
دے)، او تړلے هم دے؛ خو جزوي ترتيب نۀ دے، ځکه ضد متناظر نۀ دے۔
مثلاً، \LR{$01 \preccurlyeq 10$} او \LR{$10 \preccurlyeq 01$}، خو
\LR{$01 \neq 10$}۔
\end{ex}
\olpseganchor{OLP-0016-B010}{end}

\olpseganchor{OLP-0016-B011}{start}
\begin{ex}
د سټونو په يوه سټ باندې د \LR{$\subseteq$} اړيکه يو مهم جزوي ترتيب دے۔
دا په عمومي توګه خطي ترتيب نۀ دے، ځکه که \LR{$a \neq b$} وي او
\LR{$\Pow{\{a, b\}} = \{\emptyset, \{a\}, \{b\}, \{a,b\}\}$}
وګورو، نو وينو چې \LR{$\{a\} \nsubseteq \{b\}$}،
\LR{$\{a\} \neq \{b\}$} او \LR{$\{b\} \nsubseteq \{a\}$}۔
\end{ex}
\olpseganchor{OLP-0016-B011}{end}

\olpseganchor{OLP-0016-B012}{start}
\begin{ex}
د \emph{بې له باقي وېشل کېدو} اړيکه داسې جزوي ترتيب راکوي چې خطي
ترتيب نۀ دے۔ د صحيح عددونو \LR{$n$} او \LR{$m$} دپاره، \LR{$n \mid m$} په دې
معنا ليکو چې \LR{$n$}، \LR{$m$} (بې له باقي) وېشي؛ يعنې هله او يوازې هله
چې داسې صحيح عدد \LR{$k$} شته چې \LR{$m = kn$} وي۔ په \LR{$\Nat$} باندې دا جزوي
ترتيب دے، خو خطي ترتيب نۀ دے: مثلاً، \LR{$2 \nmid 3$} او \LR{$3 \nmid 2$}
هم۔ که وېشمنتيا په \LR{$\Int$} باندې اړيکه وګڼو، يوازې مقدماتي ترتيب
دے، ځکه ضد متناظره نۀ ده: \LR{$1 \mid -1$} او \LR{$-1 \mid 1$}، خو
\LR{$1 \neq -1$}۔
\end{ex}
\olpseganchor{OLP-0016-B012}{end}

\olpseganchor{OLP-0016-B013}{start}
\begin{ex}
د لړيو په سټ \LR{$A^*$} باندې د \emph{غځونې} اړيکه دا ده:
\LR{$s \sqsubseteq s'$} هله او يوازې هله چې \LR{$s = \emptyseq$} (تشه لړۍ)
وي، \LR{$s = s'$} وي، يا \LR{$s = \tuple{s_1, \dots, s_n}$} او
\LR{$s' = \tuple{s_1, \dots, s_n, s_{n+1}, \dots, s_m}$} وي۔ که
\LR{$s \sqsubseteq s'$} وي، دا هم وايو چې \LR{$s$} د \LR{$s'$} \emph{پيلنۍ برخه}
ده۔ په \LR{$A^*$} باندې د غځونې اړيکه جزوي ترتيب دے، خو خطي ترتيب نۀ
دے؛ مثلاً، که \LR{$a \neq b$} وي، نو \LR{$ab \not\sqsubseteq ba$} او
\LR{$ba \not\sqsubseteq ab$}۔
\end{ex}
\olpseganchor{OLP-0016-B013}{end}

\olpseganchor{OLP-0016-B014}{start}
\begin{defn}[سخت ترتيب]
\emph{سخت ترتيب} هغه اړيکه ده چې نفي انعکاسي، نامتناظره او انتقالي وي۔
\end{defn}
\olpseganchor{OLP-0016-B014}{end}

\olpseganchor{OLP-0016-B015}{start}
\begin{defn}[سخت خطي ترتيب]\ollabel{def:strictlinearorder}
هغه سخت ترتيب چې تړلے هم وي، \emph{سخت کلي ترتيب} يا
\emph{سخت خطي ترتيب} بللے شي۔
\end{defn}
\olpseganchor{OLP-0016-B015}{end}

\olpseganchor{OLP-0016-B016}{start}
\begin{ex}
\LR{$\le$} هغه خطي ترتيب دے چې له سخت خطي ترتيب \LR{$<$} سره سمون خوري۔
\LR{$\subseteq$} هغه جزوي ترتيب دے چې له سخت ترتيب \LR{$\subsetneq$} سره
سمون خوري۔
\end{ex}
\olpseganchor{OLP-0016-B016}{end}

\olpseganchor{OLP-0016-B017}{start}
په \LR{$A$} باندې هر سخت ترتيب \LR{$R$} د قطر \LR{$\Id{A}$} په ورزياتولو، يعنې
د ټولو جوړو \LR{$\tuple{x, x}$} په ورزياتولو، جزوي ترتيب ګرځولے شو۔
(دې ته د \LR{$R$} \emph{انعکاسي بندښت} وايي۔) برعکس، له جزوي ترتيبه
په پيل کولو، د \LR{$\Id{A}$} په لرې کولو سخت ترتيب ترلاسه کولے شو۔
لاندې دوه نتيجې دا خبره دقيقه کوي۔
\olpseganchor{OLP-0016-B017}{end}

\olpseganchor{OLP-0016-B018}{start}
\begin{prop}\ollabel{prop:stricttopartial}
که \LR{$R$} په \LR{$A$} باندې سخت ترتيب وي، نو \LR{$R^+ = R \cup \Id{A}$}
جزوي ترتيب دے۔ سربېره پر دې، که \LR{$R$} سخت خطي ترتيب وي، نو \LR{$R^+$}
خطي ترتيب دے۔
\end{prop}
\olpseganchor{OLP-0016-B018}{end}

\olpseganchor{OLP-0016-B019}{start}
\begin{proof}
فرض کړئ \LR{$R$} سخت ترتيب دے؛ يعنې \LR{$R \subseteq A^2$} او \LR{$R$} نفي
انعکاسي، نامتناظر او انتقالي دے۔ \LR{$R^+ = R \cup \Id{A}$} کېږدئ۔
بايد وښيو چې \LR{$R^+$} انعکاسي، ضد متناظر او انتقالي دے۔
\olpseganchor{OLP-0016-B019}{end}

\olpseganchor{OLP-0016-B020}{start}
ښکاره ده چې \LR{$R^+$} انعکاسي دے، ځکه د هر \LR{$x \in A$} دپاره
\LR{$\tuple{x, x} \in \Id{A} \subseteq R^+$}۔
\olpseganchor{OLP-0016-B020}{end}

\olpseganchor{OLP-0016-B021}{start}
د دې د ښودلو دپاره چې \LR{$R^+$} ضد متناظر دے، د تناقض له لارې د ثبوت
دپاره فرض کړئ \LR{$R^+xy$} او \LR{$R^+yx$}، خو \LR{$x \neq y$}۔ ځکه چې
\LR{$\tuple{x,y} \in R \cup \Id{A}$}، خو
\LR{$\tuple{x, y} \notin \Id{A}$}، نو خامخا \LR{$\tuple{x, y} \in R$}
لرو، يعنې \LR{$Rxy$}۔ همدارنګه، \LR{$Ryx$}۔ خو دا له دې فرضيې سره تناقض لري
چې \LR{$R$} نامتناظر دے۔
\olpseganchor{OLP-0016-B021}{end}

\olpseganchor{OLP-0016-B022}{start}
د انتقالي والي د ثابتولو دپاره، فرض کړئ \LR{$R^+xy$} او \LR{$R^+yz$}۔
که \LR{$\tuple{x, y} \in R$} او \LR{$\tuple{y,z} \in R$} دواړه وي، نو
\LR{$\tuple{x, z} \in R$}، ځکه \LR{$R$} انتقالي دے۔ که داسې نۀ وي، نو يا
\LR{$\tuple{x, y} \in \Id{A}$}، يعنې \LR{$x = y$}، يا
\LR{$\tuple{y, z} \in \Id{A}$}، يعنې \LR{$y = z$}۔ په لومړي حالت کښې د
فرضيې له مخې \LR{$R^+yz$} لرو، او \LR{$x = y$}، نو \LR{$R^+xz$}۔ په دويم
حالت کښې هم همداسې ده۔ په هر حالت کښې \LR{$R^+xz$}، نو \LR{$R^+$} انتقالي
هم دے۔
\olpseganchor{OLP-0016-B022}{end}

\olpseganchor{OLP-0016-B023}{start}
د ``سربېره پر دې'' برخې په اړه، فرض کړئ \LR{$R$} تړلے هم دے۔ نو د ټولو
\LR{$x \neq y$} دپاره يا \LR{$Rxy$} يا \LR{$Ryx$}؛ يعنې يا
\LR{$\tuple{x, y} \in R$} يا \LR{$\tuple{y, x} \in R$}۔ ځکه چې
\LR{$R \subseteq R^+$}، دا خبره د \LR{$R^+$} دپاره هم سمه پاتې کېږي۔ نو
\LR{$R^+$} تړلے هم دے۔
\end{proof}
\olpseganchor{OLP-0016-B023}{end}

\olpseganchor{OLP-0016-B024}{start}
\begin{prop}\ollabel{prop:partialtostrict}
که \LR{$R$} په \LR{$A$} باندې جزوي ترتيب وي، نو
\LR{$R^- = R \setminus \Id{A}$} سخت ترتيب دے۔ سربېره پر دې، که \LR{$R$}
خطي ترتيب وي، نو \LR{$R^-$} سخت خطي ترتيب دے۔
\end{prop}
\olpseganchor{OLP-0016-B024}{end}

\olpseganchor{OLP-0016-B025}{start}
\begin{proof}
دا د تمرين په توګه پرېښودل شوے دے۔
\end{proof}
\olpseganchor{OLP-0016-B025}{end}

\olpseganchor{OLP-0016-B026}{start}
\begin{prob}
د \olref[sfr][rel][ord]{prop:partialtostrict} ثبوت ورکړئ۔
\end{prob}
\olpseganchor{OLP-0016-B026}{end}

\olpseganchor{OLP-0016-B027}{start}
لاندې ساده نتيجه ثابتوي چې سخت خطي ترتيبونه د غړو له مخې د
برابرۍ په شان يو خاصيت پوره کوي:
\olpseganchor{OLP-0016-B027}{end}

\olpseganchor{OLP-0016-B028}{start}
\begin{prop}\ollabel{prop:extensionality-strictlinearorders}
که \LR{$<$} په \LR{$A$} باندې سخت خطي ترتيب وي، نو:
\[
  (\forall a, b \in A)((\forall x \in A)(x < a \liff x < b) \lif a = b).
\]
\end{prop}
\olpseganchor{OLP-0016-B028}{end}

\olpseganchor{OLP-0016-B029}{start}
\begin{proof}
فرض کړئ \LR{$(\forall x \in A)(x < a \liff x < b)$}۔ که \LR{$a < b$} وي،
نو \LR{$a < a$}، چې د \LR{$<$} له نفي انعکاسي والي سره تناقض لري؛ نو
\LR{$a \nless b$}۔ بالکل په همدې ډول \LR{$b \nless a$}۔ نو \LR{$a = b$}، ځکه
\LR{$<$} تړلے دے۔
\end{proof}
\olpseganchor{OLP-0016-B029}{end}

\olpseganchor{OLP-0017-B004}{start}
\olfileid{sfr}{rel}{grp}
\olsection{ګرافونه}
\olpseganchor{OLP-0017-B004}{end}

\olpseganchor{OLP-0017-B005}{start}
\emph{ګراف} داسې شکل دے چې ټکي—چې ``غوټې'' يا ``څوکې'' (د
``څوکه'' جمع) بللے شي—پکښې د څنډو په ذريعه سره نښتي وي۔ ګرافونه
په منفصل رياضي او کمپيوټرپوهنه کښې هر ځاے کارېدونکې وسيلې دي۔
د اړيکو او جوړښتونو د ښودلو او انځورولو دپاره ډېر ګټور دي؛ له
محسوسو څيزونو، لکه د بېلو ډولونو له شبکو، نيولې تر مجردو جوړښتونو،
لکه د پرېکړو تر ممکنو پايلو پورې۔ په علمي ليکنو کښې د ګرافونو ډېر
بېل ډولونه شته۔ مثلاً، په دې سره توپير کوي چې څنډې لوری لري که نه،
نوم نښې لري که نه، آيا له يوې غوټې بېرته هماغې غوټې ته څنډې کېدلے
شي، د هماغو غوټو ترمنځ څو څنډې کېدلے شي، او داسې نور۔
\emph{لوری لرونکي ګرافونه} له اړيکو سره يو خاص تړاو لري۔
\olpseganchor{OLP-0017-B005}{end}

\olpseganchor{OLP-0017-B006}{start}
\begin{defn}[لوری لرونکے ګراف]
\emph{لوری لرونکے ګراف} \LR{$G = \tuple{V, E}$} د \emph{څوکو} يو سټ
\LR{$V$} او د \emph{څنډو} يو سټ \LR{$E \subseteq V^2$} دے۔
\end{defn}
\olpseganchor{OLP-0017-B006}{end}

\olpseganchor{OLP-0017-B007}{start}
\begin{explain}
زمونږ د تعريف له مخې، ګراف يوازې يو سټ او په همدغه سټ باندې يوه
اړيکه ده۔ البته، د ګرافونو په اړه د خبرو پر وخت، په شکل کښې د
ښودلو تمه ئې طبيعي ده: ګراف داسې رسمولے شو چې دوه څوکې \LR{$v_1$} او
\LR{$v_2$} په يو غشي هله او يوازې هله سره ونښلوو چې
\LR{$\tuple{v_1, v_2} \in E$} وي۔ د يوازې اړيکې او ګراف ترمنځ توپير
يوازې دا دے چې ګراف د څوکو سټ ټاکي؛ يعنې ګراف ګوښې څوکې هم
لرلے شي۔ خو مهمه خبره دا ده چې په سټ \LR{$X$} باندې هره اړيکه \LR{$R$} د
يو لوری لرونکي ګراف \LR{$\tuple{X, R}$} په توګه کتلے شو۔ برعکس، يو
لوری لرونکے ګراف \LR{$\tuple{V, E}$} د اړيکې \LR{$E \subseteq V^2$} په
توګه کتلے شو چې سټ \LR{$V$} پکښې په څرګنده ټاکل شوے وي۔
\end{explain}
\olpseganchor{OLP-0017-B007}{end}

\olpseganchor{OLP-0017-B008}{start}
\begin{ex}
هغه ګراف \LR{$\tuple{V, E}$} چې
\LR{$V = \{1, 2, 3, 4\}$} او
\LR{$E = \{\tuple{1,1}, \allowbreak \tuple{1, 2}, \allowbreak \tuple{1, 3},
\allowbreak \tuple{2, 3}\}$} وي، داسې ښکاري:
\begin{align*}
& \begin{tikzpicture}[->,node distance=2cm]
  \node[draw,circle] (A) {\LR{$1$}};
  \node[draw,circle] (B) [right of=A] {\LR{$2$}};
  \node[draw,circle] (C) [below of=B] {\LR{$3$}};
  \node[draw,circle] (D) [right of=B] {\LR{$4$}};
  \draw (A) to [loop above]  (A);
  \draw (A) to  (B);
  \draw (A) to  (C);
  \draw (B) to  (C);
  \end{tikzpicture}
\intertext{\RL{دا له هغه ګراف \LR{$\tuple{V', E}$} څخه بېل دے چې
  \LR{$V' = \{1, 2, 3\}$} وي، او داسې ښکاري:}}
& \begin{tikzpicture}[->,node distance=2cm]
  \node[draw,circle] (A) {\LR{$1$}};
  \node[draw,circle] (B) [right of=A] {\LR{$2$}};
  \node[draw,circle] (C) [below of=B] {\LR{$3$}};
  \draw (A) to [loop above]  (A);
  \draw (A) to  (B);
  \draw (A) to  (C);
  \draw (B) to  (C);
  \end{tikzpicture}
\end{align*}
\end{ex}
\olpseganchor{OLP-0017-B008}{end}

\olpseganchor{OLP-0017-B009}{start}
\begin{prob}
  په سټ \LR{$\{1, 2, 3, 4\}$} باندې د کوچني يا برابر کېدو اړيکه
  \LR{$\le$} د ګراف په توګه وګورئ او ورسره سم شکل رسم کړئ۔
\end{prob}
\olpseganchor{OLP-0017-B009}{end}

\olpseganchor{OLP-0018-B004}{start}
\olfileid{sfr}{rel}{tre}
\olsection{ونې}
\olpseganchor{OLP-0018-B004}{end}

\olpseganchor{OLP-0018-B005}{start}
د جزوي ترتيب يو خاص ډول چې د منطق په ټولو برخو کښې مهم رول لري،
\emph{ونه} ده۔ متناهي ونې د منطق په لومړنيو برخو کښې راځي:
مثلاً، فارمولونه په نحوي ونه کښې د خپلو برخو په بېلولو پوهېدلے
شو، او په ډېرو اشتقاق نظامونو کښې اشتقاقونه هم د
متناهي ونو بڼه خپلوي۔
%
نامتناهي ونې د بياني منطق او د لومړۍ درجې منطق د بشپړتيا د قضيو
په ثبوت کښې لا راڅرګندېږي، او په ټول رياضي منطق کښې کارول کېږي۔
\olpseganchor{OLP-0018-B005}{end}

\olpseganchor{OLP-0018-B006}{start}
د سټونو د نظريې د ونې مفهوم د ګرافونو په نظريه کښې د ونې له
مفهوم سره نژدې تړاو لري۔ دا د يوې (متناهي) ونې انځور دے:
\olpseganchor{OLP-0018-B006}{end}

\olpseganchor{OLP-0018-B007}{start}
\begin{center}
\begin{tikzpicture}[nodes={draw, circle}, -]
\node{\LR{$r$}} [grow'=up]
    child { node {\LR{$a$}}
        child { node {\LR{$c$}} }
        child { node {\LR{$d$}} }
        child { node {\LR{$e$}} }
    }
    child { node {\LR{$b$}} };
\end{tikzpicture}
\end{center}
\olpseganchor{OLP-0018-B007}{end}

\olpseganchor{OLP-0018-B008}{start}
تر ټولو لاندې غوټه \LR{$r$} ريښه ده۔ له \LR{$r$} پرته د هرې غوټې سمدستي
لاندې پوره يوه مور غوټه شته۔ که له غوټې \LR{$x$} څخه د څنډو په تعقيب
پورته تلو سره غوټې \LR{$y$} ته رسېدلے شو، نو د \LR{$x$} او \LR{$y$} ترمنځ دا
اړيکه داسې ګڼلے شو چې \LR{$x$} د \LR{$y$} يو \emph{نيکۀ} دے۔
\olpseganchor{OLP-0018-B008}{end}

\olpseganchor{OLP-0018-B009}{start}
په يوه ونه کښې د نيکۀ اړيکه سخت جزوي ترتيب دے۔ دا د سټونو د
نظريې تعريف ته لاره هواروي۔ د هغه د بيانولو دپاره دوه مفهومونو ته
اړتيا لرو۔ په سټ \LR{$A$} کښې، چې د \LR{$\le$} په ذريعه جزوي مرتب شوے وي،
\emph{تر ټولو کوچنے غړے} هغه غړے \LR{$x \in A$} دے چې د
هر \LR{$y \in A$} دپاره \LR{$x \le y$} ولرو۔ يو سټ د \LR{$\le$} په ذريعه
\emph{ښۀ مرتب} دے که هر ناتش فرعي سټ ئې تر ټولو کوچنے غړے ولري۔
\olpseganchor{OLP-0018-B009}{end}

\olpseganchor{OLP-0018-B010}{start}
\begin{defn}[ونه]
\emph{ونه} داسې جوړه \LR{$T = \tuple{A, \le}$} ده چې \LR{$A$} يو سټ وي،
او \LR{$\le$} په \LR{$A$} باندې داسې جزوي ترتيب وي چې يوازينے تر ټولو
کوچنے غړے \LR{$r \in A$} ولري (چې \emph{ريښه} بللے شي)، او د هر
\LR{$x \in A$} دپاره سټ \LR{$\Setabs{y}{y \le x}$} د \LR{$\le$} په ذريعه
ښۀ مرتب وي۔
\end{defn}
\olpseganchor{OLP-0018-B010}{end}

\olpseganchor{OLP-0018-B011}{start}
\begin{defn}[ځايناستي]
فرض کړئ \LR{$T = \tuple{A, \le}$} يوه ونه ده۔ که \LR{$x,y \in A$}،
\LR{$x < y$}، او داسې \LR{$z \in A$} نۀ وي چې \LR{$x < z < y$} وي، نو وايو
چې \LR{$y$} د \LR{$x$} \emph{ځايناستے} دے۔
\end{defn}
\olpseganchor{OLP-0018-B011}{end}

\olpseganchor{OLP-0018-B012}{start}
د \LR{$x \in A$} ځايناستي د هغې \emph{اولادونه} هم بللے شي۔ که \LR{$y$}
د \LR{$x$} ځايناستے وي، نو \LR{$x$} د \LR{$y$} \emph{مخکينے} يا \emph{مور}
بولو۔
\olpseganchor{OLP-0018-B012}{end}

\olpseganchor{OLP-0018-B013}{start}
\begin{prop}
که \LR{$\tuple{A,\le}$} يوه ونه وي، نو له ريښې پرته د هر \LR{$x \in A$}
له يوه څخه زيات مخکيني نۀ شي کېدلے۔
\end{prop}
\olpseganchor{OLP-0018-B013}{end}

\olpseganchor{OLP-0018-B014}{start}
\begin{proof}
  فرض کړئ \LR{$y_1 < x$}، \LR{$y_2 < x$} او \LR{$y_1 \neq y_2$}۔ نو
  \LR{$\{y_1, y_2\} \subseteq \Setabs{z}{z<x}$}۔ ځکه چې
  \LR{$\Setabs{z}{z<x}$} د \LR{$\le$} په ذريعه ښۀ مرتب دے، د دې فرعي سټ
  \LR{$\{y_1, y_2\}$} تر ټولو کوچنے غړے لري، چې ښکاره ده يا به \LR{$y_1$}
  وي يا \LR{$y_2$}۔ نو يا \LR{$y_1 \le y_2$} يا \LR{$y_2 \le y_1$}۔ فرض مو
  کړے و چې \LR{$y_1 \neq y_2$}، نو په حقيقت کښې يا \LR{$y_1 < y_2$} يا
  \LR{$y_2 < y_1$}۔ ځکه چې \LR{$y_1 < x$} او \LR{$y_2 < x$} مو هم فرض کړي وو،
  دا هم لرو چې يا \LR{$y_1 < y_2 < x$} يا \LR{$y_2 < y_1 < x$}۔ نو
  \LR{$y_1$} او \LR{$y_2$} دواړه د \LR{$x$} مخکيني نۀ شي کېدلے۔
\end{proof}
\olpseganchor{OLP-0018-B014}{end}

\olpseganchor{OLP-0018-B015}{start}
\begin{defn}
ونه \LR{$T = \tuple{A, \le}$} \emph{نامتناهي} بللے شي که \LR{$A$} نامتناهي
سټ وي، او که نه نو \emph{متناهي} ده۔ که په \LR{$T$} کښې هر \LR{$x \in A$}
يوازې متناهي شمېر ځايناستي ولري، نو وايو چې \LR{$T$}
\emph{متناهي څانګېدونکې} ده۔
\end{defn}
\olpseganchor{OLP-0018-B015}{end}

\olpseganchor{OLP-0018-B016}{start}
\begin{defn}[څانګې]
که ونه \LR{$T = \tuple{A, \le}$} راکړل شي، نو د \LR{$T$} يوه \emph{څانګه}
په \LR{$T$} کښې يوه اعظمي ځنځير ده؛ يعنې داسې سټ \LR{$B \subseteq A$} چې
د هر \LR{$x, y \in B$} دپاره يا \LR{$x \le y$} يا \LR{$y \le x$} وي، او د
هر \LR{$z \in X \setminus B$} دپاره داسې \LR{$u \in B$} شته چې نۀ
\LR{$z \le u$} وي او نۀ \LR{$u \le z$}۔
%
د \LR{$T$} د ټولو څانګو د سټ د ښودلو دپاره \LR{$[T]$} کاروو۔
\end{defn}
\olpseganchor{OLP-0018-B016}{end}

\olpseganchor{OLP-0018-B017}{start}
\begin{ex}
د متناهي څانګېدونکې ونې يوه مشهوره بېلګه د \LR{$0$} ګانو او \LR{$1$} ګانو
د متناهي لړيو \emph{نامتناهي دوه ګونې ونه} ده۔ کله کله په
\LR{$\{0,1\}^*$} يا \LR{$\Bin^*$} ښودل کېږي، او د غځونې د اړيکې
\LR{$\sqsubseteq$} په ذريعه مرتب شوې ده (مثلاً، \LR{$101 \sqsubseteq 101101$})۔
ځکه چې هر دوه ګونے توريز تسلسل په پاے کښې د \LR{$0$} يا \LR{$1$} په
ورزياتولو تل غځولے شو، دا ونه نامتناهي ډېر غړي لري: هر غړے \LR{$s$}
پوره دوه ځايناستي لري، \LR{$s0$} او \LR{$s1$}۔ ريښه ئې تشه لړۍ \LR{$\emptyseq$} ده۔
\end{ex}
\olpseganchor{OLP-0018-B017}{end}

\olpseganchor{OLP-0018-B018}{start}
\begin{ex}
په يو څه زياته عمومي توګه، د طبيعي عددونو د متناهي لړيو سټ
\LR{$\Nat^*$} د غځونې له اړيکې \LR{$\sqsubseteq$} سره هم يوه ونه ده۔
ښکاره ده چې دا متناهي څانګېدونکې نۀ ده: هر \LR{$s \in \Nat^*$}
نامتناهي ډېر ځايناستي \LR{$sn$} لري، د هر \LR{$n \in \Nat$} دپاره يو۔
هر \LR{$A \subseteq \Nat^*$} چې د \LR{$\sqsubseteq$} لاندې تړلے وي، د
\LR{$\Nat^*$} يوه \emph{فرعي ونه} ده۔ (يعنې \LR{$A$} داسې دے چې که
\LR{$s \in A$} او \LR{$s' \sqsubseteq s$} وي، نو \LR{$s' \in A$} هم وي۔)
ټولې متناهي ونې د \LR{$\Nat^*$} د متناهي فرعي ونو په توګه ښودل کېدلے شي۔
\end{ex}
\olpseganchor{OLP-0018-B018}{end}

\olpseganchor{OLP-0018-B019}{start}
\begin{prop}[د کونيګ مرستندويه قضيه]
که \LR{$T = \tuple{A,\le}$} متناهي څانګېدونکې نامتناهي ونه وي، نو
\LR{$T$} يوه نامتناهي څانګه لري۔
\end{prop}
\olpseganchor{OLP-0018-B019}{end}

\olpseganchor{OLP-0018-B020}{start}
د کونيګ د مرستندويې قضيې يوه خاصه بڼه چې د محاسبه کېدنې په نظريه
کښې ډېره کارېږي، \emph{د کونيګ کمزورې مرستندويه قضيه} بللے شي:
د \LR{$\{0,1\}^*$} هره نامتناهي فرعي ونه يوه نامتناهي څانګه لري۔
\olpseganchor{OLP-0018-B020}{end}

\olpseganchor{OLP-0019-B004}{start}
\olfileid{sfr}{rel}{ops}
\olsection{په اړيکو باندې عمليات}
\olpseganchor{OLP-0019-B004}{end}

\olpseganchor{OLP-0019-B005}{start}
اړيکې بدلول يا سره يوځاے کول ډېر وخت ګټور وي۔ په
\olref[sfr][rel][ord]{prop:stricttopartial} کښې مو د اړيکو
\emph{اتحاد} په پام کښې نيولے و؛ دا يوازې د دوو اړيکو اتحاد دے
چې د جوړو د سټونو په توګه ګڼل شوې وي۔ همدارنګه، په
\olref[sfr][rel][ord]{prop:partialtostrict} کښې مو د اړيکو نسبي
تفاضل په پام کښې نيولے و۔ لاندې ځينې نور عمليات دي چې په اړيکو
باندې ئې کولے شو۔
\olpseganchor{OLP-0019-B005}{end}

\olpseganchor{OLP-0019-B006}{start}
\begin{defn}\ollabel{relationoperations}
فرض کړئ \LR{$R$} او \LR{$S$} اړيکې دي، او \LR{$A$} هر يو سټ دے۔
\olpseganchor{OLP-0019-B006}{end}

\olpseganchor{OLP-0019-B007}{start}
د \LR{$R$} \emph{معکوس} دا دے:
\LR{$R^{-1} = \Setabs{\tuple{y, x}}{\tuple{x, y} \in R}$}۔
\olpseganchor{OLP-0019-B007}{end}

\olpseganchor{OLP-0019-B008}{start}
د \LR{$R$} او \LR{$S$} \emph{نسبي حاصل ضرب} دا دے:
\LR{$(R \mid S) = \{\tuple{x, z} : \exists y(Rxy \land Syz)\}$}۔
\olpseganchor{OLP-0019-B008}{end}

\olpseganchor{OLP-0019-B009}{start}
تر \LR{$A$} پورې د \LR{$R$} \emph{محدودونه} دا ده:
\LR{$\funrestrictionto{R}{A}= R \cap A^2$}۔
\olpseganchor{OLP-0019-B009}{end}

\olpseganchor{OLP-0019-B010}{start}
په \LR{$A$} باندې د \LR{$R$} \emph{تطبيق} دا دے:
\LR{$\funimage{R}{A} = \{y : (\exists x \in A)Rxy\}$}
\end{defn}
\olpseganchor{OLP-0019-B010}{end}

\olpseganchor{OLP-0019-B011}{start}
\begin{ex}
فرض کړئ \LR{$S \subseteq \Int^2$} په \LR{$\Int$} باندې د ځايناستي اړيکه
ده؛ يعنې \LR{$S = \Setabs{\tuple{x, y} \in \Int^2}{x + 1 = y}$}،
نو \LR{$Sxy$} هله او يوازې هله چې \LR{$x + 1 = y$} وي۔
\olpseganchor{OLP-0019-B011}{end}

\olpseganchor{OLP-0019-B012}{start}
\LR{$S^{-1}$} په \LR{$\Int$} باندې د مخکيني اړيکه ده، يعنې:
\LR{$\Setabs{\tuple{x,y}\in\Int^2}{x -1 =y}$}۔
\olpseganchor{OLP-0019-B012}{end}

\olpseganchor{OLP-0019-B013}{start}
\LR{$S\mid S$} دا دے:
\LR{$ \Setabs{\tuple{x,y}\in\Int^2}{x + 2 =y}$}
\olpseganchor{OLP-0019-B013}{end}

\olpseganchor{OLP-0019-B014}{start}
\LR{$\funrestrictionto{S}{\Nat}$} په \LR{$\Nat$} باندې د ځايناستي اړيکه ده۔
\olpseganchor{OLP-0019-B014}{end}

\olpseganchor{OLP-0019-B015}{start}
\LR{$\funimage{S}{\{1,2,3\}}$} دا دے: \LR{$\{2, 3, 4\}$}۔
\end{ex}
\olpseganchor{OLP-0019-B015}{end}

\olpseganchor{OLP-0019-B016}{start}
\begin{defn}[انتقالي بندښت]فرض کړئ \LR{$R \subseteq A^2$} يوه دوه ځاييزه اړيکه ده۔
\olpseganchor{OLP-0019-B016}{end}

\olpseganchor{OLP-0019-B017}{start}
د \LR{$R$} \emph{انتقالي بندښت} دا دے:
\LR{$R^+ = \bigcup_{0 < n \in \Nat} R^n$}؛ دلته په بازګشتي ډول
\LR{$R^1 = R$} او \LR{$R^{n+1} = R^n \mid R$} تعريفوو۔
\olpseganchor{OLP-0019-B017}{end}

\olpseganchor{OLP-0019-B018}{start}
د \LR{$R$} \emph{انعکاسي انتقالي بندښت} دا دے:
\LR{$R^* = R^+ \cup \Id{A}$}۔
\end{defn}
\olpseganchor{OLP-0019-B018}{end}

\olpseganchor{OLP-0019-B019}{start}
\begin{ex}
د ځايناستي اړيکه \LR{$S \subseteq \Int^2$} واخلئ۔ \LR{$S^2xy$} هله او
يوازې هله چې \LR{$x + 2 = y$} وي، \LR{$S^3xy$} هله او يوازې هله چې
\LR{$x + 3 = y$} وي، او داسې نور۔ نو \LR{$S^+xy$} هله او يوازې هله چې
د کوم \LR{$n \geq 1$} دپاره \LR{$x + n = y$} وي۔ په نورو ټکو، \LR{$S^+xy$}
هله او يوازې هله چې \LR{$x < y$} وي، او \LR{$S^*xy$} هله او يوازې هله چې
\LR{$x \le y$} وي۔
\end{ex}
\olpseganchor{OLP-0019-B019}{end}

\olpseganchor{OLP-0019-B020}{start}
\begin{prob}
وښيئ چې د \LR{$R$} انتقالي بندښت په رښتيا انتقالي دے۔
\end{prob}
\olpseganchor{OLP-0019-B020}{end}

\olpseganchor{OLP-0020-B004}{start}
\olchapter{sfr}{fun}{تابعې}
\olpseganchor{OLP-0020-B004}{end}











\olpseganchor{OLP-0020-B010}{start}
%
\olpseganchor{OLP-0020-B010}{end}

\olpseganchor{OLP-0021-B004}{start}
\olfileid{sfr}{fun}{bas}
\olsection{بنسټونه}
\olpseganchor{OLP-0021-B004}{end}

\olpseganchor{OLP-0021-B005}{start}
\begin{explain}
\emph{تابع} داسې نگاشت دے چې د يوه ورکړل شوي سټ هر غړے
د کوم ورکړل شوي (بل) سټ له يوه ټاکلي غړي سره تړي۔
د بېلګې په توګه، د \LR{$1$} ورزياتولو عمل يوه تابع ټاکي: هر عدد~\LR{$n$}
له يوه يوازيني عدد~\LR{$n+1$} سره تړل کېږي۔
\olpseganchor{OLP-0021-B005}{end}
  
\olpseganchor{OLP-0021-B006}{start}
په عمومي ډول، تابعې جوړې، درې ګونې او داسې نور د ننوتونو په توګه
اخيستلے شي او يو ډول وت بېرته ورکوي۔ له بنسټيز حسابه ډېرې تابعې
پېژنو۔ د بېلګې په توګه، جمع او ضرب تابعې دي۔ دوې دوه عددونه
اخلي او درېيم عدد بېرته ورکوي۔
\olpseganchor{OLP-0021-B006}{end}

\olpseganchor{OLP-0021-B007}{start}
په دې رياضيکي، تجريدي معنٰی، تابع يو \emph{تور بکس} دے:
يوازې دا خبره اهميت لري چې کوم وت له کوم ننوت سره تړل شوے دے،
نۀ دا چې وت په کومه طريقه حسابېږي۔
\end{explain}
\olpseganchor{OLP-0021-B007}{end}

\olpseganchor{OLP-0021-B008}{start}
\begin{defn}[تابع]
\emph{تابع} \LR{$f \colon A \to B$} د~\LR{$A$} د هر غړي
د~\LR{$B$} له يوه غړي سره تړون دے۔
\olpseganchor{OLP-0021-B008}{end}

\olpseganchor{OLP-0021-B009}{start}
\LR{$A$} د~\LR{$f$} \emph{د تعريف ساحه} او \LR{$B$} د~\LR{$f$} \emph{هدف سټ} بولو۔
د~\LR{$A$} غړي د~\LR{$f$} ننوتونه يا \emph{ارګومېنټونه} بلل کېږي،
او د~\LR{$B$} هغه غړے چې~\LR{$f$} ئې له ارګومېنټ~\LR{$x$} سره تړي،
د ارګومېنټ~\LR{$x$} دپاره \emph{د~\LR{$f$} قيمت} بلل کېږي او په~\LR{$f(x)$} ليکل کېږي۔
\olpseganchor{OLP-0021-B009}{end}

\olpseganchor{OLP-0021-B010}{start}
د~\LR{$f$} \emph{د قيمتونو سټ} \LR{$\ran{f}$} د هدف سټ هغه فرعي سټ دے چې
د~\LR{$f$} له هغو قيمتونو جوړ دے چې د کوم ارګومېنټ دپاره ترلاسه کېږي؛
\LR{$\ran{f} = \Setabs{f(x)}{x \in A}$}۔
\end{defn}
\olpseganchor{OLP-0021-B010}{end}

\olpseganchor{OLP-0021-B011}{start}
په \olref{fig:function} کښې انځوريز شکل د تابعو په پوهېدو کښې مرسته
کولے شي۔ کيڼې خواته بيضوي شکل د تابعې \emph{د تعريف ساحه} ښيي؛
ښي خوا ته بيضوي شکل د تابعې \emph{هدف سټ} ښيي؛ او غشے د تعريف
په ساحه کښې له يوه \emph{ارګومېنټ} څخه په هدف سټ کښې ورته
\emph{قيمت} ته ځي۔
\olpseganchor{OLP-0021-B011}{end}

\olpseganchor{OLP-0021-B012}{start}
\begin{figure}
  \olasset{../upstream/assets/diagrams/function.tikz}
  \caption{تابع د يوه سټ د هر غړي د بل سټ له
    يوه غړي سره تړون دے۔ غشے د تعريف په ساحه کښې
    له يوه ارګومېنټ څخه په هدف سټ کښې ورته قيمت ته ځي۔}
  \ollabel{fig:function}
\end{figure}
\olpseganchor{OLP-0021-B012}{end}

\olpseganchor{OLP-0021-B013}{start}
\begin{ex}
ضرب د طبيعي عددونو جوړې د ننوتونو په توګه اخلي او طبيعي عددونه
د وتونو په توګه ورکوي؛ نو له \LR{$\Nat \times \Nat$} (د تعريف له ساحې)
څخه \LR{$\Nat$} (هدف سټ) ته ځي۔ د قيمتونو سټ هم \LR{$\Nat$} دے، ځکه چې
هر \LR{$n \in \Nat$} د \LR{$n \times 1$} په توګه ترلاسه کېږي۔
\end{ex}
\olpseganchor{OLP-0021-B013}{end}

\olpseganchor{OLP-0021-B014}{start}
\begin{ex}
ضرب تابع ده، ځکه چې هر ننوت---د طبيعي عددونو هره جوړه---له يوه
يوازيني وت سره تړي: \LR{$\times \colon \Nat^2 \to \Nat$}۔ خو له يوه
طبيعي عدد~\LR{$n$} سره داسې حقيقي عدد~\LR{$x$} تړل چې \LR{$x^2 = n$} وي،
تابع نۀ جوړوي، ځکه چې هر مثبت صحيح عدد~\LR{$n$} دوه مربع جذرونه لري:
\LR{$\sqrt{n}$} او~\LR{$-\sqrt{n}$}۔ که يوازې غير منفي (اصلي) مربع جذر بېرته ورکړو،
تابع ترې جوړولے شو: \LR{$\sqrt{\phantom{X}} \colon \Nat \to \Real$}۔
\end{ex}
\olpseganchor{OLP-0021-B014}{end}

\olpseganchor{OLP-0021-B015}{start}
\begin{ex}
هغه اړيکه چې د يوه ټولګي هر زده کوونکے د هغه له وروستۍ درجې سره
تړي، تابع ده---يو زده کوونکے په هماغه ټولګي کښې دوه بېلې وروستۍ
درجې نۀ شي اخيستلے۔ هغه اړيکه چې د يوه ټولګي هر زده کوونکے
د هغه له والدينو سره تړي، تابع نۀ ده: زده کوونکي کېدے شي هېڅ،
يا دوه، يا تر دوو زيات والدين ولري۔
\end{ex}
\olpseganchor{OLP-0021-B015}{end}

\olpseganchor{OLP-0021-B016}{start}
\begin{explain}
تابعې داسې تعريفولے شو چې په يوه کره طريقه د هر ممکن ارګومېنټ
دپاره د تابعې قيمت وټاکو۔ د دې بېلې لارې دا دي چې يو فورمول
ورکړو، د قيمت د حسابولو طريقه بيان کړو، يا د هر ارګومېنټ قيمت
په فهرست کښې وليکو۔ تابعې چې په هره طريقه تعريفوو، بايد ډاډمن
شو چې د هر ارګومېنټ دپاره يو، او يوازې يو، قيمت ټاکو۔
\end{explain}
\olpseganchor{OLP-0021-B016}{end}


\olpseganchor{OLP-0021-B017}{start}
\begin{ex}
فرض کړئ \LR{$f \colon \Nat \to \Nat$} داسې تعريف شوې ده چې \LR{$f(x) = x+1$}
وي۔ دا تعريف~\LR{$f$} د داسې تابعې په توګه ټاکي چې طبيعي عددونه اخلي
او طبيعي عددونه ورکوي۔ تعريف وايي چې د ورکړل شوي طبيعي عدد~\LR{$x$}
دپاره به~\LR{$f$} د هغه ځايناستے~\LR{$x+1$} ورکړي۔ په دې حالت کښې هدف
سټ \LR{$\Nat$} د~\LR{$f$} د قيمتونو سټ نۀ دے، ځکه چې طبيعي عدد~\LR{$0$} د
هېڅ طبيعي عدد ځايناستے نۀ دے۔ د~\LR{$f$} د قيمتونو سټ د ټولو مثبتو
صحيح عددونو سټ، \LR{$\Int^{+}$}، دے۔
\end{ex}
\olpseganchor{OLP-0021-B017}{end}

\olpseganchor{OLP-0021-B018}{start}
\begin{ex}\ollabel{examplefunext}
فرض کړئ \LR{$g \colon \Nat \to \Nat$} داسې تعريف شوې ده چې \LR{$g(x) = x+2-1$}
وي۔ دا وايي چې~\LR{$g$} يوه تابع ده چې طبيعي عددونه اخلي او طبيعي
عددونه ورکوي۔ د ورکړل شوي طبيعي عدد دپاره به~\LR{$g$} د~\LR{$x$} د
ځايناستي د ځايناستي مخکينے عدد، يعنې \LR{$x+1$}، ورکړي۔
\end{ex}
\olpseganchor{OLP-0021-B018}{end}

\olpseganchor{OLP-0021-B019}{start}
\begin{explain}
اوس مو دوه تابعې، \LR{$f$} او \LR{$g$}، له بېلو \emph{تعريفونو} سره په پام
کښې ونيولې۔ خو دا \emph{هماغه يوه تابع} ده۔ ځکه چې د هر طبيعي
عدد~\LR{$n$} دپاره \LR{$f(n) = n+1 = n+2-1 = g(n)$} لرو۔ په بله وينا:
د \LR{$f$} او~\LR{$g$} تعريفونه مو په بېلو معادلو هماغه يو نگاشت ټاکي۔ نو
په ضمني ډول د تابعو د قيمتونو له مخې د برابرۍ په يوه اصل تکيه کوو:
\[
  \text{\RL{که }}\forall x\, f(x) = g(x)\text{\RL{، نو }}f = g
\]
په دې شرط چې د \LR{$f$} او~\LR{$g$} د تعريف ساحه او هدف سټ يو شان وي۔
\end{explain}
\olpseganchor{OLP-0021-B019}{end}

\olpseganchor{OLP-0021-B020}{start}
\begin{ex}
تابعې د حالتونو په وېشلو هم تعريفولے شو۔ د بېلګې په توګه،
\LR{$h \colon \Nat \to \Nat$} داسې تعريفولے شو:
\[
h(x) =
\begin{cases}
  \frac{x}{2} & \text{\RL{که \LR{$x$} جفت وي}} \\
  \frac{x+1}{2} & \text{\RL{که \LR{$x$} طاق وي۔}}
\end{cases}
\]
ځکه چې هر طبيعي عدد يا جفت دے يا طاق، د دې تابعې وت به تل
طبيعي عدد وي۔ يوازې دا په ياد وساتئ چې که تابع د حالتونو په
وېشلو تعريفوئ، هر ممکن ننوت بايد په دقيقه توګه په يوه حالت کښې
راشي۔ ځينې وختونه د دې ثبوت ته اړتيا وي چې حالتونه ټول امکانات
رانغاړي او يو له بل سره ګډ نۀ دي۔
\end{ex}
\olpseganchor{OLP-0021-B020}{end}

\olpseganchor{OLP-0022-B004}{start}
\olfileid{sfr}{fun}{kin}
\olsection{د تابعو ډولونه}
\olpseganchor{OLP-0022-B004}{end}

\olpseganchor{OLP-0022-B005}{start}
\begin{explain}
د تابعو د هغو ځينو ډولونو دپاره چې ډېر ورسره مخ کېږو، د يو ډول
ډلبندۍ پېژندل به ګټور وي۔
\olpseganchor{OLP-0022-B005}{end}

\olpseganchor{OLP-0022-B006}{start}
په پيل کښې داسې تابعې په پام کښې نيول غواړو چې د هدف سټ هر
غړے د تابعې يو قيمت وي۔ دې تابعو ته پر هدف سټ بشپړه ويل کېږي،
او د \olref{fig:surjective} په شان ئې ښودلے شو۔
\olpseganchor{OLP-0022-B006}{end}

\olpseganchor{OLP-0022-B007}{start}
\begin{figure}
  \olasset{../upstream/assets/diagrams/surjective.tikz}
  \caption{پر هدف سټ بشپړه تابع د هدف سټ هر غړے د قيمت په توګه لري۔}
  \ollabel{fig:surjective}
\end{figure}
\end{explain}
\olpseganchor{OLP-0022-B007}{end}

\olpseganchor{OLP-0022-B008}{start}
\begin{defn}[پر هدف سټ بشپړه تابع]
تابع \LR{$f \colon A \rightarrow B$} هله او يوازې هله \emph{پر هدف سټ بشپړه}
ده چې \LR{$B$} د~\LR{$f$} د قيمتونو سټ هم وي؛ يعنې د هر \LR{$y \in B$}
دپاره لږ تر لږه يو \LR{$x \in A$} وي چې~\LR{$f(x) = y$} وي، يا په نښو:
\[
  (\forall y \in B)(\exists x \in A)f(x) = y.
\]
داسې تابع له \LR{$A$} څخه \LR{$B$} ته پر هدف سټ بشپړه تابع بولو۔
\end{defn}
\olpseganchor{OLP-0022-B008}{end}

\olpseganchor{OLP-0022-B009}{start}
\begin{explain}
که ښودل غواړئ چې \LR{$f$} يوه پر هدف سټ بشپړه تابع ده، نو بايد وښيئ چې د~\LR{$f$}
د هدف سټ هر څيز د کوم ننوت \LR{$x$} دپاره د \LR{$f(x)$} قيمت دے۔
\olpseganchor{OLP-0022-B009}{end}

\olpseganchor{OLP-0022-B010}{start}
پام وکړئ چې هره تابع يوه پر هدف سټ بشپړه تابع \emph{رامنځته کوي}۔ ځکه،
د ورکړل شوې تابعې \LR{$f \colon A \to B$} دپاره، فرض کړئ
\LR{$f' \colon A \to \ran{f}$} په \LR{$f'(x) = f(x)$} تعريف شوې ده۔ ځکه
چې \LR{$\ran{f}$} د \LR{$\Setabs{f(x) \in B}{x \in A}$} په توګه \emph{تعريف}
شوے دے، دا تضمين شوے دے چې تابع \LR{$f'$} به پر هدف سټ بشپړه تابع وي۔
\end{explain}
\olpseganchor{OLP-0022-B010}{end}

\olpseganchor{OLP-0022-B011}{start}
\begin{explain}
اوس، هره تابع هر ممکن ننوت له يوه يوازيني وت سره تړي۔ خو داسې
تابعې هم شته چې بېل ننوتونه هېڅکله له يو شان وت سره نۀ تړي۔ دې
تابعو ته يو پر يو ويل کېږي، او د \olref{fig:injective} په شان
ئې ښودلے شو۔
\begin{figure}
  \olasset{../upstream/assets/diagrams/injective.tikz}
  \caption{يو پر يو تابع هېڅکله دوه بېل ارګومېنټونه له يوه قيمت سره نۀ تړي۔}
  \ollabel{fig:injective}
\end{figure}
\end{explain}
\olpseganchor{OLP-0022-B011}{end}

\olpseganchor{OLP-0022-B012}{start}
\begin{defn}[يو پر يو تابع]
تابع \LR{$f \colon A \rightarrow B$} هله او يوازې هله \emph{يو پر يو}
ده چې د هر \LR{$y \in B$} دپاره تر يوه زيات داسې \LR{$x \in A$} نۀ وي
چې~\LR{$f(x) = y$} وي۔ داسې تابع له \LR{$A$} څخه~\LR{$B$} ته يو پر يو تابع بولو۔
\end{defn}
\olpseganchor{OLP-0022-B012}{end}

\olpseganchor{OLP-0022-B013}{start}
\begin{explain}
که ښودل غواړئ چې \LR{$f$} يوه يو پر يو تابع ده، بايد وښيئ چې د~\LR{$f$}
د تعريف د ساحې د هر دوو غړو \LR{$x$} او \LR{$y$} دپاره،
که \LR{$f(x)=f(y)$} وي، نو \LR{$x=y$} وي۔
\end{explain}
\olpseganchor{OLP-0022-B013}{end}

\olpseganchor{OLP-0022-B014}{start}
\begin{ex}
ثابته تابع \LR{$f\colon \Nat \to \Nat$} چې په \LR{$f(x) = 1$} ورکړل شوې،
نۀ يو پر يو ده او نۀ پر هدف سټ بشپړه۔
\olpseganchor{OLP-0022-B014}{end}

\olpseganchor{OLP-0022-B015}{start}
د عينيت تابع \LR{$f\colon \Nat \to \Nat$} چې په \LR{$f(x) = x$} ورکړل شوې،
هم يو پر يو ده او هم پر هدف سټ بشپړه۔
\olpseganchor{OLP-0022-B015}{end}

\olpseganchor{OLP-0022-B016}{start}
د ځايناستي تابع \LR{$f \colon \Nat \to \Nat$} چې په \LR{$f(x) = x+1$} ورکړل
شوې، يو پر يو ده خو پر هدف سټ بشپړه نۀ ده۔
\olpseganchor{OLP-0022-B016}{end}
  
\olpseganchor{OLP-0022-B017}{start}
تابع \LR{$f \colon \Nat \to \Nat$} چې داسې تعريف شوې ده:
\[
  f(x) =
  \begin{cases}
    \frac{x}{2} & \text{\RL{که \LR{$x$} جفت وي}} \\
    \frac{x+1}{2} & \text{\RL{که \LR{$x$} طاق وي۔}}
  \end{cases}
\]
پر هدف سټ بشپړه ده، خو يو پر يو نۀ ده۔
\end{ex}
\olpseganchor{OLP-0022-B017}{end}

\olpseganchor{OLP-0022-B018}{start}
\begin{explain}
ډېر وخت داسې تابعې په پام کښې نيول غواړو چې هم يو پر يو وي
او هم پر هدف سټ بشپړه۔ دا تابعې دوه اړخيزه يو پر يو بولو۔ دوې په
\olref{fig:bijective} کښې ښودل شوې تابعې ته ورته دي۔ دوه اړخيزې يو پر يو تابعې
کله کله \emph{يو پر يو مطابقتونه} هم بلل کېږي، ځکه چې د هدف سټ
غړي د تعريف د ساحې له غړو سره په يوازيني ډول جوړوي۔
\begin{figure}
  \olasset{../upstream/assets/diagrams/bijective.tikz}
  \caption{دوه اړخيزه يو پر يو تابع د هدف سټ غړي د تعريف د ساحې له غړو سره په يوازيني ډول جوړوي۔}
  \ollabel{fig:bijective}
\end{figure}
\end{explain}
\olpseganchor{OLP-0022-B018}{end}

\olpseganchor{OLP-0022-B019}{start}
\begin{defn}[دوه اړخيزه يو پر يو تابع]
تابع \LR{$f \colon A \to B$} هله او يوازې هله \emph{دوه اړخيزه يو پر يو} ده چې
هم پر هدف سټ بشپړه وي او هم يو پر يو۔ داسې تابع له \LR{$A$} څخه~\LR{$B$}
ته (يا د \LR{$A$} او~\LR{$B$} ترمنځ) دوه اړخيزه يو پر يو تابع بولو۔
\end{defn}
\olpseganchor{OLP-0022-B019}{end}

\olpseganchor{OLP-0023-B004}{start}
\olfileid{sfr}{fun}{rel}
\olpseganchor{OLP-0023-B004}{end}

\olpseganchor{OLP-0023-B005}{start}
\olsection{تابعې د اړيکو په توګه}
\olpseganchor{OLP-0023-B005}{end}

\olpseganchor{OLP-0023-B006}{start}
\begin{explain}
هغه تابع چې د~\LR{$A$} غړي د~\LR{$B$} له غړو سره
تړي، په ښکاره د \LR{$A$} او~\LR{$B$} ترمنځ يوه اړيکه ټاکي؛ يعنې هغه اړيکه
چې د \LR{$x$} او \LR{$y$} ترمنځ هله او يوازې هله وي چې \LR{$f(x) = y$} وي۔
په حقيقت کښې، که د بېلګې په توګه د رياضياتو د بنسټيزو توکو
په کمولو کښې دلچسپي لرو، ان تابع~\LR{$f$} له دې اړيکې سره، يعنې
د جوړو له يوه سټ سره، \emph{يوه ګڼلے} شو۔ بيا دا پوښتنه راپورته
کېږي: کومې اړيکې په دې ډول تابعې ټاکي؟
\end{explain}
\olpseganchor{OLP-0023-B006}{end}

\olpseganchor{OLP-0023-B007}{start}
\begin{defn}[د تابعې ګراف]فرض کړئ \LR{$f\colon A \to B$} يوه تابع ده۔
د~\LR{$f$} \emph{ګراف} هغه اړيکه \LR{$R_f \subseteq A \times B$} ده چې داسې
تعريف شوې ده:
\[
R_f = \Setabs{\tuple{x,y}}{f(x) = y}.
\]
\end{defn}
\olpseganchor{OLP-0023-B007}{end}

\olpseganchor{OLP-0023-B008}{start}
\begin{explain}
د غړو له مخې د برابرۍ په اصل، د تابعې ګراف په يوازيني ډول ټاکل
کېږي۔ سربېره پر دې، د سټونو د غړو له مخې د برابرۍ اصل سملاسي
د تابعو د قيمتونو له مخې د برابرۍ ضمني اصل هم توجيه کوي: که د
\LR{$f$} او~\LR{$g$} د تعريف ساحه او هدف سټ يو شان وي او په ټولو قيمتونو
کښې سره موافقې وي، نو هماغه يوه تابع ده۔
\olpseganchor{OLP-0023-B008}{end}

\olpseganchor{OLP-0023-B009}{start}
همدارنګه، که يوه اړيکه «د تابعې خاصيت ولري»، نو د يوې تابعې ګراف ده۔
\end{explain}
\olpseganchor{OLP-0023-B009}{end}

\olpseganchor{OLP-0023-B010}{start}
\begin{prop}\ollabel{prop:graph-function}
فرض کړئ \LR{$R \subseteq A \times B$} داسې ده چې:
\begin{enumerate}
\item که \LR{$Rxy$} او \LR{$Rxz$} وي، نو \LR{$y = z$} وي؛ او
\item د هر \LR{$x \in A$} دپاره داسې \LR{$y \in B$} شته چې \LR{$\tuple{x,
y} \in R$} وي۔
\end{enumerate}
نو \LR{$R$} د هغې تابعې \LR{$f\colon A \to B$} ګراف ده چې په دې ډول
تعريف شوې ده: \LR{$f(x) = y$} هله او يوازې هله چې \LR{$Rxy$} وي۔
\end{prop}
\olpseganchor{OLP-0023-B010}{end}

\olpseganchor{OLP-0023-B011}{start}
\begin{proof}
فرض کړئ داسې \LR{$y$} شته چې \LR{$Rxy$} وي۔ که يو بل \LR{$z \neq y$} هم واے
چې \LR{$Rxz$} واے، پر~\LR{$R$} لګېدلے شرط به مات شوے و۔ نو که داسې \LR{$y$}
شته چې \LR{$Rxy$} وي، دا \LR{$y$} يوازينے دے؛ له دې کبله~\LR{$f$} په سمه توګه
تعريف شوې ده۔ په ښکاره، \LR{$R_f = R$}۔
\end{proof}
\olpseganchor{OLP-0023-B011}{end}

\olpseganchor{OLP-0023-B012}{start}
\begin{explain}
هره تابع \LR{$f\colon A \to B$} ګراف لري؛ يعنې د همدې د تعريف ساحې
او هدف سټ ترمنځ يوه اړيکه، چې د \LR{$A \times B$} فرعي سټ دے او په \LR{$f(x) = y$}
تعريف شوې ده۔ بل خوا، هره اړيکه~\LR{$R
\subseteq A \times B$} چې په \olref{prop:graph-function} کښې ورکړل
شوي خاصيتونه ولري، د يوې تابعې~\LR{$f \colon A \to B$} ګراف ده۔ د
تابعو او د هغو د ګرافونو د دې نږدې تړاو له کبله، تابع يوازې د
هغې د ګراف په توګه ګڼلے شو۔ په بله وينا، تابعې له ځينو اړيکو،
يعنې د مرتبو څوګونو له ځينو سټونو، سره يوې ګڼلے شو۔
خو پام وکړئ چې د دې «يو ګڼلو» مطلب
هماغه دے چې په \olref[sfr][rel][ref]{sec} کښې و: دا د تابعو د
مابعدالطبيعي ماهيت په اړه ادعا نۀ ده، بلکې دا مشاهده ده چې
تابعې د ځينو سټونو په توګه \emph{ګڼل} اسانتيا لري۔ د دې اسانتيا
يو لامل دا دے چې اوس پر تابعو هم هغو ته ورته عمليات په
پام کښې نيولے شو چې پر اړيکو مو وکړل (وګورئ
\olref[sfr][rel][ops]{sec})۔ په ځانګړي ډول:
\end{explain}
\olpseganchor{OLP-0023-B012}{end}

\olpseganchor{OLP-0023-B013}{start}
\begin{defn}\ollabel{defn:funimage}
فرض کړئ \LR{$f \colon A \to B$} يوه تابع ده او \LR{$C\subseteq A$}۔
\olpseganchor{OLP-0023-B013}{end}

\olpseganchor{OLP-0023-B014}{start}
تر~\LR{$C$} پورې د~\LR{$f$} \emph{محدودونه} هغه تابع
\LR{$\funrestrictionto{f}{C}\colon C \to B$} ده چې د ټولو \LR{$x \in C$}
دپاره په \LR{$(\funrestrictionto{f}{C})(x) = f(x)$} تعريف شوې ده۔ په بله
وينا، \LR{$\funrestrictionto{f}{C} = \Setabs{\tuple{x, y} \in R_f}{x \in
C}$}۔
\olpseganchor{OLP-0023-B014}{end}

\olpseganchor{OLP-0023-B015}{start}
په~\LR{$C$} باندې د~\LR{$f$} \emph{تطبيق} دا دے: \LR{$\funimage{f}{C} =
\Setabs{f(x)}{x \in C}$}۔ دې ته د~\LR{$f$} لاندې د~\LR{$C$} \emph{انځور} هم وايو۔
\end{defn}
\olpseganchor{OLP-0023-B015}{end}

\olpseganchor{OLP-0023-B016}{start}
\begin{explain}
له دې تعريفونو څرګندېږي چې د هرې تابعې~\LR{$f$} دپاره
\LR{$\ran{f} = \funimage{f}{\dom{f}}$}۔
په \olref[sfr][rel][ops]{sec} کښې د
تعريفونو او له اړيکو سره د تابعو د يو ګڼلو له مخې، ورته
مفهومونه مو تمه لرله۔ دلته د تابعې محدودونه يوازې د ننوت
مختصات محدودوي؛ د اړيکې مخکښ تعريف دواړه مختصات محدودوي۔
خو دوه نور عمليات---معکوسونه او نسبي
حاصل ضربونه---لږ زيات تفصيل غواړي۔ هغه به په \olref[inv]{sec}
او \olref[cmp]{sec} کښې وړاندې کړو۔
\end{explain}
\olpseganchor{OLP-0023-B016}{end}

\olpseganchor{OLP-0024-B004}{start}
\olfileid{sfr}{fun}{inv}
\olsection{د تابعو معکوسونه}
\olpseganchor{OLP-0024-B004}{end}

\olpseganchor{OLP-0024-B005}{start}
\begin{explain}
تابعې د نگاشتونو په توګه ګڼو۔ نو د تابعو په اړه يوه څرګنده
پوښتنه دا ده چې ايا نگاشت «سرچپه» کېدے شي۔ د بېلګې په توګه،
د ځايناستي تابع \LR{$f(x) = x + 1$} په دې معنٰی سرچپه کېدے شي چې
تابع \LR{$g(y) = y - 1$} د \LR{$f$} کړے کار «بېرته اړوي»۔
\olpseganchor{OLP-0024-B005}{end}

\olpseganchor{OLP-0024-B006}{start}
خو بايد احتياط وکړو۔ که څه هم د~\LR{$g$} تعريف يوه تابع \LR{$\Int \to \Int$}
ټاکي، يوه \emph{تابع} \LR{$\Nat \to \Nat$} نۀ ټاکي، ځکه چې
\LR{$g(0) \notin \Nat$}۔ نو په ساده حالتونو کښې هم دا خبره پوره څرګنده
نۀ ده چې تابع سرچپه کېدے شي که نۀ؛ دا د تعريف په ساحه او هدف
سټ پورې تړلې کېدے شي۔
\olpseganchor{OLP-0024-B006}{end}

\olpseganchor{OLP-0024-B007}{start}
د تابعې د معکوس مفهوم دا خبره لا کره کوي۔
\end{explain}
\olpseganchor{OLP-0024-B007}{end}

\olpseganchor{OLP-0024-B008}{start}
\begin{defn}
تابع \LR{$g \colon B \to A$} د تابعې \LR{$f \colon A \to B$} \emph{معکوس}
ده که د ټولو \LR{$x \in A$} او \LR{$y \in B$} دپاره \LR{$f(g(y)) = y$} او
\LR{$g(f(x)) = x$} وي۔
\end{defn}
\olpseganchor{OLP-0024-B008}{end}

\olpseganchor{OLP-0024-B009}{start}
که \LR{$f$} معکوس~\LR{$g$} ولري، ډېر وخت د~\LR{$g$} پر ځاے \LR{$f^{-1}$} ليکو۔
\olpseganchor{OLP-0024-B009}{end}

\olpseganchor{OLP-0024-B010}{start}
\begin{explain}
اوس به معلومه کړو چې تابعې کله معکوسونه لري۔ د \LR{$f\colon A \to B$}
د معکوس دپاره يو مناسب نوماند \LR{$g\colon B \to A$} دے چې «په دې
تعريف شوے» وي:
\[
g(y) = \text{\RL{«هغه يوازينے» \LR{$x$} چې \LR{$f(x) = y$} وي۔}}
\]
خو د «په دې تعريف شوے» (او «هغه يوازينے») شاوخوا احتياطي نښې
اشاره کوي چې دا تعريف نۀ دے۔ لږ تر لږه، تل او په بشپړ عموميت
سره کار نۀ ورکوي۔ ځکه د دې دپاره چې دا تعريف يوه تابع وټاکي،
بايد يو او يوازې يو~\LR{$x$} وي چې \LR{$f(x) = y$} وي---د~\LR{$g$} وت بايد په
يوازيني ډول ټاکل شوے وي۔ سربېره پر دې، د هر \LR{$y \in B$} دپاره
بايد ټاکل شوے وي۔ که داسې \LR{$x_1$} او \LR{$x_2 \in A$} وي چې
\LR{$x_1 \neq x_2$} خو \LR{$f(x_1) = f(x_2)$} وي، نو د
\LR{$y = f(x_1) = f(x_2)$} دپاره به \LR{$g(y)$} په يوازيني ډول نۀ وي
ټاکل شوے۔ او که هېڅ داسې~\LR{$x$} نۀ وي چې \LR{$f(x) = y$} وي، نو \LR{$g(y)$}
به بېخي نۀ وي ټاکل شوے۔ په بله وينا، د~\LR{$g$} د تعريف کېدو دپاره
بايد~\LR{$f$} هم يو پر يو وي او هم پر هدف سټ بشپړه۔
\olpseganchor{OLP-0024-B010}{end}

\olpseganchor{OLP-0024-B011}{start}
راځئ ګام په ګام لاړ شو۔ پوښتنه به په دوو برخو ووېشو: د ورکړل
شوې تابعې~\LR{$f\colon A \to B$} دپاره، کله داسې تابع \LR{$g\colon B \to A$}
شته چې \LR{$g(f(x)) = x$} وي؟ داسې~\LR{$g$} د \LR{$f$} کړے کار «بېرته اړوي»
او د~\LR{$f$} \emph{کيڼ معکوس} بلل کېږي۔ دويمه دا چې کله داسې تابع
\LR{$h\colon B \to A$} شته چې \LR{$f(h(y)) = y$} وي؟ داسې~\LR{$h$} د~\LR{$f$}
\emph{ښي معکوس} بلل کېږي---\LR{$f$} د~\LR{$h$} کړے کار «بېرته اړوي»۔
\end{explain}
\olpseganchor{OLP-0024-B011}{end}

\olpseganchor{OLP-0024-B012}{start}
\begin{prop}
که د تعريف ساحه ناتشه وي او \LR{$f\colon A \to B$} يوه يو پر يو تابع وي، نو د~\LR{$f$} يو
\emph{کيڼ معکوس}~\LR{$g\colon B \to A$} شته چې د ټولو \LR{$x \in A$}
دپاره \LR{$g(f(x)) = x$} وي۔
\textbf{د ژباړې د نسخې سمون (OLFUN-001):} په انګرېزي سرچينه کښې
د ناتشې تعريف ساحې شرط نۀ و؛ خو لاندې ثبوت يو غړے ترې ټاکي۔
\end{prop}
\olpseganchor{OLP-0024-B012}{end}

\olpseganchor{OLP-0024-B013}{start}
\begin{proof}
فرض کړئ \LR{$f\colon A \to B$} يوه يو پر يو تابع ده۔ يو \LR{$y \in B$}
په پام کښې ونيسئ۔ که \LR{$y \in \ran{f}$} وي، نو داسې \LR{$x \in A$} شته
چې \LR{$f(x) = y$} وي۔ ځکه چې \LR{$f$} يو پر يو ده، يوازې يو داسې
\LR{$x \in A$} شته۔ بيا داسې تعريفولے شو: \LR{$g(y) = x$}؛ يعنې \LR{$g(y)$}
«هغه يوازينے» \LR{$x \in A$} دے چې \LR{$f(x) = y$} وي۔ که
\LR{$y \notin \ran{f}$} وي، له هر يوه~\LR{$a \in A$} سره ئې تړلے شو۔ نو
يو \LR{$a \in A$} ټاکلے او \LR{$g\colon B \to A$} داسې تعريفولے شو:
\[
g(y) = \begin{cases}
    x & \text{\RL{که \LR{$f(x) = y$} وي}}\\
    a & \text{\RL{که \LR{$y \notin \ran{f}$} وي۔}}
\end{cases}
\]
دا د ټولو \LR{$y \in B$} دپاره تعريف شوې ده، ځکه چې د هر داسې
\LR{$y \in \ran{f}$} دپاره په دقيقه توګه يو \LR{$x \in A$} شته چې
\LR{$f(x) = y$} وي۔ د تعريف له مخې، که \LR{$y = f(x)$} وي، نو \LR{$g(y) = x$}
وي؛ يعنې \LR{$g(f(x)) = x$}۔
\end{proof}
\olpseganchor{OLP-0024-B013}{end}

\olpseganchor{OLP-0024-B014}{start}
\begin{prob}
وښيئ چې که \LR{$f\colon A \to B$} کيڼ معکوس~\LR{$g$} ولري، نو~\LR{$f$}
يو پر يو ده۔
\end{prob}
\olpseganchor{OLP-0024-B014}{end}

\olpseganchor{OLP-0024-B015}{start}
\begin{prop}
    که \LR{$f\colon A \to B$} يوه پر هدف سټ بشپړه تابع وي، نو د~\LR{$f$}
    يو \emph{ښي معکوس}~\LR{$h\colon B \to A$} شته چې د ټولو~\LR{$y \in B$}
    دپاره \LR{$f(h(y)) = y$} وي۔
\end{prop}
\olpseganchor{OLP-0024-B015}{end}

\olpseganchor{OLP-0024-B016}{start}
\begin{proof}
فرض کړئ \LR{$f\colon A \to B$} يوه پر هدف سټ بشپړه تابع ده۔ يو \LR{$y \in B$}
په پام کښې ونيسئ۔ ځکه چې~\LR{$f$} پر هدف سټ بشپړه ده، داسې \LR{$x_y \in A$}
شته چې \LR{$f(x_y) = y$} وي۔ بيا داسې تعريفولے شو: \LR{$h(y) = x_y$}؛
يعنې د هر \LR{$y \in B$} دپاره يو داسې \LR{$x \in A$} ټاکو چې \LR{$f(x) = y$}
وي؛ ځکه چې~\LR{$f$} پر هدف سټ بشپړه ده، تل لږ تر لږه يو د ټاکلو دپاره
شته۔\footnote{ځکه چې \LR{$f$} پر هدف سټ بشپړه ده، د هر~\LR{$y \in B$}
دپاره سټ \LR{$\Setabs{x}{f(x) = y}$} تش نۀ دے۔ زموږ د~\LR{$h$} تعريف غواړي
چې له دې هر سټ څخه يو يوازينے \LR{$x$} وټاکو۔ دا چې دا کار تل ممکن
دے، په حقيقت کښې څرګنده خبره نۀ ده---د دې ټاکنو امکان په ساده
ډول د يوه بديهي اصل په توګه فرض کېږي۔ په بله وينا، دا قضيه
د ټاکنې په نامه بديهي اصل فرض کوي؛ دا مسئله به
په لنډه پرېږدو۔ خو په ډېرو ځانګړو حالتونو کښې، لکه
چې \LR{$A = \Nat$} وي يا متناهي وي، يا \LR{$f$} يوه دوه اړخيزه يو پر يو تابع وي،
د ټاکنې بديهي اصل ته اړتيا نۀ وي۔ (په هغه ځانګړي حالت کښې چې
\LR{$f$} دوه اړخيزه يو پر يو وي، د هر \LR{$y \in B$} دپاره سټ
\LR{$\Setabs{x}{f(x) = y}$} په دقيقه توګه يو غړے لري؛ نو د
انتخاب دپاره څو بديلونه نشته۔)}
د تعريف له مخې، که \LR{$x = h(y)$} وي، نو \LR{$f(x) = y$} وي؛ يعنې
د هر \LR{$y \in B$} دپاره \LR{$f(h(y)) = y$}۔
\end{proof}
\olpseganchor{OLP-0024-B016}{end}

\olpseganchor{OLP-0024-B017}{start}
\begin{prob}
وښيئ چې که \LR{$f\colon A \to B$} ښي معکوس~\LR{$h$} ولري، نو~\LR{$f$}
پر هدف سټ بشپړه ده۔
\end{prob}
\olpseganchor{OLP-0024-B017}{end}

\olpseganchor{OLP-0024-B018}{start}
\begin{explain}
  د مخکيني ثبوت د مفکورو په يوځاے کولو اوس دا ترلاسه کوو چې
  هره دوه اړخيزه يو پر يو تابع معکوس لري؛ يعنې يوه واحده تابع شته چې
  هم د~\LR{$f$} کيڼ معکوس دے او هم ښي معکوس۔
\end{explain}
\olpseganchor{OLP-0024-B018}{end}

\olpseganchor{OLP-0024-B019}{start}
\begin{prop}\ollabel{prop:bijection-inverse}
که \LR{$f\colon A \to B$} يوه دوه اړخيزه يو پر يو تابع وي، داسې تابع
\LR{$f^{-1}\colon B \to A$} شته چې د ټولو \LR{$x \in A$} دپاره
\LR{$f^{-1}(f(x)) = x$} او د ټولو \LR{$y \in B$} دپاره \LR{$f(f^{-1}(y)) = y$} وي۔
\end{prop}
\olpseganchor{OLP-0024-B019}{end}

\olpseganchor{OLP-0024-B020}{start}
\begin{proof}
تمرين۔
\end{proof}
\olpseganchor{OLP-0024-B020}{end}

\olpseganchor{OLP-0024-B021}{start}
\begin{prob}
\olref[sfr][fun][inv]{prop:bijection-inverse} ثابته کړئ۔ بايد
\LR{$f^{-1}$} تعريف کړئ، وښيئ چې تابع ده، او وښيئ چې د~\LR{$f$} معکوس
ده؛ يعنې د ټولو \LR{$x \in A$} او \LR{$y \in B$} دپاره \LR{$f^{-1}(f(x)) = x$}
او \LR{$f(f^{-1}(y)) = y$} وي۔
\end{prob}
\olpseganchor{OLP-0024-B021}{end}

\olpseganchor{OLP-0024-B022}{start}
\begin{explain}
د معکوسونو د ترلاسه کولو يوه لږه زياته عمومي لار هم شته۔ په
\olref[kin]{sec} کښې مو وليدل چې هره تابع \LR{$f$} د ټولو \LR{$x \in A$}
دپاره د \LR{$f'(x) = f(x)$} په ټاکلو سره يوه پر هدف سټ بشپړه تابع
\LR{$f' \colon A \to \ran{f}$} رامنځته کوي۔ په ښکاره، که~\LR{$f$}
يو پر يو وي، نو~\LR{$f'$} دوه اړخيزه يو پر يو ده، او له دې کبله د
\olref{prop:bijection-inverse} له مخې يوازينے معکوس لري۔ د نښو
په کارولو کښې په يو ډېر واړه تسامح سره، کله کله د \LR{$f'$} معکوس
يوازې «د~\LR{$f$} معکوس» بولو۔
\end{explain}
\olpseganchor{OLP-0024-B022}{end}

\olpseganchor{OLP-0024-B023}{start}
\begin{prop}\ollabel{prop:left-right}%
  وښيئ چې که \LR{$f\colon A \to B$} کيڼ معکوس~\LR{$g$} او ښي معکوس~\LR{$h$}
  ولري، نو \LR{$h = g$}۔
\end{prop}
\olpseganchor{OLP-0024-B023}{end}

\olpseganchor{OLP-0024-B024}{start}
\begin{proof}
  تمرين۔
\end{proof}
\olpseganchor{OLP-0024-B024}{end}

\olpseganchor{OLP-0024-B025}{start}
\begin{prob}
  \olref[sfr][fun][inv]{prop:left-right} ثابته کړئ۔
\end{prob}
\olpseganchor{OLP-0024-B025}{end}

\olpseganchor{OLP-0024-B026}{start}
\begin{prop}\ollabel{prop:inverse-unique}
هره تابع~\LR{$f$} تر يوه زيات معکوس نۀ لري۔
\end{prop}
\olpseganchor{OLP-0024-B026}{end}

\olpseganchor{OLP-0024-B027}{start}
\begin{proof}
  فرض کړئ \LR{$g$} او \LR{$h$} دواړه د~\LR{$f$} معکوسونه دي۔ نو په ځانګړي ډول
  \LR{$g$} د~\LR{$f$} کيڼ معکوس او \LR{$h$} ښي معکوس دے۔ د
  \olref{prop:left-right} له مخې، \LR{$g = h$}۔
\end{proof}
\olpseganchor{OLP-0024-B027}{end}

\olpseganchor{OLP-0025-B004}{start}
\olfileid{sfr}{fun}{cmp}
\olsection{د تابعو ترکيب}
\olpseganchor{OLP-0025-B004}{end}

\olpseganchor{OLP-0025-B005}{start}
\begin{explain}
په \olref[inv]{sec} کښې مو وليدل چې
د يوې دوه اړخيزې يو پر يو تابعې~\LR{$f$} معکوس~\LR{$f^{-1}$} پخپله يوه تابع ده۔ پر
تابعو يو بل عمل ترکيب دے: موږ دوه تابعې، \LR{$f$} او~\LR{$g$}، سره
ترکيبولے او نوې تابع تعريفولے شو؛ يعنې لومړے~\LR{$f$} او بيا~\LR{$g$}
تطبيق کړو۔ البته، دا يوازې هغه وخت ممکنه ده چې د قيمتونو سټونه
او د تعريف ساحې سره برابر راشي؛ يعنې د~\LR{$f$} د قيمتونو سټ بايد د~\LR{$g$}
د تعريف د ساحې فرعي سټ وي۔ پر تابعو
دا عمل، پر اړيکو د نسبي حاصل ضرب له هغه عمل سره ورته دے چې په
\olref[rel][ops]{sec} کښې راغلے دے۔
\olpseganchor{OLP-0025-B005}{end}

\olpseganchor{OLP-0025-B006}{start}
انځوريز شکل د ترکيب د مفکورې په تشريح کښې مرسته کولے شي۔ په
\olref{fig:composition} کښې دوه تابعې \LR{$f \colon A \to B$} او
\LR{$g \colon B \to C$} او د هغو ترکيب~\LR{$(\comp{f}{g})$} ښيو۔ تابع
\LR{$(\comp{f}{g}) \colon A \to C$} د~\LR{$A$} هر غړے د~\LR{$C$} له يوه
غړي سره تړي۔ دا چې د \LR{$A$} يو غړے
د~\LR{$C$} له کوم غړي سره تړل کېږي، داسې ئې ټاکو:
د ورکړل شوي ننوت \LR{$x \in A$} دپاره، لومړے تابع \LR{$f$} پر~\LR{$x$} تطبيق
کړئ، چې يو \LR{$f(x) = y \in B$} به ورکړي؛ بيا تابع \LR{$g$} پر~\LR{$y$} تطبيق
کړئ، چې يو \LR{$g(f(x)) = g(y) = z \in C$} به ورکړي۔
\begin{figure}
  \olasset[2\olphotowidth]{../upstream/assets/diagrams/composition.tikz}
  \caption{د دوو تابعو \LR{$f$} او~\LR{$g$} ترکيب \LR{$g \circ f$}۔}
  \ollabel{fig:composition}
\end{figure}
\end{explain}
\olpseganchor{OLP-0025-B006}{end}

\olpseganchor{OLP-0025-B007}{start}
\begin{defn}[ترکيب]
فرض کړئ \LR{$f\colon A \to B$} او \LR{$g\colon B \to C$} تابعې دي۔
له~\LR{$g$} سره د \LR{$f$} \emph{ترکيب} دا دے: \LR{$\comp{f}{g} \colon A \to C$}،
چې \LR{$(\comp{f}{g})(x) = g(f(x))$} وي۔
\end{defn}
\olpseganchor{OLP-0025-B007}{end}

\olpseganchor{OLP-0025-B008}{start}
\begin{ex}
تابعې \LR{$f(x) = x + 1$} او \LR{$g(x) = 2x$} په پام کښې ونيسئ۔ ځکه چې
\LR{$(\comp{f}{g})(x) = g(f(x))$}، نو د هر ننوت~\LR{$x$} دپاره بايد لومړے
د هغه ځايناستے واخلئ، بيا پايله په دوو ضرب کړئ۔ نو د هغو ترکيب
په \LR{$(\comp{f}{g})(x) = 2(x+1)$} ورکړل شوے دے۔
\end{ex}
\olpseganchor{OLP-0025-B008}{end}

\olpseganchor{OLP-0025-B009}{start}
\begin{prob}
وښيئ چې که \LR{$f \colon A \to B$} او \LR{$g \colon B \to C$} دواړه
يو پر يو وي، نو \LR{$\comp{f}{g}\colon A \to C$} هم يو پر يو ده۔
\end{prob}
\olpseganchor{OLP-0025-B009}{end}

\olpseganchor{OLP-0025-B010}{start}
\begin{prob}
وښيئ چې که \LR{$f \colon A \to B$} او \LR{$g \colon B \to C$} دواړه
پر هدف سټ بشپړه وي، نو \LR{$\comp{f}{g}\colon A \to C$} هم پر هدف سټ بشپړه ده۔
\end{prob}
\olpseganchor{OLP-0025-B010}{end}

\olpseganchor{OLP-0025-B011}{start}
\begin{prob}
فرض کړئ \LR{$f \colon A \to B$} او \LR{$g \colon B \to C$}۔ وښيئ چې د
\LR{$\comp{f}{g}$} ګراف \LR{$R_f \mid R_g$} دے۔
\end{prob}
\olpseganchor{OLP-0025-B011}{end}

\olpseganchor{OLP-0026-B004}{start}
\olfileid{sfr}{fun}{par}
\olpseganchor{OLP-0026-B004}{end}

\olpseganchor{OLP-0026-B005}{start}
\olsection{جزوي تابعې}
\olpseganchor{OLP-0026-B005}{end}

\olpseganchor{OLP-0026-B006}{start}
\begin{explain}
کله کله ګټوره وي چې د تابعې په تعريف کښې نرمي راولو، څو د ټولو
ممکنو ننوتونو دپاره د تابعې د وت تعريف کېدل لازم نۀ وي۔ داسې
نگاشتونه \emph{جزوي تابعې} بلل کېږي۔
\end{explain}
\olpseganchor{OLP-0026-B006}{end}

\olpseganchor{OLP-0026-B007}{start}
\begin{defn}
\emph{جزوي تابع} \LR{$f \colon A \pto B$} داسې نگاشت دے چې د~\LR{$A$} له
هر غړي سره د~\LR{$B$} تر يوه زيات غړے نۀ
تړي۔ که \LR{$f$} له \LR{$x \in A$} سره د~\LR{$B$} يو غړے وتړي، وايو چې \LR{$f(x)$}
\emph{تعريف شوے} دے، او که داسې نۀ وي، \emph{تعريف نۀ شوے} دے۔
که \LR{$f(x)$} تعريف شوے وي، \LR{$f(x) \fdefined$} ليکو؛ کنه،
\LR{$f(x) \fundefined$} ليکو۔ د جزوي تابعې~\LR{$f$} \emph{د تعريف ساحه}
د~\LR{$A$} هغه فرعي سټ دے چې تابع پرې تعريف شوې ده؛ يعنې
\LR{$\dom{f} = \Setabs{x \in A}{f(x) \fdefined}$}۔
\end{defn}
\olpseganchor{OLP-0026-B007}{end}

\olpseganchor{OLP-0026-B008}{start}
\begin{ex}
هره تابع \LR{$f\colon A \to B$} يوه جزوي تابع هم ده۔ هغه جزوي تابعې
چې پر~\LR{$A$} هرځاے تعريف شوې وي---يعنې هغه څه چې تر دې دمه مو
يوازې تابع بلل---\emph{هرځاے تعريف شوې} تابعې هم بلل کېږي۔
\end{ex}
\olpseganchor{OLP-0026-B008}{end}

\olpseganchor{OLP-0026-B009}{start}
\begin{ex}
جزوي تابع \LR{$f \colon \Real \pto \Real$} چې په \LR{$f(x) = 1/x$} ورکړل
شوې، د \LR{$x = 0$} دپاره تعريف شوې نۀ ده، او په نورو ټولو ځايونو
کښې تعريف شوې ده۔
\end{ex}
\olpseganchor{OLP-0026-B009}{end}

\olpseganchor{OLP-0026-B010}{start}
\begin{prob}
د ورکړل شوې \LR{$f\colon A \pto B$} دپاره، جزوي تابع \LR{$g\colon B \pto
A$} داسې تعريف کړئ: د هر \LR{$y \in B$} دپاره، که يو يوازينے
\LR{$x \in A$} وي چې \LR{$f(x) = y$} وي، نو \LR{$g(y) = x$}؛ کنه،
\LR{$g(y) \fundefined$}۔ وښيئ چې که \LR{$f$} يو پر يو وي، نو د ټولو
\LR{$x \in \dom{f}$} دپاره \LR{$g(f(x)) = x$}، او د ټولو \LR{$y \in \ran{f}$}
دپاره \LR{$f(g(y)) = y$} وي۔
\end{prob}
\olpseganchor{OLP-0026-B010}{end}

\olpseganchor{OLP-0026-B011}{start}
\begin{defn}[د جزوي تابعې ګراف]
فرض کړئ \LR{$f\colon A \pto B$} يوه جزوي تابع ده۔ د~\LR{$f$} \emph{ګراف}
هغه اړيکه \LR{$R_f \subseteq A \times B$} ده چې داسې تعريف شوې ده:
\[
R_f = \Setabs{\tuple{x,y}}{f(x) = y}.
\]
\end{defn}
\olpseganchor{OLP-0026-B011}{end}

\olpseganchor{OLP-0026-B012}{start}
\begin{prop}
فرض کړئ \LR{$R \subseteq A \times B$} دا خاصيت لري چې هرکله \LR{$Rxy$}
او \LR{$Rxy'$} وي، نو \LR{$y = y'$} وي۔ بيا \LR{$R$} د هغې جزوي تابعې
\LR{$f\colon A \pto B$} ګراف ده چې داسې تعريف شوې ده: که داسې \LR{$y$}
وي چې \LR{$Rxy$} وي، نو \LR{$f(x) = y$}؛ کنه، \LR{$f(x) \fundefined$}۔ که \LR{$R$}
\emph{مسلسله} هم وي، يعنې د هر \LR{$x \in A$} دپاره داسې \LR{$y \in B$}
وي چې \LR{$Rxy$} وي، نو \LR{$f$} هرځاے تعريف شوې ده۔
\end{prop}
\olpseganchor{OLP-0026-B012}{end}

\olpseganchor{OLP-0026-B013}{start}
\begin{proof}
فرض کړئ داسې \LR{$y$} شته چې \LR{$Rxy$} وي۔ که يو بل \LR{$y' \neq y$} هم واے
چې \LR{$Rxy'$} واے، پر \LR{$R$} لګېدلے شرط به مات شوے و۔ نو که داسې \LR{$y$}
شته چې \LR{$Rxy$} وي، هغه \LR{$y$} يوازينے دے؛ له دې کبله \LR{$f$} په سمه
توګه تعريف شوې ده۔ په ښکاره، \LR{$R_f = R$}، او که~\LR{$R$} مسلسله وي،
نو \LR{$f$} هرځاے تعريف شوې ده۔
\end{proof}
\olpseganchor{OLP-0026-B013}{end}

\olpseganchor{OLP-0027-B004}{start}
\olchapter{sfr}{siz}{د سټونو اندازه}
\olpseganchor{OLP-0027-B004}{end}

\olpseganchor{OLP-0027-B005}{start}
\begin{editorial}
دا باب شمېرنې، د شمېر وړ والے او د شمېر نۀ وړ والے څېړي۔
ځينې برخې دوه بڼې لري: يوه زياته ابتدايي بڼه، چې شمېرنې د
لړونو يا له \LR{$\PosInt$} څخه د هغو تابعو په توګه اخلي چې پر هدف
سټ بشپړې وي؛ او يوه زياته تجريدي بڼه، چې شمېرنې له \LR{$\Nat$} سره
د دوه اړخيزو يو پر يو تابعو په توګه تعريفوي۔
\end{editorial}
\olpseganchor{OLP-0027-B005}{end}





















\olpseganchor{OLP-0027-B016}{start}
\begin{editorial}
لاندې \olref[sfr][siz][enm-alt]{sec}،
\olref[sfr][siz][nen-alt]{sec} او \olref[sfr][siz][red-alt]{sec}
د \olref[sfr][siz][enm]{sec}، \olref[sfr][siz][nen]{sec} او
\olref[sfr][siz][red]{sec} بديلې بڼې دي چې ټيم بټن د خپل
«د سټونو خلاصه نظريه» متن کښې د کارولو دپاره برابرې کړې دي۔
دوې لږې زياتې پرمختللې دي او د شمېر وړ والي يو بېل تعريف
کاروي چې د سټونو د نظريې په سياق کښې زيات مناسب دے؛ يعنې
له \LR{$\Nat$} يا د هغې له پيلنۍ برخې سره دوه اړخيزه يو پر يو تابع
لرل، د دې پر ځاے چې په لړ کښې د راوړلو وړ وي يا له \LR{$\PosInt$}
څخه د يوې پر هدف سټ بشپړې تابعې د قيمتونو سټ وي۔
\end{editorial}
\olpseganchor{OLP-0027-B016}{end}

\olpseganchor{OLP-0028-B004}{start}
\olfileid{sfr}{siz}{int}
\olpseganchor{OLP-0028-B004}{end}

\olpseganchor{OLP-0028-B005}{start}
\olsection{پېژندنه}
\olpseganchor{OLP-0028-B005}{end}

\olpseganchor{OLP-0028-B006}{start}
کله چې ګېورګ کانتور په 1870مو کلونو کښې د سټونو نظريې ته وده ورکړه،
د هغه يو مقصد دا و چې د نامتناهي مجموعې مفکوره د منلو وړ کړي---
هغه څه چې د منځنيو پېړيو پوهانو به فعلي نامتناهي والے بللو۔ د
دې يوه مهمه برخه د بېلو سټونو د \emph{اندازې} په اړه د هغه
څېړنه وه۔ که \LR{$a$}، \LR{$b$} او \LR{$c$} ټول بېل وي، نو په وجداني ډول
سټ \LR{$\{a, b, c\}$} له \LR{$\{a, b\}$} څخه \emph{لوے} دے۔ خو د
نامتناهي سټونو په اړه څه وايو؟ ايا ټول د يو بل هومره لوے دي؟
معلومېږي چې داسې نۀ ده۔
\olpseganchor{OLP-0028-B006}{end}

\olpseganchor{OLP-0028-B007}{start}
دلته لومړۍ مهمه مفکوره د شمېرنې ده۔ هر متناهي سټ د هغه د ټولو
غړو په ليکلو په لړ کښې راوړلے شو۔ د ځينو
نامتناهي سټونو ټول غړي هم په لړ کښې راوړلے شو، که
اجازه ورکړو چې لړ پخپله نامتناهي وي۔ داسې سټونه د شمېر وړ
بلل کېږي۔ د کانتور حيرانوونکې پايله، چې د دې باب تر پايه به
ئې پوره وپېژنو، دا وه چې ځينې نامتناهي سټونه د شمېر وړ نۀ دي۔
\olpseganchor{OLP-0028-B007}{end}

\olpseganchor{OLP-0029-B004}{start}
\olfileid{sfr}{siz}{enm}
\olpseganchor{OLP-0029-B004}{end}

\olpseganchor{OLP-0029-B005}{start}
\olsection{شمېرنې او د شمېر وړ سټونه}
\olpseganchor{OLP-0029-B005}{end}

\olpseganchor{OLP-0029-B006}{start}
\begin{editorial}
  دا برخه د سټونو شمېرنې څېړي او له \LR{$\PosInt$} څخه د پر هدف سټ
  بشپړو تابعو په توګه ئې تعريفوي۔ د هغو لوستونکو دپاره چې
  لږ رياضيکي شاليد لري، خبرې ګام په ګام وړاندې کوي۔ يوه
  بديله، لنډه بڼه په \olref[enm-alt]{sec} کښې ورکړل شوې، چې
  شمېرنې په بېل ډول تعريفوي: له \LR{$\Nat$} (يا د هغې له پيلنۍ
  برخې) سره د دوه اړخيزو يو پر يو تابعو په توګه۔
\end{editorial}
\olpseganchor{OLP-0029-B006}{end}

\olpseganchor{OLP-0029-B007}{start}
\begin{explain}
مخکښې مو د سټونو بېلګې د هغو د غړو په
ليکلو ورکړې دي۔ راځئ په زياته عمومي توګه بحث وکړو چې د يوه
سټ غړي څنګه او کله په لړ کښې راوړلے شو، که هغه سټ
نامتناهي هم وي۔
\end{explain}
\olpseganchor{OLP-0029-B007}{end}

\olpseganchor{OLP-0029-B008}{start}
\begin{defn}[شمېرنه، په غير صوري ډول]
په غير صوري ډول، د سټ~\LR{$A$} \emph{شمېرنه} د~\LR{$A$} د
غړو يو (کېدے شي نامتناهي) لړ دے، داسې چې د \LR{$A$}
هر غړے د لړ په يوه متناهي مقام کښې راشي۔ که \LR{$A$} شمېرنه
ولري، نو \LR{$A$} ته \emph{د شمېر وړ} ويل کېږي۔
\end{defn}
\olpseganchor{OLP-0029-B008}{end}

\olpseganchor{OLP-0029-B009}{start}
\begin{explain}
د شمېرنو په اړه څو ټکي:
\begin{enumerate}
\item يوازې هغه لړونه شمېرنې ګڼو چې پيل ولري او له لومړني
  غړي پرته د هر غړي سمدستي
  مخکښې يو يوازينے غړے وي۔ په بله وينا، د لړ د لومړي
  غړي او هر بل غړي ترمنځ
  يوازې متناهي شمېر غړي وي۔ په ځانګړي ډول، دا معنٰی لري چې
  د شمېرنې هر غړے متناهي مقام لري: لومړے غړے
  په مقام~\LR{$1$} کښې وي، دويم په مقام~\LR{$2$} کښې، او داسې نور۔
\item د هماغه سټ~\LR{$A$} بېلې شمېرنې لرلے شو چې د
  غړو د راتلو په ترتيب کښې سره بېلې وي:
  \LR{$4$}، \LR{$1$}، \LR{$25$}، \LR{$16$}،~\LR{$9$} د لومړيو پنځو مربع عددونو (د سټ)
  هومره ښه شمېرنه ده لکه \LR{$1$}، \LR{$4$}، \LR{$9$}، \LR{$16$}،~\LR{$25$}۔
\item تکراري شمېرنې هم شمېرنې دي: \LR{$1$}، \LR{$1$}، \LR{$2$}، \LR{$2$}، \LR{$3$}،
  \LR{$3$}،~\dots{} د هماغه سټ شمېرنه ده چې \LR{$1$}، \LR{$2$}، \LR{$3$}،~\dots{}
  ئې شمېرنه کوي۔
\item کله چې يوه شمېرنه ټاکو، ترتيب او تکرار \emph{اهميت لري}:
  د مثبتو صحيح عددونو شمېرنه په \LR{$1$}، \LR{$2$}، \LR{$3$}، \LR{$1$}، \dots{}
  پيلولے شو، خو نمونه هغه وخت په اسانه ښکاري چې په معمولي
  ډول ئې په \LR{$1$}، \LR{$2$}، \LR{$3$}، \LR{$4$}،~\dots وشمېرو۔
\item شمېرنې بايد پيل ولري: \dots، \LR{$3$}، \LR{$2$}، \LR{$1$} د مثبتو صحيح
  عددونو شمېرنه نۀ ده، ځکه چې لومړے غړے نۀ لري۔ د دې
  د ليدلو دپاره چې دا له غير صوري تعريفه څنګه راځي، له ځانه
  وپوښتئ: «عدد 76 د لړ په کوم مقام کښې راځي؟»
\item لاندې لړ د مثبتو صحيح عددونو شمېرنه نۀ ده:
  \LR{$1$}، \LR{$3$}، \LR{$5$}، \dots، \LR{$2$}، \LR{$4$}، \LR{$6$}، \dots\@ ستونزه دا ده چې
  جفت عددونه د متناهي مقامونو پر ځاے په \LR{$\infty + 1$}،
  \LR{$\infty + 2$}، \LR{$\infty + 3$} مقامونو کښې راځي۔
\item تش سټ د شمېر وړ دے: تش لړ ئې شمېري!{}
\end{enumerate}
\end{explain}
\olpseganchor{OLP-0029-B009}{end}

\olpseganchor{OLP-0029-B010}{start}
\begin{prop}
  که \LR{$A$} شمېرنه ولري، بې له تکراره شمېرنه هم لري۔
\end{prop}
\olpseganchor{OLP-0029-B010}{end}

\olpseganchor{OLP-0029-B011}{start}
\begin{proof}
  فرض کړئ \LR{$A$} شمېرنه \LR{$x_1$}، \LR{$x_2$}، \dots{} لري، چې هر \LR{$x_i$}
  د~\LR{$A$} يو غړے دے۔ له شمېرنې څخه د تکراري
  غړو په لرې کولو تکرارونه لرې کولے شو۔ د
  بېلګې په توګه، شمېرنه داسې نوې شمېرنې ته اړولے شو چې \LR{$x_i$}
  پکښې هله وليکو چې د~\LR{$A$} داسې غړے وي چې په \LR{$x_1$}،
  \dots، \LR{$x_{i-1}$} کښې نۀ وي، يا \LR{$x_i$} له لړه لرې کړو که له
  مخکښې په \LR{$x_1$}، \dots،~\LR{$x_{i-1}$} کښې راغلے وي۔
\end{proof}
\olpseganchor{OLP-0029-B011}{end}

\olpseganchor{OLP-0029-B012}{start}
وروستے استدلال ښيي چې د شمېرنو او د شمېر وړ سټونو د ښه
پوهېدو او د هغو په اړه د خبرو د ثابتولو دپاره يو زيات کره
تعريف ته اړتيا لرو۔ لاندې تعريف دا اړتيا پوره کوي۔
\olpseganchor{OLP-0029-B012}{end}

\olpseganchor{OLP-0029-B013}{start}
\begin{defn}[شمېرنه، په صوري ډول]
د سټ \LR{$A \neq \emptyset$} \emph{شمېرنه} هره پر هدف سټ بشپړه
تابع \LR{$f \colon \PosInt \to A$} ده۔
\end{defn}
\olpseganchor{OLP-0029-B013}{end}

\olpseganchor{OLP-0029-B014}{start}
\begin{explain}
راځئ ځان ډاډمن کړو چې صوري تعريف او هغه غير صوري تعريف چې
يو ممکن نامتناهي لړ کاروي، همارزښته دي۔ لومړے، له \LR{$\PosInt$}
څخه سټ~\LR{$A$} ته هره پر هدف سټ بشپړه تابع د~\LR{$A$} شمېرنه کوي۔ داسې
تابع د پاسني غير صوري تعريف له مخې يوه شمېرنه ټاکي: لړ
\LR{$f(1)$}، \LR{$f(2)$}، \LR{$f(3)$}، \dots۔ ځکه چې \LR{$f$} پر هدف سټ بشپړه ده،
دا تضمين شوے چې د~\LR{$A$} هر غړے به د کوم~\LR{$n \in \PosInt$}
دپاره د~\LR{$f(n)$} قيمت وي۔ نو د \LR{$A$} هر غړے په لړ کښې
په کوم متناهي مقام راځي۔ ځکه چې کېدے شي تابع يو پر يو
نۀ وي، لړ تکرار لرلے شي؛ خو دا د منلو وړ دي، لکه چې پورته
ياد شول۔
\olpseganchor{OLP-0029-B014}{end}

\olpseganchor{OLP-0029-B015}{start}
بل خوا، که داسې لړ راکړل شوے وي چې د~\LR{$A$} ټول غړي
شمېري، يوه پر هدف سټ بشپړه تابع \LR{$f\colon \PosInt \to A$} داسې
تعريفولے شو چې \LR{$f(n)$} د لړ \LR{$n$}م غړے وټاکو، يا که
\LR{$n$}م غړے نۀ وي، نو د لړ وروستے غړے وټاکو۔
يوازې په هغه حالت کښې له دې څخه پر هدف سټ بشپړه تابع نۀ
جوړېږي چې \LR{$A$} تش وي او له دې کبله لړ هم تش وي۔ نو هر ناتش
لړ يوه پر هدف سټ بشپړه تابع \LR{$f\colon \PosInt \to A$} ټاکي۔
\end{explain}
\olpseganchor{OLP-0029-B015}{end}

\olpseganchor{OLP-0029-B016}{start}
\begin{defn}
  \ollabel{defn:enumerable}
  سټ~\LR{$A$} هله او يوازې هله د شمېر وړ دے چې تش وي يا شمېرنه ولري۔
\end{defn}
\olpseganchor{OLP-0029-B016}{end}

\olpseganchor{OLP-0029-B017}{start}
\begin{ex}
د مثبتو صحيح عددونو (\LR{$\PosInt$}) شمېرونکې تابع يوازې د عينيت
هغه تابع ده چې په \LR{$f(n) = n$} ورکړل شوې ده۔ د طبيعي عددونو
\LR{$\Nat$} شمېرونکې تابع \LR{$g(n) = n - 1$} ده۔
\end{ex}
\olpseganchor{OLP-0029-B017}{end}

\olpseganchor{OLP-0029-B018}{start}
\begin{ex}
تابعې \LR{$f\colon \PosInt \to \PosInt$} او \LR{$g \colon \PosInt \to
\PosInt$} چې داسې ورکړل شوې دي:
\begin{align*}
f(n) & = 2n \text{\RL{ او}}\\
g(n) & = 2n - 1
\end{align*}
په ترتيب سره د جفتو مثبتو صحيح عددونو او د طاقو مثبتو صحيح
عددونو شمېرنه کوي۔ خو له دوو هېڅ يوه د \LR{$\PosInt$} شمېرنه نۀ
ده، ځکه چې هېڅ يوه پر هدف سټ بشپړه نۀ ده۔
\end{ex}
\olpseganchor{OLP-0029-B018}{end}

\olpseganchor{OLP-0029-B019}{start}
\begin{prob}
د مثبتو مربع عددونو \LR{$1$}، \LR{$4$}، \LR{$9$}، \LR{$16$}، \dots يوه شمېرنه تعريف کړئ۔
\end{prob}
\olpseganchor{OLP-0029-B019}{end}

\olpseganchor{OLP-0029-B020}{start}
\begin{ex}
تابع \LR{$f(n) = (-1)^{n} \lceil \frac{(n-1)}{2}\rceil$} (چې
\LR{$\lceil x \rceil$} د \emph{پاسني صحيح عدد} تابع ښيي، يعنې \LR{$x$}
پورته تر ټولو نږدې صحيح عدد ته ګردوي) د صحيح عددونو~\LR{$\Int$}
د سټ شمېرنه کوي۔ پام وکړئ چې \LR{$f$} د مثبتو او منفي صحيح عددونو
ترمنځ په تګ راتګ سره د \LR{$\Int$} قيمتونه څنګه راوباسي:
\[
\begin{array}{c c c c c c c c}
f(1) & f(2) & f(3) & f(4) & f(5) & f(6) & f(7) & \dots \\ \\
- \lceil \tfrac{0}{2} \rceil & \lceil \tfrac{1}{2}\rceil & - \lceil \tfrac{2}{2} \rceil & \lceil \tfrac{3}{2} \rceil & - \lceil \tfrac{4}{2} \rceil  & \lceil \tfrac{5}{2}
\rceil & - \lceil \tfrac{6}{2} \rceil & \dots \\ \\
0 & 1 & -1 & 2 & -2 & 3 & -3 & \dots
\end{array}
\]
\textbf{د ژباړې د نسخې سمون (OLSIZ-001):} د انګرېزي جدول سرليک
او د فارمول قطار f(7) لري، خو د قيمتونو قطار ئې اړين قيمت -3
پرېښے و؛ دلته د هماغه فارمول له مخې ورزيات شوے دے۔
\LR{$f$} د حالتونو له مخې داسې تعريف شوې هم ګڼلے شئ:
\[
f(n) = \begin{cases}
  0 & \text{\RL{که \LR{$n = 1$} وي}}\\
  n/2 & \text{\RL{که \LR{$n$} جفت وي}}\\
  -(n-1)/2 & \text{\RL{که \LR{$n$} طاق او \LR{$>1$} وي}}
  \end{cases}
\]
\end{ex}
\olpseganchor{OLP-0029-B020}{end}

\olpseganchor{OLP-0029-B021}{start}
\begin{prob}
  وښيئ چې که \LR{$A$} او \LR{$B$} د شمېر وړ وي، نو \LR{$A \cup B$} هم
  داسې دے۔ د دې دپاره فرض کړئ چې پر هدف سټ بشپړه تابعې
  \LR{$f\colon \PosInt \to A$} او \LR{$g\colon \PosInt \to B$} شته؛
  يوه پر هدف سټ بشپړه تابع~\LR{$h\colon \PosInt \to A \cup B$}
  تعريف کړئ او ثابته کړئ چې پر هدف سټ بشپړه ده۔ هغه حالتونه
  هم په پام کښې ونيسئ چې \LR{$A$} يا~\LR{$B = \emptyset$} وي۔
\end{prob}
\olpseganchor{OLP-0029-B021}{end}
  
\olpseganchor{OLP-0029-B022}{start}
\begin{prob}
  وښيئ چې که \LR{$B \subseteq A$} او \LR{$A$} د شمېر وړ وي، نو~\LR{$B$}
  هم داسې دے۔ د دې دپاره فرض کړئ چې يوه پر هدف سټ بشپړه
  تابع \LR{$f\colon \PosInt \to A$} شته۔ يوه پر هدف سټ بشپړه
  تابع~\LR{$g\colon \PosInt \to B$} تعريف کړئ او ثابته کړئ چې
  پر هدف سټ بشپړه ده۔ که \LR{$B = \emptyset$} وي، څه پېښېږي؟
\end{prob}
\olpseganchor{OLP-0029-B022}{end}
    
\olpseganchor{OLP-0029-B023}{start}
\begin{prob}
  پر \LR{$n$} د استقرا په وسيله وښيئ چې که \LR{$A_1$}، \LR{$A_2$}، \dots،
  \LR{$A_n$} ټول د شمېر وړ وي، نو \LR{$A_1 \cup \dots \cup A_n$}
  هم داسې دے۔ دا حقيقت فرضولے شئ چې که دوه سټونه \LR{$A$} او~\LR{$B$}
  د شمېر وړ وي، نو~\LR{$A \cup B$} هم داسې دے۔
\end{prob}
\olpseganchor{OLP-0029-B023}{end}

\olpseganchor{OLP-0029-B024}{start}
که څه هم د يوه سټ د غړو د لړ کولو پر وخت
ښايي له \LR{$1$}م غړي څخه شمېر پيلول زيات طبيعي
ښکاري، رياضي پوهان د څيزونو د شمېرلو دپاره طبيعي عددونه~\LR{$\Nat$}
کارول خوښوي۔ دوے د لړ د \LR{$0$}م، \LR{$1$}م، \LR{$2$}م، او داسې نورو
غړو خبرې کوي۔ له دې سره سم، شمېرنه له
\LR{$\Nat$} څخه~\LR{$A$} ته د يوې پر هدف سټ بشپړې تابعې په توګه تعريفولے
شو۔ البته، دواړه تعريفونه همارزښته دي۔
\olpseganchor{OLP-0029-B024}{end}

\olpseganchor{OLP-0029-B025}{start}
\begin{prop}\ollabel{prop:enum-shift}
  يوه پر هدف سټ بشپړه تابع \LR{$f\colon \PosInt \to A$} هله او يوازې هله
  شته چې يوه پر هدف سټ بشپړه تابع \LR{$g\colon \Nat \to A$} شته وي۔
\end{prop}
\olpseganchor{OLP-0029-B025}{end}

\olpseganchor{OLP-0029-B026}{start}
\begin{proof}
  د ورکړل شوې پر هدف سټ بشپړې تابعې \LR{$f\colon \PosInt \to A$} دپاره،
  د ټولو \LR{$n \in \Nat$} دپاره \LR{$g(n) = f(n+1)$} تعريفولے شو۔ په
  اسانه ښکاري چې \LR{$g\colon \Nat \to A$} يوه پر هدف سټ بشپړه تابع
  ده۔ برعکس، د ورکړل شوې پر هدف سټ بشپړې تابعې \LR{$g\colon \Nat \to A$}
  دپاره، \LR{$f(n) = g(n-1)$} تعريف کړئ۔
\end{proof}
\olpseganchor{OLP-0029-B026}{end}

\olpseganchor{OLP-0029-B027}{start}
له دې لاندې پايله ترلاسه کېږي:
\olpseganchor{OLP-0029-B027}{end}

\olpseganchor{OLP-0029-B028}{start}
\begin{cor}\ollabel{cor:enum-nat}
سټ \LR{$A$} هله او يوازې هله د شمېر وړ دے چې تش وي يا يوه
پر هدف سټ بشپړه تابع \LR{$f\colon \Nat \to A$} شته وي۔
\end{cor}
\olpseganchor{OLP-0029-B028}{end}

\olpseganchor{OLP-0029-B029}{start}
پورته مو بحث وکړ چې د سټ~\LR{$A$} د غړو لړ په
بې تکراره لړ اړول کېدے شي۔ دا د شمېرنو دپاره هم رښتيا ده، خو
په کره ډول ئې بيانول او ثابتول لږ ګران دي۔ هره تابع
\LR{$f\colon \PosInt \to A$} بايد د ټولو \LR{$n \in \PosInt$} دپاره
تعريف شوې وي۔ که په~\LR{$A$} کښې يوازې متناهي شمېر غړي وي،
نو په ښکاره داسې تابع نۀ شو لرلے چې د~\LR{$\PosInt$} په نامتناهي
شمېر غړو تعريف شوې وي، د~\LR{$A$} ټول
غړي د قيمتونو په توګه واخلي، خو هېڅکله يو قيمت دوه
ځله وانخلي۔ په دې حالت کښې، يعنې کله چې بې تکراره لړ متناهي
وي، بايد د~\LR{$f$} دپاره د تعريف يوه بېله ساحه وټاکو، چې يوازې
متناهي شمېر غړي ولري۔ د تکرارونو نۀ لرل دا معنٰی
لري چې \LR{$f$} بايد يو پر يو وي۔ ځکه چې پر هدف سټ بشپړه هم ده،
نو د کوم متناهي سټ \LR{$\{1, \dots, n\}$} يا \LR{$\PosInt$} او~\LR{$A$}
ترمنځ يوه دوه اړخيزه يو پر يو تابع لټوو۔
\olpseganchor{OLP-0029-B029}{end}

\olpseganchor{OLP-0029-B030}{start}
\begin{prop}\ollabel{prop:enum-bij}
که \LR{$f\colon \PosInt \to A$} يوه پر هدف سټ بشپړه تابع وي (يعنې
د~\LR{$A$} شمېرنه وي)، نو يوه دوه اړخيزه يو پر يو تابع \LR{$g\colon Z \to A$}
شته چې \LR{$Z$} يا~\LR{$\PosInt$} وي يا د کوم~\LR{$n \in \PosInt$} دپاره
\LR{$\{1, \dots, n\}$} وي۔
\end{prop}
\olpseganchor{OLP-0029-B030}{end}

\olpseganchor{OLP-0029-B031}{start}
\begin{proof}
  تابع \LR{$g$} په بازګشتي ډول تعريفوو: فرض کړئ \LR{$g(1) = f(1)$}۔ که
  \LR{$g(i)$} له مخکښې تعريف شوے وي، نو \LR{$g(i+1)$} د \LR{$f(1)$}،
  \LR{$f(2)$}، \dots{} لومړے هغه قيمت وټاکئ چې تر اوسه په \LR{$g(1)$}،
  \dots، \LR{$g(i)$} کښې نۀ وي، که داسې قيمت شته وي۔ که \LR{$A$} يوازې
  \LR{$n$} غړي ولري، نو \LR{$g(1)$}، \dots، \LR{$g(n)$} ټول تعريف
  شوي دي، او په دې ډول مو يوه تابع \LR{$g\colon \{1, \dots, n\}
  \to A$} تعريف کړې ده۔ که \LR{$A$} نامتناهي شمېر غړي
  ولري، نو د هر \LR{$i$} دپاره بايد د \LR{$f(1)$}، \LR{$f(2)$}، \dots په
  شمېرنه کښې د~\LR{$A$} يو داسې غړے وي چې تر اوسه په
  \LR{$g(1)$}، \dots، \LR{$g(i)$} کښې نۀ وي۔ په دې حالت کښې مو يوه
  تابع \LR{$g\colon \PosInt \to A$} تعريف کړې ده۔
\olpseganchor{OLP-0029-B031}{end}
  
\olpseganchor{OLP-0029-B032}{start}
  تابع \LR{$g$} پر هدف سټ بشپړه ده، ځکه چې د~\LR{$A$} هر غړے په \LR{$f(1)$}،
  \LR{$f(2)$}، \dots{} کښې دے (ځکه چې \LR{$f$} پر هدف سټ بشپړه ده) او له
  دې کبله به بالاخره د کوم~\LR{$i$} دپاره د~\LR{$g(i)$} قيمت شي۔ دا
  يو پر يو هم ده، ځکه چې که داسې \LR{$j < i$} واے چې
  \LR{$g(j) = g(i)$} واے، نو \LR{$g(i)$} به له مخکښې په \LR{$g(1)$}،
  \dots، \LR{$g(i-1)$} کښې و، چې زموږ د~\LR{$g$} له تعريف سره په ټکر کښې دے۔
\end{proof}
\olpseganchor{OLP-0029-B032}{end}

\olpseganchor{OLP-0029-B033}{start}
\begin{cor}\ollabel{cor:enum-nat-bij}
سټ \LR{$A$} هله او يوازې هله د شمېر وړ دے چې تش وي يا يوه
دوه اړخيزه يو پر يو تابع \LR{$f\colon N \to A$} شته وي، چې يا \LR{$N = \Nat$}
وي يا د کوم \LR{$n \in \Nat$} دپاره \LR{$N = \{0, \dots, n\}$} وي۔
\end{cor}
\olpseganchor{OLP-0029-B033}{end}

\olpseganchor{OLP-0029-B034}{start}
\begin{proof}
\LR{$A$} هله او يوازې هله د شمېر وړ دے چې \LR{$A$} تش وي يا يوه
پر هدف سټ بشپړه \LR{$f\colon \PosInt \to A$} شته وي۔ د
\olref{prop:enum-bij} له مخې، وروستۍ خبره هله او يوازې هله
سمه ده چې يوه دوه اړخيزه يو پر يو تابع~\LR{$f\colon Z \to A$} شته وي،
چې \LR{$Z = \PosInt$} يا د کوم \LR{$n \in \PosInt$} دپاره
\LR{$Z = \{1, \dots, n\}$} وي۔ د \olref{prop:enum-shift} د ثبوت
د هماغه استدلال له مخې، دا بيا هله او يوازې هله سمه ده چې
يوه دوه اړخيزه يو پر يو تابع \LR{$g\colon N \to A$} شته وي، چې يا \LR{$N = \Nat$}
وي يا \LR{$N = \{0, \dots, n-1\}$} وي۔
\end{proof}
\olpseganchor{OLP-0029-B034}{end}

\olpseganchor{OLP-0029-B035}{start}
\begin{prob}
  د \olref[sfr][siz][enm]{defn:enumerable} له مخې، سټ \LR{$A$} هله
  او يوازې هله د شمېر وړ دے چې \LR{$A = \emptyset$} وي يا يوه
  پر هدف سټ بشپړه \LR{$f\colon \PosInt \to A$} شته وي۔ «د شمېر وړ
  سټ» په کره ډول داسې هم تعريفولے شو: سټ هله او يوازې هله
  د شمېر وړ دے چې يوه يو پر يو تابع \LR{$g\colon A \to \PosInt$}
  شته وي۔ وښيئ چې تعريفونه همارزښته دي؛ يعنې وښيئ چې يوه
  يو پر يو تابع \LR{$g\colon A \to \PosInt$} هله او يوازې هله
  شته چې يا \LR{$A = \emptyset$} وي يا يوه پر هدف سټ بشپړه
  \LR{$f\colon \PosInt \to A$} شته وي۔
\end{prob}
\olpseganchor{OLP-0029-B035}{end}

\olpseganchor{OLP-0030-B004}{start}
\olfileid{sfr}{siz}{zigzag}
\olpseganchor{OLP-0030-B004}{end}

\olpseganchor{OLP-0030-B005}{start}
\olsection{د کانتور د زيګزاګ طريقه}
\olpseganchor{OLP-0030-B005}{end}

\olpseganchor{OLP-0030-B006}{start}
\begin{explain}
مخکښې مو ځينې «اسانې» شمېرنې په پام کښې نيولې دي۔ اوس به
يو لږ ګران کار په پام کښې ونيسو۔ د طبيعي عددونو د جوړو سټ
په پام کښې ونيسئ، چې په
\olref[set][pai]{sec} کښې مو داسې تعريف کړے و:
\[
\Nat \times \Nat = \Setabs{\tuple{n,m}}{n,m \in \Nat}
\]
دا مرتبې جوړې په يوه \emph{جدولي ترتيب} کښې داسې برابرولے شو:
\[
\begin{array}{ c | c | c | c | c | c}
& \mathbf 0 & \mathbf 1 & \mathbf 2 & \mathbf 3 & \dots \\
\hline
\mathbf 0 & \tuple{0,0} & \tuple{0,1} & \tuple{0,2} & \tuple{0,3} & \dots \\
\hline
\mathbf 1 & \tuple{1,0} & \tuple{1,1} & \tuple{1,2} & \tuple{1,3} & \dots \\
\hline
\mathbf 2 & \tuple{2,0} & \tuple{2,1} & \tuple{2,2} & \tuple{2,3} & \dots \\
\hline
\mathbf 3 & \tuple{3,0} & \tuple{3,1} & \tuple{3,2} & \tuple{3,3} & \dots \\
\hline
\vdots & \vdots & \vdots & \vdots & \vdots & \ddots\\
\end{array}
\]
په ښکاره، په \LR{$\Nat \times \Nat$} کښې هره مرتبه جوړه به په
دې جدول کښې په دقيقه توګه يو ځل راشي۔ په ځانګړي ډول،
\LR{$\tuple{n,m}$} به په \LR{$n$}م قطار او \LR{$m$}م ستون کښې راشي۔ خو د
داسې جدول غړي په يوه «يو بُعدي» لړ کښې څنګه برابر کړو؟ په
لاندې جدول کښې نمونه د دې کار يوه لار ښيي، که څه هم البته
ډېرې نورې لارې هم شته:
\[
\begin{array}{ c | c | c | c | c | c | c}
& \mathbf 0 & \mathbf 1 & \mathbf 2 & \mathbf 3 & \mathbf 4 &\dots \\
\hline
\mathbf 0 & 0  & 1& 3 & 6& 10 &\ldots \\
\hline
\mathbf 1 &2 & 4& 7 & 11 & \dots &\ldots \\
\hline
\mathbf 2 & 5 & 8 & 12 & \ldots & \dots&\ldots \\
\hline
\mathbf 3 & 9 & 13 & \ldots & \ldots & \dots & \ldots \\
\hline
\mathbf 4 & 14 & \ldots & \ldots & \ldots & \dots & \ldots \\
\hline
\vdots & \vdots & \vdots & \vdots & \vdots&\ldots & \ddots\\
\end{array}
\]\noindent
دا نمونه \emph{د کانتور د زيګزاګ طريقه} بلل کېږي۔ دا د
\LR{$\Nat \times \Nat$} شمېرنه داسې کوي:
\[
\tuple{0,0}, \tuple{0,1}, \tuple{1,0}, \tuple{0,2}, \tuple{1,1},
\tuple{2,0}, \tuple{0,3}, \tuple{1,2}, \tuple{2,1}, \tuple{3,0}, \dots
\]
او له دې لاندې خبره ثابتېږي:
\end{explain}
\olpseganchor{OLP-0030-B006}{end}

\olpseganchor{OLP-0030-B007}{start}
\begin{prop}\ollabel{natsquaredenumerable}
\LR{$\Nat \times \Nat$} يو د شمېر وړ سټ دے۔
\end{prop}
\olpseganchor{OLP-0030-B007}{end}

\olpseganchor{OLP-0030-B008}{start}
\begin{proof}
فرض کړئ \LR{$f \colon \Nat \to \Nat\times\Nat$} هر \LR{$k \in \Nat$}
له هغې مرتبې څوګونې \LR{$\tuple{n,m} \in \Nat \times \Nat$} سره
تړي چې \LR{$k$} د کانتور د زيګزاګ په جدول کښې د \LR{$n$}م قطار او
\LR{$m$}م ستون قيمت وي۔
\end{proof}
\olpseganchor{OLP-0030-B008}{end}

\olpseganchor{OLP-0030-B009}{start}
\begin{explain}
دا تخنيک په ښه ډول عمومي کېدے هم شي۔ د بېلګې په توګه، د
طبيعي عددونو د مرتبو درې ګونو د سټ د شمېرنې دپاره ئې
کارولے شو؛ يعنې:
\[
\Nat \times \Nat \times \Nat = \Setabs{\tuple{n,m,k}}{n,m,k \in \Nat}
\]
\LR{$\Nat \times \Nat \times \Nat$} له \LR{$\Nat$} سره د \LR{$\Nat \times \Nat$}
د کارتيزي حاصل ضرب په توګه ګڼو؛ يعنې:
\[
\Nat^3 = (\Nat \times \Nat) \times \Nat =
\Setabs{\tuple{\tuple{n,m},k}}{n, m, k
  \in \Nat }
\]
نو د يوه محور د \LR{$\Nat$} په شمېرنې او د بل محور د \LR{$\Nat^2$}
په شمېرنې په نښه کولو سره، د يوه جدول په وسيله د \LR{$\Nat^3$}
شمېرنه کولے شو:
\[
\begin{array}{ c | c | c | c | c | c}
& \mathbf 0 & \mathbf 1 & \mathbf 2 & \mathbf 3 & \dots \\
\hline
\mathbf{\tuple{0,0}} & \tuple{0,0,0} & \tuple{0,0,1} & \tuple{0,0,2} & \tuple{0,0,3} & \dots \\
\hline
\mathbf{\tuple{0,1}} & \tuple{0,1,0} & \tuple{0,1,1} & \tuple{0,1,2} & \tuple{0,1,3} & \dots \\
\hline
\mathbf{\tuple{1,0}} & \tuple{1,0,0} & \tuple{1,0,1} & \tuple{1,0,2} & \tuple{1,0,3} & \dots \\
\hline
\mathbf{\tuple{0,2}} & \tuple{0,2,0} & \tuple{0,2,1} & \tuple{0,2,2} & \tuple{0,2,3} & \dots\\
\hline
\vdots & \vdots & \vdots & \vdots & \vdots & \ddots \\
\end{array}
\]
نو د کانتور د زيګزاګ طريقې ته ورته طريقې په کارولو سره د~\LR{$\Nat^3$}
شمېرنه هم ترلاسه کولے شو۔ او همداسې مخ په وړاندې تللے شو،
چې د هر طبيعي عدد \LR{$n$} دپاره د \LR{$\Nat^n$} شمېرنې ترلاسه کړو۔ نو لرو:
\end{explain}
\olpseganchor{OLP-0030-B009}{end}

\olpseganchor{OLP-0030-B010}{start}
\begin{prop}
د هر \LR{$n \in \Nat$} دپاره، \LR{$\Nat^n$} يو د شمېر وړ سټ دے۔
\end{prop}
\olpseganchor{OLP-0030-B010}{end}

\olpseganchor{OLP-0030-B011}{start}
\begin{prob}\label{sfr:siz:zigzag:prob:posint-n}
وښيئ چې د هر \LR{$n \in \Nat$} دپاره، \LR{$(\PosInt)^n$} يو د شمېر وړ سټ دے۔
\end{prob}
\olpseganchor{OLP-0030-B011}{end}

\olpseganchor{OLP-0030-B012}{start}
\begin{prob}\label{sfr:siz:zigzag:prob:posint-star}
  وښيئ چې \LR{$(\PosInt)^*$} يو د شمېر وړ سټ دے۔
  \cref{sfr:siz:zigzag:prob:posint-n} فرضولے شئ۔
\end{prob}
\olpseganchor{OLP-0030-B012}{end}

\olpseganchor{OLP-0031-B004}{start}
\olfileid{sfr}{siz}{pai}
\olsection{جوړوونکې تابعې او کوډونه}
\olpseganchor{OLP-0031-B004}{end}

\olpseganchor{OLP-0031-B005}{start}
\begin{explain}
د کانتور د زيګزاګ طريقه د \LR{$\Nat^n$} د شمېر وړ والے په انځوريز
ډول څرګندوي۔ خو راځئ پر هغه جدول تمرکز وکړو چې \LR{$\Nat^2$} ښيي۔
په جدول کښې د زيګزاګ کرښه تعقيبول او مقامونه شمېرل راښيي چې
\LR{$\tuple{1,2}$} له عدد~\LR{$7$} سره تړل شوې ده۔ خو ښه به وي که دا
په لا نېغه توګه حسابولے شو۔ يعنې ښه به وي چې د زيګزاګ شمېرنې
\emph{معکوس} تابع \LR{$g\colon \Nat^2 \to \Nat$} په لاس کښې ولرو،
داسې چې:
\[
g(\tuple{0,0}) = 0, \;
g(\tuple{0,1}) = 1, \;
g(\tuple{1,0}) = 2, \; \dots,  
g(\tuple{1,2}) = 7, \; \dots
\]
په دې سره په دقيق ډول حسابولے شو چې \LR{$\tuple{n, m}$} به زموږ
په شمېرنه کښې چېرته راشي۔
\olpseganchor{OLP-0031-B005}{end}

\olpseganchor{OLP-0031-B006}{start}
په حقيقت کښې، په دوو مشاهدو سره \LR{$g$} په نېغه توګه تعريفولے شو۔
لومړے: که د \LR{$n$}م قطار او \LR{$m$}م ستون قيمت~\LR{$v$} وي، نو د
\LR{$(n+1)$}م قطار او \LR{$(m-1)$}م ستون قيمت \LR{$v + 1$} دے۔ دويم: زموږ
د شمېرنې لومړے قطار له مثلثي عددونو جوړ دے چې له \LR{$0$}، \LR{$1$}،
\LR{$3$}، \LR{$6$} او داسې نورو پيلېږي۔ \LR{$k$}م مثلثي عدد د هغو طبيعي
عددونو جمع ده چې \LR{$< k$} دي، او په \LR{$k(k+1)/2$} حسابېږي۔ د دې دوو
مشاهدو په يوځاے کولو لاندې تابع په پام کښې ونيسئ:
\[
  g(n,m) = \frac{(n+m+1)(n+m)}{2} + n
\]
ډېر وخت د \LR{$g(\tuple{n, m})$} پر ځاے يوازې \LR{$g(n, m)$} ليکو، ځکه
چې سترګو ته اسانه دے۔ دا لومړے د \LR{$(n+m)^\text{\RL{th}}$} مثلثي عدد
ټاکل او بيا \LR{$n$} ورزياتول درته وايي۔ او دا جدول په هماغه طريقه
ډکوي چې غواړو ئې۔ نو په ځانګړي ډول، جوړه \LR{$\tuple{1, 2}$} له
\LR{$\frac{4 \times 3}{2} + 1 = 7$} سره تړل کېږي۔
\olpseganchor{OLP-0031-B006}{end}

\olpseganchor{OLP-0031-B007}{start}
دا تابع \LR{$g$} د جوړو د يوه سټ د شمېرنې \emph{معکوس} ده۔ داسې
تابعې \emph{جوړوونکې تابعې} بلل کېږي۔
\end{explain}
\olpseganchor{OLP-0031-B007}{end}

\olpseganchor{OLP-0031-B008}{start}
\begin{defn}[جوړوونکې تابع]
  تابع \LR{$f\colon A \times B \to \Nat$} يوه حسابي
  \emph{جوړوونکې تابع} ده که \LR{$f$} يو پر يو وي۔ دا هم وايو چې
  \LR{$f$} د \LR{$A \times B$} \emph{کوډول} کوي، او \LR{$f(x,y)$} د
  \LR{$\tuple{x,y}$} \emph{کوډ} دے۔
\end{defn}
\olpseganchor{OLP-0031-B008}{end}

\olpseganchor{OLP-0031-B009}{start}
\begin{explain}
جوړوونکې تابعې د بېلګې په توګه د طبيعي عددونو د جوړو د
کوډولو دپاره کارولے شو؛ يا په بله وينا، د غړو هره \emph{جوړه}
په يوه \emph{واحد} عدد ښودلے شو۔ د جوړوونکې تابعې په معکوس
عدد بېرته لوستلے شو؛ يعنې معلومولے شو چې کومه جوړه ښيي۔
\end{explain}
\olpseganchor{OLP-0031-B009}{end}

\olpseganchor{OLP-0031-B010}{start}
\begin{prob}
د ټولو غير منفي ناطقو عددونو د سټ يوه شمېرنه ورکړئ۔
\end{prob}
\olpseganchor{OLP-0031-B010}{end}

\olpseganchor{OLP-0031-B011}{start}
\begin{prob}
وښيئ چې \LR{$\Rat$} د شمېر وړ دے۔ په ياد ولرئ چې هر ناطق عدد
د کسر \LR{$z/m$} په توګه ليکل کېدے شي، چې \LR{$z \in \Int$} او
\LR{$m \in \Nat^+$} وي۔
\end{prob}
\olpseganchor{OLP-0031-B011}{end}

\olpseganchor{OLP-0031-B012}{start}
\begin{prob}
د \LR{$\Bin^*$} يوه شمېرنه تعريف کړئ۔
\end{prob}
\olpseganchor{OLP-0031-B012}{end}

\olpseganchor{OLP-0031-B013}{start}
\begin{prob}
د منطق له پېژندنيز کورسه په ياد ولرئ چې د رښتياوالي هر ممکن
جدول د رښتياوالي يوه تابع څرګندوي۔ په بله وينا، د کوم~\LR{$k$}
دپاره د رښتياوالي تابعې ټولې تابعې له \LR{$\Bin^k \to \Bin$} څخه
دي۔ ثابته کړئ چې د رښتياوالي د ټولو تابعو سټ د شمېر وړ دے۔
\end{prob}
\olpseganchor{OLP-0031-B013}{end}

\olpseganchor{OLP-0031-B014}{start}
\begin{prob}
وښيئ چې د يوه خپل سري نامتناهي د شمېر وړ سټ د ټولو
متناهي فرعي سټونو سټ د شمېر وړ دے۔
\end{prob}
\olpseganchor{OLP-0031-B014}{end}

\olpseganchor{OLP-0031-B015}{start}
\begin{prob}
د \LR{$\Nat$} يوه فرعي سټ ته هله \emph{د متناهي متمم لرونکے} ويل
کېږي چې متمم ئې په \LR{$\Nat$} کښې متناهي سټ وي؛ يعنې
\LR{$A \subseteq \Nat$} هله د متناهي متمم لرونکے دے چې
\LR{$\Nat\setminus A$} متناهي وي۔ \textbf{د ژباړې د نسخې سمون
(OLSIZ-002):} د انګرېزي سرچينې په تعريف کښې د متناهي فرعي سټ او
طبيعي عددونو د سټ تر منځ اړيکه پاتې وه؛ ورپسې هم‌ارز فارمول ښيي
چې مطلب په طبيعي عددونو کښې متناهي متمم دے۔ فرض کړئ \LR{$I$} هغه سټ دے چې
غړي ئې په دقيقه توګه د \LR{$\Nat$} متناهي او د متناهي
متمم لرونکي فرعي سټونه دي۔ وښيئ چې \LR{$I$} د شمېر وړ دے۔
\end{prob}
\olpseganchor{OLP-0031-B015}{end}

\olpseganchor{OLP-0031-B016}{start}
\begin{prob}
وښيئ چې د د شمېر وړ سټونو د شمېر وړ اتحاد
د شمېر وړ دے۔ يعنې، هرکله چې \LR{$A_1$}، \LR{$A_2$}، \dots{} سټونه
وي او هر \LR{$A_i$} د شمېر وړ وي، نو د هغو ټولو اتحاد
\LR{$\bigcup_{i=1}^\infty A_i$} هم د شمېر وړ دے۔ [يادونه:
دا ګران دے!]
\end{prob}
\olpseganchor{OLP-0031-B016}{end}

\olpseganchor{OLP-0031-B017}{start}
\begin{prob}
فرض کړئ \LR{$f \colon A \times B \to \Nat$} يوه خپله سري جوړوونکې
تابع ده او د جوړو سټ ناتش دے۔ وښيئ چې د~\LR{$f$} معکوس پر خپلې
د تعريف ساحې دوه اړخيزه يو پر يو تابع ده چې د قيمتونو سټ ئې
\LR{$A \times B$} دے، او د هغې له ساحې څخه بهر د يوې ثابتې جوړې
په ورکولو د طبيعي عددونو پر ټولو غړو غځېدے شي۔
\textbf{د ژباړې د نسخې سمون (PSSIZ-001):} انګرېزي تمرين د
معکوس تابع د تعريف ساحه د اصلي تابعې له قيمتونو سټ څخه ټولو
طبيعي عددونو ته نۀ وه غځولې، او د تش جوړو سټ حالت ئې هم نۀ و بېل کړے۔
\end{prob}
\olpseganchor{OLP-0031-B017}{end}

\olpseganchor{OLP-0031-B018}{start}
\begin{prob}
داسې تابع وټاکئ چې \LR{$\Nat^3$} کوډ کړي۔
\end{prob}
\olpseganchor{OLP-0031-B018}{end}

\olpseganchor{OLP-0032-B004}{start}
\olfileid{sfr}{siz}{pai-alt}
\olpseganchor{OLP-0032-B004}{end}

\olpseganchor{OLP-0032-B005}{start}
\olsection{يوه بله جوړوونکې تابع}
\olpseganchor{OLP-0032-B005}{end}

\olpseganchor{OLP-0032-B006}{start}
\begin{explain}
د \LR{$\Nat^2$} داسې نورې شمېرنې هم شته چې د هغو د معکوس تابعو
معلومول اسانه کوي۔ يوه ئې دا ده۔ د شمېرنې پر جدولي ترتيب د
انځورولو پر ځاے، د مثبتو صحيح عددونو له داسې لړ څخه پيل وکړئ
چې په سر کښې له تشو ځايونو سره تړلے وي۔ وانګېړئ چې دا ځايونه
يو په بل پسې په لاندې ډول په جوړو \LR{$\tuple{n,m}$} ډکوو۔ لومړے
له هغو جوړو پيل وکړئ چې په لومړي مقام کښې~\LR{$0$} لري (يعنې جوړې
\LR{$\tuple{0,m}$})؛ لومړۍ جوړه (يعنې \LR{$\tuple{0,0}$}) په لومړي تش
مقام کښې کېږدئ، بيا يو تش مقام پرېږدئ، دويمه جوړه (يعنې
\LR{$\tuple{0,1}$}) په ورپسې تش مقام کښې کېږدئ، بيا يو مقام پرېږدئ،
او همداسې مخکښې لاړ شئ۔
\textbf{د ژباړې د نسخې سمون (PSSIZ-002):} انګرېزي نثر دويمه
جوړه په تېروتنه د درېيمې جوړې په نښه ښيي؛ لاندې اصلي جدول ئې
سمه، دويمه جوړه څرګندوي۔ زموږ د شمېرنې (نيمګړے) پيل اوس داسې دے:
\[\small
\begin{array}{@{}c c c c c c c c c c c@{}}
\mathbf 1 & \mathbf 2 & \mathbf 3 & \mathbf 4 & \mathbf 5 & \mathbf 6 & \mathbf 7 & \mathbf 8 & \mathbf 9 & \mathbf{10} & \dots \\ \\
\tuple{0,0} &  & \tuple{0,1} &  & \tuple{0,2} &  & \tuple{0,3} & & \tuple{0,4} &  & \dots \\
\end{array}
\]
همدا کار د \LR{$\tuple{1,m}$} جوړو دپاره په هغو مقامونو کښې تکرار
کړئ چې لا تش پاتې دي؛ بيا هم هر بل تش مقام پرېږدئ:
\[\small
\begin{array}{@{}c c c c c c c c c c c@{}}
\mathbf 1 & \mathbf 2 & \mathbf 3 & \mathbf 4 & \mathbf 5 & \mathbf 6 & \mathbf 7 & \mathbf 8 & \mathbf 9 & \mathbf{10} & \dots \\ \\
\tuple{0,0} & \tuple{1,0} & \tuple{0,1} &  & \tuple{0,2} & \tuple{1,1} & 
\tuple{0,3} & & \tuple{0,4} &  \tuple{1,2} & \dots \\
\end{array}
\]
په همدې طريقه \LR{$\tuple{2,m}$}، \LR{$\tuple{3,m}$} او داسې نورې جوړې
ور دننه کړئ۔
\textbf{د ژباړې د نسخې سمون (OLSIZ-003):} انګرېزي نثر د دوو
پرله‌پسې قطارونو پر ځاے د يوه قطار نښه دوه ځله تکراروي؛ بشپړ
اصلي جدول د پرله‌پسې قطارونو لوستل تاييدوي۔ نو زموږ بشپړه شمېرنه
داسې پيلېږي:
\[\small
\begin{array}{@{}cc c c c c c c c c c@{}}
\mathbf 1 & \mathbf 2 & \mathbf 3 & \mathbf 4 & \mathbf 5 & \mathbf 6 & \mathbf 7 & \mathbf 8 & \mathbf 9 & \mathbf{10} & \dots \\ \\
\tuple{0,0} & \tuple{1,0} & \tuple{0,1} & \tuple{2,0}  & \tuple{0,2} & 
\tuple{1,1} & \tuple{0,3} & \tuple{3,0}  & \tuple{0,4} &  \tuple{1,2} & \dots \\
\end{array}
\]
که د پورته جدول حجرې د همدې شمېرنې له مخې وشمېرو، د زيګزاګ
پاکه کرښه به و نۀ مومو، بلکې دا ترتيب به ومومو:
\[
\begin{array}{ c | c | c | c | c | c | c | c }
& \mathbf 0 & \mathbf 1 & \mathbf 2 & \mathbf 3 & \mathbf 4 & \mathbf 5 & \dots \\
\hline
\mathbf 0 & 1 & 3 & 5 & 7 & 9 & 11 & \dots \\
\hline
\mathbf 1 & 2 & 6 & 10 & 14 & 18 & \dots & \dots \\
\hline
\mathbf 2 & 4 & 12 & 20 & 28 & \dots & \dots & \dots \\
\hline
\mathbf 3 & 8 & 24 & 40 & \dots & \dots & \dots & \dots \\
\hline
\mathbf 4 & 16 & 48 & \dots & \dots & \dots & \dots & \dots \\
\hline
\mathbf 5 & 32 & \dots & \dots & \dots & \dots & \dots & \dots \\
\hline
\vdots & \vdots & \vdots & \vdots & \vdots & \vdots & \vdots & \ddots\\
\end{array}
\]
وينو چې په~\LR{$0$} قطار کښې جوړې زموږ د شمېرنې په طاق شمېره
لرونکو مقامونو کښې دي؛ يعنې جوړه \LR{$\tuple{0,m}$} په \LR{$2m+1$} مقام
کښې ده۔ د دويم قطار جوړې، \LR{$\tuple{1,m}$}، په هغو مقامونو کښې دي
چې شمېر ئې د يوه طاق عدد دوه چنده دے، په ځانګړي ډول
\LR{$2 \cdot (2m+1)$}؛ د درېيم قطار جوړې، \LR{$\tuple{2,m}$}، په هغو
مقامونو کښې دي چې شمېر ئې د يوه طاق عدد څلور چنده،
\LR{$4 \cdot (2m+1)$}، دے؛ او همداسې نور۔ په هر قطار کښې د
\LR{$(2m+1)$} ضربي عوامل، \LR{$1$}، \LR{$2$}، \LR{$4$}، \LR{$8$}، \dots، په دقيقه توګه
د~\LR{$2$} طاقتونه دي: \LR{$1= 2^0$}، \LR{$2 = 2^1$}، \LR{$4 = 2^2$}،
\LR{$8 = 2^3$}، \dots\@ په حقيقت کښې، اړوند توان ښودونکے تل د پام
وړ جوړې لومړے غړے وي۔ نو د جوړې \LR{$\tuple{n,m}$} دپاره عامل
\LR{$2^n$} دے۔ په دې توګه عمومي فارمول لاس ته راځي:
\LR{$2^n \cdot (2m+1)$}۔ خو دا له جوړو څخه \emph{مثبتو} صحيح عددونو
ته نقشه ده؛ يعنې \LR{$\tuple{0,0}$} مقام~\LR{$1$} لري۔ که وغواړو له
مقام~\LR{$0$} څخه پيل وکړو، بايد له پايلې څخه~\LR{$1$} منفي کړو۔ نو دا
لاس ته راوړو:
\end{explain}
\olpseganchor{OLP-0032-B006}{end}

\olpseganchor{OLP-0032-B007}{start}
\begin{ex}
تابع \LR{$h\colon \Nat^2 \to \Nat$} چې په لاندې ډول ورکړل شوې ده،
\[
h(n,m) = 2^n (2m+1) - 1
\]
د طبيعي عددونو د جوړو د سټ~\LR{$\Nat^2$} يوه جوړوونکې تابع ده۔
\end{ex}
\olpseganchor{OLP-0032-B007}{end}

\olpseganchor{OLP-0032-B008}{start}
\begin{explain}
له دې سره سم، زموږ د \LR{$\Nat^2$} په دويمه شمېرنه کښې جوړه
\LR{$\tuple{0,0}$} کوډ \LR{$h(0,0) = 2^0(2\cdot 0+1) - 1 = 0$} لري؛
\LR{$\tuple{1,2}$} کوډ \LR{$2^{1} \cdot (2 \cdot 2 + 1) - 1 = 2
\cdot 5 - 1 = 9$} لري؛ او \LR{$\tuple{2,6}$} کوډ \LR{$2^{2} \cdot (2
\cdot 6 + 1) - 1 = 51$} لري۔
\end{explain}
\olpseganchor{OLP-0032-B008}{end}

\olpseganchor{OLP-0032-B009}{start}
ځينې وخت د طبيعي عددونو د جوړو~\LR{$\Nat^2$} کوډول دومره بس وي،
بې له دې چې کوډونه پر هدف سټ بشپړ وي۔ د داسې کوډولو معکوسې
تابعې يوازې جزوي تابعې وي۔
\olpseganchor{OLP-0032-B009}{end}

\olpseganchor{OLP-0032-B010}{start}
\begin{ex}
تابع \LR{$j\colon \Nat^2 \to \Nat^+$} چې په لاندې ډول ورکړل شوې ده،
\[
j(n,m) = 2^n3^m
\]
يوه يو پر يو تابع \LR{$\Nat^2 \to \Nat$} ده۔
\end{ex}
\olpseganchor{OLP-0032-B010}{end}

\olpseganchor{OLP-0033-B004}{start}
\olfileid{sfr}{siz}{nen}
\olpseganchor{OLP-0033-B004}{end}

\olpseganchor{OLP-0033-B005}{start}
\olsection{د شمېر نۀ وړ سټونه}
\olpseganchor{OLP-0033-B005}{end}

\olpseganchor{OLP-0033-B006}{start}
\begin{editorial}
  دا برخه په \olref[enm]{sec} کښې د ورکړل شوي تعريف په کارولو د
  \LR{$\Bin^\omega$} او \LR{$\Pow{\PosInt}$} د شمېر نۀ وړ والے ثابتوي۔ دا
  د \olref[enm-alt]{sec} له بڼې څخه لږه ابتدايي او لږه زياته
  تفصيلي جوړه شوې ده۔
\end{editorial}
\olpseganchor{OLP-0033-B006}{end}

\olpseganchor{OLP-0033-B007}{start}
ځينې سټونه، لکه د مثبتو صحيح عددونو سټ~\LR{$\PosInt$}، نامتناهي دي۔
تر اوسه مو د نامتناهي سټونو داسې بېلګې ليدلې چې ټولې
د شمېر وړ وې۔ خو داسې نامتناهي سټونه هم شته چې دا خاصيت نۀ
لري۔ داسې سټونه \emph{د شمېر نۀ وړ} بلل کېږي۔
\olpseganchor{OLP-0033-B007}{end}

\olpseganchor{OLP-0033-B008}{start}
تر هر څه لومړے ښايي همدا خبره لا حيرانوونکې وي چې د شمېر نۀ وړ
سټونه شته۔ د هر د شمېر وړ سټ~\LR{$A$} دپاره يوه پر هدف سټ بشپړه
تابع \LR{$f \colon \PosInt \to A$} شته۔ که يو سټ د شمېر نۀ وړ وي،
داسې تابع نۀ شته۔ يعنې هېڅ هغه تابع چې د~\LR{$\PosInt$} بې‌شمېره
غړي پر~\LR{$A$} نقشوي، ټول~\LR{$A$} نۀ شي پوره کولے۔ نو په~\LR{$A$}
کښې د بې‌شمېره مثبتو صحيح عددونو په پرتله «زيات» غړي شته۔
\olpseganchor{OLP-0033-B008}{end}

\olpseganchor{OLP-0033-B009}{start}
څنګه به ثابتوو چې يو سټ د شمېر نۀ وړ دے؟ بايد وښيئ چې داسې
پر هدف سټ بشپړه تابع وجود نۀ شي لرلے۔ په مساوي ډول، بايد وښيئ
چې د~\LR{$A$} غړي په يو اړخيز نامتناهي لړ کښې نۀ شي شمېرل کېدے۔ د
دې تر ټولو ښه لار دا ده چې وښيئ د~\LR{$A$} د غړي هر لړ بايد
لږ تر لږه يو غړے پرېږدي؛ يا دا چې هېڅ تابع
\LR{$f\colon \PosInt \to A$} پر هدف سټ بشپړه نۀ شي کېدے۔ دا د کانتور
په \emph{قطري طريقه} کولے شو۔ د~\LR{$A$} د غړي يو لړ، د بېلګې
په ډول \LR{$x_1$}، \LR{$x_2$}، \dots، چې راکړل شي، موږ د~\LR{$A$} يو بل غړے
جوړوو چې د خپل جوړښت له مخې په هېڅ ډول په هغه لړ کښې نۀ شي راتلے۔
\olpseganchor{OLP-0033-B009}{end}

\olpseganchor{OLP-0033-B010}{start}
زموږ لومړۍ بېلګه د ټولو هغو بې‌تشې نامتناهي تسلسلونو سټ
\LR{$\Bin^\omega$} دے چې له \LR{$0$}ګانو او \LR{$1$}ګانو جوړ دي۔
\olpseganchor{OLP-0033-B010}{end}

\olpseganchor{OLP-0033-B011}{start}
\begin{thm}
\ollabel{thm:nonenum-bin-omega}
\LR{$\Bin^\omega$}~د شمېر نۀ وړ دے۔
\end{thm}
\olpseganchor{OLP-0033-B011}{end}

\olpseganchor{OLP-0033-B012}{start}
\begin{proof}
د تناقض دپاره فرض کړئ چې \LR{$\Bin^\omega$} د شمېر وړ دے؛ يعنې
فرض کړئ چې د~\LR{$\Bin^\omega$} د ټولو غړي يو لړ \LR{$s_{1}$}،
\LR{$s_{2}$}، \LR{$s_{3}$}، \LR{$s_{4}$}، \dots{} شته۔ له دې هر يو \LR{$s_i$} پخپله
د \LR{$0$}ګانو او~\LR{$1$}ګانو يو نامتناهي تسلسل دے۔ د دې لړ د \LR{$i$}م تسلسل
\LR{$j$}م غړے به \LR{$s_i(j)$} وبولو۔ نو \LR{$i$}م تسلسل~\LR{$s_i$} دا دے:
\[
s_i(1), s_i(2), s_i(3), \dots
\]
\olpseganchor{OLP-0033-B012}{end}

\olpseganchor{OLP-0033-B013}{start}
دا لړ، او په کښې د هر تسلسل~\LR{$s_i$} غړي، په يوه جدولي ترتيب کښې
اېښودلے شو:
\[
\begin{array}{c|c|c|c|c|c}
& 1 & 2 & 3 & 4 & \dots \\\hline
1 & \mathbf{s_{1}(1)} & s_{1}(2) & s_{1}(3) & s_1(4) & \dots \\\hline
2 & s_{2}(1)& \mathbf{s_{2}(2)} & s_2(3) & s_2(4) & \dots \\\hline
3 & s_{3}(1)& s_{3}(2) & \mathbf{s_3(3)} & s_3(4) & \dots \\\hline
4 & s_{4}(1)& s_{4}(2) & s_4(3) & \mathbf{s_4(4)} & \dots \\\hline
\vdots & \vdots & \vdots & \vdots & \vdots & \mathbf{\ddots}
\end{array}
\]
د څنګ پر خوا نښې په لړ کښې د تسلسلونو شمېرې \LR{$s_1$}، \LR{$s_2$}،
\dots راکوي؛ د پاسنۍ خوا شمېرې د هر يوه تسلسل غړي په
نښه کوي۔ د بېلګې په ډول، \LR{$s_{1}(1)$} د هغه عدد نوم دے چې، که \LR{$0$}
وي يا~\LR{$1$}، په تسلسل~\LR{$s_{1}$} کښې لومړے غړے دے؛ او همداسې
نور۔
\olpseganchor{OLP-0033-B013}{end}

\olpseganchor{OLP-0033-B014}{start}
اوس د \LR{$0$}ګانو او \LR{$1$}ګانو يو داسې نامتناهي تسلسل
\LR{$\overline{s}$} جوړوو چې په هېڅ ډول په دې لړ کښې نۀ شي راتلے۔ د
\LR{$\overline{s}$} تعريف به په لړ \LR{$s_1$}، \LR{$s_2$}، \dots پورې تړلے وي۔
د \LR{$0$}ګانو او \LR{$1$}ګانو د نامتناهي تسلسلونو هر نامتناهي لړ يو داسې
نامتناهي تسلسل~\LR{$\overline{s}$} راپيدا کوي چې په ډاډه توګه په هغه
لړ کښې نۀ راځي۔
\olpseganchor{OLP-0033-B014}{end}

\olpseganchor{OLP-0033-B015}{start}
د \LR{$\overline{s}$} د تعريفولو دپاره د هغه ټول غړي ټاکو؛
يعنې د ټولو \LR{$n \in \PosInt$} دپاره \LR{$\overline{s}(n)$} ټاکو۔ دا کار
د پورته جدولي ترتيب پر قطر د لاندې تللو (ځکه ورته «قطري طريقه»
وايو)، بيا د هر \LR{$1$} په \LR{$0$} او د هر \LR{$0$} په~\LR{$1$} بدلولو کوو۔ په لا
مجرد ډول، \LR{$\overline{s}(n)$} د دې له مخې \LR{$0$} يا \LR{$1$} تعريفوو چې د
قطر \LR{$n$}م غړے، \LR{$s_n(n)$}، \LR{$1$} دے که \LR{$0$}:
\[
\overline{s}(n) =
\begin{cases}
1 & \text{\RL{که \LR{$s_{n}(n) = 0$} وي}}\\
0 & \text{\RL{که \LR{$s_{n}(n) = 1$} وي}}.
\end{cases}
\]
که فارمولونه مو د حالتونو تر تعريف زيات خوښېږي، دا هم تعريفولے
شئ چې \LR{$\overline{s}(n) = 1 - s_n(n)$}۔
\olpseganchor{OLP-0033-B015}{end}

\olpseganchor{OLP-0033-B016}{start}
په څرګند ډول \LR{$\overline{s}$} د \LR{$0$}ګانو او \LR{$1$}ګانو يو نامتناهي
تسلسل دے، ځکه دا يوازې د هغو \LR{$0$}ګانو او \LR{$1$}ګانو مخالف تسلسل دے
چې زموږ د جدولي ترتيب پر قطر راځي۔ نو \LR{$\overline{s}$} د
~\LR{$\Bin^\omega$} يو غړے دے۔ خو دا په لړ \LR{$s_1$}، \LR{$s_2$}،
\dots{} کښې نۀ شي راتلے۔ ولې؟
\olpseganchor{OLP-0033-B016}{end}

\olpseganchor{OLP-0033-B017}{start}
دا په لړ کښې لومړے تسلسل~\LR{$s_1$} نۀ شي کېدے، ځکه په لومړي
غړے کښې له \LR{$s_1$} څخه توپير لري۔ \LR{$s_1(1)$} چې هر څه وي،
موږ \LR{$\overline{s}(1)$} د هغه مخالف تعريف کړے دے۔ دا په لړ کښې
دويم تسلسل نۀ شي کېدے، ځکه \LR{$\overline{s}$} په دويم غړي کښې له
\LR{$s_2$} څخه توپير لري: که \LR{$s_2(2)$}، \LR{$0$} وي، \LR{$\overline{s}(2)$}،
\LR{$1$} دے، او برعکس۔ او همداسې نور۔
\olpseganchor{OLP-0033-B017}{end}

\olpseganchor{OLP-0033-B018}{start}
په لا دقيقه توګه: که \LR{$\overline{s}$} په لړ کښې وے، يو \LR{$k$} به
داسې وے چې \LR{$\overline{s} = s_{k}$}۔ دوه تسلسلونه هله او يوازې
هله يو شان وي چې په هر مقام کښې سره سمون ولري؛ يعنې د هر~\LR{$n$}
دپاره \LR{$\overline{s}(n) = s_{k}(n)$} وي۔ نو په ځانګړي ډول چې
\LR{$n = k$} واخلو، بايد \LR{$\overline{s}(k) = s_{k}(k)$} وي۔ \LR{$s_k(k)$}
يا \LR{$0$} دے يا~\LR{$1$}۔ که \LR{$0$} وي، \LR{$\overline{s}(k)$} بايد~\LR{$1$} وي؛
همداسې مو \LR{$\overline{s}$} تعريف کړے دے۔ خو که \LR{$s_k(k) = 1$} وي،
بيا هم د \LR{$\overline{s}$} د تعريف له مخې \LR{$\overline{s}(k) = 0$}۔
په هر حالت کښې \LR{$\overline{s}(k) \neq s_{k}(k)$}۔
\olpseganchor{OLP-0033-B018}{end}

\olpseganchor{OLP-0033-B019}{start}
موږ د~\LR{$\Bin^\omega$} د غړي د يوه لړ \LR{$s_1$}، \LR{$s_2$}،
\dots{} په فرضولو پيل وکړ۔ له دې لړ څخه مو تسلسل~\LR{$\overline{s}$}
جوړ کړ چې ثابته مو کړه په لړ کښې نۀ شي راتلے۔ خو که ټول \LR{$s_i$}
د \LR{$0$}ګانو او \LR{$1$}ګانو تسلسلونه وي، دا خامخا د \LR{$0$}ګانو او
\LR{$1$}ګانو يو تسلسل دے؛ يعنې \LR{$\overline{s} \in \Bin^\omega$}۔ دا په
ځانګړي ډول ښيي چې د~\LR{$\Bin^\omega$} د \emph{ټولو} غړي هېڅ
لړ نۀ شي کېدے، ځکه د هر داسې لړ دپاره يو داسې تسلسل
~\LR{$\overline{s}$} هم جوړولے شو چې په ډاډه توګه په لړ کښې نۀ راځي؛
نو دا فرض چې په~\LR{$\Bin^\omega$} کښې د ټولو تسلسلونو يو لړ شته،
تناقض ته رسېږي۔
\end{proof}
\olpseganchor{OLP-0033-B019}{end}

\olpseganchor{OLP-0033-B020}{start}
\begin{explain}
د ثبوت دې طريقې ته «قطرول» ويل کېږي، ځکه د~\LR{$\overline{s}$} د
تعريف دپاره د جدولي ترتيب قطر کاروي۔ په قطرولو کښې د جدول شتون
لازم نۀ دے: په ورته مفکوره ښودلے شو چې سټونه د شمېر وړ نۀ
دي، ان هغه وخت چې هېڅ جدول او رښتينے قطر نۀ وي۔
\end{explain}
\olpseganchor{OLP-0033-B020}{end}

\olpseganchor{OLP-0033-B021}{start}
\begin{thm}
\ollabel{thm:nonenum-pownat}
\LR{$\Pow{\PosInt}$} د شمېر وړ نۀ دے۔
\end{thm}
\olpseganchor{OLP-0033-B021}{end}

\olpseganchor{OLP-0033-B022}{start}
\begin{proof}
په هماغه طريقه مخکښې ځو: ښيو چې د~\LR{$\PosInt$} د فرعي سټونو د هر
لړ دپاره د \LR{$\PosInt$} يو داسې فرعي سټ شته چې په هغه لړ کښې نۀ
شي راتلے۔ فرض کړئ دا د~\LR{$\PosInt$} د فرعي سټونو يو ورکړل شوے لړ دے:
\[
Z_{1}, Z_{2}, Z_{3}, \dots
\]
اوس يو سټ \LR{$\overline{Z}$} داسې تعريفوو چې د هر \LR{$n \in \PosInt$}
دپاره، \LR{$n \in \overline{Z}$} هله او يوازې هله چې \LR{$n \notin Z_{n}$}:
\[
\overline{Z} = \Setabs{n \in \PosInt}{n \notin Z_n}
\]
\LR{$\overline{Z}$} په څرګند ډول د مثبتو صحيح عددونو يو سټ دے، ځکه
د فرض له مخې هر~\LR{$Z_n$} داسې دے؛ او له دې امله
\LR{$\overline{Z} \in \Pow{\PosInt}$}۔ خو \LR{$\overline{Z}$} په لړ کښې نۀ
شي راتلے۔ د دې د ښودلو دپاره به ثابته کړو چې د هر
\LR{$k \in \PosInt$} دپاره \LR{$\overline{Z} \neq Z_k$}۔
\olpseganchor{OLP-0033-B022}{end}

\olpseganchor{OLP-0033-B023}{start}
نو \LR{$k \in \PosInt$} خپله سرے ونيسئ۔ موږ \LR{$\overline{Z}$} داسې تعريف
کړے چې د هر \LR{$n \in \PosInt$} دپاره، \LR{$n \in \overline{Z}$} هله او
يوازې هله چې \LR{$n \notin Z_n$}۔ په ځانګړي ډول چې \LR{$n=k$} واخلو،
\LR{$k \in \overline{Z}$} هله او يوازې هله چې \LR{$k \notin Z_k$}۔ خو دا
ښيي چې \LR{$\overline{Z} \neq Z_k$}، ځکه \LR{$k$} د يوه سټ غړے
دے خو د بل نۀ دے، او له دې امله \LR{$\overline{Z}$} او \LR{$Z_k$} بېل
غړي لري۔ ځکه چې \LR{$k$} خپله سرے و، \LR{$\overline{Z}$} په لړ
\LR{$Z_1$}، \LR{$Z_2$}، \dots کښې نۀ دے۔
\end{proof}
\olpseganchor{OLP-0033-B023}{end}

\olpseganchor{OLP-0033-B024}{start}
\begin{explain}
تېر ثبوت د قطر نوم وانۀ خيست، خو که داسې ئې انځور کړئ، قطري
طريقه په کښې ليدلے شئ: وانګېړئ چې سټونه \LR{$Z_1$}، \LR{$Z_2$}، \dots
په يوه جدولي ترتيب کښې ليکل شوي، چې هر غړے~\LR{$j \in Z_i$}
په~\LR{$j$}م ستون کښې ليکل کېږي۔ فرض کړئ په هغه لړ کښې لومړي څلور
سټونه \LR{$\{1,2,3,\dots\}$}، \LR{$\{2, 4, 6, \dots\}$}، \LR{$\{1,2,5\}$} او
\LR{$\{3,4,5,\dots\}$} دي۔ نو جدولي ترتيب به داسې پيل شي:
\[
\begin{array}{r@{}rrrrrrr}
  Z_1 = \{ & \mathbf{1}, & 2, & 3, & 4, & 5, & 6, & \dots\}\\
  Z_2 = \{ &  & \mathbf{2}, &  & 4, &  & 6, & \dots\}\\
  Z_3 = \{ & 1, & 2, &  &  & 5\phantom{,} &  & \}\\
  Z_4 = \{ &  &  & 3, & \mathbf{4}, & 5, & 6, & \dots\}\\
  \vdots & & & & & \ddots
\end{array}
\]
بيا \LR{$\overline{Z}$} هغه سټ دے چې د قطر پر لاندې تلو جوړېږي: هغه
شمېرې پرېږدو چې پر قطر راځي، او هغه \LR{$j$}ګانې ور شاملوو چې د جدول
په \LR{$j$}م قطار/ستون کښې تش ځاے لري۔ په پورته حالت کښې به \LR{$1$} او
\LR{$2$} پرېږدو،~\LR{$3$} به ور شاملوو،~\LR{$4$} به پرېږدو، او همداسې نور۔
\end{explain}
\olpseganchor{OLP-0033-B024}{end}

\olpseganchor{OLP-0033-B025}{start}
\begin{prob}
په قطري استدلال وښيئ چې \LR{$\Pow{\Nat}$} د شمېر نۀ وړ دے۔
\end{prob}
\olpseganchor{OLP-0033-B025}{end}

\olpseganchor{OLP-0033-B026}{start}
\begin{prob}\label{sfr:siz:nen:prob:f-posint}
په څرګند قطري استدلال وښيئ چې د تابعو سټ
\LR{$f \colon \PosInt \to \PosInt$} د شمېر نۀ وړ دے۔ يعنې وښيئ
چې که \LR{$f_1$}، \LR{$f_2$}، \dots د تابعو يو لړ وي او هره
\LR{$f_i\colon \PosInt \to \PosInt$} وي، نو يوه
\LR{$\overline{f}\colon \PosInt \to \PosInt$} شته چې په دې لړ کښې نۀ ده۔
\end{prob}
\olpseganchor{OLP-0033-B026}{end}

\olpseganchor{OLP-0034-B004}{start}
\olfileid{sfr}{siz}{red}
\olpseganchor{OLP-0034-B004}{end}

\olpseganchor{OLP-0034-B005}{start}
\olsection{راکمونه}
\olpseganchor{OLP-0034-B005}{end}

\olpseganchor{OLP-0034-B006}{start}
\begin{editorial}
  دا برخه په راکمونه د شمېر نۀ وړ والے ثابتوي، او پايلې ئې د
  \olref[nen]{sec} له پايلو سره سمون خوري۔ يوه بله، لږه لنډه
  بڼه چې د \olref[nen-alt]{sec} له پايلو سره سمون خوري، په
  \olref[red-alt]{sec} کښې ورکړل شوې ده۔
\end{editorial}
\olpseganchor{OLP-0034-B006}{end}

\olpseganchor{OLP-0034-B007}{start}
موږ په قطري استدلال وښودل چې \LR{$\Pow{\PosInt}$} د شمېر نۀ وړ
دے۔ له دې مخکښې مو ثابته کړې وه چې \LR{$\Bin^\omega$}، يعنې د
\LR{$0$}ګانو او \LR{$1$}ګانو د ټولو نامتناهي تسلسلونو سټ،
د شمېر نۀ وړ دے۔ دا د \LR{$\Pow{\PosInt}$} د د شمېر نۀ وړ
کېدو د ثابتولو يوه بله لار ده: وښيئ چې \emph{که \LR{$\Pow{\PosInt}$}
د شمېر وړ وي، نو \LR{$\Bin^\omega$} هم د شمېر وړ دے}۔ ځکه
پوهېږو چې \LR{$\Bin^\omega$} د شمېر وړ نۀ دے، \LR{$\Pow{\PosInt}$}
هم داسې نۀ شي کېدے۔ دې ته د يوې مسئلې بلې ته \emph{راکمونه}
ويل کېږي؛ په دې حالت کښې د \LR{$\Bin^\omega$} د شمېرلو مسئله د
\LR{$\Pow{\PosInt}$} د شمېرلو مسئلې ته راکموو۔ د وروستۍ مسئلې يو حل،
يعنې د \LR{$\Pow{\PosInt}$} يوه شمېرنه، به د لومړۍ مسئلې حل، يعنې د
\LR{$\Bin^\omega$} يوه شمېرنه، راکړي۔
\olpseganchor{OLP-0034-B007}{end}

\olpseganchor{OLP-0034-B008}{start}
د يوه سټ~\LR{$B$} د شمېرلو مسئله څنګه د يوه سټ~\LR{$A$} د شمېرلو مسئلې
ته راکموو؟ د~\LR{$A$} د يوې شمېرنې د~\LR{$B$} په شمېرنه د اړولو يوه لار
ورکوو۔ د دې تر ټولو اسانه لار د يوې پر هدف سټ بشپړه تابع
\LR{$f\colon A \to B$} تعريفول دي۔ که \LR{$x_1$}، \LR{$x_2$}، \dots{} د~\LR{$A$}
شمېرنه وي، نو \LR{$f(x_1)$}، \LR{$f(x_2)$}، \dots{} به د~\LR{$B$} شمېرنه وي۔
زموږ په حالت کښې د يوې پر هدف سټ بشپړې تابع
\LR{$f\colon \Pow{\PosInt} \to \Bin^\omega$} په لټه کښې يو۔
\olpseganchor{OLP-0034-B008}{end}

\olpseganchor{OLP-0034-B009}{start}
\begin{prob}
وښيئ چې که يوه يو پر يو تابع \LR{$g\colon B \to A$} وي او
\LR{$B$}~د شمېر نۀ وړ وي، نو \LR{$A$} هم همداسې دے۔ دا په ښودلو
وکړئ چې څنګه د~\LR{$g$} په کارولو د~\LR{$A$} يوه شمېرنه د~\LR{$B$} په شمېرنه
اړولے شئ۔
\end{prob}
\olpseganchor{OLP-0034-B009}{end}

\olpseganchor{OLP-0034-B010}{start}
\begin{proof}[په راکمونه د {\olref[nen]{thm:nonenum-pownat}} ثبوت]
فرض کړئ چې \LR{$\Pow{\PosInt}$} د شمېر وړ وے، او له دې امله د
هغه يوه شمېرنه \LR{$Z_{1}$}، \LR{$Z_{2}$}، \LR{$Z_{3}$}، \dots وے۔
\olpseganchor{OLP-0034-B010}{end}

\olpseganchor{OLP-0034-B011}{start}
تابع \LR{$f \colon \Pow{\PosInt} \to \Bin^\omega$} داسې تعريف کړئ
چې \LR{$f(Z)$} هغه تسلسل \LR{$s$} وي چې \LR{$s(n) = 1$} هله او يوازې هله
چې \LR{$n \in Z$} وي، او په بل حالت کښې \LR{$s(n) = 0$} وي۔
\textbf{د ژباړې د نسخې سمون (OLSIZ-004):} په انګرېزي نثر کښې
د تسلسل بې تړاوه شاخص د خپل سري ورودي سټ له نښې سره بدل شو،
څو د تابع تعريف روښانه وي۔ دا په څرګند ډول تابع تعريفوي، ځکه
کله چې \LR{$Z \subseteq \PosInt$} وي، هر \LR{$n \in \PosInt$} يا د \LR{$Z$}
يو غړے وي يا نۀ وي۔ د بېلګې په ډول، د مثبتو جفتو
عددونو سټ \LR{$2\PosInt = \{2, 4, 6, \dots\}$} تسلسل
\LR{$010101\dots$} ته نقشېږي، تش سټ \LR{$0000\dots$} ته او پخپله سټ
\LR{$\PosInt$}، \LR{$1111\dots$} ته نقشېږي۔
\olpseganchor{OLP-0034-B011}{end}

\olpseganchor{OLP-0034-B012}{start}
دا پر هدف سټ بشپړه هم ده: د \LR{$0$}ګانو او \LR{$1$}ګانو هر تسلسل د مثبتو
صحيح عددونو له يوه سټ سره سمون خوري، يعنې له هغه سټ سره چې غړي
ئې هغه صحيح عددونه دي کوم چې د تسلسل د~\LR{$1$} لرونکو مقامونو سره
سمون خوري۔ په لا دقيقه توګه، فرض کړئ \LR{$s \in \Bin^\omega$}۔
\LR{$Z \subseteq \PosInt$} په لاندې ډول تعريف کړئ:
\[
Z = \Setabs{n \in \PosInt}{s(n) = 1}
\]
بيا \LR{$f(Z) = s$}، لکه چې د~\LR{$f$} تعريف ته په کتلو تاييدولے شو۔
\olpseganchor{OLP-0034-B012}{end}

\olpseganchor{OLP-0034-B013}{start}
اوس دا لړ په پام کښې ونيسئ:
\[
f(Z_1), f(Z_2), f(Z_3), \dots
\]
ځکه چې \LR{$f$} پر هدف سټ بشپړه ده، د \LR{$\Bin^\omega$} هر غړے بايد د
کوم ورودي دپاره د~\LR{$f$} د قيمت په توګه راشي، او نو بايد په لړ کښې
راشي۔ له دې امله دا لړ بايد ټول~\LR{$\Bin^\omega$} وشمېري۔
\olpseganchor{OLP-0034-B013}{end}

\olpseganchor{OLP-0034-B014}{start}
نو که \LR{$\Pow{\PosInt}$} د شمېر وړ وے، \LR{$\Bin^\omega$} به
د شمېر وړ وے۔ خو \LR{$\Bin^\omega$} د شمېر نۀ وړ دے
(\olref[nen]{thm:nonenum-bin-omega})۔ له دې امله
\LR{$\Pow{\PosInt}$} د شمېر نۀ وړ دے۔
\end{proof}
\olpseganchor{OLP-0034-B014}{end}

\olpseganchor{OLP-0034-B015}{start}
\begin{explain}
د راکمونې په لوري کښې ګډوډېدل اسانه دي۔ د بېلګې په ډول، يوه
پر هدف سټ بشپړه تابع \LR{$g \colon \Bin^\omega \to B$} دا \emph{نۀ}
ثابتوي چې \LR{$B$} د شمېر نۀ وړ دے۔ (تابع
\LR{$g \colon \Bin^\omega \to \Bin$} په پام کښې ونيسئ چې په
\LR{$g(s) = s(1)$} تعريف شوې؛ دا تابع د \LR{$0$}ګانو او \LR{$1$}ګانو يو تسلسل
د هغه لومړي غړے ته نقشوي۔ دا پر هدف سټ بشپړه ده، ځکه ځينې
تسلسلونه په \LR{$0$} او ځينې په \LR{$1$} پيلېږي۔ خو \LR{$\Bin$} متناهي دے۔)
دا هم په ياد ولرئ چې تابع~\LR{$f$} بايد پر هدف سټ بشپړه وي، ګنې
استدلال مخکښې نۀ ځي: بيا به دا ډاډ نۀ وي چې \LR{$f(x_1)$}، \LR{$f(x_2)$}،
\dots{} د~\LR{$B$} ټول غړي لري۔ د بېلګې په ډول،
\[
h(n) = \underbrace{000\dots0}_{\text{\RL{\LR{$n$} \LR{$0$}ګانې}}}111\dots
\]
تابع \LR{$h\colon \PosInt \to \Bin^\omega$} تعريفوي، خو \LR{$\PosInt$}
د شمېر وړ دے۔
\textbf{د ژباړې د نسخې سمون (OLSIZ-005):} اصلي ښودل شوے حاصل
متناهي توريز تسلسل و، نو د يادې تابع د هدف غړے نۀ و۔ دلته تر
لومړنيو صفرونو وروسته د يوونو نامتناهي لکۍ ورزياته شوې؛ د راکمونې
د لوري په اړه پايله هماغه پاتې کېږي۔
\end{explain}
\olpseganchor{OLP-0034-B015}{end}

\olpseganchor{OLP-0034-B016}{start}
\begin{prob}
په راکمونې استدلال وښيئ چې د مثبتو صحيح عددونو د جوړو د
\emph{سټونو} سټ د شمېر نۀ وړ دے۔
\end{prob}
\olpseganchor{OLP-0034-B016}{end}

\olpseganchor{OLP-0034-B017}{start}
\begin{prob}\label{sfr:siz:red:prob:nat-nat}
  په راکمونې استدلال وښيئ چې د ټولو تابعو \LR{$f\colon \Nat \to \Nat$}
  سټ~\LR{$X$} د شمېر نۀ وړ دے۔ (لارښوونه: له \LR{$X$} څخه
  ~\LR{$\Bin^\omega$} ته يوه پر هدف سټ بشپړه تابع ورکړئ۔)
\end{prob}
\olpseganchor{OLP-0034-B017}{end}

\olpseganchor{OLP-0034-B018}{start}
\begin{prob}
وښيئ چې \LR{$\Nat^\omega$}، يعنې د طبيعي عددونو د نامتناهي
تسلسلونو سټ، په راکمونې استدلال د شمېر نۀ وړ دے۔
\end{prob}
\olpseganchor{OLP-0034-B018}{end}

\olpseganchor{OLP-0034-B019}{start}
\begin{prob}
فرض کړئ \LR{$P$} له مثبتو صحيح عددونو څخه سټ~\LR{$\{0\}$} ته د تابعو سټ
دے، او \LR{$Q$} له مثبتو صحيح عددونو څخه سټ~\LR{$\{0\}$} ته د
\emph{جزوي} تابعو سټ دے۔ وښيئ چې \LR{$P$}~د شمېر وړ دے او \LR{$Q$}
نۀ دے۔ (لارښوونه: د \LR{$\Bin^\omega$} د شمېرلو مسئله د~\LR{$Q$} شمېرلو
ته راکمه کړئ۔)
\end{prob}
\olpseganchor{OLP-0034-B019}{end}

\olpseganchor{OLP-0034-B020}{start}
\begin{prob}
فرض کړئ \LR{$S$} له مثبتو صحيح عددونو څخه سټ~\LR{$\{0,1\}$} ته د ټولو
پر هدف سټ بشپړه تابعو سټ دے؛ يعنې \LR{$S$} له ټولو پر هدف سټ بشپړه
~\LR{$f\colon \PosInt \to \Bin$} څخه جوړ دے۔ وښيئ چې \LR{$S$}
د شمېر نۀ وړ دے۔
\textbf{د ژباړې د نسخې سمون (PSSIZ-003):} په انګرېزي نثر کښې
د دوه غړيز سټ نښه له رياضي حالت څخه بهر وه؛ دلته د نورو سټونو
په شان په رياضي حالت کښې ليکل شوې ده۔
\end{prob}
\olpseganchor{OLP-0034-B020}{end}

\olpseganchor{OLP-0034-B021}{start}
\begin{prob}
وښيئ چې د ټولو حقيقي عددونو سټ~\LR{$\Real$} د شمېر نۀ وړ دے۔
\end{prob}
\olpseganchor{OLP-0034-B021}{end}

\olpseganchor{OLP-0035-B004}{start}
\olfileid{sfr}{siz}{equ}
\olpseganchor{OLP-0035-B004}{end}

\olpseganchor{OLP-0035-B005}{start}
\olsection{همشمېر والی}
\olpseganchor{OLP-0035-B005}{end}

\olpseganchor{OLP-0035-B006}{start}
موږ د سټونو د «اندازې» يوه شهودي مفکوره لرو چې د متناهي سټونو
دپاره سمه کار کوي۔ خو د نامتناهي سټونو څه؟ که وغواړو د \emph{هرې}
اندازې د دوو سټونو د اندازو د پرتله کولو يوه صوري لار جوړه کړو،
ښه به وي چې لومړے تعريف کړو سټونه کله يوه اندازه لري۔ دا ئې د
فرېګه وينا ده:
\begin{quote}
  که يو خدمتګار غواړي ډاډه شي چې پر مېز ئې د پليټونو هومره چاړې
  په دقيقه توګه اېښې دي، د دواړو له شمېرلو پرته هم دا کار کولے شي:
  يوازې د هر پليټ ښي لوري ته يوه چاړه کېږدي، داسې چې پر مېز هره
  چاړه د کوم پليټ ښي لوري ته پرته وي۔ په دې توګه پليټونه او چاړې
  په ځانګړې توګه يو له بل سره، او په رښتيا هم د هماغې مکاني اړيکې
  له لارې، تړل کېږي۔ \citep[\S70]{Frege1884}
\end{quote}
د دې عبارت بصيرت په يوه صوري تعريف کښې څرګندولے شو:
\olpseganchor{OLP-0035-B006}{end}

\olpseganchor{OLP-0035-B007}{start}
\begin{defn}\ollabel{comparisondef}
  \LR{$A$} له \LR{$B$} سره \emph{همشمېره} دے، چې \LR{$\cardeq{A}{B}$} ئې ليکو،
  هله او يوازې هله چې يوه دوه اړخيزه يو پر يو تابع \LR{$f \colon A \to B$} وي۔
\end{defn}
\olpseganchor{OLP-0035-B007}{end}

\olpseganchor{OLP-0035-B008}{start}
\begin{prop}\ollabel{equinumerosityisequi}
همشمېر والی يوه د معادلتوب اړيکه ده۔
\end{prop}
\olpseganchor{OLP-0035-B008}{end}

\olpseganchor{OLP-0035-B009}{start}
\begin{proof} 
بايد وښيو چې همشمېر والی انعکاسي، متناظر او انتقالي دے۔ فرض
کړئ \LR{$A, B$} او \LR{$C$} سټونه دي۔
\olpseganchor{OLP-0035-B009}{end}

\olpseganchor{OLP-0035-B010}{start}
\emph{انعکاسيتوب۔} عيني نقشه \LR{$\Id{A} \colon A \to A$}، چې د ټولو
\LR{$x \in A$} دپاره \LR{$\Id{A} (x) = x$} وي، يوه دوه اړخيزه يو پر يو تابع ده۔ نو
\LR{$\cardeq{A}{A}$}۔
\olpseganchor{OLP-0035-B010}{end}

\olpseganchor{OLP-0035-B011}{start}
\emph{تناظر۔} فرض کړئ \LR{$\cardeq{A}{B}$}؛ يعنې يوه دوه اړخيزه يو پر يو تابع
\LR{$f\colon A \to B$} شته۔ ځکه چې \LR{$f$} دوه اړخيزه يو پر يو ده، د هغې معکوس
\LR{$f^{-1}$} شته او هغه هم دوه اړخيزه يو پر يو دے۔ له دې امله
\LR{$f^{-1}\colon B \to A$} يوه دوه اړخيزه يو پر يو تابع ده، نو \LR{$\cardeq{B}{A}$}۔
\olpseganchor{OLP-0035-B011}{end}

\olpseganchor{OLP-0035-B012}{start}
\emph{انتقاليتوب۔} فرض کړئ چې \LR{$\cardeq{A}{B}$} او
\LR{$\cardeq{B}{C}$}؛ يعنې دوه اړخيزې يو پر يو تابعې \LR{$f\colon A \to B$} او
\LR{$g\colon B \to C$} شته۔ بيا ترکيب \LR{$\comp{f}{g}\colon A \to C$}
دوه اړخيزه يو پر يو دے، نو \LR{$\cardeq{A}{C}$}۔
\end{proof}
\olpseganchor{OLP-0035-B012}{end}

\olpseganchor{OLP-0035-B013}{start}
\begin{prop}
که \LR{$\cardeq{A}{B}$} وي، نو \LR{$A$} هله او يوازې هله د شمېر وړ
دے چې \LR{$B$} داسې وي۔
\end{prop}
\olpseganchor{OLP-0035-B013}{end}

\olpseganchor{OLP-0035-B014}{start}
\begin{editorial}
لاندې ثبوت که \olref[enm]{sec} شامل وي،
\olref[enm]{defn:enumerable} کاروي، او که نۀ وي
\olref[enm-alt]{defn:enumerable} کاروي۔
\end{editorial}
\olpseganchor{OLP-0035-B014}{end}

\olpseganchor{OLP-0035-B015}{start}
\begin{proof}
فرض کړئ \LR{$\cardeq{A}{B}$}، نو کومه دوه اړخيزه يو پر يو تابع \LR{$f \colon A
\to B$} شته، او فرض کړئ چې \LR{$A$} د شمېر وړ دے۔

  بيا يا \LR{$A = \emptyset$} دے، يا يوه پر هدف سټ بشپړه تابع
  \LR{$g\colon \PosInt \to A$} شته۔ که \LR{$A = \emptyset$} وي، نو
  \LR{$B = \emptyset$} هم دے (ګنې يو غړے~\LR{$y \in B$} به وے،
  خو هېڅ \LR{$x \in A$} به نۀ وے چې \LR{$f(x) = y$})۔
  \textbf{د ژباړې د نسخې سمون (OLSIZ-006):} د انګرېزي ثبوت په
  دواړو مشروطو بڼو کښې د تش حالت معادله ناسمې تابعې ته منسوبه
  وه؛ دلته هغه دوه اړخيزې تابعې ته سمه شوې چې د سټونو همشمېر
  والی شاهدي کوي۔ له بلې خوا، که \LR{$g\colon \PosInt \to A$}
  پر هدف سټ بشپړه وي، نو \LR{$\comp{g}{f} \colon \PosInt \to B$}
  پر هدف سټ بشپړه ده۔ د دې د ليدلو دپاره \LR{$y \in B$} ونيسئ۔ ځکه
  چې \LR{$f$} پر هدف سټ بشپړه ده، يو \LR{$x \in A$} شته چې \LR{$f(x) = y$}۔ ځکه
  چې \LR{$g$} پر هدف سټ بشپړه ده، يو \LR{$n \in \PosInt$} شته چې \LR{$g(n) =
  x$}۔ له دې امله
\[
(\comp{g}{f})(n) = f(g(n)) = f(x) = y
\]
او نو \LR{$\comp{g}{f}$} پر هدف سټ بشپړه ده۔ موږ لرو چې
\LR{$\comp{g}{f}$} د~\LR{$B$} يوه شمېرنه ده، نو \LR{$B$}~د شمېر وړ دے۔
\olpseganchor{OLP-0035-B015}{end}

\olpseganchor{OLP-0035-B016}{start}
که \LR{$B$} د شمېر وړ وي، د~\LR{$f$} پر ځاے د دوه اړخيزه يو پر يو تابع
\LR{$f^{-1}\colon B \to A$} په کارولو د استدلال په تکرارولو مومو چې
\LR{$A$} د شمېر وړ دے۔
\end{proof}
\olpseganchor{OLP-0035-B016}{end}

\olpseganchor{OLP-0035-B017}{start}
\begin{prob}
وښيئ چې که \LR{$\cardeq{A}{C}$} او \LR{$\cardeq{B}{D}$} وي، او
\LR{$A \cap B = C \cap D = \emptyset$} وي، نو
\LR{$\cardeq{A \cup B}{C \cup D}$}۔
\end{prob}
\olpseganchor{OLP-0035-B017}{end}

\olpseganchor{OLP-0035-B018}{start}
\begin{prob}
وښيئ چې که \LR{$A$} نامتناهي او د شمېر وړ وي، نو
\LR{$\cardeq{A}{\Nat}$}۔
\end{prob}
\olpseganchor{OLP-0035-B018}{end}

\olpseganchor{OLP-0036-B004}{start}
\olfileid{sfr}{siz}{car}
\olpseganchor{OLP-0036-B004}{end}

\olpseganchor{OLP-0036-B005}{start}
\olsection{د بېلو اندازو سټونه او د کانتور قضيه}
\olpseganchor{OLP-0036-B005}{end}

\olpseganchor{OLP-0036-B006}{start}
\begin{explain}
موږ دا مفکوره په دقيقه وينا کښې وړاندې کړې چې دوه سټونه يوه
اندازه لري۔ دا مفکوره هم په دقيقه وينا کښې وړاندې کولے شو چې يو
سټ له بل څخه کوچنے دے۔ زموږ د «له ... څخه کوچنے (يا همشمېره)
دے» تعريف به د سټونو تر منځ د دوه اړخيزه يو پر يو تابع پر ځاے، له لومړي
سټ څخه دويم ته يوه يو پر يو تابع وغواړي۔ که داسې تابع وي، د
لومړي سټ اندازه د دويم سټ له اندازې څخه کوچنۍ يا ورسره مساوي ده۔
په شهودي ډول، له يوه سټ څخه بل ته يو پر يو تابع دا ډاډ راکوي
چې د تابعې د قيمتونو سټ لږ تر لږه د تعريف د ساحې هومره
غړي لري، ځکه د تعريف د ساحې هېڅ دوه غړي د
قيمتونو د سټ يوه غړے ته نۀ نقشېږي۔
\end{explain}
\olpseganchor{OLP-0036-B006}{end}

\olpseganchor{OLP-0036-B007}{start}
\begin{defn}
\LR{$A$} له~\LR{$B$} څخه \emph{لوے نۀ دے}، چې \LR{$\cardle{A}{B}$} ئې ليکو،
هله او يوازې هله چې يوه يو پر يو تابع \LR{$f \colon A \to B$} وي۔
\end{defn}
\olpseganchor{OLP-0036-B007}{end}

\olpseganchor{OLP-0036-B008}{start}
څرګنده ده چې دا يوه انعکاسي او انتقالي اړيکه ده، خو متناظره نۀ
ده (دا د تمرين په توګه پرېښوول شوې ده)۔ يوه داسې مفکوره هم
معرفي کولے شو چې وايي يو سټ له بل څخه (په ټينګه) کوچنے دے۔
\olpseganchor{OLP-0036-B008}{end}

\olpseganchor{OLP-0036-B009}{start}
\begin{defn}
\LR{$A$} له~\LR{$B$} څخه \emph{کوچنے دے}، چې \LR{$\cardless{A}{B}$} ئې ليکو،
هله او يوازې هله چې يوه يو پر يو تابع~\LR{$f\colon A \to B$} وي،
خو هېڅ دوه اړخيزه يو پر يو تابع~\LR{$g\colon A \to B$} نۀ وي؛ يعنې
\LR{$\cardle{A}{B}$} او \LR{$\cardneq{A}{B}$} وي۔
\end{defn}
\olpseganchor{OLP-0036-B009}{end}

\olpseganchor{OLP-0036-B010}{start}
څرګنده ده چې دا اړيکه ناانعکاسي او انتقالي ده۔ (دا د تمرين په
توګه پرېښوول شوې ده۔) په دې نښه ويلے شو چې يو سټ \LR{$A$} هله او
يوازې هله د شمېر وړ دے چې \LR{$\cardle{A}{\Nat}$} وي، او \LR{$A$}
هله او يوازې هله د شمېر نۀ وړ دے چې
\LR{$\cardless{\Nat}{A}$} وي۔
\textbf{د ژباړې د نسخې سرچينيز شرط (PSSIZ-004):} وروستے
هم‌ارز بيان دا فرض کاروي چې هر نامتناهي سټ د طبيعي عددونو يوه
يو پر يو کاپي لري؛ دا د انتخاب له اصل پرته په عمومي سټ‌تيورۍ
کښې نۀ ثابتېږي۔ په دې سره دا بيا ويلے شو:
%
\olref[sfr][siz][nen-alt]{thm:nonenum-pownat}
د دې مشاهدې په توګه چې
\LR{$\cardless{\Nat}{\Pow{\Nat}}$}۔ په
حقيقت کښې، \citet{Cantor1892} ثابته کړه چې دا وروستۍ خبره
\emph{بشپړه عمومي} ده:
\olpseganchor{OLP-0036-B010}{end}

\olpseganchor{OLP-0036-B011}{start}
\begin{thm}[کانتور]\ollabel{thm:cantor}
د هر سټ \LR{$A$} دپاره \LR{$\cardless{A}{\Pow{A}}$}۔
\end{thm}
\olpseganchor{OLP-0036-B011}{end}

\olpseganchor{OLP-0036-B012}{start}
\begin{proof}
نقشه \LR{$f(x) = \{x\}$} يوه يو پر يو تابع
\LR{$f \colon A \to \Pow{A}$} ده، ځکه که \LR{$x \neq y$} وي، نو د غړو له
مخې د برابرۍ په سبب \LR{$\{x\} \neq \{y\}$} هم وي، او نو
\LR{$f(x) \neq f(y)$}۔ له دې امله \LR{$\cardle{A}{\Pow{A}}$} لرو۔
\olpseganchor{OLP-0036-B012}{end}

\olpseganchor{OLP-0036-B013}{start}
\begin{editorial}
که \olref[nen]{sec} موجوده وي، تفصيلي ثبوت وړاندې کوو؛ ګنې يو
چټک ثبوت وړاندې کوو چې د \olref[nen-alt]{sec} سره سمون خوري۔
\end{editorial}
\olpseganchor{OLP-0036-B013}{end}

\olpseganchor{OLP-0036-B014}{start}
% 
  اوس به وښيو چې هېڅ پر هدف سټ بشپړه تابع~\LR{$g\colon A \to
  \Pow{A}$} نۀ شي کېدے، نو يوه دوه اړخيزه يو پر يو تابع خو لا پرېږده؛
  او له دې امله \LR{$\cardneq{A}{\Pow{A}}$}۔ ځکه فرض کړئ چې
  \LR{$g\colon A \to \Pow{A}$}۔ دا چې \LR{$g$} بشپړه ده، هر \LR{$x \in A$} د
  يوه فرعي سټ \LR{$g(x) \subseteq A$} سره نقشېږي۔ موږ ښودلے شو چې
  \LR{$g$} پر هدف سټ بشپړه نۀ شي کېدے۔ د دې دپاره يو فرعي سټ
  ~\LR{$\overline{A} \subseteq A$} تعريفوو چې د تعريف له مخې د~\LR{$g$}
  په قيمتونو سټ کښې نۀ شي راتلے۔ کېږدئ:
  \[
  \overline{A} = \Setabs{x \in A}{x \notin g(x)}.
  \]
  ځکه چې \LR{$g(x)$} د هر \LR{$x \in A$} دپاره تعريف شوے، \LR{$\overline{A}$}
  په څرګند ډول د~\LR{$A$} يو سم تعريف شوے فرعي سټ دے۔ خو دا د~\LR{$g$}
  د قيمتونو په سټ کښې نۀ شي راتلے۔ يو خپل سری \LR{$x \in A$} ونيسئ؛
  موږ به وښيو چې \LR{$\overline{A} \neq g(x)$}۔ که \LR{$x \in g(x)$} وي،
  نو \LR{$x \notin g(x)$} نۀ پوره کوي، او د~\LR{$\overline{A}$} د تعريف له
  مخې \LR{$x \notin \overline{A}$} لرو۔ که \LR{$x \in \overline{A}$} وي،
  د~\LR{$\overline{A}$} تعريفي خاصيت بايد پوره کړي؛ يعنې \LR{$x \in A$}
  او \LR{$x \notin g(x)$} وي۔ دا چې \LR{$x$} خپل سری و، دا ښيي چې د هر
  \LR{$x \in A$} دپاره، \LR{$x \in g(x)$} هله او يوازې هله چې
  \LR{$x \notin \overline{A}$}، او نو \LR{$g(x) \neq \overline{A}$}۔
  \textbf{د ژباړې د نسخې سمون (OLSIZ-007):} انګرېزي ثبوت په
  دې عمومي پايله کښې ورودي د ټول سټ پر ځاے يوازې قطري فرعي سټ
  ته محدود کړے و؛ دلته ئې د مخکښې ټاکل شوي خپل سري ورودي ساحه
  راوګرځول شوه۔ په بله وينا، \LR{$\overline{A}$} د~\LR{$g$} په قيمتونو
  سټ کښې نۀ شي راتلے، او دا له دې فرض سره تناقض لري چې~\LR{$g$} پر
  هدف سټ بشپړه ده۔{}
\end{proof}
\olpseganchor{OLP-0036-B014}{end}

\olpseganchor{OLP-0036-B015}{start}
\begin{explain}
د
  \olref{thm:cantor} ثبوت د
  \olref[nen]{thm:nonenum-pownat} له ثبوت سره پرتله کول ګټور دي۔
  هلته مو وښودل چې د~\LR{$\PosInt$} د فرعي سټونو د هر لړ \LR{$Z_1$}،
  \LR{$Z_2$}، \dots دپاره د عددونو يو سټ~\LR{$\overline{Z}$} جوړولے شو
  چې په ډاډه توګه په لړ کښې نۀ وي۔ په لړ کښې نۀ و، ځکه چې د هر
  \LR{$n \in \PosInt$} دپاره، \LR{$n \in Z_n$} هله او يوازې هله چې
  \LR{$n \notin \overline{Z}$}۔ په دې توګه تل کوم عدد د \LR{$Z_n$} او
  \LR{$\overline{Z}$} له يوه څخه غړے وي خو د بل نۀ وي۔ دلته
  هماغه مفکوره تعقيبوو، خو شاخصونه~\LR{$n$} اوس د~\LR{$\PosInt$} پر ځاے
  د~\LR{$A$} غړي دي۔ سټ \LR{$\overline{A}$} داسې تعريف شوے چې د
  هر \LR{$x \in A$} دپاره له~\LR{$g(x)$} څخه بېل وي، ځکه \LR{$x \in g(x)$}
  هله او يوازې هله چې \LR{$x \notin \overline{A}$}۔ بيا هم تل د~\LR{$A$}
  يو غړے شته چې له \LR{$g(x)$} او \LR{$\overline{A}$} څخه د يوه
  غړے وي خو د بل نۀ وي۔ او لکه څنګه چې له دې امله
  \LR{$\overline{Z}$} په لړ \LR{$Z_1$}، \LR{$Z_2$}، \dots کښې نۀ شي راتلے،
  \LR{$\overline{A}$} د~\LR{$g$} په قيمتونو سټ کښې نۀ شي راتلے۔
\olpseganchor{OLP-0036-B015}{end}

\olpseganchor{OLP-0036-B016}{start}
د
  \olref{thm:cantor} ثبوت د
  \olref[nen-alt]{thm:nonenum-pownat} له ثبوت سره پرتله کول ګټور
  دي۔ هلته مو وښودل چې د~\LR{$\Nat$} د فرعي سټونو د هر لړ \LR{$N_0$}،
  \LR{$N_1$}، \LR{$N_2$}، \dots دپاره د عددونو يو سټ~\LR{$D$} جوړولے شو چې په
  ډاډه توګه په لړ کښې نۀ وي۔ په لړ کښې نۀ و، ځکه چې د هر
  \LR{$n \in \Nat$} دپاره \LR{$n \in N_n$} هله او يوازې هله چې
  \LR{$n \notin D$}۔ دلته هماغه مفکوره تعقيبوو، خو شاخصونه~\LR{$n$} اوس
  د~\LR{$\Nat$} پر ځاے د~\LR{$A$} غړي دي۔ سټ \LR{$D$} داسې تعريف
  شوے چې د هر \LR{$x \in A$} دپاره له~\LR{$g(x)$} څخه بېل وي، ځکه
  \LR{$x \in g(x)$} هله او يوازې هله چې \LR{$x \notin D$}۔
\olpseganchor{OLP-0036-B016}{end}

\olpseganchor{OLP-0036-B017}{start}
دا ثبوت د رسل د پاراډوکس له ثبوت،
\olref[sfr][set][rus]{thm:russells-paradox}، سره هم د پرتله کولو
وړ دے۔ په حقيقت کښې، د کانتور قضيه د رسل د خپل پاراډوکس الهام و۔
\end{explain}
\olpseganchor{OLP-0036-B017}{end}

\olpseganchor{OLP-0036-B018}{start}
\begin{prob}
  وښيئ چې د هر سټ~\LR{$A$} دپاره هېڅ يو پر يو تابع
  \LR{$g\colon \Pow{A} \to A$} نۀ شي کېدے۔ لارښوونه: فرض کړئ
  \LR{$g\colon \Pow{A} \to A$} يو پر يو ده۔
  \LR{$D = \Setabs{g(B)}{B \subseteq A \text{\RL{ او }} g(B) \notin B}$}
  په پام کښې ونيسئ۔ کېږدئ \LR{$x = g(D)$}۔ د تناقض د راايستلو دپاره
  دا حقيقت وکاروئ چې \LR{$g$} يو پر يو ده۔
\end{prob}
\olpseganchor{OLP-0036-B018}{end}

\olpseganchor{OLP-0037-B004}{start}
\olfileid{sfr}{siz}{sb}
\olpseganchor{OLP-0037-B004}{end}

\olpseganchor{OLP-0037-B005}{start}
\olsection{د اندازې مفکوره او شروډر–برنشتاين}
\olpseganchor{OLP-0037-B005}{end}

\olpseganchor{OLP-0037-B006}{start}
\begin{explain}
دا يوه شهودي مفکوره ده: که \LR{$A$} له \LR{$B$} څخه لوے نۀ وي او \LR{$B$} له
\LR{$A$} څخه لوے نۀ وي، نو \LR{$A$} او \LR{$B$} همشمېره دي۔ رښتيا دا چې که
دا مفکوره \emph{ناسمه} وے، موږ به په ډېره ګرانه توجيه کړے شوے
چې زموږ تعريف شوې د همشمېر والي مفکوره د سټونو تر منځ د «اندازو»
له پرتله کولو سره څه اړيکه لري!{} خو له نېکه مرغه، شهودي مفکوره
سمه ده۔ دا د شروډر–برنشتاين په قضيې توجيه کېږي۔
\end{explain}
\olpseganchor{OLP-0037-B006}{end}

\olpseganchor{OLP-0037-B007}{start}
\begin{thm}[شروډر–برنشتاين]
	\ollabel{thm:schroder-bernstein}
	که \LR{$\cardle{A}{B}$} او \LR{$\cardle{B}{A}$} وي،
	نو \LR{$\cardeq{A}{B}$}۔
\end{thm}
\olpseganchor{OLP-0037-B007}{end}

\olpseganchor{OLP-0037-B008}{start}
\begin{explain}
په بله وينا، که له \LR{$A$} څخه~\LR{$B$} ته يوه يو پر يو تابع، او له
\LR{$B$} څخه~\LR{$A$} ته يوه يو پر يو تابع وي، نو له \LR{$A$} څخه~\LR{$B$} ته
يوه دوه اړخيزه يو پر يو تابع شته۔
\olpseganchor{OLP-0037-B008}{end}

\olpseganchor{OLP-0037-B009}{start}
خو د دې پايلې ثبوت په رښتيا خورا \emph{ګران} دے۔ په حقيقت کښې،
که څه هم کانتور دا پايله بيان کړه، نورو ثابته کړه۔\footnote{د
تاريخ په اړه د نورو معلوماتو دپاره، د بېلګې په ډول
\citet[pp.~165--6]{Potter2004} وګورئ۔}
%
اوس مهال بايد دا په باور ومنئ۔
\olpseganchor{OLP-0037-B009}{end}

\olpseganchor{OLP-0037-B010}{start}
له نېکه مرغه، شروډر–برنشتاين \emph{سمه} ده او دا فکر مو تاييدوي
چې زموږ تعريف شوې اړيکې، يعنې \LR{$\cardeq{A}{B}$} او
\LR{$\cardle{A}{B}$}، له «اندازې» سره څه اړيکه لري۔ پر دې سربېره، د
شروډر–برنشتاين قضيه ډېره \emph{ګټوره} ده۔ د دوو همشمېره سټونو
تر منځ د دوه اړخيزه يو پر يو تابع تصورول ګرانېدلای شي۔ د شروډر–برنشتاين
قضيه اجازه راکوي چې پرتله په حالتونو ووېشو، څو يوازې له لومړي
څخه دويم ته د يو پر يو تابع، او بيا برعکس، تصور وکړو۔
\end{explain}
\olpseganchor{OLP-0037-B010}{end}

\olpseganchor{OLP-0038-B004}{start}
\olfileid{sfr}{siz}{enm-alt}
\olpseganchor{OLP-0038-B004}{end}

\olpseganchor{OLP-0038-B005}{start}
\olsection{شمېرنې او د شمېر وړ سټونه}
\olpseganchor{OLP-0038-B005}{end}

\olpseganchor{OLP-0038-B006}{start}
\begin{editorial}
  دا برخه شمېرنې له \LR{$\Nat$} (پيلنيو ټوټو) سره د دوه اړخيزو يو
  پر يو تابعو په توګه تعريفوي، لکه په سټ‌تيورۍ کښې چې کېږي۔ نو
  د \olref[enm]{sec} له تعريفونو سره لږ ټکر لري، او هلته ټولې
  بېلګې تکراروي۔ دا د هغې برخې په پرتله لږه لنډه هم ده۔
\end{editorial}
\olpseganchor{OLP-0038-B006}{end}

\olpseganchor{OLP-0038-B007}{start}
يو متناهي سټ يوازې د هغه د غړي په شمېرلو ټاکلے شو۔ دا
هغه وخت کوو چې يو سټ داسې تعريفوو:
\[
  A = \{a_1, a_2, \ldots, a_n\}.
\]
که فرض کړو چې غړي \LR{$a_1$}، \dots، \LR{$a_n$} ټول سره بېل دي،
دا د \LR{$A$} او لومړنيو \LR{$n$} طبيعي عددونو \LR{$0$}، \dots، \LR{$n-1$} تر منځ
يوه دوه اړخيزه يو پر يو تابع راکوي۔ برعکس، ځکه هر متناهي سټ يوازې متناهي
شمېر غړي لري، هر متناهي سټ په داسې سمون کښې اېښودلے
شو۔ په بله وينا، که \LR{$A$} متناهي وي، د \LR{$A$} او
\LR{$\{0, \dots, n-1\}$} تر منځ يوه دوه اړخيزه يو پر يو تابع شته، چې \LR{$n$} د~\LR{$A$}
د غړي شمېر دے۔
\olpseganchor{OLP-0038-B007}{end}

\olpseganchor{OLP-0038-B008}{start}
که د نامتناهي سټونو ځينې ډولونه ومنو، نو ځينې نامتناهي سټونه به
هم شمېرل کېدے شي۔ دا په دې وينا دقيقولے شو چې يو نامتناهي سټ د
هغه او ټول~\LR{$\Nat$} تر منځ د دوه اړخيزه يو پر يو تابع په وسيله شمېرل کېږي۔
\olpseganchor{OLP-0038-B008}{end}

\olpseganchor{OLP-0038-B009}{start}
\begin{defn}[شمېرنه، سټ‌تيوريکي]
د يوه سټ \LR{$A$} \emph{شمېرنه} يوه دوه اړخيزه يو پر يو تابع ده چې د قيمتونو
سټ ئې \LR{$A$} او د تعريف ساحه ئې يا د طبيعي عددونو يو پيلنی سټ
\LR{$\{0, 1, \ldots, n\}$} {يا} د طبيعي عددونو ټول سټ~\LR{$\Nat$} وي۔
\end{defn}
\olpseganchor{OLP-0038-B009}{end}

\olpseganchor{OLP-0038-B010}{start}
\begin{explain}
د \emph{شمېرنې} د ويي د دې کارونې تر شا يو شهودي بنسټ شته۔ دا
ويل چې موږ يو سټ \LR{$A$} شمېرلے دے، په دې معنا دي چې يوه
دوه اړخيزه يو پر يو تابع \LR{$f$} شته چې پرې د سټ \LR{$A$} غړي شمېرلے شو۔ \LR{$0$}م غړے
\LR{$f(0)$} دے، لومړے \LR{$f(1)$}، \ldots، \LR{$n$}م \LR{$f(n)$}،
\ldots۔\footnote{هو، موږ له \LR{$0$} څخه شمېرو۔ البته له~\LR{$1$} څخه هم
پيلولے شو۔ په دې سره به لوے توپير رانۀ شي۔ يوازې به مو~\LR{$\Nat$}
په~\LR{$\PosInt$} بدلولے۔} د دې دليل د لاندې خبرې په ورزياتولو لا
روښانه کېدے شي:
\end{explain}
\olpseganchor{OLP-0038-B010}{end}

\olpseganchor{OLP-0038-B011}{start}
\begin{defn}
  \ollabel{defn:enumerable}
  يو سټ~\LR{$A$} هله او يوازې هله د شمېر وړ دے چې يا
  \LR{$A = \emptyset$} وي يا د~\LR{$A$} يوه شمېرنه وي۔ وايو چې \LR{$A$} هله
  او يوازې هله د شمېر نۀ وړ دے چې \LR{$A$} د شمېر وړ نۀ وي۔
\end{defn}
\olpseganchor{OLP-0038-B011}{end}

\olpseganchor{OLP-0038-B012}{start}
\begin{explain}
نو يو سټ هله او يوازې هله د شمېر وړ دے چې تش وي يا ئې
غړي په يوه شمېرنه شمېرلے شئ۔
\end{explain}
\olpseganchor{OLP-0038-B012}{end}

\olpseganchor{OLP-0038-B013}{start}
\begin{ex}
هغه تابع چې طبيعي عددونه شمېري يوازې عيني تابع
\LR{$\Id{\Nat} \colon \Nat \to \Nat$} ده چې په
\LR{$\Id{\Nat}(n) = n$} ورکړل شوې۔ هغه تابع چې \emph{مثبت} طبيعي
عددونه شمېري، \LR{$\Nat^+ = \Nat \setminus \{0\}$}، تابع
\LR{$g(n) = n + 1$} ده؛ يعنې راتلونکې تابع۔
\end{ex}
\olpseganchor{OLP-0038-B013}{end}

\olpseganchor{OLP-0038-B014}{start}
\begin{prob}
وښيئ چې يو سټ \LR{$A$} هله او يوازې هله د شمېر وړ دے چې يا
\LR{$A = \emptyset$} وي يا يوه پر هدف سټ بشپړه تابع \LR{$f\colon \Nat \to A$}
وي۔ وښيئ چې \LR{$A$} هله او يوازې هله د شمېر وړ دے چې يوه
يو پر يو تابع \LR{$g\colon A \to \Nat$} وي۔
\end{prob}
\olpseganchor{OLP-0038-B014}{end}

\olpseganchor{OLP-0038-B015}{start}
\begin{ex}
تابعې \LR{$f\colon \Nat \to \Nat$} او \LR{$g \colon \Nat \to \Nat$} چې
په ترتيب سره په لاندې ډول ورکړل شوې دي:
\begin{align*}
  f(n) & = 2n \text{\RL{ او}}\\
  g(n) & = 2n+1
\end{align*}
جفت طبيعي عددونه او طاق طبيعي عددونه شمېري۔ خو له دواړو يوه هم
پر هدف سټ بشپړه نۀ ده، نو يوه هم د \LR{$\Nat$} شمېرنه نۀ ده۔
\end{ex}
\olpseganchor{OLP-0038-B015}{end}

\olpseganchor{OLP-0038-B016}{start}
\begin{prob}
د مربع عددونو \LR{$1$}، \LR{$4$}، \LR{$9$}، \LR{$16$}، \dots يوه شمېرنه تعريف کړئ۔
\end{prob}
\olpseganchor{OLP-0038-B016}{end}

\olpseganchor{OLP-0038-B017}{start}
\begin{ex}
فرض کړئ \LR{$\lceil x \rceil$} د \emph{پاسني صحيح عدد} تابع ده، چې
\LR{$x$} تر ټولو نږدې صحيح عدد ته پورته ګردوي۔ بيا تابع
\LR{$f \colon \Nat \to \Int$} چې په لاندې ډول ورکړل شوې ده:
\[
  f(n) = (-1)^{n} \left\lceil\tfrac{n}{2}\right\rceil
\]
د صحيح عددونو سټ~\LR{$\Int$} په لاندې ډول شمېري:
\[
\begin{array}{c c c c c c c c}
f(0) & f(1) & f(2) & f(3) & f(4) & f(5) & f(6) & \dots \\ \\
\big\lceil \tfrac{0}{2} \big\rceil & -\big\lceil \tfrac{1}{2}\big\rceil &  \big\lceil \tfrac{2}{2} \big\rceil & -\big\lceil \tfrac{3}{2} \big\rceil & \big\lceil \tfrac{4}{2} \big\rceil  & -\big\lceil \tfrac{5}{2}\big\rceil & \big\lceil \tfrac{6}{2} \big\rceil & \dots \\ \\
0 & -1 & 1 & -2 & 2 & -3 & 3& \dots
\end{array}
\]
پام وکړئ چې \LR{$f$} څنګه د مثبتو او منفي صحيح عددونو تر منځ په
«ټوپونو» د~\LR{$\Int$} قيمتونه پيدا کوي۔ \LR{$f$} په حالتونو هم داسې
تعريفولے شئ:
\[
f(n) = \begin{cases}
  \frac{n}{2} & \text{\RL{که \LR{$n$} جفت وي}}\\
  -\frac{n+1}{2} & \text{\RL{که \LR{$n$} طاق وي}}
  \end{cases}
\]
\end{ex}
\olpseganchor{OLP-0038-B017}{end}

\olpseganchor{OLP-0038-B018}{start}
\begin{prob}
وښيئ چې که \LR{$A$} او \LR{$B$} د شمېر وړ وي، نو \LR{$A \cup B$} هم دے۔
\end{prob}
\olpseganchor{OLP-0038-B018}{end}

\olpseganchor{OLP-0038-B019}{start}
\begin{prob}
پر \LR{$n$} په استقرا وښيئ چې که \LR{$A_1$}، \LR{$A_2$}، \dots، \LR{$A_n$} ټول
د شمېر وړ وي، نو \LR{$A_1 \cup \dots \cup A_n$} هم دے۔
\end{prob}
\olpseganchor{OLP-0038-B019}{end}

\olpseganchor{OLP-0039-B004}{start}
\olfileid{sfr}{siz}{nen-alt}
\olpseganchor{OLP-0039-B004}{end}

\olpseganchor{OLP-0039-B005}{start}
\olsection{د شمېر نۀ وړ سټونه}
\olpseganchor{OLP-0039-B005}{end}

\olpseganchor{OLP-0039-B006}{start}
\begin{editorial}
  دا برخه په \olref[enm-alt]{sec} کښې د تعريفونو په کارولو، يعنې
  له \LR{$\PosInt$} څخه د پر هدف سټ بشپړې تابع پر ځاے له~\LR{$\Nat$} سره
  د دوه اړخيزې يو پر يو تابع په غوښتلو، د \LR{$\Bin^\omega$} او
  \LR{$\Pow{\Nat}$} د شمېر نۀ وړ والے ثابتوي۔
\end{editorial}
\olpseganchor{OLP-0039-B006}{end}

\olpseganchor{OLP-0039-B007}{start}
\begin{explain}
د طبيعي عددونو سټ \LR{$\Nat$} نامتناهي دے۔ دا په څرګند ډول
د شمېر وړ هم دے۔ خو د پام وړ حقيقت دا دے چې
\emph{د شمېر نۀ وړ} سټونه، يعنې هغه سټونه چې د شمېر وړ
نۀ دي، شته (\olref[sfr][siz][enm-alt]{defn:enumerable} وګورئ)۔
\olpseganchor{OLP-0039-B007}{end}

\olpseganchor{OLP-0039-B008}{start}
دا ښايي حيرانوونکې وي۔ بالاخره، دا ويل چې \LR{$A$} د شمېر نۀ وړ
دے، دا ويل دي چې \emph{هېڅ} دوه اړخيزه يو پر يو تابع
\LR{$f \colon \Nat \to A$} نۀ شته؛ يعنې هېڅ هغه تابع چې د~\LR{$\Nat$}
بې‌شمېره غړي پر~\LR{$A$} نقشوي، ټول~\LR{$A$} نۀ شي پوره کولے۔
نو که \LR{$A$} د شمېر نۀ وړ وي، د طبيعي عددونو په پرتله د~\LR{$A$}
«زيات» غړي شته۔
\olpseganchor{OLP-0039-B008}{end}

\olpseganchor{OLP-0039-B009}{start}
د دې د ثابتولو دپاره چې يو سټ د شمېر نۀ وړ دے، بايد وښيئ چې هېڅ
مناسبه دوه اړخيزه يو پر يو تابع نۀ شي موجودېدلای۔ د دې تر ټولو ښه لار دا ده
چې وښيئ د~\LR{$A$} د غړي د شمېرلو هره هڅه بايد لږ تر لږه يو
غړے پرېږدي؛ دا ښيي چې هېڅ تابع \LR{$f\colon \Nat \to A$}
پر هدف سټ بشپړه نۀ ده۔ د دې د ثابتولو يوه عمومي تګلاره د کانتور
\emph{قطري طريقه} ده۔ د~\LR{$A$} د غړي يو لړ، د بېلګې په
ډول \LR{$x_1$}، \LR{$x_2$}، \dots، چې راکړل شي، موږ د~\LR{$A$} يو بل
غړے جوړوو چې د خپل جوړښت له مخې په هېڅ ډول په هغه لړ
کښې نۀ شي راتلے۔
\olpseganchor{OLP-0039-B009}{end}

\olpseganchor{OLP-0039-B010}{start}
خو دا هر څه په بېلګه تر ټولو ښه پوهېدل کېږي۔ نو زموږ لومړۍ
بېلګه د \LR{$0$}ګانو او \LR{$1$}ګانو د ټولو نامتناهي توريزو تسلسلونو سټ
~\LR{$\Bin^\omega$} دے۔ (`\LR{$\Bin$}' د دوه‌ګوني دپاره دے، او يوازې ئې د
دوه غړيز سټ \LR{$\{0,1\}$} په توګه ګڼلے شو۔) خو اوس مهال دا لږ نرم بيان بايد هېڅ ګډوډي جوړه نۀ کړي۔
\end{explain}
\olpseganchor{OLP-0039-B010}{end}

\olpseganchor{OLP-0039-B011}{start}
\begin{thm}
\ollabel{thm:nonenum-bin-omega}
\LR{$\Bin^\omega$}~د شمېر نۀ وړ دے۔
\end{thm}
\olpseganchor{OLP-0039-B011}{end}

\olpseganchor{OLP-0039-B012}{start}
\begin{proof}
د \LR{$\Bin^\omega$} د يوه فرعي سټ هره شمېرنه په پام کښې ونيسئ۔ نو
کوم لړ \LR{$s_{0}$}، \LR{$s_{1}$}، \LR{$s_{2}$}، \dots{} لرو چې هر \LR{$s_n$} د
\LR{$0$}ګانو او~\LR{$1$}ګانو يو نامتناهي توريز تسلسل دے۔ فرض کړئ
\LR{$s_n(m)$} په دې لړ کښې د \LR{$n$}م توريز تسلسل \LR{$m$}م رقم دے۔
\textbf{د ژباړې د نسخې سمون (OLSIZ-008):} انګرېزي نثر د تسلسل
او رقم د شاخصونو رولونه سرچپه بيان کړي وو؛ لاندې اصلي جدول لومړے
شاخص د قطار او دويم شاخص د ستون په توګه ټاکي۔ نو اوس خپل لړ د يوه جدولي
ترتيب په توګه ګڼلے شو، چې \LR{$s_n(m)$} په \LR{$n$}م قطار او \LR{$m$}م ستون
کښې اېښوول شوے دے:
\[
\begin{array}{c|c|c|c|c|c}
& 0 & 1 & 2 & 3 & \dots \\\hline
0 & \mathbf{s_{0}(0)} & s_{0}(1) & s_{0}(2) & s_0(3) & \dots \\\hline
1 & s_{1}(0)& \mathbf{s_{1}(1)} & s_1(2) & s_1(3) & \dots \\\hline
2 & s_{2}(0)& s_{2}(1) & \mathbf{s_2(2)} & s_2(3) & \dots \\\hline
3 & s_{3}(0)& s_{3}(1) & s_3(2) & \mathbf{s_3(3)} & \dots \\\hline
\vdots & \vdots & \vdots & \vdots & \vdots & \mathbf{\ddots}
\end{array}
\]
اوس د \LR{$0$}ګانو او \LR{$1$}ګانو يو داسې نامتناهي توريز تسلسل \LR{$d$}
جوړوو چې په دې لړ کښې نۀ وي۔ دا د هغه د هرې خانې په ټاکلو؛
يعنې د ټولو \LR{$n \in \Nat$} دپاره د \LR{$d(n)$} په ټاکلو، کوو۔ په شهودي
ډول، د پورته جدول پر قطر لاندې ځو (ځکه ورته «قطري طريقه» وايو)،
او بيا هر \LR{$1$} په \LR{$0$} او هر \LR{$0$} په~\LR{$1$} بدلوو۔
\textbf{د ژباړې د نسخې سمون (OLSIZ-009):} انګرېزي نثر د دويم
بدلون پر ځاے لومړے بدلون بيا تکرار کړے و؛ لاندې اصلي حالت‌تعريف
د صفر او يو دواړه اړخونه ټاکي۔ په لا مجرد ډول، \LR{$d(n)$} د دې له
مخې \LR{$0$} يا \LR{$1$} تعريفوو چې د قطر \LR{$n$}م غړے، \LR{$s_n(n)$}،
\LR{$1$} دے که \LR{$0$}؛ يعنې:
\[
d(n) =
\begin{cases}
1 & \text{\RL{که \LR{$s_{n}(n) = 0$} وي}}\\
0 & \text{\RL{که \LR{$s_{n}(n) = 1$} وي}}
\end{cases}
\]
په څرګند ډول \LR{$d \in \Bin^\omega$}، ځکه دا د \LR{$0$}ګانو او \LR{$1$}ګانو
يو نامتناهي توريز تسلسل دے۔ خو موږ \LR{$d$} داسې جوړ کړے چې د هر
\LR{$n \in \Nat$} دپاره \LR{$d(n) \neq s_n(n)$}۔ يعنې \LR{$d$} له \LR{$s_n$} څخه
په خپل \LR{$n$}م مقام کښې توپير لري۔ نو د هر \LR{$n\in \Nat$} دپاره
\LR{$d \neq s_n$}۔ ځکه \LR{$d$} په لړ \LR{$s_0$}، \LR{$s_1$}، \LR{$s_2$}،
\dots کښې نۀ شي راتلے۔
\olpseganchor{OLP-0039-B012}{end}

\olpseganchor{OLP-0039-B013}{start}
موږ وښودل چې د \LR{$\Bin^\omega$} د کوم فرعي سټ هره خپله سرې شمېرنه
به د \LR{$\Bin^\omega$} کوم غړے پرېږدي۔ نو د سټ
\LR{$\Bin^\omega$} هېڅ شمېرنه نۀ شته؛ يعنې \LR{$\Bin^\omega$}
د شمېر نۀ وړ دے۔
\end{proof}
\olpseganchor{OLP-0039-B013}{end}

\olpseganchor{OLP-0039-B014}{start}
\begin{explain}
د ثبوت دې طريقې ته «قطرول» ويل کېږي، ځکه د~\LR{$d$} د تعريف دپاره
د جدول قطر کاروي۔ خو قطرول د جدول شتون ته اړتيا نۀ لري۔ په
حقيقت کښې، په ورته مفکوره ښودلے شو چې ځينې سټونه
د شمېر نۀ وړ دي، ان چې هېڅ جدول او رښتينے قطر نۀ وي۔
لاندې پايله ښيي چې څنګه۔
\end{explain}
\olpseganchor{OLP-0039-B014}{end}

\olpseganchor{OLP-0039-B015}{start}
\begin{thm}
\ollabel{thm:nonenum-pownat}
\LR{$\Pow{\Nat}$} د شمېر وړ نۀ دے۔
\end{thm}
\olpseganchor{OLP-0039-B015}{end}

\olpseganchor{OLP-0039-B016}{start}
\begin{proof}
په هماغه طريقه مخکښې ځو: ښيو چې د~\LR{$\Nat$} د فرعي سټونو هر لړ
د \LR{$\Nat$} کوم فرعي سټ پرېږدي۔ نو فرض کړئ چې د \LR{$\Nat$} د فرعي
سټونو کوم لړ \LR{$N_0, N_1, N_2, \ldots$} لرو۔ سټ \LR{$D$} داسې تعريفوو:
\LR{$n \in D$} هله او يوازې هله چې \LR{$n \notin N_{n}$}:
\[
D = \Setabs{n \in \Nat}{n \notin N_n}
\]
په څرګند ډول \LR{$D\subseteq \Nat$}۔ خو \LR{$D$} په لړ کښې نۀ شي راتلے۔
بالاخره، د جوړښت له مخې \LR{$n \in D$} هله او يوازې هله چې
\LR{$n\notin N_n$}؛ نو د هر \LR{$n \in \Nat$} دپاره \LR{$D \neq N_n$}۔
\end{proof}
\olpseganchor{OLP-0039-B016}{end}

\olpseganchor{OLP-0039-B017}{start}
\begin{explain}
تېر ثبوت د قطر نوم وانۀ خيست۔ بيا هم که داسې ئې انځور کړئ، قطري
طريقه په کښې ليدلے شئ: وانګېړئ چې سټونه \LR{$N_0$}، \LR{$N_1$}، \dots په
جدولي ترتيب کښې ليکل شوي، چې \LR{$N_n$} په \LR{$n$}م قطار کښې داسې ليکو
چې \LR{$m$} په \LR{$m$}م ستون کښې هله او يوازې هله ليکو که \LR{$m \in N_n$}۔
د بېلګې په ډول، فرض کړئ په هغه لړ کښې لومړي څلور سټونه
\LR{$\{0,1,2,\dots\}$}، \LR{$\{1, 3, 5, \dots\}$}، \LR{$\{0,1,4\}$} او
\LR{$\{2,3,4,\dots\}$} دي؛ نو زموږ جدولي ترتيب به داسې پيل شي:
\[
\begin{array}{r@{}rrrrrrr}
  N_0 = \{ & \mathbf{0}, & 1, & 2, & & & & \dots\}\\
  N_1 = \{ &  & \mathbf{1}, &  & 3, &  & 5, & \dots\}\\
  N_2 = \{ & 0, & 1, &  &  & 4\phantom{,} &  & \}\\
  N_3 = \{ &  &  & 2, & \mathbf{3}, & 4, & & \dots\}\\
  &\vdots & & & & & & \ddots\phantom{\}}
  \end{array}
\]
بيا \LR{$D$} هغه سټ دے چې پر قطر د لاندې تلو له لارې تر لاسه کېږي:
\LR{$n \in D$} هله او يوازې هله اېږدو چې \LR{$n$} پر قطر \emph{نۀ} وي۔
نو په پورته حالت کښې \LR{$0$} او \LR{$1$} پرېږدو،~\LR{$2$} شاملوو،~\LR{$3$}
پرېږدو، او همداسې نور۔
\end{explain}
\olpseganchor{OLP-0039-B017}{end}

\olpseganchor{OLP-0039-B018}{start}
\begin{prob}
په څرګند قطري استدلال وښيئ چې د ټولو تابعو
\LR{$f \colon \Nat \to \Nat$} سټ د شمېر نۀ وړ دے۔ يعنې وښيئ
چې که \LR{$f_1$}، \LR{$f_2$}، \dots د تابعو يو لړ وي او هره
\LR{$f_i\colon \Nat \to \Nat$} وي، نو کومه \LR{$g \colon \Nat \to
\Nat$} شته چې په دې لړ کښې نۀ ده۔
\end{prob}
\olpseganchor{OLP-0039-B018}{end}

\olpseganchor{OLP-0040-B004}{start}
\olfileid{sfr}{siz}{red-alt}
\olpseganchor{OLP-0040-B004}{end}

\olpseganchor{OLP-0040-B005}{start}
\olsection{راکمونه}
\olpseganchor{OLP-0040-B005}{end}

\olpseganchor{OLP-0040-B006}{start}
\begin{editorial}
  دا برخه په راکمونه د شمېر نۀ وړ والے ثابتوي، او پايلې ئې د
  \olref[nen-alt]{sec} له پايلو سره سمون خوري۔ يوه بله، لږه
  تفصيلي بڼه چې د \olref[nen]{sec} له پايلو سره سمون خوري، په
  \olref[red]{sec} کښې ورکړل شوې ده۔
\end{editorial}
\olpseganchor{OLP-0040-B006}{end}

\olpseganchor{OLP-0040-B007}{start}
موږ په يوه قطري استدلال وښودل چې \LR{$\Bin^\omega$} د شمېر نۀ وړ
دے۔ په ورته قطري استدلال مو وښودل چې \LR{$\Pow{\Nat}$}
د شمېر نۀ وړ دے۔ خو د \LR{$\Pow{\Nat}$} د د شمېر نۀ وړ
کېدو يوه بله لار دا ده: وښيئ چې \emph{که \LR{$\Pow{\Nat}$}
د شمېر وړ وي، نو \LR{$\Bin^\omega$} هم د شمېر وړ دے}۔ ځکه
چې پوهېږو \LR{$\Bin^\omega$} د شمېر نۀ وړ دے، دا به وښيي چې
\LR{$\Pow{\Nat}$} هم همداسې دے۔
\olpseganchor{OLP-0040-B007}{end}

\olpseganchor{OLP-0040-B008}{start}
دې ته د يوې مسئلې بلې ته \emph{راکمونه} ويل کېږي۔ په دې حالت
کښې د \LR{$\Bin^\omega$} د شمېرلو مسئله د \LR{$\Pow{\Nat}$} د شمېرلو
مسئلې ته راکموو۔ د وروستۍ مسئلې يو حل، يعنې د \LR{$\Pow{\Nat}$} يوه
شمېرنه، به د لومړۍ مسئلې حل، يعنې د \LR{$\Bin^\omega$} يوه شمېرنه،
راکړي۔
\olpseganchor{OLP-0040-B008}{end}

\olpseganchor{OLP-0040-B009}{start}
د يوه سټ~\LR{$B$} د شمېرلو مسئله د يوه سټ~\LR{$A$} د شمېرلو مسئلې ته د
راکمولو دپاره، د~\LR{$A$} د يوې شمېرنې د~\LR{$B$} په شمېرنه د اړولو يوه
لار ورکوو۔ د دې تر ټولو اسانه لار د يوې پر هدف سټ بشپړه تابع تابع
\LR{$f\colon A \to B$} تعريفول دي۔ که \LR{$x_1$}، \LR{$x_2$}، \dots{} د~\LR{$A$}
شمېرنه وي، نو \LR{$f(x_1)$}، \LR{$f(x_2)$}، \dots{} به د~\LR{$B$} شمېرنه وي۔
زموږ په حالت کښې د يوې پر هدف سټ بشپړه تابع تابع
\LR{$f\colon \Pow{\Nat} \to \Bin^\omega$} په لټه کښې يو۔
\olpseganchor{OLP-0040-B009}{end}

\olpseganchor{OLP-0040-B010}{start}
\begin{prob}
وښيئ چې که يوه يو پر يو تابع \LR{$g\colon B \to A$} وي او
\LR{$B$}~د شمېر نۀ وړ وي، نو \LR{$A$} هم همداسې دے۔ دا په ښودلو
وکړئ چې څنګه د~\LR{$g$} په کارولو د~\LR{$A$} يوه شمېرنه د~\LR{$B$} په شمېرنه
اړولے شئ۔
\end{prob}
\olpseganchor{OLP-0040-B010}{end}

\olpseganchor{OLP-0040-B011}{start}
\begin{proof}[په راکمونه د {\olref[nen-alt]{thm:nonenum-pownat}} ثبوت]
فرض کړئ چې \LR{$\Pow{\Nat}$} د شمېر وړ وے، او له دې امله د هغه
يوه شمېرنه \LR{$N_{1}$}، \LR{$N_{2}$}، \LR{$N_{3}$}، \dots وے۔
\olpseganchor{OLP-0040-B011}{end}

\olpseganchor{OLP-0040-B012}{start}
تابع \LR{$f \colon \Pow{\Nat} \to \Bin^\omega$} داسې تعريف کړئ چې
\LR{$f(N)$} هغه تسلسل \LR{$s$} وي چې \LR{$s(n) = 1$} هله او يوازې هله چې
\LR{$n \in N$} وي، او په بل حالت کښې \LR{$s(n) = 0$} وي۔
\textbf{د ژباړې د نسخې سمون (OLSIZ-010):} په انګرېزي نثر کښې
د تسلسل بې تړاوه شاخص د خپل سري ورودي سټ له نښې سره بدل شو، څو
د تابع تعريف روښانه وي۔
\olpseganchor{OLP-0040-B012}{end}

\olpseganchor{OLP-0040-B013}{start}
دا په څرګند ډول تابع تعريفوي، ځکه کله چې \LR{$N \subseteq \Nat$} وي،
هر \LR{$n \in \Nat$} يا د~\LR{$N$} يو غړے وي يا نۀ وي۔ د بېلګې
په ډول، د طبيعي جفتو عددونو سټ
\LR{$2\Nat = \Setabs{2n}{n \in \Nat} = \{0,2, 4, 6, \dots\}$}
تسلسل \LR{$1010101\dots$} ته نقشېږي، \LR{$\emptyset$} سټ \LR{$0000\dots$} ته او پخپله
سټ \LR{$\Nat$}، \LR{$1111\dots$} ته نقشېږي۔
\olpseganchor{OLP-0040-B013}{end}

\olpseganchor{OLP-0040-B014}{start}
دا پر هدف سټ بشپړه هم ده: د \LR{$0$}ګانو او \LR{$1$}ګانو هر تسلسل د طبيعي
عددونو له يوه سټ سره سمون خوري، يعنې له هغه سټ سره چې غړي ئې
هغه طبيعي عددونه دي کوم چې د تسلسل د~\LR{$1$} لرونکو مقامونو سره
سمون خوري۔ په لا دقيقه توګه، که \LR{$s \in \Bin^\omega$} وي، نو
\LR{$N \subseteq \Nat$} داسې تعريف کړئ:
\[
N = \Setabs{n \in \Nat}{s(n) = 1}
\]
بيا \LR{$f(N) = s$}، لکه چې د~\LR{$f$} تعريف ته په کتلو تاييدولے شو۔
\olpseganchor{OLP-0040-B014}{end}

\olpseganchor{OLP-0040-B015}{start}
اوس دا لړ په پام کښې ونيسئ:
\[
f(N_1), f(N_2), f(N_3), \dots
\]
ځکه چې \LR{$f$} پر هدف سټ بشپړه ده، د \LR{$\Bin^\omega$} هر غړے بايد د کوم
ورودي دپاره د~\LR{$f$} د قيمت په توګه راشي، او نو بايد په لړ کښې راشي۔
له دې امله دا لړ بايد ټول~\LR{$\Bin^\omega$} وشمېري۔
\olpseganchor{OLP-0040-B015}{end}

\olpseganchor{OLP-0040-B016}{start}
نو که \LR{$\Pow{\Nat}$} د شمېر وړ وے، \LR{$\Bin^\omega$} به
د شمېر وړ وے۔ خو \LR{$\Bin^\omega$} د شمېر نۀ وړ دے
(\olref[nen-alt]{thm:nonenum-bin-omega})۔ له دې امله
\LR{$\Pow{\Nat}$} د شمېر نۀ وړ دے۔
\end{proof}
\olpseganchor{OLP-0040-B016}{end}

\olpseganchor{OLP-0040-B017}{start}
%\begin{explain}
%د راکمونې په لوري کښې ګډوډېدل اسانه دي۔ د بېلګې په ډول، يوه
%پر هدف سټ بشپړه تابع $g \colon \Bin^\omega \to X$ دا \emph{نۀ}
%ثابتوي چې $X$ د شمېر نۀ وړ دے۔ (تابع $g
%\colon \Bin^\omega \to \Bin$ په پام کښې ونيسئ چې په
%$g(s) = s(1)$ تعريف شوې؛ دا تابع د $0$ګانو او $1$ګانو يو تسلسل
%د هغه لومړي غړے ته نقشوي۔ دا پر هدف سټ بشپړه ده، ځکه ځينې
%تسلسلونه په $0$ او ځينې په $1$ پيلېږي۔ خو $\Bin$ متناهي دے۔)
%دا هم په ياد ولرئ چې تابع~$f$ بايد پر هدف سټ بشپړه وي، ګنې
%استدلال مخکښې نۀ ځي: بيا به دا ډاډ نۀ وي چې $f(x_1)$، $f(x_2)$،
%\dots{} د~$Y$ ټول غړي لري۔ د بېلګې په ډول، تابع
%$h\colon \Nat \to
%\Bin^\omega$ په لاندې ډول تعريف کړئ:
%\[
%h(n) = \underbrace{000\dots0}_{\text{$n$ $0$ګانې}}
%\]
%دا يوه تابع ده، خو $\Nat$ د شمېر وړ دے۔
%\end{explain}
\olpseganchor{OLP-0040-B017}{end}

\olpseganchor{OLP-0040-B018}{start}
\begin{prob}\label{sfr:siz:red-alt:prob:nat-nat}
  په راکمونې استدلال وښيئ چې د ټولو تابعو \LR{$f\colon \Nat \to \Nat$}
  سټ~\LR{$X$} د شمېر نۀ وړ دے۔ (لارښوونه: له \LR{$X$} څخه
  ~\LR{$\Bin^\omega$} ته يوه پر هدف سټ بشپړه تابع ورکړئ۔)
\end{prob}
\olpseganchor{OLP-0040-B018}{end}

\olpseganchor{OLP-0040-B019}{start}
\begin{prob}
په راکمونې استدلال وښيئ چې د طبيعي عددونو د جوړو د
\emph{سټونو} سټ، يعنې \LR{$\Pow{\Nat \times \Nat}$}،
د شمېر نۀ وړ دے۔
\end{prob}
\olpseganchor{OLP-0040-B019}{end}

\olpseganchor{OLP-0040-B020}{start}
\begin{prob}
په راکمونې استدلال وښيئ چې \LR{$\Nat^\omega$}، يعنې د طبيعي عددونو
د نامتناهي تسلسلونو سټ، د شمېر نۀ وړ دے۔
\end{prob}
\olpseganchor{OLP-0040-B020}{end}

\olpseganchor{OLP-0040-B021}{start}
%\begin{prob}
%فرض کړئ $P$ له $\Nat$ څخه سټ~$\{0\}$ ته د تابعو سټ دے، او $Q$
%له مثبتو صحيح عددونو څخه سټ~$\{0\}$ ته د \emph{جزوي} تابعو سټ
%دے۔ وښيئ چې $P$~د شمېر وړ دے او $Q$ نۀ دے۔ (لارښوونه: د
%$\Bin^\omega$ د شمېرلو مسئله د~$Q$ شمېرلو ته راکمه کړئ۔)
%\end{prob}
\olpseganchor{OLP-0040-B021}{end}

\olpseganchor{OLP-0040-B022}{start}
\begin{prob}
فرض کړئ \LR{$S$} له \LR{$\Nat$} څخه سټ~\LR{$\{0,1\}$} ته د ټولو
پر هدف سټ بشپړې تابعې سټ دے؛ يعنې \LR{$S$} له ټولو پر هدف سټ بشپړې تابعې
~\LR{$f \colon \Nat \to \Bin$} څخه جوړ دے۔ وښيئ چې \LR{$S$}
د شمېر نۀ وړ دے۔
\end{prob}
\olpseganchor{OLP-0040-B022}{end}

\olpseganchor{OLP-0040-B023}{start}
\begin{prob}
وښيئ چې د ټولو حقيقي عددونو سټ~\LR{$\Real$} د شمېر نۀ وړ دے۔
\end{prob}
\olpseganchor{OLP-0040-B023}{end}

\olpseganchor{OLP-0041-B004}{start}
\olchapter{sfr}{arith}{د شمېر نظامونو سټي جوړونه}
\olpseganchor{OLP-0041-B004}{end}

\olpseganchor{OLP-0041-B005}{start}
\begin{editorial}
په دې څپرکي کښې مواد د ساده سټ نظريې په دننه کښې د شمېر نظامونو
جوړښت وړاندې کوي۔ دا د ټم بټن له \emph{اوپن سېټ تيوري} متن څخه
اخېستل شوي دي۔
\end{editorial}
\olpseganchor{OLP-0041-B005}{end}

\olpseganchor{OLP-0042-B004}{start}
\olfileid{sfr}{arith}{int}
\olsection{له \LR{$\Nat$} څخه \LR{$\Int$} ته}
\olpseganchor{OLP-0042-B004}{end}

\olpseganchor{OLP-0042-B005}{start}
دا دوه بنسټيزې خبرې په پام کښې ونيسئ:
	\begin{enumerate}
		\item هر صحيح عدد د \LR{$n - m$} په بڼه ليکل کېدے شي، چې \LR{$n, m \in\Nat$} وي۔
		\item په عبارت \LR{$n - m$} کښې زېرمه شوي معلومات په هماغه ډول په مرتبه جوړه \LR{$\tuple{n, m}$} کښې زېرمه کېدے شي۔
	\end{enumerate}
موږ لا پوهېږو چې د طبيعي عددونو مرتبې جوړې د \LR{$\Nat^2$} غړي دي۔ او دا فرضوو چې په \LR{$\Nat$} پوهېږو۔ نو پر دغو دوو خبرو ولاړ يو ساده وړانديز دا دے: \emph{صحيح عددونه دې د طبيعي عددونو د مرتبو جوړو په توګه وګڼو}۔
\olpseganchor{OLP-0042-B005}{end}

\olpseganchor{OLP-0042-B006}{start}
په حقيقت کښې دا وړانديز ډېر ساده دے۔ ښکاره ده چې غواړو \LR{$0- 2 = 4 - 6$} وي۔ خو په څرګند ډول \LR{$\tuple{0, 2 }\neq \tuple{4, 6}$} دے۔ نو په ساده ډول نۀ شو ويلے چې \LR{$\Nat^2$} د صحيح عددونو سټ دے۔
\olpseganchor{OLP-0042-B006}{end}

\olpseganchor{OLP-0042-B007}{start}
د مخکښې ستونزې په عمومي کولو، موږ لاندې خاصيت غواړو:
	$$a - b = c - d \text{\RL{ هله او يوازې هله چې }}a + d = c + b$$
(دا بايد څرګنده وي چې صحيح عددونه په همدې ډول \emph{چلېدل} غواړي: يوازې دواړو خواوو ته \LR{$b$} او \LR{$d$} ورزيات کړئ۔) د دې چلند د يقيني کولو اسانه لار دا ده چې د مرتبو جوړو تر منځ د هم ارزښتۍ يوه اړيکه \LR{$\Intequiv$} داسې تعريف کړو:
$$\tuple{ a, b } \Intequiv \tuple{c, d}\text{\RL{ هله او يوازې هله چې }}a + d = c + b$$
اوس بايد وښيو چې دا د هم ارزښتۍ اړيکه ده۔
\begin{prop} \LR{$\Intequiv$} د هم ارزښتۍ اړيکه ده۔
\end{prop}
	\begin{proof}
	بايد وښيو چې \LR{$\Intequiv$} انعکاسي، تناظري او انتقالي ده۔
\olpseganchor{OLP-0042-B007}{end}
	
\olpseganchor{OLP-0042-B008}{start}
	\emph{انعکاسيت:} په څرګند ډول \LR{$\tuple{a, b} \Intequiv \tuple{a, b}$}، ځکه چې \LR{$a + b = b + a$}۔
\olpseganchor{OLP-0042-B008}{end}
	
\olpseganchor{OLP-0042-B009}{start}
	\emph{تناظر:} فرض کړئ \LR{$\tuple{a, b} \Intequiv \tuple{c, d}$}، نو \LR{$a + d = c + b$}۔ بيا \LR{$c + b = a + d$}، نو \LR{$\tuple{c, d} \Intequiv \tuple{a, b}$}۔
\olpseganchor{OLP-0042-B009}{end}
	
\olpseganchor{OLP-0042-B010}{start}
	\emph{انتقاليت:} فرض کړئ \LR{$\tuple{a, b} \Intequiv \tuple{c, d}\Intequiv \tuple{m, n}$}۔ نو \LR{$a + d = c + b$} او \LR{$c + n = m + d$}۔ له دې امله \LR{$a + d + c + n = c + b + m + d$}، او نو \LR{$a + n = m + b$}۔ په پايله کښې \LR{$\tuple{a, b } \Intequiv \tuple{m, n}$}۔		
\end{proof}
\olpseganchor{OLP-0042-B010}{end}

\olpseganchor{OLP-0042-B011}{start}
اوس په دې د هم ارزښتۍ اړيکه د هم ارزښتۍ ټولګي اخيستلے شو:
\olpseganchor{OLP-0042-B011}{end}

\olpseganchor{OLP-0042-B012}{start}
\begin{defn}
		صحيح عددونه د \LR{$\Intequiv$} له مخې د طبيعي عددونو د مرتبو جوړو د هم ارزښتۍ ټولګي دي؛ يعنې \LR{$\Int = \equivclass{\Nat^2}{\Intequiv}$}۔
\end{defn}
\olpseganchor{OLP-0042-B012}{end}

\olpseganchor{OLP-0042-B013}{start}
د دې ټاکلي تعريف په اړه ښايي بېلابېل \emph{فلسفي} غبرګونونه وي۔ خو د هغو غبرګونونو تر څېړلو مخکښې، ښه ده چې ځينې فني جزئيات پر مخ يوسو۔
\olpseganchor{OLP-0042-B013}{end}

\olpseganchor{OLP-0042-B014}{start}
چې صحيح عددونه مو تعريف کړل، پر هغو به بنسټيزې تابعې او اړيکې هم تعريفول غواړو۔ \LR{$\equivrep{m, n}{\Intequiv}$} دې د \LR{$\Intequiv$} له مخې هغه د هم ارزښتۍ ټولګے وي چې \LR{$\tuple{m, n}$} ئې يو غړے دے۔\footnote{يادونه: په \olref[sfr][rel][eqv]{def:equivalenceclass} کښې د معرفي شوې نښې په کارولو به مو د همدې څيز دپاره \LR{$\equivrep{\tuple{m,n}}{\Intequiv}$} ليکل۔ خو هغه لوستل لږ ګران دي۔} يعنې:
	$$\equivrep{m, n}{\Intequiv} = \Setabs{\tuple{a, b} \in \Nat^2}{\tuple{a, b}\Intequiv \tuple{m, n}}$$
نو اوس لاندې تعريفونه وړاندې کوو:
	\begin{align*}
		\equivrep{a, b}{\Intequiv} + \equivrep{c, d}{\Intequiv} &= \equivrep{a + c, b + d}{\Intequiv}\\
		\equivrep{a, b}{\Intequiv} \times \equivrep{c, d}{\Intequiv} &= \equivrep{a c + b  d, a  d + b c}{\Intequiv}\\
		\equivrep{a, b}{\Intequiv} \leq \equivrep{c, d}{\Intequiv} &\text{\RL{ هله او يوازې هله چې }}a + d \leq b + c
	\end{align*}	
(لکه چې دود دے، د اصولو د اسانه لوستلو دپاره `\LR{$ab$}' د `\LR{$(a \times b)$}' پر ځاے کاروم۔) اوس بايد ډاډمن شو چې دا تعريفونه لکه څنګه چې \emph{بايد}، هماغسې چلېږي۔ د دې مطلب په تفصيل بيانول او بيا ئې ارزول ډېر اوږد کار دے؛ جزئيات \olref[check]{sec} ته پرېږدو۔ خو لنډه خبره دا ده: هر څه سم کار کوي!
\olpseganchor{OLP-0042-B014}{end}

\olpseganchor{OLP-0042-B015}{start}
يو وروستے کار پاتې دے۔ موږ صحيح عددونه د طبيعي عددونو په
کارولو جوړ کړي دي۔ خو د دې معنا به دا وي چې طبيعي عددونه \emph{پخپله
صحيح عددونه نۀ دي}۔ د دې فلسفي ارزښت ته به په \olref[ref]{sec}
کښې ورستانۀ شو۔ خو په سوچه فني کچه به يوې داسې لارې ته اړتيا لرو
چې طبيعي عددونه د صحيح عددونو \emph{په توګه} وکارولے شو۔ مفکوره
ډېره اسانه ده: د هر \LR{$n \in \Nat$} دپاره ټاکو چې
\LR{$n_\Int = \equivrep{n, 0}{\Intequiv}$}۔ بايد تاييد کړو چې دا تعريف
سم چلند کوي؛ يعنې د هر \LR{$m, n \in \Nat$} دپاره
\begin{align*}
	(m + n)_\Int &= m_\Int + n_\Int\\
	(m \times n)_\Int &= m_\Int \times n_\Int\\
	m \leq n &\liff m_\Int \leq n_\Int
\end{align*}
خو دا ټول پوره ساده دي۔ د بېلګې په ډول، د دې د ښودلو دپاره چې
دويم خاصيت سم دے، د طبيعي عددونو له چلنده په ساده ډول کار اخلو او
لاندې استدلال کوو:
%\begin{proof}
%	د طبيعي عددونو له چلنده په کار اخيستلو، په څرګند ډول:
	\begin{align*}
%		(m + n)_\Int  = \equivrep{m + n, 0}{\Intequiv} = \equivrep{m + n, 0 + 0}{\Intequiv} = \equivrep{m, 0}{\Intequiv} + \equivrep{n, 0}{\Intequiv} = m_\Int + n_\Int\\
		(m \times n)_\Int  &= \equivrep{m \times n, 0}{\Intequiv} \\
		&= \equivrep{m \times n + 0 \times 0, m \times 0 + 0 \times n}{\Intequiv} \\
		&= \equivrep{m, 0}{\Intequiv} \times \equivrep{n, 0}{\Intequiv} \\
		&= m_\Int \times n_\Int
	\end{align*}
%\end{proof}
د نورو دوو شرطونو تاييد د تمرين په توګه پرېږدو۔
\begin{prob}
	وښيئ چې د هر \LR{$m, n \in \Nat$} دپاره \LR{$(m + n)_\Int = m_\Int + n_\Int$} او \LR{$m \leq n \liff m_\Int \leq n_\Int$}۔
\end{prob}
\olpseganchor{OLP-0042-B015}{end}

\olpseganchor{OLP-0043-B004}{start}
\olfileid{sfr}{arith}{rat}
\olsection{له \LR{$\Int$} څخه \LR{$\Rat$} ته}
\olpseganchor{OLP-0043-B004}{end}

\olpseganchor{OLP-0043-B005}{start}
همدا اوس مو وليدل چې په ساده سټ نظريه څنګه له طبيعي عددونو څخه
صحيح عددونه جوړېږي۔ اوس به په ډېرې ورته لارې له صحيح عددونو څخه
ناطق عددونه جوړول وګورو۔ زموږ لومړنۍ خبرې دا دي:
\begin{enumerate}
	\item هر ناطق عدد د \LR{$\nicefrac{i}{j}$} په بڼه ليکل کېدے شي،
	چې \LR{$i$} او \LR{$j$} دواړه صحيح عددونه وي خو \LR{$j$} صفر نۀ وي۔
	\item په عبارت \LR{$\nicefrac{i}{j}$} کښې زېرمه شوي معلومات په هماغه
	ډول په مرتبه جوړه \LR{$\tuple{ i, j}$} کښې زېرمه کېدے شي۔
\end{enumerate}
ښکاره لار به دا وي چې ناطق عددونه د \LR{$\Int \times
(\Int \setminus \{0_\Int\})$} څخه د اخيستل شوو مرتبو جوړو \emph{په
توګه} وګڼو۔ خو لکه مخکښې، دا به لږ ډېر ساده وي، ځکه موږ غواړو
\LR{$\nicefrac{3}{2} = \nicefrac{6}{4}$} وي، خو \LR{$\tuple{ 3, 2}\neq
\tuple{ 6, 4}$}۔ په عمومي ډول لاندې خاصيت غواړو:
\[
	\nicefrac{a}{b} = \nicefrac{c}{d} \text{\RL{ هله او يوازې هله چې }} a \times d = b \times c
\]
د دې د ترلاسه کولو دپاره پر \LR{$\Int \times (\Int \setminus
\{0_\Int\})$} د هم ارزښتۍ يوه اړيکه داسې تعريفوو:
\[
	\tuple{ a, b }\Ratequiv \tuple{ c, d} \text{\RL{ هله او يوازې هله چې }}a \times d = b \times c
\]
بايد وارزوو چې دا د هم ارزښتۍ اړيکه ده۔ دا د \LR{$\Intequiv$} حالت ته
ډېر ورته دے، او موږ ئې د تمرين په توګه پرېږدو۔
\begin{prob}
وښيئ چې \LR{$\Ratequiv$} د هم ارزښتۍ اړيکه ده۔
\end{prob}
خو دا موږ ته د لاندې خبرې اجازه راکوي:
\begin{defn}
ناطق عددونه د \LR{$\Ratequiv$} له مخې د صحيح عددونو د جوړو د هم ارزښتۍ
ټولګي دي (چې دويم غړے ئې صفر نۀ وي)۔ يعنې \LR{$\Rat =
\equivclass{(\Int \times (\Int\setminus \{0_\Int\}))}{\Ratequiv}$}۔
\end{defn}
\olpseganchor{OLP-0043-B005}{end}

\olpseganchor{OLP-0043-B006}{start}
لکه د صحيح عددونو په شان، ځينې بنسټيزې عمليې هم تعريفول غواړو۔
چې \LR{$\equivrep{i,j}{\Ratequiv}$} د \LR{$\Ratequiv$} له مخې هغه د هم
ارزښتۍ ټولګے وي چې \LR{$\tuple{i, j}$} ئې يو غړے دے، وايو:
\begin{align*}
	\equivrep{a, b}{\Ratequiv} + \equivrep{c, d}{\Ratequiv} &= \equivrep{ad + bc,  bd}{\Ratequiv}\\
	\equivrep{a, b}{\Ratequiv} \times \equivrep{c, d}{\Ratequiv} &= \equivrep{a  c, b d}{\Ratequiv}.
\intertext{\RL{پر دې ناطقو عددونو د \LR{$r \leq s$} د تعريفولو دپاره دا حقيقت کاروو چې
\LR{$r \le s$} هله او يوازې هله چې \LR{$s - r$} منفي نۀ وي؛ يعنې \LR{$r - s$} د
\LR{$\nicefrac{i}{j}$} په بڼه ليکل کېدے شي چې \LR{$i$} نامنفي او \LR{$j$}~مثبت وي۔
\textbf{د ژباړې د نسخې سمون (OLARI-001):} په انګرېزي نثر کښې د وروستي توپير ترتيب سرچپه ليکل شوے؛ لاندې معادله د نامنفي توپير سم ترتيب ښيي:}}
	\equivrep{a, b}{\Ratequiv} \leq \equivrep{c, d}{\Ratequiv} &\text{\RL{ هله او يوازې هله چې }}
	\equivrep{c, d}{\Ratequiv} - \equivrep{a, b}{\Ratequiv} =
	\equivrep{i_\Int, j_\Int}{\Ratequiv}
\end{align*}
د ځينو \LR{$i \in \Nat$} او \LR{$0 \neq j \in \Nat$} دپاره۔
\olpseganchor{OLP-0043-B006}{end}

\olpseganchor{OLP-0043-B007}{start}
بيا بايد وارزوو چې دا تعريفونه لکه څنګه چې \emph{بايد}، هماغسې
چلېږي؛ او دا کار \olref[check]{sec} ته پرېږدو۔ خو رښتيا هم سم
چلېږي!{} په پاے کښې صحيح عددونه د ناطقو عددونو \emph{په توګه}
د کارولو يوه لار غواړو؛ نو د هر \LR{$i \in \Int$} دپاره ټاکو چې
\LR{$i_\Rat = \equivrep{i, 1_\Int}{\Ratequiv}$}۔ يو ځل بيا په
\olref[check]{sec} کښې ارزوو چې دا ټول سم چلند کوي۔
\olpseganchor{OLP-0043-B007}{end}

\olpseganchor{OLP-0043-B008}{start}
\begin{prob}
وښيئ چې د هر \LR{$i, j \in \Int$} دپاره \LR{$(i + j)_\Rat = i_\Rat+ j_\Rat$}
او \LR{$(i \times j)_\Rat = i_\Rat \times j_\Rat$} او
\LR{$i \leq j \liff i_\Rat \leq j_\Rat$}۔
\end{prob}
\olpseganchor{OLP-0043-B008}{end}

\olpseganchor{OLP-0044-B004}{start}
\olfileid{sfr}{arith}{real}
\olsection{د حقيقي عددونو کرښه}
\olpseganchor{OLP-0044-B004}{end}

\olpseganchor{OLP-0044-B005}{start}
بل ګام دا ښودل دي چې له ناطقو عددونو څخه حقيقي عددونه څنګه
جوړېږي۔ تر هغې مخکښې بايد پوه شو چې د حقيقي عددونو
\emph{ځانګړتيا} څه ده۔
\olpseganchor{OLP-0044-B005}{end}

\olpseganchor{OLP-0044-B006}{start}
حقيقي عددونه ناطقو عددونو ته ډېر ورته چلند کوي۔ (په فني ډول، دواړه
د \emph{مرتبو ميدانونو} بېلګې دي؛ د دې د تعريف دپاره
\olref[check]{orderedfield} وګورئ۔) اوس که مو د
\olref[sfr][siz][]{chap} تمرينونه حل کړي وي، نو پوهېږئ چې حقيقي
عددونه له ناطقو عددونو په کلکه زيات دي؛ يعنې
\LR{$\cardless{\Rat}{\Real}$}۔ دا خبره لومړے کانتور ثابته کړه۔ خو شاوخوا
دوه نيم زره کاله کېږي چې غير ناطق عددونه، يعنې هغه حقيقي عددونه چې
ناطق نۀ دي، پېژندل شوي دي۔ په رښتيا:
%\footnote{اوس: ``$\sqrt{2}$'' مثبت او منفي دواړو جذرونو ته يو شان اشاره کوي؛ خو په راتلونکو څو مخونو کښې به دا هېر کړم او يوازې مثبت جذر به په پام کښې نيسم۔}
\begin{thm}\ollabel{root2irrational}
\LR{$\sqrt{2}$} ناطق نۀ دے؛ يعنې \LR{$\sqrt{2} \notin \Rat$}
\end{thm}
\olpseganchor{OLP-0044-B006}{end}

\olpseganchor{OLP-0044-B007}{start}
\begin{proof}
د تناقض دپاره فرض کړئ چې \LR{$\sqrt{2}$} ناطق دے۔ نو د ځينو طبيعي
عددونو \LR{$m$} او \LR{$n$} دپاره \LR{$\sqrt{2} = \nicefrac{m}{n}$}۔ په حقيقت
کښې \LR{$m$} او \LR{$n$} داسې ټاکلے شو چې کسر نور نۀ شي ساده کېدے۔ په بيا
ترتيبولو \LR{$m^{2} = 2n^{2}$}۔ له دې ځايه ثبوت په دوو لارو بشپړولے شو:
\olpseganchor{OLP-0044-B007}{end}

\olpseganchor{OLP-0044-B008}{start}
\emph{لومړے، په هندسي ډول} (د ټېننباوم په پيروۍ)۔\footnote{دا
ثبوت \cite{Conway2006} روايت کړے دے۔} دا مربعګانې په پام کښې ونيسئ:
\begin{center}	
	\begin{tikzpicture}
		\draw[thick] (0,0) rectangle (3,3);
		\draw[thick, fill=red!50] (0,0) rectangle (2.3,2.3);
		\draw[thick, fill=yellow!50] (0.7,0.7) rectangle (3, 3);
		\draw[thick, fill=orange!50] (0.7,0.7) rectangle (2.3, 2.3);
		\draw[<->] (4, 0.7)--(4, 3);
		\node at (4.25, 1.85) (n) {\LR{$n$}};
		\draw[<->] (5, 0)--(5, 3);
		\node at (5.25, 1.5) (m) {\LR{$m$}};
	\end{tikzpicture}
\end{center}
ځکه چې \LR{$m^2 = 2n^2$}، هغه سيمه چې د ضلعې \LR{$n$} دواړه مربعګانې
يو پر بل راځي، له هغې سيمې سره مساوي مساحت لري چې له دغو دواړو
مربعګانو هېڅ يو نۀ ده پوښلې؛ يعنې د نارنجي مربع مساحت د دواړو
بې رنګه مربعګانو د مساحتونو له مجموعې سره مساوي دے۔ نو که د نارنجي
مربع ضلعه \LR{$p$} او د هر بې رنګه مربع ضلعه \LR{$q$} وي، نو \LR{$p^2 = 2q^2$}۔
خو اوس \LR{$\sqrt{2} = \nicefrac{p}{q}$}، چې \LR{$p < m$}، \LR{$q < n$} او
\LR{$p, q \in \Nat$} دي۔ دا له دې خبرې سره تناقض لري چې \LR{$m$} او \LR{$n$}
تر ټولو واړه ممکن عددونه ټاکل شوي وو۔
\olpseganchor{OLP-0044-B008}{end}
		
\olpseganchor{OLP-0044-B009}{start}
\emph{دويم، په صوري ډول۔} ځکه چې \LR{$m^{2} = 2n^{2}$}، نو \LR{$m$} جفت
دے۔ (دا ښودل اسانه دي چې که \LR{$x$} طاق وي، نو \LR{$x^2$} هم طاق وي۔)
نو د کوم \LR{$r \in \Nat$} دپاره \LR{$m = 2r$}۔ په بيا ترتيبولو
\LR{$2r^2 = n^2$}،
		%				\begin{align*}
		%					(2q)^{2} &= 2n^{2}\\
		%					4q^{2} &= 2n^{2}\\
		%					2q^{2} &= n^{2}
		%				\end{align*}
نو \LR{$n$} هم جفت دے۔ په دې توګه \LR{$m$} او \LR{$n$} دواړه جفت دي، او له دې
امله کسر \LR{$\nicefrac{m}{n}$} نور \emph{ساده کېدے} شي۔ تناقض!
\end{proof}
\olpseganchor{OLP-0044-B009}{end}

\olpseganchor{OLP-0044-B010}{start}
په ضمن کښې دا انځوريز ثبوت موږ ته اجازه راکوي چې د \emph{د OpenLogic اړونده سرچينيزه برخه}\nobreakspace\LR{\href{https://github.com/OpenLogicProject/OpenLogic/blob/9620cc73f9c8e0ad003c514a5d3748f29611c4c0/content/history/set-theory/mythology.tex\#L7}{\latinfont\scriptsize [OLP]}} مواد بيا وګورو۔ ټېننباوم (1927--2006) بشپړ معاصر رياضي پوه و؛ خو ثبوت بې شکه ښکلے، پوره دقيق او پر هندسي پوهېدنه ولاړ دے!
\olpseganchor{OLP-0044-B010}{end}

\olpseganchor{OLP-0044-B011}{start}
په هر حال، حقيقي عددونه تر ناطقو عددونو «ډېر پراخ» دي۔ په يوه معنا په ناطقو عددونو کښې «تشيالونه» شته چې حقيقي عددونه ئې ډکوي۔ وايرشټراس وپېژندله چې دا د حقيقي عددونو يو خاصيت بيانوي، کوم چې هغوے له ناطقو عددونو بېلوي؛ يعنې د بشپړتيا خاصيت: \emph{د حقيقي عددونو هر ناتش سټ چې پاسنی حد لري، تر ټولو لږ پاسنی حد هم لري۔}
\olpseganchor{OLP-0044-B011}{end}

\olpseganchor{OLP-0044-B012}{start}
دا ليدل اسانه دي چې ناطق عددونه د بشپړتيا خاصيت نۀ لري۔ د بېلګې په ډول له \LR{$\sqrt{2}$} څخه د وړو ناطقو عددونو دا سټ په پام کښې ونيسئ:
\[
	\Setabs{p \in \Rat}{p^2 < 2 \text{\RL{ يا }}p < 0}
\]
دا په ناطقو عددونو کښې پاسنی حد لري؛ د بېلګې په ډول د دې ټول
غړي له \LR{$3$} څخه واړه دي۔ \textbf{د ژباړې د نسخې سمون
(OLARI-002):} په انګرېزي سرچينه کښې د غړي له نښان وروسته يو بې
ځايه «کم تر» نښان پاتې و؛ دلته حذف شوے دے۔ خو د دې تر ټولو لږ
پاسنی حد څه دے؟ غواړو ووايو `\LR{$\sqrt{2}$}'؛ خو همدا اوس مو وليدل
چې \LR{$\sqrt{2}$} \emph{ناطق نۀ} دے۔ او له \LR{$\sqrt{2}$} څخه لوے تر ټولو
لږ ناطق عدد نشته۔ نو دا سټ پاسنی حد لري خو تر ټولو لږ پاسنی حد نۀ
لري۔ له دې امله ناطق عددونه د بشپړتيا خاصيت نۀ لري۔
\olpseganchor{OLP-0044-B012}{end}

\olpseganchor{OLP-0044-B013}{start}
د دې پر خلاف، پيوستار «له معنوي پلوه بايد» د بشپړتيا خاصيت ولري۔ موږ يوازې دا نۀ غواړو چې \LR{$\sqrt{2}$} يو حقيقي عدد وي؛ غواړو د ناطقو عددونو د کرښې ټول «تشيالونه» ډک کړو۔ په رښتيا غواړو چې پخپله پيوستار هېڅ «تشيال» ونۀ لري۔ همدا هغه څه دي چې د بشپړتيا له لارې به ئې ترلاسه کړو۔
\olpseganchor{OLP-0044-B013}{end}

\olpseganchor{OLP-0045-B003}{start}
\olfileid{sfr}{arith}{cuts}
\olsection{له \LR{$\Rat$} څخه \LR{$\Real$} ته}
\olpseganchor{OLP-0045-B003}{end}

\olpseganchor{OLP-0045-B004}{start}
په اصل کښې د بشپړتيا خاصيت ښيي چې د حقيقي عددونو د کرښې هر ټکی
\LR{$\alpha$} دا کرښه په بشپړ ډول پر دوو برخو وېشي: هغه عددونه چې
\LR{$\alpha$} ئې تر ټولو لږ پاسنی حد دے، او هغه عددونه چې \LR{$\alpha$} ئې
تر ټولو لوے لاندينی حد دے۔ ډېډېکېنډ وړانديز وکړ چې له ناطقو عددونو
څخه د حقيقي عددونو د \emph{جوړولو} دپاره، حقيقي عددونه يوازې د هغو
\emph{پرېکونونو} په توګه وګڼو چې ناطق عددونه وېشي۔ يعنې \LR{$\sqrt{2}$}
له هغه \emph{پرېکون} سره يو ګڼو چې ناطق عددونه \LR{$< \sqrt{2}$} له
ناطق عددونو \LR{$> \sqrt{2}$} څخه بېلوي۔
\olpseganchor{OLP-0045-B004}{end}

\olpseganchor{OLP-0045-B005}{start}
راځئ دا خبره سمه منظمه کړو۔ که ناطق عددونه پر دوو برخو پرې کړو،
نو يوازې د هغې \emph{لاندېنۍ} برخې په کتلو خپل وېش په يوازيني ډول
پېژندلے شو۔ نو په دقيق ډول لاندې تعريف وړاندې کوو:
\olpseganchor{OLP-0045-B005}{end}

\olpseganchor{OLP-0045-B006}{start}
\begin{defn}[پرېکون]
يو \emph{پرېکون} \LR{$\alpha$} د ناطقو عددونو هره ناتشه خاصه پيلنۍ
برخه ده چې تر ټولو لوے غړے نۀ لري۔ يعنې \LR{$\alpha$} هله او يوازې هله
پرېکون دے چې:
\begin{enumerate}
	\item \emph{ناتش، خاص}: \LR{$\emptyset \neq \alpha \subsetneq \Rat$}
	\item \emph{پيلنۍ}: د ټولو \LR{$p,q \in \Rat$} دپاره: که \LR{$p < q \in \alpha$} وي، نو \LR{$p \in \alpha$}
	\item \emph{تر ټولو لوے غړے نۀ لري}: د هر \LR{$p \in \alpha$} دپاره داسې \LR{$q \in \alpha$} شته چې \LR{$p < q$} وي
\end{enumerate}
بيا \LR{$\Real$} د پرېکونونو سټ دے۔
\end{defn}
\olpseganchor{OLP-0045-B006}{end}

\olpseganchor{OLP-0045-B007}{start}
نو اوس ويلے شو چې \LR{$\sqrt{2} = \Setabs{p \in \Rat}{p^2 < 2\text{\RL{
يا }}p < 0}$}۔ البته بايد وارزوو چې دا په رښتيا يو \emph{پرېکون}
دے، خو دا کار \olref[check]{sec} ته پرېږدو۔
\olpseganchor{OLP-0045-B007}{end}

\olpseganchor{OLP-0045-B008}{start}
لکه مخکښې، چې ځينې څيزونه مو تعريف کړل، بيا بايد پر هغو بنسټيزې
تابعې او اړيکې تعريف کړو۔ په يوه اسانه تعريف پيل کوو:
\begin{align*}
	\alpha \leq \beta \text{\RL{ هله او يوازې هله چې }}\alpha \subseteq \beta
\end{align*}
د مرتبې دا تعريف موږ ته اجازه راکوي چې مرکزي پايله \emph{بيان}
کړو: د پرېکونونو سټ د بشپړتيا خاصيت لري۔ په بشپړ ډول د بيان بڼه
داسې ده۔ که \LR{$S$} د پرېکونونو يو ناتش سټ وي چې پاسنی حد لري، نو
\LR{$S$} تر ټولو لږ پاسنی حد لري۔ په لا تفصيل: يو پرېکون \LR{$\lambda$} شته
چې د \LR{$S$} پاسنی حد دے، يعنې \LR{$(\forall \alpha \in S)\alpha \subseteq
\lambda$}، او \LR{$\lambda$} تر ټولو لږ داسې پرېکون دے، يعنې
\LR{$(\forall \beta \in \Real)((\forall \alpha \in S)\alpha \subseteq \beta \lif \lambda \subseteq \beta)$}۔ اوس د پايلې ثبوت دا دے:
\olpseganchor{OLP-0045-B008}{end}

\olpseganchor{OLP-0045-B009}{start}
\begin{thm}\ollabel{realcompleteness}
د پرېکونونو سټ د بشپړتيا خاصيت لري۔
\end{thm}
\olpseganchor{OLP-0045-B009}{end}

\olpseganchor{OLP-0045-B010}{start}
\begin{proof}
فرض کړئ \LR{$S$} د پرېکونونو هر ناتش سټ دے چې پاسنی حد لري۔ وټاکئ
\LR{$\lambda = \bigcup S$}۔
%\Setabs{p \in \Rat}{(\exists \alpha \in S)p \in \alpha}$$
لومړے ادعا کوو چې \LR{$\lambda$} يو پرېکون دے:
\begin{enumerate}
\item ځکه چې \LR{$S$} ناتش دے، لږ تر لږه يو پرېکون په \LR{$S$} کښې دے، نو
\LR{$\emptyset \neq \lambda$}۔ \textbf{د ژباړې د نسخې سمون
(OLARI-003):} انګرېزي سرچينه دا ګام په تېروتنه د پاسني حد له فرض
سره تړي؛ اړين دليل د همدې سټ ناتش والے دے۔ ځکه چې \LR{$S$} د پرېکونونو سټ
دے، \LR{$\lambda \subseteq \Rat$}۔ ځکه چې \LR{$S$} پاسنی حد لري، کوم
\LR{$p \in \Rat$} د هر پرېکون \LR{$\alpha \in S$} څخه بهر دے۔ نو
\LR{$p\notin \lambda$}، او له دې امله \LR{$\lambda \subsetneq \Rat$}۔
\item فرض کړئ \LR{$p < q \in \lambda$}۔ نو داسې \LR{$\alpha \in S$} شته
چې \LR{$q \in \alpha$}۔ ځکه چې \LR{$\alpha$} يو پرېکون دے،
\LR{$p \in \alpha$}۔ نو \LR{$p \in \lambda$}۔
\item فرض کړئ \LR{$p \in \lambda$}۔ نو داسې \LR{$\alpha \in S$} شته چې
\LR{$p \in \alpha$}۔ ځکه چې \LR{$\alpha$} يو پرېکون دے، داسې \LR{$q \in
\alpha$} شته چې \LR{$p < q$}۔ نو \LR{$q \in \lambda$}۔
\end{enumerate}
په دې سره ادعا ثابته شوه۔ برسېره پر دې، په څرګند ډول
\LR{$(\forall \alpha \in S)\alpha \subseteq \bigcup S = \lambda$}؛
يعنې \LR{$\lambda$} د \LR{$S$} يو پاسنی حد دے۔ اوس فرض کړئ
\LR{$\beta \in \mathbb{R}$} هم پاسنی حد دے؛ يعنې
\LR{$(\forall \alpha \in S)\alpha \subseteq \beta$}۔ د هر
\LR{$p \in \Rat$} دپاره، که \LR{$p \in \lambda$} وي، نو داسې
\LR{$\alpha \in S$} شته چې \LR{$p \in \alpha$}، او په دې توګه
\LR{$p \in \beta$}۔ په عمومي کولو \LR{$\lambda \subseteq \beta$}۔ نو
\LR{$\lambda$} د \LR{$S$} \emph{تر ټولو لږ} پاسنی حد دے۔
\end{proof}
\olpseganchor{OLP-0045-B010}{end}

\olpseganchor{OLP-0045-B011}{start}
نو موږ داسې ډېر څيزونه لرو چې د بشپړتيا خاصيت پوره کوي۔ د دې د
ويلو يوه لار دا ده: زموږ په پرېکونونو کښې هېڅ «تشيال» نشته۔
(نو د ناطقو عددونو پر ځاے د حقيقي عددونو نور «پرېکونونه» اخيستل
به هېڅ په زړه پورې نوي څيزونه رانۀ کړي۔)
\olpseganchor{OLP-0045-B011}{end}

\olpseganchor{OLP-0045-B012}{start}
بيا بايد پر حقيقي عددونو ځينې عمليې تعريف کړو۔ ناطق عددونه په
حقيقي عددونو کښې په دې ټاکنه ورګډوو چې د هر \LR{$p \in \Rat$} دپاره
\LR{$p_\Real = \Setabs{q \in \Rat}{q < p}$}۔ بيا تعريفوو:
\begin{align*}
	\alpha + \beta &= \Setabs{p + q}{p \in \alpha \land q \in \beta}\\
%	\alpha - \beta &\defis \Setabs{p - q}{p \in \alpha \land q \in \Rat \setminus \beta}\\
	\alpha \times \beta &=
	\Setabs{p \times q}{0 \leq p \in \alpha \land 0 \leq q \in \beta} \cup 0^\mathbb{R} & \text{\RL{که }}\alpha, \beta \geq 0_\Real
%	\alpha \div \beta &\defis
%	\Setabs{\nicefrac{p}{q}}{p \in \alpha \land q \in \Rat \setminus \beta}  & \text{که }\alpha \geq 0_\Real, \beta > 0_\Real
\end{align*}
\textbf{د ژباړې د نسخې سمون (OLARI-004):} سرچينې ناطق عددونه
حقيقي عددونو ته په لاندې‌ليک نښه ورګډ کړي، خو په دې اتحاد کښې د
حقيقي صفر نښه په پورته‌ليک ليکي۔ د فورمول د برابرۍ دپاره اصلي نښه
ساتل شوې؛ دلته ئې ټاکل شوې معنا د حقيقي صفر پرېکون دے۔ د ضرب د
نورو حالتونو د سمبالولو دپاره لومړے وټاکئ: %چې $0_\Real \times \alpha = 0_\Real = \alpha \times 0_\Real$، او ورزيات کړئ:
\begin{align*}
	-\alpha &= \Setabs{p - q}{p < 0 \land q \notin \alpha}
\end{align*}
او بيا وټاکئ:
\begin{align*}
	\alpha \times \beta &\defis
	\begin{cases}
		\mathord{-}\alpha \times \mathord{-}\beta &\text{\RL{که }}\alpha < 0_\Real\text{\RL{ او }}\beta < 0_\Real\\
		\mathord{-}(\mathord{-}\alpha \times \beta) &\text{\RL{که }}\alpha < 0_\Real \text{\RL{ او }}\beta > 0_\Real\\
		\mathord{-}(\alpha \times \mathord{-}\beta) &\text{\RL{که }}\alpha > 0_\Real \text{\RL{ او }}\beta < 0_\Real
	\end{cases}
\end{align*}
\textbf{د ژباړې د نسخې سمون (OLARI-005):} فعاله انګرېزي سرچينه
هغه د ضرب حالتونه نۀ تعريفوي چې يو عامل صفر او بل منفي وي؛ د صفر
قاعده يوازې په غېر فعاله تبصره کښې پاتې ده۔ نو تر هغې قاعدې پرته
پورته ټوټه‌ييز تعريف بشپړ ضرب نۀ تعريفوي۔ بيا بايد وارزوو چې دا هر
تعريف تل يو پرېکون ورکوي۔ په پاے کښې بايد په يوه اسانه (خو اوږده)
څرګندونه کښې وښيو چې په دې ډول تعريف شوي پرېکونونه کټ مټ لکه څنګه
چې بايد، هماغسې چلېږي۔ خو دا ټول \olref[check]{sec} ته پرېږدو۔
\olpseganchor{OLP-0045-B012}{end}

\olpseganchor{OLP-0046-B004}{start}
\olfileid{sfr}{arith}{ref}
\olsection{ځينې فلسفي غورونه}
\olpseganchor{OLP-0046-B004}{end}

\olpseganchor{OLP-0046-B005}{start}
فني خبرې تر دې ځايه بس دي۔ خو هغو څه ترلاسه کړل؟
\olpseganchor{OLP-0046-B005}{end}

\olpseganchor{OLP-0046-B006}{start}
ښه، دا خبره تقريباً له اختلافه پاکه ده چې هغو موږ ته څه {ښکلې}
سوچه رياضي راکړه۔ برسېره پر دې، ځينې ژورې تصوري لاسته راوړنې هم
وې۔ دا پېژندنه ژوره وه چې د بشپړتيا خاصيت د حقيقي او ناطقو عددونو
تر منځ بنسټيز توپير بيانوي۔ همدارنګه د ډېډېکېنډ د پرېکونونو په توګه
د حقيقي عددونو ښکاره جوړښت د تحليل موضوع ته کلک بنسټ ورکوي۔ پوهېږو
چې د \emph{بشپړ مرتب ميدان} تصور بې تناقضه دے، ځکه پرېکونونه کټ
مټ همداسې يو ميدان جوړوي۔
\olpseganchor{OLP-0046-B006}{end}

\olpseganchor{OLP-0046-B007}{start}
له دې ټولو سره، د دغو لاسته راوړنو په اړه بايد څو تحفظات څرګند کړو۔
\olpseganchor{OLP-0046-B007}{end}

\olpseganchor{OLP-0046-B008}{start}
لومړے، دا څرګنده نۀ ده چې د پرېکونونو په توګه د حقيقي عددونو ګڼل
د هغو د پېژندل شوو (ښايي نامتناهي) اعشاري پراختياوو په توګه تر ګڼلو
څومره \emph{زيات} دقيق دي۔ د حقيقي عددونو دا وروستے «جوړښت» د
کوشي تسلسلونو له لارې د حقيقي عددونو جوړښت ته څه ورته والے لري؛ خو
په حقيقت کښې رياضي پوهان د اوولسمې پېړۍ له پيله په بنسټيز ډول پرې
پوهېدل (\olref[cauchy]{sec} وګورئ)۔ په دقت کښې رښتينے زياتوالے له
دې پېژندنې راغے چې حقيقي عددونه د بشپړتيا خاصيت لري؛ يوازې دا وړتيا
چې حقيقي عددونه د ځانګړو سټونو په توګه جوړ کړو، ښايي پخپله دومره
په زړه پورې نۀ وي۔
\olpseganchor{OLP-0046-B008}{end}

\olpseganchor{OLP-0046-B009}{start}
دا لا هم لږ څرګنده ده چې د صحيح يا ناطقو عددونو (ډېرې اسانه) سټي
جوړونه په دغو ډګرونو کښې دقت زياتوي۔ دلته يوه ساده خبره د يادولو وړ
ده۔ صحيح عددونه مو د طبيعي عددونو د مرتبو جوړو د هم ارزښتۍ ټولګيو
په توګه \emph{جوړ} کړل، بيا مو ناطق عددونه د صحيح عددونو د مرتبو
جوړو د هم ارزښتۍ ټولګيو په توګه جوړ کړل، او بيا مو حقيقي عددونه د
ناطقو عددونو د سټونو په توګه جوړ کړل؛ خو سمدستي دا جوړښتونه
\emph{هېروو}۔ په ځانګړي ډول هېڅوک به کله هم ونۀ غواړي چې په يوه
رياضياتي ثبوت کښې دا جوړښتونه \emph{وکاروي} (البته د دې ثبوت پرته
چې جوړښتونه لکه څنګه چې ټاکل شوي وو، هماغسې چلېږي)۔ د يوه حقيقي
عدد په اړه نېغ په نېغه خبرې کول د طبيعي عددونو د سټونو د سټونو د
سټونو د سټونو د سټونو د سټونو د کوم سټ په اړه تر خبرو ډېر اسانه دي۔
%(ځکه حقيقي عددونه د ناطقو عددونو د سټونو په توګه جوړ شول؛ او دا د صحيح عددونو د مرتبو جوړو د هم ارزښتۍ ټولګيو، يعنې د صحيح عددونو د سټونو د سټونو د سټونو په توګه جوړ شول؛ او داسې نور۔)
%2/3 = [2,3]
%	يعنې {<2,3>, ... }
%	{ {{2},{2,3}}, ... }
%
%حقيقي عددونه
%	د ناطقو عددونو سټونه دي
%
%ناطق عددونه
%	د صحيح عددونو د سټونو د سټونو سټونه دي
%
%صحيح عددونه
%	د طبيعي عددونو د سټونو د سټونو سټونه دي
\olpseganchor{OLP-0046-B009}{end}
	
\olpseganchor{OLP-0046-B010}{start}
تر ټولو زيات شک په دې دے چې دا تعريفونه موږ ته وايي چې صحيح، ناطق
يا حقيقي عددونه په \emph{ماوراءالطبيعتي معنا څه دي}۔ يعنې په دې شک
دے چې حقيقي عددونه (مثلاً) ځينې ځانګړي سټونه (د سټونو د سټونو
\ldots) \emph{دي}۔ د دې نظر پر وړاندې ستر خنډ دا دے چې جوړښت په
ډېرو بېلابېلو لارو کېدے شو۔ د حقيقي عددونو په حالت کښې ځينې په
رښتيا په زړه پورې بېلابېل جوړښتونه شته (\olref[cauchy]{sec}
وګورئ)۔ خو د بېلابېلو جوړښتونو د ترلاسه کولو يوه ډېره ساده لار دا
ده: لکه په \olref[sfr][rel][ref]{sec} کښې، مرتبې جوړې مو لږ په
بله توګه تعريفولے شوې؛ که مو د مرتبې جوړې دا بديل تصور کارولے وے،
زموږ جوړښتونو به کټ مټ همدومره سم کار کړے وے، خو په پاے کښې به مو
بېل څيزونه ترلاسه کړي وو۔ له دې امله د صحيح، ناطقو او حقيقي عددونو
ډېر سيال سټي جوړښتونه شته۔ اوس دا ادعا چې صحيح عددونه (مثلاً)
\emph{دا} سټونه دي، نۀ \emph{هغه}، يوازې خپلسرې (او شرموونکې) به
وي۔ (لکه په \olref[sfr][rel][ref]{sec} کښې، دا د هغه استدلال يوه
بېلګه ده چې \citealt{Benacerraf1965} مشهور کړ۔)
\olpseganchor{OLP-0046-B010}{end}

\olpseganchor{OLP-0046-B011}{start}
يو بل ټکی هم د راپورته کولو وړ دے: زموږ په جوړښتونو کښې يو ډېر
\emph{عجيب} څيز شته۔ موږ په طبيعي عددونو پيل وکړ۔ بيا صحيح عددونه
جوړوو، او «د صحيح عددونو \LR{$0$}»، يعنې
\LR{$ \equivrep{0,0}{\Intequiv}$}، جوړوو۔ خو
\LR{$0 \neq \equivrep{0,0}{\Intequiv}$}۔ په حقيقت کښې زموږ د جوړښتونو
له مخې \emph{هېڅ} طبيعي عدد صحيح عدد نۀ دے۔ خو دا د شهودي پوهې
سخت خلاف ښکاري۔ همدارنګه په \olref[sfr][set][imp]{sec} کښې مو له
ډېر استدلال پرته ادعا وکړه چې \LR{$\Nat \subseteq \Rat$}۔ که جوړښتونه
موږ ته په دقيق ډول وايي چې عددونه \emph{څه} دي، نو دا ادعا په
څرګند ډول ناسمه وه۔
\olpseganchor{OLP-0046-B011}{end}

\olpseganchor{OLP-0046-B012}{start}
نو چې يو ګام شاته ودرېږو، کوم ځاے ته رسېږو؟ په ساده سټ نظريه کښې
په کار کولو او طبيعي عددونو ته په ورکړل شوو څيزونو کتلو، صحيح، ناطق
او حقيقي عددونه د ځينو سټونو په توګه \emph{کارولے} شو۔ په دې معنا
د دغو څيزونو نظريې په سټ نظريه کښې \emph{ورګډولے} شو۔ خو د دې
ورګډولو فلسفي ارزښت دومره ساده نۀ دے۔
\olpseganchor{OLP-0046-B012}{end}

\olpseganchor{OLP-0046-B013}{start}
البته دا هېڅ د بحث وروستۍ خبره نۀ ده!{} ټکی يوازې دا دے: د دې د
ښودلو دپاره چې د حقيقي عددونو سټي جوړونه ژور فلسفي ارزښت \emph{لري}،
يو څه نور \emph{فلسفي} استدلال ته اړتيا ده۔
\olpseganchor{OLP-0046-B013}{end}

\olpseganchor{OLP-0046-B014}{start}
%برسېره پر دې، د دې ټول څپرکي په اوږدو کښې مو طبيعي عددونه يوازې
%د راکړل شوو څيزونو په توګه ګڼلي دي۔ خو له هغو سره څه کولے شو؟
%دا د بل څپرکي موضوع ده۔
\olpseganchor{OLP-0046-B014}{end}

\olpseganchor{OLP-0047-B003}{start}
\olfileid{sfr}{arith}{check}
\olsection{مرتبې حلقې او ميدانونه}
\olpseganchor{OLP-0047-B003}{end}

\olpseganchor{OLP-0047-B004}{start}
په دې ټول څپرکي کښې مو ادعا وکړه چې ځينې تعريفونه «لکه څنګه چې
بايد» هماغسې چلېږي۔ په دې فني ضميمه کښې به روښانه کړو چې مطلب مو
څه دے، او (لنډه خاکه به ورکړو چې څنګه) وښيو تعريفونه «سم» چلېږي۔
\textbf{د ژباړې د نسخې بشپړتيا يادونه (OLARI-006):} سرچينه د حلقې
قوانين ارزوي، خو تر هغو مخکښې دا نۀ ثابتوي چې د صحيح او ناطقو
عددونو پر ټولګيو تعريف شوې عمليې د استازي له انتخابه خپلواکې دي؛
لاندې محاسبې دا اړين ښه‌تعريف‌والے له مخکښې فرضوي۔
\olpseganchor{OLP-0047-B004}{end}

\olpseganchor{OLP-0047-B005}{start}
په \olref[int]{sec} کښې مو پر \LR{$\Int$} جمع او ضرب تعريف کړل۔ غواړو
وښيو چې دا عمليې، لکه څنګه چې تعريف شوې دي، \LR{$\Int$} ته هغه جوړښت
ورکوي چې موږ ئې «غواړو»۔ په ځانګړي ډول دا جوړښت يوه تبادلي حلقه ده۔
\olpseganchor{OLP-0047-B005}{end}

\olpseganchor{OLP-0047-B006}{start}
\begin{defn}
	يوه \emph{تبادلي حلقه} يو سټ \LR{$S$} دے چې ځانګړي غړي \LR{$0$} او \LR{$1$} او عمليې \LR{$+$} او \LR{$\times$} لري، او دا اته فارمولونه پوره کوي:
	\begin{align*}
		\emph{اتحاديت}&&a + (b+ c) & = (a + b) + c \\
		&& (a \times b) \times c & = a \times (b\times c)\\
		\emph{تبادليت}&&a + b &= b+ a  \\
		&&  a \times b&= b\times a\\
		\emph{عيني غړي}&&a + 0 &= a \\
		&& a \times 1 &= a\\
		\emph{جمعي معکوس}&&(\exists b\in S)0&=a + b\\
		\emph{توزيعيت}&&a \times (b+ c ) &= (a \times b) + (a \times c)
	\end{align*}
	په ضمني ډول دا ټول په هغو کلي مقداري نښو تړل شوي چې ساحه ئې تر \LR{$S$} محدوده ده۔ او پام وکړئ چې دلته د \LR{$0$} او~\LR{$1$} غړي لازماً په همدې نومونو طبيعي عددونه نۀ دي۔
\end{defn}
\olpseganchor{OLP-0047-B006}{end}

\olpseganchor{OLP-0047-B007}{start}
نو د دې د ارزولو دپاره چې صحيح عددونه تبادلي حلقه جوړوي، يوازې
دا اته شرطونه ارزول غواړو۔ د شرطونو ثابتول {ګران} نۀ دي، خو لږ
اوږد کار دے۔ د بېلګې په ډول د جمع په حالت کښې \emph{اتحاديت}
داسې ثابتوو:
\olpseganchor{OLP-0047-B007}{end}

\olpseganchor{OLP-0047-B008}{start}
\begin{proof}
\LR{$i, j, k \in \Int$} ثابت ونيسئ۔ نو داسې \LR{$a_1, b_1, a_2, b_2, a_3, b_3 \in
\Nat$} شته چې \LR{$i = \equivrep{a_1, b_1}{}$} او \LR{$j =
\equivrep{a_2,b_2}{}$} او \LR{$k = \equivrep{a_3, b_3}{}$}۔ (د لوستلو
د اسانتيا دپاره د ``\LR{$\equivrep{x, y}{\Intequiv}$}'' پر ځاے
``\LR{$\equivrep{x, y}{}$}'' ليکو؛ په دې ټوله برخه کښې به همداسې کوو۔)
اوس:
\begin{align*}
	i + (j + k) &= \equivrep{a_1, b_1}{}+(\equivrep{a_2, b_2}{} + \equivrep{a_3, b_3}{}) \\
	&= \equivrep{a_1,  b_1}{} + \equivrep{a_2+a_3, b_2+b_3}{}\\
	&= \equivrep{a_1 + (a_2 + a_3), b_1 + (b_2 + b_3)}{}\\
	&= \equivrep{(a_1 + a_2) + a_3, (b_1 + b_2) + b_3}{}\\
	&= \equivrep{a_1 + a_2, b_1 + b_2}{} + \equivrep{a_3, b_3}{}\\
	&= (\equivrep{a_1, b_1}{} + \equivrep{a_2, b_2}{}) + \equivrep{a_3, b_3}{}\\
	&= (i+j) + k
\end{align*}
او د \LR{$\Nat$} پر جمع له چلنده ازاد کار اخلو۔
\end{proof}
\olpseganchor{OLP-0047-B008}{end}

\olpseganchor{OLP-0047-B009}{start}
په همدې ډول \emph{جمعي معکوس} داسې ثابتوو:
\olpseganchor{OLP-0047-B009}{end}

\olpseganchor{OLP-0047-B010}{start}
\begin{proof}
\LR{$i \in \Int$} ثابت ونيسئ، نو د ځينو \LR{$a,b \in \Nat$} دپاره
\LR{$i = \equivrep{a,b}{}$}۔ وټاکئ \LR{$j = \equivrep{b,a}{} \in \Int$}۔
د طبيعي عددونو له چلنده په کار اخيستلو
\LR{$(a+b) + 0 = 0 + (a+b)$}، نو د تعريف له مخې
\LR{$\tuple{a+b, b+a} \sim_\Int \tuple{0,0}$}، او له دې امله
\LR{$\equivrep{a+b, b+a}{} = \equivrep{0, 0}{} = 0_\Int$}۔ نو اوس
\LR{$i + j = \equivrep{a,b}{}+\equivrep{b,a}{}=\equivrep{a+b, b+a}{}= \equivrep{0,
0}{} = 0_\Int$}۔
\end{proof}
\olpseganchor{OLP-0047-B010}{end}

\olpseganchor{OLP-0047-B011}{start}
او د \emph{توزيعيت} ثبوت دا دے:
\olpseganchor{OLP-0047-B011}{end}

\olpseganchor{OLP-0047-B012}{start}
\begin{proof}
لکه پورته، \LR{$i = \equivrep{a_1, b_1}{}$}،
\LR{$j = \equivrep{a_2,b_2}{}$} او \LR{$k = \equivrep{a_3, b_3}{}$} ثابت
ونيسئ۔ اوس:
\begin{align*}
	i \times (j + k)
	&= \equivrep{a_1, b_1}{} \times (\equivrep{a_2,b_2}{} + \equivrep{a_3, b_3}{})\\
	&= \equivrep{a_1, b_1}{} \times \equivrep{a_2 + a_3,b_2+b_3}{}\\
	&= \equivrep{a_1  (a_2 + a_3) + b_1  (b_2+b_3), a_1  (b_2 + b_3) + b_1 (a_2 + a_3)}{}\\
	&= \equivrep{a_1 a_2 + a_1a_3 + b_1 b_2+b_1b_3, a_1 b_2 + a_1b_3 + a_2b_1 + a_3b_1}{}\\		
%		&= \equivrep{a_1 a_2 + b_1b_2 + a_1a_3 + b_1 b_2+b_1b_3, a_1 b_2 + a_1b_3 + a_2b_1 + a_3b_1}{}\\		
	&= \equivrep{a_1a_2 + b_1b_2, a_1b_2 + a_2b_1}{} + \equivrep{a_1a_3 + b_1b_3, a_1b_3 + a_3b_1}{}\\
	&= (\equivrep{a_1, b_1}{} \times \equivrep{a_2,b_2}{}) + (\equivrep{a_1, b_1}{} \times  \equivrep{a_3, b_3}{})\\
	&= (i \times j) + (i \times k)
\end{align*}
\end{proof}
\olpseganchor{OLP-0047-B012}{end}

\olpseganchor{OLP-0047-B013}{start}
د پاتې پنځو شرطونو ثبوت د تمرين په توګه پرېږدو۔ چې دا مو وکړل،
ښودلي مو دي چې \LR{$\Int$} يوه تبادلي حلقه جوړوي؛ يعنې جمع او ضرب
(لکه څنګه چې تعريف شوي) هماغسې چلېږي چې بايد۔
\olpseganchor{OLP-0047-B013}{end}

\olpseganchor{OLP-0047-B014}{start}
\begin{prob}
ثابت کړئ چې \LR{$\Int$} يوه تبادلي حلقه ده۔
\end{prob}
\olpseganchor{OLP-0047-B014}{end}

\olpseganchor{OLP-0047-B015}{start}
خو زموږ کار لا ختم نۀ دے۔ پر \LR{$\Int$} مو له جمع او ضرب سره يوه
ترتيبي اړيکه \LR{$\leq$} هم تعريف کړې، او بايد وارزوو چې دا لکه څنګه
چې بايد، هماغسې چلېږي۔ په لا تفصيل بايد وښيو چې \LR{$\Int$} يوه
\emph{مرتبه} حلقه جوړوي۔\footnote{له
	\olref[sfr][rel][ord]{def:linearorder} څخه په ياد ولرئ چې کلي
	ترتيب هغه اړيکه ده چې انعکاسي، انتقالي، ضدتناظري او وصل شوې وي۔
	د ترتيبي اړيکو په متن کښې وصل والے کله کله \emph{درېګونې} بلل
	کېږي، ځکه د هر \LR{$a$} او \LR{$b$} دپاره \LR{$a \leq b \lor a
	= b \lor a \geq b$} لرو۔}
\olpseganchor{OLP-0047-B015}{end}

\olpseganchor{OLP-0047-B016}{start}
\begin{defn}
يوه \emph{مرتبه حلقه} هغه تبادلي حلقه ده چې يوه کلي ترتيبي اړيکه
\LR{$\leq$} هم لري، داسې چې:
\begin{align*}
	a \leq b &\lif a + c \leq b + c\\
	(a \leq b \land 0 \leq c) &\lif a \times c \leq b \times c
\end{align*}
\end{defn}
\olpseganchor{OLP-0047-B016}{end}

\olpseganchor{OLP-0047-B017}{start}
\begin{prob}
ثابت کړئ چې \LR{$\Int$} يوه مرتبه حلقه ده۔
\end{prob}
\olpseganchor{OLP-0047-B017}{end}

\olpseganchor{OLP-0047-B018}{start}
لکه مخکښې، دا ښودل چې جوړ شوے~\LR{$\Int$} يوه مرتبه حلقه ده، اوږد
خو معمولي کار دے۔ دا تاسو ته پرېږدو۔
\olpseganchor{OLP-0047-B018}{end}

\olpseganchor{OLP-0047-B019}{start}
په دې سره د صحيح عددونو کار بشپړېږي۔ خو اوس د ناطقو عددونو په اړه
ورته خبرې ښودل غواړو۔ په ځانګړي ډول بايد وښيو چې زموږ د \LR{$+$}،
\LR{$\times$} او \LR{$\leq$} د ورکړل شوو تعريفونو له مخې ناطق عددونه يو
مرتب \emph{ميدان} جوړوي:
\begin{defn}\ollabel{orderedfield}
يو \emph{مرتب ميدان} هغه مرتبه حلقه ده چې دا شرط هم پوره کوي:
\begin{align*}
	\emph{ضربي معکوس}& & (\forall a \in S \setminus \{0\})(\exists b \in S) a\times b& = 1
\end{align*}
\textbf{د ژباړې د نسخې بشپړتيا يادونه (OLARI-008):} د ميدان په
دوديز تعريف کښې جمعي او ضربي عيني غړي بايد بېل وي؛ سرچينې دا شرط
نۀ دے بيان کړے، نو ليکل شوے تعريف صفر حلقه هم په تشه توګه مني۔
\end{defn}
\olpseganchor{OLP-0047-B019}{end}

\olpseganchor{OLP-0047-B020}{start}
چې مو وښودل \LR{$\Int$} يوه مرتبه حلقه جوړوي، بيا دا ښودل اسانه خو
اوږد دي چې \LR{$\Rat$} يو مرتب ميدان جوړوي۔
\olpseganchor{OLP-0047-B020}{end}

\olpseganchor{OLP-0047-B021}{start}
\begin{prob}
ثابت کړئ چې \LR{$\Rat$} يو مرتب ميدان دے۔
\end{prob}
\olpseganchor{OLP-0047-B021}{end}

\olpseganchor{OLP-0047-B022}{start}
د صحيح او ناطقو عددونو له سمبالولو وروسته يوازې حقيقي عددونه پاتې
دي۔ په ځانګړي ډول بايد وښيو چې \LR{$\Real$} يو \emph{بشپړ} مرتب ميدان
جوړوي، يعنې يو مرتب ميدان چې د بشپړتيا خاصيت لري۔
\olref[cuts]{realcompleteness} ثابته کړه چې \LR{$\Real$} د بشپړتيا خاصيت
لري۔ خو لا د دې ارزولو (ستړے) کار پاتې دے چې \LR{$\Real$} يو مرتب ميدان
دے۔ \textbf{د ژباړې د نسخې سمون (OLARI-007):} په انګرېزي جمله
کښې د «ستړي» او «ارزولو» تر منځ اړين نوم پاتې شوے و؛ دلته جمله د
همدې کار په څرګند نوم بشپړه شوې ده۔
\olpseganchor{OLP-0047-B022}{end}

\olpseganchor{OLP-0047-B023}{start}
مخکښې له دې چې \emph{هغه} اوږد تمرين پيل کړو، بايد ځينې نورې
«سمدستي» خبرې وارزوو۔ د بېلګې په ډول بايد ډاډ ولرو چې د هر دوو
پرېکونونو \LR{$\alpha$} او \LR{$\beta$} دپاره، تعريف شوے \LR{$\alpha + \beta$}
په رښتيا يو \emph{پرېکون} دے۔ د دې حقيقت ثبوت دا دے:
\olpseganchor{OLP-0047-B023}{end}

\olpseganchor{OLP-0047-B024}{start}
\begin{proof}	
ځکه چې \LR{$\alpha$} او \LR{$\beta$} دواړه پرېکونونه دي،
\LR{$\alpha + \beta = \Setabs{p + q}{p \in \alpha \land q \in \beta}$}
د \LR{$\Rat$} يو ناتش خاص فرعي سټ دے۔ اوس فرض کړئ د ځينو
\LR{$p \in \alpha$} او \LR{$q \in \beta$} دپاره \LR{$x < p + q$}۔ بيا
\LR{$x - p < q$}، نو \LR{$x - p \in \beta$}، او
\LR{$x = p + (x - p) \in \alpha + \beta$}۔ نو \LR{$\alpha + \beta$} د
\LR{$\Rat$} يوه پيلنۍ برخه ده۔ په پاے کښې، د هر
\LR{$p + q \in \alpha + \beta$} دپاره، ځکه چې \LR{$\alpha$} او \LR{$\beta$}
دواړه پرېکونونه دي، داسې \LR{$p_1 \in \alpha$} او \LR{$q_1 \in \beta$} شته
چې \LR{$p < p_1$} او \LR{$q < q_1$}؛ نو
\LR{$p + q < p_1 + q_1 \in \alpha + \beta$}؛ نو
\LR{$\alpha + \beta$} تر ټولو لوے غړے نۀ لري۔
\end{proof}
\olpseganchor{OLP-0047-B024}{end}

\olpseganchor{OLP-0047-B025}{start}
ورته هڅې به درته اجازه وکړي چې وارزوئ \LR{$\alpha - \beta$}،
\LR{$\alpha \times \beta$} او \LR{$\alpha \div \beta$} پرېکونونه دي (په
وروستي حالت کښې هغه حالت نۀ شمېرو چې \LR{$\beta$} صفر پرېکون وي)۔ خو
يو ځل بيا دا يوازې تاسو ته پرېږدو۔ \textbf{د ژباړې د نسخې بشپړتيا
يادونه (OLARI-009):} د پرېکونونو په برخه کښې د تفريق او وېش
فارمولونه غېر فعالې تبصرې دي، نو فعاله سرچينه هغه دوه عمليې نۀ
تعريفوي چې دا دعوه او راتلونکے ميدان تمرين ورته اړتيا لري۔
\olpseganchor{OLP-0047-B025}{end}

\olpseganchor{OLP-0047-B026}{start}
\begin{prob}
ثابت کړئ چې \LR{$\Real$} يو مرتب ميدان دے۔
\end{prob}
\olpseganchor{OLP-0047-B026}{end}

\olpseganchor{OLP-0047-B027}{start}
خو يوه وړه نيمګړې خبره لا سمول غواړي۔ په \olref[cuts]{sec} کښې
مو وړانديز وکړ چې \LR{$\sqrt{2} =
\Setabs{p \in \Rat}{p < 0 \text{\RL{ يا }}p^2 < 2}$} اخيستلے شو۔ خو
بايد وښيو چې دا سټ يو \emph{پرېکون} دے۔ د دې حقيقت ثبوت دا دے:
\olpseganchor{OLP-0047-B027}{end}

\olpseganchor{OLP-0047-B028}{start}
\begin{proof}
په څرګند ډول دا د ناطقو عددونو يوه ناتشه خاصه پيلنۍ برخه ده؛ نو
دومره ښودل بس دي چې تر ټولو لوے غړے نۀ لري۔ په ځانګړي ډول دا ښودل
بس دي چې که \LR{$p$} يو مثبت ناطق عدد وي، \LR{$p^2 < 2$} وي او
\LR{$q = \frac{2p+2}{p+2}$} وي، نو \LR{$p < q$} او \LR{$q^2 < 2$} دواړه دي۔
د \LR{$p < q$} د ليدلو دپاره يوازې پام وکړئ:
\begin{align*}
	p^2 &< 2\\
	p^2 + 2p &< 2 + 2p\\
	p(p + 2) &< 2 + 2p\\
	p &< \tfrac{2+2p}{p+2} = q
\end{align*}
د \LR{$q^2 < 2$} د ليدلو دپاره يوازې پام وکړئ:
\begin{align*}
	p^2 &< 2\\
	2p^2 + 4p + 2 &< p^2 + 4p+ 4\\
	4p^2 + 8p + 4 &< 2(p^2 + 4p + 4)\\
	(2p+2)^2 & <2(p+2)^2\\
	\tfrac{(2p+2)^2}{(p+2)^2} &< 2\\
	q^2 &< 2
\end{align*}
\end{proof}
\olpseganchor{OLP-0047-B028}{end}

\olpseganchor{OLP-0048-B004}{start}
\olfileid{sfr}{arith}{cauchy}
\olpseganchor{OLP-0048-B004}{end}

\olpseganchor{OLP-0048-B005}{start}
\olsection{ضميمه: حقيقي عددونه د کوشي د لړيو په توګه}
\olpseganchor{OLP-0048-B005}{end}

\olpseganchor{OLP-0048-B006}{start}
په \olref[cuts]{sec} کښې مو حقيقي عددونه د ډېډېکېنډ د پرېکونونو په
توګه جوړ کړل۔ په دې برخه کښې يوه بله جوړونه تشريح کوو۔ دا د کوشي پر
هغه تعريف ولاړه ده چې اوس ئې د کوشي لړۍ بولو؛ خو له دې تعريف څخه د
حقيقي عددونو د \emph{جوړولو} کار د نولسمې پېړۍ نورو ليکوالانو، په
ځانګړې توګه وايرشټراس، هاينې، مېراے او کانټور، کړے دے۔ (د ښۀ تاريخ
دپاره \citeauthor{OConnorRobertson:RN} \citeyear{OConnorRobertson:RN}
وګورئ۔)
\olpseganchor{OLP-0048-B006}{end}

\olpseganchor{OLP-0048-B007}{start}
مخکښې له دې چې نولسمې پېړۍ ته ورسېږو، د سايمن سټېوين (1548--1620)
خبره د پام وړ ده۔ په لنډ ډول، سټېوين پوه شو چې هر حقيقي عدد د هغه د
اعشاري پراختيا له مخې ګڼلے شو۔ نو آن يو غير ناطق عدد لکه \LR{$\sqrt{2}$}
هم يوه ښۀ اعشاري پراختيا لري، چې پيل ئې داسې دے:
\[
	1.41421356237\ldots
\]
په سټ نظريه کښې د اعشاري پراختياوو نمونه جوړول ډېر اسان دي: يوازې
داسې تابعې \LR{$d \colon \Nat \to \Nat$} ئې وګڼئ چې \LR{$d(n)$} هغه \LR{$n$}م
اعشاري ځاے وي چې موږ ئې څېړو۔ بيا لږ بدلون ته اړتيا لرو، څو د حقيقي
عدد هغه برخه هم سمبال کړو چې له اعشاري ټکي مخکښې راځي (دلته يوازې
\LR{$1$})۔ يو بل بدلون، يعنې د هم ارزښتۍ اړيکه، هم پکار ده څو د بېلګې په
توګه يقيني کړو چې \LR{$0.999\ldots = 1$}۔ خو په سټ نظريه کښې د سټېوين په
لار د حقيقي عددونو بشپړه دقيقه جوړونه وړاندې کول ګران نۀ دي۔
\olpseganchor{OLP-0048-B007}{end}

\olpseganchor{OLP-0048-B008}{start}
سټېوين زموږ اصلي موضوع نۀ ده۔ (د سټېوين په اړه نور معلومات په
\citealt{KatzKatz2012} کښې وګورئ۔) خو دا ډېره ورته خبره وګورئ۔ د
\LR{$\sqrt{2}$} اعشاري پراختيا په نېغه د څېړلو پر ځاے، د \LR{$\sqrt{2}$} د
لا کره پراختياوو په اخيستلو د هغو ناطقو اټکلونو يوه \emph{لړۍ}
کتلے شو چې کره‌والے ئې پرله‌پسې زياتېږي:
\[
1, 1.4, 1.414, 1.4142, 1.41421,\ldots
\]
دا نظر چې حقيقي عددونه د «پرله‌پسې ښو اټکلونو» له لارې ګڼلے شو، د
يوې بلې پوهېدنې بنسټ راکوي (هغو پوهېدنو ته ورته چې \LR{$\Rat$} مو له
\LR{$\Int$}، يا \LR{$\Int$} مو له \LR{$\Nat$} څخه جوړول)۔ بنسټيزې خبرې دا دي:
\begin{enumerate}
	\item هر حقيقي عدد د يوې (کېدے شي نامتناهي) اعشاري پراختيا په
	بڼه ليکل کېدے شي۔
	\item په يوه (کېدے شي نامتناهي) اعشاري پراختيا کښې زېرمه شوي
	معلومات په هماغه ډول د ناطقو عددونو په يوه لړۍ کښې زېرمه کېدے شي۔
	\item د ناطقو عددونو لړۍ له \LR{$\Nat$} څخه \LR{$\Rat$} ته يوه تابع ګڼلے
	شو؛ يوازې \LR{$f(n)$} په لړۍ کښې \LR{$n$}م ناطق عدد وټاکئ۔
\end{enumerate}
البته له \LR{$\Nat$} څخه \LR{$\Rat$} ته هره \emph{تابع} حقيقي عدد نۀ راکوي۔
د بېلګې په توګه دا تابع وګورئ:
\[
	f(n) =\begin{cases}
			1 & \text{\RL{که }}n\text{\RL{ طاق وي}}\\
			0 &\text{\RL{که }}n\text{\RL{ جفت وي}}
		\end{cases}
\]
اصلي اندېښنه دا ده چې د \LR{$0,1,0,1,0,1,0,\ldots$} لړۍ هېڅ حقيقي
عدد ته «نژدې کېدونکې» نۀ ښکاري۔ نو د دې د يقيني کولو دپاره چې هغه
لړۍ څېړو چې کوم حقيقي عدد ته ورنژدې کېږي، بايد خپل پام هغو لړيو ته
محدود کړو چې کوم \emph{حد} لري۔
\olpseganchor{OLP-0048-B008}{end}

\olpseganchor{OLP-0048-B009}{start}
موږ د حد له نظر سره په \emph{د OpenLogic اړونده سرچينيزه برخه}\nobreakspace\LR{\href{https://github.com/OpenLogicProject/OpenLogic/blob/9620cc73f9c8e0ad003c514a5d3748f29611c4c0/content/history/set-theory/limits.tex\#L11}{\latinfont\scriptsize [OLP]}} کښې مخکښې هم
مخامخ شوي يو۔ خو کټ مټ هماغه پخوانے تعريف نشو کارولے۔ هلته عبارت
«\LR{$(\forall \epsilon>0)$}» په پټه پر حقيقي عددونو کميت‌ايښودنه
درلوده؛ او موږ د حقيقي عددونو پر تابعو حدونه څېړل؛ نو د هغه تعريف
کارول به دا وي چې حقيقي عددونه لا له مخکښې ومنو، حال دا چې موږ کټ
مټ د هغو \emph{جوړول} غواړو۔ نېکمرغي ده چې د حد له يوه ډېر نژدې
نظر سره کار کولے شو۔
\begin{defn}\ollabel{def:CauchySequence}
  يوه تابع \LR{$f: \Nat \to \Rat$} هله او يوازې هله د \emph{کوشي لړۍ}
  ده چې د هر مثبت \LR{$\epsilon \in \Rat$} دپاره
  \LR{$(\exists \ell \in \Nat)(\forall m, n > \ell)|f(m) - f(n)| <
  \epsilon$} وي۔
\end{defn}
\olpseganchor{OLP-0048-B009}{end}

\olpseganchor{OLP-0048-B010}{start}
د حد عمومي نظر د پخوا په شان دے: که د کره‌والي يوه ټاکلې کچه
(چې په~\LR{$\epsilon$} اندازه کېږي) وغواړئ، د کتلو يوه «سيمه» شته
(هره ورودي چې له~\LR{$\ell$} څخه لويه وي)۔ دا هم په اسانۍ ليدل کېږي چې
زموږ لړۍ \LR{$1$}، \LR{$1.4$}، \LR{$1.414$}، \LR{$1.4142$}، \LR{$1.41421$}\ldots حد لري:
که غواړئ \LR{$\sqrt{2}$} د \LR{$\nicefrac{1}{10^{n}}$} له خطا دننه اټکل
کړئ، نو يوازې له \LR{$n$}م غړي وروسته هر غړے وګورئ۔
\olpseganchor{OLP-0048-B010}{end}

\olpseganchor{OLP-0048-B011}{start}
نو څرګند فکر به دا وي چې حقيقي عدد هماغه هره د کوشي لړۍ \emph{ده}۔
خو لکه د \LR{$\Int$} او \LR{$\Rat$} په جوړونو کښې، دا خبره به بېخي ساده وي:
د هر ورکړل شوي حقيقي عدد دپاره څو بېلې د کوشي لړۍ هماغه حقيقي عدد
ښيي۔ د دې د ليدلو يوه اسانه لار داسې ده۔ يوه د کوشي لړۍ~\LR{$f$} واخلئ
او \LR{$g$} کټ مټ د~\LR{$f$} په شان تابع تعريف کړئ، يوازې دومره توپير چې
\LR{$g(0)\neq f(0)$}۔ ځکه دواړه لړۍ له لومړي عدد وروسته په هر ځاے کښې
سره برابرې دي، نو په پاے کښې به ووايو چې د
\olref{def:CauchySequence} په معنا يو حد لري، او بايد د هماغه حقيقي
عدد «تعريفوونکې» وګڼل شي۔ نو په رښتيا بايد دا د کوشي لړۍ هماغه يو
حقيقي عدد وګڼو۔
\olpseganchor{OLP-0048-B011}{end}

\olpseganchor{OLP-0048-B012}{start}
په پايله کښې بيا بايد د کوشي پر لړيو د هم ارزښتۍ يوه اړيکه تعريف
کړو او حقيقي عددونه د هم ارزښتۍ له ټولګيو سره يو وګڼو۔
\textbf{د ژباړې د نسخې سمون (OLARI-010):} انګرېزي سرچينه په دې
ځاے «د هم ارزښتۍ اړيکې» وايي، خو جوړېدونکي حقيقي عددونه د اړيکې
\emph{ټولګي} دي؛ د لاندې تعريف او فورمول سره سم دلته هماغه ټولګي
ليکل شوي۔ لومړے داسې تابع پکار ده چې په حد کښې \LR{$0$} ته ورنژدې شي۔
د هرې تابع \LR{$h : \Nat \to \Rat$} دپاره ووايئ چې \emph{\LR{$h$} \LR{$0$} ته
ورنژدې کېږي} هله او يوازې هله چې د هر مثبت \LR{$\epsilon \in \Rat$}
دپاره \LR{$(\exists \ell \in \Nat)(\forall n > \ell)|h(n)| <
\epsilon$} وي۔\footnote{دا د \LR{$\lim_{x
\mathord{\rightarrow}\infty}f(x) = 0$} له تعريف سره پرتله کړئ چې په
\emph{د OpenLogic اړونده سرچينيزه برخه}\nobreakspace\LR{\href{https://github.com/OpenLogicProject/OpenLogic/blob/9620cc73f9c8e0ad003c514a5d3748f29611c4c0/content/history/set-theory/limits.tex\#L11}{\latinfont\scriptsize [OLP]}} کښې راغلے۔} برسېره پر دې، چې \LR{$f$} او
\LR{$g$} د \LR{$\Nat \to \Rat$} تابعې وي، \LR{$(f-g)(n) = f(n) - g(n)$}
وټاکئ۔ اوس تعريف کړئ:
\[
	f \Realequiv g \text{\RL{ هله او يوازې هله چې \LR{$(f-g)$} \LR{$0$} ته ورنژدې کېږي}}.
\]
بايد وارزوو چې \LR{$\Realequiv$} د هم ارزښتۍ اړيکه ده؛ او رښتيا هم ده۔
بيا، که وغواړو، حقيقي عددونه د \LR{$\Nat \to \Rat$} د ټولو کوشي
لړيو هغه د هم ارزښتۍ ټولګي تعريفولے شو چې د \LR{$\Realequiv$} له مخې
جوړېږي۔
\olpseganchor{OLP-0048-B012}{end}

\olpseganchor{OLP-0048-B013}{start}
\begin{prob}
د هر \LR{$n$} دپاره \LR{$f(n) = 0$} وټاکئ۔ او \LR{$g(n) = \frac{1}{(n+1)^2}$}
وټاکئ۔ وښيئ چې دواړه د کوشي لړۍ دي او په رښتيا د دواړو تابعو حد
\LR{$0$} دے، نو \LR{$f \sim_\Real g$} هم دے۔
\textbf{د ژباړې د نسخې سمون (OLARI-011):} دلته سرچينه د تازه
تعريف شوې نښې پر ځاے بېله نښه کاروي؛ د تمرين معنا د هماغې تازه تعريف شوې اړيکې له
مخې هم ارزښت دے۔
\end{prob}
\olpseganchor{OLP-0048-B013}{end}

\olpseganchor{OLP-0048-B014}{start}
له دې وروسته، د معمول په شان، \LR{$\equivrep{f}{\Realequiv}$} د هغه د
هم ارزښتۍ ټولګي دپاره ليکو چې \LR{$f$} ئې يو غړے دے۔ خو د لوست
د اسانولو دپاره به په راتلونکو برخو کښې لاندې‌ليک پرېږدو او يوازې
\LR{$\equivrep{f}{}$} ليکو۔ دا هم ټاکو چې د هر \LR{$q \in \Rat$} دپاره
\LR{$q_{\Real} = \equivrep{c_{q}}{}$} وي، چې \LR{$c_{q}$} ثابته تابع ده او
د هر \LR{$n \in \Nat$} دپاره \LR{$c_q(n) = q$} وي۔ بيا پر حقيقي عددونو
بنسټيزې اړيکې او عمليې داسې تعريفوو، د بېلګې په توګه:
\begin{align*}
	\equivrep{f}{} + \equivrep{g}{} &= 	\equivrep{(f + g)}{} \\
	\equivrep{f}{} \times 	\equivrep{g}{} &= \equivrep{(f \times g)}{}
\end{align*}
چې \LR{$(f + g)(n) = f(n) + g(n)$} او \LR{$(f \times g)(n) = f(n) \times
g(n)$}۔ البته بايد دا هم وارزوو چې که \LR{$f$} او \LR{$g$} د کوشي لړۍ وي، نو
هره يوه \LR{$(f + g)$}، \LR{$(f-g)$} او \LR{$(f\times g)$} هم د کوشي لړۍ ده؛ او
داسې ده، خو ثبوت ئې تاسو ته پرېږدو۔
\olpseganchor{OLP-0048-B014}{end}

\olpseganchor{OLP-0048-B015}{start}
په پاے کښې د مرتبې يو مفهوم تعريفوو۔ ووايئ \LR{$\equivrep{f}{}$}
\emph{مثبت} دے هله او يوازې هله چې هم \LR{$\equivrep{f}{}\neq 0_\Real$}
وي او هم \LR{$(\exists \ell \in \Nat)(\forall n > \ell)0 < f(n)$} وي۔
\textbf{د ژباړې د نسخې سمون (OLARI-012):} سرچينه په دې شرط کښې
ناطق صفر ليکي، خو کيڼه خوا د حقيقي عدد ټولګے دے او همدلته پورته
ثابته ناطقه تابع حقيقي عدد ته ورګډه شوې؛ نو د ډول له مخې سم صفر
حقيقي صفر دے۔ بيا ووايئ
\LR{$\equivrep{f}{} < \equivrep{g}{}$} هله او يوازې هله چې
\LR{$\equivrep{(g - f)}{}$} مثبت وي۔ بايد وارزوو چې دا سم تعريف شوے
دے (يعنې د هم ارزښتۍ له ټولګي څخه د «نماينده» تابعې پر ټاکنه پورې
اړه نۀ لري)۔
%\begin{align*}
%	\equivrep{f}{} \leq \equivrep{g}{} \text{ iff }&\equivrep{f}{} = \equivrep{g}{} \text{ or }\\
%	&(\exists \ell \in \Nat)(\forall n \in \Nat)(\ell < n \lif f(n) \leq g(n))
%\end{align*}
خو له دې ارزونې وروسته په ډېرې اسانۍ ښودل کېدے شي چې دا تعريفونه
سم الجبري خاصيتونه راکوي؛ يعنې:
\begin{thm}\ollabel{thm:cauchyorderedfield}
د کوشي د لړيو د هم ارزښتۍ ټولګي يو مرتب ميدان جوړوي۔
\end{thm}
\textbf{د ژباړې د نسخې سمون (OLARI-013):} انګرېزي سرچينه دلته او
د بشپړتيا په راتلونکې قضيه کښې په لنډه «د کوشي لړۍ» وايي؛ په لفظي
ډول يوازې خامې لړۍ ميدان نۀ جوړوي۔ دا پايلې د تعريف شوې اړيکې له مخې
د هغو د هم ارزښتۍ ټولګيو په اړه دي، لکه مخکښې تعريف شول۔
\olpseganchor{OLP-0048-B015}{end}

\olpseganchor{OLP-0048-B016}{start}
\begin{proof}
تمرين۔
\end{proof}
\olpseganchor{OLP-0048-B016}{end}

\olpseganchor{OLP-0048-B017}{start}
\begin{prob}
ثابته کړئ چې د کوشي د لړيو د هم ارزښتۍ ټولګي يو مرتب ميدان جوړوي۔
\end{prob}
\olpseganchor{OLP-0048-B017}{end}

\olpseganchor{OLP-0048-B018}{start}
دا ثابتول لږ ګران دي چې په دې ډول جوړ شوي حقيقي عددونه د بشپړتيا
خاصيت لري، نو ثبوت ئې وړاندې کوو۔
\olpseganchor{OLP-0048-B018}{end}

\olpseganchor{OLP-0048-B019}{start}
\begin{thm}
د کوشي د لړيو د هم ارزښتۍ د ټولګيو هر ناتش سټ چې پاسنی حد لري،
تر ټولو لږ پاسنی حد هم لري۔
\end{thm}
\olpseganchor{OLP-0048-B019}{end}

\olpseganchor{OLP-0048-B020}{start}
\begin{proof}[د ثبوت خاکه] فرض کړئ \LR{$S$} د کوشي د لړيو هر داسې ناتش
سټ دے چې پاسنی حد لري۔ نو داسې \LR{$p \in \Rat$} شته چې \LR{$p_{\Real}$} د
\LR{$S$} پاسنی حد دے۔ يو \LR{$r \in S$} واخلئ؛ بيا داسې \LR{$q \in \Rat$} شته
چې \LR{$q_{\Real} < r$} وي۔ نو که د \LR{$S$} تر ټولو لږ پاسنی حد موجود وي،
د \LR{$q_\Real$} او \LR{$p_\Real$} تر منځ (دواړه سرونه پکې شامل) دے۔
\textbf{د ژباړې د نسخې سمون (OLARI-014):} له يوه خپل‌سري حقيقي
پاسني حد څخه د ناطق پاسني حد، او د ټاکلي غړي څخه لاندې د ناطق
حد شتون د ناطقو عددونو د ارشيميدي/کثافت خاصيت ته اړتيا
لري۔ سرچينه دا منځنی دليل نۀ وړاندې کوي؛ لاندې نيمه‌کولو جوړونه
همدا ناطق سرونه فرضوي۔
\olpseganchor{OLP-0048-B020}{end}

\olpseganchor{OLP-0048-B021}{start}
موږ به تر ټولو لږ پاسني حد ته په يوه وخت له لاندې او پورته څخه
ورنژدې شو۔ په ځانګړې توګه دوه تابعې \LR{$f, g \colon \Nat \to \Rat$}
داسې تعريفوو چې موخه مو دا وي چې \LR{$f$} له پورته او \LR{$g$} له لاندې څخه
تر ټولو لږ پاسني حد ته ورنژدې شي۔ په دې تعريف پيل کوو:
\begin{align*}
	f(0) &= p \\
	g(0) &= q
\end{align*}
بيا، چې \LR{$a_n = \frac{f(n) + g(n)}{2}$} وي، وټاکئ:\footnote{دا يو
بازګشتي تعريف دے۔ خو تر اوسه مو هېڅ دليل نۀ دے ورکړے چې
بازګشتي تعريفونه سم دي۔}
\begin{align*}
	f(n+1) &=
	\begin{cases}
		a_n &\text{\RL{که }}(\forall h \in S)\equivrep{h}{} \leq (a_n)_\Real\\
		f(n)&\text{\RL{که نۀ}}
	\end{cases}\\
	g(n+1) &=
	\begin{cases}
		a_n &\text{\RL{که }}(\exists h \in S)\equivrep{h}{} \geq (a_n)_\Real\\
	 	g(n) &\text{\RL{که نۀ}}
	\end{cases}
\end{align*}
\LR{$f$} او \LR{$g$} دواړه د کوشي لړۍ دي۔ (دا په نسبتاً اسانۍ ارزول کېدے
شي، خو د تمرين په توګه ئې پرېږدو۔) پام وکړئ چې تابع \LR{$(f-g)$} \LR{$0$} ته
ورنژدې کېږي، ځکه د \LR{$f$} او \LR{$g$} تر منځ توپير په هر ګام کښې نيم کېږي۔
نو \LR{$\equivrep{f}{} = \equivrep{g}{}$}۔
\olpseganchor{OLP-0048-B021}{end}

\olpseganchor{OLP-0048-B022}{start}
لومړے ښيو چې \LR{$\equivrep{f}{}$} د \LR{$S$} پاسنی حد دے، يعنې
\LR{$(\forall h \in S)\equivrep{h}{} \leq \equivrep{f}{}$}۔
(په لار کښې به \olref{thm:cauchyorderedfield} کاروو۔) يو \LR{$h \in S$}
واخلئ او د نقيض دپاره فرض کړئ چې
\LR{$\equivrep{f}{} < \equivrep{h}{}$}، نو
\LR{$0_\Real < \equivrep{(h-f)}{}$}۔ ځکه \LR{$f$} يوه يکنواخته کموونکې د
کوشي لړۍ ده، داسې \LR{$n \in \Nat$} شته چې
\LR{$\equivrep{(c_{f(n)} - f)}{} < \equivrep{(h-f)}{}$}۔ نو:
\[
	(f(n))_\Real = \equivrep{c_{f(n)}}{} < \equivrep{f}{} + \equivrep{(h-f)}{} = \equivrep{h}{},
\]
خو دا له دې جوړښتي حقيقت سره په ټکر کښې دے چې
\LR{$\equivrep{h}{} \leq (f(n))_\Real$}۔
\olpseganchor{OLP-0048-B022}{end}

\olpseganchor{OLP-0048-B023}{start}
اوس ښيو چې \LR{$\equivrep{f}{} = \equivrep{g}{}$} د \LR{$S$} \emph{تر ټولو
لږ} پاسنی حد دے۔ نو \LR{$j$} دې هره د کوشي لړۍ وي او فرض کړئ
\LR{$\equivrep{j}{} < \equivrep{g}{}$}۔ لکه پورته (د دې حقيقت په کارولو
چې \LR{$g$} \emph{زياتېدونکې} ده) داسې \LR{$n \in \Nat$} شته چې
\LR{$\equivrep{j}{} < (g(n))_\Real$}۔ خو د جوړونې له مخې داسې \LR{$h \in S$}
شته چې \LR{$(g(n))_\Real \leq \equivrep{h}{}$}، نو
\LR{$\equivrep{j}{} < \equivrep{h}{}$} او له دې امله
\LR{$\equivrep{j}{}$} د \LR{$S$} پاسنی حد نۀ دے۔
\end{proof}
\olpseganchor{OLP-0048-B023}{end}

\olpseganchor{OLP-0049-B004}{start}
\olchapter{sfr}{infinite}{نامتناهي سټونه}
\olpseganchor{OLP-0049-B004}{end}

\olpseganchor{OLP-0049-B005}{start}
\begin{editorial}
د نامتناهي سټونو دا څپرکے د ټيم بټن له \emph{اوپن سېټ تيوري}
څخه اخيستل شوے دے۔
\end{editorial}
\olpseganchor{OLP-0049-B005}{end}

\olpseganchor{OLP-0050-B003}{start}
\olfileid{sfr}{infinite}{hilbert}
\olsection{د هيلبرټ هوټل}
\olpseganchor{OLP-0050-B003}{end}

\olpseganchor{OLP-0050-B004}{start}
د طبيعي عددونو سټ په څرګند ډول نامتناهي دے۔ نو که نۀ غواړو طبيعي
عددونه له مخکښې \emph{ومنـو}، زموږ لومړے ګام بايد دا وي چې نامتناهي
سټ په داسې اصطلاحاتو کښې وپېژنو چې خپله د طبيعي عددونو يادولو ته
اړتيا ونۀ لري۔ دا يوه ښۀ لار ده چې هيلبرټ په 1924 کال کښې په يوه
وينا کښې وړاندې کړه۔ هغه غواړي چې داسې تصور وکړو:
\begin{quote}
[\ldots] يو هوټل چې متناهي شمېر کوټې لري۔ په دې هره کوټه کښې دې
کټ مټ يو مېلمه اوسېږي۔ که مېلمانه اوس خپلې کوټې په کوم ډول سره
بدلې کړي، [خو] داسې چې په هره کوټه کښې بيا هم تر يوه زيات کس نۀ
وي، نو هېڅ کوټه به خالي نشي، او د هوټل خاوند په دې توګه د نوي
رارسېدونکي مېلمه دپاره نوے ځاے نشي پيدا کولے [\ldots
\textparagraph \ldots]
\olpseganchor{OLP-0050-B004}{end}

\olpseganchor{OLP-0050-B005}{start}
اوس ټاکو چې هوټل به نامتناهي ډېرې شمېرل شوې کوټې \LR{$1$}، \LR{$2$}، \LR{$3$}،
\LR{$4$}، \LR{$5$}، \dots ولري او په هره کښې کټ مټ يو مېلمه وي۔ کله چې يو
نوے مېلمه راشي، خاوند ته يوازې دا پکار دي چې هر پخوانے مېلمه هغې
کوټې ته ولېږي چې شمېر ئې يو زيات وي؛ بيا کوټه~\LR{$1$} د نوي رارسېدونکي
مېلمه دپاره خالي کېږي۔
	\begin{center}
		\begin{tikzpicture}[scale = .75]
		\foreach \x in {1, 2, 3, 4, 5, 6, 7, 8, 9}
		{
			\node (\x a) at (\x, 1) {\small{\x}};
			\node (\x b) at (\x, 2) {\small{\x}};
		}
		\node (dotsa) at (10, 1) {\small{\ldots}};
		\node (dotsb) at (10, 2) {\small{\ldots}};
		\draw[->] (1b)--(2a);
		\draw[->] (2b)--(3a);
		\draw[->] (3b)--(4a);
		\draw[->] (4b)--(5a);
		\draw[->] (5b)--(6a);
		\draw[->] (6b)--(7a);
		\draw[->] (7b)--(8a);
		\draw[->] (8b)--(9a);
		\draw[->] (9b)--(dotsa);
		\draw (1,1) circle (.4);
		\end{tikzpicture}
	\end{center}
	(په \citealt[730]{EwaldSieg2013} کښې خپور شوے؛ ژباړه زموږ ده)
\end{quote}
مهم ټکی دا دے چې د هيلبرټ هوټل نامتناهي ډېرې کوټې لري؛ او د هغه
تشريح د دې خبرې تعريف ګڼلے شو چې دا څه معنا لري۔ په رښتيا، د
ډېډېکېنډ لار همدا وه (البته دلته په لوے زماني بې‌ځايه‌والي سره
وړاندې شوې؛ د ډېډېکېنډ تعريف د \citeyear{Dedekind1888} کال دے):
\olpseganchor{OLP-0050-B005}{end}

\olpseganchor{OLP-0050-B006}{start}
\begin{defn}\ollabel{defn:DedekindInfinite}
يو سټ \LR{$A$} هله او يوازې هله \emph{د ډېډېکېنډ له مخې نامتناهي} دے
چې له~\LR{$A$} څخه د~\LR{$A$} يو حقيقي فرعي سټ ته کومه يو پر يو تابع موجوده
وي۔ يعنې داسې \LR{$o \in A$} او يو پر يو تابع \LR{$f \colon A \to A$} شته
چې \LR{$o \notin \ran{f}$} وي۔
\end{defn}
\olpseganchor{OLP-0050-B006}{end}

\olpseganchor{OLP-0051-B003}{start}
\olfileid{sfr}{infinite}{dedekind}
\olsection{د ډېډېکېنډ الجبرونه}
\olpseganchor{OLP-0051-B003}{end}

\olpseganchor{OLP-0051-B004}{start}
موږ يوازې دا نۀ غواړو چې طبيعي عددونه نامتناهي وي؛ غواړو چې ځينې
(الجبري) خاصيتونه هم ولري: بايد د جمع، ضرب او نورو عملياتو له مخې
سم چلند وکړي۔
\olpseganchor{OLP-0051-B004}{end}

\olpseganchor{OLP-0051-B005}{start}
د ډېډېکېنډ نظر دا و چې د \emph{تالي تابعې} مفهوم بنسټيز وګڼي، او
بيا عددونه د هغو څيزونو په توګه وپېژني چې لاندې خاصيتونه لري:
\begin{enumerate}
	\item يو عدد، \LR{$0$}، شته چې د هېڅ عدد تالي نۀ دے
	\\يعنې \LR{$0 \notin \ran{s}$}
	\\يعنې \LR{$\forall x\ s(x) \neq 0$}
	\item بېل عددونه بېل تالي لري
	\\يعنې \LR{$s$} يوه يو پر يو تابع ده
	\\يعنې \LR{$\forall x \forall y (s(x) = s(y) \lif x = y)$}
	\item\ollabel{repeatedapplication} هر عدد پر \LR{$0$} د تالي تابعې له
	بيا بيا تطبيق څخه ترلاسه کېږي۔
\end{enumerate}
لومړي دوه شرطونه د لومړۍ مرتبې د منطق په کارولو په اسانۍ سمبالېږي
(پورته وګورئ)۔ خو يوازې د لومړۍ مرتبې په منطق سره
\olref{repeatedapplication} نشو سمبالولے۔ د ډېډېکېنډ ستره بريا دا وه
چې شرط \olref{repeatedapplication} ئې په سټي بڼه داسې بيا بيان کړ:
\begin{enumerate}
	\item[3\LR{$'$}.] طبيعي عددونه تر ټولو کوچنے سټ دے چې \emph{د تالي
	تابعې له مخې بند} دے: يعنې که \LR{$s$} د سټ پر هر غړے تطبيق
	کړو، د سټ يو بل غړے ترلاسه کوو۔
\end{enumerate}
خو دا خبره بايد ورو ورو په تفصيل وليکو۔
\olpseganchor{OLP-0051-B005}{end}

\olpseganchor{OLP-0051-B006}{start}
\begin{defn}\ollabel{Closure}
	د هرې تابع \LR{$f$} دپاره، سټ \LR{$X$} د \LR{$f$} له مخې \emph{بند} {هله او يوازې هله چې}
	\LR{$(\forall x \in X)f(x) \in X$} وي۔ اوس د هر \LR{$o$} دپاره تعريف کړئ:
	$$\closureofunder{f}{o} = \bigcap\Setabs{X}{o \in X\text{\RL{ او }}X
\text{\RL{ د \LR{$f$} له مخې بند دے}}}$$
\end{defn}
\textbf{د ژباړې د نسخې سمون (OLINF-001):} د «هرې تابع» دا بيان
په پټه دا فرضوي چې تابع د يوه چاپېريال سټ له ځان څخه ځان ته ده او
ټول څېړل شوي سټونه ئې د ډومېن دننه دي۔ که دا شرط ونۀ منل شي، د
ډومېن څخه بهر د غړي د تابعې قيمت تعريف نۀ لري او د رېنج او ټاکلي
غړي اتحاد هم لازماً د تابعې له مخې بند نۀ وي۔ لاندې کارونې همدا
ځان-ته-تابع حالت لري۔
\olpseganchor{OLP-0051-B006}{end}

\olpseganchor{OLP-0051-B007}{start}
نو \LR{$\closureofunder{f}{o}$} د ټولو هغو سټونو اشتراک دے چې د \LR{$f$} له
مخې بند دي او \LR{$o$} ئې يو غړے دے۔ په وجداني ډول، نو
\LR{$\closureofunder{f}{o}$} تر ټولو \emph{کوچنے} داسې د \LR{$f$} له مخې بند
سټ دے چې \LR{$o$} ئې يو غړے دے۔ راتلونکې پايله دا وجداني خبره
دقيقه کوي؛
\begin{lem}\ollabel{closureproperties}
	د هرې تابع \LR{$f$} او هر \LR{$o \in A$} دپاره:
	\begin{enumerate}
		\item\ollabel{closurehaselem} \LR{$o \in \closureofunder{f}{o}$}؛ او
		\item\ollabel{closureclosed} \LR{$\closureofunder{f}{o}$} د \LR{$f$} له مخې بند دے؛ او
		\item\ollabel{closuresmallest} که \LR{$X$} د \LR{$f$} له مخې بند وي او \LR{$o \in
		X$} وي، نو \LR{$\closureofunder{f}{o} \subseteq X$}
	\end{enumerate}
\end{lem}
\olpseganchor{OLP-0051-B007}{end}

\olpseganchor{OLP-0051-B008}{start}
\begin{proof}
پام وکړئ چې لږ تر لږه يو داسې سټ شته چې د \LR{$f$} له مخې بند دے او \LR{$o$}
ئې يو غړے دے، يعنې \LR{$\ran{f}\cup \{o\}$}۔ نو
\LR{$\closureofunder{f}{o}$}، چې د دې \emph{ټولو} سټونو اشتراک دے، شتون
لري۔ اوس بايد \olref{closurehaselem}--\olref{closuresmallest} وارزوو۔
\olpseganchor{OLP-0051-B008}{end}

\olpseganchor{OLP-0051-B009}{start}
د \olref{closurehaselem} په اړه: \LR{$o \in \closureofunder{f}{o}$}،
ځکه دا د هغو سټونو اشتراک دے چې په ټولو کښې \LR{$o$} يو غړے دے۔
\olpseganchor{OLP-0051-B009}{end}

\olpseganchor{OLP-0051-B010}{start}
د \olref{closureclosed} په اړه: فرض کړئ \LR{$x \in
\closureofunder{f}{o}$}۔ نو که \LR{$o \in X$} او \LR{$X$} د \LR{$f$} له مخې بند
وي، \LR{$x \in X$}؛ او ځکه \LR{$X$} د \LR{$f$} له مخې بند دے، اوس \LR{$f(x) \in X$}۔
نو \LR{$f(x) \in \closureofunder{f}{o}$}۔
\olpseganchor{OLP-0051-B010}{end}

\olpseganchor{OLP-0051-B011}{start}
د \olref{closuresmallest} په اړه: په بشپړ عموميت کښې، که \LR{$X
\in C$} وي، نو \LR{$\bigcap C \subseteq X$}۔
\end{proof}
\olpseganchor{OLP-0051-B011}{end}

\olpseganchor{OLP-0051-B012}{start}
د دې په کارولو ويلے شو:
\olpseganchor{OLP-0051-B012}{end}

\olpseganchor{OLP-0051-B013}{start}
\begin{defn}
د \emph{ډېډېکېنډ الجبر} يو سټ \LR{$A$}، يوه تابع \LR{$f \colon A \to A$} او
يو \LR{$o \in A$} دي چې لاندې شرطونه پوره کوي:
	\begin{enumerate}
		\item \ollabel{ded:proper} \LR{$o \notin \ran{f}$}
		\item \ollabel{ded:injection} \LR{$f$} يوه يو پر يو تابع ده
		\item \ollabel{ded:closure} \LR{$A = \closureofunder{f}{o}$}
	\end{enumerate}
\end{defn}
\olpseganchor{OLP-0051-B013}{end}

\olpseganchor{OLP-0051-B014}{start}
ځکه \LR{$A = \closureofunder{f}{o}$} دے، زموږ پخوانۍ پايله وايي چې \LR{$A$}
تر ټولو کوچنے داسې د \LR{$f$} له مخې بند سټ دے چې \LR{$o$} ئې يو غړے
دے۔ په څرګند ډول هر ډېډېکېنډ الجبر د ډېډېکېنډ له مخې نامتناهي دے؛
يوازې د تعريف \olref{ded:proper} او \olref{ded:injection} مادې
وګورئ۔ خو تر دې په زړه پورې حقيقت دا دے چې هر د ډېډېکېنډ له مخې
نامتناهي سټ په ډېډېکېنډ الجبر بدلېدے شي۔
\olpseganchor{OLP-0051-B014}{end}

\olpseganchor{OLP-0051-B015}{start}
\begin{thm}\ollabel{thm:DedekindInfiniteAlgebra}
که کوم د ډېډېکېنډ له مخې نامتناهي سټ وي، نو يو ډېډېکېنډ الجبر شته۔
\end{thm}
\olpseganchor{OLP-0051-B015}{end}

\olpseganchor{OLP-0051-B016}{start}
\begin{proof}
فرض کړئ \LR{$D$} د ډېډېکېنډ له مخې نامتناهي دے۔ نو يوه يو پر يو تابع
\LR{$g \colon D \to D$} او يو غړے \LR{$o \in D \setminus \ran{g}$} شته۔
اوس \LR{$A = \closureofunder{g}{o}$} وټاکئ؛ د
\olref{closureproperties} له مخې \LR{$A$} شته او \LR{$o \in A$}۔
\LR{$f = \funrestrictionto{g}{A}$} وټاکئ۔ ښيو چې \LR{$A, f, o$} يو
ډېډېکېنډ الجبر جوړوي۔
\olpseganchor{OLP-0051-B016}{end}

\olpseganchor{OLP-0051-B017}{start}
د \olref{ded:proper} په اړه: \LR{$o \notin \ran{g}$} او \LR{$\ran{f}
\subseteq \ran{g}$}، نو \LR{$o\notin \ran{f}$}۔
\olpseganchor{OLP-0051-B017}{end}

\olpseganchor{OLP-0051-B018}{start}
د \olref{ded:injection} په اړه: \LR{$g$} پر \LR{$D$} يوه يو پر يو تابع ده؛
نو \LR{$f \subseteq g$} هم بايد يوه يو پر يو تابع وي۔
\olpseganchor{OLP-0051-B018}{end}

\olpseganchor{OLP-0051-B019}{start}
د \olref{ded:closure} په اړه: د \olref{closureproperties} له مخې
\LR{$A$} د \LR{$g$} له مخې بند دے؛ نو په لا قوي ډول \LR{$A$} د \LR{$f$} له مخې بند
دے۔ په دې توګه د \olref{closureproperties} له مخې
\LR{$\closureofunder{f}{o} \subseteq A$}۔ ځکه \LR{$\closureofunder{f}{o}$}
هم د \LR{$f$} له مخې بند دے او \LR{$f = \funrestrictionto{g}{A}$}، نو
\LR{$\closureofunder{f}{o}$} د \LR{$g$} له مخې بند دے۔ بيا د
\olref{closureproperties} له مخې
\LR{$A \subseteq \closureofunder{f}{o}$}۔
\end{proof}
\olpseganchor{OLP-0051-B019}{end}

\olpseganchor{OLP-0052-B004}{start}
\olfileid{sfr}{infinite}{induction}
\olsection[رياضي استقرا]{د ډېډېکېنډ الجبرونه او رياضي استقرا}
\olpseganchor{OLP-0052-B004}{end}

\olpseganchor{OLP-0052-B005}{start}
اوس مهمه خبره دا ده چې يو ډېډېکېنډ الجبر---بلکې \emph{هر}
ډېډېکېنډ الجبر---د طبيعي عددونو د بدل په توګه کار کولے شي۔ دا د
لاندې اسانې نتيجې له برکته ده:
\olpseganchor{OLP-0052-B005}{end}

\olpseganchor{OLP-0052-B006}{start}
\begin{thm}[رياضي استقرا]\ollabel{thm:dedinfiniteinduction}
فرض کړئ \LR{$N, s, o$} يو ډېډېکېنډ الجبر جوړوي۔ بيا د هر سټ \LR{$X$} دپاره:
	\begin{center}
		که \LR{$o \in X$} او \LR{$(\forall {n} \in N \cap X){s}({n}) \in X$} وي، {نو} \LR{$N \subseteq X$}۔
	\end{center}
\end{thm}
\olpseganchor{OLP-0052-B006}{end}

\olpseganchor{OLP-0052-B007}{start}
\begin{proof}
د ډېډېکېنډ الجبر د تعريف له مخې \LR{$N = \closureofunder{s}{o}$}۔ اوس
که هم \LR{${o} \in X$} او هم \LR{$(\forall {n} \in N)(n \in X \lif
{s}({n}) \in X)$} وي، نو \LR{$N = \closureofunder{s}{o} \subseteq X$}۔
\end{proof}
\olpseganchor{OLP-0052-B007}{end}

\olpseganchor{OLP-0052-B008}{start}
ځکه استقرا د طبيعي عددونو ځانګړنه ده، مقصد دا دے۔ هر د ډېډېکېنډ
له مخې نامتناهي سټ راکړل شي، ترې يو ډېډېکېنډ الجبر جوړولے شو او دا
الجبر د طبيعي عددونو د بدل په توګه کارولے شو۔
\olpseganchor{OLP-0052-B008}{end}

\olpseganchor{OLP-0052-B009}{start}
دا منو چې \olref{thm:dedinfiniteinduction} استقرا په \emph{سټي}
اصطلاحاتو کښې بيانوي۔ خو دا اصل په اسانۍ په هغو اصطلاحاتو کښې هم
بيانولے شو چې ښايي ډېرې اشنا وي:
\olpseganchor{OLP-0052-B009}{end}

\olpseganchor{OLP-0052-B010}{start}
\begin{cor}\ollabel{natinductionschema}
فرض کړئ \LR{$N, s, o$} يو ډېډېکېنډ الجبر جوړوي۔ بيا د هر فارمول
\LR{$\phi(x)$} دپاره، چې پاراميټرونه لرلے شي:
\begin{center}
	که \LR{$\phi(o)$} او \LR{$(\forall {n} \in N)(\phi(n)\lif
	\phi({s}({n})))$} وي، {نو} \LR{$(\forall n \in N)\phi(n)$}
\end{center}
\end{cor}
\olpseganchor{OLP-0052-B010}{end}

\olpseganchor{OLP-0052-B011}{start}
\begin{proof}
\LR{$X = \Setabs{n \in N}{\phi(n)}$} وټاکئ، او اوس
\olref{thm:dedinfiniteinduction} وکاروئ۔
\end{proof}
\olpseganchor{OLP-0052-B011}{end}

\olpseganchor{OLP-0052-B012}{start}
په دې پايله کښې مو وويل چې يو فارمول «پاراميټرونه لري»۔ په لنډه
معنا ئې دا ده چې د هر ډول څيزونو \LR{$c_1, \ldots, c_k$} دپاره له
\LR{$\phi(x, c_1, \ldots, c_k)$} سره کار کولے شو۔ په لا دقيق ډول، پايله
بې له دې چې «پاراميټرونه» ياد کړو، داسې بيانولے شو۔ د هر فارمول
\LR{$\phi(x, v_1, \ldots, v_k)$} دپاره، چې ټول ازاد متغيرونه ئې ښودل
شوي وي، لرو:
	\begin{align*}
			\forall v_1 \ldots \forall v_k((&\phi(o, v_1,\ldots, v_k) \land {}\\
			&	(\forall x \in N)(\phi(x,v_1, \ldots, v_k) \lif \phi(s(x), v_1,\ldots, v_k))) \lif {}\\
			&\hspace{3em} (\forall x \in N)\phi(x, v_1,\ldots, v_k))
	\end{align*}
په څرګند ډول، د «پاراميټرونو لرلو» خبره لوستل ډېر اسانه کولے شي۔
(په \emph{د OpenLogic اړونده سرچينيزه برخه}\nobreakspace\LR{\href{https://github.com/OpenLogicProject/OpenLogic/blob/9620cc73f9c8e0ad003c514a5d3748f29611c4c0/content/set-theory/set-theory.tex\#L7}{\latinfont\scriptsize [OLP]}} کښې به دا طريقه ډېره وکاروو۔)
\olpseganchor{OLP-0052-B012}{end}

\olpseganchor{OLP-0052-B013}{start}
بېرته د ډېډېکېنډ الجبرونو خبرې ته راځو: هر ډېډېکېنډ الجبر راکړل
شي، د جمع، ضرب او طاقت معمول حسابي تابعې هم تعريفولے شو۔ خو دا
اسان کار نۀ دے او د \emph{بازګشتي تعريف} طريقه کاروي۔ دا هغه طريقه
ده چې ډېر وروسته، په ډېر عمومي چوکاټ کښې، معرفي او توجيه کوو۔
(مينه‌وال ښايي له \emph{د OpenLogic اړونده سرچينيزه برخه}\nobreakspace\LR{\href{https://github.com/OpenLogicProject/OpenLogic/blob/9620cc73f9c8e0ad003c514a5d3748f29611c4c0/content/set-theory/ord-arithmetic/ord-arithmetic.tex\#L8}{\latinfont\scriptsize [OLP]}} وروسته دې
موضوع ته بيا راشي، يا د \citealt[pp.~95--8]{Potter2004} په څېر
کومه بله شننه ولولي۔) خو چې \LR{$N, s, o$} يو ډېډېکېنډ الجبر جوړوي، په
پاے کښې به لاندې ټاکنې کولے شو:
\begin{align*}
	{a} + {o} &= {a} & & & {a} \times {o} &= {o} & & & {a}^{o} &= s(o)\\
{a} + {s}({b}) &= {s}({a}+{b}) &&& {a} \times {s}({b}) &= ({a}\times {b}) + {a}  & & & {a}^{{s}({b})} &= {a}^{b} \times {a}
\end{align*}
او ښودلے شو چې دا تابعې لکه څنګه چې هيله کېږي، هماغسې چلند کوي۔
\olpseganchor{OLP-0052-B013}{end}

\olpseganchor{OLP-0053-B003}{start}
\olfileid{sfr}{infinite}{dedekindsproof}
\olpseganchor{OLP-0053-B003}{end}

\olpseganchor{OLP-0053-B004}{start}
\olsection[د ډېډېکېنډ «ثبوت»]{د نامتناهي سټ د شتون په اړه د
ډېډېکېنډ «ثبوت»}
\olpseganchor{OLP-0053-B004}{end}

\olpseganchor{OLP-0053-B005}{start}
په دې څپرکي کښې مو د طبيعي عددونو سټي شننه د ډېډېکېنډ د الجبرونو
له مخې وړاندې کړه۔ په \olref[arith][ref]{sec} کښې مو (له نورو
خبرو سره) د تحليل د سټي جوړونې پر فلسفي اهميت غور وکړ۔ اوس بايد د
هغه څه اهميت وارزوو چې دلته مو ترلاسه کړل۔
\olpseganchor{OLP-0053-B005}{end}

\olpseganchor{OLP-0053-B006}{start}
د \olref[sfr][arith][]{chap} په ټوله موده کښې مو طبيعي عددونه
راکړل شوي وګڼل او په څرګند ډول مو ترې صحيح، ناطق او حقيقي عددونه
جوړ کړل۔ په دې څپرکي کښې مو د طبيعي عددونو کوم څرګند جوړښت نۀ دے
وړاندې کړے۔ يوازې مو وښودل چې \emph{هر د ډېډېکېنډ له مخې نامتناهي
سټ راکړل شي}، داسې سټ تعريفولے شو چې کټ مټ لکه څنګه چې غواړو
\LR{$\Nat$} وچلېږي، هماغسې به وچلېږي۔
\olpseganchor{OLP-0053-B006}{end}

\olpseganchor{OLP-0053-B007}{start}
نو څرګنده ده چې دا ادعا نشو کولے چې کومې ماوراءالطبيعتي پوښتنې ته
مو ځواب ورکړے، لکه دا چې \emph{طبيعي عددونه کوم څيزونه دي}۔ خو دا
ښۀ کار دے۔ په \olref[sfr][arith][ref]{sec} کښې مو ټينګار کړے و چې
د ډېډېکېنډ د پرېکونونو د سټ په توګه د \LR{$\Real$} تعريف بايد يوه اسانه
ټاکنه وګڼو، نۀ يوه \emph{موندنه}۔ مهمه مشاهده دا ده چې د
ډېډېکېنډ پرېکونونه د حقيقي عددونو بنسټيز رياضي خاصيتونه لري۔ دلته
هم همدا خبره ده: مهمه مشاهده دا ده چې \emph{هر} ډېډېکېنډ الجبر د
طبيعي عددونو بنسټيز رياضي خاصيتونه لري۔ (په رښتيا، ډېډېکېنډ دا
ټکے په دې ثبوت لا پياوړے کړ چې ټول ډېډېکېنډ الجبرونه
\emph{آيزومورف} دي (\citeyear[Theorems
132--3]{Dedekind1888})۔ نو
دا هېڅ د حيرانۍ خبره نۀ ده چې اوسني ډېر «جوړښت‌پال» ډېډېکېنډ خپل
مخکښ بولي۔)
\olpseganchor{OLP-0053-B007}{end}

\olpseganchor{OLP-0053-B008}{start}
برسېره پر دې، موږ وښودل چې د طبيعي عددونو نظريه څنګه په ساده سټ
نظريه کښې ورګډېږي؛ دا ساده نظريه خپله لا هم يو څه غېر رسمي ده، خو
(لکه چې ښکاري) طبيعي عددونه له مخکښې راکړل شوي نۀ ګڼي۔ نو کېدے شي
د ډېډېکېنډ د خپلې سترې پروژې د پوره کولو پر لار يو، چې هغه داسې
تشريح کړې وه:
\begin{quote}
	په ساينس کښې هېڅ د ثبوت وړ خبره بايد بې له ثبوته ونۀ منل شي۔ که
	څه هم دا غوښتنه معقوله ښکاري، زۀ دا نشم منلے چې آن د تر ټولو ساده
	ساينس د بنسټ ايښودلو په وروستيو طريقو کښې پوره شوې ده؛ يعنې د منطق
	په هغې برخې کښې چې د عددونو نظريه څېړي۔ کله چې حساب (الجبر،
	تحليل) يوازې د منطق يوه برخه بولم، مقصد مې دا دے چې د عدد مفهوم
	د مکان او زمان له تصورونو يا شهود څخه بشپړ خپلواک ګڼم---بلکې دا
	د فکر د خالصو قوانينو نېغه پايله ګڼم۔
	\citep[preface]{Dedekind1888}
\end{quote}
د ډېډېکېنډ زړه‌وره مفکوره دا ده۔ همدا اوس مو وښودل چې يوازې د
(ساده) سټ نظريې په کارولو طبيعي عددونه څنګه جوړېږي۔ په
\olref[sfr][arith][]{chap} کښې مو وليدل چې د طبيعي عددونو او يوې
اندازې سټ نظريې په راکړه کېدو سره حقيقي عددونه څنګه جوړېږي۔ نو
ښايي «حساب (الجبر، تحليل)» «يوازې د منطق يوه برخه» وګرځي (د
ډېډېکېنډ د «منطق» کلمې په پراخه معنا)۔
\olpseganchor{OLP-0053-B008}{end}

\olpseganchor{OLP-0053-B009}{start}
نظر همدا دے۔ خو يوه شېبه تم شئ۔ زموږ د ډېډېکېنډ الجبر جوړول (چې
د طبيعي عددونو بدل دے) د يوه د ډېډېکېنډ له مخې نامتناهي سټ د شتون
پر شرط ولاړ دي۔ (يوازې \olref[sfr][infinite][dedekind]{thm:DedekindInfiniteAlgebra}
ته بېرته وګورئ۔) تر څو چې د ډېډېکېنډ له مخې د نامتناهي سټ شتون د
«منطق» يا «د فکر د خالصو قوانينو» له لارې ثابت نشي، پروژه درېږي۔
\olpseganchor{OLP-0053-B009}{end}

\olpseganchor{OLP-0053-B010}{start}
نو، \emph{ايا} د ډېډېکېنډ له مخې د نامتناهي سټ شتون «د فکر د
خالصو قوانينو» له لارې ثابتېدے شي؟ د ډېډېکېنډ هڅه دا وه:
\begin{quote}
  زما د فکرونو خپله نړۍ، يعنې د ټولو هغو څيزونو بشپړتيا \LR{$S$} چې زما
  د فکر موضوع کېدے شي، نامتناهي ده۔ ځکه که \LR{$s$} د~\LR{$S$} يو غړے وي،
  نو دا فکر \LR{$s'$} چې~\LR{$s$} زما د فکر موضوع کېدے شي، خپله د~\LR{$S$} يو
  غړے دے۔ که دا د غړي~\LR{$s$} يو تصوير \LR{$\phi(s)$} وګڼو، نو
  \dots~\LR{$S$} [د ډېډېکېنډ له مخې] نامتناهي دے، او همدا بايد ثابت
  شوي وے۔
	\citep[\S66]{Dedekind1888}
\end{quote}
داسې خبره د يوه داسې کتاب په منځ کښې ليدل رښتيا هم ډېره حيرانوونکې
ده چې زياته برخه ئې له بشپړو دقيقو رياضي ثبوتونو جوړه ده۔ دوه
تبصرې د پام وړ دي۔
\olpseganchor{OLP-0053-B010}{end}

\olpseganchor{OLP-0053-B011}{start}
لومړۍ: دا «ثبوت» هغه بڼه نږدې هېڅ نۀ لري چې اوس به ئې «رياضي»
وبولو۔ دا د نفسياتي څيزونو (فکرونو) خبره کوي، او هغه هم يوازې
\emph{ممکن} څيزونه۔
\olpseganchor{OLP-0053-B011}{end}

\olpseganchor{OLP-0053-B012}{start}
دويمه: لږ تر لږه لکه څنګه چې موږ د ډېډېکېنډ الجبرونه وړاندې کړي،
دا «ثبوت» يوه څرګنده فني نيمګړتيا لري۔ که د ډېډېکېنډ استدلال بريالے
وي، يوازې دا ثابتوي چې نامتناهي ډېر څيزونه (په ځانګړې توګه
نامتناهي ډېر فکرونه) شته۔ خو ډېډېکېنډ بايد دا دليل هم راکړي چې~\LR{$S$}
ولې يو يوازينے \emph{سټ} وګڼو چې نامتناهي ډېر غړي لري،
نۀ دا چې~\LR{$S$} \emph{څه څيزونه} (په جمع بڼه) وګڼو۔
\olpseganchor{OLP-0053-B012}{end}

\olpseganchor{OLP-0053-B013}{start}
دا حقيقت چې ډېډېکېنډ دلته تشه نۀ ليدله، ښايي وښيي چې د هغه د
«بشپړتيا» کلمه کټ مټ زموږ د «سټ» کلمې له کارونې سره برابره نۀ
ده۔\footnote{په رښتيا، د دې د باور دپاره نور دليلونه هم لرو؛
\citet[p.~23]{Potter2004} وګورئ۔} خو دا به ډېره د حيرانۍ خبره نۀ
وي۔ هغه پروژه چې په تېرو دوو څپرکيو کښې مو مخ ته يوړه---د طبيعي
عددونو «جوړول»، او له هغو څخه د صحيح، حقيقي او ناطقو عددونو
«جوړول»---ټوله په ساده ډول شوې ده۔ کله مو چې کوم سټ ته اړتيا لرله،
هماغه سټ مو منلے؛ خو د دې په اړه مو ډېر عمومي اصلونه نۀ دي ټاکلي
چې کټ مټ کوم سټونه کله شتون لري۔ خو پوهېږو چې \emph{ځينې} عمومي
اصلونه پکار دي، ځکه له هغو پرته به د رسل په پاراډوکس کښې ولوېږو۔
\olpseganchor{OLP-0053-B013}{end}

\olpseganchor{OLP-0053-B014}{start}
اوس د دې وخت دے چې له خپلې ساده‌ګۍ څخه ووځو۔
\olpseganchor{OLP-0053-B014}{end}

\olpseganchor{OLP-0054-B003}{start}
\olfileid{sfr}{infinite}{card-sb}
\olpseganchor{OLP-0054-B003}{end}

\olpseganchor{OLP-0054-B004}{start}
\olsection{ضميمه: د شروډر--برنشتاين ثبوت}
\olpseganchor{OLP-0054-B004}{end}

\olpseganchor{OLP-0054-B005}{start}
له ساده سټ نظريې څخه د وتلو مخکښې يو وروستے ساده (خو ډېر فني!)
ثبوت څېړو۔ دا د شروډر--برنشتاين ثبوت دے
(\olref[sfr][siz][sb]{thm:schroder-bernstein}): که
\LR{$\cardle{A}{B}$} او \LR{$\cardle{B}{A}$} وي، نو \LR{$\cardeq{A}{B}$}؛ يعنې که
يو پر يو تابعې \LR{$f \colon A \to B$} او \LR{$g \colon B \to A$} راکړل
شي، نو يوه دوه اړخيزه يو پر يو تابع \LR{$h \colon A \to B$} شته۔
\olpseganchor{OLP-0054-B005}{end}

\olpseganchor{OLP-0054-B006}{start}
په دې څپرکي کښې مو د ډېډېکېنډ د \emph{بندښتونو} مفهوم وکاراوه۔ په
رښتيا، ډېډېکېنډ په همدې مفهوم د شروډر--برنشتاين يو ښکلی ثبوت ورکړ،
چې دلته ئې وړاندې کوو۔ که لږ بېله خو په اصل کښې ورته شننه غواړئ،
دا ثبوت په نژدې ډول د \citet[pp.~157--8]{Potter2004} پيروي کوي۔
په انټرنېټ کښې لږه پلټنه به هم درته وښيي چې دا---لکه د
\LR{$\sqrt{2}$} غير ناطق‌والے---هغه قضيه ده چې \emph{ډېر} په زړه پورې
او بېل ثبوتونه لري۔
\olpseganchor{OLP-0054-B006}{end}

\olpseganchor{OLP-0054-B007}{start}
د \olref[sfr][infinite][dedekind]{Closure} په شان نښه وکاروئ او د
هر سټ \LR{$B$} او تابعې~\LR{$f$} دپاره وټاکئ:
\[
\Closureofunder{f}{B} = \bigcap \Setabs{X}{B \subseteq X
\text{\RL{ او \LR{$X$} د \LR{$f$} له مخې بند دے}}}
\]
په دې ډول تعريف شوے \LR{$\Closureofunder{f}{B}$} تر ټولو کوچنے داسې د
\LR{$f$} له مخې بند سټ دے چې~\LR{$B$} رانغاړي، يعنې:
\olpseganchor{OLP-0054-B007}{end}

\olpseganchor{OLP-0054-B008}{start}
\begin{lem}\ollabel{Closureprops}
	د هرې تابع \LR{$f$} او هر \LR{$B$} دپاره:
	\begin{enumerate}
		\item\ollabel{Closurehaselem} \LR{$B \subseteq \Closureofunder{f}{B}$}؛ او
		\item\ollabel{Closureclosed} \LR{$\Closureofunder{f}{B}$} د \LR{$f$} له مخې بند دے؛ او
		\item\ollabel{Closuresmallest} که \LR{$X$} د \LR{$f$} له مخې بند وي او \LR{$B
		\subseteq X$} وي، نو \LR{$\Closureofunder{f}{B} \subseteq X$}۔
	\end{enumerate}
\end{lem}
\olpseganchor{OLP-0054-B008}{end}

\olpseganchor{OLP-0054-B009}{start}
\begin{proof}
کټ مټ لکه په \olref[sfr][infinite][dedekind]{closureproperties} کښې۔
\end{proof}
\olpseganchor{OLP-0054-B009}{end}

\olpseganchor{OLP-0054-B010}{start}
شروډر--برنشتاين ته د رسېدو دپاره يو وروستے حقيقت پکار دے:
\olpseganchor{OLP-0054-B010}{end}

\olpseganchor{OLP-0054-B011}{start}
\begin{prop}\ollabel{sbhelper}
که \LR{$A \subseteq B \subseteq C$} او \LR{$A \approx C$} وي، نو \LR{$\cardeq{\cardeq{A}{B}}{C}$}۔
\end{prop}
\olpseganchor{OLP-0054-B011}{end}

\olpseganchor{OLP-0054-B012}{start}
\begin{proof}
يوه دوه اړخيزه يو پر يو تابع \LR{$f \colon C \to A$} راکړل شي، وټاکئ
\LR{$F = \Closureofunder{f}{C \setminus B}$} او له ډومېن \LR{$C$} سره تابع
\LR{$g$} داسې تعريف کړئ:
\[
	g(x) =
	\begin{cases}
			f(x) &\text{\RL{که \LR{$x \in F$} وي}}\\
			x & \text{\RL{که نۀ}}
	\end{cases}
\]
ښيو چې \LR{$g$} له \LR{$C \to B$} يوه دوه اړخيزه يو پر يو تابع ده؛ له دې به راووځي
چې \LR{$\comp{f^{-1}}{g} \colon A \to B$} يوه دوه اړخيزه يو پر يو تابع ده، او
ثبوت بشپړېږي۔
\olpseganchor{OLP-0054-B012}{end}

\olpseganchor{OLP-0054-B013}{start}
لومړے ادعا کوو چې که \LR{$x \in F$} خو \LR{$y\notin F$} وي، نو
\LR{$g(x) \neq g(y)$}۔ د نقيض دپاره فرض کړئ چې برابر دي؛ نو
\LR{$y = g(y) = g(x) = f(x)$}۔ ځکه \LR{$x \in F$} او د
\olref{Closureprops} له مخې \LR{$F$} د \LR{$f$} له مخې بند دے،
\LR{$y = f(x) \in F$} لرو، چې ټکر دے۔
\olpseganchor{OLP-0054-B013}{end}

\olpseganchor{OLP-0054-B014}{start}
اوس فرض کړئ \LR{$g(x) = g(y)$}۔ د پورته خبرې له مخې، \LR{$x \in F$} هله او
يوازې هله چې \LR{$y \in F$} وي۔ که \LR{$x, y \in F$} وي، نو
\LR{$f(x) = g (x) = g(y) = f(y)$}؛ او ځکه \LR{$f$} يوه دوه اړخيزه يو پر يو تابع ده،
\LR{$x = y$}۔ که \LR{$x, y \notin F$} وي، نو \LR{$x = g(x) = g(y) = y$}۔ نو
\LR{$g$} يوه يو پر يو تابع ده۔
\olpseganchor{OLP-0054-B014}{end}

\olpseganchor{OLP-0054-B015}{start}
دا پاتې ده چې وښيو \LR{$\ran{g} = B$}۔ يو \LR{$x \in B \subseteq C$}
وټاکئ۔ که \LR{$x \notin F$} وي، نو \LR{$g(x) = x$}۔ که \LR{$x \in F$} وي، نو د
کوم \LR{$y \in F$} دپاره \LR{$x = f(y)$}؛ ځکه که داسې نۀ وي، نو
\LR{$F \setminus \{x\}$} به د \LR{$f$} له مخې بند وي او \LR{$C\setminus B$} به
وغځوي، چې د \olref{Closureprops} له مخې ناشونې ده؛ اوس
\LR{$g(y) = f(y) = x$}۔
\end{proof}
\olpseganchor{OLP-0054-B015}{end}

\olpseganchor{OLP-0054-B016}{start}
په پاے کښې د اصلي پايلې ثبوت دا دے۔ ياد کړئ چې د تابعې \LR{$h$} او سټ
\LR{$D$} دپاره \LR{$\funimage{h}{D} = \Setabs{h(x)}{x \in D}$} تعريفوو۔
\olpseganchor{OLP-0054-B016}{end}

\olpseganchor{OLP-0054-B017}{start}
\begin{proof}[د شروډر--برنشتاين ثبوت] فرض کړئ \LR{$f \colon A \to B$}
او \LR{$g \colon B \to A$} يو پر يو تابعې دي۔ ځکه
\LR{$\funimage{f}{A} \subseteq B$}، نو
\LR{$\funimage{g}{\funimage{f}{A}} \subseteq g[B] \subseteq A$}۔
برسېره پر دې، \LR{$\comp{f}{g} \colon A \to
\funimage{g}{\funimage{f}{A}}$} يوه يو پر يو تابع ده، ځکه \LR{$g$} او
\LR{$f$} دواړه همداسې دي؛ او په رښتيا \LR{$\comp{f}{g}$} يوه
دوه اړخيزه يو پر يو تابع ده، يوازې په دې سبب چې مقابل ډومېن مو همداسې تعريف
کړے۔ نو \LR{$\cardeq{\funimage{g}{\funimage{f}{A}}}{A}$}، او په دې توګه
د \olref{sbhelper} له مخې يوه دوه اړخيزه يو پر يو تابع \LR{$h \colon A \to
\funimage{g}{B}$} شته۔ برسېره پر دې، \LR{$g^{-1}$} يوه دوه اړخيزه يو پر يو تابع
\LR{$\funimage{g}{B} \to B$} ده۔ نو \LR{$\comp{h}{g^{-1}} \colon A \to B$}
يوه دوه اړخيزه يو پر يو تابع ده۔
\end{proof}
\olpseganchor{OLP-0054-B017}{end}

\olpseganchor{OLP-0055-B004}{start}
\olpart{pl}{بياني منطق}
\olpseganchor{OLP-0055-B004}{end}

\olpseganchor{OLP-0055-B005}{start}
\begin{editorial}
  په دې برخه کښې د کلاسيکي بياني منطق مواد دي۔ لومړے څپرکے
  نسبتاً بنسټيز دے او يوازې تعريفونه او پايلې راغونډوي؛ ډېر ثبوتونه
  نۀ دي بشپړ شوي، بلکې د تمرينونو په توګه پرېښودل شوي دي۔ د ثبوت د
  نظامونو او د بشپړتيا د قضيې مواد د لومړۍ درجې منطق له برخې څخه
  راوړل شوي، چې د لومړۍ درجې منطق ټېګ پکې ناسم ټاکل شوے دے۔ په دې
  توګه د
  محمولونو، اصطلاحاتو او سورونو اړوند هر څه ترې وځي، او د
  جوړښتونه~\LR{$\Struct{M}$} پر ځاے د
  ارزښت ټاکنې~\LR{$\pAssign{v}$} خبره راځي۔
\olpseganchor{OLP-0055-B005}{end}

\olpseganchor{OLP-0055-B006}{start}
  پلان دا دے چې دا برخه لا په تفصيل پراخه شي، او نورې موضوعګانې او
  پايلې، لکه د رښتياوالي تابعوي بشپړتيا، هم ورزياتې شي۔
\end{editorial}
\olpseganchor{OLP-0055-B006}{end}



\olpseganchor{OLP-0055-B008}{start}
\tagfalse{FOL}
\olpseganchor{OLP-0055-B008}{end}



\olpseganchor{OLP-0055-B010}{start}
%
  

\olpseganchor{OLP-0055-B010}{end}

\olpseganchor{OLP-0055-B011}{start}
%
  

\olpseganchor{OLP-0055-B011}{end}

\olpseganchor{OLP-0055-B012}{start}
%
  

\olpseganchor{OLP-0055-B012}{end}

\olpseganchor{OLP-0055-B013}{start}
%
  

\olpseganchor{OLP-0055-B013}{end}



\olpseganchor{OLP-0055-B015}{start}
\tagtrue{FOL}
\olpseganchor{OLP-0055-B015}{end}

\olpseganchor{OLP-0056-B004}{start}
\olchapter{pl}{syn}{نحو او معنٰی پوهنه}
\olpseganchor{OLP-0056-B004}{end}

\olpseganchor{OLP-0056-B005}{start}
\begin{editorial}
  دا يوازې د تعريفونو ډېره لنډه لنډيزه ده۔ بايد پراخه شي، څو د
  فارمولونو پر جوړښت د استقرا د ثبوتونو اسانه پېژندنه او ډېرې نورې
  بېلګې پکې راشي۔
\end{editorial}
\olpseganchor{OLP-0056-B005}{end}

\olpseganchor{OLP-0057-B004}{start}
\olfileid{pl}{syn}{int}
\olpseganchor{OLP-0057-B004}{end}

\olpseganchor{OLP-0057-B005}{start}
\olsection{پېژندنه}
\olpseganchor{OLP-0057-B005}{end}

\olpseganchor{OLP-0057-B006}{start}
بياني منطق له هغو فارمولونه سره کار لري چې له
بياني متغيرونه څخه د \LR{$\lnot$}، \LR{$\land$}، \LR{$\lor$}، \LR{$\lif$} او
\LR{$\liff$} بياني نښلوونکو په کارولو جوړېږي۔ د تصور له مخې،
بياني متغير~\LR{$p$} د داسې يوې جملې يا بيان ځاے نيسي
چې رښتونے يا دروغژن وي۔ هر کله چې په فارمول کښې د
بياني متغير ``د رښتياوالي قيمت'' وټاکل شي، نو د
بياني نښلوونکو په کارولو ترې د جوړو شويو ټولو فارمولونه د
رښتياوالي قيمت هم ټاکل کېږي۔ وايو چې بياني منطق \emph{د رښتياوالي
  تابعوي} دے، ځکه چې معنٰی پوهنه ئې د رښتياوالي د قيمتونو د تابعو
له مخې ورکول کېږي۔ په تېره بيا، په بياني منطق کښې دا نۀ څېړو چې
رښتياوالے او دروغوالے نور څنګه ټاکل کېږي؛ مثلاً دا چې يو څه د يوازې
اتفاقي رښتياوالي پر ځاے بالضروره رښتيا دي، يا څوک پرې پوهېږي چې
رښتيا دي، يا اوس رښتيا دي نۀ دا چې پخوا رښتيا وو يا وروسته به رښتيا
وي۔ موږ يوازې د رښتيا (\LR{$\True$}) او دروغو (\LR{$\False$}) دوه قيمتونه په
پام کښې نيسو، نو دا امکان له بحثه وباسو چې يو بيان دې نۀ رښتونے وي
نۀ دروغژن، يا دې نيم رښتونے وي۔ موږ يوازې هغو نښلوونکو ته هم پام
کوو چې ترې د جوړ شوي فارمول د رښتياوالي قيمت د هغې د برخو د
رښتياوالي په قيمتونو پوره ټاکل کېږي، نۀ مثلاً د هغې په معنٰی۔ په
تېره بيا، دا لانجمنه ده چې په انګرېزۍ کښې د شرطيو رښتياوالي قيمت په
دې معنٰی تابعوي دے که نۀ۔ مادي شرطي~\LR{$\lif$} تابعوي دے؛ نور منطقونه
له هغو شرطيو سره کار کوي چې د رښتياوالي تابعوي نۀ وي۔
\olpseganchor{OLP-0057-B006}{end}

\olpseganchor{OLP-0057-B007}{start}
د رښتياوالي د تابعوي بياني منطق د نظريې او مابعدمنطق د ودې دپاره
بايد لومړے د هغه د عبارتونو نحو او معنٰی پوهنه تعريف کړو۔ موږ به له
بياني متغيرونه څخه د نښلوونکو په کارولو د فارمولونه
د جوړولو يوه لار بيان کړو۔ نور تعريفونه هم ممکن دي۔ نور نظامونه به
بېلې نښې غوره کړي، د بنسټيزو نښلوونکو بېل سټونه به وټاکي او قوسونه
به په بېله توګه وکاروي، يا به ئې هېڅ و نۀ کاروي، لکه په تش په نامه
پولنډۍ نښه‌اېښودنه کښې۔ خو په ټولو لارو کښې يو ګډ شے دا دے چې د
جوړېدو قاعدې د فارمولونه سټ \emph{په استقرايي ډول} تعريفوي۔ که
کار په سمه توګه وشي، هر عبارت د جوړېدو د قاعدو له مخې په اصل کښې
يوازې په يوه لار ترلاسه کېدے شي۔ هغه استقرايي تعريف چې
\emph{په يکتا ډول لوستل کېدونکي} عبارتونه ورکوي، موږ ته اجازه راکوي
چې په هماغه طريقه، يعنې په استقرايي تعريف، دغو عبارتونو ته معنٰی
ورکړو۔
\olpseganchor{OLP-0057-B007}{end}

\olpseganchor{OLP-0057-B008}{start}
عبارتونو ته معنٰی ورکول د معنٰی پوهنې ډګر دے۔ د بياني منطق په
معنٰی پوهنه کښې مرکزي مفهوم په ارزښت ټاکنه کښې صدق دے۔
ارزښت ټاکنه~\LR{$\pAssign{v}$} بياني متغيرونه ته د
رښتياوالي \LR{$\True$} او \LR{$\False$} قيمتونه ورکوي۔ هر ارزښت ټاکنه د هر
فارمول~\LR{$!A$} دپاره د رښتياوالي قيمت \LR{$\pValue{v}(!A)$} ټاکي۔
فارمول په ارزښت ټاکنه~\LR{$\pAssign{v}$} کښې هله او يوازې هله
صادق دے چې \LR{$\pValue{v}(!A) = \True$} وي؛ دا په \LR{$\pSat{v}{!A}$} ليکو۔
دا اړيکه د~\LR{$!A$} پر جوړښت د استقرا له مخې هم تعريفېدے شي: د منطقي
نښلوونکو د رښتياوالي تابعې کاروو، څو مثلاً د \LR{$!A \land !B$} صدق د
\LR{$!A$} او~\LR{$!B$} د صدق يا نۀ صدق له مخې تعريف کړو۔
\olpseganchor{OLP-0057-B008}{end}

\olpseganchor{OLP-0057-B009}{start}
د جملو د صدق د اړيکې \LR{$\pSat{v}{!A}$} پر بنسټ بيا د تاوتولوژۍ،
استلزام او د صدق وړتيا بنسټيز معنايي مفهومونه تعريفولے شو۔
فارمول هله تاوتولوژي، \LR{$\Entails !A$}، ده چې هر ارزښت ټاکنه
ئې صادق کړي، يعنې د هر~\LR{$\pAssign{v}$} دپاره
\LR{$\pValue{v}(!A) = \True$} وي۔ د فارمولونه يو سټ هله هغې ته
استلزام ورکوي، \LR{$\Gamma \Entails !A$}، چې هر هغه ارزښت ټاکنه چې د
~\LR{$\Gamma$} ټول فارمولونه صادق کړي، \LR{$!A$} هم صادق کړي۔ او د
فارمولونه يو سټ هله د صدق وړ دے چې کوم ارزښت ټاکنه ئې ټول
فارمولونه په يو وخت صادق کړي۔ ځکه چې فارمولونه په استقرايي
ډول تعريف شوي او صدق بيا د فارمولونه پر جوړښت د استقرا له مخې
تعريف شوے، موږ استقرا کارولے شو چې د خپلې معنٰی پوهنې خاصيتونه ثابت
او تعريف شوي معنايي مفهومونه يو له بل سره وتړو۔
\olpseganchor{OLP-0057-B009}{end}

\olpseganchor{OLP-0058-B004}{start}
\olfileid{pl}{syn}{fml}
\olpseganchor{OLP-0058-B004}{end}

\olpseganchor{OLP-0058-B005}{start}
\olsection{بياني فارمولونه}
\olpseganchor{OLP-0058-B005}{end}

\olpseganchor{OLP-0058-B006}{start}
د بياني منطق فارمولونه له \emph{بياني متغيرونه}، د دروغوالي له
  بياني ثابت~\LR{$\lfalse$}
    او د رښتياوالي له بياني ثابت~\LR{$\ltrue$} څخه د \emph{منطقي
  نښلوونکو} په کارولو جوړېږي۔
\olpseganchor{OLP-0058-B006}{end}

\olpseganchor{OLP-0058-B007}{start}
\begin{enumerate}
\item د بياني متغيرونه \LR{$\Obj p_0$}،
  \LR{$\Obj p_1$}، \dots يو شمېرېدونکے لامتناهي سټ~\LR{$\PVar$}
\item د دروغوالی بياني ثابت~\LR{$\lfalse$}۔
\item د رښتياوالی بياني ثابت~\LR{$\ltrue$}۔
\item منطقي نښلوونکي:
  \startycommalist
  \ycomma \LR{$\lnot$} (نفي)%
  \ycomma \LR{$\land$} (عطف)%
  \ycomma \LR{$\lor$} (فصل)%
  \ycomma \LR{$\lif$} (شرطي)%
  \ycomma \LR{$\liff$} (دوه اړخيز شرطي)%
\item د وقف نښې: (، ) او کامه۔
\end{enumerate}
\olpseganchor{OLP-0058-B007}{end}

\olpseganchor{OLP-0058-B008}{start}
د بياني منطق دا ژبه په \LR{$\Lang L_0$} ښيو۔
\olpseganchor{OLP-0058-B008}{end}

\olpseganchor{OLP-0058-B009}{start}
۔
\olpseganchor{OLP-0058-B009}{end}



\olpseganchor{OLP-0058-B011}{start}
% Alternate symbols
\olpseganchor{OLP-0058-B011}{end}

\olpseganchor{OLP-0058-B012}{start}
\begin{intro}
کېدے شي تاسو له هغو اصطلاحاتو او نښو سره بلد يئ چې زموږ له پورته
کارولو سره توپير لري۔ د منطق متنونه، او ښوونکي، عموماً د ``نفي''
دپاره \LR{$\sim$}، \LR{$\neg$} او~!، او د ``عطف'' دپاره \LR{$\wedge$}، \LR{$\cdot$} او
\LR{$\&$} کاروي۔ د ``شرطي'' يا ``استلزام'' عامې نښې \LR{$\rightarrow$}،
\LR{$\Rightarrow$} او \LR{$\supset$} دي۔
د ``دوه اړخيز شرطي''، ``دوه اړخيز استلزام''
  يا ``(مادي) هم ارزښت'' نښې \LR{$\leftrightarrow$}،
  \LR{$\Leftrightarrow$} او \LR{$\equiv$} دي۔
د \LR{$\lfalse$} نښې ته په بېلو ځايونو کښې
  ``دروغوالے''، ``فالسوم''، ``محال'' يا ``ښکته'' ويل کېږي۔
د \LR{$\ltrue$} نښې ته په بېلو ځايونو کښې
  ``رښتياوالے''، ``ويروم'' يا ``پورته'' ويل کېږي۔
\end{intro}
\olpseganchor{OLP-0058-B012}{end}

\olpseganchor{OLP-0058-B013}{start}
\begin{defn}[فارمول]
\ollabel{defn:formulas}
د بياني منطق د \emph{فارمولونه} سټ~\LR{$\Frm[L_0]$} په استقرايي ډول
داسې تعريفېږي:
\begin{enumerate}
\item \LR{$\lfalse$} يو اټومي فارمول دے۔
\olpseganchor{OLP-0058-B013}{end}

\olpseganchor{OLP-0058-B014}{start}
\item \LR{$\ltrue$} يو اټومي فارمول دے۔
\olpseganchor{OLP-0058-B014}{end}

\olpseganchor{OLP-0058-B015}{start}
\item هر بياني متغير~\LR{$\Obj p_i$} يو اټومي
  فارمول دے۔
\olpseganchor{OLP-0058-B015}{end}

\olpseganchor{OLP-0058-B016}{start}
\item که \LR{$!A$} يو فارمول وي، نو \LR{$\lnot !A$} هم
  يو فارمول دے۔
\olpseganchor{OLP-0058-B016}{end}

\olpseganchor{OLP-0058-B017}{start}
\item که \LR{$!A$} او \LR{$!B$} فارمولونه وي، نو \LR{$(!A \land
  !B)$} يو فارمول دے۔
\olpseganchor{OLP-0058-B017}{end}

\olpseganchor{OLP-0058-B018}{start}
\item که \LR{$!A$} او \LR{$!B$} فارمولونه وي، نو \LR{$(!A \lor !B)$}
  يو فارمول دے۔
\olpseganchor{OLP-0058-B018}{end}

\olpseganchor{OLP-0058-B019}{start}
\item که \LR{$!A$} او \LR{$!B$} فارمولونه وي، نو \LR{$(!A \lif !B)$}
  يو فارمول دے۔
\olpseganchor{OLP-0058-B019}{end}

\olpseganchor{OLP-0058-B020}{start}
\item که \LR{$!A$} او \LR{$!B$} فارمولونه وي، نو \LR{$(!A \liff !B)$}
  يو فارمول دے۔
\olpseganchor{OLP-0058-B020}{end}

\olpseganchor{OLP-0058-B021}{start}
\item بل هېڅ شے فارمول نۀ دے۔
\end{enumerate}
\end{defn}
\olpseganchor{OLP-0058-B021}{end}

\olpseganchor{OLP-0058-B022}{start}
\begin{explain}
د فارمولونه تعريف يو \emph{استقرايي تعريف} دے۔ په اصل کښې موږ
د فارمولونه سټ په نامتناهي ډېرو پړاوونو کښې جوړوو۔ په لومړني
پړاو کښې ټول اټومي فارمولونه فارمولونه ګڼو؛ دا د تعريف له لومړنيو
څو حالتونو سره سمون لري، يعنې د
\LR{$\ltrue$}، %
\LR{$\lfalse$}، %
\LR{$\Obj p_i$} حالتونه۔ په دې توګه ``اټومي فارمول'' د همدې بڼې هر
فارمول ته وايي۔
\olpseganchor{OLP-0058-B022}{end}

\olpseganchor{OLP-0058-B023}{start}
د تعريف نور حالتونه له هغو فارمولونه څخه د نويو فارمولونه د
جوړولو قاعدې راکوي چې له مخکښې جوړ شوي وي۔ په دويم پړاو کښې له دغو
قاعدو څخه په کار اخيستو له اټومي فارمولونه څخه فارمولونه
جوړولے شو۔ په درېيم پړاو کښې له اټومي فارمولونو او په دويم پړاو کښې
له ترلاسه شويو فارمولونو څخه نوي فارمولونه جوړوو، او همداسې مخکښې
ځو۔ فارمول هر هغه شے دے چې بالاخره په داسې کوم پړاو کښې جوړ
شي، او بل هېڅ شے نۀ دے۔
\end{explain}
\olpseganchor{OLP-0058-B023}{end}

\olpseganchor{OLP-0058-B024}{start}
کله چې د دوه ځايه نښلوونکي~\LR{$\ast$} په کارولو له \LR{$!B$} او \LR{$!C$} څخه
جوړ شوے فارمول \LR{$(!B \ast !C)$} ليکو، ډېر ځله د قوسونو تر ټولو بهرنی
جوړ پرېږدو او يوازې~\LR{$!B \ast !C$} ليکو۔
\olpseganchor{OLP-0058-B024}{end}

\olpseganchor{OLP-0058-B034}{start}
\begin{defn}[نحوي يوشانوالی]
د \LR{$\ident$} نښه د نښو د توريزو لړيو ترمنځ نحوي يوشانوالی څرګندوي؛
يعنې \LR{$!A \ident !B$} هله او يوازې هله چې \LR{$!A$} او \LR{$!B$} د نښو د يو
شان اوږدوالي توريزې لړۍ وي او په هر ځاے کښې يوه نښه ولري۔
\end{defn}
\olpseganchor{OLP-0058-B034}{end}

\olpseganchor{OLP-0058-B035}{start}
د \LR{$\ident$} نښې دواړو خواوو ته د يو ځاے کولو له لارې ترلاسه شوې
توريزې لړۍ هم راتلے شي۔ مثلاً \LR{$!A \ident (!B \lor !C)$} دا معنٰی لري:
د نښو توريزه لړۍ~\LR{$!A$} هماغه توريزه لړۍ ده چې په همدې ترتيب د پرانيستي
قوس، د \LR{$!B$} توريزې لړۍ، د \LR{$\lor$} نښې، د~\LR{$!C$} توريزې لړۍ او د تړلي
قوس په يو ځاے کولو ترلاسه کېږي۔ که داسې وي، نو پوهېږو چې د \LR{$!A$}
لومړۍ نښه پرانيستے قوس دے، \LR{$!A$} د \LR{$!B$} توريزه لړۍ د فرعي توريزې
لړۍ په توګه لري چې له دويمې نښې پيلېږي، ورپسې \LR{$\lor$} راځي، او داسې نور۔
\olpseganchor{OLP-0058-B035}{end}

\olpseganchor{OLP-0059-B004}{start}
\olfileid{pl}{syn}{pre}
\olpseganchor{OLP-0059-B004}{end}

\olpseganchor{OLP-0059-B005}{start}
\olsection{لومړنۍ پايلې}
\olpseganchor{OLP-0059-B005}{end}

\olpseganchor{OLP-0059-B006}{start}
\begin{thm}[\emph{د فارمولونه دپاره د استقرا اصل}]
  \ollabel{thm:induction}  
  که کوم خاصيت~\LR{$P$} د ټولو اټومي فارمولونه دپاره صدق کوي او داسې
  وي چې
  \begin{enumerate}
    \item هر کله چې د~\LR{$!A$} دپاره صدق کوي، د \LR{$\lnot !A$}
      دپاره هم صدق کوي؛
    \item هر کله چې د \LR{$!A$} او~\LR{$!B$} دپاره صدق کوي، د
      \LR{$(!A \land !B)$} دپاره هم صدق کوي؛
    \item هر کله چې د \LR{$!A$} او~\LR{$!B$} دپاره صدق کوي، د
      \LR{$(!A \lor !B)$} دپاره هم صدق کوي؛
    \item هر کله چې د \LR{$!A$} او~\LR{$!B$} دپاره صدق کوي، د
      \LR{$(!A \lif !B)$} دپاره هم صدق کوي؛
    \item هر کله چې د \LR{$!A$} او~\LR{$!B$} دپاره صدق کوي، د
      \LR{$(!A \liff !B)$} دپاره هم صدق کوي؛
  \end{enumerate}
  نو \LR{$P$} د ټولو فارمولونه دپاره صدق کوي۔
\end{thm}
\olpseganchor{OLP-0059-B006}{end}

\olpseganchor{OLP-0059-B007}{start}
\begin{proof}
  \LR{$S$} دې د هغو ټولو فارمولونه مجموعه وي چې خاصيت~\LR{$P$} لري۔ په
  څرګنده \LR{$S \subseteq \Frm[L_0]$}۔ \LR{$S$} د
  \olref[fml]{defn:formulas} ټول شرطونه پوره کوي: ټول اټومي
  فارمولونه لري او د منطقي عملګرونه له مخې بند دے۔ \LR{$\Frm[L_0]$} د
  دې ډول تر ټولو وړوکے ټولګے دے، نو \LR{$\Frm[L_0] \subseteq S$}۔ له دې
  امله \LR{$\Frm[L_0] = S$} او هر فارمول خاصيت~\LR{$P$} لري۔
\end{proof}
\olpseganchor{OLP-0059-B007}{end}

\olpseganchor{OLP-0059-B008}{start}
\begin{prop}\ollabel{prop:balanced}
   په~\LR{$\Frm[L_0]$} کښې هر فارمول \emph{متوازن} دے؛ يعنې د کيڼو
   قوسونو شمېر ئې د ښيو قوسونو له شمېر سره برابر دے۔
\end{prop}
\olpseganchor{OLP-0059-B008}{end}

\olpseganchor{OLP-0059-B009}{start}
\begin{prob}
د \olref[pl][syn][pre]{prop:balanced} ثبوت ورکړئ۔
\end{prob}
\olpseganchor{OLP-0059-B009}{end}

\olpseganchor{OLP-0059-B010}{start}
\begin{prop} \ollabel{prop:noinit}
د فارمول هېڅ خاصه پيلنۍ برخه فارمول نۀ ده۔
\end{prop}
\olpseganchor{OLP-0059-B010}{end}

\olpseganchor{OLP-0059-B011}{start}
\begin{prob}
  د \olref[pl][syn][pre]{prop:noinit} ثبوت ورکړئ۔
\end{prob}
\olpseganchor{OLP-0059-B011}{end}

\olpseganchor{OLP-0059-B012}{start}
\begin{prop}[يکتا لوستونتيا]
په \LR{$ \Frm[L_0]$} کښې هر فارمول~\LR{$!A$} په لاندې بڼو کښې د يوې په
توګه پوره يوه تجزيه لري:
\begin{enumerate}
\item \LR{$\lfalse$}۔
\olpseganchor{OLP-0059-B012}{end}

\olpseganchor{OLP-0059-B013}{start}
\item \LR{$\ltrue$}۔
\olpseganchor{OLP-0059-B013}{end}

\olpseganchor{OLP-0059-B014}{start}
\item د کوم \LR{$\Obj p_n \in  \PVar$} دپاره \LR{$\Obj p_n$}۔
\olpseganchor{OLP-0059-B014}{end}
  
\olpseganchor{OLP-0059-B015}{start}
\item د کوم فارمول~\LR{$!B$} دپاره \LR{$\lnot !B$}۔
\olpseganchor{OLP-0059-B015}{end}

\olpseganchor{OLP-0059-B016}{start}
\item د کومو فارمولونه \LR{$!B$} او~\LR{$!C$} دپاره
  \LR{$(!B \land !C)$}۔
\olpseganchor{OLP-0059-B016}{end}

\olpseganchor{OLP-0059-B017}{start}
\item د کومو فارمولونه \LR{$!B$} او~\LR{$!C$} دپاره
  \LR{$(!B \lor !C)$}۔
\olpseganchor{OLP-0059-B017}{end}

\olpseganchor{OLP-0059-B018}{start}
\item د کومو فارمولونه \LR{$!B$} او~\LR{$!C$} دپاره
  \LR{$(!B \lif !C)$}۔
\olpseganchor{OLP-0059-B018}{end}

\olpseganchor{OLP-0059-B019}{start}
\item د کومو فارمولونه \LR{$!B$} او~\LR{$!C$} دپاره
  \LR{$(!B \liff !C)$}۔
\end{enumerate}
سربېره پر دې، دا تجزيه \emph{يکتا} ده۔
\end{prop}
\olpseganchor{OLP-0059-B019}{end}

\olpseganchor{OLP-0059-B020}{start}
\begin{proof}
پر \LR{$!A$} استقرا کوو۔ مثلاً فرض کړئ چې \LR{$!A$} د \LR{$(!B \lif !C)$} او
\LR{$(!B' \lif !C')$} په توګه دوه بېلې لوستنې لري۔ نو \LR{$!B$} او \LR{$!B'$}
بايد يو شان وي، ګنې يو به د بل خاصه پيلنۍ برخه وي؛ نو که د \LR{$!A$}
دواړه لوستنې بېلې وي، علت ئې خامخا دا دے چې \LR{$!C$} او \LR{$!C'$} د نښو د
هماغې يوې لړۍ بېلې لوستنې دي، چې د استقرايي فرض له مخې ناشونې ده۔
\end{proof}
\olpseganchor{OLP-0059-B020}{end}


\olpseganchor{OLP-0059-B021}{start}
\begin{defn}[يو شان ځايناستي]
که \LR{$!A$} او \LR{$!B$} فارمولونه وي او \LR{$\Obj p_i$} يو بياني متغير وي، نو \LR{$\Subst{!A}{!B}{\Obj p_i}$} هغه پايله ښيي چې په
\LR{$!A$} کښې د \LR{$\Obj p_i$} هر راتګ د \LR{$!B$} په يوه راتګ بدل شي؛ همدارنګه،
د \LR{$\Obj p_1$}، \dots،~\LR{$\Obj p_n$} پر ځاے په يو وخت د فارمولونه
\LR{$!B_1$}، \dots،~\LR{$!B_n$} ځايناستي په
\LR{$\SSubst{!A}{\subst{!B_1}{\Obj p_1},\dots,\subst{!B_n}{\Obj p_n}}$}
ښودل کېږي۔
\end{defn}
\olpseganchor{OLP-0059-B021}{end}

\olpseganchor{OLP-0059-B022}{start}
\begin{prob} د لاندې پنځو فارمولونه په هر يوه کښې وټاکئ چې آيا
 فارمول د \( \Subst{!A}{!B}{\Obj p_i} \) د ځايناستي په توګه
 څرګندېدے شي، چې \( !A \) په دې کښې يو وي: (الف) \( \Obj p_0 \)؛
 (ب) \( ( \lnot \Obj p_0 \land \Obj p_1) \)؛ او (ج)
 \( ( ( \lnot \Obj p_0 \lif \Obj p_1 ) \land \Obj p_2 ) \)۔ په هر
 حالت کښې اړونده ځايناستي په ګوته کړئ۔
  \begin{enumerate}
    \item \( \Obj p_1 \) 
    \item \( ( \lnot \Obj p_0 \land \Obj p_0 ) \)
    \item \( ( ( \Obj p_0 \lor \Obj p_1 ) \land \Obj p_2 )  \)
    \item \( \lnot ( ( \Obj p_0 \lif \Obj p_1 ) \land \Obj p_2 ) \)
    \item \( (( \lnot ( \Obj p_0 \lif \Obj p_1 ) \lif ( \Obj p_0 \lor \Obj p_1 )) \land \lnot ( \Obj p_0 \land \Obj p_1 )) \)
  \end{enumerate}
\end{prob}
\olpseganchor{OLP-0059-B022}{end}

\olpseganchor{OLP-0059-B023}{start}
\begin{prob}
د \LR{$\Subst{!A}{!B}{p}$} يو له رياضي پلوه کره تعريف د استقرا له مخې
ورکړئ۔
\end{prob}
\olpseganchor{OLP-0059-B023}{end}

\olpseganchor{OLP-0060-B004}{start}
\olfileid{pl}{syn}{fseq}
\olpseganchor{OLP-0060-B004}{end}

\olpseganchor{OLP-0060-B005}{start}
\olsection{جوړښتي لړۍ}
\olpseganchor{OLP-0060-B005}{end}

\olpseganchor{OLP-0060-B006}{start}
په استقرايي تعريف سره د فارمولونه تعريفول، او د استقرا له لارې
د فارمولونه د خاصيتونو د ثابتولو بشپړوونکې طريقه، ښکلې او اغېزمنه
لار ده۔ خو کله کله دا هم ګټوره وي چې د فارمولونه جوړول له بېخه
پورته، ګام په ګام وڅېړو۔ دلته د \emph{جوړښتي لړۍ} په مفهوم همدا کار
کوو۔
\olpseganchor{OLP-0060-B006}{end}

\olpseganchor{OLP-0060-B007}{start}
\begin{defn}[د فارمولونو جوړښتي لړۍ]
\ollabel{defn:fseq-frm}
د ژبې~\LR{$\Lang L_0$} د نښو د توريزو لړيو يوه متناهي لړۍ
\LR{$\tuple{!A_0,\dotsc,!A_n}$} د \LR{$!A$} دپاره \emph{جوړښتي لړۍ} ده که
\LR{$!A \ident !A_n$} وي او د هر \LR{$i \leq n$} دپاره يا \LR{$!A_i$} اټومي
فارمول وي، يا داسې \LR{$j,k < i$} موجود وي چې له لاندې حالتونو څخه يو
صدق وکړي:
\begin{enumerate}
    \item \LR{$!A_i \ident \lnot !A_j$}۔%
    \item \LR{$!A_i \ident (!A_j \land !A_k)$}۔%
    \item \LR{$!A_i \ident (!A_j \lor !A_k)$}۔%
    \item \LR{$!A_i \ident (!A_j \lif !A_k)$}۔%
    \item \LR{$!A_i \ident (!A_j \liff !A_k)$}۔%
\end{enumerate}
\end{defn}
\olpseganchor{OLP-0060-B007}{end}

\olpseganchor{OLP-0060-B008}{start}
\begin{ex}
\[
\tuple{
    \Obj p_0,
    \Obj p_1,
    (\Obj p_1 \land \Obj p_0),
    \lnot (\Obj p_1 \land \Obj p_0)
}
\]
د \LR{$\lnot (\Obj p_1 \land \Obj p_0)$} يوه جوړښتي لړۍ ده، او دا هم ده:
\[
\tuple{
    \Obj p_0,
    \Obj p_1,
    \Obj p_0,
    (\Obj p_1 \land \Obj p_0),
    (\Obj p_0 \lif \Obj p_1),
    \lnot (\Obj p_1 \land \Obj p_0)
}.
\]
%
لکه چې له دويمې بېلګې ښکاري، په جوړښتي لړيو کښې ``بې‌ګټې توکي''
هم راتلے شي: داسې فارمولونه چې زياتي وي يا په جوړونه کښې ونډه نۀ
لري۔
\end{ex}
\olpseganchor{OLP-0060-B008}{end}

\olpseganchor{OLP-0060-B009}{start}
\begin{prop}
\ollabel{prop:formed}
په~\LR{$\Frm[L_0]$} کښې هر فارمول~\LR{$!A$} يوه جوړښتي لړۍ لري۔
\end{prop}
\olpseganchor{OLP-0060-B009}{end}

\olpseganchor{OLP-0060-B010}{start}
\begin{proof}
فرض کړئ \LR{$!A$} اټومي دے۔ نو لړۍ~\LR{$\tuple{!A}$} د~\LR{$!A$} جوړښتي لړۍ ده۔
%
اوس فرض کړئ چې \LR{$!B$} او~\LR{$!C$} په ترتيب سره جوړښتي لړۍ
\LR{$\tuple{!B_0,\dotsc,!B_n}$} او \LR{$\tuple{!C_0,\dotsc,!C_m}$} لري۔
%
\begin{enumerate}
  \item که \LR{$!A \ident \lnot !B$} وي، نو
      \LR{$\tuple{!B_0,\dotsc,!B_n,\lnot !B_n}$}
      د~\LR{$!A$} جوړښتي لړۍ ده۔
  \item که \LR{$!A \ident (!B \land !C)$} وي، نو
      \LR{$\tuple{!B_0,\dotsc,!B_n,!C_0,\dotsc,!C_m,(!B_n \land !C_m)}$}
      د~\LR{$!A$} جوړښتي لړۍ ده۔
  \item که \LR{$!A \ident (!B \lor !C)$} وي، نو
      \LR{$\tuple{!B_0,\dotsc,!B_n,!C_0,\dotsc,!C_m,(!B_n \lor !C_m)}$}
      د~\LR{$!A$} جوړښتي لړۍ ده۔
  \item که \LR{$!A \ident (!B \lif !C)$} وي، نو
      \LR{$\tuple{!B_0,\dotsc,!B_n,!C_0,\dotsc,!C_m,(!B_n \lif !C_m)}$}
      د~\LR{$!A$} جوړښتي لړۍ ده۔
  \item که \LR{$!A \ident (!B \liff !C)$} وي، نو
      \LR{$\tuple{!B_0,\dotsc,!B_n,!C_0,\dotsc,!C_m,(!B_n \liff !C_m)}$}
      د~\LR{$!A$} جوړښتي لړۍ ده۔
\end{enumerate}
د فارمولونه دپاره د استقرا د اصل له مخې، هر فارمول يوه
جوړښتي لړۍ لري۔
\end{proof}
\olpseganchor{OLP-0060-B010}{end}

\olpseganchor{OLP-0060-B011}{start}
د دې عکس هم ثابتولے شو۔ دا ځکه مهم دے چې ښيي د فارمولونو د تعريف
زموږ دواړه لارې همارزښته دي: دواړه يو شان پايلې ورکوي۔ دا هم ترې
معلومېږي چې د جوړښتي لړيو پر اوږدوالي عادي استقرا کارولے شو او د
فارمولونو په اړه قضيې ثابتولے شو۔
\olpseganchor{OLP-0060-B011}{end}

\olpseganchor{OLP-0060-B012}{start}
\begin{lem}
\ollabel{lem:fseq-init}
فرض کړئ \LR{$\tuple{!A_0,\dotsc,!A_n}$} د~\LR{$!A_n$} جوړښتي لړۍ ده او
\LR{$k \leq n$} دے۔ نو \LR{$\tuple{!A_0,\dotsc,!A_k}$} د~\LR{$!A_k$} جوړښتي لړۍ
ده۔
\end{lem}
\olpseganchor{OLP-0060-B012}{end}

\olpseganchor{OLP-0060-B013}{start}
\begin{proof}
تمرين۔
\end{proof}
\olpseganchor{OLP-0060-B013}{end}

\olpseganchor{OLP-0060-B014}{start}
\begin{thm}
\ollabel{thm:fseq-frm-equiv}
\LR{$\Frm[L_0]$} د ژبې~\LR{$\Lang L_0$} د نښو د ټولو هغو توريزو لړيو سټ دے
چې يوه جوړښتي لړۍ لري۔
\end{thm}
\olpseganchor{OLP-0060-B014}{end}

\olpseganchor{OLP-0060-B015}{start}
\begin{proof}
\LR{$F$} دې د ژبې~\LR{$\Lang L_0$} د نښو د هغو ټولو توريزو لړيو سټ وي چې
جوړښتي لړۍ لري۔ په \olref[pl][syn][fseq]{prop:formed} کښې مو وليدل
چې \LR{$\Frm[L_0] \subseteq F$}؛ اوس ئې عکس ثابتوو۔
\olpseganchor{OLP-0060-B015}{end}

\olpseganchor{OLP-0060-B016}{start}
فرض کړئ \LR{$!A$} يوه جوړښتي لړۍ \LR{$\tuple{!A_0,\dotsc,!A_n}$} لري۔ پر~\LR{$n$}
د قوي استقرا له لارې ثابتوو چې \LR{$!A \in \Frm[L_0]$}۔ زموږ استقرايي
فرض دا دے چې د نښو هره هغه توريزه لړۍ چې د \LR{$m < n$} اوږدوالي
جوړښتي لړۍ لري، په~\LR{$\Frm[L_0]$} کښې ده۔ د جوړښتي لړۍ د تعريف له مخې،
يا \LR{$!A_n$} اټومي دے، يا بايد داسې \LR{$j,k < n$} موجود وي چې له لاندې
حالتونو څخه يو صدق وکړي:
\begin{enumerate}
    \item \LR{$!A_n \ident \lnot !A_j$}۔%
    \item \LR{$!A_n \ident (!A_j \land !A_k)$}۔%
    \item \LR{$!A_n \ident (!A_j \lor !A_k)$}۔%
    \item \LR{$!A_n \ident (!A_j \lif !A_k)$}۔%
    \item \LR{$!A_n \ident (!A_j \liff !A_k)$}۔%
\end{enumerate}
اوس حالت په حالت استدلال کوو۔ که \LR{$!A_n$} اټومي وي، نو
\LR{$!A_n \in \Frm[L_0]$}۔ بل حالت کښې فرض کړئ چې \LR{$!A \ident
(!A_j \land !A_k)$}۔ د \olref[pl][syn][fseq]{lem:fseq-init} له مخې،
\LR{$\tuple{!A_0,\dotsc,!A_j}$} او \LR{$\tuple{!A_0,\dotsc,!A_k}$} په ترتيب
سره د \LR{$!A_j$} او~\LR{$!A_k$} جوړښتي لړۍ دي۔ ځکه چې دا د~\LR{$!A$} د جوړښتي
لړۍ خاصې پيلنۍ فرعي لړۍ دي، د دواړو اوږدوالے له~\LR{$n$} څخه لږ دے۔ نو
د استقرايي فرض له مخې \LR{$!A_j$} او~\LR{$!A_k$} په~\LR{$\Frm[L_0]$} کښې دي؛ او
بيا د فارمول د تعريف له مخې \LR{$(!A_j \land !A_k)$} هم پکې دے۔
نور حالتونه په ورته استدلال ثابتېږي۔
\textbf{د ژباړې د نسخې سمون (OLPL-002):} سرچينه په دې يوه کرښه کښې
د نحوي يوشانوالي د نښې پر ځاے د معنايي هم ارزښت نښه کاروي، حال دا
چې تعريف، ټول نور حالتونه او ثبوت د توريزو لړيو يوشانوالی غواړي۔
دلته د هماغه تعريف نښه بېرته کارول شوې ده۔
\end{proof}
\olpseganchor{OLP-0060-B016}{end}

\olpseganchor{OLP-0061-B004}{start}
\olfileid{pl}{syn}{val}
\olpseganchor{OLP-0061-B004}{end}

\olpseganchor{OLP-0061-B005}{start}
\olsection{ارزښت ټاکنې او صدق}
\olpseganchor{OLP-0061-B005}{end}

\olpseganchor{OLP-0061-B006}{start}
\begin{defn}[ارزښت ټاکنې]
\LR{$\{\True, \False\}$} دې د رښتياوالي د دوو قيمتونو، ``رښتيا'' او
``دروغ''، سټ وي۔ د \LR{$\Lang{L_0}$} دپاره \emph{ارزښت ټاکنه} يوه
تابع~\LR{$\pAssign{v}$} ده چې د ژبې بياني متغيرونه ته يا
\LR{$\True$} ورکوي يا \LR{$\False$}؛ يعنې
\LR{$\pAssign{v} \colon \PVar \to \{\True, \False \}$}۔
\end{defn}
\olpseganchor{OLP-0061-B006}{end}

\olpseganchor{OLP-0061-B007}{start}
\begin{defn}
  چې ارزښت ټاکنه~\LR{$\pAssign{v}$} راکړل شوے وي، د ارزونې تابع
  \LR{$\pValue{v} \colon \Frm[L_0] \to \{\True, \False \}$} په استقرايي
  ډول داسې تعريف کړئ:
  \begin{align*}
    \pValue{v}(\lfalse) & = \False; \\
    \pValue{v}(\ltrue) & = \True; \\
    \pValue{v}(\Obj p_n) & = \pAssign{v}(\Obj p_n); \\
    \pValue{v}(\lnot !A) & = \begin{cases}
        \True & \text{\RL{که }} \pValue{v}(!A) = \False;\\
        \False & \text{\RL{که داسې نۀ وي۔}}
      \end{cases} \\ 
    \pValue{v}(!A \land !B) & = \begin{cases}
        \True &
        \text{\RL{که \LR{$\pValue{v}(!A) = \True$} او \LR{$\pValue{v}(!B) = \True$} وي؛}}\\
        \False &
        \text{\RL{که \LR{$\pValue{v}(!A) = \False$} يا \LR{$\pValue{v}(!B) = \False$} وي}}.
      \end{cases}\\
    \pValue{v}(!A \lor !B) & = \begin{cases}
        \True &
        \text{\RL{که \LR{$\pValue{v}(!A) = \True$} يا \LR{$\pValue{v}(!B) = \True$} وي؛}}\\
        \False &
        \text{\RL{که \LR{$\pValue{v}(!A) = \False$} او \LR{$\pValue{v}(!B) = \False$} وي}}.
      \end{cases}\\
    \pValue{v}(!A \lif !B) & = \begin{cases}
        \True &
        \text{\RL{که \LR{$\pValue{v}(!A) = \False$} يا \LR{$\pValue{v}(!B) = \True$} وي؛}}\\
        \False &
        \text{\RL{که \LR{$\pValue{v}(!A) = \True$} او \LR{$\pValue{v}(!B) = \False$} وي}}.
      \end{cases}\\
    \pValue{v}(!A \liff !B) & = \begin{cases}
        \True &
        \text{\RL{که \LR{$\pValue{v}(!A) = \pValue{v}(!B)$} وي؛}}\\
        \False &
        \text{\RL{که \LR{$\pValue{v}(!A) \neq \pValue{v}(!B)$} وي}}.
      \end{cases}
\end{align*}
\end{defn}
\olpseganchor{OLP-0061-B007}{end}

\olpseganchor{OLP-0061-B008}{start}
\begin{explain}
دا مادې د رښتياوالي له لاندې جدولونو سره سمون لري:
\begin{center}

  \begin{tabular}{|c||c|} \hline
    \LR{$!A$} & \LR{$ \lnot !A$} \\
    \hline \hline
    \LR{$\True$} & \LR{$\False$} \\
    \LR{$\False$} & \LR{$\True$} \\
    \hline
  \end{tabular}


  \begin{tabular}{|cc||c|} \hline
    \LR{$!A$} & \LR{$!B$} & \LR{$!A \land !B$} \\
    \hline \hline
    \LR{$\True$} & \LR{$\True$} & \LR{$\True$} \\
    \LR{$\True$} & \LR{$\False$} & \LR{$\False$} \\
    \LR{$\False$} & \LR{$\True$} & \LR{$\False$} \\
    \LR{$\False$} & \LR{$\False$} & \LR{$\False$} \\
    \hline
  \end{tabular}


  \begin{tabular}{|cc||c|} \hline
    \LR{$!A$} & \LR{$!B$} & \LR{$!A \lor !B$} \\
    \hline \hline
    \LR{$\True$} & \LR{$\True$} & \LR{$\True$} \\
    \LR{$\True$} & \LR{$\False$} & \LR{$\True$} \\
    \LR{$\False$} & \LR{$\True$} & \LR{$\True$} \\
    \LR{$\False$} & \LR{$\False$} & \LR{$\False$} \\
    \hline
  \end{tabular}
%
\\[1em]

  \begin{tabular}{|cc||c|} \hline
    \LR{$!A$} & \LR{$!B$} & \LR{$!A \lif !B$} \\
    \hline \hline
    \LR{$\True$} & \LR{$\True$} & \LR{$\True$} \\
    \LR{$\True$} & \LR{$\False$} & \LR{$\False$} \\
    \LR{$\False$} & \LR{$\True$} & \LR{$\True$} \\
    \LR{$\False$} & \LR{$\False$} & \LR{$\True$} \\
    \hline
  \end{tabular}


  \begin{tabular}{|cc||c|} \hline
    \LR{$!A$} & \LR{$!B$} & \LR{$!A \liff !B$} \\
    \hline \hline
    \LR{$\True$} & \LR{$\True$} & \LR{$\True$} \\
    \LR{$\True$} & \LR{$\False$} & \LR{$\False$} \\
    \LR{$\False$} & \LR{$\True$} & \LR{$\False$} \\
    \LR{$\False$} & \LR{$\False$} & \LR{$\True$} \\
    \hline
  \end{tabular}

\end{center}
\end{explain}
\olpseganchor{OLP-0061-B008}{end}

\olpseganchor{OLP-0061-B009}{start}
\begin{prob}
فرض کړئ چې \LR{$\Lang{L_0}$} ته يو درې ځايه نښلوونکے \LR{$\diamondsuit$}
ورزياتوو چې ارزونه ئې داسې ورکول شوې ده:
\begin{gather*}
  \pValue{v}(\diamondsuit( !A , !B , !C ) ) = \begin{cases}
    \pValue{v}( !B ) &
    \text{\RL{که \LR{$\pValue{v}(!A) = \True$} وي؛}}\\
    \pValue{v}( !C ) &
    \text{\RL{که \LR{$\pValue{v}(!A) = \False $} وي}}۔
  \end{cases}
\end{gather*}
د دې نښلوونکي د رښتياوالي جدول وليکئ۔
\end{prob}
\olpseganchor{OLP-0061-B009}{end}

\olpseganchor{OLP-0061-B010}{start}
\begin{thm}[ځايي ټاکنه]
\ollabel{thm:LocalDetermination} فرض کړئ \LR{$\pAssign{v_1}$} او
\LR{$\pAssign{v_2}$} داسې ارزښت ټاکنې دي چې په \LR{$!A$} کښې په راتلونکو
بياني متغيرونه يو له بل سره موافق دي؛ يعنې هر کله چې
\LR{$\Obj p_n$} په فارمول~\LR{$!A$} کښې راشي،
\LR{$\pAssign{v_1}(\Obj p_n) = \pAssign{v_2}(\Obj p_n)$} وي۔ نو
\LR{$\pValue{v_1}$} او \LR{$\pValue{v_2}$} هم پر~\LR{$!A$} موافق دي؛ يعنې
\LR{$\pValue{v_1}(!A) = \pValue{v_2}(!A)$}۔
\end{thm}
\textbf{د ژباړې د نسخې سمون (OLPL-003):} سرچينه د فرض په دويمه
وينا کښې ``په کوم فارمول کښې'' ليکي، خو قضيه هماغه له مخکښې ټاکلے
شوے فارمول ارزوي۔ ژباړه د لومړۍ وينا سره سم د هماغه فارمول متغيرونه
اخلي؛ په دې سره د قضیې مقصود شرط څرګند پاتې کېږي۔
\olpseganchor{OLP-0061-B010}{end}

\olpseganchor{OLP-0061-B011}{start}
\begin{proof}
پر \LR{$!A$} استقرا کوو۔
\end{proof}
\olpseganchor{OLP-0061-B011}{end}

\olpseganchor{OLP-0061-B012}{start}
\begin{defn}[صدق]
\ollabel{defn:satisfaction} د \emph{فارمول~\LR{$!A$} صدق په
  ارزښت ټاکنه~\LR{$\pAssign{v}$} کښې}، يعنې \LR{$\pSat{v}{!A}$}، په
  استقرايي ډول داسې تعريفولے شو۔ (د ``\LR{$\pSat{v}{!A}$} نۀ دے'' د
  څرګندولو دپاره \LR{$\pSat/{v}{!A}$} ليکو۔)
\begin{enumerate}
\item %
  \indcase{!A}{\lfalse}{\LR{$\pSat/{v}{\indfrm}$}۔}
\olpseganchor{OLP-0061-B012}{end}

\olpseganchor{OLP-0061-B013}{start}
\item %
  \indcase{!A}{\ltrue}{\LR{$\pSat{v}{\indfrm}$}۔}
\olpseganchor{OLP-0061-B013}{end}

\olpseganchor{OLP-0061-B014}{start}
\item \indcase{!A}{\Obj p_i}{\LR{$\pSat{v}{\indfrm}$}
  هله او يوازې هله چې \LR{$\pAssign{v}(\Obj p_i) = \True$} وي۔}
\olpseganchor{OLP-0061-B014}{end}

\olpseganchor{OLP-0061-B015}{start}
\item %
  \indcase{!A}{\lnot !B}{\LR{$\pSat{v}{\indfrm}$} هله او يوازې هله چې
    \LR{$\pSat/{v}{!B}$} وي۔}
\olpseganchor{OLP-0061-B015}{end}

\olpseganchor{OLP-0061-B016}{start}
\item %
  \indcase{!A}{(!B \land !C)}{\LR{$\pSat{v}{\indfrm}$} هله او يوازې هله
    چې \LR{$\pSat{v}{!B}$} او \LR{$\pSat{v}{!C}$} دواړه وي۔}
\olpseganchor{OLP-0061-B016}{end}

\olpseganchor{OLP-0061-B017}{start}
\item %
  \indcase{!A}{(!B \lor !C)}{\LR{$\pSat{v}{\indfrm}$} هله او يوازې هله
    چې \LR{$\pSat{v}{!B}$} يا \LR{$\pSat{v}{!C}$}، يا دواړه، وي۔}
\olpseganchor{OLP-0061-B017}{end}

\olpseganchor{OLP-0061-B018}{start}
\item %
  \indcase{!A}{(!B \lif !C)}{\LR{$\pSat{v}{\indfrm}$} هله او يوازې هله
    چې \LR{$\pSat/{v}{!B}$} يا \LR{$\pSat{v}{!C}$}، يا دواړه، وي۔}
\olpseganchor{OLP-0061-B018}{end}

\olpseganchor{OLP-0061-B019}{start}
\item %
  \indcase{!A}{(!B \liff !C)}{\LR{$\pSat{v}{\indfrm}$} هله او يوازې هله
    چې يا \LR{$\pSat{v}{!B}$} او \LR{$\pSat{v}{!C}$} دواړه وي، يا نۀ
    \LR{$\pSat{v}{!B}$} وي او نۀ \LR{$\pSat{v}{!C}$}۔}
\end{enumerate}
که \LR{$\Gamma$} د فارمولونه يو سټ وي، نو \LR{$\pSat{v}{\Gamma}$} هله او
يوازې هله چې د هر~\LR{$!A \in \Gamma$} دپاره \LR{$\pSat{v}{!A}$} وي۔
\end{defn}
\olpseganchor{OLP-0061-B019}{end}

\olpseganchor{OLP-0061-B020}{start}
\begin{prop}\ollabel{prop:sat-value}
  \LR{$\pSat{v}{!A}$} هله او يوازې هله چې \LR{$\pValue{v}(!A) = \True$} وي۔
\end{prop}
\olpseganchor{OLP-0061-B020}{end}

\olpseganchor{OLP-0061-B021}{start}
\begin{proof}
  پر~\LR{$!A$} استقرا کوو۔
\end{proof}
\olpseganchor{OLP-0061-B021}{end}

\olpseganchor{OLP-0061-B022}{start}
\begin{prob}
  د \olref[pl][syn][val]{prop:sat-value} ثبوت ورکړئ۔
\end{prob}
\olpseganchor{OLP-0061-B022}{end}

\olpseganchor{OLP-0062-B004}{start}
\olfileid{pl}{syn}{sem}
\olpseganchor{OLP-0062-B004}{end}

\olpseganchor{OLP-0062-B005}{start}
\olsection{معنايي مفهومونه}
\olpseganchor{OLP-0062-B005}{end}

\olpseganchor{OLP-0062-B006}{start}
لاندې معنايي مفهومونه تعريفوو:
\olpseganchor{OLP-0062-B006}{end}

\olpseganchor{OLP-0062-B007}{start}
\begin{defn} 
\begin{enumerate}
\item فارمول~\LR{$!A$} هله \emph{د صدق وړ} دے چې د کوم
  \LR{$\pAssign{v}$} دپاره \LR{$\pSat{v}{!A}$} وي؛ او هله \emph{د صدق نۀ وړ}
  دے چې د هېڅ \LR{$\pAssign{v}$} دپاره \LR{$\pSat{v}{!A}$} نۀ وي؛
\item فارمول~\LR{$!A$} هله \emph{تاوتولوژي} ده چې د ټولو
  ارزښت ټاکنې~\LR{$\pAssign{v}$} دپاره \LR{$\pSat{v}{!A}$} وي؛
\item فارمول~\LR{$!A$} هله \emph{اتفاقي} دے چې د صدق وړ وي خو
  تاوتولوژي نۀ وي؛
\item که \LR{$\Gamma$} د فارمولونه يو سټ وي، نو \LR{$\Gamma \Entails !A$}
  (``\LR{$\Gamma$}، \LR{$!A$} ته استلزام ورکوي'') هله او يوازې هله چې د هر
  هغه ارزښت ټاکنه~\LR{$\pAssign{v}$} دپاره \LR{$\pSat{v}{!A}$} وي چې
  \LR{$\pSat{v}{\Gamma}$} وي۔
\item که \LR{$\Gamma$} د فارمولونه يو سټ وي، نو \LR{$\Gamma$} هله
  \emph{د صدق وړ} دے چې داسې ارزښت ټاکنه~\LR{$\pAssign{v}$} موجود وي
  چې \LR{$\pSat{v}{\Gamma}$} وي؛ او که داسې نۀ وي، \LR{$\Gamma$}
  \emph{د صدق نۀ وړ} دے۔
\end{enumerate} 
\end{defn}
\olpseganchor{OLP-0062-B007}{end}

\olpseganchor{OLP-0062-B008}{start}
\begin{prob} 
 د لاندې څلورو فارمولونه په هر يوه کښې وټاکئ چې آيا هغه
 (الف)~د صدق وړ، (ب)~تاوتولوژي او (ج)~اتفاقي دے۔
\begin{enumerate}
  \item \( ( \Obj p_0 \lif ( \lnot \Obj p_1 \lif \lnot \Obj p_0 ) ) \).
  \item \( ( ( \Obj p_0 \land \lnot \Obj p_1 ) \lif ( \lnot \Obj p_0 \land \Obj p_2 )) \liff ( ( \Obj p_2 \lif \Obj p_0 ) \lif ( \Obj p_0 \lif \Obj p_1 )) \).
  \item \( ( \Obj p_0 \liff \Obj p_1 ) \lif ( \Obj p_2 \liff \lnot \Obj p_1 ) \).
  \item \( (( \Obj p_0 \liff ( \lnot \Obj p_1 \land \Obj p_2 )) \lor ( \Obj p_2 \lif ( \Obj p_0 \liff \Obj p_1 ))) \).
\end{enumerate}
\end{prob}
\olpseganchor{OLP-0062-B008}{end}

\olpseganchor{OLP-0062-B009}{start}
\begin{prop}
\ollabel{prop:semanticalfacts} 
\begin{enumerate} 
\item \LR{$!A$} هله او يوازې هله تاوتولوژي ده چې
  \LR{$\emptyset \Entails !A$} وي؛ 
\item که \LR{$\Gamma \Entails !A$} او \LR{$\Gamma \Entails !A \lif !B$} وي،
  نو \LR{$\Gamma \Entails !B$}؛
\item که \LR{$\Gamma$} د صدق وړ وي، نو د \LR{$\Gamma$} هر متناهي فرعي سټ هم
  د صدق وړ دے؛ 
\item \ollabel{def:monotonicity} يو اړخيز زياتوالے: که
  \LR{$\Gamma \subseteq \Delta$} او \LR{$\Gamma \Entails !A$} وي، نو
  \LR{$\Delta \Entails !A$} هم؛
\item \ollabel{def:Cut} تعدي: که \LR{$\Gamma \Entails !A$} او
  \LR{$\Delta \cup \{ !A\} \Entails !B$} وي، نو
  \LR{$\Gamma \cup \Delta \Entails !B$}۔
\end{enumerate}
\end{prop}
\olpseganchor{OLP-0062-B009}{end}

\olpseganchor{OLP-0062-B010}{start}
\begin{proof}
تمرين۔
\end{proof}
\olpseganchor{OLP-0062-B010}{end}

\olpseganchor{OLP-0062-B011}{start}
\begin{prob}
د \olref[pl][syn][sem]{prop:semanticalfacts} ثبوت ورکړئ۔
\end{prob}
\olpseganchor{OLP-0062-B011}{end}

\olpseganchor{OLP-0062-B012}{start}
\begin{prop}\ollabel{prop:entails-unsat}
  \LR{$\Gamma \Entails !A$} هله او يوازې هله چې
  \LR{$\Gamma \cup \{\lnot !A\}$} د صدق نۀ وړ وي۔
\end{prop}
\olpseganchor{OLP-0062-B012}{end}

\olpseganchor{OLP-0062-B013}{start}
\begin{proof}
تمرين۔
\end{proof}
\olpseganchor{OLP-0062-B013}{end}

\olpseganchor{OLP-0062-B014}{start}
\begin{prob}
د \olref[pl][syn][sem]{prop:entails-unsat} ثبوت ورکړئ۔
\end{prob}
\olpseganchor{OLP-0062-B014}{end}

\olpseganchor{OLP-0062-B015}{start}
\begin{thm}[د معنايي استنتاج قضيه]
  \ollabel{thm:sem-deduction} \LR{$\Gamma \Entails !A \lif !B$} هله او
  يوازې هله چې \LR{$\Gamma \cup \{!A\} \Entails !B$} وي۔
\end{thm}
\olpseganchor{OLP-0062-B015}{end}

\olpseganchor{OLP-0062-B016}{start}
\begin{proof}
تمرين۔
\end{proof}
\olpseganchor{OLP-0062-B016}{end}

\olpseganchor{OLP-0062-B017}{start}
\begin{prob}
د \olref[pl][syn][sem]{thm:sem-deduction} ثبوت ورکړئ۔
\end{prob}
\olpseganchor{OLP-0062-B017}{end}

\olpseganchor{OLP-0063-B004}{start}
\olchapter{fol}{prf}{د اشتقاق نظامونه}
\olpseganchor{OLP-0063-B004}{end}

\olpseganchor{OLP-0063-B005}{start}
\begin{editorial}
دا څپرکے د اشتقاق د نظامونو عمومي مواد راغونډوي۔ کوم درسي
کتاب چې يو ځانګړے نظام کاروي، د هغه نظام د معرفي کوونکي څپرکي په
پيل کښې د پېژندنې برخه او د اړوندې ټولکتنې برخه ورګډولے شي۔
\end{editorial}
\olpseganchor{OLP-0063-B005}{end}

\olpseganchor{OLP-0064-B004}{start}
\olfileid{fol}{prf}{int}
\olpseganchor{OLP-0064-B004}{end}

\olpseganchor{OLP-0064-B005}{start}
\olsection{پېژندنه}
\olpseganchor{OLP-0064-B005}{end}

\olpseganchor{OLP-0064-B006}{start}
منطقونه عموماً معنٰی پوهنه او اشتقاق نظام دواړه لري۔ معنٰی
پوهنه د رښتياوالي، د صدق وړتيا، اعتبار او استلزام په څېر مفهومونو
سره کار لري۔ د اشتقاق د نظامونو موخه دا ده چې د استلزام او
اعتبار د ثابتولو يوه بشپړه نحوي طريقه ورکړي۔ دا نظامونه په دې معنٰی
بشپړ نحوي دي چې په داسې يوه نظام کښې اشتقاق يو متناهي نحوي
څيز وي، چې عموماً د تړلي فارمولونه يا فارمولونه يوه لړۍ يا بل
متناهي ترتيب وي۔ د اشتقاق ښۀ نظامونه دا خاصيت لري چې د
تړلي فارمولونه يا فارمولونه د هرې ورکړل شوې لړۍ يا ترتيب
``سموالے'' په ميخانيکي توګه کتل کېدے شي۔
\olpseganchor{OLP-0064-B006}{end}

\olpseganchor{OLP-0064-B007}{start}
د لومړۍ درجې منطق تر ټولو ساده، او له تاريخي پلوه لومړني،
اشتقاق نظامونه \emph{بديهي} وو۔ په داسې يوه نظام کښې د
فارمولونه يوه لړۍ هله اشتقاق ګڼل کېږي چې په هغې کښې هر
فارمول يا د ``بديهي اصولو'' په يوه ټاکلي سټ کښې وي، يا د ټاکلو
``استنتاجي قاعدو'' له يوې څخه په کار اخيستو په لړۍ کښې له مخکښو
فارمولونه څخه ترې راووځي؛ او په ميخانيکي توګه کتل کېدے شي چې آيا
فارمول بديهي اصل دے او آيا د استنتاج له کومې قاعدې سره سم له
نورو فارمولونه څخه راوتلے دے۔ د اشتقاق بديهي نظامونه
تشريح کول اسانه دي، او له مابعدمنطقي پلوه پرې کار کول هم اسانه دي،
خو په هغو کښې اشتقاقونه لوستل او پوهېدل ګران دي او جوړول ئې هم
ګران دي۔
\olpseganchor{OLP-0064-B007}{end}

\olpseganchor{OLP-0064-B008}{start}
د اشتقاق نور نظامونه د دې موخې دپاره جوړ شوي چې د
اشتقاقونه جوړول يا له بشپړېدو وروسته پر اشتقاقونه
پوهېدل اسانه کړي۔
بېلګې ئې طبيعي استنتاج، د رښتياوالي ونې چې د تابلو ثبوتونو په نامه
هم پېژندل کېږي، او د سېکوېنټ حساب دي۔ د اشتقاق ځينې نظامونه
په ځانګړې توګه د ميخانيکي کولو دپاره طرح شوي دي؛ مثلاً د حل طريقه
په سافټوير کښې په اسانۍ پلې کېږي، خو پر اشتقاقونه ئې پوهېدل
تقريباً ناشوني دي۔ د اشتقاق دا نور نظامونه زياتره
اشتقاقونه د لړيو پر ځاے د فارمولونه د ونو په توګه ښيي۔ په
دې سره په اسانۍ ليدل کېږي چې د اشتقاق کومې برخې پر کومو
نورو برخو تکيه لري۔
\olpseganchor{OLP-0064-B008}{end}

\olpseganchor{OLP-0064-B009}{start}
نو د يوه ورکړل شوي منطق، لکه د لومړۍ درجې منطق، د اشتقاق
بېل نظامونه به دا په بېلو ډولونو روښانه کړي چې تړلے فارمول کله
\emph{قضيه} ده او د نورو له ځينو څخه د تړلے فارمول د
د اشتقاق وړ کېدو معنٰی څه ده۔ دا که په بديهي اشتقاقونه،
طبيعي استنتاجونو، سېکوېنټي اشتقاقونه، د رښتياوالي په ونو يا د
حل په ابطالونو ترسره شي، غواړو چې دا اړيکې د اعتبار او استلزام له
معنايي مفهومونو سره سمون ولري۔ د ``\LR{$!A$}~قضيه ده'' دپاره
\LR{$\Proves !A$} او د ``\LR{$!A$} له~\LR{$\Gamma$} څخه د اشتقاق وړ دے'' دپاره
``\LR{$\Gamma \Proves !A$}'' ليکو۔ \LR{$\Proves$} چې هر ډول تعريف شي، غواړو
له \LR{$\Entails$} سره سمون ولري؛ يعنې:
\begin{enumerate}
\item \LR{$\Proves !A$} هله او يوازې هله چې \LR{$\Entails !A$}
\item \LR{$\Gamma \Proves !A$} هله او يوازې هله چې \LR{$\Gamma \Entails !A$}
\end{enumerate}
د پورته دواړو د ``يوازې هله'' لوري ته \emph{سموالے} ويل کېږي۔
اشتقاق نظام هله سم دے چې د اشتقاق وړتيا استلزام يا اعتبار
تضمين کړي۔ د اشتقاق هر د کار وړ نظام بايد سم وي؛ د
اشتقاق ناسم نظامونه هېڅ ګټه نۀ لري۔ په اصل کښې د
اشتقاق ټوله موخه دا ده چې د اعتبار يا استلزام نحوي تضمين
ورکړي۔ موږ به د وړاندې شويو اشتقاق نظامونو سموالے ثابت کړو۔
\olpseganchor{OLP-0064-B009}{end}

\olpseganchor{OLP-0064-B010}{start}
د دې عکس، يعنې ``هله'' لورے، هم مهم دے: دې ته \emph{بشپړتيا} ويل
کېږي۔ د اشتقاق بشپړ نظام دومره ځواکمن وي چې هر کله \LR{$!A$}
معتبر وي، وښيي چې \LR{$!A$} قضيه ده؛ او هر کله چې
\LR{$\Gamma \Entails !A$} وي، وښيي چې \LR{$\Gamma \Proves !A$}۔ بشپړتيا ثابتول
ګران دي، او ځينې منطقونه د اشتقاق هېڅ بشپړ نظام نۀ لري۔ د
لومړۍ درجې منطق ئې لري۔ کورت ګوډل لومړے کس و چې د لومړۍ درجې منطق د
اشتقاق د يوه نظام بشپړتيا ئې په خپله ۱۹۲۹ز کال څېړنيزه
رساله کښې ثابته کړه۔
\olpseganchor{OLP-0064-B010}{end}

\olpseganchor{OLP-0064-B011}{start}
له اشتقاق نظامونو سره يو بل تړلے مفهوم \emph{سازګاري} ده۔
د تړلي فارمولونه يو سټ هله ناسازګار بلل کېږي چې هر څه ترې
اشتقاق کېدے شي؛ او که داسې نۀ وي، سازګار دے۔ ناسازګاري د صدق
د نۀ وړتيا نحوي سياله ده: لکه د صدق نۀ وړ سټونو غوندې، د
تړلي فارمولونه ناسازګار سټونه ښې نظريې نۀ جوړوي او په بنسټيز ډول
نيمګړي وي۔ د تړلي فارمولونه سازګار سټونه کېدے شي رښتوني يا ګټور نۀ
وي، خو لږ تر لږه د منطقي ګټورتيا له دې تر ټولو ټيټ بريد څخه تېرېږي۔
د اشتقاق په بېلو نظامونو کښې د تړلي فارمولونه د سټونو د
سازګارۍ ځانګړے تعريف توپير کولے شي؛ خو لکه~\LR{$\Proves$}، غواړو
سازګاري د خپل معنايي سيال، د صدق وړتيا، سره سمه وي۔ غواړو تل داسې
وي چې \LR{$\Gamma$} هله او يوازې هله سازګار وي چې د صدق وړ وي۔ دلته د
``يوازې هله'' لورے له بشپړتيا سره برابر دے، ځکه سازګاري د صدق
وړتيا تضمينوي؛ او ``هله'' لورے له سموالي سره برابر دے، ځکه د صدق
وړتيا سازګاري تضمينوي۔ په حقيقت کښې د کلاسيکي لومړۍ درجې منطق
دپاره د سموالي او بشپړتيا دواړه بڼې همارزښته دي۔
\olpseganchor{OLP-0064-B011}{end}

\olpseganchor{OLP-0065-B004}{start}
\olfileid{fol}{prf}{seq}
\olpseganchor{OLP-0065-B004}{end}

\olpseganchor{OLP-0065-B005}{start}
\olsection{د سېکوېنټ حساب}
\olpseganchor{OLP-0065-B005}{end}

\olpseganchor{OLP-0065-B006}{start}
د اشتقاق ډېر نظامونه د تړلي فارمولونه له ترتيبونو سره کار
کوي، خو د سېکوېنټ حساب له \emph{سېکوېنټونو} سره کار کوي۔ سېکوېنټ
د دې بڼې يو عبارت دے:
\[
!A_1, \dots, !A_m \Sequent !B_1, \dots, !B_n,
\]
\textbf{د ژباړې د نسخې سمون (OLPF-001):} سرچينه د سېکوېنټ دواړو
خواوو ته يو شان وروستے اندېکس کاروي، چې دواړه لړۍ په ناپامۍ يو شان
اوږدوي؛ خو ورپسې څرګندوي چې هره لړۍ په خپلواکه توګه تشه کېدے شي۔
دلته د ښۍ لړۍ وروستے اندېکس بېل شوے، څو دواړه متناهي لړۍ خپلواکه
اوږدوالی ولري۔
يعنې دا د تړلي فارمولونه د لړيو يوه جوړه ده چې د سېکوېنټ په
نښه~\LR{$\Sequent$} بېله شوې ده۔ له دواړو لړيو هره يوه تشه کېدے شي۔ د
سېکوېنټ په حساب کښې اشتقاق د سېکوېنټونو يوه ونه ده۔ تر
ټولو پاسني سېکوېنټونه ځانګړې بڼه لري، او ``لومړني سېکوېنټونه'' يا
``بديهي اصول'' بلل کېږي؛ هر بل سېکوېنټ د استنتاج د کومې قاعدې له
لارې سم له پاسه سېکوېنټونو څخه راځي۔ د استنتاج قاعدې يا په
سېکوېنټونو کښې تړلي فارمولونه سمبالوي، يعنې کيڼې يا ښۍ خوا ته ئې
ورزياتوي، ترې لرې کوي يا ئې بيا ترتيبوي؛ او يا د قاعدې په پايله کښې
يو مرکب فارمول معرفي کوي۔ مثلاً د \LR{$\LeftR{\land}$} قاعده اجازه
راکوي چې له \LR{$!A, \Gamma \Sequent \Delta$} څخه \LR{$!A \land !B, \Gamma
\Sequent \Delta$} استنتاج کړو، او \LR{$\RightR{\lif}$} اجازه راکوي چې د
هر \LR{$\Gamma$}، \LR{$\Delta$}، \LR{$!A$} او~\LR{$!B$} دپاره له \LR{$!A, \Gamma \Sequent
\Delta, !B$} څخه \LR{$\Gamma \Sequent \Delta, !A \lif !B$} استنتاج کړو۔
(په تېره بيا، \LR{$\Gamma$} او~\LR{$\Delta$} تش کېدے شي۔)
\olpseganchor{OLP-0065-B006}{end}

\olpseganchor{OLP-0065-B007}{start}
د سېکوېنټ پر حساب ولاړه \LR{$\Proves$} اړيکه داسې تعريفېږي:
\LR{$\Gamma \Proves !A$} هله او يوازې هله چې داسې کومه لړۍ \LR{$\Gamma_0$}
موجوده وي چې په \LR{$\Gamma_0$} کښې هر \LR{$!A$} په~\LR{$\Gamma$} کښې وي، او داسې
اشتقاق موجود وي چې سېکوېنټ~\LR{$\Gamma_0 \Sequent !A$} ئې په
جرړه کښې وي۔ \LR{$!A$} هله د سېکوېنټ په حساب کښې قضيه ده چې سېکوېنټ
~\LR{$\Sequent !A$} کوم اشتقاق ولري۔ مثلاً دا اشتقاق
ښيي چې \LR{$\Proves (!A \land !B) \lif !A$}:
\begin{prooftree}
  \Axiom$!A \fCenter !A$
  \RightLabel{\LeftR{\land}}
  \UnaryInf$!A \land !B \fCenter !A$
  \RightLabel{\RightR{\lif}}
  \UnaryInf$\fCenter (!A \land !B) \lif !A$
\end{prooftree}
\olpseganchor{OLP-0065-B007}{end}

\olpseganchor{OLP-0065-B008}{start}
يو سټ \LR{$\Gamma$} د سېکوېنټ په حساب کښې هله ناسازګار دے چې د
\LR{$\Gamma_0 \Sequent$} کوم اشتقاق موجود وي، چې په کښې هر
\LR{$!A \in \Gamma_0$} په~\LR{$\Gamma$} کښې وي او د سېکوېنټ ښۍ خوا تشه وي۔
د \RightR{\Weakening} قاعدې په کارولو له ناسازګار سټ څخه هر
تړلے فارمول اشتقاق کېدے شي۔
\olpseganchor{OLP-0065-B008}{end}

\olpseganchor{OLP-0065-B009}{start}
د سېکوېنټ حساب په ۱۹۳۰مو کلونو کښې ګيرهارډ ګېنتڅن جوړ کړ۔ د منظمې
او متناظرې طرحې له امله دا د اشتقاقونه د نظريې د ودې دپاره
ډېر ګټور صوري نظام دے۔ د سېکوېنټ په حساب کښې د اشتقاقونه
موندل نسبتاً اسانه دي، خو دا اشتقاقونه ډېر ځله په سختۍ لوستل
کېږي او له ثبوتونو سره ئې تړاو کله کله په اسانۍ نۀ ښکاري۔ بيا هم،
دا د اشتقاق نظامونو ډېره ښکلې لار ثابته شوې، او ډېر منطقونه
د سېکوېنټ حساب نظامونه لري۔
\olpseganchor{OLP-0065-B009}{end}

\olpseganchor{OLP-0066-B004}{start}
\olfileid{fol}{prf}{ntd}
\olpseganchor{OLP-0066-B004}{end}

\olpseganchor{OLP-0066-B005}{start}
\olsection{طبيعي استنتاج}
\olpseganchor{OLP-0066-B005}{end}

\olpseganchor{OLP-0066-B006}{start}
طبيعي استنتاج د اشتقاق داسې نظام دے چې موخه ئې د رښتيني
استدلال، په تېره بيا د رياضي پوهانو د منظم استدلال، انځورول دي۔
رښتينے استدلال د يو شمېر ``طبيعي'' بڼو له مخې مخکښې ځي۔ مثلاً د
حالتونو ثبوت موږ ته اجازه راکوي چې د يوه فصلي فرض پر بنسټ پايله
ثابته کړو، په دې چې وښيو پايله د فصل له هرې برخې څخه راځي۔ نامستقيم
ثبوت اجازه راکوي چې پايله په دې سره ثابته کړو چې نفي ئې تناقض ته
رسېږي۔ شرطي ثبوت د ``که \dots نو \dots'' شرطي ادعا په دې سره ثابتوي
چې تالي له مقدم څخه راځي۔ طبيعي استنتاج د دغو طبيعي استنتاجونو د
ځينو صوري کول دي۔ هر منطقي نښلوونکے او هر سور دوه قاعدې لري، يوه د
معرفي کولو او بله د لرې کولو قاعده، او هره يوه ئې د داسې طبيعي
استنتاج له يوې بڼې سره سمه ده۔ مثلاً \LR{$\Intro{\lif}$} له شرطي ثبوت
سره او \LR{$\Elim{\lor}$} د حالتونو له ثبوت سره سمون لري۔ يوه په تېره
ساده قاعده \LR{$\Elim{\land}$} ده، چې اجازه راکوي له \LR{$!A \land !B$} څخه
~\LR{$!A$} يا \LR{$!B$} استنتاج کړو۔
\olpseganchor{OLP-0066-B006}{end}

\olpseganchor{OLP-0066-B007}{start}
يو خاصيت چې طبيعي استنتاج د اشتقاق له نورو نظامونو څخه
بېلوي، د فرضونو کارول دي۔ په طبيعي استنتاج کښې اشتقاق د
فارمولونه يوه ونه ده۔ يو فارمول د فارمولونه د ونې په جرړه
کښې وي، او د ونې ``پاڼې'' هغه فارمولونه وي چې پايله ترې راايستل
کېږي۔ په طبيعي استنتاج کښې ځينې پاڼه-فارمولونه د اشتقاق
دننه ونډه لري، خو تر هغې ``کارول شوي وي'' چې اشتقاق پايلې ته
رسېږي۔ دا په رښتيني استدلال کښې د هغو فرضيو د معرفي کولو له دود سره
سمون لري چې يوازې لږ وخت نافذې پاتې کېږي۔ مثلاً د حالتونو په ثبوت
کښې د فصل د هرې برخې رښتياوالے فرضوو؛ په شرطي ثبوت کښې د مقدم
رښتياوالے فرضوو؛ او په نامستقيم ثبوت کښې د پايلې د نفي رښتياوالے
فرضوو۔ د فرضي فرضونو د معرفي کولو او بيا د يوه منځني ګام د ثابتولو
دپاره د هغو ختمولو دا طريقه د طبيعي استنتاج څرګنده نښه ده۔ د طبيعي
استنتاج د يوه اشتقاق په پاڼو کښې فارمولونه فرضونه بلل کېږي،
او د استنتاج ځينې قاعدې کېدے شي هغه ``ايستل'' کړي۔ مثلاً که
مو له ځينو فرضونو څخه د~\LR{$!B$} کوم اشتقاق درلود چې~\LR{$!A$} هم
پکې شامل و، نو د \LR{$\Intro{\lif}$} قاعده اجازه راکوي چې~\LR{$!A \lif !B$}
استنتاج کړو او د~\LR{$!A$} بڼې هر فرض ختم کړو۔ (د دې ثبتولو
دپاره چې کوم فرضونه په کومو استنتاجونو کښې ختمېږي، استنتاج او هغه
فرضونه چې ختموي ئې په يوه شمېره نښه کوو۔) هغه فرضونه چې د
اشتقاق په پاے کښې نۀ ايستل شوې پاتې کېږي، په ګډه د پايلې
د رښتياوالي دپاره بس دي؛ نو اشتقاق ثابتوي چې د هغه
نۀ ايستل شوي فرضونه پايلې ته استلزام ورکوي۔
\olpseganchor{OLP-0066-B007}{end}

\olpseganchor{OLP-0066-B008}{start}
د طبيعي استنتاج پر بنسټ اړيکه \LR{$\Gamma \Proves !A$} هله صدق کوي چې
داسې اشتقاق موجود وي چې \LR{$!A$} د ونې وروستۍ تړلے فارمول وي،
او هره نۀ ايستل شوې پاڼه په~\LR{$\Gamma$} کښې وي۔ \LR{$!A$} په طبيعي
استنتاج کښې هله قضيه ده چې داسې اشتقاق موجود وي چې \LR{$!A$}
وروستۍ تړلے فارمول وي او ټول فرضونه ايستل شوي وي۔ مثلاً دا
اشتقاق ښيي چې \LR{$\Proves (!A \land !B) \lif !A$}:
\begin{prooftree}
  \AxiomC{\LR{$\Discharge{!A \land !B}{1}$}}
  \RightLabel{\Elim{\land}}
  \UnaryInfC{\LR{$!A$}}
  \DischargeRule{\Intro{\lif}}{1}
  \UnaryInfC{\LR{$(!A \land !B) \lif !A$}}
\end{prooftree}
نښه~\LR{$1$} څرګندوي چې فرض \LR{$!A \land !B$} د \Intro{\lif} په استنتاج کښې
ايستل کېږي۔
\olpseganchor{OLP-0066-B008}{end}
  
\olpseganchor{OLP-0066-B009}{start}
يو سټ~\LR{$\Gamma$} هله ناسازګار دے چې په طبيعي استنتاج کښې
\LR{$\Gamma \Proves \lfalse$} وي۔ د \FalseInt{} قاعده دا يقيني کوي چې له
ناسازګار سټ څخه هر تړلے فارمول اشتقاق کېدے شي۔
\olpseganchor{OLP-0066-B009}{end}

\olpseganchor{OLP-0066-B010}{start}
د طبيعي استنتاج نظامونه ګيرهارډ ګېنتڅن او ستانيسواف ياشکوفسکي په
۱۹۳۰مو کلونو کښې جوړ کړل، او وروسته ډاګ پراوېڅ او فريډرېک فېچ لا
پراخ کړل۔ ځکه چې استنتاجونه ئې د ثبوت طبيعي لارې انځوروي، فلسفيان
ئې خوښوي۔ د فېچ جوړې شوې بڼې زياتره د منطق په پېژندنيزو درسي
کتابونو کښې کارول کېږي۔ د منطق په فلسفه کښې کله کله د طبيعي استنتاج
قاعدې د منطقي عملګرو معنٰی ورکوي، چې دې ته ``ثبوت-نظري معنايي
پوهنه'' ويل کېږي۔
\olpseganchor{OLP-0066-B010}{end}

\olpseganchor{OLP-0067-B004}{start}
\olfileid{fol}{prf}{tab}
\olpseganchor{OLP-0067-B004}{end}

\olpseganchor{OLP-0067-B005}{start}
\olsection{تابلوګانې}
\olpseganchor{OLP-0067-B005}{end}

\olpseganchor{OLP-0067-B006}{start}
د اشتقاق ډېر نظامونه د تړلي فارمولونه له ترتيبونو سره کار
کوي، خو تابلوګانې له نښه لرونکي فارمولونه سره کار کوي۔
نښه لرونکے فارمول يوه جوړه ده چې د رښتياوالي له يوې نښې
(\LR{$\True$} يا \LR{$\False$}) او تړلے فارمول څخه جوړه ده:
\[
\sFmla{\True}{!A} \text{\RL{ يا }} \sFmla{\False}{!A}.
\]
تابلو له هغو نښه لرونکي فارمولونه څخه جوړ وي چې د لاندې خوا
ته څانګېدونکې ونه کښې ترتيب شوي وي۔ دا له يو شمېر \emph{فرضونو}
پيلېږي او په هغو نښه لرونکي فارمولونه دوام کوي چې د استنتاج د کومې
قاعدې په پلي کولو تر ځان له پاسه له يوه نښه لرونکي فارمولونه څخه
راوتلي وي۔ هره قاعده اجازه راکوي چې د يوې څانګې په پاے کښې يو يا
څو نښه لرونکي فارمولونه ورزيات کړو، يا دوه نښه لرونکي فارمولونه يو د
بل تر څنګ ورزيات کړو؛ په وروستي حالت کښې څانګه په دوو وېشل کېږي او
دوه ورزيات شوي نښه لرونکي فارمولونه د دواړو څانګو سرونه جوړوي۔
\olpseganchor{OLP-0067-B006}{end}

\olpseganchor{OLP-0067-B007}{start}
پر يوه مرکب نښه لرونکے فارمول د قاعدې له پلي کولو څخه داسې
نښه لرونکي فارمولونه ورزياتېږي چې سمدستي فرعي-فارمولونه وي۔ د هرې
قاعدې يوه جوړه راځي، د دواړو نښو هرې يوې دپاره يوه قاعده۔ مثلاً د
\LR{$\TRule{\True}{\land}$} قاعده پر \LR{$\sFmla{\True}{!A \land !B}$} پلې
کېږي او اجازه راکوي چې د هرې هغې څانګې په پاے کښې دواړه
نښه لرونکي فارمولونه \LR{$\sFmla{\True}{!A}$} او~\LR{$\sFmla{\True}{!B}$}
ورزيات کړو چې \LR{$\sFmla{\True}{!A \land !B}$} لري؛ او د
\LR{$\TRule{\False}{\land}$} قاعده اجازه راکوي چې يوه څانګه د
\LR{$\sFmla{\False}{!A}$} او \LR{$\sFmla{\False}{!B}$} يو د بل تر څنګ په
ورزياتولو ووېشو۔
\textbf{د ژباړې د نسخې سمون (OLPF-002):} په سرچينه کښې د دروغې
عطف قاعدې دويم ارګومېنټ ټول مرکب فارمول دے، خو د رښتوني عطف موازي
قاعده او د لاندې بېلګې ټولې قاعدې هلته يوازې نښلوونکے اخلي۔ دلته
هماغه د عطف نښلوونکے کارول شوے، څو د قاعدې نښه په يو شان بڼه جوړه
شي۔
تابلو هله تړلے دے چې د هغه هره څانګه د نښه لرونکي فارمولونه
سره برابره جوړه \LR{$\sFmla{\True}{!A}$} او \LR{$\sFmla{\False}{!A}$} ولري۔
\olpseganchor{OLP-0067-B007}{end}

\olpseganchor{OLP-0067-B008}{start}
پر تابلوګانې ولاړه \LR{$\Proves$} اړيکه داسې تعريفېږي:
\LR{$\Gamma \Proves !A$} هله او يوازې هله چې داسې کوم متناهي سټ
~\LR{$\Gamma_0 = \{!B_1, \dots, !B_n\} \subseteq \Gamma$} موجود وي چې د
لاندې فرضونو دپاره يو تړلے تابلو موجود وي:
\[
\{\sFmla{\False}{!A}, \sFmla{\True}{!B_1}, \dots, \sFmla{\True}{!B_n}\} 
\]
مثلاً دا تړلے تابلو ښيي چې \LR{$\Proves (!A \land !B) \lif !A$}:
\begin{oltableau}
  [\sFmla{\False}{(\formula{A} \land \formula{B}) \lif \formula{A}}, just = \TAss
    [\sFmla{\True}{\formula{A} \land \formula{B}}, just = {\TRule{\False}{\lif}[1]}
      [\sFmla{\False}{\formula{A}}, just = {\TRule{\False}{\lif}[1]}
        [\sFmla{\True}{\formula{A}}, just = {\TRule{\True}{\lif}[2]}
          [\sFmla{\True}{\formula{B}}, just = {\TRule{\True}{\lif}[2]}, close]
        ]
      ]
    ]
  ]
\end{oltableau}
\olpseganchor{OLP-0067-B008}{end}

\olpseganchor{OLP-0067-B009}{start}
يو سټ \LR{$\Gamma$} د تابلو په حساب کښې هله ناسازګار دے چې د
لاندې فرضونو دپاره يو تړلے تابلو موجود وي:
\[
\{\sFmla{\True}{!B_1}, \dots, \sFmla{\True}{!B_n}\} 
\]
د کوم \LR{$!B_i \in \Gamma$} دپاره۔
\olpseganchor{OLP-0067-B009}{end}

\olpseganchor{OLP-0067-B010}{start}
تابلوګانې په ۱۹۵۰مو کلونو کښې اېورټ بېت او ياکو هېنتيکا هر يوه
په خپلواکه توګه جوړ کړل، او رېمونډ سموليان ساده او عام کړل۔ دا کارول
ډېر اسانه دي، ځکه چې د تابلو جوړول ډېره منظمه طريقه ده۔ د
تابلوګانې د منظم طبيعت له امله په کمپيوټر کښې هم په اسانۍ پلي
کېږي۔ خو تابلو زياتره په سختۍ لوستل کېږي او له ثبوتونو سره ئې
تړاو کله کله په اسانۍ نۀ ښکاري۔ دا تګلاره هم ډېره عامه ده، او ډېر
بېل منطقونه د تابلو نظامونه لري۔ تابلوګانې له موږ سره د
هغو جوړښتونه په موندلو کښې هم مرسته کوي چې ورکړل شوي
تړلي فارمولونه، يا د هغو سټونه، صادقوي: که سټ د صدق وړ وي، تړلے
تابلو به و نۀ لري؛ يعنې هر تابلو به يوه پرانيستې څانګه
ولري۔ په دې شرط چې پر هغې څانګه هره ممکنه قاعده پلې شوې وي،
صادقوونکے جوړښت له يوې پرانيستې څانګې څخه ``لوستل کېدے''
شي۔ له سېکوېنټ حساب سره هم ډېر نږدې تړاو شته: په اصل کښې يو تړلے
تابلو د سېکوېنټ په حساب کښې يو لنډ شوے اشتقاق دے چې
سرچپه ليکل شوے وي۔
\olpseganchor{OLP-0067-B010}{end}

\olpseganchor{OLP-0068-B004}{start}
\olfileid{fol}{prf}{axd}
\olpseganchor{OLP-0068-B004}{end}

\olpseganchor{OLP-0068-B005}{start}
\olsection{بديهي اشتقاقونه}
\olpseganchor{OLP-0068-B005}{end}

\olpseganchor{OLP-0068-B006}{start}
بديهي اشتقاقونه د اشتقاق تر ټولو زاړۀ او ساده منطقي
نظامونه دي۔ د هغو اشتقاقونه يوازې د تړلي فارمولونه لړۍ دي۔ د
تړلي فارمولونه يوه لړۍ هله سم اشتقاق ګڼل کېږي چې په هغې کښې
هره تړلے فارمول~\LR{$!A$} له لاندې شرطونو څخه يو پوره کړي:
\begin{enumerate}
\item \LR{$!A$} بديهي اصل وي؛ يا
\item \LR{$!A$} د تړلي فارمولونه د ورکړل شوي سټ~\LR{$\Gamma$} يو غړے
  وي؛ يا
\item \LR{$!A$} د استنتاج په کومه قاعده توجيه شوے وي۔
\end{enumerate}
د بديهي اصل کېدو دپاره \LR{$!A$} بايد د يو شمېر ټاکلو تړلے فارمول
شېماوو له يوې سره بڼه ولري۔ د بديهي اصلونو د شېماوو ډېر داسې سټونه
شته چې د لومړۍ درجې منطق دپاره د اشتقاق د قناعت وړ، سم او
بشپړ نظام ورکوي۔ ځينې د هغو نښلوونکو له مخې ترتيب شوي وي چې واک پرې
چلوي؛ مثلاً دا شېماوې
\[
!A \lif (!B \lif !A) \qquad !B \lif (!B \lor !C) \qquad (!B \land !C) \lif !B
\]
هغه عام بديهي اصول دي چې پر \LR{$\lif$}، \LR{$\lor$} او~\LR{$\land$} واک چلوي۔ د
بديهي اصولو ځينې نظامونه غواړي د اصولو شمېر تر ټولو لږ وي۔ دا چې
کوم نښلوونکي بنسټيز وټاکل شي، ان داسې نظامونه هم موندل کېدے شي چې
يوازې له يوه بديهي اصل څخه جوړ وي۔
\olpseganchor{OLP-0068-B006}{end}

\olpseganchor{OLP-0068-B007}{start}
د استنتاج قاعده يوه شرطي وينا ده چې په اشتقاق کښې د
تړلے فارمول د توجيه کېدو دپاره بس شرط ورکوي۔ مودوس پوننس يوه ډېره
عامه داسې قاعده ده: وايي که \LR{$!A$} او \LR{$!A \lif !B$} له مخکښې توجيه
شوي وي، نو \LR{$!B$} توجيه شوے دے۔ دا معنٰی لري چې په اشتقاق
کښې هغه کرښه چې تړلے فارمول~\LR{$!B$} لري، په دې شرط توجيه شوې ده چې
\LR{$!A$} او \LR{$!A \lif !B$} دواړه، د کوم تړلے فارمول~\LR{$!A$} دپاره، په
اشتقاق کښې له~\LR{$!B$} څخه مخکښې راغلي وي۔
\olpseganchor{OLP-0068-B007}{end}

\olpseganchor{OLP-0068-B008}{start}
پر بديهي اشتقاقونه ولاړه \LR{$\Proves$} اړيکه داسې تعريفېږي:
\LR{$\Gamma \Proves !A$} هله او يوازې هله چې داسې اشتقاق موجود
وي چې تړلے فارمول~\LR{$!A$} ئې وروستے فارمول وي، او \LR{$\Gamma$} په هغه
اشتقاق کښې د هغو تړلي فارمولونه سټ ونيول شي چې د پورته (۲) له
مخې توجيه شوي وي۔ \LR{$!A$} هله قضيه ده چې \LR{$!A$} داسې اشتقاق
ولري چې~\LR{$\Gamma$} پکې تش وي؛ يعنې په اشتقاق کښې هره
تړلے فارمول يا د (۱) له مخې توجيه شوې وي يا د~(۳) له مخې۔ مثلاً دا
اشتقاق ښيي چې \LR{$\Proves !A  \lif (!B \lif (!B \lor !A))$}:
\begin{derivation}
  1. & \LR{$!B \lif (!B \lor !A)$} \\
  2. & \LR{$(!B \lif (!B \lor !A)) \lif (!A  \lif (!B \lif (!B \lor !A)))$}\\
  3. & \LR{$!A  \lif (!B \lif (!B \lor !A))$}
\end{derivation}
په ۱مه کرښه تړلے فارمول د بديهي اصل \LR{$!A \lif (!A \lor !B)$} بڼه
لري، چې د \LR{$!A$} او \LR{$!B$} ونډې پکې بدلې شوې دي۔ په ۲مه کرښه جمله د
بديهي اصل \LR{$!A \lif (!B \lif !A)$} بڼه لري۔ نو دواړه کرښې توجيه شوې
دي۔ ۳مه کرښه د مودوس پوننس له مخې توجيه شوې ده: که په لنډه ئې
\LR{$!D$} وليکو، نو ۲مه کرښه د \LR{$!C \lif !D$} بڼه لري، چې \LR{$!C$} هم
\LR{$!B \lif (!B \lor !A)$}، يعنې ۱مه کرښه، ده۔
\olpseganchor{OLP-0068-B008}{end}

\olpseganchor{OLP-0068-B009}{start}
يو سټ \LR{$\Gamma$} هله ناسازګار دے چې \LR{$\Gamma \Proves \lfalse$}۔ د
بديهي اصولو يو بشپړ نظام به دا هم ثابتوي چې د هر~\LR{$!A$} دپاره
\LR{$\lfalse \lif !A$}؛ نو که \LR{$\Gamma$} ناسازګار وي، د هر~\LR{$!A$} دپاره
\LR{$\Gamma \Proves !A$}۔
\olpseganchor{OLP-0068-B009}{end}

\olpseganchor{OLP-0068-B010}{start}
د منطق د بديهي اشتقاقونه نظامونه لومړے ګوټلوب فرېګه په خپل
۱۸۷۹ز اثر \emph{بېګريفشريفت} کښې وړاندې کړل، چې له همدې امله زياتره
د معاصر منطق لومړے اثر ګڼل کېږي۔ بيا الفرېډ نارت وايټ هېډ او برټرېنډ
رسل په \emph{پرېنسيپيا ماتيماتيکا} کښې، او ډيوېډ هيلبرټ او شاګردانو
ئې په ۱۹۲۰مو کلونو کښې بشپړ کړل۔ ځکه ورته زياتره ``د فرېګه نظامونه''
يا ``د هيلبرټ نظامونه'' ويل کېږي۔ دا ډېر انعطاف منونکي دي، ځکه چې د
يو منطق دپاره بديهي نظام زياتره په اسانۍ موندل کېږي۔ دا چې
اشتقاقونه ډېر ساده جوړښت او يوازې يوه يا دوه استنتاجي قاعدې
لري، د هغو \emph{په اړه} د څيزونو ثابتول هم نسبتاً اسانه دي۔ خو په
عمل کښې ئې کارول ډېر ګران دي؛ يعنې ثبوتونه موندل او ليکل ګران دي۔
\olpseganchor{OLP-0068-B010}{end}

\olpseganchor{OLP-0069-B004}{start}
\olchapter{fol}{seq}{د سېکوېنټ حساب}
\olpseganchor{OLP-0069-B004}{end}

\olpseganchor{OLP-0069-B005}{start}
\begin{editorial}
  دا څپرکے د کلاسیک د لومړۍ درجې منطق دپاره د ګېنتڅن معياري
  سېکوېنټ حساب \texttt{LK} وړاندې کوي۔ دې ته نور مثالونه او تمرينونه
  هم پکار دي۔ د ثبوت د نظام په توګه د سېکوېنټ له حساب سره د اړوندو
  موادو د شاملولو يا اېستلو دپاره د \texttt{prfLK} ټېګ وکاروئ۔
\end{editorial}
\olpseganchor{OLP-0069-B005}{end}





\olpseganchor{OLP-0069-B008}{start}
%
  

\olpseganchor{OLP-0069-B008}{end}







\olpseganchor{OLP-0069-B012}{start}
%
  

\olpseganchor{OLP-0069-B012}{end}







\olpseganchor{OLP-0069-B016}{start}
%
  

\olpseganchor{OLP-0069-B016}{end}



\olpseganchor{OLP-0069-B018}{start}
%
  
  

\olpseganchor{OLP-0069-B018}{end}

\olpseganchor{OLP-0070-B004}{start}
\olfileid{fol}{seq}{rul}
\olpseganchor{OLP-0070-B004}{end}

\olpseganchor{OLP-0070-B005}{start}
\olsection{قاعدې او اشتقاقونه}
\olpseganchor{OLP-0070-B005}{end}

\olpseganchor{OLP-0070-B006}{start}
د راتلونکو خبرو دپاره، پرېږدئ چې \LR{$\Gamma, \Delta, \Pi, \Lambda$} د
تړلي فارمولونه متناهي لړۍ وښيي۔
\olpseganchor{OLP-0070-B006}{end}

\olpseganchor{OLP-0070-B007}{start}
\begin{defn}[سېکوېنټ]
\emph{سېکوېنټ} د لاندې بڼې څرګندنه ده:
\[
\Gamma \Sequent \Delta
\]
دلته \LR{$\Gamma$} او \LR{$\Delta$} د ژبې \LR{$\Lang L$} د تړلي فارمولونه متناهي
(او کېدے شي تشې) لړۍ دي۔ \LR{$\Gamma$} ته \emph{مقدم} او \LR{$\Delta$} ته
\emph{تالي} ويل کېږي۔
\end{defn}
\olpseganchor{OLP-0070-B007}{end}

\olpseganchor{OLP-0070-B008}{start}
\begin{explain}
د يوه سېکوېنټ تر شا شهودي مفکوره دا ده: که د مقدم ټولې تړلي فارمولونه
صدق کوي، نو د تالي لږ تر لږه يوه تړلي فارمولونه صدق کوي۔ يعنې که
\LR{$\Gamma = \tuple{!A_1, \dots, !A_m}$} او
\LR{$\Delta = \tuple{!B_1, \dots, !B_n}$} وي، نو \LR{$\Gamma \Sequent \Delta$}
هله صدق کوي چې
\[
(!A_1 \land \cdots \land !A_m) \lif (!B_1 \lor \cdots \lor
!B_n)
\]
صدق کوي۔ دوه خاص حالتونه شته: يو هغه چې \LR{$\Gamma$} تش وي، او بل هغه
چې \LR{$\Delta$}~تش وي۔ کله چې \LR{$\Gamma$} تش وي، يعنې \LR{$m = 0$}، نو \LR{$\quad
\Sequent \Delta$} هله صدق کوي چې \LR{$!B_1 \lor \dots \lor !B_n$} صدق
کوي۔ کله چې \LR{$\Delta$} تش وي، يعنې \LR{$n = 0$}، نو \LR{$\Gamma \Sequent \quad$}
هله صدق کوي چې \LR{$\lnot(!A_1 \land \dots \land !A_m)$} صدق کوي۔
موږ يو سېکوېنټ هله معتبر بولو چې ورسره سمون لرونکې تړلے فارمول
معتبره وي۔
\end{explain}
\olpseganchor{OLP-0070-B008}{end}

\olpseganchor{OLP-0070-B009}{start}
که \LR{$\Gamma$} د تړلي فارمولونه يوه لړۍ وي، نو \LR{$\Gamma, !A$} د هغې
پايلې دپاره ليکو چې~\LR{$!A$} د~\LR{$\Gamma$} ښي پاے ته ورزيات شي (او \LR{$!A,
\Gamma$} د هغې پايلې دپاره چې~\LR{$!A$} د~\LR{$\Gamma$} کيڼ پاے ته ورزيات
شي)۔ که \LR{$\Delta$} هم د تړلي فارمولونه يوه لړۍ وي، نو \LR{$\Gamma,
\Delta$} د دواړو لړيو نښلول دي۔
\olpseganchor{OLP-0070-B009}{end}

\olpseganchor{OLP-0070-B010}{start}
\begin{defn}[لومړنے سېکوېنټ]
\emph{لومړنے سېکوېنټ} داسې سېکوېنټ دے
چې له لاندې بڼو څخه يوه لري:
  \begin{enumerate}
    \item \LR{$!A \Sequent !A$}
    \item \LR{$\quad \Sequent \ltrue$}
    \item \LR{$\lfalse \Sequent \quad$}
  \end{enumerate}، او \LR{$!A$} د ژبې هره تړلے فارمول
کېدے شي۔
\end{defn}
\olpseganchor{OLP-0070-B010}{end}

\olpseganchor{OLP-0070-B011}{start}
په سېکوېنټ حساب کښې اشتقاقونه د سېکوېنټونو داسې ټاکلې ونې
دي چې تر ټولو بره سېکوېنټونه ئې لومړني سېکوېنټونه وي؛ او که يو
سېکوېنټ د يوه يا دوو نورو سېکوېنټونو لاندې راشي، بايد د استنتاج د
يوې قاعدې له مخې ترې په سمه توګه راووځي۔ د~\LR{$\Log{LK}$} قاعدې پر دوو
اصلي ډولونو وېشل کېږي: \emph{منطقي} قاعدې او \emph{جوړښتي} قاعدې۔
منطقي قاعدې د لاندې سېکوېنټ د هغې تړلے فارمول د اصلي عملګر په نوم نومول کېږي چې \LR{$!A$} او/يا \LR{$!B$} لري۔ هره قاعده
دوه بڼې لري: يوه د داسې سېکوېنټ د استنتاج دپاره چې د
منطقي عملګر لرونکې تړلے فارمول ئې کيڼ لوري ته وي، او بله د هغه
سېکوېنټ دپاره چې دا تړلے فارمول ئې ښي لوري ته وي۔
\olpseganchor{OLP-0070-B011}{end}

\olpseganchor{OLP-0071-B004}{start}
\olfileid{fol}{seq}{prl}
\olpseganchor{OLP-0071-B004}{end}

\olpseganchor{OLP-0071-B005}{start}
\olsection{بياني قاعدې}
\olpseganchor{OLP-0071-B005}{end}

\olpseganchor{OLP-0071-B006}{start}
\subsection{د \LR{$\lnot$} دپاره قاعدې}
\olpseganchor{OLP-0071-B006}{end}

\olpseganchor{OLP-0071-B007}{start}
\begin{defish}
\Axiom$ \Gamma \fCenter \Delta, !A $
\RightLabel{\LeftR{\lnot}}
\UnaryInf$ \lnot !A, \Gamma \fCenter \Delta$
\DisplayProof
\hfill
\Axiom$!A, \Gamma \fCenter \Delta$
\RightLabel{\RightR{\lnot}}
\UnaryInf$ \Gamma \fCenter \Delta, \lnot !A $
\DisplayProof
\end{defish}
\olpseganchor{OLP-0071-B007}{end}

\olpseganchor{OLP-0071-B008}{start}
\subsection{د \LR{$\land$} دپاره قاعدې}
\olpseganchor{OLP-0071-B008}{end}

\olpseganchor{OLP-0071-B009}{start}
\begin{defish}\noindent
\begin{tabular}{l}
\Axiom$ !A, \Gamma \fCenter \Delta$
\RightLabel{\LeftR{\land}}
\UnaryInf$ !A \land !B, \Gamma \fCenter \Delta$
\DisplayProof
\\[3ex]
\Axiom$!B, \Gamma \fCenter \Delta$
\RightLabel{\LeftR{\land}}
\UnaryInf$!A \land !B, \Gamma \fCenter \Delta$
\DisplayProof
\end{tabular}
\hfill
\Axiom$\Gamma \fCenter \Delta, !A$
\Axiom$ \Gamma \fCenter \Delta, !B$
\RightLabel{\RightR{\land}}
\BinaryInf$ \Gamma \fCenter \Delta, !A \land !B $
\DisplayProof
\end{defish}
\olpseganchor{OLP-0071-B009}{end}

\olpseganchor{OLP-0071-B010}{start}
\subsection{د \LR{$\lor$} دپاره قاعدې}
\olpseganchor{OLP-0071-B010}{end}

\olpseganchor{OLP-0071-B011}{start}
\begin{defish}
\Axiom$!A, \Gamma \fCenter \Delta$
\Axiom$!B, \Gamma \fCenter \Delta$
\RightLabel{\LeftR{\lor}}
\BinaryInf$!A \lor !B, \Gamma \fCenter \Delta$
\DisplayProof
\hfill
\begin{tabular}{r}
\Axiom$\Gamma \fCenter \Delta, !A$
\RightLabel{\RightR{\lor}}
\UnaryInf$ \Gamma \fCenter \Delta, !A \lor !B$
\DisplayProof
\\[3ex]
\Axiom$ \Gamma \fCenter \Delta, !B$
\RightLabel{\RightR{\lor}}
\UnaryInf$ \Gamma \fCenter \Delta, !A \lor !B$
\DisplayProof
\end{tabular}
\end{defish}
\olpseganchor{OLP-0071-B011}{end}

\olpseganchor{OLP-0071-B012}{start}
\subsection{د \LR{$\lif$} دپاره قاعدې}
\olpseganchor{OLP-0071-B012}{end}

\olpseganchor{OLP-0071-B013}{start}
\begin{defish}
\Axiom$ \Gamma \fCenter \Delta, !A$
\Axiom$ !B, \Pi \fCenter \Lambda$
\RightLabel{\LeftR{\lif}}
\BinaryInf$ !A \lif !B, \Gamma, \Pi \fCenter \Delta, \Lambda$
\DisplayProof
\hfill
\Axiom$ !A, \Gamma \fCenter \Delta, !B$
\RightLabel{\RightR{\lif}}
\UnaryInf$ \Gamma \fCenter \Delta, !A \lif !B $
\DisplayProof
\end{defish}
\olpseganchor{OLP-0071-B013}{end}

\olpseganchor{OLP-0072-B004}{start}
\olfileid{fol}{seq}{qrl}
\olpseganchor{OLP-0072-B004}{end}

\olpseganchor{OLP-0072-B005}{start}
\olsection{د سورونو قاعدې}
\olpseganchor{OLP-0072-B005}{end}

\olpseganchor{OLP-0072-B006}{start}
\subsection{د \LR{$\lforall$} دپاره قاعدې}
\olpseganchor{OLP-0072-B006}{end}

\olpseganchor{OLP-0072-B007}{start}
\begin{defish}
\Axiom$ !A(t), \Gamma \fCenter \Delta$
\RightLabel{\LeftR{\lforall}}
\UnaryInf$ \lforall[x][!A(x)],\Gamma \fCenter \Delta$
\DisplayProof
\hfill
\Axiom$ \Gamma \fCenter \Delta, !A(a) $
\RightLabel{\RightR{\lforall}}
\UnaryInf$ \Gamma \fCenter \Delta, \lforall[x][!A(x)]$
\DisplayProof
\end{defish}
\olpseganchor{OLP-0072-B007}{end}

\olpseganchor{OLP-0072-B008}{start}
په \LeftR{\lforall} کښې \LR{$t$} يوه تړلې اصطلاح ده (يعنې داسې اصطلاح
چې متغيرونه نۀ لري)۔ په \RightR{\lforall} کښې \LR{$a$}~ثابت سمبول دے
چې بايد د \RightR{\lforall} قاعدې په لاندې سېکوېنټ کښې هېڅ ځاے نۀ
وي راغلے۔ موږ \LR{$a$} د \RightR{\forall} استنتاج \emph{ځانګړے متغير}
بولو۔\footnote{موږ د “ځانګړي متغير” اصطلاح کاروو، که څه هم په
پورتنۍ قاعده کښې \LR{$a$} ثابت سمبول دے۔ دا نوم تاريخي لاملونه لري۔}
\olpseganchor{OLP-0072-B008}{end}

\olpseganchor{OLP-0072-B009}{start}
\subsection{د \LR{$\lexists$} دپاره قاعدې}
\olpseganchor{OLP-0072-B009}{end}

\olpseganchor{OLP-0072-B010}{start}
\begin{defish}
\Axiom$ !A(a), \Gamma \fCenter \Delta $
\RightLabel{\LeftR{\lexists}}
\UnaryInf$ \lexists[x][!A(x)], \Gamma \fCenter \Delta$
\DisplayProof
\hfill
\Axiom$ \Gamma \fCenter \Delta, !A(t) $
\RightLabel{\RightR{\lexists}}
\UnaryInf$ \Gamma \fCenter \Delta, \lexists[x][!A(x)]$
\DisplayProof
\end{defish}
\olpseganchor{OLP-0072-B010}{end}

\olpseganchor{OLP-0072-B011}{start}
بيا هم، \LR{$t$}~يوه تړلې اصطلاح ده، او \LR{$a$}~داسې ثابت سمبول دے چې
د \LeftR{\lexists} قاعدې په لاندې سېکوېنټ کښې نۀ راځي۔ موږ \LR{$a$} د
\LeftR{\lexists} استنتاج \emph{ځانګړے متغير} بولو۔
\olpseganchor{OLP-0072-B011}{end}

\olpseganchor{OLP-0072-B012}{start}
دا شرط چې ځانګړے متغير د \RightR{\lforall} يا
\LeftR{\lexists} استنتاج په لاندې سېکوېنټ کښې نۀ راځي،
\emph{د ځانګړي متغير شرط} بلل کېږي۔
\olpseganchor{OLP-0072-B012}{end}

\olpseganchor{OLP-0072-B013}{start}
\begin{explain}
هغه دود راياد کړئ چې کله \LR{$!A$} داسې فارمول وي چې
متغير~\LR{$x$} پکې ازاد وي، نو دا په~\LR{$!A(x)$} ليکلو ښيو۔ په همدې
چوکاټ کښې بيا \LR{$!A(t)$} د~\LR{$\Subst{!A}{t}{x}$} لنډون دے۔ نو د
\RightR{\lexists} قاعده داسې هم ليکلے شو:
\begin{prooftree}
  \Axiom$\Gamma \fCenter \Delta, \Subst{!A}{t}{x}$
  \RightLabel{\RightR{\lexists}}
  \UnaryInf$\Gamma \fCenter \Delta, \lexists[x][!A]$
\end{prooftree}
پام وکړئ چې \LR{$t$} کېدے شي له مخکښې په~\LR{$!A$} کښې راغلے وي؛ مثلاً
\LR{$!A$}~کېدے شي~\LR{$\Atom{\Obj P}{t,x}$} وي۔ له دې امله، له
~\LR{$\Gamma \Sequent \Delta, \Atom{\Obj P}{t,t}$} څخه
\LR{$\Gamma \Sequent \Delta, \lexists[x][\Atom{\Obj P}{t,x}]$} استنتاجول
د~\RightR{\lexists} سمه کارونه ده---تاسو د مقدمې د~\LR{$t$} يوه يا
زياتې پېښې په تړلي متغير~\LR{$x$} “بدلولے” شئ، او اړينه نۀ ده
چې ټولې پېښې ئې بدلې کړئ۔ خو د \RightR{\lforall} او
~\LeftR{\lexists} د ځانګړي متغير شرطونه غواړي چې
ثابت سمبول~\LR{$a$} په~\LR{$!A$} کښې نۀ راځي۔ نو د~\LR{$\RightR{\lforall}$} په
کارولو سره له \LR{$\Gamma \Sequent \Delta, \Atom{\Obj P}{a,a}$} څخه
\LR{$\Gamma \Sequent \Delta, \lforall[x][\Atom{\Obj P}{a,x}]$} په سمه
توګه نۀ شئ استنتاجولے۔
\end{explain}
\olpseganchor{OLP-0072-B013}{end}

\olpseganchor{OLP-0072-B014}{start}
\begin{explain}
په \RightR{\lexists} او \LeftR{\lforall} کښې پر اصطلاح~\LR{$t$} هېڅ
بنديز نشته۔ بل لوري ته، د \LeftR{\lexists} او
\RightR{\lforall} په قاعدو کښې د ځانګړي متغير شرط غواړي چې
ثابت سمبول~\LR{$a$} په پورتني سېکوېنټ کښې د~\LR{$!A(a)$} څخه بهر هېڅ ځاے
نۀ راځي۔ دا شرط د نظام د سموالي د يقيني کولو دپاره اړين دے؛ يعنې
نظام يوازې معتبر سېکوېنټونه اشتقاق کړي۔ له دې شرط پرته لاندې
استنتاجونه روا کېدل:
\begin{prooftree}
  \Axiom$!A(a) \fCenter !A(a)$
  \RightLabel{*\LeftR{\lexists}}
  \UnaryInf$\lexists[x][!A(x)] \fCenter !A(a)$
  \RightLabel{\RightR{\lforall}}
  \UnaryInf$\lexists[x][!A(x)] \fCenter \lforall[x][!A(x)]$
  \DisplayProof\bottomAlignProof
  \qquad
  \Axiom$!A(a) \fCenter !A(a)$
  \RightLabel{*\RightR{\lforall}}
  \UnaryInf$!A(a) \fCenter \lforall[x][!A(x)]$
  \RightLabel{\LeftR{\lexists}}
  \UnaryInf$\lexists[x][!A(x)] \fCenter \lforall[x][!A(x)]$
\end{prooftree}
خو \LR{$\lexists[x][!A(x)] \Sequent \lforall[x][!A(x)]$} معتبر نۀ دے۔
\end{explain}
\olpseganchor{OLP-0072-B014}{end}

\olpseganchor{OLP-0073-B004}{start}
\olfileid{fol}{seq}{srl}
\olpseganchor{OLP-0073-B004}{end}

\olpseganchor{OLP-0073-B005}{start}
\olsection{جوړښتي قاعدې}
\olpseganchor{OLP-0073-B005}{end}

\olpseganchor{OLP-0073-B006}{start}
موږ څو داسې قاعدو ته هم اړتيا لرو چې د يوه سېکوېنټ په کيڼ او ښي
لوري کښې د تړلي فارمولونه د بيا ترتيبولو اجازه راکړي۔ ځکه چې منطقي
قاعدې غواړي هغه تړلي فارمولونه چې قاعده پرې کار کوي، يا په بشپړ کيڼ
او يا په بشپړ ښي پاے کښې وي، نو د “تبادلې” داسې قاعدې ته اړتيا
لرو چې تړلي فارمولونه سم ځاے ته يوسي۔ کله کله دا هم مهمه وي چې دوه
يو شان تړلي فارمولونه سره يو کړو او هر لوري ته تړلے فارمول
ورزياته کړو۔
\olpseganchor{OLP-0073-B006}{end}

\olpseganchor{OLP-0073-B007}{start}
\subsection{کمزوري کول}
\olpseganchor{OLP-0073-B007}{end}

\olpseganchor{OLP-0073-B008}{start}
\begin{defish}
\Axiom$ \Gamma \fCenter \Delta $
\RightLabel{\LeftR{\Weakening}}
\UnaryInf$ !A, \Gamma \fCenter \Delta$
\DisplayProof
\hfill
\Axiom$ \Gamma \fCenter \Delta$
\RightLabel{\RightR{\Weakening}}
\UnaryInf$ \Gamma \fCenter \Delta, !A$
\DisplayProof
\end{defish}
\olpseganchor{OLP-0073-B008}{end}

\olpseganchor{OLP-0073-B009}{start}
\subsection{انقباض}
\olpseganchor{OLP-0073-B009}{end}

\olpseganchor{OLP-0073-B010}{start}
\begin{defish}
\Axiom$ !A, !A, \Gamma \fCenter \Delta $
\RightLabel{\LeftR{\Contraction}}
\UnaryInf$ !A, \Gamma \fCenter \Delta$
\DisplayProof
\hfill
\Axiom$ \Gamma \fCenter \Delta, !A, !A$
\RightLabel{\RightR{\Contraction}}
\UnaryInf$ \Gamma \fCenter \Delta, !A$
\DisplayProof
\end{defish}
\olpseganchor{OLP-0073-B010}{end}

\olpseganchor{OLP-0073-B011}{start}
\subsection{تبادله}
\olpseganchor{OLP-0073-B011}{end}

\olpseganchor{OLP-0073-B012}{start}
\begin{defish}
\Axiom$ \Gamma, !A, !B, \Pi \fCenter \Delta $
\RightLabel{\LeftR{\Exchange}}
\UnaryInf$ \Gamma, !B, !A, \Pi \fCenter \Delta$
\DisplayProof
\hfill
\Axiom$ \Gamma \fCenter \Delta, !A, !B, \Lambda$
\RightLabel{\RightR{\Exchange}}
\UnaryInf$ \Gamma \fCenter \Delta, !B, !A, \Lambda$
\DisplayProof
\end{defish}
\olpseganchor{OLP-0073-B012}{end}

\olpseganchor{OLP-0073-B013}{start}
د کمزوري کولو، انقباض او تبادلې د استنتاجونو يوه لړۍ به زياتره د
استنتاج په دوو کرښو ښوول کېږي۔
\olpseganchor{OLP-0073-B013}{end}

\olpseganchor{OLP-0073-B014}{start}
لاندې قاعده چې “پرېکون” بلل کېږي، په کره توګه اړينه نۀ ده، خو د
اشتقاقونه بيا کارول او يو ځاے کول ډېر اسانوي۔
\begin{defish}
\[
\Axiom$ \Gamma \fCenter \Delta, !A$
\Axiom$ !A, \Pi \fCenter \Lambda $
\RightLabel{\Cut}
\BinaryInf$ \Gamma, \Pi \fCenter \Delta, \Lambda$
\DisplayProof
\]
\end{defish}
\olpseganchor{OLP-0073-B014}{end}

\olpseganchor{OLP-0074-B004}{start}
\olfileid{fol}{seq}{der}
\olpseganchor{OLP-0074-B004}{end}

\olpseganchor{OLP-0074-B005}{start}
\olsection{اشتقاقونه}
\olpseganchor{OLP-0074-B005}{end}

\olpseganchor{OLP-0074-B006}{start}
\begin{explain}
موږ وويل چې لومړنے سېکوېنټ څنګه وي او د استنتاج قاعدې مو ورکړې۔
په سېکوېنټ حساب کښې اشتقاقونه له دغو څخه په استقرايي توګه
جوړېږي: هر اشتقاق يا په يوازې ځان يو لومړنے سېکوېنټ وي،
او يا له يوه يا دوو اشتقاقونه او ورپسې يوه استنتاج څخه جوړ وي۔
\end{explain}
\olpseganchor{OLP-0074-B006}{end}

\olpseganchor{OLP-0074-B007}{start}
\begin{defn}[\LR{$\Log{LK}$} اشتقاق]
د يوه سېکوېنټ~\LR{$S$} \emph{\LR{$\Log{LK}$}-اشتقاق} د سېکوېنټونو
داسې متناهي ونه ده چې لاندې شرطونه پوره کوي:
\begin{enumerate}
\item د ونې تر ټولو بره سېکوېنټونه لومړني سېکوېنټونه دي۔
\item د ونې تر ټولو لاندې سېکوېنټ~\LR{$S$} دے۔
\item له~\LR{$S$} پرته د ونې هر سېکوېنټ د استنتاج د يوې قاعدې د سمې
  کارونې مقدمه ده، او د هغې کارونې پايله په ونه کښې د هماغه
  سېکوېنټ مخامخ لاندې ولاړه وي۔
\end{enumerate}
بيا وايو چې \LR{$S$} د اشتقاق \emph{وروستنے سېکوېنټ} دے او
\LR{$S$} \emph{په \LR{$\Log{LK}$} کښې د اشتقاق وړ} دے (يا
\LR{$\Log{LK}$}-د اشتقاق وړ دے)۔
\end{defn}
\olpseganchor{OLP-0074-B007}{end}

\olpseganchor{OLP-0074-B008}{start}
\begin{ex}
هر لومړنے سېکوېنټ، مثلاً \LR{$!C \Sequent !C$}، اشتقاق دے۔
له دې څخه د يوې بلې قاعدې، مثلاً د \LR{$\LeftR{\Weakening}$} قاعدې، په
کارولو نوے اشتقاق ترلاسه کولے شو:
\begin{prooftree}
\Axiom$ \Gamma \fCenter \Delta $
\RightLabel{\LeftR{\Weakening}}
\UnaryInf$ !A, \Gamma \fCenter \Delta$
\end{prooftree}
خو قاعده عمومي مراد لري: په قاعده کښې \LR{$!A$} د هرې تړلے فارمول،
مثلاً د~\LR{$!D$}، په ځاے راوستلے شو۔ که مقدمه زموږ له لومړني سېکوېنټ
\LR{$!C \Sequent !C$} سره سمون خوري، نو \LR{$\Gamma$} او \LR{$\Delta$} دواړه
يوازې~\LR{$!C$} دي، او پايله به بيا \LR{$!D, !C \Sequent !C$} وي۔ نو لاندې
اشتقاق دے:
\begin{prooftree}
\Axiom$ !C \fCenter !C $
\RightLabel{\LeftR{\Weakening}}
\UnaryInf$ !D, !C \fCenter !C$
\end{prooftree}
اوس يوه بله قاعده، مثلاً \LR{$\LeftR{\Exchange}$}، کارولے شو چې په کيڼ
لوري کښې د دوو تړلي فارمولونه د ځاے بدلولو اجازه راکوي۔ نو لاندې هم
سم اشتقاق دے:
\begin{prooftree}
\Axiom$ !C \fCenter !C $
\RightLabel{\LeftR{\Weakening}}
\UnaryInf$ !D, !C \fCenter !C$
\RightLabel{\LeftR{\Exchange}}
\UnaryInf$ !C, !D \fCenter !C$
\end{prooftree}
د دې قاعدې په دې کارونه کښې، چې داسې ورکړل شوې وه:
\begin{prooftree}
\Axiom$ \Gamma, !A, !B, \Pi \fCenter \Delta $
\RightLabel{\LeftR{\Exchange}}
\UnaryInf$ \Gamma, !B, !A, \Pi \fCenter \Delta,$
\end{prooftree}
\LR{$\Gamma$} او \LR{$\Pi$} دواړه تش وو، \LR{$\Delta$} هم \LR{$!C$} دے، او د \LR{$!A$} او
\LR{$!B$} ونډې په همدې ترتيب \LR{$!D$} او~\LR{$!C$} لوبوي۔ نژدې په همدې طريقه
وينو چې
\begin{prooftree}
\Axiom$ !D \fCenter !D $
\RightLabel{\LeftR{\Weakening}}
\UnaryInf$ !C, !D \fCenter !D$
\end{prooftree}
هم اشتقاق دے۔ اوس دا دواړه اشتقاقونه اخلو او د
\LR{$\RightR{\land}$} په کارولو ئې يو ځاے کوو۔ هغه قاعده داسې وه:
\begin{prooftree}
\Axiom$\Gamma \fCenter \Delta, !A$
\Axiom$ \Gamma \fCenter \Delta, !B$
\RightLabel{\RightR{\land}}
\BinaryInf$ \Gamma \fCenter \Delta, !A \land !B $
\end{prooftree}
زموږ په حالت کښې مقدمې بايد د هغو اشتقاقونه له وروستيو
سېکوېنټونو سره سمون وخوري چې په مقدمو پاے ته رسېږي۔ يعنې \LR{$\Gamma$}
د \LR{$!C, !D$}، \LR{$\Delta$} تشه، \LR{$!A$} د \LR{$!C$} او \LR{$!B$} د \LR{$!D$} سره برابره
ده۔ نو که استنتاج سم وي، پايله \LR{$!C, !D \Sequent !C
\land !D$} ده۔
\begin{prooftree}
\Axiom$ !C \fCenter !C $
\RightLabel{\LeftR{\Weakening}}
\UnaryInf$ !D, !C \fCenter !C$
\RightLabel{\LeftR{\Exchange}}
\UnaryInf$ !C, !D \fCenter !C$
\Axiom$ !D \fCenter !D $
\RightLabel{\LeftR{\Weakening}}
\UnaryInf$ !C, !D \fCenter !D$
\RightLabel{\RightR{\land}}
\BinaryInf$ !C, !D \fCenter !C \land !D $
\end{prooftree}
البته مقدمې سرچپه کولے هم شو؛ بيا \LR{$!A$} به \LR{$!D$} او \LR{$!B$} به~\LR{$!C$} وي۔
\begin{prooftree}
\Axiom$ !D \fCenter !D $
\RightLabel{\LeftR{\Weakening}}
\UnaryInf$ !C, !D \fCenter !D$
\Axiom$ !C \fCenter !C $
\RightLabel{\LeftR{\Weakening}}
\UnaryInf$ !D, !C \fCenter !C$
\RightLabel{\LeftR{\Exchange}}
\UnaryInf$ !C, !D \fCenter !C$
\RightLabel{\RightR{\land}}
\BinaryInf$ !C, !D \fCenter !D \land !C $
\end{prooftree}
\end{ex}
\olpseganchor{OLP-0074-B008}{end}

\olpseganchor{OLP-0075-B004}{start}
\olfileid{fol}{seq}{pro}
\olpseganchor{OLP-0075-B004}{end}

\olpseganchor{OLP-0075-B005}{start}
\olsection{د اشتقاقونه بېلګې}
\olpseganchor{OLP-0075-B005}{end}

\olpseganchor{OLP-0075-B006}{start}
\begin{ex}
د سېکوېنټ \LR{$!A \land !B \Sequent !A$} يو
\LR{$\Log{LK}$}-اشتقاق ورکړئ۔
\olpseganchor{OLP-0075-B006}{end}

\olpseganchor{OLP-0075-B007}{start}
د اشتقاق په لاندې برخه کښې د غوښتل شوي وروستي سېکوېنټ په
ليکلو پيل کوو۔
\begin{prooftree}
\AxiomC{}
\UnaryInf$!A\land !B \fCenter !A$
\end{prooftree}
بيا بايد معلومه کړو چې د کوم ډول استنتاج لاندې سېکوېنټ دا بڼه
لرلے شي۔ دا يوه جوړښتي قاعده کېدے شي، خو غوره ده چې لومړے کومه
منطقي قاعده ولټوؤ۔ په لاندې سېکوېنټ کښې يوازينے منطقي نښلوونکے
\LR{$\land$} دے؛ نو د \LR{$\land$} قاعده لټوؤ، او ځکه چې د \LR{$\land$} نښه په
مقدم کښې راځي، د \LeftR{\land} قاعده ګورو۔
\begin{prooftree}
\AxiomC{}
\RightLabel{\LeftR{\land}}
\UnaryInf$!A\land !B \fCenter !A$
\end{prooftree}
د \LeftR{\land} استنتاج د پورتني سېکوېنټ دپاره دوه امکانونه شته:
يا \LR{$!A \Sequent !A$} او يا \LR{$!B \Sequent !A$}۔ څرګنده ده چې
\LR{$!A \Sequent !A$} يو لومړنے سېکوېنټ دے، او دا ښه خبره ده؛ خو
\LR{$!B \Sequent !A$} په عمومي توګه د اشتقاق وړ نۀ دے۔ پورتنے سېکوېنټ
ورډکوو:
\begin{prooftree}
\Axiom$!A \fCenter !A$
\RightLabel{\LeftR{\land}}
\UnaryInf$!A\land !B \fCenter !A$
\end{prooftree}
اوس د سېکوېنټ \LR{$!A \land !B \Sequent !A$} سم
\LR{$\Log{LK}$}-اشتقاق لرو۔
\end{ex}
\olpseganchor{OLP-0075-B007}{end}

\olpseganchor{OLP-0075-B008}{start}
\begin{ex}
د سېکوېنټ \LR{$\lnot !A \lor !B \Sequent !A \lif !B$} يو
\LR{$\Log{LK}$}-اشتقاق ورکړئ۔
\olpseganchor{OLP-0075-B008}{end}

\olpseganchor{OLP-0075-B009}{start}
د اشتقاق په لاندې برخه کښې د غوښتل شوي وروستي سېکوېنټ په
ليکلو پيل وکړئ۔
\begin{prooftree}
\AxiomC{}
\UnaryInf$\lnot !A \lor !B \fCenter !A \lif !B$
\end{prooftree}
د داسې منطقي قاعدې د موندلو دپاره چې دا وروستنے سېکوېنټ راکړي، د
وروستي سېکوېنټ منطقي نښلوونکي ګورو: \LR{$\lnot$}، \LR{$\lor$} او \LR{$\lif$}۔
اوس مهال يوازې \LR{$\lor$} او \LR{$\lif$} مهم دي، ځکه چې د وروستي سېکوېنټ د
تړلي فارمولونه اصلي عملګرونه دي؛ خو \LR{$\lnot$} د بل نښلوونکي د
چاپېريال دننه دے، نو وروسته به ئې وڅېړو۔ له دې امله د وروستي
استنتاج د منطقي قاعدې امکانونه \LeftR{\lor} او \RightR{\lif} دي۔
په رښتيا هره يوه غوره کولے شو، خو \RightR{\lif} غوره کوو، يوازې
د دې دپاره چې په دوو څانګو وېشل لږ وروسته کړو۔ د هغو
\RightR{\lif} استنتاجونو د بڼې له مخې چې دا لاندې سېکوېنټ ورکولے
شي، بڼه بايد داسې وي:
\begin{prooftree}
\AxiomC{}
\UnaryInf$ !A, \lnot !A \lor !B \fCenter !B $
\RightLabel{\RightR{\lif}} \UnaryInf$ \lnot !A \lor !B \fCenter !A \lif !B $
\end{prooftree}
\olpseganchor{OLP-0075-B009}{end}

\olpseganchor{OLP-0075-B010}{start}
که \LR{$\lnot !A \lor !B$} د مقدم بهرني پاے ته يوسو، د
\LeftR{\lor} قاعده کارولے شو۔ د شېما له مخې دا بايد په لاندې ډول
پر دوو پورتنيو سېکوېنټونو ووېشل شي۔
\textbf{د ژباړې د نسخې سمون (OLSQ-001):} د سرچينې په هغو څلورو
قاعدوي ونو کښې چې د مقدم دوه فارمولونه سره بدلوي، د ښي لوري د
تبادلې نښه راغلې ده؛ دلته د قاعدې له واقعي لوري سره سمه د کيڼ
لوري د تبادلې نښه کارول کېږي۔
\begin{prooftree}
\AxiomC{}
\UnaryInf$\lnot !A, !A \fCenter !B$
\AxiomC{}
\UnaryInf$!B, !A \fCenter !B$
\RightLabel{\LeftR{\lor}}
\BinaryInf$ \lnot !A \lor !B, !A \fCenter !B $
\RightLabel{\LeftR{\Exchange}}
\UnaryInf$ !A, \lnot !A \lor !B \fCenter !B $
\RightLabel{\RightR{\lif}}
\UnaryInf$ \lnot !A \lor !B \fCenter !A \lif !B $
\end{prooftree}
راياد کړئ چې موږ هڅه کوو تر لومړنيو سېکوېنټونو پورې په پورته لوري
لار پيدا کړو؛ داسې ښکاري چې ډېر ورنژدې يو!{} ښۍ څانګه له لومړني
سېکوېنټ څخه يوازې يو کمزوري کول او يوه تبادله لرې ده، او بيا
بشپړېږي:
\begin{prooftree}
\AxiomC{}
\UnaryInf$\lnot !A, !A \fCenter !B$
\Axiom$!B \fCenter !B$
\RightLabel{\LeftR{\Weakening}}
\UnaryInf$!A, !B \fCenter !B$
\RightLabel{\LeftR{\Exchange}}
\UnaryInf$!B, !A \fCenter !B$
\RightLabel{\LeftR{\lor}}
\BinaryInf$\lnot !A \lor !B, !A \fCenter !B $
\RightLabel{\LeftR{\Exchange}}
\UnaryInf$ !A, \lnot !A \lor !B \fCenter !B $
\RightLabel{\RightR{\lif}}
\UnaryInf$ \lnot !A \lor !B \fCenter !A \lif !B $
\end{prooftree}
\olpseganchor{OLP-0075-B010}{end}

\olpseganchor{OLP-0075-B011}{start}
اوس کيڼه څانګه ګورو۔ په هره تړلے فارمول کښې يوازينے منطقي
نښلوونکے د مقدم د تړلي فارمولونه د \LR{$\lnot$} نښه ده، نو د
\LeftR{\lnot} قاعدې يوه کارونه ګورو۔
\begin{prooftree}
\AxiomC{}
\UnaryInf$ !A \fCenter !B, !A$
\RightLabel{\LeftR{\lnot}}
\UnaryInf$\lnot !A, !A \fCenter !B$
\Axiom$!B \fCenter !B$
\RightLabel{\LeftR{\Weakening}}
\UnaryInf$!A, !B \fCenter !B$
\RightLabel{\LeftR{\Exchange}}
\UnaryInf$!B, !A \fCenter !B$
\RightLabel{\LeftR{\lor}}
\BinaryInf$\lnot !A \lor !B, !A \fCenter !B $
\RightLabel{\LeftR{\Exchange}}
\UnaryInf$ !A, \lnot !A \lor !B \fCenter !B $
\RightLabel{\RightR{\lif}}
\UnaryInf$ \lnot !A \lor !B \fCenter !A \lif !B $
\end{prooftree}
لکه ښۍ څانګه چې مو بشپړه کړه، له دې کيڼې څانګې څخه هم يوازې يو
کمزوري کول او يوه تبادله لرې يو۔
\begin{prooftree}
\Axiom$!A \fCenter !A$
\RightLabel{\RightR{\Weakening}}
\UnaryInf$ !A \fCenter !A, !B$
\RightLabel{\RightR{\Exchange}}
\UnaryInf$ !A \fCenter !B, !A$
\RightLabel{\LeftR{\lnot}}
\UnaryInf$\lnot !A, !A \fCenter !B$
\Axiom$!B \fCenter !B$
\RightLabel{\LeftR{\Weakening}}
\UnaryInf$!A, !B \fCenter !B$
\RightLabel{\LeftR{\Exchange}}
\UnaryInf$!B, !A \fCenter !B$
\RightLabel{\LeftR{\lor}}
\BinaryInf$\lnot !A \lor !B, !A \fCenter !B $
\RightLabel{\LeftR{\Exchange}}
\UnaryInf$ !A, \lnot !A \lor !B \fCenter !B $
\RightLabel{\RightR{\lif}}
\UnaryInf$ \lnot !A \lor !B \fCenter !A \lif !B $
\end{prooftree}
\end{ex}
\olpseganchor{OLP-0075-B011}{end}

\olpseganchor{OLP-0075-B012}{start}
\begin{ex}
د سېکوېنټ \LR{$\lnot !A \lor \lnot !B
\Sequent \lnot (!A \land !B)$} يو \LR{$\Log{LK}$}-اشتقاق ورکړئ۔
\olpseganchor{OLP-0075-B012}{end}

\olpseganchor{OLP-0075-B013}{start}
د پورته طريقو په کارولو، غوښتل شوے وروستنے سېکوېنټ په لاندې برخه
کښې ليکو۔
\begin{prooftree}
\AxiomC{}
\UnaryInf$ \lnot !A \lor \lnot !B \fCenter \lnot (!A \land !B) $
\end{prooftree}
په وروستي سېکوېنټ کښې د تړلي فارمولونه شته اصلي نښلوونکي د \LR{$\lor$}
او \LR{$\lnot$} نښې دي۔ دلته هم \LeftR{\lor} او هم \RightR{\lnot}
کار ورکوي، خو په \RightR{\lnot} پيل کوو، ځکه چې د يوې شېبې دپاره
پر دوو څانګو وېشل ځنډوي:
\begin{prooftree}
\AxiomC{}
\UnaryInf$!A \land !B, \lnot !A \lor \lnot !B \fCenter $
\RightLabel{\RightR{\lnot}}
\UnaryInf$\lnot !A \lor \lnot !B \fCenter \lnot (!A \land !B)$
\end{prooftree}
اوس د \LeftR{\land} او \LeftR{\lor} تر منځ ټاکنه لرو۔ ګورو چې د
\LeftR{\land} قاعدې په کارولو څه کېږي: يا په سېکوېنټ \LR{$!A,
\lnot !A \lor \lnot !B \Sequent \quad$} او يا په سېکوېنټ \LR{$!B, \lnot !A
\lor \lnot !B \Sequent \quad$} پيل کولے شو۔ ځکه چې اشتقاق د
\LR{$!A$} او \LR{$!B$} له پلوه متناظر دے، لومړنے غوره کوو۔
\textbf{د ژباړې د نسخې سمون (OLSQ-002):} د سرچينې په دغو دوو
امکاني پورتنيو سېکوېنټونو کښې له دويمې جملې څخه نفي پاتې شوې وه؛
دلته د وروستي سېکوېنټ او ورپسې قاعدوي ونې له مخې په دواړو جملو
کښې نفي ساتل کېږي۔
\begin{prooftree}
\AxiomC{}
\UnaryInf$!A, \lnot !A \lor \lnot !B \fCenter $
\RightLabel{\LeftR{\land}}
\UnaryInf$!A \land !B, \lnot !A \lor \lnot !B \fCenter $
\RightLabel{\RightR{\lnot}}
\UnaryInf$\lnot !A \lor \lnot !B \fCenter \lnot (!A \land !B)$
\end{prooftree}
د اشتقاق د ډکولو په دوام وينو چې له يوې ستونزې سره مخ کېږو:
\begin{prooftree}
\Axiom$!A \fCenter !A$
\RightLabel{\LeftR{\lnot}}
\UnaryInf$ \lnot !A, !A \fCenter$
\AxiomC{}
\RightLabel{?}
\UnaryInf$!A \fCenter !B$
\RightLabel{\LeftR{\lnot}}
\UnaryInf$ \lnot !B, !A \fCenter$
\RightLabel{\LeftR{\lor}}
\BinaryInf$\lnot !A \lor \lnot !B, !A \fCenter $
\RightLabel{\LeftR{\Exchange}}
\UnaryInf$!A, \lnot !A \lor \lnot !B \fCenter $
\RightLabel{\LeftR{\land}}
\UnaryInf$!A \land !B, \lnot !A \lor \lnot !B \fCenter $
\RightLabel{\RightR{\lnot}}
\UnaryInf$\lnot !A \lor \lnot !B \fCenter \lnot (!A \land !B)$
\end{prooftree}
د ښۍ څانګې پورته برخه نوره نۀ شي راکمېدلے او د جوړښتي استنتاجونو
له لارې لومړني سېکوېنټ ته نۀ شي رسول کېدے؛ نو دا سمه لار نۀ ده۔
له دې امله پورته د \LeftR{\land} قاعدې کارول تېروتنه وه۔ بېرته
مخکيني حالت ته ځو او پر ځاے ئې \LeftR{\lor} قاعده پلې کوو، نو
ترلاسه کوو:
\begin{prooftree}
\AxiomC{}
\UnaryInf$\lnot !A, !A \land !B \fCenter $
\olpseganchor{OLP-0075-B013}{end}

\olpseganchor{OLP-0075-B014}{start}
\AxiomC{}
\UnaryInf$\lnot !B, !A \land !B \fCenter $
\olpseganchor{OLP-0075-B014}{end}

\olpseganchor{OLP-0075-B015}{start}
\RightLabel{\LeftR{\lor}}
\BinaryInf$\lnot !A \lor \lnot !B, !A \land !B \fCenter $
\RightLabel{\LeftR{\Exchange}}
\UnaryInf$!A \land !B, \lnot !A \lor \lnot !B \fCenter $
\RightLabel{\RightR{\lnot}}
\UnaryInf$\lnot !A \lor \lnot !B \fCenter \lnot (!A \land !B)$
\end{prooftree}
هره څانګه د مخکښې په شان بشپړه کړو، نو ترلاسه کوو:
\begin{prooftree}
\Axiom$ !A \fCenter!A$
\RightLabel{\LeftR{\land}}
\UnaryInf$!A \land !B \fCenter !A$
\RightLabel{\LeftR{\lnot}}
\UnaryInf$\lnot !A, !A \land !B \fCenter $
\olpseganchor{OLP-0075-B015}{end}

\olpseganchor{OLP-0075-B016}{start}
\Axiom$ !B \fCenter !B$
\RightLabel{\LeftR{\land}}
\UnaryInf$!A \land !B \fCenter !B$
\RightLabel{\LeftR{\lnot}}
\UnaryInf$\lnot !B, !A \land !B \fCenter $
\olpseganchor{OLP-0075-B016}{end}

\olpseganchor{OLP-0075-B017}{start}
\RightLabel{\LeftR{\lor}}
\BinaryInf$\lnot !A \lor \lnot !B, !A \land !B \fCenter $
\RightLabel{\LeftR{\Exchange}}
\UnaryInf$!A \land !B, \lnot !A \lor \lnot !B \fCenter $
\RightLabel{\RightR{\lnot}}
\UnaryInf$\lnot !A \lor \lnot !B \fCenter \lnot (!A \land !B)$
\end{prooftree}
(په دغو ګامونو کښې د \LR{$\land$} قاعدې د \LR{$\lnot$} له قاعدو څخه لاندې
هم پلې کېدے شوې او بيا به هم سم اشتقاق ترلاسه شوے و۔)
\end{ex}
\olpseganchor{OLP-0075-B017}{end}

\olpseganchor{OLP-0075-B018}{start}
\begin{ex}
تر اوسه مو د انقباض قاعده نۀ ده کارولې، خو کله کله اړينه وي۔ دا
ئې يوه بېلګه ده۔ فرض کړئ غواړو \LR{$\quad \Sequent !A \lor \lnot !A$}
ثابت کړو۔ د \LR{$\RightR{\lor}$} سرچپه پلي کول به له دغو دوو
اشتقاقونه څخه يو راکړي:
\begin{prooftree}
\AxiomC{}
\UnaryInf$ \fCenter !A$
\RightLabel{\RightR{\lor}}
\UnaryInf$ \fCenter !A \lor \lnot !A$
\DisplayProof\qquad\bottomAlignProof
\AxiomC{}
\UnaryInf$!A \fCenter $
\RightLabel{\RightR{\lnot}}
\UnaryInf$ \fCenter \lnot !A$
\RightLabel{\RightR{\lor}}
\UnaryInf$ \fCenter !A \lor \lnot !A$
\end{prooftree}
البته له دغو څخه هېڅ يو په لومړني سېکوېنټ نۀ پاے ته رسېږي۔ چل دا
دے چې پوه شو د انقباض قاعده د تړلے فارمول دوه نقلونه سره يو
کولے شي---او کله چې د ثبوت لټون کوو، يعنې له لاندې څخه پورته ځو،
د \LR{$!A \lor \lnot !A$} يوه کاپي په مقدمه کښې ساتلے شو، مثلاً:
\begin{prooftree}
\AxiomC{}
\UnaryInf$ \fCenter !A \lor \lnot !A, !A$
\RightLabel{\RightR{\lor}}
\UnaryInf$ \fCenter !A \lor \lnot !A, !A \lor \lnot !A$
\RightLabel{\RightR{\Contraction}}
\UnaryInf$ \fCenter !A \lor \lnot !A$
\end{prooftree}
اوس \LR{$\RightR{\lor}$} دويم ځل کارولے شو او~\LR{$\lnot !A$} هم ترلاسه کوو،
چې بشپړ اشتقاق ته رسېږي۔
\begin{prooftree}
\Axiom$!A \fCenter !A$
\RightLabel{\RightR{\lnot}}
\UnaryInf$\fCenter !A, \lnot !A$
\RightLabel{\RightR{\lor}}
\UnaryInf$\fCenter !A, !A \lor \lnot !A$
\RightLabel{\RightR{\Exchange}}
\UnaryInf$ \fCenter !A \lor \lnot !A, !A$
\RightLabel{\RightR{\lor}}
\UnaryInf$ \fCenter !A \lor \lnot !A, !A \lor \lnot !A$
\RightLabel{\RightR{\Contraction}}
\UnaryInf$ \fCenter !A \lor \lnot !A$
\end{prooftree}
\end{ex}
\olpseganchor{OLP-0075-B018}{end}

\olpseganchor{OLP-0075-B019}{start}
\begin{prob}
د لاندې سېکوېنټونو اشتقاقونه ورکړئ:
\begin{enumerate}
\item \LR{$!A \land (!B \land !C) \Sequent (!A \land !B) \land !C$}.
\item \LR{$!A \lor (!B \lor !C) \Sequent (!A \lor !B) \lor !C$}.
\item \LR{$!A \lif (!B \lif !C) \Sequent !B \lif (!A \lif !C)$}.
\item \LR{$!A \Sequent \lnot\lnot !A$}.
\end{enumerate}
\end{prob}
\olpseganchor{OLP-0075-B019}{end}

\olpseganchor{OLP-0075-B020}{start}
\begin{prob}
د لاندې سېکوېنټونو اشتقاقونه ورکړئ:
\begin{enumerate}
\item \LR{$(!A \lor !B) \lif !C \Sequent !A \lif !C$}.
\item \LR{$(!A \lif !C) \land (!B \lif !C) \Sequent (!A \lor !B) \lif !C$}.
\item \LR{$\Sequent \lnot(!A \land \lnot !A)$}.
\item \LR{$!B \lif !A \Sequent \lnot !A \lif \lnot !B$}.
\item \LR{$\Sequent (!A \lif \lnot !A) \lif \lnot !A$}.
\item \LR{$\Sequent \lnot(!A \lif !B) \lif \lnot !B$}.
\item \LR{$!A \lif !C \Sequent \lnot (!A \land \lnot !C)$}.
\item \LR{$!A \land \lnot !C \Sequent \lnot (!A \lif !C)$}.
\item \LR{$!A \lor !B, \lnot !B \Sequent !A$}.
\item \LR{$\lnot !A \lor \lnot !B \Sequent \lnot(!A \land !B)$}.
\item \LR{$\Sequent (\lnot !A \land \lnot !B) \lif\lnot(!A \lor !B)$}.
\item \LR{$\Sequent \lnot(!A \lor !B) \lif (\lnot !A \land \lnot !B)$}.
\end{enumerate}
\end{prob}
\olpseganchor{OLP-0075-B020}{end}

\olpseganchor{OLP-0075-B021}{start}
\begin{prob}
د لاندې سېکوېنټونو اشتقاقونه ورکړئ:
\begin{enumerate}
\item \LR{$\lnot(!A \lif !B) \Sequent !A$}.
\item \LR{$\lnot(!A \land !B) \Sequent \lnot !A \lor \lnot !B$}.
\item \LR{$!A \lif !B \Sequent \lnot !A \lor !B$}.
\item \LR{$\Sequent \lnot \lnot !A \lif !A$}.
\item \LR{$!A \lif !B, \lnot !A \lif !B \Sequent !B$}.
\item \LR{$(!A \land !B) \lif !C \Sequent (!A \lif !C) \lor (!B \lif !C)$}.
\item \LR{$(!A \lif !B) \lif !A \Sequent !A$}.
\item \LR{$\Sequent (!A \lif !B) \lor (!B \lif !C)$}.
\end{enumerate}
(دا ټول د~\LR{$\RightR{\Contraction}$} قاعدې ته اړتيا لري۔)
\end{prob}
\olpseganchor{OLP-0075-B021}{end}

\olpseganchor{OLP-0076-B004}{start}
\olfileid{fol}{seq}{prq}
\olpseganchor{OLP-0076-B004}{end}

\olpseganchor{OLP-0076-B005}{start}
\olsection{له سورونو سره اشتقاقونه}
\olpseganchor{OLP-0076-B005}{end}

\olpseganchor{OLP-0076-B006}{start}
\begin{ex}
د سېکوېنټ \LR{$\lexists[x][\lnot !A(x)]
\Sequent \lnot \lforall[x][!A(x)]$} يو
\LR{$\Log{LK}$}-اشتقاق ورکړئ۔
\olpseganchor{OLP-0076-B006}{end}

\olpseganchor{OLP-0076-B007}{start}
کله چې له سورونو سره کار کوو، بايد ډاډمن شو چې د ځانګړي متغير شرط
نۀ ماتوو، او کله کله د دې دپاره اړ يو چې د ځينو استنتاجونو د پلي
کولو ترتيب بدل کړو۔ په عمومي توګه دا ګټوره ده چې لومړے هغه قاعدې
وڅېړو چې د ځانګړي متغير د شرط تابع دي (په بشپړ شوي ثبوت کښې به
لاندې وي)۔ همدارنګه، غوره ده لږ وړاندې وګورو او اټکل وکړو چې
لومړنے سېکوېنټ به څنګه وي۔ زموږ په حالت کښې بايد د
\LR{$!A(a) \Sequent !A(a)$} په شان وي۔ يعنې کله چې د سورونو قاعدې
“سرچپه” کوو، بايد د \LR{$\lforall$} او \LR{$\lexists$} دواړو قاعدو دپاره
هماغه اصطلاح---چې \LR{$a$} به ئې وبولو---غوره کړو۔ که د هرې قاعدې
دپاره بېله اصطلاح غوره کړو، د \LR{$!A(a) \Sequent !A(b)$} په شان څه
ته رسېږو، چې البته د اشتقاق وړ نۀ دے۔
\olpseganchor{OLP-0076-B007}{end}

\olpseganchor{OLP-0076-B008}{start}
د معمول په شان په دې ليکلو پيل کوو:
\begin{prooftree}
\AxiomC{}
\UnaryInf$\lexists[x][\lnot !A(x)] \fCenter \lnot \lforall[x][!A(x)]$
\end{prooftree}
يا د \LeftR{\exists} قاعده پلې کولے شو او يا د
\RightR{\lnot} قاعده۔ ځکه چې \LeftR{\exists} د ځانګړي متغير د
شرط تابع ده، ښه ده چې دا ژر تر ژره وڅېړو؛ نو لومړے همدا پلې کوو۔
\begin{prooftree}
\AxiomC{}
\UnaryInf$ \lnot !A(a) \fCenter \lnot \lforall[x][!A(x)]$
\RightLabel{\LeftR{\lexists}}
\UnaryInf$ \lexists[x][\lnot !A(x)] \fCenter \lnot \lforall[x][!A(x)]$
\end{prooftree}
د \LeftR{\lnot} او \RightR{\lnot} قاعدو په سرچپه پلي کولو ترلاسه
کوو:
\begin{prooftree}
\AxiomC{}
\UnaryInf$\lforall[x][!A(x)] \fCenter !A(a)$
\RightLabel{\LeftR{\lnot}}
\UnaryInf$\lnot !A(a), \lforall[x][!A(x)] \fCenter $
\RightLabel{\LeftR{\Exchange}}
\UnaryInf$\lforall[x][!A(x)], \lnot !A(a) \fCenter $
\RightLabel{\RightR{\lnot}}
\UnaryInf$ \lnot !A(a) \fCenter \lnot \lforall[x] !A(x)$
\RightLabel{\LeftR{\lexists}}
\UnaryInf$ \lexists[x] \lnot !A(x) \fCenter \lnot \lforall[x] !A(x)$
\end{prooftree}
په دې ځاے کښې يوازينے امکان د \LeftR{\forall} قاعدې پلي کول دي۔
ځکه چې دا قاعده د ځانګړي متغير د بنديز تابع نۀ ده، نو کومه ستونزه
نشته۔ راياد کړئ چې غواړو يو لومړنے سېکوېنټ، يعنې د
\LR{$!A(a) \Sequent !A(a)$} په بڼه، ترلاسه کړو؛ نو د قاعدې د کارولو پر
مهال بايد \LR{$a$} د~\LR{$!A$} د ارګومېنټ په توګه غوره کړو۔
\begin{prooftree}
\Axiom$!A(a) \fCenter !A(a)$
\RightLabel{\LeftR{\lforall}}
\UnaryInf$\lforall[x][!A(x)] \fCenter !A(a)$
\RightLabel{\LeftR{\lnot}}
\UnaryInf$\lnot !A(a), \lforall[x][!A(x)] \fCenter $
\RightLabel{\LeftR{\Exchange}}
\UnaryInf$\lforall[x][!A(x)], \lnot !A(a) \fCenter $
\RightLabel{\RightR{\lnot}}
\UnaryInf$ \lnot !A(a) \fCenter \lnot \lforall[x][!A(x)]$
\RightLabel{\LeftR{\lexists}}
\UnaryInf$ \lexists[x][ \lnot !A(x)] \fCenter \lnot \lforall[x][!A(x)]$
\end{prooftree}
په تېره بيا له سورونو سره د کار پر مهال مهمه ده چې په دې ځاے کښې
بيا وګورو چې د ځانګړي متغير شرط خو نۀ دے مات شوے۔ ځکه چې زموږ له
کارول شوو قاعدو څخه يوازې \LeftR{\exists} د ځانګړي متغير د شرط
تابع وه، او ځانګړے متغير~\LR{$a$} د هغې په لاندې سېکوېنټ (وروستي
سېکوېنټ) کښې نۀ راځي، نو دا سم اشتقاق دے۔
\end{ex}
\olpseganchor{OLP-0076-B008}{end}

\olpseganchor{OLP-0076-B009}{start}
\begin{prob}
د لاندې سېکوېنټونو اشتقاقونه ورکړئ:
\begin{enumerate}
\item \LR{$\Sequent (\lforall[x][!A(x)] \land \lforall[y][!B(y)]) \lif
\lforall[z][(!A(z) \land !B(z))]$}.
\item \LR{$\Sequent (\lexists[x][!A(x)] \lor \lexists[y][!B(y)]) \lif
\lexists[z][(!A(z) \lor !B(z))]$}.
\item \LR{$\lforall[x][(!A(x) \lif !B)] \Sequent \lexists[y][!A(y)] \lif !B$}.
\item \LR{$\lforall[x][\lnot !A(x)] \Sequent \lnot\lexists[x][!A(x)]$}.
\item \LR{$\Sequent \lnot\lexists[x][!A(x)] \lif \lforall[x][\lnot !A(x)]$}.
\item \LR{$\Sequent \lnot\lexists[x][\lforall[y][((!A(x,y) \lif \lnot
!A(y,y)) \land (\lnot !A(y,y) \lif !A(x,y)))]]$}.
\end{enumerate}
\end{prob}
\olpseganchor{OLP-0076-B009}{end}

\olpseganchor{OLP-0076-B010}{start}
\begin{prob}
د لاندې سېکوېنټونو اشتقاقونه ورکړئ:
\begin{enumerate}
\item \LR{$\Sequent \lnot\lforall[x][!A(x)] \lif \lexists[x][\lnot!A(x)]$}.
\item \LR{$(\lforall[x][!A(x)] \lif !B) \Sequent \lexists[y][(!A(y) \lif !B)]$}.
\item \LR{$\Sequent \lexists[x][(!A(x) \lif \lforall[y][!A(y)])]$}.
\end{enumerate}
(دا ټول د~\LR{$\RightR{\Contraction}$} قاعدې ته اړتيا لري۔)
\end{prob}
\olpseganchor{OLP-0076-B010}{end}

\olpseganchor{OLP-0077-B004}{start}
\begin{editorial}
  دا برخه د سېکوېنټ حساب د اشتقاق وړتيا د اړيکې او سازګارۍ
  تعريفونه راغونډوي۔
  \textbf{د ژباړې د نسخې سمون (OLSQ-003):} په سرچينه کښې د طبيعي
  استنتاج نوم راغلے و؛ دلته د څپرکي او ورپسې تعريفونو له مخې د
  سېکوېنټ حساب نوم کارول کېږي۔
\end{editorial}
\olpseganchor{OLP-0077-B004}{end}

\olpseganchor{OLP-0077-B005}{start}
\olfileid{fol}{seq}{ptn}
\olpseganchor{OLP-0077-B005}{end}

\olpseganchor{OLP-0077-B006}{start}
\olsection{د ثبوت-تيوريکي مفاهيم}
\olpseganchor{OLP-0077-B006}{end}

\olpseganchor{OLP-0077-B007}{start}
\begin{explain}
لکه څنګه چې مو يو شمېر مهم معنايي مفاهيم (اعتبار، معنايي استلزام
او د صدق وړتيا) تعريف کړي دي، اوس ورسره سم \emph{ثبوت-تيوريکي
مفاهيم} تعريفوو۔ دا مفاهيم په جوړښتونه کښې د تړلي فارمولونه
د صدق پر بنسټ نۀ، بلکې د ځينو سېکوېنټونو د د اشتقاق وړتيا يا
د اشتقاق نۀ وړتيا پر بنسټ تعريفېږي۔ دا يوه مهمه موندنه وه چې دغه
مفاهيم سره سمون خوري۔ د دې سمون منځپانګه د \emph{سموالي} او
\emph{بشپړتيا قضيه} ده۔
\end{explain}
\olpseganchor{OLP-0077-B007}{end}

\olpseganchor{OLP-0077-B008}{start}
\begin{defn}[قضيې]
تړلے فارمول~\LR{$!A$} هله \emph{قضيه} ده چې په~\LR{$\Log{LK}$} کښې د
سېکوېنټ \LR{$\quad \Sequent !A$} اشتقاق موجود وي۔ که \LR{$!A$}
قضيه وي، \LR{$\Proves !A$} ليکو؛ او که نۀ وي، \LR{$\Proves/ !A$} ليکو۔
\end{defn}
\olpseganchor{OLP-0077-B008}{end}

\olpseganchor{OLP-0077-B009}{start}
\begin{defn}[د اشتقاق وړتيا]
تړلے فارمول~\LR{$!A$} د تړلي فارمولونه له سټ~\LR{$\Gamma$} څخه
\emph{د اشتقاق وړ} ده، يعنې \LR{$\Gamma \Proves !A$}، هله او يوازې
هله چې يو متناهي فرعي سټ~\LR{$\Gamma_0 \subseteq \Gamma$} او په
\LR{$\Gamma_0$} کښې د تړلي فارمولونه يو لړ~\LR{$\Gamma_0'$} موجود وي، داسې
چې \LR{$\Log{LK}$} د \LR{$\Gamma_0' \Sequent !A$} سېکوېنټ اشتقاق کړي۔
که \LR{$!A$} له \LR{$\Gamma$} څخه د اشتقاق وړ نۀ وي، \LR{$\Gamma \Proves/ !A$}
ليکو۔
\end{defn}
\olpseganchor{OLP-0077-B009}{end}

\olpseganchor{OLP-0077-B010}{start}
د انقباض، کمزوري کولو او تبادلې د قاعدو له امله، په~\LR{$\Gamma_0'$}
کښې د تړلي فارمولونه ترتيب او شمېر اهميت نۀ لري: که سېکوېنټ
\LR{$\Gamma_0' \Sequent !A$} د اشتقاق وړ وي، نو د هر داسې
\LR{$\Gamma_0''$} دپاره \LR{$\Gamma_0'' \Sequent !A$} هم د اشتقاق وړ دے
چې د~\LR{$\Gamma_0'$} هماغه تړلي فارمولونه لري۔ مثلاً که
\LR{$\Gamma_0 = \{!B, !C\}$} وي، نو هم
\LR{$\Gamma_0' = \tuple{!B, !B, !C}$} او هم
\LR{$\Gamma_0'' = \tuple{!C, !C, !B}$} داسې لړونه دي چې يوازې د
\LR{$\Gamma_0$} تړلي فارمولونه لري۔ که د يوه لړ لرونکے سېکوېنټ
د اشتقاق وړ وي، د بل لړ لرونکے هم دے، مثلاً:
\begin{prooftree}
  \AxiomC{}
  \Deduce$!B, !B, !C \fCenter !A$
  \RightLabel{\LeftR{\Contraction}}
  \UnaryInf$!B, !C \fCenter !A$
  \RightLabel{\LeftR{\Exchange}}
  \UnaryInf$!C, !B \fCenter !A$
  \RightLabel{\LeftR{\Weakening}}
  \UnaryInf$!C, !C, !B \fCenter !A$
\end{prooftree}
له دې وروسته به وايو چې که \LR{$\Gamma_0$} د تړلي فارمولونه يو متناهي
سټ وي، نو \LR{$\Gamma_0 \Sequent !A$} هر هغه سېکوېنټ دے چې مقدم ئې
په~\LR{$\Gamma_0$} کښې د تړلي فارمولونه يو لړ وي؛ او د اړتيا په وخت کښې
انقباضونه، تبادلې او کمزوري کول په ناويلې توګه ورګډوو۔
\olpseganchor{OLP-0077-B010}{end}

\olpseganchor{OLP-0077-B011}{start}
\begin{defn}[سازګاري]
د جملو سټ~\LR{$\Gamma$} هله او يوازې هله \emph{ناسازګار} دے چې يو
متناهي فرعي سټ~\LR{$\Gamma_0 \subseteq \Gamma$} موجود وي، داسې چې
\LR{$\Log{LK}$} د \LR{$\Gamma_0 \Sequent \quad$} سېکوېنټ اشتقاق کړي۔
که \LR{$\Gamma$} ناسازګار نۀ وي، يعنې که د هر متناهي
\LR{$\Gamma_0 \subseteq \Gamma$} دپاره \LR{$\Log{LK}$} د
\LR{$\Gamma_0 \Sequent \quad$} سېکوېنټ اشتقاق نۀ کړي، نو هغه
\emph{سازګار} بولو۔
\end{defn}
\olpseganchor{OLP-0077-B011}{end}

\olpseganchor{OLP-0077-B012}{start}
\begin{prop}[انعکاسيت]
\ollabel{prop:reflexivity}
که \LR{$!A \in \Gamma$} وي، نو \LR{$\Gamma \Proves !A$}۔
\end{prop}
\olpseganchor{OLP-0077-B012}{end}

\olpseganchor{OLP-0077-B013}{start}
\begin{proof}
لومړنے سېکوېنټ \LR{$!A \Sequent !A$} د اشتقاق وړ دے، او
\LR{$\{!A\} \subseteq \Gamma$}۔
\end{proof}
\olpseganchor{OLP-0077-B013}{end}
  
\olpseganchor{OLP-0077-B014}{start}
\begin{prop}[يکنواختي]
\ollabel{prop:monotonicity}
که \LR{$\Gamma \subseteq \Delta$} او \LR{$\Gamma \Proves !A$} وي، نو
\LR{$\Delta \Proves !A$}۔
\end{prop}
\olpseganchor{OLP-0077-B014}{end}

\olpseganchor{OLP-0077-B015}{start}
\begin{proof}
فرض کړئ \LR{$\Gamma \Proves !A$}، يعنې يو متناهي
\LR{$\Gamma_0 \subseteq \Gamma$} شته چې
\LR{$\Gamma_0 \Sequent !A$} د اشتقاق وړ دے۔ ځکه چې
\LR{$\Gamma \subseteq \Delta$}، نو \LR{$\Gamma_0$} د~\LR{$\Delta$} هم يو متناهي
فرعي سټ دے۔ له دې امله د \LR{$\Gamma_0 \Sequent !A$} اشتقاق
هم \LR{$\Delta \Proves !A$} ښيي۔
\end{proof}
\olpseganchor{OLP-0077-B015}{end}

\olpseganchor{OLP-0077-B016}{start}
\begin{prop}[انتقاليت]
\ollabel{prop:transitivity}
که \LR{$\Gamma \Proves !A$} او \LR{$\{!A\} \cup \Delta \Proves !B$} وي، نو
\LR{$\Gamma \cup \Delta \Proves !B$}۔
\end{prop}
\olpseganchor{OLP-0077-B016}{end}

\olpseganchor{OLP-0077-B017}{start}
\begin{proof}
که \LR{$\Gamma \Proves !A$} وي، يو متناهي
\LR{$\Gamma_0 \subseteq \Gamma$} او د
\LR{$\Gamma_0 \Sequent !A$} اشتقاق~\LR{$\pi_0$} شته۔ که
\LR{$\{!A\} \cup \Delta \Proves !B$} وي، نو د کوم متناهي فرعي سټ
\LR{$\Delta_0 \subseteq \Delta$} دپاره د
\LR{$!A, \Delta_0 \Sequent !B$} اشتقاق~\LR{$\pi_1$} شته۔ لاندې
اشتقاق وګورئ:
\begin{prooftree}
  \AxiomC{}
  \RightLabel{\LR{$\pi_0$}}
  \Deduce$\Gamma_0 \fCenter !A$
  \AxiomC{}
  \RightLabel{\LR{$\pi_1$}}
  \Deduce$!A, \Delta_0 \fCenter !B$
  \RightLabel{\Cut}
  \BinaryInf$\Gamma_0, \Delta_0 \fCenter !B$
\end{prooftree}
ځکه چې \LR{$\Gamma_0 \cup \Delta_0 \subseteq \Gamma \cup \Delta$}، نو
دا \LR{$\Gamma \cup \Delta \Proves !B$} ښيي۔
\end{proof}
\olpseganchor{OLP-0077-B017}{end}

\olpseganchor{OLP-0077-B018}{start}
پام وکړئ، په ځانګړې توګه د دې معنٰی دا ده چې که
\LR{$\Gamma \Proves !A$} او \LR{$!A \Proves !B$} وي، نو
\LR{$\Gamma \Proves !B$}۔ همدارنګه ترې دا هم راځي چې که
\LR{$!A_1, \dots, !A_n \Proves !B$} وي او د هر~\LR{$i$} دپاره
\LR{$\Gamma \Proves !A_i$} وي، نو \LR{$\Gamma \Proves !B$}۔
\olpseganchor{OLP-0077-B018}{end}

\olpseganchor{OLP-0077-B019}{start}
\begin{prop}
\ollabel{prop:incons}
\LR{$\Gamma$} هله او يوازې هله ناسازګار دے چې د هرې جملې~\LR{$!A$} دپاره
\LR{$\Gamma \Proves {!A}$} وي۔
\end{prop}
\olpseganchor{OLP-0077-B019}{end}

\olpseganchor{OLP-0077-B020}{start}
\begin{proof}
تمرين۔
\end{proof}
\olpseganchor{OLP-0077-B020}{end}

\olpseganchor{OLP-0077-B021}{start}
\begin{prob}
\olref[fol][seq][ptn]{prop:incons} ثابت کړئ۔
\end{prob}
\olpseganchor{OLP-0077-B021}{end}

\olpseganchor{OLP-0077-B022}{start}
\begin{prop}[متناهي شاهد]
  \ollabel{prop:proves-compact}
  \begin{enumerate}
  \item که \LR{$\Gamma \Proves !A$} وي، نو يو متناهي فرعي سټ
    \LR{$\Gamma_0 \subseteq \Gamma$} شته چې \LR{$\Gamma_0 \Proves !A$}۔
  \item که د~\LR{$\Gamma$} هر متناهي فرعي سټ سازګار وي، نو \LR{$\Gamma$}
    سازګار دے۔
  \end{enumerate}
\end{prop}
\olpseganchor{OLP-0077-B022}{end}

\olpseganchor{OLP-0077-B023}{start}
\begin{proof}
  \begin{enumerate}
    \item که \LR{$\Gamma \Proves !A$} وي، نو يو متناهي فرعي سټ
      \LR{$\Gamma_0 \subseteq \Gamma$} شته، داسې چې د
      \LR{$\Gamma_0 \Sequent !A$} سېکوېنټ اشتقاق لري۔ له دې امله
      \LR{$\Gamma_0 \Proves !A$}۔
    \item که \LR{$\Gamma$} ناسازګار وي، نو يو متناهي فرعي سټ
      \LR{$\Gamma_0 \subseteq \Gamma$} شته، داسې چې \LR{$\Log{LK}$} د
      \LR{$\Gamma_0 \Sequent \quad$} سېکوېنټ اشتقاق کړي۔ خو بيا
      \LR{$\Gamma_0$} د~\LR{$\Gamma$} يو ناسازګار متناهي فرعي سټ دے۔
  \end{enumerate}
\end{proof}
\olpseganchor{OLP-0077-B023}{end}

\olpseganchor{OLP-0078-B004}{start}
\olfileid{fol}{seq}{prv}
\olpseganchor{OLP-0078-B004}{end}

\olpseganchor{OLP-0078-B005}{start}
\olsection{د اشتقاق وړتيا او سازګاري}
\olpseganchor{OLP-0078-B005}{end}

\olpseganchor{OLP-0078-B006}{start}
اوس به د د اشتقاق وړتيا اړيکې يو شمېر خاصيتونه ثابت کړو۔ دا
خاصيتونه په خپلواک ډول هم په زړه پورې دي، خو هر يو به د بشپړتيا د
قضيې په ثبوت کښې ونډه ولري۔
\olpseganchor{OLP-0078-B006}{end}

\olpseganchor{OLP-0078-B007}{start}
\begin{prop}\ollabel{prop:provability-contr}
  که \LR{$\Gamma \Proves !A$} وي او \LR{$\Gamma \cup \{!A\}$} ناسازګار وي،
  نو \LR{$\Gamma$} ناسازګار دے۔
\end{prop}
\olpseganchor{OLP-0078-B007}{end}

\olpseganchor{OLP-0078-B008}{start}
\begin{proof}
داسې متناهي \LR{$\Gamma_0$} او \LR{$\Gamma_1 \subseteq \Gamma$} شته چې
\LR{$\Log{LK}$} د \LR{$\Gamma_0 \Sequent !A$} او
\LR{$!A, \Gamma_1 \Sequent \quad$} سېکوېنټونه اشتقاق کړي۔
د \LR{$\Gamma_0 \Sequent !A$} د \LR{$\Log{LK}$}-اشتقاق نوم~\LR{$\pi_0$}
او د \LR{$\Gamma_1, !A \Sequent \quad$} د
\LR{$\Log{LK}$}-اشتقاق نوم~\LR{$\pi_1$} کېږدو۔ بيا لاندې
اشتقاق کولے شو:
\begin{prooftree}
\AxiomC{}
\RightLabel{\LR{$\pi_0$}}
\Deduce$ \Gamma_0 \fCenter !A $
\AxiomC{}
\RightLabel{\LR{$\pi_1$}}
\Deduce$!A, \Gamma_1 \fCenter $
\RightLabel{\Cut}
\BinaryInf$ \Gamma_0,\Gamma_1 \fCenter $
\end{prooftree}
\olpseganchor{OLP-0078-B008}{end}

\olpseganchor{OLP-0078-B009}{start}
ځکه چې \LR{$\Gamma_0 \subseteq \Gamma$} او
\LR{$\Gamma_1 \subseteq \Gamma$}، نو
\LR{$\Gamma_0 \cup \Gamma_1 \subseteq \Gamma$}؛ له دې امله \LR{$\Gamma$}
ناسازګار دے۔
\end{proof}
\olpseganchor{OLP-0078-B009}{end}

\olpseganchor{OLP-0078-B010}{start}
\begin{prop}
\ollabel{prop:prov-incons}
\LR{$\Gamma \Proves !A$} هله او يوازې هله چې
\LR{$\Gamma \cup \{\lnot !A\}$} ناسازګار وي۔
\end{prop}
\olpseganchor{OLP-0078-B010}{end}

\olpseganchor{OLP-0078-B011}{start}
\begin{proof}
لومړے فرض کړئ \LR{$\Gamma \Proves !A$}، يعنې د
\LR{$\Gamma \Sequent !A$} اشتقاق~\LR{$\pi_0$} شته۔ د
\LR{$\LeftR{\lnot}$} قاعدې په ورزياتولو د
\LR{$\lnot !A, \Gamma \Sequent \quad$} اشتقاق ترلاسه کوو؛
يعنې \LR{$\Gamma \cup \{\lnot !A\}$} ناسازګار دے۔
\olpseganchor{OLP-0078-B011}{end}

\olpseganchor{OLP-0078-B012}{start}
که \LR{$\Gamma \cup \{\lnot !A\}$} ناسازګار وي، نو د
\LR{$\lnot !A, \Gamma \Sequent \quad$} اشتقاق~\LR{$\pi_1$} شته۔
لاندې د \LR{$\Gamma \Sequent !A$} اشتقاق دے:
\begin{prooftree}
  \Axiom$!A \fCenter !A$
  \RightLabel{\RightR{\lnot}}
  \UnaryInf$\fCenter !A, \lnot !A$
  \AxiomC{}
  \RightLabel{\LR{$\pi_1$}}
  \Deduce$\lnot !A, \Gamma \fCenter$
  \RightLabel{\Cut}
  \BinaryInf$\Gamma \fCenter !A$
\end{prooftree}
\end{proof}
\olpseganchor{OLP-0078-B012}{end}

\olpseganchor{OLP-0078-B013}{start}
\begin{prob}
ثابت کړئ چې \LR{$\Gamma \Proves \lnot !A$} هله او يوازې هله چې
\LR{$\Gamma \cup \{!A\}$} ناسازګار وي۔
\end{prob}
\olpseganchor{OLP-0078-B013}{end}

\olpseganchor{OLP-0078-B014}{start}
\begin{prop}\ollabel{prop:explicit-inc}
  که \LR{$\Gamma \Proves !A$} او \LR{$\lnot !A \in \Gamma$} وي، نو \LR{$\Gamma$}
  ناسازګار دے۔
\end{prop}
\olpseganchor{OLP-0078-B014}{end}

\olpseganchor{OLP-0078-B015}{start}
\begin{proof}
  فرض کړئ \LR{$\Gamma \Proves !A$} او \LR{$\lnot !A \in \Gamma$}۔ بيا د کوم
  سېکوېنټ \LR{$\Gamma_0 \Sequent !A$} اشتقاق~\LR{$\pi$} شته۔
  سېکوېنټ \LR{$\lnot !A, \Gamma_0 \Sequent \quad$} هم د اشتقاق وړ دے:
  \begin{prooftree}
    \AxiomC{}
    \RightLabel{\LR{$\pi$}}
    \Deduce$\Gamma_0 \fCenter !A$
    \Axiom$!A \fCenter !A$
    \RightLabel{\LeftR{\lnot}}
    \UnaryInf$\lnot !A, !A \fCenter$
    \RightLabel{\LeftR{\Exchange}}
    \UnaryInf$!A, \lnot !A \fCenter$
    \RightLabel{\Cut}
    \BinaryInf$\Gamma_0, \lnot !A \fCenter$
  \end{prooftree}
  ځکه چې \LR{$\lnot !A \in \Gamma$} او \LR{$\Gamma_0 \subseteq \Gamma$}،
  نو دا ښيي چې \LR{$\Gamma$} ناسازګار دے۔
\end{proof}
\olpseganchor{OLP-0078-B015}{end}

\olpseganchor{OLP-0078-B016}{start}
\begin{prop}\ollabel{prop:provability-exhaustive}
  که \LR{$\Gamma \cup \{!A\}$} او \LR{$\Gamma \cup \{\lnot !A\}$} دواړه
  ناسازګار وي، نو \LR{$\Gamma$} ناسازګار دے۔
\end{prop}
\olpseganchor{OLP-0078-B016}{end}

\olpseganchor{OLP-0078-B017}{start}
\begin{proof}
داسې متناهي سټونه \LR{$\Gamma_0 \subseteq \Gamma$} او
\LR{$\Gamma_1 \subseteq \Gamma$} او داسې
\LR{$\Log{LK}$}-اشتقاقونه \LR{$\pi_0$} او \LR{$\pi_1$} شته چې په همدې ترتيب
د \LR{$!A, \Gamma_0 \Sequent \quad$} او
\LR{$\lnot !A, \Gamma_1 \Sequent \quad$} اشتقاقونه دي۔ بيا لاندې
اشتقاق کولے شو:
\begin{prooftree}
\AxiomC{}
\RightLabel{\LR{$\pi_0$}}
\Deduce$ !A, \Gamma_0 \fCenter $
\RightLabel{\RightR{\lnot}}
\UnaryInf$ \Gamma_0 \fCenter \lnot !A$
\AxiomC{}
\RightLabel{\LR{$\pi_1$}}
\Deduce$\lnot !A, \Gamma_1 \fCenter  $
\RightLabel{\Cut}
\BinaryInf$ \Gamma_0, \Gamma_1 \fCenter $
\end{prooftree}
ځکه چې \LR{$\Gamma_0 \subseteq \Gamma$} او \LR{$\Gamma_1 \subseteq \Gamma$}،
نو \LR{$\Gamma_0 \cup \Gamma_1 \subseteq \Gamma$}۔ له دې امله \LR{$\Gamma$}
ناسازګار دے۔
\end{proof}
\olpseganchor{OLP-0078-B017}{end}

\olpseganchor{OLP-0079-B005}{start}
\olfileid{fol}{seq}{ppr}
\olpseganchor{OLP-0079-B005}{end}

\olpseganchor{OLP-0079-B006}{start}
\olsection{د اشتقاق وړتيا او قضييز نښلوونکي}
\olpseganchor{OLP-0079-B006}{end}

\olpseganchor{OLP-0079-B007}{start}
\begin{explain}
  موږ ثابتوو چې د سېکوېنټ حساب د د اشتقاق وړتيا اړيکه~\LR{$\Proves$}
  دومره پياوړې ده چې د قضييزو نښلوونکو په اړه ځينې بنسټيز حقيقتونه
  ثابت کړي، لکه \LR{$!A \land !B \Proves !A$} او
  \LR{$!A, !A \lif !B \Proves !B$} (مودوس پونېنس)۔ دغه حقيقتونه د
  بشپړتيا د قضيې د ثبوت دپاره اړين دي۔
\end{explain}
\olpseganchor{OLP-0079-B007}{end}

\olpseganchor{OLP-0079-B008}{start}
\begin{prop}\ollabel{prop:provability-land}
  \begin{enumerate}
  \item \ollabel{prop:provability-land-left} هم
    \LR{$!A \land !B \Proves !A$} او هم
    \LR{$!A \land !B \Proves !B$}۔
  \item \ollabel{prop:provability-land-right}
    \LR{$!A, !B \Proves !A \land !B$}۔
  \end{enumerate}
\end{prop}
\olpseganchor{OLP-0079-B008}{end}

\olpseganchor{OLP-0079-B009}{start}
\begin{proof}
  \begin{enumerate}
  \item دواړه سېکوېنټونه \LR{$!A \land !B \Sequent !A$} او
    \LR{$!A \land !B \Sequent !B$} د اشتقاق وړ دي:
    \begin{prooftree}
      \Axiom$!A \fCenter !A$
      \RightLabel{\LeftR{\land}}
      \UnaryInf$!A \land !B \fCenter !A$
      \DisplayProof\qquad\bottomAlignProof
      \Axiom$!B \fCenter !B$
      \RightLabel{\LeftR{\land}}
      \UnaryInf$!A \land !B \fCenter !B$
    \end{prooftree}
    \item دا د سېکوېنټ \LR{$!A, !B \Sequent !A \land !B$} يو
    اشتقاق دے۔
    \textbf{د ژباړې د نسخې سمون (OLSQ-004):} د سرچينې په ونه کښې
    د عطف-ښۍ قاعدې دواړو مقدمو يو شان مقدم نۀ درلود؛ دلته د
    جوړښتي استنتاجونو دوه ګونې کرښې دواړه مقدمې له پايلې سره سموي۔
    \begin{prooftree}
      \Axiom$!A \fCenter !A$
      \doubleLine
      \UnaryInf$!A, !B \fCenter !A$
      \Axiom$!B \fCenter !B$
      \doubleLine
      \UnaryInf$!A, !B \fCenter !B$
      \RightLabel{\RightR{\land}}
      \BinaryInf$!A, !B \fCenter !A \land !B$
    \end{prooftree}
  \end{enumerate}
\end{proof}
\olpseganchor{OLP-0079-B009}{end}

\olpseganchor{OLP-0079-B010}{start}
\begin{prop}\ollabel{prop:provability-lor}
  \begin{enumerate}
  \item \LR{$!A \lor !B, \lnot !A, \lnot !B$} ناسازګار دے۔
  \item هم \LR{$!A \Proves !A \lor !B$} او هم
    \LR{$!B \Proves !A \lor !B$}۔
  \end{enumerate}
\end{prop}
\olpseganchor{OLP-0079-B010}{end}

\olpseganchor{OLP-0079-B011}{start}
\begin{proof}
  \begin{enumerate}
  \item موږ د سېکوېنټ
    \LR{$!A \lor !B, \lnot !A, \lnot !B \Sequent$} يو
    اشتقاق ورکوو:
    \begin{prooftree}
      \Axiom$!A \fCenter !A$
      \RightLabel{\LeftR{\lnot}}
      \UnaryInf$\lnot !A, !A \fCenter$
      \doubleLine
      \UnaryInf$!A, \lnot !A, \lnot !B \fCenter$
      \Axiom$!B \fCenter !B$
      \RightLabel{\LeftR{\lnot}}
      \UnaryInf$\lnot !B, !B \fCenter$
      \doubleLine
      \UnaryInf$!B, \lnot !A, \lnot !B \fCenter$
      \RightLabel{\LeftR{\lor}}
      \BinaryInf$ !A \lor !B, \lnot !A, \lnot !B \fCenter $
    \end{prooftree}
    (راياد کړئ چې د استنتاج دوه ګونې کرښې څو کمزوري کول، انقباضونه
    او تبادلې ښيي۔)
  \item دواړه سېکوېنټونه \LR{$!A \Sequent !A \lor !B$} او
    \LR{$!B \Sequent !A \lor !B$} اشتقاقونه لري:
    \begin{prooftree}
      \Axiom$!A \fCenter !A$
      \RightLabel{\RightR{\lor}}
      \UnaryInf$!A \fCenter !A \lor !B$
      \DisplayProof\qquad\bottomAlignProof
      \Axiom$!B \fCenter !B$
      \RightLabel{\RightR{\lor}}
      \UnaryInf$!B \fCenter !A \lor !B$
    \end{prooftree}
  \end{enumerate}
\end{proof}
\olpseganchor{OLP-0079-B011}{end}

\olpseganchor{OLP-0079-B012}{start}
\begin{prop}\ollabel{prop:provability-lif}
  \begin{enumerate}
  \item \ollabel{prop:provability-lif-left}
    \LR{$!A, !A \lif !B \Proves !B$}۔
  \item \ollabel{prop:provability-lif-right}
    هم \LR{$\lnot !A \Proves !A \lif !B$} او هم
    \LR{$!B \Proves !A \lif !B$}۔
  \end{enumerate}
\end{prop}
\olpseganchor{OLP-0079-B012}{end}

\olpseganchor{OLP-0079-B013}{start}
\begin{proof}
  \begin{enumerate}
    \item سېکوېنټ \LR{$!A \lif !B, !A \Sequent !B$} د اشتقاق وړ دے:
      \begin{prooftree}
        \Axiom$!A \fCenter !A$
        \Axiom$!B \fCenter !B$
        \RightLabel{\LeftR{\lif}}
        \BinaryInf$!A \lif !B, !A  \fCenter !B$
      \end{prooftree}
    \item دواړه سېکوېنټونه
      \LR{$\lnot !A \Sequent !A \lif !B$} او
      \LR{$!B \Sequent !A \lif !B$} د اشتقاق وړ دي:
      \begin{prooftree}
        \Axiom$!A \fCenter !A$
        \RightLabel{\LeftR{\lnot}}
        \UnaryInf$\lnot !A, !A \fCenter$
        \RightLabel{\LeftR{\Exchange}}
        \UnaryInf$!A, \lnot !A \fCenter$
        \RightLabel{\RightR{\Weakening}}
        \UnaryInf$!A, \lnot !A \fCenter !B$
        \RightLabel{\RightR{\lif}}
        \UnaryInf$\lnot !A \fCenter !A \lif !B$
        \DisplayProof\qquad\bottomAlignProof
        \Axiom$!B \fCenter !B$
        \RightLabel{\LeftR{\Weakening}}
        \UnaryInf$!A, !B \fCenter !B$
        \RightLabel{\RightR{\lif}}
        \UnaryInf$!B \fCenter !A \lif !B$
      \end{prooftree}
  \end{enumerate}
\end{proof}
\olpseganchor{OLP-0079-B013}{end}

\olpseganchor{OLP-0080-B004}{start}
\olfileid{fol}{seq}{qpr}
\olsection{د اشتقاق وړتيا او سورونه}
\olpseganchor{OLP-0080-B004}{end}

\olpseganchor{OLP-0080-B005}{start}
\begin{explain}
  د بشپړتيا قضيه دا هم غواړي چې د سېکوېنټ حساب قاعدې په دې برخه
  کښې د~\LR{$\Proves$} په اړه ثابت شوي حقيقتونه راکړي۔
\end{explain}
\olpseganchor{OLP-0080-B005}{end}

\olpseganchor{OLP-0080-B006}{start}
\begin{thm}
\ollabel{thm:strong-generalization} که \LR{$c$} داسې ثابت وي چې په
\LR{$\Gamma$} يا \LR{$!A(x)$} کښې نۀ راځي، او \LR{$\Gamma \Proves !A(c)$} وي، نو
\LR{$\Gamma \Proves \lforall[x][!A(x)]$}۔
\end{thm}
\olpseganchor{OLP-0080-B006}{end}

\olpseganchor{OLP-0080-B007}{start}
\begin{proof}
د کوم متناهي \LR{$\Gamma_0 \subseteq \Gamma$} دپاره
\LR{$\Gamma_0 \Sequent !A(c)$} د \LR{$\Log{LK}$}-اشتقاق نوم~\LR{$\pi_0$}
کېږدئ۔ د \LR{$\RightR{\lforall}$} استنتاج په ورزياتولو د
\LR{$\Gamma_0 \Sequent \lforall[x][!A(x)]$} اشتقاق ترلاسه کوو؛
ځکه \LR{$c$} په \LR{$\Gamma$} يا \LR{$!A(x)$} کښې نۀ راځي، نو د ځانګړي متغير شرط
پوره کېږي۔
\end{proof}
\olpseganchor{OLP-0080-B007}{end}

\olpseganchor{OLP-0080-B008}{start}
\begin{prop}
\ollabel{prop:provability-quantifiers}
\begin{enumerate}
\item \LR{$!A(t) \Proves \lexists[x][!A(x)]$}.
\item \LR{$\lforall[x][!A(x)] \Proves !A(t)$}.
\end{enumerate}
\end{prop}
\olpseganchor{OLP-0080-B008}{end}

\olpseganchor{OLP-0080-B009}{start}
\begin{proof}
\begin{enumerate}
\item سېکوېنټ
  \LR{$!A(t) \Sequent \lexists[x][!A(x)]$} د اشتقاق وړ دے:
  \begin{prooftree}
    \Axiom$!A(t) \fCenter !A(t)$
    \RightLabel{\RightR{\lexists}}
    \UnaryInf$!A(t) \fCenter \lexists[x][!A(x)]$
    \end{prooftree}
\item سېکوېنټ
  \LR{$\lforall[x][!A(x)] \Sequent !A(t)$} د اشتقاق وړ دے:
  \begin{prooftree}
    \Axiom$!A(t) \fCenter !A(t)$
    \RightLabel{\LeftR{\lforall}}
    \UnaryInf$\lforall[x][!A(x)] \fCenter !A(t)$
    \end{prooftree}
\end{enumerate}
\end{proof}
\olpseganchor{OLP-0080-B009}{end}

\olpseganchor{OLP-0081-B004}{start}
\olfileid{fol}{seq}{sou}
\olpseganchor{OLP-0081-B004}{end}

\olpseganchor{OLP-0081-B005}{start}
\olsection{سموالے}
\olpseganchor{OLP-0081-B005}{end}

\olpseganchor{OLP-0081-B006}{start}
\begin{explain}
د سېکوېنټ حساب په شان اشتقاق نظام هله \emph{سم} دے چې
هغه څيزونه اشتقاق نۀ کړي چې په حقيقت کښې نۀ وي۔ نو سموالے د
اشتقاق نظامونو د تضمين شوې خونديتوب يو ډول خاصيت دے۔
د دې له مخې چې کوم ثبوت-تيوريکي خاصيت تر بحث لاندې وي، مثلاً غواړو
پوه شو چې:
\begin{enumerate}
\item هر د اشتقاق وړ~\LR{$!A$} معتبر ده؛
\item که تړلے فارمول له ځينو نورو څخه د اشتقاق وړ وي، د هغو
  معنايي پايله هم ده؛
\item که د تړلي فارمولونه يو سټ ناسازګار وي، د صدق وړ نۀ دے۔
\end{enumerate}
دا د اشتقاق نظام مهم خاصيتونه دي۔ که له دغو څخه کوم يو
سم نۀ وي، اشتقاق نظام نيمګړے دے---له حده زيات څيزونه به
اشتقاق کوي۔ له دې امله د اشتقاق نظام سموالے ثابتول
تر ټولو زيات اهميت لري۔
\olpseganchor{OLP-0081-B006}{end}

\olpseganchor{OLP-0081-B007}{start}
ځکه دغه ټول ثبوت-تيوريکي خاصيتونه په سېکوېنټ حساب کښې د ځينو
سېکوېنټونو د د اشتقاق وړتيا له لارې تعريف شوي دي، د پورته
(۱)--(۳) ثابتول غواړي چې د د اشتقاق وړ سېکوېنټونو د معنايي
خاصيتونو په اړه يو څه ثابت کړو۔ لومړے به تعريف کړو چې د سېکوېنټ
\emph{معتبر} کېدل څه معنٰی لري، او بيا به وښيو چې هر
د اشتقاق وړ سېکوېنټ معتبر دے۔ وروسته (۱)--(۳) د دې پايلې د
لازمو نتيجو په توګه راځي۔
\end{explain}
\olpseganchor{OLP-0081-B007}{end}

\olpseganchor{OLP-0081-B008}{start}
\begin{defn}
جوړښت~\LR{$\Struct
  M$} يو سېکوېنټ
\LR{$\Gamma \Sequent \Delta$} هله او يوازې هله \emph{پوره کوي} چې يا
د کوم \LR{$!A \in \Gamma$} دپاره
\LR{$\Sat/{M}{!A}$} وي، يا د کوم
\LR{$!A \in \Delta$} دپاره
\LR{$\Sat{M}{!A}$} وي۔
\olpseganchor{OLP-0081-B008}{end}

\olpseganchor{OLP-0081-B009}{start}
يو سېکوېنټ هله او يوازې هله \emph{معتبر} دے چې هر
جوړښت~\LR{$\Struct
  M$} هغه پوره کړي۔
\end{defn}
\olpseganchor{OLP-0081-B009}{end}

\olpseganchor{OLP-0081-B010}{start}
\begin{thm}[سموالے]
\ollabel{thm:sequent-soundness} که \LR{$\Log{LK}$} د
\LR{$\Theta \Sequent \Xi$} سېکوېنټ اشتقاق کړي، نو
\LR{$\Theta \Sequent \Xi$} معتبر دے۔
\end{thm}
\olpseganchor{OLP-0081-B010}{end}

\olpseganchor{OLP-0081-B011}{start}
\begin{proof}
فرض کړئ \LR{$\pi$} د \LR{$\Theta \Sequent \Xi$} اشتقاق دے۔ موږ
په~\LR{$\pi$} کښې د استنتاجونو پر شمېر~\LR{$n$} استقرا کوو۔
\olpseganchor{OLP-0081-B011}{end}

\olpseganchor{OLP-0081-B012}{start}
که د استنتاجونو شمېر~\LR{$0$} وي، نو \LR{$\pi$} يوازې له يوه لومړني
سېکوېنټ څخه جوړ دے۔ هر لومړنے سېکوېنټ
\LR{$!A \Sequent !A$} په څرګند ډول معتبر دے؛ ځکه د هر
\LR{$\Struct M$} دپاره يا
\LR{$\Sat/{M}{!A}$} وي يا
\LR{$\Sat{M}{!A}$}۔
\olpseganchor{OLP-0081-B012}{end}

\olpseganchor{OLP-0081-B013}{start}
که د استنتاجونو شمېر له~0 څخه زيات وي، د تر ټولو لاندې استنتاج
د ډول له مخې حالتونه بېلوو۔ د استقرا د فرضيې له مخې منلے شو چې د
دغه استنتاج مقدمې معتبرې دي؛ ځکه د هرې مقدمې په اشتقاق
کښې د استنتاجونو شمېر له~\LR{$n$} څخه کم دے۔
\olpseganchor{OLP-0081-B013}{end}

\olpseganchor{OLP-0081-B014}{start}
لومړے هغه ممکن استنتاجونه څېړو چې يوازې يوه مقدمه لري۔
\begin{enumerate}
\item وروستنے استنتاج کمزوري کول دي۔ بيا
  \LR{$\Theta \Sequent \Xi$} يا \LR{$!A, \Gamma \Sequent \Delta$} دے
  (که وروستنے استنتاج \LeftR{\Weakening} وي) يا
  \LR{$\Gamma \Sequent \Delta, !A$} دے
  (که \RightR{\Weakening} وي)، او اشتقاق له لاندې دواړو
  بڼو څخه په يوه پاے ته رسېږي:
  \begin{prooftree}
    \AxiomC{}
    \Deduce$\Gamma \fCenter \Delta$
    \RightLabel{\LeftR{\Weakening}}
    \UnaryInf$!A, \Gamma \fCenter \Delta$
    \DisplayProof\qquad\bottomAlignProof
    \AxiomC{}
    \Deduce$\Gamma \fCenter \Delta$
    \RightLabel{\RightR{\Weakening}}
    \UnaryInf$\Gamma \fCenter \Delta, !A$
  \end{prooftree}
  د استقرا د فرضيې له مخې \LR{$\Gamma \Sequent \Delta$} معتبر دے؛ يعنې
  د هر جوړښت~\LR{$\Struct{M}$} دپاره يا کوم
  \LR{$!C \in \Gamma$} شته چې
  \LR{$\Sat/{M}{!C}$} وي، يا کوم
  \LR{$!C \in \Delta$} شته چې
  \LR{$\Sat{M}{!C}$} وي۔
\olpseganchor{OLP-0081-B014}{end}
  
\olpseganchor{OLP-0081-B015}{start}
  که د کوم \LR{$!C \in \Gamma$} دپاره
  \LR{$\Sat/{M}{!C}$} وي، نو
  \LR{$!C \in \Theta$} هم دے؛ ځکه \LR{$\Theta = !A, \Gamma$}، نو د کوم
  \LR{$!C \in \Theta$} دپاره
  \LR{$\Sat/{M}{!C}$} دے۔ همدارنګه، که د کوم
  \LR{$!C \in \Delta$} دپاره
  \LR{$\Sat{M}{!C}$} وي، ځکه
  \LR{$!C \in \Xi$}، نو د کوم \LR{$!C \in \Xi$} دپاره
  \LR{$\Sat{M}{!C}$} دے۔ له دې امله
  \LR{$\Theta \Sequent \Xi$} معتبر دے۔
\item وروستنے استنتاج \LeftR{\lnot} دے: بيا د وروستي استنتاج
  مقدمه \LR{$\Gamma \Sequent \Delta, !A$} او پايله
  \LR{$\lnot !A, \Gamma \Sequent \Delta$} ده، يعنې اشتقاق داسې
  پاے ته رسېږي:
  \begin{prooftree}
    \AxiomC{}
    \Deduce$\Gamma \fCenter \Delta, !A$
    \RightLabel{\LeftR{\lnot}}
    \UnaryInf$\lnot !A, \Gamma \fCenter \Delta$
  \end{prooftree}
  او \LR{$\Theta = \lnot !A, \Gamma$}، په داسې حال کښې چې
  \LR{$\Xi = \Delta$}۔
\olpseganchor{OLP-0081-B015}{end}

\olpseganchor{OLP-0081-B016}{start}
  د استقرا فرضيه وايي چې \LR{$\Gamma \Sequent \Delta, !A$} معتبر دے؛
  يعنې د هر \LR{$\Struct{M}$} دپاره يا
  (الف) د کوم \LR{$!C \in \Gamma$} دپاره
  \LR{$\Sat/{M}{!C}$}، يا
  (ب) د کوم \LR{$!C \in \Delta$} دپاره
  \LR{$\Sat{M}{!C}$}، يا
  (ج) \LR{$\Sat{M}{!A}$} وي۔ غواړو وښيو چې
  \LR{$\Theta \Sequent \Xi$} هم معتبر دے۔ فرض کړئ
  \LR{$\Struct{M}$} جوړښت دے۔ که (الف) سم وي، نو داسې
  \LR{$!C \in \Gamma$} شته چې
  \LR{$\Sat/{M}{!C}$}، خو
  \LR{$!C \in \Theta$} هم دے۔ که (ب) سم وي، نو داسې
  \LR{$!C \in \Delta$} شته چې
  \LR{$\Sat{M}{!C}$}، خو
  \LR{$!C \in \Xi$} هم دے۔ په پاے کښې، که
  \LR{$\Sat{M}{!A}$} وي، نو
  \LR{$\Sat/{M}{\lnot !A}$} دے۔ ځکه
  \LR{$\lnot !A \in \Theta$}، نو داسې \LR{$!C \in \Theta$} شته چې
  \LR{$\Sat/{M}{!C}$} وي۔ له دې امله
  \LR{$\Theta \Sequent \Xi$} معتبر دے۔
\item وروستنے استنتاج \RightR{\lnot} دے: تمرين۔
\item وروستنے استنتاج \LeftR{\land} دے: دوه بڼې شته؛
  \LR{$!A \land !B$} په کيڼ لوري کښې د مقدمې له \LR{$!A$} يا \LR{$!B$} څخه
  استنتاج کېدے شي۔ په لومړي حالت کښې \LR{$\pi$} داسې پاے ته رسېږي:
  \begin{prooftree}
    \AxiomC{}
    \Deduce$!A, \Gamma \fCenter \Delta$
    \RightLabel{\LeftR{\land}}
    \UnaryInf$!A \land !B, \Gamma \fCenter \Delta$
  \end{prooftree}
  او \LR{$\Theta = !A \land !B, \Gamma$}، په داسې حال کښې چې
  \LR{$\Xi = \Delta$}۔ يو
  جوړښت~\LR{$\Struct
    M$} په پام کښې ونيسئ۔ ځکه د
  استقرا د فرضيې له مخې \LR{$!A, \Gamma \Sequent \Delta$} معتبر دے،
  يا (الف) \LR{$\Sat/{M}{!A}$}،
  يا (ب) د کوم \LR{$!C \in \Gamma$} دپاره
  \LR{$\Sat/{M}{!C}$}، يا
  (ج) د کوم \LR{$!C \in \Delta$} دپاره
  \LR{$\Sat{M}{!C}$} وي۔
  په (الف) حالت کښې
  \LR{$\Sat/{M}{!A \land !B}$}، نو داسې
  \LR{$!C \in \Theta$} (يعنې \LR{$!A \land !B$}) شته چې
  \LR{$\Sat/{M}{!C}$} وي۔ په (ب) حالت کښې داسې
  \LR{$!C \in \Gamma$} شته چې
  \LR{$\Sat/{M}{!C}$}، او
  \LR{$!C \in \Theta$} هم دے۔ په (ج) حالت کښې داسې
  \LR{$!C \in \Delta$} شته چې
  \LR{$\Sat{M}{!C}$}، او ځکه
  \LR{$\Xi = \Delta$}، \LR{$!C \in \Xi$} هم دے۔ نو په هر حالت کښې
  \LR{$\Struct M$} د
  \LR{$!A \land !B, \Gamma \Sequent \Delta$} سېکوېنټ پوره کوي۔
  \textbf{د ژباړې د نسخې سمون (OLSQ-005):} د سرچينې په ورپسې
  پايله کښې عطف لرونکې مقدمه پاتې شوې وه؛ دلته د همدغه ثابت شوي
  بشپړ سېکوېنټ اعتبار ليکل کېږي۔
  ځکه \LR{$\Struct{M}$} اختياري و،
  \LR{$!A \land !B, \Gamma \Sequent \Delta$} معتبر دے۔ هغه حالت چې
  \LR{$!A \land !B$} له \LR{$!B$} څخه استنتاج کېږي په همدې ډول، د \LR{$!A$} پر
  ځاے د \LR{$!B$} په راوړلو، حل کېږي۔
\item وروستنے استنتاج \RightR{\lor} دے: دوه بڼې شته؛
  \LR{$!A \lor !B$} په ښي لوري کښې د مقدمې له \LR{$!A$} يا \LR{$!B$} څخه
  استنتاج کېدے شي۔ په لومړي حالت کښې \LR{$\pi$} داسې پاے ته رسېږي:
  \begin{prooftree}
    \AxiomC{}
    \Deduce$\Gamma \fCenter \Delta, !A$
    \RightLabel{\RightR{\lor}}
    \UnaryInf$\Gamma \fCenter \Delta, !A \lor !B$
  \end{prooftree}
  اوس \LR{$\Theta = \Gamma$} او \LR{$\Xi = \Delta, !A \lor !B$} دي۔ يو
  جوړښت~\LR{$\Struct{M}$} په پام کښې ونيسئ۔ ځکه
  \LR{$\Gamma \Sequent \Delta, !A$} معتبر دے، يا
  (الف) \LR{$\Sat{M}{!A}$}،
  يا (ب) د کوم \LR{$!C \in \Gamma$} دپاره
  \LR{$\Sat/{M}{!C}$}، يا
  (ج) د کوم \LR{$!C \in \Delta$} دپاره
  \LR{$\Sat{M}{!C}$} وي۔
  په (الف) حالت کښې
  \LR{$\Sat{M}{!A \lor !B}$}۔ په (ب) حالت کښې
  داسې \LR{$!C \in \Gamma$} شته چې
  \LR{$\Sat/{M}{!C}$}۔ په (ج) حالت کښې داسې
  \LR{$!C \in \Delta$} شته چې
  \LR{$\Sat{M}{!C}$}۔ نو په هر حالت کښې
  \LR{$\Struct{M}$} د
  \LR{$\Gamma \Sequent \Delta, !A \lor !B$}، يعنې
  \LR{$\Theta \Sequent \Xi$}، پوره کوي۔ ځکه
  \LR{$\Struct{M}$} اختياري و،
  \LR{$\Theta \Sequent \Xi$} معتبر دے۔ هغه حالت چې \LR{$!A \lor !B$} له
  \LR{$!B$} څخه استنتاج کېږي په همدې ډول، د \LR{$!A$} پر ځاے د \LR{$!B$} په
  راوړلو، حل کېږي۔
\item وروستنے استنتاج \RightR{\lif} دے: بيا \LR{$\pi$} داسې پاے ته رسېږي:
  \begin{prooftree}
    \AxiomC{}
    \Deduce$!A, \Gamma \fCenter \Delta, !B$
    \RightLabel{\RightR{\lif}}
    \UnaryInf$\Gamma \fCenter \Delta, !A \lif !B$
  \end{prooftree}
  بيا هم د استقرا فرضيه وايي چې مقدمه معتبره ده؛ غواړو وښيو چې
  پايله هم معتبره ده۔ يو اختياري
  \LR{$\Struct{M}$} ونيسئ۔ ځکه
  \LR{$!A, \Gamma \Sequent \Delta, !B$} معتبر دے، لږ تر لږه يو له دغو
  حالتونو څخه سم دے: (الف)
  \LR{$\Sat/{M}{!A}$}، (ب)
  \LR{$\Sat{M}{!B}$}، (ج) د کوم
  \LR{$~!C \in \Gamma$} دپاره
  \LR{$\Sat/{M}{!C}$}، يا (د) د کوم
  \LR{$!C \in \Delta$} دپاره
  \LR{$\Sat{M}{!C}$}۔
  په (الف) او (ب) حالتونو کښې
  \LR{$\Sat{M}{!A \lif !B}$}، نو داسې
  \LR{$!C \in \Delta, !A \lif !B$} شته چې
  \LR{$\Sat{M}{!C}$} وي۔ په (ج) حالت کښې د کوم
  \LR{$!C \in \Gamma$} دپاره
  \LR{$\Sat/{M}{!C}$} دے۔ په (د) حالت کښې د
  کوم \LR{$!C \in \Delta$} دپاره
  \LR{$\Sat{M}{!C}$} دے۔ په هر حالت کښې
  \LR{$\Struct{M}$} د
  \LR{$\Gamma \Sequent \Delta, !A \lif !B$} سېکوېنټ پوره کوي۔ ځکه
  \LR{$\Struct{M}$} اختياري و،
  \LR{$\Gamma \Sequent \Delta, !A \lif !B$} معتبر دے۔
  %
\item وروستنے استنتاج \LeftR{\lforall} دے: بيا داسې
  فارمول~\LR{$!A(x)$} او تړلې اصطلاح~\LR{$t$} شته چې \LR{$\pi$} داسې
  پاے ته رسېږي:
  \begin{prooftree}
    \AxiomC{}
    \Deduce$!A(t), \Gamma \fCenter \Delta$
    \RightLabel{\LeftR{\lforall}}
    \UnaryInf$\lforall[x][!A(x)], \Gamma \fCenter \Delta$
  \end{prooftree}
  غواړو وښيو چې پايله
  \LR{$\lforall[x][!A(x)], \Gamma \Sequent \Delta$} معتبره ده۔ يو
  جوړښت~\LR{$\Struct M$} په پام کښې ونيسئ۔ ځکه چې مقدمه
  \LR{$!A(t), \Gamma \Sequent \Delta$} معتبره ده، يا
  (الف) \LR{$\Sat/{M}{!A(t)}$}، (ب) د کوم \LR{$!C \in \Gamma$} دپاره
  \LR{$\Sat/{M}{!C}$}، يا (ج) د کوم \LR{$!C \in \Delta$} دپاره
  \LR{$\Sat{M}{!C}$} دے۔ په (الف) حالت کښې، د
  \olref[syn][sem]{prop:quant-terms} له مخې، که
  \LR{$\Sat{M}{\lforall[x][!A(x)]}$} وي، نو \LR{$\Sat{M}{!A(t)}$} دے۔ ځکه
  \LR{$\Sat/{M}{!A(t)}$}، نو
  \LR{$\Sat/{M}{\lforall[x][!A(x)]}$}۔ په (ب) او (ج) حالتونو کښې
  \LR{$\Struct{M}$} هم \LR{$\lforall[x][!A(x)], \Gamma \Sequent \Delta$}
  پوره کوي۔ ځکه \LR{$\Struct M$} اختياري و،
  \LR{$\lforall[x][!A(x)], \Gamma \Sequent \Delta$} معتبر دے۔
\item وروستنے استنتاج \RightR{\lexists} دے: تمرين۔
\item وروستنے استنتاج \RightR{\lforall} دے: بيا داسې
  فارمول~\LR{$!A(x)$} او ثابت سمبول~\LR{$a$} شته چې \LR{$\pi$} داسې
  پاے ته رسېږي:
  \begin{prooftree}
    \AxiomC{}
    \Deduce$\Gamma \fCenter \Delta, !A(a)$
    \RightLabel{\RightR{\lforall}}
    \UnaryInf$\Gamma \fCenter \Delta, \lforall[x][!A(x)]$
  \end{prooftree}
  دلته د ځانګړي متغير شرط پوره کېږي؛ يعنې \LR{$a$} په \LR{$!A(x)$}،
  \LR{$\Gamma$} يا \LR{$\Delta$} کښې نۀ راځي۔ د استقرا د فرضيې له مخې د
  وروستي استنتاج مقدمه معتبره ده۔ بايد وښيو چې پايله هم معتبره ده؛
  يعنې د هر جوړښت~\LR{$\Struct M$} دپاره يا
  (الف) \LR{$\Sat{M}{\lforall[x][!A(x)]}$}، يا
  (ب) د کوم \LR{$!C \in \Gamma$} دپاره \LR{$\Sat/{M}{!C}$}، يا
  (ج) د کوم \LR{$!C \in \Delta$} دپاره \LR{$\Sat{M}{!C}$} وي۔
\olpseganchor{OLP-0081-B016}{end}

\olpseganchor{OLP-0081-B017}{start}
  فرض کړئ \LR{$\Struct{M}$} يو اختياري جوړښت دے۔ که (ب) يا (ج)
  سم وي، کار بشپړ دے؛ نو فرض کړئ يو هم سم نۀ دے: د هر
  \LR{$!C \in \Gamma$} دپاره \LR{$\Sat{M}{!C}$}، او د هر
  \LR{$!C \in \Delta$} دپاره \LR{$\Sat/{M}{!C}$}۔ بايد وښيو چې (الف) سم دے،
  يعنې \LR{$\Sat{M}{\lforall[x][!A(x)]}$}۔ د
  \olref[syn][ass]{prop:sat-quant} له مخې کافي ده وښيو چې د ټولو
  متغير-ټاکنو~\LR{$s$} دپاره \LR{$\Sat{M}{!A(x)}[s]$}۔ نو يو اختياري
  متغير-ټاکنه~\LR{$s$} ونيسئ۔ داسې جوړښت~\LR{$\Struct{M'}$} په پام کښې
  ونيسئ چې د~\LR{$\Struct{M}$} په شان دے، پرته له دې چې
  \LR{$\Assign{a}{M'} = s(x)$}۔ د
  \olref[syn][ext]{cor:extensionality-sent} له مخې، د هر
  \LR{$!C \in \Gamma$} دپاره \LR{$\Sat{M'}{!C}$}؛ ځکه \LR{$a$} په~\LR{$\Gamma$} کښې
  نۀ راځي، او د هر \LR{$!C \in \Delta$} دپاره \LR{$\Sat/{M'}{!C}$}۔ خو
  مقدمه معتبره ده، نو \LR{$\Sat{M'}{!A(a)}$}۔ د
  \olref[syn][ass]{prop:sentence-sat-true} له مخې
  \LR{$\Sat{M'}{!A(a)}[s]$}، ځکه \LR{$!A(a)$} يوه جمله ده۔ اوس
  \LR{$\varAssign{s}{s}{x}$} او
  \LR{$s(x) = \Value{a}{M'}[s]$}، ځکه \LR{$\Struct{M'}$} مو همداسې تعريف
  کړے دے۔ نو \olref[syn][ext]{prop:ext-formulas} پلي کېږي او
  \LR{$\Sat{M'}{!A(x)}[s]$} ترلاسه کوو۔ ځکه \LR{$a$} په~\LR{$!A(x)$} کښې نۀ راځي،
  د \olref[syn][ext]{prop:extensionality} له مخې
  \LR{$\Sat{M}{!A(x)}[s]$}۔ ځکه \LR{$s$} اختياري و، دا مو ثابت کړل چې د ټولو
  متغير-ټاکنو دپاره \LR{$\Sat{M}{!A(x)}[s]$}۔
\item وروستنے استنتاج \LeftR{\lexists} دے: تمرين۔

\end{enumerate}
اوس هغه ممکن استنتاجونه څېړو چې دوه مقدمې لري۔
\begin{enumerate}
\item وروستنے استنتاج پرېکون دے: بيا \LR{$\pi$} داسې پاے ته رسېږي:
  \begin{prooftree}
    \AxiomC{}
    \Deduce$\Gamma \fCenter \Delta, !A$
    \AxiomC{}
    \Deduce$!A, \Pi \fCenter \Lambda$
    \RightLabel{\Cut}
    \BinaryInf$\Gamma, \Pi \fCenter \Delta, \Lambda$
  \end{prooftree}
  فرض کړئ \LR{$\Struct{M}$} جوړښت دے۔ د استقرا د فرضيې له مخې دواړه
  مقدمې معتبرې دي، نو
  \LR{$\Struct{M}$} دواړه پوره کوي۔
  دوه حالتونه بېلوو: (الف)
  \LR{$\Sat/{M}{!A}$} او (ب)
  \LR{$\Sat{M}{!A}$}۔ په (الف) حالت کښې، د دې
  دپاره چې \LR{$\Struct{M}$} کيڼه مقدمه
  پوره کړي، بايد \LR{$\Gamma \Sequent \Delta$} پوره کړي۔ بيا پايله هم
  پوره کوي۔ په (ب) حالت کښې، د دې دپاره چې
  \LR{$\Struct{M}$} ښۍ مقدمه پوره کړي،
  \textbf{د ژباړې د نسخې سمون (OLSQ-006):} سرچينې دلته د
  سېکوېنټ د نښې پر ځاے د سټ-توپير نښه لرله؛ د قاعدې له مقدمو سره
  سمه نښه راوستل کېږي۔
  بايد \LR{$\Pi \Sequent \Lambda$} پوره کړي۔ بيا هم
  \LR{$\Struct{M}$} پايله پوره کوي۔
\item وروستنے استنتاج \RightR{\land} دے۔ بيا \LR{$\pi$} داسې پاے ته رسېږي:
  \begin{prooftree}
    \AxiomC{}
    \Deduce$\Gamma \fCenter \Delta, !A$
    \AxiomC{}
    \Deduce$\Gamma \fCenter \Delta, !B$
    \RightLabel{\RightR{\land}}
    \BinaryInf$\Gamma \fCenter \Delta, !A \land !B$
  \end{prooftree}
  يو جوړښت~\LR{$\Struct
    M$} په پام کښې ونيسئ۔ که
  \LR{$\Struct{M}$} د
  \LR{$\Gamma \Sequent \Delta$} سېکوېنټ پوره کوي، کار بشپړ دے۔ نو فرض
  کړئ نۀ ئې پوره کوي۔ ځکه د استقرا د فرضيې له مخې
  \LR{$\Gamma \fCenter \Delta, !A$} معتبر دے،
  \LR{$\Sat{M}{!A}$}۔ همدارنګه، ځکه
  \LR{$\Gamma \Sequent \Delta, !B$} معتبر دے،
  \LR{$\Sat{M}{!B}$}۔ نو
  \LR{$\Sat{M}{!A \land !B}$}۔
\item وروستنے استنتاج \LeftR{\lor} دے: تمرين۔
\item وروستنے استنتاج \LeftR{\lif} دے۔ بيا \LR{$\pi$} داسې پاے ته رسېږي:
  \begin{prooftree}
    \AxiomC{}
    \Deduce$\Gamma \fCenter \Delta, !A$
    \AxiomC{}
    \Deduce$!B, \Pi \fCenter \Lambda$
    \RightLabel{\LeftR{\lif}}
    \BinaryInf$!A \lif !B, \Gamma, \Pi \fCenter \Delta, \Lambda$
  \end{prooftree}
  بيا هم يو
  جوړښت~\LR{$\Struct{M}$} په پام کښې ونيسئ او فرض کړئ چې
  \LR{$\Struct{M}$} د
  \LR{$\Gamma, \Pi \Sequent \Delta, \Lambda$} سېکوېنټ نۀ پوره کوي۔
  بايد وښيو چې
  \LR{$\Sat/{M}{!A \lif !B}$}۔ که
  \LR{$\Struct{M}$} د
  \LR{$\Gamma, \Pi \Sequent \Delta, \Lambda$} سېکوېنټ نۀ پوره کوي،
  نۀ \LR{$\Gamma \Sequent \Delta$} او نۀ \LR{$\Pi \Sequent \Lambda$} پوره
  کوي۔ ځکه \LR{$\Gamma \Sequent \Delta, !A$} معتبر دے،
  \LR{$\Sat{M}{!A}$}؛ او ځکه
  \LR{$!B, \Pi \Sequent \Lambda$} معتبر دے،
  \LR{$\Sat/{M}{!B}$}۔ خو بيا
  \LR{$\Sat/{M}{!A \lif !B}$}، او همدا مو
  غوښتل۔
\end{enumerate}
\end{proof}
\olpseganchor{OLP-0081-B017}{end}

\olpseganchor{OLP-0081-B018}{start}
\begin{prob}
د \olref[fol][seq][sou]{thm:sequent-soundness} ثبوت بشپړ کړئ۔
\end{prob}
\olpseganchor{OLP-0081-B018}{end}



\olpseganchor{OLP-0081-B020}{start}
\begin{cor}
\ollabel{cor:weak-soundness}
که \LR{$\Proves !A$} وي، نو \LR{$!A$} معتبر ده۔
\end{cor}
\olpseganchor{OLP-0081-B020}{end}

\olpseganchor{OLP-0081-B021}{start}
\begin{cor}
\ollabel{cor:entailment-soundness}
که \LR{$\Gamma \Proves !A$} وي، نو \LR{$\Gamma \Entails !A$}۔
\end{cor}
\olpseganchor{OLP-0081-B021}{end}

\olpseganchor{OLP-0081-B022}{start}
\begin{proof}
که \LR{$\Gamma \Proves !A$} وي، نو د کوم متناهي فرعي سټ
\LR{$\Gamma_0 \subseteq \Gamma$} دپاره د
\LR{$\Gamma_0 \Sequent !A$} اشتقاق شته۔ د
\olref{thm:sequent-soundness} له مخې، هر
جوړښت~\LR{$\Struct{M}$} يا کوم \LR{$!B \in \Gamma_0$}
ناسموي يا \LR{$!A$} رښتيا کوي۔ نو که
\LR{$\Sat{M}{\Gamma}$} وي، بيا
\LR{$\Sat{M}{!A}$} هم دے۔
\end{proof}
\olpseganchor{OLP-0081-B022}{end}

\olpseganchor{OLP-0081-B023}{start}
\begin{cor}
\ollabel{cor:consistency-soundness}
که \LR{$\Gamma$} د صدق وړ وي، نو سازګار دے۔
\end{cor}
\olpseganchor{OLP-0081-B023}{end}

\olpseganchor{OLP-0081-B024}{start}
\begin{proof}
موږ مقابل عکس ثابتوو۔ فرض کړئ \LR{$\Gamma$} سازګار نۀ دے۔ بيا يو
متناهي \LR{$\Gamma_0 \subseteq \Gamma$} او د
\LR{$\Gamma_0 \Sequent \quad$} اشتقاق شته۔ د
\olref{thm:sequent-soundness} له مخې
\LR{$\Gamma_0 \Sequent \quad$} معتبر دے۔ په بله وينا، د هر
جوړښت~\LR{$\Struct{M}$} دپاره داسې
\LR{$!C \in \Gamma_0$} شته چې
\LR{$\Sat/{M}{!C}$}، او ځکه
\LR{$\Gamma_0 \subseteq \Gamma$}، دغه \LR{$!C$} په~\LR{$\Gamma$} کښې هم دے۔
نو هېڅ \LR{$\Struct{M}$} د~\LR{$\Gamma$} سټ نۀ
پوره کوي، او \LR{$\Gamma$} د صدق وړ نۀ دے۔
\end{proof}
\olpseganchor{OLP-0081-B024}{end}

\olpseganchor{OLP-0082-B004}{start}
\olfileid{fol}{seq}{ide}
\olpseganchor{OLP-0082-B004}{end}

\olpseganchor{OLP-0082-B005}{start}
\olsection{له عينيت سره اشتقاقونه}
\olpseganchor{OLP-0082-B005}{end}

\olpseganchor{OLP-0082-B006}{start}
له عينيت سره اشتقاقونه اضافي لومړنيو سېکوېنټونو او
استنتاجي قاعدو ته اړتيا لري۔
\olpseganchor{OLP-0082-B006}{end}

\olpseganchor{OLP-0082-B007}{start}
\begin{defn}[د~\LR{$\eq$} دپاره لومړني سېکوېنټونه]
که \LR{$t$} تړلې اصطلاح وي، نو \LR{${} \Sequent \eq[t][t]$} يو لومړنے
سېکوېنټ دے۔
\end{defn}
\olpseganchor{OLP-0082-B007}{end}

\olpseganchor{OLP-0082-B008}{start}
د \LR{$\eq$} قاعدې دا دي (\LR{$t_1$} او \LR{$t_2$} تړلې اصطلاحګانې دي):
\olpseganchor{OLP-0082-B008}{end}

\olpseganchor{OLP-0082-B009}{start}
\begin{defish}
\Axiom$ \eq[t_1][t_2], \Gamma \fCenter \Delta, !A(t_1) $
\RightLabel{\LR{$\eq$}}
\UnaryInf$\eq[t_1][t_2], \Gamma \fCenter \Delta, !A(t_2)$
\DisplayProof
\hfill
\Axiom$\eq[t_1][t_2], \Gamma \fCenter \Delta, !A(t_2) $
\RightLabel{\LR{$\eq$}}
\UnaryInf$\eq[t_1][t_2], \Gamma  \fCenter \Delta, !A(t_1)$
\DisplayProof
\end{defish}
\olpseganchor{OLP-0082-B009}{end}

\olpseganchor{OLP-0082-B010}{start}
\begin{ex}
که \LR{$s$} او \LR{$t$} تړلې اصطلاحګانې وي، نو
\LR{$\eq[s][t], !A(s) \Proves !A(t)$}:
\begin{prooftree}
\Axiom$ !A(s) \fCenter !A(s)$
\RightLabel{\LeftR{\Weakening}}
\UnaryInf$\eq[s][t], !A(s)  \fCenter !A(s)$
\RightLabel{\LR{$\eq$}}
\UnaryInf$\eq[s][t], !A(s)  \fCenter !A(t)$
\end{prooftree}
دا به د يو شان څيزونو د تعويض د اصل يا د لايبنېڅ د قانون په نامه
اشنا وي۔
\olpseganchor{OLP-0082-B010}{end}

\olpseganchor{OLP-0082-B011}{start}
\LR{$\Log{LK}$} ثابتوي چې \LR{$\eq$} متناظره او انتقالي ده:
\begin{prooftree}
\Axiom$ \fCenter \eq[t_1][t_1] $
\RightLabel{\LeftR{\Weakening}}
\UnaryInf$ \eq[t_1][t_2] \fCenter \eq[t_1][t_1] $
\RightLabel{\LR{$\eq$}}
\UnaryInf$ \eq[t_1][t_2] \fCenter \eq[t_2][t_1]$
\DisplayProof\qquad\bottomAlignProof
\Axiom$ \eq[t_1][t_2] \fCenter \eq[t_1][t_2] $
\RightLabel{\LeftR{\Weakening}}
\UnaryInf$\eq[t_2][t_3], \eq[t_1][t_2]  \fCenter \eq[t_1][t_2] $
\RightLabel{\LR{$\eq$}}
\UnaryInf$\eq[t_2][t_3], \eq[t_1][t_2]  \fCenter \eq[t_1][t_3]$
\RightLabel{\LeftR{\Exchange}}
\UnaryInf$\eq[t_1][t_2], \eq[t_2][t_3]  \fCenter \eq[t_1][t_3]$
\end{prooftree}
په کيڼ اشتقاق کښې فارمول~\LR{$\eq[x][t_1]$} زموږ
\LR{$!A(x)$} دے۔ په ښي اړخ کښې \LR{$!A(x)$} د~\LR{$\eq[t_1][x]$} په توګه اخلو۔
\end{ex}
\olpseganchor{OLP-0082-B011}{end}

\olpseganchor{OLP-0082-B012}{start}
\begin{prob}
د لاندې سېکوېنټونو اشتقاقونه ورکړئ:
\begin{enumerate}
\item \LR{$\Sequent \lforall[x][\lforall[y][((x = y \land !A(x)) \lif !A(y))]]$}
\item \LR{$\lexists[x][!A(x)] \land \lforall[y][\lforall[z][((!A(y) \land
    !A(z)) \lif y = z)]] \Sequent 
\lexists[x][(!A(x) \land \lforall[y][(!A(y) \lif y = x)])]$}
\end{enumerate}
\end{prob}
\olpseganchor{OLP-0082-B012}{end}

\olpseganchor{OLP-0083-B004}{start}
\olfileid{fol}{seq}{sid}
\olpseganchor{OLP-0083-B004}{end}

\olpseganchor{OLP-0083-B005}{start}
\olsection{له عينيت سره سموالے}
\olpseganchor{OLP-0083-B005}{end}

\olpseganchor{OLP-0083-B006}{start}
\begin{prop}
\LR{$\Log{LK}$} د عينيت له لومړنيو سېکوېنټونو او قاعدو سره سم دے۔
\end{prop}
\olpseganchor{OLP-0083-B006}{end}

\olpseganchor{OLP-0083-B007}{start}
\begin{proof}
د \LR{${} \Sequent \eq[t][t]$} په بڼه لومړني سېکوېنټونه معتبر دي؛
ځکه د هر جوړښت~\LR{$\Struct M$} دپاره
\LR{$\Sat{M}{\eq[t][t]}$}۔ (پام وکړئ، فرض کوو چې اصطلاح~\LR{$t$} تړلې ده،
يعنې هېڅ متغير نۀ لري؛ نو متغير-ټاکنې اهميت نۀ لري۔)
\olpseganchor{OLP-0083-B007}{end}

\olpseganchor{OLP-0083-B008}{start}
فرض کړئ په اشتقاق کښې وروستنے استنتاج \LR{$=$} دے۔ بيا مقدمه
\LR{$\eq[t_1][t_2], \Gamma \Sequent \Delta, !A(t_1)$} او پايله
\LR{$\eq[t_1][t_2], \Gamma \Sequent \Delta, !A(t_2)$} ده۔ يو
جوړښت~\LR{$\Struct M$} په پام کښې ونيسئ۔ بايد وښيو چې پايله
معتبره ده؛ يعنې که \LR{$\Sat{M}{\eq[t_1][t_2]}$} او
\LR{$\Sat{M}{\Gamma}$} وي، نو يا د کوم \LR{$!C \in \Delta$} دپاره
\LR{$\Sat{M}{!C}$}، يا \LR{$\Sat{M}{!A(t_2)}$} دے۔
\olpseganchor{OLP-0083-B008}{end}

\olpseganchor{OLP-0083-B009}{start}
د استقرا د فرضيې له مخې مقدمه معتبره ده۔ د دې معنٰی دا ده چې که
\LR{$\Sat{M}{\eq[t_1][t_2]}$} او \LR{$\Sat{M}{\Gamma}$} وي، نو يا
(الف) د کوم \LR{$!C \in \Delta$} دپاره \LR{$\Sat{M}{!C}$}، يا
(ب) \LR{$\Sat{M}{!A(t_1)}$} دے۔ په (الف) حالت کښې کار بشپړ دے۔ (ب)
حالت وګورئ۔ داسې متغير-ټاکنه~\LR{$s$} ونيسئ چې
\LR{$s(x) = \Value{t_1}{M}$}۔ د
\olref[syn][ass]{prop:sentence-sat-true} له مخې
\LR{$\Sat{M}{!A(t_1)}[s]$}۔ ځکه \LR{$\varAssign{s}{s}{x}$}، د
\olref[syn][ext]{prop:ext-formulas} له مخې
\LR{$\Sat{M}{!A(x)}[s]$}۔ ځکه \LR{$\Sat{M}{\eq[t_1][t_2]}$}، نو
\LR{$\Value{t_1}{M} = \Value{t_2}{M}$}، او له دې امله
\LR{$s(x) = \Value{t_2}{M}$}۔ د
\olref[syn][ext]{prop:ext-formulas} په بيا پلي کولو
\LR{$\Sat{M}{!A(t_2)}[s]$} هم لرو۔ د
\olref[syn][ass]{prop:sentence-sat-true} له مخې
\LR{$\Sat{M}{!A(t_2)}$}۔
\end{proof}
\olpseganchor{OLP-0083-B009}{end}

\olpseganchor{OLP-0084-B004}{start}
\olchapter{fol}{ntd}{فطري استنتاج}
\olpseganchor{OLP-0084-B004}{end}

\olpseganchor{OLP-0084-B005}{start}
\begin{editorial}
دا څپرکے د ګېنتڅن/پراويڅ په طريقه د فطري استنتاج يو نظام وړاندې
کوي۔
\olpseganchor{OLP-0084-B005}{end}
  
\olpseganchor{OLP-0084-B006}{start}
د ثبوت د يوه نظام په توګه له فطري استنتاج سره اړوند مواد د شاملولو
يا ايستلو دپاره د~\texttt{prfND} ټګ وکاروئ۔
\end{editorial}
\olpseganchor{OLP-0084-B006}{end}





\olpseganchor{OLP-0084-B009}{start}
%
  

\olpseganchor{OLP-0084-B009}{end}





\olpseganchor{OLP-0084-B012}{start}
%
  

\olpseganchor{OLP-0084-B012}{end}







\olpseganchor{OLP-0084-B016}{start}
%
  

\olpseganchor{OLP-0084-B016}{end}



\olpseganchor{OLP-0084-B018}{start}
%
  
  

\olpseganchor{OLP-0084-B018}{end}

\olpseganchor{OLP-0085-B004}{start}
\olfileid{fol}{ntd}{rul}
\olpseganchor{OLP-0085-B004}{end}

\olpseganchor{OLP-0085-B005}{start}
\olsection{قاعدې او اشتقاقونه}
\olpseganchor{OLP-0085-B005}{end}

\olpseganchor{OLP-0085-B006}{start}
\begin{explain}
د فطري استنتاج له نظامونو څخه موخه دا ده چې په رياضيکي ثبوت کښې له
کارېدونکي غيررسمي استدلال سره نژدې موازي وي (همدا وجه ده چې تر يوه
حده “فطري” دے)۔ د فطري استنتاج ثبوتونه له فرضيو پيلېږي۔ بيا د
استنتاج قاعدې ورباندې کارول کېږي۔ فرضيې د~\Intro{\lnot}،
\Intro{\lif}، \Elim{\lor} او \Elim{\lexists} د استنتاج د قاعدو په وسيله “ايستل” کېږي،
او د روښانتيا دپاره د ايستل شوې فرضيې نښان د استنتاج تر
څنګ ليکل کېږي۔
\end{explain}
\olpseganchor{OLP-0085-B006}{end}

\olpseganchor{OLP-0085-B007}{start}
\begin{defn}[فرضيه]
يوه \emph{فرضيه} هره هغه تړلے فارمول ده چې د هرې څانګې په تر ټولو
پورتني ځاے کښې وي۔
\end{defn}
\olpseganchor{OLP-0085-B007}{end}

\olpseganchor{OLP-0085-B008}{start}
په فطري استنتاج کښې اشتقاقونه د تړلي فارمولونه ځانګړې ونې دي؛
تر ټولو پورتنۍ تړلي فارمولونه فرضيې وي، او که يوه تړلے فارمول تر
يوې، دوو يا دريو نورو جملو لاندې وي، بايد په سمه توګه د استنتاج له
يوې قاعدې څخه پايله شي۔ د استنتاج په سر کښې تړلي فارمولونه د هغې
\emph{مقدمې} او لاندې تړلے فارمول د استنتاج \emph{پايله} بلل کېږي۔
قاعدې په جوړو کښې راځي: د هر منطقي عملګر دپاره يوه د داخلولو او
يوه د ايستلو قاعده۔ دوئ په پايله کښې منطقي عملګر داخلوي يا
منطقي عملګر د قاعدې له يوې مقدمې څخه باسي۔ ځينې قاعدې اجازه
ورکوي چې د يوه ټاکلي
ډول فرضيه \emph{ايستل} شي۔ د دې د ښودلو دپاره چې کومه
فرضيه د کوم استنتاج په وسيله ايستل کېږي، فرضيې او استنتاج
دواړو ته نښانونه ورکوو۔ دا د فرضيې په “\LR{$\Discharge{!A}{n}$}” ليکلو
ښودل کېږي۔ \textbf{د ژباړې د نسخې سمون (OLND-001):} سرچينه په دې
ونه کښې لاندې ولاړ توکي “سېکوېنټونه” بولي؛ دلته “جملې” راوړل شوې،
ځکه خپله سرچينه ونه د جملو ونه تعريفوي او قاعدې پر جملو لګېږي۔
\olpseganchor{OLP-0085-B008}{end}

\olpseganchor{OLP-0085-B009}{start}
% Only include this sentence is on of \land, lor, lif or lnot is defined.
\olpseganchor{OLP-0085-B009}{end}

\olpseganchor{OLP-0085-B010}{start}
دا دود دے چې د ټولو منطقي عملګرونه~\LR{$\land$}، \LR{$\lor$}، \LR{$\lif$}،
    \LR{$\lnot$} او~\LR{$\lfalse$} دپاره قاعدې په پام کښې ونيول شي، که څه هم
    ځينې ئې تعريف شوي وي۔
\olpseganchor{OLP-0085-B010}{end}

\olpseganchor{OLP-0086-B004}{start}
\olfileid{fol}{ntd}{prl}
\olpseganchor{OLP-0086-B004}{end}

\olpseganchor{OLP-0086-B005}{start}
\olsection{قضيې قاعدې}
\olpseganchor{OLP-0086-B005}{end}


\olpseganchor{OLP-0086-B006}{start}
\subsection{د \LR{$\land$} دپاره قاعدې}
\olpseganchor{OLP-0086-B006}{end}

\olpseganchor{OLP-0086-B007}{start}
\begin{defish}
\AxiomC{\LR{$!A$}}
\AxiomC{\LR{$!B$}}
\RightLabel{\Intro{\land}}
\BinaryInfC{\LR{$!A \land !B$}}
\DisplayProof
\hfill
\begin{tabular}{r}
\AxiomC{\LR{$!A \land !B$}}
\RightLabel{\Elim{\land}}
\UnaryInfC{\LR{$!A$}}
\DisplayProof
\\[3ex]
\AxiomC{\LR{$!A \land !B$}}
\RightLabel{\Elim{\land}}
\UnaryInfC{\LR{$!B$}}
\DisplayProof
\end{tabular}
\end{defish}
\olpseganchor{OLP-0086-B007}{end}

\olpseganchor{OLP-0086-B008}{start}
\subsection{د \LR{$\lor$} دپاره قاعدې}
\olpseganchor{OLP-0086-B008}{end}

\olpseganchor{OLP-0086-B009}{start}
\begin{defish}
\begin{tabular}{r}
\AxiomC{\LR{$!A$}}
\RightLabel{\Intro{\lor}}
\UnaryInfC{\LR{$!A \lor !B$}}
\DisplayProof
\\[3ex]
\AxiomC{\LR{$!B$}}
\RightLabel{\Intro{\lor}}
\UnaryInfC{\LR{$!A \lor !B$}}
\DisplayProof
\end{tabular}
\hfill
\AxiomC{\LR{$!A \lor !B$}}
\AxiomC{\LR{$\Discharge{!A}{n}$}}
\DeduceC{\LR{$!C$}}
\AxiomC{\LR{$\Discharge{!B}{n}$}}
\DeduceC{\LR{$!C$}}
\DischargeRule{\Elim{\lor}}{n}
\TrinaryInfC{\LR{$!C$}}
\DisplayProof
\end{defish}
\olpseganchor{OLP-0086-B009}{end}

\olpseganchor{OLP-0086-B010}{start}
\subsection{د \LR{$\lif$} دپاره قاعدې}
\olpseganchor{OLP-0086-B010}{end}

\olpseganchor{OLP-0086-B011}{start}
\begin{defish}
\AxiomC{\LR{$\Discharge{!A}{n}$}}
\DeduceC{\LR{$!B$}}
\DischargeRule{\Intro{\lif}}{n}
\UnaryInfC{\LR{$!A \lif !B$}}
\DisplayProof
\hfill
\AxiomC{\LR{$!A \lif !B$}}
\AxiomC{\LR{$!A$}}
\RightLabel{\Elim{\lif}}
\BinaryInfC{\LR{$!B$}}
\DisplayProof
\end{defish}
\olpseganchor{OLP-0086-B011}{end}

\olpseganchor{OLP-0086-B012}{start}
\subsection{د \LR{$\lnot$} دپاره قاعدې}
\olpseganchor{OLP-0086-B012}{end}

\olpseganchor{OLP-0086-B013}{start}
\begin{defish}
\AxiomC{\LR{$\Discharge{!A}{n}$}}
\noLine
\DeduceC{\LR{$\lfalse$}}
\DischargeRule{\Intro{\lnot}}{n}
\UnaryInfC{\LR{$\lnot !A$}}
\DisplayProof
\hfill
\AxiomC{\LR{$\lnot !A$}}
\AxiomC{\LR{$!A$}}
\RightLabel{\Elim{\lnot}}
\BinaryInfC{\LR{$\lfalse$}}
\DisplayProof
\end{defish}
\olpseganchor{OLP-0086-B013}{end}

\olpseganchor{OLP-0086-B014}{start}
\subsection{د \LR{$\lfalse$} دپاره قاعدې}
\olpseganchor{OLP-0086-B014}{end}

\olpseganchor{OLP-0086-B015}{start}
\begin{defish}
\AxiomC{\LR{$\lfalse$}}
\RightLabel{\FalseInt}
\UnaryInfC{\LR{$!A$}}
\DisplayProof
\hfill
\AxiomC{\LR{$\Discharge{\lnot !A}{n}$}}
\DeduceC{\LR{$\lfalse$}}
\DischargeRule{\FalseCl}{n}
\UnaryInfC{\LR{$!A$}}
\DisplayProof
\end{defish}
\olpseganchor{OLP-0086-B015}{end}

\olpseganchor{OLP-0086-B016}{start}
پام وکړئ چې \LR{$\Intro{\lnot}$} او~\LR{$\FalseCl$} ډېر سره ورته دي: توپير
دا دے چې \LR{$\Intro{\lnot}$} يوه نفي شوې تړلے فارمول~\LR{$\lnot !A$} راوباسي،
خو \LR{$\FalseCl$} يوه مثبته تړلے فارمول~\LR{$!A$} راوباسي۔
\olpseganchor{OLP-0086-B016}{end}

\olpseganchor{OLP-0086-B017}{start}
کله چې يوه قاعده ښيي چې کومه فرضيه ايستل کېدے شي، موږ دا اجازه
بولو، شرط ئې نۀ بولو۔ مثلاً د~\LR{$\Intro{\lif}$} په قاعده کښې د
مقدمې~\LR{$!B$} په اشتقاق کښې د~\LR{$!A$} په بڼه هر شمېر فرضيې
ايستلے شو، چې صفر هم پکې شامل دے۔
\olpseganchor{OLP-0086-B017}{end}

\olpseganchor{OLP-0087-B004}{start}
\olfileid{fol}{ntd}{qrl}
\olpseganchor{OLP-0087-B004}{end}

\olpseganchor{OLP-0087-B005}{start}
\olsection{د سورونو قاعدې}
\olpseganchor{OLP-0087-B005}{end}

\olpseganchor{OLP-0087-B006}{start}
\subsection{د \LR{$\lforall$} دپاره قاعدې}
\olpseganchor{OLP-0087-B006}{end}

\olpseganchor{OLP-0087-B007}{start}
\begin{defish}
\AxiomC{\LR{$!A(a)$}}
\RightLabel{\Intro{\lforall}}
\UnaryInfC{\LR{$\lforall[x][\Atom{!A}{x}]$}}
\DisplayProof
\hfill
\AxiomC{\LR{$\lforall[x][\Atom{!A}{x}]$}}
\RightLabel{\Elim{\lforall}}
\UnaryInfC{\LR{$!A(t)$}}
\DisplayProof
\end{defish}
\olpseganchor{OLP-0087-B007}{end}

\olpseganchor{OLP-0087-B008}{start}
د~\LR{$\lforall$} په قاعدو کښې \LR{$t$} يوه تړلې اصطلاح ده (يعنې داسې اصطلاح
چې هېڅ متغير نۀ لري)، او \LR{$a$}~داسې ثابت سمبول دے چې په
پايله~\LR{$\lforall[x][!A(x)]$} يا په هرې هغې فرضيې کښې نۀ راځي چې د
مقدمې~\LR{$!A(a)$} په پای ته رسېدونکي اشتقاق کښې
نۀ ايستل شوې پاتې وي۔ موږ \LR{$a$} د~\Intro{\lforall} استنتاج
\emph{ځانګړے متغير} بولو۔\footnote{موږ د “ځانګړي متغير” اصطلاح
کاروو، که څه هم په پورتنۍ قاعده کښې \LR{$a$} يو ثابت دے۔ دا نوم تاريخي
لاملونه لري۔}
\olpseganchor{OLP-0087-B008}{end}

\olpseganchor{OLP-0087-B009}{start}
\subsection{د \LR{$\lexists$} دپاره قاعدې}
\olpseganchor{OLP-0087-B009}{end}

\olpseganchor{OLP-0087-B010}{start}
\begin{defish}
\AxiomC{\LR{$\Atom{!A}{t}$}}
\RightLabel{\Intro{\lexists}}
\UnaryInfC{\LR{$\lexists[x][\Atom{!A}{x}]$}}
\DisplayProof
\hfill
\AxiomC{\LR{$\lexists[x][\Atom{!A}{x}]$}}
\AxiomC{[\LR{$\Atom{!A}{a}$}]\LR{$^n$}}
\DeduceC{\LR{$!C$}}
\DischargeRule{\Elim{\lexists}}{n}
\BinaryInfC{\LR{$!C$}}
\DisplayProof
\end{defish}
\olpseganchor{OLP-0087-B010}{end}

\olpseganchor{OLP-0087-B011}{start}
بيا هم، \LR{$t$} يوه تړلې اصطلاح ده، او \LR{$a$} داسې ثابت سمبول دے چې
په مقدمه~\LR{$\lexists[x][!A(x)]$}، پايله~\LR{$!C$}، يا په هرې هغې فرضيې
کښې نۀ راځي چې د دواړو مقدمو په پای ته رسېدونکو اشتقاقونه کښې
نۀ ايستل شوې پاتې وي (د~\LR{$!A(a)$} فرضيو پرته)۔ موږ \LR{$a$} د
\Elim{\lexists} استنتاج \emph{ځانګړے متغير} بولو۔
\olpseganchor{OLP-0087-B011}{end}

\olpseganchor{OLP-0087-B012}{start}
هغه شرط چې په~\Intro{\lforall} يا~\Elim{\lexists} استنتاج کښې
ځانګړے متغير په پايله يا په هرې هغې فرضيې کښې نۀ راځي چې مقدمو ته
په رسېدونکو اشتقاقونه کښې نۀ ايستل شوې وي، او د
په وجودي مقدمه کښې هم نۀ راځي، \emph{د ځانګړي
متغير شرط} بلل کېږي۔ \textbf{د ژباړې د نسخې سمون (OLND-002):}
سرچينه په دې عمومي جمله کښې وايي چې ځانګړے متغير په “مقدمو” کښې
نۀ راځي، خو د دواړو قاعدو مرستندويه مقدمه همدا ځانګړے ثابت لري؛
دلته د سرچينې له مخکښې کره شرطونو سره سم پايله، نۀ ايستل شوې فرضيې
او د وجود ايستلو وجودي مقدمه ټاکل شوې دي۔
\olpseganchor{OLP-0087-B012}{end}

\olpseganchor{OLP-0087-B013}{start}
\begin{explain}
هغه دود راياد کړئ چې کله \LR{$!A$} داسې فارمول وي چې
متغير~\LR{$x$} پکې ازاد وي، نو دا په~\LR{$!A(x)$} ليکلو ښيو۔ په همدې
چوکاټ کښې بيا \LR{$!A(t)$} د~\LR{$\Subst{!A}{t}{x}$} لنډون دے۔ نو د
\LR{$\Intro\lexists$} قاعده داسې هم ليکلے شو:
\begin{prooftree}
\AxiomC{\LR{$\Subst{!A}{t}{x}$}}
\RightLabel{\Intro{\lexists}}
\UnaryInfC{\LR{$\lexists[x][!A]$}}
\end{prooftree}
پام وکړئ چې \LR{$t$} کېدے شي له مخکښې په~\LR{$!A$} کښې راغلے وي؛ مثلاً
\LR{$!A$}~کېدے شي~\LR{$\Atom{\Obj P}{t,x}$} وي۔ له دې امله، له
~\LR{$\Atom{\Obj P}{t,t}$} څخه \LR{$\lexists[x][\Atom{\Obj P}{t,x}]$}
استنتاجول د~\LR{$\Intro\lexists$} سمه کارونه ده---تاسو د مقدمې د~\LR{$t$}
يوه يا زياتې پېښې په تړلي متغير~\LR{$x$} “بدلولے” شئ، او اړينه
نۀ ده چې ټولې پېښې ئې بدلې کړئ۔ خو د~\LR{$\Intro\lforall$} او
~\LR{$\Elim\lexists$} د ځانګړي متغير شرطونه غواړي چې ثابت سمبول~\LR{$a$}
په~\LR{$!A$} کښې نۀ راځي۔ نو د~\LR{$\Intro\lforall$} په کارولو سره له
\LR{$\Atom{\Obj P}{a,a}$} څخه \LR{$\lforall[x][\Atom{\Obj P}{a,x}]$} په سمه
توګه نۀ شئ استنتاجولے۔
\end{explain}
\olpseganchor{OLP-0087-B013}{end}

\olpseganchor{OLP-0087-B014}{start}
\begin{explain}
په~\Intro{\lexists} او~\Elim{\lforall} کښې له تړلتيا پرته بل
بنديز نشته، او اصطلاح~\LR{$t$} هره تړلې اصطلاح کېدے شي؛ نو د بل شرط
اندېښنه نۀ شته۔ بل لوري ته، د~\Elim{\lexists} او
\Intro{\lforall} په قاعدو کښې د ځانګړي متغير شرط غواړي چې
ثابت سمبول~\LR{$a$} په پايله يا کومې نۀ ايستل شوې فرضيې کښې هېڅ
ځاے نۀ راځي۔ دا شرط د نظام د سموالي د يقيني کولو دپاره اړين دے؛
يعنې نظام له هغو نۀ ايستل شوې فرضيو څخه يوازې هغه تړلي فارمولونه
اشتقاق کړي چې ترې معنايي پايله کېږي۔ له دې شرط پرته لاندې به
روا و:
\begin{prooftree}
  \AxiomC{\LR{$\lexists[x][!A(x)]$}}
  \AxiomC{\LR{$\Discharge{!A(a)}{1}$}}
  \RightLabel{*\Intro{\lforall}}
  \UnaryInfC{\LR{$\lforall[x][!A(x)]$}}
  \RightLabel{\Elim{\lexists}}
  \BinaryInfC{\LR{$\lforall[x][!A(x)]$}}
\end{prooftree}
خو \LR{$\lexists[x][!A(x)] \Entails/ \lforall[x][!A(x)]$}۔
\olpseganchor{OLP-0087-B014}{end}

\olpseganchor{OLP-0087-B015}{start}
ځکه چې د سورونو د ايستلو قاعدې د متغيرونه پر ځاے يوازې د تړلو
اصطلاحاتو ايښوولو اجازه ورکوي، نو هره هغه فارمول چې د
تړلي فارمولونه له يوه سټ څخه وايستل شي، پخپله تړلے فارمول ده۔
\end{explain}
\olpseganchor{OLP-0087-B015}{end}

\olpseganchor{OLP-0088-B004}{start}
\olfileid{fol}{ntd}{der}
\olpseganchor{OLP-0088-B004}{end}

\olpseganchor{OLP-0088-B005}{start}
\olsection{اشتقاقونه}
\olpseganchor{OLP-0088-B005}{end}

\olpseganchor{OLP-0088-B006}{start}
\begin{explain}
موږ وويل چې فرضيه څه ده او د استنتاج قاعدې مو ورکړې۔ په فطري
استنتاج کښې اشتقاقونه له دغو څخه په استقرايي ډول جوړېږي: هر
اشتقاق يا يوازې يوه فرضيه وي، يا له يو، دوو يا دريو
اشتقاقونه او ورپسې له يوه سم استنتاج څخه جوړ وي۔
\end{explain}
\olpseganchor{OLP-0088-B006}{end}

\olpseganchor{OLP-0088-B007}{start}
\begin{defn}[اشتقاق]
 له فرضيو~\LR{$\Gamma$} څخه د تړلے فارمول~\LR{$!A$} يو \emph{اشتقاق}
د تړلي فارمولونه يوه متناهي ونه ده چې لاندې شرطونه پوره کوي:
\begin{enumerate}
\item د ونې تر ټولو پورتنۍ تړلي فارمولونه يا په~\LR{$\Gamma$} کښې وي يا
  په ونه کښې د يوه استنتاج په وسيله ايستل شوې وي۔
\item د ونې تر ټولو لاندې تړلے فارمول~\LR{$!A$} ده۔
\item په ونه کښې هره تړلے فارمول، پرته له لاندېنۍ جملې~\LR{$!A$}
  څخه، د استنتاج د يوې قاعدې د سمې کارونې مقدمه وي چې پايله ئې په
  ونه کښې نېغ د هغې تړلے فارمول لاندې ولاړه وي۔
\end{enumerate}
بيا وايو چې \LR{$!A$} د اشتقاق \emph{پايله} او~\LR{$\Gamma$} ئې
نۀ ايستل شوې فرضيې دي۔
\olpseganchor{OLP-0088-B007}{end}

\olpseganchor{OLP-0088-B008}{start}
که له~\LR{$\Gamma$} څخه د~\LR{$!A$} کوم اشتقاق وي، وايو چې \LR{$!A$}
له~\LR{$\Gamma$} څخه \emph{د اشتقاق وړ} دے، يا په نښو کښې:
\LR{$\Gamma \Proves !A$}۔ که د~\LR{$!A$} داسې اشتقاق وي چې هره
فرضيه پکې ايستل شوې وي، نو~\LR{$\Proves !A$} ليکو۔
\end{defn}
\olpseganchor{OLP-0088-B008}{end}

\olpseganchor{OLP-0088-B009}{start}
\begin{ex}
هره فرضيه په يوازې ځان اشتقاق ده۔ نو مثلاً \LR{$!A$} په يوازې
ځان اشتقاق دے او \LR{$!B$} هم په يوازې ځان همداسې دے۔
له دغو څخه د، مثلاً، \LR{$\Intro{\land}$} قاعدې په کارولو يو نوے
اشتقاق اخستلے شو:
\begin{prooftree}
\AxiomC{\LR{$!A$}}
\AxiomC{\LR{$!B$}}
\RightLabel{\Intro{\land}}
\BinaryInfC{\LR{$!A \land !B$}}
\end{prooftree}
دا قاعدې عامې دي: پکې \LR{$!A$} او~\LR{$!B$} په هر ډول تړلي فارمولونه، مثلاً
په~\LR{$!C$} او~\LR{$!D$}، بدلولے شو۔ بيا به پايله \LR{$!C \land !D$} وي؛ نو
\begin{prooftree}
\AxiomC{\LR{$!C$}}
\AxiomC{\LR{$!D$}}
\RightLabel{\Intro{\land}}
\BinaryInfC{\LR{$!C \land !D$}}
\end{prooftree}
يو سم اشتقاق دے۔ البته فرضيې سره اړولے هم شو، څو \LR{$!D$} د~\LR{$!A$}
او \LR{$!C$} د~\LR{$!B$} ونډه واخلي۔ له دې امله
\begin{prooftree}
\AxiomC{\LR{$!D$}}
\AxiomC{\LR{$!C$}}
\RightLabel{\Intro{\land}}
\BinaryInfC{\LR{$!D \land !C$}}
\end{prooftree}
هم يو سم اشتقاق دے۔
\olpseganchor{OLP-0088-B009}{end}

\olpseganchor{OLP-0088-B010}{start}
اوس بله قاعده، مثلاً \LR{$\Intro{\lif}$}، کارولے شو۔ دا اجازه راکوي چې
يو شرطي قضيه پايله کړو او هره هغه فرضيه ايستل کړو چې له
دغه شرطي قضيې له مقدم سره يو شان وي۔ نو لاندې دواړه سم
اشتقاقونه دي:
\begin{prooftree}
\AxiomC{\LR{$\Discharge{!C}{1}$}}
\AxiomC{\LR{$!D$}}
\RightLabel{\Intro{\land}}
\BinaryInfC{\LR{$!C \land !D$}}
\DischargeRule{\Intro{\lif}}{1}
\UnaryInfC{\LR{$!C \lif (!C \land !D)$}}
\DisplayProof\bottomAlignProof
\AxiomC{\LR{$!C$}}
\AxiomC{\LR{$\Discharge{!D}{1}$}}
\RightLabel{\Intro{\land}}
\BinaryInfC{\LR{$!C \land !D$}}
\DischargeRule{\Intro{\lif}}{1}
\UnaryInfC{\LR{$!D \lif (!C \land !D)$}}
\end{prooftree}
دوئ په ترتيب سره ښيي چې \LR{$!D \Proves !C \lif (!C \land !D)$} او
\LR{$!C \Proves !D \lif (!C \land !D)$}۔
\olpseganchor{OLP-0088-B010}{end}

\olpseganchor{OLP-0088-B011}{start}
راياد کړئ چې د فرضيو ايستل اجازه ده، شرط نۀ دے: لازم نۀ ده چې
فرضیې وباسو۔ په ځانګړي ډول، يوه قاعده هله هم کارولے شو چې فرضيې په
اشتقاق کښې موجودې نۀ وي۔ مثلاً لاندې اشتقاق روا دے، که څه
هم د ايستل کېدو دپاره هېڅ فرضيه~\LR{$!A$} نۀ شته:
\begin{prooftree}
  \AxiomC{\LR{$!B$}}
  \DischargeRule{\Intro{\lif}}{1}
  \UnaryInfC{\LR{$!A \lif !B$}}
\end{prooftree}
\end{ex}
\olpseganchor{OLP-0088-B011}{end}

\olpseganchor{OLP-0089-B004}{start}
\olfileid{fol}{ntd}{pro}
\olpseganchor{OLP-0089-B004}{end}

\olpseganchor{OLP-0089-B005}{start}
\olsection{د اشتقاقونه بېلګې}
\olpseganchor{OLP-0089-B005}{end}

\olpseganchor{OLP-0089-B006}{start}
\begin{ex}
راځئ چې د تړلے فارمول~\LR{$(!A \land !B) \lif !A$} يو اشتقاق
ورکړو۔
\olpseganchor{OLP-0089-B006}{end}

\olpseganchor{OLP-0089-B007}{start}
په پيل کښې مطلوبه پايله د اشتقاق په لاندې برخه کښې ليکو۔
\begin{prooftree}
\AxiomC{}
\UnaryInfC{\LR{$(!A\land !B) \lif !A$}}
\end{prooftree}
\olpseganchor{OLP-0089-B007}{end}

\olpseganchor{OLP-0089-B008}{start}
اوس بايد معلومه کړو چې د کوم ډول استنتاج پايله د دې بڼې
تړلے فارمول کېدے شي۔ د پايلې اصلي عملګر~\LR{$\lif$} دے، نو هڅه
کوو چې د~\Intro{\lif} قاعدې په کارولو پايلې ته ورسېږو۔ غوره ده چې
اړوندې فرضيې وليکو او د پرمختګ پر مهال د استنتاج قاعدو ته نښانونه
ورکړو، څو د ثبوت په پای کښې په اسانۍ ووينو چې ټولې فرضيې
ايستل شوې دي که نۀ۔
\begin{prooftree}
\AxiomC{\LR{$\Discharge{!A \land !B}{1}$}}
\DeduceC{\LR{$!A$}}
\DischargeRule{\Intro{\lif}}{1} 
\UnaryInfC{\LR{$(!A\land !B) \lif !A$}}
\end{prooftree}
\olpseganchor{OLP-0089-B008}{end}

\olpseganchor{OLP-0089-B009}{start}
اوس بايد له فرضيې~\LR{$!A \land !B$} څخه تر~\LR{$!A$} پورې پړاوونه ډک کړو۔
ځکه چې يوازې له يوه رابط~\LR{$\land$} سره کار لرو، بايد د~\LR{$\land$} د
ايستلو قاعده وکاروو۔ له دې څخه لاندې ثبوت لاس ته راځي:
\begin{prooftree}
\AxiomC{\LR{$\Discharge{!A \land !B}{1}$}}
\RightLabel{\Elim{\land}}
\UnaryInfC{\LR{$!A$}}
\DischargeRule{\Intro{\lif}}{1} 
\UnaryInfC{\LR{$(!A\land !B) \lif !A$}}
\end{prooftree}
اوس د~\LR{$(!A \land !B) \lif !A$} يو سم اشتقاق لرو۔
\end{ex}
\olpseganchor{OLP-0089-B009}{end}

\olpseganchor{OLP-0089-B010}{start}
\begin{ex}
اوس راځئ چې د~\LR{$(\lnot !A \lor !B) \lif (!A \lif !B)$} يو
اشتقاق ورکړو۔
\olpseganchor{OLP-0089-B010}{end}

\olpseganchor{OLP-0089-B011}{start}
په پيل کښې مطلوبه پايله د اشتقاق په لاندې برخه کښې ليکو۔
\begin{prooftree}
\AxiomC{}
\UnaryInfC{\LR{$(\lnot !A \lor !B) \lif (!A \lif !B)$}}
\end{prooftree}
د هغې منطقي قاعدې د موندلو دپاره چې دا پايله راکولے شي، په پايله
کښې منطقي رابطونو~\LR{$\lnot$}، \LR{$\lor$} او~\LR{$\lif$} ته ګورو۔ اوس يوازې د
\LR{$\lif$} لومړۍ پېښه راته مهمه ده، ځکه همدا د وروستۍ پايلې د
تړلے فارمول~اصلي عملګر دے؛ \LR{$\lnot$}، \LR{$\lor$} او د~\LR{$\lif$}
دويمه پېښه د بل رابط په ساحه کښې دي، نو وروسته به ئې وڅېړو۔ له دې
امله د~\Intro{\lif} له قاعدې پيل کوو۔ سمه کارونه ئې بايد داسې وي:
\textbf{د ژباړې د نسخې سمون (OLND-004):} سرچينه دلته د فطري
استنتاج وروستۍ جمله “په وروستي سېکوېنټ کښې” بولي؛ دلته “وروستۍ
پايله” راوړل شوې، ځکه دا ونه د جملو ثبوت دے او سېکوېنټ نۀ لري۔
\begin{prooftree}
\AxiomC{\LR{$\Discharge{\lnot !A \lor !B}{1}$}}
\DeduceC{\LR{$!A \lif !B$}}
\DischargeRule{\Intro{\lif}}{1}
\UnaryInfC{\LR{$(\lnot !A \lor !B) \lif (!A \lif !B)$}}
\end{prooftree}
له دې وروسته د دوام دوه لارې لرو۔ يا له لاندې څخه پورته کار روان
ساتلے او د~\Intro{\lif} بله کارونه لټولے شو، يا له پاسه لاندې کار
کولے او د~\Elim{\lor} قاعده کارولے شو۔ دويمه لار اخلو۔ فرضيه
\LR{$\lnot !A \lor !B$} د~\Elim{\lor} د تر ټولو کيڼې مقدمې په توګه
کاروو۔ د~\Elim{\lor} د سمې کارونې دپاره نورې دوه مقدمې بايد له
پايلې~\LR{$!A \lif !B$} سره يو شان وي، خو هره يوه په خپل وار له بلې
فرضیې، يعنې د~\LR{$\lnot !A \lor !B$} له دوو بېلندونو څخه له يوه،
رايستل کېدے شي۔ نو اشتقاق به داسې وي:
\begin{prooftree}
\AxiomC{\LR{$\Discharge{\lnot !A \lor !B}{1}$}}
\AxiomC{\LR{$\Discharge{\lnot !A}{2}$}}
\DeduceC{\LR{$!A \lif !B$}}
\AxiomC{\LR{$\Discharge{!B}{2}$}}
\DeduceC{\LR{$!A \lif !B$}}
\DischargeRule{\Elim{\lor}}{2}
\TrinaryInfC{\LR{$!A \lif !B$}}
\DischargeRule{\Intro{\lif}}{1} 
\UnaryInfC{\LR{$(\lnot !A \lor !B) \lif (!A \lif !B)$}}
\end{prooftree}
\olpseganchor{OLP-0089-B011}{end}

\olpseganchor{OLP-0089-B012}{start}
په ښي لوري دواړو څانګو کښې غواړو \LR{$!A \lif !B$}~اشتقاق کړو؛
غوره لار ئې د~\Intro{\lif} کارول دي۔
\begin{prooftree}
\AxiomC{\LR{$\Discharge{\lnot !A \lor !B}{1}$}}
\AxiomC{\LR{$\Discharge{\lnot !A}{2}, \Discharge{!A}{3}$}}
\DeduceC{\LR{$!B$}}
\DischargeRule{\Intro{\lif}}{3}
\UnaryInfC{\LR{$!A \lif !B$}}
\AxiomC{\LR{$\Discharge{!B}{2}, \Discharge{!A}{4}$}}
\DeduceC{\LR{$!B$}}
\DischargeRule{\Intro{\lif}}{4}
\UnaryInfC{\LR{$!A \lif !B$}}
\DischargeRule{\Elim{\lor}}{2}
\TrinaryInfC{\LR{$!A \lif !B$}}
\DischargeRule{\Intro{\lif}}{1} 
\UnaryInfC{\LR{$(\lnot !A \lor !B) \lif (!A \lif !B)$}}
\end{prooftree}
\olpseganchor{OLP-0089-B012}{end}

\olpseganchor{OLP-0089-B013}{start}
د اشتقاق په دوو تشو برخو کښې له~\LR{$\lnot !A$} او~\LR{$!A$} څخه،
او له~\LR{$!A$} او~\LR{$!B$} څخه، د~\LR{$!B$} اشتقاقونه غواړو۔ لومړۍ برخه
اخلو۔ \LR{$\lnot !A$} او~\LR{$!A$} د~\Elim{\lnot} دوه مقدمې دي:
\begin{prooftree}
\AxiomC{\LR{$\Discharge{\lnot !A}{2}$}}
\AxiomC{\LR{$\Discharge{!A}{3}$}}
\RightLabel{\Elim{\lnot}}
\BinaryInfC{\LR{$\lfalse$}}
\DeduceC{\LR{$!B$}}
\end{prooftree}
د~\FalseInt په کارولو \LR{$!B$} د پايلې په توګه اخستلے او څانګه بشپړولے
شو۔
\begin{prooftree}
\AxiomC{\LR{$\Discharge{\lnot !A \lor !B}{1}$}}
\AxiomC{\LR{$\Discharge{\lnot !A}{2}$}}
\AxiomC{\LR{$\Discharge{!A}{3}$}}
\RightLabel{\Elim{\lnot}}
\BinaryInfC{\LR{$\lfalse$}}
\RightLabel{\FalseInt}
\UnaryInfC{\LR{$!B$}}
\DischargeRule{\Intro{\lif}}{3}
\UnaryInfC{\LR{$!A \lif !B$}}
\AxiomC{\LR{$\Discharge{!B}{2}, \Discharge{!A}{4}$}}
\DeduceC{\LR{$!B$}}
\DischargeRule{\Intro{\lif}}{4}
\UnaryInfC{\LR{$!A \lif !B$}}
\DischargeRule{\Elim{\lor}}{2}
\TrinaryInfC{\LR{$!A \lif !B$}}
\DischargeRule{\Intro{\lif}}{1} 
\UnaryInfC{\LR{$(\lnot !A \lor !B) \lif (!A \lif !B)$}}
\end{prooftree}
\textbf{د ژباړې د نسخې سمون (OLND-003):} سرچينه په دې بشپړه ونه
کښې له نفي شوې مقدمې او مثبتې مقدمې څخه کاذبۍ ته استنتاج د کاذبۍ
د داخلولو په نښان ښيي؛ د مخکښې جزوي ونې او ټاکلې قاعدې سره سم دلته
د نفي ايستلو نښان راوړل شوے دے۔
\olpseganchor{OLP-0089-B013}{end}

\olpseganchor{OLP-0089-B014}{start}
اوس تر ټولو ښۍ څانګې ته ګورو۔ دلته بايد پوه شو چې د اشتقاق
تعريف \emph{د فرضيو ايستل روا بولي} خو \emph{لازمي ئې نۀ بولي}۔ په
بله وينا، که له فرضيو~\LR{$!A$} او~\LR{$!B$} څخه د يوې په کارولو \LR{$!B$} راوباسو
او بله و نۀ کاروو، دا روا ده۔ له~\LR{$!B$} څخه د~\LR{$!B$}~اشتقاق کول
څرګند دي: \LR{$!B$} په يوازې ځان داسې اشتقاق دے او هېڅ استنتاج
نۀ غواړي۔ نو فرضيه~\LR{$!A$} ساده ړنګولے شو۔
\begin{prooftree}
\AxiomC{\LR{$\Discharge{\lnot !A \lor !B}{1}$}}
\AxiomC{\LR{$\Discharge{\lnot !A}{2}$}}
\AxiomC{\LR{$\Discharge{!A}{3}$}}
\RightLabel{\Elim{\lnot}}
\BinaryInfC{\LR{$\lfalse$}}
\RightLabel{\FalseInt}
\UnaryInfC{\LR{$!B$}}
\DischargeRule{\Intro{\lif}}{3}
\UnaryInfC{\LR{$!A \lif !B$}}
\AxiomC{\LR{$\Discharge{!B}{2}$}}
\RightLabel{\Intro{\lif}}
\UnaryInfC{\LR{$!A \lif !B$}}
\DischargeRule{\Elim{\lor}}{2}
\TrinaryInfC{\LR{$!A \lif !B$}}
\DischargeRule{\Intro{\lif}}{1} 
\UnaryInfC{\LR{$(\lnot !A \lor !B) \lif (!A \lif !B)$}}
\end{prooftree}
پام وکړئ چې په بشپړ اشتقاق کښې د تر ټولو ښي اړخ
\Intro{\lif} استنتاج په حقيقت کښې هېڅ فرضيه نۀ باسي۔
\end{ex}
\olpseganchor{OLP-0089-B014}{end}

\olpseganchor{OLP-0089-B015}{start}
\begin{ex}
تر اوسه مو د~\FalseCl{} قاعدې ته اړتيا نۀ ده پېښه شوې۔ دا ځکه
ځانګړې ده چې داسې فرضيه ايستل راکوي چې د قاعدې د پايلې فرعي
فارمول نۀ وي۔ دا له~\FalseInt{} قاعدې سره نژدې تړلې ده۔ په
حقيقت کښې \FalseInt{} د~\FalseCl{} قاعدې ځانګړے حالت دے---د
“شهودي منطق” په نامه يو منطق شته چې يوازې \FalseInt{} پکې روا دے۔
کله چې بله هېڅ لار کار نۀ کوي، \FalseCl{} وروستۍ وسيله ده۔ مثلاً
فرض کړئ غواړو \LR{$!A \lor \lnot !A$}~اشتقاق کړو۔ زموږ عادي تګلاره
به دا وي چې د~\LR{$!A \lor \lnot !A$} د~\LR{$\Intro{\lor}$} په کارولو
اشتقاق کړو۔ خو دا
غواړي چې له هېڅ فرضیې پرته يا~\LR{$!A$} يا~\LR{$\lnot !A$}~اشتقاق کړو،
او دا نۀ کېږي۔ \FalseCl{} مو ژغوري!
\begin{prooftree}
  \AxiomC{\LR{$\Discharge{\lnot(!A \lor \lnot !A)}{1}$}}
  \DeduceC{\LR{$\lfalse$}}
  \DischargeRule{\FalseCl}{1}
  \UnaryInfC{\LR{$!A \lor \lnot !A$}}
\end{prooftree}
اوس له~\LR{$\lnot(!A \lor \lnot !A)$} څخه د~\LR{$\lfalse$} يو
اشتقاق لټوو۔ ځکه \LR{$\lfalse$} د~\LR{$\Elim{\lnot}$} پايله ده، نو
کېدے شي دا هڅه وکړو:
\begin{prooftree}
  \AxiomC{\LR{$\Discharge{\lnot(!A \lor \lnot !A)}{1}$}}
  \DeduceC{\LR{$\lnot !A$}}
  \AxiomC{\LR{$\Discharge{\lnot(!A \lor \lnot !A)}{1}$}}
  \DeduceC{\LR{$!A$}}
  \RightLabel{\Elim{\lnot}}
  \BinaryInfC{\LR{$\lfalse$}}
  \DischargeRule{\FalseCl}{1}
  \UnaryInfC{\LR{$!A \lor \lnot !A$}}
\end{prooftree}
د~\LR{$\lnot !A$} د اشتقاق موندلو تګلاره د~\LR{$\Intro{\lnot}$}
کارونه غواړي:
\begin{prooftree}
  \AxiomC{\LR{$\Discharge{\lnot(!A \lor \lnot !A)}{1}, \Discharge{!A}{2}$}}
  \DeduceC{\LR{$\lfalse$}}
  \DischargeRule{\Intro{\lnot}}{2}
  \UnaryInfC{\LR{$\lnot !A$}}
  \AxiomC{\LR{$\Discharge{\lnot(!A \lor \lnot !A)}{1}$}}
  \DeduceC{\LR{$!A$}}
  \RightLabel{\Elim{\lnot}}
  \BinaryInfC{\LR{$\lfalse$}}
  \DischargeRule{\FalseCl}{1}
  \UnaryInfC{\LR{$!A \lor \lnot !A$}}
\end{prooftree}
دلته \LR{$\lfalse$} په اسانۍ اخستلے شو: \LR{$\Elim{\lnot}$} پر
\LR{$\lnot(!A \lor \lnot !A)$} او~\LR{$!A \lor \lnot !A$} کاروو؛ دويمه
جمله زموږ له نوې فرضيې~\LR{$!A$} څخه د~\LR{$\Intro{\lor}$} په وسيله راځي:
\begin{prooftree}
  \AxiomC{\LR{$\Discharge{\lnot(!A \lor \lnot !A)}{1}$}}
  \AxiomC{\LR{$\Discharge{!A}{2}$}}
  \RightLabel{\Intro{\lor}}
  \UnaryInfC{\LR{$!A \lor \lnot !A$}}
  \RightLabel{\Elim{\lnot}}
  \BinaryInfC{\LR{$\lfalse$}}
  \DischargeRule{\Intro{\lnot}}{2}
  \UnaryInfC{\LR{$\lnot !A$}}
  \AxiomC{\LR{$\Discharge{\lnot(!A \lor \lnot !A)}{1}$}}
  \DeduceC{\LR{$!A$}}
  \RightLabel{\Elim{\lnot}}
  \BinaryInfC{\LR{$\lfalse$}}
  \DischargeRule{\FalseCl}{1}
  \UnaryInfC{\LR{$!A \lor \lnot !A$}}
\end{prooftree}
په ښي لوري کښې همدا تګلاره کاروو، خو \LR{$!A$} د~\FalseCl په وسيله
اخلو:
\begin{prooftree}
  \AxiomC{\LR{$\Discharge{\lnot(!A \lor \lnot !A)}{1}$}}
  \AxiomC{\LR{$\Discharge{!A}{2}$}}
  \RightLabel{\Intro{\lor}}
  \UnaryInfC{\LR{$!A \lor \lnot !A$}}
  \RightLabel{\Elim{\lnot}}
  \BinaryInfC{\LR{$\lfalse$}}
  \DischargeRule{\Intro{\lnot}}{2}
  \UnaryInfC{\LR{$\lnot !A$}}
  \AxiomC{\LR{$\Discharge{\lnot(!A \lor \lnot !A)}{1}$}}
  \AxiomC{\LR{$\Discharge{\lnot !A}{3}$}}
  \RightLabel{\Intro{\lor}}
  \UnaryInfC{\LR{$!A \lor \lnot !A$}}
  \RightLabel{\Elim{\lnot}}
  \BinaryInfC{\LR{$\lfalse$}}
  \DischargeRule{\FalseCl}{3}
  \UnaryInfC{\LR{$!A$}}
  \RightLabel{\Elim{\lnot}}
  \BinaryInfC{\LR{$\lfalse$}}
  \DischargeRule{\FalseCl}{1}
  \UnaryInfC{\LR{$!A \lor \lnot !A$}}
\end{prooftree}
\end{ex}
\olpseganchor{OLP-0089-B015}{end}

\olpseganchor{OLP-0089-B016}{start}
\begin{prob}
داسې اشتقاقونه ورکړئ چې لاندې وښيي:
\begin{enumerate}
\item \LR{$!A \land (!B \land !C) \Proves (!A \land !B) \land !C$}.
\item \LR{$!A \lor (!B \lor !C) \Proves (!A \lor !B) \lor !C$}.
\item \LR{$!A \lif (!B \lif !C) \Proves !B \lif (!A \lif !C)$}.
\item \LR{$!A \Proves \lnot\lnot !A$}.
\end{enumerate}
\end{prob}
\olpseganchor{OLP-0089-B016}{end}

\olpseganchor{OLP-0089-B017}{start}
\begin{prob}
داسې اشتقاقونه ورکړئ چې لاندې وښيي:
\begin{enumerate}
\item \LR{$(!A \lor !B) \lif !C \Proves !A \lif !C$}.
\item \LR{$(!A \lif !C) \land (!B \lif !C) \Proves (!A \lor !B) \lif !C$}.
\item \LR{$\Proves \lnot(!A \land \lnot !A)$}.
\item \LR{$!B \lif !A \Proves \lnot !A \lif \lnot !B$}.
\item \LR{$\Proves (!A \lif \lnot !A) \lif \lnot !A$}.
\item \LR{$\Proves \lnot(!A \lif !B) \lif \lnot !B$}.
\item \LR{$!A \lif !C \Proves \lnot (!A \land \lnot !C)$}.
\item \LR{$!A \land \lnot !C \Proves \lnot (!A \lif !C)$}.
\item \LR{$!A \lor !B, \lnot !B \Proves !A$}.
\item \LR{$\lnot !A \lor \lnot !B \Proves \lnot(!A \land !B)$}.
\item \LR{$\Proves (\lnot !A \land \lnot !B) \lif\lnot(!A \lor !B)$}.
\item \LR{$\Proves \lnot(!A \lor !B) \lif (\lnot !A \land \lnot !B)$}.
\end{enumerate}
\end{prob}
\olpseganchor{OLP-0089-B017}{end}

\olpseganchor{OLP-0089-B018}{start}
\begin{prob}
داسې اشتقاقونه ورکړئ چې لاندې وښيي:
\begin{enumerate}
\item \LR{$\lnot(!A \lif !B) \Proves !A$}.
\item \LR{$\lnot(!A \land !B) \Proves \lnot !A \lor \lnot !B$}.
\item \LR{$!A \lif !B \Proves \lnot !A \lor !B$}.
\item \LR{$\Proves \lnot \lnot !A \lif !A$}.
\item \LR{$!A \lif !B, \lnot !A \lif !B \Proves !B$}.
\item \LR{$(!A \land !B) \lif !C \Proves (!A \lif !C) \lor (!B \lif !C)$}.
\item \LR{$(!A \lif !B) \lif !A \Proves !A$}.
\item \LR{$\Proves (!A \lif !B) \lor (!B \lif !C)$}.
\end{enumerate}
(دا ټول د~\LR{$\FalseCl$} قاعده غواړي۔)
\end{prob}
\olpseganchor{OLP-0089-B018}{end}

\olpseganchor{OLP-0090-B004}{start}
\olfileid{fol}{ntd}{prq}
\olpseganchor{OLP-0090-B004}{end}

\olpseganchor{OLP-0090-B005}{start}
\olsection{له سورونو سره اشتقاقونه}
\olpseganchor{OLP-0090-B005}{end}

\olpseganchor{OLP-0090-B006}{start}
\begin{ex}
له سورونو سره د کار پر مهال بايد ډاډ واخلو چې د ځانګړي متغير شرط
نۀ ماتوو، او کله کله د دې دپاره د ځينو استنتاجونو د عملي کولو ترتيب
بدلول را بدلول غواړي۔ په عمومي توګه ګټه کوي چې لومړے د هغو قاعدو
د حل هڅه وکړو چې د ځانګړي متغير شرط ورباندې لګېږي (په بشپړ ثبوت
کښې به دوئ لاندې راځي)۔
\olpseganchor{OLP-0090-B006}{end}

\olpseganchor{OLP-0090-B007}{start}
راځئ وګورو چې د فارمول~\LR{$\lexists[x][\lnot !A(x)] \lif \lnot
\lforall[x][!A(x)]$} يو اشتقاق څنګه ورکوو۔ د معمول په شان
داسې پيل کوو:
\begin{prooftree}
\AxiomC{}
\UnaryInfC{\LR{$\lexists[x][\lnot !A(x)]\lif \lnot \lforall[x][!A(x)]$}}
\end{prooftree}
لومړے ليکو چې د~\Intro{\lif} د قاعدې په وسيله د وروستي پړاو د
توجيه کولو دپاره څه پکار دي۔
\begin{prooftree}
\AxiomC{\LR{$\Discharge{\lexists[x][\lnot !A(x)]}{1}$}}
\DeduceC{\LR{$\lnot \lforall[x][!A(x)]$}}
\DischargeRule{\Intro{\lif}}{1}
\UnaryInfC{\LR{$\lexists[x][\lnot !A(x)]\lif \lnot \lforall[x][!A(x)]$}}
\end{prooftree}
ځکه پر~\LR{$\lnot \lforall[x][!A(x)]$} د کارولو دپاره څرګنده قاعده نۀ
شته، اشتقاق داسې برابروو چې د~\Elim{\lexists} قاعده
وکارولے شو۔ دلته بايد د ځانګړي متغير شرط ته پام وکړو او داسې ثابت
وټاکو چې په~\LR{$\lexists[x][!A(x)]$} يا د هغې په هرې اړوندې فرضيې کښې
نۀ راځي۔ (خو ځکه چې هېڅ ثابت سمبولونه نۀ دي راغلي، هر انتخاب سم
دے۔)
\begin{prooftree}
\AxiomC{\LR{$\Discharge{\lexists[x][\lnot !A(x)]}{1}$}}
\AxiomC{\LR{$\Discharge{\lnot !A(a)}{2}$}}
\DeduceC{\LR{$\lnot \lforall[x][!A(x)]$}}
\DischargeRule{\Elim{\lexists}}{2}
\BinaryInfC{\LR{$\lnot \lforall[x][!A(x)]$}}
\DischargeRule{\Intro{\lif}}{1}
\UnaryInfC{\LR{$\lexists[x][\lnot !A(x)] \lif \lnot \lforall[x][!A(x)]$}}
\end{prooftree}
د~\LR{$\lnot \lforall[x][!A(x)]$} د راايستلو دپاره د~\Intro{\lnot}
قاعده کارول غواړو: دا غواړي چې يو تناقض راوباسو، او کېدے شي
\LR{$\lforall[x][!A(x)]$} د يوې اضافي فرضيې په توګه وکاروو۔ البته په دې
تناقض کښې هغه فرضيه~\LR{$\lnot !A(a)$} هم راتلے شي چې د
\Elim{\lexists} استنتاج به ئې وباسي۔ داسې ئې برابروو:
\begin{prooftree}
\AxiomC{\LR{$\Discharge{\lexists[x][\lnot !A(x)]}{1}$}}
\AxiomC{\LR{$\Discharge{\lnot !A(a)}{2}, \Discharge{\lforall[x][!A(x)]}{3}$}}
\DeduceC{\LR{$\lfalse$}}
\DischargeRule{\Intro{\lnot}}{3}
\UnaryInfC{\LR{$\lnot \lforall[x][!A(x)]$}}
\DischargeRule{\Elim{\lexists}}{2}
\BinaryInfC{\LR{$\lnot \lforall[x][!A(x)]$}}
\DischargeRule{\Intro{\lif}}{1}
\UnaryInfC{\LR{$\lexists[x][\lnot !A(x)]\lif \lnot \lforall[x][!A(x)]$}}
\end{prooftree}
داسې ښکاري چې تناقض ته نژدې يو۔ تر ټولو اسانه قاعده
\Elim{\lforall} ده چې د ځانګړي متغير هېڅ شرط نۀ لري۔ ځکه د عمومي
سور لرونکي~\LR{$x$} پر ځاے هره تړلې اصطلاح کارولے شو، ښه ده چې هماغه
~\LR{$a$} وکاروو څو تناقض ته ورسېږو۔
\begin{prooftree}
\AxiomC{\LR{$\Discharge{\lexists[x][\lnot !A(x)]}{1}$}}
\AxiomC{\LR{$\Discharge{\lnot !A(a)}{2}$}}
\AxiomC{\LR{$\Discharge{\lforall[x][!A(x)]}{3}$}}
\RightLabel{\Elim{\lforall}}
\UnaryInfC{\LR{$!A(a)$}}
\RightLabel{\Elim{\lnot}}
\BinaryInfC{\LR{$\lfalse$}}
\DischargeRule{\Intro{\lnot}}{3}
\UnaryInfC{\LR{$\lnot \lforall[x][!A(x)]$}}
\DischargeRule{\Elim{\lexists}}{2}
\BinaryInfC{\LR{$\lnot \lforall[x][!A(x)]$}}
\DischargeRule{\Intro{\lif}}{1}
\UnaryInfC{\LR{$\lexists[x][\lnot !A(x)]\lif \lnot \lforall[x][!A(x)]$}}
\end{prooftree}
\olpseganchor{OLP-0090-B007}{end}

\olpseganchor{OLP-0090-B008}{start}
په ځانګړي ډول له سورونو سره د کار پر مهال اوس دا بيا کتل مهم دي
چې د ځانګړي متغير شرط نۀ دے مات شوے۔ دلته دغه شرط يوازې پر
\Elim{\exists} لګېدۀ، او ځانګړے متغير~\LR{$a$} په هغو فرضيو کښې نۀ
راځي چې استنتاج ورپورې تړلے دے؛ نو دا سم اشتقاق دے۔
\end{ex}
\olpseganchor{OLP-0090-B008}{end}

\olpseganchor{OLP-0090-B009}{start}
\begin{ex}
کله کله له نورو فارمولونه څخه فارمول راايستلے شو۔ په دغو
حالتونو کښې کېدے شي نۀ ايستل شوې فرضيې ولرو۔ د خپلې وروستۍ موخې
تر څنګ د فرضيو حساب ساتل هم مهم دي۔
\olpseganchor{OLP-0090-B009}{end}

\olpseganchor{OLP-0090-B010}{start}
راځئ وګورو چې له فرضيو~\LR{$\lexists[x][(!A(x) \land !B(x))]$} او
\LR{$\lforall[x][(!B(x) \lif !C(x,b))]$} څخه د فارمول
\LR{$\lexists[x][!C(x,b)]$} يو اشتقاق څنګه ورکوو۔ د معمول په
شان پايله لاندې ليکو۔
\begin{prooftree}
\AxiomC{}
\UnaryInfC{\LR{$\lexists[x][!C(x,b)]$}}
\end{prooftree}
\olpseganchor{OLP-0090-B010}{end}

\olpseganchor{OLP-0090-B011}{start}
د کار دپاره دوه مقدمې لرو۔ د لومړۍ د کارولو، يعنې له
\LR{$\lexists[x][(!A(x) \land !B(x))]$} څخه د
\LR{$\lexists[x][!C(x, b)]$} د اشتقاق موندلو دپاره به د
\Elim{\lexists} قاعده کاروو۔ ځکه د ځانګړي متغير شرط لري، دا قاعده
لومړۍ کاروو۔ لاندې لاس ته راځي:
\begin{prooftree}
\AxiomC{\LR{$\lexists[x][(!A(x) \land !B(x))]$}}
\AxiomC{\LR{$\Discharge{!A(a) \land !B(a)}{1}$}}
\DeduceC{\LR{$\lexists[x][!C(x,b)]$}}
\DischargeRule{\Elim{\lexists}}{1}
\BinaryInfC{\LR{$\lexists[x][!C(x,b)]$}}
\end{prooftree}
هغه دوه فرضيې چې ورسره کار کوو، \LR{$!B$} شريکوي۔ ښايي اوس د
\Elim{\land} په کارولو د~\LR{$!B(a)$} بېلول ګټور وي۔
\begin{prooftree}
\AxiomC{\LR{$\lexists[x][(!A(x) \land !B(x)])$}}
\AxiomC{\LR{$\Discharge{!A(a) 
\land !B(a)}{1}$}}
\RightLabel{\Elim{\land}}
\UnaryInfC{\LR{$!B(a)$}}
\DeduceC{\LR{$\lexists[x][!C(x,b)]$}}
\DischargeRule{\Elim{\lexists}}{1}
\BinaryInfC{\LR{$\lexists[x][!C(x,b)]$}}
\end{prooftree}
\olpseganchor{OLP-0090-B011}{end}

\olpseganchor{OLP-0090-B012}{start}
د کار دپاره دويمه فرضيه~\LR{$\lforall[x][(!B(x) \lif !C(x,b))]$} ده۔
ځکه د ځانګړي متغير شرط نۀ لري، د~\Elim{\lforall} په کارولو د
ثابت سمبول~\LR{$a$} په وسيله \LR{$x$} بېلګه کولے او~\LR{$!B(a) \lif !C(a, b)$}
اخستلے شو۔ اوس \LR{$!B(a) \lif !C(a,b)$} او~\LR{$!B(a)$} دواړه لرو۔ بل ګام
بايد د~\Elim{\lif} نېغه کارونه وي۔
\begin{prooftree}
\AxiomC{\LR{$\lexists[x][(!A(x) \land !B(x))]$}}
\AxiomC{\LR{$\lforall[x][(!B(x) \lif !C(x,b))]$}}
\RightLabel{\Elim{\lforall}}
\UnaryInfC{\LR{$!B(a) \lif !C(a,b)$}}
\AxiomC{\LR{$\Discharge{!A(a) 
\land !B(a)}{1}$}}
\RightLabel{\Elim{\land}}
\UnaryInfC{\LR{$!B(a)$}}
\RightLabel{\Elim{\lif}}
\BinaryInfC{\LR{$!C(a,b)$}}
\DeduceC{\LR{$\lexists[x][!C(x,b)]$}}
\DischargeRule{\Elim{\lexists}}{1}
\insertBetweenHyps{\hspace{-5em}}
\BinaryInfC{\LR{$\lexists[x][!C(x,b)]$}}
\end{prooftree}
موخې ته ډېر نژدې يو!{} د~\Intro{\lexists} يوه کارونه موخې ته مو
رسوي۔
\begin{prooftree}
\AxiomC{\LR{$\lexists[x][(!A(x) \land !B(x))]$}}
\AxiomC{\LR{$\lforall[x][(!B(x) \lif !C(x,b))]$}}
\RightLabel{\Elim{\lforall}}
\UnaryInfC{\LR{$!B(a) \lif !C(a,b)$}}
\AxiomC{\LR{$\Discharge{!A(a) 
\land !B(a)}{1}$}}
\RightLabel{\Elim{\land}}
\UnaryInfC{\LR{$!B(a)$}}
\RightLabel{\Elim{\lif}}
\BinaryInfC{\LR{$!C(a,b)$}}
\RightLabel{\Intro{\lexists}}
\UnaryInfC{\LR{$\lexists[x][!C(x,b)]$}}
\DischargeRule{\Elim{\lexists}}{1}
\insertBetweenHyps{\hspace{-5em}}
\BinaryInfC{\LR{$\lexists[x][!C(x,b)]$}}
\end{prooftree}
ځکه په هر پړاو کښې مو ډاډ اخستے چې د ځانګړي متغير شرطونه نۀ دي
مات شوي، باور کولے شو چې دا سم اشتقاق دے۔
\end{ex}
\olpseganchor{OLP-0090-B012}{end}

\olpseganchor{OLP-0090-B013}{start}
\begin{ex}
له فرضيو~\LR{$\lforall[x][!A(x)] \lif \lexists[y][!B(y)]$} او
\LR{$\lnot\lexists[y][!B(y)]$} څخه د فارمول
\LR{$\lnot\lforall[x][!A(x)]$} يو اشتقاق ورکړئ۔ د معمول په شان
موخې فارمول لاندې ليکو۔
\begin{prooftree}
\AxiomC{}
\UnaryInfC{\LR{$\lnot\lforall[x][!A(x)]$}}
\end{prooftree}
د اشتقاق وروستۍ کرښه نفي ده، نو \Intro{\lnot} کارول غواړو۔
دا غواړي چې معلومه کړو تناقض څنګه اشتقاق کړو۔
\begin{prooftree}
\AxiomC{\LR{$\Discharge{\lforall[x][!A(x)]}{1}$}}
\DeduceC{\LR{$\lfalse$}}
\DischargeRule{\Intro{\lnot}}{1}
\UnaryInfC{\LR{$\lnot\lforall[x][!A(x)]$}}
\end{prooftree}
تر دې ځايه ښه يو۔ \Elim{\lforall} کارولے شو، خو څرګنده نۀ ده چې
موخې ته به مو ورسوي۔ پر ځاے ئې خپله يوه فرضيه کاروو۔
\LR{$\lforall[x][!A(x)] \lif \lexists[y][!B(y)]$} له
\LR{$\lforall[x][!A(x)]$} سره د~\Elim{\lif} قاعدې کارول راکوي۔
\begin{prooftree}
\AxiomC{\LR{$\lforall[x][!A(x)] \lif \lexists[y][!B(y)]$}}
\AxiomC{\LR{$\Discharge{\lforall[x][!A(x)]}{1}$}}
\RightLabel{\Elim{\lif}}
\BinaryInfC{\LR{$\lexists[y][!B(y)]$}}
\DeduceC{\LR{$\lfalse$}}
\DischargeRule{\Intro{\lnot}}{1}
\UnaryInfC{\LR{$\lnot\lforall[x][!A(x)]$}}
\end{prooftree}
اوس د کار دپاره يوه وروستۍ فرضيه لرو، او ښکاري چې د~\Elim{\lnot}
په کارولو تناقض ته مو رسولے شي۔
\begin{prooftree}
\AxiomC{\LR{$\lnot\lexists[y][!B(y)]$}}
\AxiomC{\LR{$\lforall[x][!A(x)] \lif \lexists[y][!B(y)]$}}
\AxiomC{\LR{$\Discharge{\lforall[x][!A(x)]}{1}$}}
\RightLabel{\Elim{\lif}}
\BinaryInfC{\LR{$\lexists[y][!B(y)]$}}
\RightLabel{\Elim{\lnot}}
\BinaryInfC{\LR{$\lfalse$}}
\DischargeRule{\Intro{\lnot}}{1}
\UnaryInfC{\LR{$\lnot\lforall[x][!A(x)]$}}
\end{prooftree}
\end{ex}
\olpseganchor{OLP-0090-B013}{end}

\olpseganchor{OLP-0090-B014}{start}
\begin{prob}
داسې اشتقاقونه ورکړئ چې لاندې وښيي:
\begin{enumerate}
\item \LR{$\Proves (\lforall[x][!A(x)] \land \lforall[y][!B(y)]) \lif
\lforall[z][(!A(z) \land !B(z))]$}.
\item \LR{$\Proves (\lexists[x][!A(x)] \lor \lexists[y][!B(y)]) \lif
\lexists[z][(!A(z) \lor !B(z))]$}.
\item \LR{$\lforall[x][(!A(x) \lif !B)] \Proves \lexists[y][!A(y)] \lif !B$}.
\item \LR{$\lforall[x][\lnot !A(x)] \Proves \lnot\lexists[x][!A(x)]$}.
\item \LR{$\Proves \lnot\lexists[x][!A(x)] \lif \lforall[x][\lnot !A(x)]$}.
\item \LR{$\Proves \lnot\lexists[x][\lforall[y][((!A(x,y) \lif \lnot
!A(y,y)) \land (\lnot !A(y,y) \lif !A(x,y)))]]$}.
\end{enumerate}
\end{prob}
\olpseganchor{OLP-0090-B014}{end}

\olpseganchor{OLP-0090-B015}{start}
\begin{prob}
داسې اشتقاقونه ورکړئ چې لاندې وښيي:
\begin{enumerate}
\item \LR{$\Proves \lnot\lforall[x][!A(x)] \lif \lexists[x][\lnot!A(x)]$}.
\item \LR{$(\lforall[x][!A(x)] \lif !B) \Proves \lexists[y][(!A(y) \lif !B)]$}.
\item \LR{$\Proves \lexists[x][(!A(x) \lif \lforall[y][!A(y)])]$}.
\end{enumerate}
(دا ټول د~\LR{$\FalseCl$} قاعده غواړي۔)
\end{prob}
\olpseganchor{OLP-0090-B015}{end}

\olpseganchor{OLP-0091-B004}{start}
\olfileid{fol}{ntd}{ptn}
\olpseganchor{OLP-0091-B004}{end}

\olpseganchor{OLP-0091-B005}{start}
\olsection{ثبوت-تيوريکي مفاهيم}
\olpseganchor{OLP-0091-B005}{end}

\olpseganchor{OLP-0091-B006}{start}
\begin{editorial}
دا برخه د فطري استنتاج د اثبات‌پذيرۍ اړيکې او سازګارۍ تعريفونه
راټولوي۔
\end{editorial}
\olpseganchor{OLP-0091-B006}{end}

\olpseganchor{OLP-0091-B007}{start}
\begin{explain}
لکه څنګه چې مو يو شمېر مهم معنايي مفاهيم (اعتبار، معنايي استلزام،
د صدق وړتيا) تعريف کړل، اوس ورسره سم \emph{ثبوت-تيوريکي مفاهيم}
تعريفوو۔ دا مفاهيم په جوړښتونه کښې د تړلي فارمولونه د پوره
کېدو له لارې نۀ، بلکې له نورو څخه د ځينو تړلي فارمولونه د
د اشتقاق وړتيا يا د اشتقاق نۀ وړتيا له لارې تعريفېږي۔ دا مهمه
موندنه وه چې دغه مفاهيم سره سمون خوري۔ د دې سمون منځپانګه د
\emph{سموالي} او \emph{بشپړتيا قضيې} دي۔
\end{explain}
\olpseganchor{OLP-0091-B007}{end}


\olpseganchor{OLP-0091-B008}{start}
\begin{defn}[قضيې]
يوه تړلے فارمول~\LR{$!A$} هله \emph{قضيه} ده چې په فطري استنتاج کښې
د~\LR{$!A$} داسې اشتقاق وي چې ټولې فرضيې پکې ايستل
شوې وي۔ که \LR{$!A$} قضيه وي، \LR{$\Proves !A$} ليکو، او که نۀ وي
\LR{$\Proves/ !A$} ليکو۔
\end{defn}
\olpseganchor{OLP-0091-B008}{end}

\olpseganchor{OLP-0091-B009}{start}
\begin{defn}[د اشتقاق وړتيا]
تړلے فارمول~\LR{$!A$} د تړلي فارمولونه له سټ~\LR{$\Gamma$} څخه،
\LR{$\Gamma \Proves !A$}، هله \emph{د اشتقاق وړ} ده چې داسې
اشتقاق وي چې پايله ئې~\LR{$!A$} وي او هره فرضيه پکې يا
ايستل شوې وي يا په~\LR{$\Gamma$} کښې وي۔ که \LR{$!A$} له~\LR{$\Gamma$}
څخه د اشتقاق وړ نۀ وي، \LR{$\Gamma \Proves/ !A$} ليکو۔
\end{defn}
\olpseganchor{OLP-0091-B009}{end}

\olpseganchor{OLP-0091-B010}{start}
\begin{defn}[سازګاري]
د تړلي فارمولونه يو سټ~\LR{$\Gamma$} هله او يوازې هله \emph{ناسازګار}
دے چې \LR{$\Gamma \Proves \lfalse$} وي۔ که \LR{$\Gamma$} ناسازګار نۀ وي،
يعنې \LR{$\Gamma \Proves/ \lfalse$} وي، نو \emph{سازګار} ئې بولو۔
\end{defn}
\olpseganchor{OLP-0091-B010}{end}

\olpseganchor{OLP-0091-B011}{start}
\begin{prop}[انعکاسيت]
\ollabel{prop:reflexivity}
که \LR{$!A \in \Gamma$} وي، نو \LR{$\Gamma \Proves !A$}۔
\end{prop}
\olpseganchor{OLP-0091-B011}{end}

\olpseganchor{OLP-0091-B012}{start}
\begin{proof}
فرضيه~\LR{$!A$} په يوازې ځان د~\LR{$!A$} داسې اشتقاق دے چې هره
نۀ ايستل شوې فرضيه ئې (يعنې \LR{$!A$}) په~\LR{$\Gamma$} کښې ده۔
\end{proof}
\olpseganchor{OLP-0091-B012}{end}
  

\olpseganchor{OLP-0091-B013}{start}
\begin{prop}[يکنواختي]
\ollabel{prop:monotonicity}
که \LR{$\Gamma \subseteq \Delta$} او \LR{$\Gamma \Proves !A$} وي، نو
\LR{$\Delta \Proves !A$}۔
\end{prop}
\olpseganchor{OLP-0091-B013}{end}

\olpseganchor{OLP-0091-B014}{start}
\begin{proof}
له~\LR{$\Gamma$} څخه د~\LR{$!A$} هر اشتقاق له~\LR{$\Delta$} څخه هم د~\LR{$!A$}
اشتقاق دے۔
\end{proof}
\olpseganchor{OLP-0091-B014}{end}

\olpseganchor{OLP-0091-B015}{start}
\begin{prop}[انتقاليت]
\ollabel{prop:transitivity}
که \LR{$\Gamma \Proves !A$} او \LR{$\{!A\} \cup \Delta \Proves !B$} وي، نو
\LR{$\Gamma \cup \Delta \Proves !B$}۔
\end{prop}
\olpseganchor{OLP-0091-B015}{end}

\olpseganchor{OLP-0091-B016}{start}
\begin{proof}
که \LR{$\Gamma \Proves !A$} وي، د~\LR{$!A$} داسې اشتقاق~\LR{$\delta_0$}
شته چې ټولې نۀ ايستل شوې فرضيې ئې په~\LR{$\Gamma$} کښې دي۔ که
\LR{$\{!A\} \cup \Delta \Proves !B$} وي، د~\LR{$!B$} داسې
اشتقاق~\LR{$\delta_1$} شته چې ټولې نۀ ايستل شوې فرضيې ئې
په~\LR{$\{!A\} \cup \Delta$} کښې دي۔ اوس لاندې وګورئ:
\begin{prooftree}
  \AxiomC{\LR{$\Delta, \Discharge{!A}{1}$}}
  \RightLabel{\LR{$\delta_1$}}
  \DeduceC{\LR{$!B$}}
  \DischargeRule{\Intro{\lif}}{1}
  \UnaryInfC{\LR{$!A \lif !B$}}
  \AxiomC{\LR{$\Gamma$}}
  \RightLabel{\LR{$\delta_0$}}
  \DeduceC{\LR{$!A$}}
  \RightLabel{\Elim{\lif}}
  \BinaryInfC{\LR{$!B$}}
\end{prooftree}
اوس ټولې نۀ ايستل شوې فرضيې په~\LR{$\Gamma \cup \Delta$} کښې دي؛
نو دا \LR{$\Gamma \cup \Delta \Proves !B$} ښيي۔
\end{proof}
\olpseganchor{OLP-0091-B016}{end}

\olpseganchor{OLP-0091-B017}{start}
کله چې \LR{$\Gamma = \{!A_1, !A_2, \ldots, !A_k\}$} يو متناهي سټ وي،
د~\LR{$\Gamma \Proves !B$} دپاره لنډه نښه
\LR{$!A_1,!A_2,\ldots,!A_k \Proves !B$} کارولے شو؛ په ځانګړي ډول
\LR{$!A \Proves !B$} دا معنٰی لري چې \LR{$\{!A\} \Proves !B$}۔
\olpseganchor{OLP-0091-B017}{end}

\olpseganchor{OLP-0091-B018}{start}
پام وکړئ که \LR{$\Gamma \Proves !A$} او \LR{$!A \Proves !B$} وي، نو
\LR{$\Gamma \Proves !B$}۔ همدارنګه، که
\LR{$!A_1, \dots, !A_n \Proves !B$} او د هر~\LR{$i$} دپاره
\LR{$\Gamma \Proves !A_i$} وي، نو \LR{$\Gamma \Proves !B$}۔
\olpseganchor{OLP-0091-B018}{end}

\olpseganchor{OLP-0091-B019}{start}
\begin{prop}
\ollabel{prop:incons}
لاندې شرطونه معادل دي۔
    \begin{enumerate}
      \item \( \Gamma \) ناسازګار دے۔
      \item د هرې تړلے فارمول~\( {!A} \) دپاره \( \Gamma \Proves {!A} \)۔
      \item د کومې تړلے فارمول~\( {!A} \) دپاره \( \Gamma \Proves {!A} \) او \( \Gamma \Proves \lnot {!A} \)۔
    \end{enumerate}
\end{prop}
\olpseganchor{OLP-0091-B019}{end}

\olpseganchor{OLP-0091-B020}{start}
\begin{proof}
تمرين۔
\end{proof}
\olpseganchor{OLP-0091-B020}{end}

\olpseganchor{OLP-0091-B021}{start}
\begin{prob}
\olref[fol][ntd][ptn]{prop:incons} ثابت کړئ۔
\end{prob}
\olpseganchor{OLP-0091-B021}{end}



\olpseganchor{OLP-0091-B023}{start}
\begin{prop}[تراکم]
  \ollabel{prop:proves-compact}
  \begin{enumerate}
  \item که \LR{$\Gamma \Proves !A$} وي، نو داسې متناهي فرعي سټ
    \LR{$\Gamma_0 \subseteq \Gamma$} شته چې \LR{$\Gamma_0 \Proves !A$} وي۔
  \item که د~\LR{$\Gamma$} هر متناهي فرعي سټ سازګار وي، نو \LR{$\Gamma$}
    سازګار دے۔
  \end{enumerate}
\end{prop}
\olpseganchor{OLP-0091-B023}{end}

\olpseganchor{OLP-0091-B024}{start}
\begin{proof}
  \begin{enumerate}
    \item که \LR{$\Gamma \Proves !A$} وي، نو له~\LR{$\Gamma$} څخه د~\LR{$!A$}
      کوم اشتقاق~\LR{$\delta$} شته۔ \LR{$\Gamma_0$} د~\LR{$\delta$} د
      نۀ ايستل شوې فرضيو سټ ونيسئ۔ ځکه هر اشتقاق متناهي
      دے، \LR{$\Gamma_0$} يوازې متناهي ډېرې تړلي فارمولونه لرلے شي۔ نو
      \LR{$\delta$} له متناهي~\LR{$\Gamma_0 \subseteq \Gamma$} څخه د~\LR{$!A$}
      اشتقاق دے۔
    \item دا د~(۱) د هغه ځانګړي حالت نقيض عکس دے چې
      \LR{$!A \ident \lfalse$} وي۔
  \end{enumerate}
\end{proof}
\olpseganchor{OLP-0091-B024}{end}

\olpseganchor{OLP-0092-B004}{start}
\olfileid{fol}{ntd}{prv}
\olpseganchor{OLP-0092-B004}{end}
      
\olpseganchor{OLP-0092-B005}{start}
\olsection{د اشتقاق وړتيا او سازګاري}
\olpseganchor{OLP-0092-B005}{end}

\olpseganchor{OLP-0092-B006}{start}
اوس به د د اشتقاق وړتيا اړيکې يو شمېر خاصيتونه ثابت کړو۔ دوئ په
خپلواکه توګه په زړه پورې دي، خو هر يو به د بشپړتيا قضيې په ثبوت
کښې ونډه ولري۔
\olpseganchor{OLP-0092-B006}{end}

\olpseganchor{OLP-0092-B007}{start}
\begin{prop}\ollabel{prop:provability-contr}
  که \LR{$\Gamma \Proves !A$} وي او \LR{$\Gamma \cup \{!A\}$} ناسازګار وي،
  نو \LR{$\Gamma$} ناسازګار دے۔
\end{prop}
\olpseganchor{OLP-0092-B007}{end}

\olpseganchor{OLP-0092-B008}{start}
\begin{proof}
له~\LR{$\Gamma$} څخه د~\LR{$!A$} اشتقاق~\LR{$\delta_1$} او له
\LR{$\Gamma \cup \{!A\}$} څخه د~\LR{$\lfalse$} اشتقاق~\LR{$\delta_2$}
ونوموئ۔ بيا داسې اشتقاق کولے شو:
\begin{prooftree}
\AxiomC{\LR{$\Gamma, \Discharge{!A}{1}$}}
\RightLabel{\LR{$\delta_2$}}
\DeduceC{\LR{$\lfalse$}}
\DischargeRule{\Intro{\lnot}}{1}
\UnaryInfC{\LR{$\lnot !A$}}
\AxiomC{\LR{$\Gamma$}}
\RightLabel{\LR{$\delta_1$}}
\DeduceC{\LR{$!A$}}
\RightLabel{\Elim{\lnot}}
\BinaryInfC{\LR{$\lfalse$}}
\end{prooftree}
په نوي اشتقاق کښې فرضيه~\LR{$!A$}~ايستل شوې ده، نو دا
له~\LR{$\Gamma$} څخه اشتقاق دے۔
\end{proof}
\olpseganchor{OLP-0092-B008}{end}

\olpseganchor{OLP-0092-B009}{start}
\begin{prop}
\ollabel{prop:prov-incons}
\LR{$\Gamma \Proves !A$} هله او يوازې هله چې
\LR{$\Gamma \cup \{\lnot !A\}$} ناسازګار وي۔
\end{prop}
\olpseganchor{OLP-0092-B009}{end}

\olpseganchor{OLP-0092-B010}{start}
\begin{proof}
لومړے فرض کړئ \LR{$\Gamma \Proves !A$}؛ يعنې د~\LR{$!A$} داسې
اشتقاق~\LR{$\delta_0$} شته چې نۀ ايستل شوې فرضيې ئې
\LR{$\Gamma$} دي۔ له~\LR{$\Gamma \cup \{\lnot !A\}$} څخه د~\LR{$\lfalse$} يو
اشتقاق داسې اخلو:
\begin{prooftree}
  \AxiomC{\LR{$\lnot !A$}}
  \AxiomC{\LR{$\Gamma$}}
  \RightLabel{\LR{$\delta_0$}}
  \DeduceC{\LR{$!A$}}
  \RightLabel{\Elim{\lnot}}
  \BinaryInfC{\LR{$\lfalse$}}
\end{prooftree}
\olpseganchor{OLP-0092-B010}{end}

\olpseganchor{OLP-0092-B011}{start}
اوس فرض کړئ \LR{$\Gamma \cup \{\lnot !A\}$} ناسازګار دے، او
\LR{$\delta_1$} له~\LR{$\Gamma \cup \{\lnot !A\}$} کښې له نۀ ايستل شوې
فرضيو څخه د~\LR{$\lfalse$} اړوند اشتقاق دے۔ يوازې له~\LR{$\Gamma$}
څخه د~\LR{$!A$} يو اشتقاق د~\LR{$\FalseCl$} په کارولو اخلو:
\begin{prooftree}
  \AxiomC{\LR{$\Gamma, \Discharge{\lnot !A}{1}$}}
  \RightLabel{\LR{$\delta_1$}}
  \DeduceC{\LR{$\lfalse$}}
  \DischargeRule{\FalseCl}{1}
  \UnaryInfC{\LR{$!A$}}
\end{prooftree}
\textbf{د ژباړې د نسخې سمون (OLND-005):} سرچينه په دې ونه کښې
له ايستلو نښان سره د همدې قاعدې يو دويم ښي نښان هم ږدي؛ زيات نښان
لرې شوے او د فرضيې ايستل او د قاعدې نوم په ټاکلي ګډ نښان کښې ساتل
شوي دي۔
\end{proof}
\olpseganchor{OLP-0092-B011}{end}

\olpseganchor{OLP-0092-B012}{start}
\begin{prob}
ثابت کړئ چې \LR{$\Gamma \Proves \lnot !A$} هله او يوازې هله چې
\LR{$\Gamma \cup \{!A\}$} ناسازګار وي۔
\end{prob}
\olpseganchor{OLP-0092-B012}{end}

\olpseganchor{OLP-0092-B013}{start}
\begin{prop}\ollabel{prop:explicit-inc}
  که \LR{$\Gamma \Proves !A$} او \LR{$\lnot !A \in \Gamma$} وي، نو \LR{$\Gamma$}
  ناسازګار دے۔
\end{prop}
\olpseganchor{OLP-0092-B013}{end}

\olpseganchor{OLP-0092-B014}{start}
\begin{proof}
  فرض کړئ \LR{$\Gamma \Proves !A$} او \LR{$\lnot !A \in \Gamma$}۔ بيا
  له~\LR{$\Gamma$} څخه د~\LR{$!A$} کوم اشتقاق~\LR{$\delta$} شته۔ د
  \LR{$\Elim{\lnot}$} قاعدې دا ساده کارونه وګورئ:
  \begin{prooftree}
    \AxiomC{\LR{$\lnot !A$}}
    \AxiomC{\LR{$\Gamma$}}
    \RightLabel{\LR{$\delta$}}
    \DeduceC{\LR{$!A$}}
    \RightLabel{\Elim{\lnot}}
    \BinaryInfC{\LR{$\lfalse$}}
  \end{prooftree}
  ځکه \LR{$\lnot !A \in \Gamma$} او ټولې نۀ ايستل شوې فرضيې
  په~\LR{$\Gamma$} کښې دي، دا \LR{$\Gamma \Proves \lfalse$} ښيي۔
\end{proof}
\olpseganchor{OLP-0092-B014}{end}

\olpseganchor{OLP-0092-B015}{start}
\begin{prop}\ollabel{prop:provability-exhaustive}
  که \LR{$\Gamma \cup \{!A\}$} او \LR{$\Gamma \cup \{\lnot !A\}$} دواړه
  ناسازګار وي، نو \LR{$\Gamma$} ناسازګار دے۔
\end{prop}
\olpseganchor{OLP-0092-B015}{end}

\olpseganchor{OLP-0092-B016}{start}
\begin{proof}
په ترتيب سره له~\LR{$\Gamma \cup \{ !A \}$} څخه د~\LR{$\lfalse$} او له
\LR{$\Gamma \cup \{ \lnot !A \}$} څخه د~\LR{$\lfalse$}~اشتقاقونه
\LR{$\delta_1$} او~\LR{$\delta_2$} شته۔ بيا داسې اشتقاق کولے شو:
\begin{prooftree}
\AxiomC{\LR{$\Gamma, \Discharge{\lnot !A}{2}$}}
\RightLabel{\LR{$\delta_2$}}
\DeduceC{\LR{$\lfalse$}}
\DischargeRule{\Intro{\lnot}}{2}
\UnaryInfC{\LR{$\lnot \lnot !A$}}
\AxiomC{\LR{$\Gamma, \Discharge{!A}{1}$}}
\RightLabel{\LR{$\delta_1$}}
\DeduceC{\LR{$\lfalse$}}
\DischargeRule{\Intro{\lnot}}{1}
\UnaryInfC{\LR{$\lnot !A$}}
\RightLabel{\Elim{\lnot}}
\BinaryInfC{\LR{$\lfalse$}}
\end{prooftree}
ځکه فرضيې~\LR{$!A$} او~\LR{$\lnot !A$}~ايستل شوې دي، دا يوازې
له~\LR{$\Gamma$} څخه د~\LR{$\lfalse$} يو اشتقاق دے۔ نو \LR{$\Gamma$}
ناسازګار دے۔
\end{proof}
\olpseganchor{OLP-0092-B016}{end}

\olpseganchor{OLP-0093-B005}{start}
\olfileid{fol}{ntd}{ppr}
\olpseganchor{OLP-0093-B005}{end}

\olpseganchor{OLP-0093-B006}{start}
\olsection{د اشتقاق وړتيا او قضيې رابطونه}
\olpseganchor{OLP-0093-B006}{end}

\olpseganchor{OLP-0093-B007}{start}
\begin{explain}
  موږ ثابتوو چې د فطري استنتاج د د اشتقاق وړتيا اړيکه~\LR{$\Proves$}
  دومره پياوړې ده چې د قضيې رابطونو په اړه ځينې بنسټيز حقيقتونه،
  لکه \LR{$!A \land !B \Proves !A$} او
  \LR{$!A, !A \lif !B \Proves !B$} (قياسي استنتاج)، ثابت کړي۔ دا
  حقيقتونه د بشپړتيا قضيې د ثبوت دپاره پکار دي۔
\end{explain}
\olpseganchor{OLP-0093-B007}{end}

\olpseganchor{OLP-0093-B008}{start}
\begin{prop}\ollabel{prop:provability-land}
  \begin{enumerate}
  \item \ollabel{prop:provability-land-left} دواړه
    \LR{$!A \land !B \Proves !A$} او \LR{$!A \land !B \Proves !B$}۔
  \item \ollabel{prop:provability-land-right}
    \LR{$!A, !B \Proves !A \land !B$}۔
  \end{enumerate}
\end{prop}
\olpseganchor{OLP-0093-B008}{end}

\olpseganchor{OLP-0093-B009}{start}
\begin{proof}
  \begin{enumerate}
  \item دواړه داسې اشتقاق کولے شو:
    \begin{prooftree}
      \AxiomC{\LR{$!A \land !B$}}
      \RightLabel{\Elim{\land}}
      \UnaryInfC{\LR{$!A$}}
      \DisplayProof\qquad\bottomAlignProof
      \AxiomC{\LR{$!A \land !B$}}
      \RightLabel{\Elim{\land}}
      \UnaryInfC{\LR{$!B$}}
    \end{prooftree}
  \item داسې ئې اشتقاق کولے شو:
    \begin{prooftree}
      \AxiomC{\LR{$!A$}}
      \AxiomC{\LR{$!B$}}
      \RightLabel{\Intro{\land}}
      \BinaryInfC{\LR{$!A \land !B$}}
    \end{prooftree}
  \end{enumerate}
\end{proof}
\olpseganchor{OLP-0093-B009}{end}


\olpseganchor{OLP-0093-B010}{start}
\begin{prop}\ollabel{prop:provability-lor}
  \begin{enumerate}
  \item \LR{$!A \lor !B, \lnot !A, \lnot !B$} ناسازګار دے۔
  \item دواړه \LR{$!A \Proves !A \lor !B$} او \LR{$!B \Proves !A \lor !B$}۔
  \end{enumerate}
\end{prop}
\olpseganchor{OLP-0093-B010}{end}

\olpseganchor{OLP-0093-B011}{start}
\begin{proof}
  \begin{enumerate}
  \item لاندې اشتقاق وګورئ:
    \begin{prooftree}
      \AxiomC{\LR{$!A \lor !B$}}
      \AxiomC{\LR{$\lnot !A$}}
      \AxiomC{\LR{$\Discharge{!A}{1}$}}
      \RightLabel{\Elim{\lnot}}
      \BinaryInfC{\LR{$\lfalse$}}
      \AxiomC{\LR{$\lnot !B$}}
      \AxiomC{\LR{$\Discharge{!B}{1}$}}
      \RightLabel{\Elim{\lnot}}
      \BinaryInfC{\LR{$\lfalse$}}
      \DischargeRule{\Elim{\lor}}{1}
      \TrinaryInfC{\LR{$\lfalse$}}
    \end{prooftree}
    دا له نۀ ايستل شوې فرضيو~\LR{$!A \lor !B$}، \LR{$\lnot !A$} او
    \LR{$\lnot !B$} څخه د~\LR{$\lfalse$} يو اشتقاق دے۔
  \item دواړه داسې اشتقاق کولے شو:
    \begin{prooftree}
      \AxiomC{\LR{$!A$}}
      \RightLabel{\Intro{\lor}}
      \UnaryInfC{\LR{$!A \lor !B$}}
      \DisplayProof\qquad\bottomAlignProof
      \AxiomC{\LR{$!B$}}
      \RightLabel{\Intro{\lor}}
      \UnaryInfC{\LR{$!A \lor !B$}}
    \end{prooftree}
  \end{enumerate}
\end{proof}
\olpseganchor{OLP-0093-B011}{end}

\olpseganchor{OLP-0093-B012}{start}
\begin{prop}\ollabel{prop:provability-lif}
  \begin{enumerate}
  \item \ollabel{prop:provability-lif-left}
    \LR{$!A, !A \lif !B \Proves !B$}۔
  \item \ollabel{prop:provability-lif-right}
    دواړه \LR{$\lnot !A \Proves !A \lif !B$} او
    \LR{$!B \Proves !A \lif !B$}۔
  \end{enumerate}
\end{prop}
\olpseganchor{OLP-0093-B012}{end}

\olpseganchor{OLP-0093-B013}{start}
\begin{proof}
  \begin{enumerate}
  \item داسې ئې اشتقاق کولے شو:
    \begin{prooftree}
      \AxiomC{\LR{$!A \lif !B$}}
      \AxiomC{\LR{$!A$}}
      \RightLabel{\Elim{\lif}}
      \BinaryInfC{\LR{$!B$}}
    \end{prooftree}
\olpseganchor{OLP-0093-B013}{end}
    
\olpseganchor{OLP-0093-B014}{start}
  \item دا د لاندې دوو اشتقاقونه په وسيله ښودل کېږي:
    \begin{prooftree}
      \AxiomC{\LR{$\lnot !A$}}
      \AxiomC{\LR{$\Discharge{!A}{1}$}}
      \RightLabel{\Elim{\lnot}}
      \BinaryInfC{\LR{$\lfalse$}}
      \RightLabel{\FalseInt}
      \UnaryInfC{\LR{$!B$}}
      \DischargeRule{\Intro{\lif}}{1}
      \UnaryInfC{\LR{$!A \lif !B$}}
      \DisplayProof\qquad\bottomAlignProof
      \AxiomC{\LR{$!B$}}
      \RightLabel{\Intro{\lif}}
      \UnaryInfC{\LR{$!A \lif !B$}}
    \end{prooftree}
    پام وکړئ چې \LR{$\Intro{\lif}$} د فرضيې~\LR{$!A$}~ايستل کولے
    شي، خو لازمه نۀ ده چې و ئې باسي۔
  \end{enumerate}
\end{proof}
\olpseganchor{OLP-0093-B014}{end}

\olpseganchor{OLP-0094-B005}{start}
\olfileid{fol}{ntd}{qpr}
\olpseganchor{OLP-0094-B005}{end}

\olpseganchor{OLP-0094-B006}{start}
\olsection{د اشتقاق وړتيا او سورونه}
\olpseganchor{OLP-0094-B006}{end}

\olpseganchor{OLP-0094-B007}{start}
\begin{explain}
  د بشپړتيا قضيه دا هم غواړي چې د فطري استنتاج قاعدې د~\LR{$\Proves$}
  په اړه هغه حقيقتونه ورکړي چې په دې برخه کښې ثابتېږي۔
\end{explain}
\olpseganchor{OLP-0094-B007}{end}

\olpseganchor{OLP-0094-B008}{start}
\begin{thm}
\ollabel{thm:strong-generalization} که \LR{$c$} داسې ثابت وي چې په
\LR{$\Gamma$} يا~\LR{$!A(x)$} کښې نۀ راځي، او \LR{$\Gamma \Proves !A(c)$} وي، نو
\LR{$\Gamma \Proves \lforall[x][!A(x)]$}۔
\end{thm}
\olpseganchor{OLP-0094-B008}{end}

\olpseganchor{OLP-0094-B009}{start}
\begin{proof}
\LR{$\delta$} له~\LR{$\Gamma$} څخه د~\LR{$!A(c)$} يو اشتقاق ونيسئ۔ د
\Intro{\lforall} استنتاج په زياتولو د~\LR{$\lforall[x][!A(x)]$} يو
اشتقاق اخلو۔ ځکه \LR{$c$} په~\LR{$\Gamma$} يا~\LR{$!A(x)$} کښې نۀ راځي،
د ځانګړي متغير شرط پوره کېږي۔
\end{proof}
\olpseganchor{OLP-0094-B009}{end}

\olpseganchor{OLP-0094-B010}{start}
\begin{prop}
\ollabel{prop:provability-quantifiers}
\begin{enumerate}
\item \LR{$!A(t) \Proves \lexists[x][!A(x)]$}.
\olpseganchor{OLP-0094-B010}{end}

\olpseganchor{OLP-0094-B011}{start}
\item \LR{$\lforall[x][!A(x)] \Proves !A(t)$}.
\end{enumerate}
\end{prop}
\olpseganchor{OLP-0094-B011}{end}

\olpseganchor{OLP-0094-B012}{start}
\begin{proof}
\begin{enumerate}
\item لاندې له~\LR{$!A(t)$} څخه د~\LR{$\lexists[x][!A(x)]$} يو
اشتقاق دے:
\begin{prooftree}
\AxiomC{\LR{$!A(t)$}}
\RightLabel{\Intro{\lexists}}
\UnaryInfC{\LR{$\lexists[x][!A(x)]$}}
\end{prooftree}
\olpseganchor{OLP-0094-B012}{end}

\olpseganchor{OLP-0094-B013}{start}
\item لاندې له~\LR{$\lforall[x][!A(x)]$} څخه د~\LR{$!A(t)$} يو
اشتقاق دے:
\begin{prooftree}
\AxiomC{\LR{$\lforall[x][!A(x)]$}}
\RightLabel{\Elim{\lforall}}
\UnaryInfC{\LR{$!A(t)$}}
\end{prooftree}
\end{enumerate}
\end{proof}
\olpseganchor{OLP-0094-B013}{end}

\olpseganchor{OLP-0095-B004}{start}
\olfileid{fol}{ntd}{sou}
\olpseganchor{OLP-0095-B004}{end}
      
\olpseganchor{OLP-0095-B005}{start}
\olsection{سموالے}
\olpseganchor{OLP-0095-B005}{end}

\olpseganchor{OLP-0095-B006}{start}
\begin{explain}
د فطري استنتاج په شان اشتقاق نظام هله \emph{سم} دے چې
هغه څيزونه اشتقاق نۀ کړي چې په حقيقت کښې ترې نۀ لازمېږي۔
نو سموالے د اشتقاق نظامونو د تضمين شوې خونديتوب يو ډول
خاصيت دے۔ د دې له مخې چې کوم ثبوت-تيوريکي خاصيت تر بحث لاندې وي،
مثلاً غواړو پوه شو چې:
\begin{enumerate}
\item هره د اشتقاق وړ~تړلے فارمول معتبره ده؛
\item که تړلے فارمول له ځينو نورو څخه د اشتقاق وړ وي، د هغو
  معنايي پايله هم ده؛
\item که د تړلي فارمولونه يو سټ ناسازګار وي، د صدق وړ نۀ دے۔
\end{enumerate}
دا د اشتقاق نظام مهم خاصيتونه دي۔ که له دغو څخه کوم يو
سم نۀ وي، اشتقاق نظام نيمګړے دے---له حده زيات څيزونه به
اشتقاق کوي۔ له دې امله د اشتقاق نظام سموالے ثابتول
تر ټولو زيات اهميت لري۔
\end{explain}
\olpseganchor{OLP-0095-B006}{end}

\olpseganchor{OLP-0095-B007}{start}
\begin{thm}[سموالے]
\ollabel{thm:soundness}
که \LR{$!A$} له نۀ ايستل شوې فرضيو~\LR{$\Gamma$} څخه د اشتقاق وړ وي،
نو \LR{$\Gamma \Entails !A$}۔
\end{thm}
\olpseganchor{OLP-0095-B007}{end}

\olpseganchor{OLP-0095-B008}{start}
\begin{proof}
فرض کړئ \LR{$\delta$} د~\LR{$!A$} يو اشتقاق دے۔ موږ په~\LR{$\delta$}
کښې د استنتاجونو پر شمېر استقرا کوو۔
\olpseganchor{OLP-0095-B008}{end}

\olpseganchor{OLP-0095-B009}{start}
د استقرا د بنسټ دپاره هغه حالت ثابتوو چې د استنتاجونو شمېر~\LR{$0$}
وي۔ په دې حالت کښې \LR{$\delta$} يوازې له يوې تړلے فارمول~\LR{$!A$} څخه
جوړ دے، يعنې يوه فرضيه ده۔ دغه فرضيه نۀ ايستل شوې ده؛ ځکه
فرضیې يوازې د استنتاجونو په وسيله ايستل کېدے شي، او دلته
هېڅ استنتاج نشته۔ نو هر
جوړښت~\LR{$\Struct{M}$}
چې د ثبوت ټولې نۀ ايستل شوې فرضيې پوره کوي، \LR{$!A$} هم پوره کوي۔
\olpseganchor{OLP-0095-B009}{end}

\olpseganchor{OLP-0095-B010}{start}
اوس د استقرا پړاو اخلو۔ فرض کړئ \LR{$\delta$}، \LR{$n$} استنتاجونه لري۔
د تر ټولو لاندې استنتاج مقدمه يا مقدمې د فرعي-اشتقاقونه په
وسيله اشتقاق شوې دي، او په هر يوه کښې تر~\LR{$n$} کم استنتاجونه
دي۔ د استقرا فرضيه منو: د تر ټولو لاندې استنتاج مقدمې د هغو
فرعي-اشتقاقونه له نۀ ايستل شوې فرضيو څخه لازمېږي چې په
دغو مقدمو پاے ته رسېږي۔ بايد وښيو چې پايله~\LR{$!A$} د ټول ثبوت له
نۀ ايستل شوې فرضيو څخه لازمېږي۔
\olpseganchor{OLP-0095-B010}{end}

\olpseganchor{OLP-0095-B011}{start}
د تر ټولو لاندې استنتاج د ډول له مخې حالتونه بېلوو۔ لومړے هغه
ممکن استنتاجونه څېړو چې يوازې يوه مقدمه لري۔
\olpseganchor{OLP-0095-B011}{end}

\olpseganchor{OLP-0095-B012}{start}
\begin{enumerate}
\item فرض کړئ وروستنے استنتاج \Intro{\lnot} دے: اشتقاق
  دا بڼه لري:
  \begin{prooftree}
    \AxiomC{\LR{$\Gamma, \Discharge{!A}{n}$}}
    \RightLabel{\LR{$\delta_1$}}
    \DeduceC{\LR{$\lfalse$}}
    \DischargeRule{\Intro{\lnot}}{n}
    \UnaryInfC{\LR{$\lnot !A$}}
  \end{prooftree}
  د استقرا د فرضيې له مخې \LR{$\lfalse$} د~\LR{$\delta_1$} له
  نۀ ايستل شوې فرضيو~\LR{$\Gamma \cup \{!A\}$} څخه لازمېږي۔ يو
  جوړښت~\LR{$\Struct{M}$}
  ونيسئ۔ بايد وښيو چې که
  \LR{$\Sat{M}{\Gamma}$} وي، نو
  \LR{$\Sat{M}{\lnot !A}$} دے۔ د خلف په توګه
  فرض کړئ \LR{$\Sat{M}{\Gamma}$}، خو
  \LR{$\Sat/{M}{\lnot !A}$}، يعنې
  \LR{$\Sat{M}{!A}$}۔ له دې به
  \LR{$\Sat{M}{\Gamma \cup \{!A\}}$} لازم شي۔
  دا زموږ د استقرا د فرضيې خلاف دے۔ نو
  \LR{$\Sat{M}{\lnot !A}$}۔
\olpseganchor{OLP-0095-B012}{end}

\olpseganchor{OLP-0095-B013}{start}
\item وروستنے استنتاج \Elim{\land} دے: دوه بڼې لري؛ له مقدمې
  \LR{$!A \land !B$} څخه يا~\LR{$!A$} يا~\LR{$!B$} راايستل کېدے شي۔ لومړے حالت
  ونيسئ۔ اشتقاق~\LR{$\delta$} داسې ښکاري:
  \begin{prooftree}
    \AxiomC{\LR{$\Gamma$}}
    \RightLabel{\LR{$\delta_1$}}
    \DeduceC{\LR{$!A \land !B$}}
    \RightLabel{\Elim{\land}}
    \UnaryInfC{\LR{$!A$}}
  \end{prooftree}
  د استقرا د فرضيې له مخې \LR{$!A \land !B$} د~\LR{$\delta_1$} له
  نۀ ايستل شوې فرضيو~\LR{$\Gamma$} څخه لازمېږي۔
  يو جوړښت~\LR{$\Struct{M}$} ونيسئ۔
  بايد وښيو چې که
  \LR{$\Sat{M}{\Gamma}$} وي، نو
  \LR{$\Sat{M}{!A}$} دے۔ فرض کړئ
  \LR{$\Sat{M}{\Gamma}$}۔ زموږ د استقرا د
  فرضيې (\LR{$\Gamma \Entails !A \land !B$}) له مخې پوهېږو چې
  \LR{$\Sat{M}{!A \land !B}$}۔ د تعريف له مخې
  \LR{$\Sat{M}{!A \land !B}$} هله او يوازې هله
  چې \LR{$\Sat{M}{!A}$} او
  \LR{$\Sat{M}{!B}$} دواړه وي۔ (هغه حالت چې له
  \LR{$!A \land !B$} څخه~\LR{$!B$} راايستل کېږي، په ورته ډول څېړل کېږي۔)
\olpseganchor{OLP-0095-B013}{end}
  
\olpseganchor{OLP-0095-B014}{start}
\item وروستنے استنتاج \Intro{\lor} دے: دوه بڼې لري؛
  \LR{$!A \lor !B$} يا له مقدمې~\LR{$!A$} يا له مقدمې~\LR{$!B$} څخه راايستل
  کېدے شي۔ لومړے حالت ونيسئ۔ اشتقاق دا بڼه لري:
  \begin{prooftree}
    \AxiomC{\LR{$\Gamma$}}
    \RightLabel{\LR{$\delta_1$}}
    \DeduceC{\LR{$!A$}}
    \RightLabel{\Intro{\lor}}
    \UnaryInfC{\LR{$!A \lor !B$}}
  \end{prooftree}
  د استقرا د فرضيې له مخې \LR{$!A$} د~\LR{$\delta_1$} له
  نۀ ايستل شوې فرضيو~\LR{$\Gamma$} څخه لازمېږي۔ يو
  جوړښت~\LR{$\Struct{M}$}
  ونيسئ۔ بايد وښيو چې که
  \LR{$\Sat{M}{\Gamma}$} وي، نو
  \LR{$\Sat{M}{!A \lor !B}$} دے۔ فرض کړئ
  \LR{$\Sat{M}{\Gamma}$}؛ بيا
  \LR{$\Sat{M}{!A}$}، ځکه
  \LR{$\Gamma \Entails !A$} (د استقرا فرضيه)۔ نو
  \LR{$\Sat{M}{!A \lor !B}$} هم بايد وي۔
  (هغه حالت چې \LR{$!A \lor !B$} له~\LR{$!B$} څخه راايستل کېږي، په ورته ډول
  څېړل کېږي۔)
\olpseganchor{OLP-0095-B014}{end}
  
\olpseganchor{OLP-0095-B015}{start}
\item وروستنے استنتاج \Intro{\lif} دے: \LR{$!A \lif !B$} له داسې فرعي
  ثبوت څخه راايستل کېږي چې فرضيه ئې~\LR{$!A$} او پايله ئې~\LR{$!B$} ده،
  يعنې:
  \begin{prooftree}
    \AxiomC{\LR{$\Gamma, \Discharge{!A}{n}$}}
    \RightLabel{\LR{$\delta_1$}}
    \DeduceC{\LR{$!B$}}
    \DischargeRule{\Intro{\lif}}{n}
    \UnaryInfC{\LR{$!A \lif !B$}}
  \end{prooftree}
  د استقرا د فرضيې له مخې \LR{$!B$} د~\LR{$\delta_1$} له
  نۀ ايستل شوې فرضيو څخه لازمېږي، يعنې
  \LR{$\Gamma \cup \{!A\} \Entails !B$}۔ يو
  جوړښت~\LR{$\Struct{M}$}
  ونيسئ۔ د~\LR{$\delta$}~نۀ ايستل شوې فرضيې يوازې~\LR{$\Gamma$} دي؛ ځکه
  \LR{$!A$} په وروستي استنتاج کښې ايستل شوې ده۔ نو بايد وښيو
  چې \LR{$\Gamma \Entails !A \lif !B$}۔ د خلف په توګه فرض کړئ چې د کوم
  جوړښت~\LR{$\Struct{M}$}
  دپاره \LR{$\Sat{M}{\Gamma}$}، خو
  \LR{$\Sat/{M}{!A \lif !B}$} وي۔ نو
  \LR{$\Sat{M}{!A}$} او
  \LR{$\Sat/{M}{!B}$}۔ خو د فرضيې له مخې
  \LR{$!B$} د~\LR{$\Gamma \cup \{!A\}$} معنايي پايله ده، يعنې
  \LR{$\Sat{M}{!B}$}، او دا تناقض دے۔ نو
  \LR{$\Gamma \Entails !A \lif !B$}۔
\olpseganchor{OLP-0095-B015}{end}

\olpseganchor{OLP-0095-B016}{start}
\item وروستنے استنتاج \FalseInt دے: دلته \LR{$\delta$} داسې پاے ته رسېږي:
  \begin{prooftree}
    \AxiomC{\LR{$\Gamma$}}
    \RightLabel{\LR{$\delta_1$}}
    \DeduceC{\LR{$\lfalse$}}
    \RightLabel{\FalseInt}
    \UnaryInfC{\LR{$!A$}}
  \end{prooftree}
  د استقرا د فرضيې له مخې \LR{$\Gamma \Entails \lfalse$}۔ بايد وښيو چې
  \LR{$\Gamma \Entails !A$}۔ فرض کړئ داسې نۀ وي؛ بيا د کوم
  \LR{$\Struct{M}$} دپاره
  \LR{$\Sat{M}{\Gamma}$} او
  \LR{$\Sat/{M}{!A}$} دي۔ خو
  کاذبي هېڅ جوړښت يا ارزونه نۀ پوره کوي؛ يعنې تل
  \LR{$\Sat/{M}{\lfalse}$} وي؛ نو له دې به
  \LR{$\Gamma \Entails/ \lfalse$} لازم شي، چې د استقرا د فرضيې خلاف دے۔
\olpseganchor{OLP-0095-B016}{end}

\olpseganchor{OLP-0095-B017}{start}
\item وروستنے استنتاج \FalseCl دے: تمرين۔
\olpseganchor{OLP-0095-B017}{end}

\olpseganchor{OLP-0095-B018}{start}

\item وروستنے استنتاج \Intro{\lforall} دے: بيا \LR{$\delta$} دا بڼه لري:
  \begin{prooftree}
    \AxiomC{\LR{$\Gamma$}}
    \RightLabel{\LR{$\delta_1$}}
    \DeduceC{\LR{$!A(a)$}}
    \RightLabel{\Intro{\lforall}}
    \UnaryInfC{\LR{$\lforall[x][!A(x)]$}}
  \end{prooftree}
  د استقرا د فرضيې له مخې مقدمه~\LR{$!A(a)$} د~نۀ ايستل شوې
  فرضيو~\LR{$\Gamma$} معنايي پايله ده۔ داسې جوړښت~\LR{$\Struct{M}$} ونيسئ
  چې \LR{$\Sat{M}{\Gamma}$} وي۔ بايد وښيو چې
  \LR{$\Sat{M}{\lforall[x][!A(x)]}$}۔ ځکه
  \LR{$\lforall[x][!A(x)]$} يوه تړلے فارمول ده، د دې معنٰی دا ده چې
  د هرې متغير-ګومارنې~\LR{$s$} دپاره بايد
  \LR{$\Sat{M}{!A(x)}[s]$} وښيو
  (\olref[syn][ass]{prop:sat-quant})۔ ځکه \LR{$\Gamma$} ټول له جملو جوړ
  دے، د هر~\LR{$!B \in \Gamma$} دپاره
  \LR{$\Sat{M}{!B}[s]$} د~\olref[syn][sat]{defn:satisfaction} له مخې
  دے۔ \LR{$\Struct{M'}$} داسې جوړښت ونيسئ چې د~\LR{$\Struct{M}$} په شان وي،
  خو \LR{$\Assign{a}{M'} = s(x)$} وي۔ ځکه \LR{$a$} په~\LR{$\Gamma$} کښې نۀ راځي،
  د~\olref[syn][ext]{cor:extensionality-sent} له مخې
  \LR{$\Sat{M'}{\Gamma}$}۔ ځکه \LR{$\Gamma \Entails !A(a)$}،
  \LR{$\Sat{M'}{!A(a)}$}۔ ځکه \LR{$!A(a)$} يوه تړلے فارمول ده، د
  \olref[syn][ass]{prop:sentence-sat-true} له مخې
  \LR{$\Sat{M'}{!A(a)}[s]$}۔ د
  \olref[syn][ext]{prop:ext-formulas} له مخې
  \LR{$\Sat{M'}{!A(x)}[s]$} هله او يوازې هله چې
  \LR{$\Sat{M'}{!A(a)}$} (ياد ولرئ \LR{$!A(a)$} هماغه
  \LR{$\Subst{!A(x)}{a}{x}$} دے)۔ نو \LR{$\Sat{M'}{!A(x)}[s]$}۔ ځکه \LR{$a$}
  په~\LR{$!A(x)$} کښې نۀ راځي، د
  \olref[syn][ext]{prop:extensionality} له مخې
  \LR{$\Sat{M}{!A(x)}[s]$}۔ خو \LR{$s$} خپل‌سری متغير-ګومارنه وه؛ نو
  \LR{$\Sat{M}{\lforall[x][!A(x)]}$}۔
\olpseganchor{OLP-0095-B018}{end}
  
\olpseganchor{OLP-0095-B019}{start}
\item وروستنے استنتاج \Intro{\lexists} دے: تمرين۔
\olpseganchor{OLP-0095-B019}{end}

\olpseganchor{OLP-0095-B020}{start}
\item وروستنے استنتاج \Elim{\forall} دے: تمرين۔

\end{enumerate}
\olpseganchor{OLP-0095-B020}{end}

\olpseganchor{OLP-0095-B021}{start}
اوس هغه ممکن استنتاجونه څېړو چې څو مقدمې لري:
\Elim{\lor}، \Intro{\land}، \Elim{\lnot}،
\Elim{\lif} او~\Elim{\lexists}۔
\textbf{د ژباړې د نسخې سمون (OLND-006):} سرچينې د څو-مقدمې
قاعدو له دې بشپړ لړ څخه د نفي ايستلو قاعده غورځولې ده، سره له دې
چې لاندې تمرين ئې هم راوړي؛ دلته بېرته شامله شوه۔
\begin{enumerate}
\olpseganchor{OLP-0095-B021}{end}

\olpseganchor{OLP-0095-B022}{start}
\item وروستنے استنتاج \Intro{\land} دے۔ \LR{$!A \land !B$} له
  مقدمو~\LR{$!A$} او~\LR{$!B$} څخه راايستل کېږي او \LR{$\delta$} دا بڼه لري:
  \begin{prooftree}
    \AxiomC{\LR{$\Gamma_1$}}
    \RightLabel{\LR{$\delta_1$}}
    \DeduceC{\LR{$!A$}}
    \AxiomC{\LR{$\Gamma_2$}}
    \RightLabel{\LR{$\delta_2$}}
    \DeduceC{\LR{$!B$}}
    \RightLabel{\Intro{\land}}
    \BinaryInfC{\LR{$!A \land !B$}}
  \end{prooftree}
  د استقرا د فرضيې له مخې \LR{$!A$} د~\LR{$\delta_1$} له
  نۀ ايستل شوې فرضيو~\LR{$\Gamma_1$} څخه، او~\LR{$!B$} د~\LR{$\delta_2$} له
  نۀ ايستل شوې فرضيو~\LR{$\Gamma_2$} څخه لازمېږي۔ د~\LR{$\delta$}
  نۀ ايستل شوې فرضيې \LR{$\Gamma_1 \cup \Gamma_2$} دي؛ نو بايد
  \LR{$\Gamma_1 \cup \Gamma_2 \Entails !A \land !B$} وښيو۔ داسې
  جوړښت~\LR{$\Struct{M}$}
  ونيسئ چې
  \LR{$\Sat{M}{\Gamma_1 \cup \Gamma_2}$} وي۔
  ځکه \LR{$\Sat{M}{\Gamma_1}$}، او
  \LR{$\Gamma_1 \Entails !A$}، نو
  \LR{$\Sat{M}{!A}$}۔ همدارنګه، ځکه
  \LR{$\Sat{M}{\Gamma_2}$} او
  \LR{$\Gamma_2 \Entails !B$}، نو
  \LR{$\Sat{M}{!B}$}۔ دواړه يو ځاے
  \LR{$\Sat{M}{!A \land !B}$} راکوي۔
\olpseganchor{OLP-0095-B022}{end}
  
\olpseganchor{OLP-0095-B023}{start}
\item وروستنے استنتاج \Elim{\lor} دے: تمرين۔
\olpseganchor{OLP-0095-B023}{end}

\olpseganchor{OLP-0095-B024}{start}
\item وروستنے استنتاج \Elim{\lif} دے۔ \LR{$!B$} له مقدمو
  \LR{$!A \lif !B$} او~\LR{$!A$} څخه راايستل کېږي۔ اشتقاق~\LR{$\delta$}
  داسې ښکاري:
  \begin{prooftree}
    \AxiomC{\LR{$\Gamma_1$}}
    \RightLabel{\LR{$\delta_1$}}
    \DeduceC{\LR{$!A \lif !B$}}
    \AxiomC{\LR{$\Gamma_2$}}
    \RightLabel{\LR{$\delta_2$}}
    \DeduceC{\LR{$!A$}}
    \RightLabel{\Elim{\lif}}
    \BinaryInfC{\LR{$!B$}}
  \end{prooftree}
  د استقرا د فرضيې له مخې \LR{$!A \lif !B$} د~\LR{$\delta_1$} له
  نۀ ايستل شوې فرضيو~\LR{$\Gamma_1$} څخه، او~\LR{$!A$} د~\LR{$\delta_2$} له
  نۀ ايستل شوې فرضيو~\LR{$\Gamma_2$} څخه لازمېږي۔ يو
  جوړښت~\LR{$\Struct{M}$}
  ونيسئ۔ بايد وښيو چې که
  \LR{$\Sat{M}{\Gamma_1 \cup \Gamma_2}$} وي، نو
  \LR{$\Sat{M}{!B}$} دے۔ فرض کړئ
  \LR{$\Sat{M}{\Gamma_1 \cup \Gamma_2}$}۔
  ځکه \LR{$\Gamma_1 \Entails !A \lif !B$}، نو
  \LR{$\Sat{M}{!A \lif !B}$}۔ ځکه
  \LR{$\Gamma_2 \Entails !A$}، نو
  \LR{$\Sat{M}{!A}$}۔ له دې څخه
  \LR{$\Sat{M}{!B}$} لازمېږي (ځکه که
  \LR{$\Sat/{M}{!B}$} وے، نو له
  \LR{$\Sat{M}{!A}$} سره به
  \LR{$\Sat/{M}{!A \lif !B}$} وے، چې د
  \LR{$\Sat{M}{!A \lif !B}$} خلاف دے)۔
\olpseganchor{OLP-0095-B024}{end}

\olpseganchor{OLP-0095-B025}{start}
\item وروستنے استنتاج \Elim{\lnot} دے: تمرين۔
\olpseganchor{OLP-0095-B025}{end}
  
\olpseganchor{OLP-0095-B026}{start}
\item وروستنے استنتاج \Elim{\lexists} دے: تمرين۔
\end{enumerate}
\end{proof}
\olpseganchor{OLP-0095-B026}{end}

\olpseganchor{OLP-0095-B027}{start}
\begin{prob}
د~\olref[fol][ntd][sou]{thm:soundness} ثبوت بشپړ کړئ۔
\end{prob}
\olpseganchor{OLP-0095-B027}{end}



\olpseganchor{OLP-0095-B029}{start}
\begin{cor}
\ollabel{cor:weak-soundness}
که \LR{$\Proves !A$} وي، نو \LR{$!A$} معتبره ده۔
\end{cor}
\olpseganchor{OLP-0095-B029}{end}

\olpseganchor{OLP-0095-B030}{start}
\begin{cor}
\ollabel{cor:consistency-soundness}
که \LR{$\Gamma$} د صدق وړ وي، نو سازګار دے۔
\end{cor}
\olpseganchor{OLP-0095-B030}{end}

\olpseganchor{OLP-0095-B031}{start}
\begin{proof}
نقيض عکس ثابتوو۔ فرض کړئ \LR{$\Gamma$} سازګار نۀ دے۔ بيا
\LR{$\Gamma \Proves \lfalse$}، يعنې له~\LR{$\Gamma$} څخه د
نۀ ايستل شوې فرضيو پر بنسټ د~\LR{$\lfalse$} يو اشتقاق شته۔
د~\olref{thm:soundness} له مخې هر
جوړښت~\LR{$\Struct{M}$}
چې \LR{$\Gamma$} پوره کوي، بايد \LR{$\lfalse$} هم پوره کړي۔ ځکه د هر
جوړښت~\LR{$\Struct{M}$}
دپاره \LR{$\Sat/{M}{\lfalse}$}، نو هېڅ
\LR{$\Struct{M}$}، \LR{$\Gamma$} نۀ شي پوره
کولے؛ يعنې \LR{$\Gamma$} د صدق وړ نۀ دے۔
\end{proof}
\olpseganchor{OLP-0095-B031}{end}

\olpseganchor{OLP-0096-B004}{start}
\olfileid{fol}{ntd}{ide}
\olpseganchor{OLP-0096-B004}{end}

\olpseganchor{OLP-0096-B005}{start}
\olsection{له عينيت سره اشتقاقونه}
\olpseganchor{OLP-0096-B005}{end}

\olpseganchor{OLP-0096-B006}{start}
له عينيت سره اشتقاقونه اضافي استنتاجي قاعدو ته اړتيا لري۔
\olpseganchor{OLP-0096-B006}{end}

\olpseganchor{OLP-0096-B007}{start}
\begin{defish}
\AxiomC{}
\RightLabel{\Intro{\eq}}
\UnaryInfC{\LR{$\eq[t][t]$}}
\DisplayProof
\hfill
\begin{tabular}{r}
\AxiomC{\LR{$\eq[t_1][t_2]$}}
\AxiomC{\LR{$!A(t_1)$}}
\RightLabel{\Elim{\eq}}
\BinaryInfC{\LR{$!A(t_2)$}}
\DisplayProof
\\[3ex]
\AxiomC{\LR{$\eq[t_1][t_2]$}}
\AxiomC{\LR{$!A(t_2)$}}
\RightLabel{\Elim{\eq}}
\BinaryInfC{\LR{$!A(t_1)$}}
\DisplayProof
\end{tabular}
\end{defish}
\olpseganchor{OLP-0096-B007}{end}

\olpseganchor{OLP-0096-B008}{start}
په پورته قاعدو کښې \LR{$t$}، \LR{$t_1$} او~\LR{$t_2$} تړلې اصطلاحګانې دي۔
د~\Intro{\eq} قاعده راکوي چې د~\LR{$\eq[t][t]$} په بڼه هره د يووالي
جمله نېغ په نېغه، له هېڅ فرضيې پرته، اشتقاق کړو۔
\olpseganchor{OLP-0096-B008}{end}

\olpseganchor{OLP-0096-B009}{start}
\begin{ex}
که \LR{$s$} او~\LR{$t$} تړلې اصطلاحګانې وي، نو
\LR{$!A(s), \eq[s][t] \Proves !A(t)$}:
\begin{prooftree}
\AxiomC{\LR{$\eq[s][t]$}}
\AxiomC{\LR{$!A(s)$}}
\RightLabel{\LR{$\Elim{\eq}$}}
\BinaryInfC{\LR{$!A(t)$}}
\end{prooftree}
دا به د «يو شان څيزونو د تعويض د اصل» يا د لايبنېڅ د قانون په نامه
اشنا وي۔
\end{ex}
\olpseganchor{OLP-0096-B009}{end}

\olpseganchor{OLP-0096-B010}{start}
\begin{prob}
ثابت کړئ چې \LR{$=$} متناظره او انتقالي دواړه ده؛ يعنې د
\LR{$\lforall[x][\lforall[y][(\eq[x][y] \lif
    \eq[y][x])]]$} او~\LR{$\lforall[x][\lforall[y][\lforall[z]((\eq[x][y]
    \land \eq[y][z]) \lif \eq[x][z])]]$}
اشتقاقونه ورکړئ۔
\end{prob}
\olpseganchor{OLP-0096-B010}{end}

\olpseganchor{OLP-0096-B011}{start}
\begin{ex}
موږ دا تړلے فارمول~اشتقاق کوو:
\begin{align*}
& \lforall[x][\lforall[y][((!A(x) \land !A(y)) \lif \eq[x][y])]]
\intertext{\RL{له دې تړلے فارمول څخه:}}
& \lexists[x][\lforall[y][(!A(y) \lif \eq[y][x])]]
\end{align*}
اشتقاق له لاندې څخه پورته جوړوو:
\begin{prooftree}
\AxiomC{\LR{$\lexists[x][\lforall[y][(!A(y) \lif \eq[y][x])]]
\quad \Discharge{!A(a) \land !A(b)}{1}$}}
\DeduceC{\LR{$\eq[a][b]$}}
\DischargeRule{\Intro{\lif}}{1}
\UnaryInfC{\LR{$((!A(a) \land !A(b)) \lif \eq[a][b])$}}
\RightLabel{\Intro{\lforall}}
\UnaryInfC{\LR{$\lforall[y][((!A(a) \land !A(y)) \lif \eq[a][y])]$}}
\RightLabel{\Intro{\lforall}}
\UnaryInfC{\LR{$\lforall[x][\lforall[y][((!A(x) \land !A(y)) \lif \eq[x][y])]]$}}
\end{prooftree}
اوس بايد اصلي فرضيه وکاروو: ځکه دا وجودي فارمول ده، د
منځنۍ پايلې~\LR{$\eq[a][b]$} د~اشتقاق کولو دپاره
\Elim{\lexists} کاروو۔
\begin{prooftree}
\AxiomC{\LR{$\lexists[x][\lforall[y][(!A(y) \lif \eq[y][x])]]$}}
\AxiomC{\LR{$\Discharge{\lforall[y][(!A(y) \lif \eq[y][c]])}{2}$}}
\noLine
\UnaryInfC{\LR{$\Discharge{!A(a) \land !A(b)}{1}$}}
\DeduceC{\LR{$\eq[a][b]$}}
\DischargeRule{\Elim{\lexists}}{2}
\BinaryInfC{\LR{$\eq[a][b]$}}
\DischargeRule{\Intro{\lif}}{1}
\UnaryInfC{\LR{$((!A(a) \land !A(b)) \lif \eq[a][b])$}}
\RightLabel{\Intro{\lforall}}
\UnaryInfC{\LR{$\lforall[y][((!A(a) \land !A(y)) \lif \eq[a][y])]$}}
\RightLabel{\Intro{\lforall}}
\UnaryInfC{\LR{$\lforall[x][\lforall[y][((!A(x) \land !A(y)) \lif \eq[x][y])]]$}}
\end{prooftree}
د پاس په ښي لوري کښې فرعي-اشتقاق د خپلو فرضيو په کارولو
بشپړېږي، څو \LR{$\eq[a][c]$} او~\LR{$\eq[b][c]$} وښيي۔ دا دوه جلا
اشتقاقونه غواړي۔ د~\LR{$\eq[a][c]$}~اشتقاق داسې دے:
\begin{prooftree}
\AxiomC{\LR{$\Discharge{\lforall[y][(!A(y) \lif \eq[y][c]])}{2}$}}
\RightLabel{\Elim{\lforall}}
\UnaryInfC{\LR{$!A(a) \lif \eq[a][c]$}}
\AxiomC{\LR{$\Discharge{!A(a) \land !A(b)}{1}$}}
\RightLabel{\Elim{\land}}
\UnaryInfC{\LR{$!A(a)$}}
\RightLabel{\Elim{\lif}}
\BinaryInfC{\LR{$\eq[a][c]$}}
\end{prooftree}
له~\LR{$\eq[a][c]$} او~\LR{$\eq[b][c]$} څخه
\LR{$\eq[a][b]$} د~\Elim{\eq} په وسيله اشتقاق کوو۔
\end{ex}
\olpseganchor{OLP-0096-B011}{end}

\olpseganchor{OLP-0096-B012}{start}
\begin{prob}
د لاندې فارمولونه~اشتقاقونه ورکړئ:
\begin{enumerate}
\item \LR{$\lforall[x][\lforall[y][((\eq[x][y] \land !A(x)) \lif !A(y))]]$}
\item \LR{$\lexists[x][!A(x)] \land \lforall[y][\lforall[z][((!A(y) \land
    !A(z)) \lif \eq[y][z])]] \lif \lexists[x][(!A(x) \land
  \lforall[y][(!A(y) \lif \eq[y][x])])]$}
\end{enumerate}
\end{prob}
\olpseganchor{OLP-0096-B012}{end}

\olpseganchor{OLP-0097-B004}{start}
\olfileid{fol}{ntd}{sid}
\olpseganchor{OLP-0097-B004}{end}

\olpseganchor{OLP-0097-B005}{start}
\olsection{له عينيت سره سموالے}
\olpseganchor{OLP-0097-B005}{end}

\olpseganchor{OLP-0097-B006}{start}
\begin{prop}
د~\LR{$\eq$} له قاعدو سره فطري استنتاج سم دے۔
\end{prop}
\olpseganchor{OLP-0097-B006}{end}

\olpseganchor{OLP-0097-B007}{start}
\begin{proof}
د~\LR{$\eq[t][t]$} په بڼه هره فارمول معتبره ده؛ ځکه د هر
جوړښت~\LR{$\Struct M$} دپاره \LR{$\Sat{M}{\eq[t][t]}$}۔ (پام وکړئ
چې اصطلاح~\LR{$t$} تړلې فرضوو، يعنې هېڅ متغير نۀ لري؛ نو
متغير-ګومارنې اغېز نۀ لري۔)
\olpseganchor{OLP-0097-B007}{end}

\olpseganchor{OLP-0097-B008}{start}
فرض کړئ په کوم اشتقاق کښې وروستنے استنتاج \Elim{\eq} دے،
يعنې اشتقاق لاندې بڼه لري:
\begin{prooftree}
  \AxiomC{\LR{$\Gamma_1$}}
  \RightLabel{\LR{$\delta_1$}}
  \DeduceC{\LR{$\eq[t_1][t_2]$}}
  \AxiomC{\LR{$\Gamma_2$}}
  \RightLabel{\LR{$\delta_2$}}
  \DeduceC{\LR{$!A(t_1)$}}
  \RightLabel{\Elim{\eq}}
  \BinaryInfC{\LR{$!A(t_2)$}}
\end{prooftree}
مقدمې~\LR{$\eq[t_1][t_2]$} او~\LR{$!A(t_1)$} په خپل وار له
نۀ ايستل شوې فرضيو~\LR{$\Gamma_1$} او~\LR{$\Gamma_2$} څخه
اشتقاق شوې دي۔ غواړو وښيو چې \LR{$!A(t_2)$} له
\LR{$\Gamma_1 \cup \Gamma_2$} څخه لازمېږي۔ داسې
جوړښت~\LR{$\Struct{M}$} ونيسئ چې
\LR{$\Sat{M}{\Gamma_1 \cup \Gamma_2}$} وي۔ د استقرا د فرضيې له مخې
\LR{$\Sat{M}{!A(t_1)}$} او~\LR{$\Sat{M}{\eq[t_1][t_2]}$}۔ نو
\LR{$\Value{t_1}{M} = \Value{t_2}{M}$}۔ \LR{$s$} هره متغير-ګومارنه ونيسئ،
او \LR{$m = \Value{t_1}{M} = \Value{t_2}{M}$}۔ د
\olref[fol][syn][ext]{prop:ext-formulas} له مخې
\LR{$\Sat{M}{!A(t_1)}[s]$} هله او يوازې هله چې
\LR{$\Sat{M}{!A(x)}[\Subst{s}{m}{x}]$}، او دا هله او يوازې هله چې
\LR{$\Sat{M}{!A(t_2)}[s]$}۔ ځکه \LR{$\Sat{M}{!A(t_1)}$}، نو
\LR{$\Sat{M}{!A(t_2)}$}۔
\end{proof}
\olpseganchor{OLP-0097-B008}{end}

\olpseganchor{OLP-0098-B004}{start}
\olchapter{fol}{tab}{تابلوګانې}
\olpseganchor{OLP-0098-B004}{end}

\olpseganchor{OLP-0098-B005}{start}
\begin{editorial}
دا څپرکے يو نښه لرونکے تحليلي تابلو نظام وړاندې کوي۔
\olpseganchor{OLP-0098-B005}{end}
  
\olpseganchor{OLP-0098-B006}{start}
د ثبوت د يوه نظام په توګه له تابلو سره اړوند مواد د شاملولو يا
ايستلو دپاره د~\texttt{prfTab} ټګ وکاروئ۔
\textbf{د ژباړې د نسخې سمون (OLTAB-001):} سرچينه دلته په تېروتنه
فطري استنتاج يادوي؛ د څپرکي، ټګ او ټول وارد شوي موادو له مخې مطلب
تابلو نظام دے، نو نوم ئې دلته سم شوے دے۔
\end{editorial}
\olpseganchor{OLP-0098-B006}{end}





\olpseganchor{OLP-0098-B009}{start}
%
  

\olpseganchor{OLP-0098-B009}{end}





\olpseganchor{OLP-0098-B012}{start}
%
  

\olpseganchor{OLP-0098-B012}{end}







\olpseganchor{OLP-0098-B016}{start}
%
  

\olpseganchor{OLP-0098-B016}{end}



\olpseganchor{OLP-0098-B018}{start}
%
  
  

\olpseganchor{OLP-0098-B018}{end}

\olpseganchor{OLP-0099-B004}{start}
\olfileid{fol}{tab}{rul}
\olpseganchor{OLP-0099-B004}{end}

\olpseganchor{OLP-0099-B005}{start}
\olsection{قاعدې او تابلوګانې}
\olpseganchor{OLP-0099-B005}{end}

\olpseganchor{OLP-0099-B006}{start}
تابلو په جوړښت کښې د~تړلے فارمول د رښتيا يا
دروغ کېدو د ممکنو لارو منظمه پلټنه ده۔ د تابلو بنسټيز توکي
نښه لرونکي فارمولونه دي: تړلي فارمولونه له يوه رښتياوالي ارزښت
«نښې» سره، چې يا~\LR{$\True$} وي يا~\LR{$\False$}۔ دا نښه لرونکي فارمولونه
په داسې ونه کښې ترتيبېږي چې لاندې خوا ته غځېږي۔
\olpseganchor{OLP-0099-B006}{end}

\olpseganchor{OLP-0099-B007}{start}
\begin{defn}
  \emph{نښه لرونکے فارمول} له يوه رښتياوالي ارزښت او~تړلے فارمول
  څخه جوړه ده، يعنې له دې دواړو څخه يوه:
  \[
  \sFmla{\True}{!A} \text{\RL{ يا }} \sFmla{\False}{!A}.
  \]
\end{defn}
\olpseganchor{OLP-0099-B007}{end}

\olpseganchor{OLP-0099-B008}{start}
په شهودي ډول \LR{$\sFmla{\True}{!A}$} داسې لوستلے شو چې «کېدے شي
\LR{$!A$} رښتيا وي»، او~\LR{$\sFmla{\False}{!A}$} داسې چې «کېدے شي
\LR{$!A$} دروغ وي» (په کوم جوړښت کښې)۔
\olpseganchor{OLP-0099-B008}{end}

\olpseganchor{OLP-0099-B009}{start}
په ونه کښې هر نښه لرونکے فارمول يا \emph{فرض} دے (چې د ونې په
تر ټولو پاسنۍ برخه کښې ليکل کېږي)، يا تر ځان له پاسه له کوم
نښه لرونکے فارمول څخه د يوې استنتاجي قاعدې په وسيله اخستل شوے
دے۔ د مخکښې فارمول د هر ممکن اصلي عملګر دپاره دوه
قاعدې شته: يوه د هغه حالت دپاره چې نښه~\LR{$\True$} وي، او بله د هغه
حالت دپاره چې نښه~\LR{$\False$} وي۔ ځينې قاعدې ونې ته څانګې ورکوي او
ځينې يوازې څانګې ته نښه لرونکي فارمولونه ورزياتوي۔ يوه قاعده
کېدے شي (او ډېر ځله بايد) سمدستي پر مخکښ نښه لرونکے فارمول نۀ،
بلکې له ريښې څخه د قاعدې د پلي کېدو تر ځاے پورې په څانګه کښې پر هر
نښه لرونکے فارمول پلې شي۔
\olpseganchor{OLP-0099-B009}{end}

\olpseganchor{OLP-0099-B010}{start}
يوه څانګه هله \emph{تړلې} ده چې
\LR{$\sFmla{\True}{!A}$} او~\LR{$\sFmla{\False}{!A}$} دواړه ولري۔
تړلے تابلو هغه دے چې هره څانګه ئې تړلې وي۔ د شهودي تعبير له
مخې هره څانګه يو ګډ امکان بيانوي، خو
\LR{$\sFmla{\True}{!A}$} او~\LR{$\sFmla{\False}{!A}$} يو ځاے ممکن نۀ دي۔
په بله وينا، که يوه څانګه تړلې وي، هغه امکان رد شوے دے چې څانګه
ئې بيانوي۔ په ځانګړي ډول، د دې معنٰی دا ده چې تړلے تابلو هغه
ټول امکانات ردوي چې د~\LR{$\sFmla{\True}{!A}$} په بڼه هر فرض په يو وخت
رښتيا او د~\LR{$\sFmla{\False}{!A}$} په بڼه هر فرض دروغ کړي۔
\olpseganchor{OLP-0099-B010}{end}

\olpseganchor{OLP-0099-B011}{start}
\emph{د~\LR{$!A$} دپاره} تړلے تابلو داسې تړلے تابلو دے چې
ريښه ئې~\LR{$\sFmla{\False}{!A}$} وي۔ که داسې تړلے تابلو موجود
وي، د~\LR{$!A$} د دروغ کېدو ټول امکانات رد شوي؛ يعنې \LR{$!A$} بايد په هر
جوړښت کښې رښتيا وي۔
\olpseganchor{OLP-0099-B011}{end}

\olpseganchor{OLP-0100-B004}{start}
\olfileid{fol}{tab}{prl}
\olpseganchor{OLP-0100-B004}{end}

\olpseganchor{OLP-0100-B005}{start}
\olsection{قضيوي قاعدې}
\olpseganchor{OLP-0100-B005}{end}

\olpseganchor{OLP-0100-B006}{start}
\subsection{د~\LR{$\lnot$} قاعدې}
\olpseganchor{OLP-0100-B006}{end}

\olpseganchor{OLP-0100-B007}{start}
\begin{defish}
\AxiomC{\sFmla{\True}{\lnot !A}}
\RightLabel{\TRule{\True}{\lnot}}
\UnaryInfC{\sFmla{\False}{!A}}
\DisplayProof
\hfill
\AxiomC{\sFmla{\False}{\lnot !A}}
\RightLabel{\TRule{\False}{\lnot}}
\UnaryInfC{\sFmla{\True}{!A}}
\DisplayProof
\end{defish}
\olpseganchor{OLP-0100-B007}{end}

\olpseganchor{OLP-0100-B008}{start}
\subsection{د~\LR{$\land$} قاعدې}
\olpseganchor{OLP-0100-B008}{end}

\olpseganchor{OLP-0100-B009}{start}
\begin{defish}\noindent
\AxiomC{\sFmla{\True}{!A \land !B}}
\RightLabel{\TRule{\True}{\land}}
\UnaryInfC{\sFmla{\True}{!A}}
\noLine
\UnaryInfC{\sFmla{\True}{!B}}
\DisplayProof
\hfill
\AxiomC{\sFmla{\False}{!A \land !B}}
\RightLabel{\TRule{\False}{\land}}
\UnaryInfC{\LR{$\sFmla{\False}{!A} \quad \mid \quad \sFmla{\False}{!B}$}}
\DisplayProof
\end{defish}
\olpseganchor{OLP-0100-B009}{end}

\olpseganchor{OLP-0100-B010}{start}
\subsection{د~\LR{$\lor$} قاعدې}
\olpseganchor{OLP-0100-B010}{end}

\olpseganchor{OLP-0100-B011}{start}
\begin{defish}
\AxiomC{\sFmla{\True}{!A \lor !B}}
\RightLabel{\TRule{\True}{\lor}}
\UnaryInfC{\LR{$\sFmla{\True}{!A} \quad \mid \quad \sFmla{\True}{!B}$}}
\DisplayProof
\hfill
\AxiomC{\sFmla{\False}{!A \lor !B}}
\RightLabel{\TRule{\False}{\lor}}
\UnaryInfC{\sFmla{\False}{!A}}
\noLine
\UnaryInfC{\sFmla{\False}{!B}}
\DisplayProof
\end{defish}
\olpseganchor{OLP-0100-B011}{end}

\olpseganchor{OLP-0100-B012}{start}
\subsection{د~\LR{$\lif$} قاعدې}
\olpseganchor{OLP-0100-B012}{end}

\olpseganchor{OLP-0100-B013}{start}
\begin{defish}
\AxiomC{\sFmla{\True}{!A \lif !B}}
\RightLabel{\TRule{\True}{\lif}}
\UnaryInfC{\LR{$\sFmla{\False}{!A} \quad \mid \quad \sFmla{\True}{!B}$}}
\DisplayProof
\hfill
\AxiomC{\sFmla{\False}{!A \lif !B}}
\RightLabel{\TRule{\False}{\lif}}
\UnaryInfC{\sFmla{\True}{!A}}
\noLine
\UnaryInfC{\sFmla{\False}{!B}}
\DisplayProof
\end{defish}
\olpseganchor{OLP-0100-B013}{end}

\olpseganchor{OLP-0100-B014}{start}
\subsection{د پرېکون قاعده}
\olpseganchor{OLP-0100-B014}{end}

\olpseganchor{OLP-0100-B015}{start}
\begin{defish}
\AxiomC{}
\RightLabel{\Cut}
\UnaryInfC{\LR{$\sFmla{\True}{!A} \quad \mid \quad \sFmla{\False}{!A}$}}
\DisplayProof
\end{defish}
\olpseganchor{OLP-0100-B015}{end}

\olpseganchor{OLP-0100-B016}{start}
د~\Cut{} قاعده پر مخکښ نښه لرونکے فارمول نۀ «پلې کېږي»؛ بلکې
راکوي چې په تابلو کښې هره څانګه په دوو ووېشو: يوه څانګه
\LR{$\sFmla{\True}{!A}$} او بله~\LR{$\sFmla{\False}{!A}$} لري۔ دا اړينه نۀ
ده---د نښه لرونکي فارمولونه هر هغه سټ چې تړلے تابلو ولري،
له~\Cut پرته هم يو تړلے تابلو لري---خو راکوي چې تابلوګانې په
اسانه طريقه سره يو ځاے کړو۔
\olpseganchor{OLP-0100-B016}{end}

\olpseganchor{OLP-0101-B004}{start}
\olfileid{fol}{tab}{qrl}
\olpseganchor{OLP-0101-B004}{end}

\olpseganchor{OLP-0101-B005}{start}
\olsection{د سورونو قاعدې}
\olpseganchor{OLP-0101-B005}{end}

\olpseganchor{OLP-0101-B006}{start}
\subsection{د~\LR{$\lforall$} قاعدې}
\olpseganchor{OLP-0101-B006}{end}

\olpseganchor{OLP-0101-B007}{start}
\begin{defish}
\AxiomC{\sFmla{\True}{\lforall[x][!A(x)]}}
\RightLabel{\TRule{\True}{\forall}}
\UnaryInfC{\sFmla{\True}{!A(t)}}
\DisplayProof
\hfill
\AxiomC{\sFmla{\False}{\lforall[x][!A(x)]}}
\RightLabel{\TRule{\False}{\lforall}}
\UnaryInfC{\sFmla{\False}{!A(a)}}
\DisplayProof
\end{defish}
\olpseganchor{OLP-0101-B007}{end}

\olpseganchor{OLP-0101-B008}{start}
په~\TRule{\True}{\lforall} کښې \LR{$t$} تړلې اصطلاح ده (يعنې متغيرونه
نۀ لري)۔ په~\TRule{\False}{\lforall} کښې \LR{$a$} داسې ثابت سمبول
دے چې بايد د~\TRule{\False}{\lforall} تر قاعدې پاس په څانګه کښې
هېڅ ځاے نۀ وي راغلے۔ \LR{$a$} د~\TRule{\False}{\forall} استنتاج
\emph{ځانګړے متغير} بولو۔\footnote{د «ځانګړي متغير» اصطلاح کاروو،
سره له دې چې په پورته قاعده کښې \LR{$a$} يو ثابت سمبول دے۔ دا
تاريخي لاملونه لري۔}
\olpseganchor{OLP-0101-B008}{end}

\olpseganchor{OLP-0101-B009}{start}
\subsection{د~\LR{$\lexists$} قاعدې}
\olpseganchor{OLP-0101-B009}{end}

\olpseganchor{OLP-0101-B010}{start}
\begin{defish}
\AxiomC{\sFmla{\True}{\lexists[x][!A(x)]}}
\RightLabel{\TRule{\True}{\lexists}}
\UnaryInfC{\sFmla{\True}{!A(a)}}
\DisplayProof
\hfill
\AxiomC{\sFmla{\False}{\lexists[x][!A(x)]}}
\RightLabel{\TRule{\False}{\lexists}}
\UnaryInfC{\sFmla{\False}{!A(t)}}
\DisplayProof
\end{defish}
\olpseganchor{OLP-0101-B010}{end}

\olpseganchor{OLP-0101-B011}{start}
بيا هم \LR{$t$} تړلې اصطلاح ده، او~\LR{$a$} داسې ثابت سمبول دے چې د
\TRule{\True}{\lexists} تر قاعدې پاس په څانګه کښې نۀ راځي۔
\LR{$a$} د~\TRule{\True}{\lexists} استنتاج \emph{ځانګړے متغير} بولو۔
\olpseganchor{OLP-0101-B011}{end}

\olpseganchor{OLP-0101-B012}{start}
هغه شرط چې ځانګړے متغير د~\TRule{\False}{\lforall} يا
\TRule{\True}{\lexists} استنتاج تر پاسه په څانګه کښې نۀ راځي،
\emph{د ځانګړي متغير شرط} بلل کېږي۔
\olpseganchor{OLP-0101-B012}{end}

\olpseganchor{OLP-0101-B013}{start}
\begin{explain}
هغه قرارداد را ياد کړئ چې کله \LR{$!A$} داسې فارمول وي چې
متغير~\LR{$x$} پکې ازاد وي، دا د~\LR{$!A(x)$} په ليکلو ښيو۔ په همدې
سياق کښې \LR{$!A(t)$} د~\LR{$\Subst{!A}{t}{x}$} لنډه بڼه ده۔ نو د
\TRule{\False}{\lexists} قاعده داسې هم ليکلے شو:
\begin{prooftree}
  \AxiomC{\sFmla{\False}{\lexists[x][!A]}}
  \RightLabel{\TRule{\False}{\lexists}}
  \UnaryInfC{\sFmla{\False}{\Subst{!A}{t}{x}}}
\end{prooftree}
پام وکړئ چې \LR{$t$} له مخکښې په~\LR{$!A$} کښې راتلے شي؛ مثلاً \LR{$!A$} کېدے
شي \LR{$\Atom{\Obj P}{t,x}$} وي۔ نو له
\LR{$ \sFmla{\False}{\lexists[x][\Atom{\Obj P}{t,x}]}$} څخه
\LR{$\sFmla{\False}{\Atom{\Obj P}{t,t}}$} استنتاجول د
\TRule{\False}{\lexists} سمه کارونه ده۔ خو په
\TRule{\False}{\lforall} او~\TRule{\True}{\lexists} کښې د ځانګړي
متغير شرطونه غواړي چې ثابت سمبول~\LR{$a$} په~\LR{$!A$} کښې نۀ راځي۔ نو له
\LR{$\sFmla{\False}{\lforall[x][\Atom{\Obj P}{a,x}]}$}
څخه د~\LR{$\TRule{\False}{\lforall}$} په کارولو
\LR{$\sFmla{\False}{\Atom{\Obj P}{a,a}}$} په سمه توګه نۀ شئ استنتاجولے۔
\end{explain}
\olpseganchor{OLP-0101-B013}{end}

\olpseganchor{OLP-0101-B014}{start}
\begin{explain}
په~\TRule{\True}{\lforall} او~\TRule{\False}{\lexists} کښې د
تړلې اصطلاح~\LR{$t$} پر انتخاب نور بنديز نشته۔ بل لوري ته، د
\TRule{\True}{\lexists} او~\TRule{\False}{\lforall} په قاعدو کښې
د ځانګړي متغير شرط غواړي چې ثابت سمبول~\LR{$a$} د اړوند استنتاج تر
پاسه په څانګو کښې هېڅ ځاے نۀ راځي۔ دا شرط د نظام د سموالي د
يقيني کولو دپاره اړين دے۔ له دې شرط پرته لاندې به د
\LR{$\lexists[x][\formula{A}(x)] \lif \lforall[x][\formula{A}(x)]$}
دپاره يو تړلے تابلو و:
\begin{center}
\begin{tableau}{}
  [\sFmla{\False}{\lexists[x][\formula{A}(x)] \lif \lforall[x][\formula{A}(x)]}, just=\TAss
    [\sFmla{\True}{\lexists[x][\formula{A}(x)]},
      just={\TRule{\False}{\lif}[1]}
      [\sFmla{\False}{\lforall[x][\formula{A}(x)]},
        just={\TRule{\False}{\lif}[1]}
        [\sFmla{\True}{\formula{A}(a)},
          just={\TRule{\True}{\lexists}[2]}
          [\sFmla{\False}{\formula{A}(a)},
            just={\TRule{\False}{\lforall}[3]}, close]
        ]
      ]
    ]
  ]
\end{tableau}
\end{center}
خو \LR{$\lexists[x][\formula{A}(x)] \lif
\lforall[x][\formula{A}(x)]$} معتبر نۀ دے۔
\end{explain}
\olpseganchor{OLP-0101-B014}{end}

\olpseganchor{OLP-0102-B004}{start}
\olfileid{fol}{tab}{der}
\olpseganchor{OLP-0102-B004}{end}

\olpseganchor{OLP-0102-B005}{start}
\olsection{تابلوګانې}
\olpseganchor{OLP-0102-B005}{end}

\olpseganchor{OLP-0102-B006}{start}
\begin{explain}
و مو ويل چې فرض څه دے، او د استنتاج قاعدې مو ورکړې۔ له دغو څخه
تابلوګانې په استقرايي ډول جوړېږي: هر تابلو يا يوه يوازېنۍ
څانګه ده چې له يوه يا زياتو فرضونو جوړه وي، يا پر يوه څانګه د
استنتاج د کومې قاعدې په پلي کولو له بل تابلو څخه اخستل کېږي۔
\end{explain}
\olpseganchor{OLP-0102-B006}{end}

\olpseganchor{OLP-0102-B007}{start}
\begin{defn}[تابلو]
د فرضونو~\LR{$\sFmla{S_1}{!A_1}$}، \dots،
\LR{$\sFmla{S_n}{!A_n}$} دپاره (چې هر~\LR{$S_i$} يا~\LR{$\True$} وي
يا~\LR{$\False$}) تابلو د~نښه لرونکي فارمولونه يوه متناهي ونه
ده چې لاندې شرطونه پوره کوي:
\begin{enumerate}
\item د ونې تر ټولو پاسني \LR{$n$}~نښه لرونکي فارمولونه
  \LR{$\sFmla{S_i}{!A_i}$} دي، يو تر بل لاندې۔
\item په ونه کښې هر نښه لرونکے فارمول چې له فرضونو څخه نۀ وي، تر
  ځان پاس په څانګه کښې پر نښه لرونکے فارمول د يوې استنتاجي قاعدې
  له سمې کارونې څخه راوتلے وي۔
\end{enumerate}
د~تابلو يوه څانګه هله او يوازې هله \emph{تړلې} ده چې
\LR{$\sFmla{\True}{!A}$} او~\LR{$\sFmla{\False}{!A}$} دواړه ولري؛ که نۀ وي
\emph{پرانيستې} ده۔ هغه تابلو چې هره څانګه ئې تړلې وي، د
خپلو فرضونو د سټ \emph{تړلے تابلو} دے۔ که تابلو تړلے
نۀ وي، يعنې لږ تر لږه يوه پرانيستې څانګه ولري، نو \emph{پرانيستے}
دے۔
\end{defn}
\olpseganchor{OLP-0102-B007}{end}

\olpseganchor{OLP-0102-B008}{start}
\begin{ex}
  د فرضونو هر سټ په يوازې ځان تابلو دے، خو په عام ډول به
  تړلے نۀ وي۔ (په څرګند ډول يوازې هله تړلے دے چې فرضونه له مخکښې
  د~نښه لرونکي فارمولونه يوه جوړه
  \LR{$\sFmla{\True}{!A}$} او~\LR{$\sFmla{\False}{!A}$} ولري۔)
\olpseganchor{OLP-0102-B008}{end}

\olpseganchor{OLP-0102-B009}{start}
  له يوه تابلو (پرانيستي يا تړلي) څخه د کومې استنتاجي قاعدې
  په پلي کولو يو نوے، لوے تابلو اخستلے شو۔ قاعده په هغه کښې پر
  نښه لرونکے فارمول~\LR{$!A$} پلې کوو۔ دا قاعده به د هرې هغې څانګې په
  پاے کښې يو يا زيات نښه لرونکي فارمولونه ورزيات کړي چې د~\LR{$!A$} هغه
  پېښه لري چې قاعده مو پرې پلې کړې ده۔
\olpseganchor{OLP-0102-B009}{end}

\olpseganchor{OLP-0102-B010}{start}
  مثلاً فرض~\LR{$\sFmla{\True}{!A \land \lnot !A}$} ونيسئ۔ لاندې هغه
  (پرانيستے) تابلو دے چې يوازې له همدې فرض څخه جوړ دے:
  \begin{center}
    \begin{tableau}{}
      [\sFmla{\True}{\formula{A} \land \lnot \formula{A}}, just=\TAss]
    \end{tableau}{}
  \end{center}
  د~\LR{$\TRule{\True}{\land}$} قاعده پر فرض پلې کوو او ترې نوے
  تابلو اخلو۔ دا قاعده راکوي چې تابلو ته دوه نوې کرښې،
  \LR{$\sFmla{\True}{!A}$} او~\LR{$\sFmla{\True}{\lnot !A}$}، ورزياتې کړو:
  \begin{center}
    \begin{tableau}{}
      [\sFmla{\True}{\formula{A} \land \lnot \formula{A}}, just=\TAss
        [\sFmla{\True}{\formula{A}}, just={\TRule{\True}{\land}[1]},
          [\sFmla{\True}{\lnot \formula{A}}, just={\TRule{\True}{\land}[1]}
          ]
        ]
      ]
    \end{tableau}{}
  \end{center}
  کله چې تابلوګانې ليکو، پلې شوې قاعدې په ښي لوري کښې ثبتوو
  (مثلاً \LR{$\TRule{\True}{\land} 1$} دا معنٰی لري چې پر دې کرښه
  نښه لرونکے فارمول د~\LR{$\TRule{\True}{\land}$} قاعدې د هغې کارونې
  پايله ده چې د~\LR{$1$} پر کرښه نښه لرونکے فارمول ته شوې ده)۔ دا نوے
  تابلو اوس نور نښه لرونکي فارمولونه لري، خو يوازې پر يوه
  (\LR{$\sFmla{\True}{\lnot !A}$}) قاعده پلې کولے شو (په دې حالت کښې
  \LR{$\TRule{\True}{\lnot}$} قاعده)۔ له دې څخه لاندې تړلے
  تابلو راځي:
  \begin{center}
    \begin{tableau}{}
      [\sFmla{\True}{\formula{A} \land \lnot \formula{A}}, just=\TAss
        [\sFmla{\True}{\formula{A}}, just={\TRule{\True}{\land}[1]}
          [\sFmla{\True}{\lnot \formula{A}}, just={\TRule{\True}{\land}[1]}
            [\sFmla{\False}{\formula{A}}, just={\TRule{\True}{\lnot}[3]}, close]
          ]
        ]
      ]
    \end{tableau}{}
  \end{center}
\end{ex}
\olpseganchor{OLP-0102-B010}{end}

\olpseganchor{OLP-0103-B004}{start}
\olfileid{fol}{tab}{pro}
\olpseganchor{OLP-0103-B004}{end}

\olpseganchor{OLP-0103-B005}{start}
\olsection{د~تابلوګانې بېلګې}
\olpseganchor{OLP-0103-B005}{end}

\olpseganchor{OLP-0103-B006}{start}
\begin{ex}
راځئ د~تړلے فارمول~\LR{$(!A \land !B) \lif !A$} دپاره يو تړلے
تابلو ومومو۔
\olpseganchor{OLP-0103-B006}{end}

\olpseganchor{OLP-0103-B007}{start}
لومړے اړوند فرض د~تابلو په سر کښې ليکو۔
\begin{oltableau}
  [\sFmla{\False}{(\formula{A} \land \formula{B}) \lif \formula{A}},
    just = \TAss]
\end{oltableau}
\olpseganchor{OLP-0103-B007}{end}

\olpseganchor{OLP-0103-B008}{start}
يوازې يو فرض شته، نو يوازې يو نښه لرونکے فارمول دے چې قاعده پرې
پلې کولے شو۔ (د هر نښه لرونکے فارمول دپاره تل تر ټولو زياته يوه
قاعده پلې کېدے شي: دا د~تړلے فارمول د اړوندې نښې او
اصلي عملګر قاعده ده۔) په دې حالت کښې بايد
\LR{$\TRule{\False}{\lif}$} پلې کړو۔
\begin{oltableau}
  [\sFmla{\False}{(\formula{A} \land \formula{B}) \lif \formula{A}},
    checked, just = \TAss
    [\sFmla{\True}{\formula{A} \land \formula{B}},
      just={\TRule{\False}{\lif}[1]}
      [\sFmla{\False}{\formula{A}}, just={\TRule{\False}{\lif}[1]}]
    ]
  ]
\end{oltableau}
د دې حساب ساتلو دپاره چې پر کومو نښه لرونکي فارمولونه مو اړوندې
قاعدې پلې کړې دي، د جملې تر څنګ د سم نښان ليکو۔ خو د سم نښان
\emph{يوازې} هله وليکئ چې قاعده پر ټولو پرانيستو څانګو پلې شوې
وي۔ کله چې پر کوم نښه لرونکے فارمول په هره پرانيستې څانګه کښې
اړونده قاعده پلې شي، بيا ورته د بېرته ورګرځېدو او د قاعدې د بيا
پلي کولو اړتيا نۀ لرو۔ په دې حالت کښې يوازې يوه څانګه ده، نو قاعده
يوازې يو ځل پلې کول غواړي۔ (پام وکړئ چې د سم نښانونه يوازې د
تابلوګانو د جوړولو اسانتيا ده او په رسمي ډول د تابلوګانو د نحوې
برخه نۀ ده۔)
\olpseganchor{OLP-0103-B008}{end}

\olpseganchor{OLP-0103-B009}{start}
يو نوے نښه لرونکے فارمول شته چې قاعده پرې پلې کولے شو: پر
\LR{$2$}~کرښه \LR{$\sFmla{\True}{!A \land !B}$}۔ د
\LR{$\TRule{\True}{\land}$} قاعدې له پلي کولو څخه دا راځي:
\begin{oltableau}
  [\sFmla{\False}{(\formula{A} \land \formula{B}) \lif \formula{A}},
    checked, just = \TAss
    [\sFmla{\True}{\formula{A} \land \formula{B}},
      just={\TRule{\False}{\lif}[1]}, checked
      [\sFmla{\False}{\formula{A}}, just={\TRule{\False}{\lif}[1]}
        [\sFmla{\True}{\formula{A}}, just={\TRule{\True}{\land}[2]}
          [\sFmla{\True}{\formula{B}}, just={\TRule{\True}{\land}[2]}, close
          ]
        ]
      ]
    ]
  ]
\end{oltableau}
ځکه اوس څانګه \LR{$\sFmla{\True}{!A}$} (پر~\LR{$4$} کرښه) او
\LR{$\sFmla{\False}{!A}$} (پر~\LR{$3$} کرښه) دواړه لري، څانګه تړلې ده۔
ځکه همدا يوازېنۍ څانګه ده، تابلو تړلے دے۔ د
\LR{$(!A \land !B) \lif !A$} دپاره مو يو تړلے تابلو وموند۔
\end{ex}
\olpseganchor{OLP-0103-B009}{end}

\olpseganchor{OLP-0103-B010}{start}
\begin{ex}
اوس راځئ د~\LR{$(\lnot !A \lor !B) \lif (!A
\lif !B)$} دپاره يو تړلے تابلو ومومو۔
\olpseganchor{OLP-0103-B010}{end}

\olpseganchor{OLP-0103-B011}{start}
له اړوند فرض څخه پيل کوو:
\begin{oltableau}
  [\sFmla{\False}{(\lnot \formula{A} \lor \formula{B}) \lif
      (\formula{A} \lif \formula{B})}, just=\TAss]
\end{oltableau}
په دې تابلو کښې د يوازېني نښه لرونکے فارمول
اصلي عملګر~\LR{$\lif$} او نښه~\LR{$\False$} ده، نو د
\LR{$\TRule{\False}{\lif}$} قاعده پرې پلې کوو او دا اخلو:
\begin{oltableau}
  [\sFmla{\False}{(\lnot \formula{A} \lor \formula{B})
      \lif (\formula{A} \lif \formula{B})}, just=\TAss, checked
    [\sFmla{\True}{\lnot \formula{A} \lor \formula{B}},
      just={\TRule{\False}{\lif}[1]}
      [\sFmla{\False}{(\formula{A} \lif \formula{B})},
        just={\TRule{\False}{\lif}[1]}
      ]
    ]
  ]
\end{oltableau}
اوس دا ټاکنه لرو چې پر~\LR{$2$} کرښه~\LR{$\TRule{\True}{\lor}$} پلې کړو
که پر~\LR{$3$} کرښه \LR{$\TRule{\False}{\lif}$}۔ په حقيقت کښې د ترتيب
ټاکنه توپير نۀ کوي، تر هغې چې د هر نښه لرونکے فارمول اړونده قاعده
په هره څانګه کښې پلې شي۔ نو لومړۍ اخلو۔ د
\LR{$\TRule{\True}{\lor}$} قاعده تابلو ته څانګې ورکوي، او د
قاعدې دوه پايلې به هغه نوي نښه لرونکي فارمولونه وي چې دوو نويو
څانګو ته ورزياتېږي۔ له دې څخه دا راځي:
\begin{oltableau}
  [\sFmla{\False}{(\lnot \formula{A} \lor \formula{B}) \lif
      (\formula{A} \lif \formula{B})}, just=\TAss, checked
    [\sFmla{\True}{\lnot \formula{A} \lor \formula{B}},
      just={\TRule{\False}{\lif}[1]}, checked
      [\sFmla{\False}{(\formula{A} \lif \formula{B})},
        just={\TRule{\False}{\lif}[1]}
        [\sFmla{\True}{\lnot \formula{A}}, just={\TRule{\True}{\lor}[2]}]
        [\sFmla{\True}{\formula{B}}, just={\TRule{\True}{\lor}[2]}]
      ]
    ]
  ]
\end{oltableau}
لا مو پر~\LR{$3$} کرښه د~\LR{$\TRule{\False}{\lif}$} قاعده نۀ ده پلې
کړې؛ اوس ئې کوو۔ د وخت د سپما دپاره ئې پر دواړو څانګو پلې کوو۔
را ياد کړئ چې د کوم نښه لرونکے فارمول تر څنګ يوازې هله د سم نښان
ليکو چې اړونده قاعده مو په هره پرانيستې څانګه کښې پلې کړې وي۔ نو
غوره ده چې قاعده د هرې هغې څانګې په پاے کښې پلې کړو چې اړوند
نښه لرونکے فارمول لري۔ په دې توګه به په بېلابېلو څانګو کښې لاندې
هماغه نښه لرونکے فارمول ته بېرته نۀ ورګرځو۔
\begin{oltableau}
  [\sFmla{\False}{(\lnot \formula{A} \lor \formula{B}) \lif
      (\formula{A} \lif \formula{B})}, just=\TAss, checked
    [\sFmla{\True}{\lnot \formula{A} \lor \formula{B}},
      just={\TRule{\False}{\lif}[1]}, checked
      [\sFmla{\False}{(\formula{A} \lif \formula{B})},
        just={\TRule{\False}{\lif}[1]}, checked
        [\sFmla{\True}{\lnot \formula{A}}, just={\TRule{\True}{\lor}[2]}
          [\sFmla{\True}{\formula{A}}, just={\TRule{\False}{\lif}[3]}
            [\sFmla{\False}{\formula{B}}, just={\TRule{\False}{\lif}[3]}]
          ]
        ]
        [\sFmla{\True}{\formula{B}}, just={\TRule{\True}{\lor}[2]}
          [\sFmla{\True}{\formula{A}}, just={\TRule{\False}{\lif}[3]}
            [\sFmla{\False}{\formula{B}}, just={\TRule{\False}{\lif}[3]}, close]
          ]
        ]
      ]
    ]
  ]
\end{oltableau}
اوس ښۍ څانګه تړلې ده۔ په کيڼه څانګه کښې لا هم پر~\LR{$4$} کرښه د
\LR{$\TRule{\True}{\lnot}$} قاعده پلې کولے شو۔ له دې څخه
\LR{$\sFmla{\False}{!A}$} راځي او کيڼه څانګه تړل کېږي:
\begin{oltableau}
  [\sFmla{\False}{(\lnot \formula{A} \lor \formula{B}) \lif
      (\formula{A} \lif \formula{B})}, just=\TAss, checked
    [\sFmla{\True}{\lnot \formula{A} \lor \formula{B}},
      just={\TRule{\False}{\lif}[1]}, checked
      [\sFmla{\False}{(\formula{A} \lif \formula{B})},
        just={\TRule{\False}{\lif}[1]}, checked
        [\sFmla{\True}{\lnot \formula{A}}, just={\TRule{\True}{\lor}[2]}
          [\sFmla{\True}{\formula{A}}, just={\TRule{\False}{\lif}[3]}
            [\sFmla{\False}{\formula{B}}, just={\TRule{\False}{\lif}[3]}
              [\sFmla{\False}{\formula{A}},
                just={\TRule{\True}{\lnot}[4]}, close]
            ]
          ]
        ]
        [\sFmla{\True}{\formula{B}}, just={\TRule{\True}{\lor}[2]}
          [\sFmla{\True}{\formula{A}}, just={\TRule{\False}{\lif}[3]}
            [\sFmla{\False}{\formula{B}},
              just={\TRule{\False}{\lif}[3]}, close]
          ]
        ]
      ]
    ]
  ]
\end{oltableau}
\end{ex}
\olpseganchor{OLP-0103-B011}{end}

\olpseganchor{OLP-0103-B012}{start}
\begin{ex}
د فرضونو په توګه د هر شمېر نښه لرونکي فارمولونه دپاره
تابلوګانې ورکولے شو۔ ډېر ځله له يوې زياتې داسې قاعدې پلي کول هم
اړين وي چې څانګې ورکوي؛ او په عام ډول تابلو هر شمېر څانګې
لرلے شي۔ مثلاً د
\LR{$\{\sFmla{\True}{!A \lor (!B \land !C)}, \sFmla{\False}{(!A \lor !B) \land (!A
  \lor !C)}\}$} دپاره تابلو ونيسئ۔ لومړے پر لومړي فرض د
\LR{$\TRule{\True}{\lor}$} قاعده پلې کوو:
\begin{oltableau}
  [\sFmla{\True}{\formula{A} \lor (\formula{B} \land \formula{C})},
    just=\TAss, checked
    [\sFmla{\False}{(\formula{A} \lor \formula{B}) \land
        (\formula{A} \lor \formula{C})}, just=\TAss
      [\sFmla{\True}{\formula{A}}, just={\TRule{\True}{\lor}[1]}]
      [\sFmla{\True}{\formula{B} \land \formula{C}},
        just={\TRule{\True}{\lor}[1]}]
    ]
  ]
\end{oltableau}
اوس پر~\LR{$2$} کرښه د~\LR{$\TRule{\False}{\land}$} قاعده پلې کولے
شو۔ دا په دواړو څانګو کښې په يو وخت کوو، نو پر~\LR{$2$} کرښه د سم نښان
وهلے شو:
\begin{oltableau}
  [\sFmla{\True}{\formula{A} \lor (\formula{B} \land \formula{C})},
    just=\TAss, checked
    [\sFmla{\False}{(\formula{A} \lor \formula{B}) \land
        (\formula{A} \lor \formula{C})}, just=\TAss, checked
      [\sFmla{\True}{\formula{A}}, just={\TRule{\True}{\lor}[1]}
        [\sFmla{\False}{\formula{A} \lor \formula{B}},
          just={\TRule{\False}{\land}[2]}]
        [\sFmla{\False}{\formula{A} \lor \formula{C}},
          just={\TRule{\False}{\land}[2]}]
      ]
      [\sFmla{\True}{\formula{B} \land \formula{C}},
        just={\TRule{\True}{\lor}[1]}
        [\sFmla{\False}{\formula{A} \lor \formula{B}},
          just={\TRule{\False}{\land}[2]}]
        [\sFmla{\False}{\formula{A} \lor \formula{C}},
          just={\TRule{\False}{\land}[2]}]
      ]
    ]
  ]
\end{oltableau}
اوس \LR{$\TRule{\False}{\lor}$} پر ټولو هغو څانګو پلې کولے شو چې
\LR{$!A \lor !B$} لري:
\begin{oltableau}
  [\sFmla{\True}{\formula{A} \lor (\formula{B} \land \formula{C})},
    just=\TAss, checked
    [\sFmla{\False}{(\formula{A} \lor \formula{B}) \land
        (\formula{A} \lor \formula{C})}, just=\TAss, checked
      [\sFmla{\True}{\formula{A}}, just={\TRule{\True}{\lor}[1]}
        [\sFmla{\False}{\formula{A} \lor \formula{B}},
          just={\TRule{\False}{\land}[2]}, checked
          [\sFmla{\False}{\formula{A}}, just={\TRule{\False}{\lor}[4]}
            [\sFmla{\False}{\formula{B}},
              just={\TRule{\False}{\lor}[4]}, close]
          ]
        ]
        [\sFmla{\False}{\formula{A} \lor \formula{C}},
          just={\TRule{\False}{\land}[2]}]
      ]
      [\sFmla{\True}{\formula{B} \land \formula{C}},
        just={\TRule{\True}{\lor}[1]}
        [\sFmla{\False}{\formula{A} \lor \formula{B}},
          just={\TRule{\False}{\land}[2]}, checked
          [\sFmla{\False}{\formula{A}}, just={\TRule{\False}{\lor}[4]}
            [\sFmla{\False}{\formula{B}},
              just={\TRule{\False}{\lor}[4]}]
          ]
        ]
        [\sFmla{\False}{\formula{A} \lor \formula{C}},
          just={\TRule{\False}{\land}[2]}]
      ]
    ]
  ]
\end{oltableau}
اوس تر ټولو کيڼه څانګه تړلې ده۔ راځئ
\LR{$\TRule{\False}{\lor}$} پر~\LR{$!A \lor !C$} پلې کړو:
\begin{oltableau}
  [\sFmla{\True}{\formula{A} \lor (\formula{B} \land \formula{C})},
    just=\TAss, checked
    [\sFmla{\False}{(\formula{A} \lor \formula{B}) \land
        (\formula{A} \lor \formula{C})}, just=\TAss, checked
      [\sFmla{\True}{\formula{A}}, just={\TRule{\True}{\lor}[1]}
        [\sFmla{\False}{\formula{A} \lor \formula{B}},
          just={\TRule{\False}{\land}[2]}, checked
          [\sFmla{\False}{\formula{A}},
            just={\TRule{\False}{\lor}[4]}
            [\sFmla{\False}{\formula{B}},
              just={\TRule{\False}{\lor}[4]}, close]
          ]
        ]
        [\sFmla{\False}{\formula{A} \lor \formula{C}},
          just={\TRule{\False}{\land}[2]}, checked
          [\sFmla{\False}{\formula{A}},
            just={\TRule{\False}{\lor}[4]}, move by=2
            [\sFmla{\False}{\formula{C}},
              just={\TRule{\False}{\lor}[4]}, close]
          ]
        ]
      ]
      [\sFmla{\True}{\formula{B} \land \formula{C}},
        just={\TRule{\True}{\lor}[1]}
        [\sFmla{\False}{\formula{A} \lor \formula{B}},
          just={\TRule{\False}{\land}[2]}, checked
          [\sFmla{\False}{\formula{A}},
            just={\TRule{\False}{\lor}[4]}
            [\sFmla{\False}{\formula{B}},
              just={\TRule{\False}{\lor}[4]}
            ]
          ]
        ]
        [\sFmla{\False}{\formula{A} \lor \formula{C}},
          just={\TRule{\False}{\land}[2]}, checked
          [\sFmla{\False}{\formula{A}},
            just={\TRule{\False}{\lor}[4]}, move by=2
            [\sFmla{\False}{\formula{C}},
              just={\TRule{\False}{\lor}[4]}]
          ]
        ]
      ]
    ]
  ]
\end{oltableau}
پام وکړئ چې د روښانتيا دپاره مو د
\LR{$\TRule{\False}{\lor}$} د دويم ځل کارونې پايله لاندې يوړه۔ په دې
حالت کښې دې ته اړتيا نۀ وه؛ ځکه توجيهونه به يو شان وو۔
\olpseganchor{OLP-0103-B012}{end}

\olpseganchor{OLP-0103-B013}{start}
دوه څانګې پرانيستې پاتې دي، او پر~\LR{$3$} کرښه
\LR{$\sFmla{\True}{!B \land !C}$} لا بې د سم نښانه دے۔ پر هغې
\LR{$\TRule{\True}{\land}$} پلې کوو او تړلے تابلو اخلو:
\begin{oltableau}
  [\sFmla{\True}{\formula{A} \lor (\formula{B} \land \formula{C})},
    just=\TAss, checked
    [\sFmla{\False}{(\formula{A} \lor \formula{B}) \land
        (\formula{A} \lor \formula{C})}, just=\TAss, checked
      [\sFmla{\True}{\formula{A}}, just={\TRule{\True}{\lor}[1]}
        [\sFmla{\False}{\formula{A} \lor \formula{B}},
          just={\TRule{\False}{\land}[2]}, checked
          [\sFmla{\False}{\formula{A}},
            just={\TRule{\False}{\lor}[4]}
            [\sFmla{\False}{\formula{B}},
              just={\TRule{\False}{\lor}[4]}, close]
          ]
        ]
        [\sFmla{\False}{\formula{A} \lor \formula{C}},
          just={\TRule{\False}{\land}[2]}, checked
          [\sFmla{\False}{\formula{A}},
            just={\TRule{\False}{\lor}[4]}
            [\sFmla{\False}{\formula{C}},
              just={\TRule{\False}{\lor}[4]}, close]
          ]
        ]
      ]
      [\sFmla{\True}{\formula{B} \land \formula{C}},
        just={\TRule{\True}{\lor}[1]}, checked
        [\sFmla{\False}{\formula{A} \lor \formula{B}},
          just={\TRule{\False}{\land}[2]}, checked
          [\sFmla{\False}{\formula{A}},
            just={\TRule{\False}{\lor}[4]}
            [\sFmla{\False}{\formula{B}},
              just={\TRule{\False}{\lor}[4]}
              [\sFmla{\True}{\formula{B}},
                just={\TRule{\True}{\land}[3]}
                [\sFmla{\True}{\formula{C}},
                  just={\TRule{\True}{\land}[3]},close
                ]
              ]
            ]
          ]
        ]
        [\sFmla{\False}{\formula{A} \lor \formula{C}},
          just={\TRule{\False}{\land}[2]}, checked
          [\sFmla{\False}{\formula{A}},
            just={\TRule{\False}{\lor}[4]}
            [\sFmla{\False}{\formula{C}},
              just={\TRule{\False}{\lor}[4]}
              [\sFmla{\True}{\formula{B}},
                just={\TRule{\True}{\land}[3]}
                [\sFmla{\True}{\formula{C}},
                  just={\TRule{\True}{\land}[3]},close
                ]
              ]
            ]
          ]
        ]
      ]
    ]
  ]
\end{oltableau}
\olpseganchor{OLP-0103-B013}{end}

\olpseganchor{OLP-0103-B014}{start}
د پرتلې دپاره، لاندې د فرضونو د همدې سټ يو تړلے تابلو دے
چې قاعدې پکې په بل ترتيب پلې شوې دي:
\begin{oltableau}
  [\sFmla{\True}{\formula{A} \lor (\formula{B} \land \formula{C})},
    just=\TAss, checked
    [\sFmla{\False}{(\formula{A} \lor \formula{B}) \land
        (\formula{A} \lor \formula{C})}, just=\TAss, checked
      [\sFmla{\False}{\formula{A} \lor \formula{B}},
        just={\TRule{\False}{\land}[2]}, checked
        [\sFmla{\False}{\formula{A}},
          just={\TRule{\False}{\lor}[3]}
          [\sFmla{\False}{\formula{B}},
            just={\TRule{\False}{\lor}[3]}
            [\sFmla{\True}{\formula{A}},
              just={\TRule{\True}{\lor}[1]},close
            ]
            [\sFmla{\True}{\formula{B} \land \formula{C}},
              just={\TRule{\True}{\lor}[1]}, checked
              [\sFmla{\True}{\formula{B}},
                just={\TRule{\True}{\land}[6]}
                [\sFmla{\True}{\formula{C}},
                  just={\TRule{\True}{\land}[6]},close
                ]
              ]
            ]
          ]
        ]
      ]
      [\sFmla{\False}{\formula{A} \lor \formula{C}},
        just={\TRule{\False}{\land}[2]}, checked
        [\sFmla{\False}{\formula{A}},
          just={\TRule{\False}{\lor}[3]}
          [\sFmla{\False}{\formula{C}},
            just={\TRule{\False}{\lor}[3]}
            [\sFmla{\True}{\formula{A}},
              just={\TRule{\True}{\lor}[1]},close
            ]
            [\sFmla{\True}{\formula{B} \land \formula{C}},
              just={\TRule{\True}{\lor}[1]}, checked
              [\sFmla{\True}{\formula{B}},
                just={\TRule{\True}{\land}[6]}
                [\sFmla{\True}{\formula{C}},
                  just={\TRule{\True}{\land}[6]},close
                ]
              ]
            ]
          ]
        ]
      ]
    ]
  ]
\end{oltableau}
\end{ex}
\olpseganchor{OLP-0103-B014}{end}

\olpseganchor{OLP-0103-B015}{start}
\begin{prob}
د لاندې دپاره تړلي تابلوګانې ورکړئ:
\begin{enumerate}
\item \LR{$\sFmla{\True}{!A \land (!B \land !C)}, \sFmla{\False}{(!A \land !B) \land !C}$}.
\item \LR{$\sFmla{\True}{!A \lor (!B \lor !C)}, \sFmla{\False}{(!A \lor !B) \lor !C}$}.
\item \LR{$\sFmla{\True}{!A \lif (!B \lif !C)}, \sFmla{\False}{!B \lif (!A \lif !C)}$}.
\item \LR{$\sFmla{\True}{!A}, \sFmla{\False}{\lnot\lnot !A}$}.
\end{enumerate}
\end{prob}
\olpseganchor{OLP-0103-B015}{end}

\olpseganchor{OLP-0103-B016}{start}
\begin{prob}
د لاندې دپاره تړلي تابلوګانې ورکړئ:
\begin{enumerate}
\item \LR{$\sFmla{\True}{(!A \lor !B) \lif !C}, \sFmla{\False}{!A \lif !C}$}.
\item \LR{$\sFmla{\True}{(!A \lif !C) \land (!B \lif !C)}, \sFmla{\False}{(!A \lor !B) \lif !C}$}.
\item \LR{$\sFmla{\False}{\lnot(!A \land \lnot !A)}$}.
\item \LR{$\sFmla{\True}{!B \lif !A}, \sFmla{\False}{\lnot !A \lif \lnot !B}$}.
\item \LR{$\sFmla{\False}{(!A \lif \lnot !A) \lif \lnot !A}$}.
\item \LR{$\sFmla{\False}{\lnot(!A \lif !B) \lif \lnot !B}$}.
\item \LR{$\sFmla{\True}{!A \lif !C}, \sFmla{\False}{\lnot (!A \land \lnot !C)}$}.
\item \LR{$\sFmla{\True}{!A \land \lnot !C}, \sFmla{\False}{\lnot (!A \lif !C)}$}.
\item \LR{$\sFmla{\True}{!A \lor !B}, \sFmla{\True}{\lnot !B}, \sFmla{\False}{!A}$}.
\textbf{د ژباړې د نسخې سمون (OLTAB-002):} سرچينې دوه رښتوني فرضونه
د يو نښه لرونکي فارمول په دننه کښې سره يو ځاے کړي وو؛ دلته هر فرض
خپله رښتيا نښه لري، لکه د مسئلې په اړوند استلزام کښې۔
\item \LR{$\sFmla{\True}{\lnot !A \lor \lnot !B}, \sFmla{\False}{\lnot(!A \land !B)}$}.
\item \LR{$\sFmla{\False}{(\lnot !A \land \lnot !B) \lif\lnot(!A \lor !B)}$}.
\item \LR{$\sFmla{\False}{\lnot(!A \lor !B) \lif (\lnot !A \land \lnot !B)}$}.
\end{enumerate}
\end{prob}
\olpseganchor{OLP-0103-B016}{end}

\olpseganchor{OLP-0103-B017}{start}
\begin{prob}
د لاندې دپاره تړلي تابلوګانې ورکړئ:
\begin{enumerate}
\item \LR{$\sFmla{\True}{\lnot(!A \lif !B)}, \sFmla{\False}{!A}$}.
\item \LR{$\sFmla{\True}{\lnot(!A \land !B)}, \sFmla{\False}{\lnot !A \lor \lnot !B}$}.
\item \LR{$\sFmla{\True}{!A \lif !B}, \sFmla{\False}{\lnot !A \lor !B}$}.
\item \LR{$\sFmla{\False}{\lnot \lnot !A \lif !A}$}.
\item \LR{$\sFmla{\True}{!A \lif !B}, \sFmla{\True}{\lnot !A \lif !B}, \sFmla{\False}{!B}$}.
\item \LR{$\sFmla{\True}{(!A \land !B) \lif !C}, \sFmla{\False}{(!A \lif !C) \lor (!B \lif !C)}$}.
\item \LR{$\sFmla{\True}{(!A \lif !B) \lif !A}, \sFmla{\False}{!A}$}.
\item \LR{$\sFmla{\False}{(!A \lif !B) \lor (!B \lif !C)}$}.
\end{enumerate}
\end{prob}
\olpseganchor{OLP-0103-B017}{end}

\olpseganchor{OLP-0104-B004}{start}
\olfileid{fol}{tab}{prq}
\olpseganchor{OLP-0104-B004}{end}

\olpseganchor{OLP-0104-B005}{start}
\olsection{له سورونو سره تابلوګانې}
\olpseganchor{OLP-0104-B005}{end}

\olpseganchor{OLP-0104-B006}{start}
\begin{ex}
له سورونو سره د کار پر مهال بايد ډاډ واخلو چې د ځانګړي متغير
شرط نۀ ماتوو، او کله کله د دې دپاره د ځينو استنتاجونو د پلي کولو
ترتيب بدلول را بدلول غواړي۔ په عام ډول ګټه کوي چې لومړے د هغو
قاعدو د حل هڅه وکړو چې د ځانګړي متغير شرط ورباندې لګېږي (په بشپړ
تابلو کښې به دوئ پاس راځي)۔
\olpseganchor{OLP-0104-B006}{end}

\olpseganchor{OLP-0104-B007}{start}
راځئ وګورو چې د~تړلے فارمول
\LR{$\lexists[x][\lnot !A(x)] \lif \lnot \lforall[x][!A(x)]$}
دپاره تابلو څنګه ورکوو۔ د معمول په شان د فرض له ثبتولو
پيل کوو:
\begin{oltableau}
[\sFmla{\False}{\lexists[x][\lnot \formula{A}(x)]\lif \lnot
    \lforall[x][\formula{A}(x)]}, just=\TAss]
\end{oltableau}
ځکه اصلي عملګر~\LR{$\lif$} دے، نو
\LR{$\TRule{\False}{\lif}$} پلې کوو:
\begin{oltableau}
[\sFmla{\False}{\lexists[x][\lnot \formula{A}(x)]\lif \lnot
    \lforall[x][\formula{A}(x)]}, just=\TAss, checked
  [\sFmla{\True}{\lexists[x][\lnot \formula{A}(x)]},
    just={\TRule{\False}{\lif}[1]}
    [\sFmla{\False}{\lnot\lforall[x][\formula{A}(x)]},
      just={\TRule{\False}{\lif}[1]}]
  ]
]
\end{oltableau}
راتلونکې د کار کرښه~\LR{$2$} ده۔ \LR{$\TRule{\True}{\lexists}$} کاروو۔
دا يو نوے ثابت سمبول غواړي؛ ځکه لا هېڅ ثابت سمبولونه نۀ دي
راغلي، هر يو ټاکلے شو، مثلاً~\LR{$a$}۔
\begin{oltableau}
[\sFmla{\False}{\lexists[x][\lnot \formula{A}(x)]\lif \lnot
    \lforall[x][\formula{A}(x)]}, just=\TAss, checked
  [\sFmla{\True}{\lexists[x][\lnot \formula{A}(x)]},
    just={\TRule{\False}{\lif}[1]}, checked
    [\sFmla{\False}{\lnot\lforall[x][\formula{A}(x)]},
      just={\TRule{\False}{\lif}[1]}
      [\sFmla{\True}{\lnot\formula{A}(a)}, just={\TRule{\True}{\lexists}[2]}]
    ]
  ]
]
\end{oltableau}
اوس پر~\LR{$3$} کرښه \LR{$\TRule{\False}{\lnot}$} پلې کوو:
\begin{oltableau}
[\sFmla{\False}{\lexists[x][\lnot \formula{A}(x)]\lif \lnot
    \lforall[x][\formula{A}(x)]}, just=\TAss, checked
  [\sFmla{\True}{\lexists[x][\lnot \formula{A}(x)]},
    just={\TRule{\False}{\lif}[1]}, checked
    [\sFmla{\False}{\lnot\lforall[x][\formula{A}(x)]},
      just={\TRule{\False}{\lif}[1]}, checked
      [\sFmla{\True}{\lnot\formula{A}(a)}, just={\TRule{\True}{\lexists}[2]}
        [\sFmla{\True}{\lforall[x][\formula{A}(x)]},
          just = {\TRule{\False}{\lnot}[3]}
        ]
      ]
    ]
  ]
]
\end{oltableau}
پر~\LR{$4$} کرښه د~\LR{$\TRule{\True}{\lnot}$} او ورپسې پر~\LR{$5$} کرښه د
\LR{$\TRule{\True}{\lforall}$} په پلي کولو تړلے تابلو اخلو۔
\begin{oltableau}
[\sFmla{\False}{\lexists[x][\lnot \formula{A}(x)]\lif \lnot
    \lforall[x][\formula{A}(x)]}, just=\TAss, checked
  [\sFmla{\True}{\lexists[x][\lnot \formula{A}(x)]},
    just={\TRule{\False}{\lif}[1]}, checked
    [\sFmla{\False}{\lnot\lforall[x][\formula{A}(x)]},
      just={\TRule{\False}{\lif}[1]}, checked
      [\sFmla{\True}{\lnot\formula{A}(a)}, just={\TRule{\True}{\lexists}[2]}
        [\sFmla{\True}{\lforall[x][\formula{A}(x)]},
          just = {\TRule{\False}{\lnot}[3]}
          [\sFmla{\False}{\formula{A}(a)}, just={\TRule{\True}{\lnot}[4]}
            [\sFmla{\True}{\formula{A}(a)},
              just={\TRule{\True}{\lforall}[5]}, close
            ]
          ]
        ]
      ]
    ]
  ]
]
\end{oltableau}
\end{ex}
\olpseganchor{OLP-0104-B007}{end}

\olpseganchor{OLP-0104-B008}{start}
\begin{ex}
راځئ وګورو چې د لاندې سټ دپاره تابلو څنګه ورکوو:
\[
\sFmla{\False}{\lexists[x][!C(x,b)]},
\sFmla{\True}{\lexists[x][(!A(x) \land !B(x))]},
\sFmla{\True}{\lforall[x][(!B(x) \lif !C(x,b))]}.
\]
د معمول په شان له فرضونو پيل کوو:
\begin{oltableau}
  [\sFmla{\False}{\lexists[x][\formula{C}(x,b)]}, just=\TAss
    [\sFmla{\True}{\lexists[x][(\formula{A}(x) \land \formula{B}(x))]},
      just=\TAss
      [\sFmla{\True}{\lforall[x][(\formula{B}(x) \lif \formula{C}(x,b))]},
        just=\TAss]
    ]
  ]
\end{oltableau}
تل بايد لومړے هغه قاعده پلې کړو چې د ځانګړي متغير شرط لري؛ په
دې حالت کښې هغه پر~\LR{$2$} کرښه \LR{$\TRule{\True}{\lexists}$} ده۔ ځکه
فرضونه ثابت سمبول~\LR{$b$} لري، بايد بل يو وکاروو؛ بيا~\LR{$a$} اخلو۔
\begin{oltableau}
  [\sFmla{\False}{\lexists[x][\formula{C}(x,b)]}, just=\TAss
    [\sFmla{\True}{\lexists[x][(\formula{A}(x) \land \formula{B}(x))]},
      just=\TAss, checked
      [\sFmla{\True}{\lforall[x][(\formula{B}(x) \lif \formula{C}(x,b))]},
        just=\TAss
        [\sFmla{\True}{\formula{A}(a) \land \formula{B}(a)},
          just={\TRule{\True}{\lexists}[2]}]
      ]
    ]
  ]
\end{oltableau}
که اوس پر~\LR{$1$} کرښه \LR{$\TRule{\False}{\lexists}$} يا پر~\LR{$3$} کرښه
\LR{$\TRule{\True}{\lforall}$} پلې کړو، بايد وټاکو چې د~\LR{$x$} پر ځاے
کومه اصطلاح~\LR{$t$} کښېږدو۔ ځکه دا قاعدې د ځانګړي متغير شرط نۀ لري،
هره خوښه اصطلاح ټاکلے شو۔ په ځينو حالتونو کښې ښايي قاعده څو ځله
له بېلابېلو~\LR{$t$}s سره پلې کول هم لازم شي۔ خو د عمومي قاعدې په توګه
ګټه کوي چې له هغو اصطلاحګانو څخه يوه واخلو چې له مخکښې په
تابلو کښې راغلې وي---دلته \LR{$a$} او~\LR{$b$}---او په دې حالت کښې
اټکل کولے شو چې له~\LR{$a$} سره د تړلې څانګې لاس ته راتلل ډېر ممکن دي۔
\begin{oltableau}
  [\sFmla{\False}{\lexists[x][\formula{C}(x,b)]}, just=\TAss
    [\sFmla{\True}{\lexists[x][(\formula{A}(x) \land \formula{B}(x))]},
      just=\TAss, checked
      [\sFmla{\True}{\lforall[x][(\formula{B}(x) \lif \formula{C}(x,b))]}, just=\TAss
        [\sFmla{\True}{\formula{A}(a) \land \formula{B}(a)}, just={\TRule{\True}{\lexists}[2]}
          [\sFmla{\False}{\formula{C}(a,b)}, just={\TRule{\False}{\lexists}[1]}
            [\sFmla{\True}{\formula{B}(a) \lif \formula{C}(a,b)},
              just={\TRule{\True}{\lforall}[3]}
            ]
          ]
        ]
      ]
    ]
  ]
\end{oltableau}
پر~\LR{$1$} او~\LR{$3$} کرښو نښه لرونکي فارمولونه ته د سم نښان نۀ ورکوو؛
ځکه ښايي بيا ئې وکاروو۔ اوس پر~\LR{$4$} کرښه
\LR{$\TRule{\True}{\land}$} پلې کړئ:
\begin{oltableau}
  [\sFmla{\False}{\lexists[x][\formula{C}(x,b)]}, just=\TAss
    [\sFmla{\True}{\lexists[x][(\formula{A}(x) \land \formula{B}(x))]},
      just=\TAss, checked
      [\sFmla{\True}{\lforall[x][(\formula{B}(x) \lif \formula{C}(x,b))]}, just=\TAss
        [\sFmla{\True}{\formula{A}(a) \land \formula{B}(a)},
          just={\TRule{\True}{\lexists}[2]}, checked
          [\sFmla{\False}{\formula{C}(a,b)}, just={\TRule{\False}{\lexists}[1]}
            [\sFmla{\True}{\formula{B}(a) \lif \formula{C}(a,b)},
              just={\TRule{\True}{\lforall}[3]}
              [\sFmla{\True}{\formula{A}(a)}, just={\TRule{\True}{\land}[4]}
                [\sFmla{\True}{\formula{B}(a)}, just={\TRule{\True}{\land}[4]}
                ]
              ]
            ]
          ]
        ]
      ]
    ]
  ]
\end{oltableau}
که اوس پر~\LR{$6$} کرښه \LR{$\TRule{\True}{\lif}$} پلې کړو،
تابلو تړل کېږي:
\begin{oltableau}
  [\sFmla{\False}{\lexists[x][\formula{C}(x,b)]}, just=\TAss
    [\sFmla{\True}{\lexists[x][(\formula{A}(x) \land \formula{B}(x))]},
      just=\TAss, checked
      [\sFmla{\True}{\lforall[x][(\formula{B}(x) \lif \formula{C}(x,b))]},
        just=\TAss
        [\sFmla{\True}{\formula{A}(a) \land \formula{B}(a)},
          just={\TRule{\True}{\lexists}[2]}, checked
          [\sFmla{\False}{\formula{C}(a,b)},
            just={\TRule{\False}{\lexists}[1]}
            [\sFmla{\True}{\formula{B}(a) \lif \formula{C}(a,b)},
              just={\TRule{\True}{\lforall}[3]}, checked
              [\sFmla{\True}{\formula{A}(a)},
                just={\TRule{\True}{\land}[4]}
                [\sFmla{\True}{\formula{B}(a)},
                  just={\TRule{\True}{\land}[4]}
                  [\sFmla{\False}{\formula{B}(a)},
                    just={\TRule{\True}{\lif}[6]},close
                  ]
                  [\sFmla{\True}{\formula{C}(a,b)},
                    just={\TRule{\True}{\lif}[6]},close
                  ]
                ]
              ]
            ]
          ]
        ]
      ]
    ]
  ]
\end{oltableau}
\olpseganchor{OLP-0104-B008}{end}


\olpseganchor{OLP-0104-B009}{start}
\end{ex}
\olpseganchor{OLP-0104-B009}{end}

\olpseganchor{OLP-0104-B010}{start}
\begin{ex}
د لاندې سټ دپاره تابلو جوړوو:
\[
\sFmla{\True}{\lforall[x][!A(x)]}, \sFmla{\True}{\lforall[x][!A(x)] 
  \lif \lexists[y][!B(y)]}, \sFmla{\True}{\lnot\lexists[y][!B(y)]}.
\]
د معمول په شان فرضونه ليکو:
\begin{oltableau}
  [\sFmla{\True}{\lforall[x][\formula{A}(x)]}, just=\TAss
    [\sFmla{\True}{\lforall[x][\formula{A}(x)] \lif
        \lexists[y][\formula{B}(y)]}, just=\TAss
      [\sFmla{\True}{\lnot\lexists[y][\formula{B}(y)]}, just=\TAss
      ]
    ]
  ]
\end{oltableau}
لومړے پر~\LR{$3$} کرښه د~\LR{$\TRule{\True}{\lnot}$} قاعده پلې کوو۔
د «تل لومړے د ځانګړي متغير شرط لرونکې قاعدې پلې کړئ» قاعدې يوه
پايله دا ده: «بې د ځانګړي متغير شرطه د سورونو قاعدې تر اړتيا پورې
وځنډوئ۔» همدارنګه هغه قاعدې هم وځنډوئ چې وېش رامنځته کوي۔
\begin{oltableau}
  [\sFmla{\True}{\lforall[x][\formula{A}(x)]}, just=\TAss
    [\sFmla{\True}{\lforall[x][\formula{A}(x)] \lif
        \lexists[y][\formula{B}(y)]}, just=\TAss
      [\sFmla{\True}{\lnot\lexists[y][\formula{B}(y)]}, just=\TAss, checked
        [\sFmla{\False}{\lexists[y][\formula{B}(y)]}, just={\TRule{\True}{\lnot}[3]}]
      ]
    ]
  ]
\end{oltableau}
نوې \LR{$4$}~کرښه \LR{$\TRule{\False}{\lexists}$} غواړي، چې د ځانګړي
متغير له شرط پرته د سور قاعده ده۔ نو دا ځنډوو او پر~\LR{$2$} کرښه
\LR{$\TRule{\True}{\lif}$} کاروو۔
\begin{oltableau}
  [\sFmla{\True}{\lforall[x][\formula{A}(x)]}, just=\TAss
    [\sFmla{\True}{\lforall[x][\formula{A}(x)] \lif
        \lexists[y][\formula{B}(y)]}, just=\TAss, checked
      [\sFmla{\True}{\lnot\lexists[y][\formula{B}(y)]}, just=\TAss, checked
        [\sFmla{\False}{\lexists[y][\formula{B}(y)]},
          just={\TRule{\True}{\lnot}[3]},
          [\sFmla{\False}{\lforall[x][\formula{A}(x)]}, just={\TRule{\True}{\lif}[2]}]
          [\sFmla{\True}{\lexists[y][\formula{B}(y)]}, just={\TRule{\True}{\lif}[2]}]
        ]
      ]
    ]
  ]
\end{oltableau}
دواړه نوي نښه لرونکي فارمولونه د ځانګړي متغير د شرط لرونکې
قاعدې غواړي، نو راتلونکے ګام بايد همدا وي:
\begin{oltableau}
  [\sFmla{\True}{\lforall[x][\formula{A}(x)]}, just=\TAss
    [\sFmla{\True}{\lforall[x][\formula{A}(x)] \lif
        \lexists[y][\formula{B}(y)]}, just=\TAss, checked
      [\sFmla{\True}{\lnot\lexists[y][\formula{B}(y)]}, just=\TAss, checked
        [\sFmla{\False}{\lexists[y][\formula{B}(y)]},
          just={\TRule{\True}{\lnot}[3]}
          [\sFmla{\False}{\lforall[x][\formula{A}(x)]}, just={\TRule{\True}{\lif}[2]},checked
            [\sFmla{\False}{\formula{A}(b)}, just={\TRule{\False}{\lforall}[5]}]
          ]
          [\sFmla{\True}{\lexists[y][\formula{B}(y)]}, just={\TRule{\True}{\lif}[2]},checked
            [\sFmla{\True}{\formula{B}(c)}, just={\TRule{\True}{\lexists}[5]}]
          ]
        ]
      ]
    ]
  ]
\end{oltableau}
د څانګو د تړلو دپاره بايد پر~\LR{$1$} او~\LR{$4$} کرښو
نښه لرونکي فارمولونه وکاروو۔ د هغو اړوندې قاعدې
(\TRule{\True}{\lforall} او~\TRule{\False}{\lexists}) د ځانګړي
متغير شرط نۀ لري، نو هرې مناسبې اصطلاحګانې ټاکلے شو۔ په دې حالت
کښې په خپل وار همدا \LR{$b$} او~\LR{$c$} دي۔
\textbf{د ژباړې د نسخې سمون (OLTAB-003):} سرچينه دلته د
دروغ-وجودي نښه لرونکي فارمول د کرښې شمېر درې ښيي، خو درېيمه کرښه
د نفي مقدمه ده؛ اړونده دروغ-وجودي قاعده او وروستۍ ونه دواړه څلورمه
کرښه کاروي، نو اشاره سمه شوه۔
\begin{oltableau}
  [\sFmla{\True}{\lforall[x][\formula{A}(x)]}, just=\TAss
    [\sFmla{\True}{\lforall[x][\formula{A}(x)] \lif
        \lexists[y][\formula{B}(y)]}, just=\TAss, checked
      [\sFmla{\True}{\lnot\lexists[y][\formula{B}(y)]}, just=\TAss, checked
        [\sFmla{\False}{\lexists[y][\formula{B}(y)]},
          just={\TRule{\True}{\lnot}[3]}
          [\sFmla{\False}{\lforall[x][\formula{A}(x)]},
            just={\TRule{\True}{\lif}[2]},checked
            [\sFmla{\False}{\formula{A}(b)},
              just={\TRule{\False}{\lforall}[5]}
              [\sFmla{\True}{\formula{A}(b)},
                just={\TRule{\True}{\lforall}[1]},close
              ]
            ]
          ]
          [\sFmla{\True}{\lexists[y][\formula{B}(y)]},
            just={\TRule{\True}{\lif}[2]},checked
            [\sFmla{\True}{\formula{B}(c)},
              just={\TRule{\True}{\lexists}[5]}
              [\sFmla{\False}{\formula{B}(c)},
                just={\TRule{\False}{\lexists}[4]},close
              ]
            ]
          ]
        ]
      ]
    ]
  ]
\end{oltableau}
\end{ex}
\olpseganchor{OLP-0104-B010}{end}

\olpseganchor{OLP-0104-B011}{start}
\begin{prob}
د لاندې دپاره تړلي تابلوګانې ورکړئ:
\begin{enumerate}
\item \LR{$\sFmla{\False}{(\lforall[x][!A(x)] \land \lforall[y][!B(y)])
\lif \lforall[z][(!A(z) \land !B(z))]}$}.
\item \LR{$\sFmla{\False}{(\lexists[x][!A(x)] \lor \lexists[y][!B(y)])
\lif \lexists[z][(!A(z) \lor !B(z))]}$}.
\item \LR{$\sFmla{\True}{\lforall[x][(!A(x) \lif !B)]},
\sFmla{\False}{\lexists[y][!A(y)] \lif !B}$}.
\item \LR{$\sFmla{\True}{\lforall[x][\lnot !A(x)]},
\sFmla{\False}{\lnot\lexists[x][!A(x)]}$}.
\item \LR{$\sFmla{\False}{\lnot\lexists[x][!A(x)] \lif \lforall[x][\lnot
!A(x)]}$}.
\item \LR{$\sFmla{\False}{\lnot\lexists[x][\lforall[y][((!A(x,y) \lif
\lnot !A(y,y)) \land (\lnot !A(y,y) \lif !A(x,y)))]]}$}.
\end{enumerate}
\end{prob}
\olpseganchor{OLP-0104-B011}{end}

\olpseganchor{OLP-0104-B012}{start}
\begin{prob}
د لاندې دپاره تړلي تابلوګانې ورکړئ:
\begin{enumerate}
\item \LR{$\sFmla{\False}{\lnot\lforall[x][!A(x)] \lif \lexists[x][\lnot!A(x)]}$}.
\item \LR{$\sFmla{\True}{(\lforall[x][!A(x)] \lif !B)}, \sFmla{\False}{\lexists[y][(!A(y) \lif !B)]}$}.
\item \LR{$\sFmla{\False}{\lexists[x][(!A(x) \lif \lforall[y][!A(y)])]}$}.
\end{enumerate}
\end{prob}
\olpseganchor{OLP-0104-B012}{end}

\olpseganchor{OLP-0105-B004}{start}
\olfileid{fol}{tab}{ptn}
\olpseganchor{OLP-0105-B004}{end}
      
\olpseganchor{OLP-0105-B005}{start}
\olsection{ثبوت-تيوريکي مفاهيم}
\olpseganchor{OLP-0105-B005}{end}

\olpseganchor{OLP-0105-B006}{start}
\begin{editorial}
دا برخه د تابلوګانو د اثبات‌پذيرۍ اړيکې او سازګارۍ تعريفونه
راټولوي۔
\end{editorial}
\olpseganchor{OLP-0105-B006}{end}

\olpseganchor{OLP-0105-B007}{start}
\begin{explain}
لکه څنګه چې مو يو شمېر مهم معنايي مفاهيم (اعتبار، استلزام او د
صدق وړتيا) تعريف کړل، اوس ورته \emph{ثبوت-تيوريکي مفاهيم}
تعريفوو۔ دا مفاهيم په جوړښتونه کښې د تړلي فارمولونه د پوره
کېدو له مخې نۀ، بلکې د ځينو تړلو تابلوګانې د موجوديت له مخې
تعريفېږي۔ دا مهمه موندنه وه چې دغه مفاهيم سره سمون خوري۔ د دې
سمون منځپانګه د \emph{سموالي} او \emph{بشپړتيا قضيې} دي۔
\end{explain}
\olpseganchor{OLP-0105-B007}{end}


\olpseganchor{OLP-0105-B008}{start}
\begin{defn}[قضيې]
يوه تړلے فارمول~\LR{$!A$} هله \emph{قضيه} ده چې د
\LR{$\sFmla{\False}{!A}$} دپاره تړلے تابلو موجود وي۔ که \LR{$!A$}
قضيه وي، \LR{$\Proves !A$} ليکو، او که نۀ وي \LR{$\Proves/ !A$} ليکو۔
\end{defn}
\olpseganchor{OLP-0105-B008}{end}

\olpseganchor{OLP-0105-B009}{start}
\begin{defn}[د اشتقاق وړتيا]
تړلے فارمول~\LR{$!A$} د تړلي فارمولونه له سټ~\LR{$\Gamma$} څخه،
\LR{$\Gamma \Proves !A$}، هله \emph{د اشتقاق وړ} ده چې داسې متناهي
سټ~\LR{$\{!B_1, \dots, !B_n\} \subseteq \Gamma$} او د لاندې سټ دپاره
تړلے تابلو موجود وي:
\[
\{
  \sFmla{\False}{!A},
  \sFmla{\True}{!B_1}, \dots,
  \sFmla{\True}{!B_n}
  \}.
\]
که \LR{$!A$} له~\LR{$\Gamma$} څخه د اشتقاق وړ نۀ وي،
\LR{$\Gamma \Proves/ !A$} ليکو۔
\end{defn}
\olpseganchor{OLP-0105-B009}{end}

\olpseganchor{OLP-0105-B010}{start}
\begin{defn}[سازګاري]
د تړلي فارمولونه يو سټ~\LR{$\Gamma$} هله او يوازې هله \emph{ناسازګار}
دے چې داسې متناهي سټ~\LR{$\{!B_1, \dots, !B_n\} \subseteq \Gamma$}
او د لاندې سټ دپاره تړلے تابلو موجود وي:
\[
\{
  \sFmla{\True}{!B_1}, \dots,
  \sFmla{\True}{!B_n}  
  \}.
\]
که \LR{$\Gamma$} ناسازګار نۀ وي، \emph{سازګار} ئې بولو۔
\end{defn}
\olpseganchor{OLP-0105-B010}{end}

\olpseganchor{OLP-0105-B011}{start}
\begin{prop}[انعکاسيت]
\ollabel{prop:reflexivity}
که \LR{$!A \in \Gamma$} وي، نو \LR{$\Gamma \Proves !A$}۔
\end{prop}
\olpseganchor{OLP-0105-B011}{end}

\olpseganchor{OLP-0105-B012}{start}
\begin{proof}
که \LR{$!A \in \Gamma$} وي، نو \LR{$\{!A\}$} د~\LR{$\Gamma$} يو متناهي فرعي سټ
دے، او لاندې تابلو تړلے دے:
\begin{oltableau}
  [\sFmla{\False}{\formula{A}}, just = \TAss
    [\sFmla{\True}{\formula{A}}, just = \TAss,close]
  ]
\end{oltableau}
\end{proof}
\olpseganchor{OLP-0105-B012}{end}
  

\olpseganchor{OLP-0105-B013}{start}
\begin{prop}[يکنواختي]
\ollabel{prop:monotonicity}
که \LR{$\Gamma \subseteq \Delta$} او \LR{$\Gamma \Proves !A$} وي، نو
\LR{$\Delta \Proves !A$}۔
\end{prop}
\olpseganchor{OLP-0105-B013}{end}

\olpseganchor{OLP-0105-B014}{start}
\begin{proof}
د~\LR{$\Gamma$} هر متناهي فرعي سټ د~\LR{$\Delta$} هم متناهي فرعي سټ دے۔
\end{proof}
\olpseganchor{OLP-0105-B014}{end}

\olpseganchor{OLP-0105-B015}{start}
\begin{prop}[انتقاليت]
\ollabel{prop:transitivity}
که \LR{$\Gamma \Proves !A$} او
\LR{$\{!A\} \cup \Delta \Proves !B$} وي، نو
\LR{$\Gamma \cup \Delta \Proves !B$}۔
\end{prop}
\olpseganchor{OLP-0105-B015}{end}

\olpseganchor{OLP-0105-B016}{start}
\begin{proof}
که \LR{$\{!A\} \cup \Delta \Proves !B$} وي، نو داسې متناهي فرعي سټ
\LR{$\Delta_0 = \{!C_1, \dots, !C_n\} \subseteq \Delta$} شته چې
\begin{align*}
\{\sFmla{\False}{!B}, & \sFmla{\True}{!A}, \sFmla{\True}{!C_1},
\dots, \sFmla{\True}{!C_n}\}
\intertext{\RL{تړلے تابلو لري۔ که \LR{$\Gamma \Proves !A$} وي، نو
  داسې \LR{$!D_1$}, \dots, \LR{$!D_m \in \Gamma$} شته چې}}
\{\sFmla{\False}{!A}, & \sFmla{\True}{!D_1},
\dots, \sFmla{\True}{!D_m}\}
\end{align*}
تړلے تابلو لري۔
\textbf{د ژباړې د نسخې سمون (OLTAB-004):} سرچينه په منځنۍ جمله
کښې د ډي-لړ يوازې وروستۍ جمله د ګاما فرعي سټ بولي؛ دلته د لړ
ټولې جملې د ګاما غړي ښودل شوې، لکه د انتقاليت تعريف چې غواړي۔
\olpseganchor{OLP-0105-B016}{end}

\olpseganchor{OLP-0105-B017}{start}
اوس د لاندې فرضونو تابلو ونيسئ:
\[
\sFmla{\False}{!B},
\sFmla{\True}{!C_1}, \dots, \sFmla{\True}{!C_n},
\sFmla{\True}{!D_1}, \dots, \sFmla{\True}{!D_m}.
\]
پر~\LR{$!A$} د~\Cut{} قاعده پلې کړئ۔ دا دوه څانګې جوړوي: يوه پکې
\LR{$\sFmla{\True}{!A}$}، او بله پکې \LR{$\sFmla{\False}{!A}$} لري۔ نو په
يوه څانګه کښې دا ټول موجود دي:
\[
\{\sFmla{\False}{!B}, \sFmla{\True}{!A},
\sFmla{\True}{!C_1}, \dots, \sFmla{\True}{!C_n}\}
\]
ځکه د دغو فرضونو دپاره تړلے تابلو شته، هغه له دې څانګې سره
نښلولے شو؛ هره هغه څانګه تړل کېږي چې له
\LR{$\sFmla{\True}{!A}$} څخه تېرېږي۔ په بله څانګه کښې دا ټول موجود دي:
\[
\{\sFmla{\False}{!A}, \sFmla{\True}{!D_1}, \dots,
\sFmla{\True}{!D_m}\}
\]
نو بل لورے هم بشپړولے او تړلے تابلو اخستلے شو۔ دا
\LR{$\Gamma \cup \Delta \Proves !B$} ښيي۔
\end{proof}
\olpseganchor{OLP-0105-B017}{end}

\olpseganchor{OLP-0105-B018}{start}
پام وکړئ د دې يوه ځانګړې پايله دا ده: که \LR{$\Gamma \Proves !A$} او
\LR{$!A \Proves !B$} وي، نو \LR{$\Gamma \Proves !B$}۔ له دې څخه دا هم
لازمېږي چې که \LR{$!A_1,\dots,!A_n \Proves !B$} او د هر~\LR{$i$} دپاره
\LR{$\Gamma \Proves !A_i$} وي، نو \LR{$\Gamma \Proves !B$}۔
\olpseganchor{OLP-0105-B018}{end}

\olpseganchor{OLP-0105-B019}{start}
\begin{prop}
\ollabel{prop:incons}
\LR{$\Gamma$} هله او يوازې هله ناسازګار دے چې د هرې
تړلے فارمول~\LR{$!A$} دپاره \LR{$\Gamma \Proves !A$} وي۔
\end{prop}
\olpseganchor{OLP-0105-B019}{end}

\olpseganchor{OLP-0105-B020}{start}
\begin{proof}
تمرين۔
\end{proof}
\olpseganchor{OLP-0105-B020}{end}

\olpseganchor{OLP-0105-B021}{start}
\begin{prob}
\olref[fol][tab][ptn]{prop:incons} ثابت کړئ۔
\end{prob}
\olpseganchor{OLP-0105-B021}{end}



\olpseganchor{OLP-0105-B023}{start}
\begin{prop}[تراکم]
  \ollabel{prop:proves-compact}
  \begin{enumerate}
  \item که \LR{$\Gamma \Proves !A$} وي، نو داسې متناهي فرعي سټ
    \LR{$\Gamma_0 \subseteq \Gamma$} شته چې \LR{$\Gamma_0 \Proves !A$} وي۔
  \item که د~\LR{$\Gamma$} هر متناهي فرعي سټ سازګار وي، نو \LR{$\Gamma$}
    سازګار دے۔
  \end{enumerate}
\end{prop}
\olpseganchor{OLP-0105-B023}{end}

\olpseganchor{OLP-0105-B024}{start}
\begin{proof}
  \begin{enumerate}
    \item که \LR{$\Gamma \Proves !A$} وي، نو داسې متناهي فرعي سټ
      \LR{$\Gamma_0 = \{!B_1, \dots, !B_n\}$} او د لاندې دپاره تړلے
      تابلو شته:
      \[
      \{\sFmla{\False}{!A}, \sFmla{\True}{!B_1}, \dots, \sFmla{\True}{!B_n}\}
      \]
      همدا تابلو، \LR{$\Gamma_0 \Proves !A$} هم ښيي۔
    \item که \LR{$\Gamma$} ناسازګار وي، نو د کوم متناهي فرعي سټ
      \LR{$\Gamma_0 = \{!B_1, \dots, !B_n\}$} دپاره د لاندې يو تړلے
      تابلو شته:
      \[
      \{\sFmla{\True}{!B_1}, \dots, \sFmla{\True}{!B_n}\}
      \]
      دا تړلے تابلو ښيي چې \LR{$\Gamma_0$} ناسازګار دے۔
  \end{enumerate}
\end{proof}
\olpseganchor{OLP-0105-B024}{end}

\olpseganchor{OLP-0106-B004}{start}
\olfileid{fol}{tab}{prv}
\olpseganchor{OLP-0106-B004}{end}

\olpseganchor{OLP-0106-B005}{start}
\olsection{د اشتقاق وړتيا او سازګاري}
\olpseganchor{OLP-0106-B005}{end}

\olpseganchor{OLP-0106-B006}{start}
اوس د~د اشتقاق وړتيا اړيکې يو شمېر خاصيتونه ثابتوو۔ دوئ په
خپلواک ډول هم په زړه پورې دي، خو هر يو به د بشپړتيا قضیې په ثبوت
کښې ونډه ولري۔
\olpseganchor{OLP-0106-B006}{end}

\olpseganchor{OLP-0106-B007}{start}
\begin{prop}\ollabel{prop:provability-contr}
  که \LR{$\Gamma \Proves !A$} او \LR{$\Gamma \cup \{!A\}$} ناسازګار وي،
  نو \LR{$\Gamma$} ناسازګار دے۔
\end{prop}
\olpseganchor{OLP-0106-B007}{end}

\olpseganchor{OLP-0106-B008}{start}
\begin{proof}
داسې متناهي
\LR{$\Gamma_0 = \{!B_1, \dots, !B_n\}$} او
\LR{$\Gamma_1 =\{!C_1, \dots, !C_m\} \subseteq \Gamma$} شته چې
  \begin{align*}
    \{\sFmla{\False}{!A}, &
    \sFmla{\True}{!B_1}, \dots, \sFmla{\True}{!B_n}\} \\
    \{\sFmla{\True}{!A}, &
    \sFmla{\True}{!C_1}, \dots, \sFmla{\True}{!C_m}\}
  \end{align*}
  تړلي تابلوګانې لري۔ پر~\LR{$!A$} د~\Cut{} قاعدې په کارولو دا په
  يوه تړلي تابلو کښې يو ځاے کولے شو چې د
  \LR{$\Gamma_0 \cup \Gamma_1$} ناسازګاري ښيي۔ ځکه
  \LR{$\Gamma_0 \subseteq \Gamma$} او~\LR{$\Gamma_1 \subseteq \Gamma$}،
  نو \LR{$\Gamma_0 \cup \Gamma_1 \subseteq \Gamma$}؛ له دې امله
  \LR{$\Gamma$} ناسازګار دے۔
\textbf{د ژباړې د نسخې سمون (OLTAB-005):} سرچينه د دويم متناهي
سټ په تعريف کښې لړ په اېن شاخص پاے ته رسوي، خو په تابلو کښې ئې
وروستے غړے د اېم شاخص لري؛ دلته د لړ وروستے شاخص يو شان شوے دے۔
\end{proof}
\olpseganchor{OLP-0106-B008}{end}

\olpseganchor{OLP-0106-B009}{start}
\begin{prop}
\ollabel{prop:prov-incons}
\LR{$\Gamma \Proves !A$} هله او يوازې هله چې
\LR{$\Gamma \cup \{\lnot !A\}$} ناسازګار وي۔
\end{prop}
\olpseganchor{OLP-0106-B009}{end}

\olpseganchor{OLP-0106-B010}{start}
\begin{proof}
لومړے فرض کړئ \LR{$\Gamma \Proves !A$}، يعنې د لاندې دپاره تړلے
تابلو شته:
\[
\{\sFmla{\False}{!A},
\sFmla{\True}{!B_1}, \dots, \sFmla{\True}{!B_n}\}
\]
د~\LR{$\TRule{\True}{\lnot}$} قاعدې په کارولو دا د لاندې دپاره په تړلي
تابلو اړولے شو:
\[
\{\sFmla{\True}{\lnot !A},
\sFmla{\True}{!B_1}, \dots, \sFmla{\True}{!B_n}\}.
\]
\olpseganchor{OLP-0106-B010}{end}

\olpseganchor{OLP-0106-B011}{start}
بل لوري ته، که د وروستي سټ دپاره تړلے تابلو وي، د لومړي سټ
په تړلي تابلو ئې اړولے شو: هغه هر فارمول ړنګوو چې پر لومړي
فرض~\LR{$\sFmla{\True}{\lnot !A}$} د~\TRule{\True}{\lnot} له کارونې
راوتلے وي، خپله دغه فرض هم ړنګوو، او فرض
\LR{$\sFmla{\False}{!A}$} ورزياتوو۔ ځکه که کومه څانګه مخکښې د
\TRule{\True}{\lnot} د هغې پايلې له امله تړلې وه چې پر
\LR{$\sFmla{\True}{\lnot !A}$} پلې شوې، يعنې
\LR{$\sFmla{\False}{!A}$}، نو په نوي تابلو کښې اړونده څانګه هم
تړلې ده۔ که په زوړ تابلو کښې کومه څانګه د فرض
\LR{$\sFmla{\True}{\lnot !A}$} او~\LR{$\sFmla{\False}{\lnot !A}$} دواړو له
امله تړلې وه، پر~\LR{$\sFmla{\False}{\lnot !A}$} د
\LR{$\TRule{\False}{\lnot}$} په پلي کولو او د
\LR{$\sFmla{\True}{!A}$} په اخستلو ئې په تړلې څانګه اړولے شو۔ دا څانګه
ځکه تړل کېږي چې \LR{$\sFmla{\False}{!A}$} مو د فرض په توګه ورزيات کړے
دے۔
\end{proof}
\olpseganchor{OLP-0106-B011}{end}

\olpseganchor{OLP-0106-B012}{start}
\begin{prob}
ثابت کړئ چې \LR{$\Gamma \Proves \lnot !A$} هله او يوازې هله چې
\LR{$\Gamma \cup \{!A\}$} ناسازګار وي۔
\end{prob}
\olpseganchor{OLP-0106-B012}{end}

\olpseganchor{OLP-0106-B013}{start}
\begin{prop}\ollabel{prop:explicit-inc}
  که \LR{$\Gamma \Proves !A$} او~\LR{$\lnot !A \in \Gamma$} وي، نو
  \LR{$\Gamma$} ناسازګار دے۔
\end{prop}
\olpseganchor{OLP-0106-B013}{end}

\olpseganchor{OLP-0106-B014}{start}
\begin{proof}
  فرض کړئ \LR{$\Gamma \Proves !A$} او~\LR{$\lnot !A \in \Gamma$}۔ بيا داسې
  \LR{$!B_1$}، \dots، \LR{$!B_n \in \Gamma$} شته چې
  \[
  \{\sFmla{\False}{!A}, \sFmla{\True}{!B_1}, \dots, \sFmla{\True}{!B_n}\}
  \]
  تړلے تابلو لري۔ فرض~\sFmla{\False}{!A} په
  \sFmla{\True}{\lnot !A} بدل کړئ، او د فرضونو وروسته هغه
  \sFmla{\False}{!A} پايله ورننباسئ چې پر
  \sFmla{\True}{\lnot !A} د~\TRule{\True}{\lnot} له کارونې
  راځي۔ په زوړ تابلو کښې هره تړلے فارمول چې د~\LR{$1$} کرښې په
  حواله توجيه شوې وه، اوس د~\LR{$n+1$} کرښې په حواله توجيه کېږي۔ نو که
  زوړ تابلو تړلے و، نوے تابلو هم تړلے دے۔ دا د~\LR{$\Gamma$} ناسازګاري
  ښيي؛ ځکه ټول فرضونه په~\LR{$\Gamma$} کښې دي۔
  \textbf{د ژباړې د نسخې سمون (OLTAB-006):} سرچينه د «داسې چې»
  وروسته يو بې دليله تېک شاتګی ږدي؛ دلته لرې شوے دے۔
  \textbf{د ژباړې د نسخې سمون (OLTAB-007):} سرچينه وايي د رښتونې
  نفي قاعده پر دروغې جملې پلې کېږي؛ قاعده په حقيقت کښې پر رښتونې
  نفي شوې جملې پلې کېږي او دروغ مثبت فارمول راوباسي، لکه د همدې
  ثبوت په راتلونکې جمله کښې۔
\end{proof}
\olpseganchor{OLP-0106-B014}{end}

\olpseganchor{OLP-0106-B015}{start}
\begin{prop}\ollabel{prop:provability-exhaustive}
  که \LR{$\Gamma \cup \{!A\}$} او~\LR{$\Gamma \cup \{\lnot !A\}$} دواړه
  ناسازګار وي، نو \LR{$\Gamma$} ناسازګار دے۔
\end{prop}
\olpseganchor{OLP-0106-B015}{end}

\olpseganchor{OLP-0106-B016}{start}
\begin{proof}
  که داسې \LR{$!B_1$}، \dots، \LR{$!B_n \in \Gamma$} او
  \LR{$!C_1$}، \dots، \LR{$!C_m \in \Gamma$} وي چې
  \begin{align*}
    \{\sFmla{\True}{!A}, &
    \sFmla{\True}{!B_1}, \dots, \sFmla{\True}{!B_n}\} \text{\RL{ او}}\\
    \{\sFmla{\True}{\lnot !A}, &
    \sFmla{\True}{!C_1}, \dots, \sFmla{\True}{!C_m}\}
  \end{align*}
  دواړه تړلي تابلوګانې ولري، يو داسې ګډ تابلو جوړولے شو
  چې د~\LR{$\Gamma$} ناسازګاري ښيي۔ د فرضونو په توګه
  \LR{$\sFmla{\True}{!B_1}$}، \dots، \LR{$\sFmla{\True}{!B_n}$} له
  \LR{$\sFmla{\True}{!C_1}$}، \dots، \LR{$\sFmla{\True}{!C_m}$} سره يو
  ځاے کاروو، او بيا د~\Cut{} قاعده پلې کوو۔ دا دوه څانګې ورکوي:
  يوه په~\LR{$\sFmla{\True}{!A}$} او بله په~\LR{$\sFmla{\False}{!A}$} پيلېږي۔
\olpseganchor{OLP-0106-B016}{end}
  
\olpseganchor{OLP-0106-B017}{start}
  په کيڼ لوري کښې د لومړي تابلو د فرضونو لاندې برخه ورزياته
  کړئ۔ دلته د قاعدې هره کارونه لا هم سمه ده؛ ځکه د لومړي
  تابلو هر فرض، د~\LR{$\sFmla{\True}{!A}$} په شمول، موجود دے۔
  نو د~\LR{$\sFmla{\True}{!A}$} لاندې هره څانګه تړل کېږي۔
\olpseganchor{OLP-0106-B017}{end}
  
\olpseganchor{OLP-0106-B018}{start}
  په ښي لوري کښې د دويم تابلو د فرض لاندې برخه ورزياته کړئ،
  خو د~\LR{$\sFmla{\True}{\lnot !A}$} پر فرض د
  \LR{$\TRule{\True}{\lnot}$} د هرې کارونې پايلې ترې لرې کړئ۔ پر
  \LR{$\sFmla{\True}{\lnot !A}$} د~\LR{$\TRule{\True}{\lnot}$} پايله
  \LR{$\sFmla{\False}{!A}$} ده، چې بيا هم موجوده ده؛ ځکه د ګډ
  تابلو په ښي لوري کښې د~\Cut{} قاعدې پايله ده۔
\olpseganchor{OLP-0106-B018}{end}

\olpseganchor{OLP-0106-B019}{start}
  که په دويم تابلو کښې کومه څانګه د فرض
  \LR{$\sFmla{\True}{\lnot !A}$} (چې په ګډ تابلو کښې نور فرض نۀ
  دے) او~\LR{$\sFmla{\False}{\lnot !A}$} دواړو له امله تړلې وه، پر
  \LR{$\sFmla{\False}{\lnot !A}$} د~\LR{$\TRule{\False}{\lnot}$} په پلي کولو
  \LR{$\sFmla{\True}{!A}$} اخستلے شو۔ اوس په ګډ تابلو کښې اړونده
  څانګه هم تړل کېږي؛ ځکه د~\Cut{} قاعدې ښۍ پايله،
  \LR{$\sFmla{\False}{!A}$}، پکې ده۔ که په دويم تابلو کښې يوه
  څانګه د بل کوم لامل له امله تړلې وه، په ګډ تابلو کښې
  اړونده څانګه هم تړلې ده؛ ځکه په زوړ، دويم تابلو کښې پر
  څانګه له~\LR{$\sFmla{\True}{\lnot !A}$} پرته هر بل نښه لرونکي فارمولونه
  د ګډ تابلو پر اړوندې څانګې هم راځي۔
  \end{proof}
\olpseganchor{OLP-0106-B019}{end}

\olpseganchor{OLP-0107-B005}{start}
\olfileid{fol}{tab}{ppr}
\olpseganchor{OLP-0107-B005}{end}

\olpseganchor{OLP-0107-B006}{start}
\olsection{د اشتقاق وړتيا او قضيوي رابطونه}
\olpseganchor{OLP-0107-B006}{end}

\olpseganchor{OLP-0107-B007}{start}
\begin{explain}
  ثابتوو چې د تابلوګانو د اشتقاق وړتيا اړيکه~\LR{$\Proves$} دومره
  پياوړې ده چې د قضيوي رابطونو په اړه ځينې بنسټيز حقيقتونه ثابت
  کړي؛ لکه \LR{$!A \land !B \Proves !A$} او
  \LR{$!A, !A \lif !B \Proves !B$} (مودوس پوننس)۔ دا حقيقتونه د
  بشپړتيا قضیې د ثبوت دپاره پکار دي۔
\end{explain}
\olpseganchor{OLP-0107-B007}{end}

\olpseganchor{OLP-0107-B008}{start}
\begin{prop}\ollabel{prop:provability-land}
  \begin{enumerate}
  \item \ollabel{prop:provability-land-left} دواړه
    \LR{$!A \land !B \Proves !A$} او~\LR{$!A \land !B \Proves !B$}۔
  \item \ollabel{prop:provability-land-right}
    \LR{$!A, !B \Proves !A \land !B$}۔
  \end{enumerate}
\end{prop}
\olpseganchor{OLP-0107-B008}{end}

\olpseganchor{OLP-0107-B009}{start}
\begin{proof}
  \begin{enumerate}
  \item د~\LR{$\{\sFmla{\False}{!A}, \sFmla{\True}{!A \land !B}\}$} او
    \LR{$\{\sFmla{\False}{!B}, \sFmla{\True}{!A \land !B}\}$} دواړو
    دپاره تړلي تابلوګانې شته:
    \begin{oltableau}{}
      [\sFmla{\False}{\formula{A}}, just=\TAss
        [\sFmla{\True}{\formula{A} \land \formula{B}}, just=\TAss
          [\sFmla{\True}{\formula{A}},just={\TRule{\True}{\land}[2]}
            [\sFmla{\True}{\formula{B}},just={\TRule{\True}{\land}[2]}, close
            ]
          ]
        ]
      ]
    \end{oltableau}
    \begin{oltableau}{}
      [\sFmla{\False}{\formula{B}}, just=\TAss
        [\sFmla{\True}{\formula{A} \land \formula{B}}, just=\TAss
          [\sFmla{\True}{\formula{A}},just={\TRule{\True}{\land}[2]}
            [\sFmla{\True}{\formula{B}},just={\TRule{\True}{\land}[2]}, close
            ]
          ]
        ]
      ]
    \end{oltableau}
    \item د~\LR{$\{\sFmla{\True}{!A},
      \sFmla{\True}{!B}, \sFmla{\False}{!A \land !B}\}$} دپاره
      لاندې تړلے تابلو دے:
      \begin{oltableau}
        [\sFmla{\False}{\formula{A} \land \formula{B}}, just = \TAss
          [\sFmla{\True}{\formula{A}}, just = \TAss
            [\sFmla{\True}{\formula{B}}, just=\TAss
              [\sFmla{\False}{\formula{A}}, just = {\TRule{\False}{\land}[1]}, close]
              [\sFmla{\False}{\formula{B}}, just = {\TRule{\False}{\land}[1]}, close]
            ]
          ]
        ]
      \end{oltableau}
  \end{enumerate}
  \textbf{د ژباړې د نسخې سمون (OLTAB-008):} سرچينه په لومړيو
  څلورو تابلوګانو کښې د اتو نښه لرونکو فارمولونو د نښې او فارمول
  تر منځ جلا کونکے قوس غورځوي؛ دلته د تعريف شوې دوه-ارګومېنټه
  د \texttt{\string\sFmla} دوه-ارګومېنټه بڼه بېرته راوړل شوې ده۔
\end{proof}
\olpseganchor{OLP-0107-B009}{end}

\olpseganchor{OLP-0107-B010}{start}
\begin{prop}\ollabel{prop:provability-lor}
  \begin{enumerate}
  \item \LR{$\{!A \lor !B, \lnot !A, \lnot !B\}$} ناسازګار دے۔
  \item دواړه \LR{$!A \Proves !A \lor !B$} او
    \LR{$!B \Proves !A \lor !B$}۔
  \end{enumerate}
\end{prop}
\olpseganchor{OLP-0107-B010}{end}

\olpseganchor{OLP-0107-B011}{start}
\begin{proof}
  \begin{enumerate}
  \item د~\LR{$\{\sFmla{\True}{!A \lor !B},
    \sFmla{\True}{\lnot !A}, \sFmla{\True}{\lnot !B}\}$} تړلے
    تابلو ورکوو:
    \begin{oltableau}
      [\sFmla{\True}{\formula{A} \lor \formula{B}}, just = \TAss
        [\sFmla{\True}{\lnot \formula{A}}, just = \TAss
          [\sFmla{\True}{\lnot \formula{B}}, just = \TAss
            [\sFmla{\False}{\formula{A}}, just = {\TRule{\True}{\lnot}[2]}
              [\sFmla{\False}{\formula{B}}, just = {\TRule{\True}{\lnot}[3]}
                [\sFmla{\True}{\formula{A}}, just = {\TRule{\True}{\lor}[1]}, close]
                [\sFmla{\True}{\formula{B}}, just = {\TRule{\True}{\lor}[1]}, close]
              ]
            ]
          ]
        ]
      ]
    \end{oltableau}
  \item د~\LR{$\{\sFmla{\False}{!A \lor !B}, \sFmla{\True}{!A}\}$} او
    \LR{$\{\sFmla{\False}{!A \lor !B}, \sFmla{\True}{!B}\}$} دواړو
    دپاره تړلي تابلوګانې شته:
    \begin{oltableau}{}
      [\sFmla{\False}{\formula{A} \lor \formula{B}}, just=\TAss
        [\sFmla{\True}{\formula{A}}, just=\TAss
          [\sFmla{\False}{\formula{A}},just={\TRule{\False}{\lor}[1]}
            [\sFmla{\False}{\formula{B}},just={\TRule{\False}{\lor}[1]}, close
            ]
          ]
        ]
      ]
    \end{oltableau}
    \begin{oltableau}{}
      [\sFmla{\False}{\formula{A} \lor \formula{B}}, just=\TAss
        [\sFmla{\True}{\formula{B}}, just=\TAss
          [\sFmla{\False}{\formula{A}},just={\TRule{\False}{\lor}[1]}
            [\sFmla{\False}{\formula{B}},just={\TRule{\False}{\lor}[1]}, close
            ]
          ]
        ]
      ]
    \end{oltableau}
  \end{enumerate}
\end{proof}
\olpseganchor{OLP-0107-B011}{end}

\olpseganchor{OLP-0107-B012}{start}
\begin{prop}\ollabel{prop:provability-lif}
  \begin{enumerate}
  \item \ollabel{prop:provability-lif-left}
    \LR{$!A, !A \lif !B \Proves !B$}۔
  \item \ollabel{prop:provability-lif-right}
    دواړه \LR{$\lnot !A \Proves !A \lif !B$} او
    \LR{$!B \Proves !A \lif !B$}۔
  \end{enumerate}
\end{prop}
\olpseganchor{OLP-0107-B012}{end}

\olpseganchor{OLP-0107-B013}{start}
\begin{proof}
  \begin{enumerate}
    \item د~\LR{$\{\sFmla{\False}{!B}, \sFmla{\True}{!A \lif !B},
      \sFmla{\True}{!A}\}$} يو تړلے تابلو شته:
      \begin{oltableau}
        [\sFmla{\False}{\formula{B}}, just=\TAss
          [\sFmla{\True}{\formula{A} \lif \formula{B}}, just=\TAss
            [\sFmla{\True}{\formula{A}}, just=\TAss
              [\sFmla{\False}{\formula{A}}, just = {\TRule{\True}{\lif}[2]}, close]
              [\sFmla{\True}{\formula{B}}, just = {\TRule{\True}{\lif}[2]}, close]
            ]
          ]
        ]
      \end{oltableau}
    \item د~\LR{$\{\sFmla{\False}{!A \lif !B},
      \sFmla{\True}{\lnot !A}\}$} او
      \LR{$\{\sFmla{\False}{!A \lif !B}, \sFmla{\True}{!B}\}$} دواړو
      دپاره تړلي تابلوګانې شته:
      \begin{oltableau}
        [\sFmla{\False}{\formula{A} \lif \formula{B}}, just = \TAss
          [\sFmla{\True}{\lnot \formula{A}}, just = \TAss
            [\sFmla{\True}{\formula{A}}, just = {\TRule{\False}{\lif}[1]}
              [\sFmla{\False}{\formula{B}}, just = {\TRule{\False}{\lif}[1]}
                [\sFmla{\False}{\formula{A}}, just = {\TRule{\True}{\lnot}[2]}, close]
              ]
            ]
          ]
        ]
      \end{oltableau}
      \begin{oltableau}
        [\sFmla{\False}{\formula{A} \lif \formula{B}}, just = \TAss
          [\sFmla{\True}{\formula{B}}, just = \TAss
            [\sFmla{\True}{\formula{A}}, just = {\TRule{\False}{\lif}[1]}
              [\sFmla{\False}{\formula{B}}, just = {\TRule{\False}{\lif}[1]},close
              ]
            ]
          ]
        ]
      \end{oltableau}
  \end{enumerate}
\end{proof}
\olpseganchor{OLP-0107-B013}{end}

\olpseganchor{OLP-0108-B004}{start}
\olfileid{fol}{tab}{qpr}
\olsection{د اشتقاق وړتيا او سورونه}
\olpseganchor{OLP-0108-B004}{end}

\olpseganchor{OLP-0108-B005}{start}
\begin{explain}
  د بشپړتيا قضيه دا هم غواړي چې د تابلو قاعدې د~\LR{$\Proves$} په اړه
  هغه حقيقتونه ورکړي چې په دې برخه کښې ثابتېږي۔
\end{explain}
\olpseganchor{OLP-0108-B005}{end}

\olpseganchor{OLP-0108-B006}{start}
\begin{thm}
\ollabel{thm:strong-generalization} که \LR{$c$} داسې ثابت وي چې په
\LR{$\Gamma$} يا~\LR{$!A(x)$} کښې نۀ راځي، او \LR{$\Gamma \Proves !A(c)$} وي،
نو \LR{$\Gamma \Proves \lforall[x][!A(x)]$}۔
\end{thm}
\olpseganchor{OLP-0108-B006}{end}

\olpseganchor{OLP-0108-B007}{start}
\begin{proof}
فرض کړئ \LR{$\Gamma \Proves !A(c)$}، يعنې داسې
\LR{$!B_1$}، \dots، \LR{$!B_n \in \Gamma$} او د لاندې دپاره تړلے
تابلو شته:
\begin{align*}
\{ \sFmla{\False}{!A(c)}, &
\sFmla{\True}{!B_1}, \dots, \sFmla{\True}{!B_n} \}.
\intertext{\RL{بايد وښيو چې د لاندې دپاره هم تړلے تابلو شته:}}
\{ \sFmla{\False}{\lforall[x][!A(x)]}, &
\sFmla{\True}{!B_1}, \dots, \sFmla{\True}{!B_n} \}.
\end{align*}
تړلے تابلو واخلئ، لومړے فرض ئې په
\LR{$\sFmla{\False}{\lforall[x][!A(x)]}$} بدل کړئ، او د فرضونو وروسته
\LR{$\sFmla{\False}{!A(c)}$} ورننباسئ۔
\begin{center}
\begin{tableau}{not line numbering}
  [\sFmla{\False}{\formula{A}(c)}
    [\sFmla{\True}{\formula{B}_1}
      [\vdots
      [\sFmla{\True}{\formula{B}_n}, 
        [][][]]]]]
\end{tableau}\qquad
\begin{tableau}{not line numbering}
  [\sFmla{\False}{\lforall[x][\formula{A}(x)]}
    [\sFmla{\True}{\formula{B}_1}
      [\vdots
        [\sFmla{\True}{\formula{B}_n},
          [\sFmla{\False}{\formula{A}(c)}
        [][][]]]]]]
\end{tableau}
\end{center}
تابلو لا هم تړلے دے؛ ځکه ټولې هغه تړلي فارمولونه چې مخکښې د فرضونو
په توګه موجودې وې، اوس هم د~تابلو په سر کښې موجودې دي۔
ننويستل شوې کرښه د~\LR{$\TRule{\False}{\lforall}$} قاعدې د سمې کارونې
پايله ده؛ ځکه ثابت سمبول~\LR{$c$} په~\LR{$!B_1$}، \dots، \LR{$!B_n$} يا
\LR{$\lforall[x][!A(x)]$} کښې نۀ راځي، يعنې په نوي تابلو کښې د
ننويستل شوې کرښې پاس نۀ راځي۔
\end{proof}
\olpseganchor{OLP-0108-B007}{end}

\olpseganchor{OLP-0108-B008}{start}
\begin{prop}
\ollabel{prop:provability-quantifiers}
\begin{enumerate}
\item \LR{$!A(t) \Proves \lexists[x][!A(x)]$}.
\item \LR{$\lforall[x][!A(x)] \Proves !A(t)$}.
\end{enumerate}
\end{prop}
\olpseganchor{OLP-0108-B008}{end}

\olpseganchor{OLP-0108-B009}{start}
\begin{proof}
\begin{enumerate}
\item د
  \LR{$\sFmla{\False}{\lexists[x][!A(x)]}, \sFmla{\True}{!A(t)}$} دپاره
  تړلے تابلو دا دے:
  \begin{oltableau}
    [\sFmla{\False}{\lexists[x][\formula{A}(x)]}, just = \TAss
      [\sFmla{\True}{\formula{A}(t)}, just = \TAss
        [\sFmla{\False}{\formula{A}(t)}, just = {\TRule{\False}{\lexists}[1]}, close]]]
    \end{oltableau}
\item د
  \LR{$\sFmla{\False}{\formula{A}(t)}, \sFmla{\True}{\lforall[x][!A(x)]}, $} دپاره
  تړلے تابلو دا دے:
  \begin{oltableau}
    [\sFmla{\False}{\formula{A}(t)}, just = \TAss
      [\sFmla{\True}{\lforall[x][\formula{A}(x)]}, just = \TAss
        [\sFmla{\True}{\formula{A}(t)}, just = {\TRule{\True}{\lforall}[2]}, close]]]
    \end{oltableau}
\end{enumerate}
\end{proof}
\olpseganchor{OLP-0108-B009}{end}

\olpseganchor{OLP-0109-B004}{start}
\olfileid{fol}{tab}{sou}
\olpseganchor{OLP-0109-B004}{end}
      
\olpseganchor{OLP-0109-B005}{start}
\olsection{سموالے}
\olpseganchor{OLP-0109-B005}{end}

\olpseganchor{OLP-0109-B006}{start}
\begin{explain}
د تابلوګانو په شان اشتقاق نظام هله \emph{سم} دے چې
هغه څيزونه اشتقاق نۀ کړي چې په حقيقت کښې نۀ وي۔ نو سموالے د
اشتقاق نظامونو د تضمين شوې خونديتوب يو ډول خاصيت دے۔
د دې له مخې چې کوم ثبوت-تيوريکي خاصيت تر بحث لاندې وي، مثلاً غواړو
پوه شو چې:
\begin{enumerate}
\item هر د اشتقاق وړ~\LR{$!A$} معتبر دے؛
\item که تړلے فارمول له ځينو نورو څخه د اشتقاق وړ وي، د هغو
  معنايي پايله هم ده؛
\item که د تړلي فارمولونه يو سټ ناسازګار وي، د صدق وړ نۀ دے۔
\end{enumerate}
دا د اشتقاق نظام مهم خاصيتونه دي۔ که له دغو څخه کوم يو
سم نۀ وي، اشتقاق نظام نيمګړے دے---له حده زيات څيزونه به
اشتقاق کوي۔ له دې امله د اشتقاق نظام سموالے ثابتول
تر ټولو زيات اهميت لري۔
\olpseganchor{OLP-0109-B006}{end}

\olpseganchor{OLP-0109-B007}{start}
ځکه دغه ټول ثبوت-تيوريکي خاصيتونه د يو يا بل ډول د تړلو
تابلوګانې له لارې تعريف شوي دي، د پورته (۱)--(۳) ثابتول غواړي
چې د تړلو تابلوګانې د معنايي خاصيتونو په اړه يو څه ثابت کړو۔
لومړے به تعريف کړو چې د يو نښه لرونکے فارمول پوره کېدل په يو
جوړښت کښې څه معنٰی لري؛ بيا به وښيو چې که يو تابلو تړلے وي،
هېڅ جوړښت ئې ټول فرضونه نۀ پوره کوي۔ وروسته (۱)--(۳) د دې پايلې
د لازمو نتيجو په توګه راځي۔
\end{explain}
\olpseganchor{OLP-0109-B007}{end}

\olpseganchor{OLP-0109-B008}{start}
\begin{defn}
جوړښت~\LR{$\Struct{M}$}
يو نښه لرونکے فارمول~\LR{$\sFmla{\True}{!A}$} هله او يوازې هله
\emph{پوره کوي} چې \LR{$\Sat{M}{!A}$}، او
\LR{$\sFmla{\False}{!A}$} هله او يوازې هله پوره کوي چې
\LR{$\Sat/{M}{!A}$}۔
\LR{$\Struct{M}$} د نښه لرونکي فارمولونه يو
سټ~\LR{$\Gamma$} هله او يوازې هله پوره کوي چې هر
\LR{$\sFmla{S}{!A} \in \Gamma$} پوره کړي۔ \LR{$\Gamma$} هله \emph{د صدق وړ}
دے چې کوم جوړښت ئې پوره کړي،
او که داسې نۀ وي \emph{د صدق وړ نۀ} دے۔
\end{defn}
\olpseganchor{OLP-0109-B008}{end}

\olpseganchor{OLP-0109-B009}{start}
\begin{thm}[سموالے]
  \ollabel{thm:tableau-soundness}
  که \LR{$\Gamma$} يو تړلے تابلو ولري، نو \LR{$\Gamma$} د صدق وړ نۀ دے۔
\end{thm}
\olpseganchor{OLP-0109-B009}{end}

\olpseganchor{OLP-0109-B010}{start}
\begin{proof}
د يو تابلو څانګه هله او يوازې هله د صدق وړ بولو چې پر هغې
د نښه لرونکي فارمولونه سټ د صدق وړ وي؛ او تابلو هله د صدق وړ
بولو چې لږ تر لږه يوه د صدق وړ څانګه ولري۔
\olpseganchor{OLP-0109-B010}{end}

\olpseganchor{OLP-0109-B011}{start}
دا خبره ښيو: د استنتاج د کومې قاعدې په وسيله د يو د صدق وړ
تابلو پراخول تل يو د صدق وړ تابلو ورکوي۔ دا به قضيه
ثابته کړي: هر تړلے تابلو د~\LR{$\Gamma$} يوازې له فرضونو څخه جوړ
تابلو ته د استنتاج د قاعدو په پلي کولو تر لاسه کېږي۔ نو که
\LR{$\Gamma$} د صدق وړ واے، د هغې هر تابلو به د صدق وړ و۔ خو تړلے
تابلو په څرګند ډول د صدق وړ نۀ دے: هره څانګه دواړه
\LR{$\sFmla{\True}{!A}$} او \LR{$\sFmla{\False}{!A}$} لري، او هېڅ جوړښت
نۀ شي کولے چې \LR{$!A$} هم پوره او هم نۀ پوره کړي۔
\olpseganchor{OLP-0109-B011}{end}

\olpseganchor{OLP-0109-B012}{start}
فرض کړئ يو د صدق وړ تابلو لرو، يعنې داسې تابلو چې لږ
تر لږه يوه د صدق وړ څانګه لري۔ د استنتاج قاعده يا يوې څانګې ته
نښه لرونکي فارمولونه ورزياتوي يا يوه څانګه په دوو وېشي۔ که په
تابلو کښې کومه د صدق وړ څانګه وي چې اړوند قاعدوي تطبيق ئې نۀ
پراخوي، هغه په پراخ شوي تابلو کښې هم د صدق وړ پاتې کېږي؛ نو
پراخ شوے تابلو د صدق وړ دے۔ ځکه يوازې هغه حالت څېړل پکار دي چې
قاعده پر يوې د صدق وړ څانګې پلي کېږي۔
\olpseganchor{OLP-0109-B012}{end}

\olpseganchor{OLP-0109-B013}{start}
\LR{$\Gamma$} پر هغې څانګې د نښه لرونکي فارمولونه سټ وبولئ، او
\LR{$\sFmla{S}{!A} \in \Gamma$} هغه نښه لرونکے فارمول وبولئ چې قاعده
پرې پلي کېږي۔ که قاعده څانګه نۀ وېشي، بايد وښيو چې پراخه شوې څانګه،
يعنې \LR{$\Gamma$} له قاعدوي پايلو سره، لا هم د صدق وړ ده۔ که قاعده
څانګه وېشي، بايد وښيو چې له دوو ترلاسه شوو څانګو لږ تر لږه يوه د
صدق وړ ده۔
\olpseganchor{OLP-0109-B013}{end}
  
\olpseganchor{OLP-0109-B014}{start}
لومړے هغه ممکن استنتاجونه څېړو چې څانګه نۀ وېشي۔
\begin{enumerate}
\item څانګه پر \LR{$\sFmla{\True}{\lnot !B} \in \Gamma$} د
  \LR{$\TRule{\True}{\lnot}$} په پلي کولو پراخېږي۔ بيا پراخه شوې څانګه
  د \LR{$\Gamma \cup \{\sFmla{\False}{!B}\}$} نښه لرونکي فارمولونه لري۔
  فرض کړئ \LR{$\Sat{M}{\Gamma}$}۔ په ځانګړي ډول
  \LR{$\Sat{M}{\lnot !B}$}۔ نو
  \LR{$\Sat/{M}{!B}$}، يعنې
  \LR{$\Struct{M}$}،
  \LR{$\sFmla{\False}{!B}$} پوره کوي۔
\item څانګه پر \LR{$\sFmla{\False}{\lnot !B} \in \Gamma$} د
  \LR{$\TRule{\False}{\lnot}$} په پلي کولو پراخېږي: تمرين۔
\item څانګه پر \LR{$\sFmla{\True}{!B \land !C} \in \Gamma$} د
  \LR{$\TRule{\True}{\land}$} په پلي کولو پراخېږي؛ له دې څخه پر څانګه
  دوه نوي نښه لرونکي فارمولونه ترلاسه کېږي: \LR{$\sFmla{\True}{!B}$} او
  \LR{$\sFmla{\True}{!C}$}۔ فرض کړئ
  \LR{$\Sat{M}{\Gamma}$}، او په ځانګړي ډول
  \LR{$\Sat{M}{!B \land !C}$}۔ بيا
  \LR{$\Sat{M}{!B}$} او
  \LR{$\Sat{M}{!C}$}۔ يعنې
  \LR{$\Struct{M}$} دواړه
  \LR{$\sFmla{\True}{!B}$} او \LR{$\sFmla{\True}{!C}$} پوره کوي۔
\item څانګه پر \LR{$\sFmla{\False}{!B \lor !C} \in \Gamma$} د
  \LR{$\TRule{\False}{\lor}$} په پلي کولو پراخېږي: تمرين۔
\item څانګه پر \LR{$\sFmla{\False}{!B \lif !C} \in \Gamma$} د
  \LR{$\TRule{\False}{\lif}$} په پلي کولو پراخېږي۔ له دې څخه پر څانګه
  دوه نوي نښه لرونکي فارمولونه ترلاسه کېږي: \LR{$\sFmla{\True}{!B}$} او
  \LR{$\sFmla{\False}{!C}$}۔ فرض کړئ
  \LR{$\Sat{M}{\Gamma}$}، او په ځانګړي ډول
  \LR{$\Sat/{M}{!B \lif !C}$}۔ بيا
  \LR{$\Sat{M}{!B}$} او
  \LR{$\Sat/{M}{!C}$}۔ يعنې
  \LR{$\Struct{M}$} دواړه
  \LR{$\sFmla{\True}{!B}$} او \LR{$\sFmla{\False}{!C}$} پوره کوي۔
\olpseganchor{OLP-0109-B014}{end}

\olpseganchor{OLP-0109-B015}{start}
%
\textbf{د ژباړې د نسخې سمون (OLTAB-009):} سرچينې په دواړو
کلي-کميت ټاکونکو قاعدوي حالتونو کښې د قاعدې په مقدمه کښې د بي ميتا-فارمول
ليکلے و، خو پايله او ټول ورپسې معنايي استدلال د اې ميتا-فارمول
کاروي۔ دلته اې په هر حالت کښې يو شان ساتل کېږي۔
\item څانګه پر \LR{$\sFmla{\True}{\lforall[x][!A(x)]} \in \Gamma$} د
  \LR{$\TRule{\True}{\lforall}$} په پلي کولو پراخېږي۔ له دې څخه پر څانګه
  يو نوے نښه لرونکے فارمول~\LR{$\sFmla{\True}{!A(t)}$} ترلاسه کېږي۔
  فرض کړئ \LR{$\Sat{M}{\Gamma}$}، او په ځانګړي ډول
  \LR{$\Sat{M}{\lforall[x][!A(x)]}$}۔ د
  \olref[syn][sem]{prop:quant-terms} له مخې \LR{$\Sat{M}{!A(t)}$}۔ نو
  \LR{$\Struct{M}$}، \LR{$\sFmla{\True}{!A(t)}$} پوره کوي۔
\olpseganchor{OLP-0109-B015}{end}

\olpseganchor{OLP-0109-B016}{start}
\item څانګه پر \LR{$\sFmla{\False}{\lforall[x][!A(x)]} \in \Gamma$} د
  \LR{$\TRule{\False}{\lforall}$} په پلي کولو پراخېږي۔ له دې څخه پر
  څانګه يو نوے نښه لرونکے فارمول~\LR{$\sFmla{\False}{!A(a)}$} ترلاسه
  کېږي، چې \LR{$a$} داسې ثابت سمبول دے چې په~\LR{$\Gamma$} کښې نۀ راځي۔
  ځکه \LR{$\Gamma$} د صدق وړ دے، داسې \LR{$\Struct{M}$} شته چې
  \LR{$\Sat{M}{\Gamma}$}، او په ځانګړي ډول
  \LR{$\Sat/{M}{\lforall[x][!A(x)]}$}۔ بايد وښيو چې
  \LR{$\Gamma \cup \{\sFmla{\False}{!A(a)}\}$} د صدق وړ دے۔ د دې دپاره
  مناسب~\LR{$\Struct{M'}$} داسې تعريفوو۔
\olpseganchor{OLP-0109-B016}{end}

\olpseganchor{OLP-0109-B017}{start}
  د \olref[syn][ass]{prop:sat-quant} له مخې
  \LR{$\Sat/{M}{\lforall[x][!A(x)]}$} هله او يوازې هله چې د کوم~\LR{$s$}
  دپاره \LR{$\Sat/{M}{!A(x)}[s]$}۔ اوس \LR{$\Struct{M'}$} د~\LR{$\Struct{M}$} په
  شان وبولئ، خو \LR{$\Assign{a}{M'} = s(x)$}۔ د
  \olref[syn][ext]{cor:extensionality-sent} له مخې، د هر
  \LR{$\sFmla{\True}{!C} \in \Gamma$} دپاره \LR{$\Sat{M'}{!C}$}، او د هر
  \LR{$\sFmla{\False}{!C} \in \Gamma$} دپاره \LR{$\Sat/{M'}{!C}$}؛ ځکه
  \LR{$a$} په~\LR{$\Gamma$} کښې نۀ راځي۔
\olpseganchor{OLP-0109-B017}{end}

\olpseganchor{OLP-0109-B018}{start}
  د \olref[syn][ext]{prop:extensionality} له مخې
  \LR{$\Sat/{M'}{!A(x)}[s]$}۔ د \olref[syn][ext]{prop:ext-formulas} له مخې
  \LR{$\Sat/{M'}{!A(a)}[s]$}۔ ځکه \LR{$!A(a)$} يوه تړلے فارمول ده، د
  \olref[syn][ass]{prop:sentence-sat-true} له مخې
  \LR{$\Sat/{M'}{!A(a)}$}؛ يعنې \LR{$\Struct{M'}$}،
  \LR{$\sFmla{\False}{!A(a)}$} پوره کوي۔
\item څانګه پر \LR{$\sFmla{\True}{\lexists[x][!B(x)]} \in \Gamma$} د
  \LR{$\TRule{\True}{\lexists}$} په پلي کولو پراخېږي: تمرين۔
\item څانګه پر \LR{$\sFmla{\False}{\lexists[x][!B(x)]} \in \Gamma$} د
  \LR{$\TRule{\False}{\lexists}$} په پلي کولو پراخېږي: تمرين۔
\end{enumerate}
اوس هغه ممکن استنتاجونه څېړو چې څانګه وېشي۔
\begin{enumerate}
\item څانګه پر \LR{$\sFmla{\False}{!B \land !C} \in \Gamma$} د
  \LR{$\TRule{\False}{\land}$} په پلي کولو پراخېږي؛ له دې څخه دوه څانګې
  ترلاسه کېږي: کيڼه څانګه د~\LR{$\sFmla{\False}{!B}$} له لارې او ښۍ
  څانګه د~\LR{$\sFmla{\False}{!C}$} له لارې دوام کوي۔ فرض کړئ
  \LR{$\Sat{M}{\Gamma}$}، او په ځانګړي ډول
  \LR{$\Sat/{M}{!B \land !C}$}۔ بيا يا
  \LR{$\Sat/{M}{!B}$} يا
  \LR{$\Sat/{M}{!C}$}۔ په لومړي حالت کښې
  \LR{$\Struct{M}$}،
  \LR{$\sFmla{\False}{!B}$} پوره کوي؛ يعنې
  \LR{$\Struct{M}$} د کيڼې څانګې فارمولونه
  پوره کوي۔ په دويم حالت کښې
  \LR{$\Struct{M}$}،
  \LR{$\sFmla{\False}{!C}$} پوره کوي؛ يعنې
  \LR{$\Struct{M}$} د ښۍ څانګې فارمولونه
  پوره کوي۔
\item څانګه پر \LR{$\sFmla{\True}{!B \lor !C} \in \Gamma$} د
  \LR{$\TRule{\True}{\lor}$} په پلي کولو پراخېږي: تمرين۔
\item څانګه پر \LR{$\sFmla{\True}{!B \lif !C} \in \Gamma$} د
  \LR{$\TRule{\True}{\lif}$} په پلي کولو پراخېږي: تمرين۔
\item څانګه د~\Cut{} په وسيله پراخېږي۔ له دې څخه دوه څانګې ترلاسه
  کېږي: يوه \LR{$\sFmla{\True}{!B}$} او بله \LR{$\sFmla{\False}{!B}$} لري۔
  ځکه \LR{$\Sat{M}{\Gamma}$} او يا
  \LR{$\Sat{M}{!B}$} يا
  \LR{$\Sat/{M}{!B}$}،
  \LR{$\Struct{M}$} يا کيڼه يا ښۍ څانګه پوره
  کوي۔
\end{enumerate}
\end{proof}
\olpseganchor{OLP-0109-B018}{end}

\olpseganchor{OLP-0109-B019}{start}
\begin{prob}
د~\olref[fol][tab][sou]{thm:tableau-soundness} ثبوت بشپړ کړئ۔
\end{prob}
\olpseganchor{OLP-0109-B019}{end}



\olpseganchor{OLP-0109-B021}{start}
\begin{cor}
\ollabel{cor:weak-soundness}
که \LR{$\Proves !A$} وي، نو \LR{$!A$} معتبر دے۔
\end{cor}
\olpseganchor{OLP-0109-B021}{end}

\olpseganchor{OLP-0109-B022}{start}
\begin{cor}
\ollabel{cor:entailment-soundness}
که \LR{$\Gamma \Proves !A$} وي، نو \LR{$\Gamma \Entails !A$}۔
\end{cor}
\olpseganchor{OLP-0109-B022}{end}

\olpseganchor{OLP-0109-B023}{start}
\begin{proof}
  که \LR{$\Gamma \Proves !A$} وي، نو د ځينو
  \LR{$!B_1$}، \dots، \LR{$!B_n \in \Gamma$} دپاره
  \LR{$\{\sFmla{\False}{!A}, \sFmla{\True}{!B_1}, \dots,
  \sFmla{\True}{!B_n}\}$} يو تړلے تابلو لري۔ د
  \olref{thm:tableau-soundness} له مخې، هر
  جوړښت~\LR{$\Struct{M}$}
  يا کوم \LR{$!B_i$} ناسموي يا \LR{$!A$} رښتيا کوي۔ نو که
  \LR{$\Sat{M}{\Gamma}$} وي، بيا
  \LR{$\Sat{M}{!A}$} هم دے۔
\end{proof}
\olpseganchor{OLP-0109-B023}{end}

\olpseganchor{OLP-0109-B024}{start}
\begin{cor}
\ollabel{cor:consistency-soundness}
که \LR{$\Gamma$} د صدق وړ وي، نو سازګار دے۔
\end{cor}
\olpseganchor{OLP-0109-B024}{end}

\olpseganchor{OLP-0109-B025}{start}
\begin{proof}
مقابل عکس ثابتوو۔ فرض کړئ \LR{$\Gamma$} سازګار نۀ دے۔ بيا
\LR{$!B_1$}، \dots، \LR{$!B_n \in \Gamma$} شته او د
\LR{$\{\sFmla{\True}{!B_1}, \dots, \sFmla{\True}{!B_n}\}$} دپاره يو
تړلے تابلو شته۔ د \olref{thm:tableau-soundness} له مخې، داسې
\LR{$\Struct{M}$} نشته چې د ټولو
\LR{$i=1$}، \dots،~\LR{$n$} دپاره \LR{$\Sat{M}{!B_i}$} وي۔
خو بيا \LR{$\Gamma$} د صدق وړ نۀ دے۔
\end{proof}
\olpseganchor{OLP-0109-B025}{end}

\olpseganchor{OLP-0110-B004}{start}
\olfileid{fol}{tab}{ide}
\olpseganchor{OLP-0110-B004}{end}

\olpseganchor{OLP-0110-B005}{start}
\olsection{له عينيت سره تابلوګانې}
\olpseganchor{OLP-0110-B005}{end}

\olpseganchor{OLP-0110-B006}{start}
له عينيت سره تابلوګانې د استنتاج اضافي قاعدې غواړي۔
د~\LR{$\eq$} قاعدې دا دي (\LR{$t$}، \LR{$t_1$} او \LR{$t_2$} تړلي ترمونه دي):
\olpseganchor{OLP-0110-B006}{end}

\olpseganchor{OLP-0110-B007}{start}
\begin{defish}
\AxiomC{}
\RightLabel{\LR{$\eq$}}
\UnaryInfC{\sFmla{\True}{\eq[t][t]}}
\DisplayProof
\hfill
\AxiomC{\sFmla{\True}{\eq[t_1][t_2]}}
\noLine
\UnaryInfC{\sFmla{\True}{!A(t_1)}}
\RightLabel{\LR{$\TRule{\True}{\eq}$}}
\UnaryInfC{\sFmla{\True}{!A(t_2)}}
\DisplayProof
\hfill
\AxiomC{\sFmla{\True}{\eq[t_1][t_2]}}
\noLine
\UnaryInfC{\sFmla{\False}{!A(t_1)}}
\RightLabel{\LR{$\TRule{\False}{\eq}$}}
\UnaryInfC{\sFmla{\False}{!A(t_2)}}
\DisplayProof
\end{defish}
پام وکړئ چې له نورو ټولو قاعدو سره په توپير کښې،
\LR{$\TRule{\True}{\eq}$} او \LR{$\TRule{\False}{\eq}$} غواړي چې \emph{دوه}
نښه لرونکي فارمولونه له وړاندې پر څانګه وي؛ يعنې دواړه
\LR{$\sFmla{\True}{\eq[t_1][t_2]}$} او \LR{$\sFmla{S}{!A(t_1)}$}۔
\olpseganchor{OLP-0110-B007}{end}

\olpseganchor{OLP-0110-B008}{start}
\begin{ex}
که \LR{$s$} او \LR{$t$} تړلي ترمونه وي، نو
\LR{$\eq[s][t], !A(s) \Proves !A(t)$}:
\begin{oltableau}
  [\sFmla{\False}{\formula{A}(t)}, just = \TAss
    [\sFmla{\True}{\eq[s][t]}, just = \TAss
      [\sFmla{\True}{\formula{A}(s)}, just = \TAss
        [\sFmla{\True}{\formula{A}(t)}, just={\TRule{\True}{\eq}[2, 3]}, close]
      ]
    ]
  ]
\end{oltableau}
کېدے شي دا د يو شان څيزونو د تعويض د اصل يا د لايبنېڅ د قانون په
نوم درته اشنا وي۔
\olpseganchor{OLP-0110-B008}{end}

\olpseganchor{OLP-0110-B009}{start}
تابلوګانې ثابتوي چې \LR{$\eq$} متناظره ده، يعنې
\LR{$\eq[s_1][s_2] \Proves \eq[s_2][s_1]$}:
\begin{oltableau}
  [\sFmla{\False}{\eq[s_2][s_1]}, just = \TAss
    [\sFmla{\True}{\eq[s_1][s_2]}, just = \TAss
      [\sFmla{\True}{\eq[s_1][s_1]}, just = {\LR{$\eq$}}
        [\sFmla{\True}{\eq[s_2][s_1]}, just = {\TRule{\True}{\eq}[2, 3]}, close]
      ]
    ]
  ]
\end{oltableau}
دلته کرښه~\LR{$2$} د~\LR{$\TRule{\True}{\eq}$} لومړنے اړين
فارمول~\LR{$\sFmla{\True}{\eq[s_1][s_2]}$} دے۔ کرښه~\LR{$3$} دويم
اړين فارمول دے، چې د~\LR{$\sFmla{\True}{!A(s_1)}$} بڼه لري---\LR{$!A(x)$}
د~\LR{$\eq[x][s_1]$} په توګه وګڼئ؛ بيا \LR{$!A(s_1)$}،
\LR{$\eq[s_1][s_1]$} او \LR{$!A(s_2)$}، \LR{$\eq[s_2][s_1]$} کېږي۔
\textbf{د ژباړې د نسخې سمون (OLTAB-010):} سرچينې کرښه درې د اې
د دويم ترم بڼه بللې وه، خو د قاعدې په دې تطبيق کښې هغه د اې د
لومړي ترم بڼه ده، لکه ورپسې تشريح چې هم ښيي۔
\olpseganchor{OLP-0110-B009}{end}

\olpseganchor{OLP-0110-B010}{start}
تابلوګان دا هم ثابتوي چې \LR{$\eq$} انتقالي ده، يعنې
\LR{$\eq[s_1][s_2], \eq[s_2][s_3] \Proves \eq[s_1][s_3]$}:
\begin{oltableau}
  [\sFmla{\False}{\eq[s_1][s_3]}, just = \TAss
    [\sFmla{\True}{\eq[s_1][s_2]}, just = \TAss
      [\sFmla{\True}{\eq[s_2][s_3]}, just = \TAss
        [\sFmla{\True}{\eq[s_1][s_3]}, just = {\TRule{\True}{\eq}[3, 2]}, close]
      ]
    ]
  ]
\end{oltableau}
په دې تابلو کښې د~\LR{$\TRule{\True}{\eq}$} لومړنے اړين
فارمول کرښه~\LR{$3$}، \LR{$\sFmla{\True}{\eq[s_2][s_3]}$} ده
(\LR{$s_2$} د~\LR{$t_1$} او \LR{$s_3$} د~\LR{$t_2$} ونډه لري)۔ دويم اړين فارمول چې
د~\LR{$\sFmla{\True}{!A(s_2)}$} بڼه لري، کرښه~\LR{$2$} ده۔ دلته \LR{$!A(x)$}
د~\LR{$\eq[s_1][x]$} په توګه وګڼئ؛ په دې سره \LR{$!A(s_2)$}،
\LR{$\eq[s_1][s_2]$} (يعنې کرښه~\LR{$2$}) او \LR{$!A(s_3)$} په پايله کښې
فارمول~\LR{$\eq[s_1][s_3]$} کېږي۔
\textbf{د ژباړې د نسخې سمون (OLTAB-011):} سرچينې د اې د دويم
ترم د بڼې د پراخولو پر ځاے د قاعدې عمومي برابري ليکلې وه؛ خو
د همدلته ټاکل شوي اې له مخې اړوند فارمول د لومړي او دويم ترم
برابري ده۔
\end{ex}
\olpseganchor{OLP-0110-B010}{end}

\olpseganchor{OLP-0110-B011}{start}
\begin{prob}
د لاندې دپاره تړلي تابلوګانې ورکړئ:
\begin{enumerate}
\item \LR{$\sFmla{\False}{\lforall[x][\lforall[y][((x = y \land !A(x))
      \lif !A(y))]]}$}
\item \LR{$\sFmla{\False}{\lexists[x][(!A(x) \land
    \lforall[y][(!A(y) \lif y = x)])]}$},\\
  \LR{$\sFmla{\True}{\lexists[x][!A(x)] \land
    \lforall[y][\lforall[z][((!A(y) \land !A(z)) \lif y = z)]]}$}
\end{enumerate}
\end{prob}
\olpseganchor{OLP-0110-B011}{end}

\olpseganchor{OLP-0111-B004}{start}
\olfileid{fol}{tab}{sid}
\olpseganchor{OLP-0111-B004}{end}

\olpseganchor{OLP-0111-B005}{start}
\olsection{له عينيت سره سموالے}
\olpseganchor{OLP-0111-B005}{end}

\olpseganchor{OLP-0111-B006}{start}
\begin{prop}
د عينيت له قاعدو سره تابلوګانې سم دي: هېڅ تړلے تابلو
د صدق وړ نۀ دے۔
\end{prop}
\olpseganchor{OLP-0111-B006}{end}

\olpseganchor{OLP-0111-B007}{start}
\begin{proof}
يوازې بايد د پخوا په شان وښيو چې که يو تابلو د صدق وړ
څانګه ولري، نو د~\LR{$\eq$} له قاعدو څخه د يوې په پلي کولو ترلاسه شوې
څانګه هم د صدق وړ ده۔ \LR{$\Gamma$} پر څانګه د نښه لرونکو فارمولونه
سټ وبولئ، او \LR{$\Struct{M}$} داسې جوړښت وبولئ چې \LR{$\Gamma$}
پوره کوي۔
\olpseganchor{OLP-0111-B007}{end}

\olpseganchor{OLP-0111-B008}{start}
فرض کړئ څانګه د~\LR{$\eq$} په کارولو، يعنې د نښه لرونکي
فارمول~\LR{$\sFmla{\True}{\eq[t][t]}$} په زياتولو پراخېږي۔ په
څرګند ډول \LR{$\Sat{M}{\eq[t][t]}$}؛ نو \LR{$\Struct{M}$}،
\LR{$\Gamma \cup \{\sFmla{\True}{\eq[t][t]}\}$} هم پوره کوي۔
\olpseganchor{OLP-0111-B008}{end}

\olpseganchor{OLP-0111-B009}{start}
که څانګه د~\LR{$\TRule{\True}{\eq}$} په کارولو پراخېږي، موږ يو نښه
لرونکے فارمول~\LR{$\sFmla{\True}{!A(t_2)}$} ورزياتوو، خو \LR{$\Gamma$}
دواړه~\LR{$\sFmla{\True}{\eq[t_1][t_2]}$} او
\LR{$\sFmla{\True}{!A(t_1)}$} لري۔
\textbf{د ژباړې د نسخې سمون (OLTAB-012):} سرچينې د رښتوني-عينيت
قاعدې پايلې ته عمومي S نښه ورکړې وه؛ د قاعدې، دواړو مقدمو او
لاندې معنايي ثبوت له مخې دا رښتونې نښه ده۔
نو \LR{$\Sat{M}{\eq[t_1][t_2]}$} او \LR{$\Sat{M}{!A(t_1)}$} لرو۔ داسې
متغير-ټاکنه~\LR{$s$} واخلئ چې \LR{$s(x) = \Value{t_1}{M}$}۔ د
\olref[syn][ass]{prop:sentence-sat-true} له مخې
\LR{$\Sat{M}{!A(t_1)}[s]$}۔ ځکه \LR{$\varAssign{s}{s}{x}$}، د
\olref[syn][ext]{prop:ext-formulas} له مخې \LR{$\Sat{M}{!A(x)}[s]$}۔
ځکه \LR{$\Sat{M}{\eq[t_1][t_2]}$}، نو
\LR{$\Value{t_1}{M} = \Value{t_2}{M}$}، او ځکه
\LR{$s(x) = \Value{t_2}{M}$}۔ د
\olref[syn][ext]{prop:ext-formulas} په بيا پلي کولو
\LR{$\Sat{M}{!A(t_2)}[s]$} هم لرو۔ د
\olref[syn][ass]{prop:sentence-sat-true} له مخې
\LR{$\Sat{M}{!A(t_2)}$}۔ د~\LR{$\TRule{\False}{\eq}$} حالت هم په همدې ډول
حل کېږي۔
\end{proof}
\olpseganchor{OLP-0111-B009}{end}

\olpseganchor{OLP-0112-B004}{start}
\olchapter{fol}{axd}{بديهي اشتقاقونه}
\olpseganchor{OLP-0112-B004}{end}

\olpseganchor{OLP-0112-B005}{start}
\begin{editorial}
  لا تر اوسه دا هڅه نۀ ده شوې چې ډاډ ترلاسه شي چې د دې څپرکي مواد
  د هغو بېلابېلو ټاګونو درناوے کوي چې ښيي کوم نښلوونکي او کميت
  ټاکونکي بنسټيز يا تعريف شوي دي: له~\LR{$\liff$} پرته، چې تعريف شوے
  ګڼل کېږي، ټول بنسټيز ګڼل شوي دي۔ که د لومړۍ درجې منطق ټاګ رښتونے
  وي، له کميت ټاکونکو سره بڼه جوړوو؛ که نۀ وي، بې له هغو بڼه جوړوو۔
\end{editorial}
\olpseganchor{OLP-0112-B005}{end}





\olpseganchor{OLP-0112-B008}{start}
%
  

\olpseganchor{OLP-0112-B008}{end}



\olpseganchor{OLP-0112-B010}{start}
%
  

\olpseganchor{OLP-0112-B010}{end}





\olpseganchor{OLP-0112-B013}{start}
%
  

\olpseganchor{OLP-0112-B013}{end}





\olpseganchor{OLP-0112-B016}{start}
%
  

\olpseganchor{OLP-0112-B016}{end}

\olpseganchor{OLP-0113-B004}{start}
\olfileid{fol}{axd}{rul}
\olpseganchor{OLP-0113-B004}{end}

\olpseganchor{OLP-0113-B005}{start}
\olsection{قاعدې او اشتقاقونه}
\olpseganchor{OLP-0113-B005}{end}

\olpseganchor{OLP-0113-B006}{start}
\begin{explain}
  بديهي اشتقاقونه ښايي د منطق د اشتقاق تر ټولو ساده
  نظام وي۔ اشتقاق يوازې د فارمولونه يوه لړۍ ده۔ د
  اشتقاق ګڼل کېدو دپاره، په لړۍ کښې هر فارمول بايد
  يا د کوم بديهي اصل يوه بېلګه وي، يا د استنتاج د کومې قاعدې له مخې
  له يوه يا زياتو هغو فارمولونه څخه راشي چې په لړۍ کښې ترې وړاندې
  دي۔ اشتقاق خپل وروستے فارمول، اشتقاق کوي۔
\end{explain}
\olpseganchor{OLP-0113-B006}{end}

\olpseganchor{OLP-0113-B007}{start}
\begin{defn}[د اشتقاق وړتيا]
که \LR{$\Gamma$} د~\LR{$\Lang L$} د فارمولونه يو سټ وي، نو له~\LR{$\Gamma$}
څخه يو \emph{اشتقاق} د فارمولونه يوه متناهي لړۍ
\LR{$!A_1$}، \dots،~\LR{$!A_n$} ده، چې د هر \LR{$i \le n$} دپاره له لاندې څخه
يو شرط پوره کېږي:
\begin{enumerate}
\item \LR{$!A_i \in \Gamma$}؛ يا
\item \LR{$!A_i$} يو بديهي اصل دے؛ يا
\item \LR{$!A_i$} د استنتاج د کومې قاعدې له مخې له کوم \LR{$!A_j$} (او
  \LR{$!A_k$}) څخه راځي، چې \LR{$j < i$} (او \LR{$k < i$})۔
\end{enumerate}
\end{defn}
\olpseganchor{OLP-0113-B007}{end}

\olpseganchor{OLP-0113-B008}{start}
دا چې سم اشتقاق څه ګڼل کېږي، په دې پورې اړه لري چې موږ
کومې استنتاجي قاعدې منو (او البته کوم څيزونه بديهي اصول ګڼو)۔
استنتاجي قاعده يوه که-نو وينا ده چې راته وايي د ځينو شرطونو لاندې
په اشتقاق کښې يو ګام~\LR{$!A_i$} سم استنتاجي ګام دے۔
\textbf{د ژباړې د نسخې سمون (OLAX-001):} سرچينې په دې يو ځای کښې
د اې-آی د ميتا-فارمول نښه پرېښې وه؛ تعريف او د همدې پاراګراف نورې
ټولې کارونې هغه د فارمول په توګه ليکي، نو نښه دلته بېرته راوستل شوه۔
\olpseganchor{OLP-0113-B008}{end}

\olpseganchor{OLP-0113-B009}{start}
\begin{defn}[د استنتاج قاعده]
د \emph{استنتاج قاعده} يو بس شرط ورکوي چې له~\LR{$\Gamma$} څخه په
اشتقاق کښې څه شی سم استنتاجي ګام ګڼل کېږي۔
\end{defn}
\olpseganchor{OLP-0113-B009}{end}

\olpseganchor{OLP-0113-B010}{start}
د بېلګې په توګه، ځکه چې هره يو-غړيزه لړۍ~\LR{$!A$} چې
\LR{$!A \in \Gamma$} وي، په څرګند ډول اشتقاق ګڼل کېږي، لاندې
ډېره ساده استنتاجي قاعده کېدے شي:
\begin{quote}
  که \LR{$!A \in \Gamma$} وي، نو \LR{$!A$} له~\LR{$\Gamma$} څخه په هر
  اشتقاق کښې تل سم استنتاجي ګام دے۔
\end{quote}
همدارنګه که \LR{$!A$} له بديهي اصولو څخه يو وي، نو يوازې \LR{$!A$} هم
اشتقاق دے؛ ځکه دا هم يوه استنتاجي قاعده ده:
\begin{quote}
  که \LR{$!A$} بديهي اصل وي، نو \LR{$!A$} سم استنتاجي ګام دے۔
\end{quote}
خبره هغه وخت لا په زړه پورې شي چې استنتاجي قاعده هغو فارمولونه
ته رجوع کوي چې تر څېړل شوي ګام وړاندې راغلي وي۔ لاندې قاعده
\emph{مودوس پوننس} بلل کېږي:
\begin{quote}
  که \LR{$!B \lif !A$} او \LR{$!B$} په اشتقاق کښې وړاندې راغلي وي،
  نو~\LR{$!A$} سم استنتاجي ګام دے۔
\end{quote}
که همدا يوازينۍ استنتاجي قاعده وي، نو د اشتقاق پورته تعريف
مو داسې کېږي: \LR{$!A_1$}، \dots،~\LR{$!A_n$} هله او يوازې هله
اشتقاق دے چې د هر \LR{$i \le n$} دپاره له لاندې څخه يو شرط
پوره کېږي:
\begin{enumerate}
\item \LR{$!A_i \in \Gamma$}؛ يا
\item \LR{$!A_i$} يو بديهي اصل دے؛ يا
\item د کوم \LR{$j < i$} دپاره \LR{$!A_j$}، \LR{$!B \lif !A_i$} وي، او د کوم
  \LR{$k < i$} دپاره \LR{$!A_k$}،~\LR{$!B$} وي۔
\end{enumerate}
وروستے بند وايي چې \LR{$!A_i$} له~\LR{$!A_j$} (\LR{$!B \lif !A_i$}) او
\LR{$!A_k$} (\LR{$!B$}) څخه د مودوس پوننس په وسيله راځي۔ که له~\LR{$1$} څخه
تر~\LR{$n$} پورې تېرېدے شو، او هر ځل داسې فارمول~\LR{$!A_i$} ومومو
چې يا په~\LR{$\Gamma$} کښې وي، يا بديهي اصل وي، يا استنتاجي قاعده راته
وايي چې سم استنتاجي ګام دے، نو ټوله لړۍ سم اشتقاق ګڼل کېږي۔
\olpseganchor{OLP-0113-B010}{end}

\olpseganchor{OLP-0113-B011}{start}
\begin{defn}[د اشتقاق وړتيا]
يو فارمول~\LR{$!A$} هله له~\LR{$\Gamma$} څخه \emph{د اشتقاق وړ} دے،
چې \LR{$\Gamma \Proves !A$} ليکو، چې له~\LR{$\Gamma$} څخه داسې
اشتقاق وي چې په~\LR{$!A$} پاے ته رسېږي۔
\end{defn}
\olpseganchor{OLP-0113-B011}{end}

\olpseganchor{OLP-0113-B012}{start}
\begin{defn}[قضيې]
فارمول~\LR{$!A$} هله \emph{قضيه} ده چې له تش سټ څخه د~\LR{$!A$} يو
اشتقاق وي۔ که \LR{$!A$} قضيه وي \LR{$\Proves !A$}، او که نۀ وي
\LR{$\Proves/ !A$} ليکو۔
\end{defn}
\olpseganchor{OLP-0113-B012}{end}

\olpseganchor{OLP-0114-B004}{start}
\olfileid{fol}{axd}{prp}
\olpseganchor{OLP-0114-B004}{end}

\olpseganchor{OLP-0114-B005}{start}
\olsection{د قضيې نښلوونکو دپاره بديهي اصول او قاعدې}
\olpseganchor{OLP-0114-B005}{end}

\olpseganchor{OLP-0114-B006}{start}
\begin{defn}[بديهي اصول]
د قضيې نښلوونکو د بديهي اصولو سټ~\LR{$\PAx$} د لاندې ټولو بڼو
فارمولونه لري:
\begin{align}
  & (!A \land !B) \lif !A \ollabel{ax:land1}\\
  & (!A \land !B) \lif !B \ollabel{ax:land2}\\
  & !A \lif (!B \lif (!A \land !B)) \ollabel{ax:land3}\\
  & !A \lif (!A \lor !B) \ollabel{ax:lor1}\\
  & !A \lif (!B \lor !A) \ollabel{ax:lor2}\\
  & (!A \lif !C) \lif ((!B \lif !C) \lif ((!A \lor !B) \lif !C)) \ollabel{ax:lor3}\\
  & !A \lif (!B \lif !A) \ollabel{ax:lif1}\\
  & (!A \lif (!B \lif !C)) \lif ((!A \lif !B) \lif (!A \lif !C)) \ollabel{ax:lif2}\\
  & (!A \lif !B) \lif ((!A \lif \lnot !B) \lif \lnot !A) \ollabel{ax:lnot1}\\
  & \lnot !A \lif (!A \lif !B) \ollabel{ax:lnot2}\\
  & \ltrue \ollabel{ax:ltrue}\\
  & \lfalse \lif !A \ollabel{ax:lfalse1}\\
  & (!A \lif \lfalse) \lif \lnot !A \ollabel{ax:lfalse2}\\
  & \lnot\lnot !A \lif !A \ollabel{ax:dne}
\end{align}
\end{defn}
\olpseganchor{OLP-0114-B006}{end}

\olpseganchor{OLP-0114-B007}{start}
\begin{defn}[مودوس پوننس]
  که \LR{$!B$} او \LR{$!B \lif !A$} له وړاندې په اشتقاق کښې راغلي
  وي، نو \LR{$!A$} سم استنتاجي ګام دے۔
\end{defn}
\olpseganchor{OLP-0114-B007}{end}

\olpseganchor{OLP-0114-B008}{start}
د مودوس پوننس قاعده به په ``\MP'' لنډوو۔
\olpseganchor{OLP-0114-B008}{end}

\olpseganchor{OLP-0115-B004}{start}
\olfileid{fol}{axd}{qua}
\olpseganchor{OLP-0115-B004}{end}

\olpseganchor{OLP-0115-B005}{start}
\olsection{د کميت ټاکونکو دپاره بديهي اصول او قاعدې}
\olpseganchor{OLP-0115-B005}{end}

\olpseganchor{OLP-0115-B006}{start}
\begin{defn}[د کميت ټاکونکو بديهي اصول]
پر کميت ټاکونکو واکمن \emph{بديهي اصول} د لاندې ټولو بڼو بېلګې دي:
\begin{align}
\ollabel{ax:q1} & \lforall[x][!B] \lif !B(t), \\
\ollabel{ax:q2} & !B(t) \lif \lexists[x][!B].
\end{align}
د هر تړلي ترم~\LR{$t$} دپاره۔
\end{defn}
\olpseganchor{OLP-0115-B006}{end}

\olpseganchor{OLP-0115-B007}{start}
\begin{defn}[د کميت ټاکونکو قاعدې]
\item که \LR{$!B \lif !A(a)$} له وړاندې په اشتقاق کښې راغلے وي
  او \LR{$a$} په~\LR{$\Gamma$} يا~\LR{$!B$} کښې نۀ راځي، نو
  \LR{$!B \lif \lforall[x][!A(x)]$} سم استنتاجي ګام دے۔
\item که \LR{$!A(a) \lif !B$} له وړاندې په اشتقاق کښې راغلے وي
  او \LR{$a$} په~\LR{$\Gamma$} يا~\LR{$!B$} کښې نۀ راځي، نو
  \LR{$\lexists[x][!A(x)] \lif !B$} سم استنتاجي ګام دے۔
\end{defn}
\olpseganchor{OLP-0115-B007}{end}

\olpseganchor{OLP-0115-B008}{start}
له دغو دواړو هره يوه به په ``\QR'' لنډوو۔
\olpseganchor{OLP-0115-B008}{end}

\olpseganchor{OLP-0116-B004}{start}
\olfileid{fol}{axd}{pro}
\olpseganchor{OLP-0116-B004}{end}

\olpseganchor{OLP-0116-B005}{start}
\olsection{د اشتقاقونه بېلګې}
\olpseganchor{OLP-0116-B005}{end}


\olpseganchor{OLP-0116-B006}{start}
\begin{ex}
فرض کړئ غواړو \LR{$(\lnot !D \lor !E) \lif (!D \lif !E)$} ثابت کړو۔
دا په څرګند ډول زموږ د هېڅ بديهي اصل بېلګه نۀ ده، نو د هغې د
اشتقاق کولو دپاره د~\MP{} قاعده کارول پکار دي۔ زموږ يوازينۍ
قاعده مودوس پوننس ده، چې \LR{$!A$} او \LR{$!A \lif !B$} راکړل شي، د~\LR{$!B$}
توجيه راکوي۔ يوه تګلاره دا ده چې \olref[prp]{ax:lor3} داسې وکاروو
چې \LR{$!A$}، \LR{$\lnot !D$}؛ \LR{$!B$}، \LR{$!E$}؛ او \LR{$!C$}، \LR{$!D \lif !E$} وي؛ يعنې
دا بېلګه:
\[
(\lnot !D \lif (!D \lif !E)) \lif ((!E \lif (!D \lif !E)) \lif ((\lnot
!D \lor !E) \lif (!D \lif !E))).
\]
ولې؟ د مودوس پوننس دوه تطبيقونه وروستۍ برخه راکوي، او همدا مو
غوښتې ده۔ په اسانۍ وينو چې \LR{$\lnot !D \lif (!D \lif !E)$} د
\olref[prp]{ax:lnot2}، او \LR{$!E \lif (!D \lif !E)$} د
\olref[prp]{ax:lif1} بېلګه ده۔ نو زموږ اشتقاق دا دے:
\begin{derivation}
  1. &  \LR{$\lnot !D \lif (!D \lif !E)$} & \olref[prp]{ax:lnot2} \\
  2. & \LR{$(\lnot !D \lif (!D \lif !E)) \lif {}$}\\
  &\qquad \LR{$((!E \lif (!D \lif !E)) \lif ((\lnot !D \lor !E) \lif (!D \lif !E)))$} & \olref[prp]{ax:lor3}\\
  3. & \LR{$(!E \lif (!D \lif !E)) \lif ((\lnot !D \lor !E) \lif (!D \lif !E))$} & 1, 2, \MP\\
  4. & \LR{$!E \lif (!D \lif !E)$} & \olref[prp]{ax:lif1}\\
  5. & \LR{$(\lnot !D \lor !E) \lif (!D \lif !E)$} & 3, 4, \MP
\end{derivation}
\end{ex}
\olpseganchor{OLP-0116-B006}{end}

\olpseganchor{OLP-0116-B007}{start}
\begin{ex}\ollabel{ex:identity}
راځئ د~\LR{$!D \lif !D$} يو اشتقاق ولټوو۔ دا د بديهي اصل بېلګه
نۀ ده، نو په~\MP{} ئې اشتقاق کول پکار دي۔
\olref[prp]{ax:lif1} د~\LR{$!A \lif !B$} په بڼه يو بديهي اصل دے چې
\MP{} پرې پلي کېدے شي۔ د ګټې دپاره، البته، هغه~\LR{$!B$} چې \MP{} ئې په
دې حالت کښې د سم ګام په توګه توجيه کوي بايد~\LR{$!D \lif !D$} وي؛ ځکه
همدا مو غوښتل چې اشتقاق ئې کړو۔ په دې سره \LR{$!A$} هم بايد \LR{$!D$} وي؛
يعنې د \olref[prp]{ax:lif1} دې بېلګې ته کتلے شو:
\[
!D \lif (!D \lif !D)
\]
د~\MP{} د پلي کولو دپاره به د دويمې اړوندې مقدمې، يعنې~\LR{$!A$}،
توجيه هم پکار وي۔ خو زموږ په حالت کښې هغه~\LR{$!D$} دے، او \LR{$!D$} په
يوازې ځان اشتقاق کولے نۀ شو۔ نو بله تګلاره پکار ده۔
\olpseganchor{OLP-0116-B007}{end}

\olpseganchor{OLP-0116-B008}{start}
بل بديهي اصل چې يوازې~\LR{$\lif$} لري \olref[prp]{ax:lif2} دے، يعنې
\[
(!A \lif (!B \lif !C)) \lif ((!A \lif !B) \lif (!A \lif !C))
\]
د~\MP{} په دوه ځله پلي کولو وروستي دننني شرطي ته رسېدے شو۔ بيا دا
معنٰی لري چې د \olref[prp]{ax:lif2} داسې بېلګه غواړو چې
\LR{$!A \lif !C$} پکې \LR{$!D \lif !D$}، يعنې زموږ موخه فارمول وي۔
نو \LR{$!A$} او \LR{$!C$} دواړه~\LR{$!D$} دي۔ \LR{$!B$} څنګه وټاکو چې دواړه
\LR{$!A \lif (!B \lif !C)$} او \LR{$!A \lif !B$}، يعنې زموږ په حالت کښې
\LR{$!D \lif (!B \lif !D)$} او \LR{$!D \lif !B$}، هم د اشتقاق وړ وي؟
له دغو لومړنے، هر~\LR{$!B$} چې وټاکو، لا د \olref[prp]{ax:lif1} بېلګه
ده۔ او که \LR{$!B$}، \LR{$(!D \lif !D)$} وي، نو \LR{$!D \lif !B$} به هم د
\olref[prp]{ax:lif1} بله بېلګه وي۔ نو زموږ اشتقاق دا دے:
\begin{derivation}
  1. & \LR{$!D \lif ((!D \lif !D) \lif !D)$} & \olref[prp]{ax:lif1}\\
  2. & \LR{$(!D \lif ((!D \lif !D) \lif !D)) \lif {}$}\\
  & \qquad \LR{$((!D \lif (!D \lif !D)) \lif (!D \lif !D))$} & \olref[prp]{ax:lif2}\\
  3. & \LR{$(!D \lif (!D \lif !D)) \lif (!D \lif !D)$} & 1, 2, \MP\\
  4. & \LR{$!D \lif (!D \lif !D)$} & \olref[prp]{ax:lif1}\\
  5. & \LR{$!D \lif !D$} & 3, 4, \MP
\end{derivation}
\end{ex}
\olpseganchor{OLP-0116-B008}{end}

\olpseganchor{OLP-0116-B009}{start}
\begin{ex}\ollabel{ex:chain}
کله کله غواړو وښيو چې له ځينو نورو فارمولونه~\LR{$\Gamma$} څخه د کوم
فارمول يو اشتقاق شته۔ مثلاً راځئ وښيو چې له
\LR{$\Gamma = \{!A \lif !B, !B \lif !C\}$} څخه
\LR{$!A \lif !C$}، اشتقاق کولے شو۔
\begin{derivation}
  1. & \LR{$!A \lif !B$} & \Hyp\\
  2. & \LR{$!B \lif !C$} & \Hyp\\
  3. & \LR{$(!B \lif !C) \lif (!A \lif (!B \lif !C))$} & \olref[prp]{ax:lif1} \\
  4. & \LR{$!A \lif (!B \lif !C)$} & 2, 3, \MP\\
  5. & \LR{$(!A \lif (!B \lif !C)) \lif {}$}\\
  & \qquad \LR{$((!A \lif !B) \lif (!A \lif !C))$} & \olref[prp]{ax:lif2}\\
  6. & \LR{$((!A \lif !B) \lif (!A \lif !C))$} & 4, 5, \MP\\
  7. & \LR{$!A \lif !C$} & 1, 6, \MP
\end{derivation}
هغه کرښې چې ``\Hyp'' (د ``فرضیې'' دپاره) پرې ليکل شوي دي، ښيي چې
په هغې کرښه فارمول د~\LR{$\Gamma$} يو غړے دے۔
\end{ex}
\olpseganchor{OLP-0116-B009}{end}

\olpseganchor{OLP-0116-B010}{start}
\begin{prop}
  \ollabel{prop:chain} که \LR{$\Gamma \Proves !A \lif !B$} او \LR{$\Gamma
  \Proves !B \lif !C$} وي، نو \LR{$\Gamma \Proves !A \lif !C$}
\end{prop}
\olpseganchor{OLP-0116-B010}{end}

\olpseganchor{OLP-0116-B011}{start}
\begin{proof}
  فرض کړئ \LR{$\Gamma \Proves !A \lif !B$} او \LR{$\Gamma \Proves !B \lif
  !C$}۔ بيا له~\LR{$\Gamma$} څخه د~\LR{$!A \lif !B$} يو اشتقاق شته؛
  او له~\LR{$\Gamma$} څخه د~\LR{$!B \lif !C$} يو اشتقاق هم شته۔ د
  هغو په يو بل پسې نښلولو ئې په يو اشتقاق کښې يو ځاے کړئ۔
  اوس د مخکښې بېلګې د~اشتقاق ۳--۷ کرښې ورزياتې کړئ۔ دا له~\LR{$\Gamma$} څخه د
  \LR{$!A \lif !C$} يو اشتقاق دے---او همدا د نوي
  اشتقاق وروستۍ کرښه ده۔ پام وکړئ چې د ۴مې او ۷مې کرښې
  توجيه لا هم سمه پاتې کېږي که د ۱مې کرښې حواله د~\LR{$!A \lif !B$}
  د~اشتقاق د وروستۍ کرښې په حواله، او د ۲مې کرښې حواله د
  \LR{$!B \lif !C$} د~اشتقاق د وروستۍ کرښې په حواله بدله شي۔
\end{proof}
\olpseganchor{OLP-0116-B011}{end}

\olpseganchor{OLP-0116-B012}{start}
\begin{prob} 
له بديهي اصولو څخه د اشتقاقونه په وړاندې کولو وښيئ چې لاندې
خبرې سمې دي:
\begin{enumerate} 
\item \LR{$(!A \land !B) \lif (!B \land !A)$}
\item \LR{$((!A \land !B) \lif !C) \lif (!A \lif (!B \lif !C))$}
\item \LR{$\lnot(!A \lor !B) \lif \lnot !A$}
\end{enumerate} 
\end{prob}
\olpseganchor{OLP-0116-B012}{end}

\olpseganchor{OLP-0117-B004}{start}
\olfileid{fol}{axd}{prq}
\olpseganchor{OLP-0117-B004}{end}

\olpseganchor{OLP-0117-B005}{start}
\olsection{له کميت ټاکونکو سره اشتقاقونه}
\olpseganchor{OLP-0117-B005}{end}

\olpseganchor{OLP-0117-B006}{start}
\begin{ex}
راځئ د~\LR{$(\lforall[x][!A(x)] \land
\lforall[y][!B(y)]) \lif \lforall[x][(!A(x) \land !B(x))]$} يو
اشتقاق ورکړو۔
\olpseganchor{OLP-0117-B006}{end}

\olpseganchor{OLP-0117-B007}{start}
لومړے پام وکړئ چې
\begin{align*}
  (\lforall[x][!A(x)] \land \lforall[y][!B(y)]) & \lif \lforall[x][!A(x)]\\
  \intertext{\RL{د \olref[prp]{ax:land1} يوه بېلګه ده، او}}
  \lforall[x][!A(x)] & \lif !A(a) \\
  \intertext{\RL{د \olref[qua]{ax:q1} بېلګه ده۔ نو د \olref[pro]{prop:chain} له مخې پوهېږو چې}}
  (\lforall[x][!A(x)] \land \lforall[y][!B(y)]) & \lif !A(a) \\
  \intertext{\RL{د اشتقاق وړ دے۔ همدارنګه، ځکه}}
  (\lforall[x][!A(x)] \land \lforall[y][!B(y)]) & \lif \lforall[y][!B(y)] \qquad\text{\RL{او}}\\
    \lforall[y][!B(y)] & \lif !B(a)\\
    \intertext{\RL{په ترتيب د \olref[prp]{ax:land2} او \olref[qua]{ax:q1} بېلګې دي،}}
    (\lforall[x][!A(x)] \land \lforall[y][!B(y)]) & \lif !B(a)\\
    \intertext{\RL{د \olref[pro]{prop:chain} په وسيله د ثبوت وړ دے۔ د \olref[prp]{ax:land3} د يوې مناسبې بېلګې او د~\MP{} د دوو تطبيقونو په کارولو وينو چې}}
    (\lforall[x][!A(x)] \land \lforall[y][!B(y)]) & \lif (!A(a) \land !B(a))\\
    \intertext{\RL{د ثبوت وړ دے۔ اوس \QR{} پلي کولے شو او دا ترلاسه کوو:}}
(\lforall[x][!A(x)] \land \lforall[y][!B(y)]) & \lif \lforall[x][(!A(x) \land !B(x))].
\end{align*}
\end{ex}
\olpseganchor{OLP-0117-B007}{end}

\olpseganchor{OLP-0118-B004}{start}
\olfileid{fol}{axd}{ptn}
\olpseganchor{OLP-0118-B004}{end}

\olpseganchor{OLP-0118-B005}{start}
\olsection{ثبوت-تيوريکي مفاهيم}
\olpseganchor{OLP-0118-B005}{end}

\olpseganchor{OLP-0118-B006}{start}
\begin{explain}
لکه څنګه چې مو يو شمېر مهم معنايي مفاهيم
(اعتبار، معنايي لازميت او د صدق وړتيا)
تعريف کړي دي، اوس ورسره سم \emph{ثبوت-تيوريکي مفاهيم} تعريفوو۔
دا د جوړښتونه په دننه کښې د تړلي فارمولونه پوره کېدو ته په
رجوع نۀ، بلکې د ځينو فارمولونو د اشتقاق وړتيا يا
د اشتقاق نۀ وړتيا ته په رجوع تعريفېږي۔ دا يوه مهمه موندنه وه چې
دغه مفاهيم يو شان راځي۔ د هغو يو شان راتلل د \emph{سموالي} او
\emph{بشپړوالي د قضيو} منځپانګه ده۔
\end{explain}
\olpseganchor{OLP-0118-B006}{end}

\olpseganchor{OLP-0118-B007}{start}
\begin{defn}[د اشتقاق وړتيا]
فارمول~\LR{$!A$} هله له~\LR{$\Gamma$} څخه \emph{د اشتقاق وړ} دے،
چې \LR{$\Gamma \Proves !A$} ليکو، چې له~\LR{$\Gamma$} څخه داسې
اشتقاق وي چې په~\LR{$!A$} پاے ته رسېږي۔
\end{defn}
\olpseganchor{OLP-0118-B007}{end}

\olpseganchor{OLP-0118-B008}{start}
\begin{defn}[قضيې]
فارمول~\LR{$!A$} هله \emph{قضيه} ده چې له تش سټ څخه د~\LR{$!A$} يو
اشتقاق وي۔ که \LR{$!A$} قضيه وي \LR{$\Proves !A$}، او که نۀ وي
\LR{$\Proves/ !A$} ليکو۔
\end{defn}
\olpseganchor{OLP-0118-B008}{end}

\olpseganchor{OLP-0118-B009}{start}
\begin{defn}[سازګاري]
د فارمولونه يو سټ~\LR{$\Gamma$} هله او يوازې هله \emph{سازګار} دے
چې \LR{$\Gamma\Proves/ \lfalse$}؛ که داسې نۀ وي \emph{ناسازګار} دے۔
\end{defn}
\olpseganchor{OLP-0118-B009}{end}

\olpseganchor{OLP-0118-B010}{start}
\begin{prop}[انعکاس]
\ollabel{prop:reflexivity}
که \LR{$!A \in \Gamma$} وي، نو \LR{$\Gamma \Proves !A$}۔
\end{prop}
\olpseganchor{OLP-0118-B010}{end}

\olpseganchor{OLP-0118-B011}{start}
\begin{proof}
  يوازې فارمول~\LR{$!A$} له~\LR{$\Gamma$} څخه د~\LR{$!A$} يو
  اشتقاق دے۔
\end{proof}
\olpseganchor{OLP-0118-B011}{end}

\olpseganchor{OLP-0118-B012}{start}
\begin{prop}[زياتېدنه]
\ollabel{prop:monotonicity}
که \LR{$\Gamma \subseteq \Delta$} او \LR{$\Gamma \Proves !A$} وي، نو
\LR{$\Delta \Proves !A$}۔
\end{prop}
\olpseganchor{OLP-0118-B012}{end}

\olpseganchor{OLP-0118-B013}{start}
\begin{proof}
له~\LR{$\Gamma$} څخه د~\LR{$!A$} هر اشتقاق له~\LR{$\Delta$} څخه هم د~\LR{$!A$}
يو اشتقاق دے۔
\end{proof}
\olpseganchor{OLP-0118-B013}{end}

\olpseganchor{OLP-0118-B014}{start}
\begin{prop}[انتقال]
\ollabel{prop:transitivity}
که \LR{$\Gamma \Proves !A$} او \LR{$\{!A\} \cup \Delta \Proves !B$} وي،
نو \LR{$\Gamma \cup \Delta \Proves !B$}۔
\end{prop}
\olpseganchor{OLP-0118-B014}{end}

\olpseganchor{OLP-0118-B015}{start}
\begin{proof}
  فرض کړئ \LR{$\{!A\} \cup \Delta \Proves !B$}۔ بيا له
  \LR{$\{!A\} \cup \Delta$} څخه يو اشتقاق
  \LR{$!B_1$}، \dots، \LR{$!B_l = !B$} شته۔ په دې اشتقاق کښې ځينې
  ګامونه به د هغې قاعدې له مخې سم وي چې کومې مخکښې کرښې
  \LR{$!B_i = !A$} ته رجوع کوي۔ د فرض له مخې له~\LR{$\Gamma$} څخه د~\LR{$!A$}
  يو اشتقاق شته؛ يعنې داسې اشتقاق
  \LR{$!A_1$}، \dots، \LR{$!A_k = !A$} چې هر \LR{$!A_i$} يا بديهي اصل وي، يا
  د~\LR{$\Gamma$} يو غړے وي، يا د استنتاج د قاعدې له مخې سم
  وي۔ اوس دا لړۍ په پام کښې ونيسئ:
  \[
  !A_1, \dots, !A_k = !A, !B_1, \dots, !B_l = !B.
  \]
  دا له~\LR{$\Gamma \cup \Delta$} څخه د~\LR{$!B$} سم اشتقاق دے؛ ځکه
  هر \LR{$!B_i = !A$} اوس د هماغې قاعدې له مخې توجيه کېږي چې
  \LR{$!A_k = !A$} توجيه کوي۔
  \textbf{د ژباړې د نسخې سمون (OLAX-002):} سرچينې په وروستۍ جمله
  کښې يوازې له بي-آی څخه د ميتا-فارمول نښه پرېښې وه؛ هماغه کرښه
  پورته او په نښلول شوې لړۍ کښې په نښه لرلې بڼه تعريف شوې ده۔
\end{proof}
\olpseganchor{OLP-0118-B015}{end}

\olpseganchor{OLP-0118-B016}{start}
پام وکړئ چې په ځانګړي ډول دا معنٰی لري: که
\LR{$\Gamma \Proves !A$} او \LR{$!A \Proves !B$} وي، نو
\LR{$\Gamma \Proves !B$}۔ له دې څخه دا هم راځي چې که
\LR{$!A_1, \dots, !A_n \Proves !B$} او د هر~\LR{$i$} دپاره
\LR{$\Gamma \Proves !A_i$} وي، نو \LR{$\Gamma \Proves !B$}۔
\olpseganchor{OLP-0118-B016}{end}

\olpseganchor{OLP-0118-B017}{start}
\begin{prop}
\ollabel{prop:incons}
\LR{$\Gamma$} هله او يوازې هله ناسازګار دے چې د هر~\LR{$!A$} دپاره
\LR{$\Gamma \Proves !A$} وي۔
\end{prop}
\olpseganchor{OLP-0118-B017}{end}

\olpseganchor{OLP-0118-B018}{start}
\begin{proof}
تمرين۔
\end{proof}
\olpseganchor{OLP-0118-B018}{end}

\olpseganchor{OLP-0118-B019}{start}
\begin{prob}
\olref[fol][axd][ptn]{prop:incons} ثابت کړئ۔
\end{prob}
\olpseganchor{OLP-0118-B019}{end}



\olpseganchor{OLP-0118-B021}{start}
\begin{prop}[تراکم]
\ollabel{prop:proves-compact}
  \begin{enumerate}
  \item که \LR{$\Gamma \Proves !A$} وي، نو داسې متناهي فرعي سټ
    \LR{$\Gamma_0 \subseteq \Gamma$} شته چې \LR{$\Gamma_0 \Proves !A$}۔
  \item که د~\LR{$\Gamma$} هر متناهي فرعي سټ سازګار وي، نو \LR{$\Gamma$}
    سازګار دے۔
  \end{enumerate}
\end{prop}
\olpseganchor{OLP-0118-B021}{end}

\olpseganchor{OLP-0118-B022}{start}
\begin{proof}
  \begin{enumerate}
    \item که \LR{$\Gamma \Proves !A$} وي، نو د فارمولونه داسې متناهي
      لړۍ \LR{$!A_1$}، \dots،~\LR{$!A_n$} شته چې \LR{$!A \ident !A_n$}؛ او هر
      \LR{$!A_i$} يا منطقي بديهي اصل وي، يا د~\LR{$\Gamma$} يو غړے
      وي، يا د مودوس پوننس له مخې له مخکښو فارمولونه څخه راځي۔
      \LR{$\Gamma_0$} د هغو \LR{$!A_i$} سټ ونيسئ چې په~\LR{$\Gamma$} کښې دي۔ بيا
      همدا اشتقاق له~\LR{$\Gamma_0$} څخه هم يو اشتقاق
      دے، نو \LR{$\Gamma_0 \Proves !A$}۔
    \item دا د ځانګړي حالت~\LR{$!A \ident \lfalse$} دپاره د~(۱) مقابل
      عکس دے۔
\end{enumerate}
\end{proof}
\olpseganchor{OLP-0118-B022}{end}

\olpseganchor{OLP-0119-B005}{start}
\olfileid{fol}{axd}{ded}
\olpseganchor{OLP-0119-B005}{end}

\olpseganchor{OLP-0119-B006}{start}
\olsection{د استنتاج قضيه}
\olpseganchor{OLP-0119-B006}{end}

\olpseganchor{OLP-0119-B007}{start}
لکه چې مو وليدل، په بديهي نظام کښې اشتقاقونه ورکول ستونزمن
دي، او اشتقاقونه موندل هم ګرانېږي۔ د فارمولونه اوږدې لړۍ
په رښتيا ليکلو پر ځاے، زياتره دا استدلال اسانه وي چې داسې
اشتقاقونه شته؛ د دې دپاره څو ساده پايلې کاروو۔ درې داسې پايلې
مو لا ثابتې کړې دي: \olref[ptn]{prop:reflexivity} وايي کله چې
\LR{$!A \in \Gamma$} پوهېږو، تل \LR{$\Gamma \Proves !A$} ويلے شو۔
\olref[ptn]{prop:monotonicity} وايي که \LR{$\Gamma \Proves !A$} وي، نو
\LR{$\Gamma \cup \{!B\} \Proves !A$} هم دے۔ او
\olref[ptn]{prop:transitivity} ښيي که \LR{$\Gamma \Proves !A$} او
\LR{$!A \Proves !B$} وي، نو \LR{$\Gamma \Proves !B$}۔ دا د مودوس پوننس يوه
بله ساده، ``ميتا'' بڼه ده:
\olpseganchor{OLP-0119-B007}{end}

\olpseganchor{OLP-0119-B008}{start}
\begin{prop}
  \ollabel{prop:mp} که \LR{$\Gamma \Proves !A$} او
  \LR{$\Gamma \Proves !A \lif !B$} وي، نو \LR{$\Gamma \Proves !B$}۔
\end{prop}
\olpseganchor{OLP-0119-B008}{end}

\olpseganchor{OLP-0119-B009}{start}
\begin{proof}
  موږ \LR{$\{!A, !A \lif !B\} \Proves !B$} لرو:
  \begin{derivation}
    1. & \LR{$!A$} & فرض\\
    2. & \LR{$!A \lif !B$} & فرض\\
    3. & \LR{$!B$} & 1, 2, مودوس پوننس
  \end{derivation}
  د \olref[ptn]{prop:transitivity} له مخې \LR{$\Gamma \Proves !B$}۔
\end{proof}
\olpseganchor{OLP-0119-B009}{end}

\olpseganchor{OLP-0119-B010}{start}
په دې برخه کښې تر ټولو مهمه پايله چې کاروو ئې د استنتاج قضيه ده:
\olpseganchor{OLP-0119-B010}{end}

\olpseganchor{OLP-0119-B011}{start}
\begin{thm}[د استنتاج قضيه]
\ollabel{thm:deduction-thm} \LR{$\Gamma \cup \{!A\} \Proves !B$} هله
او يوازې هله چې \LR{$\Gamma \Proves !A \lif !B$}۔
\end{thm}
\olpseganchor{OLP-0119-B011}{end}

\olpseganchor{OLP-0119-B012}{start}
\begin{proof}
د ``که'' لوری سمدستي دے۔ که \LR{$\Gamma \Proves !A \lif !B$} وي، د
\olref[ptn]{prop:monotonicity} له مخې
\LR{$\Gamma \cup \{!A\} \Proves !A \lif !B$} هم دے۔ همدارنګه د
\olref[ptn]{prop:reflexivity} له مخې
\LR{$\Gamma \cup \{!A\} \Proves !A$}۔ نو د \olref{prop:mp} له مخې
\LR{$\Gamma \cup \{!A\} \Proves !B$}۔
\olpseganchor{OLP-0119-B012}{end}

\olpseganchor{OLP-0119-B013}{start}
د ``يوازې که'' لوري دپاره له~\LR{$\Gamma \cup \{!A\}$} څخه د~\LR{$!B$} د
اشتقاق پر اوږدوالي استقرا کوو۔
\olpseganchor{OLP-0119-B013}{end}

\olpseganchor{OLP-0119-B014}{start}
د استقرا د بنسټ دپاره دعوه د اوږدوالي~\LR{$1$} د هر اشتقاق
دپاره ثابتوو۔ له~\LR{$\Gamma \cup \{!A\}$} څخه د~\LR{$!B$} د اوږدوالي~\LR{$1$}
اشتقاق يوازې له~\LR{$!B$} څخه جوړ وي؛ او که سم وي، يا
\LR{$!B \in \Gamma \cup \{!A\}$} وي يا \LR{$!B$} بديهي اصل وي۔
\textbf{د ژباړې د نسخې سمون (OLAX-003):} سرچينې په دې غړيتوب کښې
له اې-متا فارمول څخه مخکښې يوازې د غړيتوب نښه ليکلې وه؛ د جملې
فاعل دلته د بشپړ فارمول په دننه کښې راوستل شو۔
که \LR{$!B \in \Gamma$} وي يا بديهي اصل وي، نو \LR{$\Gamma \Proves !B$}۔
همدارنګه د \olref[prp]{ax:lif1} له مخې
\LR{$\Gamma \Proves !B \lif (!A \lif !B)$}، او \olref{prop:mp} موږ ته
\LR{$\Gamma \Proves !A \lif !B$} راکوي۔ که \LR{$!B \in \{ !A\}$} وي، نو
\LR{$\Gamma \Proves !A \lif !B$}؛ ځکه وروستۍ تړلے فارمول
\LR{$!A \lif !B$} د~\LR{$!A \lif !A$} هماغه جمله ده، او په
\olref[pro]{ex:identity} کښې مو هغه اشتقاق کړې ده۔
\olpseganchor{OLP-0119-B014}{end}

\olpseganchor{OLP-0119-B015}{start}
د استقرا په پړاو کښې فرض کړئ له~\LR{$\Gamma \cup \{!A\}$} څخه د~\LR{$!B$}
يو اشتقاق په داسې ګام~\LR{$!B$} پاے ته رسېږي چې په مودوس پوننس
توجيه شوے دے۔ (که په مودوس پوننس توجيه شوے نۀ وي، نو
\LR{$!B \in \Gamma$}، \LR{$!B \ident !A$}، يا \LR{$!B$} بديهي اصل دے، او د استقرا
د بنسټ هماغه استدلال پرې پلي کېږي۔) بيا د کوم فارمول~\LR{$!C$}
دپاره په اشتقاق کښې ځينې مخکښې ګامونه \LR{$!C \lif !B$} او
\LR{$!C$} دي؛ يعنې \LR{$\Gamma \cup \{!A\} \Proves !C \lif !B$} او
\LR{$\Gamma \cup \{!A\} \Proves !C$}، او اړوند اشتقاقونه لنډ دي،
نو د استقرا فرضيه پرې پلي کېږي۔ ځکه دواړه دا لرو:
\begin{align*}
  & \Gamma \Proves !A \lif (!C \lif !B); \\
  & \Gamma \Proves !A \lif !C.
\end{align*}
خو د \olref[prp]{ax:lif2} له مخې دا هم لرو:
\[
\Gamma \Proves (!A \lif (!C \lif !B)) \lif
((!A\lif !C)  \lif (!A \lif !B)),
\]
او د \olref{prop:mp} دوه تطبيقونه، لکه چې پکار ده،
\LR{$\Gamma \Proves !A \lif !B$} راکوي۔
\end{proof}
\olpseganchor{OLP-0119-B015}{end}

\olpseganchor{OLP-0119-B016}{start}
پام وکړئ چې \olref[prp]{ax:lif1} او \olref[prp]{ax:lif2} په عين
دې موخه ټاکل شوي وو چې د استنتاج قضيه سمه وي۔
\olpseganchor{OLP-0119-B016}{end}

\olpseganchor{OLP-0119-B017}{start}
لاندې د~د اشتقاق وړتيا په اړه ځينې ګټور حقيقتونه دي چې د تمرينونو
په توګه ئې پرېږدو۔
\olpseganchor{OLP-0119-B017}{end}

\olpseganchor{OLP-0119-B018}{start}
\begin{prop}
  \ollabel{prop:derivfacts}
\begin{enumerate}
\item \LR{$\Proves (!A \lif !B) \lif ((!B \lif !C)
  \lif (!A \lif !C))$}؛ \ollabel{derivfacts:a}
\textbf{د ژباړې د نسخې سمون (OLAX-004):} سرچينې د لومړۍ دعوې د
بهرني وروستي قوس د تړلو نښه پرېښې وه؛ دلته د پرانيستو قوسونو سره
سمه يوه تړونکې نښه بېرته راوستل شوې ده۔
\item که \LR{$\Gamma \cup \{ \lnot !A\}
  \Proves \lnot !B$} وي، نو \LR{$\Gamma \cup \{ !B\} \Proves
  !A$} (مقابل عکس)؛ \ollabel{derivfacts:b}
\item \LR{$\{ !A, \lnot!A\} \Proves
    !B$} (له کاذبۍ هر څه، انفجار)؛ \ollabel{derivfacts:c}
\item \LR{$\{ \lnot\lnot!A\} \Proves
  !A$} (دوه ګونې نفي حذف)؛\ollabel{derivfacts:d}
\item که \LR{$\Gamma \Proves \lnot\lnot!A$} وي، نو \LR{$\Gamma \Proves
  !A$}؛\ollabel{derivfacts:e}
\end{enumerate}
\end{prop}
\olpseganchor{OLP-0119-B018}{end}

\olpseganchor{OLP-0119-B019}{start}
\begin{prob}
  \olref[fol][axd][ded]{prop:derivfacts} ثابت کړئ۔
\end{prob}
\olpseganchor{OLP-0119-B019}{end}

\olpseganchor{OLP-0120-B004}{start}
\olfileid{fol}{axd}{ddq}
\olpseganchor{OLP-0120-B004}{end}

\olpseganchor{OLP-0120-B005}{start}
\olsection{له کميت ټاکونکو سره د استنتاج قضيه}
\olpseganchor{OLP-0120-B005}{end}

\olpseganchor{OLP-0120-B006}{start}
\begin{thm}[د استنتاج قضيه]
\ollabel{thm:deduction-thm-q} که
\LR{$\Gamma \cup \{!A\} \Proves !B$} وي، نو
\LR{$\Gamma \Proves !A \lif !B$}۔
\end{thm}
\olpseganchor{OLP-0120-B006}{end}

\olpseganchor{OLP-0120-B007}{start}
\begin{proof}
يو ځل بيا له~\LR{$\Gamma \cup \{!A\}$} څخه د~\LR{$!B$} د اشتقاق
پر اوږدوالي استقرا کوو۔
\olpseganchor{OLP-0120-B007}{end}

\olpseganchor{OLP-0120-B008}{start}
د استقرا د بنسټ ثبوت د
\olref[ded]{thm:deduction-thm} په ثبوت کښې له بنسټ سره يو شان دے۔
\olpseganchor{OLP-0120-B008}{end}

\olpseganchor{OLP-0120-B009}{start}
د استقرا په پړاو کښې بيا فرض کړئ له~\LR{$\Gamma \cup \{!A\}$} څخه د
\LR{$!B$} اشتقاق په داسې ګام~\LR{$!B$} پاے ته رسېږي چې په استنتاجي
قاعده توجيه شوے دے۔ که استنتاجي قاعده مودوس پوننس وي، د
\olref[ded]{thm:deduction-thm} د ثبوت په شان مخکښې ځو۔ که قاعده
\QR{} وي، پوهېږو چې \LR{$!B \ident !C \lif \lforall[x][!D(x)]$} او د
\LR{$!C \lif !D(a)$} بڼې فارمول په اشتقاق کښې له مخکښې
راغلے دے، چې \LR{$a$} په~\LR{$!C$}، \LR{$!A$} يا~\LR{$\Gamma$} کښې نۀ راځي۔ نو دا لرو:
\begin{align*}
  \Gamma \cup \{!A\} & \Proves !C \lif !D(a),\\
  \intertext{\RL{او د استقرا فرضيه پلي کېږي؛ يعنې دا لرو:}}
    \Gamma & \Proves !A \lif (!C \lif !D(a)).\\
  \intertext{\RL{له دې څخه}}
  & \Proves (!A \lif (!C \lif !D(a))) \lif ((!A \land !C) \lif !D(a))\\
  \intertext{\RL{او مودوس پوننس سره ترلاسه کوو:}}
  \Gamma & \Proves (!A \land !C) \lif !D(a).\\
  \intertext{\RL{ځکه د ځانګړي متغير شرط لا هم پلي کېږي، دې اشتقاق ته په \QR{} توجيه شوے ګام ورزياتولے شو، او ترلاسه کوو:}}
    \Gamma & \Proves (!A \land !C) \lif \lforall[x][!D(x)].\\
    \intertext{\RL{دا هم لرو:}}
    & \Proves ((!A \land !C) \lif \lforall[x][!D(x)]) \lif (!A \lif (!C \lif \lforall[x][!D(x)])),\\
    \intertext{\RL{نو د مودوس پوننس له مخې،}}
    \Gamma & \Proves !A \lif (!C \lif \lforall[x][!D(x)]),
\end{align*}
\textbf{د ژباړې د نسخې سمون (OLAX-005):} سرچينې د پورته ماقبل
وروستۍ بديهي بېلګه کښې د بهرني شرطي يو تړونکے قوس پرېښے و؛ هغه
قوس دلته د بشپړې شېما دپاره بېرته راوستل شو۔
يعنې \LR{$\Gamma \Proves !A \lif !B$}۔
\textbf{د ژباړې د نسخې سمون (OLAX-006):} سرچينې وروستۍ پايله يوازې
د بي د ثبوت په توګه ليکلې وه؛ خو د قضیې موخه او سمدستي مخکښې کرښه
د اې-که-بي ثبوت ورکوي، نو هماغه بشپړه پايله راوستل شوې ده۔
\olpseganchor{OLP-0120-B009}{end}

\olpseganchor{OLP-0120-B010}{start}
هغه حالت د تمرين په توګه پرېږدو چې \LR{$!B$} د~\QR{} قاعدې په وسيله
توجيه کېږي، خو د~\LR{$\lexists[x][!D(x)] \lif !C$} بڼه لري۔
\end{proof}
\olpseganchor{OLP-0120-B010}{end}

\olpseganchor{OLP-0120-B011}{start}
\begin{prob}
  د~\olref[fol][axd][ddq]{thm:deduction-thm-q} ثبوت بشپړ کړئ۔
\end{prob}
\olpseganchor{OLP-0120-B011}{end}

\olpseganchor{OLP-0121-B004}{start}
\olfileid{fol}{axd}{prv}
\olpseganchor{OLP-0121-B004}{end}

\olpseganchor{OLP-0121-B005}{start}
\olsection{د اشتقاق وړتيا او سازګاري}
\olpseganchor{OLP-0121-B005}{end}

\olpseganchor{OLP-0121-B006}{start}
اوس به د~د اشتقاق وړتيا اړيکې يو شمېر خاصيتونه ثابت کړو۔ دا په
خپلواکه توګه په زړه پورې دي، خو هر يو به د بشپړوالي د قضیې په
ثبوت کښې ونډه ولري۔
\olpseganchor{OLP-0121-B006}{end}

\olpseganchor{OLP-0121-B007}{start}
\begin{prop}\ollabel{prop:provability-contr}
  که \LR{$\Gamma \Proves !A$} او \LR{$\Gamma \cup \{!A\}$} ناسازګار وي، نو
  \LR{$\Gamma$} ناسازګار دے۔
\end{prop}
\olpseganchor{OLP-0121-B007}{end}

\olpseganchor{OLP-0121-B008}{start}
\begin{proof}
  که \LR{$\Gamma \cup \{!A\}$} ناسازګار وي، نو
  \LR{$\Gamma \cup \{!A\} \Proves \lfalse$}۔ د
  \olref[ptn]{prop:reflexivity} له مخې د هر \LR{$!B \in \Gamma$} دپاره
  \LR{$\Gamma \Proves !B$}۔ د فرض له مخې \LR{$\Gamma \Proves !A$} هم دے،
  نو د هر \LR{$!B \in \Gamma \cup \{!A\}$} دپاره
  \LR{$\Gamma \Proves !B$}۔ د \olref[ptn]{prop:transitivity} له مخې
  \LR{$\Gamma \Proves \lfalse$}؛ يعنې \LR{$\Gamma$} ناسازګار دے۔
\end{proof}
\olpseganchor{OLP-0121-B008}{end}

\olpseganchor{OLP-0121-B009}{start}
\begin{prop}
\ollabel{prop:prov-incons}
\LR{$\Gamma \Proves !A$} هله او يوازې هله چې
\LR{$\Gamma \cup \{\lnot !A\}$} ناسازګار وي۔
\end{prop}
\olpseganchor{OLP-0121-B009}{end}

\olpseganchor{OLP-0121-B010}{start}
\begin{proof}
لومړے فرض کړئ \LR{$\Gamma \Proves !A$}۔ د
\olref[ptn]{prop:monotonicity} له مخې
\LR{$\Gamma \cup \{\lnot !A\} \Proves !A$}۔ د
\olref[ptn]{prop:reflexivity} له مخې
\LR{$\Gamma \cup \{\lnot !A\} \Proves \lnot !A$}۔ د
\olref[prp]{ax:lnot2} له مخې
\LR{$\Proves \lnot !A \lif (!A \lif \lfalse)$} هم لرو۔ نو د
\olref[ded]{prop:mp} په دوه ځله پلي کولو
\LR{$\Gamma \cup \{\lnot !A\} \Proves \lfalse$} لرو۔
\olpseganchor{OLP-0121-B010}{end}

\olpseganchor{OLP-0121-B011}{start}
اوس فرض کړئ \LR{$\Gamma \cup \{\lnot !A\}$} ناسازګار دے؛ يعنې
\LR{$\Gamma \cup \{\lnot !A\} \Proves \lfalse$}۔ د استنتاج د قضیې له
مخې \LR{$\Gamma \Proves \lnot !A \lif \lfalse$}۔ د
\olref[prp]{ax:lfalse2} له مخې
\LR{$\Gamma \Proves (\lnot !A \lif \lfalse) \lif \lnot\lnot !A$}، نو
د \olref[ded]{prop:mp} له مخې \LR{$\Gamma \Proves \lnot\lnot !A$}۔
ځکه \LR{$\Gamma \Proves \lnot\lnot !A \lif !A$}
(\olref[prp]{ax:dne})، د \olref[ded]{prop:mp} په بيا پلي کولو
\LR{$\Gamma \Proves !A$} لرو۔
\end{proof}
\olpseganchor{OLP-0121-B011}{end}

\olpseganchor{OLP-0121-B012}{start}
\begin{prob}
ثابته کړئ چې \LR{$\Gamma \Proves \lnot !A$} هله او يوازې هله چې
\LR{$\Gamma \cup \{!A\}$} ناسازګار وي۔
\end{prob}
\olpseganchor{OLP-0121-B012}{end}

\olpseganchor{OLP-0121-B013}{start}
\begin{prop}\ollabel{prop:explicit-inc}
  که \LR{$\Gamma \Proves !A$} او \LR{$\lnot !A \in \Gamma$} وي، نو
  \LR{$\Gamma$} ناسازګار دے۔
\end{prop}
\olpseganchor{OLP-0121-B013}{end}

\olpseganchor{OLP-0121-B014}{start}
\begin{proof}
  د \olref[prp]{ax:lnot2} له مخې
  \LR{$\Gamma \Proves \lnot !A \lif (!A \lif \lfalse)$}۔ د
  \olref[ded]{prop:mp} په دوه ځله پلي کولو
  \LR{$\Gamma \Proves \lfalse$}۔
\end{proof}
\olpseganchor{OLP-0121-B014}{end}


\olpseganchor{OLP-0121-B015}{start}
\begin{prop}\ollabel{prop:provability-exhaustive}
  که \LR{$\Gamma \cup \{!A\}$} او \LR{$\Gamma \cup \{\lnot !A\}$} دواړه
  ناسازګار وي، نو \LR{$\Gamma$} ناسازګار دے۔
\end{prop}
\olpseganchor{OLP-0121-B015}{end}

\olpseganchor{OLP-0121-B016}{start}
\begin{proof}
تمرين۔
\end{proof}
\olpseganchor{OLP-0121-B016}{end}

\olpseganchor{OLP-0121-B017}{start}
\begin{prob}
  \olref[fol][axd][prv]{prop:provability-exhaustive} ثابت کړئ۔
\end{prob}
\olpseganchor{OLP-0121-B017}{end}

\olpseganchor{OLP-0122-B004}{start}
\olfileid{fol}{axd}{ppr}
\olpseganchor{OLP-0122-B004}{end}

\olpseganchor{OLP-0122-B005}{start}
\olsection{د اشتقاق وړتيا او د قضیې نښلوونکي}
\olpseganchor{OLP-0122-B005}{end}

\olpseganchor{OLP-0122-B006}{start}
\begin{explain}
  ثابتوو چې د بديهي استنتاج د~\LR{$\Proves$} د اشتقاق وړتيا اړيکه
  دومره پياوړې ده چې د قضيې د نښلوونکو ځينې بنسټيز حقيقتونه ثابت
  کړي؛ لکه \LR{$!A \land !B \Proves !A$} او
  \LR{$!A, !A \lif !B \Proves !B$} (مودوس پوننس)۔ دغه حقيقتونه د
  بشپړوالي د قضیې د ثبوت دپاره پکار دي۔
\end{explain}
\olpseganchor{OLP-0122-B006}{end}

\olpseganchor{OLP-0122-B007}{start}
\begin{prop}\ollabel{prop:provability-land}
  \begin{enumerate}
  \item \ollabel{prop:provability-land-left} دواړه
    \LR{$!A \land !B \Proves !A$} او \LR{$!A \land !B \Proves !B$}
  \item \ollabel{prop:provability-land-right}
    \LR{$!A, !B \Proves !A \land !B$}۔
  \end{enumerate}
\end{prop}
\olpseganchor{OLP-0122-B007}{end}

\olpseganchor{OLP-0122-B008}{start}
\begin{proof}
  \begin{enumerate}
    \item د \olref[prp]{ax:land1} او \olref[prp]{ax:land2} څخه په
      مودوس پوننس۔
      \textbf{د ژباړې د نسخې سمون (OLAX-007):} سرچينې د دواړو
      جلا عطف-ايستنو دپاره لومړني بديهي اصل ته دوه ځله حواله کړې وه؛
      دويمه دعوه د دويم عطف-ايستنې بديهي اصل کاروي۔
  \item د \olref[prp]{ax:land3} څخه د مودوس پوننس په دوه تطبيقونو۔
  \end{enumerate}
\end{proof}
\olpseganchor{OLP-0122-B008}{end}


\olpseganchor{OLP-0122-B009}{start}
\begin{prop}\ollabel{prop:provability-lor}
  \begin{enumerate}
  \item \LR{$!A \lor !B, \lnot !A, \lnot !B$} ناسازګار دے۔
  \item دواړه \LR{$!A \Proves !A \lor !B$} او
    \LR{$!B \Proves !A \lor !B$}۔
  \end{enumerate}
\end{prop}
\olpseganchor{OLP-0122-B009}{end}

\olpseganchor{OLP-0122-B010}{start}
\begin{proof}
  \begin{enumerate}
  \item له \olref[prp]{ax:lnot2} څخه
    \LR{$\Proves \lnot !A \lif (!A \lif \lfalse)$} او
    \LR{$\Proves \lnot !B \lif (!B \lif \lfalse)$} ترلاسه کوو۔ نو د
    استنتاج د قضیې له مخې \LR{$\{\lnot !A\} \Proves !A \lif \lfalse$}
    او \LR{$\{\lnot !B\} \Proves !B \lif \lfalse$} لرو۔ له
    \olref[prp]{ax:lor3} څخه
    \LR{$\{\lnot !A, \lnot !B\} \Proves (!A \lor !B) \lif \lfalse$}
    ترلاسه کوو۔ د استنتاج د قضیې له مخې
    \LR{$\{!A \lor !B, \lnot !A, \lnot !B\} \Proves \lfalse$}۔
    \textbf{د ژباړې د نسخې سمون (OLAX-008):} سرچينې دلته د نفي
    لومړني بديهي اصل ته حواله کړې وه؛ ښودل شوې انفجاري شېماوې د نفي
    دويم بديهي اصل بېلګې دي۔
  \item له \olref[prp]{ax:lor1} او \olref[prp]{ax:lor2} څخه په
    مودوس پوننس۔
  \end{enumerate}
\end{proof}
\olpseganchor{OLP-0122-B010}{end}

\olpseganchor{OLP-0122-B011}{start}
\begin{prop}\ollabel{prop:provability-lif}
  \begin{enumerate}
  \item \ollabel{prop:provability-lif-left}
    \LR{$!A, !A \lif !B \Proves !B$}۔
  \item \ollabel{prop:provability-lif-right}
    دواړه \LR{$\lnot !A \Proves !A \lif !B$} او
    \LR{$!B \Proves !A \lif !B$}۔
  \end{enumerate}
\end{prop}
\olpseganchor{OLP-0122-B011}{end}

\olpseganchor{OLP-0122-B012}{start}
\begin{proof}
  \begin{enumerate}
  \item دا اشتقاق کولے شو:
    \begin{derivation}
      1. & \LR{$!A$} & \Hyp\\
      2. & \LR{$!A \lif !B$} & \Hyp\\
      3. & \LR{$!B$} & 1, 2, \MP
    \end{derivation}
\olpseganchor{OLP-0122-B012}{end}
    
\olpseganchor{OLP-0122-B013}{start}
  \item په ترتيب د \olref[prp]{ax:lnot2} او
    \olref[prp]{ax:lif1}، او د استنتاج د قضیې له مخې۔
  \end{enumerate}
\end{proof}
\olpseganchor{OLP-0122-B013}{end}

\olpseganchor{OLP-0123-B005}{start}
\olfileid{fol}{axd}{qpr}
\olpseganchor{OLP-0123-B005}{end}

\olpseganchor{OLP-0123-B006}{start}
\olsection{د اشتقاق وړتيا او کميت ټاکونکي}
\olpseganchor{OLP-0123-B006}{end}

\olpseganchor{OLP-0123-B007}{start}
\begin{explain}
  د بشپړوالي قضيه دا هم غواړي چې بديهي استنتاجونه د~\LR{$\Proves$} په
  اړه په دې برخه کښې ثابت شوي حقيقتونه ورکړي۔
\end{explain}
\olpseganchor{OLP-0123-B007}{end}

\olpseganchor{OLP-0123-B008}{start}
\begin{thm}
\ollabel{thm:strong-generalization} که \LR{$c$} داسې ثابت سمبول وي چې
په~\LR{$\Gamma$} يا~\LR{$!A(x)$} کښې نۀ راځي، او \LR{$\Gamma \Proves !A(c)$} وي،
نو \LR{$\Gamma \Proves \lforall[x][!A(x)]$}۔
\end{thm}
\olpseganchor{OLP-0123-B008}{end}

\olpseganchor{OLP-0123-B009}{start}
\begin{proof}
د استنتاج د قضیې له مخې \LR{$\Gamma \Proves \ltrue \lif !A(c)$}۔ ځکه
\LR{$c$} په~\LR{$\Gamma$} يا~\LR{$\top$} کښې نۀ راځي، ترلاسه کوو چې
\LR{$\Gamma \Proves \ltrue \lif \lforall[x][!A(x)]$}۔ ځکه رښتيا يو
بديهي اصل دے، د ميتا-مودوس پوننس له مخې
\LR{$\Gamma \Proves \lforall[x][!A(x)]$}۔
\textbf{د ژباړې د نسخې سمون (OLAX-009):} سرچينې وروستے ګام د
استنتاج قضیې ته يو ځل بيا منسوب کړے و؛ هغه قضيه يوازې د رښتيا فرض
ورزياتوي، حال دا چې دلته رښتيا بديهي اصل او مودوس پوننس فرض بېرته
لرې کوي۔
\end{proof}
\olpseganchor{OLP-0123-B009}{end}

\olpseganchor{OLP-0123-B010}{start}
\begin{prop}
\ollabel{prop:provability-quantifiers}
\begin{enumerate}
\item \LR{$!A(t) \Proves \lexists[x][!A(x)]$}.
\olpseganchor{OLP-0123-B010}{end}

\olpseganchor{OLP-0123-B011}{start}
\item \LR{$\lforall[x][!A(x)] \Proves !A(t)$}.
\end{enumerate}
\end{prop}
\olpseganchor{OLP-0123-B011}{end}

\olpseganchor{OLP-0123-B012}{start}
\begin{proof}
\begin{enumerate}
\item د \olref[qua]{ax:q2} او د استنتاج د قضیې له مخې۔
\item د \olref[qua]{ax:q1} او د استنتاج د قضیې له مخې۔
\end{enumerate}
\end{proof}
\olpseganchor{OLP-0123-B012}{end}

\olpseganchor{OLP-0124-B004}{start}
\olfileid{fol}{axd}{sou}
\olpseganchor{OLP-0124-B004}{end}
      
\olpseganchor{OLP-0124-B005}{start}
\olsection{سموالے}
\olpseganchor{OLP-0124-B005}{end}

\olpseganchor{OLP-0124-B006}{start}
\begin{explain}
د بديهي استنتاج په شان اشتقاق نظام هله \emph{سم} دے چې
هغه څيزونه اشتقاق نۀ کړي چې په حقيقت کښې نۀ صدق کېږي۔ نو
سموالے د اشتقاق نظامونو د تضمين شوې خونديتوب يو ډول خاصيت
دے۔ د دې له مخې چې کوم ثبوت-تيوريکي خاصيت تر بحث لاندې وي، مثلاً
غواړو پوه شو چې:
\begin{enumerate}
\item هر د اشتقاق وړ~\LR{$!A$} معتبر دے؛
\item که \LR{$!A$} له ځينو نورو~\LR{$\Gamma$} څخه د اشتقاق وړ وي، د هغو
  معنايي پايله هم ده؛
\item که د فارمولونه يو سټ~\LR{$\Gamma$} ناسازګار وي، د صدق وړ نۀ
  دے۔
\end{enumerate}
دا د اشتقاق نظام مهم خاصيتونه دي۔ که له دغو څخه کوم يو
سم نۀ وي، اشتقاق نظام نيمګړے دے---له حده زيات څيزونه به
اشتقاق کوي۔ له دې امله د اشتقاق نظام سموالے ثابتول
تر ټولو زيات اهميت لري۔
\end{explain}
\olpseganchor{OLP-0124-B006}{end}

\olpseganchor{OLP-0124-B007}{start}
\begin{prop}
  که \LR{$!A$} يو بديهي اصل وي، نو
  د هر جوړښت~\LR{$\Struct{M}$} او متغير-ګومارنې~\LR{$s$} دپاره \LR{$\Sat{M}{!A}[s]$}۔
\end{prop}
\olpseganchor{OLP-0124-B007}{end}

\olpseganchor{OLP-0124-B008}{start}
\begin{proof}
بايد تصديق کړو چې ټول بديهي اصول معتبر دي۔ د بېلګې
  په توګه د~\olref[qua]{ax:q1} حالت دا دے: فرض کړئ \LR{$t$} په~\LR{$!A$}
  کښې د~\LR{$x$} دپاره ازاد وي، او
  \LR{$\Sat{M}{\lforall[x][!A]}[s]$} ومنئ۔ بيا د پوره کېدو د تعريف له
  مخې، د هرې \LR{$\varAssign{s'}{s}{x}$} دپاره \LR{$\Sat{M}{!A}[s']$} هم
  دے، او په ځانګړي ډول دا هله سم دے چې
  \LR{$s'(x) = \Value{t}{M}[s]$} وي۔ د
  \olref[syn][ext]{prop:ext-formulas} له مخې
  \LR{$\Sat{M}{\Subst{!A}{t}{x}}[s]$}۔ دا ښيي چې
  \LR{$\Sat{M}{(\lforall[x][!A] \lif \Subst{!A}{t}{x})}[s]$}۔
\end{proof}
\olpseganchor{OLP-0124-B008}{end}

\olpseganchor{OLP-0124-B009}{start}
\begin{thm}[سموالے]
\ollabel{thm:soundness}
که \LR{$\Gamma \Proves !A$} وي، نو \LR{$\Gamma \Entails !A$}۔
\end{thm}
\olpseganchor{OLP-0124-B009}{end}

\olpseganchor{OLP-0124-B010}{start}
\begin{proof}
له~\LR{$\Gamma$} څخه د~\LR{$!A$} د اشتقاق پر اوږدوالي استقرا کوو۔
که داسې هېڅ ګام نۀ وي چې په استنتاج توجيه شوے وي، نو په
اشتقاق کښې ټول فارمولونه يا د بديهي اصولو بېلګې دي يا
په~\LR{$\Gamma$} کښې دي۔ د مخکښې قضیې له مخې ټول بديهي اصول
معتبر دي؛ نو که \LR{$!A$} بديهي اصل وي،
\LR{$\Gamma \Entails !A$}۔ که \LR{$!A \in \Gamma$} وي، نو په څرګنده
\LR{$\Gamma \Entails !A$}۔
\olpseganchor{OLP-0124-B010}{end}

\olpseganchor{OLP-0124-B011}{start}
که د~\LR{$!A$} د اشتقاق وروستنے ګام په مودوس پوننس توجيه شوے
وي، نو په اشتقاق کښې فارمولونه~\LR{$!B$} او~\LR{$!B \lif !A$}
شته، او د استقرا فرضيه د اشتقاق پر هغو برخو پلي کېږي چې په
دغو فارمولونه پاے ته رسېږي (ځکه په هره برخه کښې لږ تر لږه يو
په استنتاج توجيه شوے ګام کم دے)۔ نو د استقرا د فرضيې له مخې
\LR{$\Gamma \Entails !B$} او \LR{$\Gamma \Entails !B \lif !A$}۔ بيا د
\olref[syn][sem]{thm:sem-deduction}
له مخې \LR{$\Gamma \Entails !A$}۔
\olpseganchor{OLP-0124-B011}{end}

\olpseganchor{OLP-0124-B012}{start}
اوس فرض کړئ وروستنے ګام په~\QR توجيه شوے دے۔ بيا هغه
ګام د~\LR{$!C \lif \lforall[x][B(x)]$} بڼه لري، او ترې مخکښې ګام
\LR{$!C \lif !B(c)$} دے، چې \LR{$c$} په~\LR{$\Gamma$}، \LR{$!C$} يا
\LR{$\lforall[x][B(x)]$} کښې نۀ راځي۔ د استقرا د فرضيې له مخې
\LR{$\Gamma \Entails !C \lif !B(c)$}۔ د
\olref[syn][sem]{thm:sem-deduction} له مخې
\LR{$\Gamma \cup \{!C\} \Entails !B(c)$}۔
\olpseganchor{OLP-0124-B012}{end}

\olpseganchor{OLP-0124-B013}{start}
داسې يو جوړښت~\LR{$\Struct{M}$} په پام کښې ونيسئ چې
\LR{$\Sat{M}{\Gamma \cup \{!C\}}$} وي۔ بايد وښيو چې
\LR{$\Sat{M}{\lforall[x][!B(x)]}$}۔ ځکه
\LR{$\lforall[x][!B(x)]$} يوه تړلے فارمول ده، د دې معنٰی دا ده چې د
هرې متغير-ګومارنې~\LR{$s$} دپاره بايد \LR{$\Sat{M}{!B(x)}[s]$} وښيو
(\olref[syn][ass]{prop:sat-quant})۔ ځکه \LR{$\Gamma \cup \{!C\}$} ټول
له جملو جوړ دے، د هر~\LR{$!D \in \Gamma$} دپاره
\LR{$\Sat{M}{!D}[s]$} د~\olref[syn][sat]{defn:satisfaction} له مخې دے۔
\LR{$\Struct{M'}$} داسې جوړښت ونيسئ چې د~\LR{$\Struct{M}$} په شان وي، خو
\LR{$\Assign{c}{M'} = s(x)$} وي۔ ځکه \LR{$c$} په~\LR{$\Gamma$} يا~\LR{$!C$} کښې نۀ
راځي، د~\olref[syn][ext]{cor:extensionality-sent} له مخې
\LR{$\Sat{M'}{\Gamma \cup \{!C\}}$}۔ ځکه
\LR{$\Gamma \cup \{!C\} \Entails !B(c)$}، نو \LR{$\Sat{M'}{!B(c)}$}۔
\textbf{د ژباړې د نسخې سمون (OLAX-010):} سرچينې په دې پوره کېدو
ادعا کښې د بي-سي د ميتا-فارمول نښه پرېښې وه؛ په مخکښې لازميت او
ورپسې ټولو استعمالونو کښې همدا نښه موجوده ده، نو دلته بېرته
راوستل شوه۔ ځکه \LR{$!B(c)$} يوه تړلے فارمول ده، د
\olref[syn][ass]{prop:sentence-sat-true} له مخې
\LR{$\Sat{M'}{!B(c)}[s]$}۔ د~\olref[syn][ext]{prop:ext-formulas} له مخې
\LR{$\Sat{M'}{!B(x)}[s]$} هله او يوازې هله چې
\LR{$\Sat{M'}{!B(c)}[s]$} (ياد ولرئ چې \LR{$!B(c)$} هماغه
\LR{$\Subst{!B(x)}{c}{x}$} دے)۔ نو \LR{$\Sat{M'}{!B(x)}[s]$}۔ ځکه \LR{$c$} په
\LR{$!B(x)$} کښې نۀ راځي، د~\olref[syn][ext]{prop:extensionality} له
مخې \LR{$\Sat{M}{!B(x)}[s]$}۔ خو \LR{$s$} اختياري متغير-ګومارنه وه؛ نو
\LR{$\Sat{M}{\lforall[x][!B(x)]}$}۔ له دې امله
\LR{$\Gamma \cup \{!C\} \Entails \lforall[x][!B(x)]$}۔ د
\olref[syn][sem]{thm:sem-deduction} له مخې
\LR{$\Gamma \Entails !C \lif \lforall[x][!B(x)]$}۔
\olpseganchor{OLP-0124-B013}{end}

\olpseganchor{OLP-0124-B014}{start}
هغه حالت د تمرين په توګه پرېږدو چې \LR{$!A$} په~\QR{} توجيه کېږي، خو
د~\LR{$\lexists[x][!B(x)] \lif !C$} بڼه لري۔
\end{proof}
\olpseganchor{OLP-0124-B014}{end}

\olpseganchor{OLP-0124-B015}{start}
\begin{prob}
  د~\olref[fol][axd][sou]{thm:soundness} ثبوت بشپړ کړئ۔
\end{prob}
\olpseganchor{OLP-0124-B015}{end}

\olpseganchor{OLP-0124-B016}{start}
\begin{cor}
\ollabel{cor:weak-soundness}
که \LR{$\Proves !A$} وي، نو \LR{$!A$} معتبر دے۔
\end{cor}
\olpseganchor{OLP-0124-B016}{end}

\olpseganchor{OLP-0124-B017}{start}
\begin{cor}
\ollabel{cor:consistency-soundness}
که \LR{$\Gamma$} د صدق وړ وي، نو سازګار دے۔
\end{cor}
\olpseganchor{OLP-0124-B017}{end}

\olpseganchor{OLP-0124-B018}{start}
\begin{proof}
نقيض عکس ثابتوو۔ فرض کړئ \LR{$\Gamma$} سازګار نۀ دے۔ بيا
\LR{$\Gamma \Proves \lfalse$}؛ يعنې له~\LR{$\Gamma$} څخه د~\LR{$\lfalse$} يو
اشتقاق شته۔ د~\olref{thm:soundness} له مخې هر
جوړښت~\LR{$\Struct{M}$}
چې \LR{$\Gamma$} پوره کوي، بايد \LR{$\lfalse$} هم پوره کړي۔ ځکه د هر
جوړښت~\LR{$\Struct{M}$} دپاره \LR{$\Sat/{M}{\lfalse}$}، نو هېڅ
\LR{$\Struct{M}$}، \LR{$\Gamma$} نۀ شي پوره
کولے؛ يعنې \LR{$\Gamma$} د صدق وړ نۀ دے۔
\end{proof}
\olpseganchor{OLP-0124-B018}{end}

\olpseganchor{OLP-0125-B004}{start}
\olfileid{fol}{axd}{ide}
\olpseganchor{OLP-0125-B004}{end}

\olpseganchor{OLP-0125-B005}{start}
\olsection{له عينيت سره اشتقاقونه}
\olpseganchor{OLP-0125-B005}{end}

\olpseganchor{OLP-0125-B006}{start}
په اشتقاقونه کښې د~\LR{$\eq$} د شاملولو دپاره يوازې نوې بديهي
شېماوې ورزياتوو۔ د~اشتقاق او~\LR{$\Proves$} تعريف هماغسې پاتې
کېږي؛ يوازې نوي بديهي اصول هم منو۔
\olpseganchor{OLP-0125-B006}{end}

\olpseganchor{OLP-0125-B007}{start}
\begin{defn}[د عينيت بديهي اصول]
\begin{align}
\ollabel{ax:id1} & \eq[t][t], \\
\ollabel{ax:id2} & \eq[t_1][t_2] \lif (!B(t_1) \lif
  !B(t_2)),
\end{align}
د هر تړلي ترم~\LR{$t$}، \LR{$t_1$} او~\LR{$t_2$} دپاره۔
\end{defn}
\olpseganchor{OLP-0125-B007}{end}

\olpseganchor{OLP-0125-B008}{start}
\begin{prop}\ollabel{prop:sound}
  بديهي اصول~\olref{ax:id1} او~\olref{ax:id2} معتبر دي۔
\end{prop}
\olpseganchor{OLP-0125-B008}{end}

\olpseganchor{OLP-0125-B009}{start}
\begin{proof}
  تمرين۔
\end{proof}
\olpseganchor{OLP-0125-B009}{end}

\olpseganchor{OLP-0125-B010}{start}
\begin{prob}
\olref[fol][axd][ide]{prop:sound} ثابت کړئ۔
\end{prob}
\olpseganchor{OLP-0125-B010}{end}

\olpseganchor{OLP-0125-B011}{start}
\begin{prop}\ollabel{prop:iden1}
 د هر تړلي ترم~\LR{$t$} او سټ~\LR{$\Gamma$} دپاره
 \LR{$\Gamma \Proves \eq[t][t]$}۔
 \textbf{د ژباړې د نسخې سمون (OLAX-011):} سرچينې دلته «هر ترم»
 ويلي وو، خو د يووالي لومړے بديهي اصل يوازې پر تړلو ترمونو شېما
 ده؛ هماغه
 لازم قيد بېرته څرګند شو۔
\end{prop}
\olpseganchor{OLP-0125-B011}{end}

\olpseganchor{OLP-0125-B012}{start}
\begin{prop}\ollabel{prop:iden2}
  که \LR{$\Gamma \Proves !A(t_1)$} او \LR{$\Gamma \Proves
  \eq[t_1][t_2]$} وي، نو \LR{$\Gamma \Proves !A(t_2)$}، په دې شرط چې
  دواړه ښودل شوي ترمونه تړلي وي۔
  \textbf{د ژباړې د نسخې سمون (OLAX-012):} سرچينې په دې قضيه کښې
  د تړلو ترمونو شرط نۀ و بيان کړے، سره له دې چې ثبوت ئې نېغ په
  نېغه د يووالي دويم بديهي اصل کاروي او هغه شېما يوازې تړلي
  ترمونه مني۔
\end{prop}
\olpseganchor{OLP-0125-B012}{end}

\olpseganchor{OLP-0125-B013}{start}
\begin{proof}
دا فارمول
\[
(\eq[t_1][t_2] \lif (!A(t_1) \lif !A(t_2)))
\]
د~\olref{ax:id2} يوه بېلګه ده۔ پايله د~\MP په دوه پلي کولو سره
ترلاسه کېږي۔
\end{proof}
\olpseganchor{OLP-0125-B013}{end}

\olpseganchor{OLP-0126-B004}{start}
\olchapter{fol}{com}{د بشپړتيا قضيه}
\olpseganchor{OLP-0126-B004}{end}







\olpseganchor{OLP-0126-B008}{start}
%
  

\olpseganchor{OLP-0126-B008}{end}





\olpseganchor{OLP-0126-B011}{start}
%
  

\olpseganchor{OLP-0126-B011}{end}







\olpseganchor{OLP-0126-B015}{start}
%
  

\olpseganchor{OLP-0126-B015}{end}

\olpseganchor{OLP-0127-B004}{start}
\olfileid{fol}{com}{int}
\olpseganchor{OLP-0127-B004}{end}

\olpseganchor{OLP-0127-B005}{start}
\olsection{سريزه}
\olpseganchor{OLP-0127-B005}{end}

\olpseganchor{OLP-0127-B006}{start}
د بشپړتيا قضيه د منطق په اړه له تر ټولو بنسټيزو پايلو څخه يوه ده۔
دا په دوو بڼو وړاندې کېږي، او موږ به د دواړو همارزښت ثابت کړو۔ په
لومړۍ بڼه کښې دا زموږ د معنٰی پوهنې د پايلې او د اشتقاق
نظام ترمنځ د اړيکې په اړه يو بنسټيز حقيقت وايي: که
تړلے فارمول~\LR{$!A$} د ځينو تړلي فارمولونه~\LR{$\Gamma$} معنايي پايله
وي، نو داسې اشتقاق هم شته چې
\LR{$\Gamma \Proves !A$} ثابتوي۔ په دې توګه، د اشتقاق نظام
هغومره پياوړے دے څومره چې د هغو خبرو له ثابتولو پرته کېدے شي چې
په رښتيا نۀ لازمېږي۔
\olpseganchor{OLP-0127-B006}{end}

\olpseganchor{OLP-0127-B007}{start}
په دويمه بڼه کښې دا د مدل د شتون د پايلې په توګه ويل کېدے شي: د
تړلي فارمولونه هر سازګار سټ د صدق وړ دے۔ سازګاري يو ثبوت-تيوريکي
مفهوم دے: دا وايي چې زموږ د اشتقاق نظام د يو ټاکلي ډول
اشتقاقونه نۀ شي جوړولے۔ خو دا څوک تضمينوي چې يوازې ځکه چې
له~\LR{$\Gamma$} څخه د يو ټاکلي ډول هېڅ اشتقاقونه نشته، خامخا
داسې جوړښت~\LR{$\Struct{M}$}؟ د بشپړتيا د
قضيې تر لومړي ثبوت وړاندې---بلکې د هغو اشتقاق نظامونو تر
رامنځته کېدو هم وړاندې چې اوس ئې لرو---ستر جرمن رياضي‌پوه ډېوېډ
هيلبرټ دا نظر درلود چې د رياضي نظريو سازګاري د هغو څيزونو شتون
تضمينوي چې نظريې ئې په اړه دي۔ هغه ګوتلوب فرېګه ته په يو ليک کښې
داسې وويل:
\begin{quote}
  که په خپل‌سري ډول ورکړل شوي بديهي اصول له خپلو ټولو پايلو سره يو
  بل نۀ نقضوي، نو رښتوني دي او هغه څيزونه چې بديهي اصول ئې تعريفوي
  شتون لري۔ زما دپاره همدا د رښتياوالي او شتون معيار دے۔
\end{quote}
فرېګه سخت مخالفت وکړ۔ د بشپړتيا د قضيې دويمه بڼه ښيي چې هيلبرټ
لږ تر لږه په دې معنٰی سم و چې که بديهي اصول سازګار وي، نو داسې
\emph{کوم}
جوړښت شته چې ټول رښتوني کوي۔
\olpseganchor{OLP-0127-B007}{end}

\olpseganchor{OLP-0127-B008}{start}
دا يوازيني لاملونه نۀ دي چې د بشپړتيا قضيه---يا لا کره، د هغې
ثبوت---مهمه ده۔ دا يو شمېر مهمې پايلې لري چې ځينې به ئې جلا وڅېړو۔
د بېلګې په توګه، ځکه هر اشتقاق چې
\LR{$\Gamma \Proves !A$} ښيي متناهي دے، نو له~\LR{$\Gamma$} څخه يوازې
متناهي ډېر تړلي فارمولونه کارولے شي۔ له دې امله د بشپړتيا له قضيې
څخه لازمېږي چې که~\LR{$!A$} د~\LR{$\Gamma$} معنايي پايله وي، نو لا له وړاندې
د~\LR{$\Gamma$} د کوم متناهي فرعي سټ معنايي پايله ده۔ دې ته
\emph{د متناهي شاهد خاصيت} ويل کېږي۔ په همارز ډول، که د~\LR{$\Gamma$}
هر متناهي فرعي سټ سازګار وي، نو په خپله~\LR{$\Gamma$} هم سازګار دے۔
\olpseganchor{OLP-0127-B008}{end}

\olpseganchor{OLP-0127-B009}{start}
که څه هم د متناهي شاهد قضيه د بشپړتيا له قضيې څخه د
اشتقاقونه د منځنۍ لارې په تېرېدو لازمېږي، د بشپړتيا د قضيې
\emph{خپله ثبوت} هم نېغ په نېغه د هغې د ثابتولو دپاره کارول کېدے
شي۔ ځکه ثبوت د تړلي فارمولونه داسې سټ اخلي چې يو ټاکلے خاصيت---
سازګاري---لري، او له دې سټه داسې جوړښت جوړوي چې ټاکلي
خاصيتونه لري (په دې حالت کښې، سټ پوره کوي)۔ کابو همدا جوړونه د
متناهي شاهد خاصيت د نېغ ثبوت دپاره کارېږي، که د سازګارو سټونو پر
ځاے د تړلي فارمولونه له «په متناهي ډول د صدق وړ» سټونو څخه پيل
وکړو۔
دا جوړونه نورې پايلې هم ورکوي؛ د بېلګې په توګه، د
  تړلي فارمولونه هر د صدق وړ سټ يو متناهي يا شمېرېدونکے لامتناهي مدل
  لري۔ (دې پايلې ته د لوېنهايم--سکولم قضيه ويل کېږي۔) په عمومي ډول،
  د تړلي فارمولونه له سټونو څخه د جوړښتونه جوړول په منطق کښې
  ډېر ځله، او کله کله حتٰی په فلسفه کښې هم، کارول کېږي۔
\olpseganchor{OLP-0127-B009}{end}

\olpseganchor{OLP-0128-B004}{start}
\olfileid{fol}{com}{out}
\olpseganchor{OLP-0128-B004}{end}

\olpseganchor{OLP-0128-B005}{start}
\olsection{د ثبوت لنډيز}
\olpseganchor{OLP-0128-B005}{end}

\olpseganchor{OLP-0128-B006}{start}
د بشپړتيا د قضيې ثبوت يو څه پېچلے دے، او په لومړي لوست کښې ترې
لار ورکېدل اسانه دي۔ نو راځئ د ثبوت لنډيز وړاندې کړو۔ لومړے ګام د
نظر د زاويې داسې بدلون دے چې د ثبوت لاره راښيي۔ که بشپړتيا داسې
وګڼو چې «هر کله \LR{$\Gamma \Entails !A$} وي، نو
\LR{$\Gamma \Proves !A$} وي»، ښايي د کومې مفکورې موندل هم ګران وي:
ځکه د~\LR{$\Gamma \Proves !A$} د ښودلو دپاره بايد اشتقاق
ومومو، او داسې نۀ ښکاري چې فرضيه~\LR{$\Gamma \Entails !A$} په دې کښې
څه مرسته راسره کوي۔ د ځينو ثبوتي نظامونو دپاره اشتقاق
نېغ جوړېدے شي، خو موږ لږه بله لار نيسو۔ د نظر اړين بدلون دا دے:
بشپړتيا داسې هم ويل کېدے شي چې «که~\LR{$\Gamma$} سازګار وي، نو د صدق
وړ دے»۔ ښايي د~\LR{$\Gamma$} معلومات له دې فرضيې سره يو ځاے وکاروو چې
هغه سازګار دے، او داسې جوړښت
جوړ کړو چې په~\LR{$\Gamma$} کښې هر
تړلے فارمول پوره کړي۔ بالاخره موږ پوهېږو
چې د کوم ډول جوړښت په لټه کښې
يو: هغه چې هماغسې وي لکه~\LR{$\Gamma$} چې ئې بيانوي!{}
\olpseganchor{OLP-0128-B006}{end}

\olpseganchor{OLP-0128-B007}{start}
که~\LR{$\Gamma$} يوازې اتومي
  تړلي فارمولونه ولري، نو د هغې دپاره
مدل جوړول اسانه دي۔ فرض کړئ ټولې اتومي تړلي فارمولونه
  د~\LR{$\Atom{P}{a_1,\dots,a_n}$} بڼه لري، چې \LR{$a_i$} ګانې
  ثابت سمبولونه دي۔ يوازې دومره پکار دي چې
%
  د تعريف ساحه~\LR{$\Domain{M}$} او د~\LR{$P$} دپاره داسې ګومارنه پيدا کړو چې
  \LR{$\Sat{M}{\Atom{P}{a_1,\ldots,a_n}}$} وي۔ دا ډېره ګرانه نۀ ده:
  \LR{$\Domain{M} = \Nat$}، \LR{$\Assign{\Obj c_i}{M} = i$} وټاکئ، او د هر
  \LR{$\Atom{P}{a_1,\ldots,a_n} \in \Gamma$} دپاره جوړه
  \LR{$\tuple{k_1, \dots, k_n}$} په~\LR{$\Assign{P}{M}$} کښې واچوئ، چې
  \LR{$k_i$} د ثابت سمبول~\LR{$a_i$} شمېر دے (يعنې
  \LR{$a_i \ident \Obj c_{k_i}$})۔
\olpseganchor{OLP-0128-B007}{end}

\olpseganchor{OLP-0128-B008}{start}
اوس فرض کړئ~\LR{$\Gamma$} کوم داسې فارمول~\LR{$\lnot !B$} لري چې
\LR{$!B$} اتومي دے۔ ښايي اندېښنه وکړو چې د
\LR{$\Struct{M}$} جوړونه د~\LR{$\lnot !B$} د
رښتوني کولو له امکان سره ټکر کوي۔ خو دلته د~\LR{$\Gamma$} سازګاري
کار راځي: که \LR{$\lnot !B \in \Gamma$} وي، نو
\LR{$!B \notin \Gamma$}، کنه~\LR{$\Gamma$} به ناسازګار و۔ او که
\LR{$!B \notin \Gamma$} وي، نو زموږ د
\LR{$\Struct{M}$} د جوړونې له مخې
\LR{$\Sat/{M}{!B}$}؛ نو
\LR{$\Sat{M}{\lnot !B}$}۔ تر دې ځايه هر څه سم
دي۔
\olpseganchor{OLP-0128-B008}{end}

\olpseganchor{OLP-0128-B009}{start}
که~\LR{$\Gamma$} پېچلي، نااتومي فارمولونه ولري نو څه؟ فرض کړئ
\LR{$!A \land !B$} لري۔ د دې د رښتوني کولو دپاره بايد داسې پر مخ لاړ
شو لکه دواړه~\LR{$!A$} او~\LR{$!B$} چې په~\LR{$\Gamma$} کښې وي۔ او که
\LR{$!A \lor !B \in \Gamma$} وي، بايد لږ تر لږه يو ئې رښتونے کړو؛
يعنې داسې پر مخ لاړ شو لکه يو ئې چې په~\LR{$\Gamma$} کښې وي۔
\olpseganchor{OLP-0128-B009}{end}

\olpseganchor{OLP-0128-B010}{start}
دا لاندې مفکوره راښيي: په~\LR{$\Gamma$} کښې نور فارمولونه داسې
ورزياتوو چې (الف)~ترې جوړېدونکے سټ سازګار وساتو، (ب)~ډاډ ترلاسه
کړو چې د هرې ممکنې تړلے فارمول~\LR{$!A$} دپاره يا~\LR{$!A$} په
جوړ شوي سټ کښې وي يا~\LR{$\lnot !A$}، او (ج)~که هر کله
\LR{$!A \land !B$} په سټ کښې وي، دواړه~\LR{$!A$} او~\LR{$!B$} هم پکښې وي؛ که
\LR{$!A \lor !B$} په سټ کښې وي، لږ تر لږه يو له~\LR{$!A$} او~\LR{$!B$} څخه هم
پکښې وي؛ او داسې نور۔ \textbf{د ژباړې د نسخې سمون (OLCOM-001):}
سرچينې په شرط~(ب) کښې «اتومي جمله» ليکلې وه، خو همدا پاراګراف دا
شرط د بشپړ سټ تعريف بولي او ورپسې رسمي تعريف ئې د هرې جملې دپاره
ټاکي؛ نو ناسمه «اتومي» ځانګړنه لرې شوه۔ موږ دې کار ته (که اړينه
وي، تر ابده) دوام
ورکوو۔ د ټولو داسې زياتو شوو فارمولونه سټ~\LR{$\Gamma^*$} وبولئ۔
بيا به پورته جوړونه داسې
جوړښت~\LR{$\Struct{M}$}
راکړي چې په استقرا سره ثابتولے شو په~\LR{$\Gamma^*$} کښې ټولې جملې
پوره کوي، او ځکه چې \LR{$\Gamma \subseteq \Gamma^*$}، په~\LR{$\Gamma$}
کښې ټولې جملې هم پوره کوي۔ څرګندېږي چې د (الف) او~(ب) تضمين بس
دے۔ د جملو هغه سټ چې (ب) پوره کوي \emph{بشپړ} بلل کېږي۔ نو زموږ
کار به دا وي چې سازګار سټ~\LR{$\Gamma$} يو سازګار او بشپړ
سټ~\LR{$\Gamma^*$} ته وغځوو۔
\olpseganchor{OLP-0128-B010}{end}

\olpseganchor{OLP-0128-B011}{start}
%
په دې پلان کښې يوه ستونزه شته: که
\LR{$\lexists[x][!A(x)] \in \Gamma$} وي، غواړو په دې بهير کښې کوم
ثابت سمبول~\LR{$c$} وټاکو او~\LR{$!A(c)$} ورزيات کړو۔ خو څنګه پوهېږو چې
دا کار تل کېدے شي؟ ښايي زموږ په ژبه کښې يوازې څو ثابت سمبولونه
وي، او د هر يوه دپاره \LR{$\lnot !A(c) \in \Gamma$} ولرو۔ ورسره
\LR{$!A(c)$} هم نۀ شو زياتولے، ځکه سټ به ناسازګار کړي؛ بيا به نۀ پوهېدو
چې~\LR{$\Struct{M}$} بايد~\LR{$!A(c)$} رښتونے کړي که~\LR{$\lnot !A(c)$}۔ لا دا
هم کېدے شي چې~\LR{$\Gamma$} يوازې د داسې ژبې جملې ولري چې هېڅ ثابت
سمبول نۀ لري (د بېلګې په توګه، د سټونو د تيورۍ ژبه)۔
\olpseganchor{OLP-0128-B011}{end}

\olpseganchor{OLP-0128-B012}{start}
د دې ستونزې حل دا دے چې په پيل کښې بېخي نامتناهي ډېر ثابتونه، او
ورسره داسې جملې، زياتې کړو چې په سم ډول ئې له کميت ټاکونکو سره
وتړي۔ (البته بايد تصديق کړو چې دا کار ناسازګاري نۀ شي رامنځته
کولے۔)
\olpseganchor{OLP-0128-B012}{end}

\olpseganchor{OLP-0128-B013}{start}
زموږ لومړنۍ جوړونه هغه وخت ښه کار کوي چې په اتومي جملو کښې يوازې
ثابت سمبولونه ولرو۔ خو ژبه ښايي تابعې هم ولري۔ په هغه
حالت کښې ښايي پر~\LR{$\Nat$} د سمې تابعې موندل، چې دغو تابعې
ته ئې وګومارو او هر څه سم شي، ستونزمن وي۔ نو بله چاره دا ده: د
\LR{$\Obj c_i$} د تفسير دپاره د~\LR{$i$} د کارولو پر ځاے، د ثابت سمبولونه
خپله سټ د تعريف ساحه وټاکئ۔ بيا~\LR{$\Struct M$} هر ثابت سمبول خپل
ځان ته ګومارلے شي: \LR{$\Assign{\Obj c_i}{M} = \Obj c_i$}۔ خو ولې دا
کار تر پايه نۀ رسوو: د ژبې ټول \emph{ترمونه}~\LR{$\Domain{M}$} وټاکئ!{}
په دې صورت کښې تابعې ته د تابعو څرګنده ګومارنه شته (دغه
تابعې ترمونه د ارګومېنټونو په توګه اخلي او ترمونه د قيمتونو په توګه
لري): تابع~\LR{$\Obj f^n_i$} ته هغه تابع ګومارو چې \LR{$n$} ترمونه
\LR{$t_1$}، \dots،~\LR{$t_n$} د ننوت په توګه واخلي او ترم
\LR{$\Obj f^n_i(t_1, \dots, t_n)$} د قيمت په توګه ورکړي۔
\olpseganchor{OLP-0128-B013}{end}

\olpseganchor{OLP-0128-B014}{start}
د معما وروستۍ ټوټه دا ده چې له~\LR{$\eq$} سره څه وکړو۔
پريديکات~\LR{$\eq$} ثابت تفسير لري:
\LR{$\Sat{M}{\eq[t][t']}$} هله او يوازې هله چې
\LR{$\Value{t}{M} = \Value{t'}{M}$} وي۔ اوس که هر څه داسې وټاکو چې د
ترم~\LR{$t$} قيمت په خپله~\LR{$t$} وي، نو دا جوړښت به د
\LR{$\eq[t][t']$} په بڼه \emph{هېڅ} جمله رښتونې نۀ کړي، مګر دا چې
\LR{$t$} او~\LR{$t'$} يو او هماغه ترم وي۔ څرګنده ده چې دا ستونزه ده، ځکه د
تابعې لرونکې ژبې کابو هره په زړه پورې تيوري به داسې جملې
\LR{$\eq[t][t']$} د قضيو په توګه ولري چې~\LR{$t$} او~\LR{$t'$} يو ترم نۀ وي (د
بېلګې په توګه، د حساب په تيوريو کښې:
\LR{$\eq[(\Obj 0+ \Obj 0)][\Obj 0]$})۔ د دې ستونزې د هوارولو دپاره د
\LR{$\Struct M$} تعريف ساحه بدلوو: په~\LR{$\Domain{M}$} کښې د څيزونو په
توګه د ترمونو پر ځاے د ترمونو سټونه کاروو، او هر سټ ټول هغه ترمونه
لري چې د~\LR{$\Gamma$} جملې ئې برابر ګڼل غواړي۔ نو د بېلګې په توګه، که
\LR{$\Gamma$} د حساب تيوري وي، له دغو سټونو څخه يو به~\LR{$\Obj 0$}،
\LR{$(\Obj 0 + \Obj 0)$}، \LR{$(\Obj 0 \times \Obj 0)$}، او داسې نور ولري۔
همدا سټ به~\LR{$\Obj 0$} ته ګومارو، او څرګنده به شي چې همدا سټ په دې
کښې د ټولو ترمونو قيمت هم دے، لکه د~\LR{$(\Obj 0 + \Obj 0)$} قيمت۔
نو په دې بدله شوې جوړښت کښې به جمله
\LR{$\eq[(\Obj 0+ \Obj 0)][\Obj 0]$} رښتونې وي۔
\olpseganchor{OLP-0128-B014}{end}

\olpseganchor{OLP-0128-B015}{start}
نو دا کارونه به کوو۔ لومړے د بشپړ سازګارو سټونو خاصيتونه
څېړو؛ په ځانګړي ډول ثابتوو چې بشپړ سازګار سټ
\LR{$!A \land !B$} هله او يوازې هله لري چې دواړه~\LR{$!A$} او~\LR{$!B$} ولري،
او~\LR{$!A \lor !B$} هله او يوازې هله لري چې لږ تر لږه يو ئې ولري؛
او داسې نور (\olref[ccs]{prop:ccs})۔ بيا د جملو
  «مشبوع» سټونه تعريفوو او څېړو۔ مشبوع سټ داسې شرطي فارمولونه لري
  چې هره کميت ټاکونکې تړلے فارمول د هغې له بېلګو سره تړي
  (\olref[hen]{defn:henkin-exp})۔ ښيو چې هر سازګار سټ~\LR{$\Gamma$}
  تل مشبوع سټ~\LR{$\Gamma'$} ته غځېدے شي (\olref[hen]{lem:henkin})۔ که
  يو سټ سازګار، مشبوع او بشپړ وي، دا خاصيت هم لري چې
  \LR{$\lexists[x][!A(x)]$} هله او يوازې هله لري چې د کوم
    تړلي ترم~\LR{$t$} دپاره~\LR{$!A(t)$} ولري او
  \LR{$\lforall[x][!A(x)]$} هله او يوازې هله لري چې د
    ټولو تړلو ترمونو~\LR{$t$} دپاره~\LR{$!A(t)$} ولري
  (\olref[hen]{prop:saturated-instances})۔
بيا مشبوع سازګار سټ
\LR{$\Gamma'$} اخلو، او ښيو چې داسې
مشبوع، سازګار او بشپړ
سټ~\LR{$\Gamma^*$} ته غځېدے شي (\olref[lin]{lem:lindenbaum})۔ همدا
سټ~\LR{$\Gamma^*$} د خپل ترمي مدل~\LR{$\Struct M(\Gamma^*)$} د تعريف دپاره کاروو۔
د ترمي مدل تعريف ساحه د تړلو ترمونو سټ دے، او د هغه
  د پريديکاتونه تفسير د هغو اتومي
  تړلي فارمولونه له مخې ټاکل کېږي چې په~\LR{$\Gamma^*$} کښې دي
(\olref[mod]{defn:termmodel})۔ د مشبوع، بشپړو
سازګارو سټونو خاصيتونه کاروو، څو وښيو چې په رښتيا
\LR{$\Sat{M(\Gamma^*)}{!A}$} هله او
يوازې هله چې \LR{$!A \in \Gamma^*$}
(\olref[mod]{lem:truth})؛ نو په ځانګړي ډول
\LR{$\Sat{M(\Gamma^*)}{\Gamma}$}۔
په پاے کښې دا هم څېړو چې که~\LR{$\Gamma$}،~\LR{$\eq$} هم ولري،
  ترمي مدل څنګه تعريف کړو (\olref[ide]{defn:term-model-factor})،
  او ښيو چې هغه~\LR{$\Gamma^*$} پوره کوي (\olref[ide]{lem:truth})۔
\olpseganchor{OLP-0128-B015}{end}

\olpseganchor{OLP-0129-B005}{start}
\olfileid{fol}{com}{ccs}
\olpseganchor{OLP-0129-B005}{end}
      
\olpseganchor{OLP-0129-B006}{start}
\olsection{د تړلي فارمولونه بشپړ سازګار سټونه}
\olpseganchor{OLP-0129-B006}{end}

\olpseganchor{OLP-0129-B007}{start}
\begin{defn}[بشپړ سټ]
\ollabel{def:complete-set} د تړلي فارمولونه يو سټ~\LR{$\Gamma$} هله او
يوازې هله \emph{بشپړ} دے چې د هرې تړلے فارمول~\LR{$!A$}
دپاره يا~\LR{$!A \in \Gamma$} وي يا~\LR{$\lnot !A \in \Gamma$}۔
\end{defn}
\olpseganchor{OLP-0129-B007}{end}

\olpseganchor{OLP-0129-B008}{start}
\begin{explain}
د جملو بشپړ سټونه هېڅ پوښتنه بې ځوابه نۀ پرېږدي۔ د هرې
تړلے فارمول~\LR{$!A$} دپاره~\LR{$\Gamma$} «وايي» چې~\LR{$!A$} رښتونے دے که
کاذب۔ د بشپړ سټونو ارزښت د بشپړتيا د قضيې له ثبوت څخه هم
اوړي۔ د بېلګې په توګه، هغه تيوري چې بشپړ او د بديهي اصولو
له لارې څرګندېدونکې وي، تل د پرېکړې وړ وي۔
\end{explain}
\olpseganchor{OLP-0129-B008}{end}

\olpseganchor{OLP-0129-B009}{start}
\begin{explain}
بشپړ سازګار سټونه د بشپړتيا په ثبوت کښې مهم دي، ځکه
تضمينولے شو چې د تړلي فارمولونه هر سازګار سټ~\LR{$\Gamma$} په
بشپړ سازګار سټ~\LR{$\Gamma^*$} کښې شامل دے۔ بشپړ
سازګار سټ د هرې تړلے فارمول~\LR{$!A$} دپاره يا~\LR{$!A$} لري يا د هغې
نفي~\LR{$\lnot !A$}، خو دواړه نۀ لري۔ دا په ځانګړي ډول د
اتومي تړلي فارمولونه دپاره هم سم دي؛
نو له بشپړ سازګار سټ څخه په داسې ژبه کښې چې
ثابت سمبولونه په مناسبه توګه ورزيات شوي وي، داسې
جوړښت جوړولے شو چې پکښې د
پريديکاتونه تفسير د هغو
اتومي تړلي فارمولونه له
مخې تعريفېږي چې په~\LR{$\Gamma^*$} کښې دي۔ بيا ښودل کېدے شي چې دا
جوړښت په~\LR{$\Gamma^*$} کښې ټولې
تړلي فارمولونه (او ځکه په~\LR{$\Gamma$} کښې ټولې) رښتونې کوي۔ د دې
وروستي حقيقت ثبوت دا غواړي چې
\LR{$\lnot !A \in \Gamma^*$} هله او يوازې هله چې
\LR{$!A \notin \Gamma^*$}، او~\LR{$(!A \lor !B) \in \Gamma^*$} هله او
يوازې هله چې \LR{$!A \in \Gamma^*$} يا~\LR{$!B \in \Gamma^*$}؛ او داسې
نور۔
\end{explain}
\olpseganchor{OLP-0129-B009}{end}

\olpseganchor{OLP-0129-B010}{start}
په لاندې بحث کښې به ډېر ځله د~\LR{$\Proves$} د انعکاسيت، يکنواختۍ او
انتقاليت خاصيتونه په ناڅرګند ډول کاروو (وګورئ
%
  \tagrefs{prfSC/{fol:seq:ptn:sec},prfND/{fol:ntd:ptn:sec},prfAX/{fol:axd:ptn:sec},prfTab/{fol:tab:ptn:sec}})۔
\olpseganchor{OLP-0129-B010}{end}

\olpseganchor{OLP-0129-B011}{start}
\begin{prop}
\ollabel{prop:ccs}
فرض کړئ~\LR{$\Gamma$} بشپړ او سازګار دے۔ بيا:
\begin{enumerate}
\item \ollabel{prop:ccs-prov-in} که \LR{$\Gamma \Proves !A$} وي، نو
  \LR{$!A \in \Gamma$}۔
\olpseganchor{OLP-0129-B011}{end}

\olpseganchor{OLP-0129-B012}{start}
\item \ollabel{prop:ccs-and} \LR{$!A \land !B \in \Gamma$}
  هله او يوازې هله چې دواړه \LR{$!A \in \Gamma$} او
  \LR{$!B \in \Gamma$} وي۔
\olpseganchor{OLP-0129-B012}{end}

\olpseganchor{OLP-0129-B013}{start}
\item \ollabel{prop:ccs-or} \LR{$!A \lor !B \in \Gamma$} هله
  او يوازې هله چې يا~\LR{$!A \in \Gamma$} وي يا~\LR{$!B \in \Gamma$}۔
\olpseganchor{OLP-0129-B013}{end}

\olpseganchor{OLP-0129-B014}{start}
\item \ollabel{prop:ccs-if} \LR{$!A \lif !B \in \Gamma$} هله
  او يوازې هله چې يا~\LR{$!A \notin \Gamma$} وي يا~\LR{$!B \in \Gamma$}۔
\end{enumerate}
\end{prop}
\olpseganchor{OLP-0129-B014}{end}

\olpseganchor{OLP-0129-B015}{start}
\begin{proof}
په لاندې ټولو برخو کښې فرض کوو چې~\LR{$\Gamma$} بشپړ او
سازګار دے۔
\begin{enumerate}
\item که \LR{$\Gamma \Proves !A$} وي، نو \LR{$!A \in \Gamma$}۔
\olpseganchor{OLP-0129-B015}{end}

\olpseganchor{OLP-0129-B016}{start}
فرض کړئ \LR{$\Gamma \Proves !A$}۔ د نقيض په توګه فرض کړئ چې
\LR{$!A \notin \Gamma$}۔ ځکه~\LR{$\Gamma$} بشپړ دے، نو
\LR{$\lnot !A \in \Gamma$}۔ د
%
  \tagrefs{prfSC/{fol:seq:prv:prop:explicit-inc},
    prfND/{fol:ntd:prv:prop:explicit-inc},
    prfAX/{fol:axd:prv:prop:explicit-inc},
    prfTab/{fol:tab:prv:prop:explicit-inc}}
له مخې~\LR{$\Gamma$} ناسازګار دے۔ دا له دې فرضيې سره تناقض دے چې
\LR{$\Gamma$} سازګار دے۔ نو \LR{$!A \notin \Gamma$} نۀ شي کېدے، او ځکه
\LR{$!A \in \Gamma$}۔
\olpseganchor{OLP-0129-B016}{end}

\olpseganchor{OLP-0129-B017}{start}
\item %
%
\LR{$!A \land !B \in \Gamma$} هله او يوازې هله چې دواړه
\LR{$!A \in \Gamma$} او~\LR{$!B \in \Gamma$} وي:
\olpseganchor{OLP-0129-B017}{end}

\olpseganchor{OLP-0129-B018}{start}
د مخامخ لوري دپاره، فرض کړئ \LR{$!A \land !B \in \Gamma$}۔ بيا د
%
  \tagrefs{prfSC/{fol:seq:ppr:prop:provability-land},
    prfND/{fol:ntd:ppr:prop:provability-land},
    prfAX/{fol:axd:ppr:prop:provability-land},
    prfTab/{fol:tab:ppr:prop:provability-land}}، د توکي~(۱) له
مخې \LR{$\Gamma \Proves !A$} او~\LR{$\Gamma \Proves !B$}۔ د
\olref{prop:ccs-prov-in} له مخې \LR{$!A \in \Gamma$} او
\LR{$!B \in \Gamma$}، لکه څنګه چې پکار وو۔
\olpseganchor{OLP-0129-B018}{end}

\olpseganchor{OLP-0129-B019}{start}
د عکس لوري دپاره، \LR{$!A \in \Gamma$} او~\LR{$!B \in \Gamma$} ونيسئ۔ د
%
  \tagrefs{prfSC/{fol:seq:ppr:prop:provability-land},
    prfND/{fol:ntd:ppr:prop:provability-land},
    prfAX/{fol:axd:ppr:prop:provability-land},
    prfTab/{fol:tab:ppr:prop:provability-land}}، د توکي~(۲) له
مخې \LR{$\Gamma \Proves !A \land !B$}۔ د
\olref{prop:ccs-prov-in} له مخې \LR{$!A \land !B \in \Gamma$}۔
\olpseganchor{OLP-0129-B019}{end}

\olpseganchor{OLP-0129-B020}{start}
\item %
%
لومړے ښيو چې که \LR{$!A \lor !B \in \Gamma$} وي، نو يا
\LR{$!A \in \Gamma$} وي يا~\LR{$!B \in \Gamma$}۔ فرض کړئ
\LR{$!A \lor !B \in \Gamma$}، خو \LR{$!A \notin \Gamma$} او
\LR{$!B \notin \Gamma$}۔ ځکه~\LR{$\Gamma$} بشپړ دے، نو
\LR{$\lnot !A \in \Gamma$} او~\LR{$\lnot !B \in \Gamma$}۔ د
%
  \tagrefs{prfSC/{fol:seq:ppr:prop:provability-lor},
    prfND/{fol:ntd:ppr:prop:provability-lor},
    prfAX/{fol:axd:ppr:prop:provability-lor},
    prfTab/{fol:tab:ppr:prop:provability-lor}}، د توکي~(۱) له
مخې~\LR{$\Gamma$} ناسازګار دے؛ دا تناقض دے۔ نو يا
\LR{$!A \in \Gamma$} وي يا~\LR{$!B \in \Gamma$}۔
\olpseganchor{OLP-0129-B020}{end}

\olpseganchor{OLP-0129-B021}{start}
د عکس لوري دپاره، فرض کړئ يا~\LR{$!A \in \Gamma$} وي يا
\LR{$!B \in \Gamma$}۔ د
%
  \tagrefs{prfSC/{fol:seq:ppr:prop:provability-lor},
    prfND/{fol:ntd:ppr:prop:provability-lor},
    prfAX/{fol:axd:ppr:prop:provability-lor},
    prfTab/{fol:tab:ppr:prop:provability-lor}}، د توکي~(۲) له
مخې \LR{$\Gamma \Proves !A \lor !B$}۔ د
\olref{prop:ccs-prov-in} له مخې \LR{$!A \lor !B \in \Gamma$}، لکه
څنګه چې پکار وو۔
\olpseganchor{OLP-0129-B021}{end}

\olpseganchor{OLP-0129-B022}{start}
\item %
%
د مخامخ لوري دپاره، فرض کړئ \LR{$!A \lif !B \in \Gamma$}، او د نقيض
په توګه فرض کړئ چې \LR{$!A \in \Gamma$} او~\LR{$!B \notin \Gamma$}۔ د
دغو فرضونو له مخې \LR{$!A \lif !B \in \Gamma$} او
\LR{$!A \in \Gamma$}۔ د
%
  \tagrefs{prfSC/{fol:seq:ppr:prop:provability-lif},
    prfND/{fol:ntd:ppr:prop:provability-lif},
    prfAX/{fol:axd:ppr:prop:provability-lif},
    prfTab/{fol:tab:ppr:prop:provability-lif}}، د توکي~(۱) له
مخې \LR{$\Gamma \Proves !B$}۔ خو بيا د
\olref{prop:ccs-prov-in} له مخې \LR{$!B \in \Gamma$}؛ دا له
\LR{$!B \notin \Gamma$} سره تناقض دے۔
\olpseganchor{OLP-0129-B022}{end}

\olpseganchor{OLP-0129-B023}{start}
د عکس لوري دپاره، لومړے هغه حالت وڅېړئ چې
\LR{$!A \notin \Gamma$} وي۔ ځکه~\LR{$\Gamma$} بشپړ دے، نو
\LR{$\lnot !A \in \Gamma$}۔ د
%
  \tagrefs{prfSC/{fol:seq:ppr:prop:provability-lif},
    prfND/{fol:ntd:ppr:prop:provability-lif},
    prfAX/{fol:axd:ppr:prop:provability-lif},
    prfTab/{fol:tab:ppr:prop:provability-lif}}، د توکي~(۲) له
مخې \LR{$\Gamma \Proves !A \lif !B$}۔ بيا د
\olref{prop:ccs-prov-in} له مخې \LR{$!A \lif !B \in \Gamma$}، لکه
څنګه چې پکار وو۔
\olpseganchor{OLP-0129-B023}{end}

\olpseganchor{OLP-0129-B024}{start}
اوس هغه حالت وڅېړئ چې \LR{$!B \in \Gamma$} وي۔ بيا د
%
  \tagrefs{prfSC/{fol:seq:ppr:prop:provability-lif},
    prfND/{fol:ntd:ppr:prop:provability-lif},
    prfTab/{fol:tab:ppr:prop:provability-lif},
    prfAX/{fol:axd:ppr:prop:provability-lif}}، يو ځل بيا د
توکي~(۲) له مخې \LR{$\Gamma \Proves !A \lif !B$}۔ د
\olref{prop:ccs-prov-in} له مخې \LR{$!A \lif !B \in \Gamma$}۔
\end{enumerate}
\end{proof}
\olpseganchor{OLP-0129-B024}{end}

\olpseganchor{OLP-0129-B025}{start}
\begin{prob}
د~\olref[fol][com][ccs]{prop:ccs} ثبوت بشپړ کړئ۔
\end{prob}
\olpseganchor{OLP-0129-B025}{end}

\olpseganchor{OLP-0130-B004}{start}
\olfileid{fol}{com}{hen}
\olpseganchor{OLP-0130-B004}{end}

\olpseganchor{OLP-0130-B005}{start}
\olsection{د هنکين غځونه}
\olpseganchor{OLP-0130-B005}{end}

\olpseganchor{OLP-0130-B006}{start}
\begin{explain}
د بشپړتيا د قضيې د ثابتولو د ستونزې يوه برخه دا ده چې هغه مدل چې
له بشپړ سازګار سټ~\LR{$\Gamma$} څخه ئې جوړوو بايد په~\LR{$\Gamma$} کښې ټول
کميت ټاکونکي لرونکي فارمولونه رښتوني کړي۔ د دې د تضمين دپاره د
ليون هنکين يو چل کاروو۔ په اصل کښې دا چل د ژبې په لامتناهي ډېرو
ثابت سمبولونه غځول او د هر داسې فارمول~\LR{$!A(x)$} دپاره چې يو
ازاد متغير لري، د لاندې بڼې يو فارمول ورزياتول دي:
\LR{$\lexists[x][!A(x)] \lif !A(c)$}،
چې~\LR{$c$} يو له نوو ثابت سمبولونه څخه دے۔ کله چې هغه جوړښت
جوړوو چې~\LR{$\Gamma$} پوره کوي، دا به تضمين کړي چې د نوو ثابت سمبولونو
په منځ کښې هره
رښتونې وجودي جمله يو شاهد ولري۔
\end{explain}
\olpseganchor{OLP-0130-B006}{end}

\olpseganchor{OLP-0130-B007}{start}
\begin{prop}
\ollabel{prop:lang-exp}
که~\LR{$\Gamma$} په~\LR{$\Lang L$} کښې سازګار وي او~\LR{$\Lang L'$} له~\LR{$\Lang L$}
څخه د نوو ثابت سمبولونه د~شمېرېدونکے لامتناهي سټ~\LR{$\Obj d_0$}،
\LR{$\Obj d_1$}، \dots په ورزياتولو ترلاسه شي، نو~\LR{$\Gamma$} په~\LR{$\Lang L'$}
کښې سازګار دے۔
\end{prop}
\olpseganchor{OLP-0130-B007}{end}

\olpseganchor{OLP-0130-B008}{start}
\begin{defn}[مشبوع سټ]
د ژبې~\LR{$\Lang {L}$} د فارمولونه يو سټ~\LR{$\Gamma$} هله او يوازې هله
\emph{مشبوع} دے چې د هر داسې فارمول~\LR{$!A(x) \in \Frm[L]$}
دپاره چې يو ازاد متغير~\LR{$x$} لري، داسې ثابت سمبول
\LR{$c \in \Lang{L}$} موجود وي چې
\LR{$\lexists[x][!A(x)] \lif !A(c) \in \Gamma$}۔
\end{defn}
\olpseganchor{OLP-0130-B008}{end}

\olpseganchor{OLP-0130-B009}{start}
لاندې تعريف به د ورپسې قضيې په ثبوت کښې وکارول شي۔
\olpseganchor{OLP-0130-B009}{end}

\olpseganchor{OLP-0130-B010}{start}
\begin{defn}
\ollabel{defn:henkin-exp}
فرض کړئ~\LR{$\Lang L'$} هماغسې وي لکه په~\olref{prop:lang-exp} کښې۔
د~\LR{$\Lang L'$} د ټولو هغو فارمولونه~\LR{$!A_i(x_i)$} يوه شمېرنه
\LR{$!A_0(x_0)$}، \LR{$!A_1(x_1)$}، \dots وټاکئ چې يو متغير~(\LR{$x_i$}) پکښې
ازاد راځي۔ تړلي فارمولونه~\LR{$!D_n$} پر~\LR{$n$} د استقرا له لارې تعريفوو۔
\olpseganchor{OLP-0130-B010}{end}

\olpseganchor{OLP-0130-B011}{start}
\LR{$c_0$} دې د هغو~\LR{$\Obj d_i$} په منځ کښې لومړنے ثابت سمبول وي چې
مونږ~\LR{$\Lang{L}$} ته ورزيات کړي دي او په~\LR{$!A_0(x_0)$} کښې نۀ راځي۔
په دې فرض چې~\LR{$!D_0$}، \dots،~\LR{$!D_{n-1}$} له مخکښې تعريف شوي دي،
\LR{$c_n$} دې د نوو ثابت سمبولونه~\LR{$\Obj d_i$} په منځ کښې لومړنے هغه وي
چې نۀ په~\LR{$!D_0$}، \dots،~\LR{$!D_{n-1}$} کښې راځي او نۀ په~\LR{$!A_n(x_n)$}
کښې۔
\olpseganchor{OLP-0130-B011}{end}

\olpseganchor{OLP-0130-B012}{start}
اوس~\LR{$!D_{n}$} دا فارمول وټاکئ: \LR{$\lexists[x_{n}][!A_{n}(x_{n})] \lif
  !A_{n}(c_{n})$}۔
\end{defn}
\olpseganchor{OLP-0130-B012}{end}

\olpseganchor{OLP-0130-B013}{start}
\begin{lem}
\ollabel{lem:henkin}
هر سازګار سټ~\LR{$\Gamma$} يوه مشبوع سازګار سټ~\LR{$\Gamma'$} ته غځېدے
شي۔
\end{lem}
\olpseganchor{OLP-0130-B013}{end}

\olpseganchor{OLP-0130-B014}{start}
\begin{proof}
د ژبې~\LR{$\Lang{L}$} د جملو يو سازګار سټ~\LR{$\Gamma$} راکړل
شوے دے؛ ژبه د نوو ثابت سمبولونه د~شمېرېدونکے لامتناهي سټ په
ورزياتولو وغځوئ، څو~\LR{$\Lang{L'}$} جوړ شي۔ د~\olref{prop:lang-exp}
له مخې~\LR{$\Gamma$} لا هم په بډايه ژبه کښې سازګار دے۔ همدارنګه،
\LR{$!D_i$} دې هماغسې وي لکه په~\olref{defn:henkin-exp} کښې۔ وټاکئ:
\begin{align*}
\Gamma_0 & = \Gamma \\
\Gamma_{n+1} & = \Gamma_n \cup \{!D_n \}
\end{align*}
يعنې~\LR{$\Gamma_{n+1} = \Gamma \cup \{ !D_0, \dots, !D_n \}$}، او
\LR{$\Gamma' = \bigcup_{n} \Gamma_n$} وټاکئ۔ \LR{$\Gamma'$} په څرګند ډول
مشبوع دے۔
\olpseganchor{OLP-0130-B014}{end}

\olpseganchor{OLP-0130-B015}{start}
که~\LR{$\Gamma'$} ناسازګار وے، نو د کوم~\LR{$n$} دپاره به~\LR{$\Gamma_n$}
ناسازګار وے (تمرين: څرګنده کړئ چې ولې)۔ نو د~\LR{$\Gamma'$} د سازګارۍ
د ښودلو دپاره دا بسنه کوي چې پر~\LR{$n$} د استقرا له لارې وښيو چې هر
سټ~\LR{$\Gamma_n$} سازګار دے۔
\olpseganchor{OLP-0130-B015}{end}

\olpseganchor{OLP-0130-B016}{start}
د استقرا بنسټ يوازې دا ادعا ده چې~\LR{$\Gamma_0 = \Gamma$} سازګار دے،
او دا د قضيې فرضيه ده۔ د استقرا د ګام دپاره، فرض کړئ چې
\LR{$\Gamma_{n}$} سازګار دے، خو~\LR{$\Gamma_{n+1} =
\Gamma_n \cup \{!D_n\}$} ناسازګار دے۔ در په ياد کړئ چې~\LR{$!D_n$} دا
دے:
\LR{$\lexists[x_{n}][!A_{n}(x_n)] \lif !A_{n}(c_{n})$}،
چې~\LR{$!A_n(x_n)$} د~\LR{$\Lang{L'}$} داسې فارمول دے چې يوازې متغير
\LR{$x_n$} ئې ازاد دے۔ د~\LR{$c_n$} د غوره کولو د طريقې له مخې (وګورئ
\olref{defn:henkin-exp})، \LR{$c_n$} نۀ په~\LR{$!A_n(x_n)$} کښې راځي او نۀ
په~\LR{$\Gamma_n$} کښې۔
\olpseganchor{OLP-0130-B016}{end}

\olpseganchor{OLP-0130-B017}{start}
که~\LR{$\Gamma_n \cup \{!D_n\}$} ناسازګار وي، نو
\LR{$\Gamma_n \Proves \lnot !D_n$}، او ځکه لاندې دواړه سم دي:

\[
\Gamma_n \Proves \lexists[x_n][!A_n(x_n)]
\qquad
\Gamma_n \Proves \lnot !A_n(c_n)
\]
ځکه چې~\LR{$c_n$} نۀ په~\LR{$\Gamma_n$} کښې راځي او نۀ په~\LR{$!A_n(x_n)$}
کښې، د قوي تعميم قضيه
\tagrefs{prfAX/{fol:axd:qpr:thm:strong-generalization},prfSC/{fol:seq:qpr:thm:strong-generalization},prfND/{fol:ntd:qpr:thm:strong-generalization},prfTab/{fol:tab:qpr:thm:strong-generalization}}
پلې کېږي۔ له \LR{$\Gamma_n \Proves \lnot !A_n(c_n)$}
څخه دا ترلاسه کوو:
\LR{$\Gamma_n \Proves \lforall[x_n][\lnot !A_n(x_n)]$}۔
نو دواړه دا لرو:
\LR{$\Gamma_n \Proves \lexists[x_n][!A_n(x_n)]$} او
\LR{$\Gamma_n \Proves \lforall[x_n][\lnot !A_n(x_n)]$}؛
نو~\LR{$\Gamma_n$} پخپله ناسازګار دے۔
(پام وکړئ چې
\LR{$\lforall[x_n][\lnot !A_n(x_n)] \Proves
  \lnot\lexists[x_n][!A_n(x_n)]$}۔)
تناقض: فرض دا و چې~\LR{$\Gamma_n$} سازګار دے۔ نو
\LR{$\Gamma_n \cup \{ !D_n\}$} سازګار دے۔
\end{proof}
\olpseganchor{OLP-0130-B017}{end}

\olpseganchor{OLP-0130-B018}{start}
\begin{explain}
اوس به وښيو چې مشبوع، سازګار او \emph{بشپړ} سټونه دا
خاصيت لري چې يوه عمومي کميت ټاکونکې تړلے فارمول
  هله او يوازې هله لري چې د هغې ټولې بېلګې ولري
  او يوه وجودي کميت ټاکونکې تړلے فارمول هله او
  يوازې هله لري چې لږ تر لږه د هغې يوه بېلګه ولري۔ له دې څخه به
د دې ښودلو دپاره کار واخلو چې هغه جوړښت چې له مشبوع،
سازګار او بشپړ سټ څخه ئې جوړوو، د هغه ټولې کميت ټاکونکې جملې
رښتونې کوي۔
\end{explain}
\olpseganchor{OLP-0130-B018}{end}

\olpseganchor{OLP-0130-B019}{start}
\begin{prop}\ollabel{prop:saturated-instances}
فرض کړئ~\LR{$\Gamma$} بشپړ، سازګار او مشبوع دے۔
\begin{enumerate}
\item \LR{$\lexists[x][!A(x)] \in \Gamma$} هله او يوازې هله
  چې لږ تر لږه د يوه تړلي ترم~\LR{$t$} دپاره~\LR{$!A(t) \in \Gamma$} وي۔
\item \LR{$\lforall[x][!A(x)] \in \Gamma$} هله او يوازې هله
  چې د ټولو تړلو ترمونو~\LR{$t$} دپاره~\LR{$!A(t) \in \Gamma$} وي۔
\end{enumerate}
\end{prop}
\olpseganchor{OLP-0130-B019}{end}

\olpseganchor{OLP-0130-B020}{start}
\begin{proof}
  \begin{enumerate}
  \item %
    لومړے فرض کړئ چې
      \LR{$\lexists[x][!A(x)] \in \Gamma$}۔ ځکه~\LR{$\Gamma$} مشبوع دے، نو
      د کوم ثابت سمبول~\LR{$c$} دپاره
      \LR{$(\lexists[x][!A(x)] \lif !A(c)) \in \Gamma$}۔ د
      \tagrefs{prfAX/{fol:axd:ppr:prop:provability-lif},prfSC/{fol:seq:ppr:prop:provability-lif},prfND/{fol:ntd:ppr:prop:provability-lif},prfTab/{fol:tab:ppr:prop:provability-lif}}، د توکي~(۱)،
      او \olref[ccs]{prop:ccs}\olref[ccs]{prop:ccs-prov-in} له مخې
      \LR{$!A(c) \in \Gamma$}۔
\olpseganchor{OLP-0130-B020}{end}

\olpseganchor{OLP-0130-B021}{start}
    د بل لوري دپاره مشبوع والی لازم نۀ دے: فرض کړئ
    \LR{$!A(t) \in \Gamma$}۔ بيا د
      \tagrefs{prfAX/{fol:axd:qpr:prop:provability-quantifiers},prfSC/{fol:seq:qpr:prop:provability-quantifiers},prfND/{fol:ntd:qpr:prop:provability-quantifiers},prfTab/{fol:tab:qpr:prop:provability-quantifiers}}، د توکي~(۱) له مخې
    \LR{$\Gamma \Proves \lexists[x][!A(x)]$}۔ د
    \olref[ccs]{prop:ccs}\olref[ccs]{prop:ccs-prov-in} له مخې
    \LR{$\lexists[x][!A(x)] \in \Gamma$}۔
  \item %
    فرض کړئ چې د ټولو تړلو ترمونو~\LR{$t$}
      دپاره~\LR{$!A(t) \in \Gamma$}۔ د نقيض په توګه فرض کړئ چې
      \LR{$\lforall[x][!A(x)] \notin \Gamma$}۔ ځکه~\LR{$\Gamma$} بشپړ دے،
      نو~\LR{$\lnot\lforall[x][!A(x)] \in \Gamma$}۔ د مشبوع والي له مخې،
      د کوم ثابت سمبول~\LR{$c$} دپاره
      \LR{$(\lexists[x][\lnot !A(x)] \lif \lnot !A(c)) \in
        \Gamma$}۔ د فرض له مخې، ځکه چې~\LR{$c$} يو تړلے ترم دے، نو
      \LR{$!A(c) \in \Gamma$}۔ خو له دې څخه به~\LR{$\Gamma$} ناسازګار شي۔
      (تمرين: هغه اشتقاق ورکړئ چې ښيي
      \[
      \lnot \lforall[x][!A(x)],
      \lexists[x][\lnot !A(x)] \lif \lnot
        !A(c), !A(c)
      \]
      ناسازګار دے۔)
\olpseganchor{OLP-0130-B021}{end}

\olpseganchor{OLP-0130-B022}{start}
      د عکس لوري دپاره مشبوع والی لازم نۀ دے: فرض کړئ
      \LR{$\lforall[x][!A(x)] \in \Gamma$}۔ بيا د
      \tagrefs{prfAX/{fol:axd:qpr:prop:provability-quantifiers},prfSC/{fol:seq:qpr:prop:provability-quantifiers},prfND/{fol:ntd:qpr:prop:provability-quantifiers},prfTab/{fol:tab:qpr:prop:provability-quantifiers}}،
      د توکي~(۲) له مخې~\LR{$\Gamma \Proves !A(t)$}۔ د
      \olref[ccs]{prop:ccs} له مخې~\LR{$!A(t) \in \Gamma$}۔
  \end{enumerate}
\end{proof}
\olpseganchor{OLP-0130-B022}{end}

\olpseganchor{OLP-0131-B004}{start}
\olfileid{fol}{com}{lin}
\olpseganchor{OLP-0131-B004}{end}

\olpseganchor{OLP-0131-B005}{start}
\olsection{د لينډنباوم لمه}
\olpseganchor{OLP-0131-B005}{end}

\olpseganchor{OLP-0131-B006}{start}
\begin{explain}
اوس يوه مرستندويه قضيه ثابتوو چې ښيي د تړلي فارمولونه هر سازګار
سټ د جملو په داسې سټ کښې شامل دے چې يوازې سازګار نۀ، بلکې
بشپړ هم وي۔ ثبوت په هر ځل د يوې تړلے فارمول په زياتولو
کار کوي، او په هر پړاو کښې تضمينوي چې سټ سازګار پاتې شي۔ دا ځکه
کوو چې د هر~\LR{$!A$} دپاره يا~\LR{$!A$} يا~\LR{$\lnot !A$} په کوم پړاو کښې
ورزيات شي۔ بيا د دې جوړونې د ټولو پړاوونو اتحاد يا~\LR{$!A$} لري يا
ئې نفي~\LR{$\lnot !A$}، او په دې توګه بشپړ دے۔ دا سازګار هم دے، ځکه
په هر پړاو کښې پام کوو چې ناسازګاري رامنځته نۀ کړو۔
\end{explain}
\olpseganchor{OLP-0131-B006}{end}

\olpseganchor{OLP-0131-B007}{start}
\begin{lem}[د لينډنباوم لمه]
\ollabel{lem:lindenbaum} په ژبه~\LR{$\Lang{L}$} کښې هر سازګار
سټ~\LR{$\Gamma$} داسې غځېدے شي چې يو بشپړ او سازګار
سټ~\LR{$\Gamma^*$} شي۔
\end{lem}
\olpseganchor{OLP-0131-B007}{end}

\olpseganchor{OLP-0131-B008}{start}
\begin{proof}
فرض کړئ~\LR{$\Gamma$} سازګار دے۔ د~\LR{$\Lang L$} د ټولو تړلي فارمولونه يوه
شمېرنه~\LR{$!A_0$}، \LR{$!A_1$}، \dots{} وي۔ \LR{$\Gamma_0 = \Gamma$} تعريف
کړئ، او
\[
\Gamma_{n+1} =
\begin{cases}
\Gamma_n \cup \{ !A_n \} & \textrm{که \LR{$\Gamma_n \cup \{!A_n\}$}
  سازګار وي؛} \\
\Gamma_n \cup \{ \lnot !A_n \} & \textrm{په بل صورت کښې۔}
\end{cases}
\]
\LR{$\Gamma^* = \bigcup_{n \geq 0} \Gamma_n$} وټاکئ۔
\olpseganchor{OLP-0131-B008}{end}

\olpseganchor{OLP-0131-B009}{start}
هر~\LR{$\Gamma_n$} سازګار دے: \LR{$\Gamma_0$} د تعريف او فرض له مخې سازګار
دے۔ که~\LR{$\Gamma_{n+1} = \Gamma_n \cup \{!A_n\}$} وي، لامل ئې دا
دے چې وروستنے سټ سازګار دے۔ که دا سازګار نۀ وي، نو
\LR{$\Gamma_{n+1} = \Gamma_n \cup \{\lnot !A_n\}$}۔ بايد تصديق کړو
چې~\LR{$\Gamma_n \cup \{\lnot !A_n\}$} سازګار دے۔ فرض کړئ چې سازګار
نۀ دے۔ بيا \emph{دواړه} \LR{$\Gamma_n \cup \{!A_n\}$} او
\LR{$\Gamma_n \cup \{\lnot !A_n\}$} ناسازګار دي۔ دا په دې معنا ده چې
\LR{$\Gamma_n$} به د
%
  \tagrefs{prfAX/{fol:axd:prv:prop:provability-exhaustive},
    prfSC/{fol:seq:prv:prop:provability-exhaustive},
    prfND/{fol:ntd:prv:prop:provability-exhaustive},
    prfTab/{fol:tab:prv:prop:provability-exhaustive}}
له مخې ناسازګار وے، چې دا د استقرا له فرضيې سره تناقض دے۔
\olpseganchor{OLP-0131-B009}{end}

\olpseganchor{OLP-0131-B010}{start}
د هر~\LR{$n$} او هر~\LR{$i < n$} دپاره، \LR{$\Gamma_i \subseteq \Gamma_n$}۔
دا پر~\LR{$n$} د ساده استقرا له لارې ترلاسه کېږي۔ د~\LR{$n=0$} دپاره هېڅ
\LR{$i < 0$} نشته، نو ادعا په خپله سمه ده۔ د استقرا د ګام دپاره فرض
کړئ چې ادعا د~\LR{$n$} دپاره سمه ده۔ ښيو چې که~\LR{$i < n+1$} وي، نو
\LR{$\Gamma_i \subseteq \Gamma_{n+1}$}۔ د جوړونې له مخې يا
\LR{$\Gamma_{n+1} = \Gamma_n \cup \{!A_n\}$} ده يا
\LR{$= \Gamma_n \cup \{\lnot !A_n\}$}۔ نو
\LR{$\Gamma_n \subseteq \Gamma_{n+1}$}۔ که~\LR{$i < n+1$} وي، نو د
استقرا د فرضيې له مخې~\LR{$\Gamma_i \subseteq \Gamma_n$} (که~\LR{$i < n$}
وي)، يا دا ساده حقيقت لرو چې~\LR{$\Gamma_n \subseteq \Gamma_n$}
(که~\LR{$i = n$} وي)۔ د~\LR{$\subseteq$} د انتقاليت له مخې
\LR{$\Gamma_i \subseteq \Gamma_{n+1}$} ترلاسه کوو۔
\olpseganchor{OLP-0131-B010}{end}

\olpseganchor{OLP-0131-B011}{start}
له دې څخه څرګندېږي چې~\LR{$\Gamma^*$} سازګار دے۔ لامل دا دے: فرض کړئ
\LR{$\Gamma' \subseteq \Gamma^*$} متناهي دے۔
\textbf{د ژباړې د نسخې سمون (OLCOM-002):} سرچينې د تش متناهي
فرعي سټ حالت نۀ و بېل کړے، حال دا چې «تر ټولو لوے» پړاو يوازې د
نا تش سټ دپاره ټاکل کېدے شي؛ د تش حالت څرګنده سازګاري ورزياته
شوه۔ نو فرض کړئ چې دا فرعي سټ نا تش دے۔ هر~\LR{$!B \in \Gamma'$} د
کوم~\LR{$i$} دپاره په~\LR{$\Gamma_i$} کښې هم دے۔ \LR{$n$} دې د دغو شاخصونو تر
ټولو لوے وي۔ ځکه که~\LR{$i \le n$} وي نو
\LR{$\Gamma_i \subseteq \Gamma_n$}، هر~\LR{$!B \in \Gamma'$} هم
\LR{$\in \Gamma_n$} دے؛ يعنې~\LR{$\Gamma' \subseteq \Gamma_n$}، او
\LR{$\Gamma_n$} سازګار دے۔ نو هر متناهي فرعي سټ
\LR{$\Gamma' \subseteq \Gamma^*$} سازګار دے۔ د
%
  \tagrefs{prfAX/{fol:axd:ptn:prop:proves-compact},
    prfSC/{fol:seq:ptn:prop:proves-compact},
    prfND/{fol:ntd:ptn:prop:proves-compact},
    prfTab/{fol:tab:ptn:prop:proves-compact}} له مخې، \LR{$\Gamma^*$}
سازګار دے۔
\olpseganchor{OLP-0131-B011}{end}

\olpseganchor{OLP-0131-B012}{start}
د~\LR{$\Frm[L]$} هره تړلے فارمول په هغه لړ کښې راځي چې د~\LR{$\Gamma^*$}
د تعريف دپاره کارول شوے دے۔ که~\LR{$!A_n \notin \Gamma^*$} وي، لامل
ئې دا دے چې~\LR{$\Gamma_n \cup \{!A_n\}$} ناسازګار و۔ خو بيا
\LR{$\lnot !A_n \in \Gamma^*$}، نو~\LR{$\Gamma^*$} بشپړ دے۔
\end{proof}
\olpseganchor{OLP-0131-B012}{end}

\olpseganchor{OLP-0132-B004}{start}
\olfileid{fol}{com}{mod}
\olpseganchor{OLP-0132-B004}{end}

\olpseganchor{OLP-0132-B005}{start}
\olsection{د يو مدل جوړونه}
\olpseganchor{OLP-0132-B005}{end}

\olpseganchor{OLP-0132-B006}{start}
\begin{explain}
اوس مهال مو له~\LR{$\eq$} سره کار نشته؛ يعنې يوازې دا ښودل
  غواړو چې د~تړلي فارمولونه يو سازګار سټ~\LR{$\Gamma$} چې~\LR{$\eq$} نۀ لري،
  د صدق وړ دے۔ لومړے~\LR{$\Gamma$} يوه سازګار، بشپړ او مشبوع
  سټ~\LR{$\Gamma^*$} ته غځوو۔ په دې حالت کښې د مدل
  \LR{$\Struct{M(\Gamma^*)}$} تعريف ساده دے: د~\LR{$\Lang{L'}$} د تړلو
  ترمونو سټ د تعريف ساحه ګرځوو۔ هر ثابت سمبول پخپله ځان ته
  ورګومارو، او په عمومي ډول دا باوري کوو چې د هر تړلي ترم~\LR{$t$}
  دپاره~\LR{$\Value{t}{M(\Gamma^*)} = t$}۔ پريديکاتونه داسې تفسيروو
  چې يوه اتومي تړلے فارمول په~\LR{$\Struct{M(\Gamma^*)}$} کښې هله او
  يوازې هله رښتونې وي چې په~\LR{$\Gamma^*$} کښې وي۔ دا به څرګنده توګه
  په~\LR{$\Gamma^*$} کښې ټولې اتومي تړلي فارمولونه
  په~\LR{$\Struct{M(\Gamma^*)}$} کښې رښتونې کړي۔ نورې جملې هم رښتونې
  دي، په دې شرط چې هغه~\LR{$\Gamma^*$} چې ترې پيل کوو، سازګار، بشپړ او
  مشبوع وي۔
\end{explain}
\olpseganchor{OLP-0132-B006}{end}

\olpseganchor{OLP-0132-B007}{start}
%
\begin{defn}[ترمي مدل]
\ollabel{defn:termmodel}
فرض کړئ~\LR{$\Gamma^*$} په ژبه~\LR{$\Lang L$} کښې د~تړلي فارمولونه يو
بشپړ، سازګار او مشبوع سټ دے۔ د~\LR{$\Gamma^*$}
\emph{ترمي مدل}~\LR{$\Struct M(\Gamma^*)$} هغه جوړښت ده چې په
لاندې ډول تعريفېږي:
\begin{enumerate}
\item د د تعريف ساحه~\LR{$\Domain{M(\Gamma^*)}$} د~\LR{$\Lang L$} د ټولو
  تړلو ترمونو سټ دے۔
\item د ثابت سمبول~\LR{$c$} تفسير پخپله~\LR{$c$} دے:
  \LR{$\Assign{c}{M(\Gamma^*)} = c$}۔
\item تابع~\LR{$f$} ته هغه تابع ګومارل کېږي چې تړلي ترمونه
  \LR{$t_1$}، \dots،~\LR{$t_n$} د ارګومېنټونو په توګه واخلي او تړلے
  ترم~\LR{$f(t_1, \dots, t_n)$} د قيمت په توګه ولري:
\[
\Assign{f}{M(\Gamma^*)}(t_1, \dots, t_n) = f(t_1,\dots, t_n)
\]
\item که~\LR{$R$} يو \LR{$n$}-مقامي پريديکات وي، نو
  \[
  \tuple{t_1, \dots,
  t_n} \in \Assign{R}{M(\Gamma^*)} \text{\RL{ هله او يوازې هله چې }} \Atom{R}{t_1, \dots,
    t_n} \in \Gamma^*.
  \]
\end{enumerate}
\end{defn}

\olpseganchor{OLP-0132-B007}{end}

\olpseganchor{OLP-0132-B008}{start}
%
اوس به وګورو چې په رښتيا~\LR{$\Value{t}{M(\Gamma^*)} = t$} لرو۔
\olpseganchor{OLP-0132-B008}{end}

\olpseganchor{OLP-0132-B009}{start}
\begin{lem}
\ollabel{lem:val-in-termmodel} فرض کړئ~\LR{$\Struct M(\Gamma^*)$} د
\olref{defn:termmodel} ترمي مدل دے؛ بيا د هر تړلي ترم دپاره
\LR{$\Value{t}{M(\Gamma^*)} = t$}۔
\end{lem}
\begin{proof}
 ثبوت پر~\LR{$t$} د استقرا له لارې دے۔ د بنسټ حالت، چې~\LR{$t$}
 يو ثابت سمبول وي، نېغ په نېغه د ترمي مدل له تعريف څخه ترلاسه
 کېږي۔ د استقرا د ګام دپاره فرض کړئ~\LR{$t_1, \ldots, t_n$} داسې تړلي
 ترمونه دي چې~\LR{$\Value{t_i}{M(\Gamma^*)} = t_i$}، او~\LR{$f$} يوه
 \LR{$n$}-مقامي تابع ده۔ بيا
\begin{align*}
\Value{f(t_1,\ldots,t_n)}{M(\Gamma^*)} &= \Assign{f}{M(\Gamma^*)}(\Value{t_1}
 {M(\Gamma^*)},
\ldots, \Value{t_n}{M(\Gamma^*)}) \\
&= \Assign{f}{M(\Gamma^*)}(t_1, \dots, t_n) \\
&= f(t_1,\dots, t_n),
\end{align*}
او نو د استقرا له مخې دا د هر تړلي ترم~\LR{$t$} دپاره سم دي۔
\end{proof}
\textbf{د ژباړې د نسخې سمون (OLCOM-008):} په سرچينه کښې د دې
ټاکنيز شرط درېيم (د لومړۍ درجې منطق د نۀ ټاکل کېدو) تش ارګومېنټ پاتې شوے و؛
دلته په څرګنده توګه ورزيات شو، څو راتلونکے شرط د هغه پر ځاے وانۀ
خيستل شي۔

\olpseganchor{OLP-0132-B009}{end}

\olpseganchor{OLP-0132-B010}{start}
%
\begin{explain}
  جوړښت~\LR{$\Struct{M}$} ښائي يوه په وجودي ډول
    کميت ټاکونکې تړلے فارمول~\LR{$\lexists[x][!A(x)]$} رښتونې کړي،
    بې له دې چې کومه بېلګه~\LR{$!A(t)$} ئې رښتونې کړې وي۔
  جوړښت~\LR{$\Struct{M}$} ښائي د يوې په عمومي
    ډول کميت ټاکونکې تړلے فارمول~\LR{$\lforall[x][!A(x)]$} ټولې
    بېلګې~\LR{$!A(t)$} رښتونې کړي، خو پخپله~\LR{$\lforall[x][!A(x)]$}
    رښتونې نۀ کړي۔ لامل دا دے چې په عمومي ډول د~\LR{$\Domain{M}$}
  هر غړے د کوم تړلي ترم قيمت نۀ وي
  (\LR{$\Struct{M}$} ښائي پوښلے نۀ وي)۔ همدا لامل دے چې د صدق اړيکه د
  متغير-ګومارنو له لارې تعريفېږي۔ خو زموږ د ترمي مدل
  \LR{$\Struct{M(\Gamma^*)}$} دپاره دا اړينه نۀ ده، ځکه هغه پوښلے دے۔
  راتلونکې نتيجه همدا وايي۔
\end{explain}
\olpseganchor{OLP-0132-B010}{end}

\olpseganchor{OLP-0132-B011}{start}
\begin{prop}
\ollabel{prop:quant-termmodel}
فرض کړئ~\LR{$\Struct M(\Gamma^*)$} د~\olref{defn:termmodel} ترمي مدل دے۔
\begin{enumerate}
\item \LR{$\Sat{M(\Gamma^*)}{\lexists[x][!A(x)]}$} هله او يوازې هله چې
  لږ تر لږه د يوه تړلي ترم~\LR{$t$} دپاره~\LR{$\Sat{M(\Gamma^*)}{!A(t)}$} وي۔
\item \LR{$\Sat{M(\Gamma^*)}{\lforall[x][!A(x)]}$} هله او يوازې هله چې
  د ټولو تړلو ترمونو~\LR{$t$} دپاره~\LR{$\Sat{M(\Gamma^*)}{!A(t)}$} وي۔
\end{enumerate}
\end{prop}
\olpseganchor{OLP-0132-B011}{end}

\olpseganchor{OLP-0132-B012}{start}
\begin{proof}
\begin{enumerate}
\item %
  د~\olref[syn][ass]{prop:sat-quant} له مخې،
    \LR{$\Sat{M(\Gamma^*)}{\lexists[x][!A(x)]}$} هله او يوازې هله چې لږ
    تر لږه د يوې متغير-ګومارنې~\LR{$s$} دپاره
    \LR{$\Sat{M(\Gamma^*)}{!A(x)}[s]$} وي۔ ځکه
    \LR{$\Domain{M(\Gamma^*)}$} د~\LR{$\Lang{L}$} له تړلو ترمونو جوړ دے، دا
    هله او يوازې هله کېږي چې لږ تر لږه يو تړلے ترم~\LR{$t$} وي، داسې چې
    \LR{$s(x) = t$} او~\LR{$\Sat{M(\Gamma^*)}{!A(x)}[s]$}۔ د
    \olref[fol][syn][ext]{prop:ext-formulas} له مخې، چې~\LR{$s(x) = t$}
    وي، \LR{$\Sat{M(\Gamma^*)}{!A(x)}[s]$} هله او يوازې هله چې
    \LR{$\Sat{M(\Gamma^*)}{!A(t)}[s]$} وي۔ د
    \olref[fol][syn][ass]{prop:sentence-sat-true} له مخې،
    \LR{$\Sat{M(\Gamma^*)}{!A(t)}[s]$} هله او يوازې هله چې
    \LR{$\Sat{M(\Gamma^*)}{!A(t)}$} وي، ځکه~\LR{$!A(t)$} يوه جمله ده۔
\item %
  د~\olref[syn][ass]{prop:sat-quant} له مخې،
    \LR{$\Sat{M(\Gamma^*)}{\lforall[x][!A(x)]}$} هله او يوازې هله چې د
    هرې متغير-ګومارنې~\LR{$s$} دپاره~\LR{$\Sat{M(\Gamma^*)}{!A(x)}[s]$} وي۔
    راياد کړئ چې~\LR{$\Domain{M(\Gamma^*)}$} د~\LR{$\Lang{L}$} له تړلو
    ترمونو جوړ دے؛ نو د هر تړلي ترم~\LR{$t$} دپاره~\LR{$s(x) = t$} داسې يوه
    متغير-ګومارنه ده، او د هرې متغير-ګومارنې دپاره~\LR{$s(x)$} کوم تړلے
    ترم~\LR{$t$} دے۔ د~\olref[fol][syn][ext]{prop:ext-formulas} له مخې،
    چې~\LR{$s(x) = t$} وي، \LR{$\Sat{M(\Gamma^*)}{!A(x)}[s]$} هله او يوازې
    هله چې~\LR{$\Sat{M(\Gamma^*)}{!A(t)}[s]$} وي۔ د
    \olref[fol][syn][ass]{prop:sentence-sat-true} له مخې،
    \LR{$\Sat{M(\Gamma^*)}{!A(t)}[s]$} هله او يوازې هله چې
    \LR{$\Sat{M(\Gamma^*)}{!A(t)}$} وي، ځکه~\LR{$!A(t)$} يوه جمله ده۔
\end{enumerate}
\end{proof}

\olpseganchor{OLP-0132-B012}{end}



\olpseganchor{OLP-0132-B014}{start}
\begin{lem}[د صدق لمه]
\ollabel{lem:truth} فرض کړئ~\LR{$!A$}،~\LR{$\eq$} نۀ لري۔ بيا
\LR{$\Sat{M(\Gamma^*)}{!A}$} هله او يوازې هله چې \LR{$!A \in \Gamma^*$} وي۔
\end{lem}
\olpseganchor{OLP-0132-B014}{end}

\olpseganchor{OLP-0132-B015}{start}
\begin{proof}
دواړه لوري په يو وخت، او پر~\LR{$!A$} د استقرا له لارې ثابتوو۔
\begin{enumerate}
\item \indcase{!A}{\lfalse}
  {\LR{$\Sat/{M(\Gamma^*)}{\lfalse}$}
    د صدق د تعريف له مخې۔ له بلې خوا، ځکه~\LR{$\Gamma^*$} سازګار دے،
    نو~\LR{$\lfalse \notin \Gamma^*$}۔}
\olpseganchor{OLP-0132-B015}{end}

\olpseganchor{OLP-0132-B016}{start}
\item \indcase{!A}{\ltrue}
  {\LR{$\Sat{M(\Gamma^*)}{\ltrue}$}
    د صدق د تعريف له مخې۔ له بلې خوا، ځکه~\LR{$\Gamma^*$} سازګار او
    بشپړ دے، او~\LR{$\Gamma^* \Proves \ltrue$}، نو
    \LR{$\ltrue \in \Gamma^*$}۔}
\olpseganchor{OLP-0132-B016}{end}

\olpseganchor{OLP-0132-B017}{start}
\item \indcase{!A}{R(t_1, \dots, t_n)}
  {\LR{$\Sat{M(\Gamma^*)}{\Atom{R}{t_1, \dots, t_n}}$} هله او يوازې هله چې \LR{$\tuple{t_1,
      \dots, t_n} \in \Assign{R}{M(\Gamma^*)}$} وي (د صدق د تعريف
    له مخې)، او دا هله او يوازې هله چې
    \LR{$R(t_1, \dots, t_n) \in \Gamma^*$} وي (د~\LR{$\Struct
    M(\Gamma^*)$} د جوړونې له مخې)۔}
\olpseganchor{OLP-0132-B017}{end}

\olpseganchor{OLP-0132-B018}{start}
\item %
  \indcase{!A}{\lnot !B}
    {\LR{$\Sat{M(\Gamma^*)}{\indfrm}$} هله او يوازې هله چې
    \LR{$\Sat/{M(\Gamma^*)}{!B}$}
    وي (د صدق د تعريف له مخې)۔ د استقرا د فرضيې له مخې،
    \LR{$\Sat/{M(\Gamma^*)}{!B}$}
    هله او يوازې هله چې~\LR{$!B \notin \Gamma^*$} وي۔ ځکه~\LR{$\Gamma^*$}
    سازګار او بشپړ دے، \LR{$!B \notin \Gamma^*$} هله او يوازې
    هله چې~\LR{$\lnot !B \in \Gamma^*$} وي۔}
\olpseganchor{OLP-0132-B018}{end}

\olpseganchor{OLP-0132-B019}{start}
\item %
%
  \indcase{!A}{!B \land
    !C}{\LR{$\Sat{M(\Gamma^*)}{\indfrm}$}
    هله او يوازې هله چې دواړه
    \LR{$\Sat{M(\Gamma^*)}{!B}$} او
    \LR{$\Sat{M(\Gamma^*)}{!C}$} ولرو
    (د صدق د تعريف له مخې)، او دا هله او يوازې هله چې دواړه
    \LR{$!B \in \Gamma^*$} او~\LR{$!C \in \Gamma^*$} وي (د استقرا د فرضيې
    له مخې)۔ د~\olref[ccs]{prop:ccs}\olref[ccs]{prop:ccs-and} له
    مخې، دا هله او يوازې هله کېږي چې~\LR{$(!B \land !C) \in \Gamma^*$}۔}
\olpseganchor{OLP-0132-B019}{end}

\olpseganchor{OLP-0132-B020}{start}
\item %
%
  \indcase{!A}{!B \lor !C}
    {\LR{$\Sat{M(\Gamma^*)}{\indfrm}$}
    هله او يوازې هله چې
    \LR{$\Sat{M(\Gamma^*)}{!B}$} يا
    \LR{$\Sat{M(\Gamma^*)}{!C}$} وي
    (د صدق د تعريف له مخې)، او دا هله او يوازې هله چې
    \LR{$!B \in \Gamma^*$} يا~\LR{$!C \in \Gamma^*$} وي (د استقرا د فرضيې
    له مخې)۔ دا هله او يوازې هله کېږي چې
    \LR{$(!B \lor !C) \in \Gamma^*$} وي (د
    \olref[ccs]{prop:ccs}\olref[ccs]{prop:ccs-or} له مخې)۔}
\olpseganchor{OLP-0132-B020}{end}

\olpseganchor{OLP-0132-B021}{start}
\item %
%
  \indcase{!A}{!B \lif !C}
    {\LR{$\Sat{M(\Gamma^*)}{\indfrm}$}
    هله او يوازې هله چې~\LR{$\Sat/{M(\Gamma^*)}{!B}$} يا~\LR{$\Sat{M (\Gamma^*)}{!C}$} وي (د
    صدق د تعريف له مخې)، او دا هله او يوازې هله چې
    \LR{$!B \notin \Gamma^*$} يا~\LR{$!C \in \Gamma^*$} وي (د استقرا د فرضيې
    له مخې)۔ دا هله او يوازې هله کېږي چې
    \LR{$(!B \lif !C) \in \Gamma^*$} وي (د
    \olref[ccs]{prop:ccs}\olref[ccs]{prop:ccs-if} له مخې)۔}
\olpseganchor{OLP-0132-B021}{end}

\olpseganchor{OLP-0132-B022}{start}
%
\item %
  %
    \indcase{!A}{\lforall[x][!B(x)]}{\LR{$\Sat{M(\Gamma^*)}{\indfrm}$} هله او يوازې هله چې
      د ټولو تړلو ترمونو~\LR{$t$} دپاره~\LR{$\Sat{M(\Gamma^*)}{!B(t)}$} وي
      (\olref{prop:quant-termmodel})۔ د استقرا د فرضيې له مخې، دا
      هله او يوازې هله چې د ټولو تړلو ترمونو~\LR{$t$} دپاره
      \LR{$!B(t) \in \Gamma^*$} وي؛ د
      \olref[hen]{prop:saturated-instances} له مخې بيا دا هله او
      يوازې هله چې~\LR{$\lforall[x][!B(x)] \in \Gamma^*$} وي۔
      \textbf{د ژباړې د نسخې سمون (OLCOM-003):} د سرچينې په
      وروستۍ فورمول کښې د استقرا له اوسني حالت سره ناسمه
      فورمول-نښه راغلې؛ دلته د همدې حالت نښه کارول شوې۔}
\olpseganchor{OLP-0132-B022}{end}

\olpseganchor{OLP-0132-B023}{start}
\item %
  %
    \indcase{!A}{\lexists[x][!B(x)]}{\LR{$\Sat{M(\Gamma^*)}{\indfrm}$} هله او يوازې هله چې
      لږ تر لږه د يوه تړلي ترم~\LR{$t$} دپاره
      \LR{$\Sat{M(\Gamma^*)}{!B(t)}$} وي
      (\olref{prop:quant-termmodel})۔ د استقرا د فرضيې له مخې، دا
      هله او يوازې هله چې لږ تر لږه د يوه تړلي ترم~\LR{$t$} دپاره
      \LR{$!B(t) \in \Gamma^*$} وي۔ د
      \olref[hen]{prop:saturated-instances} له مخې بيا دا هله او
      يوازې هله چې~\LR{$\lexists[x][!B(x)] \in \Gamma^*$} وي۔}

\end{enumerate}
\end{proof}
\olpseganchor{OLP-0132-B023}{end}

\olpseganchor{OLP-0133-B004}{start}
\olfileid{fol}{com}{ide}
\olsection{عينيت}
\olpseganchor{OLP-0133-B004}{end}

\olpseganchor{OLP-0133-B005}{start}
\begin{explain}
په مخکښې برخه کښې د ترمي مدل ورکړل شوې جوړونه د هغو سټونو
\LR{$\Gamma$} دپاره د اول مرتبې د منطق د بشپړوالي ثابتولو ته بس ده
چې~\LR{$\eq$} نۀ لري۔ ترمي مدل هر داسې~\LR{$!A \in \Gamma^*$} رښتونے کوي
چې~\LR{$\eq$} نۀ لري (او په دې توګه ټول~\LR{$!A \in \Gamma$} هم)۔ خو که
\LR{$\eq$} موجود وي، دا کار نۀ کوي۔ لامل دا دے چې بيا ښائي~\LR{$\Gamma^*$}
داسې يوه تړلے فارمول~\LR{$\eq[t][t']$} ولري، خو په ترمي مدل کښې د هر
ترم قيمت پخپله هماغه ترم دے۔ نو که~\LR{$t$} او~\LR{$t'$} بېل ترمونه وي، په
ترمي مدل کښې ئې قيمتونه، يعنې په ترتيب سره~\LR{$t$} او~\LR{$t'$}، هم بېل دي؛
ځکه~\LR{$\eq[t][t']$} ناسم دے۔ خو دا ستونزه د هغې جوړونې په کارولو
هوارولے شو چې «خارج قسمت جوړونه» بللے شي۔
\end{explain}
\olpseganchor{OLP-0133-B005}{end}

\olpseganchor{OLP-0133-B006}{start}
\begin{defn}
  فرض کړئ~\LR{$\Gamma^*$} په~\LR{$\Lang L$} کښې د جملو يو سازګار او
  بشپړ سټ دے۔ د~\LR{$\Lang L$} د تړلو ترمونو پر سټ اړيکه
  \LR{$\approx$} داسې تعريفوو:
  \[
  t \approx t' \text{\RL{\quad هله او يوازې هله چې \quad}} \eq[t][t'] \in \Gamma^*
  \]
\end{defn}
\olpseganchor{OLP-0133-B006}{end}

\olpseganchor{OLP-0133-B007}{start}
\begin{prop}
\ollabel{prop:approx-equiv}
اړيکه~\LR{$\approx$} لاندې خاصيتونه لري:
\begin{enumerate}
\item \LR{$\approx$} انعکاسي ده۔
\item \LR{$\approx$} متناظره ده۔
\item \LR{$\approx$} انتقالي ده۔
\item که~\LR{$t \approx t'$} وي،~\LR{$f$} يوه تابع وي، او
  \LR{$t_1$}، \dots،~\LR{$t_{i-1}$}، \LR{$t_{i+1}$}، \dots،~\LR{$t_n$} تړلي ترمونه وي، نو
\[
\Atom{f}{t_1,\dots, t_{i-1}, t, t_{i+1}, \dots, t_n} \approx
\Atom{f}{t_1,\dots, t_{i-1}, t', t_{i+1}, \dots, t_n}.
\]
\item که~\LR{$t \approx t'$} وي،~\LR{$R$} يو پريديکات وي، او
  \LR{$t_1$}، \dots،~\LR{$t_{i-1}$}، \LR{$t_{i+1}$}، \dots،~\LR{$t_n$} تړلي ترمونه وي، نو
\begin{multline*}
\Atom{R}{t_1,\dots, t_{i-1}, t, t_{i+1}, \dots, t_n} \in \Gamma^* \text{\RL{ هله او يوازې هله چې }} \\
\Atom{R}{t_1,\dots, t_{i-1}, t', t_{i+1}, \dots, t_n} \in \Gamma^*.
\end{multline*}
\end{enumerate}
\end{prop}
\olpseganchor{OLP-0133-B007}{end}

\olpseganchor{OLP-0133-B008}{start}
\begin{proof}
ځکه~\LR{$\Gamma^*$} سازګار او بشپړ دے، نو
\LR{$\eq[t][t'] \in \Gamma^*$} هله او يوازې هله چې
\LR{$\Gamma^* \Proves \eq[t][t']$} وي۔ نو د لاندې خبرو ښودل کافي دي:
\begin{enumerate}
\item د هر تړلي ترم~\LR{$t$} دپاره~\LR{$\Gamma^* \Proves \eq[t][t]$}۔
\item که~\LR{$\Gamma^* \Proves \eq[t][t']$} وي، نو
  \LR{$\Gamma^* \Proves \eq[t'][t]$}۔
\item که~\LR{$\Gamma^* \Proves \eq[t][t']$} او
  \LR{$\Gamma^* \Proves \eq[t'][t'']$} وي، نو
  \LR{$\Gamma^* \Proves \eq[t][t'']$}۔
\item که~\LR{$\Gamma^* \Proves \eq[t][t']$} وي، نو
\[
\Gamma^* \Proves
\eq[\Atom{f}{t_1,\dots,t_{i-1},t,t_{i+1},\dots,t_n}][\Atom{f}{t_1,\dots,t_{i-1},t',t_{i+1},\dots,t_n}]
\]
\textbf{د ژباړې د نسخې سمون (OLCOM-004):} د سرچينې په لومړي
تابعي ترم کښې يوه تکراري کامه وه؛ دلته لرې شوې ده۔ دا د هرې
\LR{$n$}-مقامي تابع~\LR{$f$} او تړلو ترمونو \LR{$t_1$}، \dots،
\LR{$t_{i-1}$}، \LR{$t_{i+1}$}، \dots،~\LR{$t_n$} دپاره سم دي۔
\item که~\LR{$\Gamma^* \Proves \eq[t][t']$} او
\LR{$\Gamma^* \Proves
\Atom{R}{t_1,\dots,t_{i-1},t,t_{i+1},\dots,t_n}$} وي، نو
\LR{$\Gamma^* \Proves \Atom{R}{t_1,\dots,t_{i-1},t',t_{i+1},\dots,t_n}$}
د هر \LR{$n$}-مقامي پريديکات~\LR{$R$} او تړلو ترمونو \LR{$t_1$}، \dots،
\LR{$t_{i-1}$}، \LR{$t_{i+1}$}، \dots،~\LR{$t_n$} دپاره۔
\end{enumerate}
\end{proof}
\olpseganchor{OLP-0133-B008}{end}

\olpseganchor{OLP-0133-B009}{start}
\begin{prob}
د~\olref[fol][com][ide]{prop:approx-equiv} ثبوت بشپړ کړئ۔
\end{prob}
\olpseganchor{OLP-0133-B009}{end}

\olpseganchor{OLP-0133-B010}{start}
\begin{defn}
فرض کړئ~\LR{$\Gamma^*$} په يوه ژبه~\LR{$\Lang L$} کښې يو سازګار او
بشپړ سټ دے،~\LR{$t$} يو تړلے ترم دے، او~\LR{$\approx$} د مخکښې
تعريف اړيکه ده۔ بيا:
\[
\equivrep{t}{\approx} = \Setabs{t'}{t'\in \Trm[L], t \approx t'}
\]
او~\LR{$\equivclass{\Trm[L]}{\approx} = \Setabs{\equivrep{t}{\approx}}{t \in \Trm[L]}$}۔
\end{defn}
\olpseganchor{OLP-0133-B010}{end}

\olpseganchor{OLP-0133-B011}{start}
\begin{defn}
\ollabel{defn:term-model-factor}
فرض کړئ~\LR{$\Struct M = \Struct M(\Gamma^*)$} د~\LR{$\Gamma^*$} هغه ترمي
مدل دے چې په~\olref[mod]{defn:termmodel} کښې راغلے۔ بيا
\LR{$\Struct{\equivclass{M}{\approx}}$} لاندې جوړښت ده:
\begin{enumerate}
\item \LR{$\Domain{\equivclass{M}{\approx}} = \equivclass{\Trm[L]}{\approx}$}۔
\item \LR{$\Assign{c}{\equivclass{M}{\approx}} = \equivrep{c}{\approx}$}۔
\item \LR{$\Assign{f}{\equivclass{M}{\approx}}(\equivrep{t_1}{\approx}, \dots,
  \equivrep{t_n}{\approx}) = \equivrep{\Atom{f}{t_1,\dots, t_n}}{\approx}$}۔
\item \LR{$\tuple{\equivrep{t_1}{\approx}, \dots, \equivrep{t_n}{\approx}} \in
  \Assign{R}{\equivclass{M}{\approx}}$} هله او يوازې هله چې
  \LR{$\Sat{M}{\Atom{R}{t_1,\dots, t_n}}$} وي؛ يعنې هله او يوازې هله چې
  \LR{$\Atom{R}{t_1,\dots, t_n} \in \Gamma^*$} وي۔
\end{enumerate}
\end{defn}
\olpseganchor{OLP-0133-B011}{end}

\olpseganchor{OLP-0133-B012}{start}
\begin{explain}
پام وکړئ چې~\LR{$\Assign{f}{\equivclass{M}{\approx}}$} او
\LR{$\Assign{R}{\equivclass{M}{\approx}}$} مو د
\LR{$\equivclass{\Trm[L]}{\approx}$} پر عناصرو داسې تعريف کړي چې هغوی
مو د~\LR{$\equivrep{t}{\approx}$} په بڼه ياد کړي؛ يعنې د
\emph{استازو}~\LR{$t \in \equivrep{t}{\approx}$} له لارې۔ بايد باوري
کړو چې دا تعريفونه د دې استازو پر انتخاب پورې تړلي نۀ دي؛ يعنې د
نورو داسې انتخابونو~\LR{$t'$} دپاره چې هماغه د هم ارزښتۍ ټولګي ټاکي
(\LR{$\equivrep{t}{\approx} = \equivrep{t'}{\approx}$})، تعريفونه هماغه
پايله ورکړي۔ د بېلګې په توګه، که~\LR{$R$} يو يو-مقامي پريديکات
وي، د تعريف وروستۍ فقره وايي چې
\LR{$\equivrep{t}{\approx} \in \Assign{R}{\equivclass{M}{\approx}}$} هله
او يوازې هله چې~\LR{$\Sat{M}{\Atom{R}{t}}$} وي۔ که د کوم بل ترم~\LR{$t'$} دپاره
چې~\LR{$t \approx t'$} وي،~\LR{$\Sat/{M}{\Atom{R}{t'}}$}، نو تعريف به وغواړي
چې~\LR{$\equivrep{t'}{\approx} \notin \Assign{R}{\equivclass{M}{\approx}}$}۔
\textbf{د ژباړې د نسخې سمون (OLCOM-005):} د سرچينې په منفي صدق
فورمول کښې د لومړي استازي نښه راغلې وه؛ د جملې د «بل ترم» او
ورپسې ټولګي له مخې دلته د بل استازي نښه کارول شوې ده۔ که
\LR{$t \approx t'$} وي، نو~\LR{$\equivrep{t}{\approx} =
\equivrep{t'}{\approx}$}؛ خو نۀ شو کولے چې هم
\LR{$\equivrep{t}{\approx} \in \Assign{R}{\equivclass{M}{\approx}}$} او هم
\LR{$\equivrep{t}{\approx} \notin \Assign{R}{\equivclass{M}{\approx}}$}
ولرو۔ خو~\olref{prop:approx-equiv} تضمينوي چې دا نۀ شي پېښېدلے۔
\end{explain}
\olpseganchor{OLP-0133-B012}{end}

\olpseganchor{OLP-0133-B013}{start}
\begin{prop}
\LR{$\Struct{\equivclass{M}{\approx}}$} په ښه ډول تعريف شوے دے؛ يعنې که
\LR{$t_1$}، \dots،~\LR{$t_n$}، \LR{$t_1'$}، \dots،~\LR{$t_n'$} تړلي ترمونه وي، او د هرې
اړوندې جوړې دپاره~\LR{$t_i \approx t_i'$} وي، نو
\begin{enumerate}
\item \LR{$\equivrep{\Atom{f}{t_1,\dots, t_n}}{\approx} =
    \equivrep{\Atom{f}{t_1',\dots, t_n'}}{\approx}$}؛ يعنې
  \[
  \Atom{f}{t_1,\dots, t_n} \approx \Atom{f}{t_1',\dots, t_n'}
  \]
  او
\item \LR{$\Sat{M}{\Atom{R}{t_1,\dots, t_n}}$} هله او يوازې هله چې
  \LR{$\Sat{M}{\Atom{R}{t_1',\dots, t_n'}}$} وي؛ يعنې
  \[
    \Atom{R}{t_1,\dots, t_n} \in \Gamma^* \text{\RL{ هله او يوازې هله چې }}
    \Atom{R}{t_1',\dots, t_n'} \in \Gamma^*.
  \]
\end{enumerate}
\end{prop}
\olpseganchor{OLP-0133-B013}{end}

\olpseganchor{OLP-0133-B014}{start}
\begin{proof}
له~\olref{prop:approx-equiv} څخه پر~\LR{$n$} د استقرا له لارې ترلاسه کېږي۔
\end{proof}
\olpseganchor{OLP-0133-B014}{end}

\olpseganchor{OLP-0133-B015}{start}
د ترمي مدل د حالت په څېر، د صدق لمې تر ثبوت مخکښې لاندې لمې ته
اړتيا لرو۔
\olpseganchor{OLP-0133-B015}{end}

\olpseganchor{OLP-0133-B016}{start}
\begin{lem}
\ollabel{lem:val-in-termmodel-factored} فرض کړئ~\LR{$\Struct M = \Struct M
 (\Gamma^*)$}؛ بيا~\LR{$\Value{t}{\equivclass{M}{\approx}} = \equivrep{t}
 {\approx}$}۔
\end{lem}
\begin{proof}
ثبوت د~\olref[mod]{lem:val-in-termmodel} له ثبوت سره ورته دے۔
\end{proof}
\olpseganchor{OLP-0133-B016}{end}

\olpseganchor{OLP-0133-B017}{start}
\begin{prob}
د~\olref[fol][com][ide]{lem:val-in-termmodel-factored} ثبوت بشپړ کړئ۔
\end{prob}
\olpseganchor{OLP-0133-B017}{end}

\olpseganchor{OLP-0133-B018}{start}
\begin{lem}
\ollabel{lem:truth}
د ټولو جملو~\LR{$!A$} دپاره~\LR{$\Sat{\equivclass{M}{\approx}}{!A}$} هله او
يوازې هله چې~\LR{$!A \in \Gamma^*$} وي۔
\end{lem}
\olpseganchor{OLP-0133-B018}{end}

\olpseganchor{OLP-0133-B019}{start}
\begin{proof}
پر~\LR{$!A$} د استقرا له لارې، کټ مټ د~\olref[mod]{lem:truth} د ثبوت
په څېر۔ يوازې هغه حالت زيات پام غواړي چې~\LR{$!A \ident \eq[t][t']$} وي۔
\begin{align*}
\Sat{\equivclass{M}{\approx}}{\eq[t][t']} & \text{\RL{ هله او يوازې هله چې }} \equivrep{t}{\approx} = \equivrep{t'}{\approx}
\text{\RL{ وي (د \LR{$\Struct{\equivclass{M}{\approx}}$} د تعريف له مخې)}}\\
& \text{\RL{ هله او يوازې هله چې }} t \approx t' \text{\RL{ وي (د \LR{$\equivrep{t}{\approx}$} د تعريف له مخې)}}\\
& \text{\RL{ هله او يوازې هله چې }} \eq[t][t'] \in \Gamma^* \text{\RL{ وي (د \LR{$\approx$} د تعريف له مخې)۔}}
\end{align*}
\end{proof}
\olpseganchor{OLP-0133-B019}{end}

\olpseganchor{OLP-0133-B020}{start}
\begin{digress}
پام وکړئ: سره له دې چې~\LR{$\Struct{M(\Gamma^*)}$} تل د شمېر وړ او
نامتناهي دے،~\LR{$\Struct{\equivclass{M}{\approx}}$} ښائي متناهي وي، ځکه
کېدلے شي چې د~\LR{$\equivrep{t}{\approx}$} ټولګي يوازې متناهي شمېر وي۔
دا تمه کېدونکې ده، ځکه ښائي~\LR{$\Gamma$} داسې تړلي فارمولونه ولري چې هره
هغه جوړښت چې دوى پکښې رښتونې وي، بايد متناهي وي۔ د بېلګې
په توګه،~\LR{$\lforall[x][\lforall[y][x = y]]$} يوه سازګاره تړلے فارمول
ده، خو يوازې په هغو جوړښتونه کښې رښتونې کېږي چې د تعريف ساحه
ئې کټ مټ يو غړے لري۔
\end{digress}
\olpseganchor{OLP-0133-B020}{end}

\olpseganchor{OLP-0134-B004}{start}
\olfileid{fol}{com}{cth}
\olpseganchor{OLP-0134-B004}{end}

\olpseganchor{OLP-0134-B005}{start}
\olsection{د بشپړتيا قضيه}
\olpseganchor{OLP-0134-B005}{end}

\olpseganchor{OLP-0134-B006}{start}
\begin{explain}
راځئ خپلې پايلې سره يوځاے کړو: د بشپړتيا قضيې ته رسېږو۔
\end{explain}
\olpseganchor{OLP-0134-B006}{end}

\olpseganchor{OLP-0134-B007}{start}
\begin{thm}[د بشپړتيا قضيه]
\ollabel{thm:completeness}
فرض کړئ~\LR{$\Gamma$} د~تړلي فارمولونه يو سټ دے۔ که~\LR{$\Gamma$} سازګار وي،
نو د صدق وړ دے۔
\end{thm}
\olpseganchor{OLP-0134-B007}{end}

\olpseganchor{OLP-0134-B008}{start}
\begin{proof}
فرض کړئ~\LR{$\Gamma$} سازګار دے۔ د
  \olref[hen]{lem:henkin} له مخې، يو مشبوع سازګار سټ
  \LR{$\Gamma' \supseteq \Gamma$} شته۔ د~\olref[lin]{lem:lindenbaum} له
مخې، داسې~\LR{$\Gamma^* \supseteq \Gamma'$} شته چې
سازګار او بشپړ دے۔

  ځکه~\LR{$\Gamma' \subseteq \Gamma^*$}، نو د هرې
  فارمول~\LR{$!A(x)$} دپاره~\LR{$\Gamma^*$} په لاندې بڼه يوه تړلے فارمول
  لري:
  \LR{$\lexists[x][!A(x)] \lif !A(c)$}
  او په دې توګه~\LR{$\Gamma^*$} مشبوع دے۔ که~\LR{$\Gamma$}،~\LR{$\eq$} نۀ لري،
  نو د~\olref[mod]{lem:truth} له مخې،
\LR{$\Sat{M(\Gamma^*)}{!A}$} هله او
يوازې هله چې~\LR{$!A \in \Gamma^*$} وي۔ له دې څخه په تېره بيا دا ترلاسه
کېږي چې د ټولو~\LR{$!A \in \Gamma$} دپاره
\LR{$\Sat{M(\Gamma^*)}{!A}$}؛ نو
\LR{$\Gamma$} د صدق وړ دے۔
  که~\LR{$\Gamma$}،~\LR{$\eq$} ولري، نو د~\olref[ide]{lem:truth} له مخې د ټولو
  تړلي فارمولونه~\LR{$!A$} دپاره
  \LR{$\Sat{\equivclass{M}{\approx}}{!A}$} هله او يوازې هله چې
  \LR{$!A \in \Gamma^*$} وي۔ په تېره بيا، د ټولو~\LR{$!A \in \Gamma$} دپاره
  \LR{$\Sat{\equivclass{M}{\approx}}{!A}$}؛ نو~\LR{$\Gamma$} د صدق وړ دے۔
\end{proof}
\olpseganchor{OLP-0134-B008}{end}

\olpseganchor{OLP-0134-B009}{start}
\begin{cor}[د بشپړتيا قضيه، دويمه بڼه]
\ollabel{cor:completeness}
د ټولو~\LR{$\Gamma$} او تړلي فارمولونه~\LR{$!A$} دپاره: که
\LR{$\Gamma \Entails !A$} وي، نو~\LR{$\Gamma \Proves !A$}۔
\end{cor}
\olpseganchor{OLP-0134-B009}{end}

\olpseganchor{OLP-0134-B010}{start}
\begin{proof}
پام وکړئ چې په~\olref{cor:completeness} او
\olref{thm:completeness} کښې~\LR{$\Gamma$} ګانې په کلي ډول کميت ټاکل شوې
دي۔ د دې دپاره چې ځان ګډوډ نۀ کړو،~\olref{thm:completeness} د يوه
بېل متغير په کارولو بيا بيانوو: د~تړلي فارمولونه د هر سټ~\LR{$\Delta$}
دپاره، که~\LR{$\Delta$} سازګار وي، نو د صدق وړ دے۔ د نقيض عکس له مخې،
که~\LR{$\Delta$} د صدق وړ نۀ وي، نو~\LR{$\Delta$} ناسازګار دے۔ دا به د لازمې
نتيجې په ثبوت کښې وکاروو۔
\olpseganchor{OLP-0134-B010}{end}

\olpseganchor{OLP-0134-B011}{start}
فرض کړئ~\LR{$\Gamma \Entails !A$}۔ بيا د
\olref[syn][sem]{prop:entails-unsat} له مخې
\LR{$\Gamma \cup \{\lnot !A\}$} د صدق وړ نۀ دے۔ که
\LR{$\Gamma \cup \{\lnot !A\}$} خپل~\LR{$\Delta$} ونيسو، د
\olref{thm:completeness} مخکښې بڼه راکوي چې
\LR{$\Gamma \cup \{\lnot !A\}$} ناسازګار دے۔ د
%
  \tagrefs{prfAX/{fol:axd:prv:prop:prov-incons},
    prfSC/{fol:seq:prv:prop:prov-incons},
    prfND/{fol:ntd:prv:prop:prov-incons},
    prfTab/{fol:tab:prv:prop:prov-incons}},
له مخې~\LR{$\Gamma \Proves !A$}۔
\end{proof}
\olpseganchor{OLP-0134-B011}{end}

\olpseganchor{OLP-0134-B012}{start}
\begin{prob}
\olref[fol][com][cth]{cor:completeness} وکاروئ چې
\olref[fol][com][cth]{thm:completeness} ثابت کړئ، او په دې توګه
وښيئ چې د بشپړتيا قضيې دواړه بيانونه هم ارزښته دي۔
\end{prob}
\olpseganchor{OLP-0134-B012}{end}



\olpseganchor{OLP-0134-B014}{start}
\begin{prob}
د دې دپاره چې د~اشتقاق نظام بشپړ وي، قاعدې ئې بايد دومره
پياوړې وي چې هر د صدق وړ نۀ سټ ناسازګار ثابت کړي۔ د بشپړتيا د ثبوت
دپاره د~اشتقاق کومې قاعدې اړينې وې؟ آيا له دې قاعدو څخه
کومه يوه د ثبوت په هېڅ ځاے کښې نۀ ده کارول شوې؟ د دې پوښتنو د ځواب
دپاره داسې يو لړ يا شکل جوړ کړئ چې وښيي د~اشتقاق کومې
قاعدې په هغو کومو پايلو کښې کارول شوې چې د
\olref[fol][com][cth]{thm:completeness} ثبوت ته رسېږي۔ په دې ثبوتونو
کښې د قاعدو هر ضمني استعمال هم خامخا ياد کړئ۔
\end{prob}
\olpseganchor{OLP-0134-B014}{end}

\olpseganchor{OLP-0135-B004}{start}
\olfileid{fol}{com}{com}
\olpseganchor{OLP-0135-B004}{end}

\olpseganchor{OLP-0135-B005}{start}
\olsection{د متناهي شاهد قضيه}
\olpseganchor{OLP-0135-B005}{end}

\olpseganchor{OLP-0135-B006}{start}
د بشپړتيا قضيې يوه مهمه پايله د متناهي شاهد قضيه ده۔ د متناهي
شاهد قضيه وايي چې که د~تړلي فارمولونه د يوه سټ هر
\emph{متناهي} فرعي سټ د صدق وړ وي، نو ټول سټ د صدق وړ دے، حتٰی
که سټ پخپله نامتناهي وي۔ دا هېڅ څرګنده خبره نۀ ده۔ لږ تر لږه په
لومړي نظر داسې هېڅ څۀ نۀ ښکاري چې د~تړلي فارمولونه د داسې نامتناهي
سټونو امکان رد کړي چې متناقض وي، خو تناقض ئې، په اصطلاح، يوازې له
نامتناهي شمېر څخه راولاړېږي۔ د متناهي شاهد قضيه وايي چې داسې حالت
ردېدلے شي: د~تړلي فارمولونه داسې د صدق وړ نۀ نامتناهي سټونه نشته
چې هر متناهي فرعي سټ ئې د صدق وړ وي۔ د بشپړتيا قضيې په څېر، د دې
يوه بڼه هم له استلزام سره تړاو لري: که د~تړلي فارمولونه يو نامتناهي
سټ کوم څۀ لازموي، نو لا له وړاندې ئې يو متناهي فرعي سټ لازموي۔
\olpseganchor{OLP-0135-B006}{end}

\olpseganchor{OLP-0135-B007}{start}
\begin{defn}
  د~فارمولونه يو سټ~\LR{$\Gamma$} هله او يوازې هله
  \emph{په متناهي ډول د صدق وړ} دے چې هر متناهي
  \LR{$\Gamma_0 \subseteq \Gamma$} د صدق وړ وي۔
\end{defn}
\olpseganchor{OLP-0135-B007}{end}

\olpseganchor{OLP-0135-B008}{start}
\begin{thm}[د متناهي شاهد قضيه]
\ollabel{thm:compactness}
د هرې د جملو ټولګې~\LR{$\Gamma$} او هرې جملې~\LR{$!A$} دپاره لاندې خبرې سمې دي:
\begin{enumerate}
  \item \LR{$\Gamma \Entails !A$} هله او يوازې هله چې يو متناهي
    \LR{$\Gamma_0 \subseteq \Gamma$} وي، داسې چې
    \LR{$\Gamma_0 \Entails !A$}۔
  \item \LR{$\Gamma$} هله او يوازې هله د صدق وړ دے چې په متناهي ډول
    د صدق وړ وي۔
\end{enumerate}
\end{thm}
\olpseganchor{OLP-0135-B008}{end}

\olpseganchor{OLP-0135-B009}{start}
\begin{proof}
موږ (2) ثابتوو۔ که~\LR{$\Gamma$} د صدق وړ وي، نو داسې
جوړښت~\LR{$\Struct{M}$}
شته چې د ټولو~\LR{$!A \in \Gamma$} دپاره
\LR{$\Sat{M}{!A}$}۔ څرګنده ده چې همدا
\LR{$\Struct{M}$} د~\LR{$\Gamma$} هر متناهي فرعي
سټ هم پوره کوي؛ نو~\LR{$\Gamma$} په متناهي ډول د صدق وړ دے۔
\olpseganchor{OLP-0135-B009}{end}

\olpseganchor{OLP-0135-B010}{start}
اوس فرض کړئ~\LR{$\Gamma$} په متناهي ډول د صدق وړ دے۔ بيا هر متناهي
فرعي سټ~\LR{$\Gamma_0 \subseteq \Gamma$} د صدق وړ دے۔ د سموالي له مخې
(%
  \tagrefs{prfAX/{fol:axd:sou:cor:consistency-soundness},
    prfSC/{fol:seq:sou:cor:consistency-soundness},
    prfND/{fol:ntd:sou:cor:consistency-soundness},
    prfTab/{fol:tab:sou:cor:consistency-soundness}})،
هر متناهي فرعي سټ سازګار دے۔ بيا د لاندې پايلې له مخې~\LR{$\Gamma$}
پخپله هم بايد سازګار وي:
%
  \tagrefs{prfAX/{fol:axd:ptn:prop:proves-compact},
    prfSC/{fol:seq:ptn:prop:proves-compact},
    prfND/{fol:ntd:ptn:prop:proves-compact},
    prfTab/{fol:tab:ptn:prop:proves-compact}}.
د بشپړتيا له مخې (\olref[cth]{thm:completeness})، ځکه~\LR{$\Gamma$}
سازګار دے، نو د صدق وړ دے۔
\end{proof}
\olpseganchor{OLP-0135-B010}{end}

\olpseganchor{OLP-0135-B011}{start}
\begin{prob}
د~\olref[fol][com][com]{thm:compactness} برخه (1) ثابته کړئ۔
\end{prob}
\olpseganchor{OLP-0135-B011}{end}



\olpseganchor{OLP-0135-B013}{start}
%
\begin{ex}
د يوې تيورۍ~\LR{$\Gamma$} په هر مدل~\LR{$\Struct{M}$} کښې، هر ترم~\LR{$t$} خامخا
د~\LR{$\Domain{M}$} يو غړے ټاکي۔ آيا دا هم تضمينولے شو چې د
\LR{$\Domain{M}$} هر غړے کوم ترم ټاکلے وي؟ په بله وينا، آيا داسې
تيورۍ~\LR{$\Gamma$} شته چې ټول مدلونه ئې پوښلي وي؟ د متناهي شاهد قضيه
ښيي چې که~\LR{$\Gamma$} نامتناهي مدلونه ولري، داسې نۀ ده۔ دا داسې
ليدلے شو: فرض کړئ~\LR{$\Struct{M}$} د~\LR{$\Gamma$} يو نامتناهي مدل دے، او
\LR{$c$} داسې يو ثابت سمبول دے چې د~\LR{$\Gamma$} په ژبه کښې نشته۔ فرض
کړئ~\LR{$\Delta$} د ټولو هغو جملو~\LR{$\eq/[c][t]$} سټ دے چې~\LR{$t$} د~\LR{$\Gamma$}
په ژبه~\LR{$\Lang{L}$} کښې يو تړلے ترم وي؛ يعنې
  \[
  \Delta = \Setabs{\eq/[c][t]}{t \in \Trm[L]}.
  \]
د~\LR{$\Gamma \cup \Delta$} يو متناهي فرعي سټ د
\LR{$\Gamma' \cup \Delta'$} په بڼه ليکل کېدے شي، چې
\LR{$\Gamma' \subseteq \Gamma$} او~\LR{$\Delta' \subseteq \Delta$} وي۔ ځکه
\LR{$\Delta'$} متناهي دے، يوازې متناهي ډېر ترمونه لرلے شي۔ داسې
\LR{$a \in \Domain{M}$} واخلئ چې د~\LR{$\Domain{M}$} يو غړے وي او له
هغوى څخه ئې هېڅ يو نۀ ټاکي؛ او~\LR{$\Struct{M'}$} داسې جوړښت
واخلئ چې کټ مټ د~\LR{$\Struct{M}$} په څېر وي، خو
\LR{$\Assign{c}{M'} = a$} هم ولري۔ ځکه په~\LR{$\Delta'$} کښې د راغلو ټولو
\LR{$t$} دپاره~\LR{$a \neq \Value{t}{M}$}، نو~\LR{$\Sat{M'}{\Delta'}$}۔ ځکه
\LR{$\Sat{M}{\Gamma}$}،~\LR{$\Gamma' \subseteq \Gamma$}، او~\LR{$c$} په~\LR{$\Gamma$}
کښې نۀ راځي، نو~\LR{$\Sat{M'}{\Gamma'}$} هم۔ دواړه سره يوځاے راکوي چې
د~\LR{$\Gamma \cup \Delta$} د هر متناهي فرعي سټ
\LR{$\Gamma' \cup \Delta'$} دپاره~\LR{$\Sat{M'}{\Gamma' \cup \Delta'}$}۔
نو د~\LR{$\Gamma \cup \Delta$} هر متناهي فرعي سټ د صدق وړ دے۔ د
متناهي شاهد قضيې له مخې،~\LR{$\Gamma \cup \Delta$} پخپله د صدق وړ دے۔
نو داسې مدلونه شته چې~\LR{$\Sat{M}{\Gamma \cup \Delta}$}۔ هر داسې
\LR{$\Struct{M}$} د~\LR{$\Gamma$} مدل دے، خو پوښلے نۀ دے، ځکه د~\LR{$\Lang{L}$}
د ټولو ترمونو~\LR{$t$} دپاره
\LR{$\Value{c}{M} \neq \Value{t}{M}$}۔
\end{ex}
\olpseganchor{OLP-0135-B013}{end}

\olpseganchor{OLP-0135-B014}{start}
\begin{ex}
داسې يوه ژبه~\LR{$\Lang{L}$} په پام کښې ونيسئ چې
پريديکات~\LR{$<$}، ثابت سمبولونه~\LR{$\Obj{0}$} او~\LR{$\Obj{1}$}، او
تابعې~\LR{$+$}،~\LR{$\times$} او~\LR{$-$} ولري۔ فرض کړئ~\LR{$\Gamma$} په دې
ژبه کښې د ټولو هغو تړلي فارمولونه سټ دے چې په
جوړښت~\LR{$\Struct{Q}$} کښې رښتونې دي، چې د تعريف ساحه ئې
\LR{$\Rat$} او تفسيرونه ئې څرګند دي۔~\LR{$\Gamma$} د~\LR{$\Lang{L}$} د ټولو هغو
تړلي فارمولونه سټ دے چې د ناطقو عددونو په اړه رښتونې دي۔ څرګنده ده
چې په~\LR{$\Rat$} (او حتٰی په~\LR{$\Real$}) کښې داسې عددونه~\LR{$r$} نشته چې
له~\LR{$0$} څخه لوی، خو د ټولو~\LR{$k \in \PosInt$} دپاره له~\LR{$1/k$} څخه
کوچني وي۔ که داسې عدد موجود واے، يو \emph{بې حده کوچنے عدد} به
و: له صفره بېل، خو بې حده کوچنے۔ د متناهي شاهد قضيه د دې ښودلو
دپاره کارېدلے شي چې د~\LR{$\Gamma$} داسې مدلونه شته چې بې حده کوچني
عددونه پکښې موجود دي۔ زموږ په ژبه کښې د تقسيم تابع نشته
(پر صفر تقسيم ناتعريف دے، او تابعې بايد په بشپړو تابعو
تفسير شي)۔ خو بيا هم څرګندولے شو چې~\LR{$r < 1/k$}، ځکه دا هله او يوازې
هله سم دي چې~\LR{$r \cdot k < 1$}۔ اوس~\LR{$c$} يو نوے ثابت سمبول ونيسئ او
\LR{$\Delta$} داسې وټاکئ:
\[
\{\Obj{0} < c\}
\cup \Setabs{ c \times \num k < \Obj{1} }{k \in \PosInt}
\]
(چې~\LR{$\num{k} = (\Obj{1} + (\Obj{1} + \dots + (\Obj{1} +
\Obj{1})\dots))$} کښې \LR{$k$} ځله~\LR{$\Obj{1}$} راغلے)۔ د~\LR{$\Delta$} د هر
متناهي فرعي سټ~\LR{$\Delta_0$} دپاره داسې~\LR{$K$} شته چې په~\LR{$\Delta_0$}
کښې د ټولو تړلي فارمولونه~\LR{$c \times \num{k} < \Obj{1}$} دپاره
\LR{$k < K$} وي۔ که~\LR{$\Struct{Q}$} داسې~\LR{$\Struct{Q'}$} ته پراخه کړو چې
\LR{$\Assign{c}{Q'} = 1/K$} ولري، نو د هر متناهي
\LR{$\Gamma_0 \subseteq \Gamma$} دپاره
\LR{$\Sat{Q'}{\Gamma_0 \cup \Delta_0}$} لرو؛ نو
\LR{$\Gamma \cup \Delta$} په متناهي ډول د صدق وړ دے (تمرين: دا په
تفصيل ثابت کړئ)۔ د متناهي شاهد قضيې له مخې،
\LR{$\Gamma \cup \Delta$} د صدق وړ دے۔ د~\LR{$\Gamma \cup \Delta$} هر
مدل~\LR{$\Struct{S}$} يو بې حده کوچنے عدد لري؛ يعنې~\LR{$\Assign{c}{S}$}۔
\end{ex}

\olpseganchor{OLP-0135-B014}{end}

\olpseganchor{OLP-0135-B015}{start}
\begin{prob}
د حساب په معياري مدل~\LR{$\Struct{N}$} کښې داسې هېڅ
غړے~\LR{$k \in \Domain{N}$} نشته چې هره صيغه~\LR{$\num{n} < x$}
پوره کړي (چې~\LR{$\num{n}$} هغه~\LR{$\Obj{0}^{\prime\dots\prime}$} دے چې
\LR{$n$} دانې~\LR{$\prime$} لري)۔ د متناهي شاهد قضيه وکاروئ او وښيئ چې د
حساب په ژبه کښې هغه~تړلي فارمولونه چې د حساب په معياري
مدل~\LR{$\Struct{N}$} کښې رښتونې دي، په داسې جوړښت~\LR{$\Struct{N'}$}
کښې هم رښتونې دي چې يو داسې غړے لري چې په رښتيا هره
صيغه~\LR{$\num{n} < x$} پوره کوي۔
\end{prob}
\olpseganchor{OLP-0135-B015}{end}

\olpseganchor{OLP-0135-B016}{start}

\begin{ex}
پوهېږو چې له~عينيت سره د اول مرتبې منطق دا څرګندولے شي چې
د~د تعريف ساحه کچه بايد لږ تر لږه ټاکلې اندازه ولري: جمله
\LR{$!A_{\ge n}$} (چې وايي «لږ تر لږه~\LR{$n$} بېل څيزونه شته») يوازې په
هغو جوړښتونو کښې رښتونې ده چې~\LR{$\Domain{M}$} لږ تر لږه~\LR{$n$} څيزونه
ولري۔ نو که داسې وټاکو:
  \[
  \Delta = \Setabs{!A_{\ge n}}{n \ge 1}
  \]
بيا د~\LR{$\Delta$} هر مدل بايد نامتناهي وي۔ په دې توګه يوې تيورۍ ته
د~\LR{$\Delta$} په ورزياتولو تضمينولے شو چې يوازې نامتناهي مدلونه
ولري: د~\LR{$\Gamma \cup \Delta$} مدلونه کټ مټ د~\LR{$\Gamma$} نامتناهي
مدلونه دي۔
\olpseganchor{OLP-0135-B016}{end}

\olpseganchor{OLP-0135-B017}{start}
نو د اول مرتبې منطق نامتناهي‌والے څرګندولے شي۔ خو د متناهي شاهد
قضيه ښيي چې متناهي‌والے نۀ شي څرګندولے۔ ځکه فرض کړئ د جملو کوم
سټ~\LR{$\Lambda$} په ټولو او يوازې متناهي جوړښتونه کښې رښتونے
وي۔ بيا~\LR{$\Delta \cup \Lambda$} په متناهي ډول د صدق وړ دے۔ ولې؟
فرض کړئ
\LR{$\Delta' \cup \Lambda' \subseteq \Delta \cup \Lambda$} متناهي دے،
چې~\LR{$\Delta' \subseteq \Delta$} او~\LR{$\Lambda' \subseteq \Lambda$} وي۔
که لومړے فرعي سټ خالي وي، د اندېکس ارزښت يو ټاکو؛ او که خالي نۀ وي،
\LR{$n$} دې تر ټولو لوی هغه عدد وي چې~\LR{$!A_{\ge n} \in \Delta'$}۔
\textbf{د ژباړې د نسخې سمون (OLCOM-006):} سرچينه سمدستي تر ټولو
لوی اندېکس ټاکي، خو خالي متناهي فرعي سټ داسې اندېکس نۀ لري؛ دلته
خالي حالت جلا ونيول شو۔ ځکه~\LR{$\Lambda$} په ټولو متناهي جوړښتونو کښې
رښتونے دے، داسې مدل~\LR{$\Struct{M}$} لري چې متناهي ډېر، خو
\LR{$\ge n$}،~غړي ولري۔ خو بيا
\LR{$\Sat{M}{\Delta' \cup \Lambda'}$}۔ د متناهي شاهد قضيې له مخې،
\LR{$\Delta \cup \Lambda$} يو نامتناهي مدل لري، او دا له هغې فرضيې سره
تناقض لري چې~\LR{$\Lambda$} يوازې په متناهي جوړښتونه کښې رښتونے
دے۔
\end{ex}

\olpseganchor{OLP-0135-B017}{end}

\olpseganchor{OLP-0136-B004}{start}
\olfileid{fol}{com}{cpd}
\olpseganchor{OLP-0136-B004}{end}

\olpseganchor{OLP-0136-B005}{start}
\olsection{د متناهي شاهد قضيې نېغ ثبوت}
\olpseganchor{OLP-0136-B005}{end}

\olpseganchor{OLP-0136-B006}{start}
د بشپړتيا قضيې له کارولو پرته، د بشپړتيا قضيې د ثبوت د هماغو
نظرونو په کارولو، د متناهي شاهد قضيه نېغ په نېغه ثابتولے شو۔ د
بشپړتيا قضيې په ثبوت کښې مو د~تړلي فارمولونه له يوه سازګار
سټ~\LR{$\Gamma$} څخه پيل وکړ، هغه مو د~تړلي فارمولونه يوه
سازګار، مشبوع او بشپړ سټ~\LR{$\Gamma^*$} ته
پراخ کړ، او بيا مو وښودله چې له~\LR{$\Gamma^*$} څخه په جوړ شوي
ترمي
  مدل~\LR{$\Struct{M(\Gamma^*)}$}
کښې د~\LR{$\Gamma$} ټولې تړلي فارمولونه رښتونې دي؛ نو~\LR{$\Gamma$} د صدق وړ دے۔
\olpseganchor{OLP-0136-B006}{end}

\olpseganchor{OLP-0136-B007}{start}
همدا طريقه د دې ښودلو دپاره کارولے شو چې د جملو يو په متناهي ډول
د صدق وړ سټ، د صدق وړ دے۔ يوازې د هغو پايلو اړوندې بڼې ثابتول
پکار دي چې د صدق لمې ته رسېږي، او پکښې «سازګار» په «په متناهي
ډول د صدق وړ» بدلوو۔
\olpseganchor{OLP-0136-B007}{end}

\olpseganchor{OLP-0136-B008}{start}
\begin{prop}
\ollabel{prop:fsat-ccs}
فرض کړئ~\LR{$\Gamma$}،~بشپړ او په متناهي ډول د صدق وړ دے۔ بيا:
\begin{enumerate}
\item \LR{$(!A \land !B) \in \Gamma$}
  هله او يوازې هله چې دواړه \LR{$!A \in \Gamma$} او \LR{$!B \in \Gamma$} وي۔
\olpseganchor{OLP-0136-B008}{end}

\olpseganchor{OLP-0136-B009}{start}
\item \LR{$(!A \lor !B) \in \Gamma$} هله او يوازې هله چې
  يا~\LR{$!A \in \Gamma$} يا~\LR{$!B \in \Gamma$} وي۔
\olpseganchor{OLP-0136-B009}{end}

\olpseganchor{OLP-0136-B010}{start}
\item \LR{$(!A \lif !B) \in \Gamma$} هله او يوازې هله چې
  يا~\LR{$!A \notin \Gamma$} يا~\LR{$!B \in \Gamma$} وي۔
\end{enumerate}
\end{prop}
\olpseganchor{OLP-0136-B010}{end}

\olpseganchor{OLP-0136-B011}{start}
\begin{prob}
\olref[fol][com][cpd]{prop:fsat-ccs} ثابت کړئ۔~\LR{$\Proves$} مۀ کاروئ۔
\end{prob}
\olpseganchor{OLP-0136-B011}{end}



\olpseganchor{OLP-0136-B013}{start}
%
\begin{lem}
\ollabel{lem:fsat-henkin} هر په متناهي ډول د صدق وړ سټ~\LR{$\Gamma$}
يو مشبوع، په متناهي ډول د صدق وړ سټ~\LR{$\Gamma'$} ته غځېدلے شي۔
\end{lem}

\olpseganchor{OLP-0136-B013}{end}

\olpseganchor{OLP-0136-B014}{start}
\begin{prob}
\olref[fol][com][cpd]{lem:fsat-henkin} ثابت کړئ۔ (اشاره: مهم ګام
دا ښودل دي چې که~\LR{$\Gamma_n$} په متناهي ډول د صدق وړ وي، نو
\LR{$\Gamma_n \cup \{!D_n\}$} هم همداسې دے؛ بې له دې چې
اشتقاقونه يا سازګارۍ ته مراجعه وشي۔)
\end{prob}
\olpseganchor{OLP-0136-B014}{end}

\olpseganchor{OLP-0136-B015}{start}

\begin{prop}\ollabel{prop:fsat-instances}
فرض کړئ~\LR{$\Gamma$} بشپړ، په متناهي ډول د صدق وړ، او مشبوع دے۔
\begin{enumerate}
\item \LR{$\lexists[x][!A(x)] \in \Gamma$} هله او يوازې هله چې
لږ تر لږه د يوه تړلي ترم~\LR{$t$} دپاره~\LR{$!A(t) \in \Gamma$} وي۔
\item \LR{$\lforall[x][!A(x)] \in \Gamma$} هله او يوازې هله چې
د ټولو تړلو ترمونو~\LR{$t$} دپاره~\LR{$!A(t) \in \Gamma$} وي۔
\end{enumerate}
\end{prop}

\olpseganchor{OLP-0136-B015}{end}

\olpseganchor{OLP-0136-B016}{start}
\begin{prob}
\olref[fol][com][cpd]{prop:fsat-instances} ثابت کړئ۔
\end{prob}
\olpseganchor{OLP-0136-B016}{end}

\olpseganchor{OLP-0136-B017}{start}
\begin{lem}
\ollabel{lem:fsat-lindenbaum} هر په متناهي ډول د صدق وړ سټ~\LR{$\Gamma$}
يو بشپړ او په متناهي ډول د صدق وړ سټ~\LR{$\Gamma^*$} ته
غځېدلے شي۔
\end{lem}
\olpseganchor{OLP-0136-B017}{end}

\olpseganchor{OLP-0136-B018}{start}
\begin{prob}
\olref[fol][com][cpd]{lem:fsat-lindenbaum} ثابت کړئ۔ (اشاره: مهم
ګام دا ښودل دي چې که~\LR{$\Gamma_n$} په متناهي ډول د صدق وړ وي، نو يا
\LR{$\Gamma_n \cup \{!A_n\}$} يا~\LR{$\Gamma_n \cup \{\lnot !A_n\}$} په
متناهي ډول د صدق وړ دے۔)
\end{prob}
\olpseganchor{OLP-0136-B018}{end}



\olpseganchor{OLP-0136-B020}{start}
\begin{thm}[د متناهي شاهد خاصيت]
\ollabel{thm:compactness-direct} \LR{$\Gamma$} هله او يوازې هله د صدق
وړ دے چې په متناهي ډول د صدق وړ وي۔
\end{thm}
\olpseganchor{OLP-0136-B020}{end}

\olpseganchor{OLP-0136-B021}{start}
\begin{proof}
که~\LR{$\Gamma$} د صدق وړ وي، نو داسې
جوړښت~\LR{$\Struct{M}$}
شته چې د ټولو~\LR{$!A \in \Gamma$} دپاره
\LR{$\Sat{M}{!A}$} وي۔ څرګنده ده چې همدا
\LR{$\Struct{M}$} د~\LR{$\Gamma$} هر متناهي
فرعي سټ هم پوره کوي؛ نو~\LR{$\Gamma$} په متناهي ډول د صدق وړ دے۔
\olpseganchor{OLP-0136-B021}{end}

\olpseganchor{OLP-0136-B022}{start}
اوس فرض کړئ~\LR{$\Gamma$} په متناهي ډول د صدق وړ دے۔ د
  \olref{lem:fsat-henkin} له مخې، داسې په متناهي ډول د صدق وړ
  مشبوع سټ~\LR{$\Gamma' \supseteq \Gamma$} شته۔ د
\olref{lem:fsat-lindenbaum} له مخې،
\LR{$\Gamma'$} يو بشپړ او په متناهي ډول د
صدق وړ سټ~\LR{$\Gamma^*$} ته غځېدلے شي، او~\LR{$\Gamma^*$} لا
هم مشبوع دے۔ که عينيت موجود نۀ وي، د
ترمي مدل~\LR{$\Struct{M(\Gamma^*)}$}؛ او که عينيت موجود وي، د مخکښې
برخې خارج قسمت شوے ترمي مدل د
\olref[mod]{defn:termmodel} د جوړونې پر بنسټ جوړ کړئ۔
پام وکړئ چې~\olref[mod]{prop:quant-termmodel} د~\LR{$\Gamma^*$} پر
سازګارۍ (يا، په دې حساب، پر بشپړ والي يا مشبوع والي) تکيه
نۀ کوله؛ يوازې پر دې ئې تکيه کوله چې~\LR{$\Struct{M(\Gamma^*)}$} پوښلے
دے۔ د عينيت
په حالت کښې د هم ارزښتۍ او سازګارتيا شرطونه هم نېغ له متناهي صدق
وړتيا څخه ترلاسه کېږي: د هر شرط ناکامي به يو د صدق وړ نۀ متناهي
فرعي سټ ورکړي۔ د صدق لمې ثبوت
(\olref[mod]{lem:truth}) هغه وخت هم سم پر مخ ځي چې
\olref[ccs]{prop:ccs} او
\olref[hen]{prop:saturated-instances} ته مراجعې په ترتيب سره
\olref{prop:fsat-ccs} او~\olref{prop:fsat-instances} ته په مراجعو
بدلې کړو۔
\textbf{د ژباړې د نسخې سمون (OLCOM-009):} په سرچينه کښې د دې
ټاکنيز شرط درېيم تش ارګومېنټ پاتې شوے و؛ دلته په څرګنده توګه
ورزيات شو، څو ورپسې متن د هغه پر ځاے وانۀ خيستل شي۔ د عينيت په حالت کښې د برابرۍ استقرايي حالت د خارج قسمت
له تعريف څخه نېغ ترلاسه کېږي۔
\textbf{د ژباړې د نسخې سمون (OLCOM-007):} د سرچينې د اول مرتبې
نېغ ثبوت يوازې عادي ترمي مدل يادوي، چې د عينيت لرونکې جملې نۀ
پوښي؛ دلته د عينيت دپاره اړينه خارج قسمت جوړونه او د هغې متناهي
صدق-وړتيايي دليل څرګند شو۔
\end{proof}
\olpseganchor{OLP-0136-B022}{end}

\olpseganchor{OLP-0136-B023}{start}
\begin{prob}
د صدق لمې~(\olref[fol][com][mod]{lem:truth}) بشپړ ثبوت په هغه بڼه
وليکئ چې د~\olref[fol][com][cpd]{thm:compactness-direct} د ثبوت
دپاره اړينه ده۔
\end{prob}
\olpseganchor{OLP-0136-B023}{end}

\olpseganchor{OLP-0137-B004}{start}
\olfileid{fol}{com}{dls}
\olsection{د لوېنهايم--سکولم قضيه}
\olpseganchor{OLP-0137-B004}{end}

\olpseganchor{OLP-0137-B005}{start}
د لوېنهايم--سکولم قضيه وايي چې که يوه تيوري نامتناهي مدل ولري،
نو داسې يو مدل هم لري چې تر ډېره شمېرېدونکے لامتناهي وي۔ د دې خبرې يوه
سملاسي پايله دا ده چې د اول مرتبې منطق دا نۀ شي څرګندولے چې د
جوړښت کچه د شمېر نۀ وړ ده: هره تړلے فارمول يا د
تړلي فارمولونه هر سټ چې په ټولو د شمېر نۀ وړ جوړښتونه
کښې رښتونے وي، په کوم د شمېر وړ جوړښت کښې هم رښتونے دے۔
\olpseganchor{OLP-0137-B005}{end}

\olpseganchor{OLP-0137-B006}{start}
\begin{thm}
\ollabel{thm:downward-ls} که~\LR{$\Gamma$} سازګار وي، نو يو
د شمېر وړ مدل لري؛ يعنې په داسې جوړښت کښې د صدق وړ
دے چې د تعريف ساحه ئې يا متناهي يا شمېرېدونکے لامتناهي وي۔
\end{thm}
\olpseganchor{OLP-0137-B006}{end}

\olpseganchor{OLP-0137-B007}{start}
\begin{proof}
که~\LR{$\Gamma$} سازګار وي، د بشپړتيا قضيې ثبوت چې کوم
جوړښت~\LR{$\Struct M$} ورکوي، د هغه د تعريف ساحه~\LR{$\Domain{M}$}
د ژبې~\LR{$\Lang L$} د ترمونو له سټ څخه لويه نۀ ده۔ نو~\LR{$\Struct M$} تر
ډېره شمېرېدونکے لامتناهي دے۔
\end{proof}
\olpseganchor{OLP-0137-B007}{end}

\olpseganchor{OLP-0137-B008}{start}
\begin{thm}
\ollabel{noidentity-ls} که~\LR{$\Gamma$} د عينيت پرته د اول مرتبې د
منطق په ژبه کښې د~تړلي فارمولونه يو سازګار سټ وي، نو يو
شمېرېدونکے لامتناهي مدل لري؛ يعنې په داسې جوړښت کښې د صدق وړ
دے چې د تعريف ساحه ئې نامتناهي او د شمېر وړ وي۔
\end{thm}
\olpseganchor{OLP-0137-B008}{end}

\olpseganchor{OLP-0137-B009}{start}
\begin{proof}
که~\LR{$\Gamma$} سازګار وي او داسې هېڅ جمله نۀ لري چې عينيت پکښې
راغلے وي، نو د بشپړتيا قضيې ثبوت چې کوم
جوړښت~\LR{$\Struct M$} ورکوي، د هغه د تعريف ساحه~\LR{$\Domain{M}$}
کټ مټ د ژبې~\LR{$\Lang L'$} د ترمونو سټ دے۔ نو~\LR{$\Struct{M}$}
شمېرېدونکے لامتناهي دے، ځکه~\LR{$\Trm[L']$} هم داسې دے۔
\end{proof}
\olpseganchor{OLP-0137-B009}{end}

\olpseganchor{OLP-0137-B010}{start}
\begin{ex}[د سکولم پاراډوکس]
د زرملو--فرېنکل سټ تيوري~\LR{$\Log{ZFC}$} يو ډېر پياوړے چوکاټ دے چې
کابو ټولې رياضيکي خبرې پکښې څرګندېدلے شي، د سټونو د کچو په اړه
حقيقتونه هم۔ نو د بېلګې په توګه،~\LR{$\Log{ZFC}$} ثابتولے شي چې د
حقيقي عددونو سټ~\LR{$\Real$}،~د شمېر نۀ وړ دے؛ د کانټور قضيه
ثابتولے شي چې د هر سټ تواني سټ پخپله له هغه سټ څخه لوی دے، او داسې
نور۔ که~\LR{$\Log{ZFC}$} سازګاره وي، ټول مدلونه ئې نامتناهي دي؛ او
سربېره پر دې، ټول داسې غړي لري چې تيوري ئې په اړه وايي
چې~د شمېر نۀ وړ دي، لکه هغه عنصر چې د~\LR{$\Log{ZFC}$} دا قضيه
رښتونې کوي چې د طبيعي عددونو تواني سټ موجود دے۔ د
لوېنهايم--سکولم قضيې له مخې،~\LR{$\Log{ZFC}$}،~د شمېر وړ مدلونه هم
لري: داسې مدلونه چې «د شمېر نۀ وړ» سټونه لري، خو پخپله
د شمېر وړ دي۔
\end{ex}
\olpseganchor{OLP-0137-B010}{end}

\olpseganchor{OLP-0138-B004}{start}
\olpart{fol}{د لومړۍ درجې منطق}
\olpseganchor{OLP-0138-B004}{end}

\olpseganchor{OLP-0138-B005}{start}
\begin{editorial}
  په دې برخه کښې د لومړۍ درجې منطق مابعدمنطق تر بشپړتيا پورې
  څېړل کېږي۔ اوس مهال دا برخه د بياني منطق پر جلا بيان تکيه نۀ
  کوي؛ هر څۀ پکښې ثابتېږي۔ که د \texttt{FOL} ټېګ ناسم وي، د سرچينې
  فايلونه به د کميت ټاکونکو اړوند مواد وباسي (او
  ``جوړښت'' به په ``ارزښت ټاکنه''، \LR{$\Struct{M}$} به په
  \LR{$\pAssign{v}$}، او داسې نورو بدل کړي)۔ په حقيقت کښې، د بياني منطق
  د برخې زياتره مواد همدا د لومړۍ درجې منطق مواد دي چې د \texttt{FOL}
  ټېګ پکښې بند کړے شوے دے۔
\olpseganchor{OLP-0138-B005}{end}

\olpseganchor{OLP-0138-B006}{start}
  که د بياني منطق برخه هم شامله وي، نو ډېر تکرار ترې جوړېږي۔ خو
  پلان دا دے چې دا برخه د بياني منطق مواد په پام کښې نيولے شي (او
  مخکښې څېړل شوي مواد وباسي، او همدارنګه ثبوتونه د بياني برخې
  اړوندو ځايونو ته په حوالو لنډ کړي)۔
\end{editorial}
\olpseganchor{OLP-0138-B006}{end}











\olpseganchor{OLP-0138-B012}{start}
%
  

\olpseganchor{OLP-0138-B012}{end}

\olpseganchor{OLP-0138-B013}{start}
%
  

\olpseganchor{OLP-0138-B013}{end}

\olpseganchor{OLP-0138-B014}{start}
%
  

\olpseganchor{OLP-0138-B014}{end}

\olpseganchor{OLP-0138-B015}{start}
%
  

\olpseganchor{OLP-0138-B015}{end}

\olpseganchor{OLP-0139-B004}{start}
\olchapter{fol}{int}{د لومړۍ درجې منطق پېژندنه}
\olpseganchor{OLP-0139-B004}{end}

\olpseganchor{OLP-0140-B004}{start}
\olfileid{fol}{int}{fol}
\olpseganchor{OLP-0140-B004}{end}

\olpseganchor{OLP-0140-B005}{start}
\olsection{د لومړۍ درجې منطق}
\olpseganchor{OLP-0140-B005}{end}

\olpseganchor{OLP-0140-B006}{start}
تاسو غالباً د صوري منطق په لومړۍ پېژندنه کښې د لومړۍ درجې له منطق
سره بلد يئ۔\footnote{په حقيقت کښې موږ لږ و ډېر دا فرضوو چې تاسو
ورسره بلد يئ!{} که نۀ يئ، نو يو لا ابتدايي درسي کتاب، لکه
\emph{\texttt{forall x}} \citep{Magnus2021}، بيا کتلے شئ۔} ښايي تاسو دا د
``کميت ټاکونکو منطق'' يا ``پريديکاتي منطق'' په نامه پېژنئ۔ د لومړۍ
درجې منطق تر هر څه مخکښې يوه صوري ژبه ده۔ يعنې يو ټاکلے لغتځونډ
لري، او عبارتونه ئې د همدې لغتځونډ تارونه دي۔ خو هر تار جايز نۀ
دے۔ د جايزو عبارتونو بېلابېل ډولونه شته: ترمونه، فارمولونه او
تړلي فارمولونه۔ موږ په اصل کښې د لومړۍ درجې منطق له تړلي فارمولونه
سره لېوال يو: هغوئ د انګرېزي ژبې د جملو صوري سيال راکوي، او د هغوئ
په باب هغه پوښتنې کولے شو چې منطقپوه عموماً ورسره لېوال وي۔ مثلاً:
\begin{itemize}
    \item ايا \LR{$!B$} له~\LR{$!A$} څخه په منطقي ډول لازمېږي؟
    \item ايا \LR{$!A$} په منطقي ډول رښتونې، په منطقي ډول دروغجنه، که
اتفاقي ده؟
    \item ايا \LR{$!A$} او \LR{$!B$} هم‌ارزښته دي؟
\end{itemize}
\olpseganchor{OLP-0140-B006}{end}

\olpseganchor{OLP-0140-B007}{start}
دا پوښتنې په بنسټيز ډول د لومړۍ درجې منطق د تړلي فارمولونه د
``معنٰی'' په باب دي۔ مثلاً يو فيلسوف به دا پوښتنه، چې ايا \LR{$!B$} له
~\LR{$!A$} څخه په منطقي ډول لازمېږي، داسې وڅېړي: ايا کوم داسې حالت شته
چې \LR{$!A$} پکښې رښتونې خو~\LR{$!B$} دروغجنه وي (\LR{$!B$} له~\LR{$!A$} څخه نۀ
لازمېږي)، که هر هغه حالت چې \LR{$!A$} رښتونې کوي~\LR{$!B$} هم رښتونې کوي
(\LR{$!B$} له~\LR{$!A$} څخه لازمېږي)؟ خو لا راته نۀ دي ويل شوي چې ``حالت''
څۀ دے---دا د \emph{معنٰی پوهنې} کار دے۔ د لومړۍ درجې منطق معنٰی
پوهنه د فيلسوف د ``حالت'' د وجداني مفکورې، او هم---چې دا مهمه
ده---په يوه حالت کښې د تړلے فارمول~\LR{$!A$} د \emph{رښتونې کېدو}
رياضيکي دقيق مدل راکوي۔ هغه رياضيکي دقيق مدل چې جوړوو ئې
جوړښت بولو۔ هغه اړيکه چې ``پکښې رښتونې'' په دقيق ډول
بيانوي، د \emph{صدق} اړيکه بلل کېږي۔ نو موږ به د
تړلي فارمولونه~\LR{$!A$} او جوړښتونه~\LR{$\Struct{M}$} دپاره دا تعريفوو
چې ``\LR{$!A$} په~\LR{$\Struct{M}$} کښې صدق کوي'' (په سمبولونو کښې:
\LR{$\Sat{M}{!A}$})۔ له دې وروسته نور معنايي اصطلاحات، لکه ``ترې
لازمېږي'' يا ``په منطقي ډول رښتونې ده''، هم په دقيق ډول تعريفولے
شو۔ دا تعريفونه به ممکنه کړي چې په رياضيکي دقت سره، مثلاً، فيصله
وکړو چې ايا
\LR{$\lforall[x][(!A(x) \lif !B(x))], \lexists[x][!A(x)] \Entails \lexists[x][!B(x)]$}
ده که نۀ۔ ځواب به طبعاً ``هو'' وي۔
\textbf{د ژباړې د نسخې سمون (OLFOL-001):} په سرچينه کښې د دې
استلزام د لومړۍ عمومي جملې تړونکى قوس تر وروستۍ پايلې پورې ځنډېدلے
او کميت ټاکونکى ئې په ناسم ډول پر ټول استلزام خور کړے و؛ دلته قوس
د شرطي فارمول سمدستي وروسته تړل شو۔ که تاسو مخکښې د
انګرېزي جملو د لومړۍ درجې منطق په سمبولونو اړولو زده کړه کړې وي، دا
به مثلاً د ``ټولې مېږي حشرات دي، مېږي شته، نو حشرات شته'' د
سمبولونو بڼه وپېژنئ۔ دا په څرګند ډول معتبر استدلال دے؛ نو زموږ د
صوري ژبې د ``ترې لازمېږي'' رياضيکي مدل هم بايد همدا ځواب ورکړي۔
\olpseganchor{OLP-0140-B007}{end}

\olpseganchor{OLP-0140-B008}{start}
بله موضوع چې غالباً د صوري منطق له لومړۍ پېژندنې څخه مو ياده ده،
\emph{اشتقاقونه} دي۔ که مو د صوري منطق لومړے کورس لوستلے وي،
ښوونکي به درباندې د داسې اشتقاقونه د موندلو تمرين کړے وي؛
ښايي حتٰی داسې اشتقاق مو موندلے وي چې پورته استلزام
ښيي۔ د اشتقاقونه د ورکولو ډېرې بېلابېلې لارې شته: ښايي تاسو
``فطري استنتاج'' يا ``د رښتياوالي ونې'' کارولې وي، خو نورې ډېرې
لارې هم شته۔ د اشتقاق د نظامونو موخه داسې وسيلې برابرول دي
چې د منطقپوهانو پورته پوښتنې پرې ځواب شي: مثلاً د فطري استنتاج داسې
اشتقاق چې \LR{$\lforall[x][(!A(x) \lif !B(x))]$} او
\LR{$\lexists[x][!A(x)]$} ئې مقدمې وي او \LR{$\lexists[x][!B(x)]$} ئې پايله
(وروستۍ کرښه) وي، دا \emph{تصديقوي} چې \LR{$\lexists[x][!B(x)]$} په
منطقي ډول له \LR{$\lforall[x][(!A(x) \lif !B(x))]$} او
\LR{$\lexists[x][!A(x)]$} څخه لازمېږي۔ په دې پاراګراف کښې هم د
OLFOL-001 له مخې د هماغې عمومي جملې تړونکى قوس او د پايلې زيات قوس
سم شوي دي۔
\olpseganchor{OLP-0140-B008}{end}

\olpseganchor{OLP-0140-B009}{start}
خو دا ولې؟ په ظاهره د اشتقاق نظامونه له معنٰی پوهنې سره
هېڅ تړاو نۀ لري: د يوه صوري اشتقاق ورکول يوازې دا غواړي چې
سمبولونه د ځينو قاعدو له مخې سره منظم شي؛ هلته د ``حالتونو'' يا
``پکښې رښتونې'' هېڅ يادونه نشته۔ د اشتقاق د نظامونو او
معنٰی پوهنې تر منځ تړاو بايد په مابعدمنطقي څېړنه ثابت شي۔ پکار دا
ده چې مثلاً په رياضيکي ډول ثابته کړو: له مقدمو
\LR{$\lforall[x][(!A(x) \lif !B(x))]$} او \LR{$\lexists[x][!A(x)]$} څخه د
\LR{$\lexists[x][!B(x)]$} صوري اشتقاق هله او يوازې هله ممکن دے
چې \LR{$\lforall[x][(!A(x) \lif !B(x))]$} او
\LR{$\lexists[x][!A(x)]$} په ګډه \LR{$\lexists[x][!B(x)]$} لازم کړي۔ د
OLFOL-001 له مخې دلته هم د تکرار شوې شېما دوه ناسمه قوسونه سم شوي
دي۔ خو تر دې مخکښې چې دا کار وشي، د نحو او معنٰی پوهنې د تعريفونو
د سمولو دپاره ډېر زيارکښ کار کول پکار دي۔
\olpseganchor{OLP-0140-B009}{end}

\olpseganchor{OLP-0141-B004}{start}
\olfileid{fol}{int}{syn}
\olpseganchor{OLP-0141-B004}{end}

\olpseganchor{OLP-0141-B005}{start}
\olsection{نحو}
\olpseganchor{OLP-0141-B005}{end}

\olpseganchor{OLP-0141-B006}{start}
لومړے بايد دا په دقيق ډول وټاکو چې د سمبولونو کوم تارونه د لومړۍ
درجې منطق تړلي فارمولونه ګڼل کېږي۔ دا کار به وروسته وکړو؛ اوس به
يوازې د بېلګو له مخې مخکښې ځو۔ د لومړۍ درجې منطق بنسټيزې جوړښتي
خښتې---يعنې لغتځونډ---په دوو برخو وېشل کېږي۔ لومړۍ برخه هغه
سمبولونه دي چې په وسيله ئې ځانګړي څيزونه ټاکو يا د هغوئ په باب
ځانګړې خبرې کوو۔ څيزونه په ثابت سمبولونه ټاکو، او بيا د ټاکلو شوو
څيزونو په باب په پريديکاتونه څه وايو۔ مثلاً ښايي~\LR{$\Obj a$} د
ثابت سمبول په توګه د يوه څيز د ټاکلو دپاره وکاروو، او بيا د
تړلے فارمول~\LR{$\Atom{\Obj P}{\Obj a}$} په وسيله د هغه په باب څه
ووايو۔ که د ``\LR{$\Obj a$}'' او ``\LR{$\Obj P$}'' معناوې په ذهن کښې ولرئ،
\LR{$\Atom{\Obj P}{\Obj a}$} د انګرېزي ژبې د يوې جملې په توګه لوستلے
شئ (او غالباً مو د صوري منطق د لومړۍ زده کړې پر وخت همداسې کړي
دي)۔ کله چې د لومړۍ درجې منطق داسې ساده تړلي فارمولونه ولرو، د
لغتځونډ په دويمه برخه، يعنې منطقي سمبولونو (نښلوونکو او کميت
ټاکونکو)، ترې نورې پېچلې جملې جوړولے شو۔ نو مثلاً داسې عبارتونه
جوړولے شو لکه \LR{$(\Atom{\Obj P}{\Obj a} \land \Atom{\Obj Q}{\Obj b})$}
يا~\LR{$\lexists[\Obj x][\Atom{\Obj P}{\Obj x}]$}۔
\olpseganchor{OLP-0141-B006}{end}

\olpseganchor{OLP-0141-B007}{start}
د معنٰی پوهنې د دقيقو تعريفونو او د هغو د اشتقاق د نظامونو
د قاعدو د ورکولو دپاره چې سخته مابعدمنطقي څېړنه ئې غواړي، تر هر څه
مخکښې بايد په دقيق ډول تعريف کړو چې د لومړۍ درجې منطق
تړلے فارمول څه ګڼل کېږي۔ بنسټيزه مفکوره دومره اسانه ده: ځينې ساده
تړلي فارمولونه يوازې له پريديکاتونه او ثابت سمبولونه څخه جوړولے
شو، لکه~\LR{$\Atom{\Obj P}{\Obj a}$}۔ بيا له دغو څخه د نښلوونکو او کميت
ټاکونکو په وسيله پېچلې جملې جوړوو۔ خو هغه قواعد په دقيق ډول کوم دي
چې له مخې ئې د پېچلو تړلي فارمولونه جوړول راته جايز دي؟ دا بايد
وټاکل شي، کنه ``د لومړۍ درجې منطق تړلے فارمول'' مو په کافي دقت نۀ
دے تعريف کړے۔ څو مسئلې شته۔ لومړۍ دا چې سم تارونه د تړلي فارمولونه
په توګه وشمېرو۔ دويمه دا چې کار په داسې ډول وکړو چې د
\emph{ټولو} تړلي فارمولونه په باب رياضيکي ثبوتونه ورکولے شو۔ په پاے
کښې بايد پر تړلي فارمولونه د ځينو لومړنيو عملياتو دقيق تعريفونه هم
ورکړو، لکه ``په~\LR{$!A$} کښې هر~\LR{$\Obj x$} په~\LR{$\Obj b$} بدل کړه''۔ ستونزه
دا ده چې د لومړۍ درجې منطق کميت ټاکونکي او متغيرونه دا کار
بشپړ څرګند نۀ پرېږدي۔ مثلاً ايا
\LR{$\lexists[\Obj x][\Atom{\Obj P}{\Obj a}]$} بايد تړلے فارمول وګڼل
شي؟ او \LR{$\lexists[\Obj x][\lexists[\Obj x][\Atom{\Obj P}{\Obj x}]]$}
څنګه؟ ``په \LR{$(\Atom{\Obj P}{\Obj x} \land \lexists[\Obj x][\Atom{\Obj P}{\Obj x}])$}
کښې \LR{$\Obj x$} په~\LR{$\Obj b$} بدل کړه'' پايله بايد څه وي؟
\olpseganchor{OLP-0141-B007}{end}

\olpseganchor{OLP-0142-B004}{start}
\olfileid{fol}{int}{fml}
\olpseganchor{OLP-0142-B004}{end}

\olpseganchor{OLP-0142-B005}{start}
\section{فارمولونه}
\olpseganchor{OLP-0142-B005}{end}

\olpseganchor{OLP-0142-B006}{start}
د لومړۍ درجې منطق د تړلي فارمولونه د سخت تعريف او د متغيرونه
له کارولو څخه راولاړو شوو مسئلو د هوارولو دپاره به لاندې لاره
کاروو۔ لومړے د عبارتونو يو \emph{بل} سټ تعريفوو: فارمولونه۔ له
دې وروسته پکښې د متغيرونه ونډه څېړلے شو---او د ځينو نورو
مفکورو، يعنې ``ازادو'' او ``تړلو'' متغيرونه، پر بنسټ تعريفولے
شو چې تړلے فارمول څه ده (يعنې داسې فارمول چې ازاد
متغيرونه نۀ لري)۔ دا کار يوازې د ``تړلے فارمول'' د تعريف د
اسانه کولو دپاره نۀ کوو؛ بلکې د صدق د معنايي مفهوم د تعريف په
طريقه کښې به هم ډېر اړين وي۔
\olpseganchor{OLP-0142-B006}{end}

\olpseganchor{OLP-0142-B007}{start}
راځئ ``فارمول'' د لومړۍ درجې د يوې ساده ژبې دپاره تعريف کړو:
داسې ژبه چې يوازې يو پريديکات~\LR{$\Obj P$}، يو
ثابت سمبول~\LR{$\Obj a$}، او يوازې منطقي سمبولونه \LR{$\lnot$}، \LR{$\land$}
او~\LR{$\lexists$} لري۔ زموږ بشپړ تعريفونه به ډېر عام وي: نامتناهي ډېر
پريديکاتونه او ثابت سمبولونه به جايز کړو۔ په حقيقت کښې هغه
تابعې به هم څېړو چې له ثابت سمبولونه او متغيرونه سره
يوځاے ``ترمونه'' جوړوي۔ اوس دپاره~\LR{$\Obj a$} او متغيرونه زموږ يوازيني
ترمونه دي۔ البته نامتناهي ډېرو متغيرونه ته اړتيا لرو۔ په رسمي
ډول به \LR{$\Obj v_0$}، \LR{$\Obj v_1$}، \dots، د متغيرونو په توګه کاروو۔
\olpseganchor{OLP-0142-B007}{end}

\olpseganchor{OLP-0142-B008}{start}
\begin{defn}
د \emph{فارمولونه} سټ~\LR{$\Frm$} داسې تعريفېږي:
\begin{enumerate}
\item\ollabel{fmls-atom} \LR{$\Atom{\Obj P}{\Obj a}$} او \LR{$\Atom{\Obj
  P}{\Obj v_i}$} فارمولونه دي (\LR{$i \in \Nat$})۔
\olpseganchor{OLP-0142-B008}{end}

\olpseganchor{OLP-0142-B009}{start}
\item \ollabel{fmls-not}که \LR{$!A$} فارمول وي، نو \LR{$\lnot !A$}
  فارمول ده۔
\olpseganchor{OLP-0142-B009}{end}

\olpseganchor{OLP-0142-B010}{start}
\item که \LR{$!A$} او \LR{$!B$} فارمولونه وي، نو \LR{$(!A \land
  !B)$} فارمول ده۔
\olpseganchor{OLP-0142-B010}{end}

\olpseganchor{OLP-0142-B011}{start}
\item \ollabel{fmls-ex}که \LR{$!A$} فارمول او \LR{$x$}
  متغير وي، نو \LR{$\lexists[x][!A]$} فارمول ده۔
\olpseganchor{OLP-0142-B011}{end}

\olpseganchor{OLP-0142-B012}{start}
\item \ollabel{fmls-limit}بل هېڅ څيز فارمول نۀ دے۔
\end{enumerate}
\end{defn}
\olpseganchor{OLP-0142-B012}{end}

\olpseganchor{OLP-0142-B013}{start}
\olref{fmls-atom} راته وايي چې د هر \LR{$i \in \Nat$} دپاره
\LR{$\Atom{\Obj P}{\Obj a}$} او \LR{$\Atom{\Obj P}{\Obj v_i}$}
فارمولونه دي۔ دا تش په نامه \emph{اټومي} فارمولونه دي۔ هغوئ
د پيل دپاره څه راکوي۔ نور شرطونه له هغو څخه چې لا مو
جوړ کړي، د نويو فارمولونه د جوړولو لارې راکوي۔ نو مثلاً د
\olref{fmls-not} له مخې \LR{$\lnot \Atom{\Obj P}{\Obj v_2}$}
فارمول ده، ځکه \LR{$\Atom{\Obj P}{\Obj v_2}$} د
\olref{fmls-atom} له مخې لا فارمول ده۔ بيا د
\olref{fmls-ex} له مخې \LR{$\lexists[\Obj v_2][\lnot \Atom{\Obj P}{\Obj v_2}]$}
بله فارمول ده، او همداسې مخکښې۔ \olref{fmls-limit} راته وايي
چې \emph{يوازې} هغه تارونه فارمولونه ګڼل کېږي چې په همدې طريقه
ئې جوړولے شو۔ په ځانګړي ډول
\LR{$\lexists[\Obj v_0][\Atom{\Obj P}{\Obj a}]$} او
\LR{$\lexists[\Obj v_0][\lexists[\Obj v_0][\Atom{\Obj P}{\Obj a}]]$}
په رښتيا فارمولونه ګڼل کېږي، خو \LR{$(\lnot \Atom{\Obj P}{\Obj a})$}
نۀ ګڼل کېږي، ځکه دباندې قوسونه ئې زيات دي۔
\olpseganchor{OLP-0142-B013}{end}

\olpseganchor{OLP-0142-B014}{start}
د فارمولونه د تعريف دا طريقه \emph{استقرايي تعريف} بلل کېږي،
او اجازه راکوي چې د استقرايي ثبوت په هغه بڼه د فارمولونه په باب
څه ثابت کړو چې \emph{جوړښتي استقرا} بلل کېږي۔ دا دواړه په عمومي
ډول په \emph{د OpenLogic اړونده سرچينيزه برخه}\nobreakspace\LR{\href{https://github.com/OpenLogicProject/OpenLogic/blob/9620cc73f9c8e0ad003c514a5d3748f29611c4c0/content/methods/induction/inductive-definitions.tex\#L11}{\latinfont\scriptsize [OLP]}} او \emph{د OpenLogic اړونده سرچينيزه برخه}\nobreakspace\LR{\href{https://github.com/OpenLogicProject/OpenLogic/blob/9620cc73f9c8e0ad003c514a5d3748f29611c4c0/content/methods/induction/structural-induction.tex\#L11}{\latinfont\scriptsize [OLP]}}
کښې څېړل شوي دي؛ وروسته ثبوتونو ته له ننوتلو مخکښې ئې بيا وګورئ۔
بنسټيزه مفکوره دا ده: که غواړئ ثابته کړئ چې يو څۀ د ټولو
فارمولونه دپاره رښتيا دي، لومړے ښيئ چې د اټومي فارمولونه
دپاره رښتيا دي؛ بيا ښيئ چې \emph{که} د هر فارمول~\LR{$!A$}
(او~\LR{$!B$}) دپاره رښتيا وي، د \LR{$\lnot !A$}، \LR{$(!A \land !B)$} او
\LR{$\lexists[x][!A]$} دپاره \emph{هم} رښتيا دي۔ مثلاً په دې سره ثابتېږي
چې دا د \LR{$\lexists[\Obj v_2][\lnot \Atom{\Obj P}{\Obj v_2}]$} دپاره
رښتيا دي: له لومړۍ برخې پوهېږئ چې د اټومي
فارمول~\LR{$\Atom{\Obj P}{\Obj v_2}$} دپاره رښتيا دي۔ بيا له دويمې
برخې ترلاسه کوئ چې د \LR{$\lnot \Atom{\Obj P}{\Obj v_2}$} دپاره رښتيا
دي، او ورپسې بيا د \LR{$\lexists[\Obj v_2][\lnot \Atom{\Obj P}{\Obj v_2}]$}
په خپله دپاره رښتيا دي۔ ځکه چې ټول فارمولونه په استقرايي ډول له
اټومي فارمولونه څخه جوړېږي، دا طريقه د هر فارمول دپاره کار کوي۔
\olpseganchor{OLP-0142-B014}{end}

\olpseganchor{OLP-0143-B004}{start}
\olfileid{fol}{int}{sat}
\olpseganchor{OLP-0143-B004}{end}

\olpseganchor{OLP-0143-B005}{start}
\olsection{صدق}
\olpseganchor{OLP-0143-B005}{end}

\olpseganchor{OLP-0143-B006}{start}
کله چې پوه شو فارمولونه څه دي، سمدستي د لومړۍ درجې منطق معنٰی
پوهنې ته وړاندې تللے شو: دلته بنسټيز تعريف د جوړښت دے۔
زموږ د ساده ژبې دپاره جوړښت~\LR{$\Struct M$} يوازې درې برخې
لري: يو ناتش سټ~\LR{$\Domain{M}$} چې \emph{د تعريف ساحه} بلل کېږي، هغه
څيز چې~\LR{$\Obj a$} ئې په~\LR{$\Struct M$} کښې ټاکي، او هغه څيزونه چې
\LR{$\Obj P$} ورباندې په~\LR{$\Struct M$} کښې رښتونے دے۔ د~\LR{$\Obj a$} ټاکلے
څيز په~\LR{$\Assign{\Obj a}{M}$}، او د هغو څيزونو سټ چې~\LR{$\Obj P$}
ورباندې رښتونے دے په~\LR{$\Assign{\Obj P}{M}$} ښودل کېږي۔
جوړښت~\LR{$\Struct{M}$} همدا درې څيزونه لري: \LR{$\Domain{M}$}،
\LR{$\Assign{\Obj a}{M} \in \Domain{M}$} او
\LR{$\Assign{\Obj P}{M} \subseteq \Domain{M}$}۔ عمومي حالت به لا پېچلے
وي، ځکه ډېر پريديکاتونه او ثابت سمبولونه به وي،
پريديکاتونه له يوه څخه زيات ځايونه لرلے شي، او تابعې به
هم وي۔ \textbf{د ژباړې د نسخې سمون (OLFOL-002):} سرچينې دلته په
تېروتنه کښې ويلي چې ثابت سمبولونه له يوه څخه زيات ځايونه لرلے شي؛ د
څوځاييزوالي خاصيت د پريديکاتونو دے، لکه د همدې څپرکي دوه‌ځاييزه
اړيکه۔
\olpseganchor{OLP-0143-B006}{end}

\olpseganchor{OLP-0143-B007}{start}
دا د هغو فارمولونه د صدق د تعريف دپاره بس ده چې متغيرونه
نۀ لري۔ مفکوره دا ده چې داسې استقرايي تعريف ورکړو چې د
فارمولونه زموږ د تعريف طريقه منعکس کړي۔ ټاکو چې يو اټومي فارمول
کله په~\LR{$\Struct{M}$} کښې صدق کوي؛ بيا مثلاً دا ټاکو چې \LR{$\lnot !A$}
کله په~\LR{$\Struct{M}$} کښې صدق کوي، د دې پر بنسټ چې~\LR{$!A$} په~\LR{$\Struct{M}$} کښې صدق
کوي که نۀ۔ مثلاً داسې تعريفولے شو:
\begin{enumerate}
  \item \LR{$\Atom{\Obj P}{\Obj a}$} په~\LR{$\Struct{M}$} کښې هله صدق کوي
  چې \LR{$\Assign{\Obj a}{M} \in \Assign{\Obj P}{M}$} وي۔
  \item \LR{$\lnot !A$} په~\LR{$\Struct{M}$} کښې هله صدق کوي چې \LR{$!A$} په
  ~ \LR{$\Struct{M}$} کښې صدق نۀ کوي۔
  \item \LR{$(!A \land !B)$} په~\LR{$\Struct{M}$} کښې هله صدق کوي چې \LR{$!A$}
  په~\LR{$\Struct{M}$} کښې صدق وکړي، او \LR{$!B$} هم په~\LR{$\Struct{M}$} کښې
  صدق وکړي۔
\end{enumerate}
فرض کړئ \LR{$\Domain{M} = \{0, 1, 2\}$}، \LR{$\Assign{\Obj a}{M} = 1$}، او
\LR{$\Assign{\Obj P}{M} = \{1, 2\}$} دي۔ دا تعريف راته وايي چې
\LR{$\Atom{\Obj P}{\Obj a}$} په~\LR{$\Struct{M}$} کښې صدق کوي (ځکه
\LR{$\Assign{\Obj a}{M} =  1 \in \{1,2\} = \Assign{\Obj P}{M}$})۔
ورپسې راته وايي چې \LR{$\lnot \Atom{\Obj P}{\Obj a}$} په~\LR{$\Struct{M}$}
کښې صدق نۀ کوي؛ نو په خپل وار \LR{$\lnot\lnot \Atom{\Obj P}{\Obj a}$}
صدق کوي او \LR{$(\lnot \Atom{\Obj P}{\Obj a} \land \Atom{\Obj P}{\Obj a})$}
صدق نۀ کوي، او همداسې مخکښې۔
\olpseganchor{OLP-0143-B007}{end}

\olpseganchor{OLP-0143-B008}{start}
ستونزه هغه وخت پيدا کېږي چې د کميت ټاکونکو دپاره تعريف ورکول
غواړو: زړه مو غواړي داسې ووايو چې
``\LR{$\lexists[\Obj v_0][\Atom{\Obj P}{\Obj v_0}]$} هله صدق کوي چې
\LR{$\Atom{\Obj P}{\Obj v_0}$} صدق وکړي''۔ خو جوړښت~\LR{$\Struct{M}$}
راته نۀ وايي چې له متغيرونه سره څه وکړو۔ په اصل کښې غواړو
ووايو چې \LR{$\Atom{\Obj P}{\Obj v_0}$} \emph{د~\LR{$\Obj v_0$} د کوم قيمت
دپاره} صدق کوي۔ د دې د دقيق کولو دپاره داسې لارې ته اړتيا لرو چې
د~\LR{$\Domain{M}$} غړي يوازې~\LR{$\Obj a$} ته نۀ، بلکې~\LR{$\Obj v_0$}
ته هم وګومارو۔ د همدې دپاره د متغير \emph{ګومارنې} ورپېژنو۔
د متغير ګومارنه يوازې داسې تابع~\LR{$s$} ده چې متغيرونه
د~\LR{$\Domain{M}$} غړي ته وړي (زموږ په بېلګه کښې له~\LR{$0$}، \LR{$1$}
يا~\LR{$2$} څخه يوه ته)۔ \textbf{د ژباړې د نسخې سمون (OLFOL-003):}
سرچينې دلته يو، دوه او درې ليکلي وو، خو پورته ټاکلې د تعريف ساحه
صفر، يو او دوه لري؛ نو لړ د هماغې ساحې غړو ته سم شو۔ ځکه مخکښې نۀ
پوهېږو چې په فارمول کښې به کوم متغيرونه راشي، دا نۀ شو
محدودولے چې~\LR{$s$} کومو متغيرونه ته قيمتونه وټاکي۔ ساده حل دا
دے چې وغواړو~\LR{$s$} \emph{ټولو} متغيرونه~\LR{$\Obj v_0$}،
\LR{$\Obj v_1$}، \dots\@ ته قيمتونه وګوماري۔ موږ به يوازې د اړينو هغو
قيمتونه کاروو۔
\olpseganchor{OLP-0143-B008}{end}

\olpseganchor{OLP-0143-B009}{start}
د دې پر ځاے چې د فارمولونه صدق يوازې د جوړښت په نسبت
تعريف کړو، دا به د جوړښت~\LR{$\Struct{M}$} \emph{او} د
متغير ګومارنې~\LR{$s$} دواړو په نسبت تعريف کړو، او په لنډ ډول به
\LR{$\Sat{M}{!A}[s]$} ليکو۔ اوس به زموږ تعريف د هغو اټومي فارمولونه
دپاره يو اضافي شرط هم ولري چې متغيرونه لري:
\begin{enumerate}
  \item \LR{$\Sat{M}{\Atom{\Obj P}{\Obj a}}[s]$} هله او يوازې هله چې
  \LR{$\Assign{\Obj a}{M} \in \Assign{\Obj P}{M}$} وي۔
  \item \LR{$\Sat{M}{\Atom{\Obj P}{\Obj v_i}}[s]$} هله او يوازې هله چې
  \LR{$s(\Obj v_i) \in \Assign{\Obj P}{M}$} وي۔
  \item \LR{$\Sat{M}{\lnot !A}[s]$} هله او يوازې هله چې
  \LR{$\Sat{M}{!A}[s]$} نۀ وي۔
  \item \LR{$\Sat{M}{(!A \land !B)}[s]$} هله او يوازې هله چې
  \LR{$\Sat{M}{!A}[s]$} او \LR{$\Sat{M}{!B}[s]$} دواړه وي۔
\end{enumerate}
ښه، په دې سره يوه ستونزه هواره شوه: اوس ويلے شو چې~\LR{$\Struct{M}$} د
\LR{$s(\Obj v_0)$} قيمت دپاره \LR{$\Atom{\Obj P}{\Obj v_0}$} کله پوره کوي۔
د \LR{$\lexists[\Obj v_0][\Atom{\Obj P}{\Obj v_0}]$} د تعريف د سمولو
دپاره يو بل کار پکار دے: غواړو
\LR{$\Sat{M}{\lexists[\Obj v_0][\Atom{\Obj P}{\Obj v_0}]}[s]$} هله او
يوازې هله وي چې د~\LR{$\Obj v_0$} د قيمت د ګومارلو د \emph{کومې} لارې
\LR{$s'$} دپاره \LR{$\Sat{M}{\Atom{\Obj P}{\Obj v_0}}[s']$} وي۔ خو~\LR{$\Obj v_0$}
ته ګومارلے قيمت لازماً هغه قيمت نۀ دے چې~\LR{$s(\Obj v_0)$} ئې ټاکي۔
د دې دپاره يوه نښه ورپېژنو: که \LR{$m \in \Domain{M}$} وي، نو
\LR{$\Subst{s}{m}{\Obj v_0}$} هغه ګومارنه ده چې (له~\LR{$\Obj v_0$} پرته د
ټولو متغيرونه دپاره) کټ مټ د~\LR{$s$} په شان وي، خو~\LR{$\Obj v_0$} ته
\LR{$m$} ګوماري۔ اوس زموږ تعريف داسې کېدے شي:
\begin{enumerate}\setcounter{enumi}{4}
  \item \LR{$\Sat{M}{\lexists[\Obj v_i][!A]}[s]$} هله او يوازې هله چې
  د کوم~\LR{$m \in \Domain{M}$} دپاره
  \LR{$\Sat{M}{!A}[\Subst{s}{m}{\Obj v_i}]$} وي۔
\end{enumerate}
ايا دا کار کوي؟ فرض کړئ د ټولو~\LR{$i \in \Nat$} دپاره
\LR{$s(\Obj v_i) = 0$} ټاکو۔
\LR{$\Sat{M}{\lexists[\Obj v_0][\Atom{\Obj P}{\Obj v_0}]}[s]$} هله او
يوازې هله چې داسې~\LR{$m \in \Domain{M}$} وي چې
\LR{$\Sat{M}{\Atom{\Obj P}{\Obj v_0}}[\Subst{s}{m}{\Obj v_0}]$} وي۔ او
داسې غړے شته: \LR{$m = 1$} يا \LR{$m = 2$} ټاکلے شو۔ پام وکړئ دا حتٰی هغه
وخت رښتيا دي چې~\LR{$s$} پخپله~\LR{$\Obj v_0$} ته ګومارلے قيمت
\LR{$s(\Obj v_0)$}---دلته~\LR{$0$}---کار نۀ کوي۔
\LR{$\Sat{M}{\Atom{\Obj P}{\Obj v_0}}[\Subst{s}{1}{\Obj v_0}]$} لرو، خو
\LR{$\Sat{M}{\Atom{\Obj P}{\Obj v_0}}[s]$} نۀ لرو۔
\olpseganchor{OLP-0143-B009}{end}

\olpseganchor{OLP-0143-B010}{start}
که دا مغشوش او دروند ښکاري، همداسې دے۔ خو دا زياته پېچلتيا د ټولو
فارمولونه د صدق د دقيق، استقرايي تعريف دپاره اړينه ده، او د
معنايي مفاهيمو د دقيق تعريف دپاره ورته څه پکار دي۔ نورې لارې هم
شته، خو ټولې په همدې کچه ښکلې (يا ناشاعرانه) دي۔
\olpseganchor{OLP-0143-B010}{end}

\olpseganchor{OLP-0144-B004}{start}
\olfileid{fol}{int}{snt}
\olpseganchor{OLP-0144-B004}{end}

\olpseganchor{OLP-0144-B005}{start}
\olsection{تړلي فارمولونه}
\olpseganchor{OLP-0144-B005}{end}

\olpseganchor{OLP-0144-B006}{start}
ښه، اوس د جوړښتونه او فارمولونه دپاره د صدق (``پکښې
رښتونې'') د تعريف يوه طرحه لرو۔ خو دا يو اضافي څيز---د
متغير ګومارنه---غواړي، حال دا چې زموږ غوښتنه د
تړلي فارمولونه تعريف و۔ له ګومارنو څنګه ځان خلاصوو، او تړلي فارمولونه
څه دي؟
\olpseganchor{OLP-0144-B006}{end}

\olpseganchor{OLP-0144-B007}{start}
غالباً د صوري منطق له لومړۍ پېژندنې څخه مو د متغيرونه او کميت
ټاکونکو د اړيکې بحث ياد وي۔ له کميت ټاکونکي څخه تل يو متغير
وروسته راځي، او بيا د تړلے فارمول په هغه برخه کښې چې کميت ټاکونکے
ورباندې پلي کېږي (د هغه ``ساحه'')، متغير د هماغه کميت ټاکونکي
له خوا ``تړلے'' ګڼو۔ په فارمولونه کښې دا شرط نۀ و چې هر
متغير دې ورته برابر کميت ټاکونکے ولري؛ بې له برابر کميت
ټاکونکي متغيرونه ``ازاد'' يا ``ناتړلي'' دي۔ موږ به ټول هغه
فارمولونه تړلي فارمولونه بولو چې ازاد متغيرونه نۀ لري۔
\olpseganchor{OLP-0144-B007}{end}

\olpseganchor{OLP-0144-B008}{start}
په فارمول~\LR{$!A$} کښې د متغير د يوې پېښې د تړل کېدو، د
هغې د تړونکي کميت ټاکونکي، او د هغې د ازاد پاتې کېدو وجداني مفکوره
بيا هم ګرانه نۀ ده۔ ښايي د دې د ازمېښت کومه طريقه مو زده کړې وي،
چې ښايي د قوسونو شمېرل هم پکښې وي۔ خو موږ بايد پر دقيق تعريف ټينګار
وکړو---او ځکه چې فارمولونه مو په استقرا تعريف کړي، په
فارمول~\LR{$!A$} کښې د متغير~\LR{$x$} ازادې او تړلې پېښې هم په
استقرا تعريفولے شو۔ مثلاً زموږ د ساده شوې ژبې دپاره داسې کېدے شي:
\begin{enumerate}
  \item که \LR{$!A$} اټومي وي، د~\LR{$x$} ټولې پېښې پکښې ازادې دي (يعنې په
  ~ \LR{$\Atom{\Obj P}{x}$} کښې د~\LR{$x$} پېښه ازاده ده)۔
  \item که \LR{$!A$} د~\LR{$\lnot !B$} بڼه ولري، نو په~\LR{$\lnot !B$} کښې د~\LR{$x$}
  يوه پېښه هله او يوازې هله ازاده ده چې په~\LR{$!B$} کښې د~\LR{$x$} ورته
  پېښه ازاده وي (يعنې په~\LR{$!B$} کښې د متغيرونو ازادې پېښې کټ مټ په
  ~ \LR{$\lnot !B$} کښې ورته پېښې دي)۔
  \item که \LR{$!A$} د~\LR{$(!B \land !C)$} بڼه ولري، نو په~\LR{$(!B \land !C)$}
  کښې د~\LR{$x$} يوه پېښه هله او يوازې هله ازاده ده چې په~\LR{$!B$} يا~\LR{$!C$}
  کښې د~\LR{$x$} ورته پېښه ازاده وي۔
  \item که \LR{$!A$} د~\LR{$\lexists[x][!B]$} بڼه ولري، نو په~\LR{$!A$} کښې د~\LR{$x$}
  هېڅ پېښه ازاده نۀ ده؛ او که د~\LR{$\lexists[y][!B]$} بڼه ولري، چې~\LR{$y$}
  له~\LR{$x$} څخه بېل متغير وي، نو په~\LR{$\lexists[y][!B]$} کښې د~\LR{$x$}
  يوه پېښه هله او يوازې هله ازاده ده چې په~\LR{$!B$} کښې د~\LR{$x$} ورته
  پېښه ازاده وي۔
\end{enumerate}
\olpseganchor{OLP-0144-B008}{end}

\olpseganchor{OLP-0144-B009}{start}
کله چې د متغيرونو د ازادو او تړلو پېښو دقيق تعريف ولرو، په ساده
ډول ويلے شو: تړلے فارمول هر هغه فارمول ده چې د
متغيرونه ازادې پېښې نۀ لري۔
\olpseganchor{OLP-0144-B009}{end}

\olpseganchor{OLP-0145-B004}{start}
\olfileid{fol}{int}{sem}
\olpseganchor{OLP-0145-B004}{end}

\olpseganchor{OLP-0145-B005}{start}
\olsection{معنايي مفاهيم}
\olpseganchor{OLP-0145-B005}{end}

\olpseganchor{OLP-0145-B006}{start}
پورته مو يادونه وکړه چې کله ګورو ايا \LR{$\Sat{M}{!A}[s]$} صدق کوي، د
اسانتيا دپاره~\LR{$s$} ته پرېږدو چې ټولو متغيرونه ته قيمتونه
وګوماري، خو يوازې هغه قيمتونه کارېږي چې په~\LR{$!A$} کښې
متغيرونه ته ئې ګوماري۔ په حقيقت کښې، يوازې په~\LR{$!A$} کښې د
\emph{ازادو} متغيرونو قيمتونه اهميت لري۔ البته ځکه چې موږ احتياط
کوو، دا حقيقت به ثابتوو۔ ځکه چې تړلي فارمولونه ازاد متغيرونه نۀ
لري، نو په جوړښت کښې د هغوئ د صدق دپاره~\LR{$s$} هېڅ اهميت نۀ
لري۔ نو کله چې~\LR{$!A$} تړلے فارمول وي، \LR{$\Sat{M}{!A}$} داسې تعريفولے
شو چې معنٰی ئې ``د ټولو~\LR{$s$} دپاره \LR{$\Sat{M}{!A}[s]$}'' ده؛ او تصادفاً
دا هله او يوازې هله رښتيا دي چې لږ تر لږه د يوه~\LR{$s$} دپاره
\LR{$\Sat{M}{!A}[s]$} وي۔ د فارمولونه د صدق د کار وړ تعريف دپاره د
متغير ګومارنې ورپېژندل پکار دي، خو د تړلي فارمولونه صدق د
متغير له ګومارنو څخه خپلواک دے۔
\olpseganchor{OLP-0145-B006}{end}

\olpseganchor{OLP-0145-B007}{start}
کله چې د ``\LR{$\Sat{M}{!A}$}'' تعريف ولرو، د لومړۍ درجې منطق د
تړلي فارمولونه په باب پوهېږو چې ``حالت'' او ``پکښې رښتونې'' څه معنٰی
لري۔ د تړلي فارمولونه دپاره د~\LR{$\Sat{M}{!A}$} د تعريف پر بنسټ بيا د
اعتبار، استلزام او صدق وړتيا بنسټيز معنايي مفاهيم تعريفولے شو۔ يوه
جمله، \LR{$\Entails !A$}، هله معتبره ده چې هره جوړښت ئې پوره
کړي۔ د تړلي فارمولونه د يوه سټ له خوا، \LR{$\Gamma \Entails !A$}،
هله لازمېږي چې هره هغه جوړښت چې په~\LR{$\Gamma$} کښې ټولې
تړلي فارمولونه پوره کوي،~\LR{$!A$} هم پوره کړي۔ او د تړلي فارمولونه يو سټ
هله د صدق وړ دے چې کومه جوړښت پکښې ټولې تړلي فارمولونه په
يو وخت پوره کړي۔
\olpseganchor{OLP-0145-B007}{end}

\olpseganchor{OLP-0145-B008}{start}
ځکه چې فارمولونه په استقرا تعريف شوي، او صدق بيا د فارمولونه
پر جوړښت په استقرا تعريفېږي، د خپلې معنٰی پوهنې د خاصيتونو د ثبوت
او د تعريف شوو معنايي مفاهيمو د تړلو دپاره استقرا کارولے شو۔ ځينې
دغه خاصيتونه به راټول او ثابت کړو؛ يوه وجه دا ده چې هر يو پخپله په
زړه پورې دے، خو اصلي وجه دا ده چې وروسته د معنٰی پوهنې او د
اشتقاق د نظامونو د اړيکې په څېړنه کښې ډېر راپکارېږي۔ د دې
کار دپاره به ځينې داسې نحوي مفاهيم او عمليات هم په دقيق ډول (يعنې
په استقرا) تعريفول پکار وي چې لا مو نۀ دي ياد کړي۔
\olpseganchor{OLP-0145-B008}{end}

\olpseganchor{OLP-0146-B004}{start}
\olfileid{fol}{int}{sub}
\olpseganchor{OLP-0146-B004}{end}

\olpseganchor{OLP-0146-B005}{start}
\olsection{ځايناستي}
\olpseganchor{OLP-0146-B005}{end}

\olpseganchor{OLP-0146-B006}{start}
يوه بېلګه څېړو څو څرګنده شي چې خبرې څنګه سره تړلې دي، او د نحو او
معنٰی پوهنې وده زموږ د وروستيو پرمختللو څېړنو بنسټ څنګه جوړوي۔ زموږ
د اشتقاق نظامونه بايد اجازه راکړي چې
\LR{$\Atom{\Obj P}{\Obj a}$} له
\LR{$\lforall[\Obj v_0][\Atom{\Obj P}{\Obj v_0}]$} څخه اشتقاق کړو۔
\textbf{د ژباړې د نسخې سمون (OLFOL-004):} د سرچينې په عمومي مقدمه
کښې اټومي فارمول په ناسم ډول د عمومي کميت ټاکونکي له دويم ارګومېنټ څخه
بهر ايستل شوے و؛ دلته اټومي فارمول د کميت ټاکونکي بشپړه بدنه ده،
لکه د همدې پاراګراف عمومي فارمول۔ ښايي وغواړو دا
حتٰی د استنتاج د قاعدې په توګه بيان کړو۔ خو د دې دپاره بايد په تر
ټولو عامو اصطلاحاتو کښې ئې ويلے شو: يوازې د~\LR{$\Obj P$}، \LR{$\Obj a$} او
\LR{$\Obj v_0$} دپاره نۀ، بلکې د هر فارمول~\LR{$!A$}، ترم~\LR{$t$} او
متغير~\LR{$x$} دپاره۔ (ياد مو وي چې ثابت سمبولونه ترمونه دي، خو
له ثابت سمبولونه او تابعې څخه جوړ شوي لا پېچلي ترمونه به هم
څېړو۔) نو غواړو داسې څه ووايو: ``هر کله چې مو
\LR{$\lforall[x][!A(x)]$} اشتقاق کړے وي، د~\LR{$!A(t)$} په استنتاج کښې په
حقه يئ---يعنې هغه پايله چې~\LR{$\lforall[x]$} لرې او~\LR{$x$} په~\LR{$t$} بدل
کړي''۔ خو ``\LR{$x$} په~\LR{$t$} بدلول'' په دقيق ډول څه معنٰی لري؟ د~\LR{$!A(x)$}
او~\LR{$!A(t)$} تر منځ اړيکه څه ده؟ ايا دا تل کار کوي؟
\olpseganchor{OLP-0146-B006}{end}

\olpseganchor{OLP-0146-B007}{start}
د دې د دقيق کولو دپاره د \emph{ځايناستي} عمل تعريفوو۔ ځايناستي په
اصل کښې نازکه ده، ځکه په~\LR{$!A$} کښې ټول~\LR{$x$} ګانې په ساده ډول په~\LR{$t$}
نۀ شو بدلولے، او هر~\LR{$t$} د هر~\LR{$x$} پر ځاے نۀ شي راتلے۔ دا به بيا په
استقرايي تعريفونو حل کوو۔ خو چې دا کار وشي، د استنتاج د قاعدې داسې
ټاکل چې ``له~\LR{$\lforall[x][!A(x)]$} څخه~\LR{$!A(t)$} استنتاج کړه'' په دقيق
تعريف بدلېږي۔ سربېره پر دې به وښودلے شو چې دا په دې معنٰی د
استنتاج ښه قاعده ده چې \LR{$\lforall[x][!A(x)]$}،~\LR{$!A(t)$} لازموي۔ خو د
دې د ثبوت دپاره بيا داسې څه ثابتول پکار دي چې په لومړي نظر مو دې
پوښتنې ته هڅوي: ``دا ولې کوو؟'' له~\LR{$\lforall[x][!A(x)]$} څخه د
~\LR{$!A(t)$} لازمېدل پر دې حقيقت ولاړ دي چې \LR{$\Sat{M}{!A(t)}$} صدق کوي
که نۀ، دا يوازې د ترم~\LR{$t$} په قيمت پورې تړلې ده؛ يعنې که~\LR{$m$} د
~\LR{$\Domain{M}$} هر هغه غړے وټاکو چې~\LR{$t$} ئې ټاکي، نو
\LR{$\Sat{M}{!A(t)}[s]$} هله او يوازې هله چې
\LR{$\Sat{M}{!A(x)}[\Subst{s}{m}{x}]$} وي۔ دا حتٰی هغه وخت هم صدق کوي
چې~\LR{$t$} متغيرونه ولري، خو بايد د پايلې په دقيق بيان کښې احتياط
وکړو۔
\olpseganchor{OLP-0146-B007}{end}

\olpseganchor{OLP-0147-B004}{start}
\olfileid{fol}{int}{mod}
\olpseganchor{OLP-0147-B004}{end}

\olpseganchor{OLP-0147-B005}{start}
\olsection{مدلونه او تيورۍ}
\olpseganchor{OLP-0147-B005}{end}

\olpseganchor{OLP-0147-B006}{start}
کله چې د لومړۍ درجې منطق نحو او معنٰی پوهنه تعريف کړو، د
جوړښتونه او معنايي مفاهيمو د خاصيتونو څېړل پيلولے شو۔ د
اشتقاق نظامونه هم تعريف او وڅېړلے شو۔ د تړلي فارمولونه د
يوه سټ په باب پوښتلے شو: کومې جوړښتونه د هغه سټ ټولې
تړلي فارمولونه رښتونې کوي؟ د تړلي فارمولونه يو سټ~\LR{$\Gamma$} چې راکړل
شي، هغه جوړښت~\LR{$\Struct{M}$} چې هغوئ پوره کوي د~\LR{$\Gamma$}
\emph{مدل} بلل کېږي۔ ښايي له~\LR{$\Gamma$} څخه پيل او د هغې مدلونه
ولټوو---هغوئ څنګه دي؟ بايد څومره لوے يا واړه وي؟ خو ښايي له يوې
جوړښت يا د جوړښتونه له يوې ټولګې څخه هم پيل او
پوښتنه وکړو: کومې تړلي فارمولونه پکښې رښتونې دي؟ ايا داسې
تړلي فارمولونه شته چې دا جوړښتونه په دې معنٰی
\emph{ځانګړې} کړي چې هغوئ، او يوازې هغوئ، پکښې رښتونې وي؟ دا ډول
پوښتنې د \emph{مدل تيورۍ} ډګر دے۔ همدغه پوښتنې د \emph{بديهي
طريقې} بنسټ هم جوړوي: د جوړښتونه د يوې ټولګې بيانول د
تړلي فارمولونه په يوه سټ، يعنې د يوې تيورۍ په بديهي اصولو۔ دا د دې
مشاهدې له امله ممکنېږي چې کټ مټ هغه تړلي فارمولونه د بديهي اصولو په
ټولو مدلونو کښې رښتونې دي چې د لومړۍ درجې منطق کښې له هغو بديهي
اصولو څخه لازمېږي۔
\olpseganchor{OLP-0147-B006}{end}

\olpseganchor{OLP-0147-B007}{start}
د يوې ډېرې ساده بېلګې په توګه مخکښ‌مرتبې وګورئ۔ مخکښ‌مرتبه پر کوم
سټ~\LR{$A$} داسې اړيکه~\LR{$R$} ده چې انعکاسي او انتقالي دواړه وي۔ سټ~\LR{$A$}
چې دوه‌ځاييزه اړيکه~\LR{$R \subseteq A \times A$} ورباندې وي، کټ مټ
هغه څه دي چې د يوه دوه‌ځاييز اړوند سمبول~\LR{$\Obj P$} لرونکې د لومړۍ
درجې ژبې د جوړښت دپاره ئې غواړو: \LR{$\Domain{M} = A$} او
\LR{$\Assign{\Obj P}{M} = R$} ټاکو۔ ځکه~\LR{$R$} مخکښ‌مرتبه ده، انعکاسي او
انتقالي ده، او د لومړۍ درجې منطق د تړلي فارمولونه داسې سټ~\LR{$\Gamma$}
موندلے شو چې دا خبره وايي:
\begin{align*}
  & \lforall[\Obj v_0][\Atom{\Obj P}{\Obj v_0,\Obj v_0}]\\
  & \lforall[\Obj v_0][\lforall[\Obj v_1][\lforall[\Obj v_2][((\Atom{\Obj P}{\Obj v_0,\Obj v_1} \land \Atom{\Obj P}{\Obj v_1,\Obj v_2}) \lif \Atom{\Obj P}{\Obj v_0,\Obj v_2})]]]
\end{align*}
دا تړلي فارمولونه يوازې د ``د هر~\LR{$x$} دپاره~\LR{$Rxx$}'' (\LR{$R$} انعکاسي
ده) او ``هر کله چې~\LR{$Rxy$} او~\LR{$Ryz$} وي، نو~\LR{$Rxz$} هم وي'' (\LR{$R$}
انتقالي ده) سمبولي بڼې دي۔ وينو چې جوړښت~\LR{$\Struct{M}$} د
دغو دوو تړلي فارمولونه~\LR{$\Gamma$} هله او يوازې هله مدل دے چې~\LR{$R$}
(يعنې~\LR{$\Assign{\Obj P}{M}$}) پر~\LR{$A$} (يعنې~\LR{$\Domain{M}$}) مخکښ‌مرتبه
وي۔ په بل عبارت، د~\LR{$\Gamma$} مدلونه کټ مټ مخکښ‌مرتبې دي۔ د ټولو
مخکښ‌مرتبو هر هغه خاصيت چې د~\LR{$\Obj P$} يوازيني پريديکات لرونکې
د لومړۍ درجې ژبه کښې بيانېدے شي (لکه پورته انعکاسيت او انتقاليت)،
په~\LR{$\Gamma$} کښې له دوو تړلي فارمولونه څخه لازمېږي، او برعکس۔ نو د
~\LR{$\Gamma$} د مدلونو په باب چې څه ثابتولے شو، د ټولو مخکښ‌مرتبو په
باب مو ثابت کړي دي۔
\olpseganchor{OLP-0147-B007}{end}

\olpseganchor{OLP-0147-B008}{start}
د هرې ځانګړې تيورۍ او د مدلونو د ټولګي (لکه~\LR{$\Gamma$} او ټولې
مخکښ‌مرتبې) دپاره به په اړونده د لومړۍ درجې ژبه کښې د بيانېدونکو او
نۀ بيانېدونکو څيزونو په باب په زړه پورې پوښتنې وي۔ د
جوړښتونه ځينې داسې خاصيتونه شته چې د ټولو ژبو او د مدلونو د
ټولګيو دپاره په زړه پورې دي، يعنې د د تعريف ساحه د کچې اړوند خاصيتونه۔
مثلاً د هر~\LR{$n \in \PosInt$} دپاره تل بيانولے شو چې د تعريف ساحه کټ مټ
\LR{$n$}~غړي لري۔ د نامتناهي ډېرو تړلي فارمولونه د يوه سټ په
کارولو دا هم بيانولے شو چې د تعريف ساحه نامتناهي ده۔ خو دا نۀ شو
بيانولے چې د تعريف ساحه متناهي يا د شمېر نۀ وړ ده۔ د لومړۍ
درجې ژبو د محدوديتونو دا پايلې د متناهي شاهد او
د لوېنهايم--سکولم له قضيو څخه راځي۔
\olpseganchor{OLP-0147-B008}{end}

\olpseganchor{OLP-0148-B004}{start}
\olfileid{fol}{int}{scp}
\olpseganchor{OLP-0148-B004}{end}

\olpseganchor{OLP-0148-B005}{start}
\olsection{سموالے او بشپړتيا}
\olpseganchor{OLP-0148-B005}{end}

\olpseganchor{OLP-0148-B006}{start}
د لومړۍ درجې منطق دپاره به د اشتقاق نظامونه هم ورپېژنو۔
منطقپوهانو ډېر د اشتقاق نظامونه جوړ کړي دي، خو هغوئ ټول د
تړلي فارمولونه تر منځ هماغه د د اشتقاق وړتيا اړيکه تعريفوي۔ وايو
چې~\LR{$\Gamma$}،~\LR{$!A$} اشتقاق کوي، \LR{$\Gamma \Proves !A$}، که د يوه
ټاکلي، په دقيق ډول تعريف شوي ډول اشتقاق موجود وي۔
اشتقاقونه تل د سمبولونو متناهي ترتيبونه دي---ښايي د
تړلي فارمولونه لړ وي، يا کوم لا پېچلے جوړښت۔ د اشتقاق د
نظامونو موخه داسې وسيله برابرول دي چې وټاکي ايا کوم سټ~\LR{$\Gamma$}
تړلے فارمول لازموي که نۀ۔ د دې موخې د پوره کولو دپاره بايد
\LR{$\Gamma \Entails !A$} هله او يوازې هله وي چې \LR{$\Gamma \Proves !A$}
وي۔
\olpseganchor{OLP-0148-B006}{end}

\olpseganchor{OLP-0148-B007}{start}
که \LR{$\Gamma \Proves !A$} وي خو \LR{$\Gamma \Entails !A$} نۀ وي، زموږ د
اشتقاق نظام به ډېر پياوړے وي او له حده زيات څه به ثابتوي۔
هغه خاصيت چې له~\LR{$\Gamma \Proves !A$} څخه \LR{$\Gamma \Entails !A$}
لازمېږي، \emph{سموالے} بلل کېږي، او د اشتقاق د هر ښه نظام
لږ تر لږه غوښتنه ده۔ له بلې خوا، که \LR{$\Gamma \Entails !A$} وي خو
\LR{$\Gamma \Proves !A$} نۀ وي، نو زموږ د اشتقاق نظام ډېر کمزورے
دے او کافي څه نۀ ثابتوي۔ هغه خاصيت چې له~\LR{$\Gamma \Entails !A$} څخه
\LR{$\Gamma \Proves !A$} لازمېږي، \emph{بشپړتيا} بلل کېږي۔ سموالے
عموماً نسبتاً اسانه ثابتېږي (د اشتقاقونه پر جوړښت د استقرا له
لارې، ځکه هغوئ په استقرا تعريف شوي دي)۔ بشپړتيا ثابتول ګران دي۔
\olpseganchor{OLP-0148-B007}{end}

\olpseganchor{OLP-0148-B008}{start}
سموالے او بشپړتيا څو مهمې پايلې لري۔ که د تړلي فارمولونه يو
سټ~\LR{$\Gamma$} کوم تناقض (لکه~\LR{$!A \land \lnot !A$}) اشتقاق کړي،
\emph{ناسازګار} بلل کېږي۔ ناسازګار~\LR{$\Gamma$} هېڅ مدل نۀ شي لرلے؛ د
صدق وړ نۀ دے۔ له بشپړتيا څخه ئې عکس لازمېږي: هر هغه~\LR{$\Gamma$} چې
ناسازګار نۀ وي---يا، لکه موږ چې به وايو، \emph{سازګار} وي---يو مدل
لري۔ په حقيقت کښې دا له بشپړتيا سره هم‌ارزښته ده، او موږ به د
بشپړتيا همدا بڼه په اصل کښې ثابتوو۔ دا ژوره او ښايي حيرانوونکې
پايله ده: يوازې دا چې له~\LR{$\Gamma$} څخه~\LR{$!A \land \lnot !A$} نۀ شئ
ثابتولے، تضمينوي چې داسې جوړښت شته چې کټ مټ د~\LR{$\Gamma$} د
بيان په شان ده۔ نو بشپړتيا دې پوښتنې ته ځواب ورکوي: د
تړلي فارمولونه کوم سټونه مدلونه لري؟ ځواب: ټول او يوازې سازګار
سټونه۔
\olpseganchor{OLP-0148-B008}{end}

\olpseganchor{OLP-0148-B009}{start}
د سموالي او بشپړتيا قضيې دوه مهمې پايلې لري: د متناهي شاهد قضيه او
د لوېنهايم--سکولم قضيه۔ دا د مدلونو په تيورۍ کښې مهمې پايلې دي
او د ډېرو په زړه پورې پايلو د ثبوت دپاره کارېدے شي۔ دوه مو لا يادې
کړې دي: د لومړۍ درجې منطق دا نۀ شي بيانولے چې د جوړښت
د تعريف ساحه متناهي يا د شمېر نۀ وړ ده۔
\olpseganchor{OLP-0148-B009}{end}

\olpseganchor{OLP-0148-B010}{start}
په تاريخي ډول د دې ټولو کارونو سم معلومول ډېر وخت ونيو---د لومړۍ
درجې منطق نحو او معنٰی پوهنه څنګه تعريف شي، د اشتقاق ښه
نظامونه څنګه تعريف شي، د هغوئ سموالے او بشپړتيا څنګه ثابته شي، او
د لومړۍ درجې ژبو د بيانېدونکو او نۀ بيانېدونکو څيزونو توپير څنګه
روښانه شي۔ اوس پوهېږو چې دا هر څه څنګه وکړو، خو په ټولو جزئياتو
تېرېدل لا هم مغشوشوونکي او ستومانوونکي کېدے شي۔ بيا هم مهم دي، ځکه
دلته د لومړۍ درجې منطق د صوري ژبې دپاره جوړې شوې طريقې په منطق،
کمپيوټر ساينس او ژبپوهنه کښې هر ځاے پلي کېږي۔ نو د جزئياتو زيار په
اوږده موده کښې ګټه کوي۔
\olpseganchor{OLP-0148-B010}{end}

\olpseganchor{OLP-0149-B004}{start}
\olchapter{fol}{syn}{د لومړۍ درجې منطق نحوه}
\olpseganchor{OLP-0149-B004}{end}

\olpseganchor{OLP-0150-B004}{start}
\olfileid{fol}{syn}{itx}
\olpseganchor{OLP-0150-B004}{end}

\olpseganchor{OLP-0150-B005}{start}
\olsection{پېژندنه}
\olpseganchor{OLP-0150-B005}{end}

\olpseganchor{OLP-0150-B006}{start}
د لومړۍ درجې منطق د تيورۍ او ميتاتيورۍ د پراختيا دپاره بايد تر هر څه
مخکښې د دې منطق د عبارتونو نحوه او معنٰی پوهنه تعريف کړو۔ د لومړۍ درجې
منطق عبارتونه ترمونه او فارمولونه دي۔ ترمونه له متغيرونه،
ثابت سمبولونه او تابعې څخه جوړېږي۔ فارمولونه بيا له
پريديکاتونه او ترمونو څخه جوړېږي (په دې ډول تر ټولو واړۀ، يعنې
``اټومي'' فارمولونه جوړېږي)، او بيا له اټومي فارمولونه څخه د
منطقي تړونکو او کميت ټاکونکو په وسيله نور پېچلي فارمولونه جوړولے شو۔ د
جوړېدو د قاعدو د ټاکلو ډېرې بېلابېلې لارې شته؛ موږ يوازې يوه ممکنه لاره
وړاندې کوو۔ نور نظامونه بېل سمبولونه غوره کوي، د تړونکو بېلې ټولګې د
لومړنيو په توګه ټاکي، او قوسونه په بله طريقه کاروي (يا، لکه د تش په نامه
پولنډۍ ليکنې په حالت کښې، ئې بيخي نۀ کاروي)۔ خو ټولې لارې په دې کښې
شريکې دي چې د جوړېدو قواعد د ترمونو او فارمولونه ټولګه
\emph{په استقرايي ډول} تعريفوي۔ که دا کار سم وشي، نو هر عبارت د جوړېدو
د قواعدو له مخې په بنسټيز ډول يوازې په يوه طريقه تر لاسه کېږي۔ هغه
استقرايي تعريف چې \emph{يوازېنۍ لوستنې وړ} عبارتونه رامنځته کوي، موږ ته
دا امکان راکوي چې همدې طريقې---يعنې استقرايي تعريف---سره دغو عبارتونو ته
معناوې هم ورکړو۔
\olpseganchor{OLP-0150-B006}{end}

\olpseganchor{OLP-0151-B004}{start}
\olfileid{fol}{syn}{fol}
\olpseganchor{OLP-0151-B004}{end}

\olpseganchor{OLP-0151-B005}{start}
\olsection{د لومړۍ درجې ژبې}
\olpseganchor{OLP-0151-B005}{end}


\olpseganchor{OLP-0151-B006}{start}
د لومړۍ درجې د منطق عبارتونه له يوه بنسټيز لغت څخه جوړېږي چې
\emph{متغيرونه}، \emph{ثابت سمبولونه}، \emph{پريديکاتونه}
او کله ناکله \emph{تابعې} لري۔ له دغو سمبولونو سره، د منطقي
نښلوونکو، سورو او د وقف د نښو لکه قوسونو او کامو په کارولو،
\emph{ترمونه} او \emph{فارمولونه} جوړېږي۔
\olpseganchor{OLP-0151-B006}{end}

\olpseganchor{OLP-0151-B007}{start}
\begin{explain}
په غير رسمي ډول، پريديکاتونه د ځانګړنو او اړيکو نومونه دي،
ثابت سمبولونه د انفرادي څيزونو نومونه دي، او تابعې د نقشو
نومونه دي۔ دا ټول، له عينيت~\LR{$\eq$} پرته، \emph{غير منطقي
سمبولونه} دي او په ګډه يوه ژبه جوړوي۔ هره د لومړۍ درجې ژبه~\LR{$\Lang L$}
د خپلو غير منطقي سمبولونو له مخې ټاکل کېږي۔ په تر ټولو عام حالت کښې،
\LR{$\Lang L$} د هر ډول سمبولونه په بې نهايته شمېر لري۔
\end{explain}
\olpseganchor{OLP-0151-B007}{end}

\olpseganchor{OLP-0151-B008}{start}
په عام حالت کښې موږ په لومړۍ درجې منطق کښې لاندې سمبولونه کاروو:
\olpseganchor{OLP-0151-B008}{end}

\olpseganchor{OLP-0151-B009}{start}
\begin{enumerate}
\item منطقي سمبولونه
\begin{enumerate}
\item منطقي نښلوونکي:
  \startycommalist
  \ycomma \LR{$\lnot$} (نفي)%
  \ycomma \LR{$\land$} (عطف)%
  \ycomma \LR{$\lor$} (فصل)%
  \ycomma \LR{$\lif$} (شرطي)%
  \ycomma \LR{$\liff$} (دوه اړخيز شرطي)%
  \ycomma \LR{$\lforall$} (عمومي سور)%
  \ycomma \LR{$\lexists$} (وجودي سور).
\item د دروغوالی~\LR{$\lfalse$} لپاره بياني ثابت۔
\item د رښتياوالی~\LR{$\ltrue$} لپاره بياني ثابت۔
\item دوه ځايه عينيت~\LR{$\eq$}۔
\item د متغيرونه يو شمېرېدونکے لامتناهي سټ: \LR{$\Obj v_0$}، \LR{$\Obj v_1$}، \LR{$\Obj
  v_2$}، \dots
\end{enumerate}
\item غير منطقي سمبولونه، چې د لومړۍ درجې منطق \emph{معياري
  ژبه} جوړوي
\begin{enumerate}
\item د هر \LR{$n>0$} لپاره د \LR{$n$}-ځايه پريديکاتونه يو شمېرېدونکے لامتناهي سټ: \LR{$\Obj
  A^n_0$}، \LR{$\Obj A^n_1$}، \LR{$\Obj A^n_2$}، \dots
\item د ثابت سمبولونه يو شمېرېدونکے لامتناهي سټ: \LR{$\Obj c_0$}، \LR{$\Obj c_1$}، \LR{$\Obj
  c_2$}، \dots۔
\item د هر \LR{$n>0$} لپاره د \LR{$n$}-ځايه تابعې يو شمېرېدونکے لامتناهي سټ:
  \LR{$\Obj f^n_0$}، \LR{$\Obj f^n_1$}، \LR{$\Obj f^n_2$}، \dots
\end{enumerate}
\item د وقف نښې: پرانيستے قوس (، تړلے قوس )، او کامه۔
\end{enumerate}
\olpseganchor{OLP-0151-B009}{end}

\olpseganchor{OLP-0151-B010}{start}
زموږ ډېری تعريفونه او پايلې به د لومړۍ درجې منطق د بشپړې معياري ژبې
لپاره بيانېږي۔ خو د کارونې له ډول سره سم ژبه يوازې څو پريديکاتونه،
ثابت سمبولونه او تابعې ته هم محدودولے شو۔
\olpseganchor{OLP-0151-B010}{end}

\olpseganchor{OLP-0151-B011}{start}
\begin{ex}
د حساب ژبه~\LR{$\Lang L_A$} يو دوه ځايه پريديکات~\LR{$<$}، يو
ثابت سمبول~\LR{$\Obj 0$}، يو يو ځايه تابع~\LR{$\prime$}، او دوه دوه
ځايه تابعې~\LR{$+$} او~\LR{$\times$} لري۔
\end{ex}
\olpseganchor{OLP-0151-B011}{end}

\olpseganchor{OLP-0151-B012}{start}
\begin{ex}
د سټ تيورۍ ژبه~\LR{$\Lang L_Z$} يوازې يو دوه ځايه پريديکات~\LR{$\in$}
لري۔
\end{ex}
\olpseganchor{OLP-0151-B012}{end}

\olpseganchor{OLP-0151-B013}{start}
\begin{ex}
د ترتيبونو ژبه~\LR{$\Lang L_\le$} يوازې يو دوه ځايه پريديکات~\LR{$\le$}
لري۔
\end{ex}
\olpseganchor{OLP-0151-B013}{end}

\olpseganchor{OLP-0151-B014}{start}
دا بيا هم يوازې دودونه دي: په رسمي ډول دا سمبولونه مستعار نومونه دي؛
د بېلګې په توګه، \LR{$<$}، \LR{$\in$} او \LR{$\le$} د \LR{$\Obj A^2_0$}، \LR{$\Obj 0$} د
\LR{$\Obj c_0$}، \LR{$\prime$} د \LR{$\Obj f^1_0$}، \LR{$+$} د \LR{$\Obj f^2_0$}، او \LR{$\times$} د
\LR{$\Obj f^2_1$} مستعار نومونه دي۔
\olpseganchor{OLP-0151-B014}{end}

\olpseganchor{OLP-0151-B015}{start}
.
\olpseganchor{OLP-0151-B015}{end}



\olpseganchor{OLP-0151-B017}{start}
% Alternate symbols
\olpseganchor{OLP-0151-B017}{end}

\olpseganchor{OLP-0151-B018}{start}
\begin{intro}
کېدے شي تاسو له هغو اصطلاحاتو او سمبولونو سره بلد يئ چې زموږ له پورته
کارولو سره توپير لري۔ د منطق متنونه، او ښوونکي، عموماً د ``نفي''
دپاره \LR{$\sim$}، \LR{$\neg$} او~!، او د ``عطف'' دپاره \LR{$\wedge$}، \LR{$\cdot$} او
\LR{$\&$} کاروي۔ د ``شرطي'' يا ``استلزام'' عام سمبولونه \LR{$\rightarrow$}،
\LR{$\Rightarrow$} او \LR{$\supset$} دي۔
د ``دوه اړخيز شرطي''، ``دوه اړخيز استلزام'' يا
  ``(مادي) هم ارزښت'' سمبولونه \LR{$\leftrightarrow$}،
  \LR{$\Leftrightarrow$} او \LR{$\equiv$} دي۔
د \LR{$\lfalse$} سمبول ته په بېلو ځايونو کښې
  ``دروغوالے''، ``فالسوم''، ``محال'' يا ``ښکته'' ويل کېږي۔
د \LR{$\ltrue$} سمبول ته په بېلو ځايونو کښې
  ``رښتياوالے''، ``ويروم'' يا ``پورته'' ويل کېږي۔
\olpseganchor{OLP-0151-B018}{end}

\olpseganchor{OLP-0151-B019}{start}
دود دا دے چې د لاتيني الفبا د پيل واړه توري (لکه \LR{$a$}، \LR{$b$}، \LR{$c$}) د
ثابت سمبولونه لپاره (چې کله ناکله نومونه بلل کېږي)، او د الفبا د پای
واړه توري (لکه \LR{$x$}، \LR{$y$}، \LR{$z$}) د متغيرونه لپاره وکارول شي۔ سورونه
له متغيرونه سره يوځای راځي، لکه \LR{$x$}؛ د ليکدود په بديلونو کښې د
عمومي سور لپاره \LR{$\forall x$}، \LR{$(\forall x)$}، \LR{$(x)$}، \LR{$\Pi x$}،
\LR{$\bigwedge_x$}، او د وجودي سور لپاره \LR{$\exists x$}، \LR{$(\exists
x)$}، \LR{$(Ex)$}، \LR{$\Sigma x$}، \LR{$\bigvee_x$} شامل دي۔
\end{intro}
\olpseganchor{OLP-0151-B019}{end}

\olpseganchor{OLP-0151-B020}{start}
\begin{explain}
موږ ټول بياني عملګرونه او دواړه سورونه د ژبې بنسټيز سمبولونه ګڼلے شو۔
يا د بنسټيزو سمبولونو يوه کوچنۍ ټولګه غوره کولے او نور منطقي عملګرونه
تعريف شوي ګڼلے شو۔ د بوليني عملګرونو په هغو ټولګو کښې چې ``د رښتياوالي
د تابع له مخې بشپړې'' دي، \LR{$\{ \lnot, \lor \}$}، \LR{$\{ \lnot, \land \}$}
او \LR{$\{ \lnot, \lif\}$} شامل دي---په بياني ځواک کښې د بشپړې لومړۍ
درجې ژبې دپاره له دې هره ټولګه له هر يوه سور سره يوځای کېدے شي۔
\olpseganchor{OLP-0151-B020}{end}

\olpseganchor{OLP-0151-B021}{start}
کېدے شي تاسو له دوو نورو منطقي عملګرونه سره هم اشنا يئ: د شېفر
کرښه~\LR{$|$} (چې د هنري شېفر په نوم نومول شوې ده)، او د پيرس
غشی~\LR{$\downarrow$}، چې د کواین خنجر هم بلل کېږي۔ کله چې دې عملګرونو ته په
ترتيب سره د ``نېنډ'' او ``نور'' معمولې لوستنې ورکړل شي، هر يو ئې په
يوازې ځان د رښتياوالي د تابع له مخې بشپړ وي۔
\end{explain}
\olpseganchor{OLP-0151-B021}{end}

\olpseganchor{OLP-0152-B004}{start}
\olfileid{fol}{syn}{frm}
\olpseganchor{OLP-0152-B004}{end}

\olpseganchor{OLP-0152-B005}{start}
\olsection{ترمونه او فارمولونه}
\olpseganchor{OLP-0152-B005}{end}

\olpseganchor{OLP-0152-B006}{start}
يوه د لومړۍ درجې ژبه~\LR{$\Lang L$} چې راکړل شي، د~\LR{$\Lang L$} له بنسټيز
لغت څخه جوړ شوي عبارتونه تعريفولے شو۔ په دغو کښې په ځانګړي ډول
\emph{ترمونه} او \emph{فارمولونه} شامل دي۔
\olpseganchor{OLP-0152-B006}{end}

\olpseganchor{OLP-0152-B007}{start}
\begin{defn}[ترمونه]
\ollabel{defn:terms}
د~\LR{$\Lang L$} د \emph{ترمونو} سټ~\LR{$\Trm[L]$} په استقرايي ډول داسې
تعريفېږي:
\begin{enumerate}
\item هر متغير يو ترم دے۔
\item د~\LR{$\Lang L$} هر ثابت سمبول يو ترم دے۔
\item که \LR{$f$} يو \LR{$n$}-ځايه تابع وي او \LR{$t_1$}، \dots، \LR{$t_n$}
  ترمونه وي، نو \LR{$\Atom{f}{t_1, \ldots, t_n}$} يو ترم دے۔
\item 
  بل هېڅ څيز ترم نۀ دے۔
\end{enumerate}
هغه ترم چې هېڅ متغيرونه نۀ لري، \emph{تړلے ترم} بلل کېږي۔
\end{defn}
\olpseganchor{OLP-0152-B007}{end}

\olpseganchor{OLP-0152-B008}{start}
\begin{explain}
ثابت سمبولونه زموږ د ژبې په مشخصاتو او د ترمونو په تعريف کښې د
سمبولونو د يوې جلا ډلې په توګه راځي، خو د صفر ځايه تابعې په
توګه هم ور شاملېدے شول۔ بيا د ترمونو د تعريف دويم شرط ته اړتيا نۀ
وه۔ يوازې دا بايد ومنو چې که \LR{$n = 0$} وي، نو
\LR{$\Atom{f}{t_1, \ldots, t_n}$} پخپله يوازې \LR{$f$} معنا لري۔
\end{explain}
\olpseganchor{OLP-0152-B008}{end}

\olpseganchor{OLP-0152-B009}{start}
\begin{defn}[فارمولونه]
\ollabel{defn:formulas}
د ژبې~\LR{$\Lang L$} د \emph{فارمولونه} سټ~\LR{$\Frm[L]$} په استقرايي
ډول داسې تعريفېږي:
\begin{enumerate}
\item \LR{$\lfalse$} يو اټومي فارمول دے۔
\olpseganchor{OLP-0152-B009}{end}

\olpseganchor{OLP-0152-B010}{start}
\item \LR{$\ltrue$} يو اټومي فارمول دے۔
\olpseganchor{OLP-0152-B010}{end}

\olpseganchor{OLP-0152-B011}{start}
\item که \LR{$R$} د~\LR{$\Lang L$} يو \LR{$n$}-ځايه پريديکات وي او \LR{$t_1$}، \dots،
  \LR{$t_n$} د~\LR{$\Lang L$} ترمونه وي، نو \LR{$\Atom{R}{t_1,\ldots, t_n}$} يو
  اټومي فارمول دے۔
\olpseganchor{OLP-0152-B011}{end}

\olpseganchor{OLP-0152-B012}{start}
\item که \LR{$t_1$} او \LR{$t_2$} د~\LR{$\Lang L$} ترمونه وي، نو
  \LR{$\Atom{\eq}{t_1, t_2}$} يو اټومي فارمول دے۔
\olpseganchor{OLP-0152-B012}{end}

\olpseganchor{OLP-0152-B013}{start}
\item که \LR{$!A$} فارمول وي، نو \LR{$\lnot !A$} هم
  فارمول دے۔
\olpseganchor{OLP-0152-B013}{end}

\olpseganchor{OLP-0152-B014}{start}
\item که \LR{$!A$} او \LR{$!B$} فارمولونه وي، نو \LR{$(!A \land
  !B)$} فارمول دے۔
\olpseganchor{OLP-0152-B014}{end}

\olpseganchor{OLP-0152-B015}{start}
\item که \LR{$!A$} او \LR{$!B$} فارمولونه وي، نو \LR{$(!A \lor !B)$}
  فارمول دے۔
\olpseganchor{OLP-0152-B015}{end}

\olpseganchor{OLP-0152-B016}{start}
\item که \LR{$!A$} او \LR{$!B$} فارمولونه وي، نو \LR{$(!A \lif !B)$}
  فارمول دے۔
\olpseganchor{OLP-0152-B016}{end}

\olpseganchor{OLP-0152-B017}{start}
\item که \LR{$!A$} او \LR{$!B$} فارمولونه وي، نو \LR{$(!A \liff !B)$}
  فارمول دے۔
\olpseganchor{OLP-0152-B017}{end}

\olpseganchor{OLP-0152-B018}{start}
\item که \LR{$!A$} فارمول او \LR{$x$} متغير وي، نو
  \LR{$\lforall[x][!A]$} فارمول دے۔
\olpseganchor{OLP-0152-B018}{end}

\olpseganchor{OLP-0152-B019}{start}
\item که \LR{$!A$} فارمول او \LR{$x$} متغير وي، نو
  \LR{$\lexists[x][!A]$} فارمول دے۔
\olpseganchor{OLP-0152-B019}{end}

\olpseganchor{OLP-0152-B020}{start}
\item بل هېڅ څيز فارمول نۀ دے۔
\end{enumerate}
\end{defn}
\olpseganchor{OLP-0152-B020}{end}

\olpseganchor{OLP-0152-B021}{start}
\begin{explain}
د ترمونو د سټ او د فارمولونه د سټ تعريفونه \emph{استقرايي
تعريفونه} دي۔ په اصل کښې موږ د فارمولونه سټ په نامتناهي ډېرو
پړاوونو کښې جوړوو۔ په لومړني پړاو کښې ټول اټومي فارمولونه فارمولونه
ګڼو؛ دا د تعريف له لومړنيو څو حالتونو سره سمون لري، يعنې د
\LR{$\ltrue$}، %
\LR{$\lfalse$}، %
\LR{$\Atom{R}{t_1,\dots,t_n}$} او \LR{$\Atom{\eq}{t_1,t_2}$} حالتونه۔ نو
``اټومي فارمول'' د همدې بڼو هر فارمول ته وايي۔
\olpseganchor{OLP-0152-B021}{end}

\olpseganchor{OLP-0152-B022}{start}
د تعريف نور حالتونه له هغو فارمولونه څخه چې لا جوړ شوي وي، د نويو
فارمولونه د جوړولو قاعدې راکوي۔ په دويم پړاو کښې د دغو قاعدو په
کارولو له اټومي فارمولونه څخه فارمولونه جوړولے شو۔ په درېيم
پړاو کښې له اټومي فارمولونو او په دويم پړاو کښې له ترلاسه شوو فارمولونو
څخه نوي فارمولونه جوړوو، او همداسې مخکښې۔ فارمول هر هغه څيز
دے چې بالاخره په داسې يوه پړاو کښې جوړ شي، او بل هېڅ څيز نۀ دے۔
\end{explain}
\olpseganchor{OLP-0152-B022}{end}

\olpseganchor{OLP-0152-B023}{start}
د دود له مخې، \LR{$\eq$} د خپلو مدخلو تر منځ ليکو او قوسونه پرېږدو:
\LR{$\eq[t_1][t_2]$} د \LR{$\Atom{\eq}{t_1,t_2}$} لنډون دے۔ همدارنګه
\LR{$\lnot \Atom{\eq}{t_1,t_2}$} د \LR{$\eq/[t_1][t_2]$} په بڼه لنډوو۔ کله
چې د دوه ځايه نښلوونکي~\LR{$\ast$} په کارولو له \LR{$!B$} او \LR{$!C$} څخه جوړ
شوے فارمول \LR{$(!B \ast !C)$} ليکو، ډېر ځله تر ټولو دباندې جوړه قوسونه
پرېږدو او يوازې~\LR{$!B \ast !C$} ليکو۔
\olpseganchor{OLP-0152-B023}{end}

\olpseganchor{OLP-0152-B024}{start}
\begin{intro}
د منطق ځينې متنونه د دې دپاره چې
\LR{$\lexists[x][!A]$} %
او %
\LR{$\lforall[x][!A]$} %
فارمولونه
وګڼل شي، دا شرط لګوي چې متغير~\LR{$x$} بايد په~\LR{$!A$} کښې راشي۔
که دا شرط ونۀ لګوئ، کومه ستونزه نۀ پيدا کېږي او کار اسانه کېږي۔
\end{intro}
\olpseganchor{OLP-0152-B024}{end}

\olpseganchor{OLP-0152-B036}{start}
کله چې د يوې ځانګړې کارونې په ژبه کښې کار کوو، ډېر ځله دوه ځايه
پريديکاتونه او تابعې د اړوندو ترمونو تر منځ ليکو؛ مثلاً د
حساب په ژبه کښې \LR{$t_1 < t_2$} او \LR{$(t_1 + t_2)$}، او د سټ تيورۍ په ژبه
کښې \LR{$t_1 \in t_2$} ليکو۔ د حساب په ژبه کښې د ناستي تابع د دود له مخې
آن د خپل مدخل \emph{وروسته} ليکل کېږي:~\LR{$t'$}۔ خو په رسمي ډول دا په
ترتيب سره يوازې د \LR{$\Atom{\Obj A^2_0}{t_1, t_2}$}،
\LR{$\Obj f^2_0(t_1, t_2)$}، \LR{$\Atom{\Obj A^2_0}{t_1, t_2}$} او
\LR{$\Obj f^1_0(t)$} دوديز لنډونونه دي۔
\textbf{د ژباړې د نسخې سمون (OLFOL-006):} سرچينه د ناستي تابع په
رسمي سمبول کښې د څيز نښه پرېږدي؛ دلته هغه د همدې فهرست له سمبوليزم
سره سمه ور زياته شوې ده۔
\olpseganchor{OLP-0152-B036}{end}

\olpseganchor{OLP-0152-B037}{start}
\begin{defn}[نحوي يوشانوالے]
سمبول \LR{$\ident$} د سمبولونو د تارونو تر منځ نحوي يوشانوالے څرګندوي؛
يعنې \LR{$!A \ident !B$} هله او يوازې هله چې \LR{$!A$} او \LR{$!B$} د يو شان
اوږدوالي لرونکي د سمبولونو تارونه وي او په هر ځاے کښې يو شان سمبول
ولري۔
\end{defn}
\olpseganchor{OLP-0152-B037}{end}

\olpseganchor{OLP-0152-B038}{start}
د \LR{$\ident$} سمبول دواړو خواوو ته هغه تارونه هم ليکل کېدے شي چې په نښلولو
ترلاسه شوي وي؛ مثلاً \LR{$!A \ident (!B \lor !C)$} دا معنا لري: د
سمبولونو تار~\LR{$!A$} له هغه تار سره يو شان دے چې په همدې ترتيب د پرانيستي
قوس، تار \LR{$!B$}، سمبول \LR{$\lor$}، تار~\LR{$!C$} او تړلي قوس له نښلولو ترلاسه
کېږي۔ که دا حالت وي، نو پوهېږو چې د~\LR{$!A$} لومړے سمبول پرانيستے قوس
دے، \LR{$!A$} د دويم سمبول له ځايه \LR{$!B$} د يوه فرعي تار په توګه لري، له
هغه فرعي تار وروسته~\LR{$\lor$} راځي، او همداسې نور۔
\olpseganchor{OLP-0152-B038}{end}

\olpseganchor{OLP-0152-B039}{start}
څرنګه چې ترمونه او فارمولونه د استقرايي تعريفونو له لارې له بنسټيزو
اجزاوو څخه جوړېږي، د هغو په باب د څه ثابتولو دپاره لاندې استقرايي
اصول کارولے شو۔
\olpseganchor{OLP-0152-B039}{end}

\olpseganchor{OLP-0152-B040}{start}
\begin{lem}[\emph{پر ترمونو د استقرا اصل}]
    \ollabel{lem:trmind}
    فرض کړئ \LR{$\Lang L$} د لومړۍ درجې يوه ژبه ده۔
    که يو خاصيت~\LR{$P$} داسې وي چې
    %
    \begin{enumerate}
        \item د هر متغير~\LR{$v$} دپاره صدق کوي؛
        %
        \item د~\LR{$\Lang L$} د هر ثابت سمبول~\LR{$a$} دپاره صدق کوي؛ او
        %
        \item کله چې د \LR{$t_1$}،~\dots، \LR{$t_n$} دپاره صدق کوي او \LR{$f$}
        د~\LR{$\Lang L$} يو \LR{$n$}-ځايه تابع وي، د
        \LR{$f(t_1,\dotsc,t_n)$} دپاره هم صدق کوي،
    \end{enumerate}
    (په دې فرض چې \LR{$t_1$}،~\dots، \LR{$t_n$} د~\LR{$\Lang{L}$} ترمونه دي)،
    نو \LR{$P$} د~\LR{$\Trm[L]$} د هر ترم دپاره صدق کوي۔
\end{lem}
\olpseganchor{OLP-0152-B040}{end}

\olpseganchor{OLP-0152-B041}{start}
\begin{prob}
    \olref[fol][syn][frm]{lem:trmind} ثابته کړئ۔
\end{prob}
\olpseganchor{OLP-0152-B041}{end}

\olpseganchor{OLP-0152-B042}{start}
\begin{lem}[\emph{پر فارمولونه د استقرا اصل}]
    \ollabel{thm:frmind}
    فرض کړئ \LR{$\Lang L$} د لومړۍ درجې يوه ژبه ده۔
    که يو خاصيت~\LR{$P$} د ټولو اټومي فارمولونه دپاره صدق کوي او داسې
    وي چې
    %
    \begin{enumerate}
        \item کله چې د~\LR{$!A$} دپاره صدق کوي، د
            \LR{$\lnot !A$} دپاره هم صدق کوي؛
        \item کله چې د \LR{$!A$} او~\LR{$!B$} دپاره صدق کوي، د
            \LR{$(!A \land !B)$} دپاره هم صدق کوي؛
        \item کله چې د \LR{$!A$} او~\LR{$!B$} دپاره صدق کوي، د
            \LR{$(!A \lor !B)$} دپاره هم صدق کوي؛
        \item کله چې د \LR{$!A$} او~\LR{$!B$} دپاره صدق کوي، د
            \LR{$(!A \lif !B)$} دپاره هم صدق کوي؛
        \item کله چې د \LR{$!A$} او~\LR{$!B$} دپاره صدق کوي، د
            \LR{$(!A \liff !B)$} دپاره هم صدق کوي؛
        \item کله چې د~\LR{$!A$} دپاره صدق کوي، د
            \LR{$\lexists[x][!A]$} دپاره هم صدق کوي؛
        \item کله چې د~\LR{$!A$} دپاره صدق کوي، د
            \LR{$\lforall[x][!A]$} دپاره هم صدق کوي؛
  \end{enumerate}
  (په دې فرض چې \LR{$!A$} او \LR{$!B$} د~\LR{$\Lang{L}$} فارمولونه دي)،
  نو \LR{$P$} د~\LR{$\Frm[L]$} د ټولو فارمولونو دپاره صدق کوي۔
\end{lem}
\olpseganchor{OLP-0152-B042}{end}

\olpseganchor{OLP-0153-B004}{start}
\olfileid{fol}{syn}{unq}
\olpseganchor{OLP-0153-B004}{end}

\olpseganchor{OLP-0153-B005}{start}
\olsection{يوازينے لوست}
\olpseganchor{OLP-0153-B005}{end}

\olpseganchor{OLP-0153-B006}{start}
\begin{explain}
موږ چې فارمولونه په کومه طريقه تعريف کړي دي، هغه دا تضمينوي چې هر
فارمول يو \emph{يوازينے لوست} لري؛ يعنې د فارمولونه د جوړولو
زموږ د قاعدو له مخې ئې د جوړولو په اصل کښې يوازې يوه طريقه شته، او د
``تفسيرولو'' ئې هم يوازې يوه طريقه شته۔ که داسې نۀ واے، مبهم
فارمولونه به مو لرل؛ يعنې داسې فارمولونه چې له يوه څخه زيات
لوستونه يا تفسيرونه لري---او څرګنده ده چې موږ ترې ډډه کول غواړو۔ خو
تر دې مهمه دا ده چې له دې خاصيت پرته به زموږ ډېری راتلونکي تعريفونه
او ثبوتونه کار ورنۀ کړي۔
\textbf{د ژباړې د نسخې سمون (OLFOL-014):} سرچينه په همدې جمله کښې
د «تفسير» د انګريزي کلمې املا ناسمه ليکي؛ ژباړه ئې سمه معنا راوړي۔
\olpseganchor{OLP-0153-B006}{end}

\olpseganchor{OLP-0153-B007}{start}
د دې خبرې د روښانولو ښايي غوره لار دا وي چې وګورو، که د فارمولونه
د جوړولو داسې ناسمې قاعدې مو ورکړې واے چې يوازينے لوست نۀ تضمينوي،
څه به پېښ شوي وو۔ مثلاً کېدے شي د نښلوونکو د جوړولو په قاعدو کښې مو
قوسونه هېر کړي او دا مو روا ګڼلي واے:
\begin{quote}
که \LR{$!A$} او \LR{$!B$} فارمولونه وي، نو \LR{$!A \lif !B$} هم فارمول دے۔
\end{quote}
له يوه اټومي فارمول \LR{$!D$} څخه په پيلولو به دې قاعدې د \LR{$!D \lif !D$}
جوړول را ته روا کړي وو۔ له دې او \LR{$!D$} څخه به مو \LR{$!D \lif !D \lif
!D$} ترلاسه کړے و۔ خو دا په دوو طريقو جوړېدے شي:
\begin{enumerate}
\item \LR{$!D$} د \LR{$!A$} په توګه، او \LR{$!D \lif !D$} د \LR{$!B$} په توګه اخلو۔
\item \LR{$!A$} د \LR{$!D \lif !D$} په توګه اخلو، او \LR{$!B$} پخپله~\LR{$!D$} دے۔
\end{enumerate}
په همدې ډول د فارمول~\LR{$!D \lif !D \lif !D$} د ``لوستلو'' دوه
طريقې شته۔ دا د \LR{$!B \lif !C$} بڼه لري چې \LR{$!B$} پکښې \LR{$!D$} او \LR{$!C$}
پکښې \LR{$!D \lif !D$} دے، خو \emph{همدارنګه} د \LR{$!B \lif !C$} بڼه لري
چې \LR{$!B$} پکښې \LR{$!D \lif !D$} او \LR{$!C$} پکښې~\LR{$!D$} دے۔
\olpseganchor{OLP-0153-B007}{end}

\olpseganchor{OLP-0153-B008}{start}
که داسې وشي، زموږ تعريفونه به تل کار نۀ کوي۔ مثلاً کله چې د يوه فارمول
اصلي عملګر تعريفوو، وايو: د \LR{$!B \lif !C$} بڼه لرونکي فارمول
اصلي عملګر د~\LR{$\lif$} همدا ښودل شوے مورد دے۔ خو که فارمول \LR{$!D
\lif !D \lif !D$} له \LR{$!B \lif !C$} سره په پورته يادو شوو دوو بېلو
طريقو برابرولے شو، نو په يوه حالت کښې د \LR{$\lif$} لومړے مورد او په بل
حالت کښې ئې دويم مورد د اصلي عملګر په توګه ترلاسه کوو۔ خو موږ
غواړو چې اصلي عملګر د فارمول يوه \emph{تابع} وي؛ يعنې هر
فارمول بايد کټ مټ د اصلي عملګر يو مورد ولري۔
\end{explain}
\olpseganchor{OLP-0153-B008}{end}

\olpseganchor{OLP-0153-B009}{start}
\begin{lem}
په فارمول~\LR{$!A$} کښې د کيڼو او ښيو قوسونو شمېر سره برابر دے۔
\end{lem}
\olpseganchor{OLP-0153-B009}{end}

\olpseganchor{OLP-0153-B010}{start}
\begin{proof}
دا د \LR{$!A$} د جوړېدو پر طريقه په استقرا ثابتوو۔ دې ته دوه څيزونه پکار
دي: (الف) لومړے بايد ثابته کړو چې ټول اټومي فارمولونه ياد خاصيت لري
(د استقرا بنسټ)۔ (ب) بيا بايد ثابته کړو چې کله له راکړل شوو فارمولونو
څخه نوي فارمولونه جوړوو، که زاړه فارمولونه ياد خاصيت ولري، نو نوي ئې
هم لري۔
\olpseganchor{OLP-0153-B010}{end}

\olpseganchor{OLP-0153-B011}{start}
\LR{$l(!A)$} دې په~\LR{$!A$} کښې د کيڼو قوسونو شمېر وي او \LR{$r(!A)$} دې د ښيو
قوسونو شمېر وي؛ \LR{$l(t)$} او \LR{$r(t)$} دې په همدې ډول په يوه ترم~\LR{$t$} کښې
د کيڼو او ښيو قوسونو شمېرونه وي۔
\olpseganchor{OLP-0153-B011}{end}

\olpseganchor{OLP-0153-B012}{start}
\begin{prob}
    ثابته کړئ چې د هر ترم~\LR{$t$} دپاره \LR{$l(t) = r(t)$}۔
\end{prob}
\olpseganchor{OLP-0153-B012}{end}

\olpseganchor{OLP-0153-B013}{start}
\begin{enumerate}
\item \indcase{!A}{\lfalse}{\LR{$\indfrm$} \LR{$0$} کيڼ او \LR{$0$}
    ښي قوسونه لري۔}
\olpseganchor{OLP-0153-B013}{end}

\olpseganchor{OLP-0153-B014}{start}
\item \indcase{!A}{\ltrue}{\LR{$\indfrm$} \LR{$0$} کيڼ او \LR{$0$}
    ښي قوسونه لري۔}
\olpseganchor{OLP-0153-B014}{end}

\olpseganchor{OLP-0153-B015}{start}
\item \indcase{!A}{\Atom{R}{t_1,\dots,t_n}}{\LR{$l(\indfrm) = 1 + l(t_1) +
  \dots + l(t_n) = 1 + r(t_1) + \dots + r(t_n) = r(\indfrm)$}۔ دلته
  هغه حقيقت کاروو چې د تمرين په توګه پرېښودل شوے دے: د هر ترم~\LR{$t$}
  دپاره \LR{$l(t) = r(t)$}۔}
\olpseganchor{OLP-0153-B015}{end}

\olpseganchor{OLP-0153-B016}{start}
\item \indcase{!A}{\eq[t_1][t_2]}{\LR{$l(\indfrm) = l(t_1) + l(t_2) =
  r(t_1) + r(t_2) = r(\indfrm)$}۔}
\olpseganchor{OLP-0153-B016}{end}

\olpseganchor{OLP-0153-B017}{start}
\item \indcase{!A}{\lnot !B}{د استقرايي فرض له مخې
    \LR{$l(!B) = r(!B)$}۔ نو \LR{$l(\indfrm) = l(!B) = r(!B) =
    r(\indfrm)$}۔}
\olpseganchor{OLP-0153-B017}{end}

\olpseganchor{OLP-0153-B018}{start}
\item \indcase{!A}{(!B \ast !C)}{د استقرايي فرض له مخې \LR{$l(!B) =
  r(!B)$} او \LR{$l(!C) = r(!C)$}۔ نو \LR{$l(\indfrm) = 1 + l(!B) + l(!C) =
  1 + r(!B) + r(!C) = r(\indfrm)$}۔}
\olpseganchor{OLP-0153-B018}{end}

\olpseganchor{OLP-0153-B019}{start}
\item \indcase{!A}{\lforall[x][!B]}{د استقرايي فرض له مخې
    \LR{$l(!B) = r(!B)$}۔ نو \LR{$l(\indfrm) = l(!B) = r(!B) =
    r(\indfrm)$}۔}
\olpseganchor{OLP-0153-B019}{end}

\olpseganchor{OLP-0153-B020}{start}
% Print case for \lexists if prvEx and not prvAll

\olpseganchor{OLP-0153-B020}{end}

\olpseganchor{OLP-0153-B021}{start}
% Just say similarly if prvEx and prvAll
\item \indcase{!A}{\lexists[x][!B]}{په همدې ډول۔}
\end{enumerate}
\end{proof}
\olpseganchor{OLP-0153-B021}{end}

\olpseganchor{OLP-0153-B022}{start}
\begin{defn}[حقيقي مخکښ]
د سمبولونو تار \LR{$!B$} د سمبولونو د تار~\LR{$!A$} \emph{حقيقي مخکښ} دے، که
له \LR{$!B$} سره د سمبولونو د يوه ناتش تار نښلول~\LR{$!A$} ترلاسه کړي۔
\end{defn}
\olpseganchor{OLP-0153-B022}{end}

\olpseganchor{OLP-0153-B023}{start}
\begin{lem}\ollabel{lem:no-prefix}
که \LR{$!A$} فارمول وي او \LR{$!B$} د \LR{$!A$} حقيقي مخکښ وي، نو \LR{$!B$}
فارمول نۀ دے۔
\end{lem}
\olpseganchor{OLP-0153-B023}{end}

\olpseganchor{OLP-0153-B024}{start}
\begin{proof}
تمرين۔
\end{proof}
\olpseganchor{OLP-0153-B024}{end}

\olpseganchor{OLP-0153-B025}{start}
\begin{prob}
\olref[fol][syn][unq]{lem:no-prefix} ثابته کړئ۔
\end{prob}
\olpseganchor{OLP-0153-B025}{end}

\olpseganchor{OLP-0153-B026}{start}
\begin{prop}
\ollabel{prop:unique-atomic}
که \LR{$!A$} يو اټومي فارمول وي، نو له لاندې شرطونو څخه يو، او يوازې
يو، پوره کوي۔
\begin{enumerate}
\item \LR{$!A \ident \lfalse$}۔
\item \LR{$!A \ident \ltrue$}۔
\item \LR{$!A \ident \Atom{R}{t_1,\dots,t_n}$}، چې \LR{$R$} يو \LR{$n$}-ځايه
  پريديکات، \LR{$t_1$}، \dots، \LR{$t_n$} ترمونه، او هر يو \LR{$R$}،
  \LR{$t_1$}، \dots، \LR{$t_n$} په يوازيني ډول ټاکل شوي دي۔
\item \LR{$!A \ident \eq[t_1][t_2]$}، چې \LR{$t_1$} او \LR{$t_2$} په يوازيني ډول
  ټاکل شوي ترمونه دي۔
\end{enumerate}
\end{prop}
\olpseganchor{OLP-0153-B026}{end}

\olpseganchor{OLP-0153-B027}{start}
\begin{proof}
تمرين۔
\end{proof}
\olpseganchor{OLP-0153-B027}{end}

\olpseganchor{OLP-0153-B028}{start}
\begin{prob}
\olref[fol][syn][unq]{prop:unique-atomic} ثابته کړئ۔ (لارښوونه: د
ترمونو دپاره د \olref[fol][syn][unq]{lem:no-prefix} يوه بڼه بيان او
ثابته کړئ۔)
\end{prob}
\olpseganchor{OLP-0153-B028}{end}

\olpseganchor{OLP-0153-B029}{start}
\begin{prop}[يوازينے لوست]
هر فارمول له لاندې شرطونو څخه يو، او يوازې يو، پوره کوي۔
\begin{enumerate}
\item \LR{$!A$} اټومي دے۔
\olpseganchor{OLP-0153-B029}{end}

\olpseganchor{OLP-0153-B030}{start}
\item \LR{$!A$} د \LR{$\lnot !B$} بڼه لري۔
\olpseganchor{OLP-0153-B030}{end}

\olpseganchor{OLP-0153-B031}{start}
\item \LR{$!A$} د \LR{$(!B \land !C)$} بڼه لري۔
\olpseganchor{OLP-0153-B031}{end}

\olpseganchor{OLP-0153-B032}{start}
\item \LR{$!A$} د \LR{$(!B \lor !C)$} بڼه لري۔
\olpseganchor{OLP-0153-B032}{end}

\olpseganchor{OLP-0153-B033}{start}
\item \LR{$!A$} د \LR{$(!B \lif !C)$} بڼه لري۔
\olpseganchor{OLP-0153-B033}{end}

\olpseganchor{OLP-0153-B034}{start}
\item \LR{$!A$} د \LR{$(!B \liff !C)$} بڼه لري۔
\olpseganchor{OLP-0153-B034}{end}

\olpseganchor{OLP-0153-B035}{start}
\item \LR{$!A$} د \LR{$\lforall[x][!B]$} بڼه لري۔
\olpseganchor{OLP-0153-B035}{end}

\olpseganchor{OLP-0153-B036}{start}
\item \LR{$!A$} د \LR{$\lexists[x][!B]$} بڼه لري۔
\end{enumerate}
سربېره پر دې، په هر حالت کښې \LR{$!B$}، يا \LR{$!B$} او \LR{$!C$}، په يوازيني ډول
ټاکل شوي دي۔ يعنې مثلاً د \LR{$!B$}، \LR{$!C$} او \LR{$!B'$}، \LR{$!C'$} داسې بېلې جوړې
نشته چې \LR{$!A$} په يو وخت د دواړو بڼه ولري:
\LR{$(!B \lif !C)$} او \LR{$(!B' \lif !C')$}۔
\end{prop}
\olpseganchor{OLP-0153-B036}{end}

\olpseganchor{OLP-0153-B037}{start}
\begin{proof}
د جوړولو قاعدې غواړي چې که فارمول اټومي نۀ وي، بايد په
پرانيستي قوس~(، \LR{$\lnot$}، يا په يوه سور پيل شي۔ بل
خوا، هر فارمول چې له لاندې سمبولونو څخه په يوه پيلېږي، بايد
اټومي وي: پريديکات، تابع، ثابت سمبول
، \LR{$\lfalse$}، \LR{$\ltrue$}۔
\olpseganchor{OLP-0153-B037}{end}

\olpseganchor{OLP-0153-B038}{start}
نو په اصل کښې يوازې دا ښودل پکار دي چې که \LR{$!A$} د \LR{$(!B \ast !C)$}
او هم د \LR{$(!B' \mathbin{\ast'} !C')$} بڼه لري، نو \LR{$!B \ident !B'$}،
\LR{$!C \ident !C'$}، او \LR{$\ast = {\ast'}$}۔
\olpseganchor{OLP-0153-B038}{end}

\olpseganchor{OLP-0153-B039}{start}
نو فرض کړئ چې هم \LR{$!A \ident (!B \ast !C)$} او هم \LR{$!A \ident (!B'
\mathbin{\ast'} !C')$}۔ اوس يا \LR{$!B \ident !B'$} دے، يا نۀ دے۔ که وي،
نو څرګنده ده چې \LR{$\ast = {\ast'}$} او \LR{$!C \ident !C'$}، ځکه بيا دا د
\LR{$!A$} هغه فرعي تارونه دي چې له يوه ځايه پيلېږي او يو شان اوږدوالے لري۔
بل حالت \LR{$!B \not\ident !B'$} دے۔ څرنګه چې \LR{$!B$} او \LR{$!B'$} دواړه د
\LR{$!A$} هغه فرعي تارونه دي چې له يوه ځايه پيلېږي، يو بايد د بل حقيقي
مخکښ وي۔ خو دا د \olref{lem:no-prefix} له مخې ناشونې ده۔
\end{proof}
\olpseganchor{OLP-0153-B039}{end}

\olpseganchor{OLP-0154-B004}{start}
\olfileid{fol}{syn}{mai}
\olpseganchor{OLP-0154-B004}{end}

\olpseganchor{OLP-0154-B005}{start}
\olsection{د فارمول اصلي عملګر}
\olpseganchor{OLP-0154-B005}{end}

\olpseganchor{OLP-0154-B006}{start}
\begin{explain}
ډېر ځله د هغه وروستي عملګر په باب خبرې کول ګټور وي چې د
فارمول~\LR{$!A$} په جوړولو کښې کارول شوے وي۔ دا عملګر د~\LR{$!A$}
\emph{اصلي عملګر} بلل کېږي۔ په شهودي ډول دا د~\LR{$!A$} ``تر ټولو دباندې''
عملګر دے۔ مثلاً د \LR{$\lnot !A$} اصلي عملګر \LR{$\lnot$} دے، د \LR{$(!A \lor
!B)$} اصلي عملګر \LR{$\lor$} دے، او همداسې نور۔
\end{explain}
\olpseganchor{OLP-0154-B006}{end}


\olpseganchor{OLP-0154-B007}{start}
\begin{defn}[اصلي عملګر]
\ollabel{def:main-op}
د فارمول~\LR{$!A$} \emph{اصلي عملګر} په لاندې ډول تعريفېږي:
\begin{enumerate}
\item \indcase*{!A}{!A}{\LR{$\indfrm$} هېڅ اصلي عملګر نۀ لري۔}
\olpseganchor{OLP-0154-B007}{end}

\olpseganchor{OLP-0154-B008}{start}
\item \indcase{!A}{\lnot !B}{د \LR{$\indfrm$} اصلي عملګر
  سمبول~\LR{$\lnot$} دے۔}
\olpseganchor{OLP-0154-B008}{end}

\olpseganchor{OLP-0154-B009}{start}
\item \indcase{!A}{(!B \land !C)}{د \LR{$\indfrm$}
  اصلي عملګر سمبول~\LR{$\land$} دے۔}
\olpseganchor{OLP-0154-B009}{end}

\olpseganchor{OLP-0154-B010}{start}
\item \indcase{!A}{(!B \lor !C)}{د \LR{$\indfrm$}
  اصلي عملګر سمبول~\LR{$\lor$} دے۔}
\olpseganchor{OLP-0154-B010}{end}

\olpseganchor{OLP-0154-B011}{start}
\item \indcase{!A}{(!B \lif !C)}{د \LR{$\indfrm$}
  اصلي عملګر سمبول~\LR{$\lif$} دے۔}
\olpseganchor{OLP-0154-B011}{end}

\olpseganchor{OLP-0154-B012}{start}
\item \indcase{!A}{(!B \liff !C)}{د \LR{$\indfrm$}
  اصلي عملګر سمبول~\LR{$\liff$} دے۔}
\olpseganchor{OLP-0154-B012}{end}

\olpseganchor{OLP-0154-B013}{start}
\item \indcase{!A}{\lforall[x][!B]}{د \LR{$\indfrm$}
  اصلي عملګر سمبول~\LR{$\lforall$} دے۔}
\olpseganchor{OLP-0154-B013}{end}

\olpseganchor{OLP-0154-B014}{start}
\item \indcase{!A}{\lexists[x][!B]}{د \LR{$\indfrm$}
  اصلي عملګر سمبول~\LR{$\lexists$} دے۔}
\end{enumerate}
\end{defn}
\olpseganchor{OLP-0154-B014}{end}

\olpseganchor{OLP-0154-B015}{start}
په هر حالت کښې زموږ موخه په فارمول کښې د اصلي عملګر همدا ښودل
شوے ځانګړے \emph{مورد} دے۔ مثلاً فارمول \LR{$((!D \lif !E) \lif (!E
\lif !D))$} د \LR{$(!B \lif !C)$} بڼه لري، چې \LR{$!B$} پکښې \LR{$(!D \lif !E)$}
او \LR{$!C$} پکښې \LR{$(!E \lif !D)$} دے؛ نو د \LR{$\lif$} دويم مورد ئې
اصلي عملګر دے۔
\olpseganchor{OLP-0154-B015}{end}

\olpseganchor{OLP-0154-B016}{start}
\begin{explain}
دا د هغې تابع يو \emph{بازګشتي} تعريف دے چې غير اټومي
فارمولونه د هغو د اصلي عملګر موردونو ته لېږي۔ څرنګه چې
فارمولونه په استقرايي ډول تعريف شوي دي، هر فارمول~\LR{$!A$} په
\olref{def:main-op} کښې له ورکړل شوو حالتونو څخه يو پوره کوي۔ دا
تضمينوي چې د هر غير اټومي فارمول~\LR{$!A$} دپاره اصلي عملګر موجود دے۔ څرنګه چې هر فارمول له دغو شرطونو څخه
يوازې يو پوره کوي، او په هر حالت کښې هغه کوچني فارمولونه په
يوازيني ډول ټاکل شوي دي چې \LR{$!A$} ترې جوړ شوے، د~\LR{$!A$} د اصلي عملګر مورد هم يوازينے دے؛ نو په رښتيا مو يوه تابع تعريف کړې ده۔
\end{explain}
\olpseganchor{OLP-0154-B016}{end}

\olpseganchor{OLP-0154-B017}{start}
موږ فارمولونه په
\olref{tab:main-op} کښې د هغو د اصلي عملګر د سمبول له مخې په
ورکړل شوو نومونو يادوو۔

\olpseganchor{OLP-0154-B017}{end}

\olpseganchor{OLP-0154-B018}{start}
\begin{table}[!h]
\centering
\begin{tabular}{c | c | c}
اصلي عملګر & د فارمول ډول & بېلګه\\
\hline
هېڅ & اټومي (فارمول) &
\LR{$\lfalse$},
\LR{$\ltrue$},
\LR{$\Atom{R}{t_1, \dots, t_n}$},
\LR{$\eq[t_1][t_2]$}\\
\LR{$\lnot$} & نفي & \LR{$\lnot !A$} \\
\LR{$\land$} & عطف & \LR{$(!A \land !B)$} \\
\LR{$\lor$} & فصل & \LR{$(!A \lor !B)$} \\
\LR{$\lif$} & شرطي & \LR{$(!A \lif !B$}) \\
\LR{$\liff$} & دوه اړخيز شرطي & \LR{$(!A \liff !B)$} \\
\LR{$\lforall[][]$} & عمومي (فارمول)& \LR{$\lforall[x][!A]$} \\
\LR{$\lexists[][]$} & وجودي (فارمول)& \LR{$\lexists[x][!A]$}
\end{tabular}
\caption{اصلي عملګر او د فارمولونه نومونه}
\ollabel{tab:main-op}
\end{table}
\textbf{د ژباړې د نسخې سمون (OLFOL-007):} سرچينه د عطف او فصل د
بېلګو تړونکي قوسونه د رياضي له چاپېريال څخه بهر ږدي؛ دلته دواړه قوسونه
له خپلو فارمولونو سره يو ځاے شوي دي۔
\olpseganchor{OLP-0154-B018}{end}

\olpseganchor{OLP-0155-B004}{start}
\olfileid{fol}{syn}{sbf}
\olpseganchor{OLP-0155-B004}{end}

\olpseganchor{OLP-0155-B005}{start}
\olsection{فرعي فارمولونه}
\olpseganchor{OLP-0155-B005}{end}

\olpseganchor{OLP-0155-B006}{start}
\begin{explain}
ډېر ځله د هغو فارمولونه په باب خبرې کول ګټور وي چې يو راکړل شوے
فارمول ``جوړوي''۔ دې ته د هغه \emph{فرعي فارمولونه} وايو۔ هر
فارمول~\LR{$!A$} د خپل ځان فرعي فارمول هم ګڼل کېږي؛ د \LR{$!A$} هر هغه
فرعي فارمول چې پخپله \LR{$!A$} نۀ وي، \emph{حقيقي فرعي فارمول} دے۔
\end{explain}
\olpseganchor{OLP-0155-B006}{end}

\olpseganchor{OLP-0155-B007}{start}
\begin{defn}[سمدستي فرعي فارمول]
که \LR{$!A$} فارمول وي، د \LR{$!A$} \emph{سمدستي فرعي فارمولونه} په
استقرايي ډول داسې تعريفېږي:
\begin{enumerate}
\item اټومي فارمولونه سمدستي فرعي فارمولونه نۀ لري۔
\olpseganchor{OLP-0155-B007}{end}

\olpseganchor{OLP-0155-B008}{start}
\item \indcase{!A}{\lnot !B}{د \LR{$\indfrm$} يوازينے سمدستي
    فرعي فارمول همدا~\LR{$!B$} دے۔}
\olpseganchor{OLP-0155-B008}{end}

\olpseganchor{OLP-0155-B009}{start}
\item \indcase{!A}{(!B \ast !C)}{د \LR{$\indfrm$} سمدستي
  فرعي فارمولونه همدا \LR{$!B$} او \LR{$!C$} دي (\LR{$\ast$} له دوه ځايه
  نښلوونکو څخه هر يو کېدے شي)۔}
\olpseganchor{OLP-0155-B009}{end}

\olpseganchor{OLP-0155-B010}{start}
\item \indcase{!A}{\lforall[x][!B]}{د \LR{$\indfrm$} يوازينے
    سمدستي فرعي فارمول همدا~\LR{$!B$} دے۔}
\olpseganchor{OLP-0155-B010}{end}

\olpseganchor{OLP-0155-B011}{start}
\item \indcase{!A}{\lexists[x][!B]}{د \LR{$\indfrm$} يوازينے
    سمدستي فرعي فارمول همدا~\LR{$!B$} دے۔}
\end{enumerate}
\end{defn}
\olpseganchor{OLP-0155-B011}{end}

\olpseganchor{OLP-0155-B012}{start}
\begin{defn}[حقيقي فرعي فارمول]
که \LR{$!A$} فارمول وي، د \LR{$!A$} \emph{حقيقي فرعي فارمولونه} په
بازګشتي ډول داسې تعريفېږي:
\begin{enumerate}
\item اټومي فارمولونه حقيقي فرعي فارمولونه نۀ لري۔
\olpseganchor{OLP-0155-B012}{end}

\olpseganchor{OLP-0155-B013}{start}
\item \indcase{!A}{\lnot !B}{د \LR{$\indfrm$} حقيقي
    فرعي فارمولونه همدا~\LR{$!B$} او د~\LR{$!B$} ټول حقيقي فرعي فارمولونه دي۔}
\olpseganchor{OLP-0155-B013}{end}

\olpseganchor{OLP-0155-B014}{start}
\item \indcase{!A}{(!B \ast !C)}{د \LR{$\indfrm$} حقيقي
  فرعي فارمولونه همدا \LR{$!B$}، \LR{$!C$}، د \LR{$!B$} ټول حقيقي فرعي فارمولونه
  او د~\LR{$!C$} ټول حقيقي فرعي فارمولونه دي۔}
\olpseganchor{OLP-0155-B014}{end}

\olpseganchor{OLP-0155-B015}{start}
\item \indcase{!A}{\lforall[x][!B]}{د \LR{$\indfrm$} حقيقي
    فرعي فارمولونه همدا~\LR{$!B$} او د~\LR{$!B$} ټول حقيقي
    فرعي فارمولونه دي۔}
\olpseganchor{OLP-0155-B015}{end}

\olpseganchor{OLP-0155-B016}{start}
\item \indcase{!A}{\lexists[x][!B]}{د \LR{$\indfrm$} حقيقي
    فرعي فارمولونه همدا~\LR{$!B$} او د~\LR{$!B$} ټول حقيقي
    فرعي فارمولونه دي۔}
\end{enumerate}
\end{defn}
\olpseganchor{OLP-0155-B016}{end}

\olpseganchor{OLP-0155-B017}{start}
\begin{defn}[فرعي فارمول]
د \LR{$!A$} فرعي فارمولونه پخپله \LR{$!A$} او د هغه ټول حقيقي
فرعي فارمولونه دي۔
\end{defn}
\olpseganchor{OLP-0155-B017}{end}

\olpseganchor{OLP-0155-B018}{start}
\begin{explain}
پام وکړئ چې سمدستي فرعي فارمولونه او حقيقي فرعي فارمولونه مو په
لږ توپير تعريف کړي دي۔ په لومړي حالت کښې مو د~\LR{$!A$} د هرې ممکنې بڼې
دپاره د هغه سمدستي فرعي فارمولونه نېغ په نېغه تعريف کړي دي۔ دا د
حالتونو له مخې يو څرګند تعريف دے، او حالتونه ئې د فارمولونه د سټ
له استقرايي تعريف سره سمون لري۔ په دويم حالت کښې مو هم د ټولو
فارمولونه د سټ د تعريف طريقه تعقيب کړې ده، خو په هر حالت کښې مو د
کوچنيو فارمولونه \LR{$!B$} او \LR{$!C$} د حقيقي فرعي فارمولونه تر څنګ،
پخپله دا فارمولونه هم ور شامل کړي دي۔ په دې سره تعريف \emph{بازګشتي}
کېږي۔ په عمومي ډول، پر يوه استقرايي تعريف شوي سټ (زموږ په حالت کښې
فارمولونه) د تابع تعريف هله بازګشتي وي چې د تابع د تعريف حالتونه
پخپله هماغه تابع وکاروي۔ خو د ښه تعريف دپاره بايد ډاډه شو چې د تابع
قيمتونه تل يوازې د هغو مدخلو دپاره کاروو چې تر تعريفېدونکې مدخل
``مخکښې'' راځي---زموږ په حالت کښې، د \LR{$(!B \ast !C)$} د ``حقيقي
فرعي فارمول'' د تعريف پر وخت يوازې د ``مخکښو'' فارمولونه \LR{$!B$}
او \LR{$!C$} حقيقي فرعي فارمولونه کاروو۔
\end{explain}
\olpseganchor{OLP-0155-B018}{end}

\olpseganchor{OLP-0155-B019}{start}
\begin{prop}
\ollabel{prop:subfrm-trans}
فرض کړئ \LR{$!B$} د \LR{$!A$} يو فرعي فارمول او \LR{$!C$} د \LR{$!B$} يو فرعي فارمول
دے۔ بيا \LR{$!C$} د \LR{$!A$} يو فرعي فارمول دے۔ په بله وينا، د فرعي فارمول
اړيکه انتقالي ده۔
\end{prop}
\olpseganchor{OLP-0155-B019}{end}

\olpseganchor{OLP-0155-B020}{start}
\begin{prob}
\olref[fol][syn][sbf]{prop:subfrm-trans} ثابته کړئ۔
\end{prob}
\olpseganchor{OLP-0155-B020}{end}

\olpseganchor{OLP-0155-B021}{start}
\begin{prop}
\ollabel{prop:count-subfrms}
فرض کړئ \LR{$!A$} داسې فارمول دے چې \LR{$n$} نښلوونکي او سورونه لري۔ بيا \LR{$!A$}
تر ډېره \LR{$2n+1$} فرعي فارمولونه لري۔
\end{prop}
\olpseganchor{OLP-0155-B021}{end}

\olpseganchor{OLP-0155-B022}{start}
\begin{prob}
\olref[fol][syn][sbf]{prop:count-subfrms} ثابته کړئ۔
\end{prob}
\olpseganchor{OLP-0155-B022}{end}

\olpseganchor{OLP-0156-B004}{start}
\olfileid{fol}{syn}{fseq}
\olpseganchor{OLP-0156-B004}{end}

\olpseganchor{OLP-0156-B005}{start}
\olsection{جوړښتي لړۍ}
\olpseganchor{OLP-0156-B005}{end}

\olpseganchor{OLP-0156-B006}{start}
په استقرايي تعريف سره د فارمولونه تعريفول، او د استقرا له لارې د
فارمولونه د خاصيتونو د ثابتولو بشپړوونکې طريقه، ښکلې او اغېزمنه لار
ده۔ خو کله کله دا هم ګټوره وي چې د فارمولونه جوړول له بېخه پورته،
ګام په ګام وڅېړو۔ دلته د \emph{جوړښتي لړۍ} په مفهوم همدا کار کوو۔
%
د دې د ښودلو دپاره چې ترمونه او فارمولونه څنګه له خپلو استقرايي
تعريفونو پرته په دې طريقه معرفي کېدے شي، لومړے د يوې ژبې~\LR{$\Lang L$}
له سمبولونو څخه د جوړ شوي خپل سري تار مفهوم معرفي کوو۔
\olpseganchor{OLP-0156-B006}{end}

\olpseganchor{OLP-0156-B007}{start}
\begin{defn}[تارونه]
\ollabel{defn:string}
فرض کړئ \LR{$\Lang L$} د لومړۍ درجې يوه ژبه ده۔ \emph{د~\LR{$\Lang L$} تار}
د~\LR{$\Lang L$} د سمبولونو يوه متناهي لړۍ ده۔ چې ژبه~\LR{$\Lang L$} له
چاپېريال څخه په څرګند ډول ټاکلې وي، د~\LR{$\Lang L$} تار ته به ډېر ځله
يوازې \emph{تار} وايو۔
\end{defn}
\olpseganchor{OLP-0156-B007}{end}

\olpseganchor{OLP-0156-B008}{start}
\begin{ex}
د لومړۍ درجې د هرې ژبې \LR{$\Lang L$} دپاره ټول
\LR{$\Lang L$}-فارمولونه د~\LR{$\Lang L$} تارونه دي، خو عکس ئې سم نۀ دے۔
مثلاً \[)(\Obj v_0\lif\lexists\] د~\LR{$\Lang L$} يو تار دے، خو د~\LR{$\Lang L$}
فارمول نۀ دے۔
\end{ex}
\olpseganchor{OLP-0156-B008}{end}

\olpseganchor{OLP-0156-B009}{start}
\begin{defn}[د ترمونو جوړښتي لړۍ]
\ollabel{defn:fseq-trm}
د~\LR{$\Lang L$} د تارونو متناهي لړۍ \LR{$\tuple{t_0,\dotsc,t_n}$} د يوه
ترم \LR{$t$} \emph{جوړښتي لړۍ} ده، که \LR{$t \ident t_n$} وي او د هر \LR{$i \leq
n$} دپاره يا \LR{$t_i$} متغير يا ثابت سمبول وي، او يا~\LR{$\Lang L$}
يو \LR{$k$}-ځايه تابع~\LR{$f$} ولري او داسې \LR{$m_1,\dotsc,m_k < i$}
موجود وي چې \LR{$t_i \ident f(t_{m_1},\dotsc,t_{m_k})$}۔ چې توپير ئې اړين
وي، د ترمونو جوړښتي لړۍ به \emph{د ترم جوړښتي لړۍ} بولو۔
\textbf{د ژباړې د نسخې سمون (OLFOL-008):} سرچينه د ټاکلي ځاييز والي
په پرتله په تېروتنه يو زيات شاخص او يو زيات مدخل شمېري؛ دلته دواړه
لړونه د تابع له ټاکلي ځاييز والي سره برابر شوي دي۔
\end{defn}
\olpseganchor{OLP-0156-B009}{end}

\olpseganchor{OLP-0156-B010}{start}
\begin{ex}
لړۍ
\[
    \tuple{\Obj c_0, \Obj v_0, \Atom{\Obj f^2_0}{\Obj c_0, \Obj v_0}, \Atom{\Obj f^1_0}{\Atom{\Obj f^2_0}{\Obj c_0, \Obj v_0}}}
\]
د ترم \LR{$\Atom{\Obj f^1_0}{\Atom{\Obj f^2_0}{\Obj c_0, \Obj v_0}}$}
يوه جوړښتي لړۍ ده، او دا هم ده:
\[
    \tuple{\Obj v_0, \Obj c_0, \Atom{\Obj f^2_0}{\Obj c_0, \Obj v_0}, \Atom{\Obj f^1_0}{\Atom{\Obj f^2_0}{\Obj c_0, \Obj v_0}}}.
\]
\end{ex}
\olpseganchor{OLP-0156-B010}{end}

\olpseganchor{OLP-0156-B011}{start}
\begin{defn}[د فارمولونو جوړښتي لړۍ]
\ollabel{defn:fseq-frm}
د~\LR{$\Lang L$} د تارونو متناهي لړۍ \LR{$\tuple{!A_0,\dotsc,!A_n}$} د~\LR{$!A$}
\emph{جوړښتي لړۍ} ده، که \LR{$!A \ident !A_n$} وي او د هر \LR{$i \leq n$}
دپاره يا \LR{$!A_i$} يو اټومي فارمول وي، يا داسې \LR{$j,k < i$} او
متغير~\LR{$x$} موجود وي چې له لاندې حالتونو څخه يو صدق وکړي:
\begin{enumerate}
    \item \LR{$!A_i \ident \lnot !A_j$}۔%
    \item \LR{$!A_i \ident (!A_j \land !A_k)$}۔%
    \item \LR{$!A_i \ident (!A_j \lor !A_k)$}۔%
    \item \LR{$!A_i \ident (!A_j \lif !A_k)$}۔%
    \item \LR{$!A_i \ident (!A_j \liff !A_k)$}۔%
    \item \LR{$!A_i \ident \lforall[x][!A_j]$}۔%
    \item \LR{$!A_i \ident \lexists[x][!A_j]$}۔%
\end{enumerate}
چې توپير ئې اړين وي، د فارمولونو جوړښتي لړۍ به \emph{د فارمول
جوړښتي لړۍ} بولو۔
\end{defn}
\olpseganchor{OLP-0156-B011}{end}

\olpseganchor{OLP-0156-B012}{start}
\begin{ex}
\[
\tuple{
    \Atom{\Obj A^1_0}{\Obj v_0},
    \Atom{\Obj A^1_1}{\Obj c_1},
    (\Atom{\Obj A^1_1}{\Obj c_1} \land \Atom{\Obj A^1_0}{\Obj v_0}),
    \lexists[\Obj v_0][(\Atom{\Obj A^1_1}{\Obj c_1} \land \Atom{\Obj A^1_0}{\Obj v_0})]
}
\]
د \LR{$\lexists[\Obj v_0][(\Atom{\Obj A^1_1}{\Obj c_1} \land \Atom{\Obj
A^1_0}{\Obj v_0})]$} يوه جوړښتي لړۍ ده، او دا هم ده:
\begin{multline*}
\tuple{
    \Atom{\Obj A^1_0}{\Obj v_0},
    \Atom{\Obj A^1_1}{\Obj c_1},
    (\Atom{\Obj A^1_1}{\Obj c_1} \land \Atom{\Obj A^1_0}{\Obj v_0}),
    \Atom{\Obj A^1_1}{\Obj c_1},\\
    \lforall[\Obj v_1][\Atom{\Obj A^1_0}{\Obj v_0}],
    \lexists[\Obj v_0][(\Atom{\Obj A^1_1}{\Obj c_1} \land \Atom{\Obj A^1_0}{\Obj v_0})]
}.
\end{multline*}
%
لکه چې له دويمې بېلګې ښکاري، په جوړښتي لړيو کښې ``بې‌ګټې توکي''
هم راتلے شي: داسې فارمولونه چې زياتي وي يا په جوړونه کښې ونډه نۀ
لري۔
\end{ex}
\olpseganchor{OLP-0156-B012}{end}

\olpseganchor{OLP-0156-B013}{start}
\begin{prop}\ollabel{prop:formed}
په~\LR{$\Frm[L]$} کښې هر فارمول~\LR{$!A$} يوه جوړښتي لړۍ لري۔
\end{prop}
\olpseganchor{OLP-0156-B013}{end}

\olpseganchor{OLP-0156-B014}{start}
\begin{proof}
فرض کړئ \LR{$!A$} اټومي دے۔ نو لړۍ~\LR{$\tuple{!A}$} د~\LR{$!A$} جوړښتي لړۍ ده۔
%
اوس فرض کړئ چې \LR{$!B$} او~\LR{$!C$} په ترتيب سره جوړښتي لړۍ
\LR{$\tuple{!B_0,\dotsc,!B_n}$} او \LR{$\tuple{!C_0,\dotsc,!C_m}$} لري۔
%
\begin{enumerate}
  \item که \LR{$!A \ident \lnot !B$} وي، نو
      \LR{$\tuple{!B_0,\dotsc,!B_n,\lnot !B_n}$}
      د~\LR{$!A$} جوړښتي لړۍ ده۔
  \item که \LR{$!A \ident (!B \land !C)$} وي، نو
      \LR{$\tuple{!B_0,\dotsc,!B_n,!C_0,\dotsc,!C_m,(!B_n \land !C_m)}$}
      د~\LR{$!A$} جوړښتي لړۍ ده۔
  \item که \LR{$!A \ident (!B \lor !C)$} وي، نو
      \LR{$\tuple{!B_0,\dotsc,!B_n,!C_0,\dotsc,!C_m,(!B_n \lor !C_m)}$}
      د~\LR{$!A$} جوړښتي لړۍ ده۔
  \item که \LR{$!A \ident (!B \lif !C)$} وي، نو
      \LR{$\tuple{!B_0,\dotsc,!B_n,!C_0,\dotsc,!C_m,(!B_n \lif !C_m)}$}
      د~\LR{$!A$} جوړښتي لړۍ ده۔
  \item که \LR{$!A \ident (!B \liff !C)$} وي، نو
      \LR{$\tuple{!B_0,\dotsc,!B_n,!C_0,\dotsc,!C_m,(!B_n \liff !C_m)}$}
      د~\LR{$!A$} جوړښتي لړۍ ده۔
  \item که \LR{$!A \ident \lforall[x][!B]$} وي، نو
      \LR{$\tuple{!B_0,\dotsc,!B_n,\lforall[x][!B_n]}$}
      د~\LR{$!A$} جوړښتي لړۍ ده۔
  \item که \LR{$!A \ident \lexists[x][!B]$} وي، نو
      \LR{$\tuple{!B_0,\dotsc,!B_n,\lexists[x][!B_n]}$}
      د~\LR{$!A$} جوړښتي لړۍ ده۔
\end{enumerate}
پر فارمولونه د استقرا د اصل له مخې، هر فارمول يوه جوړښتي لړۍ
لري۔
\end{proof}
\olpseganchor{OLP-0156-B014}{end}

\olpseganchor{OLP-0156-B015}{start}
د دې عکس هم ثابتولے شو۔ دا ځکه مهم دے چې ښيي د فارمولونو د تعريف
زموږ دواړه لارې هم ارزښته دي: دواړه يو شان پايلې ورکوي۔ دا هم ترې
معلومېږي چې د جوړښتي لړيو پر اوږدوالي عادي استقرا کارولے او د
فارمولونو په باب قضيې ثابتولے شو۔
\olpseganchor{OLP-0156-B015}{end}

\olpseganchor{OLP-0156-B016}{start}
\begin{lem}
\ollabel{lem:fseq-init}
فرض کړئ \LR{$\tuple{!A_0,\dotsc,!A_n}$} د~\LR{$!A_n$} جوړښتي لړۍ ده او \LR{$k
\leq n$} دے۔ نو \LR{$\tuple{!A_0,\dotsc,!A_k}$} د~\LR{$!A_k$} جوړښتي لړۍ ده۔
\end{lem}
\olpseganchor{OLP-0156-B016}{end}

\olpseganchor{OLP-0156-B017}{start}
\begin{proof}
تمرين۔
\end{proof}
\olpseganchor{OLP-0156-B017}{end}

\olpseganchor{OLP-0156-B018}{start}
\begin{prob}
\olref[fol][syn][fseq]{lem:fseq-init} ثابته کړئ۔
\end{prob}
\olpseganchor{OLP-0156-B018}{end}

\olpseganchor{OLP-0156-B019}{start}
\begin{thm}
\ollabel{thm:fseq-frm-equiv}
\LR{$\Frm[L]$} د~\LR{$\Lang L$} د ټولو هغو تارونو~\LR{$!A$} سټ دے چې د~\LR{$!A$} دپاره
د فارمول جوړښتي لړۍ موجوده وي۔
\end{thm}
\olpseganchor{OLP-0156-B019}{end}

\olpseganchor{OLP-0156-B020}{start}
\begin{proof}
\LR{$F$} دې د ژبې~\LR{$\Lang L$} د سمبولونو د ټولو هغو تارونو سټ وي چې جوړښتي
لړۍ لري۔ په
\olref[fol][syn][fseq]{prop:formed} کښې مو وليدل چې \LR{$\Frm[L] \subseteq
F$}؛ نو اوس ئې عکس ثابتوو۔
\olpseganchor{OLP-0156-B020}{end}

\olpseganchor{OLP-0156-B021}{start}
فرض کړئ \LR{$!A$} يوه جوړښتي لړۍ \LR{$\tuple{!A_0,\dotsc,!A_n}$} لري۔ پر~\LR{$n$}
د قوي استقرا له لارې ثابتوو چې \LR{$!A \in \Frm[L]$}۔ زموږ استقرايي فرض
دا دے چې د سمبولونو هر هغه تار چې د وروستي شاخص \LR{$m < n$} لرونکې
جوړښتي لړۍ لري، په~\LR{$\Frm[L]$} کښې دے۔ د جوړښتي لړۍ د تعريف له مخې، يا
\LR{$!A \ident !A_n$} وي او وروستے تار اټومي وي، يا بايد داسې \LR{$j,k < n$}
موجود وي چې له لاندې حالتونو څخه يو صدق وکړي:
\begin{enumerate}
    \item \LR{$!A \ident \lnot !A_j$}۔%
    \item \LR{$!A \ident (!A_j \land !A_k)$}۔%
    \item \LR{$!A \ident (!A_j \lor !A_k)$}۔%
    \item \LR{$!A \ident (!A_j \lif !A_k)$}۔%
    \item \LR{$!A \ident (!A_j \liff !A_k)$}۔%
    \item \LR{$!A \ident \lforall[x][!A_j]$}۔%
    \item \LR{$!A \ident \lexists[x][!A_j]$}۔%
\end{enumerate}
اوس حالت په حالت استدلال کوو۔ که \LR{$!A$} اټومي وي، نو
\LR{$!A_n \in \Frm[L]$}۔ بل حالت کښې فرض کړئ چې
\LR{$!A \ident (!A_j \land !A_k)$}۔ د
\olref[fol][syn][fseq]{lem:fseq-init} له مخې،
\LR{$\tuple{!A_0,\dotsc,!A_j}$} او \LR{$\tuple{!A_0,\dotsc,!A_k}$} په ترتيب
سره د \LR{$!A_j$} او~\LR{$!A_k$} جوړښتي لړۍ دي۔ څرنګه چې دا د~\LR{$!A$} د جوړښتي
لړۍ حقيقي پيلنۍ فرعي لړۍ دي، د دواړو وروستے شاخص له~\LR{$n$} څخه کم دے۔
نو د استقرايي فرض له مخې \LR{$!A_j$} او~\LR{$!A_k$} په~\LR{$\Frm[L]$} کښې دي؛ او
بيا د فارمول د تعريف له مخې \LR{$(!A_j \land !A_k)$} هم پکې دے۔ نور
حالتونه په ورته استدلال ثابتېږي۔
\textbf{د ژباړې د نسخې سمونونه (OLFOL-009--011):} سرچينه دلته د
نحوي يوشانوالي پر ځاے د معنايي هم ارزښت نښه، د روانې ژبې پر ځاے د
بياني ژبې شاخص، او د وروستي شاخص پر ځاے د لړۍ اوږدوالے يادوي۔ دلته
ثبوت د همدې برخې له تعريف او د قوي استقرا له متغير سره برابر شوے دے۔
\end{proof}
\olpseganchor{OLP-0156-B021}{end}

\olpseganchor{OLP-0156-B022}{start}
د ترمونو جوړښتي لړۍ د فارمولونه د جوړښتي لړيو په شان خاصيتونه لري۔
\olpseganchor{OLP-0156-B022}{end}

\olpseganchor{OLP-0156-B023}{start}
\begin{prop}
\ollabel{prop:fseq-trm-equiv}
\LR{$\Trm[L]$} د~\LR{$\Lang L$} د ټولو هغو تارونو \LR{$t$} سټ دے چې د~\LR{$t$} دپاره
د ترم جوړښتي لړۍ موجوده وي۔
\end{prop}
\olpseganchor{OLP-0156-B023}{end}

\olpseganchor{OLP-0156-B024}{start}
\begin{proof}
تمرين۔
\end{proof}
\olpseganchor{OLP-0156-B024}{end}

\olpseganchor{OLP-0156-B025}{start}
\begin{prob}
\olref[fol][syn][fseq]{prop:fseq-trm-equiv} ثابته کړئ۔
لارښوونه: د
\olref[fol][syn][fseq]{thm:fseq-frm-equiv} په ثبوت کښې کارول شوې
طريقې ته ورته تګلاره وکاروئ۔
\end{prob}
\olpseganchor{OLP-0156-B025}{end}

\olpseganchor{OLP-0156-B026}{start}
په جوړښتي لړيو کښې دوه ډوله ``بې‌ګټې توکي'' راتلے شي: تکراري غړي، او
هغه غړي چې د فارمول يا ترم په جوړولو کښې ونډه نۀ لري۔ دواړه د تر ټولو
لنډو جوړښتي لړيو په کتلو لرې کولے شو۔
\olpseganchor{OLP-0156-B026}{end}

\olpseganchor{OLP-0156-B027}{start}
\begin{defn}[تر ټولو لنډې جوړښتي لړۍ]
\ollabel{defn:minimal-fseq}
د فارمول~\LR{$!A$} جوړښتي لړۍ \LR{$\tuple{!A_0, \ldots, !A_n}$} د~\LR{$!A$}
\emph{تر ټولو لنډه جوړښتي لړۍ} ده، که د~\LR{$!A$} هره بله جوړښتي
لړۍ~\LR{$s$} داسې وي چې د~\LR{$s$} اوږدوالے له~\LR{$n+1$} سره برابر يا ترې زيات وي۔
\olpseganchor{OLP-0156-B027}{end}

\olpseganchor{OLP-0156-B028}{start}
په همدې ډول، د ترم~\LR{$t$} جوړښتي لړۍ \LR{$\tuple{t_0, \ldots, t_n}$} د~\LR{$t$}
\emph{تر ټولو لنډه جوړښتي لړۍ} ده، که د~\LR{$t$} هره بله جوړښتي
لړۍ~\LR{$s$} داسې وي چې د~\LR{$s$} اوږدوالے له~\LR{$n+1$} سره برابر يا ترې زيات وي۔
\end{defn}
\olpseganchor{OLP-0156-B028}{end}

\olpseganchor{OLP-0156-B029}{start}
پام وکړئ چې يو فارمول يا ترم له يوې څخه زياتې تر ټولو لنډې جوړښتي لړۍ
لرلے شي، خو په ټولو کښې به کټ مټ يو شان تارونه وي۔
\olpseganchor{OLP-0156-B029}{end}

\olpseganchor{OLP-0156-B030}{start}
\begin{prop}
\ollabel{prop:subformula-equivs}
لاندې شرطونه هم ارزښته دي:
\begin{enumerate}
    \item \LR{$!B$} د~\LR{$!A$} يو فرعي-فارمول دے۔
    \item \LR{$!B$} د~\LR{$!A$} په هره جوړښتي لړۍ کښې راځي۔
    \item \LR{$!B$} د~\LR{$!A$} په کومه تر ټولو لنډه جوړښتي لړۍ کښې راځي۔
\end{enumerate}
\end{prop}
\olpseganchor{OLP-0156-B030}{end}

\olpseganchor{OLP-0156-B031}{start}
\begin{proof}
تمرين۔
\end{proof}
\olpseganchor{OLP-0156-B031}{end}

\olpseganchor{OLP-0156-B032}{start}
\begin{prob}
\olref[fol][syn][fseq]{prop:subformula-equivs} ثابته کړئ۔
\end{prob}
\olpseganchor{OLP-0156-B032}{end}

\olpseganchor{OLP-0156-B033}{start}
\begin{history}
جوړښتي لړۍ ريمونډ سمولين په خپل درسي کتاب
\emph{First-Order Logic} \citep{Smullyan1968} کښې معرفي کړې۔
د جوړښتي لړيو نور خاصيتونه \citet{Zuckerman1973} ثابت کړل۔
\end{history}
\olpseganchor{OLP-0156-B033}{end}

\olpseganchor{OLP-0157-B004}{start}
\olfileid{fol}{syn}{fvs}
\olpseganchor{OLP-0157-B004}{end}

\olpseganchor{OLP-0157-B005}{start}
\olsection{ازاد متغيرونه او تړلي فارمولونه}
\olpseganchor{OLP-0157-B005}{end}

\olpseganchor{OLP-0157-B006}{start}
\begin{defn}[د متغير ازاد موردونه]
\ollabel{defn:free-occ}
په فارمول کښې د متغير \emph{ازاد} موردونه په استقرايي
ډول داسې تعريفېږي:
\begin{enumerate}
\item \indcase*{!A}{\LR{$!A$} اټومي دے}{په \LR{$\indfrm$} کښې د
  متغير ټول موردونه ازاد دي۔}
\olpseganchor{OLP-0157-B006}{end}

\olpseganchor{OLP-0157-B007}{start}
\item \indcase{!A}{\lnot !B}{د \LR{$\indfrm$} د متغير
  ازاد موردونه کټ مټ د \LR{$!B$} ازاد موردونه دي۔}
\olpseganchor{OLP-0157-B007}{end}

\olpseganchor{OLP-0157-B008}{start}
\item \indcase{!A}{(!B \ast !C)}{د \LR{$\indfrm$} د متغير ازاد
  موردونه په \LR{$!B$} کښې ازاد موردونه او ورسره په~\LR{$!C$} کښې ازاد موردونه
  دي۔}
\olpseganchor{OLP-0157-B008}{end}

\olpseganchor{OLP-0157-B009}{start}
\item \indcase{!A}{\lforall[x][!B]}{په \LR{$\indfrm$} کښې د
  متغير ازاد موردونه په~\LR{$!B$} کښې ټول ازاد موردونه دي، د~\LR{$x$}
  له موردونو پرته۔}
\olpseganchor{OLP-0157-B009}{end}

\olpseganchor{OLP-0157-B010}{start}
\item \indcase{!A}{\lexists[x][!B]}{په \LR{$\indfrm$} کښې د
  متغير ازاد موردونه په~\LR{$!B$} کښې ټول ازاد موردونه دي، د~\LR{$x$}
  له موردونو پرته۔}
\end{enumerate}
\end{defn}
\olpseganchor{OLP-0157-B010}{end}

\olpseganchor{OLP-0157-B011}{start}
\begin{defn}[تړلي متغيرونه]
په يوه فارمول~\LR{$!A$} کښې د متغير يو مورد هله \emph{تړلے} دے
چې ازاد نۀ وي۔
\end{defn}
\olpseganchor{OLP-0157-B011}{end}

\olpseganchor{OLP-0157-B012}{start}
\begin{prob}
په \olref[fol][syn][fvs]{defn:free-occ} پسې د تړلو متغيرونو د موردونو
يو استقرايي تعريف ورکړئ۔
\end{prob}
\olpseganchor{OLP-0157-B012}{end}

\olpseganchor{OLP-0157-B013}{start}
\begin{defn}[ساحه]
که \LR{$\lforall[x][!B]$} په يوه فارمول~\LR{$!A$} کښې د يوه
  فرعي فارمول مورد وي، نو په~\LR{$!A$} کښې د~\LR{$!B$} اړوند مورد د
  \LR{$\lforall[x]$} د اړوند مورد \emph{ساحه} بلل کېږي۔
  د \LR{$\lexists[x]$} دپاره هم همدا شان۔
\olpseganchor{OLP-0157-B013}{end}

\olpseganchor{OLP-0157-B014}{start}
که \LR{$!B$} په~\LR{$!A$} کښې د سور د يوه مورد
\LR{$\lforall[x]$} يا
    \LR{$\lexists[x]$} ساحه وي، نو په~\LR{$!B$} کښې د
\LR{$x$} ازاد موردونه په \LR{$\lforall[x][!B]$} او
\LR{$\lexists[x][!B]$} کښې تړلي دي۔ وايو چې دا
موردونه د ياد شوي سور د مورد \emph{له خوا تړل شوي} دي۔
\end{defn}
\olpseganchor{OLP-0157-B014}{end}

\olpseganchor{OLP-0157-B015}{start}
\begin{ex}
لاندې فارمول په پام کښې ونيسئ:
\[
\lexists[\Obj v_0][\underbrace{\Atom{\Obj A^2_0}{\Obj v_0,\Obj v_1}}_{!B}]
\]
\LR{$!B$} د \LR{$\lexists[\Obj v_0]$} ساحه ده۔ سور په \LR{$!B$} کښې د \LR{$\Obj v_0$}
مورد تړي، خو د \LR{$\Obj v_1$} مورد نۀ تړي۔ نو په دې حالت کښې \LR{$\Obj v_1$}
ازاد متغير دے۔
\olpseganchor{OLP-0157-B015}{end}


\olpseganchor{OLP-0157-B016}{start}
اوس وينو چې په يوه زيات پېچلي فارمول~\LR{$!A$} کښې دا څنګه کار کوي:
\[
\lforall[\Obj v_0][\underbrace{(\Atom{\Obj A^1_0}{\Obj v_0} \lif
    \Atom{\Obj A^2_0}{\Obj v_0, \Obj v_1})}_{!B}] \lif \lexists[\Obj
  v_1][\underbrace{(\Atom{\Obj A^2_1}{\Obj v_0, \Obj v_1} \lor \lforall[\Obj v_0][\overbrace{\lnot \Atom{\Obj A^1_1}{\Obj v_0}}^{!D}])}_{!C}]
\]
\LR{$!B$} د لومړي \LR{$\lforall[\Obj v_0]$} ساحه، \LR{$!C$} د
\LR{$\lexists[\Obj v_1]$} ساحه، او \LR{$!D$} د دويم \LR{$\lforall[\Obj v_0]$} ساحه
ده۔ لومړے \LR{$\lforall[\Obj v_0]$} په~\LR{$!B$} کښې د \LR{$\Obj v_0$} موردونه
تړي، \LR{$\lexists[\Obj v_1]$} په \LR{$!C$} کښې د \LR{$\Obj v_1$} مورد تړي، او دويم
\LR{$\lforall[\Obj v_0]$} په~\LR{$!D$} کښې د \LR{$\Obj v_0$} مورد تړي۔ د \LR{$\Obj v_1$}
لومړے مورد او د \LR{$\Obj v_0$} څلورم مورد په~\LR{$!A$} کښې ازاد دي۔ د
\LR{$\Obj v_0$} وروستے مورد په \LR{$!D$} کښې ازاد، خو په \LR{$!C$} او~\LR{$!A$} کښې
تړلے دے۔
\end{ex}
\olpseganchor{OLP-0157-B016}{end}

\olpseganchor{OLP-0157-B017}{start}
\begin{defn}[جمله]
فارمول~\LR{$!A$} هله او يوازې هله يوه 
\emph{تړلے فارمول} ده چې د متغيرونه هېڅ ازاد مورد ونۀ لري۔
\end{defn}
\olpseganchor{OLP-0157-B017}{end}

\olpseganchor{OLP-0157-B018}{start}
% add examples!
\olpseganchor{OLP-0157-B018}{end}

\olpseganchor{OLP-0158-B004}{start}
\olfileid{fol}{syn}{sub}
\olpseganchor{OLP-0158-B004}{end}

\olpseganchor{OLP-0158-B005}{start}
\olsection{ځايناستي}
\olpseganchor{OLP-0158-B005}{end}

\olpseganchor{OLP-0158-B006}{start}
\begin{defn}[په ترم کښې ځايناستي]
\LR{$\Subst{s}{t}{x}$}، يعنې په \LR{$s$} کښې د~\LR{$x$} د هر مورد پر ځاے د \LR{$t$}
د \emph{ځايناستي} پايله، په بازګشتي ډول داسې تعريفوو:
\begin{enumerate}
\item \indcase{s}{c}{\LR{$\Subst{\indfrm}{t}{x}$} پخپله \LR{$s$} دے۔}
\olpseganchor{OLP-0158-B006}{end}

\olpseganchor{OLP-0158-B007}{start}
\item \indcase{s}{y}{که \LR{$y$} يو متغير او \LR{$y \not\ident x$} وي، نو
  \LR{$\Subst{\indfrm}{t}{x}$} هم پخپله~\LR{$s$} دے۔}
\olpseganchor{OLP-0158-B007}{end}

\olpseganchor{OLP-0158-B008}{start}
\item \indcase{s}{x}{\LR{$\Subst{\indfrm}{t}{x}$} همدا~\LR{$t$} دے۔}
\olpseganchor{OLP-0158-B008}{end}

\olpseganchor{OLP-0158-B009}{start}
\item\indcase{s}{\Atom{f}{t_1, \dots, t_n}}{\LR{$\Subst{\indfrmp}{t}{x}$}
  همدا \LR{$\Atom{f}{\Subst{t_1}{t}{x}, \dots, \Subst{t_n}{t}{x}}$} دے۔}
\end{enumerate}
\end{defn}
\olpseganchor{OLP-0158-B009}{end}

\olpseganchor{OLP-0158-B010}{start}
\begin{defn}
يو ترم~\LR{$t$} په \LR{$!A$} کښې د \LR{$x$} دپاره هله \emph{ازاد} دے چې
په \LR{$!A$} کښې د~\LR{$x$} هېڅ ازاد مورد د داسې سور په ساحه کښې نۀ راځي چې
په~\LR{$t$} کښې کوم متغير تړي۔
\end{defn}
\olpseganchor{OLP-0158-B010}{end}

\olpseganchor{OLP-0158-B011}{start}
\begin{ex} ~
\begin{enumerate}
\item \LR{$\Obj v_8$} په \LR{$\lexists[\Obj
  v_3]\Atom{\Obj A^2_4}{\Obj v_3,\Obj v_1}$} کښې د \LR{$\Obj v_1$} دپاره
  ازاد دے۔
\olpseganchor{OLP-0158-B011}{end}

\olpseganchor{OLP-0158-B012}{start}
\item \LR{$\Obj f^2_1(\Obj v_1, \Obj v_2)$} په
  \LR{$\lforall[\Obj v_2]\Atom{\Obj A^2_4}{\Obj v_0,\Obj v_2}$} کښې د
  \LR{$\Obj v_0$} دپاره \emph{ازاد نۀ} دے۔
\end{enumerate}
\end{ex}
\olpseganchor{OLP-0158-B012}{end}

\olpseganchor{OLP-0158-B013}{start}
\begin{defn}[په فارمول کښې ځايناستي]
که \LR{$!A$} فارمول، \LR{$x$}~متغير، او \LR{$t$}~داسې ترم وي چې
په~\LR{$!A$} کښې د~\LR{$x$} دپاره ازاد وي، نو
\LR{$\Subst{!A}{t}{x}$} په~\LR{$!A$} کښې د~\LR{$x$} د ټولو ازادو موردونو پر ځاے
د~\LR{$t$} د ځايناستي پايله ده۔
\begin{enumerate}
\item \indcase{!A}{\lfalse}{\LR{$\Subst{\indfrm}{t}{x}$}
    همدا \LR{$\lfalse$} دے۔}
\olpseganchor{OLP-0158-B013}{end}

\olpseganchor{OLP-0158-B014}{start}
\item \indcase{!A}{\ltrue}{\LR{$\Subst{\indfrm}{t}{x}$}
    همدا \LR{$\ltrue$} دے۔}
\olpseganchor{OLP-0158-B014}{end}

\olpseganchor{OLP-0158-B015}{start}
\item \indcase{!A}{\Atom{P}{t_1,\dots,
    t_n}}{\LR{$\Subst{\indfrm}{t}{x}$} همدا
    \LR{$\Atom{P}{\Subst{t_1}{t}{x},
    \dots, \Subst{t_n}{t}{x}}$} دے۔}
\olpseganchor{OLP-0158-B015}{end}

\olpseganchor{OLP-0158-B016}{start}
\item \indcase{!A}{\eq[t_1][t_2]}{\LR{$\Subst{\indfrmp}{t}{x}$} همدا
  \LR{$\Subst{t_1}{t}{x} = \Subst{t_2}{t}{x}$} دے۔}
\olpseganchor{OLP-0158-B016}{end}

\olpseganchor{OLP-0158-B017}{start}
\item \indcase{!A}{\lnot !B}{\LR{$\Subst{\indfrmp}{t}{x}$}
    همدا \LR{$\lnot \Subst{!B}{t}{x}$} دے۔}
\olpseganchor{OLP-0158-B017}{end}

\olpseganchor{OLP-0158-B018}{start}
\item \indcase{!A}{(!B \land
    !C)}{\LR{$\Subst{\indfrmp}{t}{x}$} همدا \LR{$(\Subst{!B}{t}{x} \land
    \Subst{!C}{t}{x})$} دے۔}
\olpseganchor{OLP-0158-B018}{end}

\olpseganchor{OLP-0158-B019}{start}
\item \indcase{!A}{(!B \lor
    !C)}{\LR{$\Subst{\indfrmp}{t}{x}$} همدا \LR{$(\Subst{!B}{t}{x} \lor
    \Subst{!C}{t}{x})$} دے۔}
\olpseganchor{OLP-0158-B019}{end}

\olpseganchor{OLP-0158-B020}{start}
\item \indcase{!A}{(!B \lif
    !C)}{\LR{$\Subst{\indfrmp}{t}{x}$} همدا \LR{$(\Subst{!B}{t}{x} \lif
    \Subst{!C}{t}{x})$} دے۔}
\olpseganchor{OLP-0158-B020}{end}

\olpseganchor{OLP-0158-B021}{start}
\item \indcase{!A}{(!B \liff
    !C)}{\LR{$\Subst{\indfrmp}{t}{x}$} همدا \LR{$(\Subst{!B}{t}{x} \liff
    \Subst{!C}{t}{x})$} دے۔}
\olpseganchor{OLP-0158-B021}{end}

\olpseganchor{OLP-0158-B022}{start}
\item 
\indcase{!A}{\lforall[y][!B]}{که \LR{$y$} له \LR{$x$} څخه بېل متغير وي، نو
    \LR{$\Subst{\indfrmp}{t}{x}$} همدا
    \LR{$\lforall[y][\Subst{!B}{t}{x}]$} دے؛ که نۀ وي، نو
    \LR{$\Subst{\indfrmp}{t}{x}$} پخپله \LR{$\indfrm$} دے۔}
\olpseganchor{OLP-0158-B022}{end}

\olpseganchor{OLP-0158-B023}{start}
\item 
\indcase{!A}{\lexists[y][!B]}{که \LR{$y$} له \LR{$x$} څخه بېل متغير وي، نو
    \LR{$\Subst{\indfrmp}{t}{x}$} همدا
    \LR{$\lexists[y][\Subst{!B}{t}{x}]$} دے؛ که نۀ وي، نو
    \LR{$\Subst{\indfrmp}{t}{x}$} پخپله \LR{$\indfrm$} دے۔}
\end{enumerate}
\end{defn}
\olpseganchor{OLP-0158-B023}{end}

\olpseganchor{OLP-0158-B024}{start}
\begin{explain}
پام وکړئ چې ځايناستي کېدے شي بې‌اغېزې وي: که \LR{$x$} په \LR{$!A$} کښې هېڅ
نۀ وي راغلے، نو \LR{$\Subst{!A}{t}{x}$} پخپله~\LR{$!A$} دے۔
\olpseganchor{OLP-0158-B024}{end}

\olpseganchor{OLP-0158-B025}{start}
دا محدوديت چې \LR{$t$} بايد په~\LR{$!A$} کښې د~\LR{$x$} دپاره ازاد وي، د
لاندې حالتونو د مخنيوي دپاره اړين دے۔ که \LR{$!A \ident
\lexists[y][x < y]$} او \LR{$t \ident y$} وي، نو \LR{$\Subst{!A}{t}{x}$} به
\LR{$\lexists[y][y < y]$} شي۔ په دې حالت کښې ازاد متغير \LR{$y$} د ځايناستي
پر مهال د سور \LR{$\lexists[y]$} له خوا «راګېرېږي»، او دا ناغوښتې پايله
ده۔ مثلاً غواړو چې هر کله \LR{$\lforall[x][!B]$} صدق کوي،
\LR{$\Subst{!B}{t}{x}$} هم صدق وکړي۔ خو
\LR{$\lforall[x][\lexists[y][x < y]]$} په پام کښې ونيسئ (دلته \LR{$!B$} همدا
\LR{$\lexists[y][x < y]$} دے)۔ دا د بېلګې په توګه د طبيعي عددونو په باب
رښتونې جمله ده: د هر عدد~\LR{$x$} دپاره تر هغه لوے يو عدد~\LR{$y$} شته۔ که
مو~\LR{$y$} د~\LR{$x$} د ممکنې ځايناستۍ په توګه روا ګڼلے واے، پايله به مو
\LR{$\Subst{!B}{y}{x} \ident \lexists[y][y < y]$} وه، چې دروغ ده۔ د دې
مخنيوے په دې کوو چې په~\LR{$t$} کښې هېڅ ازاد متغير بايد په~\LR{$!A$} کښې د
کوم سور له خوا تړلے نۀ شي۔
\end{explain}
\olpseganchor{OLP-0158-B025}{end}

\olpseganchor{OLP-0158-B026}{start}
د پېچلې نښې د مخنيوي دپاره ډېر ځله لاندې دود کاروو: که \LR{$!A$} داسې
فارمول وي چې متغير~\LR{$x$} پکښې ازاد راتلے شي، د دې د
ښودلو دپاره~\LR{$!A(x)$} هم ليکو۔ چې موخه مو کوم \LR{$!A$} او~\LR{$x$} دي څرګنده
وي، او \LR{$t$} يو ترم وي (چې فرض کېږي په \LR{$!A(x)$} کښې د \LR{$x$} دپاره ازاد
دے)، نو \LR{$!A(t)$} د \LR{$\Subst{!A}{t}{x}$} لنډون ليکو۔ مثلاً ويلے شو:
«\LR{$!A(t)$} د~\LR{$\lforall[x][!A(x)]$} يو مورد بولو۔» له دې څخه موخه دا
ده چې که \LR{$!A$} هر فارمول، \LR{$x$}~متغير، او \LR{$t$}~داسې ترم وي
چې په~\LR{$!A$} کښې د~\LR{$x$} دپاره ازاد وي، نو \LR{$\Subst{!A}{t}{x}$} د
\LR{$\lforall[x][!A]$} يو مورد دے۔
\olpseganchor{OLP-0158-B026}{end}

\olpseganchor{OLP-0159-B004}{start}
\olchapter{fol}{sem}{د لومړۍ درجې منطق معنٰی پوهنه}
\olpseganchor{OLP-0159-B004}{end}

\olpseganchor{OLP-0160-B004}{start}
\olfileid{fol}{syn}{its}
\olpseganchor{OLP-0160-B004}{end}

\olpseganchor{OLP-0160-B005}{start}
\olsection{پېژندنه}
\olpseganchor{OLP-0160-B005}{end}

\olpseganchor{OLP-0160-B006}{start}
عبارتونو ته معنا ورکول د معنٰی پوهنې کار دے۔ د معنٰی پوهنې مرکزي
مفهوم په جوړښت کښې د صدق مفهوم دے۔ جوړښت د ژبې
بنسټيزو اجزاوو ته معنا ورکوي: د تعريف ساحه د څيزونو يو ناتش سټ دے۔
سورونه داسې تفسيرېږي چې پر همدې ساحه ګرځي، ثابت سمبولونه د ساحې
له غړو سره ګومارل کېږي، تابعې له د تعريف ساحه څخه
همدې ساحې ته له تابعو سره ګومارل کېږي، او پريديکاتونه پر
د تعريف ساحه له اړيکو سره ګومارل کېږي۔ د تعريف ساحه او د بنسټيز لغت
ګومارنې په ګډه جوړښت جوړوي۔ متغيرونه په فارمولونه
کښې راتلے شي، او د معنٰی پوهنې د ورکولو دپاره بايد هغو ته هم د
د تعريف ساحه غړي وګومارو---دا د متغير ګومارنه ده۔ په پای کښې
د صدق اړيکه دا ټول سره يوځاے کوي۔ فارمول په
جوړښت~\LR{$\Struct{M}$} کښې د متغير ګومارنې~\LR{$s$} په نسبت
صدق کولے شي، چې \LR{$\Sat{M}{!A}[s]$} ليکل کېږي۔ دا اړيکه هم د~\LR{$!A$} پر
جوړښت په استقرا تعريفېږي: د منطقي نښلوونکو د رښتياوالي جدولونه
کاروو، څو مثلاً د \LR{$(!A \land !B)$} صدق د \LR{$!A$} او~\LR{$!B$} د صدق يا نه
صدق له مخې تعريف کړو۔ بيا څرګندېږي چې که فارمول~\LR{$!A$}
تړلے فارمول وي، يعنې ازاد متغيرونه ونۀ لري، د متغير ګومارنه
اړونده نۀ ده؛ نو ويلے شو چې تړلي فارمولونه په ساده ډول په
جوړښتونه کښې صدق کوي يا نۀ کوي۔
\olpseganchor{OLP-0160-B006}{end}

\olpseganchor{OLP-0160-B007}{start}
د تړلي فارمولونه دپاره د صدق د اړيکې \LR{$\Sat{M}{!A}$} پر بنسټ بيا د
اعتبار، معنايي استلزام، او د صدق وړتيا بنسټيز معنايي مفهومونه
تعريفولے شو۔ تړلے فارمول هله معتبره ده، \LR{$\Entails !A$}، چې هر
جوړښت پکښې صدق وکړي۔ د تړلي فارمولونه يو سټ هله هغه لازموي،
\LR{$\Gamma \Entails !A$}، چې هر هغه جوړښت چې د~\LR{$\Gamma$} په ټولو
تړلي فارمولونه کښې صدق کوي، په~\LR{$!A$} کښې هم صدق وکړي۔ او د
تړلي فارمولونه يو سټ هله د صدق وړ دے چې کوم جوړښت په هغه کښې
په يو وخت په ټولو تړلي فارمولونه کښې صدق وکړي۔ څرنګه چې فارمولونه
په استقرايي ډول تعريف شوي، او صدق بيا د فارمولونه پر جوړښت په
استقرا تعريفېږي، استقرا کارولے شو څو د خپلې معنٰی پوهنې خاصيتونه ثابت
او تعريف شوي معنايي مفهومونه سره وتړو۔
\olpseganchor{OLP-0160-B007}{end}

\olpseganchor{OLP-0161-B004}{start}
\olfileid{fol}{syn}{str}
\olpseganchor{OLP-0161-B004}{end}

\olpseganchor{OLP-0161-B005}{start}
\olsection{د لومړۍ درجې ژبو جوړښتونه}
\olpseganchor{OLP-0161-B005}{end}

\olpseganchor{OLP-0161-B006}{start}
\begin{explain}
د لومړۍ درجې ژبې پخپله \emph{بې‌تفسيره} دي: له ثابت سمبولونه،
تابعې او پريديکاتونه سره کومه ځانګړې معنا نۀ وي تړلې۔
معنا د  \emph{جوړښت} په ټاکلو ورکول کېږي۔
هغه \emph{د تعريف ساحه} ټاکي؛ يعنې هغه څيزونه چې ثابت سمبولونه ئې
ټاکي، تابعې پرې عمل کوي، او سورونه پرې ګرځي۔ سربېره پر دې،
دا هم ټاکي چې کوم ثابت سمبولونه کوم څيزونه ټاکي، تابع
څيزونه څنګه نورو څيزونو ته لېږي، او پريديکاتونه پر کومو څيزونو
صدق کوي۔ جوړښتونه په منطق کښې د \emph{معنايي} مفهومونو بنسټ
دي، لکه لازميت، اعتبار، او د صدق وړتيا۔ په ليکنو کښې ورته په بېلابېل
ډول «جوړښتونه»، «تفسيرونه» يا «مدلونه» ويل کېږي۔
\end{explain}
\olpseganchor{OLP-0161-B006}{end}

\olpseganchor{OLP-0161-B007}{start}
\begin{defn}[جوړښتونه]
د لومړۍ درجې منطق د يوې ژبې \LR{$\Lang{L}$} دپاره 
\emph{جوړښت}~\LR{$\Struct M$} له لاندې اجزاوو جوړ دے:
\begin{enumerate}
\item \emph{د تعريف ساحه:} يو ناتش سټ، \LR{$\Domain M$}
\item \emph{د ثابت سمبولونه تفسير:} د \LR{$\Lang{L}$} د هر
  ثابت سمبول~\LR{$c$} دپاره يو غړے \LR{$\Assign{c}{M} \in \Domain M$}
\item \emph{د پريديکاتونه تفسير:} د \LR{$\Lang{L}$} د هر \LR{$n$}-ځايه
  پريديکات~\LR{$R$} دپاره (له \LR{$\eq$} پرته) يوه \LR{$n$}-ځايه اړيکه
  \LR{$\Assign{R}{M} \subseteq \Domain{M}^n$}
\item \emph{د تابعې تفسير:} د \LR{$\Lang{L}$} د هر \LR{$n$}-ځايه
  تابع~\LR{$f$} دپاره يوه \LR{$n$}-ځايه تابع \LR{$\Assign{f}{M}
  \colon \Domain{M}^n \to \Domain{M}$}
\end{enumerate}
\end{defn}
\olpseganchor{OLP-0161-B007}{end}

\olpseganchor{OLP-0161-B008}{start}
\begin{ex}
د حساب د ژبې دپاره جوړښت~\LR{$\Struct M$} له يوه سټ؛ د
\LR{$\Domain M$} له يوه غړي \LR{$\Assign{\Obj 0}{M}$}، چې د
ثابت سمبول~\LR{$\Obj 0$} تفسير دے؛ له يوې يو ځايه تابع
\LR{$\Assign{\Obj \prime}{M} \colon \Domain{M} \to \Domain M$}؛ له دوو
دوه ځايه تابعو \LR{$\Assign{\Obj +}{M}$} او \LR{$\Assign{\Obj
  \times}{M}$}، چې دواړه \LR{$\Domain M^2 \to \Domain M$} دي؛ او له يوې
دوه ځايه اړيکې \LR{$\Assign{\Obj <}{M} \subseteq \Domain{M}^2$} جوړ دے۔
\olpseganchor{OLP-0161-B008}{end}

\olpseganchor{OLP-0161-B009}{start}
د داسې جوړښت يوه څرګنده بېلګه دا ده:
\begin{enumerate}
\item \LR{$\Domain N = \Nat$}
\item \LR{$\Assign{\Obj 0}{N} = 0$}
\item د هر \LR{$n \in \Nat$} دپاره \LR{$\Assign{\Obj \prime}{N}(n) = n + 1$}
\item د هر \LR{$n, m \in \Nat$} دپاره \LR{$\Assign{\Obj +}{N}(n, m) = n + m$}
\item د هر \LR{$n, m \in \Nat$} دپاره \LR{$\Assign{\Obj \times}{N}(n, m) = n\cdot m$}
\item \LR{$\Assign{\Obj <}{N} = \Setabs{\tuple{n, m}}{n \in \Nat, m \in
  \Nat, n < m}$}
\end{enumerate}
د \LR{$\Lang L_A$} دپاره په دې ډول تعريف شوے جوړښت~\LR{$\Struct N$} د
\emph{حساب معياري مدل} بلل کېږي، ځکه د~\LR{$\Lang L_A$} غير منطقي ثابتونه
کټ مټ هغسې تفسيروي چې تمه ئې کوئ۔
\olpseganchor{OLP-0161-B009}{end}

\olpseganchor{OLP-0161-B010}{start}
خو د~\LR{$\Lang L_A$} دپاره ډېر نور ممکن جوړښتونه هم شته۔ مثلاً د
ساحې په توګه د~\LR{$\Nat$} پر ځاے د صحيح عددونو سټ~\LR{$\Int$} اخيستلے، او د
\LR{$\Obj 0$}، \LR{$\Obj \prime$}، \LR{$\Obj +$}، \LR{$\Obj \times$} او \LR{$\Obj <$} تفسيرونه
ورسره سم تعريفولے شو۔ خو د~\LR{$\Lang L_A$} داسې جوړښتونه هم تعريفولے شو
چې له عددونو سره هېڅ لرې اړيکه هم نۀ لري۔
\end{ex}
\olpseganchor{OLP-0161-B010}{end}

\olpseganchor{OLP-0161-B011}{start}
\begin{ex}
د سټ تيورۍ د ژبې~\LR{$\Lang L_Z$} يو جوړښت~\LR{$\Struct M$} يوازې يوه ساحه
او يوه دوه ځايه اړيکه غواړي۔ نو په فني لحاظ مثلاً د خلکو سټ او اړيکه
«\LR{$x$} له \LR{$y$} څخه مشر دے» هم د \LR{$\Lang L_Z$} دپاره جوړښت
کېدے شي، او هم~\LR{$\Nat$} له \LR{$n \ge m$} سره د \LR{$n, m \in \Nat$} دپاره۔
\textbf{د ژباړې د نسخې سمون (OLFOL-012):} سرچينه د اړيکې د ځاييز
والي په عبارت کښې ناسمه نښلوونکې کرښه ږدي؛ دلته مطلب «يوه دوه ځايه
اړيکه» لوستل شوے دے۔
\olpseganchor{OLP-0161-B011}{end}

\olpseganchor{OLP-0161-B012}{start}
د \LR{$\Lang L_Z$} يو په ځانګړي ډول په زړه پورې جوړښت، چې د ساحې
غړي ئې په رښتيا سټونه او د \LR{$\Obj \in$} تفسير ئې په رښتيا
اړيکه «\LR{$x$} د~\LR{$y$} غړے دے» وي، د \emph{موروثي متناهي
سټونو} جوړښت~\LR{$\Struct{{HF}}$} دے:
\begin{enumerate}
\item \LR{$\Domain{{{HF}}} = \emptyset \cup \Pow{\emptyset} \cup
  \Pow{\Pow{\emptyset}} \cup \Pow{\Pow{\Pow{\emptyset}}} \cup \dots$}؛
\item \LR{$\Assign{\Obj \in}{{{HF}}} = \Setabs{\tuple{x, y}}{x, y \in
  \Domain{{{HF}}}, x \in y}$}۔
\end{enumerate}
\end{ex}
\olpseganchor{OLP-0161-B012}{end}

\olpseganchor{OLP-0161-B013}{start}
\begin{digress}
دا چې د څه شي د جوړښت ګڼلو دپاره کوم شرطونه ږدو، زموږ پر
منطق اغېز کوي۔ مثلاً د تشو ساحو نۀ منل، د سورو لرونکو جملو د صدق د
معمول حساب له مخې، دا تضمينوي چې
\LR{$\lexists[x][(!A(x) \lor \lnot !A(x))]$} معتبر دے---يعنې يو منطقي صدق
دے۔ او دا شرط چې ټول ثابت سمبولونه بايد د ساحې کوم څيز ته اشاره وکړي،
تضمينوي چې وجودي تعميم د استنتاج يوه سمه بڼه ده:
\LR{$!A(a)$}، نو \LR{$\lexists[x][!A(x)]$}۔ که مو نومونو ته د ساحې بهر يا هېڅ
څيز ته اشاره کول روا ګڼلي واے، د \emph{ازاد منطق} پر لور به تللي وو؛
په هغه کښې وجودي تعميم يوه اضافي مقدمې ته اړتيا لري:
\LR{$!A(a)$} او \LR{$\lexists[x][\eq[x][a]]$}، نو \LR{$\lexists[x][!A(x)]$}۔
\end{digress}
\olpseganchor{OLP-0161-B013}{end}

\olpseganchor{OLP-0162-B004}{start}
\olfileid{fol}{syn}{cov}
\olpseganchor{OLP-0162-B004}{end}

\olpseganchor{OLP-0162-B005}{start}
\olsection{پوښل شوي جوړښتونه د لومړۍ درجې ژبو دپاره}
\olpseganchor{OLP-0162-B005}{end}

\olpseganchor{OLP-0162-B006}{start}
\begin{explain}
په ياد ولرئ چې يو ترم هله \emph{تړلے} دے چې هېڅ متغيرونه ونۀ لري۔
\end{explain}
\olpseganchor{OLP-0162-B006}{end}

\olpseganchor{OLP-0162-B007}{start}
\begin{defn}[د تړلو ترمونو قيمت]
که \LR{$t$} د ژبې~\LR{$\Lang L$} يو تړلے ترم او \LR{$\Struct M$} د~\LR{$\Lang L$} يو
جوړښت وي، \emph{قيمت}~\LR{$\Value{t}{M}$} ئې په لاندې ډول
تعريفېږي:
\begin{enumerate}
\item که \LR{$t$} يوازې ثابت سمبول~\LR{$c$} وي، نو
  \LR{$\Value{c}{M} = \Assign{c}{M}$}۔
\item که \LR{$t$} د \LR{$\Atom{f}{t_1, \ldots, t_n}$} بڼه ولري، نو
  \[
  \Value{t}{M} = \Assign{f}{M}(\Value{t_1}{M}, \ldots,
  \Value{t_n}{M}).
  \]
\end{enumerate}
\end{defn}
\olpseganchor{OLP-0162-B007}{end}

\olpseganchor{OLP-0162-B008}{start}
\begin{defn}[پوښل شوے جوړښت]
يو جوړښت هله \emph{پوښل شوے} دے چې د ساحې هر غړے د کوم تړلي
ترم قيمت وي۔
\end{defn}
\olpseganchor{OLP-0162-B008}{end}

\olpseganchor{OLP-0162-B009}{start}
\begin{ex}
فرض کړئ~\LR{$\Lang L$} هغه ژبه ده چې ثابت سمبولونه \LR{$\Obj{zero}$}،
\LR{$\Obj{one}$}، \LR{$\Obj{two}$}، \dots، دوه ځايه پريديکات~\LR{$<$}، او
دوه ځايه تابعې \LR{$+$} او \LR{$\times$} لري۔ د~\LR{$\Lang L$}
جوړښت~\LR{$\Struct M$} دې هغه وي چې ساحه ئې
\LR{$\Domain M = \{0, 1, 2, \ldots \}$} او ګومارنې ئې
\LR{$\Assign{\Obj{zero}}{M} = 0$}، \LR{$\Assign{\Obj{one}}{M} = 1$}،
\LR{$\Assign{\Obj{two}}{M} = 2$}، او همداسې نورې دي۔ د دوه ځايه اړيکې د
سمبول \LR{$<$} دپاره، سټ \LR{$\Assign{<}{M}$} د ټولو هغو جوړو
\LR{$\tuple{c_1, c_2} \in \Domain{M}^2$} سټ دے چې \LR{$c_1$} له~\LR{$c_2$} څخه
کوچنے وي: مثلاً \LR{$\tuple{1, 3} \in \Assign{<}{M}$}، خو
\LR{$\tuple{2, 2} \notin \Assign{<}{M}$}۔ د دوه ځايه تابع \LR{$+$}
دپاره \LR{$\Assign{+}{M}$} په معمول ډول تعريف کړئ---مثلاً
\LR{$\Assign{+}{M}(2,3)$} عدد~\LR{$5$} ته لېږي؛ د دوه ځايه
تابع~\LR{$\times$} دپاره هم همداسې وکړئ۔ نو د \LR{$\Obj{four}$}
قيمت پخپله~\LR{$4$} دے، او د \LR{$\times(\Obj{two},
+(\Obj{three},\Obj{zero}))$} قيمت (يا په منځ‌ليکي نښه کښې،
\LR{$\Obj{two} \times (\Obj{three} + \Obj{zero})$}) دا دے:
\begin{multline*}
\Value{\times(\Obj{two}, +(\Obj{three},\Obj{zero}))}{M} \\
\begin{aligned}
& =\Assign{\times}{M}(\Value{\Obj{two}}{M}, \Value{+(\Obj{three}, \Obj{zero})}{M})\\
& = \Assign{\times}{M}(\Value{\Obj{two}}{M}, \Assign{+}{M}(\Value{\Obj{three}}{M},
\Value{\Obj{zero}}{M})) \\
& = \Assign{\times}{M}(\Assign{\Obj{two}}{M}, \Assign{+}{M}(\Assign{\Obj{three}}{M},
\Assign{\Obj{zero}}{M})) \\
& = \Assign{\times}{M}(2, \Assign{+}{M}(3, 0)) \\
& = \Assign{\times}{M}(2, 3) \\
& = 6
\end{aligned}
\end{multline*}
\textbf{د ژباړې د نسخې سمون (OLFOL-013):} د سرچينې د محاسبې په
لومړۍ کرښه کښې د برابرۍ نښه د بلې کرښې له پيل سره تکرار شوې؛ دلته
لومړۍ تکراري نښه لرې شوې ده۔
\end{ex}
\olpseganchor{OLP-0162-B009}{end}

\olpseganchor{OLP-0162-B010}{start}
\begin{prob}
ايا \LR{$\Struct N$}، يعنې د حساب معياري مدل، پوښل شوے دے؟ څرګنده ئې کړئ۔
\end{prob}
\olpseganchor{OLP-0162-B010}{end}

\olpseganchor{OLP-0163-B004}{start}
\olfileid{fol}{syn}{sat}
\olpseganchor{OLP-0163-B004}{end}

\olpseganchor{OLP-0163-B005}{start}
\olsection{په  جوړښت کښې د
   فارمول صدق}
\olpseganchor{OLP-0163-B005}{end}

\olpseganchor{OLP-0163-B006}{start}
\begin{explain}
هغه بنسټيز مفهومونه چې يوې خوا ته عبارتونه لکه ترمونه او
فارمولونه، او بلې خوا ته جوړښتونه سره تړي، د يوه ترم
\emph{قيمت} او د فارمول \emph{صدق} دي۔ په غير رسمي ډول،
د يوه ترم قيمت د جوړښت يو غړے دے---که ترم
يوازې يو ثابت وي، قيمت ئې هغه څيز دے چې جوړښت له هغه
ثابت سره ګومارلے وي؛ او که د تابعې په کارولو جوړ شوے وي، د
هغه قيمت د ثابتونو له قيمتونه او هغو تابعو څخه حسابېږي چې
جوړښت په ترم کښې تابعو ته ګومارلې وي۔ فارمول په يوه
جوړښت کښې هله \emph{صدق کوي} چې اړوندو ته ورکړل شوے تفسير
فارمول د جوړښت په ساحه کښې رښتونے کړي۔ د صدق دا مفهوم په
استقرا ټاکل کېږي: د جوړښت مشخصات نېغ په نېغه وايي چې اټومي
فارمولونه کله صدق کوي، او دا چې يو پېچلے فارمول کله صدق کوي،
د هغه له اصلي نښلوونکي يا سور او د هغه د سمدستي فرعي فارمولونه له
صدق يا نه صدق سره سم تعريفوو۔
\olpseganchor{OLP-0163-B006}{end}

\olpseganchor{OLP-0163-B007}{start}
دلته د سورو حالت لږ نازک دے، ځکه د سور لرونکي فارمول سمدستي
فرعي فارمول يو ازاد متغير لري، او جوړښتونه د
متغيرونه قيمتونه نۀ ټاکي۔ د دې ستونزې د حل دپاره د
\emph{متغير ګومارنې} هم معرفي کوو، او صدق يوازې د
جوړښت په نسبت نۀ، بلکې د جوړښت او متغير
ګومارنې په نسبت تعريفوو۔
\end{explain}
\olpseganchor{OLP-0163-B007}{end}

\olpseganchor{OLP-0163-B008}{start}
\begin{defn}[د متغير ګومارنه]
د جوړښت~\LR{$\Struct{M}$} دپاره \emph{د متغير ګومارنه}~\LR{$s$} هغه
تابع ده چې هر متغير د~\LR{$\Domain M$} کوم غړے ته لېږي؛
يعنې \LR{$s\colon \Var \to \Domain M$}۔
\end{defn}
\olpseganchor{OLP-0163-B008}{end}

\olpseganchor{OLP-0163-B009}{start}
\begin{explain}
جوړښت هر ثابت سمبول ته قيمت، او د متغير ګومارنه
هر متغير ته يو قيمت ورکوي۔ خو غواړو چې له دغو څخه جوړ شوي ترمونه هم
د د تعريف ساحه غړي ونوموي۔ د دې دپاره د ترمونو قيمت په
استقرا تعريفوو۔ د ثابت سمبولونه او متغيرونو قيمت هماغه دے چې
جوړښت يا د متغير ګومارنه ئې ټاکي؛ د زيات پېچلو ترمونو قيمت
د هغو تابعو په بازګشتي کارولو حسابېږي چې جوړښت له
تابعې سره ګومارلې وي۔
\end{explain}
\olpseganchor{OLP-0163-B009}{end}

\olpseganchor{OLP-0163-B010}{start}
\begin{defn}[د ترمونو قيمت]
که \LR{$t$} د ژبې~\LR{$\Lang L$} يو ترم، \LR{$\Struct M$} د~\LR{$\Lang L$} يو
جوړښت، او \LR{$s$} د~\LR{$\Struct M$} دپاره متغير ګومارنه وي،
\emph{قيمت}~\LR{$\Value{t}{M}[s]$} ئې په لاندې ډول تعريفېږي:
\begin{enumerate}
\item \indcase{t}{c}{\LR{$\Value{\indfrm}{M}[s] = \Assign{\indcomplex}{M}$}۔}
\item \indcase{t}{x}{\LR{$\Value{\indfrm}{M}[s] = s(\indcomplex)$}۔}
\item \indcase{t}{\Atom{f}{t_1, \ldots, t_n}}{
\[
\Value{\indfrm}{M}[s] = \Assign{f}{M}(\Value{t_1}{M}[s], \ldots,
\Value{t_n}{M}[s]).
\]}
\end{enumerate}
\end{defn}
\olpseganchor{OLP-0163-B010}{end}

\olpseganchor{OLP-0163-B011}{start}
\begin{defn}[د~\LR{$x$} بديله ګومارنه]
که \LR{$s$} د جوړښت~\LR{$\Struct M$} دپاره متغير ګومارنه وي،
د~\LR{$\Struct M$} هره هغه متغير ګومارنه~\LR{$s'$} چې له~\LR{$s$} سره تر
ډېره يوازې د~\LR{$x$} په ګومارل شوي څيز کښې توپير لري، د~\LR{$s$}
\emph{د~\LR{$x$} بديله ګومارنه} بلل کېږي۔ که \LR{$s'$} د~\LR{$s$} د~\LR{$x$} بديله
ګومارنه وي، \LR{$\varAssign{s'}{s}{x}$} ليکو۔
\end{defn}
\olpseganchor{OLP-0163-B011}{end}

\olpseganchor{OLP-0163-B012}{start}
\begin{explain}
پام وکړئ چې د يوې ګومارنې~\LR{$s$} د~\LR{$x$} بديله ګومارنه اړ نۀ ده چې
په رښتيا~\LR{$x$} ته بېل څيز وګوماري۔ په حقيقت کښې هره ګومارنه د خپل ځان
د~\LR{$x$} بديله ګومارنه هم ګڼل کېږي۔
\end{explain}
\olpseganchor{OLP-0163-B012}{end}

\olpseganchor{OLP-0163-B013}{start}
\begin{defn}
  که \LR{$s$} د جوړښت~\LR{$\Struct M$} دپاره متغير ګومارنه
  او \LR{$m \in \Domain{M}$} وي، نو ګومارنه~\LR{$\Subst{s}{m}{x}$} هغه د
  متغير ګومارنه ده چې داسې تعريفېږي:
  \[\Subst{s}{m}{x}(y) = \begin{cases}
    m & \text{\RL{که }} y \ident x\\
    s(y) & \text{\RL{که نۀ وي}}.
  \end{cases}\]
\end{defn}
\olpseganchor{OLP-0163-B013}{end}

\olpseganchor{OLP-0163-B014}{start}
په بله وينا، \LR{$\Subst{s}{m}{x}$} د~\LR{$s$} د~\LR{$x$} هغه ځانګړې بديله ګومارنه
ده چې د ساحې غړے~\LR{$m$} له~\LR{$x$} سره ګوماري، او له~\LR{$x$} پرته
متغيرونه ته هماغه څيزونه ګوماري چې \LR{$s$} ئې ګوماري۔
\olpseganchor{OLP-0163-B014}{end}

\olpseganchor{OLP-0163-B015}{start}
\begin{defn}[صدق]
\ollabel{defn:satisfaction}
په جوړښت~\LR{$\Struct M$} کښې د فارمول~\LR{$!A$} صدق د
متغير ګومارنې~\LR{$s$} په نسبت، چې په سمبولونو کښې
\LR{$\Sat{M}{!A}[s]$} دے، په بازګشتي ډول داسې تعريفېږي۔ (موږ
\LR{$\Sat/{M}{!A}[s]$} د «\LR{$\Sat{M}{!A}[s]$} نۀ دے» په معنا ليکو۔)
\begin{enumerate}
\item %
  \indcase{!A}{\lfalse}{\LR{$\Sat/{M}{\indfrm}[s]$}۔}
\olpseganchor{OLP-0163-B015}{end}

\olpseganchor{OLP-0163-B016}{start}
\item %
  \indcase{!A}{\ltrue}{\LR{$\Sat{M}{\indfrm}[s]$}۔}
\olpseganchor{OLP-0163-B016}{end}

\olpseganchor{OLP-0163-B017}{start}
\item \indcase{!A}{\Atom{R}{t_1, \dots, t_n}}{\LR{$\Sat{M}{\indfrm}[s]$}
  هله او يوازې هله چې \LR{$\langle \Value{t_1}{M}[s], \dots,
  \Value{t_n}{M}[s] \rangle \in \Assign{R}{M}$}۔}
\olpseganchor{OLP-0163-B017}{end}

\olpseganchor{OLP-0163-B018}{start}
\item \indcase{!A}{\eq[t_1][t_2]}{\LR{$\Sat{M}{\indfrm}[s]$} هله او
  يوازې هله چې \LR{$\Value{t_1}{M}[s] = \Value{t_2}{M}[s]$}۔}
\olpseganchor{OLP-0163-B018}{end}

\olpseganchor{OLP-0163-B019}{start}
\item %
  \indcase{!A}{\lnot !B}{\LR{$\Sat{M}{\indfrm}[s]$} هله او يوازې هله چې
    \LR{$\Sat/{M}{!B}[s]$}۔}
\olpseganchor{OLP-0163-B019}{end}

\olpseganchor{OLP-0163-B020}{start}
\item %
  \indcase{!A}{(!B \land !C)}{\LR{$\Sat{M}{\indfrm}[s]$} هله او يوازې
    هله چې \LR{$\Sat{M}{!B}[s]$} او \LR{$\Sat{M}{!C}[s]$}۔}
\olpseganchor{OLP-0163-B020}{end}

\olpseganchor{OLP-0163-B021}{start}
\item %
  \indcase{!A}{(!B \lor !C)}{\LR{$\Sat{M}{\indfrm}[s]$} هله او يوازې
    هله چې \LR{$\Sat{M}{!B}[s]$} يا \LR{$\Sat{M}{!C}[s]$} (يا دواړه)۔}
\olpseganchor{OLP-0163-B021}{end}

\olpseganchor{OLP-0163-B022}{start}
\item %
  \indcase{!A}{(!B \lif !C)}{\LR{$\Sat{M}{\indfrm}[s]$} هله او يوازې
    هله چې \LR{$\Sat/{M}{!B}[s]$} يا \LR{$\Sat{M}{!C}[s]$} (يا دواړه)۔}
\olpseganchor{OLP-0163-B022}{end}

\olpseganchor{OLP-0163-B023}{start}
\item %
  \indcase{!A}{(!B \liff !C)}{\LR{$\Sat{M}{\indfrm}[s]$} هله او يوازې
    هله چې يا دواړه \LR{$\Sat{M}{!B}[s]$} او \LR{$\Sat{M}{!C}[s]$} وي، يا
    نۀ \LR{$\Sat{M}{!B}[s]$} وي او نۀ \LR{$\Sat{M}{!C}[s]$}۔}
\olpseganchor{OLP-0163-B023}{end}

\olpseganchor{OLP-0163-B024}{start}
\item %
  \indcase{!A}{\lforall[x][!B]}{\LR{$\Sat{M}{\indfrm}[s]$} هله او يوازې
    هله چې د هر غړے~\LR{$m \in \Domain M$} دپاره
    \LR{$\Sat{M}{!B}[\Subst{s}{m}{x}]$}۔}
\olpseganchor{OLP-0163-B024}{end}

\olpseganchor{OLP-0163-B025}{start}
\item %
  \indcase{!A}{\lexists[x][!B]}{\LR{$\Sat{M}{\indfrm}[s]$} هله او يوازې
  هله چې لږ تر لږه د يوه غړے~\LR{$m \in \Domain M$} دپاره
  \LR{$\Sat{M}{!B}[\Subst{s}{m}{x}]$}۔}
\end{enumerate}
\end{defn}
\olpseganchor{OLP-0163-B025}{end}

\olpseganchor{OLP-0163-B026}{start}
\begin{explain}
د متغير ګومارنې په وروستي
دوو شرطونو کښې مهمې
دي۔ د \LR{$\lforall[x][!B(x)]$} صدق داسې نۀ شو تعريفولے:
«د ټولو \LR{$m \in \Domain{M}$} دپاره \LR{$\Sat{M}{!B(m)}$}۔»
د \LR{$\lexists[x][!B(x)]$} صدق داسې نۀ شو تعريفولے: «لږ تر لږه د يوه
\LR{$m \in \Domain{M}$} دپاره \LR{$\Sat{M}{!B(m)}$}۔» لامل دا دے چې که
\LR{$m \in \Domain M$} وي، هغه د ژبې سمبول نۀ دے؛ نو \LR{$!B(m)$}
فارمول نۀ دے (يعنې \LR{$\Subst{!B}{m}{x}$} ناتعريف دے)۔ دا هم نۀ
شو فرضولے چې داسې ثابت سمبولونه يا ترمونه لرو چې د~\LR{$\Struct{M}$} هر
غړے ونوموي، ځکه د جوړښتونه په تعريف کښې هېڅ شرط دا
نۀ غواړي۔ په معياري ژبه کښې د ثابت سمبولونه سټ شمېرېدونکے لامتناهي دے؛
نو که \LR{$\Domain{M}$} د شمېر وړ نۀ وي، د هر څيز د نومولو دپاره
کافي ثابت سمبولونه هم نشته۔
\olpseganchor{OLP-0163-B026}{end}

\olpseganchor{OLP-0163-B027}{start}
دا ستونزه د متغير ګومارنو په معرفي کولو حلوو؛ په هغو سره
متغيرونه نېغ په نېغه د ساحې له غړي سره تړلے شو۔ بيا د دې
پر ځاے چې مثلاً ووايو
\LR{$\lexists[x][!B(x)]$} په
\LR{$\Struct M$} کښې هله او يوازې هله صدق کوي چې
لږ تر لږه د يوه \LR{$m \in \Domain{M}$} دپاره
\LR{$\Sat{M}{!B(m)}$} وي، وايو چې دا په~\LR{$\Struct M$} کښې د~\LR{$s$}
\emph{په نسبت} هله او يوازې هله صدق کوي چې \LR{$!B(x)$} د
\LR{$\Subst{s}{m}{x}$} په نسبت لږ تر لږه د يوه
\LR{$m \in \Domain M$} دپاره صدق وکړي۔
\end{explain}
\olpseganchor{OLP-0163-B027}{end}

\olpseganchor{OLP-0163-B028}{start}
\begin{ex}
فرض کړئ \LR{$\Lang{L} = \{a, b, f, R\}$}، چې \LR{$a$} او \LR{$b$}
ثابت سمبولونه، \LR{$f$}~دوه ځايه تابع، او \LR{$R$}~دوه ځايه
پريديکات دي۔ لاندې تعريف شوے جوړښت~\LR{$\Struct{M}$} په پام
کښې ونيسئ:
\begin{enumerate}
\item \LR{$\Domain M = \{1, 2, 3, 4\}$}
\item \LR{$\Assign{a}{M} = 1$}
\item \LR{$\Assign{b}{M} = 2$}
\item که \LR{$x+y \le 3$} وي، \LR{$\Assign{f}{M}(x, y) = x+y$}، او که نۀ وي
  \LR{$= 3$}۔
\item \LR{$\Assign{R}{M} = \{\tuple{1, 1}, \tuple{1, 2}, \tuple{2, 3}, \tuple{2, 4}\}$}
\end{enumerate}
تابع \LR{$s(x) = 1$}، چې هر متغير ته \LR{$1 \in \Domain{M}$} ګوماري،
د~\LR{$\Struct{M}$} دپاره د متغير ګومارنه ده۔
\olpseganchor{OLP-0163-B028}{end}

\olpseganchor{OLP-0163-B029}{start}
بيا
\begin{align*}
\Value{f(a,b)}{M}[s] & = \Assign{f}{M}(\Value{a}{M}[s], \Value{b}{M}[s]).
\intertext{\RL{څرنګه چې \LR{$a$} او \LR{$b$} ثابت سمبولونه دي،
  \LR{$\Value{a}{M}[s] = \Assign{a}{M} = 1$} او
  \LR{$\Value{b}{M}[s] = \Assign{b}{M} = 2$}۔ نو}}
\Value{f(a,b)}{M}[s] & = \Assign{f}{M}(1, 2) = 1+2 = 3.
\intertext{\RL{د \LR{$f(f(a,b),a)$} د قيمت د حسابولو دپاره بايد دا وګورو:}}
\Value{f(f(a,b),a)}{M}[s] & = \Assign{f}{M}(\Value{f(a, b)}{M}[s],
\Value{a}{M}[s]) = \Assign{f}{M}(3, 1) = 3,
\intertext{\RL{ځکه \LR{$3+1 > 3$}۔ څرنګه چې \LR{$s(x) = 1$} او
  \LR{$\Value{x}{M}[s] = s(x)$}، دا هم لرو:}}
\Value{f(f(a,b),x)}{M}[s] & = \Assign{f}{M}(\Value{f(a, b)}{M}[s],
\Value{x}{M}[s]) = \Assign{f}{M}(3, 1) = 3,
\end{align*}
\olpseganchor{OLP-0163-B029}{end}

\olpseganchor{OLP-0163-B030}{start}
يو اټومي فارمول~\LR{$R(t_1, t_2)$} هله صدق کوي چې د مدخلو د قيمتونو
جوړه، يعنې \LR{$\tuple{\Value{t_1}{M}[s],
  \Value{t_2}{M}[s]}$}، د~\LR{$\Assign{R}{M}$} غړے وي۔ نو مثلاً
\LR{$\Sat{M}{R(b,f(a,b))}[s]$} لرو، ځکه
\LR{$\tuple{\Value{b}{M}, \Value{f(a,b)}{M}} = \tuple{2, 3} \in
\Assign{R}{M}$}؛ خو \LR{$\Sat/{M}{R(x, f(a,b))}[s]$}، ځکه
\LR{$\tuple{1, 3} \notin \Assign{R}{M}$}۔
\olpseganchor{OLP-0163-B030}{end}

\olpseganchor{OLP-0163-B031}{start}
د دې د معلومولو دپاره چې يو غير اټومي فارمول~\LR{$!A$} صدق کوي که نۀ، د
استقرايي تعريف هغه شرط کاروئ چې پر اصلي نښلوونکي عملي کېږي۔ مثلاً د
\LR{$R(a, a) \lif (R(b, x) \lor R(x, b))$} اصلي نښلوونکے~\LR{$\lif$} دے، او
\begin{align*}
  & \Sat{M}{R(a, a) \lif (R(b, x) \lor R(x, b))}[s] \text{\RL{ هله او يوازې هله چې }}\\
  & \qquad
  \Sat/{M}{R(a,a)}[s] \text{\RL{ يا }} \Sat{M}{R(b, x) \lor R(x, b)}[s]
  \intertext{\RL{څرنګه چې \LR{$\Sat{M}{R(a,a)}[s]$} (ځکه
    \LR{$\tuple{1,1} \in \Assign{R}{M}$})، ځواب لا نۀ شو ټاکلے او لومړے
    بايد معلومه کړو چې \LR{$\Sat{M}{R(b, x) \lor R(x, b)}[s]$} دے که نۀ:}}
  & \Sat{M}{R(b, x) \lor R(x, b)}[s] \text{\RL{ هله او يوازې هله چې }}\\
  & \qquad \Sat{M}{R(b,x)}[s] \text{\RL{ يا }} \Sat{M}{R(x, b)}[s]
  \intertext{\RL{او دا حالت سم دے، ځکه \LR{$\Sat{M}{R(x, b)}[s]$}
    (ځکه \LR{$\tuple{1,2} \in \Assign{R}{M}$})۔}}
\end{align*}
\olpseganchor{OLP-0163-B031}{end}

\olpseganchor{OLP-0163-B032}{start}
په ياد ولرئ چې د~\LR{$s$} د~\LR{$x$} بديله ګومارنه هغه د متغير ګومارنه ده چې
له~\LR{$s$} سره تر ډېره يوازې د~\LR{$x$} په ګومارنه کښې توپير لري۔ د
\LR{$\Domain{M}$} د هر غړے دپاره د~\LR{$s$} يوه د~\LR{$x$} بديله ګومارنه
شته:
\begin{align*}
  s_1 & = \Subst{s}{1}{x}, &
  s_2 & = \Subst{s}{2}{x},\\
  s_3 & = \Subst{s}{3}{x}, &
  s_4 & = \Subst{s}{4}{x}.
\end{align*}
نو مثلاً \LR{$s_2(x) = 2$}، او له~\LR{$x$} پرته د هر متغير~\LR{$y$} دپاره
\LR{$s_2(y) = s(y) = 1$}۔ د جوړښت~\LR{$\Struct{M}$} دپاره دا د~\LR{$s$} د~\LR{$x$}
ټولې بديلې ګومارنې دي، ځکه \LR{$\Domain{M} = \{1, 2, 3, 4\}$}۔ په ځانګړي
ډول پام وکړئ چې \LR{$s_1 = s$} (هره~\LR{$s$} تل د خپل ځان د~\LR{$x$} بديله ګومارنه
هم وي)۔
\olpseganchor{OLP-0163-B032}{end}

\olpseganchor{OLP-0163-B033}{start}
د دې د معلومولو دپاره چې د وجودي سور لرونکے
  فارمول~\LR{$\lexists[x][!A(x)]$} صدق کوي که نۀ، بايد معلومه کړو چې
  لږ تر لږه د يوه \LR{$m \in \Domain M$} دپاره
  \LR{$\Sat{M}{!A(x)}[\Subst{s}{m}{x}]$} دے که نۀ۔ نو
  \[
  \Sat{M}{\lexists[x][(R(b,x) \lor R(x,b))]}[s],
  \]
  ځکه \LR{$\Sat{M}{R(b,x) \lor R(x, b)}[\Subst{s}{1}{x}]$}
  (\LR{$\Subst{s}{3}{x}$} هم شرط پوره کوي)۔ خو
  \[
  \Sat/{M}{\lexists[x][(R(b,x) \land R(x,b))]}[s]
  \]
  ځکه هر \LR{$m \in \Domain{M}$} چې واخلو،
  \LR{$\Sat/{M}{R(b,x) \land R(x,b)}[\Subst{s}{m}{x}]$}۔
\olpseganchor{OLP-0163-B033}{end}

\olpseganchor{OLP-0163-B034}{start}
د دې د معلومولو دپاره چې د عمومي سور لرونکے
  فارمول~\LR{$\lforall[x][!A(x)]$} صدق کوي که نۀ، بايد معلومه کړو چې
  د ټولو \LR{$m \in \Domain M$} دپاره
  \LR{$\Sat{M}{!A(x)}[\Subst{s}{m}{x}]$} دے که نۀ۔ نو
  \[
  \Sat{M}{\lforall[x][(R(x,a) \lif R(a,x))]}[s],
  \]
  ځکه د ټولو \LR{$m \in \Domain M$} دپاره
  \LR{$\Sat{M}{R(x,a) \lif R(a,x)}[\Subst{s}{m}{x}]$}۔ د \LR{$m = 1$} دپاره
  \LR{$\Sat{M}{R(a,x)}[\Subst{s}{1}{x}]$} لرو، نو تالي رښتونے دے؛ او د
  \LR{$m = 2$}، \LR{$3$} او~\LR{$4$} دپاره
  \LR{$\Sat/{M}{R(x,a)}[\Subst{s}{m}{x}]$}، نو مقدم دروغ دے۔ خو
  \[
  \Sat/{M}{\lforall[x][(R(a,x) \lif R(x,a))]}[s]
  \]
  ځکه \LR{$\Sat/{M}{R(a,x) \lif R(x,a)}[\Subst{s}{2}{x}]$}
  (ځکه \LR{$\Sat{M}{R(a, x)}[\Subst{s}{2}{x}]$} او
  \LR{$\Sat/{M}{R(x, a)}[\Subst{s}{2}{x}]$})۔
\olpseganchor{OLP-0163-B034}{end}





\olpseganchor{OLP-0163-B037}{start}
د يو زيات پېچلي حالت دپاره دا فارمول په پام کښې ونيسئ:
\[
\lforall[x][(R(a,x) \lif \lexists[y][R(x,y)])].
\]
څرنګه چې \LR{$\Sat/{M}{R(a,x)}[\Subst{s}{3}{x}]$} او
\LR{$\Sat/{M}{R(a,x)}[\Subst{s}{4}{x}]$}، هغه په زړه پورې حالتونه چې د
شرطي د تالي په باب فکر پکښې پکار دے يوازې \LR{$m = 1$} او \LR{$m = 2$} دي۔
ايا \LR{$\Sat{M}{\lexists[y][R(x,y)]}[\Subst{s}{1}{x}]$}؟ دا هله صدق کوي
چې لږ تر لږه يو \LR{$n \in \Domain M$} داسې وي چې
\LR{$\Sat{M}{R(x,y)}[\Subst{\Subst{s}{1}{x}}{n}{y}]$}۔ په حقيقت کښې که
\LR{$n = 1$} واخلو، نو
\LR{$\Subst{\Subst{s}{1}{x}}{n}{y} = \Subst{s}{1}{y} = s$}۔ څرنګه چې
\LR{$s(x) = 1$}، \LR{$s(y) = 1$}، او \LR{$\tuple{1,1} \in \Assign{R}{M}$}، ځواب هو
دے۔
\olpseganchor{OLP-0163-B037}{end}

\olpseganchor{OLP-0163-B038}{start}
د دې د معلومولو دپاره چې
\LR{$\Sat{M}{\lexists[y][R(x,y)]}[\Subst{s}{2}{x}]$}، بايد د
متغير ګومارنې \LR{$\Subst{\Subst{s}{2}{x}}{n}{y}$} وګورو۔ دلته
د \LR{$n = 1$} دپاره دا ګومارنه~\LR{$s_2 = \Subst{s}{2}{x}$} ده؛ او
\LR{$R(x,y)$} پکښې صدق نۀ کوي (\LR{$s_2(x) = 2$}، \LR{$s_2(y) = 1$}، او
\LR{$\tuple{2,1}\notin \Assign{R}{M}$})۔ خو
\LR{$\Subst{\Subst{s}{2}{x}}{3}{y} = \Subst{s_2}{3}{y}$} په پام کښې
ونيسئ۔ \LR{$\Sat{M}{R(x,y)}[\Subst{s_2}{3}{y}]$}، ځکه
\LR{$\tuple{2,3} \in \Assign{R}{M}$}؛ نو
\LR{$\Sat{M}{\lexists[y][R(x,y)]}[s_2]$}۔
\olpseganchor{OLP-0163-B038}{end}

\olpseganchor{OLP-0163-B039}{start}
نو د ټولو \LR{$m \in \Domain M$} دپاره يا
\LR{$\Sat/{M}{R(a,x)}[\Subst{s}{m}{x}]$} (که \LR{$m = 3$} يا \LR{$4$} وي)، يا
\LR{$\Sat{M}{\lexists[y][R(x,y)]}[\Subst{s}{m}{x}]$} (که \LR{$m = 1$} يا \LR{$2$}
وي)؛ نو
\[
\Sat{M}{\lforall[x][(R(a,x) \lif \lexists[y][R(x,y)])]}[s].
\]
بل خوا،
\[
\Sat/{M}{\lexists[x][(R(a,x) \land \lforall[y][R(x,y)])]}[s].
\]
\LR{$\Sat{M}{R(a,x)}[\Subst{s}{m}{x}]$} يوازې د \LR{$m = 1$} او \LR{$m = 2$}
دپاره لرو۔ خو د~\LR{$m$} د دواړو دغو قيمتونو دپاره داسې
\LR{$n \in \Domain M$} شته، يعنې \LR{$n = 4$}، چې
\LR{$\Sat/{M}{R(x,y)}[\Subst{\Subst{s}{m}{x}}{n}{y}]$}؛ نو د \LR{$m = 1$} او
\LR{$m = 2$} دپاره
\LR{$\Sat/{M}{\lforall[y][R(x,y)]}[\Subst{s}{m}{x}]$}۔ په ټوله کښې داسې
\LR{$m \in \Domain M$} نشته چې
\LR{$\Sat{M}{R(a,x) \land \lforall[y][R(x,y)]}[\Subst{s}{m}{x}]$}۔
\textbf{د ژباړې د نسخې سمونونه (OLFOL-015--020):} سرچينه په بېلو
ځايونو کښې د وجودي شرط د صدق څرګندونه پرېږدي، د اړوند له تفسير سره
د متغير ګومارنه نښلوي، په فارمول کښې زياته کامه او زيات قوسونه ږدي،
له يوه عددي مساوات څخه متغير غورځوي، او د وروستي عمومي حالت متغير
بدلوي۔ دلته هر مورد د همدې تعريف او ټاکلي جوړښت له مخې سم شوے دے۔
\end{ex}
\olpseganchor{OLP-0163-B039}{end}

\olpseganchor{OLP-0163-B046}{start}
\begin{prob}
فرض کړئ \LR{$\Lang L = \{c, f, A\}$}، چې يو ثابت سمبول، يوه يو ځايه
تابع او يو دوه ځايه پريديکات لري، او
جوړښت~\LR{$\Struct{M}$} داسې ورکړل شوے دے:
\begin{enumerate}
\item \LR{$\Domain M = \{1, 2, 3\}$}
\item \LR{$\Assign{c}{M} = 3$}
\item \LR{$\Assign{f}{M}(1) = 2, \Assign{f}{M}(2) = 3, \Assign{f}{M}(3) = 2$}
\item \LR{$\Assign{A}{M} = \{\tuple{1, 2}, \tuple{2, 3}, \tuple{3, 3}\}$}
\end{enumerate}
(الف) د ټولو متغيرونه~\LR{$v$} دپاره \LR{$s(v) = 1$} وټاکئ۔ معلومه کړئ
چې آيا
\[
\Sat{M}{\lexists[x][(A(f(z), c) \lif \lforall[y][(A(y, x) \lor A(f(y),
      x))])]}[s]
\]
صدق کوي که نۀ، او لامل ئې څرګند کړئ۔
\olpseganchor{OLP-0163-B046}{end}

\olpseganchor{OLP-0163-B047}{start}
(ب) يو بېل جوړښت او متغير ګومارنه ورکړئ چې فارمول پکښې
صدق ونۀ کړي۔
\end{prob}
\olpseganchor{OLP-0163-B047}{end}

\olpseganchor{OLP-0164-B004}{start}
\olfileid{fol}{syn}{ass}
\olpseganchor{OLP-0164-B004}{end}

\olpseganchor{OLP-0164-B005}{start}
\olsection{د متغيرونو ګومارنې}
\olpseganchor{OLP-0164-B005}{end}

\olpseganchor{OLP-0164-B006}{start}
\begin{explain}
د متغير ګومارنه~\LR{$s$} \emph{هر} متغير ته يو ارزښت ورکوي---او
متغيرونه بې‌شمېره دي۔ البته دا اړينه نۀ ده۔ موږ له متغير
ګومارنو څخه يوازې د اسانتيا دپاره غواړو چې ټولو متغيرونه ته
ارزښتونه وګوماري۔ د يوه ترم~\LR{$t$} ارزښت، او دا چې فارمول~\LR{$!A$} د
ګومارنې~\LR{$s$} په نسبت په جوړښت کښې صدق کوي که نۀ، يوازې په هغو
ارزښتونو پورې تړلي دي چې~\LR{$s$} ئې د~\LR{$t$} متغيرونه او د~\LR{$!A$} ازادو
متغيرونه ته ګوماري۔ د لاندې دوو قضيو منځپانګه همدا ده۔ د «پورې
تړلي دي» مانا دقيقه کولو دپاره ښيو چې هرې دوې هغه متغير ګومارنې چې
د~\LR{$t$} پر ټولو متغيرونو سره موافقې وي، هغه ته يو شان ارزښت ورکوي؛ او
که دوې متغير ګومارنې د~\LR{$!A$} پر ټولو ازادو متغيرونو سره موافقې وي، نو
\LR{$!A$} د يوې په نسبت هله صدق کوي چې د بلې په نسبت هم صدق وکړي۔
\end{explain}
\olpseganchor{OLP-0164-B006}{end}

\olpseganchor{OLP-0164-B007}{start}
\begin{prop}\ollabel{prop:valindep}
که په ترم~\LR{$t$} کښې متغيرونه د \LR{$x_1$}، \dots،~\LR{$x_n$} په منځ کښې وي،
او \LR{$s_1(x_i) = s_2(x_i)$} د \LR{$i = 1$}، \dots،~\LR{$n$} دپاره وي، نو
\LR{$\Value{t}{M}[s_1] = \Value{t}{M}[s_2]$}۔
\end{prop}
\olpseganchor{OLP-0164-B007}{end}

\olpseganchor{OLP-0164-B008}{start}
\begin{proof}
د~\LR{$t$} پر پېچلتيا د استقرا له لارې۔ په بنسټيز حالت کښې~\LR{$t$} يا
ثابت سمبول وي يا له \LR{$x_1$}، \dots،~\LR{$x_n$} متغيرونو څخه يو وي۔ که
\LR{$t = c$} وي، نو
\LR{$\Value{t}{M}[s_1] = \Assign{c}{M} = \Value{t}{M}[s_2]$}۔ که
\LR{$t = x_i$} وي، د قضيې د فرض له مخې \LR{$s_1(x_i) = s_2(x_i)$}؛ نو
\LR{$\Value{t}{M}[s_1] = s_1(x_i) = s_2(x_i) = \Value{t}{M}[s_2]$}۔
\olpseganchor{OLP-0164-B008}{end}

\olpseganchor{OLP-0164-B009}{start}
د استقرايي ګام دپاره فرض کړئ چې
\LR{$t = \Atom{f}{t_1, \dots, t_k}$}، او ادعا د
\LR{$t_1$}، \dots، \LR{$t_k$} دپاره صدق کوي۔ نو
\begin{align*}
  \Value{t}{M}[s_1] & = \Value{\Atom{f}{t_1, \dots, t_k}}{M}[s_1] \\
  & = \Assign{f}{M}(\Value{t_1}{M}[s_1], \dots, \Value{t_k}{M}[s_1]).
\intertext{\RL{د \LR{$j = 1$}، \dots،~\LR{$k$} دپاره، د~\LR{$t_j$} متغيرونه د
    \LR{$x_1$}، \dots،~\LR{$x_n$} په منځ کښې دي۔ د استقرايي فرض له مخې
    \LR{$\Value{t_j}{M}[s_1] = \Value{t_j}{M}[s_2]$}۔ نو}}
\Value{t}{M}[s_1] & = \Value{\Atom{f}{t_1, \dots, t_k}}{M}[s_1] \\
  & = \Assign{f}{M}(\Value{t_1}{M}[s_1], \dots, \Value{t_k}{M}[s_1]) \\
  & = \Assign{f}{M}(\Value{t_1}{M}[s_2], \dots, \Value{t_k}{M}[s_2]) \\
  & = \Value{\Atom{f}{t_1, \dots, t_k}}{M}[s_2] = \Value{t}{M}[s_2].
\end{align*}
\end{proof}
\olpseganchor{OLP-0164-B009}{end}

\olpseganchor{OLP-0164-B010}{start}
\begin{prop}\ollabel{prop:satindep}
که په~\LR{$!A$} کښې ازاد متغيرونه د \LR{$x_1$}، \dots،~\LR{$x_n$} په منځ کښې
وي، او \LR{$s_1(x_i) = s_2(x_i)$} د \LR{$i = 1$}، \dots،~\LR{$n$} دپاره وي، نو
\LR{$\Sat{M}{!A}[s_1]$} هله او يوازې هله چې \LR{$\Sat{M}{!A}[s_2]$}۔
\end{prop}
\olpseganchor{OLP-0164-B010}{end}

\olpseganchor{OLP-0164-B011}{start}
\begin{proof}
د~\LR{$!A$} پر پېچلتيا استقرا کاروو۔ په بنسټيز حالت کښې، چې~\LR{$!A$} اټومي
وي، \LR{$!A$} له لاندې څخه يو کېدے شي:
\LR{$\ltrue$}،
\LR{$\lfalse$}،
د يوه \LR{$k$}-ځايه اړوند~\LR{$R$} او ترمونو \LR{$t_1$}، \dots،~\LR{$t_k$} دپاره
\LR{$\Atom{R}{t_1, \dots, t_k}$}، يا د ترمونو~\LR{$t_1$} او~\LR{$t_2$} دپاره
\LR{$\eq[t_1][t_2]$}۔ په وروستيو دوو حالتونو کښې يوازې د
دوه اړخيز شرطي مخامخ لور ثابتوو، ځکه د شاتني لور ثبوت ئې هم‌شکله
دے۔
\olpseganchor{OLP-0164-B011}{end}

\olpseganchor{OLP-0164-B012}{start}
\begin{enumerate}
\item %
  \indcase{!A}{\ltrue}{دواړه \LR{$\Sat{M}{\indfrm}[s_1]$} او
    \LR{$\Sat{M}{\indfrm}[s_2]$} صدق کوي۔}
\olpseganchor{OLP-0164-B012}{end}

\olpseganchor{OLP-0164-B013}{start}
\item %
  \indcase{!A}{\lfalse}{دواړه \LR{$\Sat/{M}{!A}[s_1]$} او
    \LR{$\Sat/{M}{!A}[s_2]$}۔}
\olpseganchor{OLP-0164-B013}{end}

\olpseganchor{OLP-0164-B014}{start}
\item
  \indcase{!A}{\Atom{R}{t_1, \ldots, t_k}}{فرض کړئ
    \LR{$\Sat{M}{\indfrm}[s_1]$}۔ نو
    \[
    \langle \Value{t_1}{M}[s_1], \ldots, \Value{t_k}{M}[s_1] \rangle
    \in \Assign{R}{M}.
    \]
    د \LR{$i = 1$}، \dots،~\LR{$k$} دپاره، د \olref{prop:valindep} له مخې
    \LR{$\Value{t_i}{M}[s_1] = \Value{t_i}{M}[s_2]$}۔ نو دا هم لرو چې
    \LR{$\langle \Value{t_1}{M}[s_2], \ldots, \Value{t_k}{M}[s_2] \rangle
    \in \Assign{R}{M}$}؛ ځکه \LR{$\Sat{M}{\indfrm}[s_2]$}۔}
\item
  \indcase{!A}{\eq[t_1][t_2]}{فرض کړئ \LR{$\Sat{M}{\indfrm}[s_1]$}۔
    بيا \LR{$\Value{t_1}{M}[s_1] = \Value{t_2}{M}[s_1]$}۔ نو
    \begin{align*}
      \Value{t_1}{M}[s_2] & = \Value{t_1}{M}[s_1]
      & \text{\RL{(د \olref{prop:valindep} له مخې)}} \\
      & = \Value{t_2}{M}[s_1]
      & \text{\RL{(ځکه چې \LR{$\Sat{M}{\eq[t_1][t_2]}[s_1]$})}}\\
      &= \Value{t_2}{M}[s_2]
      & \text{\RL{(د \olref{prop:valindep} له مخې)،}}
    \end{align*}
    نو \LR{$\Sat{M}{\eq[t_1][t_2]}[s_2]$}۔}
\end{enumerate}
\textbf{د ژباړې د نسخې سمون (OLFOL-021):} سرچينه د اټومي اړوند د
حالت په وروستي ټوپل کښې لومړي ترم ته په ناسمه توګه د عمومي شاخص
نښه ورکوي؛ دلته هغه د ټوپل له لومړي شاخص سره سم شوے دے۔
\olpseganchor{OLP-0164-B014}{end}

\olpseganchor{OLP-0164-B015}{start}
اوس فرض کړئ چې د~\LR{$!A$} څخه د لږ پېچليو ټولو فارمولونه~\LR{$!B$} دپاره
\LR{$\Sat{M}{!B}[s_1]$} هله او يوازې هله چې \LR{$\Sat{M}{!B}[s_2]$}۔ استقرايي
ګام د~\LR{$!A$} د اصلي چلوونکي له مخې په حالتونو وېشل کېږي۔ په هر
حالت کښې يوازې د دوه اړخيز شرطي مخامخ لور ثابتوو؛ د شاتني لور
ثبوت ئې هم‌شکله دے۔ د سورو له حالتونو پرته، په ټولو حالتونو کښې
استقرايي فرض د~\LR{$!A$} پر فرعي فارمولونه~\LR{$!B$} کاروو۔ د~\LR{$!B$} ازاد
متغيرونه د~\LR{$!A$} د ازادو متغيرونو په منځ کښې دي۔ نو که~\LR{$s_1$} او~\LR{$s_2$}
د~\LR{$!A$} پر ازادو متغيرونو سره موافقې وي، د~\LR{$!B$} پر هغو هم موافقې دي،
او استقرايي فرض پر~\LR{$!B$} کارول کېدے شي۔
\olpseganchor{OLP-0164-B015}{end}

\olpseganchor{OLP-0164-B016}{start}
\begin{enumerate}
\item %
  %
    \indcase{!A}{\lnot !B}{که \LR{$\Sat{M}{\indfrm}[s_1]$} وي، نو
      \LR{$\Sat/{M}{!B}[s_1]$}؛ د استقرايي فرض له مخې
      \LR{$\Sat/{M}{!B}[s_2]$}، ځکه \LR{$\Sat{M}{\indfrm}[s_2]$}۔}
\olpseganchor{OLP-0164-B016}{end}

\olpseganchor{OLP-0164-B017}{start}
\item %
  %
    \indcase{!A}{!B \land !C}{که \LR{$\Sat{M}{\indfrm}[s_1]$} وي، نو
      \LR{$\Sat{M}{!B}[s_1]$} او \LR{$\Sat{M}{!C}[s_1]$}؛ د استقرايي فرض له
      مخې \LR{$\Sat{M}{!B}[s_2]$} او \LR{$\Sat{M}{!C}[s_2]$}۔ ځکه
      \LR{$\Sat{M}{\indfrm}[s_2]$}۔}
\olpseganchor{OLP-0164-B017}{end}

\olpseganchor{OLP-0164-B018}{start}
\item %
  %
    \indcase{!A}{!B \lor !C}{که \LR{$\Sat{M}{\indfrm}[s_1]$} وي، نو
      \LR{$\Sat{M}{!B}[s_1]$} يا \LR{$\Sat{M}{!C}[s_1]$}۔ د استقرايي فرض له
      مخې \LR{$\Sat{M}{!B}[s_2]$} يا \LR{$\Sat{M}{!C}[s_2]$}؛ نو
      \LR{$\Sat{M}{\indfrm}[s_2]$}۔}
\olpseganchor{OLP-0164-B018}{end}

\olpseganchor{OLP-0164-B019}{start}
\item %
  %
    \indcase{!A}{!B \lif !C}{که \LR{$\Sat{M}{\indfrm}[s_1]$} وي، نو
    \LR{$\Sat/{M}{!B}[s_1]$} يا \LR{$\Sat{M}{!C}[s_1]$}۔ د استقرايي فرض له مخې
    \LR{$\Sat/{M}{!B}[s_2]$} يا \LR{$\Sat{M}{!C}[s_2]$}؛ نو
    \LR{$\Sat{M}{\indfrm}[s_2]$}۔}
\olpseganchor{OLP-0164-B019}{end}

\olpseganchor{OLP-0164-B020}{start}
\item %
  %
    \indcase{!A}{!B \liff !C}{که \LR{$\Sat{M}{\indfrm}[s_1]$} وي، نو يا
      \LR{$\Sat{M}{!B}[s_1]$} او \LR{$\Sat{M}{!C}[s_1]$} دواړه، يا
      \LR{$\Sat/{M}{!B}[s_1]$} او \LR{$\Sat/{M}{!C}[s_1]$} دواړه صدق نۀ کوي۔
      د استقرايي فرض له مخې، يا \LR{$\Sat{M}{!B}[s_2]$} او
      \LR{$\Sat{M}{!C}[s_2]$} دواړه، يا \LR{$\Sat/{M}{!B}[s_2]$} او
      \LR{$\Sat/{M}{!C}[s_2]$} دواړه صدق نۀ کوي۔ په هر حالت کښې
      \LR{$\Sat{M}{\indfrm}[s_2]$}۔}
\olpseganchor{OLP-0164-B020}{end}

\olpseganchor{OLP-0164-B021}{start}
\item %
  %
    \indcase{!A}{\lexists[x][!B]}{که \LR{$\Sat{M}{\indfrm}[s_1]$} وي، داسې
      \LR{$m \in \Domain{M}$} شته چې \LR{$\Sat{M}{!B}[\Subst{s_1}{m}{x}]$}۔
      \LR{$s_1' = \Subst{s_1}{m}{x}$} او \LR{$s_2' =
      \Subst{s_2}{m}{x}$} ونيسئ۔ د~\LR{$!B$} ازاد متغيرونه د \LR{$x_1$}،
      \dots، \LR{$x_n$} او~\LR{$x$} په منځ کښې دي۔ \LR{$s_1'(x_i) = s_2'(x_i)$}،
      ځکه \LR{$s_1'$} او~\LR{$s_2'$} په ترتيب سره د~\LR{$s_1$} او~\LR{$s_2$}
      \LR{$x$}-ډولونه دي، او د فرض له مخې \LR{$s_1(x_i) = s_2(x_i)$}۔
      د~\LR{$s_1'$} او~\LR{$s_2'$} د تعريف له مخې
      \LR{$s_1'(x) = s_2'(x) = m$}۔ نو استقرايي فرض پر~\LR{$!B$} او
      \LR{$s_1'$}، \LR{$s_2'$} کارول کېږي، او \LR{$\Sat{M}{!B}[s_2']$} لرو۔ څرنګه
      چې \LR{$s_2' = \Subst{s_2}{m}{x}$}، داسې \LR{$m \in \Domain{M}$} شته چې
      \LR{$\Sat{M}{!B}[\Subst{s_2}{m}{x}]$}؛ ځکه
      \LR{$\Sat{M}{\indfrm}[s_2]$}۔}
\olpseganchor{OLP-0164-B021}{end}

\olpseganchor{OLP-0164-B022}{start}
\item %
  %
    \indcase{!A}{\lforall[x][!B]}{که \LR{$\Sat{M}{\indfrm}[s_1]$} وي، نو
      د هر \LR{$m \in \Domain{M}$} دپاره
      \LR{$\Sat{M}{!B}[\Subst{s_1}{m}{x}]$}۔ غواړو وښيو چې د هر
      \LR{$m \in \Domain{M}$} دپاره
      \LR{$\Sat{M}{!B}[\Subst{s_2}{m}{x}]$} هم صدق کوي۔ يو خپل‌سري
      \LR{$m \in \Domain{M}$} ونيسئ، او
      \LR{$s_1' = \Subst{s_1}{m}{x}$} او \LR{$s_2' =
      \Subst{s_2}{m}{x}$} په پام کښې ونيسئ۔
      \LR{$\Sat{M}{!B}[s_1']$} لرو۔ د~\LR{$!B$} ازاد متغيرونه د \LR{$x_1$}،
      \dots، \LR{$x_n$} او~\LR{$x$} په منځ کښې دي۔
      \LR{$s_1'(x_i) = s_2'(x_i)$}، ځکه \LR{$s_1'$} او~\LR{$s_2'$} په ترتيب سره
      د~\LR{$s_1$} او~\LR{$s_2$} \LR{$x$}-ډولونه دي، او د فرض له مخې
      \LR{$s_1(x_i) = s_2(x_i)$}۔ د~\LR{$s_1'$} او~\LR{$s_2'$} د تعريف له مخې
      \LR{$s_1'(x) = s_2'(x) = m$}۔ نو استقرايي فرض پر~\LR{$!B$} او
      \LR{$s_1'$}، \LR{$s_2'$} کارول کېږي، او \LR{$\Sat{M}{!B}[s_2']$} لرو۔ دا د هر
      \LR{$m \in \Domain{M}$} دپاره صدق کوي، يعنې د ټولو
      \LR{$m \in \Domain{M}$} دپاره
      \LR{$\Sat{M}{!B}[\Subst{s_2}{m}{x}]$}؛ نو
      \LR{$\Sat{M}{\indfrm}[s_2]$}۔}
\end{enumerate}
\textbf{د ژباړې د نسخې سمون (OLFOL-022):} سرچينه د عمومي سور په
حالت کښې دواړه بدلې شوې ګومارنې له بې‌شاخصه ګومارنې څخه جوړوي؛ دلته
په ترتيب سره د لومړۍ او دويمې ګومارنې سم شاخصونه کارول شوي دي۔
د استقرا له مخې پايله اخلو چې هر کله د~\LR{$!A$} ازاد متغيرونه د
\LR{$x_1$}، \dots، \LR{$x_n$} په منځ کښې وي او
\LR{$s_1(x_i)=s_2(x_i)$} د \LR{$i = 1$}، \dots،~\LR{$n$} دپاره وي، نو
\LR{$\Sat{M}{!A}[s_1]$} هله او يوازې هله چې \LR{$\Sat{M}{!A}[s_2]$}۔
\end{proof}
\olpseganchor{OLP-0164-B022}{end}



\olpseganchor{OLP-0164-B024}{start}
\begin{explain}
تړلي فارمولونه ازاد متغيرونه نۀ لري؛ نو هرې دوې متغير ګومارنې د هرې
جملې ټولو (صفر) ازادو متغيرونو ته يو شان څيزونه ګوماري۔ له دې کبله
اوس ثابته شوې قضيه وايي چې تړلے فارمول د متغير ګومارنې په نسبت په
جوړښت کښې صدق کوي که نۀ، له هغې ګومارنې څخه په بشپړ ډول خپلواک دي۔
دا حقيقت ثبتوو۔ همدا د لاندې تعريف بنسټ دے، چې د متغير ګومارنې له
يادولو پرته په جوړښت کښې د تړلے فارمول صدق تعريفوي۔
\end{explain}
\olpseganchor{OLP-0164-B024}{end}

\olpseganchor{OLP-0164-B025}{start}
\begin{cor}
\ollabel{cor:sat-sentence}
که~\LR{$!A$} تړلے فارمول او~\LR{$s$} يوه متغير ګومارنه وي، نو
\LR{$\Sat{M}{!A}[s]$} هله او يوازې هله چې د هرې متغير ګومارنې~\LR{$s'$} دپاره
\LR{$\Sat{M}{!A}[s']$}۔
\end{cor}
\olpseganchor{OLP-0164-B025}{end}

\olpseganchor{OLP-0164-B026}{start}
\begin{proof}
  \LR{$s'$} هره متغير ګومارنه ونيسئ۔ څرنګه چې~\LR{$!A$} تړلے فارمول ده، ازاد
  متغيرونه نۀ لري؛ ځکه هره متغير ګومارنه~\LR{$s'$} په څرګند ډول د~\LR{$!A$}
  ټولو ازادو متغيرونو ته هماغه څيزونه ګوماري چې~\LR{$s$} ئې ګوماري۔ نو د
  \olref{prop:satindep} شرط پوره کېږي، او
  \LR{$\Sat{M}{!A}[s]$} هله او يوازې هله چې \LR{$\Sat{M}{!A}[s']$}۔
\end{proof}
\olpseganchor{OLP-0164-B026}{end}

\olpseganchor{OLP-0164-B027}{start}
\begin{defn}
\ollabel{defn:satisfaction}
که~\LR{$!A$} تړلے فارمول وي، وايو چې جوړښت~\LR{$\Struct M$}
په~\LR{$!A$} کښې \emph{صدق کوي}، \LR{$\Sat{M}{!A}$}، هله او يوازې هله چې د
ټولو متغير ګومارنو~\LR{$s$} دپاره \LR{$\Sat{M}{!A}[s]$}۔
\end{defn}
\olpseganchor{OLP-0164-B027}{end}

\olpseganchor{OLP-0164-B028}{start}
که \LR{$\Sat{M}{!A}$} وي، په ساده ډول دا هم وايو چې \emph{\LR{$!A$}
په~\LR{$\Struct{M}$} کښې رښتيا ده۔} د صدق مفهوم په طبيعي ډول له انفرادي
تړلي فارمولونه څخه د تړلي فارمولونه سټونو ته غځېږي۔
\olpseganchor{OLP-0164-B028}{end}

\olpseganchor{OLP-0164-B029}{start}
\begin{defn}
\ollabel{defn:sat}
که~\LR{$\Gamma$} د تړلي فارمولونه يو سټ وي، وايو چې
جوړښت~\LR{$\Struct M$} په~\LR{$\Gamma$} کښې \emph{صدق کوي}،
\LR{$\Sat{M}{\Gamma}$}، هله او يوازې هله چې د ټولو
\LR{$!A \in \Gamma$} دپاره \LR{$\Sat{M}{!A}$}۔
\end{defn}
\textbf{د ژباړې د نسخې سمون (OLFOL-023):} سرچينه د دې تعريف په
پيل کښې د جملو د سټ نښه بې‌ځايه دوه ځله راوړي؛ دلته تکرار لرې شوے
دے۔
\olpseganchor{OLP-0164-B029}{end}

\olpseganchor{OLP-0164-B030}{start}
\begin{prop}\ollabel{prop:sentence-sat-true}
  فرض کړئ \LR{$\Struct{M}$} جوړښت، \LR{$!A$} تړلے فارمول، او~\LR{$s$}
  يوه متغير ګومارنه ده۔ \LR{$\Sat{M}{!A}$} هله او يوازې هله چې
  \LR{$\Sat{M}{!A}[s]$}۔
\end{prop}
\olpseganchor{OLP-0164-B030}{end}

\olpseganchor{OLP-0164-B031}{start}
\begin{proof}
تمرین۔
\end{proof}
\olpseganchor{OLP-0164-B031}{end}

\olpseganchor{OLP-0164-B032}{start}
\begin{prob}
\olref[fol][syn][ass]{prop:sentence-sat-true} ثابته کړئ۔
\end{prob}
\olpseganchor{OLP-0164-B032}{end}

\olpseganchor{OLP-0164-B033}{start}
\begin{prop}\ollabel{prop:sat-quant}
  فرض کړئ \LR{$!A(x)$} يوازې~\LR{$x$} ازاد لري، او~\LR{$\Struct M$}
  جوړښت دے۔ نو:
  \begin{enumerate}
    \item \LR{$\Sat{M}{\lexists[x][!A(x)]}$} هله او يوازې هله چې
      لږ تر لږه د يوې متغير ګومارنې~\LR{$s$} دپاره
      \LR{$\Sat{M}{!A(x)}[s]$}۔
    \item \LR{$\Sat{M}{\lforall[x][!A(x)]}$} هله او يوازې هله چې
      د ټولو متغير ګومارنو~\LR{$s$} دپاره \LR{$\Sat{M}{!A(x)}[s]$}۔
\end{enumerate}
\end{prop}
\olpseganchor{OLP-0164-B033}{end}

\olpseganchor{OLP-0164-B034}{start}
\begin{proof}
تمرین۔
\end{proof}
\olpseganchor{OLP-0164-B034}{end}

\olpseganchor{OLP-0164-B035}{start}
\begin{prob}
\olref[fol][syn][ass]{prop:sat-quant} ثابته کړئ۔
\end{prob}
\olpseganchor{OLP-0164-B035}{end}

\olpseganchor{OLP-0164-B036}{start}
\begin{prob}
\DeclareRobustCommand{\VDash}{\mathrel{||}\joinrel\Relbar}
فرض کړئ~\LR{$\Lang L$} له تابعې پرته ژبه ده۔ د
جوړښت~\LR{$\Struct M$}، ثابت سمبول~\LR{$c$}، او
\LR{$a \in \Domain M$} په لرلو سره، \LR{$\Struct M[a/c]$} هغه جوړښت
تعريف کړئ چې کټ مټ د~\LR{$\Struct M$} په شان دے، خو
\LR{$\Assign{c}{M[a/c]} = a$}۔ د تړلي فارمولونه~\LR{$!A$} دپاره
\LR{$\Struct M \VDash !A$} داسې تعريف کړئ:
\begin{enumerate}
\item %
  \indcase{!A}{\lfalse}{\LR{$\Struct{M} \VDash \indfrm$} نۀ۔}
\olpseganchor{OLP-0164-B036}{end}

\olpseganchor{OLP-0164-B037}{start}
\item %
  \indcase{!A}{\ltrue}{\LR{$\Struct{M} \VDash \indfrm$}۔}
\olpseganchor{OLP-0164-B037}{end}

\olpseganchor{OLP-0164-B038}{start}
\item \indcase{!A}{\Atom{R}{d_1, \dots, d_n}}{\LR{$\Struct M \VDash \indfrm$}
  هله او يوازې هله چې
  \LR{$\langle \Assign{d_1}{M}, \dots, \Assign{d_n}{M} \rangle \in
  \Assign{R}{M}$}۔}
\olpseganchor{OLP-0164-B038}{end}

\olpseganchor{OLP-0164-B039}{start}
\item \indcase{!A}{\eq[d_1][d_2]}{\LR{$\Struct M \VDash \indfrm$} هله او
  يوازې هله چې \LR{$\Assign{d_1}{M} = \Assign{d_2}{M}$}۔}
\olpseganchor{OLP-0164-B039}{end}

\olpseganchor{OLP-0164-B040}{start}
\item %
  \indcase{!A}{\lnot !B}{\LR{$\Struct{M} \VDash \indfrm$} هله او يوازې
    هله چې \LR{$\Struct{M} \VDash {!B}$} نۀ وي۔}
\olpseganchor{OLP-0164-B040}{end}

\olpseganchor{OLP-0164-B041}{start}
\item %
  \indcase{!A}{(!B \land !C)}{\LR{$\Struct{M} \VDash
    \indfrm$} هله او يوازې هله چې \LR{$\Struct{M} \VDash!B$} او
    \LR{$\Struct{M} \VDash !C$}۔}
\olpseganchor{OLP-0164-B041}{end}

\olpseganchor{OLP-0164-B042}{start}
\item %
  \indcase{!A}{(!B \lor !C)}{\LR{$\Struct{M} \VDash \indfrm$} هله او
    يوازې هله چې \LR{$\Struct{M} \VDash !B$} يا
    \LR{$\Struct{M} \VDash !C$} (يا دواړه)۔}
\olpseganchor{OLP-0164-B042}{end}

\olpseganchor{OLP-0164-B043}{start}
\item %
  \indcase{!A}{(!B \lif !C)}{\LR{$\Struct{M} \VDash \indfrm$} هله او
    يوازې هله چې \LR{$\Struct {M} \VDash !B$} نۀ وي يا
    \LR{$\Struct M \VDash !C$} وي (يا دواړه)۔}
\olpseganchor{OLP-0164-B043}{end}

\olpseganchor{OLP-0164-B044}{start}
\item %
  \indcase{!A}{(!B \liff !C)}{\LR{$\Struct{M} \VDash {\indfrm}$} هله او
    يوازې هله چې يا \LR{$\Struct{M} \VDash {!B}$} او
    \LR{$\Struct{M} \VDash {!C}$} دواړه وي، يا نۀ
    \LR{$\Struct{M} \VDash {!B}$} وي او نۀ \LR{$\Struct{M} \VDash {!C}$}۔}
\olpseganchor{OLP-0164-B044}{end}

\olpseganchor{OLP-0164-B045}{start}
\item %
  \indcase{!A}{\lforall[x][!B]}{\LR{$\Struct{M} \VDash {\indfrm}$} هله او
    يوازې هله چې د ټولو \LR{$a \in \Domain{M}$} دپاره
    \LR{$\Struct{M[a/c]} \VDash \Subst{!B}{c}{x}$}، په دې شرط چې~\LR{$c$}
    په~\LR{$!B$} کښې نۀ وي راغلے۔}
\olpseganchor{OLP-0164-B045}{end}

\olpseganchor{OLP-0164-B046}{start}
\item %
  \indcase{!A}{\lexists[x][!B]}{\LR{$\Struct{M} \VDash
  {\indfrm}$} هله او يوازې هله چې داسې \LR{$a \in \Domain M$} وي چې
  \LR{$\Struct{M[a/c]} \VDash \Subst{!B}{c}{x}$}، په دې شرط چې~\LR{$c$}
  په~\LR{$!B$} کښې نۀ وي راغلے۔}
\end{enumerate}
فرض کړئ \LR{$x_1$}، \dots، \LR{$x_n$} په~\LR{$!A$} کښې ټول ازاد متغيرونه،
\LR{$c_1$}، \dots، \LR{$c_n$} هغه ثابت سمبولونه دي چې په~\LR{$!A$} کښې نۀ دي،
\LR{$a_1$}، \dots، \LR{$a_n \in \Domain M$}، او \LR{$s(x_i) = a_i$}۔
\olpseganchor{OLP-0164-B046}{end}

\olpseganchor{OLP-0164-B047}{start}
وښيئ چې \LR{$\Sat{M}{!A}[s]$} هله او يوازې هله چې
\LR{$\Struct M[a_1/c_1,\dots,a_n/c_n]
\VDash \Subst{\Subst{!A}{c_1}{x_1}\dots}{c_n}{x_n}$}۔
\olpseganchor{OLP-0164-B047}{end}

\olpseganchor{OLP-0164-B048}{start}
(دا مسئله ښيي چې د لومړۍ درجې منطق دپاره داسې معنٰی پوهنه ورکول
ممکن دي چې له متغير ګومارنو پرته کار وکړي۔)
\end{prob}
\olpseganchor{OLP-0164-B048}{end}

\olpseganchor{OLP-0164-B049}{start}
\begin{prob}
فرض کړئ~\LR{$f$} يو داسې تابعي سمبول دے چې په~\LR{$!A(x,y)$} کښې نۀ وي۔ وښيئ
چې داسې جوړښت~\LR{$\Struct{M}$} شته چې
\LR{$\Sat{M}{\lforall[x][\lexists[y][!A(x,y)]]}$} هله او يوازې هله چې داسې
\LR{$\Struct M'$} شته چې \LR{$\Sat{M'}{\lforall[x][!A(x,f(x))]}$}۔
\olpseganchor{OLP-0164-B049}{end}

\olpseganchor{OLP-0164-B050}{start}
(دا مسئله د سکولم د قضيې په نامه د پېژندل شوې قضيې يو ځانګړے حالت
دے؛ \LR{$\lforall[x][!A(x,f(x))]$} ته د
\LR{$\lforall[x][\lexists[y][!A(x,y)]]$} \emph{د سکولم نورماله بڼه} ويل
کېږي۔)
\end{prob}
\olpseganchor{OLP-0164-B050}{end}

\olpseganchor{OLP-0165-B004}{start}
\olfileid{fol}{syn}{ext}
\olpseganchor{OLP-0165-B004}{end}

\olpseganchor{OLP-0165-B005}{start}
\olsection{امتداديتوب}
\olpseganchor{OLP-0165-B005}{end}

\olpseganchor{OLP-0165-B006}{start}
\begin{explain}
امتداديتوب، چې کله کله اړوندتيا هم بلل کېږي، په غير رسمي ډول داسې
څرګندولے شو: يوازېني عوامل چې د فارمول~\LR{$!A$} پر صدق د
جوړښت~\LR{$\Struct M$} او د متغير ګومارنې~\LR{$s$} په نسبت
اغېز لري، د د تعريف ساحه کچه او هغه ګومارنې دي چې~\LR{$\Struct M$} او~\LR{$s$} ئې
د ژبې هغو اجزاوو ته کوي چې په رښتيا په~\LR{$!A$} کښې راځي۔
\olpseganchor{OLP-0165-B006}{end}

\olpseganchor{OLP-0165-B007}{start}
د امتداديتوب يوه سمدستي پايله دا ده: که دوه
جوړښتونه~\LR{$\Struct M$} او~\LR{$\Struct M'$} د ژبې پر ټولو هغو اجزاوو
سره موافق وي چې په جمله~\LR{$!A$} کښې راځي، او يوه ساحه ولري، نو
\LR{$\Struct M$} او~\LR{$\Struct M'$} بايد پر دې هم موافق وي چې پخپله~\LR{$!A$}
رښتيا ده که نۀ۔
\end{explain}
\olpseganchor{OLP-0165-B007}{end}

\olpseganchor{OLP-0165-B008}{start}
\begin{prop}[امتداديتوب]
  \ollabel{prop:extensionality}
  فرض کړئ~\LR{$!A$} فارمول، او~\LR{$\Struct M_1$} او~\LR{$\Struct M_2$} داسې
  جوړښتونه دي چې \LR{$\Domain{M_1} = \Domain{M_2}$}، او~\LR{$s$} پر
  \LR{$\Domain{M_1} = \Domain{M_2}$} يوه متغير ګومارنه ده۔ که په~\LR{$!A$} کښې
  د راتلونکو هر ثابت سمبول~\LR{$c$}، اړوند سمبول~\LR{$R$}، او
  تابع~\LR{$f$} دپاره
  \LR{$\Assign{c}{M_1} = \Assign{c}{M_2}$}،
  \LR{$\Assign{R}{M_1}=\Assign{R}{M_2}$}، او
  \LR{$\Assign{f}{M_1} = \Assign{f}{M_2}$} وي، نو
  \LR{$\Sat{M_1}{!A}[s]$} هله او يوازې هله چې \LR{$\Sat{M_2}{!A}[s]$}۔
\end{prop}
\olpseganchor{OLP-0165-B008}{end}

\olpseganchor{OLP-0165-B009}{start}
\begin{proof}
  لومړے (پر~\LR{$t$} د استقرا له لارې) ثابته کړئ چې د هر ترم دپاره
  \LR{$\Value{t}{M_1}[s] = \Value{t}{M_2}[s]$}۔ بيا قضيه پر~\LR{$!A$} د استقرا
  له لارې ثابته کړئ، او د استقرا په بنسټ کښې (چې~\LR{$!A$} اټومي وي) له
  همدې ثابتې شوې ادعا څخه کار واخلئ۔
\end{proof}
\olpseganchor{OLP-0165-B009}{end}

\olpseganchor{OLP-0165-B010}{start}
\begin{prob}
د \olref[fol][syn][ext]{prop:extensionality} ثبوت په تفصيل بشپړ کړئ۔
\end{prob}
\olpseganchor{OLP-0165-B010}{end}

\olpseganchor{OLP-0165-B011}{start}
\begin{cor}[د تړلي فارمولونه دپاره امتداديتوب]
  \ollabel{cor:extensionality-sent}
  فرض کړئ~\LR{$!A$} تړلے فارمول، او~\LR{$\Struct{M_1}$}، \LR{$\Struct{M_2}$} لکه
  په \olref{prop:extensionality} کښې وي۔ نو \LR{$\Sat{M_1}{!A}$} هله او
  يوازې هله چې \LR{$\Sat{M_2}{!A}$}۔
\end{cor}
\olpseganchor{OLP-0165-B011}{end}

\olpseganchor{OLP-0165-B012}{start}
\begin{proof}
له \olref{prop:extensionality} او \olref[ass]{cor:sat-sentence} څخه
لاس ته راځي۔
\end{proof}
\olpseganchor{OLP-0165-B012}{end}

\olpseganchor{OLP-0165-B013}{start}
پر دې سربېره، د يوه ترم ارزښت، او دا چې جوړښت په
فارمول کښې صدق کوي که نۀ، يوازې د هغه د فرعي ترمونو په
ارزښتونو پورې تړلي دي۔
\olpseganchor{OLP-0165-B013}{end}

\olpseganchor{OLP-0165-B014}{start}
\begin{prop}\ollabel{prop:ext-terms}
فرض کړئ~\LR{$\Struct M$} جوړښت، \LR{$t$} او~\LR{$t'$} ترمونه، او~\LR{$s$} يوه
متغير ګومارنه ده۔ نو \LR{$\Value{\Subst{t}{t'}{x}}{M}[s] =
\Value{t}{M}[\Subst{s}{\Value{t'}{M}[s]}{x}]$}۔
\end{prop}
\olpseganchor{OLP-0165-B014}{end}

\olpseganchor{OLP-0165-B015}{start}
\begin{proof}
پر~\LR{$t$} د استقرا له لارې۔
\begin{enumerate}
\item که~\LR{$t$} يو ثابت وي، يعنې \LR{$t\ident c$}، نو
  \LR{$\Subst{t}{t'}{x} = c$}، او \LR{$\Value{c}{M}[s] = \Assign{c}{M} =
  \Value{c}{M}[\Subst{s}{\Value{t'}{M}[s]}{x}]$}۔
\olpseganchor{OLP-0165-B015}{end}

\olpseganchor{OLP-0165-B016}{start}
\item که~\LR{$t$} له~\LR{$x$} پرته بل متغير وي، يعنې \LR{$t \ident y$}، نو
  \LR{$\Subst{t}{t'}{x} = y$}، او \LR{$\Value{y}{M}[s] =
  \Value{y}{M}[\Subst{s}{\Value{t'}{M}[s]}{x}]$}، ځکه
  \LR{$\varAssign{s}{\Subst{s}{\Value{t'}{M}[s]}{x}}{x}$}۔
\olpseganchor{OLP-0165-B016}{end}

\olpseganchor{OLP-0165-B017}{start}
\item که \LR{$t \ident x$} وي، نو \LR{$\Subst{t}{t'}{x} = t'$}۔ خو د
  \LR{$\Subst{s}{\Value{t'}{M}[s]}{x}$} د تعريف له مخې
  \LR{$\Value{x}{M}[\Subst{s}{\Value{t'}{M}[s]}{x}] = \Value{t'}{M}[s]$}۔
\olpseganchor{OLP-0165-B017}{end}

\olpseganchor{OLP-0165-B018}{start}
\item که \LR{$t \ident \Atom{f}{t_1,\dots,t_n}$} وي، نو:
\begin{multline*}
  \Value{\Subst{t}{t'}{x}}{M}[s]  \\
  \begin{aligned}[b]
& = \Value{\Atom{f}{\Subst{t_1}{t'}{x}, \dots, \Subst{t_n}{t'}{x}}}{M}[s]\\
& \qquad    \text{\RL{د \LR{$\Subst{t}{t'}{x}$} د تعريف له مخې}}\\
& = \Assign{f}{M}(\Value{\Subst{t_1}{t'}{x}}{M}[s], \dots,
    \Value{\Subst{t_n}{t'}{x}}{M}[s])\\
    & \qquad  \text{\RL{د \LR{$\Value{\Atom{f}{\dots}}{M}[s]$} د تعريف له مخې}}\\
& = \Assign{f}{M}(\Value{t_1}{M}[\Subst{s}{\Value{t'}{M}[s]}{x}], \dots,
   \Value{t_n}{M}[\Subst{s}{\Value{t'}{M}[s]}{x}])\\
& \qquad    \text{\RL{د استقرايي فرض له مخې}}\\
& = \Value{t}{M}[\Subst{s}{\Value{t'}{M}[s]}{x}]
    \text{\RL{د \LR{$\Value{\Atom{f}{\dots}}{M}[\Subst{s}{\Value{t'}{M}[s]}{x}]$} د تعريف له مخې}}
  \end{aligned}
\end{multline*}
\textbf{د ژباړې د نسخې سمون (OLFOL-024):} سرچينه د دې محاسبې د
لومړۍ کرښې په پای او د دويمې کرښې په پيل کښې د مساوات نښه دوه ځله
راوړي؛ دلته زياته نښه لرې شوې ده۔
\end{enumerate}
\end{proof}
\olpseganchor{OLP-0165-B018}{end}

\olpseganchor{OLP-0165-B019}{start}
\begin{prop}\ollabel{prop:ext-formulas}
فرض کړئ~\LR{$\Struct M$} جوړښت، \LR{$!A$} فارمول، \LR{$t'$} يو ترم،
او~\LR{$s$} يوه متغير ګومارنه ده۔ نو
\LR{$\Sat{M}{\Subst{!A}{t'}{x}}[s]$} هله او يوازې هله چې
\LR{$\Sat{M}{!A}[\Subst{s}{\Value{t'}{M}[s]}{x}]$}۔
\end{prop}
\olpseganchor{OLP-0165-B019}{end}

\olpseganchor{OLP-0165-B020}{start}
\begin{proof}
تمرین۔
\end{proof}
\olpseganchor{OLP-0165-B020}{end}

\olpseganchor{OLP-0165-B021}{start}
\begin{prob}
\olref[fol][syn][ext]{prop:ext-formulas} ثابته کړئ۔
\end{prob}
\olpseganchor{OLP-0165-B021}{end}

\olpseganchor{OLP-0165-B022}{start}
\begin{explain}
  د \cref{fol:syn:ext:prop:ext-terms,fol:syn:ext:prop:ext-formulas}
  مطلب دا دے۔ فرض کړئ يو ترم~\LR{$t$} يا فارمول~\LR{$!A$}، او بل ترم~\LR{$t'$}
  لرو، او غواړو د~\LR{$\Subst{t}{t'}{x}$} ارزښت، يا دا معلوم کړو چې
  \LR{$\Subst{!A}{t'}{x}$} د متغير ګومارنې~\LR{$s$} په نسبت په
  جوړښت~\LR{$\Struct M$} کښې صدق کوي که نۀ۔ يوې لارې کښې لومړے
  تعويض ترسره کولے، او بيا د~\LR{$\Struct{M}$} او~\LR{$s$} په نسبت ارزښت يا
  صدق څېړلے شو۔ بلې لارې کښې لومړے په~\LR{$\Struct{M}$} کښې د~\LR{$s$} په نسبت
  د~\LR{$t'$} ارزښت~\LR{$m = \Value{t'}{M}[s]$} معلومولے، بيا د
  متغير ګومارنه په~\LR{$\Subst{s}{m}{x}$} بدلولے، او ورپسې په
  \LR{$\Struct{M}$} کښې د~\LR{$\Subst{s}{m}{x}$} په نسبت د~\LR{$t$} ارزښت، يا
  \LR{$\Sat{M}{!A}[\Subst{s}{m}{x}]$} څېړلے شو۔
  \Cref{fol:syn:ext:prop:ext-terms,fol:syn:ext:prop:ext-formulas}
  تضمينوي چې په دواړو لارو کښې ځواب يو شان وي۔
\end{explain}
\olpseganchor{OLP-0165-B022}{end}

\olpseganchor{OLP-0166-B004}{start}
\olfileid{fol}{syn}{sem}
\olsection{معنايي مفهومونه}
\olpseganchor{OLP-0166-B004}{end}

\olpseganchor{OLP-0166-B005}{start}
\begin{explain}
د لومړۍ درجې ژبو د جوړښتونه د تعريف په لرلو سره، د جملو ځينې
بنسټيز معنايي خاصيتونه او د هغو تر منځ اړيکې تعريفولے شو۔ تر ټولو
ساده ئې د جملې د \emph{اعتبار} مفهوم دے۔ يوه جمله هله معتبره ده چې
په هر جوړښت کښې صدق وکړي۔ معتبرې جملې هغه دي چې د هغو د
نامنطقي سمبولونو د تفسير له څرنګوالي پرته صدق کوي۔ ځکه معتبرې جملې
\emph{منطقي حقيقتونه} هم بلل کېږي---هغوی په هر جوړښت کښې
رښتيا دي، يعنې صدق کوي؛ نو رښتينولي ئې يوازې په هغو کښې پر راغلو
منطقي سمبولونو او د هغو پر نحوي جوړښت پورې تړلې ده، پر
نامنطقي سمبولونو يا د هغو پر تفسير نۀ۔
\end{explain}
\olpseganchor{OLP-0166-B005}{end}

\olpseganchor{OLP-0166-B006}{start}
\begin{defn}[اعتبار]
جمله~\LR{$!A$} هله \emph{معتبره} ده، \LR{$\Entails !A$}، چې د هر
جوړښت~\LR{$\Struct M$} دپاره \LR{$\Sat{M}{!A}$}۔
\end{defn}
\olpseganchor{OLP-0166-B006}{end}

\olpseganchor{OLP-0166-B007}{start}
\begin{defn}[استلزام]
د جملو سټ~\LR{$\Gamma$} هله جمله~\LR{$!A$} \emph{لازموي}،
\LR{$\Gamma \Entails !A$}، چې د هر هغه جوړښت~\LR{$\Struct M$} دپاره
چې \LR{$\Sat{M}{\Gamma}$} وي، \LR{$\Sat{M}{!A}$} هم وي۔
\end{defn}
\olpseganchor{OLP-0166-B007}{end}


\olpseganchor{OLP-0166-B008}{start}
\begin{defn}[د صدق وړتيا]
د جملو سټ~\LR{$\Gamma$} هله \emph{د صدق وړ} دے چې د کوم
جوړښت~\LR{$\Struct M$} دپاره \LR{$\Sat{M}{\Gamma}$} وي۔ که~\LR{$\Gamma$} د
صدق وړ نۀ وي، \emph{د صدق نۀ وړ} بلل کېږي۔
\end{defn}
\olpseganchor{OLP-0166-B008}{end}

\olpseganchor{OLP-0166-B009}{start}
\begin{prop}
جمله~\LR{$!A$} هله او يوازې هله معتبره ده چې د جملو د هر سټ~\LR{$\Gamma$}
دپاره \LR{$\Gamma \Entails !A$}۔
\end{prop}
\olpseganchor{OLP-0166-B009}{end}

\olpseganchor{OLP-0166-B010}{start}
\begin{proof}
د مخامخ لور دپاره، فرض کړئ~\LR{$!A$} معتبره ده، او~\LR{$\Gamma$} د جملو يو
سټ دے۔ \LR{$\Struct M$} داسې جوړښت ونيسئ چې
\LR{$\Sat{M}{\Gamma}$}۔ څرنګه چې~\LR{$!A$} معتبره ده، \LR{$\Sat{M}{!A}$}؛ نو
\LR{$\Gamma \Entails !A$}۔
\olpseganchor{OLP-0166-B010}{end}

\olpseganchor{OLP-0166-B011}{start}
د شاتني لور د عکس نقيض دپاره، فرض کړئ~\LR{$!A$} نامعتبره ده؛ نو داسې
جوړښت~\LR{$\Struct M$} شته چې \LR{$\Sat/{M}{!A}$}۔ کله چې
\LR{$\Gamma = \{ \ltrue \}$} وي، د~\LR{$\ltrue$} د اعتبار له کبله
\LR{$\Sat{M}{\Gamma}$}۔ ځکه داسې جوړښت~\LR{$\Struct M$} شته چې
\LR{$\Sat{M}{\Gamma}$} خو \LR{$\Sat/{M}{!A}$}؛ نو~\LR{$\Gamma$}، \LR{$!A$} نۀ لازموي۔
\end{proof}
\olpseganchor{OLP-0166-B011}{end}

\olpseganchor{OLP-0166-B012}{start}
\begin{prop}
\ollabel{prop:entails-unsat}
\LR{$\Gamma \Entails !A$} هله او يوازې هله چې
\LR{$\Gamma \cup \{\lnot !A\}$} د صدق نۀ وړ وي۔
\end{prop}
\olpseganchor{OLP-0166-B012}{end}

\olpseganchor{OLP-0166-B013}{start}
\begin{proof}
د مخامخ لور دپاره، فرض کړئ \LR{$\Gamma \Entails !A$}، او د عکس فرض وکړئ
چې داسې جوړښت~\LR{$\Struct M$} شته چې
\LR{$\Sat{M}{\Gamma \cup \{ \lnot !A \}}$}۔ څرنګه چې
\LR{$\Sat{M}{\Gamma}$} او \LR{$\Gamma \Entails !A$}، نو \LR{$\Sat{M}{!A}$}۔
له بلې خوا، له \LR{$\Sat{M}{\Gamma\cup \{ \lnot !A \}}$} څخه
\LR{$\Sat{M}{\lnot !A}$}؛ نو هم \LR{$\Sat{M}{!A}$} او هم \LR{$\Sat/{M}{!A}$} لرو،
چې تناقض دے۔ ځکه داسې جوړښت~\LR{$\Struct M$} نۀ شي موجودېدے؛ نو
\LR{$\Gamma \cup \{ \lnot !A \}$} د صدق نۀ وړ دے۔
\olpseganchor{OLP-0166-B013}{end}

\olpseganchor{OLP-0166-B014}{start}
د شاتني لور دپاره، فرض کړئ \LR{$\Gamma \cup \{ \lnot !A \}$} د صدق نۀ وړ
دے۔ نو د هر جوړښت~\LR{$\Struct M$} دپاره يا
\LR{$\Sat/{M}{\Gamma}$}، يا \LR{$\Sat{M}{!A}$}۔ ځکه د هر هغه
جوړښت~\LR{$\Struct M$} دپاره چې \LR{$\Sat{M}{\Gamma}$} وي،
\LR{$\Sat{M}{!A}$}؛ نو \LR{$\Gamma \Entails !A$}۔
\end{proof}
\olpseganchor{OLP-0166-B014}{end}

\olpseganchor{OLP-0166-B015}{start}
\begin{prob}
\begin{enumerate}
\item وښيئ چې \LR{$\Gamma \Entails \bot$} هله او يوازې هله چې~\LR{$\Gamma$} د
صدق نۀ وړ وي۔
\item وښيئ چې \LR{$\Gamma \cup \{!A\} \Entails \bot$} هله او يوازې هله چې
\LR{$\Gamma \Entails \lnot !A$}۔
\item فرض کړئ~\LR{$c$} په~\LR{$!A$} يا~\LR{$\Gamma$} کښې نۀ راځي۔ وښيئ چې
  \LR{$\Gamma \Entails \lforall[x][!A]$} هله او يوازې هله چې
  \LR{$\Gamma \Entails \Subst{!A}{c}{x}$}۔
\end{enumerate}
\end{prob}
\olpseganchor{OLP-0166-B015}{end}

\olpseganchor{OLP-0166-B016}{start}
\begin{prop}
که \LR{$\Gamma \subseteq \Gamma'$} او \LR{$\Gamma \Entails !A$} وي، نو
\LR{$\Gamma' \Entails !A$}۔
\end{prop}
\olpseganchor{OLP-0166-B016}{end}

\olpseganchor{OLP-0166-B017}{start}
\begin{proof}
فرض کړئ \LR{$\Gamma \subseteq \Gamma'$} او \LR{$\Gamma \Entails !A$}۔
\LR{$\Struct M$} داسې جوړښت ونيسئ چې \LR{$\Sat{M}{\Gamma'}$}؛ نو
\LR{$\Sat{M}{\Gamma}$}، او له \LR{$\Gamma \Entails !A$} څخه
\LR{$\Sat{M}{!A}$} اخلو۔ ځکه هر کله چې \LR{$\Sat{M}{\Gamma'}$} وي،
\LR{$\Sat{M}{!A}$}؛ نو \LR{$\Gamma' \Entails !A$}۔
\end{proof}
\olpseganchor{OLP-0166-B017}{end}


\olpseganchor{OLP-0166-B018}{start}
\begin{thm}[د معنايي استنتاج قضيه]
\ollabel{thm:sem-deduction}
\LR{$\Gamma \cup \{!A\} \Entails !B$} هله او يوازې هله چې
\LR{$\Gamma \Entails !A \lif !B$}۔
\end{thm}
\olpseganchor{OLP-0166-B018}{end}

\olpseganchor{OLP-0166-B019}{start}
\begin{proof}
د مخامخ لور دپاره، فرض کړئ
\LR{$\Gamma \cup \{ !A \} \Entails !B$}، او~\LR{$\Struct M$} داسې
جوړښت ونيسئ چې \LR{$\Sat{M}{\Gamma}$}۔ که \LR{$\Sat{M}{!A}$} وي، نو
\LR{$\Sat{M}{\Gamma \cup \{ !A \} }$}؛ او څرنګه چې
\LR{$\Gamma \cup \{ !A \}$}، \LR{$!B$} لازموي، \LR{$\Sat{M}{!B}$} اخلو۔ ځکه
\LR{$\Sat{M}{!A \lif !B}$}، نو \LR{$\Gamma \Entails !A \lif !B$}۔
\olpseganchor{OLP-0166-B019}{end}

\olpseganchor{OLP-0166-B020}{start}
د شاتني لور دپاره، فرض کړئ \LR{$\Gamma \Entails !A \lif !B$}، او
\LR{$\Struct M$} داسې جوړښت ونيسئ چې
\LR{$\Sat{M}{\Gamma \cup \{ !A \}}$}۔ بيا \LR{$\Sat{M}{\Gamma}$}؛ نو
\LR{$\Sat{M}{!A \lif !B}$}، او له \LR{$\Sat{M}{!A}$} څخه \LR{$\Sat{M}{!B}$}۔
ځکه هر کله چې \LR{$\Sat{M}{\Gamma \cup \{ !A \} }$} وي،
\LR{$\Sat{M}{!B}$}؛ نو \LR{$\Gamma \cup \{ !A \} \Entails !B$}۔
\end{proof}
\olpseganchor{OLP-0166-B020}{end}

\olpseganchor{OLP-0166-B021}{start}
\begin{prop}\ollabel{prop:quant-terms}
  فرض کړئ~\LR{$\Struct{M}$} جوړښت، \LR{$!A(x)$} داسې فارمول دے
  چې يو ازاد متغير~\LR{$x$} لري، او~\LR{$t$} يو تړلے ترم دے۔ نو:
  \begin{enumerate}
    \item \LR{$!A(t) \Entails \lexists[x][!A(x)]$}
    \item \LR{$\lforall[x][!A(x)] \Entails !A(t)$}
  \end{enumerate}
\end{prop}
\olpseganchor{OLP-0166-B021}{end}

\olpseganchor{OLP-0166-B022}{start}
\begin{proof}
  \begin{enumerate}
    \item %
      فرض کړئ \LR{$\Sat{M}{!A(t)}$}۔ \LR{$s$} داسې
        متغير ګومارنه ونيسئ چې \LR{$s(x) = \Value{t}{M}$}۔ څرنګه چې
        \LR{$!A(t)$} تړلے فارمول ده، \LR{$\Sat{M}{!A(t)}[s]$}۔ د
        \olref[ext]{prop:ext-formulas} له مخې
        \LR{$\Sat{M}{!A(x)}[s]$}۔ د \olref[ass]{prop:sat-quant} له مخې
        \LR{$\Sat{M}{\lexists[x][!A(x)]}$}۔
    \item %
      فرض کړئ
        \LR{$\Sat{M}{\lforall[x][!A(x)]}$}۔ \LR{$s$} داسې متغير ګومارنه ونيسئ
        چې \LR{$s(x) = \Value{t}{M}$}۔ د
        \olref[ass]{prop:sat-quant} له مخې
        \LR{$\Sat{M}{!A(x)}[s]$}۔ د \olref[ext]{prop:ext-formulas} له مخې
        \LR{$\Sat{M}{!A(t)}[s]$}۔ څرنګه چې \LR{$!A(t)$} تړلے فارمول ده، د
        \olref[ass]{prop:sentence-sat-true} له مخې
        \LR{$\Sat{M}{!A(t)}$}۔
  \end{enumerate}
\end{proof}
\olpseganchor{OLP-0166-B022}{end}

\olpseganchor{OLP-0167-B004}{start}
\olchapter{fol}{mat}{تيورۍ او د هغو مدلونه}
\olpseganchor{OLP-0167-B004}{end}

\olpseganchor{OLP-0168-B004}{start}
\olfileid{fol}{mat}{int}
\olsection{پېژندنه}
\olpseganchor{OLP-0168-B004}{end}

\olpseganchor{OLP-0168-B005}{start}
\begin{explain}
د بديهي طريقې وده د ساينس په تاريخ کښې يوه ستره لاسته راوړنه ده،
او د رياضۍ په تاريخ کښې ځانګړے اهميت لري۔ د يوه ډګر بديهي وده ډېرې
پوښتنې روښانوي: دا ډګر د څه په باب دے؟ تر ټولو بنسټيز مفاهيم ئې کوم
دي؟ هغوئ له يو بل سره څنګه تړلي دي؟ ايا د ډګر ټول مفاهيم د دغو
بنسټيزو مفاهيمو له مخې تعريفېدے شي؟ دا مفاهيم د کومو قوانينو پيروي
کوي، او بايد ئې وکړي؟
\olpseganchor{OLP-0168-B005}{end}

\olpseganchor{OLP-0168-B006}{start}
بديهي طريقه او منطق د يو بل دپاره جوړ شوي دي۔ صوري منطق د بديهي
تيوريو د جوړولو، د تيورۍ له بديهي اصولو څخه په کره ټاکلې طريقه د
قضيو د ثابتولو، او د بديهي اصولو پوره کوونکو ټولو نظامونو د خاصيتونو
د منظمې څېړنې وسيلې برابروي۔
\end{explain}
\olpseganchor{OLP-0168-B006}{end}

\olpseganchor{OLP-0168-B007}{start}
\begin{defn}
د تړلي فارمولونه يو سټ~\LR{$\Gamma$} هله او يوازې هله \emph{بند} دے چې
هر کله \LR{$\Gamma \Entails !A$} وي، نو \LR{$!A \in \Gamma$} هم وي۔ د
تړلي فارمولونه د يوه سټ~\LR{$\Gamma$} \emph{بندښت} دا دے:
\LR{$\Setabs{!A}{\Gamma \Entails !A}$}۔
\olpseganchor{OLP-0168-B007}{end}

\olpseganchor{OLP-0168-B008}{start}
وايو چې~\LR{$\Gamma$} د جملو د يوه سټ~\LR{$\Delta$} له خوا \emph{په بديهي
اصولو ټاکل شوے} دے، که~\LR{$\Gamma$} د~\LR{$\Delta$} بندښت وي۔
\end{defn}
\olpseganchor{OLP-0168-B008}{end}

\olpseganchor{OLP-0168-B009}{start}
\begin{explain}
يوه بديهي تيوري د هغو جملو سټ ګڼلے شو چې د هغې د بديهي اصولو
سټ~\LR{$\Delta$} په بديهي اصولو ټاکلے وي۔ په بل عبارت، کله چې د لومړۍ درجې ژبه د هغه
ساينس د بنسټيزو مفاهيمو دپاره غيرمنطقي سمبولونه ولري چې غواړو په
بديهي ډول ئې وڅېړو، او ورسره د تړلي فارمولونه داسې سټ ولرو چې د
ساينس بنسټيز قوانين بيانوي، نو تيوري په دې ژبه کښې د ټولو هغو
تړلي فارمولونه په توګه انځورولے شو چې له بديهي اصولو څخه لازمېږي۔
دا لړۍ له هغو ساده بېلګو څخه، چې يوازې يو بنسټيز مفهوم او ساده
بديهي اصول لري، لکه د جزوي ترتيبونو تيوري، تر نيوټني ميخانيک غوندې
پېچلو تيوريو پورې غځېږي۔
\olpseganchor{OLP-0168-B009}{end}

\olpseganchor{OLP-0168-B010}{start}
لاندې مهم منطقي حقيقتونه دي چې دغه صوري چلند د بديهي طريقې دپاره
دومره ارزښتمن کوي۔ فرض کړئ~\LR{$\Gamma$} د يوې تيورۍ بديهي نظام، يعنې د
جملو يو سټ دے۔
\begin{enumerate}
\item په کره توګه ويلے شو چې يو بديهي نظام د جوړښتونه مطلوبه
  ټولګه کله رانغاړي۔ يعنې که د جوړښتونه يوه ټاکلې ټولګه مو په
  پام کښې وي، بديهي نظام~\LR{$\Gamma$} هغه هله او يوازې هله په برياليتوب
  رانغاړي چې مطلوب جوړښتونه کټ مټ هغه~\LR{$\Struct M$} وي چې
  \LR{$\Sat{M}{\Gamma}$} وي۔
\item کېدے شي په دې کار کښې ځکه پاتې راشو چې داسې~\LR{$\Struct M$} وي
  چې \LR{$\Sat{M}{\Gamma}$}، خو~\LR{$\Struct M$} زموږ له مطلوبو
  جوړښتونه څخه نۀ وي۔ دا ښايي دې ته مو اړ کړي چې داسې بديهي
  اصول ورزيات کړو چې په~\LR{$\Struct M$} کښې رښتوني نۀ وي۔
\item که لږ تر لږه په دې کښې بريالي يو چې~\LR{$\Gamma$} په ټولو مطلوبو
  جوړښتونه کښې رښتونے وي، نو هر کله چې
  \LR{$\Gamma \Entails !A$} وي، جمله~\LR{$!A$} هم په ټولو مطلوبو
  جوړښتونه کښې رښتونې ده۔ نو د منطقي وسيلو (لکه د
  اشتقاق طريقو) په کارولو ښودلے شو چې جملې په ټولو مطلوبو
  جوړښتونه کښې رښتونې دي: بس دا وښيو چې له بديهي اصولو څخه
  لازمېږي۔
\item کله کله مو کوم مطلوب جوړښتونه په پام کښې نۀ وي، بلکې
  له خپلو بديهي اصولو پيل کوو: ځينې بنسټيز مفاهيم اخلو چې غواړو
  ټاکلي قوانين پوره کړي، او دغه قوانين په يوه بديهي نظام کښې ثبتوو۔
  سمدستي دا تصديق کول غواړو چې بديهي اصول يو بل نۀ نقضوي؛ که ئې
  وکړي، هېڅ مفاهيم دغه قوانين نۀ شي پوره کولے او تيوري مو ناسازګاره
  جوړه کړې ده۔ د~\LR{$\Gamma$} د يوه مدل په موندلو تصديق کولے شو چې داسې
  نۀ دي شوي۔ او که تيوري مو مدلونه ولري، په منطقي طريقو ئې څېړلے او
  مدلونه هم جوړولے شو۔
\item د بديهي اصولو استقلال هم يوه مهمه پوښتنه ده۔ کېدے شي يو بديهي
  اصل په رښتيا د نورو پايله وي او له همدې امله زائد وي۔ ثابتولے شو
  چې په~\LR{$\Gamma$} کښې بديهي اصل~\LR{$!A$} زائد دے، که
  \LR{$\Gamma \setminus \{!A\} \Entails !A$} ثابت کړو۔ دا هم ثابتولے شو
  چې بديهي اصل زائد نۀ دے، که وښيو چې
  \LR{$(\Gamma \setminus \{!A\}) \cup \{\lnot !A\}$} د صدق وړ دے۔ د
  بېلګې په توګه، د هندسې د نورو بديهي اصولو څخه د موازي کرښو د اصل
  استقلال همداسې ښودل شوے و۔
\item بله مهمه پوښتنه په يوه تيورۍ کښې د مفاهيمو د تعريفېدنې په باب ده: د
  ژبې ټاکنه دا ټاکي چې د تيورۍ مدلونه له څه جوړ وي۔ خو د تيورۍ هر اړخ
  لازم نۀ دے چې په مدلونو کښې جلا استازيتوب ولري۔ مثلاً هر ترتيب
  \LR{$\le$} ورسره سم سخت ترتيب~\LR{$<$} ټاکي---چې يو راکړل شي، بل ترې تعريفولے
  شو۔ نو د داسې ترتيب لرونکې تيورۍ په مدل کښې د اړوند سخت ترتيب
  \emph{جلا} شتون لازم نۀ دے۔ په عمومي ډول، يوه اړيکه د نورو له مخې
  کله تعريفېدے شي؟ يوه اړيکه د نورو له مخې کله نۀ شي تعريفېدے (او له
  همدې امله بايد د ژبې په بنسټيزو سمبولونو کښې ورزياته شي)؟
\end{enumerate}
\end{explain}
\olpseganchor{OLP-0168-B010}{end}

\olpseganchor{OLP-0169-B004}{start}
\olfileid{fol}{mat}{exs}
\olpseganchor{OLP-0169-B004}{end}

\olpseganchor{OLP-0169-B005}{start}
\olsection{د جوړښتونه د خاصيتونو بيانول}
\olpseganchor{OLP-0169-B005}{end}

\olpseganchor{OLP-0169-B006}{start}
\begin{explain}
د تابعو او اړيکو شرطونه، يا په لا عمومي ډول دا خبره بيانول چې په
يوه جوړښت کښې تابعې او اړيکې دغه شرطونه پوره کوي، ډېر ځله ګټور او
مهم وي۔ د بېلګې په توګه، غواړو د يوې ژبې هغه جوړښتونه چې
پريديکاتونه پکښې زموږ مطلوبه ``معنٰی'' لري، له هغو جوړښتونو څخه
بېل کړو چې داسې نۀ دي۔ البته بشپړه ازادي لرو چې خپل ``مطلوب''
جوړښتونه وټاکو؛ مثلاً ويلے شو چې د پريديکات~\LR{$\le$} تفسير بايد يو
ترتيب وي، يا يوازې د~\LR{$\Lang L$} له هغو تفسيرونو سره دلچسپي لرو چې
ساحه ئې له سټونو جوړه وي او~\LR{$\Obj \in$} د ``د غړے په توګه
شتون'' په اړيکه تفسير شي۔ خو ايا دا د ژبې په تړلي فارمولونه کولے شو؟
په بل عبارت، پر جوړښت~\LR{$\Struct M$} کوم شرطونه د~\LR{$\Struct M$}
د ژبې په يوه تړلے فارمول (يا ښايي د تړلي فارمولونه په يوه سټ)
بيانولے شو؟ ځينې شرطونه نۀ شو بيانولے۔ مثلاً د~\LR{$\Lang L_A$} داسې
هېڅ جمله نشته چې يوازې په هغه جوړښت~\LR{$\Struct M$} کښې رښتونې
وي چې~\LR{$\Domain M = \Nat$} وي۔ ``ساحه يوازې طبيعي عددونه لري'' نۀ شو
بيانولے۔ خو ښايي د جوړښتونه ځينې ``جوړښتي خاصيتونه'' بيانولے
شو۔ د جوړښتونه کوم خاصيتونه په تړلي فارمولونه بيانولے شو؟ يا
په بله وينا، د جوړښتونه کومې ټولګې د هغو جوړښتونو په توګه
بيانولے شو چې يوه تړلے فارمول (يا د تړلي فارمولونه يو سټ) رښتونې
کوي؟
\end{explain}
\olpseganchor{OLP-0169-B006}{end}

\olpseganchor{OLP-0169-B007}{start}
\begin{defn}[د يوه سټ مدل]
فرض کړئ~\LR{$\Gamma$} په ژبه~\LR{$\Lang L$} کښې د تړلي فارمولونه يو سټ دے۔
وايو چې جوړښت~\LR{$\Struct M$} د~\LR{$\Gamma$} \emph{مدل دے}، که د
هر~\LR{$!A \in \Gamma$} دپاره \LR{$\Sat{M}{!A}$} وي۔
\end{defn}
\olpseganchor{OLP-0169-B007}{end}

\olpseganchor{OLP-0169-B008}{start}
\begin{ex}
جمله~\LR{$\lforall[x][x \le x]$} په~\LR{$\Struct M$} کښې هله او يوازې هله
رښتونې ده چې~\LR{$\Assign{\le}{M}$} انعکاسي اړيکه وي۔ جمله
\LR{$\lforall[x][\lforall[y][((x \le y \land y \le x) \lif x = y)]]$}
په~\LR{$\Struct M$} کښې هله او يوازې هله رښتونې ده چې
\LR{$\Assign{\le}{M}$} ضد متناظره وي۔ جمله
\LR{$\lforall[x][\lforall[y][\lforall[z][((x \le y \land y \le z)
      \lif x \le z)]]]$} په~\LR{$\Struct M$} کښې هله او يوازې هله رښتونې
ده چې~\LR{$\Assign{\le}{M}$} انتقالي وي۔ نو د لاندې سټ مدلونه
\begin{align*}
\{\quad &\lforall[x][x \le x], \\
   & \lforall[x][\lforall[y][((x \le y \land y \le
    x) \lif x = y)]], \\
   &\lforall[x][\lforall[y][\lforall[z][((x \le y
      \land y \le z) \lif x \le z)]]] \quad \}
\end{align*}
کټ مټ هغه جوړښتونه دي چې~\LR{$\Assign{\le}{M}$} پکښې انعکاسي، ضد
متناظره او انتقالي، يعنې جزوي ترتيب وي۔ له همدې امله، دا د
\emph{جزوي ترتيبونو د لومړۍ درجې تيورۍ} د بديهي اصولو په توګه
اخستلے شو۔
\end{ex}
\olpseganchor{OLP-0169-B008}{end}

\olpseganchor{OLP-0170-B004}{start}
\olfileid{fol}{mat}{the}
\olpseganchor{OLP-0170-B004}{end}

\olpseganchor{OLP-0170-B005}{start}
\olsection{د لومړۍ درجې تيوريو بېلګې}
\olpseganchor{OLP-0170-B005}{end}

\olpseganchor{OLP-0170-B006}{start}
\begin{ex}
په ژبه~\LR{$\Lang L_<$} کښې د سختو خطي ترتيبونو تيوري د لاندې سټ له خوا
په بديهي اصولو ټاکل کېږي:
\begin{align*}
\{\quad & \lforall[x][\lnot x < x], \\
& \lforall[x][\lforall[y][((x < y \lor y <
    x) \lor x = y)]], \\
& \lforall[x][\lforall[y][\lforall[z][((x < y
      \land y < z) \lif x < z)]]] \quad \}
\end{align*}
دا مطلوب جوړښتونه په بشپړ ډول رانغاړي: هر سخت خطي ترتيب د دې
بديهي نظام مدل دے؛ او برعکس، که~\LR{$R$} پر يوه سټ~\LR{$X$} خطي ترتيب وي، نو
هغه جوړښت~\LR{$\Struct M$} چې \LR{$\Domain M = X$} او
\LR{$\Assign{<}{M} = R$} لري، د دې تيورۍ مدل دے۔
\end{ex}
\olpseganchor{OLP-0170-B006}{end}

\olpseganchor{OLP-0170-B007}{start}
\begin{ex}
د ګروپونو تيوري په هغه ژبه کښې چې~\LR{$\Obj 1$} (ثابت سمبول) او~\LR{$\cdot$}
(دوه‌ځاييزه تابع) لري، د لاندې فورمولونو له خوا په بديهي
اصولو ټاکل کېږي:
\begin{align*}
& \lforall[x][\eq[(x \cdot \Obj 1)][x]]\\
& \lforall[x][\lforall[y][\lforall[z][\eq[(x \cdot (y \cdot z))][((x
          \cdot y) \cdot z)]]]]\\
& \lforall[x][\lexists[y][\eq[(x \cdot y)][\Obj 1]]]
\end{align*}
\end{ex}
\olpseganchor{OLP-0170-B007}{end}

\olpseganchor{OLP-0170-B008}{start}
\begin{ex}
د پيانو د حساب تيوري د حساب په ژبه~\LR{$\Lang L_A$} کښې د لاندې جملو له
خوا په بديهي اصولو ټاکل کېږي۔
\begin{align*}
& \lforall[x][\lforall[y][(\eq[x'][y'] \lif \eq[x][y])]]\\
& \lforall[x][\eq/[\Obj 0][x']]\\
& \lforall[x][\eq[(x + \Obj 0)][x]]\\
& \lforall[x][\lforall[y][\eq[(x + y')][(x + y)']]]\\
& \lforall[x][\eq[(x \times \Obj 0)][\Obj 0]]\\
& \lforall[x][\lforall[y][\eq[(x \times y')][((x \times y) + x)]]]\\
& \lforall[x][\lforall[y][(x < y \liff \lexists[z][\eq[(z' + x)][y])]]]\\
\intertext{\RL{او ورسره د لاندې بڼې ټولې جملې}}
& (!A(\Obj 0) \land \lforall[x][(!A(x) \lif !A(x'))]) \lif \lforall[x][!A(x)]
\end{align*}
ځکه چې د وروستۍ بڼې نامتناهي ډېرې جملې شته، دا بديهي نظام
نامتناهي دے۔ وروستۍ بڼه \emph{د استقرا شېما} بلل کېږي۔ (په اصل کښې
د استقرا شېما تر هغه لږه پېچلې ده چې دلته مو ښودلې ده۔)
\olpseganchor{OLP-0170-B008}{end}

\olpseganchor{OLP-0170-B009}{start}
وروستے بديهي اصل د~\LR{$<$} يو \emph{صريح تعريف} دے۔
\end{ex}
\olpseganchor{OLP-0170-B009}{end}

\olpseganchor{OLP-0170-B010}{start}
\begin{ex}
د خالصو سټونو تيوري د رياضۍ په بنسټونو (او فلسفه) کښې مهمه ونډه
لري۔ يو سټ هله خالص دے چې ټول غړي ئې هم خالص سټونه وي۔ نو
تش سټ خالص ګڼل کېږي، خو هغه سټ چې کوم داسې څيز ئې غړے وي
چې خپله سټ نۀ وي، خالص نۀ دے۔ په دې توګه خالص سټونه يوازې له تش سټ
څخه، بې له ``غيرسټي بنسټيزو عناصرو''، جوړ شوي سټونه دي؛ يعنې داسې
څيزونه نۀ لري چې خپله سټونه نۀ وي۔
\olpseganchor{OLP-0170-B010}{end}

\olpseganchor{OLP-0170-B011}{start}
لاندې فورمولونه د خالصو سټونو د يوې تيورۍ بديهي نظام ګڼل کېدے شي:
\begin{align*}
& \lexists[x][\lnot \lexists[y][y \in x]]\\
& \lforall[x][\lforall[y][(\lforall[z](z \in x \liff z \in y) \lif
\eq[x][y])]]\\
& \lforall[x][\lforall[y][\lexists[z][\lforall[u][(u \in z \liff
(\eq[u][x] \lor \eq[u][y]))]]]]\\
& \lforall[x][\lexists[y][\lforall[z][(z \in y \liff \lexists[u][(z \in
        u \land u \in x)])]]]\\
\intertext{\RL{او ورسره د لاندې بڼې ټولې جملې}} &
\lexists[x][\lforall[y][(y \in x \liff !A(y))]]
\end{align*}
لومړے بديهي اصل وايي چې داسې يو سټ شته چې هېڅ غړي نۀ لري
(يعنې~\LR{$\emptyset$} شته)؛ دويم وايي چې سټونه امتدادي دي؛ درېيم وايي چې د هر
دوو سټونو~\LR{$X$} او~\LR{$Y$} دپاره سټ~\LR{$\{X, Y\}$} شته؛ څلورم وايي چې د هر
سټ~\LR{$X$} دپاره سټ~\LR{$\cup X$} شته، چېرته چې~\LR{$\cup X$} د~\LR{$X$} د ټولو
عناصرو اتحاد دے۔
\olpseganchor{OLP-0170-B011}{end}

\olpseganchor{OLP-0170-B012}{start}
وروستۍ يادې شوې تړلي فارمولونه په ګډه \emph{د ساده ادراک شېما} بلل
کېږي۔ په اصل کښې وايي چې د هر~\LR{$!A(x)$} دپاره سټ
\LR{$\Setabs{x}{!A(x)}$} شته---نو په لومړي نظر يو رښتونے، ګټور او ښايي
حتا ضروري بديهي اصل ښکاري۔ ``ساده'' ځکه بلل کېږي چې په پايله کښې دا
تيوري د صدق وړ نۀ پرېږدي: که~\LR{$!A(y)$} د~\LR{$\lnot y \in y$} په توګه
واخلئ، لاندې تړلے فارمول لاس ته راځي:
\[
\lexists[x][\lforall[y][(y \in x \liff \lnot y \in y)]]
\]
او دا تړلے فارمول په هېڅ جوړښت کښې نۀ صدق کوي۔
\end{ex}
\olpseganchor{OLP-0170-B012}{end}

\olpseganchor{OLP-0170-B013}{start}
\begin{ex}
د \emph{مېرېولوژۍ} يا جزپوهنې په ډګر کښې د \emph{برخه‌والي} اړيکه
بنسټيزه اړيکه ده۔ لکه د سټونو د تيوريو په شان، د برخه‌والي داسې
تيورۍ هم شته چې د دې اړيکې بېلابېل (او کله کله يو بل سره ټکر لرونکي)
تصورات په بديهي اصولو ټاکي۔
\olpseganchor{OLP-0170-B013}{end}

\olpseganchor{OLP-0170-B014}{start}
د مېرېولوژۍ ژبه يو دوه‌ځاييز اړوند سمبول~\LR{$\Obj P$} لري، او
\LR{$\Atom{\Obj P}{x, y}$} دا ``معنٰی'' لري چې~\LR{$x$} د~\LR{$y$} يوه برخه ده۔
کله چې دا تفسير په پام کښې ولرو، د دې ژبې جوړښت د
\emph{برخه‌والي جوړښت} بلل کېږي۔ البته د يوه دوه‌ځاييز اړوند هر
جوړښت په رښتيا د دې نوم وړ نۀ دے۔ د ``برخه‌والي'' د رانغاړلو دپاره
بايد~\LR{$\Assign{\Obj P}{M}$} ځينې شرطونه پوره کړي، او دا شرطونه د
برخه‌والي د تيورۍ د بديهي اصولو په توګه ټاکلے شو۔ مثلاً برخه‌والے
پر څيزونو جزوي ترتيب دے: هر څيز د خپل ځان (که څه هم \emph{غيرحقيقي})
برخه ده؛ دوه بېل څيزونه د يو بل برخې نۀ شي کېدے؛ او د يوه څيز د
برخې برخه خپله هم د هغه څيز برخه ده۔ پام وکړئ چې په دې معنٰی ``د
يوه څيز برخه ده'' له ``فرعي سټ دے'' سره ورته ده، خو له ``عنصر دے''
سره نۀ ده ورته، ځکه عنصرول نۀ انعکاسي دي او نۀ انتقالي۔
\begin{align*}
& \lforall[x][\Atom{\Obj P}{x,x}] \\
& \lforall[x][\lforall[y][((\Part{x}{y} \land \Part{y}{x})
      \lif \eq[x][y])]] \\
& \lforall[x][\lforall[y][\lforall[z][((\Part{x}{y} \land
        \Part{y}{z}) \lif \Part{x}{z})]]]\\
\intertext{\RL{سربېره پر دې، هر دوه څيزونه مېرېولوژيکي مجموعه لري (داسې
  څيز چې دغه دواړه څيزونه ئې برخې وي او په دې لحاظ تر ټولو کوچنے وي)۔}}  &
\lforall[x][\lforall[y][\lexists[z][\lforall[u][(\Part{z}{u} \liff
        (\Part{x}{u} \land \Part{y}{u}))]]]]
\end{align*}
دا يوازې د برخه‌والي ځينې هغه بنسټيز اصول دي چې مابعدالطبيعتپوهان ئې
څېړي۔ خو نور اصول ډېر ژر داسې پېچلي شي چې تر ځينو تعريف شويو اړيکو
وړاندې ئې جوړول يا ليکل ستونزمن وي۔ مثلاً په مېرېولوژۍ کښې زياتره
دلچسپي لرونکي مابعدالطبيعتپوهان لاندې اصل هم معتبر ګڼي: هر کله چې
څيز~\LR{$x$} حقيقي برخه~\LR{$y$} ولري، نو داسې برخه~\LR{$z$} هم لري چې له~\LR{$y$} سره
هېڅ ګډه برخه ونۀ لري، او د~\LR{$y$} او~\LR{$z$} امتزاج~\LR{$x$} وي۔
\end{ex}
\olpseganchor{OLP-0170-B014}{end}

\olpseganchor{OLP-0171-B004}{start}
\olfileid{fol}{mat}{exr}
\olpseganchor{OLP-0171-B004}{end}

\olpseganchor{OLP-0171-B005}{start}
\olsection{په  جوړښت کښې د
  اړيکو بيانول}
\olpseganchor{OLP-0171-B005}{end}

\olpseganchor{OLP-0171-B006}{start}
\begin{explain}
د فارمولونه يوه اصلي ګټه دا ده چې په جوړښت~\LR{$\Struct M$}
کښې د~\LR{$\Struct M$} د ژبې~\LR{$\Lang L$} د بنسټيزو سمبولونو له مخې
خاصيتونه او اړيکې بيانوي۔ مطلب مو دا دے: د~\LR{$\Struct M$} د تعريف ساحه
د څيزونو يو سټ دے۔ ثابت سمبولونه، تابعې او پريديکاتونه
په~\LR{$\Struct M$} کښې په ترتيب سره د~\LR{$\Domain M$} په ځينو څيزونو،
پر~\LR{$\Domain M$} تابعو او پر~\LR{$\Domain M$} اړيکو تفسيرېږي۔ مثلاً که
\LR{$\Obj A^2_0$} په~\LR{$\Lang L$} کښې وي، نو~\LR{$\Struct M$} ورته
اړيکه~\LR{$R = \Assign{{\Obj A^2_0}}{M}$} ګوماري۔ بيا فارمول
\LR{$\Atom{\Obj A^2_0}{\Obj v_1, \Obj v_2}$} همدغه اړيکه په لاندې معنٰی
\emph{بيانوي}: که يوه متغير-ګومارنه~\LR{$s$}، \LR{$\Obj v_1$} له
\LR{$a \in \Domain{M}$} او~\LR{$\Obj v_2$} له~\LR{$b \in \Domain M$} سره وتړي، نو
\[
Rab \text{\RL{\quad هله او يوازې هله چې\quad}} \Sat{M}{\Atom{\Obj A^2_0}{\Obj v_1, \Obj v_2}}[s].
\]
پام وکړئ چې دلته بايد متغير-ګومارنې شاملې کړو: يوازې دا نۀ شو ويلے
چې ``\LR{$Rab$} هله او يوازې هله چې
\LR{$\Sat{M}{\Atom{\Obj A^2_0}{a, b}}$}''، ځکه~\LR{$a$} او~\LR{$b$} زموږ د ژبې
سمبولونه نۀ، بلکې د~\LR{$\Domain{M}$} غړي دي۔
\olpseganchor{OLP-0171-B006}{end}

\olpseganchor{OLP-0171-B007}{start}
موږ يوازې اټومي فارمولونه نۀ لرو، بلکې په منطقي نښلوونکو او
سورونو ئې يوځاے کولے شو؛ نو لا پېچلي فارمولونه نورې هغه اړيکې
تعريفولے شي چې نېغ په نېغه په~\LR{$\Struct M$} کښې نۀ دي شاملې۔ موږ دې
ته ګورو چې دا څنګه کېږي، او په ځانګړي ډول په جوړښت کښې کومې
اړيکې تعريفولے شو۔
\end{explain}
\olpseganchor{OLP-0171-B007}{end}

\olpseganchor{OLP-0171-B008}{start}
\begin{defn}
فرض کړئ~\LR{$!A(\Obj v_1,\dots, \Obj v_n)$} د~\LR{$\Lang L$} داسې
فارمول دے چې يوازې~\LR{$\Obj v_1$}،\dots،~\LR{$\Obj v_n$} پکښې ازاد
راځي، او~\LR{$\Struct M$} د~\LR{$\Lang L$} يو جوړښت دے۔
\LR{$!A(\Obj v_1,\dots, \Obj v_n)$} هله او يوازې هله
\emph{اړيکه}~\LR{$R \subseteq \Domain M^n$} \emph{بيانوي} چې
\[
Ra_1\dots a_n \text{\RL{\quad هله او يوازې هله چې\quad}} \Sat{M}{\Atom{!A}{\Obj
    v_1,\dots, \Obj v_n}}[s]
\]
د هرې داسې متغير-ګومارنې~\LR{$s$} دپاره چې
\LR{$s(\Obj v_i) = a_i$} (\LR{$i = 1, \dots, n$}) وي۔
\end{defn}
\olpseganchor{OLP-0171-B008}{end}

\olpseganchor{OLP-0171-B009}{start}
\begin{ex}
د حساب په معياري مدل~\LR{$\Struct N$} کښې فارمول
\LR{$\Obj v_1 < \Obj v_2 \lor \eq[\Obj v_1][\Obj v_2]$} پر~\LR{$\Nat$} د
\LR{$\le$} اړيکه بيانوي۔ فارمول~\LR{$\eq[\Obj v_2][\Obj v_1']$} د ناستي
اړيکه بيانوي؛ يعنې هغه~\LR{$R \subseteq \Nat^2$} چې~\LR{$Rnm$} هله وي چې~\LR{$m$}
د~\LR{$n$} ناستے وي۔ فارمول~\LR{$\eq[\Obj v_1][\Obj v_2']$} د مخکښ اړيکه
بيانوي۔ دواړه فارمولونه
\LR{$\lexists[\Obj v_3][(\eq/[\Obj v_3][\Obj 0] \land
  \eq[\Obj v_2][(\Obj v_1 + \Obj v_3)])]$} او~\LR{$\lexists[\Obj
  v_3][\eq[(\Obj v_1 + {\Obj v_3}')][\Obj v_2]]$} د~\LR{$\Obj <$} اړيکه
بيانوي۔ د دې معنٰی دا ده چې اړوند سمبول~\LR{$<$} په رښتيا د حساب په ژبه
کښې زائد دے؛ تعريفېدے شي۔
\textbf{د ژباړې د نسخې سمون (OLFOL-025):} سرچينې د دويم تعريفوونکي
فارمول له وروستي متغير څخه د څيز-ژبې بڼه غورځولې وه؛ دلته هغه بڼه د
همدې فارمول له نورو متغيرونو سره يو شان بېرته راوستل شوې ده۔
\end{ex}
\olpseganchor{OLP-0171-B009}{end}

\olpseganchor{OLP-0171-B010}{start}
\begin{explain}
دا مفکوره يوازې په ځانګړو جوړښتونه کښې په زړه پورې نۀ ده،
بلکې هر کله مهمه ده چې يوه ژبه د يوه يا څو مطلوبو مدلونو د بيانولو
دپاره کاروو، يعنې تيورۍ څېړو۔ دغه تيورۍ زياتره يوازې څو
پريديکاتونه د بنسټيزو سمبولونو په توګه لري، خو په هغه ساحه کښې
چې بيانوي ئې، ډېرې نورې اړيکې مهمې ونډې لري۔ که دغه نورې اړيکې د ژبې
د بنسټيزو پريديکاتونه د تفسيروونکو اړيکو له مخې په منظمه توګه
بيانېدے شي، وايو چې هغوئ په ژبه کښې \emph{تعريفولے} شو۔
\end{explain}
\olpseganchor{OLP-0171-B010}{end}

\olpseganchor{OLP-0171-B011}{start}
\begin{prob}
په~\LR{$\Lang L_A$} کښې داسې فارمولونه ومومئ چې لاندې اړيکې تعريف کړي:
\begin{enumerate}
\item \LR{$n$} د~\LR{$i$} او~\LR{$j$} تر منځ دے؛
\item \LR{$n$}، \LR{$m$} په پوره ډول وېشي (يعنې~\LR{$m$} د~\LR{$n$} مضرب دے)؛
\item \LR{$n$} اول عدد دے (يعنې له~\LR{$1$} او~\LR{$n$} پرته بل هېڅ عدد~\LR{$n$} په پوره
  ډول نۀ وېشي)۔
\end{enumerate}
\end{prob}
\olpseganchor{OLP-0171-B011}{end}

\olpseganchor{OLP-0171-B012}{start}
\begin{prob}
فرض کړئ فارمول~\LR{$!A(\Obj v_1, \Obj v_2)$} په
جوړښت~\LR{$\Struct M$} کښې اړيکه~\LR{$R \subseteq \Domain M^2$}
بيانوي۔ داسې فارمولونه ومومئ چې لاندې اړيکې بيان کړي:
\begin{enumerate}
\item د~\LR{$R$} معکوسه اړيکه~\LR{$R^{-1}$}؛
\item نسبي حاصل‌ضرب~\LR{$R \mid R$}؛
\end{enumerate}
ايا د~\LR{$R$} د انتقالي بندښت~\LR{$R^+$} د بيانولو کومه طريقه موندلے شئ؟
\end{prob}
\olpseganchor{OLP-0171-B012}{end}

\olpseganchor{OLP-0171-B013}{start}
\begin{prob}
فرض کړئ~\LR{$\Lang{L}$} هغه ژبه ده چې يوازې دوه‌ځاييز اړوند سمبول~\LR{$<$}
لري (بل هېڅ ثابت سمبولونه، تابعې يا پريديکاتونه نۀ لري---
البته له~\LR{$\eq$} پرته)۔ فرض کړئ~\LR{$\Struct{N}$} داسې جوړښت دے چې
\LR{$\Domain{N} = \Nat$}، او
\LR{$\Assign{<}{N} = \Setabs{\tuple{n,m}}{n < m}$}۔ لاندې ثابت کړئ:
\begin{enumerate}
\item \LR{$\{ 0 \}$} په~\LR{$\Struct{N}$} کښې تعريفېدے شي؛
\item \LR{$\{ 1 \}$} په~\LR{$\Struct{N}$} کښې تعريفېدے شي؛
\item \LR{$\{ 2 \}$} په~\LR{$\Struct{N}$} کښې تعريفېدے شي؛
\item د هر~\LR{$n \in \Nat$} دپاره سټ~\LR{$\{ n \}$} په~\LR{$\Struct{N}$} کښې
  تعريفېدے شي؛
\item د~\LR{$\Domain{N}$} هر متناهي فرعي سټ په~\LR{$\Struct{N}$} کښې
  تعريفېدے شي؛
\item د~\LR{$\Domain{N}$} هر د متناهي متمم لرونکے فرعي سټ په~\LR{$\Struct{N}$} کښې
  تعريفېدے شي (چېرته چې~\LR{$X \subseteq \Nat$} هله او يوازې هله
  د متناهي متمم لرونکے دے چې~\LR{$\Nat \setminus X$} متناهي وي)۔
\end{enumerate}
\end{prob}
\olpseganchor{OLP-0171-B013}{end}

\olpseganchor{OLP-0172-B004}{start}
\olfileid{fol}{mat}{set}
\olpseganchor{OLP-0172-B004}{end}

\olpseganchor{OLP-0172-B005}{start}
\olsection{د سټونو تيوري}
\olpseganchor{OLP-0172-B005}{end}

\olpseganchor{OLP-0172-B006}{start}
د رياضۍ کابو ټولې برخې د سټونو په تيورۍ کښې جوړېدے شي۔ په دې
تيورۍ کښې د رياضۍ جوړول څو کارونه غواړي۔ لومړے، د~\LR{$\in$} اړيکې دپاره
د بديهي اصولو يو سټ غواړي۔ بېلابېل بديهي نظامونه جوړ شوي دي چې کله
کله د~\LR{$\in$} په باب يو بل سره ټکر لرونکي خاصيتونه لري۔ د~\LR{$\Log{ZFC}$}
په نوم بديهي نظام، يعنې د انتخاب له بديهي اصل سره د
زرملو--فرېنکل سټ تيوري، ځانګړے مقام لري: دا تر بل هر بديهي نظام ډېر
کارول کېږي او څېړل کېږي، ځکه بديهي اصول ئې د
کابو ټولو هغو خبرو د ثبوت دپاره بس کېږي چې رياضيپوهان ئې ثابتول
غواړي۔ خو د دې له څرګندولو وړاندې بايد لومړے روښانه کړو چې د
رياضيپوهانو مطلوبې خبرې اصلًا څنګه \emph{بيانولے} شو۔ د پيل دپاره،
ژبه هېڅ ثابت سمبولونه يا تابعې نۀ لري؛ نو په لومړي نظر دا
روښانه نۀ ده چې د ځانګړو سټونو (لکه~\LR{$\emptyset$} يا~\LR{$\Nat$})، پر
سټونو د عملونو (لکه~\LR{$X \cup Y$} او~\LR{$\Pow{X}$})، يا لا د هغو جوړښتونو
په باب څنګه خبرې وکړو چې له سټونو پرته نور څيزونه هم رانغاړي، لکه
اړيکې او تابعې۔
\olpseganchor{OLP-0172-B006}{end}

\olpseganchor{OLP-0172-B007}{start}
په پيل کښې، ``عنصر دے'' يوازېنۍ مهمه اړيکه نۀ ده؛ ``فرعي سټ دے''
کابو همدومره مهمه ښکاري۔ خو ``فرعي سټ دے'' د ``عنصر دے'' له مخې
\emph{تعريفولے} شو۔ د دې دپاره بايد د سټ تيورۍ په ژبه کښې داسې
فارمول~\LR{$!A(x, y)$} ومومو چې د سټونو يوه جوړه~\LR{$\tuple{X, Y}$}
هله او يوازې هله پوره کړي چې~\LR{$X \subseteq Y$} وي۔ خو~\LR{$X$} هله او
يوازې هله د~\LR{$Y$} فرعي سټ دے چې د~\LR{$X$} ټول غړي د~\LR{$Y$} هم
غړي وي۔ نو~\LR{$\subseteq$} په لاندې فارمول تعريفولے شو:
\[
\lforall[z][(z \in x \lif z \in y)]
\]
اوس چې هر کله په يوه فارمول کښې اړيکه~\LR{$\subseteq$} کارول غواړو، د
هغې پر ځاے همدا فارمول کارولے شو (په دې شرط چې~\LR{$x$} او~\LR{$y$} په مناسب
ډول بدل او تړلے متغير~\LR{$z$} د اړتيا پر وخت بيا ونومول شي)۔ مثلاً د
سټونو امتداديتوب وايي چې که هر دوه سټونه~\LR{$x$} او~\LR{$y$} په دوه اړخيزه
توګه په يو بل کښې شامل وي، نو~\LR{$x$} او~\LR{$y$} بايد هماغه يو سټ وي۔ دا په
\LR{$\lforall[x][\lforall[y][((x \subseteq y \land y
    \subseteq x) \lif x = y)]]$} بيانېدے شي؛ يا که~\LR{$\subseteq$} په
پورته تعريف بدل کړو، په لاندې ډول:
\[
\lforall[x][\lforall[y][((\lforall[z][(z \in x \lif z \in y)] \land
    \lforall[z][(z \in y \lif z \in x)]) \lif x = y)]].
\]
دا په رښتيا د~\LR{$\Log{ZFC}$} يو بديهي اصل، يعنې ``د امتداديتوب بديهي
اصل'' دے۔
\olpseganchor{OLP-0172-B007}{end}

\olpseganchor{OLP-0172-B008}{start}
د~\LR{$\emptyset$} دپاره هېڅ ثابت سمبول نشته، خو ``\LR{$x$} تش دے'' په
\LR{$\lnot \lexists[y][y \in x]$} بيانولے شو۔ بيا ``\LR{$\emptyset$} شته''
په تړلے فارمول~\LR{$\lexists[x][\lnot \lexists[y][y \in x]]$} اوړي۔
دا د~\LR{$\Log{ZFC}$} يو بل بديهي اصل دے۔ (پام وکړئ چې د امتداديتوب
بديهي اصل دا لازموي چې يوازې يو تش سټ وي۔) هر کله چې د سټ تيورۍ په
ژبه کښې د~\LR{$\emptyset$} په باب خبرې کول غواړو، داسې به ليکو: ``داسې
يو سټ شته چې تش دے او~\dots''۔ د بېلګې په توګه، دا حقيقت چې
\LR{$\emptyset$} د هر سټ فرعي سټ دے، په لاندې ډول بيانولے شو:
\[
\lexists[x][(\lnot\lexists[y][y \in x] \land \lforall[z][x \subseteq
    z])]
\]
چېرته چې، البته، \LR{$x \subseteq z$} به بيا په خپل تعريف بدلېږي۔
\olpseganchor{OLP-0172-B008}{end}

\olpseganchor{OLP-0172-B009}{start}
د سټونو د عملونو، لکه~\LR{$X \cup Y$} او~\LR{$\Pow{X}$}، په باب د خبرو دپاره
ورته چل کارول پکار دي۔ د سټ تيورۍ په ژبه کښې د تابعې سمبولونه نشته،
خو تابعوي اړيکې~\LR{$X \cup Y = Z$} او~\LR{$\Pow{X} = Y$} په لاندې ډول
بيانولے شو:
\begin{align*}
& \lforall[u][((u \in x \lor u \in y) \liff u \in z)]\\
& \lforall[u][(u \subseteq x \liff u \in y )]
\end{align*}
ځکه د~\LR{$X \cup Y$} غړي کټ مټ هغه سټونه دي چې يا د~\LR{$X$}
غړي وي يا د~\LR{$Y$} (يا د دواړو) غړي وي، او د~\LR{$\Pow{X}$}
غړي کټ مټ
د~\LR{$X$} فرعي سټونه دي۔ خو دا اجازه نۀ راکوي چې~\LR{$x \cup y$} يا~\LR{$\Pow{x}$}
د ترمونو په توګه وکاروو: يوازې ټول هغه فارمولونه کارولے شو چې
اړيکې~\LR{$X \cup Y = Z$} او~\LR{$\Pow{X} = Y$} تعريفوي۔ په اصل کښې لا دا هم
نۀ پوهېږو چې دغه اړيکې کله هم پوره کېږي، يعنې نۀ پوهېږو چې اتحادونه
او د فرعي سټونو سټونه تل شته۔ مثلاً تړلے فارمول
\LR{$\lforall[x][\lexists[y][\Pow{x} = y]]$} د~\LR{$\Log{ZFC}$} يو بل بديهي
اصل دے (د فرعي سټونو د سټ بديهي اصل)۔
\olpseganchor{OLP-0172-B009}{end}

\olpseganchor{OLP-0172-B010}{start}
اوس د مرتبو جوړو يا تابعو په باب خبرې څنګه کوو؟ دلته بايد څرګنده کړو
چې مرتبې جوړې او تابعې څنګه د سټونو ځانګړي ډولونه ګڼلے شو۔ د مرتبې
جوړې~\LR{$\tuple{x, y}$} د تعريف يوه طريقه سټ~\LR{$\{\{x\}, \{x, y\}\}$} ده۔
خو لکه مخکښې، داسې تابع نۀ شو معرفي کولے چې دغه سټ ونوموي؛
يوازې اړيکه~\LR{$\tuple{x, y} = z$}، يعنې
\LR{$\{\{x\}, \{x, y\}\} = z$}، تعريفولے شو:
\[
\lforall[u][(u \in z \liff (\lforall[v][(v \in u \liff v = x)] \lor
  \lforall[v][(v \in u \liff (v = x \lor v = y))]))]
\]
دا وايي چې د~\LR{$z$} غړي~\LR{$u$} کټ مټ هغه سټونه دي چې يا يوازې
\LR{$x$} ئې غړے وي، يا يوازې~\LR{$x$} او~\LR{$y$} ئې غړي وي (په
بل عبارت، هغه سټونه چې يا له~\LR{$\{x\}$} سره يو شان وي يا
له~\LR{$\{x, y\}$} سره)۔ چې دا ولرو، نورې خبرې هم کولے شو؛ مثلاً
\LR{$X \times Y = Z$}:
\[
\lforall[z][(z \in Z \liff \lexists[x][\lexists[y][(x \in
        X \land y \in Y \land \tuple{x, y} = z)]])]
\]
\olpseganchor{OLP-0172-B010}{end}

\olpseganchor{OLP-0172-B011}{start}
تابعه~\LR{$f \colon X \to Y$} د اړيکې~\LR{$f(x) = y$} په توګه، يعنې د جوړو
د سټ~\LR{$\Setabs{\tuple{x,y}}{f(x) = y}$} په توګه، ګڼلے شو۔ بيا ويلے شو
چې سټ~\LR{$f$} له~\LR{$X$} څخه~\LR{$Y$} ته تابعه ده، که (الف) يوه اړيکه
\LR{$\subseteq X \times Y$} وي؛ (ب) ټوله وي، يعنې د هر~\LR{$x \in X$} دپاره
داسې~\LR{$y \in Y$} وي چې~\LR{$\tuple{x, y} \in f$}؛ او (ج) تابعوي وي، يعنې
هر کله چې~\LR{$\tuple{x, y}, \tuple{x, y'} \in f$} وي، نو~\LR{$y = y'$} وي
(ځکه د تابعې قيمتونه بايد يوازيني وي)۔ نو ``\LR{$f$} له~\LR{$X$} څخه~\LR{$Y$} ته
تابعه ده'' داسې ليکل کېدے شي:
\begin{align*}
 \lforall[u][(u \in f \lif {} & \lexists[x][\lexists[y][(x \in X \land y \in
      Y \land \tuple{x, y} = u)]])] \land {}\\
 \lforall[x][(x \in X \lif {} &
      (\lexists[y][(y \in Y \land \mathrm{maps}(f, x, y))] \land {}\\
& (\lforall[y][\lforall[y'][((\mathrm{maps}(f, x, y) \land
    \mathrm{maps}(f, x, y')) \lif y = y')]])))]
\end{align*}
چېرته چې~\LR{$\mathrm{maps}(f, x, y)$} د
\LR{$\lexists[v][(v \in f \land \tuple{x, y} = v)]$} لنډون دے (دا
فارمول ``\LR{$f(x) = y$}'' بيانوي)۔
\textbf{د ژباړې د نسخې سمون (OLFOL-026):} سرچينې په دواړو کلي
شرطونو کښې ساحوي قوسونه د شرط له پايلې وړاندې تړلي وو؛ دلته دواړه
پايلې د خپلو کميت ټاکونکو په ساحه کښې راوستل شوې دي۔
\olpseganchor{OLP-0172-B011}{end}

\olpseganchor{OLP-0172-B012}{start}
اوس دا هم ستونزمنه نۀ ده چې مثلاً بيان کړو چې~\LR{$f\colon X \to Y$}
يو پر يو ده:
\begin{multline*}
 f \colon X \to Y \land \lforall[x][\lforall[x'][((x \in X \land x' \in
     X \land {} \\
 \lexists[y][(\mathrm{maps}(f, x, y) \land \mathrm{maps}(f,
        x', y))]) \lif x = x')]]
\end{multline*}
تابعه~\LR{$f\colon X \to Y$} هله او يوازې هله يو پر يو ده چې هر
کله~\LR{$f$}، \LR{$x, x' \in X$} يوه~\LR{$y$} ته ولېږي، نو~\LR{$x = x'$} وي۔ که دا
فارمول په~\LR{$\mathrm{inj}(f, X, Y)$} لنډ کړو، لا له مخکښې په داسې مقام
کښې يو چې د سټ تيورۍ په ژبه د کانتور د قضيې غوندې نۀ معمولي خبره
بيان کړو: له~\LR{$\Pow{X}$} څخه~\LR{$X$} ته هېڅ يو پر يو تابعه نشته:
\[
\lforall[X][\lforall[Y][(\Pow{X} = Y \lif
    \lnot\lexists[f][\mathrm{inj}(f, Y, X)])]]
\]
\textbf{د ژباړې د نسخې سمون (OLFOL-027):} سرچينې دواړه کلي سوري د
مشترک قيمت د شتون له شرط وړاندې تړلي وو؛ دلته د هماغو سورو تړونکي
قوسونه د بشپړ شرطي فارمول تر پای وروسته راغلي دي۔
\olpseganchor{OLP-0172-B012}{end}

\olpseganchor{OLP-0172-B013}{start}
کېدے شي څوک فکر وکړي چې سټ تيوري يو بل داسې بديهي اصل غواړي چې د هر
تعريفوونکي خاصيت دپاره د يوه سټ شتون تضمين کړي۔ که~\LR{$!A(x)$} د سټ
تيورۍ داسې فارمول وي چې متغير~\LR{$x$} پکښې ازاد وي، لاندې تړلے فارمول
په پام کښې نيولے شو:
\[
\lexists[y][\lforall[x][(x \in y \liff !A(x))]].
\]
دا تړلے فارمول وايي چې داسې سټ~\LR{$y$} شته چې غړي ئې ټول او
يوازې هغه~\LR{$x$} دي چې~\LR{$!A(x)$} پوره کوي۔ دا شېما ``د ادراک اصل'' بلل
کېږي۔ ډېره ګټوره ښکاري؛ خو له بده مرغه ناسازګاره ده۔
\LR{$!A(x) \ident \lnot x \in x$} واخلئ؛ بيا د ادراک اصل وايي:
\[
\lexists[y][\lforall[x][(x \in y \liff x \notin x)]],
\]
يعنې د ټولو هغو سټونو د سټ شتون وايي چې د خپل ځان غړي نۀ
وي۔ داسې سټ نۀ شي موجودېدے---دا د راسل تناقض دے۔ په اصل کښې
\LR{$\Log{ZFC}$} د دې اصل يوه محدوده---او سازګاره---بڼه لري، يعنې د
بېلتون اصل:
\[
\lforall[z][\lexists[y][\lforall[x][(x \in y \liff (x \in z \land
      !A(x))]]].
\]
\olpseganchor{OLP-0172-B013}{end}

\olpseganchor{OLP-0172-B014}{start}
\begin{prob}
داسې اشتقاق ورکړئ چې لاندې وښيي، او په دې توګه ثابته کړئ
چې د ادراک اصل ناسازګار دے:
\[
\lexists[y][\lforall[x][(x \in y \liff x \notin x)]] \Proves \lfalse.
\]
ښايي لومړے د لاندې ثبوت ګټور وي:
\LR{$(A \lif \lnot A) \land (\lnot A \lif A) \Proves \lfalse$}۔
\end{prob}
\olpseganchor{OLP-0172-B014}{end}

\olpseganchor{OLP-0173-B004}{start}
\olfileid{fol}{mat}{siz}
\olpseganchor{OLP-0173-B004}{end}

\olpseganchor{OLP-0173-B005}{start}
\olsection{د جوړښتونه د کچې بيانول}
\olpseganchor{OLP-0173-B005}{end}

\olpseganchor{OLP-0173-B006}{start}
\begin{explain}
د جوړښتونو ځينې خاصيتونه د ژبې د غيرمنطقي سمبولونو له کارولو پرته
هم بيانولے شو۔ مثلاً داسې تړلي فارمولونه شته چې په جوړښت
کښې هله او يوازې هله رښتونې وي چې د تعريف ساحه، يعنې د جوړښت
ساحه، لږ تر لږه، زيات نه زيات، يا کټ مټ يو ټاکلے شمېر~\LR{$n$}
غړي ولري۔
\end{explain}
\olpseganchor{OLP-0173-B006}{end}

\olpseganchor{OLP-0173-B007}{start}
\begin{prop}
تړلے فارمول
\begin{multline*}
  !A_{\ge n} \ident \lexists[x_1][\lexists[x_2][\dots\lexists[x_n][{}\\
  \begin{aligned}
  (\eq/[x_1][x_2] \land {}
  \eq/[x_1][x_3] \land \eq/[x_1][x_4] \land \dots \land \eq/[x_1][x_n] \land {}\\
\eq/[x_2][x_3] \land \eq/[x_2][x_4] \land \dots \land {} \eq/[x_2][x_n] \land {} \\
\vdots\\
\eq/[x_{n-1}][x_n])
\end{aligned}]]]
\end{multline*}
په جوړښت~\LR{$\Struct M$} کښې هله او يوازې هله رښتونې ده چې
\LR{$\Domain M$} لږ تر لږه~\LR{$n$} غړي ولري۔ په پايله کښې،
\LR{$\Sat{M}{\lnot !A_{\ge n+1}}$} هله او يوازې هله چې~\LR{$\Domain M$}
زيات نه زيات~\LR{$n$} غړي ولري۔
\textbf{د ژباړې د نسخې سمون (OLFOL-028):} سرچينې د وجودي سورو
ټول ساحوي قوسونه د بېلوالي له شرطونو وړاندې تړلي وو او شرطونه ئې د
سورو له ساحې بهر پرېښي وو؛ دلته قوسونه د بشپړ شرط وروسته تړل شوي دي۔
\olpseganchor{OLP-0173-B007}{end}

\olpseganchor{OLP-0173-B008}{start}
\end{prop}
\olpseganchor{OLP-0173-B008}{end}

\olpseganchor{OLP-0173-B009}{start}
\begin{prop}
تړلے فارمول
\begin{multline*}
!A_{= n} \ident \lexists[x_1][\lexists[x_2][\dots\lexists[x_n][{} \\
\begin{aligned}
  (\eq/[x_1][x_2] \land {}
  \eq/[x_1][x_3] \land \eq/[x_1][x_4] \land \dots \land \eq/[x_1][x_n] \land {}\\
\eq/[x_2][x_3] \land \eq/[x_2][x_4] \land \dots \land {} \eq/[x_2][x_n] \land {} \\
\vdots\\
\eq/[x_{n-1}][x_n] \land {} \\
\lforall[y][(\eq[y][x_1] \lor \dots \lor \eq[y][x_n])])
\end{aligned}]]]
\end{multline*}
په جوړښت~\LR{$\Struct M$} کښې هله او يوازې هله رښتونې ده چې
\LR{$\Domain M$} کټ مټ~\LR{$n$} غړي ولري۔
\textbf{د ژباړې د نسخې سمون (OLFOL-029):} سرچينې دلته هم وجودي
سوري له خپل شرط وړاندې تړلي وو؛ د وروستي کلي سور تړونکي قوسونه هم
په سرچينه کښې په ناسم ترتيب او يو زيات تړونکي قوس سره راغلي وو۔
دلته ټول شرط د وجودي سورو په ساحه کښې او د کلي سور قوسونه جوړې دي۔
\end{prop}
\olpseganchor{OLP-0173-B009}{end}

\olpseganchor{OLP-0173-B010}{start}
\begin{prop}
يو جوړښت هله او يوازې هله نامتناهي دے چې د لاندې سټ مدل وي:
\[
\{!A_{\ge 1}, !A_{\ge 2}, !A_{\ge 3}, \dots \}.
\]
\end{prop}
\olpseganchor{OLP-0173-B010}{end}

\olpseganchor{OLP-0173-B011}{start}
هېڅ يوه خالصه منطقي جمله نشته چې په~\LR{$\Struct M$} کښې هله او يوازې
هله رښتونې وي چې~\LR{$\Domain M$} نامتناهي وي۔ خو له غيرمنطقي
پريديکاتونه سره داسې تړلي فارمولونه ورکولے شو چې يوازې نامتناهي
مدلونه لري (که څه هم هر نامتناهي جوړښت ئې مدل نۀ دے)۔ د
متناهي جوړښت خاصيت او د د شمېر نۀ وړ جوړښت خاصيت حتا د
تړلي فارمولونه په نامتناهي سټ هم نۀ شي بيانېدے۔ دا حقيقتونه د
متناهي شاهد او لوېنهايم--سکولم له قضيو څخه راځي۔
\olpseganchor{OLP-0173-B011}{end}

\olpseganchor{OLP-0174-B004}{start}
\olchapter{fol}{byd}{د لومړۍ درجې له منطق څخه وړاندې}
\olpseganchor{OLP-0174-B004}{end}

\olpseganchor{OLP-0174-B005}{start}
\begin{editorial}
  دا څپرکى، چې د جېرمي اېوېګېډ د منطق له يادښتونو څخه برابر شوے، د
  شته نورو منطقي نظامونو يوازې تر ټولو لنډه کتنه وړاندې کوي۔ موخه ئې دا
  ده چې په هغه کورس کښې چې دا موضوعات نۀ رانغاړي، د نورو مطالعو موضوعات
  وروښيي۔ هيله ده چې دلته ياده شوې هره موضوع به بالاخره په پرانيستې
  منطق پروژه کښې د يوې مستقلې برخې په کچه بشپړ بحث ومومي۔
\end{editorial}
\olpseganchor{OLP-0174-B005}{end}

\olpseganchor{OLP-0175-B004}{start}
\olfileid{fol}{byd}{int}
\olpseganchor{OLP-0175-B004}{end}

\olpseganchor{OLP-0175-B005}{start}
\olsection{ټوليزه کتنه}
\olpseganchor{OLP-0175-B005}{end}

\olpseganchor{OLP-0175-B006}{start}
د لومړۍ درجې منطق د پام وړ يوازينے منطقي نظام نۀ دے: د دې منطق ډېر
غځېدلي او بدل ډولونه شته۔ يو منطق په عام ډول د يوې ژبې صوري مشخصات،
اکثره، خو تل نۀ، د استنتاج يو نظام، او اکثره، خو تل نۀ، يوه ټاکل شوې
معنٰی پوهنه لري۔ خو د دې ويي تخنيکي کارونه يوه څرګنده پوښتنه راپورته
کوي: هغه منطقونه چې د لومړۍ درجې منطق نۀ دي، په وجداني يا فلسفي معنا
د ``منطق'' له ويي سره څه اړيکه لري؟ لاندې ټول بيان شوي نظامونه د يوه
يا بل ډول استدلال د مدل کولو دپاره جوړ شوي؛ آيا ويلے شو چې څه شي دوى
منطقي کوي؟
\olpseganchor{OLP-0175-B006}{end}

\olpseganchor{OLP-0175-B007}{start}
اسانه ځوابونه نۀ تر سترګو کېږي۔ د ``منطق'' ويے په بېلابېلو لارو او
بېلابېلو چاپېريالونو کښې کارول کېږي، او دا مفهوم، لکه د ``رښتيا''
مفهوم، له ډېرو فلسفي دريځونو څېړل شوے دے۔ د بېلګې په توګه، څوک ښايي
د منطقي استدلال موخه دا وګڼي چې معلومه کړي کومې ويناوې په لازمي ډول
رښتونې، له تجربې وړاندې رښتونې، د غيرمنطقي اصطلاحاتو له تفسير څخه
خپلواکې رښتونې، د خپلې بڼې له مخې رښتونې، يا د ژبني تړون له مخې
رښتونې دي؛ او له دغو مفکورو څخه هره يوه ډېر وضاحت غواړي۔ که پام يوازې
په رياضياتو کښې کارېدونکي منطق ته هم محدود کړو، د هغه د ساحې په هکله لږ
توافق شته۔ د بېلګې په توګه، رسل او وايټ‌هېډ په
\textit{پرېنسيپيا مېتېماتيکا} کښې هڅه وکړه چې رياضيات د منطق پر بنسټ،
د فرېګې له پيل کړې {\em منطقپالې} دود سره سم، جوړ کړي۔ د دوى منطقي
نظام د لوړو ډولونو د منطق يوه بڼه وه چې لاندې بيان شوي نظام ته ورته
وه۔ په پاى کښې اړ شول داسې بديهي اصول ورزيات کړي چې د ډېرو معيارونو
له مخې سوچه منطقي نۀ ښکاري (په ځانګړي ډول د لامتناهي والي بديهي اصل
او د راکمېدنې بديهي اصل)؛ خو بيا هم څوک دا دريځ لرلے شي چې د لوړې
درجې استدلال ځينې بڼې بايد منطقي ومنل شي۔ د دې پر خلاف، کواين، چې
هستپوهنه ئې ``قضيې'' د بحث د روا څيزونو په توګه نۀ مني، استدلال کوي
چې د دويمې او لوړې درجې منطق په اصل کښې د سټ تيوري ده چې د پسه پوست
ئې اغوستے؛ په نورو ټکو، هغه نظامونه چې پر پريديکاتونو کميت ټاکنه لري
سوچه منطقي نۀ دي۔
\olpseganchor{OLP-0175-B007}{end}

\olpseganchor{OLP-0175-B008}{start}
اوس دپاره غوره ده چې دا ډول فلسفي مسئلې بل وخت ته پرېږدو او لاندې
نظامونه يوازې د بېلابېلو ډولونو استدلال، که منطقي وي که بل څه، صوري
ايډيالې بڼې وګڼو۔
\olpseganchor{OLP-0175-B008}{end}

\olpseganchor{OLP-0176-B004}{start}
\olfileid{fol}{byd}{msl}
\olpseganchor{OLP-0176-B004}{end}

\olpseganchor{OLP-0176-B005}{start}
\olsection{د ډېرو ډولونو منطق}
\olpseganchor{OLP-0176-B005}{end}

\olpseganchor{OLP-0176-B006}{start}
د لومړۍ درجې په منطق کښې متغيرونه او کميت ټاکونکي پر يوه
د تعريف ساحه ګرځي۔ خو ډېر ځله د څو (يو له بله بېلو) د تعريف ساحې درلودل
ګټور وي: د بېلګې په توګه، ښايي د عددونو د تعريف ساحه، د هندسي څيزونو
د تعريف ساحه، له عددونو څخه عددونو ته د تابعو د تعريف ساحه، د ابېلي
ګروپونو د تعريف ساحه او داسې نور وغواړئ۔
\olpseganchor{OLP-0176-B006}{end}

\olpseganchor{OLP-0176-B007}{start}
د ډېرو ډولونو منطق دا ډول چوکاټ برابروي۔ پيل د ``ډولونو'' له يوه لړ
څخه کېږي---د يوه څيز ``ډول'' هغه ``د تعريف ساحه'' ښيي چې څيز بايد پکې
اوسېږي۔ بيا د هر ډول دپاره متغيرونه او کميت ټاکونکي، او (په عام
ډول) د هر ډول دپاره يو عينيت لرو۔ تابعې او اړيکې هم د هغو
څيزونو د ډولونو له مخې ``ټايپ'' کېږي چې د ارګومېنټونو په توګه ئې
اخيستلے شي۔ له دې پرته د لومړۍ درجې د منطق عادي قاعدې ساتل کېږي، او
د کميت ټاکونکو د قاعدو بڼې د هر ډول دپاره تکرارېږي۔
\olpseganchor{OLP-0176-B007}{end}

\olpseganchor{OLP-0176-B008}{start}
د بېلګې په توګه، د نړيوالو اړيکو د مطالعې دپاره ښايي داسې ژبه وټاکو
چې دوه ډوله څيزونه، فرانسوي اتباع او جرمن اتباع، ولري۔ ښايي د لومړي
ډول د څيزونو دپاره يوه يوځاييزه اړيکه، ``شراب څښي''، د دويم ډول د
څيزونو دپاره بله يوځاييزه اړيکه، ``ورسټ خوري''، او يوه دوه ځاييزه
اړيکه، ``څو ملتيزه واده شوې جوړه جوړوي''، ولرو۔ وروستۍ اړيکه دوه
ارګومېنټونه اخلي: لومړے د لومړي ډول او دويم د دويم ډول وي۔ که \LR{$a$}،
\LR{$b$}، \LR{$c$} متغيرونه پر فرانسوي اتباعو او \LR{$x$}، \LR{$y$}، \LR{$z$} متغيرونه پر
جرمن اتباعو وګرځوو، نو
\[
\lforall[a][\lforall[x][(\Atom{\Obj{MarriedTo}}{a,x} \lif
(\Atom{\Obj{DrinksWine}}{a} \lor \lnot \Atom{\Obj{EatsWurst}}{x})]]
\]
دا وايي چې که کوم فرانسوے له يوه جرمن سره وادۀ شوے وي، نو يا هغه
فرانسوے شراب څښي يا هغه جرمن ورسټ نۀ خوري۔
\textbf{د ژباړې د نسخې سمون (OLFOL-030):} سرچينې د دويم عمومي کميت
ټاکونکي د متغير پرانيستونکې نښه غورځولې وه؛ دلته هغه نښه بېرته
ايښوول شوې ده۔
\olpseganchor{OLP-0176-B008}{end}

\olpseganchor{OLP-0176-B009}{start}
د ډېرو ډولونو منطق په طبيعي ډول د لومړۍ درجې په منطق کښې ځايېدلے شي:
د ډېرو ډولونو د د تعريف ساحې ټول څيزونه په يوه د لومړۍ درجې د تعريف ساحه
کښې يوځاے کوو، د ډولونو د پېژندلو دپاره يوځاييز پريديکاتونه کاروو
او کميت ټاکونکي نسبي کوو۔ د بېلګې په توګه، د پاسنۍ بېلګې اړونده د
لومړۍ درجې ژبه به د نورو بيان شويو اړيکو تر څنګ ``\LR{$\Obj{German}$}''
او ``\LR{$\Obj{French}$}'' يوځاييز پريديکاتونه ولري، او د ډول شرطونه
به ترې لرې وي۔ يو ډول‌لرونکے کميت ټاکونکے \LR{$\lforall[x][!A]$}، چې \LR{$x$}
د جرمن ډول متغير وي، داسې ژباړل کېږي:
\[
\lforall[x][(\Atom{\Obj{German}}{x} \lif !A)].
\]
بايد داسې بديهي اصول هم ورزيات کړو چې ډولونه بېل تضمين کړي---لکه
\LR{$\lforall[x][\lnot (\Atom{\Obj{German}}{x} \land
  \Atom{\Obj{French}}{x})]$}---او هم داسې بديهي اصول چې تضمين کړي
``شراب څښي'' يوازې د \LR{$\Atom{\Obj{French}}{x}$} پريديکات پوره کوونکو
څيزونو دپاره صدق کوي، او داسې نور۔ له دې تړونونو او بديهي اصولو سره
دا ښودل ګران نۀ دي چې د ډېرو ډولونو تړلي فارمولونه د لومړۍ درجې
تړلي فارمولونه ته، او د ډېرو ډولونو اشتقاقونه د لومړۍ درجې
اشتقاقونه ته ژباړل کېږي۔ همدارنګه، د ډېرو ډولونو جوړښتونه
د لومړۍ درجې اړوندو جوړښتونه ته او برعکس ``ژباړل'' کېږي؛ نو د
ډېرو ډولونو د منطق دپاره هم د بشپړتيا قضيه لرو۔
\olpseganchor{OLP-0176-B009}{end}

\olpseganchor{OLP-0177-B004}{start}
\olfileid{fol}{byd}{sol}
\olpseganchor{OLP-0177-B004}{end}

\olpseganchor{OLP-0177-B005}{start}
\olsection{د دويمې درجې منطق}
\olpseganchor{OLP-0177-B005}{end}

\olpseganchor{OLP-0177-B006}{start}
د دويمې درجې منطق ژبه يوازې د افرادو پر د تعريف ساحه کميت ټاکنې ته
اجازه نۀ ورکوي، بلکې د هغه د تعريف ساحه پر اړيکو هم کميت ټاکل پرې کېدے
شي۔ د لومړۍ درجې يوې ژبې~\LR{$\Lang{L}$} ته د هر~\LR{$k$} دپاره داسې
متغيرونه~\LR{$R$} ورزياتېږي چې پر \LR{$k$}-ځاييزو اړيکو ګرځي، او پر دغو
متغيرونو کميت ټاکنه هم روا کېږي۔ که \LR{$R$} د يوې \LR{$k$}-ځاييزې اړيکې
متغير وي او \LR{$t_1$}، \dots،~\LR{$t_k$} عادي (د لومړۍ درجې) ترمونه
وي، نو \LR{$\Atom{R}{t_1,\dots,t_k}$} يو اټومي فارمول دے۔ له دې پرته
د فارمولونه سټ هماغسې تعريفېږي لکه د لومړۍ درجې په منطق کښې، خو د
دويمې درجې کميت ټاکنې اضافي شرطونه لري۔ پام وکړئ چې عينيت
يوازې د لومړۍ درجې د ترمونو دپاره لرو: که \LR{$R$} او \LR{$S$} د يوه شان
ځايښت~\LR{$k$} اړيکيز متغيرونه وي، \LR{$\eq[R][S]$} د لاندې عبارت د لنډون
په توګه تعريفولے شو:
\[
\lforall[x_1 \dots][\lforall[x_k][(\Atom{R}{x_1, \dots, x_k} \liff
  \Atom{S}{x_1, \dots, x_k})]].
\]
\olpseganchor{OLP-0177-B006}{end}

\olpseganchor{OLP-0177-B007}{start}
د دويمې درجې منطق قاعدې يوازې د کميت ټاکونکو قاعدې د دويمې درجې
نويو متغيرونو ته غځوي۔ خو دلته بايد لږ پام وشي چې دا متغيرونه د
ژبې~\LR{$\Lang{L}$} له پريديکاتونه، او په عام ډول د~\LR{$\Lang{L}$} له
فارمولونه سره څنګه چلند کوي۔ لږ تر لږه، اړيکيز متغيرونه ترمونه
ګڼل کېږي، نو د لاندې بڼې استنتاجونه لرو:
\[
!A(R) \Proves \lexists[R][!A(R)]
\]
خو که \LR{$\Lang{L}$} د حساب ژبه وي او د ثابتې اړيکې نښه~\LR{$<$} ولري، تمه
کېږي چې لاندې استنتاج هم معتبر وي:
\[
x < y \Proves \lexists[R][\Atom{R}{x,y}]
\]
يا د يوه ورکړل شوي فارمول~\LR{$!A$} دپاره:
\[
!A(x_1, \dots, x_k) \Proves \lexists[R][\Atom{R}{x_1,\dots,x_k}]
\]
په لا عام ډول، ښايي د لاندې بڼې استنتاجونه روا کول وغواړو:
\[
\Subst{!A}{\lambd[\vec x][!B(\vec x)]}{R} \Proves
\lexists[R][!A]
\]
دلته \LR{$\Subst{!A}{\lambd[\vec x][!B(\vec x)]}{R}$} هغه پايله ښيي چې
په~\LR{$!A$} کښې د \LR{$\Atom{R}{t_1,\dots,t_k}$} بڼې هر اټومي فارمول په
\LR{$!B(t_1, \dots, t_k)$} بدل شي۔ دا وروستۍ قاعده د {\em ادراک شېما} له
شتون سره برابره ده، يعنې د لاندې بڼې يو بديهي اصل:
\[
\lexists[R][\lforall[x_1, \dots, x_k][(!A(x_1, \dots, x_k) \liff
\Atom{R}{x_1, \dots, x_k})]],
\]
د دويمې درجې په ژبه کښې د هر فارمول~\LR{$!A$} دپاره يوه داسې بېلګه
چې \LR{$R$} پکې ازاد متغير نۀ وي۔ (تمرین: وښيئ چې که \LR{$R$} ته اجازه ورکړل
شي چې په~\LR{$!A$} کښې راشي، دا شېما ناسازګاره ده!)
\textbf{د ژباړې د نسخې سمون (OLFOL-031):} د بدلولو په تشريح کښې
سرچينې د اړيکې اټومي نښه د څيز په نښه ليکلې وه؛ دلته د اټومي اړيکې
سمه نښه بېرته ايښوول شوې ده۔
\olpseganchor{OLP-0177-B007}{end}

\olpseganchor{OLP-0177-B008}{start}
کله چې منطقپوهان د ``دويمې درجې منطق بديهي اصول'' يادوي، په عام ډول
د دويمې درجې د کميت ټاکونکو پر قاعدو او د ادراک پر شېما د لومړۍ درجې
منطق لږ تر لږه غځونه مرادوي۔ خو د دې بديهي اصولو او قاعدو د کمزورو
فرعي نظامونو مطالعه هم ډېر ځله په زړه پورې وي۔ د بېلګې په توګه، پام
وکړئ چې د ادراک بديهي شېما په خپل بشپړ عموميت کښې
\emph{ناوړاندوينيزه} ده: دا د داسې اړيکې
\LR{$\Atom{R}{x_1, \dots, x_k}$} د شتون دعوه روا کوي چې په داسې
فارمول ``تعريف'' شوې وي چې د دويمې درجې کميت ټاکونکي لري؛ او دا
کميت ټاکونکي د ټولو داسې اړيکو پر سټ ګرځي---هغه سټ چې خپله \LR{$R$} هم پکې
ده!{} د
شلمې پېړۍ د پيل شاوخوا د راسل تناقض ته يو عام غبرګون دا و چې دا ډول
تعريفونه ملامت وګڼل شي او د رياضياتو د بنسټونو په جوړولو کښې ترې ډډه
وشي۔ که په فارمول~\LR{$!A$} کښې د دويمې درجې کميت ټاکونکو کارونه منع
شي، د ادراک يوه {\em وړاندوينيزه} بڼه ترلاسه کېږي چې څه ناڅه کمزورې
ده۔
\olpseganchor{OLP-0177-B008}{end}

\olpseganchor{OLP-0177-B009}{start}
د معنايي نظر له مخې، د دويمې درجې يو جوړښت داسې ګڼلے شو چې
د ژبې دپاره د لومړۍ درجې يو جوړښت او ورسره پر د تعريف ساحه د
هغو اړيکو يو سټ لري چې د دويمې درجې کميت ټاکونکي پرې ګرځي (په لا
دقيق ډول، د هر~\LR{$k$} دپاره د ځايښت~\LR{$k$} اړيکو يو سټ شته)۔ البته که
ادراک په اشتقاق نظام کښې شامل وي، دا اضافي شرط لرو چې د
``دويمې درجې په برخه'' کښې د ادراک د بديهي اصولو د پوره کولو دپاره
کافي اړيکې وي---ګنې د اشتقاق نظام سم نۀ دے!{} د کافي اړيکو د
تضمين يوه اسانه لار دا ده چې د دويمې درجې برخه د لومړۍ درجې پر برخه
له \emph{ټولو} اړيکو جوړه وګڼو۔ داسې جوړښت ته \emph{بشپړ}
ويل کېږي، او په يوه معنا په رښتيا د ژبې ``ټاکل شوے جوړښت''
دے۔ که خپل پام بشپړو جوړښتونه ته محدود کړو، هغه څه ترلاسه کوو
چې د دويمې درجې \emph{بشپړه} معنٰی پوهنه بلل کېږي۔ په دې حالت کښې د
يوه جوړښت مشخصول يوازې د لومړۍ درجې د برخې مشخصولو ته راکمېږي،
ځکه د دويمې درجې د برخې منځپانګه په ضمني ډول ترې معلومېږي۔
\olpseganchor{OLP-0177-B009}{end}

\olpseganchor{OLP-0177-B010}{start}
په لنډيز کښې، د دويمې درجې منطق په خبرو کښې څه ابهام شته۔ د
اشتقاق نظام له مخې ښايي له لاندې دواړو څخه يو مراد وي:
\begin{enumerate}
\item د دويمې درجې يو ``لږ تر لږه'' اشتقاق نظام، له ځينو
  ادراکي بديهي اصولو سره۔
\item د دويمې درجې ``معياري'' اشتقاق نظام، له بشپړ ادراک سره۔
\end{enumerate}
د معنٰی پوهنې له مخې هم ښايي له لاندې دواړو څخه يوه د پام وړ وي:
\begin{enumerate}
\item ``کمزورې'' معنٰی پوهنه، چې جوړښت پکې د لومړۍ درجې يوه
  برخه او د ادراک د بديهي اصولو د پوره کولو دپاره کافي لويه د دويمې
  درجې برخه لري۔
\item د دويمې درجې ``معياري'' معنٰی پوهنه، چې يوازې بشپړ
  جوړښتونه پکې څېړل کېږي۔
\end{enumerate}
کله چې منطقپوهان خپل مراد اشتقاق نظام يا معنٰی پوهنه نۀ
مشخصوي، اکثره د هر لړ دويم توکے مرادوي۔ لکه چې وبه وينو، د دې معنٰی
پوهنې ګټه دا ده چې د ډېرو طبيعي رياضيکي جوړښتونو قطعي توصيفونه راکوي؛
په ورته وخت کښې اشتقاق نظام ډېر پياوړے او د دې معنٰی پوهنې
دپاره سم دے۔ نيمګړتيا ئې دا ده چې اشتقاق نظام د دې معنٰی
پوهنې دپاره \emph{بشپړ نۀ} دے؛ په حقيقت کښې د دويمې درجې د بشپړې
معنٰی پوهنې دپاره \emph{هېڅ} په اغېزمن ډول ورکړل شوے اشتقاق
نظام بشپړ نۀ دے۔ له بلې خوا به ووينو چې اشتقاق نظام د کمزورې
معنٰی پوهنې دپاره \emph{بشپړ} دے؛ نو که يوه جمله د اثبات وړ نۀ وي،
يو \emph{داسې} جوړښت شته، چې لازم نۀ دے بشپړ وي، او جمله پکې
دروغ ده۔
\olpseganchor{OLP-0177-B010}{end}

\olpseganchor{OLP-0177-B011}{start}
د دويمې درجې منطق ژبه ډېره بډايه ده۔ يوځاييزې اړيکې د د تعريف ساحه له
فرعي سټونو سره يو شان ګڼلے شو، نو په ځانګړي ډول پر دغو سټونو کميت
ټاکلے شو؛ د بېلګې په توګه، د طبيعي عددونو دپاره استقرا په يوه بديهي
اصل بيانولے شو:
\[
\lforall[R][((\Atom{R}{\Obj{0}} \land \lforall[x][(\Atom{R}{x} \lif
    \Atom{R}{x'})]) \lif \lforall[x][\Atom{R}{x}])].
\]
که د حساب ژبه د \LR{$\Obj 0, \Obj \prime, +, \times$} او~\LR{$<$} نښې ولري، د
هغوى د چلند د بيانولو دپاره لاندې بديهي اصول ورزياتولے شو:
\begin{enumerate}
\item \LR{$\lforall[x][\lnot x' = \Obj 0]$}
\item \LR{$\lforall[x][\lforall[y][(x' = y' \lif x = y)]]$}
\item \LR{$\lforall[x][(x + \Obj 0) = x]$}
\item \LR{$\lforall[x][\lforall[y][(x + y') = (x + y)']]$}
\item \LR{$\lforall[x][(x \times \Obj 0) = \Obj 0]$}
\item \LR{$\lforall[x][\lforall[y][(x \times y') = ((x \times y) + x)]]$}
\item \LR{$\lforall[x][\lforall[y][(x < y \liff \lexists[z][y = (x + z')])]]$}
\end{enumerate}
دا ښودل ګران نۀ دي چې دغه بديهي اصول د پاسني استقرا له بديهي اصل سره،
په دې شرط چې د دويمې درجې بشپړه معنٰی پوهنه کاروو، د حساب د معياري
جوړښت~\LR{$\Struct{N}$}، د حساب د معياري مدل، قطعي توصيف وړاندې
کوي۔ که کوم جوړښت~\LR{$\Struct{M}$}
دا بديهي اصول رښتوني کوي، له~\LR{$\Nat$} څخه د~\LR{$\Struct{M}$} د تعريف ساحه ته
تابع~\LR{$f$} پر~\LR{$\Nat$} د عادي بازګښت په وسيله داسې تعريف کړئ چې
\LR{$f(0) = \Assign{\Obj 0}{M}$} او \LR{$f(x+1) = \Assign{\prime}{M}(f(x))$} وي۔
پر~\LR{$\Nat$} د عادي استقرا او په~\LR{$\Struct M$} کښې د بديهي اصولو (1)
او~(2) له صدق څخه وينو چې \LR{$f$} يو پر يو ده۔ د دې د ښودلو دپاره
چې \LR{$f$} پر هدف سټ بشپړه ده، \LR{$P$} د~\LR{$\Domain{M}$} د هغو غړو سټ ونيسئ چې
د~\LR{$f$} د قيمتونو په سټ کښې دي۔ ځکه چې \LR{$\Struct M$} بشپړ دے، \LR{$P$} د
دويمې درجې په د تعريف ساحه کښې دے۔ د~\LR{$f$} له جوړونې پوهېږو چې
\LR{$\Assign{\Obj 0}{M}$} په~\LR{$P$} کښې دے او \LR{$P$} د~\LR{$\Assign{\prime}{M}$}
لاندې تړلے دے۔ دا چې د استقرا بديهي اصل په~\LR{$\Struct M$} کښې (په ځانګړي
ډول د~\LR{$P$} دپاره) صدق کوي، تضمينوي چې \LR{$P$} د~\LR{$\Struct M$} د لومړۍ درجې
له ټول د تعريف ساحه سره برابر دے۔ نو \LR{$f$} دوه اړخيزه يو پر يو تابع ده۔ دا
ښودل هم زيات ګران نۀ دي چې \LR{$f$} هومومورفيزم دے؛ پر~\LR{$\Nat$} عادي استقرا
بيا بيا کاروو۔
\textbf{د ژباړې د نسخې سمون (OLFOL-032):} په دويم بديهي اصل کښې
سرچينې د تابعې يوه بېله نښه کارولې وه، که څه هم ژبه او پاتې بحث د
پريـم نښه کاروي؛ دلته دويم اصل هماغې پريـم نښې ته يو شان شوے دے۔
\olpseganchor{OLP-0177-B011}{end}

\olpseganchor{OLP-0177-B012}{start}
د سټ تيورۍ په اصطلاحاتو کښې تابع يوازې د اړيکې يو ځانګړے ډول دے؛ د
بېلګې په توګه، يوځاييزه تابع~\LR{$f$} له داسې دوه ځاييزې اړيکې~\LR{$R$} سره يو
شان ګڼل کېدے شي چې \LR{$\lforall[x][\lexists![y][R(x,y)]]$} پوره کوي۔ نو
پر تابعو هم کميت ټاکلے شو۔ بيا د بشپړې معنٰی پوهنې په کارولو د
لامتناهي جوړښتونه ټولګے د هغو جوړښتونه~\LR{$\Struct M$} د
ټولګي په توګه تعريفولے شو چې د~\LR{$\Struct M$} له د تعريف ساحه څخه د هغه
يو مناسب فرعي سټ ته يو پر يو تابع لري:
\[
\lexists[f][(\lforall[x][\lforall[y][(\eq[f(x)][f(y)] \lif \eq[x][y])]]
  \land \lexists[y][\lforall[x][\eq/[f(x)][y]]])].
\]
بيا د دې جملې نفي د متناهي جوړښتونه ټولګے تعريفوي۔
\olpseganchor{OLP-0177-B012}{end}

\olpseganchor{OLP-0177-B013}{start}
سربېره پر دې، د خطي ترتيب تعريف ته د لاندې شرط په زياتولو د ښو
ترتيبونو ټولګے تعريفولے شو:
\[
\lforall[P][(\lexists[x][\Atom{P}{x}] \lif \lexists[x][(\Atom{P}{x}
    \land \lforall[y][(y < x \lif \lnot \Atom{P}{y})])])].
\]
دا وايي چې هر ناتش سټ تر ټولو کوچنے غړے لري، که ``سټ'' له ``يوځاييزې
اړيکې'' سره يو شان وګڼو۔ د بلې بېلګې دپاره، د ګرافونو تړلتيا داسې
بيانولے شو چې د څوکو هېڅ ناتشه بېلتون په ناتړلو برخو کښې نشته:
\[
\lnot \lexists[A][(\lexists[x][A(x)] \land \lexists[y][\lnot A(y)]
  \land \lforall[w][\lforall[z][((\Atom{A}{w} \land \lnot \Atom{A}{z})
      \lif \lnot \Atom{R}{w,z})]])].
\]
د يوې بلې بېلګې په توګه، په تمرين کښې هڅه کولے شئ د هغو متناهي
جوړښتونه ټولګے تعريف کړئ چې د تعريف ساحه ئې جفته اندازه لري۔ تر
دې هم په زړه پورې دا چې حقيقي عددونه د ناطقو عددونو لرونکي بشپړ مرتب
ميدان په توګه په قطعي ډول توصيفولے شو۔
\olpseganchor{OLP-0177-B013}{end}

\olpseganchor{OLP-0177-B014}{start}
په لنډ ډول، د دويمې درجې منطق د لومړۍ درجې له منطق څخه ډېر زيات
بيانوونکے دے۔ دا ښه خبر و؛ اوس بد خبر ته راځو۔ مخکې مو وويل چې د
دويمې درجې د بشپړې معنٰی پوهنې دپاره هېڅ بشپړ اغېزمن اشتقاق
نظام نشته۔ که ښه ئې وګڼو که بد، د لومړۍ درجې منطق ډېر خاصيتونه دلته
نشته، چې د متناهي شاهد خاصيت او د لوېنهايم--سکولم قضيې هم پکې دي۔
\olpseganchor{OLP-0177-B014}{end}

\olpseganchor{OLP-0177-B015}{start}
له بلې خوا، که څوک د دويمې درجې بشپړه معنٰی پوهنه پرېږدي او کمزورې
بڼه ئې ومني، نو د دويمې درجې لږ تر لږه اشتقاق نظام د دې معنٰی
پوهنې دپاره بشپړ دے۔ په نورو ټکو، که \LR{$\Proves$} داسې ولولو چې ``په لږ
تر لږه نظام کښې ثابتوي'' او \LR{$\Entails$} داسې چې ``په کمزورې معنٰی پوهنه
کښې معنايي استلزام لري''، ښودلے شو چې هر کله \LR{$\Gamma \Entails !A$} وي
نو \LR{$\Gamma \Proves !A$} وي۔ که څوک په اشتقاق نظام کښې ځانګړي
ادراکي بديهي اصول شاملول غواړي، بايد معنٰی پوهنه د دويمې درجې هغو
جوړښتونو ته محدوده کړي چې دا اصول پوره کوي: د بېلګې په توګه، که
\LR{$\Delta$} د ادراکي بديهي اصولو يو سټ وي (کېدے شي ټول وي)، نو که
\LR{$\Gamma \cup \Delta \Entails !A$} وي، \LR{$\Gamma \cup \Delta \Proves !A$}
هم وي۔ په ځانګړي ډول، که \LR{$!A$} د هغو ادراکي بديهي اصولو په کارولو د
اثبات وړ نۀ وي چې موږ ئې څېړو، د~\LR{$\lnot !A$} داسې يو مدل شته چې بيا هم
دغه ادراکي بديهي اصول پوره کوي۔
\olpseganchor{OLP-0177-B015}{end}

\olpseganchor{OLP-0177-B016}{start}
دا چې د بشپړتيا قضيه د کمزورې معنٰی پوهنې دپاره ولې صدق کوي، تر ټولو
اسانه ئې دا ده چې د دويمې درجې منطق په لاندې ډول د ډېرو ډولونو منطق
وګڼو۔ يو ډول د عادي ``لومړۍ درجې'' د تعريف ساحه په توګه تفسيرېږي، او
بيا د هر~\LR{$k$} دپاره د ``ځايښت~\LR{$k$} اړيکو'' د تعريف ساحه لرو۔ ژبې ته
جوړې دننه اړيکيزې نښې ``\LR{$\Atom{\Obj{true}_k}{R,x_1,\dots,x_k}$}''
ورکوو؛ مراد ئې دا دے چې \LR{$R$} د~\LR{$x_1$}، \dots،~\LR{$x_k$} دپاره صدق کوي،
چرته چې \LR{$R$} د ``\LR{$k$}-ځاييزې اړيکې'' د ډول متغير او \LR{$x_1$}،
\dots،~\LR{$x_k$} د لومړۍ درجې د ډول څيزونه دي۔
\olpseganchor{OLP-0177-B016}{end}

\olpseganchor{OLP-0177-B017}{start}
د دې يو شان ګڼلو له مخې، د دويمې درجې کمزورې معنٰی پوهنه په اصل کښې د
ډېرو ډولونو د منطق عادي معنٰی پوهنه ده؛ او مخکې مو ليدلي چې د ډېرو
ډولونو منطق د لومړۍ درجې په منطق کښې ځايېدلے شي۔ نو د دواړو خواوو
ژباړو ته په پام، د دويمې درجې منطق کمزورے تصور په رښتيا د لومړۍ درجې
منطق يوه پټه بڼه ده چې د تعريف ساحه ئې هم ``څيزونه'' او هم ``اړيکې''
لري او مناسب بديهي اصول پرې واکمن دي۔
\olpseganchor{OLP-0177-B017}{end}

\olpseganchor{OLP-0178-B004}{start}
\olfileid{fol}{byd}{hol}
\olpseganchor{OLP-0178-B004}{end}

\olpseganchor{OLP-0178-B005}{start}
\olsection{د لوړې درجې منطق}
\olpseganchor{OLP-0178-B005}{end}

\olpseganchor{OLP-0178-B006}{start}
له لومړۍ درجې منطق څخه دويمې درجې منطق ته په تېرېدو مو وکړے شول چې
په صوري ژبه کښې د لومړۍ درجې د د تعريف ساحه د څيزونو د سټونو په هکله
خبرې وکړو۔ ولې همدلته ودرېږو؟ د بېلګې په توګه، د درېيمې درجې منطق
بايد موږ ته د څيزونو د سټونو د سټونو، يا ښايي آن د هغو سټونو، په
هکله د خبرو وس راکړي چې هم څيزونه او هم د څيزونو سټونه لري۔ د څلورمې
درجې منطق بيا د دې ډول څيزونو د سټونو په هکله خبرې راکوي۔ لکه چې
اټکل مو کړے وي، دا مفکوره په خپله خوښه هر څومره تکرارولے شو۔
\olpseganchor{OLP-0178-B006}{end}

\olpseganchor{OLP-0178-B007}{start}
په عملي کار کښې د لوړې درجې منطق ډېر ځله د اړيکو پر ځاے د تابعو په
اصطلاحاتو کښې د فارمول په بڼه وړاندې کېږي۔ (له طبيعي يو شان ګڼلو
سره دا توپير بنسټيز نۀ دے۔) چې ځينې بنسټيز ``ډولونه'' \LR{$A$}، \LR{$B$}،
\LR{$C$}،~\dots راکړل شي (چې اوس به ئې ``ټايپونه'' بولو)، د لاندې شرط په
ورکولو نوي ټايپونه جوړولے شو:
\begin{quote}
که \LR{$\sigma$} او \LR{$\tau$} متناهي ټايپونه وي، \LR{$\sigma \to \tau$} هم
متناهي ټايپ دے۔
\end{quote}
ټايپونه داسې نحوي ``نښکې'' وګڼئ چې هغه څيزونه ډلبندي کوي چې په خپل
د تعريف ساحه کښې ئې غواړو؛ \LR{$\sigma \to \tau$} هغه څيزونه بيانوي چې د
ټايپ~\LR{$\sigma$} څيزونه د ټايپ~\LR{$\tau$} څيزونو ته وړونکې تابعې دي۔ د
بېلګې په توګه، ښايي د ``رښتوني'' او ``دروغ'' د رښتياوالي قيمتونو يو
ټايپ~\LR{$\Omega$} او د طبيعي عددونو يو ټايپ~\LR{$\Nat$} وغواړو۔ په دې حالت
کښې د ټايپ~\LR{$\Nat \to \Omega$} څيزونه د يوځاييزو اړيکو يا د~\LR{$\Nat$}
فرعي سټونو په توګه ګڼلے شئ؛ د ټايپ~\LR{$\Nat \to \Nat$} څيزونه له طبيعي
عددونو څخه طبيعي عددونو ته تابعې دي؛ او د ټايپ
\LR{$(\Nat \to \Nat) \to \Nat$} څيزونه ``فانکشنلونه'' دي، يعنې د لوړې
درجې داسې تابعې چې تابعې اخلي او عددونه ورکوي۔
\olpseganchor{OLP-0178-B007}{end}

\olpseganchor{OLP-0178-B008}{start}
لکه د دويمې درجې منطق، د لوړې درجې منطق هم د ډېرو ډولونو يو داسې
منطق ګڼلے شو چې د څيز د هر مطلوب ټايپ دپاره يو ډول لري۔ خو اکثره دا
لا روښانه وي چې د لوړو ټايپونو نحو له بنسټه تعريف کړو۔ د بېلګې په
توګه، د متناهي ټايپونو يو سټ په استقرايي ډول داسې تعريفولے شو:
\begin{enumerate}
\item \LR{$\Nat$} يو متناهي ټايپ دے۔
\item که \LR{$\sigma$} او \LR{$\tau$} متناهي ټايپونه وي، \LR{$\sigma
  \to \tau$} هم متناهي ټايپ دے۔
\item که \LR{$\sigma$} او \LR{$\tau$} متناهي ټايپونه وي، \LR{$\sigma \times
  \tau$} هم متناهي ټايپ دے۔
\end{enumerate}
په وجداني ډول، \LR{$\Nat$} د طبيعي عددونو ټايپ، \LR{$\sigma \to \tau$} له
\LR{$\sigma$} څخه~\LR{$\tau$} ته د تابعو ټايپ، او \LR{$\sigma \times \tau$} د هغو
جوړو ټايپ ښيي چې يو څيز ئې له~\LR{$\sigma$} او بل له~\LR{$\tau$} څخه وي۔ بيا د
ترمونو يو سټ په استقرايي ډول داسې تعريفولے شو:
\begin{enumerate}
\item د هر ټايپ~\LR{$\sigma$} دپاره د \LR{$x$}، \LR{$y$}، \LR{$z$}، \dots په څېر د
  ټايپ~\LR{$\sigma$} متغيرونو يوه زېرمه شته۔
\item \LR{$\Obj 0$} د ټايپ~\LR{$\Nat$} يو ترم دے۔
\item \LR{$\Obj S$} (ناستي) د ټايپ~\LR{$\Nat \to \Nat$} يو ترم دے۔
\item که \LR{$s$} د ټايپ~\LR{$\sigma$} ترم او \LR{$t$} د ټايپ
  \LR{$\Nat \to (\sigma \to \sigma)$} ترم وي، \LR{$\Obj{R}_{st}$} د ټايپ
  \LR{$\Nat \to \sigma$} يو ترم دے۔
\item که \LR{$s$} د ټايپ~\LR{$\tau \to \sigma$} ترم او \LR{$t$} د ټايپ~\LR{$\tau$}
  ترم وي، \LR{$s(t)$} د ټايپ~\LR{$\sigma$} يو ترم دے۔
\item که \LR{$s$} د ټايپ~\LR{$\sigma$} ترم او \LR{$x$} د ټايپ~\LR{$\tau$} متغير وي،
  \LR{$\lambd[x][s]$} د ټايپ~\LR{$\tau \to \sigma$} يو ترم دے۔
\item که \LR{$s$} د ټايپ~\LR{$\sigma$} ترم او \LR{$t$} د ټايپ~\LR{$\tau$} ترم وي،
  \LR{$\tuple{s, t}$} د ټايپ~\LR{$\sigma \times \tau$} يو ترم دے۔
\item که \LR{$s$} د ټايپ~\LR{$\sigma \times \tau$} يو ترم وي، \LR{$p_1(s)$} د
  ټايپ~\LR{$\sigma$} يو ترم او \LR{$p_2(s)$} د ټايپ~\LR{$\tau$} يو ترم دے۔
\end{enumerate}
په وجداني ډول، \LR{$\Obj{R}_{st}$} هغه تابع ښيي چې په بازګشتي ډول داسې
تعريف شوې وي:
\begin{align*}
\Obj{R}_{st}(0) & = s \\
\Obj{R}_{st}(x+1) & = t(x, R_{st}(x)),
\end{align*}
\LR{$\tuple{s, t}$} هغه جوړه ښيي چې لومړے جز ئې~\LR{$s$} او دويم جز ئې~\LR{$t$}
دے، او \LR{$p_1(s)$} او~\LR{$p_2(s)$} د~\LR{$s$} لومړے او دويم غړي
(``پروجيکشنونه'') ښيي۔ په پاى کښې، \LR{$\lambd[x][s]$} هغه تابع~\LR{$f$} ښيي
چې داسې تعريف شوې وي:
\[
f(x) = s
\]
د ټايپ~\LR{$\tau$} د هر~\LR{$x$} دپاره؛ نو توکے~(6) د ادراک يوه بڼه راکوي او
د ترمونو په وسيله د تابعو تعريفول شوني کوي۔ فارمولونه د يوه شان
ټايپ د ترمونو تر منځ له عينيت ويناوو~\LR{$\eq[s][t]$}، عادي بياني
تړونکو او د لوړو ټايپونو له کميت ټاکنې څخه جوړېږي۔ بيا د نظام بديهي
اصول هغه بنسټيزې معادلې ګڼلے شو چې پر پاس تعريف شويو ترمونو واکمنې
دي، او ورسره د کميت ټاکونکو او عينيت لرونکي منطق عادي قاعدې۔
\textbf{د ژباړې د نسخې سمون (OLFOL-033):} په سرچينه کښې د لامبډا ترم
متغير د پايلې ټايپ ته منسوب شوے و؛ دلته هغه د توکي~(6) له تعريف سره
سم د تابع د ورودي ټايپ ته منسوب شوے دے۔
\olpseganchor{OLP-0178-B008}{end}

\olpseganchor{OLP-0178-B009}{start}
که د متناهي ټايپونو نظام ته د رښتياوالي د قيمتونو ټايپ~\LR{$\Omega$}
ورزيات شي، بايد د هغه پر کارونه واکمن بديهي اصول هم شامل شي۔ په حقيقت
کښې، که لږ ځير واوسو، پېچلي فارمولونه په بشپړ ډول لرې کولے او د
ټايپ~\LR{$\Omega$} په ترمونو ئې بدلولے شو!{} بيا د ثبوت نظام هم ورسره سم
بدلېدے شي۔ پايله په اصل کښې د \emph{ټايپونو ساده تيوري} ده چې الونزو
چرچ په ۱۹۳۰مو کلونو کښې وړاندې کړې وه۔
\olpseganchor{OLP-0178-B009}{end}

\olpseganchor{OLP-0178-B010}{start}
لکه د دويمې درجې منطق، د لوړو ټايپونو د معنٰی پوهنې هم بېلابېلې بڼې
شته چې ښايي کارول ئې وغواړو۔ په بشپړه بڼه کښې د ټايپ
\LR{$\sigma \to \tau$} متغيرونه د ټايپ~\LR{$\sigma$} له څيزونو څخه د
ټايپ~\LR{$\tau$} څيزونو ته د \emph{ټولو} تابعو پر سټ ګرځي۔ لکه چې تمه
کېږي، دا معنٰی پوهنه دومره پياوړې ده چې يو بشپړ او اغېزمن
اشتقاق نظام نۀ مني۔ خو کمزورې معنٰی پوهنه هم څېړلے شو چې
جوړښت پکې د هر ټايپ~\LR{$\tau$} دپاره د غړو سټونه~\LR{$T_\tau$} او
ورسره د تطبيق، پروجيکشن او داسې نورو مناسب عملونه لري۔ که جزئيات سم
بشپړ شي، د پورته بيان شويو اشتقاق نظامونو د ډولونو دپاره د
بشپړتيا قضيې ترلاسه کولے شو۔
\olpseganchor{OLP-0178-B010}{end}

\olpseganchor{OLP-0178-B011}{start}
د لوړو ټايپونو منطق ځکه زړه راښکونکے دے چې يو داسې چوکاټ برابروي چې
د رياضياتو لويه برخه پکې په طبيعي ډول ځايولے شو: له~\LR{$\Nat$} څخه په
پيل حقيقي عددونه، مسلسل تابعې او داسې نور تعريفولے شو۔ دا منطق د
شهودي منطق په چاپېريال کښې په ځانګړي ډول زړه راښکونکے هم دے، ځکه
ټايپونه روښانه ``تعميري'' تفسيرونه لري۔ په حقيقت کښې د لوړو ټايپونو
د معنٰی پوهنې تعميري بڼې (چې د کلاسیک منطق پر ځاے پر شهودي منطق
ولاړې وي) جوړولے شو؛ دا بڼې دغه تعميري تفسيرونه ډېر ښه روښانوي او په
ډېرو لارو د خپلو کلاسیکو سيالانو په پرتله په زړه پورې دي۔
\olpseganchor{OLP-0178-B011}{end}

\olpseganchor{OLP-0179-B004}{start}
\olfileid{fol}{byd}{il}
\olpseganchor{OLP-0179-B004}{end}

\olpseganchor{OLP-0179-B005}{start}
\olsection{شهودي منطق}
\olpseganchor{OLP-0179-B005}{end}

\olpseganchor{OLP-0179-B006}{start}
د دويمې او لوړې درجې د منطق پر خلاف، د لومړۍ درجې شهودي منطق د
کلاسیکې بڼې يو محدود شوے نظام دے چې د لا ``تعميري'' ډول استدلال د مدل
کولو دپاره جوړ شوے دے۔ لاندې بېلګې ښايي ځينې بنسټيزې انګېزې روښانه کړي۔
\olpseganchor{OLP-0179-B006}{end}

\olpseganchor{OLP-0179-B007}{start}
فرض کړئ يوه ورځ يو څوک درته راشي او اعلان وکړي چې طبيعي عدد~\LR{$x$} ئې
ټاکلے دے او دا خاصيت لري: که \LR{$x$} اولي وي، د ريمان فرضيه رښتونې ده؛
او که \LR{$x$} مرکب وي، د ريمان فرضيه دروغ ده۔ دا خو لوے زېرے دے!{} د
ريمان د فرضيې رښتياوالے يا دروغوالے د رياضياتو له لويو ناحل شويو
پوښتنو څخه دے، او داسې ښکاري چې دې کس مسئله محاسبې ته راکمه کړې،
يعنې يوازې بايد معلومه شي چې يو ځانګړے عدد اولي دے که نۀ۔
\olpseganchor{OLP-0179-B007}{end}

\olpseganchor{OLP-0179-B008}{start}
د~\LR{$x$} جادويي قيمت څو دے؟ هغه ئې داسې بيانوي: \LR{$x$} هغه طبيعي عدد دے
چې که د ريمان فرضيه رښتونې وي له~\LR{$7$} سره برابر دے، او که نۀ،
له~\LR{$9$} سره۔
\olpseganchor{OLP-0179-B008}{end}

\olpseganchor{OLP-0179-B009}{start}
تاسو په غوسه خپلې پيسې بېرته غواړئ۔ له کلاسیک نظره پاسنے بيان په
رښتيا د~\LR{$x$} يو يوازينے قيمت ټاکي؛ خو څه چې تاسو ئې په اصل کښې غواړئ
د~\LR{$x$} داسې قيمت دے چې \emph{په څرګند ډول} درکړل شوے وي۔
\olpseganchor{OLP-0179-B009}{end}

\olpseganchor{OLP-0179-B010}{start}
د يوې بلې، ښايي لږې مصنوعي، بېلګې دپاره لاندې پوښتنه وګورئ۔ پوهېږو
چې يو غير ناطق عدد ناطق توان ته پورته کېدے او ناطقه پايله ورکولے شي؛
د بېلګې په توګه، \LR{$\sqrt{2}^2 = 2$}۔ دومره څرګنده نۀ ده چې يو غير ناطق
عدد \emph{غير ناطق} توان ته پورته شي او بيا هم ناطقه پايله ورکړي که
نۀ۔ لاندې قضيه مثبت ځواب ورکوي:
\olpseganchor{OLP-0179-B010}{end}

\olpseganchor{OLP-0179-B011}{start}
\begin{thm}
داسې غير ناطق عددونه \LR{$a$} او \LR{$b$} شته چې \LR{$a^b$} ناطق وي۔
\end{thm}
\olpseganchor{OLP-0179-B011}{end}

\olpseganchor{OLP-0179-B012}{start}
\begin{proof}
\LR{$\sqrt{2}^{\sqrt{2}}$} وګورئ۔ که دا ناطق وي، کار بشپړ دے: کولے شو
\LR{$a = b = \sqrt{2}$} ونيسو۔ که نۀ وي، دا غير ناطق دے۔ بيا لرو:
\[
(\sqrt{2}^{\sqrt{2}})^{\sqrt{2}} = \sqrt{2}^{\sqrt{2} \cdot
  \sqrt{2}} = \sqrt{2}^2 = 2,
\]
چې خامخا ناطق دے۔ نو په دې حالت کښې \LR{$a$} له
\LR{$\sqrt{2}^{\sqrt{2}}$} سره او \LR{$b$} له~\LR{$\sqrt 2$} سره برابر ونيسئ۔
\end{proof}
\olpseganchor{OLP-0179-B012}{end}

\olpseganchor{OLP-0179-B013}{start}
آيا دا يو معتبر ثبوت دے؟ زياتره رياضيپوهان ئې معتبر ګڼي۔ خو دلته بيا
هم څه ناڅه نۀ قانع کوونکې خبره شته: د حقيقي عددونو د يوې جوړې شتون
مو په يوه خاصيت سره ثابت کړے، بې له دې چې ويلے شو دا \emph{کومه}
جوړه ده۔ همدا پايله داسې هم ثابتېدے شي چې جوړه~\LR{$a$}، \LR{$b$} پخپله په
ثبوت کښې \emph{ورکړل شوې} وي: \LR{$a = \sqrt{3}$} او \LR{$b = \log_3 4$}
ونيسئ۔ بيا
\[
a^b = \sqrt{3}^{\log_3 4} = 3^{1/2 \cdot \log_3 4} = (3^{\log_3
  4})^{1/2} = 4^{1/2}= 2,
\]
ځکه چې \LR{$3^{\log_3 x} = x$}۔
\olpseganchor{OLP-0179-B013}{end}

\olpseganchor{OLP-0179-B014}{start}
شهودي منطق د داسې استدلال د مدل کولو دپاره جوړ شوے چې د لومړي ثبوت
په څېر ګامونه نۀ مني۔ د~\LR{$!A(x)$} پوره کوونکي~\LR{$x$} د شتون ثابتول په دې
معنا دي چې يو ځانګړے~\LR{$x$} او د هغه له خوا د~\LR{$!A$} د پوره کولو ثبوت
ورکړئ، لکه په دويم ثبوت کښې۔ دا ثابتول چې \LR{$!A$} يا \LR{$!B$} صدق کوي،
غواړي چې له دواړو څخه يو په رښتيا ثابت کړے شئ۔
\olpseganchor{OLP-0179-B014}{end}

\olpseganchor{OLP-0179-B015}{start}
په صوري ډول، د لومړۍ درجې شهودي منطق هغه وخت ترلاسه کېږي چې د لومړۍ
درجې د منطق اشتقاق نظام په يوه ټاکلې طريقه محدود کړئ۔ همداسې
د دويمې يا لوړې درجې منطق شهودي بڼې هم شته۔ له رياضيکي نظره دا يوازې
صوري استنتاجي نظامونه دي، خو لکه مخکې چې وويل شول، موخه ئې د رياضيکي
استدلال يو ډول مدل کول دي۔ څوک دا هغه استدلال ګڼلے شي چې د رياضياتو
يو ځانګړے فلسفي دريځ ئې روا بولي (لکه د بروور شهوديت)؛ څوک ئې د
رياضيکي استدلال يو لا ``ملموس'' او قانع کوونکے ډول ګڼلے شي (د بيشپ د
تعميريت په دود)؛ او دا بحث هم کېدے شي چې صوري بيان غيرصوري انګېزه په
رښتيا رانغاړي که نۀ۔ خو هر فلسفي دريځ چې ولرو، شهودي منطق د يوه صوري
وړاندې شوي منطق په توګه مطالعه کولے شو؛ او د هر دليل له مخې، ډېر
رياضيکي منطقپوهان دا مطالعه په زړه پورې ګڼي۔
\olpseganchor{OLP-0179-B015}{end}

\olpseganchor{OLP-0179-B016}{start}
د شهودي تړونکو يو غيرصوري تعميري تفسير شته چې په عام ډول د بي‌اېچ‌کې
تفسير بلل کېږي (د بروور، هايتنګ او کولموګوروف په نومونو نومول شوے)۔
دا داسې دے: د~\LR{$!A \land !B$} ثبوت د~\LR{$!A$} له يوه ثبوت او د~\LR{$!B$} له
يوه ثبوت څخه جوړه وي؛ د~\LR{$!A \lor !B$} ثبوت يا د~\LR{$!A$} ثبوت وي يا
د~\LR{$!B$}، او څرګند معلومات لرو چې کوم يو دے؛ د~\LR{$!A \lif !B$} ثبوت يوه
کړنلاره وي چې د~\LR{$!A$} ثبوت د~\LR{$!B$} په ثبوت بدلوي؛ د
\LR{$\lforall[x][!A(x)]$} ثبوت يوه کړنلاره وي چې د~\LR{$x$} د هر قيمت دپاره د
\LR{$!A(x)$} ثبوت بېرته ورکوي؛ او د~\LR{$\lexists[x][!A(x)]$} ثبوت د~\LR{$x$} له
يوه قيمت او دې ثبوت څخه جوړ وي چې هغه قيمت~\LR{$!A$} پوره کوي۔ دا تفسير
په هغو محاسبوي اصطلاحاتو هم بيانېدے شي چې د ``کري--هاورډ
آيزومورفيزم'' يا ``د ټايپونو په توګه د فارمولونه بېلګه'' بلل
کېږي: فارمول د معلوماتو يو ټاکلے ټايپ مشخصوونکے وګڼئ، او ثبوتونه
د همدې معلوماتي ټايپونو محاسبوي څيزونه وګڼئ چې راښيي اړوند
فارمول رښتونے دے۔
\olpseganchor{OLP-0179-B016}{end}

\olpseganchor{OLP-0179-B017}{start}
شهودي منطق ډېر ځله کلاسیک منطق ``د خارج وسط له اصل پرته'' ګڼل کېږي۔
لاندې قضيه دا خبره لا دقيقه کوي۔
\begin{thm}
په شهودي ډول لاندې بديهي شېمې يو له بل سره برابرې دي:
\begin{enumerate}
\item \LR{$(\lnot !A \lif \lfalse) \lif !A$}.
\item \LR{$!A \lor \lnot !A$}
\item \LR{$\lnot \lnot !A \lif !A$}
\end{enumerate}
\end{thm}
له دوو نورو څخه د هرې شېمې بېلګې ترلاسه کول په شهودي منطق کښې يو ښه
تمرين دے۔
\olpseganchor{OLP-0179-B017}{end}

\olpseganchor{OLP-0179-B018}{start}
د شهودي بياني منطق لومړني استنتاجي نظامونه، چې د بروور د شهوديت د
صوري بڼو په توګه وړاندې شوي وو، کولموګوروف، ګليوېنکو او هايتنګ په
خپلواکه توګه جوړ کړل۔ د لومړۍ درجې شهودي منطق (او د شهودي رياضياتو
د ځينو برخو) لومړۍ صوري بڼه هايتنګ وړاندې کړه۔ که څه هم يو شمېر په
کلاسیک ډول معتبرې شېمې په شهودي ډول معتبرې نۀ دي، ډېرې نورې معتبرې
دي۔
\olpseganchor{OLP-0179-B018}{end}

\olpseganchor{OLP-0179-B019}{start}
\emph{دوه‌ګونې نفي ژباړه} د کلاسیک او شهودي منطق تر منځ يوه مهمه
اړيکه بيانوي۔ دا په استقرايي ډول داسې تعريفېږي (\LR{$!A^N$} د کلاسیک
فارمول~\LR{$!A$} ``شهودي'' ژباړه وګڼئ):
\begin{align*}
!A^N & \ident \lnot\lnot !A \quad \text{\RL{د اټومي فارمولونه \LR{$!A$} دپاره}} \\
(!A \land !B)^N & \ident (!A^N \land !B^N) \\
(!A \lor !B)^N & \ident  \lnot\lnot (!A^N \lor !B^N) \\
(!A \lif !B)^N & \ident (!A^N \lif !B^N) \\
(\lforall[x][!A])^N & \ident \lforall[x][!A^N] \\
(\lexists[x][!A])^N & \ident \lnot\lnot\lexists[x][!A^N]
\end{align*}
کولموګوروف او ګليوېنکو د بياني منطق دپاره د دې ژباړې بڼې لرلې؛ د
پريديکاتي منطق دپاره ګوډل او ګېنڅن په خپلواکه توګه وړاندې کړه۔ لرو:
\olpseganchor{OLP-0179-B019}{end}

\olpseganchor{OLP-0179-B020}{start}
\begin{thm}
\begin{enumerate}
\item \LR{$!A \liff !A^N$} په کلاسیک ډول د اثبات وړ دے۔
\item که \LR{$!A$} په کلاسیک ډول د اثبات وړ وي، \LR{$!A^N$} په شهودي ډول د
  اثبات وړ دے۔
\end{enumerate}
\end{thm}
\olpseganchor{OLP-0179-B020}{end}

\olpseganchor{OLP-0179-B021}{start}
اوس لاندې خبرې اترې تصورولے شو۔ کلاسیک رياضيپوه: ``ما \LR{$!A$} ثابت
کړے!'' شهودي رياضيپوه: ``ستا ثبوت معتبر نۀ دے۔ تا په اصل کښې
\LR{$!A^N$} ثابت کړے دے۔'' کلاسیک رياضيپوه: ``ماته سمه ده!'' د کلاسیک
رياضيپوه له نظره شهودي رياضيپوه يوازې په وړو توپيرونو نښتے، ځکه دواړه
سره برابر دي۔ خو شهودي رياضيپوه ټينګار کوي چې توپير شته۔
\olpseganchor{OLP-0179-B021}{end}

\olpseganchor{OLP-0179-B022}{start}
پام وکړئ چې پاسنۍ ژباړه يوازې له سوچه منطق سره کار لري؛ دا پوښتنه نۀ
څېړي چې د کلاسیکو او شهودي رياضياتو دپاره مناسب \emph{غيرمنطقي}
بديهي اصول کوم دي يا د دواړو تر منځ څه اړيکه ده۔ خو د پاسنۍ قضيې
لاندې لږ غځونه څه ګټور معلومات راکوي:
\olpseganchor{OLP-0179-B022}{end}

\olpseganchor{OLP-0179-B023}{start}
\begin{thm}
که \LR{$\Gamma$} په کلاسیک ډول \LR{$!A$} ثابتوي، \LR{$\Gamma^N$} په شهودي ډول
\LR{$!A^N$} ثابتوي۔
\end{thm}
\olpseganchor{OLP-0179-B023}{end}

\olpseganchor{OLP-0179-B024}{start}
په نورو ټکو، که \LR{$!A$} له ځينو فرضيو څخه په کلاسیک ډول د اثبات وړ وي،
\LR{$!A^N$} د هغو له دوه‌ګونې نفي ژباړو څخه د اثبات وړ دے۔
\olpseganchor{OLP-0179-B024}{end}

\olpseganchor{OLP-0179-B025}{start}
د دې د ښودلو دپاره چې يوه جمله يا بياني فارمول په شهودي ډول
معتبر دے، يوازې يو ثبوت ورکول پکار دي۔ خو دا به څنګه ښيئ چې معتبر نۀ
دے؟ د دې دپاره داسې معنٰی پوهنې ته اړ يو چې سمه، او غوره دا چې بشپړه
وي۔ د کرېپکي معنٰی پوهنه دا اړتيا ښه پوره کوي۔
\olpseganchor{OLP-0179-B025}{end}

\olpseganchor{OLP-0179-B026}{start}
هماغه کار کولے شو چې د کلاسیک منطق دپاره مو وکړ: معنٰی پوهنه تعريف
کړو او سموالے او بشپړتيا ثابته کړو۔ خو لاندې توپير ته پام ارزښت لري۔
د کلاسیک منطق معنٰی پوهنه په يوه معنا ``څرګنده'' وه او د تړونکو په
معنا کښې ضمني پرته وه۔ که څه هم د کرېپکي د معنٰی پوهنې دپاره څه
وجداني انګېزه ورکولے شو، دا د هماغه ډول حتمي والي احساس نۀ راکوي۔
سربېره پر دې، د يوه کلاسیک جوړښت مفهوم يو طبيعي رياضيکي
مفهوم دے؛ نو يا د جوړښت مفهوم د کلاسیک د لومړۍ درجې منطق د
مطالعې وسيله ګڼلے شو، يا کلاسیک د لومړۍ درجې منطق د رياضيکي
جوړښتونه د مطالعې وسيله۔ د دې پر خلاف، د کرېپکي جوړښتونه
يوازې منطقي جوړونې ګڼل کېدے شي؛ داسې نۀ ښکاري چې خپلواک رياضيکي ارزښت
ولري۔
\olpseganchor{OLP-0179-B026}{end}

\olpseganchor{OLP-0179-B027}{start}
د بياني ژبې دپاره د کرېپکي جوړښت~\LR{$\mModel M = \tuple{W, R, V}$}
د يوه سټ~\LR{$W$}، پر~\LR{$W$} د يوه جزوي ترتيب~\LR{$R$} چې تر ټولو کوچنے
غړے ولري، او د بياني متغيرونو له يوې ``يکنواختې'' ګومارنې
څخه د~\LR{$W$} غړي ته جوړ دے۔ وجداني تصور دا دے چې د~\LR{$W$}
غړي ``نړۍګانې'' يا ``د پوهې حالتونه'' ښيي؛ غړے
\LR{$v \geq u$} د~\LR{$u$} يو ``ممکن راتلونکے حالت'' ښيي؛ او~\LR{$u$} ته ګومارل شوي
بياني متغيرونه هغه قضيې دي چې په حالت~\LR{$u$} کښې رښتونې پېژندل شوې دي۔
بيا د تحميل اړيکه~\LR{$\mSat{M}{!A}[w]$} دا اړيکه د ژبې ټولو
فارمولونه ته غځوي؛ \LR{$\mSat{M}{!A}[w]$} داسې ولولئ چې ``\LR{$!A$} په
حالت~\LR{$w$} کښې رښتونے دے''۔ اړيکه په استقرايي ډول داسې تعريفېږي:
\begin{enumerate}
\item \LR{$\mSat{M}{\Obj p_i}[w]$} هله او يوازې هله چې \LR{$\Obj p_i$} يو له
  هغو بياني متغيرونو څخه وي چې~\LR{$w$} ته ګومارل شوي۔
\item \LR{$\mSat/{M}{\lfalse}[w]$}۔
\item \LR{$\mSat{M}{(!A \land !B)}[w]$} هله او يوازې هله چې
  \LR{$\mSat{M}{!A}[w]$} او \LR{$\mSat{M}{!B}[w]$} وي۔
\item \LR{$\mSat{M}{(!A \lor !B)}[w]$} هله او يوازې هله چې
  \LR{$\mSat{M}{!A}[w]$} يا \LR{$\mSat{M}{!B}[w]$} وي۔
\item \LR{$\mSat{M}{(!A \lif !B)}[w]$} هله او يوازې هله چې هر کله
  \LR{$w' \geq w$} او \LR{$\mSat{M}{!A}[w']$} وي، نو \LR{$\mSat{M}{!B}[w']$} هم وي۔
\end{enumerate}
دا ښودل يو ښه تمرين دے چې
\LR{$\lnot (p \land q) \lif (\lnot p \lor \lnot q)$} په شهودي ډول معتبر
نۀ دے؛ د کرېپکي داسې جوړښت جوړ کړئ چې مقابله بېلګه ورکړي۔
\olpseganchor{OLP-0179-B027}{end}

\olpseganchor{OLP-0180-B004}{start}
\olfileid{fol}{byd}{mod}
\olpseganchor{OLP-0180-B004}{end}

\olpseganchor{OLP-0180-B005}{start}
\olsection{مودالي منطقونه}
\olpseganchor{OLP-0180-B005}{end}

\olpseganchor{OLP-0180-B006}{start}
د يوې شرطي جملې دا بېلګه وګورئ:
\begin{quote}
  که جېرمي په هغې کوټه کښې يوازې وي، نو هغه مست او بربنډ دے او پر
  څوکيو ګډېږي۔
\end{quote}
دا داسې شرطي وينا ده چې ښايي په مادي لحاظ رښتونې خو بيا هم
غولوونکې وي، ځکه عبارت داسې ښيي چې د مقدم او پايلې تر منځ تر دې
پياوړې اړيکه شته چې يوازې يا مقدم دروغ وي يا تالي رښتونے۔ يعنې، عبارت
دا ښيي چې دعوه يوازې په دې ځانګړې نړۍ کښې رښتونې نۀ ده (چرته چې
ښايي ځکه په تش ډول رښتونې وي چې جېرمي په کوټه کښې يوازې نۀ دے)، بلکې
که مقدم رښتونے \emph{وے} نو پايله به هم رښتونې \emph{وه}۔ په نورو
ټکو، دا وينا داسې اخيستل کېدے شي چې دعوه يوازې په دې نړۍ کښې نۀ، بلکې
په هره ``ممکنه'' نړۍ کښې رښتونې ده؛ يا دا چې د دې ځانګړې نړۍ د محض
رښتياوالي پر خلاف، \emph{لازماً} رښتونې ده۔
\olpseganchor{OLP-0180-B006}{end}

\olpseganchor{OLP-0180-B007}{start}
مودالي منطق د دې ډول لزوم د معنا کولو دپاره جوړ شوے۔ مودالي بياني
منطق له عادي بياني منطق څخه د بکس عملګر په زياتولو ترلاسه کېږي؛ يعنې
که \LR{$!A$} فارمول وي، نو \LR{$\Box !A$} هم دے۔ په وجداني ډول \LR{$\Box !A$}
وايي چې \LR{$!A$} \emph{لازماً} رښتونے، يا په هره ممکنه نړۍ کښې رښتونے
دے۔ \LR{$\Diamond !A$} په عام ډول د \LR{$\lnot \Box \lnot !A$} لنډون ګڼل کېږي
او داسې لوستل کېږي چې \LR{$!A$} \emph{ممکن} رښتونے دے۔ البته موداليت په
پريديکاتي منطق کښې هم ورزياتېدلے شي۔
\olpseganchor{OLP-0180-B007}{end}

\olpseganchor{OLP-0180-B008}{start}
کرېپکي جوړښتونه مودالي منطق ته معنٰی پوهنه ورکولے شي؛ په حقيقت
کښې کرېپکي دا معنٰی پوهنه لومړے د مودالي منطق په پام کښې نيولو سره
جوړه کړې وه۔ دلته پر جزوي ترتيبونو د محدودېدو پر ځاے په عام ډول د
``ممکنو نړۍګانو'' يو سټ~\LR{$P$} او د نړۍګانو تر منځ دوه ځاييزه
``د لاسرسي'' اړيکه \LR{$\Atom{R}{x,y}$} لرو۔ په وجداني ډول
\LR{$\Atom{R}{p,q}$} وايي چې نړۍ~\LR{$q$} له~\LR{$p$} سره سازګاره ده؛ يعنې که موږ
په نړۍ~\LR{$p$} کښې ``يو''، بايد دا امکان ومنو چې نړۍ به د~\LR{$q$} په څېر وه۔
\olpseganchor{OLP-0180-B008}{end}

\olpseganchor{OLP-0180-B009}{start}
مودالي منطق کله د ``مفهومي'' منطق په نامه يادېږي، د ``مصداقي'' منطق
پر خلاف۔ د کلاسیک منطق په څېر د يوه مصداقي منطق ټاکل شوې معنٰی پوهنه
يوازې يوې، ``واقعي'' نړۍ ته اشاره کوي؛ خو د يوه ``مفهومي'' منطق معنٰی
پوهنه پر يوې پراخې هستپوهنې تکيه کوي۔ د لزوم د جوړښت کولو تر
څنګ، موداليت د نورو ژبنيو جوړښتونو د جوړښت کولو دپاره هم
کارېدلے شي، او \LR{$\Box$} او \LR{$\Diamond$} د اړوند کارونې له مخې بيا تفسيرېږي۔
د بېلګې په توګه:
\begin{enumerate}
\item د اثبات‌وړتيا په منطق کښې \LR{$\Box !A$} داسې لوستل کېږي چې ``\LR{$!A$}
  د اثبات وړ دے'' او \LR{$\Diamond !A$} داسې چې ``\LR{$!A$} سازګار دے''۔
\item په معرفتي منطق کښې \LR{$\Box !A$} داسې لوستل کېدے شي چې ``زه په
  \LR{$!A$} پوهېږم'' يا ``زه پر \LR{$!A$} باور لرم''۔
\item په زماني منطق کښې \LR{$\Box !A$} داسې لوستل کېدے شي چې ``\LR{$!A$} تل
  رښتونے وي'' او \LR{$\Diamond !A$} داسې چې ``\LR{$!A$} کله کله رښتونے وي''۔
\end{enumerate}
\olpseganchor{OLP-0180-B009}{end}

\olpseganchor{OLP-0180-B010}{start}
غواړو منطق په هغو قاعدو او بديهي اصولو پياوړے کړو چې له موداليت سره
کار لري۔ د بېلګې په توګه، نظام~\LR{$\Log{S4}$} د بياني منطق عادي بديهي
اصول او قاعدې، او ورسره لاندې بديهي اصول لري:
\begin{align*}
& \Box (!A \lif !B) \lif (\Box !A \lif \Box !B)\\
& \Box !A \lif !A \\
& \Box !A \lif \Box \Box !A
\intertext{\RL{او دا قاعده هم: ``له \LR{$!A$} څخه \LR{$\Box !A$} پايله واخلئ''۔
  \LR{$\Log{S5}$} لاندې بديهي اصل ورزياتوي:}}
& \Diamond !A \lif \Box \Diamond !A
\end{align*}
د دې بديهي اصولو بدل ډولونه ښايي د بېلابېلو کارونو دپاره مناسب وي؛ د
بېلګې په توګه، S5 په عام ډول د منطقي لزوم د مفهوم ځانګړوونکے ګڼل
کېږي۔ ښه خبره دا ده چې اکثره داسې معنٰی پوهنه موندلے شو چې د کرېپکي
په جوړښتونه کښې د لاسرسي اړيکه په طبيعي لارو محدودوي او د هغې
په نسبت د اشتقاق نظام سم او بشپړ وي۔ د بېلګې په توګه،
\LR{$\Log{S4}$} د کرېپکي د هغو جوړښتونه له ټولګي سره سمون لري چې د
لاسرسي اړيکه ئې انعکاسي او انتقالي وي۔ \LR{$\Log{S5}$} د کرېپکي د هغو
جوړښتونه له ټولګي سره سمون لري چې د لاسرسي اړيکه ئې {\em عمومي}
وي، يعنې هره نړۍ له هرې بلې نړۍ څخه د لاسرسي وړ وي؛ نو \LR{$\Box !A$} هله
او يوازې هله صدق کوي چې \LR{$!A$} په هره نړۍ کښې صدق وکړي۔
\olpseganchor{OLP-0180-B010}{end}

\olpseganchor{OLP-0181-B004}{start}
\olfileid{fol}{byd}{oth}
\olpseganchor{OLP-0181-B004}{end}

\olpseganchor{OLP-0181-B005}{start}
\olsection{نور منطقونه}
\olpseganchor{OLP-0181-B005}{end}

\olpseganchor{OLP-0181-B006}{start}
لکه چې ښايي تر اوسه مو معلومه کړې وي، د نوي منطق جوړول ګران نۀ دي۔
تاسو هم خپله نحو جوړولے، د استنتاج يو نظام پنځولے، او ورته معنٰی پوهنه
ترتيبولے شئ۔ که غواړئ د اشتقاق نظام مو د معنٰی پوهنې په نسبت
بشپړ وي، ښايي لږ ځيرکتيا ته اړ شئ؛ او د نړۍ د قانع کولو دپاره هم څه
هڅه پکار ده چې منطق مو په رښتيا په زړه پورې دے۔ خو په بدل کښې ئې د
خپل منطق د رياضيکي او محاسبوي خاصيتونو په سپړلو ډېر ساعتونه پاکه او
خوندوره بوختيا لرلے شئ۔
\olpseganchor{OLP-0181-B006}{end}

\olpseganchor{OLP-0181-B007}{start}
وروستيو لسيزو د صوري منطقونو رښتينې چاودنه ليدلې ده۔ فزي منطق د
مبهمو خاصيتونو په هکله د استدلال د مدل کولو دپاره جوړ شوے۔ احتمالي
منطق د ناڅرګندتيا په هکله استدلال مدل کوي۔ ډيفالټ منطقونه او
نايکنواخته منطقونه د ماتېدونکو استدلالي بڼو د مدل کولو دپاره جوړ شوي؛
يعنې داسې ``معقول'' استنتاجونه چې د نويو معلوماتو په راتلو وروسته
ردېدلے شي۔ معرفتي منطقونه د پوهې په هکله استدلال، سببي منطقونه د سبب
او مسبب د اړيکو په هکله استدلال، او آن ``تکليفي'' منطقونه د اخلاقي
مکلفيتونو په هکله استدلال مدل کوي۔ دا نظامونه چې د فلسفي، رياضيکي يا
محاسبوي انګېزې له مخې معرفي شوي وي، ښايي د رياضيکي منطق، فلسفي منطق،
مصنوعي ځيرکتيا، ادراکي پوهنې يا بلې څانګې تر سرليک لاندې مطالعه شي۔
\olpseganchor{OLP-0181-B007}{end}

\olpseganchor{OLP-0181-B008}{start}
لړ همداسې اوږدېږي او امکانونه بې‌پايه ښکاري۔ ښايي هېڅکله د لايبنېڅ
دې خوب ته و نۀ رسېږو چې ټول انساني عقل محاسبې ته راکم کړو---خو دا مو
له هڅې څخه نۀ شي منع کولے۔
\olpseganchor{OLP-0181-B008}{end}

\olpseganchor{OLP-0182-B004}{start}
\olpart{mod}{مدل تيوري}
\olpseganchor{OLP-0182-B004}{end}

\olpseganchor{OLP-0182-B005}{start}
\begin{editorial}
  د مدل تيورۍ مواد نيمګړي او ازمېښتي دي۔ دا مهال دا يوازې د مدل
  تيورۍ په هکله د الډو انتونېلي له يادښتونو څخه اخيستل شوې بڼه ده، له
  هغو موضوعاتو پرته چې د لومړۍ درجې منطق په برخه کښې راغلې دي (تيورۍ،
  بشپړتيا او د متناهي شاهد خاصيت)۔ د درسي کتاب دپاره د ګټورتيا په موخه
  دا برخه لا ډېرې سريزې، انګېزې، سپړنې او تمرينونو ته اړتيا لري۔ اېنډي
  ارانا پلان لري چې په ځانګړي ډول پر دې برخه کار وکړي (\gitissue{65})۔
\end{editorial}
\olpseganchor{OLP-0182-B005}{end}

\olpseganchor{OLP-0183-B004}{start}
\olchapter{mod}{bas}{د مدل تيورۍ بنسټونه}
\olpseganchor{OLP-0183-B004}{end}















\olpseganchor{OLP-0183-B012}{start}
%
\olpseganchor{OLP-0183-B012}{end}

\olpseganchor{OLP-0184-B004}{start}
\olfileid{mod}{bas}{red}
\section{راکمونې او غځونې}
\olpseganchor{OLP-0184-B004}{end}

\olpseganchor{OLP-0184-B005}{start}
ډېر ځله د هغو ژبو، او د دغو ژبو د جوړښتونه پرتله ګټوره يا
اړينه وي چې ځينې نښې سره ګډې لري۔ تر ټولو عام حالت دا دے چې د يوې
ژبه~\LR{$\Lang{L}$} ټولې نښې د بلې ژبه~\LR{$\Lang{L'}$}
برخه هم وي، يعنې \LR{$\Lang{L} \subseteq \Lang{L'}$}۔ بيا يو
\LR{$\Lang{L}$}-جوړښت~\LR{$\Struct{M}$} تل د اضافي نښو په تفسيرولو
\LR{$\Lang{L'}$}-جوړښت ته غځېدلے شي، په داسې حال کښې چې د ګډو نښو
تفسيرونه هماغسې پاتې شي۔ له بلې خوا، له يوه
\LR{$\Lang{L'}$}-جوړښت~\LR{$\Struct{M'}$} څخه د هغو نښو د تفسيرونو په
``هېرولو'' يو \LR{$\Lang{L}$}-جوړښت ترلاسه کولے شو چې په~\LR{$\Lang{L}$} کښې
نۀ راځي۔
\olpseganchor{OLP-0184-B005}{end}

\olpseganchor{OLP-0184-B006}{start}
\begin{defn}
\ollabel{defn:reduct}
فرض کړئ \LR{$\Lang L \subseteq \Lang L'$}، \LR{$\Struct M$} يو
\LR{$\Lang L$}-جوړښت او \LR{$\Struct M'$} يو \LR{$\Lang L'$}-جوړښت
دے۔ \LR{$\Struct M$} د~\LR{$\Struct M'$} تر~\LR{$\Lang L$} پورې \emph{راکمونه}، او
\LR{$\Struct M'$} د~\LR{$\Struct M$} تر~\LR{$\Lang L'$} پورې \emph{غځونه} هله او
يوازې هله ده چې
\begin{enumerate}
\item \LR{$\Domain{M} = \Domain{M'}$}
\item د هر داسې ثابت سمبول~\LR{$c \in \Lang L$} دپاره،
  \LR{$\Assign{c}{M} = \Assign{c}{M'}$} وي۔
\item د هر داسې تابع~\LR{$f \in \Lang L$} دپاره،
  \LR{$\Assign{f}{M} = \Assign{f}{M'}$} وي۔
\item د هر داسې پريديکات~\LR{$P \in \Lang L$} دپاره،
  \LR{$\Assign{P}{M} = \Assign{P}{M'}$} وي۔
\end{enumerate}
\end{defn}
\olpseganchor{OLP-0184-B006}{end}

\olpseganchor{OLP-0184-B007}{start}
\begin{prop}
\ollabel{prop:reduct}
که يو \LR{$\Lang{L}$}-جوړښت~\LR{$\Struct{M}$} د يوه
\LR{$\Lang{L'}$}-جوړښت~\LR{$\Struct{M'}$} راکمونه وي، نو د~\LR{$\Lang{L}$}
د ټولو تړلي فارمولونه~\LR{$!A$} دپاره
\[
\Sat{M}{!A} \text{\RL{ هله او يوازې هله چې }} \Sat{M'}{!A}.
\]
\end{prop}
\olpseganchor{OLP-0184-B007}{end}

\olpseganchor{OLP-0184-B008}{start}
\begin{proof}
  تمرين۔
\end{proof}
\olpseganchor{OLP-0184-B008}{end}

\olpseganchor{OLP-0184-B009}{start}
\begin{prob}
\olref[mod][bas][red]{prop:reduct} ثابته کړئ۔
\end{prob}
\olpseganchor{OLP-0184-B009}{end}

\olpseganchor{OLP-0184-B010}{start}
\begin{defn}
که يو \LR{$\Lang{L}$}-جوړښت \LR{$\Struct{M}$} ولرو، او \LR{$\Lang{L'} =
\Lang{L} \cup \{P\}$} د~\LR{$\Lang{L}$} هغه غځونه وي چې يوازينے
\LR{$n$}-ځاييز پريديکات~\LR{$P$} ورزياتوي، او \LR{$R \subseteq \Domain{M}^n$}
يوه \LR{$n$}-ځاييزه اړيکه وي، نو د~\LR{$\Struct{M}$} هغه
غځونه~\LR{$\Struct{M'}$} چې \LR{$\Assign{P}{M'} = R$} وي، په
\LR{$\Expan{M}{R}$} ليکو۔
\end{defn}
\olpseganchor{OLP-0184-B010}{end}

\olpseganchor{OLP-0185-B004}{start}
\olfileid{mod}{bas}{sub}
\olsection{فرعي جوړښتونه}
\olpseganchor{OLP-0185-B004}{end}

\olpseganchor{OLP-0185-B005}{start}
د يوه جوړښت~\LR{$\Struct{M}$} د تعريف ساحه ښايي د بل
جوړښت~\LR{$\Struct{M'}$} فرعي سټ وي۔ خو څرګنده ده چې \LR{$\Struct{M}$} بايد
يوازې هغه وخت د~\LR{$\Struct{M'}$} يوه ``برخه'' وګڼو چې نۀ يوازې
\LR{$\Domain{M} \subseteq \Domain{M'}$} وي، بلکې \LR{$\Struct{M}$} او
\LR{$\Struct{M'}$} لږ تر لږه پر ګډې برخې~\LR{$\Domain{M}$} د ژبې د نښو په
تفسيرولو کښې هم ``سره سلا'' وي۔
\olpseganchor{OLP-0185-B005}{end}

\olpseganchor{OLP-0185-B006}{start}
\begin{defn}
\ollabel{defn:substructure}
د يوې ژبې~\LR{$\Lang L$} جوړښتونه \LR{$\Struct M$} او \LR{$\Struct M'$} راکړل شي۔
وايو \LR{$\Struct M$} د~\LR{$\Struct M'$} \emph{فرعي جوړښت}، او
\LR{$\Struct M'$} د~\LR{$\Struct M$} \emph{غځونه} ده، چې
\LR{$\Struct M \substruct \Struct M'$} ئې ليکو، هله او يوازې هله چې
\begin{enumerate}
\item \LR{$\Domain{M} \subseteq \Domain{M'}$} وي؛
\item د هر داسې ثابت \LR{$c \in \Lang L$} دپاره، \LR{$\Assign{c}{M} =
    \Assign{c}{M'}$} وي؛
\item د هر داسې \LR{$n$}-ځاييز تابع~\LR{$f \in \Lang L$} دپاره، د
  ټولو \LR{$a_1$}، \dots،~\LR{$a_n \in \Domain{M}$} دپاره
  \LR{$\Assign{f}{M}(a_1, \dots, a_n) = \Assign{f}{M'}(a_1, \dots, a_n)$}
  وي؛
\item د هر داسې \LR{$n$}-ځاييز پريديکات~\LR{$R \in \Lang L$} دپاره، د
  ټولو \LR{$a_1$}، \dots،~\LR{$a_n \in \Domain{M}$} دپاره
  \LR{$\langle a_1, \dots, a_n\rangle \in \Assign{R}{M}$} هله او يوازې
  هله چې \LR{$\langle a_1, \dots, a_n\rangle \in \Assign{R}{M'}$} وي۔
\end{enumerate}
\end{defn}
\olpseganchor{OLP-0185-B006}{end}

\olpseganchor{OLP-0185-B007}{start}
\begin{rem}
\ollabel{rem:substructure}
که ژبه ثابت يا تابعې ونۀ لري، نو هر
\LR{$N \subseteq \Domain{M}$} د~\LR{$\Struct M$} يو فرعي
جوړښت~\LR{$\Struct{N}$} ټاکي چې د تعريف ساحه ئې
\LR{$\Domain{N} = N$} دے؛ د دې دپاره \LR{$\Assign{R}{N} =
\Assign{R}{M} \cap N^n$} ږدو۔
\end{rem}
\olpseganchor{OLP-0185-B007}{end}

\olpseganchor{OLP-0185-B008}{start}
% prove something about this? Examples?
\olpseganchor{OLP-0185-B008}{end}

\olpseganchor{OLP-0186-B004}{start}
\olfileid{mod}{bas}{ove}
\olsection{له بريده اوښتنه}
\olpseganchor{OLP-0186-B004}{end}

\olpseganchor{OLP-0186-B005}{start}
\begin{thm}
\ollabel{overspill} که د جملو يو سټ~\LR{$\Gamma$} د هرې کچې متناهي
مدلونه ولري، نو يو نامتناهي مدل هم لري۔
\end{thm}
\olpseganchor{OLP-0186-B005}{end}

\olpseganchor{OLP-0186-B006}{start}
\begin{proof}
د~\LR{$\Gamma$} ژبه د شمېر وړ نويو ثابتونو \LR{$c_0$}، \LR{$c_1$}، \dots په
ورزياتولو وغځوئ او سټ \LR{$\Gamma \cup \{c_i \neq c_j : i \neq j\}$}
وګورئ۔ دا ويل چې \LR{$\Gamma$} د هرې کچې متناهي مدلونه لري، معنا ئې دا ده
چې د هر \LR{$m >0$} دپاره داسې \LR{$n\ge m$} شته چې \LR{$\Gamma$} د شمېر~\LR{$n$} يو مدل
لري۔ له دې پايله کېږي چې \LR{$\Gamma \cup \{c_i \neq c_j : i \neq j\}$}
په متناهي ډول د صدق وړ دے۔ د متناهي شاهد قضيې له مخې، \LR{$\Gamma \cup
\{c_i \neq c_j : i \neq j\}$} يو مدل~\LR{$\Struct M$} لري چې د تعريف ساحه
ئې بايد نامتناهي وي، ځکه ټولې نابرابرۍ \LR{$c_i \neq c_j$} پوره کوي۔
\end{proof}
\olpseganchor{OLP-0186-B006}{end}

\olpseganchor{OLP-0186-B007}{start}
\begin{prop}
\ollabel{inf-not-fo}
د لومړۍ درجې د هېڅ ژبې داسې جمله~\LR{$!A$} نشته چې په يوه
جوړښت~\LR{$\Struct M$} کښې هله او يوازې هله رښتونې وي چې د دغه
جوړښت د تعريف ساحه~\LR{$\Domain{M}$} نامتناهي وي۔
\end{prop}
\olpseganchor{OLP-0186-B007}{end}

\olpseganchor{OLP-0186-B008}{start}
\begin{proof}
که داسې~\LR{$!A$} موجوده وے، نفي~\LR{$\lnot !A$} به ئې په ټولو او يوازې په
متناهي جوړښتونه کښې رښتونې وه؛ نو د هرې کچې متناهي مدلونه به
ئې لرل خو نامتناهي مدل به ئې نۀ درلود، چې دا له~\olref{overspill}
سره په ټکر کښې ده۔
\end{proof}
\olpseganchor{OLP-0186-B008}{end}

\olpseganchor{OLP-0187-B004}{start}
\olfileid{mod}{bas}{iso}
\olsection{آيزومورف جوړښتونه}
\olpseganchor{OLP-0187-B004}{end}

\olpseganchor{OLP-0187-B005}{start}
د لومړۍ درجې جوړښتونه په دوو لارو يو شان کېدے شي۔ يوه لار دا
ده چې يو شان تړلي فارمولونه رښتونې کړي۔ داسې جوړښتونه
\emph{ابتدايي هم‌ارز} بولو۔ خو جوړښتونه ډېر سره توپير لرلے او بيا هم
يو شان تړلي فارمولونه رښتونې کولے شي---د بېلګې په توګه، يو ئې
د شمېر وړ کېدے شي او بل نۀ۔ لامل دا دے چې د يوه
جوړښت ډېر خاصيتونه د لومړۍ درجې په ژبو کښې نۀ شي بيانېدلے؛
يا ځکه چې ژبه پوره بډايه نۀ وي، يا د لومړۍ درجې منطق د بنسټيزو
محدوديتونو، لکه د لوېنهايم--سکولم د قضيې، له امله۔ نو بله او لا سخته
معنا چې جوړښتونه پکې يو شان کېدے شي دا ده چې له بنسټه هماغه وي:
يوازې په هغو څيزونو کښې توپير ولري چې ترې جوړ دي، خو په جوړښتي
خاصيتونو کښې نۀ۔ د دې خبرې د دقيقولو يوه لار د \emph{آيزومورفيزم}
مفهوم دے۔
\olpseganchor{OLP-0187-B005}{end}

\olpseganchor{OLP-0187-B006}{start}
\begin{defn}
\ollabel{defn:elem-equiv}
د يوې ژبه~\LR{$\Lang{L}$} دوه جوړښتونه \LR{$\Struct{M}$} او
\LR{$\Struct M'$} راکړل شي۔ وايو \LR{$\Struct{M}$} له~\LR{$\Struct M'$} سره
\emph{ابتدايي هم‌ارز} دے، چې \LR{$\Struct{M} \equiv \Struct M'$} ئې
ليکو، هله او يوازې هله چې د~\LR{$\Lang{L}$} د هرې تړلے فارمول~\LR{$!A$}
دپاره \LR{$\Sat{M}{!A}$} هله او يوازې هله چې \LR{$\Sat{M'}{!A}$} وي۔
\end{defn}
\olpseganchor{OLP-0187-B006}{end}

\olpseganchor{OLP-0187-B007}{start}
\begin{defn}
\ollabel{defn:isomorphism} د يوې ژبه~\LR{$\Lang L$} دوه
جوړښتونه \LR{$\Struct{M}$} او \LR{$\Struct M'$} راکړل شي۔ وايو
\LR{$\Struct{M}$} له~\LR{$\Struct M'$} سره \emph{آيزومورف} دے، چې
\LR{$\Struct{M} \simeq \Struct M'$} ئې ليکو، هله او يوازې هله چې داسې
تابع \LR{$h \colon \Domain{M} \to \Domain{M'}$} شته چې:
\begin{enumerate}
\item \LR{$h$} يو پر يو ده: که \LR{$h(x) = h(y)$} وي، نو \LR{$x = y$}؛
\item \LR{$h$} پر هدف سټ بشپړه ده: د هر \LR{$y \in \Domain{M'}$} دپاره داسې
  \LR{$x \in \Domain{M}$} شته چې \LR{$h(x) = y$}؛
\item \ollabel{defn:iso-const}د هر ثابت سمبول~\LR{$c$} دپاره:
  \LR{$h(\Assign{c}{M}) = \Assign{c}{M'}$}؛
\item \ollabel{defn:iso-pred}د هر \LR{$n$}-ځاييز پريديکات~\LR{$P$} دپاره:
  \[
  \tuple{a_1, \dots, a_n}\in \Assign{P}{M} \quad\text{\RL{هله او يوازې هله چې}}\quad
  \tuple{h(a_1), \dots, h(a_n)} \in \Assign{P}{M'};
  \]
\item \ollabel{defn:iso-func}د هر \LR{$n$}-ځاييز تابع~\LR{$f$} دپاره:
  \[
  h(\Assign{f}{M}(a_1, \dots, a_n)) =
  \Assign{f}{M'}(h(a_1), \dots, h(a_n)).
  \]
\end{enumerate}
\end{defn}
\olpseganchor{OLP-0187-B007}{end}

\olpseganchor{OLP-0187-B008}{start}
\begin{thm}
\ollabel{thm:isom}
که \LR{$\Struct{M} \iso \Struct M'$} وي، نو \LR{$\Struct{M} \elemequiv
\Struct{M'}$}۔
\end{thm}
\olpseganchor{OLP-0187-B008}{end}

\olpseganchor{OLP-0187-B009}{start}
\begin{proof}
فرض کړئ \LR{$h$} له~\LR{$\Struct{M}$} څخه پر~\LR{$\Struct M'$} يو آيزومورفيزم دے۔
د هرې ګومارنې~\LR{$s$} دپاره \LR{$h \circ s$} د~\LR{$h$} او~\LR{$s$} ترکيب، يعنې
په~\LR{$\Struct{M'}$} کښې هغه ګومارنه ده چې
\LR{$(h \circ s)(x) = h(s(x))$} وي۔ پر~\LR{$t$} او~\LR{$!A$} په استقرا لاندې
پياوړې دعوې ثابتولے شو:
\begin{enumerate}
  \item[a.] \LR{$h(\Value{t}{M}[s]) = \Value{t}{M'}[h\circ s]$}۔
  \item[b.] \LR{$\Sat{M}{!A}[s]$} هله او يوازې هله چې
    \LR{$\Sat{M'}{!A}[h \circ s]$}۔
\end{enumerate}
لومړۍ دعوه د~\LR{$t$} پر پېچلتيا په استقرا ثابتېږي۔
\begin{enumerate}
\item که \LR{$t \ident c$} وي، نو \LR{$\Value{c}{M}[s] = \Assign{c}{M}$} او
  \LR{$\Value{c}{M'}[h \circ s] = \Assign{c}{M'}$}۔ نو
  \LR{$h(\Value{t}{M}[s]) = h(\Assign{c}{M}) = \Assign{c}{M'}$} (د
  \olref{defn:isomorphism} له~\olref{defn:iso-const} څخه)
  \LR{$= \Value{t}{M'}[h \circ s]$}۔
\item که \LR{$t \ident x$} وي، نو \LR{$\Value{x}{M}[s] = s(x)$} او
  \LR{$\Value{x}{M'}[h \circ s] = h(s(x))$}۔ نو
  \LR{$h(\Value{x}{M}[s]) = h(s(x)) = \Value{x}{M'}[h \circ s]$}۔
\item که \LR{$t \ident f(t_1, \dots, t_n)$} وي، نو
  \begin{align*}
    \Value{t}{M}[s] & = \Assign{f}{M}(\Value{t_1}{M}[s], \dots, \Value{t_n}{M}[s]) \quad\text{\RL{او}}\\
  \Value{t}{M'}[h \circ s] & = \Assign{f}{M'}(\Value{t_1}{M'}[h \circ
    s], \dots, \Value{t_n}{M'}[h \circ s]).
  \end{align*}
  \textbf{د ژباړې د نسخې سمون (OLMOD-001):} سرچينې په دويمه کرښه کښې
  د دويم جوړښت د ترم قيمت د لومړي جوړښت د تابعې په تفسير ليکلے و؛
  دلته د هماغې کرښې د دويم جوړښت تفسير بېرته ايښوول شوے دے۔ د استقرا
  فرض دا دے چې د هر~\LR{$i$} دپاره \LR{$h(\Value{t_i}{M}[s]) =
  \Value{t_i}{M'}[h\circ s]$}۔ نو
  \begin{align}
    h(\Value{t}{M}[s])
    & = h(\Assign{f}{M}(\Value{t_1}{M}[s], \dots, \Value{t_n}{M}[s])) \notag\\
    & = \Assign{f}{M'}(h(\Value{t_1}{M}[s]), \dots,
    h(\Value{t_n}{M}[s])) \ollabel{iso-1}\\
    & = \Assign{f}{M'}(\Value{t_1}{M'}[h \circ s], \dots,
    \Value{t_n}{M'}[h \circ s]) \ollabel{iso-2}\\
    & = \Value{t}{M'}[h\circ s] \notag
  \end{align}
  \textbf{د ژباړې د نسخې سمون (OLMOD-002):} سرچينې د لومړۍ کرښې د
  بهرنۍ تابعې تړونکې قوس غورځولې وه؛ دلته هغه قوس بېرته ايښوول شوے دے۔
  دلته \olref{iso-1} د~\olref{defn:isomorphism} له
  \olref{defn:iso-func} او \olref{iso-2} د استقرا له فرض څخه پايله
  کېږي۔
\end{enumerate}
د (b) برخې ثبوت د تمرين دپاره پرېښودل شوے دے۔
\olpseganchor{OLP-0187-B009}{end}

\olpseganchor{OLP-0187-B010}{start}
که \LR{$!A$} جمله وي، ګومارنې~\LR{$s$} او~\LR{$h \circ s$} بې اغېزې دي، او
\LR{$\Sat{M}{!A}$} هله او يوازې هله چې \LR{$\Sat{M'}{!A}$} لرو۔
\end{proof}
\olpseganchor{OLP-0187-B010}{end}

\olpseganchor{OLP-0187-B011}{start}
\begin{prob}
د~\olref[mod][bas][iso]{thm:isom} د (b) برخې ثبوت په تفصيل بشپړ
کړئ۔ پام وکړئ چې د~\olref[mod][bas][iso]{defn:isomorphism} د
آيزومورفيزم ځانګړوونکي پنځه خاصيتونه هر يو چېرته کارول کېږي۔
\end{prob}
\olpseganchor{OLP-0187-B011}{end}

\olpseganchor{OLP-0187-B012}{start}
\begin{defn}
د يوه جوړښت~\LR{$\Struct{M}$} \emph{اتومورفيزم} له~\LR{$\Struct{M}$} څخه پر
خپل ځان يو آيزومورفيزم دے۔
\end{defn}
\olpseganchor{OLP-0187-B012}{end}

\olpseganchor{OLP-0187-B013}{start}
\begin{prob}
وښيئ چې د هر جوړښت~\LR{$\Struct{M}$} دپاره، که \LR{$X$} د~\LR{$\Struct{M}$} يو
تعريفېدونکی فرعي سټ او \LR{$h$} د~\LR{$\Struct{M}$} يو اتومورفيزم وي، نو
\LR{$X = \Setabs{h(x)}{x \in X}$} (يعنې \LR{$X$} د~\LR{$h$} لاندې ثابت پاتې کېږي)۔
\end{prob}
\olpseganchor{OLP-0187-B013}{end}

\olpseganchor{OLP-0188-B004}{start}
\olfileid{mod}{bas}{thm}
\olpseganchor{OLP-0188-B004}{end}

\olpseganchor{OLP-0188-B005}{start}
\section{د يوه جوړښت تيوري}
\olpseganchor{OLP-0188-B005}{end}

\olpseganchor{OLP-0188-B006}{start}
هر جوړښت~\LR{$\Struct{M}$} ځينې تړلي فارمولونه رښتونې او ځينې
دروغې کوي۔ د هغو ټولو تړلي فارمولونه سټ چې دا جوړښت ئې رښتونې کوي، د
هغه \emph{تيوري} بلل کېږي۔ دا سټ په رښتيا يوه تيوري ده، ځکه د هغې هر
معنايي استلزام بايد د هغې په ټولو مدلونو، او له دې ډلې
په~\LR{$\Struct{M}$} کښې، رښتونے وي۔
\olpseganchor{OLP-0188-B006}{end}

\olpseganchor{OLP-0188-B007}{start}
\begin{defn}
  يو جوړښت~\LR{$\Struct M$} راکړل شي۔ د~\LR{$\Struct{M}$}
  \emph{تيوري} د هغو تړلي فارمولونه سټ~\LR{$\Theory{M}$} دے چې
  په~\LR{$\Struct{M}$} کښې رښتونې وي، يعنې
  \LR{$\Theory{M} = \Setabs{!A}{\Sat{M}{!A}}$}۔
\end{defn}
\olpseganchor{OLP-0188-B007}{end}

\olpseganchor{OLP-0188-B008}{start}
د ``تيورۍ'' اصطلاح په غيرصوري ډول د هغو تړلي فارمولونه سټونو دپاره
هم کاروو چې يو ټاکلے مطلوب تفسير لري، که د استنتاج لاندې تړلي وي او
که نۀ۔
\olpseganchor{OLP-0188-B008}{end}

\olpseganchor{OLP-0188-B009}{start}
\begin{prop}
د هر~\LR{$\Struct{M}$} دپاره \LR{$\Theory{M}$} بشپړه ده۔
\end{prop}
\olpseganchor{OLP-0188-B009}{end}

\olpseganchor{OLP-0188-B010}{start}
\begin{proof}
د هرې تړلے فارمول~\LR{$!A$} دپاره يا \LR{$\Sat{M}{!A}$} وي يا
\LR{$\Sat{M}{\lnot !A}$}؛ نو يا \LR{$!A \in \Theory{M}$} وي يا
\LR{$\lnot !A \in \Theory{M}$}۔
\end{proof}
\olpseganchor{OLP-0188-B010}{end}

\olpseganchor{OLP-0188-B011}{start}
\begin{prop}\ollabel{prop:equiv}
  که د هر~\LR{$!A \in \Theory{M}$} دپاره \LR{$\Struct{N} \models !A$} وي، نو
  \LR{$\Struct{M} \elemequiv \Struct{N}$}۔
\end{prop}
\olpseganchor{OLP-0188-B011}{end}

\olpseganchor{OLP-0188-B012}{start}
\begin{proof}
دا چې د ټولو \LR{$!A \in \Theory{M}$} دپاره \LR{$\Sat{N}{!A}$} دے،
\LR{$\Theory{M} \subseteq \Theory{N}$}۔ که \LR{$\Sat{N}{!A}$} وي، نو
\LR{$\Sat/{N}{\lnot !A}$}؛ ځکه نو \LR{$\lnot !A \notin \Theory{M}$}۔ دا چې
\LR{$\Theory{M}$} بشپړه ده، \LR{$!A \in \Theory{M}$}۔ نو
\LR{$\Theory{N} \subseteq \Theory{M}$}، او له دې
\LR{$\Struct{M} \elemequiv \Struct{N}$} لرو۔
\end{proof}
\olpseganchor{OLP-0188-B012}{end}

\olpseganchor{OLP-0188-B013}{start}
\begin{rem}\ollabel{remark:R}
  \LR{$\Struct{R} = \langle\Real, <\rangle$} وګورئ: دا هغه جوړښت
  دے چې د تعريف ساحه ئې د حقيقي عددونو سټ~\LR{$\Real$} او ژبه ئې
  يوازې يو دوه ځاييز پريديکات لري چې پر حقيقي عددونو د~\LR{$<$}
  اړيکې په توګه تفسير شوے۔ څرګنده ده چې \LR{$\Struct{R}$}
  د شمېر نۀ وړ دے؛ خو دا چې \LR{$\Theory{R}$} څرګنده سازګاره ده، د
  لوېنهايم--سکولم د قضيې له مخې د شمېر وړ مدل لري، مثلاً
  \LR{$\Struct{S}$}، او د~\olref{prop:equiv} له مخې \LR{$\Struct{R}
  \equiv \Struct{S}$}۔ سربېره پر دې، دا چې \LR{$\Struct{R}$} او
  \LR{$\Struct{S}$} آيزومورف نۀ دي، ښيي چې د~\olref[iso]{thm:isom} معکوس
  په عام ډول نۀ صدق کوي۔
\end{rem}
\olpseganchor{OLP-0188-B013}{end}

\olpseganchor{OLP-0189-B004}{start}
\olfileid{mod}{bas}{pis}
\section{جزوي آيزومورفيزمونه}
\olpseganchor{OLP-0189-B004}{end}

\olpseganchor{OLP-0189-B005}{start}
\begin{defn}
  دوه جوړښتونه \LR{$\Struct{M}$} او \LR{$\Struct{N}$} راکړل شي۔
  له~\LR{$\Struct{M}$} څخه~\LR{$\Struct{N}$} ته \emph{جزوي آيزومورفيزم} يوه
  متناهي جزوي تابع~\LR{$p$} ده چې ارګومېنټونه ئې په~\LR{$\Domain M$} کښې او
  قيمتونه ئې په~\LR{$\Domain N$} کښې وي، او د~\olref[iso]{defn:isomorphism}
  د آيزومورفيزم شرطونه پر خپله د تعريف ساحه پوره کوي:
  \begin{enumerate}
  \item \LR{$p$} يو پر يو ده؛
  \item د هر ثابت سمبول~\LR{$c$} دپاره: که \LR{$p(\Assign{c}{M})$} تعريف
    شوے وي، نو \LR{$p(\Assign{c}{M}) = \Assign{c}{N}$}؛
  \item د هر \LR{$n$}-ځاييز پريديکات~\LR{$P$} دپاره: که \LR{$a_1$}،
    \dots،~\LR{$a_n$} د~\LR{$p$} د تعريف په ساحه کښې وي، نو
    \LR{$\langle a_1, \dots, a_n\rangle \in \Assign P M$} هله او يوازې
    هله چې \LR{$\langle p(a_1), \dots, p(a_n) \rangle \in \Assign P N$}؛
  \item د هر \LR{$n$}-ځاييز تابع~\LR{$f$} دپاره: که \LR{$a_1$}،
    \dots،~\LR{$a_n$} د~\LR{$p$} د تعريف په ساحه کښې وي، نو
    \LR{$p(\Assign f M (a_1, \dots,a_n)) =
    \Assign f N (p(a_1), \dots, p(a_n))$}۔
  \end{enumerate}
  د~\LR{$p$} متناهي والی په دې معنا دے چې \LR{$\dom{p}$} متناهي ده۔
\end{defn}
\olpseganchor{OLP-0189-B005}{end}

\olpseganchor{OLP-0189-B006}{start}
پام وکړئ چې تشه تابع~\LR{$\emptyset$} تل د هر دوو جوړښتونه تر منځ
يو جزوي آيزومورفيزم وي۔
\olpseganchor{OLP-0189-B006}{end}

\olpseganchor{OLP-0189-B007}{start}
\begin{defn}\ollabel{defn:partialisom}
  دوه جوړښتونه \LR{$\Struct{M}$} او \LR{$\Struct{N}$}
  \emph{په جزوي ډول آيزومورف} دي، چې \LR{$\Struct{M} \iso[p]
  \Struct{N}$} ئې ليکو، هله او يوازې هله چې د~\LR{$\Struct{M}$} او
  \LR{$\Struct{N}$} تر منځ د جزوي آيزومورفيزمونو يو ناتش سټ~\LR{$I$} وي چې د
  \emph{مخ او شا} خاصيت پوره کوي:
  \begin{enumerate}
  \item (\emph{مخ}) د هر \LR{$p \in I$} او \LR{$a \in \Domain M$} دپاره داسې
    \LR{$q \in I$} شته چې \LR{$p \subseteq q$} او~\LR{$a$} د~\LR{$q$} د تعريف په ساحه
    کښې وي؛
  \item (\emph{شا}) د هر \LR{$p \in I$} او \LR{$b \in \Domain N$} دپاره داسې
    \LR{$q \in I$} شته چې \LR{$p \subseteq q$} او~\LR{$b$} د~\LR{$q$} د قيمتونو په ساحه
    کښې وي۔
  \end{enumerate}
\end{defn}
\olpseganchor{OLP-0189-B007}{end}

\olpseganchor{OLP-0189-B008}{start}
\begin{thm}\ollabel{thm:p-isom1}
  که \LR{$\Struct{M} \iso[p] \Struct{N}$} وي او \LR{$\Struct{M}$} او
  \LR{$\Struct{N}$} د شمېر وړ وي، نو \LR{$\Struct{M} \iso \Struct{N}$}۔
\end{thm}
\olpseganchor{OLP-0189-B008}{end}

\olpseganchor{OLP-0189-B009}{start}
\begin{proof}
  دا چې \LR{$\Struct{M}$} او \LR{$\Struct{N}$} د شمېر وړ دي، فرض کړئ
  \LR{$\Domain{M} = \{a_0, a_1, \ldots \}$} او
  \LR{$\Domain{N} = \{b_0, b_1, \ldots \}$}۔ له يوه خپل سري
  \LR{$p_0 \in I$} څخه په پيل، د جزوي آيزومورفيزمونو زياتېدونکې لړۍ
  \LR{$p_0 \subseteq p_1 \subseteq p_2 \subseteq \cdots$} داسې تعريفوو:
  \begin{enumerate}
  \item که \LR{$n+1$} طاق وي، يعنې \LR{$n = 2r$}، نو د مخ خاصيت په کارولو
    داسې \LR{$p_{n+1} \in I$} ومومئ چې \LR{$p_n \subseteq p_{n+1}$} او~\LR{$a_r$}
    د~\LR{$p_{n+1}$} د تعريف په ساحه کښې وي؛
  \item که \LR{$n+1$} جفت وي، يعنې \LR{$n+1 =2r$}، نو د شا خاصيت په کارولو
    داسې \LR{$p_{n+1} \in I$} ومومئ چې \LR{$p_n \subseteq p_{n+1}$} او~\LR{$b_r$}
    د~\LR{$p_{n+1}$} د قيمتونو په ساحه کښې وي۔
  \end{enumerate}
اوس که
\[
p = \bigcup_{n\ge 0} p_n,
\]
کېږدو، نو \LR{$p$} د~\LR{$\Struct{M}$} او~\LR{$\Struct{N}$} تر منځ يو آيزومورفيزم
دے۔
\end{proof}
\olpseganchor{OLP-0189-B009}{end}

\olpseganchor{OLP-0189-B010}{start}
\begin{prob}
  په تفصيل وښيئ چې په~\olref[mod][bas][pis]{thm:p-isom1} کښې تعريف
  شوے~\LR{$p$} په رښتيا يو آيزومورفيزم دے۔
\end{prob}
\olpseganchor{OLP-0189-B010}{end}

\olpseganchor{OLP-0189-B011}{start}
\begin{thm}\ollabel{thm:p-isom2}
  فرض کړئ \LR{$\Struct{M}$} او \LR{$\Struct{N}$} د يوې سوچه اړيکيزې
  ژبه جوړښتونه دي (داسې ژبه چې يوازې
  پريديکاتونه لري او تابعې يا ثابتونه نۀ لري)۔ بيا که
  \LR{$\Struct{M} \iso[p] \Struct{N}$} وي، نو \LR{$\Struct{M} \elemequiv
  \Struct{N}$} هم وي۔
\end{thm}
\olpseganchor{OLP-0189-B011}{end}

\olpseganchor{OLP-0189-B012}{start}
\begin{proof}
  پر فارمولونه په استقرا ښودل کېږي چې که \LR{$a_1$}، \dots،~\LR{$a_n$} او
  \LR{$b_1$}، \dots،~\LR{$b_n$} داسې وي چې يو جزوي آيزومورفيزم~\LR{$p$} هر~\LR{$a_i$}
  پر~\LR{$b_i$} وړي، او \LR{$s_1(x_i) =a_i$} او \LR{$s_2(x_i) =b_i$} وي (د
  \LR{$i =1$}، \dots،~\LR{$n$} دپاره)، نو \LR{$\Sat{M}{!A}[s_1]$} هله او يوازې هله
  چې \LR{$\Sat{N}{!A}[s_2]$} وي۔ د~\LR{$n=0$} حالت
  \LR{$\Struct{M} \elemequiv \Struct{N}$} راکوي۔
\end{proof}
\olpseganchor{OLP-0189-B012}{end}

\olpseganchor{OLP-0189-B013}{start}
\begin{rem}
که تابعې موجودې وي، مخکينۍ پايله بيا هم رښتونې ده، خو د هغه
آيزومورفيزم څېړل پکار دي چې~\LR{$p$} ئې د~\LR{$a_1$}، \dots،~\LR{$a_n$} له خوا د
\LR{$\Struct{M}$} پر توليد شوي فرعي جوړښت او د~\LR{$b_1$}،
\dots،~\LR{$b_n$} له خوا د~\LR{$\Struct{N}$} پر توليد شوي فرعي جوړښت
رامنځته کوي۔
\end{rem}
\olpseganchor{OLP-0189-B013}{end}

\olpseganchor{OLP-0189-B014}{start}
مخکينۍ پايله په پړاوونو هم ``ماتېدے'' شي: په يوه فارمول کښې د
يو بل دننه کميت ټاکونکو د شمېر او د اړوندو جزوي آيزومورفيزمونو د
غځېدو د ځلونو تر منځ اړيکه ټينګوو۔
\olpseganchor{OLP-0189-B014}{end}

\olpseganchor{OLP-0189-B015}{start}
\begin{defn}
  د هر فارمول~\LR{$!A$} دپاره، د~\LR{$!A$} \emph{د کميت ټاکونکو رتبه}، چې
  \LR{$\QuantRank{!A} \in \Nat$} ئې ليکو، په بازګشتي ډول په~\LR{$!A$} کښې د
  يو بل دننه کميت ټاکونکو تر ټولو لوړه شمېره تعريفېږي۔ دوه
  جوړښتونه \LR{$\Struct{M}$} او \LR{$\Struct{N}$} \emph{\LR{$n$}-هم‌ارز} دي،
  چې \LR{$\Struct{M} \elemequiv[n] \Struct{N}$} ئې ليکو، که د کميت
  ټاکونکو د مرتبې تر~\LR{$n$} پورې پر ټولو جملو سره سلا وي۔
\end{defn}
\olpseganchor{OLP-0189-B015}{end}

\olpseganchor{OLP-0189-B016}{start}
\begin{prop}\ollabel{prop:qr-finite}
  فرض کړئ \LR{$\Lang{L}$} يوه متناهي سوچه اړيکيزه ژبه ده، يعنې
  داسې ژبه چې متناهي شمېر پريديکاتونه او ثابت سمبولونه
  لري او تابعې نۀ لري۔ بيا د هر \LR{$n \in \Nat$} دپاره، له
  منطقي هم‌ارزۍ پرته، د~\LR{$\Lang{L}$} په ژبه کښې يوازې متناهي ډېرې داسې
  د لومړۍ درجې تړلي فارمولونه شته چې د کميت ټاکونکو رتبه ئې تر~\LR{$n$}
  زياته نۀ وي۔
\end{prop}
\olpseganchor{OLP-0189-B016}{end}

\olpseganchor{OLP-0189-B017}{start}
\begin{proof}
  پر~\LR{$n$} په استقرا۔
\end{proof}
\olpseganchor{OLP-0189-B017}{end}

\olpseganchor{OLP-0189-B018}{start}
\begin{defn}
  يو جوړښت~\LR{$\Struct{M}$} راکړل شي۔ \LR{$\Domain M^{<\omega}$} د
  \LR{$\Domain{M}$} پر غړو د ټولو متناهي لړيو سټ ونيسئ۔ په~\LR{$\mathbf{a}, \mathbf{b}, \mathbf{c}, \ldots$} د غړو متناهي لړۍ ښيو۔ که
  \LR{$\mathbf{a} \in \Domain{M}^{<\omega}$} او \LR{$a \in \Domain{M}$} وي، نو
  \LR{$\mathbf{a}a$} له~\LR{$a$} سره د~\LR{$\mathbf{a}$} \emph{نښلون} ښيي۔
\end{defn}
\olpseganchor{OLP-0189-B018}{end}

\olpseganchor{OLP-0189-B019}{start}
\begin{defn}
  جوړښتونه \LR{$\Struct{M}$} او \LR{$\Struct{N}$} راکړل شي۔ د يو شان
  اوږدوالي لړيو تر منځ اړيکې
  \LR{$I_n \subseteq \Domain M^{<\omega} \times \Domain N^{<\omega}$}
  پر~\LR{$n$} په بازګښت داسې تعريفوو:
   \begin{enumerate}
   \item \LR{$I_0(\mathbf{a},\mathbf{b})$} هله او يوازې هله چې
     \LR{$\mathbf{a}$} او \LR{$\mathbf{b}$} په~\LR{$\Struct{M}$} او~\LR{$\Struct{N}$}
     کښې يو شان اټومي فارمولونه پوره کړي؛ يعنې که د دواړو لړيو
     اوږدوالی~\LR{$k$} وي، \LR{$s_1(x_i) = a_i$} او \LR{$s_2(x_i) = b_i$} وي، او
     \LR{$!A$} داسې اټومي فارمول وي چې ټول متغيرونه ئې په
     \LR{$x_1$}، \dots،~\LR{$x_k$} کښې وي، نو \LR{$\Sat{M}{!A}[s_1]$} هله او يوازې
     هله چې \LR{$\Sat{N}{!A}[s_2]$} وي۔
   \item \LR{$I_{n+1} (\mathbf{a},\mathbf{b})$} هله او يوازې هله چې د هر
     \LR{$a\in \Domain M$} دپاره داسې \LR{$b\in \Domain N$} وي چې
     \LR{$I_n(\mathbf{a}a,\mathbf{b}b)$}، او برعکس۔
   \end{enumerate}
  \textbf{د ژباړې د نسخې سمون (OLMOD-003):} سرچينې په لومړي شرط کښې
  د لړيو د اوږدوالي شاخص د اړيکو د بازګشتي شاخص سره يو شان نيولے و؛
  دلته اوږدوالی په بېل~\LR{$k$} ښودل شوے، څو لومړنۍ اړيکه د هر اوږدوالي د
  لړيو اټومي فارمولونه پرتله کړي۔
\end{defn}
\olpseganchor{OLP-0189-B019}{end}


\olpseganchor{OLP-0189-B020}{start}
\begin{defn}
  \LR{$\Struct{M} \approx_n \Struct{N}$} ليکو که
  \LR{$I_n(\emptyseq,\emptyseq)$} د~\LR{$\Struct{M}$} او~\LR{$\Struct{N}$} دپاره
  صدق وکړي (دلته \LR{$\emptyseq$} تشه لړۍ ده)۔
\end{defn}
\olpseganchor{OLP-0189-B020}{end}

\olpseganchor{OLP-0189-B021}{start}
\begin{thm}\ollabel{thm:b-n-f}
  فرض کړئ \LR{$\Lang{L}$} يوه سوچه اړيکيزه ژبه ده۔ بيا
  \LR{$I_n(\mathbf{a},\mathbf{b})$} دا لازموي چې د هر داسې~\LR{$!A$} دپاره چې
  \LR{$\QuantRank{!A} \le n$} وي، \LR{$\Sat{M}{!A}[\mathbf{a}]$} هله او يوازې
  هله چې \LR{$\Sat{N}{!A}[\mathbf{b}]$} وي (دلته هم، \LR{$\mathbf{a}$} هغه
  وخت~\LR{$!A$} پوره کوي چې هره داسې~\LR{$s$} چې \LR{$s(x_i) = a_i$} وي، \LR{$!A$} پوره
  کړي)۔ سربېره پر دې، که \LR{$\Lang{L}$} متناهي وي، معکوس هم صدق کوي۔
\end{thm}
\olpseganchor{OLP-0189-B021}{end}

\olpseganchor{OLP-0189-B022}{start}
\begin{proof}
  دا چې \LR{$I_n(\mathbf{a},\mathbf{b})$} لازموي چې \LR{$\mathbf{a}$} او
  \LR{$\mathbf{b}$} د کميت ټاکونکو تر~\LR{$n$} پورې يو شان فارمولونه پوره
  کوي، پر~\LR{$!A$} په اسانه استقرا ثابتېږي۔ د معکوس دپاره پر~\LR{$n$} استقرا
  کوو او \olref{prop:qr-finite} کاروو، چې تضمينوي د هر~\LR{$n$} دپاره د
  هغې مرتبې تر ډېره متناهي ډېر يو له بله نۀ هم‌ارز فارمولونه شته۔
\olpseganchor{OLP-0189-B022}{end}

\olpseganchor{OLP-0189-B023}{start}
  د~\LR{$n=0$} دپاره، دا فرض چې \LR{$\mathbf{a}$} او \LR{$\mathbf{b}$} يو شان بې
  کميت ټاکونکو فارمولونه پوره کوي، راکوي چې يو شان اټومي فارمولونه
  پوره کوي؛ نو \LR{$I_0(\mathbf{a},\mathbf{b})$}۔
\olpseganchor{OLP-0189-B023}{end}

\olpseganchor{OLP-0189-B024}{start}
  د~\LR{$n+1$} په حالت کښې فرض کړئ \LR{$\mathbf{a}$} او \LR{$\mathbf{b}$} د کميت
  ټاکونکو تر~\LR{$n+1$} پورې يو شان فارمولونه پوره کوي۔ د
  \LR{$I_{n+1}(\mathbf{a},\mathbf{b})$} د ښودلو دپاره کافي ده چې د هر
  \LR{$a \in \Domain M$} دپاره داسې \LR{$b \in \Domain N$} وښيو چې
  \LR{$I_n(\mathbf{a}a,\mathbf{b}b)$}؛ او د استقرا له فرضه بيا کافي ده
  وښيو چې د هر \LR{$a \in \Domain M$} دپاره داسې \LR{$b \in \Domain N$} شته
  چې \LR{$\mathbf{a}a$} او \LR{$\mathbf{b}b$} د کميت ټاکونکو تر~\LR{$n$} پورې يو
  شان فارمولونه پوره کوي۔
\olpseganchor{OLP-0189-B024}{end}

\olpseganchor{OLP-0189-B025}{start}
  يو \LR{$a \in \Domain M$} راکړل شي۔ \LR{$!T^a_n$} د رتبې تر~\LR{$n$} پورې د هغو
  فارمولونه~\LR{$!B(x,\mathbf{y})$} سټ ونيسئ چې \LR{$\mathbf{a}a$} ئې
  په~\LR{$\Struct{M}$} کښې پوره کوي؛ \LR{$!T^a_n$} له منطقي هم‌ارزۍ پرته
  متناهي دے، نو د غړو د تړون په اخيستو ئې يو د لومړۍ درجې
  فارمول ګڼلے شو۔ له دې پايله کېږي چې \LR{$\mathbf{a}$}
  \LR{$\lexists[x][!T^a_n(x,\mathbf{y})]$} پوره کوي، چې د کميت ټاکونکو
  رتبه ئې تر~\LR{$n+1$} زياته نۀ ده۔ د فرض له مخې \LR{$\mathbf{b}$} هماغه
  فارمول په~\LR{$\Struct{N}$} کښې پوره کوي؛ نو داسې
  \LR{$b \in \Domain N$} شته چې \LR{$\mathbf{b}b$}، \LR{$!T^a_n$} پوره کوي۔ په
  ځانګړي ډول، \LR{$\mathbf{b}b$} د کميت ټاکونکو تر~\LR{$n$} پورې هماغه
  فارمولونه پوره کوي چې \LR{$\mathbf{a}a$} ئې پوره کوي۔ په ورته ډول
  ښودل کېږي چې د هر \LR{$b \in \Domain N$} دپاره داسې \LR{$a\in \Domain M$}
  شته چې \LR{$\mathbf{a}a$} او \LR{$\mathbf{b}b$} د کميت ټاکونکو تر~\LR{$n$} پورې
  يو شان فارمولونه پوره کوي؛ په دې سره ثبوت بشپړېږي۔
\end{proof}
\olpseganchor{OLP-0189-B025}{end}

\olpseganchor{OLP-0189-B026}{start}
\begin{cor}\ollabel{cor:b-n-f}
  که \LR{$\Struct{M}$} او \LR{$\Struct{N}$} په يوه متناهي ژبه کښې
  سوچه اړيکيز جوړښتونه وي، نو \LR{$\Struct{M} \approx_n\Struct{N}$}
  هله او يوازې هله چې \LR{$\Struct{M} \elemequiv[n] \Struct{N}$} وي۔ په
  ځانګړي ډول، \LR{$\Struct{M} \elemequiv \Struct{N}$} هله او يوازې هله
  چې د هر~\LR{$n$} دپاره \LR{$\Struct{M} \approx_n \Struct{N}$} وي۔
\end{cor}
\olpseganchor{OLP-0189-B026}{end}

\olpseganchor{OLP-0190-B004}{start}
\olfileid{mod}{bas}{dlo}
\section{بې له سرحدي ټکو ګڼ خطي ترتيبونه}
\olpseganchor{OLP-0190-B004}{end}

\olpseganchor{OLP-0190-B005}{start}
\begin{defn}
  \emph{بې له سرحدي ټکو ګڼ خطي ترتيب} د هغې ژبه يو جوړښت
  \LR{$\Struct{M}$} دے چې يوازينے دوه ځاييز پريديکات~\LR{$<$} لري او لاندې
  جملې پوره کوي:
  \begin{enumerate}
  \item \LR{$\lforall[x][\lnot x < x]$}؛
  \item \LR{$\lforall[x][\lforall[y][\lforall[z][(x < y \lif (y < z \lif x
    <z ))]]]$}؛
  \item \LR{$\lforall[x][\lforall[y][(x< y \lor \eq[x][y] \lor y < x)]]$}؛
  \item \LR{$\lforall[x][\lexists[y][x < y]]$}؛
  \item \LR{$\lforall[x][\lexists[y][y < x]]$}؛
  \item \LR{$\lforall[x][\lforall[y][(x < y \lif \lexists[z][(x < z \land
        z < y))]]]$}۔
 \end{enumerate}
\end{defn}
\olpseganchor{OLP-0190-B005}{end}

\olpseganchor{OLP-0190-B006}{start}
\begin{thm}\ollabel{thm:cantorQ}
  هر دوه د شمېر وړ، بې له سرحدي ټکو ګڼ خطي ترتيبونه آيزومورف دي۔
\end{thm}
\olpseganchor{OLP-0190-B006}{end}

\olpseganchor{OLP-0190-B007}{start}
\begin{proof}
  فرض کړئ \LR{$\Struct{M_1}$} او \LR{$\Struct{M_2}$} د شمېر وړ، بې له
  سرحدي ټکو ګڼ خطي ترتيبونه دي، \LR{${<_1} = \Assign{<}{M_1}$} او
  \LR{${<_2} = \Assign{<}{M_2}$}، او \LR{$\PIso{I}$} د دوى تر منځ د ټولو جزوي
  آيزومورفيزمونو سټ دے۔ \LR{$\PIso{I}$} تش نۀ دے، ځکه لږ تر لږه
  \LR{$\emptyset \in \PIso{I}$}۔ ښيو چې \LR{$\PIso{I}$} د مخ او شا خاصيت پوره
  کوي۔ بيا \LR{$\Struct{M_1} \iso[p] \Struct{M_2}$}، او قضيه
  له~\olref[pis]{thm:p-isom1} څخه پايله کېږي۔
\olpseganchor{OLP-0190-B007}{end}

\olpseganchor{OLP-0190-B008}{start}
  د دې د ښودلو دپاره چې \LR{$\PIso{I}$} د مخ خاصيت پوره کوي، فرض کړئ
  \LR{$p \in \PIso{I}$} او د~\LR{$i = 1$}، \dots،~\LR{$n$} دپاره
  \LR{$p(a_i) = b_i$}، او د عموميت له زيان پرته فرض کړئ
  \LR{$a_1 <_1 a_2 <_1 \cdots <_1 a_n$}۔ يو \LR{$a \in \Domain{M_1}$} چې
  راکړل شي، \LR{$b \in \Domain{M_2}$} داسې ومومئ:
  \begin{enumerate}
  \item که \LR{$a <_1 a_1$} وي، \LR{$b \in \Domain{M_2}$} داسې ونيسئ چې
    \LR{$b <_2 b_1$}؛
  \item که \LR{$a_n <_1 a$} وي، \LR{$b \in \Domain{M_2}$} داسې ونيسئ چې
    \LR{$b_n <_2 b$}؛
 \item که د کوم~\LR{$i$} دپاره \LR{$a_i <_1 a <_1 a_{i+1}$} وي، نو
   \LR{$b \in \Domain{M_2}$} داسې ونيسئ چې \LR{$b_i <_2 b <_2 b_{i+1}$}۔
  \end{enumerate}
  د مطلوب خاصيت لرونکے~\LR{$b$} تل موندل کېدے شي، ځکه \LR{$\Struct{M_2}$} بې
  له سرحدي ټکو ګڼ خطي ترتيب دے۔ \LR{$q = p \cup \{ \langle a, b \rangle
  \}$} داسې تعريف کړئ چې \LR{$q \in \PIso{I}$} د~\LR{$p$} مطلوبه غځونه وي۔ په
  دې سره د مخ خاصيت ثابت شو۔ د شا خاصيت هم ورته دے۔ نو
  \LR{$\Struct{M_1} \iso[p] \Struct{M_2}$}؛ او د
  \olref[pis]{thm:p-isom1} له مخې \LR{$\Struct{M_1} \iso \Struct{M_2}$}۔
\end{proof}
\olpseganchor{OLP-0190-B008}{end}

\olpseganchor{OLP-0190-B009}{start}
\begin{prob}
  د~\olref[mod][bas][dlo]{thm:cantorQ} ثبوت په دې ښودلو بشپړ کړئ چې
  \LR{$\PIso{I}$} د شا خاصيت پوره کوي۔
\end{prob}
\olpseganchor{OLP-0190-B009}{end}

\olpseganchor{OLP-0190-B010}{start}
\begin{rem}
  \LR{$\Struct{S}$}، بې له سرحدي ټکو هر د شمېر وړ ګڼ خطي ترتيب،
  ونيسئ۔ بيا (د~\olref{thm:cantorQ} له مخې) \LR{$\Struct{S} \iso
  \Struct{Q}$}، چې \LR{$\Struct{Q} = (\Rat, <)$} هغه د شمېر وړ ګڼ
  خطي ترتيب دے چې د تعريف ساحه ئې د ناطقو عددونو سټ~\LR{$\Rat$} دے۔ اوس د
  \olref[thm]{remark:R} جوړښت~\LR{$\Struct{R} = (\Real, <)$} بيا وګورئ۔ مو
  وليدل چې داسې د شمېر وړ جوړښت~\LR{$\Struct{S}$} شته چې
  \LR{$\Struct{R} \elemequiv \Struct{S}$}۔ خو \LR{$\Struct{S}$}، بې له سرحدي
  ټکو يو د شمېر وړ ګڼ خطي ترتيب دے؛ نو له جوړښت~\LR{$\Struct{Q}$}
  سره آيزومورف (او ځکه ابتدايي هم‌ارز) جوړښت دے۔ د ابتدايي هم‌ارزۍ له
  انتقاليت څخه \LR{$\Struct{R} \elemequiv \Struct{Q}$} پايله کېږي۔ (دا مو
  په مستقيم ډول هم ښودلے شول: د هماغه مخ او شا دليل په وسيله مو
  \LR{$\Struct{R} \iso[p] \Struct{Q}$} ټينګولے شو۔)
\end{rem}
\olpseganchor{OLP-0190-B010}{end}

\olpseganchor{OLP-0191-B004}{start}
\olchapter{mod}{mar}{د حساب مدلونه}
\olpseganchor{OLP-0191-B004}{end}

\olpseganchor{OLP-0192-B004}{start}
\olfileid{mod}{mar}{int}
\section{سريزه}
\olpseganchor{OLP-0192-B004}{end}

\olpseganchor{OLP-0192-B005}{start}
د حساب \emph{معياري مدل} هغه جوړښت~\LR{$\Struct{N}$} دے چې
\LR{$\Domain{N} = \Nat$} وي او \LR{$\Obj{0}$}، \LR{$\prime$}، \LR{$+$}، \LR{$\times$} او \LR{$<$}
هغسې تفسير شوي وي لکه تمه چې ترې کېږي۔ يعنې \LR{$\Obj{0}$}، \LR{$0$} دے،
\LR{$\prime$} د جانشين تابع ده، \LR{$+$} د جمع په توګه او \LR{$\times$} د ضرب په
توګه د~\LR{$\Nat$} د شمېرو تفسير لري۔ په ځانګړي ډول:
\begin{align*}
  \Assign{\Obj{0}}{N} & = 0\\
  \Assign{\prime}{N}(n) & = n + 1\\
  \Assign{+}{N}(n, m) & = n + m\\
  \Assign{\times}{N}(n, m) & = nm
\end{align*}
البته د~\LR{$\Lang{L_A}$} داسې جوړښتونه هم شته چې تعريف ساحې ئې له~\LR{$\Nat$}
څخه بېلې دي۔ د بېلګې په توګه، \LR{$\Struct{M}$} داسې ونيسئ چې
\LR{$\Domain{M} = \{a\}^*$} وي (د يوازېنۍ نښې~\LR{$a$} متناهي لړۍ؛ يعنې
\LR{$\emptyset$}، \LR{$a$}، \LR{$aa$}، \LR{$aaa$}، \dots)، او تفسيرونه ئې داسې وي:
\begin{align*}
  \Assign{\Obj{0}}{M} & = \emptyset\\
  \Assign{\prime}{M}(s) & = s \concat a\\
  \Assign{+}{M}(n, m) & = a^{n + m}\\
  \Assign{\times}{M}(n, m) & = a^{nm}
\end{align*}
دا دوه جوړښتونه په دې معنا «په اصل کښې يو شان» دي چې يوازينے
توپير ئې د غړي او د تعريف ساحې په څيزونو کښې دے، نه دا چې
د تفسير د تابعو له مخې د د تعريف ساحې غړي له يو بل سره څنګه تړاو لري۔
موږ وايو چې دا دوه جوړښتونه \emph{آيزومورف} دي۔
\olpseganchor{OLP-0192-B005}{end}

\olpseganchor{OLP-0192-B006}{start}
د متناهي شاهد له قضيه څخه په اسانۍ پايله کېږي چې هره تيوري چې په
\LR{$\Struct{N}$} کښې رښتونې وي، داسې مدلونه هم لري چې له~\LR{$\Struct{N}$} سره
آيزومورف نۀ وي۔ دا ډول جوړښتونه \emph{نامعياري} بلل کېږي۔ په معياري
مدل کښې (يعنې په~\LR{$\Struct{N}$} او له هغې سره په ټولو آيزومورفو
جوړښتونه کښې) ټول غړي د معياري عددنښو په ارزښتونو
\LR{$\num{n}$} کښې راټولېږي، يعنې:
\[
\Domain{N} = \Setabs{\Value{\num{n}}{N}}{n \in \Nat}
\]
خو په نامعياري مدلونو کښې داسې نۀ ده: که \LR{$\Struct{M}$} نامعياري وي، نو
لږ تر لږه يو \LR{$x \in \Domain{M}$} شته چې د هېڅ~\LR{$n$} دپاره
\LR{$x \neq \Value{\num{n}}{M}$} وي۔
\olpseganchor{OLP-0192-B006}{end}

\olpseganchor{OLP-0192-B007}{start}
دا نامعياري عناصر په زړه پورې دي: «نامتناهي طبيعي شمېرې» دي۔
خو د هغو شتون د ناتکميلۍ پديدې هم په يوه معنا روښانوي۔ د بېلګې په
توګه د پيانو د حساب د سازګارۍ جمله، \LR{$\OCon[\Th{PA}]$}، يعنې
\LR{$\lnot \lexists[x][\OPrf[\Th{PA}](x, \gn{\lfalse})]$}، راواخلئ۔
څرنګه چې \LR{$\Th{PA}$} نه \LR{$\OCon[\Th{PA}]$} ثابتوي او نه
\LR{$\lnot \OCon[\Th{PA}]$}، نو هره يوه جمله په سازګار ډول له~\LR{$\Th{PA}$}
سره يوځاے کېدے شي۔ څرنګه چې \LR{$\Th{PA}$} سازګاره ده،
\LR{$\Sat{N}{\OCon[\Th{PA}]}$} او په پايله کښې
\LR{$\Sat/{N}{\lnot \OCon[\Th{PA}]}$} لرو۔ نو \LR{$\Struct{N}$} د
\LR{$\Th{PA} \cup \{\lnot \OCon[\Th{PA}]\}$} مدل \emph{نۀ} دے، او ټول
مدلونه ئې بايد نامعياري وي۔ د \LR{$\Th{PA} \cup
\{\lnot \OCon[\Th{PA}]\}$} مدلونه بايد داسې يو غړے ولري چې د
\LR{$\lexists[x][\OPrf[\Th{PA}](x,\gn{\lfalse})]$} د رښتينولۍ شاهد وي، يعنې
له~\LR{$\Th{PA}$} څخه د تناقض د اشتقاق ګوډل عدد وي۔ داسې
غړے معياري کېدے نۀ شي، ځکه د هر~\LR{$n$} دپاره
\LR{$\Th{PA} \Proves \lnot \OPrf[\Th{PA}](\num{n}, \gn{\lfalse})$}۔
\textbf{د ژباړې د نسخې سمون (OLMOD-004):} سرچينې د وجودي جملې په
دويم ځل ليکلو کښې د ثبوت د کوډ ارګومېنټ x غورځولے و؛ دلته د ثبوت د
نسبت دوه ارګومېنټه بڼه بېرته ليکل شوې، لکه د همدې فقرې په لومړي
تعريف او په ورپسې معياري عددنښه کښې چې راغلې ده۔
\olpseganchor{OLP-0192-B007}{end}

\olpseganchor{OLP-0193-B004}{start}
\olfileid{mod}{mar}{stm}
\olsection{د حساب معياري مدلونه}
\olpseganchor{OLP-0193-B004}{end}

\olpseganchor{OLP-0193-B005}{start}
د حساب ژبه~\LR{$\Lang{L_A}$} په څرګند ډول د شمېرو، په ځانګړي ډول د
طبيعي شمېرو، په هکله ده۔ نو «معياري مدل»~\LR{$\Struct{N}$} ځانګړے دے:
همدا هغه مدل دے چې موږ ئې په هکله خبرې کول غواړو۔ خو په منطق کښې
زموږ پام ډېر ځله يوازې جوړښتي ځانګړتياوو ته وي، او هر دوه آيزومورف
جوړښتونه دغه ځانګړتياوې ګډې لري۔ نو لږ پراخ فکر کولے شو او
هر هغه جوړښت چې له~\LR{$\Struct{N}$} سره آيزومورف وي «معياري»
وبولو۔
\olpseganchor{OLP-0193-B005}{end}

\olpseganchor{OLP-0193-B006}{start}
\begin{defn}
د~\LR{$\Lang{L_A}$} يو جوړښت \emph{معياري} دے، که له~\LR{$\Struct{N}$}
سره آيزومورف وي۔
\end{defn}
\olpseganchor{OLP-0193-B006}{end}

\olpseganchor{OLP-0193-B007}{start}
\begin{prop}
\ollabel{prop:standard-domain} که جوړښت~\LR{$\Struct{M}$} معياري
وي، نو تعريف ساحه ئې د معياري عددنښو د ارزښتونو سټ دے، يعنې:
\[
\Domain{M} = \Setabs{\Value{\num{n}}{M}}{n \in \Nat}
\]
\end{prop}
\olpseganchor{OLP-0193-B007}{end}

\olpseganchor{OLP-0193-B008}{start}
\begin{proof}
څرګنده ده چې هر \LR{$\Value{\num{n}}{M} \in \Domain{M}$} دے۔ يوازې دا
ښودل پاتې دي چې هر \LR{$x \in \Domain{M}$} د کوم~\LR{$n$} دپاره
\LR{$\Value{\num{n}}{M}$} دے۔ څرنګه چې \LR{$\Struct{M}$} معياري دے، له
\LR{$\Struct{N}$} سره آيزومورف دے۔ فرض کړئ \LR{$g\colon \Nat \to \Domain{M}$}
آيزومورفيزم وي۔ بيا \LR{$g(n) = g(\Value{\num{n}}{N}) =
\Value{\num{n}}{M}$}۔ خو د هر \LR{$x \in \Domain{M}$} دپاره داسې~\LR{$n \in
\Nat$} شته چې \LR{$g(n) = x$} وي، ځکه \LR{$g$} پر هدف سټ بشپړه ده۔
\end{proof}
\olpseganchor{OLP-0193-B008}{end}

\olpseganchor{OLP-0193-B009}{start}
\begin{explain}
که د~\LR{$\Lang{L_A}$} يو جوړښت~\LR{$\Struct{M}$} معياري وي، د~د تعريف ساحه
ټول غړي ئې د معياري عددنښو \LR{$\num{0}$}، \LR{$\num{1}$}، \LR{$\num{2}$}،
\dots په نوم يادېدے شي، يعنې د \LR{$\Obj{0}$}، \LR{$\Obj{0}'$}،
\LR{$\Obj{0}''$} او داسې نورو ترمونو په وسيله۔ البته دا معنا نۀ لري چې د
\LR{$\Domain{M}$} غړي پخپله شمېرې دي؛ يوازې دومره ده چې د هغو
په نښه کولے شو، لکه څنګه چې په~\LR{$\Domain{N}$} کښې شمېرې نښه کوو۔
\end{explain}
\olpseganchor{OLP-0193-B009}{end}

\olpseganchor{OLP-0193-B010}{start}
\begin{prob}
وښيئ چې د~\olref[mod][mar][stm]{prop:standard-domain} معکوس ناسم دے؛
يعنې داسې يو جوړښت~\LR{$\Struct{M}$} ورکړئ چې
\LR{$\Domain{M} = \Setabs{\Value{\num{n}}{M}}{n \in \Nat}$} وي، خو له
\LR{$\Struct{N}$} سره آيزومورف نۀ وي۔
\end{prob}
\olpseganchor{OLP-0193-B010}{end}

\olpseganchor{OLP-0193-B011}{start}
\begin{prop}
\ollabel{prop:thq-standard}
که \LR{$\Sat{M}{\Th{Q}}$} او \LR{$\Domain{M} =
\Setabs{\Value{\num{n}}{M}}{n \in \Nat}$} وي، نو \LR{$\Struct{M}$} معياري دے۔
\end{prop}
\olpseganchor{OLP-0193-B011}{end}

\olpseganchor{OLP-0193-B012}{start}
\begin{proof}
بايد وښيو چې \LR{$\Struct{M}$} له~\LR{$\Struct{N}$} سره آيزومورف دے۔ تابع
\LR{$g\colon \Nat \to \Domain{M}$} داسې ونيسئ: \LR{$g(n) =
\Value{\num{n}}{M}$}۔ د فرض له مخې \LR{$g$} پر هدف سټ بشپړه ده۔ دا
يو پر يو هم ده: هر کله چې \LR{$n \neq m$} وي، \LR{$\Th{Q} \Proves
\eq/[\num{n}][\num{m}]$}۔ نو له \LR{$\Sat{M}{\Th{Q}}$} څخه
\LR{$\Sat{M}{\eq/[\num{n}][\num{m}]}$}، هر کله چې \LR{$n \neq m$}۔ په پايله
کښې \LR{$n \neq m$} وي، نو \LR{$\Value{\num{n}}{M} \neq
\Value{\num{m}}{M}$}، يعنې \LR{$g(n) \neq g(m)$}۔
\olpseganchor{OLP-0193-B012}{end}

\olpseganchor{OLP-0193-B013}{start}
اوس بايد دا هم وڅېړو چې \LR{$g$} آيزومورفيزم دے۔
\begin{enumerate}
\item \LR{$g(\Assign{\Obj{0}}{N}) = g(0)$}، ځکه \LR{$\Assign{\Obj{0}}{N} = 0$}۔
  د~\LR{$g$} له تعريفه \LR{$g(0) = \Value{\num{0}}{M}$}۔ خو \LR{$\num{0}$} هماغه
  \LR{$\Obj{0}$} دے، او د داسې ترم ارزښت چې ثابت سمبول وي، هغه څه دے
  چې جوړښت ئې دغه ثابت سمبول ته ټاکي، يعنې \LR{$\Value{\Obj{0}}{M} =
  \Assign{\Obj{0}}{M}$}۔ نو \LR{$g(\Assign{\Obj{0}}{N}) =
  \Assign{\Obj{0}}{M}$}، لکه اړين چې و۔
\item \LR{$g(\Assign{\prime}{N}(n)) = g(n+1)$}، ځکه په~\LR{$\Struct{N}$} کښې
  \LR{$\prime$} پر~\LR{$\Nat$} د جانشين تابع ده۔ بيا د~\LR{$g$} له تعريفه
  \LR{$g(n+1) = \Value{\num{n+1}}{M}$}۔ خو \LR{$\num{n+1}$} له~\LR{$\num{n}'$}
  سره يو ترم دے، نو \LR{$\Value{\num{n+1}}{M} =
  \Value{\num{n}'}{M}$}۔ د ارزښت د تابعې له تعريفه دا
  \LR{$= \Assign{\prime}{M}(\Value{\num{n}}{M})$} دے۔ څرنګه چې
  \LR{$\Value{\num{n}}{M} = g(n)$}، نو
  \LR{$g(\Assign{\prime}{N}(n)) = \Assign{\prime}{M}(g(n))$}۔
\item \LR{$g(\Assign{+}{N}(n,m)) = g(n+m)$}، ځکه په~\LR{$\Struct{N}$} کښې \LR{$+$}
  د~\LR{$\Nat$} د جمع تابع ده۔ بيا د~\LR{$g$} له تعريفه
\LR{$g(n+m) = \Value{\num{n+m}}{M}$}۔
  خو \LR{$\Th{Q} \Proves \num{n+m} = (\num{n} + \num{m})$}، نو
  \LR{$\Value{\num{n+m}}{M} = \Value{\num{n}+\num{m}}{M}$}۔ د ارزښت د
  تابعې له تعريفه دا
  \LR{$= \Assign{+}{M}(\Value{\num{n}}{M},\Value{\num{m}}{M})$} دے۔
  څرنګه چې \LR{$\Value{\num{n}}{M} = g(n)$} او
  \LR{$\Value{\num{m}}{M} = g(m)$}، نو
  \LR{$g(\Assign{+}{N}(n, m)) = \Assign{+}{M}(g(n), g(m))$}۔
\item \LR{$g(\Assign{\times}{N}(n, m)) = \Assign{\times}{M}(g(n), g(m))$}:
  تمرين۔
\item \LR{$\tuple{n,m} \in \Assign{<}{N}$} هله او يوازې هله چې \LR{$n < m$}۔
  که \LR{$n < m$} وي، نو \LR{$\Th{Q} \Proves \num{n} < \num{m}$} او
  \LR{$\Sat{M}{\num{n} < \num{m}}$} هم لرو۔ نو
  \LR{$\tuple{\Value{\num{n}}{M}, \Value{\num{m}}{M}} \in
  \Assign{<}{M}$}، يعنې \LR{$\tuple{g(n), g(m)} \in \Assign{<}{M}$}۔
  که \LR{$n \not< m$} وي، نو \LR{$\Th{Q} \Proves \lnot \num{n} < \num{m}$} او
  په پايله کښې \LR{$\Sat/{M}{\num{n} < \num{m}}$}۔ نو، لکه مخکې،
  \LR{$\tuple{g(n), g(m)} \notin \Assign{<}{M}$}۔ په ګډه:
  \LR{$\tuple{n,m} \in \Assign{<}{N}$} هله او يوازې هله چې
  \LR{$\tuple{g(n), g(m)} \in \Assign{<}{M}$}۔
\end{enumerate}
\end{proof}
\olpseganchor{OLP-0193-B013}{end}

\olpseganchor{OLP-0193-B014}{start}
\begin{explain}
تابع~\LR{$g$} د~\LR{$\Lang{L_A}$} د هر بل جوړښت~\LR{$\Struct{M}$} د تعريف
ساحې ته له~\LR{$\Nat$} څخه د نقشې تر ټولو څرګنده لار ده، ځکه هر داسې
\LR{$\Struct{M}$} هغه غړي لري چې \LR{$\num{0}$}، \LR{$\num{1}$}،
\LR{$\num{2}$} او داسې نورو په نومونو يادېږي۔ نو د حيرانتيا خبره نۀ ده چې
که \LR{$\Struct{M}$} د~\LR{$\num{n}$} په باب لږ تر لږه ځينې بنسټيزې ويناوې
هماغسې رښتونې کړي لکه~\LR{$\Struct{N}$} او \LR{$g$} هم دوه اړخيزه يو په يو
وي، نو \LR{$g$} آيزومورفيزم ګرځي۔ په حقيقت کښې که \LR{$\Domain{M}$} له هغو
غړي پرته چې \LR{$\num{n}$} ئې نوموي بل څه نۀ لري، نو همدا
يوازينے آيزومورفيزم دے۔
\end{explain}
\olpseganchor{OLP-0193-B014}{end}

\olpseganchor{OLP-0193-B015}{start}
\begin{prop}
\ollabel{prop:thq-unique-iso} که \LR{$\Struct{M}$} معياري وي، نو د
\olref{prop:thq-standard} له ثبوت څخه \LR{$g$} له~\LR{$\Struct{N}$} څخه
\LR{$\Struct{M}$} ته يوازينے آيزومورفيزم دے۔
\end{prop}
\olpseganchor{OLP-0193-B015}{end}

\olpseganchor{OLP-0193-B016}{start}
\begin{proof}
فرض کړئ \LR{$h\colon \Nat \to \Domain{M}$} له~\LR{$\Struct{N}$} څخه
\LR{$\Struct{M}$} ته آيزومورفيزم وي۔ د~\LR{$n$} پر بنسټ استقرا ښيو چې \LR{$g=h$}۔
که \LR{$n=0$} وي، نو د~\LR{$g$} له تعريفه \LR{$g(0)=\Assign{\Obj{0}}{M}$}۔ خو
څرنګه چې \LR{$h$} آيزومورفيزم دے، \LR{$h(0)=h(\Assign{\Obj{0}}{N}) =
\Assign{\Obj{0}}{M}$}، نو \LR{$g(0)=h(0)$}۔
\olpseganchor{OLP-0193-B016}{end}

\olpseganchor{OLP-0193-B017}{start}
اوس د \LR{$n+1$} حالت راواخلئ:
\begin{align*}
  g(n+1) & = \Value{\num{n+1}}{M} \text{\RL{ د \LR{$g$} له تعريفه}}\\
  & = \Value{\num{n}'}{M} \text{\RL{ ځکه \LR{$\num{n+1}\ident \num{n}'$}}}\\
  & = \Assign{\prime}{M}(\Value{\num{n}}{M})
    \text{\RL{ د \LR{$\Value{t'}{M}$} له تعريفه}}\\
  & = \Assign{\prime}{M}(g(n)) \text{\RL{ د \LR{$g$} له تعريفه}}\\
  & = \Assign{\prime}{M}(h(n)) \text{\RL{ د استقرايي فرض له مخې}}\\
  & = h(\Assign{\prime}{N}(n)) \text{\RL{ ځکه \LR{$h$} آيزومورفيزم دے}}\\
  & = h(n+1)
\end{align*}
\end{proof}
\olpseganchor{OLP-0193-B017}{end}

\olpseganchor{OLP-0193-B018}{start}
\begin{explain}
د هر شمېرېدونکے لامتناهي سټ~\LR{$M$} دپاره د~\LR{$\Nat$} او~\LR{$M$} تر منځ يوه
دوه اړخيزه يو پر يو تابع شته، نو هر داسې~\LR{$M$} په بالقوه ډول د معياري
مدل~\LR{$\Struct{M}$} د تعريف ساحه کېدے شي۔ په حقيقت کښې، کله چې يو څيز
\LR{$z \in M$} او مناسبه تابع~\LR{$s$} د~\LR{$\Assign{\Obj{0}}{M}$} او
\LR{$\Assign{\prime}{M}$} په توګه وټاکئ، د \LR{$+$}، \LR{$\times$} او \LR{$<$} تفسيرونه
لا له مخکې ټاکل شوي وي۔ يوازې هغه تابعې~\LR{$s\colon M \to M \setminus
\{z\}$} په معياري مدل کښې د~\LR{$\Assign{\prime}{M}$} په توګه مناسبې دي چې
يو پر يو او پر هدف سټ بشپړه دواړه وي۔ د~\LR{$s$} د قيمتونو ساحه
\LR{$z$} لرلے نۀ شي، ځکه که داسې واے نو
\LR{$\lforall[x][\eq/[\Obj 0][x']]$} به دروغ واے۔ دا تړلے فارمول
په~\LR{$\Struct{N}$} کښې رښتونې ده، نو \LR{$\Struct{M}$} ئې هم بايد رښتونې کړي۔ تابع~\LR{$s$} بايد
يو پر يو وي، ځکه په~\LR{$\Struct{N}$} کښې د جانشين
تابع~\LR{$\Assign{\prime}{N}$} همداسې ده، او د هغې
\LR{$\Assign{\prime}{N}$} يو پر يو توب په~\LR{$\Struct{N}$} کښې د
تړلے فارمول له خوا بيانېږي۔ دا بايد
پر هدف سټ بشپړه هم وي، ځکه که نه نو داسې~\LR{$x \in M \setminus \{z\}$}
به وي چې د~\LR{$s$} د قيمتونو په ساحه کښې نۀ وي، يعنې دا تړلے فارمول
\LR{$\lforall[x][(\eq[x][\Obj 0] \lor
\lexists[y][\eq[y'][x]])]$} به په~\LR{$\Struct{M}$} کښې دروغ وي، حال دا چې
په~\LR{$\Struct{N}$} کښې رښتونې ده۔
\textbf{د ژباړې د نسخې سمون (OLMOD-005):} سرچينې په وروستۍ جمله کښې
د تابعې د «قيمتونو ساحې» پر ځاے د «تعريف ساحه» ليکلې وه، حال دا
چې تابع s پر ټول M تعريف شوې ده۔ دلته د جملې له فورمول او د
پرښتياوالي له شرط سره سم د قيمتونو ساحه ليکل شوې ده۔
\end{explain}
\olpseganchor{OLP-0193-B018}{end}

\olpseganchor{OLP-0194-B004}{start}
\olfileid{mod}{mar}{nst}
\section{نامعياري مدلونه}
\olpseganchor{OLP-0194-B004}{end}

\olpseganchor{OLP-0194-B005}{start}
\begin{explain}
د~\LR{$\Lang{L_A}$} يو جوړښت هغه وخت معياري بولو چې له
\LR{$\Struct{N}$} سره آيزومورف وي۔ که يو جوړښت له~\LR{$\Struct{N}$}
سره آيزومورف نۀ وي، نامعياري بلل کېږي۔
\end{explain}
\olpseganchor{OLP-0194-B005}{end}

\olpseganchor{OLP-0194-B006}{start}
\begin{defn}
د~\LR{$\Lang{L_A}$} يو جوړښت~\LR{$\Struct{M}$} \emph{نامعياري} دے،
که له~\LR{$\Struct{N}$} سره آيزومورف نۀ وي۔ هغه غړي
\LR{$x \in \Domain{M}$} چې د کوم \LR{$n \in \Nat$} دپاره
\LR{$\Value{\num{n}}{M}$} وي، د~\LR{$\Struct{M}$} \emph{معياري شمېرې} بلل
کېږي؛ او نورې \emph{نامعياري شمېرې} بلل کېږي۔
\end{defn}
\olpseganchor{OLP-0194-B006}{end}

\olpseganchor{OLP-0194-B007}{start}
\begin{explain}
د~\olref[stm]{prop:standard-domain} له مخې د~\LR{$\Lang{L_A}$} هر
معياري جوړښت يوازې معياري غړي لري۔ نو يو نامعياري
جوړښت بايد لږ تر لږه يو نامعياري عنصر ولري۔ په حقيقت کښې،
د يوه نامعياري غړے شتون دا تضمينوي چې جوړښت
نامعياري دے۔
\end{explain}
\olpseganchor{OLP-0194-B007}{end}

\olpseganchor{OLP-0194-B008}{start}
\begin{prop}
که د~\LR{$\Lang{L_A}$} يو جوړښت~\LR{$\Struct{M}$} يوه نامعياري شمېره
ولري، نو \LR{$\Struct{M}$} نامعياري دے۔
\end{prop}
\olpseganchor{OLP-0194-B008}{end}

\olpseganchor{OLP-0194-B009}{start}
\begin{proof}
برعکس فرض کړئ، يعنې \LR{$\Struct{M}$} معياري وي خو نامعياري شمېره~\LR{$x$}
ولري۔ \LR{$g\colon \Nat \to \Domain{M}$} يو آيزومورفيزم ونيسئ۔ په
اسانه (پر~\LR{$n$} په استقرا) ليدل کېږي چې
\LR{$g(\Value{\num{n}}{N}) = \Value{\num{n}}{M}$}۔ په بله وينا، \LR{$g$} د
\LR{$\Struct{N}$} معياري شمېرې د~\LR{$\Struct{M}$} معياري شمېرو ته وړي۔ که
\LR{$\Struct{M}$} يوه نامعياري شمېره ولري، \LR{$g$} پر هدف سټ بشپړه کېدے نۀ
شي، چې له فرض سره ټکر دے۔
\end{proof}
\olpseganchor{OLP-0194-B009}{end}

\olpseganchor{OLP-0194-B010}{start}
\begin{prob}
په ياد راوړئ چې \LR{$\Th{Q}$} لاندې بديهي اصول لري:
\begin{align*}
& \lforall[x][\lforall[y][(\eq[x'][y'] \lif \eq[x][y])]] \tag{\LR{$!Q_1$}}\\
& \lforall[x][\eq/[\Obj 0][x']] \tag{\LR{$!Q_2$}}\\
& \lforall[x][(\eq[x][\Obj 0] \lor \lexists[y][\eq[x][y']])] \tag{\LR{$!Q_3$}}
\end{align*}
داسې جوړښتونه~\LR{$\Struct{M_1}$}، \LR{$\Struct{M_2}$}، \LR{$\Struct{M_3}$}
ورکړئ چې:
\begin{enumerate}
\item \LR{$\Sat{M_1}{!Q_1}$}، \LR{$\Sat{M_1}{!Q_2}$}، \LR{$\Sat/{M_1}{!Q_3}$}؛
\item \LR{$\Sat{M_2}{!Q_1}$}، \LR{$\Sat/{M_2}{!Q_2}$}، \LR{$\Sat{M_2}{!Q_3}$}؛ او
\item \LR{$\Sat/{M_3}{!Q_1}$}، \LR{$\Sat{M_3}{!Q_2}$}، \LR{$\Sat{M_3}{!Q_3}$}۔
\end{enumerate}
څرګنده ده چې د هر يوه دپاره يوازې~\LR{$\Assign{\Obj{0}}{M_i}$} او
\LR{$\Assign{\prime}{M_i}$} ټاکل پکار دي۔
\end{prob}
\olpseganchor{OLP-0194-B010}{end}

\olpseganchor{OLP-0194-B011}{start}
\begin{explain}
د~\LR{$\Lang{L_A}$} دپاره نامعياري جوړښتونه په اسانه ټاکل کېدے
شي۔ د بېلګې په توګه، داسې جوړښت ونيسئ چې د تعريف ساحه ئې~\LR{$\Int$} وي او
ټولې غيرمنطقي نښې په معمولي ډول تفسير کړي۔ څرنګه چې منفي شمېرې د
هېڅ~\LR{$n$} دپاره د~\LR{$\num{n}$} ارزښتونه نۀ دي، دا جوړښت نامعياري دے۔
البته دا به د حساب \emph{مدل} په دې معنا نۀ وي چې هماغه جملې رښتونې
کړي چې~\LR{$\Struct{N}$} ئې رښتونې کوي۔ مثلاً
\LR{$\lforall[x][\eq/[x'][\Obj{0}]]$} دروغ ده۔ خو د متناهي شاهد له قضيې
څخه په اسانۍ ثابتولے شو چې د حساب نامعياري مدلونه شته۔
\end{explain}
\olpseganchor{OLP-0194-B011}{end}

\olpseganchor{OLP-0194-B012}{start}
\begin{prop}
فرض کړئ \LR{$\Th{TA} = \Setabs{!A}{\Sat{N}{!A}}$} د~\LR{$\Struct{N}$} تيوري
ده۔ \LR{$\Th{TA}$} يو د شمېر وړ نامعياري مدل لري۔
\end{prop}
\olpseganchor{OLP-0194-B012}{end}

\olpseganchor{OLP-0194-B013}{start}
\begin{proof}
ژبه~\LR{$\Lang{L_A}$} په يوه نوي ثابت سمبول~\LR{$c$} وغځوئ او د
تړلي فارمولونه لاندې سټ ونيسئ:
\[
\Gamma = \Th{TA} \cup \{\eq/[c][\num{0}], \eq/[c][\num{1}],
\eq/[c][\num{2}], \dots\}
\]
د~\LR{$\Gamma$} هر مدل~\LR{$\Struct{M^c}$} به داسې يو غړے
\LR{$x = \Assign{c}{M}$} ولري چې نامعياري وي، ځکه د هر~\LR{$n \in \Nat$}
دپاره \LR{$x \neq \Value{\num{n}}{M}$}۔ دا هم څرګنده ده چې
\LR{$\Sat{M^c}{\Th{TA}}$}، ځکه \LR{$\Th{TA} \subseteq \Gamma$}۔ که له
\LR{$\Struct{M^c}$} څخه يوازې د~\LR{$c$} تفسير هېر کړو او هغه د~\LR{$\Lang{L_A}$}
يو جوړښت~\LR{$\Struct{M}$} وګرځوو، تعريف ساحه ئې لا هم نامعياري
\LR{$x$} لري او \LR{$\Sat{M}{\Th{TA}}$} هم وي۔ وروستۍ خبره ځکه تضمينېږي چې
\LR{$c$} په~\LR{$\Th{TA}$} کښې نۀ راځي۔ نو کافي ده وښيو چې \LR{$\Gamma$} مدل لري۔
\olpseganchor{OLP-0194-B013}{end}

\olpseganchor{OLP-0194-B014}{start}
د متناهي شاهد قضيه کاروو۔ که د~\LR{$\Gamma$} هر متناهي فرعي سټ د صدق وړ
وي، نو \LR{$\Gamma$} هم د صدق وړ دے۔ يو متناهي فرعي سټ
\LR{$\Gamma_0 \subseteq \Gamma$} ونيسئ۔ په~\LR{$\Gamma_0$} کښې د~\LR{$\Th{TA}$}
ځينې تړلي فارمولونه او ښايي د~\LR{$\eq/[c][\num{n}]$} د بڼې متناهي ډېرې جملې وي۔
که د دې بڼې کومه جمله وي، فرض کړئ~\LR{$k$} تر ټولو لوے عدد دے چې
\LR{$\eq/[c][\num{k}] \in \Gamma_0$} وي؛ او که هېڅ داسې جمله نه وي، د
k ارزښت صفر ونيسئ۔ \LR{$\Struct{N_k}$} د~\LR{$\Struct{N}$} داسې غځونه تعريف کړئ چې
\LR{$\Assign{c}{N_k} = k+1$} وي۔ بيا \LR{$\Sat{N_k}{\Gamma_0}$}: که
\LR{$!A \in \Th{TA}$} وي، \LR{$\Sat{N_k}{!A}$}، ځکه \LR{$\Struct{N_k}$} له~\LR{$c$}
پرته په هر څه کښې د~\LR{$\Struct{N}$} په شان دے او \LR{$c$} په~\LR{$!A$} کښې نۀ
راځي۔ او د دغې بڼې هرې راغلې جملې دپاره
\LR{$\Sat{N_k}{\eq/[c][\num{n}]}$}، ځکه \LR{$n \le k$} او
\LR{$\Value{c}{N_k} = k+1$}۔ نو د~\LR{$\Gamma$} هر متناهي فرعي سټ د صدق وړ
دے۔ د متناهي شاهد له قضيې \LR{$\Gamma$} يو مدل لري؛ او څرنګه چې ژبه
شمېرېدونکې او دغه مدل نامتناهي دے، د ښکته لوېنهايم--سکولم له قضيې
يو شمېرېدونکے داسې مدل اخلو۔
\textbf{د ژباړې د نسخې سمون (OLMOD-006):} سرچينې د متناهي فرعي سټ
دپاره تر ټولو لوے شاخص يوازې هغه وخت ټاکلے و چې لږ تر لږه يوه
د c او عددنښې د نابرابرۍ جمله پکې وي؛ دلته د ټولو راغلو شاخصونو يو
پاسنی حد نيول شوے، چې د هېڅ داسې جملې حالت هم رانغاړي۔
\textbf{د ژباړې د نسخې سمون (OLMOD-007):} د سرچينې ثبوت د متناهي
شاهد له قضيې يوازې د يو مدل شتون راويست، خو قضيه شمېرېدونکے مدل
غواړي؛ وروستۍ جمله د مخکې ثابتې شوې ښکته لوېنهايم--سکولم قضيې په
کارولو دا اړين ګام څرګندوي۔
\end{proof}
\olpseganchor{OLP-0194-B014}{end}

\olpseganchor{OLP-0195-B004}{start}
\olfileid{mod}{mar}{mdq}
\section{د~\LR{$\Th{Q}$} مدلونه}
\olpseganchor{OLP-0195-B004}{end}

\olpseganchor{OLP-0195-B005}{start}
\begin{explain}
موږ پوهېږو چې داسې نامعياري جوړښتونه شته چې هماغه
تړلي فارمولونه رښتونې کوي چې~\LR{$\Struct{N}$} ئې رښتونې کوي؛ يعنې د
\LR{$\Th{TA}$} مدلونه دي۔ څرنګه چې \LR{$\Sat{N}{\Th{Q}}$}، د~\LR{$\Th{TA}$} هر
مدل د~\LR{$\Th{Q}$} هم مدل دے۔ \LR{$\Th{Q}$} له~\LR{$\Th{TA}$} څخه ډېره کمزورې
ده؛ مثلاً
\LR{$\Th{Q} \Proves/ \lforall[x][\lforall[y][\eq[(x + y)][(y+x)]]]$}۔
کمزورې تيورۍ په اسانه پوره کېږي: ډېر مدلونه لري۔ مثلاً \LR{$\Th{Q}$}
داسې مدلونه لري چې
\LR{$\lforall[x][\lforall[y][\eq[(x + y)][(y+x)]]]$} دروغوي، خو هغه د
\LR{$\Th{TA}$} يا حتا د~\LR{$\Th{PA}$} مدلونه نۀ شي کېدے۔ د~\LR{$\Th{Q}$} مدلونه
نسبتاً ساده هم دي: په څرګند ډول ئې ټاکلے شو۔
\end{explain}
\olpseganchor{OLP-0195-B005}{end}

\olpseganchor{OLP-0195-B006}{start}
\begin{ex}
\ollabel{ex:model-K-of-Q}
هغه جوړښت~\LR{$\Struct{K}$} ونيسئ چې تعريف ساحه ئې
\LR{$\Domain{K} = \Nat \cup \{a\}$} او تفسيرونه ئې دا وي:
\begin{align*}
  \Assign{\Obj{0}}{K} & = 0\\
  \Assign{\prime}{K}(x) & =
  \begin{cases}
    x+1 & \text{\RL{که \LR{$x\in \Nat$} وي}}\\
    a & \text{\RL{که \LR{$x = a$} وي}}
  \end{cases}\\
  \Assign{+}{K}(x, y) & =
  \begin{cases}
    x+y & \text{\RL{که \LR{$x$}، \LR{$y \in\Nat$} وي}}\\
    a & \text{\RL{په نورو حالاتو کښې}}
  \end{cases}\\
  \Assign{\times}{K}(x, y) & =
  \begin{cases}
    xy & \text{\RL{که \LR{$x$}، \LR{$y \in\Nat$} وي}}\\
    0 & \text{\RL{که \LR{$x = 0$} يا \LR{$y = 0$} وي}}\\
    a & \text{\RL{په نورو حالاتو کښې}}\\
  \end{cases}\\
  \Assign{<}{K} & =
  \Setabs{\tuple{x,y}}{x, y \in \Nat \text{\RL{ او }} x<y} \cup
  \Setabs{\tuple{x,a}}{x \in \Domain{K}}
\end{align*}
د~\LR{$\Sat{K}{\Th{Q}}$} ښودلو دپاره بايد ثابته کړو چې د~\LR{$\Th{Q}$} ټول
بديهي اصول په~\LR{$\Struct{K}$} کښې رښتوني دي۔ د اسانتيا دپاره
\LR{$x^\nssucc$} د~\LR{$\Assign{\prime}{K}(x)$} په ځاے ليکو (په~\LR{$\Struct{K}$}
کښې د~\LR{$x$} «جانشين»)، \LR{$x \nsplus y$} د~\LR{$\Assign{+}{K}(x,y)$} په ځاے
(په~\LR{$\Struct{K}$} کښې د~\LR{$x$} او~\LR{$y$} «مجموعه»)، \LR{$x \nstimes y$} د
\LR{$\Assign{\times}{K}(x,y)$} په ځاے (په~\LR{$\Struct{K}$} کښې د~\LR{$x$}
او~\LR{$y$} «حاصل ضرب»)، او
\LR{$x \nsless y$} د~\LR{$\tuple{x,y}\in\Assign{<}{K}$} په ځاے ليکو۔ په دې
لنډونو سره د~\LR{$\Struct{K}$} عمليات په لا څرګند ډول داسې ښودلے شو:
\[
\begin{array}{c|c}
  x & x^\nssucc \\
  \hline
  n & n+1 \\
  a & a
\end{array}
\qquad
\begin{array}{c|ccc}
  x \nsplus y & 0 & m & a \\
  \hline
  0 & 0 & m & a \\
  n & n & n+m & a \\
  a & a & a & a \\
\end{array}
\qquad
\begin{array}{c|ccc}
  x \nstimes y & 0 & m & a \\
  \hline
  0 & 0 & 0 & 0 \\
  n & 0 & nm & a \\
  a & 0 & a & a \\
\end{array}
\]
\LR{$n\nsless m$} هله او يوازې هله چې \LR{$n<m$}، د~\LR{$n$}، \LR{$m\in\Nat$}
دپاره؛ او د هر~\LR{$x\in\Domain{K}$} دپاره \LR{$x\nsless a$}۔
\olpseganchor{OLP-0195-B006}{end}

\olpseganchor{OLP-0195-B007}{start}
\LR{$\Sat{K}{\lforall[x][\lforall[y][(\eq[x'][y'] \lif \eq[x][y])]]}$}،
ځکه \LR{$\nssucc$} يو پر يو ده۔
\LR{$\Sat{K}{\lforall[x][\eq/[\Obj 0][x']]}$}، ځکه په~\LR{$\Struct{K}$} کښې
\LR{$0$} د~\LR{$\nssucc$} جانشين نۀ دے۔
\LR{$\Sat{K}{\lforall[x][(\eq[x][\Obj 0] \lor
\lexists[y][\eq[x][y']])]}$}، ځکه د هر~\LR{$n>0$} دپاره
\LR{$n=(n-1)^\nssucc$} او \LR{$a=a^\nssucc$}۔
\olpseganchor{OLP-0195-B007}{end}

\olpseganchor{OLP-0195-B008}{start}
\LR{$\Sat{K}{\lforall[x][\eq[(x + \Obj 0)][x]]}$}، ځکه
\LR{$n\nsplus0=n+0=n$} او د~\LR{$\nsplus$} له تعريفه \LR{$a\nsplus0=a$}۔
\LR{$\Sat{K}{\lforall[x][\lforall[y][\eq[(x + y')][(x + y)']]]}$} لږه
پېچلې ده۔ که \LR{$n$} او~\LR{$m$} دواړه معياري وي:
\begin{align*}
(n \nsplus m^\nssucc) & = (n+(m+1)) = (n+m)+1 = (n \nsplus m)^\nssucc
\intertext{\RL{ځکه \LR{$\nsplus$} او \LR{$^\nssucc$} پر معياري شمېرو له~\LR{$+$} او
  \LR{$\prime$} سره سلا دي۔ اوس فرض کړئ \LR{$x \in \Domain{K}$}۔ بيا}}
(x \nsplus a^\nssucc) & = (x \nsplus a) = a = a^\nssucc = (x \nsplus a)^\nssucc
\intertext{\RL{پاتې حالت دا دے چې \LR{$y \in \Domain{K}$} وي خو \LR{$x=a$}۔
  دلته هم بايد وڅېړو چې \LR{$y=n$} معياري دے يا \LR{$y=a$}:}}
(a \nsplus n^\nssucc) & = (a \nsplus (n+1)) = a = a^\nssucc = (a \nsplus n)^\nssucc\\
(a \nsplus a^\nssucc) & = (a \nsplus a) = a = a^\nssucc = (a \nsplus a)^\nssucc
\end{align*}
\textbf{د ژباړې د نسخې سمون (OLMOD-008):} سرچينې په دې وېش کښې
د y وروستے حالت b ليکلے و، خو د K د تعريف ساحه يوازې طبيعي شمېرې
او a لري؛ دلته له ورپسې معادلې سره سم وروستے حالت a ليکل شوے دے۔
البته دا له اړتيا زيات تفصيل دے۔ مثلاً څرنګه چې د هر~\LR{$z$} دپاره
\LR{$a\nsplus z=a$}، نو سمدستي \LR{$a\nsplus a^\nssucc=a$} پايله کولے شو۔
پاتې بديهي اصول هم په همدې ډول کتل کېدے شي۔
\olpseganchor{OLP-0195-B008}{end}

\olpseganchor{OLP-0195-B009}{start}
نو \LR{$\Struct{K}$} د~\LR{$\Th{Q}$} يو مدل دے۔ «جمع»~\LR{$\nsplus$} ئې تبادلي هم
ده۔ خو داسې نورې جملې شته چې په~\LR{$\Struct{N}$} کښې رښتونې او په
\LR{$\Struct{K}$} کښې دروغ دي، او برعکس۔ مثلاً \LR{$a\nsless a$}، نو
\LR{$\Sat{K}{\lexists[x][x < x]}$} او
\LR{$\Sat/{K}{\lforall[x][\lnot x<x]}$}۔ له دې څرګندېږي چې
\LR{$\Th{Q} \Proves/ \lforall[x][\lnot x < x]$}۔
\end{ex}
\olpseganchor{OLP-0195-B009}{end}

\olpseganchor{OLP-0195-B010}{start}
\begin{prob}
ثابته کړئ چې د~\olref[mod][mar][mdq]{ex:model-K-of-Q} جوړښت
\LR{$\Struct{K}$} د~\LR{$\Th{Q}$} پاتې بديهي اصول پوره کوي:
\begin{align*}
  & \lforall[x][\eq[(x \times \Obj 0)][\Obj 0]] \tag{\LR{$!Q_6$}}\\
  & \lforall[x][\lforall[y][\eq[(x \times y')][((x \times y) + x)]]] \tag{\LR{$!Q_7$}}\\
  & \lforall[x][\lforall[y][(x < y \liff \lexists[z][\eq[(z' + x)][y]])]] \tag{\LR{$!Q_8$}}
\end{align*}
داسې تړلے فارمول ومومئ چې يوازې~\LR{$\prime$} ولري، په~\LR{$\Struct{N}$}
کښې رښتونې او په~\LR{$\Struct{K}$} کښې دروغ وي۔
\end{prob}
\olpseganchor{OLP-0195-B010}{end}

\olpseganchor{OLP-0195-B011}{start}
\begin{ex}
\ollabel{ex:model-L-of-Q} هغه جوړښت~\LR{$\Struct{L}$} ونيسئ چې
تعريف ساحه ئې \LR{$\Domain{L} = \Nat \cup \{a,b\}$} او تفسيرونه
\LR{$\Assign{\prime}{L} = \nssucc$}، \LR{$\Assign{+}{L} = \nsplus$} د لاندې
جدولونو له مخې وي:
\[
\begin{array}{c|c}
  x & x^\nssucc \\
  \hline
  n & n+1 \\
  a & a\\
  b & b
\end{array}
\qquad
\begin{array}{c|ccc}
  x \nsplus y & m & a & b\\
  \hline
  n & n+m & b & a\\
  a & a & b & a\\
  b & b & b & a
\end{array}
\]
څرنګه چې \LR{$\nssucc$} يو پر يو ده، \LR{$0$} ئې د قيمتونو په ساحه کښې
نۀ دے، او له~\LR{$0$} پرته هر \LR{$x\in\Domain{L}$} پکې دے، نو
\LR{$!Q_1$}--\LR{$!Q_3$} په~\LR{$\Struct{L}$} کښې رښتوني دي۔ د هر~\LR{$x$} دپاره
\LR{$x\nsplus0=x$}، نو \LR{$!Q_4$} هم رښتونے دے۔ د~\LR{$!Q_5$} دپاره
\LR{$x\nsplus y^\nssucc$} او \LR{$(x\nsplus y)^\nssucc$} پرتله کړئ۔ که \LR{$x$}
او~\LR{$y$} دواړه معياري وي، برابر دي، ځکه بيا \LR{$\nssucc$} او~\LR{$\nsplus$}
له~\LR{$\prime$} او~\LR{$+$} سره سلا دي۔ که \LR{$x$} نامعياري او~\LR{$y$} معياري وي،
\LR{$x\nsplus y^\nssucc=x=x^\nssucc=(x\nsplus y)^\nssucc$}۔ که \LR{$x$} او
\LR{$y$} دواړه نامعياري وي، څلور حالتونه لرو:
\begin{align*}
&  a \nsplus a^\nssucc  = b = b^\nssucc = (a \nsplus a)^\nssucc\\
&  b \nsplus b^\nssucc  = a = a^\nssucc = (b \nsplus b)^\nssucc\\
&  b \nsplus a^\nssucc  = b = b^\nssucc = (b \nsplus a)^\nssucc\\
&  a \nsplus b^\nssucc  = a = a^\nssucc = (a \nsplus b)^\nssucc\\
\intertext{\RL{که \LR{$x$} معياري خو \LR{$y$} نامعياري وي، لرو:}}
&  n \nsplus a^\nssucc  = n \nsplus a = b = b^\nssucc = (n \nsplus a)^\nssucc\\
&  n \nsplus b^\nssucc  = n \nsplus b = a = a^\nssucc = (n \nsplus b)^\nssucc
\end{align*}
\textbf{د ژباړې د نسخې سمون (OLMOD-009):} سرچينې د درېيم نامعياري
حالت په ښي اړخ کښې ناتړلے y ليکلے و؛ دلته د همدې کرښې له چپ اړخ
او د جدول له مخې a بېرته ايښوول شوے دے۔
نو \LR{$\Sat{L}{!Q_5}$}۔ خو \LR{$a\nsplus0\neq0\nsplus a$}، نو
\LR{$\Sat/{L}{\lforall[x][\lforall[y][\eq[(x+y)][(y+x)]]]}$}۔
\end{ex}
\olpseganchor{OLP-0195-B011}{end}

\olpseganchor{OLP-0195-B012}{start}
\begin{prob}
د~\olref[mod][mar][mdq]{ex:model-L-of-Q} جوړښت~\LR{$\Struct{L}$} په
\LR{$\nstimes$} او~\LR{$\nsless$} سره وغځوئ، څو د~\LR{$\times$} او~\LR{$<$} تفسيرونه
وي۔ وښيئ چې جوړښت مو د~\LR{$\Th{Q}$} پاتې بديهي اصول پوره کوي:
\begin{align*}
& \lforall[x][\eq[(x \times \Obj 0)][\Obj 0]] \tag{\LR{$!Q_6$}}\\ &
  \lforall[x][\lforall[y][\eq[(x \times y')][((x \times y) + x)]]]
  \tag{\LR{$!Q_7$}}\\ & \lforall[x][\lforall[y][(x < y \liff \lexists[z][\eq[(z'+x)][y]])]] \tag{\LR{$!Q_8$}}
\end{align*}
\end{prob}
\olpseganchor{OLP-0195-B012}{end}

\olpseganchor{OLP-0195-B013}{start}
\begin{prob}
د~\olref[mod][mar][mdq]{ex:model-L-of-Q} په~\LR{$\Struct{L}$} کښې
\LR{$a^\nssucc=a$} او \LR{$b^\nssucc=b$} دي۔ ايا د~\LR{$\Th{Q}$} داسې مدل شته چې
\LR{$a^\nssucc=b$} او \LR{$b^\nssucc=a$} وي؟
\end{prob}
\olpseganchor{OLP-0195-B013}{end}

\olpseganchor{OLP-0195-B014}{start}
\begin{explain}
موږ د~\LR{$\Th{Q}$} داسې مدلونه په څرګند ډول جوړ کړل چې نامعياري
غړي ئې د معياري عناصرو «هاخوا» پراتۀ دي۔ په حقيقت کښې،
بديهي اصول همدومره غواړي۔ يو نامعياري غړے~\LR{$x$} د~\LR{${} \nsless 0$} شرط
پوره کولے نۀ شي، ځکه \LR{$\Th{Q} \Proves
\lforall[x][\lnot x<0]$} (وګورئ
\olref[inc][req][min]{lem:less-zero})۔ همدارنګه د هر~\LR{$n$} دپاره
\LR{$\Th{Q} \Proves \lforall[x][(x < \num{n}' \lif
(\eq[x][\num{0}] \lor \eq[x][\num{1}] \lor \dots \lor
\eq[x][\num{n}]))]$} (\olref[inc][req][min]{lem:less-nsucc})، نو د
هېڅ~\LR{$n>0$} دپاره \LR{$a\nsless n$} نۀ شو لرلے۔
\end{explain}
\olpseganchor{OLP-0195-B014}{end}

\olpseganchor{OLP-0196-B004}{start}
\olfileid{mod}{mar}{mpa}
\section{د~\LR{$\Th{PA}$} مدلونه}
\olpseganchor{OLP-0196-B004}{end}

\olpseganchor{OLP-0196-B005}{start}
\begin{explain}
د~\LR{$\Th{TA}$} هر نامعياري مدل د~\LR{$\Th{PA}$} هم نامعياري مدل دے۔ موږ
پوهېږو چې د~\LR{$\Th{TA}$}، او په پايله کښې د~\LR{$\Th{PA}$}، نامعياري
مدلونه شته۔ دا هم پوهېږو چې داسې نامعياري مدلونه نامعياري «شمېرې»
لري، يعنې د تعريف ساحې داسې غړي چې له ټولو معياري «شمېرو»
هاخوا دي۔ خو دا څنګه ترتيب شوي دي؟ څومره دي؟ ومو ليدل چې د کمزورې
تيورۍ~\LR{$\Th{Q}$} مدلونه حتا يوازې يوه نامعياري شمېره لرلے شي۔ خو دغه
ساده جوړښتونه د~\LR{$\Th{PA}$} يا~\LR{$\Th{TA}$} مدلونه نۀ دي۔
\olpseganchor{OLP-0196-B005}{end}

\olpseganchor{OLP-0196-B006}{start}
د~\LR{$\Th{PA}$} يا~\LR{$\Th{TA}$} د مدلونو د جوړښت د پوهېدو کلي دا ده چې
وګورو په دغو تيوريو کښې کوم حقيقتونه د اشتقاق وړ دي۔ مثلاً
\LR{$\Th{PA}$} لا له مخکې
\LR{$\lforall[x][\eq/[x][x']]$} او
\LR{$\lforall[x][\lforall[y][\eq[(x+y)][(y+x)]]]$} ثابتوي، نو هغه ساده
جوړښتونه چې دا تړلي فارمولونه دروغوي د~\LR{$\Th{PA}$} د مدلونو په توګه
ردوي۔
\olpseganchor{OLP-0196-B006}{end}

\olpseganchor{OLP-0196-B007}{start}
فرض کړئ~\LR{$\Struct{M}$} د~\LR{$\Th{PA}$} مدل دے۔ که
\LR{$\Th{PA} \Proves !A$} وي، نو \LR{$\Sat{M}{!A}$}۔ بيا د
\LR{$\Assign{\Obj{0}}{M}$} دپاره~\LR{$\nszero$}، د~\LR{$\Assign{\prime}{M}$} دپاره
\LR{$\nssucc$}، د~\LR{$\Assign{+}{M}$} دپاره~\LR{$\nsplus$}، د
\LR{$\Assign{\times}{M}$} دپاره~\LR{$\nstimes$} او د~\LR{$\Assign{<}{M}$} دپاره
\LR{$\nsless$} کاروو۔ هره تړلے فارمول~\LR{$!A$} بيا د~\LR{$\nszero$}، \LR{$\nssucc$}،
\LR{$\nsplus$}، \LR{$\nstimes$} او~\LR{$\nsless$} په هکله کوم شرط بيانوي، او که
\LR{$\Sat{M}{!A}$} وي نو هغه شرط بايد پوره شي۔ مثلاً که
\LR{$\Sat{M}{!Q_1}$}، يعنې
\LR{$\Sat{M}{\lforall[x][\lforall[y][(\eq[x'][y'] \lif \eq[x][y])]]}$}،
نو \LR{$\nssucc$} بايد يو پر يو وي۔
\end{explain}
\olpseganchor{OLP-0196-B007}{end}

\olpseganchor{OLP-0196-B008}{start}
\begin{prop}
په~\LR{$\Struct{M}$} کښې \LR{$\nsless$} يو خطي سخت ترتيب دے، يعنې:
\begin{enumerate}
\item د هېڅ~\LR{$x \in \Domain{M}$} دپاره \LR{$x \nsless x$} نۀ وي۔
\item که \LR{$x \nsless y$} او \LR{$y \nsless z$} وي، نو \LR{$x \nsless z$}۔
\item د هر \LR{$x\neq y$} دپاره يا \LR{$x\nsless y$}، يا \LR{$y\nsless x$}۔
\end{enumerate}
\end{prop}
\olpseganchor{OLP-0196-B008}{end}

\olpseganchor{OLP-0196-B009}{start}
\begin{proof}
\LR{$\Th{PA}$} دا ثابتوي:
\begin{enumerate}
\item \LR{$\lforall[x][\lnot x < x]$}
\item \LR{$\lforall[x][\lforall[y][\lforall[z][((x < y \land y < z) \lif x < z)]]]$}
\item \LR{$\lforall[x][\lforall[y][((x < y \lor y < x) \lor \eq[x][y]))]]$}
\end{enumerate}
\end{proof}
\olpseganchor{OLP-0196-B009}{end}


\olpseganchor{OLP-0196-B010}{start}
\begin{prop}
\ollabel{prop:M-discrete} \LR{$\nszero$} د~\LR{$\Domain{M}$} تر ټولو کوچنے
غړے د~\LR{$\nsless$} په ترتيب کښې دے۔ د هر~\LR{$x$} دپاره
\LR{$x\nsless x^\nssucc$}، او \LR{$x^\nssucc$} تر ټولو کوچنے
\LR{$\nsless$}-لوي غړے دے۔ له صفر څخه پرته د هر~\LR{$x$} دپاره داسې يوازينے~\LR{$y$} شته چې
\LR{$y^\nssucc=x$} وي۔ (موږ~\LR{$y$} په~\LR{$\Struct{M}$} کښې
د~\LR{$x$} «مخکښېنی» بولو او په~\LR{$^\nssucc x$} ئې ښيو۔)
\textbf{د ژباړې د نسخې سمون (OLMOD-010):} سرچينې د يوازيني
مخکښيني دعوه د هر x دپاره کړې وه، چې صفر هم رانغاړي، حال دا چې د
Q2 بديهي اصل له مخې صفر د هېڅ څيز جانشين نۀ دے۔ دلته دعوه په سم
ډول له صفر څخه پرته غړو ته محدوده شوې ده۔
\end{prop}
\olpseganchor{OLP-0196-B010}{end}

\olpseganchor{OLP-0196-B011}{start}
\begin{proof}
تمرين۔
\end{proof}
\olpseganchor{OLP-0196-B011}{end}

\olpseganchor{OLP-0196-B012}{start}
\begin{prob}
په~\LR{$\Lang{L_A}$} کښې داسې تړلي فارمولونه ومومئ چې په~\LR{$\Th{PA}$} کښې
د اشتقاق وړ وي (او ځکه په~\LR{$\Struct{N}$} کښې رښتونې وي) او په
\olref[mod][mar][mpa]{prop:M-discrete} کښې د~\LR{$\nszero$}، \LR{$\nssucc$}
او~\LR{$\nsless$} ځانګړتياوې تضمين کړي۔
\end{prob}
\olpseganchor{OLP-0196-B012}{end}

\olpseganchor{OLP-0196-B013}{start}
\begin{prop}
د~\LR{$\Struct{M}$} ټول معياري غړي د~\LR{$\nsless$} له مخې له ټولو
نامعياري غړي څخه کوچني دي۔
\end{prop}
\olpseganchor{OLP-0196-B013}{end}

\olpseganchor{OLP-0196-B014}{start}
\begin{proof}
\LR{$n$} د معياري غړے~\LR{$\Value{\num{n}}{M}$} د
\LR{$\Struct{M}$} غړي د لنډون په توګه کاروو۔ لا~\LR{$\Th{Q}$} د هر~\LR{$n\in\Nat$} دپاره
\LR{$\lforall[x][(x < \num{n}' \lif (\eq[x][\num{0}] \lor
  \eq[x][\num{1}] \lor \dots \lor \eq[x][\num{n}]))]$} ثابتوي۔
د~\LR{$\nsless \nszero$} شرط پوره کوونکي غړي نشته۔ نو که~\LR{$n$} معياري او~\LR{$x$}
نامعياري وي، \LR{$x\nsless n$} نۀ شو لرلے۔ د تعريف له مخې نامعياري
عنصر د هېڅ~\LR{$n\in\Nat$} دپاره \LR{$\Value{\num{n}}{M}$} نۀ وي، نو
\LR{$x\neq n$} هم دے۔ څرنګه چې \LR{$\nsless$} خطي ترتيب دے، بايد
\LR{$n\nsless x$} وي۔
\end{proof}
\olpseganchor{OLP-0196-B014}{end}

\olpseganchor{OLP-0196-B015}{start}
\begin{prop}
د~\LR{$\Domain{M}$} هر نامعياري غړے~\LR{$x$} د لاندې فرعي سټ يو غړے
دے:
\[
\dots ^{\nssucc\nssucc\nssucc}x \nsless ^{\nssucc\nssucc}x \nsless
^{\nssucc}x \nsless x \nsless x^{\nssucc} \nsless x^{\nssucc\nssucc}
\nsless x^{\nssucc\nssucc\nssucc} \nsless \dots
\]
موږ دا فرعي سټ د~\LR{$x$} \emph{بلاک} بولو او په~\LR{$[x]$} ئې ليکو۔ نه تر
ټولو کوچنے او نه تر ټولو لوے غړے لري۔ دا د هغو
\LR{$y\in\Domain{M}$} د سټ په توګه ځانګړې کېدے شي چې د کوم معياري~\LR{$n$}
دپاره \LR{$x\nsplus n=y$} يا \LR{$y\nsplus n=x$} وي۔
\end{prop}
\olpseganchor{OLP-0196-B015}{end}

\olpseganchor{OLP-0196-B016}{start}
\begin{proof}
داسې سټ~\LR{$[x]$} تل شته، ځکه د~\LR{$\Domain{M}$} هر غړے~\LR{$y$}
يوازينے جانشين~\LR{$y^\nssucc$} لري او که دغه غړے نامعياري وي،
يوازينے مخکښېنی~\LR{$^\nssucc y$} لري۔ د پرله‌پسې غړي \LR{$y$} او~\LR{$y^\nssucc$}
دپاره \LR{$y\nsless y^\nssucc$}، او \LR{$y^\nssucc$} د~\LR{$\nsless$} له مخې
د~\LR{$\Domain{M}$} تر ټولو کوچنے هغه غړے دے چې \LR{$y$} ترې
\LR{$\nsless$} وي۔ څرنګه چې تل
\LR{$^\nssucc y\nsless y$} او \LR{$y\nsless y^\nssucc$}، نو \LR{$[x]$} نه تر ټولو
کوچنے او نه تر ټولو لوے غړے لري۔ که \LR{$y\in[x]$} وي، نو
\LR{$x\in[y]$}، ځکه بيا يا \LR{$y^{\nssucc\dots\nssucc}=x$} يا
\LR{$x^{\nssucc\dots\nssucc}=y$}۔ که
\LR{$y^{\nssucc\dots\nssucc}=x$} وي (له~\LR{$n$} ځله~\LR{$\nssucc$} سره)، نو
\LR{$y\nsplus n=x$} او برعکس، ځکه (که~\LR{$n$} د~\LR{$\prime$} نښو شمېر وي)
\LR{$\Th{PA} \Proves \lforall[x][\eq[x^{\prime\dots\prime}][(x +
    \num{n})]]$}۔
\end{proof}
\olpseganchor{OLP-0196-B016}{end}

\olpseganchor{OLP-0196-B017}{start}
\begin{prop}
که \LR{$[x]\neq[y]$} او \LR{$x\nsless y$} وي، نو د هر \LR{$u\in[x]$} او
\LR{$v\in[y]$} دپاره \LR{$u\nsless v$}۔
\end{prop}
\olpseganchor{OLP-0196-B017}{end}

\olpseganchor{OLP-0196-B018}{start}
\begin{proof}
پام وکړئ چې \LR{$\Th{PA} \Proves
\lforall[x][\lforall[y][(x < y \lif (x' < y \lor x' = y))]]$}۔
نو که \LR{$u\nsless v$} وي او \LR{$[u]\neq[v]$}، د هر~\LR{$n$} دپاره
\LR{$u\nsplus n^\nssucc\nsless v$} هم لرو۔
\olpseganchor{OLP-0196-B018}{end}

\olpseganchor{OLP-0196-B019}{start}
هر \LR{$u\in[x]$} د~\LR{$\nsless y$} شرط پوره کوي: د فرض له مخې \LR{$x\nsless y$}۔ که
\LR{$u\nsless x$} وي، د انتقال له مخې \LR{$u\nsless y$}۔ او که
\LR{$x\nsless u$} خو \LR{$u\in[x]$} وي، نو د کوم~\LR{$n$} دپاره
\LR{$u=x\nsplus n^\nssucc$}؛ نو د هغه حقيقت له مخې چې اوس مو ثابت کړ
\LR{$u\nsless y$}۔
\olpseganchor{OLP-0196-B019}{end}

\olpseganchor{OLP-0196-B020}{start}
اوس فرض کړئ \LR{$v\in[y]$} د~\LR{$\nsless y$} شرط پوره کوي، يعنې د کوم معياري~\LR{$m$}
دپاره \LR{$v\nsplus m^\nssucc=y$}۔ بيا \LR{$v\nsless x$} کېدے نۀ شي، ګنې
\LR{$y=v\nsplus m^\nssucc\nsless x$} به و۔ همدارنګه \LR{$x\neq v$}، ګنې
\LR{$x\nsplus m^\nssucc=v\nsplus m^\nssucc=y$} او \LR{$[x]=[y]$} به وو۔
نو \LR{$x\nsless v$}۔ بيا د هر~\LR{$n$} دپاره
\LR{$x\nsplus n^\nssucc\nsless v$} هم۔ له دې امله که \LR{$x\nsless u$} او
\LR{$u\in[x]$} وي، نو \LR{$u\nsless v$}۔ که \LR{$u\nsless x$} وي، د انتقال له
مخې \LR{$u\nsless v$}۔
\olpseganchor{OLP-0196-B020}{end}

\olpseganchor{OLP-0196-B021}{start}
په پای کښې که \LR{$y\nsless v$} وي، نو \LR{$u\nsless v$}، ځکه لکه چې ومو
ښودل \LR{$u\nsless y$} او \LR{$y\nsless v$}۔
\end{proof}
\olpseganchor{OLP-0196-B021}{end}

\olpseganchor{OLP-0196-B022}{start}
\begin{cor}
که \LR{$[x]\neq[y]$} وي، نو \LR{$[x]\cap[y]=\emptyset$}۔
\end{cor}
\olpseganchor{OLP-0196-B022}{end}

\olpseganchor{OLP-0196-B023}{start}
\begin{proof}
فرض کړئ \LR{$z\in[x]$} او \LR{$x\nsless y$}۔ بيا د هر \LR{$u\in[y]$} دپاره
\LR{$z\nsless u$}۔ که \LR{$z\in[y]$} هم واے، \LR{$z\nsless z$} به مو لرلے۔
د~\LR{$y\nsless x$} حالت هم ورته دے۔
\end{proof}
\olpseganchor{OLP-0196-B023}{end}

\olpseganchor{OLP-0196-B024}{start}
\begin{explain}
دا معنا لري چې بلاکونه پخپله داسې ترتيبېدے شي چې د~\LR{$\nsless$} درناوی
وکړي: \LR{$[x]\nsless[y]$} هله او يوازې هله چې \LR{$x\nsless y$}؛ يا په
هم‌ارزه توګه، چې د هر \LR{$u\in[x]$} او \LR{$v\in[y]$} دپاره
\LR{$u\nsless v$} وي۔ څرګنده ده چې معياري بلاک~\LR{$[0]$} تر ټولو کوچنے بلاک
دے۔ له هېڅ نامعياري بلاک سره تقاطع نۀ لري، او هېڅ دوه نامعياري
بلاکونه هم يو بل نۀ پرې کوي۔ په ځانګړي ډول، د پرله‌پسې جانشينونو يا
مخکښينو په اخيستو بل بلاک ته نۀ شئ رسېدے۔
\end{explain}
\olpseganchor{OLP-0196-B024}{end}

\olpseganchor{OLP-0196-B025}{start}
\begin{prop}
که \LR{$x$} او~\LR{$y$} نامعياري وي، نو \LR{$x\nsless x\nsplus y$} او
\LR{$x\nsplus y\notin[x]$}۔
\end{prop}
\olpseganchor{OLP-0196-B025}{end}

\olpseganchor{OLP-0196-B026}{start}
\begin{proof}
که \LR{$y$} نامعياري وي، نو \LR{$y\neq\nszero$}۔
\LR{$\Th{PA} \Proves \lforall[x][\lforall[y][(y \neq \Obj{0} \lif x < (x+y))]]$}۔
\textbf{د ژباړې د نسخې سمون (OLMOD-011):} سرچينې د دې عمومي
حقيقت په ليکلو کښې y ازاد پرېښے و؛ دلته پر y هم کلي سور ورزيات
شوے، لکه دعوه او ورپسې استعمال چې غواړي۔
اوس فرض کړئ \LR{$x\nsplus y\in[x]$}۔ څرنګه چې
\LR{$x\nsless x\nsplus y$}، د کوم معياري عدد دپاره
\LR{$x\nsplus n^\nssucc=x\nsplus y$}۔ خو د جمع د اختصار قانون
\LR{$\Th{PA} \Proves \lforall[x][\lforall[y][\lforall[z][(\eq[(x+y)][(x+z)] \lif y = z)]]]$}
دے۔ له دې به د کوم معياري~\LR{$n$} دپاره \LR{$y=n^\nssucc$} پايله شي؛ خو
\LR{$y$} نامعياري فرض شوې ده۔
\end{proof}
\olpseganchor{OLP-0196-B026}{end}

\olpseganchor{OLP-0196-B027}{start}
\begin{prop}
تر ټولو کوچنے نامعياري بلاک نشته۔
\end{prop}
\olpseganchor{OLP-0196-B027}{end}

\olpseganchor{OLP-0196-B028}{start}
\begin{proof}
\LR{$\Th{PA} \Proves \lforall[x][\lexists[y][(\eq[(y+y)][x] \lor
      \eq[(y+y)'][x])]]$}؛ يعنې هر~\LR{$x$} پر~\LR{$2$} وېشل کېږي (ښايي پاتې
شونې~\LR{$1$} ولري)۔ که \LR{$x$} نامعياري وي، \LR{$y$} هم نامعياري دے۔ د مخکينۍ
قضيې له مخې \LR{$y\nsless y\nsplus y$} او \LR{$y\nsplus y\notin[y]$}۔
بيا \LR{$y\nsless(y\nsplus y)^\nssucc$} او
\LR{$(y\nsplus y)^\nssucc\notin[y]$} هم۔ خو
\LR{$x=y\nsplus y$} يا \LR{$x=(y\nsplus y)^\nssucc$}، نو \LR{$y\nsless x$} او
\LR{$y\notin[x]$}۔
\end{proof}
\olpseganchor{OLP-0196-B028}{end}

\olpseganchor{OLP-0196-B029}{start}
\begin{prop}
تر ټولو لوے بلاک نشته۔
\end{prop}
\olpseganchor{OLP-0196-B029}{end}

\olpseganchor{OLP-0196-B030}{start}
\begin{proof}
تمرين۔
\end{proof}
\olpseganchor{OLP-0196-B030}{end}

\olpseganchor{OLP-0196-B031}{start}
\begin{prob}
وښيئ چې د~\LR{$\Th{PA}$} په يوه نامعياري مدل کښې تر ټولو لوے بلاک
نشته۔
\end{prob}
\olpseganchor{OLP-0196-B031}{end}

\olpseganchor{OLP-0196-B032}{start}
\begin{prop}
\ollabel{prop:blocks-dense}
د بلاکونو ترتيب ګڼ دے۔ يعنې که \LR{$x\nsless y$} او \LR{$[x]\neq[y]$} وي،
نو داسې بلاک~\LR{$[z]$} شته چې له دواړو بېل او د هغو تر منځ وي۔
\end{prop}
\olpseganchor{OLP-0196-B032}{end}

\olpseganchor{OLP-0196-B033}{start}
\begin{proof}
فرض کړئ \LR{$x\nsless y$}۔ لکه مخکې، \LR{$x\nsplus y$} پر دوه وېشل کېږي
(ښايي پاتې شونې ولري): داسې \LR{$z\in\Domain{M}$} شته چې يا
\LR{$x \nsplus y = z \nsplus z$} يا
\LR{$x \nsplus y = (z \nsplus z)^\nssucc$}۔ عنصر~\LR{$z$} د~\LR{$x$} او~\LR{$y$}
«اوسط» دے، او \LR{$x\nsless z$} او \LR{$z\nsless y$}۔
\textbf{د ژباړې د نسخې سمون (OLMOD-012):} سرچينې د دې ثبوت په
دوو معادلو کښې ناتعريفه جمع کارولې وه؛ دلته د ټول څپرکي ټاکلې
نامعياري جمع بېرته ايښوول شوې ده۔
\end{proof}
\olpseganchor{OLP-0196-B033}{end}

\olpseganchor{OLP-0196-B034}{start}
\begin{prob}
د~\olref[mod][mar][mpa]{prop:blocks-dense} مفصل ثبوت وليکئ۔
د~\LR{$z$} د شتون د تضمين دپاره \LR{$\Th{PA}$} بايد کومه تړلے فارمول
اشتقاق کړي؟ ولې \LR{$x\nsless z$} او \LR{$z\nsless y$}، او ولې
\LR{$[x]\neq[z]$} او \LR{$[z]\neq[y]$} دي؟
\end{prob}
\olpseganchor{OLP-0196-B034}{end}

\olpseganchor{OLP-0196-B035}{start}
\begin{explain}
له دې امله، په هر شمېرېدونکي نامعياري مدل کښې نامعياري بلاکونه د
ناطقو په شان ترتيب لري: يو شمېرېدونکے لامتناهي ګڼ خطي ترتيب بې له
سرحدي ټکو جوړوي۔ ښودل کېدے شي چې هر دوه داسې شمېرېدونکے لامتناهي
ترتيبونه آيزومورف دي۔ نو د رښتيني حساب د هر دوه د شمېر وړ
نامعياري مدلونو \LR{$\Struct{M}_1$} او~\LR{$\Struct{M_2}$} هغه راکمونې چې
يوازې~\LR{$<$} او~\LR{$=$} لرونکې ژبې ته کېږي، آيزومورفې دي۔ په رښتيا،
آيزومورفيزم~\LR{$h$} داسې تعريفېدے شي: د~\LR{$\Struct{M_1}$} او
\LR{$\Struct{M_2}$} معياري برخې د معياري مدل~\LR{$\Struct{N}$} سره، او ځکه له
يو بل سره، آيزومورفې دي۔ د نامعياري برخې بلاکونه پخپله د ناطقو په
شان ترتيب شوي او ځکه آيزومورف دي؛ د بلاکونو يو آيزومورفيزم د هر
بلاک دننه د خپل‌سرو عناصرو په سره برابرولو غځېدے شي، او بيا د
\LR{$\Struct{M_1}$} کښې د~\LR{$x$} د جانشين انځور د~\LR{$\Struct{M_2}$} کښې د
\LR{$x$} د انځور جانشين وټاکل شي۔ له دې دا \emph{نۀ} راځي چې
\LR{$\mathfrak{M}_1$} او~\LR{$\mathfrak{M}_2$} د حساب په بشپړه ژبه کښې
آيزومورف دي (آيزومورفيزم تل د يوې ژبه په نسبت وي)، ځکه پر
\LR{$\Domain{M_1}$} او~\LR{$\Domain{M_2}$} د جمع او ضرب د تعريفولو نامتجانسې
لارې شته۔ (دا د واټ د يوې مشهورې قضيې څخه هم راځي چې د يوې بشپړې
تيورۍ د شمېرېدونکو مدلونو شمېر 2 کېدے نۀ شي۔)
\textbf{د ژباړې د نسخې سمون (OLMOD-013):} سرچينې د بلاکونو
شمېرېدنه له هر نامعياري مدل څخه بې له شرطه پايله کړې وه؛ دا يوازې د شمېرېدونکي مدل دپاره لازمېږي۔ دلته دغه شرط څرګند شوے
او د ورپسې جملې له شمېرېدونکو مدلونو سره برابر شوے دے۔
\end{explain}
\olpseganchor{OLP-0196-B035}{end}

\olpseganchor{OLP-0197-B004}{start}
\olfileid{mod}{mar}{cmp}
\section{د حساب محاسبه کېدونکي مدلونه}
\olpseganchor{OLP-0197-B004}{end}

\olpseganchor{OLP-0197-B005}{start}
\begin{explain}
معياري مدل~\LR{$\Struct{N}$} دوه ښې ځانګړتياوې لري۔ تعريف ساحه ئې طبيعي
شمېرې~\LR{$\Nat$} دي، يعنې عناصر ئې هماغه ډول څيزونه دي چې د حساب
په ژبه ئې په هکله خبرې کول غواړو، او معياري عددنښه~\LR{$\num{n}$} په
رښتيا~\LR{$n$} په نښه کوي۔ بله ښه ځانګړتيا دا ده چې د~\LR{$\Lang{L_A}$} د
غيرمنطقي نښو ټول تفسيرونه \emph{محاسبه کېدونکي} دي۔ د جانشين، جمع
او ضرب تابعې چې \LR{$\Assign{\prime}{N}$}، \LR{$\Assign{+}{N}$} او
\LR{$\Assign{\times}{N}$} دي، د شمېرو محاسبه کېدونکې تابعې دي (مثلاً د
ټيورينګ ماشينونو له خوا محاسبه کېدونکې، يا په ابتدايي بازګښت
تعريفېدونکې)۔ او په~\LR{$\Struct{N}$} کښې د کوچنيوالي اړيکه، يعنې
\LR{$\Assign{<}{N}$}، پرېکړه کېدونکې ده۔
\olpseganchor{OLP-0197-B005}{end}

\olpseganchor{OLP-0197-B006}{start}
د~\LR{$\Th{Q}$} او~\LR{$\Th{PA}$} په شان د حسابي تيوريو نامعياري مدلونه بايد
نامعياري عناصر ولري۔ نو تعريف ساحې ئې عموماً له~\LR{$\Nat$} سره نور
غړي هم لري۔ خو هر نامتناهي شمېرېدونکے جوړښت پر هر
شمېرېدونکے لامتناهي سټ، د~\LR{$\Nat$} په ګډون، جوړېدے شي۔ نو داسې نامعياري
مدلونه هم شته چې تعريف ساحه ئې~\LR{$\Nat$} وي۔ البته په داسې
مدل~\LR{$\Struct{M}$} کښې لږ تر لږه ځينې شمېرې خپل معمولي نقش نۀ شي
لوبولے، ځکه داسې~\LR{$k$} بايد وي چې د ټولو~\LR{$n \in \Nat$} دپاره
له~\LR{$\Value{\num{n}}{M}$} څخه بېل وي۔
\textbf{د ژباړې د نسخې سمون (OLMOD-023):} سرچينې دلته هر
شمېرېدونکے جوړښت ياد کړے و؛ خو د دې نسخې د شمېر وړتيا اصطلاح متناهي
سټونه هم رانغاړي، او متناهي جوړښت د طبيعي شمېرو پر سټ نۀ شي
جوړېدے۔ دلته د
نامتناهي والي اړين شرط څرګند شوے؛ د بحث نامعياري حسابي مدلونه له
مخکې لا نامتناهي دي۔
\end{explain}
\olpseganchor{OLP-0197-B006}{end}

\olpseganchor{OLP-0197-B007}{start}
\begin{defn}
د~\LR{$\Lang{L_A}$} يو جوړښت~\LR{$\Struct{M}$} \emph{محاسبه
کېدونکے} دے هله او يوازې هله چې \LR{$\Domain{M}=\Nat$}؛
\LR{$\Assign{\prime}{M}$}، \LR{$\Assign{+}{M}$} او~\LR{$\Assign{\times}{M}$}
محاسبه کېدونکې تابعې؛ او \LR{$\Assign{<}{M}$} پرېکړه کېدونکې اړيکه وي۔
\end{defn}
\olpseganchor{OLP-0197-B007}{end}

\olpseganchor{OLP-0197-B008}{start}
\begin{ex}\ollabel{ex:comp-model-q}
د~\olref[mdq]{ex:model-K-of-Q} جوړښت~\LR{$\Struct{K}$} په ياد راوړئ۔
تعريف ساحه ئې \LR{$\Domain{K}=\Nat\cup\{a\}$} او تفسيرونه ئې دا وو:
\begin{align*}
  \Assign{\Obj{0}}{K} & = 0\\
  \Assign{\prime}{K}(x) & =
  \begin{cases}
    x+1 & \text{\RL{که \LR{$x\in \Nat$} وي}}\\
    a & \text{\RL{که \LR{$x = a$} وي}}
  \end{cases}\\
  \Assign{+}{K}(x, y) & =
  \begin{cases}
    x+y & \text{\RL{که \LR{$x$}، \LR{$y \in\Nat$} وي}}\\
    a & \text{\RL{په نورو حالاتو کښې}}
  \end{cases}\\
  \Assign{\times}{K}(x, y) & =
  \begin{cases}
    xy & \text{\RL{که \LR{$x$}، \LR{$y \in\Nat$} وي}}\\
    0 & \text{\RL{که \LR{$x=0$} يا \LR{$y=0$} وي}}\\
    a & \text{\RL{په نورو حالاتو کښې}}\\
  \end{cases}\\
  \Assign{<}{K} & =
  \Setabs{\tuple{x,y}}{x, y \in \Nat \text{\RL{ او }} x<y} \cup
  \Setabs{\tuple{x,a}}{x \in \Domain{K}}
\end{align*}
\textbf{د ژباړې د نسخې سمون (OLMOD-014):} سرچينې د وروستي سټ په
شرط کښې n ليکلے و، حال دا چې په مرتبه جوړه کښې x راغلے دے؛ دلته
شرط د x د تعريف ساحې غړيتوب ته سم شوے، لکه د K په اصلي تعريف کښې
چې راغلے دے۔
خو \LR{$\Domain{K}$} شمېرېدونکے لامتناهي او له~\LR{$\Nat$} سره همشمېره دے۔
مثلاً \LR{$g\colon\Nat\to\Domain{K}$} چې \LR{$g(0)=a$} او د~\LR{$n>0$} دپاره
\LR{$g(n)=n-1$} وي، يوه دوه اړخيزه يو پر يو تابع ده۔ دا د~\LR{$\Th{Q}$} د نوي
مدل~\LR{$\Struct{K'}$} او~\LR{$\Struct{K}$} تر منځ پر آيزومورفيزم بدلولے شو۔
په~\LR{$\Struct{K'}$} کښې بايد د~\LR{$\Lang{L_A}$} نښو ته بېلې تابعې او
اړيکې وټاکو، ځکه د~\LR{$\Nat$} بېل غړي د معياري او نامعياري
شمېرو نقشونه لوبوي۔
\textbf{د ژباړې د نسخې سمون (OLMOD-015):} سرچينې د مثبت n دپاره
د n جمع يو نقشه ليکلې وه، چې صفر او يو د تابعې له قيمتونو بهر
پرېږدي او دوه اړخيزه يو په يو نۀ ده۔ دلته د n منفي يو نقشه ليکل
شوې؛ همدا نقشه د ورپسې ټولو لېږدول شويو تفسيرونو سره سلا ده۔
\olpseganchor{OLP-0197-B008}{end}

\olpseganchor{OLP-0197-B009}{start}
په ځانګړي ډول، اوس~\LR{$0$} د~\LR{$a$} نقش لوبوي، نه د تر ټولو کوچنۍ معياري
شمېرې۔ تر ټولو کوچنۍ معياري شمېره اوس~\LR{$1$} ده۔ نو
\LR{$\Assign{\Obj{0}}{K'}=1$} ټاکو۔ د جانشين تابع هم اوس بېله ده: يوه
معياري شمېره، يعنې \LR{$n>0$}، ورکړل شي، بيا هم~\LR{$n+1$} راګرځوي۔ خو~\LR{$0$}
اوس د~\LR{$a$} نقش لوبوي، چې خپل جانشين دے۔ نو
\LR{$\Assign{\prime}{K'}(0)=0$}۔ د جمع او ضرب دپاره هم لرو:
\begin{align*}
\Assign{+}{K'}(x, y) & =
  \begin{cases}
    x+y-1 & \text{\RL{که \LR{$x$}، \LR{$y >0$} وي}}\\
    0 & \text{\RL{په نورو حالاتو کښې}}
  \end{cases}\\
  \Assign{\times}{K'}(x, y) & =
  \begin{cases}
    1 & \text{\RL{که \LR{$x = 1$} يا \LR{$y = 1$} وي}}\\
    xy - x - y + 2 & \text{\RL{که \LR{$x$}، \LR{$y > 1$} وي}}\\
    0 & \text{\RL{په نورو حالاتو کښې}}\\
  \end{cases}
\end{align*}
او \LR{$\tuple{x,y}\in\Assign{<}{K'}$} هله او يوازې هله چې \LR{$x<y$}،
\LR{$x>0$} او \LR{$y>0$} وي، يا \LR{$y=0$} وي۔
\olpseganchor{OLP-0197-B009}{end}

\olpseganchor{OLP-0197-B010}{start}
دا ټولې تابعې د طبيعي شمېرو محاسبه کېدونکې تابعې دي او
\LR{$\Assign{<}{K'}$} پر~\LR{$\Nat$} يوه پرېکړه کېدونکې اړيکه ده؛ خو دغه
تابعې پر~\LR{$\Nat$} د جانشين، جمع او ضرب له معمولي تابعو سره يو شان نۀ
دي، او \LR{$\Assign{<}{K'}$} هم پر~\LR{$\Nat$} له معمولي~\LR{$<$} سره يو شان نۀ ده۔
\end{ex}
\olpseganchor{OLP-0197-B010}{end}

\olpseganchor{OLP-0197-B011}{start}
\begin{prob}
داسې جوړښت~\LR{$\Struct{L'}$} ورکړئ چې \LR{$\Domain{L'}=\Nat$} وي او
د~\olref[mod][mar][mdq]{ex:model-L-of-Q} له~\LR{$\Struct{L}$} سره
آيزومورف وي۔
\end{prob}
\olpseganchor{OLP-0197-B011}{end}

\olpseganchor{OLP-0197-B012}{start}
\begin{explain}
\Olref{ex:comp-model-q} ښيي چې \LR{$\Th{Q}$} پر~\LR{$\Nat$} د تعريف ساحې
محاسبه کېدونکي نامعياري مدلونه لري۔ خو لاندې پايله ښيي چې د
\LR{$\Th{PA}$} د مدلونو دپاره دا خبره رښتونې نۀ ده (او ځکه د~\LR{$\Th{TA}$}
د مدلونو دپاره هم نۀ ده)۔
\end{explain}
\olpseganchor{OLP-0197-B012}{end}

\olpseganchor{OLP-0197-B013}{start}
\begin{thm}[د ټېننباوم قضيه]
\LR{$\Struct{N}$} تر آيزومورفيزم پورې د~\LR{$\Th{PA}$} يوازينے محاسبه
کېدونکے مدل دے۔
\textbf{د ژباړې د نسخې سمون (OLMOD-016):} سرچينې «يوازينے» په
لفظي ډول ويلي وو، خو د طبيعي شمېرو پر بنسټ محاسبه کېدونکي او
آيزومورف بيا-نومول شوي نقلونه شته؛ د ټېننباوم د قضيې دقيقه پايله دا ده چې
هېڅ محاسبه کېدونکے \emph{نامعياري} مدل نشته۔ دلته «تر
آيزومورفيزم پورې» دا دقيقوالی څرګندوي۔
\end{thm}
\olpseganchor{OLP-0197-B013}{end}

\olpseganchor{OLP-0198-B004}{start}
\olchapter{mod}{int}{د منځګړيتوب قضيه}
\olpseganchor{OLP-0198-B004}{end}

\olpseganchor{OLP-0199-B004}{start}
\olfileid{mod}{int}{int}
\olpseganchor{OLP-0199-B004}{end}

\olpseganchor{OLP-0199-B005}{start}
\olsection{سريزه}
\olpseganchor{OLP-0199-B005}{end}

\olpseganchor{OLP-0199-B006}{start}
د منځګړيتوب قضيه لاندې پايله ده: فرض کړئ
\LR{$\Entails !A \lif !B$}۔ بيا داسې تړلے فارمول~\LR{$!C$} شته چې
\LR{$\Entails !A \lif !C$} او \LR{$\Entails !C \lif !B$} وي۔ سربېره پر دې،
په~\LR{$!C$} کښې هر ثابت سمبول، تابع او پريديکات (له
\LR{$\eq$} پرته) په~\LR{$!A$} او~\LR{$!B$} دواړو کښې راځي۔ تړلے فارمول~\LR{$!C$} د
\LR{$!A$} او~\LR{$!B$} يو \emph{منځګړی} بلل کېږي۔
\olpseganchor{OLP-0199-B006}{end}

\olpseganchor{OLP-0199-B007}{start}
د منځګړيتوب قضيه په خپل ذات کښې په زړه پورې ده، خو اصلي ارزښت ئې
په دې کښې دے چې د يوې تيورۍ دننه د تعريف‌وړتيا په هکله پايلې، او
هغه شرطونه پرې ثابتولے شو چې لاندې ئې د دوو سازګارو تيوريو يوځاے
کول يوه سازګاره تيوري ورکوي۔ لومړۍ پايله د بېت د تعريف‌وړتيا قضيه،
او دويمه د رابنسن د ګډې سازګارۍ قضيه بلل کېږي۔
\olpseganchor{OLP-0199-B007}{end}

\olpseganchor{OLP-0200-B004}{start}
\olfileid{mod}{int}{sep}
\olpseganchor{OLP-0200-B004}{end}

\olpseganchor{OLP-0200-B005}{start}
\olsection{د تړلي فارمولونه بېلول}
\olpseganchor{OLP-0200-B005}{end}

\olpseganchor{OLP-0200-B006}{start}
د منځګړيتوب د قضيې ثبوت ته تر تلو مخکې لږ بنسټيز کار پکار دے۔ د
\LR{$!A$} او~\LR{$!B$} منځګړی داسې تړلے فارمول~\LR{$!C$} ده چې
\LR{$!A \Entails !C$} او \LR{$!C \Entails !B$} وي۔ د نقيض عکس له مخې،
وروستۍ خبره هله او يوازې هله رښتونې ده چې
\LR{$\lnot !B \Entails \lnot !C$} وي۔ هغه تړلے فارمول~\LR{$!C$} چې دا
خاصيت ولري، ويل کېږي چې \LR{$!A$} او~\LR{$\lnot !B$} \emph{بېلوي}۔ نو د
\LR{$!A$} او~\LR{$!B$} د يو منځګړي موندل د داسې تړلے فارمول له موندلو سره
برابر دي چې \LR{$!A$} او~\LR{$\lnot !B$} بېل کړي۔ د معمول په څېر به د يوې
عمومي شوې بڼې څېړل ګټور وي: داسې جمله چې د تړلي فارمولونه دوه
\emph{سټونه} بېل کړي۔
\olpseganchor{OLP-0200-B006}{end}

\olpseganchor{OLP-0200-B007}{start}
\begin{defn}
يوه جمله~\LR{$!C$} د جملو سټونه~\LR{$\Gamma$} او~\LR{$\Delta$} هله او يوازې هله
\emph{بېلوي} چې \LR{$\Gamma \Entails !C$} او
\LR{$\Delta \Entails \lnot !C$} وي۔ که داسې هېڅ تړلے فارمول نۀ وي،
نو~\LR{$\Gamma$} او~\LR{$\Delta$} \emph{نه بېلېدونکي} دي۔
\end{defn}
\olpseganchor{OLP-0200-B007}{end}

\olpseganchor{OLP-0200-B008}{start}
د~\LR{$\Gamma$}، \LR{$\Delta$} او~\LR{$!C$} د مدلونو د ټولګيو تر منځ د شاملېدو
اړيکې لاندې ښودل شوې دي:
\begin{figure}[h]
  \centering
  \begin{tikzpicture}[node distance=2cm, auto, thick]
    \draw [rounded corners] (0,0) -- (6,0) -- (6,3) -- (0,3) --  cycle;
    \draw (1.5,1.5) circle (0.9cm);
    \draw (4.5,1.5) circle (0.9cm);
    \path node at (1.5,1.5) {\Large \LR{$\Gamma$}};
    \path node at (4.5,1.5) {\Large \LR{$\Delta$}};
    \path node at (0.4,2.6) {\Large \LR{$\formula{C}$}};
    \path node at (3.4,0.4) {\Large \LR{$\lnot \formula{C}$}};
    \draw (2.5,0) .. controls (2.5,1.5) and (3.5,1.5) .. (3.5,3);
  \end{tikzpicture}
  \caption{\LR{$!C$}، \LR{$\Gamma$} او~\LR{$\Delta$} بېلوي}
  \ollabel{fig:sep}
\end{figure}
\olpseganchor{OLP-0200-B008}{end}


\olpseganchor{OLP-0200-B009}{start}
\begin{lem}\ollabel{lem:sep1}
فرض کړئ~\LR{$\Lang{L}_0$} هغه ژبه ده چې هر هغه ثابت سمبول،
تابع او پريديکات (له~\LR{$\doteq$} پرته) لري چې په
\LR{$\Gamma$} او~\LR{$\Delta$} \emph{دواړو} کښې راځي؛ او~\LR{$\Lang{L}'_0$} د
هر~\LR{$n \ge 0$} دپاره د نامتناهي ډېرو نويو ثابت سمبولونه~\LR{$\Obj c_n$}
په ورزياتولو ترې ترلاسه شوې وي۔ بيا که~\LR{$\Gamma$} او~\LR{$\Delta$} په
\LR{$\Lang{L}_0$} کښې نه بېلېدونکي وي، په~\LR{$\Lang{L}'_0$} کښې هم نه
بېلېدونکي دي۔
\end{lem}
\olpseganchor{OLP-0200-B009}{end}

\olpseganchor{OLP-0200-B010}{start}
\begin{proof}
په غيرمستقيم ډول پر مخ ځو: د تناقض په موخه فرض کړئ چې~\LR{$\Gamma$} او
\LR{$\Delta$} په~\LR{$\Lang{L}'_0$} کښې بېلېږي۔ بيا د کوم~\LR{$!C \in
\Lang{L}_0$} دپاره \LR{$\Gamma \Entails \Subst{!C}{c}{x}$} او
\LR{$\Delta \Entails \lnot \Subst{!C}{c}{x}$} دي (دلته~\LR{$c$} يو نوی
ثابت سمبول دے؛ هغه حالت چې~\LR{$!C$} تر يوه زيات داسې نوي
ثابت سمبول ولري، ورته دے)۔ د متناهي شاهد له مخې د~\LR{$\Gamma$} يو
متناهي فرعي سټ~\LR{$\Gamma_0$} او د~\LR{$\Delta$} يو متناهي فرعي
سټ~\LR{$\Delta_0$} شته، داسې چې
\LR{$\Gamma_0 \Entails \Subst{!C}{c}{x}$} او
\LR{$\Delta_0 \Entails \lnot \Subst{!C}{c}{x}$} وي۔ پرېږدئ~\LR{$!G$} په
\LR{$\Gamma_0$} کښې د ټولو فارمولونه عطف، او~\LR{$!H$} په~\LR{$\Delta_0$}
کښې د ټولو فارمولونه عطف وي۔ بيا
\begin{align*}
  !G & \Entails \Subst{!C}{c}{x}, & !H  \Entails \lnot \Subst{!C}{c}{x}.
\end{align*}
له لومړۍ څخه د تعميم په قاعده \LR{$!G \Entails
\lforall[x][!C]$} لرو؛ او له دويمې څخه د نقيض عکس په قاعده
\LR{$\Subst{!C}{c}{x} \Entails \lnot !H$}، او له دې څخه هم
\LR{$\lforall[x][!C] \Entails \lnot !H$} ترلاسه کېږي۔
\textbf{د ژباړې د نسخې سمون (OLMOD-017):} سرچينې په وروستي استلزام
کښې ناټاکلې نښه~\LR{$\delta$} ليکلې وه؛ دلته~\LR{$!H$} ليکل شوے، ځکه همدا د
\LR{$\Delta_0$} د جملو عطف دے او د مخکنۍ کرښې نقيض عکس همدا پايله ورکوي۔
د نقيض عکس بيا کارول \LR{$!H \Entails
\lnot \lforall[x][!C]$} راکوي۔ د يکنواختۍ له مخې
\begin{align*}
  \Gamma &\Entails \lforall[x][!C], &
  \Delta & \Entails \lnot \lforall[x][!C],
\end{align*}
نو~\LR{$\lforall[x][!C]$} په~\LR{$\Lang{L}_0$} کښې~\LR{$\Gamma$} او~\LR{$\Delta$}
بېلوي۔
\end{proof}
\olpseganchor{OLP-0200-B010}{end}

\olpseganchor{OLP-0200-B011}{start}
\begin{lem}\ollabel{lem:sep2}
فرض کړئ~\LR{$\Gamma \cup \{ \lexists[x][!S] \}$} او~\LR{$\Delta$} نه
بېلېدونکي دي، او~\LR{$c$} يو نوی ثابت سمبول دے چې په~\LR{$\Gamma$}،
\LR{$\Delta$} يا~\LR{$!S$} کښې نۀ راځي۔ بيا
\LR{$\Gamma \cup \{ \lexists[x][!S], \Subst{!S}{c}{x} \}$} او
\LR{$\Delta$} هم نه بېلېدونکي دي۔
\end{lem}
\olpseganchor{OLP-0200-B011}{end}

\olpseganchor{OLP-0200-B012}{start}
\begin{proof}
د تناقض دپاره فرض کړئ چې~\LR{$!C$}،
\LR{$\Gamma \cup \{\lexists[x][!S], \Subst{!S}{c}{x}\}$} او
\LR{$\Delta$} بېلوي، خو په عين وخت کښې
\LR{$\Gamma \cup \{\lexists[x]{!S} \}$} او~\LR{$\Delta$} نه بېلېدونکي دي۔
دوه حالتونه بېلوو:
\begin{enumerate}
\item \LR{$c$} په~\LR{$!C$} کښې نۀ راځي: په دې حالت کښې
  \LR{$\Gamma \cup \{\lexists[x][!S], \lnot!C \}$} د صدق وړ دے (که
  داسې نۀ وے، \LR{$!C$} به~\LR{$\Gamma \cup \{\lexists[x][!S] \}$} او
  \LR{$\Delta$} بېل کړي وو)۔ که~\LR{$\Subst{!S}{c}{x}$} ورزيات هم شي، سټ د
  صدق وړ پاتې کېږي؛ نو په پايله کښې~\LR{$!C$} بيا هم
  \LR{$\Gamma \cup \{ \lexists[x][!S], \Subst{!S}{c}{x} \}$} او
  \LR{$\Delta$} نۀ بېلوي۔
\item \LR{$c$} په~\LR{$!C$} کښې راځي، نو~\LR{$!C$} د~\LR{$\Subst{!C}{c}{x}$} بڼه
  لري۔ بيا لرو چې
  \[
  \Gamma \cup \{ \lexists[x][!S], \Subst{!S}{c}{x}\} \Entails \Subst{!C}{c}{x},
  \]
  نو د استنتاج د قضيې او تعميم له مخې
  \LR{$\Gamma, \lexists[x][!S] \Entails \lforall[x][(!S \lif !C)]$}،
  او په پای کښې
  \LR{$\Gamma \cup \{ \lexists[x][!S] \} \Entails \lexists[x][!C]$}۔
  له بلې خوا، \LR{$\Delta \Entails \lnot \Subst{!C}{c}{x}$} او ځکه د
  تعميم له مخې~\LR{$\Delta \Entails \lnot \lexists[x][!C]$}۔ نو
  \LR{$\Gamma \cup \{\lexists[x][!S] \}$} او~\LR{$\Delta$} د بېلېدو وړ دي؛
  دا تناقض دے۔
\end{enumerate}
\end{proof}
\olpseganchor{OLP-0200-B012}{end}

\olpseganchor{OLP-0201-B004}{start}
\olfileid{mod}{int}{prf}
\olpseganchor{OLP-0201-B004}{end}

\olpseganchor{OLP-0201-B005}{start}
\olsection{د کرېګ د منځګړيتوب قضيه}
\olpseganchor{OLP-0201-B005}{end}

\olpseganchor{OLP-0201-B006}{start}
\begin{thm}[د کرېګ د منځګړيتوب قضيه]
\ollabel{thm:interpol}
که~\LR{$\Entails !A \lif !B$} وي، نو داسې تړلے فارمول~\LR{$!C$} شته چې
\LR{$\Entails !A \lif !C$} او \LR{$\Entails !C \lif !B$} وي، او په~\LR{$!C$}
کښې هر ثابت سمبول، تابع او پريديکات (له~\LR{$\eq$} پرته)
په~\LR{$!A$} او~\LR{$!B$} دواړو کښې راځي۔ تړلے فارمول~\LR{$!C$} د~\LR{$!A$} او~\LR{$!B$}
يو \emph{منځګړی} بلل کېږي۔
\end{thm}
\olpseganchor{OLP-0201-B006}{end}

\olpseganchor{OLP-0201-B007}{start}
\begin{proof}
فرض کړئ~\LR{$\Lang{L}_1$} د~\LR{$!A$} ژبه او~\LR{$\Lang{L}_2$} د~\LR{$!B$} ژبه ده۔
\LR{$\Lang{L}_0 = \Lang{L}_1 \cap \Lang{L}_2$} پرېږدئ۔ د هر
\LR{$i \in \{0, 1, 2 \}$} دپاره~\LR{$\Lang{L}'_i$} له~\LR{$\Lang{L}_i$} څخه د
نامتناهي ډېرو نويو ثابت سمبولونه~\LR{$\Obj c_0, \Obj c_1, \Obj c_2,
\dots$} په ورزياتولو ترلاسه شوې ژبه وي۔
\olpseganchor{OLP-0201-B007}{end}

\olpseganchor{OLP-0201-B008}{start}
که~\LR{$!A$} د صدق وړ نۀ وي، \LR{$\lexists[x][\eq/[x][x]]$} يو منځګړی دے۔
که~\LR{$\lnot !B$} د صدق وړ نۀ وي (او له همدې امله~\LR{$!B$} معتبره وي)،
\LR{$\lexists[x][\eq[x][x]]$} يو منځګړی دے۔ نو دا هم فرض کولے شو چې
\LR{$!A$} او~\LR{$\lnot !B$} دواړه د صدق وړ دي۔
\olpseganchor{OLP-0201-B008}{end}

\olpseganchor{OLP-0201-B009}{start}
د منځګړيتوب د قضيې د نقيض عکس د ثابتولو دپاره فرض کړئ چې د~\LR{$!A$}
او~\LR{$!B$} هېڅ منځګړی نشته۔ په بله وينا، فرض کړئ چې~\LR{$\{!A\}$} او
\LR{$\{\lnot !B\}$} په~\LR{$\Lang{L}_0$} کښې نه بېلېدونکي دي۔
\olpseganchor{OLP-0201-B009}{end}

\olpseganchor{OLP-0201-B010}{start}
موخه مو دا ده چې جوړه~\LR{$(\{ !A \}, \{\lnot!B\})$} يوې اعظمي نه
بېلېدونکې جوړې~\LR{$(\Gamma^*, \Delta^*)$} ته وغځوو۔ پرېږدئ~\LR{$!A_0$}،
\LR{$!A_1$}، \LR{$!A_2$}، \dots، د~\LR{$\Lang{L}'_1$} تړلي فارمولونه، او~\LR{$!B_0$}،
\LR{$!B_1$}، \LR{$!B_2$}، \dots، د~\LR{$\Lang{L}'_2$} تړلي فارمولونه په لړ کښې
راوړي۔
\textbf{د ژباړې د نسخې سمون (OLMOD-018):} سرچينې دلته د نا پراخو
ژبو~\LR{$\Lang{L}_1$} او~\LR{$\Lang{L}_2$} جملې په لړ کښې راوړې وې، حال دا
چې ورپسې ادعا په نويو ثابتونو لرونکو ژبو~\LR{$\Lang{L}'_1$} او
\LR{$\Lang{L}'_2$} کښې اعظمي سازګاري غواړي؛ دلته شمېرنه همدغو پراخو
ژبو ته سمه شوې ده۔
د~\LR{$n \ge 0$} دپاره د تړلي فارمولونه د سټونو دوه زياتېدونکې لړۍ
\LR{$(\Gamma_n, \Delta_n)$} داسې تعريفوو۔
\LR{$\Gamma_0 = \{ !A\}$} او~\LR{$\Delta_0 = \{\lnot !B \}$} کښېږدئ۔ که
\LR{$(\Gamma_n, \Delta_n)$} لا تعريف شوي وي، \LR{$\Gamma_{n+1}$} او
\LR{$\Delta_{n+1}$} داسې تعريف کړئ:
\begin{enumerate}
\item که~\LR{$\Gamma_n \cup \{!A_n \}$} او~\LR{$\Delta_n$} په
  \LR{$\Lang{L}'_0$} کښې نه بېلېدونکي وي، \LR{$!A_n$} په~\LR{$\Gamma_{n+1}$} کښې
  واچوئ۔ سربېره پر دې، که~\LR{$!A_n$} يو وجودي فارمول
  \LR{$\lexists[x][!S]$} وي، داسې نوی ثابت سمبول~\LR{$c$} وټاکئ چې په
  \LR{$\Gamma_n$}، \LR{$\Delta_n$}، \LR{$!A_n$} يا~\LR{$!B_n$} کښې نۀ راځي، او
  \LR{$\Subst{!S}{c}{x}$} په~\LR{$\Gamma_{n+1}$} کښې واچوئ۔
\item که~\LR{$\Gamma_{n+1}$} او~\LR{$\Delta_n \cup \{!B_n \}$} په
  \LR{$\Lang{L}'_0$} کښې نه بېلېدونکي وي، \LR{$!B_n$} په~\LR{$\Delta_{n+1}$} کښې
  واچوئ۔ سربېره پر دې، که~\LR{$!B_n$} يو وجودي فارمول
  \LR{$\lexists[x][!S]$} وي، داسې نوی ثابت سمبول~\LR{$c$} وټاکئ چې په
  \LR{$\Gamma_{n+1}$}، \LR{$\Delta_n$}، \LR{$!A_n$} يا~\LR{$!B_n$} کښې نۀ راځي، او
  \LR{$\Subst{!S}{c}{x}$} په~\LR{$\Delta_{n+1}$} کښې واچوئ۔
\end{enumerate}
په پای کښې داسې تعريف کړئ:
\begin{align*}
  \Gamma^* & = \bigcup_{n\ge 0} \Gamma_n, &
  \Delta^* & = \bigcup_{n\ge 0} \Delta_n.
\end{align*}
اوس پر~\LR{$n$} په هممهاله استقرا ثابتولے شو چې:
\begin{enumerate}
\item\ollabel{part-a} \LR{$\Gamma_n$} او~\LR{$\Delta_n$} په~\LR{$\Lang{L}'_0$}
  کښې نه بېلېدونکي دي؛
\item\ollabel{part-b} \LR{$\Gamma_{n+1}$} او~\LR{$\Delta_n$} په
  \LR{$\Lang{L}'_0$} کښې نه بېلېدونکي دي۔
\end{enumerate}
د~\olref{part-a} بنسټ له~\olref[sep]{lem:sep1} څخه ترلاسه کېږي۔ د
\olref{part-b} دپاره بايد درې حالتونه بېل کړو:
\begin{enumerate}
\item که~\LR{$\Gamma_0 \cup \{!A_0 \}$} او~\LR{$\Delta_0$} بېلېدونکي وي،
  نو~\LR{$\Gamma_1 = \Gamma_0$} او~\olref{part-b} يوازې هماغه
  \olref{part-a} ده؛
\item که~\LR{$\Gamma_1 = \Gamma_0 \cup\{ !A_0\}$} وي، نو~\LR{$\Gamma_1$}
  او~\LR{$\Delta_0$} د جوړونې له مخې نه بېلېدونکي دي؛
\item هغه حالت پاتې کېږي چې~\LR{$!A_0$} وجودي وي، نو
  \LR{$\Gamma_1 = \Gamma_0 \cup \{ \lexists[x][!S],
  \Subst{!S}{c}{x} \}$}۔ د جوړونې له مخې
  \LR{$\Gamma_0 \cup \{ \lexists[x][!S]\}$} او~\LR{$\Delta_0$} نه
  بېلېدونکي دي، نو د~\olref[sep]{lem:sep2} له مخې
  \LR{$\Gamma_0 \cup \{ \lexists[x][!S], \Subst{!S}{c}{x} \}$} او
  \LR{$\Delta_0$} هم نه بېلېدونکي دي۔
\end{enumerate}
په دې سره د پورته~\olref{part-a} او~\olref{part-b} د استقرا بنسټ
بشپړېږي۔ اوس استقرايي ګام ته ورځو۔ د~\olref{part-a} دپاره، که
\LR{$\Delta_{n+1} = \Delta_n \cup \{ !B_n \}$} وي، نو
\LR{$\Gamma_{n+1}$} او~\LR{$\Delta_{n+1}$} د جوړونې له مخې نه بېلېدونکي دي
(حتٰی چې~\LR{$!B_n$} وجودي هم وي، د~\olref[sep]{lem:sep2} له مخې)؛ او
که~\LR{$\Delta_{n+1} = \Delta_n$} وي (ځکه چې~\LR{$\Gamma_{n+1}$} او
\LR{$\Delta_n \cup \{!B_n\}$} بېلېدونکي دي)، د~\olref{part-b}
استقرايي فرضيه کاروو۔ د~\olref{part-b} د استقرايي ګام دپاره، که
\LR{$\Gamma_{n+2} = \Gamma_{n+1} \cup \{!A_{n+1} \}$} وي، نو
\LR{$\Gamma_{n+2}$} او~\LR{$\Delta_{n+1}$} د جوړونې له مخې نه بېلېدونکي دي
(حتٰی چې~\LR{$!A_{n+1}$} وجودي هم وي، د~\olref[sep]{lem:sep2} له مخې)؛
او که~\LR{$\Gamma_{n+2} = \Gamma_{n+1}$} وي، د~\olref{part-a} هغه
استقرايي حالت کاروو چې همدا اوس مو ثابت کړ۔ په دې سره پر
\olref{part-a} او~\olref{part-b} استقرا پاے ته رسېږي۔
\olpseganchor{OLP-0201-B010}{end}

\olpseganchor{OLP-0201-B011}{start}
له دې څخه پايله کېږي چې~\LR{$\Gamma^*$} او~\LR{$\Delta^*$} نه بېلېدونکي دي؛
که داسې نۀ وے، د متناهي شاهد له مخې به داسې~\LR{$n \ge 0$} وای چې
\LR{$\Gamma_n$} او~\LR{$\Delta_n$} بېلېدونکي وي، د~\olref{part-a} خلاف۔ په
ځانګړي ډول، \LR{$\Gamma^*$} او~\LR{$\Delta^*$} سازګار دي: ځکه که لومړنے يا
دويم سټ ناسازګار وے، په ترتيب سره به
\LR{$\lexists[x][\eq/[x][x]]$} يا~\LR{$\lforall[x][\eq[x][x]]$} بېل کړي وو۔
\olpseganchor{OLP-0201-B011}{end}

\olpseganchor{OLP-0201-B012}{start}
اوس ښيو چې~\LR{$\Gamma^*$} په~\LR{$\Lang{L}'_1$} کښې او په همدې ډول
\LR{$\Delta^*$} په~\LR{$\Lang{L}'_2$} کښې اعظمي سازګار دے۔ د لومړي سټ دپاره
فرض کړئ چې د کوم~\LR{$n \ge 0$} دپاره \LR{$!A_n \notin \Gamma^*$} او
\LR{$\lnot !A_n \notin \Gamma^*$} وي۔ که~\LR{$!A_n \notin \Gamma^*$} وي،
نو~\LR{$\Gamma_n \cup \{!A_n \}$} له~\LR{$\Delta_n$} څخه بېلېدونکے دے، او
ځکه داسې~\LR{$!C \in \Lang{L}'_0$} شته چې دواړه لاندې خبرې سمې وي:
\begin{align*}
  \Gamma^* & \Entails !A_n \lif !C, &
  \Delta^* & \Entails \lnot !C.
\end{align*}
همدا ډول، که~\LR{$\lnot !A_n \notin \Gamma^*$} وي، داسې
\LR{$!C' \in \Lang{L}'_0$} شته چې دواړه لاندې خبرې سمې وي:
\begin{align*}
  \Gamma^* & \Entails \lnot !A_n \lif !C', &
  \Delta^* & \Entails \lnot !C'.
\end{align*}
د بياني منطق له مخې~\LR{$\Gamma^* \Entails !C \lor !C'$} او
\LR{$\Delta^* \Entails \lnot (!C \lor !C')$}، نو~\LR{$!C \lor !C'$}،
\LR{$\Gamma^*$} او~\LR{$\Delta^*$} بېلوي۔ ورته استدلال ثابتوي چې~\LR{$\Delta^*$}
هم اعظمي دے۔
\olpseganchor{OLP-0201-B012}{end}

\olpseganchor{OLP-0201-B013}{start}
په پای کښې ښيو چې~\LR{$\Gamma^* \cap \Delta^*$} په~\LR{$\Lang{L}'_0$} کښې
اعظمي سازګار دے۔ دا په څرګنده سازګار دے، ځکه د دوو سازګارو سټونو
تقاطع ده۔ د اعظميت د ښودلو دپاره~\LR{$!S \in \Lang{L}'_0$} ونيسئ۔ اوس
\LR{$\Gamma^*$} په~\LR{$\Lang{L}'_1 \supseteq \Lang{L}'_0$} کښې اعظمي دے،
او په همدې ډول~\LR{$\Delta^*$} په
\LR{$\Lang{L}'_2 \supseteq \Lang{L}'_0$} کښې اعظمي دے۔ نو يا
\LR{$!S \in \Gamma^*$} يا~\LR{$\lnot !S \in \Gamma^*$}، او يا
\LR{$!S \in \Delta^*$} يا~\LR{$\lnot !S \in \Delta^*$} وي۔ که
\LR{$!S \in \Gamma^*$} او~\LR{$\lnot !S \in \Delta^*$} وي، نو~\LR{$!S$} به
\LR{$\Gamma^*$} او~\LR{$\Delta^*$} بېل کړي؛ او که~\LR{$\lnot !S \in \Gamma^*$}
او~\LR{$!S \in \Delta^*$} وي، نو~\LR{$\lnot !S$} به~\LR{$\Gamma^*$} او
\LR{$\Delta^*$} بېل کړي۔ له دې امله، يا
\LR{$!S \in \Gamma^* \cap \Delta^*$} يا
\LR{$\lnot !S \in \Gamma^* \cap \Delta^*$} وي، او
\LR{$\Gamma^* \cap \Delta^*$} اعظمي دے۔
\olpseganchor{OLP-0201-B013}{end}

\olpseganchor{OLP-0201-B014}{start}
څرنګه چې~\LR{$\Gamma^*$} اعظمي سازګار دے، داسې يو مدل~\LR{$\Struct{M}'_1$}
لري چې د تعريف ساحه~\LR{$\Domain{M'_1}$} ئې په هغو او يوازې هغو عناصرو
مشتمله ده چې ثابت سمبولونه تفسيروي، يعنې~\LR{$\Assign{c}{M'_1}$}؛ کټ مټ
لکه د بشپړتيا د قضيې په ثبوت کښې
(\olref[fol][com][cth]{thm:completeness})۔ په همدې ډول، \LR{$\Delta^*$}
داسې يو مدل~\LR{$\Struct{M}'_2$} لري چې د تعريف ساحه~\LR{$\Domain{M'_2}$} ئې د
ثابت سمبولونه په تفسيرونو~\LR{$\Assign{c}{M'_2}$} جوړه ده۔
\olpseganchor{OLP-0201-B014}{end}

\olpseganchor{OLP-0201-B015}{start}
\LR{$\Struct{M_1}$} له~\LR{$\Struct{M'_1}$} څخه داسې ترلاسه کړئ چې په
\LR{$\Lang{L}'_1 \setminus \Lang{L}'_0$} کښې د ثابت سمبولونه،
تابعې او پريديکاتونه تفسيرونه ترې وغورځوئ؛ او
\LR{$\Struct{M_2}$} هم په همدې ډول ترلاسه کړئ۔ بيا
\LR{$h \colon M_1 \to M_2$} چې د
\LR{$h(\Assign{c}{M'_1}) = \Assign{c}{M'_2}$} له مخې تعريف شوے، په
\LR{$\Lang{L}'_0$} کښې يو آيزومورفيزم دے، ځکه لکه چې ومو ښودل
\LR{$\Gamma^* \cap \Delta^*$} په~\LR{$\Lang{L}'_0$} کښې اعظمي سازګار دے۔
دا ځکه پايله کېږي چې هره~\LR{$\Lang{L}'_0$}-تړلے فارمول يا په
\LR{$\Gamma^*$} او~\LR{$\Delta^*$} دواړو کښې ده، يا په هېڅ يوه کښې نۀ ده:
نو~\LR{$\Assign{c}{M'_1} \in \Assign{P}{M'_1}$} هله او يوازې هله چې
\LR{$\Atom{P}{c} \in \Gamma^*$}، هله او يوازې هله چې
\LR{$\Atom{P}{c} \in \Delta^*$}، او هله او يوازې هله چې
\LR{$\Assign{c}{M'_2} \in \Assign{P}{M'_2}$} وي۔ د آيزومورفيزم نور شرطونه
هم په ورته ډول ثابتېدے شي۔
\olpseganchor{OLP-0201-B015}{end}

\olpseganchor{OLP-0201-B016}{start}
اوس به د ژبه~\LR{$\Lang{L}'_1 \cup \Lang{L}'_2$} يو
مدل~\LR{$\Struct{M}$} داسې تعريف کړو:
\begin{enumerate}
\item د تعريف ساحه~\LR{$\Domain{M}$} هماغه~\LR{$\Domain{M_2}$} دے، يعنې د ټولو
  \LR{$\Assign{c}{M'_2}$} عناصرو سټ؛
\item که يو پريديکات~\LR{$P$} په
  \LR{$\Lang{L}'_2 \setminus \Lang{L}'_1$} کښې وي، نو
  \LR{$\Assign{P}{M} = \Assign{P}{M'_2}$}؛
\item که يو پريديکات~\LR{$P$} په
  \LR{$\Lang{L}'_1\setminus \Lang{L}'_2$} کښې وي، نو
  \LR{$\Assign{P}{M} = h(\Assign{P}{M'_1})$}؛ يعنې
  \LR{$\tuple{\Assign{c_1}{M'_2}, \dots, \Assign{c_n}{M'_2}} \in
  \Assign{P}{M}$} هله او يوازې هله چې
  \LR{$\tuple{\Assign{c_1}{M'_1}, \dots,
  \Assign{c_n}{M'_1}} \in \Assign{P}{M'_1}$} وي۔
  \textbf{د ژباړې د نسخې سمون (OLMOD-019):} سرچينې په لومړۍ
  معادله کښې د~\LR{$\Lang{L}'_1$} ځانګړي پريديکات تفسير په ناسم ډول له
  \LR{$\Struct{M'_2}$} څخه لېږداوه، چې هلته اصلاً نۀ تفسيرېږي؛ دلته د
  ورپسې «يعنې» شرط سره سم د~\LR{$\Struct{M'_1}$} تفسير د~\LR{$h$} له لارې
  لېږدول شوے دے۔
\item که يو پريديکات~\LR{$P$} په~\LR{$\Lang{L}'_0$} کښې وي، نو
  \LR{$\Assign{P}{M} = \Assign{P}{M'_2} = h(\Assign{P}{M'_1})$}۔
\item د~\LR{$\Lang{L}'_1 \cup \Lang{L}'_2$} تابعې، د
  ثابت سمبولونه په ګډون، په ورته ډول اداره کېږي۔
\end{enumerate}
\textbf{د ژباړې د نسخې سمون (OLMOD-020):} سرچينې وروستے ګډ مدل
يوازې د نا پراخو ژبو~\LR{$\Lang{L}_1 \cup \Lang{L}_2$} دپاره تعريف کړ،
خو بيا ئې د نويو ثابتونو لرونکو ژبو~\LR{$\Lang{L}'_1$} او
\LR{$\Lang{L}'_2$} له مدلونو سره موافق او د~\LR{$\Gamma^* \cup \Delta^*$}
مدل وباله۔ دلته ګډه ژبه او د نښو ټول حالتونه پراخو ژبو ته سم شوي،
څو نوي ثابتونه هم تفسير ولري او وروستۍ د صدق ادعا تعريف شوې وي۔
\olpseganchor{OLP-0201-B016}{end}

\olpseganchor{OLP-0201-B017}{start}
په پای کښې پر فارمولونه په استقرا ښودل کېږي چې~\LR{$\Struct{M}$} د
\LR{$\Lang{L}'_1$} پر ټولو فارمولونه له~\LR{$\Struct{M'_1}$} سره، او د
\LR{$\Lang{L}'_2$} پر ټولو فارمولونه له~\LR{$\Struct{M'_2}$} سره موافق دے۔
په ځانګړي ډول، \LR{$\Struct{M} \Entails \Gamma^* \cup \Delta^*$}؛ نو
\LR{$\Struct{M} \Entails !A$} او~\LR{$\Struct{M} \Entails \lnot!B$}، او
\LR{$\not\Entails !A \lif !B$}۔ په دې سره د کرېګ د منځګړيتوب د قضيې
ثبوت پاے ته رسېږي۔
\end{proof}
\olpseganchor{OLP-0201-B017}{end}

\olpseganchor{OLP-0202-B004}{start}
\olfileid{mod}{int}{def}
\olpseganchor{OLP-0202-B004}{end}

\olpseganchor{OLP-0202-B005}{start}
\olsection{د تعريف‌وړتيا قضيه}
\olpseganchor{OLP-0202-B005}{end}

\olpseganchor{OLP-0202-B006}{start}
د منځګړيتوب د قضيې يو مهم استعمال د بېت د تعريف‌وړتيا قضيه ده۔ د
يوې~\LR{$n$}-ځاييزې اړيکې~\LR{$R$} د تعريفولو دپاره داسې فارمول~\LR{$!C$}
ورکولے شو چې~\LR{$n$} ازاد متغيرونه ولري او~\LR{$R$} پکې نۀ راځي۔ دا به
د~\LR{$!C$} له مخې د~\LR{$R$} يو \emph{څرګند} تعريف وي۔ بيا دا هم ويلے شو چې
په داسې ژبه کښې يوه تيوري~\LR{$\Sigma(P)$} چې~\LR{$n$}-ځاييز
پريديکات~\LR{$P$} لري، \LR{$P$} په څرګنده توګه تعريفوي که يو صوري
څرګند تعريف ولري (يا لږ تر لږه ترې لازم راشي)، يعنې
\[
\Sigma(P) \Entails \lforall[x_1][\dots
  \lforall[x_n][(\Atom{P}{x_1,\dots, x_n} \liff !C(x_1, \dots,
    x_n))]].
\]
خو څرګند تعريف د يوې اړيکې د تعريفولو---يعنې په بشپړ ډول د هغې د
ټاکلو---يوازې يوه لار ده۔ يوه تيوري داسې هم کېدے شي چې په هر مدل
کښې د~\LR{$P$} تفسير د پاتې ژبه په تفسير سره ټاکل شوے وي۔ د
تعريف‌وړتيا قضيه وايي چې هر کله يوه تيوري د~\LR{$P$} تفسير په دې ډول
ټاکي---يعنې~\LR{$P$} \emph{په ضمني توګه تعريفوي}---نو هغه په څرګنده
توګه هم تعريفوي۔
\olpseganchor{OLP-0202-B006}{end}

\olpseganchor{OLP-0202-B007}{start}
\begin{defn}
فرض کړئ~\LR{$\Lang{L}$} داسې ژبه ده چې پريديکات~\LR{$P$} نۀ
لري۔ د~\LR{$\Lang{L} \cup \{P\}$} د تړلي فارمولونه يو سټ~\LR{$\Sigma(P)$}،
\LR{$P$} هله او يوازې هله \emph{په څرګنده توګه تعريفوي} چې د
\LR{$\Lang{L}$} داسې فارمول~\LR{$!C(x_1, \dots, x_n)$} وي چې
\[
\Sigma(P) \Entails \lforall[x_1][\dots
  \lforall[x_n][(\Atom{P}{x_1,\dots, x_n} \liff !C(x_1, \dots,
    x_n))]].
\]
\end{defn}
\olpseganchor{OLP-0202-B007}{end}

\olpseganchor{OLP-0202-B008}{start}
\begin{defn}
فرض کړئ~\LR{$\Lang{L}$} داسې ژبه ده چې پريديکاتونه~\LR{$P$} او
\LR{$P'$} نۀ لري۔ د~\LR{$\Lang{L} \cup \{P\}$} د تړلي فارمولونه يو
سټ~\LR{$\Sigma(P)$}، \LR{$P$} هله او يوازې هله \emph{په ضمني توګه تعريفوي}
چې
\[
\Sigma(P) \cup \Sigma(P') \Entails \lforall[x_1][\dots
  \lforall[x_n][(\Atom{P}{x_1,\dots, x_n} \liff \Atom{P'}{x_1,\dots,
      x_n})]],
\]
چې~\LR{$\Sigma(P')$} په~\LR{$\Sigma(P)$} کښې د~\LR{$P$} پر ځاے په يو شان ډول د
\LR{$P'$} له اېښودلو څخه ترلاسه کېږي۔
\end{defn}
\olpseganchor{OLP-0202-B008}{end}

\olpseganchor{OLP-0202-B009}{start}
په بله وينا، د هر مدل~\LR{$\Struct{M}$} او
\LR{$R, R' \subseteq \Domain{M}^n$} دپاره، که دواړه
\LR{$\Expan{M}{R} \Entails \Sigma(P)$} او
\LR{$\Expan{M}{R'} \Entails \Sigma(P')$} وي، نو~\LR{$R=R'$}؛ دلته
\LR{$\Expan{M}{R}$} هغه جوړښت~\LR{$\Struct{M'}$} دے چې~\LR{$\Lang{L}$}
\LR{$\Lang{L} \cup \{P\}$} ته غځوي او~\LR{$\Assign{P}{M'} = R$} کوي، او د
\LR{$\Expan{M}{R'}$} دپاره هم همداسې ده۔
\olpseganchor{OLP-0202-B009}{end}

\olpseganchor{OLP-0202-B010}{start}
\begin{thm}[د بېت د تعريف‌وړتيا قضيه]
د~\LR{$\Lang{L} \cup\{P\}$} د جملو يو سټ~\LR{$\Sigma(P)$}، \LR{$P$}
هله او يوازې هله په ضمني توګه تعريفوي چې~\LR{$\Sigma(P)$}، \LR{$P$} په
څرګنده توګه تعريف کړي۔
\textbf{د ژباړې د نسخې سمون (OLMOD-021):} سرچينې د قضيې په بيان
کښې~\LR{$\Sigma(P)$} د فارمولونه سټ بللے و، سره له دې چې دواړه مخکني
تعريفونه او ثبوت جملې کاروي؛ او د «هله او يوازې هله» دويم «هله»
ترې پاتې و۔ دلته «جملې» او بشپړه دوه اړخيزه رابطه ليکل شوې ده۔
\end{thm}
\olpseganchor{OLP-0202-B010}{end}

\olpseganchor{OLP-0202-B011}{start}
\begin{proof}
که~\LR{$\Sigma(P)$}، \LR{$P$} په څرګنده توګه تعريف کړي، نو دواړه
\begin{align*}
  \Sigma(P) & \Entails & \lforall[x_1][\dots \lforall[x_n]
    [(\Atom{P}{x_1,\dots, x_n} \liff !C(x_1,\dots,x_n))]]\\
  \Sigma(P') & \Entails & \lforall[x_1][\dots \lforall[x_n]
    [(\Atom{P'}{x_1,\dots, x_n} \liff !C(x_1,\dots,x_n))]]
\end{align*}
دي، او پايله ترې ترلاسه کېږي۔ د بل لوري دپاره فرض کړئ چې
\LR{$\Sigma(P)$}، \LR{$P$} په ضمني توګه تعريفوي۔ لومړے ثابت سمبولونه~\LR{$c_1$}،
\dots،~\LR{$c_n$} په~\LR{$\Lang{L}$} کښې ورزياتوو۔ بيا
\[
\Sigma(P) \cup \Sigma(P') \Entails
\Atom{P}{c_1, \dots, c_n} \to  \Atom{P'}{c_1, \dots, c_n}.
\]
د متناهي شاهد له مخې داسې متناهي سټونه
\LR{$\Delta_0 \subseteq \Sigma(P)$} او~\LR{$\Delta_1 \subseteq \Sigma(P')$}
شته چې
\[
\Delta_0 \cup \Delta_1 \Entails
\Atom{P}{c_1, \dots, c_n} \to \Atom{P'}{c_1, \dots, c_n}.
\]
پرېږدئ~\LR{$!D(P)$} د ټولو هغو تړلي فارمولونه~\LR{$!A(P)$} عطف وي چې يا
\LR{$!A(P) \in \Delta_0$} يا~\LR{$!A(P') \in \Delta_1$}؛ او~\LR{$!D(P')$} د
ټولو هغو تړلي فارمولونه~\LR{$!A(P')$} عطف وي چې يا
\LR{$!A(P) \in \Delta_0$} يا~\LR{$!A(P') \in \Delta_1$}۔ بيا
\LR{$!D(P) \land !D(P') \Entails \Atom{P}{c_1, \dots, c_n} \to
\Atom{P'}{c_1,\dots,c_n}$}۔
\textbf{د ژباړې د نسخې سمون (OLMOD-022):} سرچينې په دې پايله کښې
وروستے اټومي فارمول په خامه بڼه~\LR{$P'c_1\dots c_n$} ليکلے و؛ دلته د
همدې ثبوت له ثابتې نښې سره سم~\LR{$\Atom{P'}{c_1,\dots,c_n}$} ليکل
شوے دے۔
دا داسې بيا ترتيبولے شو چې هر پريديکات د~\LR{$\Entails$} په يوه
اړخ کښې راشي:
\[
!D(P) \land \Atom{P}{c_1, \dots, c_n} \Entails
!D(P') \to \Atom{P'}{c_1, \dots, c_n}.
\]
د کرېګ د منځګړيتوب له قضيې داسې تړلے فارمول
\LR{$!C(c_1,\dots, c_n)$} شته چې~\LR{$P$} يا~\LR{$P'$} نۀ لري او:
\begin{align*}
  !D(P) \land \Atom{P}{c_1, \dots, c_n} & \Entails !C(c_1,\dots, c_n); \\
  !C(c_1,\dots, c_n) & \Entails !D(P') \to \Atom{P'}{c_1, \dots, c_n}.
\end{align*}
له دې دوو استلزامونو له لومړي څخه لرو چې
\LR{$!D(P) \Entails \Atom{P}{c_1,\dots, c_n} \lif
!C(c_1,\dots, c_n)$}۔ او له دويم څخه، ځکه چې د~\LR{$\Lang{L} \cup
\{P\}$} يو مدل~\LR{$\Expan{M}{R} \Entails !A(P)$} هله او يوازې هله چې
ورسره اړوند د~\LR{$\Lang{L} \cup \{P'\}$} مدل
\LR{$\Expan{M}{R} \models !A(P')$} وي، لرو چې
\LR{$!C(c_1,\dots, c_n) \Entails !D(P) \lif
\Atom{P}{c_1,\dots, c_n}$}؛ له دې څخه:
\[
!D(P) \Entails !C(c_1,\dots,c_n) \to \Atom{P}{c_1,\dots, c_n}.
\]
دواړه سره يوځاے کول راکوي چې
\LR{$!D(P) \Entails \Atom{P}{c_1,\dots, c_n}
\liff !C(c_1, \dots, c_n)$}، او د يکنواختۍ او تعميم له مخې هم
\[
\Sigma(P) \Entails
\lforall[x_1][\dots\lforall[x_n][(\Atom{P}{x_1,\dots, x_n} \liff
    !C(x_1,\dots, x_n))]].
\]
\end{proof}
\olpseganchor{OLP-0202-B011}{end}

\olpseganchor{OLP-0203-B004}{start}
\olchapter{mod}{lin}{د ليندستروم قضيه}
\olpseganchor{OLP-0203-B004}{end}

\olpseganchor{OLP-0204-B004}{start}
\olfileid{mod}{lin}{int}
\section{پېژندګلو}
\olpseganchor{OLP-0204-B004}{end}

\olpseganchor{OLP-0204-B005}{start}
په دې څپرکي کښې موخه دا ده چې د لومړۍ درجې منطق په اړه د ليندستروم
ځانګړنه ثابته کړو: د هغو منطقونو په ډله کښې چې ځينې نور شرطونه هم
پوره کوي، دا هغه اعظمي منطق دے چې د متناهي شاهد او د لوېنهايم--سکولم
ښکته قضیې پکښې صدق کوي
(\olref[fol][com][com]{thm:compactness} او
\olref[fol][com][dls]{thm:downward-ls})۔ لومړے بايد د منطقونو هغه عمومي
ټولګی لا په عمومي ډول ځانګړے کړو چې قضيه پرې عملي کېږي۔ بحث به مو
تر \emph{اړيکيزو} ژبو محدود وي، يعنې هغو ژبو ته چې يوازې
پريديکاتونه او فردي ثابتونه لري، خو هېڅ تابعې نۀ لري۔
\olpseganchor{OLP-0204-B005}{end}

\olpseganchor{OLP-0205-B004}{start}
\olfileid{mod}{lin}{alg}
\olsection{انتزاعي منطقونه}
\olpseganchor{OLP-0205-B004}{end}

\olpseganchor{OLP-0205-B005}{start}
\begin{defn}
يو \emph{انتزاعي منطق} جوړه~\LR{$\tuple{L, \models_L}$} ده، چې~\LR{$L$} داسې
تابع ده چې هرې ژبه~\LR{$\Lang{L}$} ته د تړلي فارمولونه يو
سټ~\LR{$L(\Lang{L})$} ورکوي، او~\LR{$\models_L$} د~\LR{$\Lang{L}$} د
جوړښتونه او د~\LR{$L(\Lang{L})$} د غړي تر منځ اړيکه ده۔
په ځانګړي ډول،~\LR{$\tuple{F, \models}$} عادي د لومړۍ درجې منطق دے؛ يعنې
\LR{$F$} هغه تابع ده چې ژبه~\LR{$\Lang{L}$} ته د~\LR{$\Lang{L}$} له
ثابتونو څخه د جوړو شوو د لومړۍ درجې تړلي فارمولونه سټ ورکوي، او
\LR{$\models$} د لومړۍ درجې منطق د پوره کولو اړيکه ده۔
\end{defn}
\olpseganchor{OLP-0205-B005}{end}

\olpseganchor{OLP-0205-B006}{start}
پام وکړئ چې د يوې ورکړې شوې ژبه د جوړښت هماغه مفهوم
کاروو چې د لومړۍ درجې منطق دپاره ئې کاروو؛ خو دا له مخکښې نۀ فرضوو
چې تړلي فارمولونه د~\LR{$\Lang{L}$} له بنسټيزو نښو څخه په معمولي ډول جوړې
شوې دي، او نه دا چې~\LR{$\models_L$} اړيکه د لومړۍ درجې منطق په شان په
بازګشتي ډول تعريف شوې ده۔ نو دا بشپړ عمومي تعريف، د بېلګې په توګه،
هغه حالت هم رانغاړي چې د~\LR{$\tuple{L,\models_L}$} تړلي فارمولونه نامتناهي
اوږده تړونونه يا بېلتونونه ولري، له~\LR{$\lexists$} او~\LR{$\lforall$} پرته نور
کميت ټاکونکي ولري (لکه «نامتناهي ډېر داسې~\LR{$x$} شته چې \dots»)، يا
ښايي د کميت ټاکونکو نامتناهي اوږده مخ‌لړۍ ولري۔ د دې د ټينګار دپاره
چې په~\LR{$L(\Lang{L})$} کښې «تړلي فارمولونه» اړينه نۀ ده د لومړۍ درجې
منطق معمولي تړلي فارمولونه وي، په دې څپرکي کښې پرې~\LR{$!E$}،~\LR{$!F$}،~\dots
متغيرونه چلوو، او~\LR{$!A$}،~\LR{$!B$}،~\dots د لومړۍ درجې معمولي
فارمولونه ته ساتو۔
\olpseganchor{OLP-0205-B006}{end}

\olpseganchor{OLP-0205-B007}{start}
\begin{defn}
\LR{$\Mod(L){!E}$} دې ټولګی~\LR{$\Setabs{\Struct{M}}{\Struct{M}
  \models_L !E}$} وښيي۔ که ژبه په څرګنده ښودل پکار وي،
\LR{$\Mod[L](L){!E}$} ليکو۔ د~\LR{$\Lang{L}$} دوه جوړښتونه~\LR{$\Struct{M}$}
او~\LR{$\Struct{N}$} په~\LR{$\tuple{L, \models_L}$} کښې \emph{ابتدايي هم‌ارز}
دي، چې~\LR{$\Struct{M} \elemequiv[L] \Struct{N}$} ئې ليکو، که د
\LR{$L(\Lang{L})$} هماغه تړلي فارمولونه په دواړو کښې رښتونې وي۔
\end{defn}
\olpseganchor{OLP-0205-B007}{end}

\olpseganchor{OLP-0205-B008}{start}
\begin{defn}
د~ژبه~\LR{$\Lang{L}$} يو انتزاعي منطق~\LR{$\tuple{L,\models_L}$}
\emph{نورمال} دے که لاندې خاصيتونه پوره کړي:
\begin{enumerate}
\item (\emph{\LR{$L$}-زياتېدنه}) د~ژبې~\LR{$\Lang{L}$} او
  \LR{$\Lang{L'}$} دپاره، که~\LR{$\Lang{L} \subseteq \Lang{L'}$} وي، نو
  \LR{$L(\Lang{L}) \subseteq L(\Lang{L'})$}۔
\item (\emph{د غځونې خاصيت}) د هر~\LR{$!E \in L(\Lang{L})$} دپاره د
  \LR{$\Lang{L}$} يو \emph{متناهي} فرعي سټ~\LR{$\Lang{L'}$} شته، داسې چې
  اړيکه~\LR{$\Struct{M} \models_L !E$} يوازې د~\LR{$\Lang{L'}$} پر
  راکمونې پورې اړه لري؛ يعنې که~\LR{$\Struct{M}$} او~\LR{$\Struct{N}$} پر
  \LR{$\Lang{L'}$} يوه راکمونه ولري، نو~\LR{$\Struct{M} \models_L !E$} هله او
  يوازې هله چې~\LR{$\Struct{N} \models_L !E$} وي۔
\item (\emph{د آيزومورفيزم خاصيت}) که~\LR{$\Struct{M} \models_L !E$} او
  \LR{$\Struct{M} \simeq \Struct{N}$} وي، نو~\LR{$\Struct{N} \models_L !E$}
  هم وي۔
\item (\emph{د نوم‌اړونې خاصيت}) اړيکه~\LR{$\models_L$} د نوم‌اړونې له
  مخې ساتل کېږي: که ژبه~\LR{$\Lang{L}'$} له~\LR{$\Lang{L}$} څخه داسې
  لاس ته راشي چې هره نښه~\LR{$P$} د هماغې ځاييزې کچې په نښه~\LR{$P'$} او هر
  ثابت~\LR{$c$} په بېل ثابت~\LR{$c'$} بدل شي، نو د هر
  جوړښت~\LR{$\Struct{M}$} او هرې تړلے فارمول~\LR{$!E$} دپاره،
  \LR{$\Struct{M} \models_L !E$} هله او يوازې هله چې
  \LR{$\Struct{M}' \models_L !E'$} وي؛ دلته~\LR{$\Struct{M}'$} د~\LR{$\Struct{M}$}
  اړوند~\LR{$\Lang{L}'$}-جوړښت او~\LR{$!E' \in L(\Lang{L}')$} د~\LR{$!E$}
  نوم‌اړول شوې جمله ده۔
\item (\emph{بولي خاصيت}) انتزاعي منطق~\LR{$\tuple{L,
  \models_L}$} د بولي رابطونو له مخې بند دے، په دې معنا چې د هر
  \LR{$!E \in L(\Lang{L})$} دپاره داسې~\LR{$!F \in L(\Lang{L})$} شته چې
  \LR{$\Struct{M} \models_L !F$} هله او يوازې هله چې
  \LR{$\Struct{M} \not\models_L !E$} وي؛ او د هر~\LR{$!E$} او~\LR{$!F$} دپاره داسې
  \LR{$!G$} شته چې~\LR{$\Mod(L){!G} = \Mod(L){!E} \cap
  \Mod(L){!F}$} وي۔ د اټومي فارمولونه او نورو رابطونو دپاره هم
  همداسې ده۔
\item (\emph{د کميت ټاکونکي خاصيت}) په~\LR{$\Lang{L}$} کښې د هر ثابت~\LR{$c$}
  او هر~\LR{$!E \in L(\Lang{L})$} دپاره داسې~\LR{$!F \in L(\Lang{L})$} شته
  چې
  \[
  \Mod[L'](L){!F} = \Setabs{\Struct{M}}{\Expan{M}{a} \in
  \Mod[L](L){!E} \text{\RL{ چې }} a \in \Domain{M}},
  \]
  چې دلته~\LR{$\Lang{L'} = \Lang{L} \setminus \{c\}$} او
  \LR{$\Expan{M}{a}$} د~\LR{$\Struct{M}$} هغه~\LR{$\Lang{L}$}-غځونه ده چې~\LR{$a$}
  ثابت~\LR{$c$} ته ورکوي۔
\item (\emph{د نسبي کولو خاصيت}) يوه تړلے فارمول
  \LR{$!E \in L(\Lang{L})$}، يو نوے~\LR{$(n+1)$}-ځاييز محمول~\LR{$R$}، او
  داسې نښې~\LR{$c_1$}، \dots،~\LR{$c_n$} راکړل شي چې په~\LR{$\Lang{L}$} کښې نۀ وي؛
  بيا يوه تړلے فارمول~\LR{$!F \in L(\Lang{L} \cup
  \{R,c_1,\ldots,c_n\})$} شته چې د~\LR{$!E$} \emph{پر
  \LR{$\Atom{R}{x, c_1, \dots c_n}$} نسبي کول} بلل کېږي، او د هر
  جوړښت~\LR{$\Struct{M}$} دپاره:
  \[
  \Expan{M}{X, b_1, \ldots, b_n} \models_L !F \text{\RL{ هله او
    يوازې هله چې }} \Struct{N} \models_L !E,
  \]
  چې~\LR{$\Struct{N}$} د~\LR{$\Struct{M}$} هغه فرعي جوړښت دے چې د تعريف ساحه ئې
  \LR{$\Domain{N} = \Setabs{a\in \Domain{M}}{\Assign{R}{M}(a, b_1, \dots, b_n)}$}
  دے (\olref[bas][sub]{rem:substructure} وګورئ)، او
  \LR{$\Expan{M}{X, b_1, \ldots, b_n}$} د~\LR{$\Struct{M}$} هغه غځونه ده چې
  \LR{$R$}،~\LR{$c_1$}، \dots،~\LR{$c_n$} په ترتيب سره په~\LR{$X$}،~\LR{$b_1,$}
  \dots،~\LR{$b_n$} تفسيروي (چې~\LR{$X \subseteq M^{n+1}$})۔
\end{enumerate}
\textbf{د ژباړې د نسخې سمون (OLMOD-024):} سرچينې د نوم‌اړونې په
شرط کښې نوم‌اړول شوے جوړښت د اصلي ژبه اړوند بللے و؛ دلته هغه د اصلي
جوړښت اړوند وبلل شو، او نوم‌اړول شوې جمله هم په څرګنده وپېژندل شوه۔
\textbf{د ژباړې د نسخې سمون (OLMOD-025):} سرچينې د نسبي کولو په شرط
کښې د نوي محمول ځاييزه کچه نۀ وه ټاکلې؛ له ورپسې تعريف ساحې او غځونې
سره سم دلته هغه يو زيات له ورکړو ثابتونو، يعنې~\LR{$(n+1)$}-ځاييز، ټاکل
شوے دے۔
\end{defn}
\olpseganchor{OLP-0205-B008}{end}

\olpseganchor{OLP-0205-B009}{start}
\begin{defn}
دوه انتزاعي منطقونه~\LR{$\tuple{L_1, \models_{L_1}}$} او
\LR{$\tuple{L_2, \models_{L_2}}$} راکړل شي۔ وايو چې دويم \emph{لږ تر
لږه د لومړي هومره څرګندوونکی} دے، چې~\LR{$\tuple{L_1, \models_{L_1}}
\leq \tuple{L_2, \models_{L_2}}$} ئې ليکو، که د هرې
ژبه~\LR{$\Lang{L}$} او هرې تړلے فارمول
\LR{$!E \in L_1(\Lang{L})$} دپاره داسې تړلے فارمول
\LR{$!F \in L_2(\Lang{L})$} شته چې~\LR{$\Mod[L](L_1){!E} =
\Mod[L](L_2){!F}$} وي۔ منطقونه~\LR{$\tuple{L_1, \models_{L_1}}$} او
\LR{$\tuple{L_2, \models_{L_2}}$} \emph{هم‌ارز} دي که
\LR{$\tuple{L_1, \models_{L_1}} \leq \tuple{L_2, \models_{L_2}}$} او
\LR{$\tuple{L_2, \models_{L_2}} \leq \tuple{L_1, \models_{L_1}}$} وي۔
\end{defn}
\olpseganchor{OLP-0205-B009}{end}

\olpseganchor{OLP-0205-B010}{start}
\begin{rem}
  د لومړۍ درجې منطق، يعنې انتزاعي منطق~\LR{$\tuple{F, \models}$}، نورمال
  دے۔ په حقيقت کښې، پورته خاصيتونه تر ډېره د لومړۍ درجې منطق دپاره
  سملاسي دي۔ يوازې دومره يادوو چې د غځونې خاصيت له مصداقيت څخه
  راځي، او پر~\LR{$\Atom{R}{x, c_1, \dots, c_n}$} د تړلے فارمول~\LR{$!E$}
  نسبي کول د هرې فرعي فارمول~\LR{$\lforall[x][!F]$} په
  \LR{$\lforall[x][(\Atom{R}{x, c_1, \dots, c_n} \lif \ !F)]$} بدلولو
  سره تر لاسه کېږي۔ سربېره پر دې، که~\LR{$\tuple{L, \models_L}$} نورمال
  وي، نو~\LR{$\tuple{F, \models} \leq \tuple{L, \models_{L}}$}، لکه چې
  دا خبره د لومړۍ درجې پر فارمولونه په استقرا ښودل کېدے شي۔ نو له
  عموميت څخه بې له کوم زيان څخه فرض کولے شو چې د لومړۍ درجې هره
  تړلے فارمول د هر نورمال منطق غړې ده۔
\end{rem}
\olpseganchor{OLP-0205-B010}{end}

\olpseganchor{OLP-0206-B004}{start}
\olfileid{mod}{lin}{lsp}
\olpseganchor{OLP-0206-B004}{end}

\olpseganchor{OLP-0206-B005}{start}
\olsection{د متناهي شاهد او لوېنهايم--سکولم خاصيتونه}
\olpseganchor{OLP-0206-B005}{end}

\olpseganchor{OLP-0206-B006}{start}
اوس د متناهي شاهد او لوېنهايم--سکولم خاصيتونه د انتزاعي منطقونو حالت
ته په څرګند ډول غځوو۔
\olpseganchor{OLP-0206-B006}{end}

\olpseganchor{OLP-0206-B007}{start}
\begin{defn}
يو انتزاعي منطق~\LR{$\tuple{L, \models_L}$} \emph{د متناهي شاهد خاصيت}
لري که د~\LR{$L(\Lang{L})$}-تړلي فارمولونه هر سټ~\LR{$\Gamma$} هغه وخت د صدق
وړ وي چې هر متناهي~\LR{$\Gamma_0 \subseteq \Gamma$} ئې د صدق وړ وي۔
\end{defn}
\olpseganchor{OLP-0206-B007}{end}

\olpseganchor{OLP-0206-B008}{start}
\begin{defn}
\LR{$\tuple{L, \models_L}$} \emph{د لوېنهايم--سکولم ښکته خاصيت} لري که
هر د صدق وړ~\LR{$\Gamma$} يو د شمېر وړ مدل ولري۔
\end{defn}
\olpseganchor{OLP-0206-B008}{end}


\olpseganchor{OLP-0206-B009}{start}
د~\olref[bas][pis]{defn:partialisom} د جزوي آيزومورفيزم مفهوم سوچه
«الجبري» دے؛ يعنې د ژبې د تړلي فارمولونه له يادونې پرته او يوازې د
جوړښتونه د~ژبه~\LR{$\Lang{L}$} د ثابتونو له مخې ورکړل شوے
دے۔ نو دا مفهوم پر انتزاعي منطقونو هم عملي کېږي۔ د لومړۍ درجې منطق
په حالت کښې له~\olref[bas][pis]{thm:p-isom2} څخه پوهېږو چې که دوه
جوړښتونه په جزوي ډول آيزومورف وي، نو ابتدايي هم‌ارز دي۔ هغه
ثبوت انتزاعي منطقونو ته نۀ غځېږي، ځکه پر فارمولونه استقرا د خپل
سري~\LR{$!E \in L(\Lang{L})$} دپاره ښايي موجوده نۀ وي؛ خو که د
لوېنهايم--سکولم خاصيت صدق وکړي، قضيه بيا هم رښتونې ده۔
\olpseganchor{OLP-0206-B009}{end}

\olpseganchor{OLP-0206-B010}{start}
\begin{thm}
\ollabel{thm:abstract-p-isom}
فرض کړئ~\LR{$\tuple{L, \models_L}$} يو نورمال منطق دے چې د
لوېنهايم--سکولم خاصيت لري۔ بيا هر دوه په جزوي ډول آيزومورف
جوړښتونه په~\LR{$\tuple{L, \models_L}$} کښې ابتدايي هم‌ارز دي۔
\end{thm}
\olpseganchor{OLP-0206-B010}{end}

\olpseganchor{OLP-0206-B011}{start}
\begin{proof}
فرض کړئ~\LR{$\Struct{M} \simeq_p \Struct{N}$}، خو د کوم~\LR{$!E$} دپاره
\LR{$\Struct{M} \models_L !E$} او~\LR{$\Struct{N} \not\models_L !E$} هم وي۔ د
آيزومورفيزم د خاصيت له مخې فرض کولے شو چې~\LR{$\Domain{M}$} او
\LR{$\Domain{N}$} سره بېل دي؛ او د غځونې د خاصيت له مخې فرض کولے شو چې
\LR{$!E \in L(\Lang{L})$} د يوې متناهي ژبه~\LR{$\Lang{L}$} دپاره
وي۔ \LR{$\mathcal{I}$} دې د~\LR{$\Struct{M}$} او~\LR{$\Struct{N}$} تر منځ د جزوي
آيزومورفيزمونو يو سټ وي؛ له عموميت څخه بې له زيان څخه دا هم فرض کړئ
چې که~\LR{$p \in \mathcal{I}$} او~\LR{$q \subseteq p$} وي، نو
\LR{$q \in \mathcal{I}$} هم وي۔
\olpseganchor{OLP-0206-B011}{end}

\olpseganchor{OLP-0206-B012}{start}
\LR{$\Domain{M}^{<\omega}$} د~\LR{$\Domain{M}$} د غړي د متناهي لړيو
سټ دے۔ \LR{$S$} دې پر~\LR{$\Domain{M}^{<\omega}$} هغه درې‌ځاييزه اړيکه وي چې
نښلون ښيي؛ يعنې که~\LR{$\mathbf{a}, \mathbf{b},
\mathbf{c} \in \Domain{M}^{<\omega}$} وي، نو
\LR{$S(\mathbf{a}, \mathbf{b}, \mathbf{c})$} هله او يوازې هله صدق کوي
چې~\LR{$\mathbf{c}$} د~\LR{$\mathbf{a}$} او~\LR{$\mathbf{b}$} نښلون وي۔ همدارنګه،
\LR{$T$} دې هغه درې‌ځاييزه اړيکه وي چې د~\LR{$b \in M$} او
\LR{$\mathbf{a}, \mathbf{c} \in \Domain{M}^{<\omega}$} دپاره
\LR{$T(\mathbf{a}, b, \mathbf{c})$} هله او يوازې هله صدق کوي چې
\LR{$\mathbf{a} = a_1, \dots a_n$} او
\LR{$\mathbf{c} = a_1, \dots a_n, b$} وي۔ دوه نوي درې‌ځاييز
پريديکاتونه~\LR{$P$} او~\LR{$Q$} واخلئ، او جوړښت~\LR{$\Struct{M}^*$}
داسې جوړ کړئ چې جهان ئې~\LR{$\Domain{M} \cup \Domain{M}^{<\omega}$} وي،
\LR{$\Struct{M}$} د فرعي جوړښت په توګه ولري، او~\LR{$P$} او~\LR{$Q$} په نښلوونکو
اړيکو~\LR{$S$} او~\LR{$T$} تفسير کړي؛ نو~\LR{$\Struct{M}^*$} د
ژبه~\LR{$\Lang{L} \cup \{ P, Q\}$} جوړښت دے۔ د اړتيا پر وخت د
جهان د دوو برخو بېل نقلونه اخلو، څو دا اتحاد سره بېل کوډ شوے وي۔
\olpseganchor{OLP-0206-B012}{end}

\olpseganchor{OLP-0206-B013}{start}
\LR{$\Domain{N}^{<\omega}$}،~\LR{$S'$}،~\LR{$T'$}،~\LR{$P'$}،~\LR{$Q'$} او
\LR{$\Struct{N}^*$} په ورته ډول تعريف کړئ۔ ځکه چې د فرض له مخې
\LR{$\Struct{M} \simeq_p \Struct{N}$}، د~\LR{$\Domain{M}^{<\omega}$} او
\LR{$\Domain{N}^{<\omega}$} تر منځ داسې اړيکه~\LR{$I$} شته چې
\LR{$I(\mathbf{a}, \mathbf{b})$} هله او يوازې هله صدق کوي چې
\LR{$\mathbf{a}$} او~\LR{$\mathbf{b}$} آيزومورف وي او د
\olref[bas][pis]{defn:partialisom} د مخ او شا شرط پوره کړي۔ اوس
\LR{$\Struct{K}$} هغه جوړښت ونيسئ چې د تعريف ساحه ئې د
\LR{$\Struct{M}^*$} او~\LR{$\Struct{N}^*$} د د تعريف ساحې سره بېل کوډ شوے اتحاد
وي،~\LR{$\Struct{M}^*$} او~\LR{$\Struct{N}^*$} د فرعي جوړښتونه په توګه
ولري، او ژبه ئې يو اضافي دوه‌ځاييز پريديکات~\LR{$R$} ولري چې په
اړيکه~\LR{$I$} تفسيرېږي، او داسې پريديکاتونه هم ولري چې د دواړو برخو
د تعريف ساحې~\LR{$\Domain{M}^*$} او~\LR{$\Domain{N}^*$} ټاکي۔
\olpseganchor{OLP-0206-B013}{end}

\olpseganchor{OLP-0206-B014}{start}
\begin{figure}[h]
  \centering
  \begin{tikzpicture}[node distance=2cm, auto, thick, >=stealth']
    \draw [rounded corners] (0,0) -- (8,0) -- (8,4) -- (0,4) --  cycle;
    \draw (2,2) circle (0.5cm);
    \draw (2,2) circle (1.25cm);
    \draw (6,2) circle (0.5cm);
    \draw (6,2) circle (1.25cm);
    \path node at (0.75,3.5) {\large \LR{$\Struct{K}$}};
    \path node at (2,2) {\large \LR{$\Struct{M}$}};
    \path node at (6,2) {\large \LR{$\Struct{N}$}};
    \path node at (3.5,1) {\large \LR{$\Struct{M}^*$}};
    \path node at (7.5,1) {\large \LR{$\Struct{N}^*$}};
    \node (Idom) at (2.8,2) {};
    \node (Irng) at (5.2,2) {};
    \draw[<->, bend left] (Idom) to node {\large \LR{$I$}} (Irng) ;
  \end{tikzpicture}
  \caption{جوړښت~\LR{$\Struct{K}$} او د هغه دننه جزوي
    آيزومورفيزم۔}
\end{figure}
\olpseganchor{OLP-0206-B014}{end}

\olpseganchor{OLP-0206-B015}{start}
مهمه کتنه دا ده چې د~جوړښت~\LR{$\Struct{K}$} په ژبه کښې يوه د
انتزاعي منطق تړلے فارمول~\LR{$!D_1$} شته چې په~\LR{$\Struct{K}$} کښې رښتونې
ده او وايي~\LR{$\Struct{M} \models_L !E$} او
\LR{$\Struct{N} \not\models_L !E$}؛ د دې دپاره د نسبي کولو خاصيت پکار
دے۔ همدارنګه، يوه د لومړۍ درجې تړلے فارمول~\LR{$!D_2$} په~\LR{$\Struct{K}$}
کښې رښتونې ده او وايي چې~\LR{$\Struct{M} \simeq_p \Struct{N}$} د جزوي
آيزومورفيزم~\LR{$I$} له لارې۔ د لوېنهايم--سکولم د خاصيت له مخې،~\LR{$!D_1$}
او~\LR{$!D_2$} په ګډه په يوه د شمېر وړ مدل~\LR{$\Struct{K}_0$} کښې
رښتونې دي؛ په دې مدل کښې په جزوي ډول آيزومورف فرعي جوړښتونه
\LR{$\Struct{M}_0$} او~\LR{$\Struct{N}_0$} شته، داسې چې
\LR{$\Struct{M}_0 \models_L !E$} او~\LR{$\Struct{N}_0 \not\models_L !E$}۔ خو
د شمېر وړ او په جزوي ډول آيزومورف جوړښتونه د
\olref[bas][pis]{thm:p-isom1} له مخې په حقيقت کښې آيزومورف دي؛ دا د
نورمالو انتزاعي منطقونو د آيزومورفيزم له خاصيت سره تناقض دے۔
\textbf{د ژباړې د نسخې سمون (OLMOD-026):} سرچينې دواړه شاهدې جملې د
لومړۍ درجې بللې وې، حال دا چې لومړۍ جمله د خپل سري انتزاعي جملې د
نسبي کولو له لارې لاس ته راځي؛ دلته لومړۍ د انتزاعي منطق جمله او
دويمه د لومړۍ درجې جمله وبلل شوه۔
\textbf{د ژباړې د نسخې سمون (OLMOD-027):} سرچينې ګډ جوړښت ته د
پخواني چپ جوړښت هماغه نوم، او بيا ښکته مدل او چپ فرعي جوړښت ته هم
يو شان نوم ورکړے و؛ دلته ګډ جوړښت او د هغه شمېرېدونکی مدل په بېلو
نښو ښودل شوي، او په انځور کښې بهرنی نوم هم ورسره سم شوے دے۔
\end{proof}
\olpseganchor{OLP-0206-B015}{end}

\olpseganchor{OLP-0207-B004}{start}
\olfileid{mod}{lin}{prf}
\olpseganchor{OLP-0207-B004}{end}

\olpseganchor{OLP-0207-B005}{start}
\olsection{د ليندستروم قضيه}
\olpseganchor{OLP-0207-B005}{end}

\olpseganchor{OLP-0207-B006}{start}
\begin{lem}
\ollabel{lem:lindstrom}
فرض کړئ~\LR{$!E \in L(\Lang{L})$} او~\LR{$\Lang{L}$} متناهي ده؛ دا هم فرض
کړئ چې داسې~\LR{$n \in \Nat$} شته چې د هر دوو جوړښتونه
\LR{$\Struct{M}$} او~\LR{$\Struct{N}$} دپاره، که~\LR{$\Struct{M} \equiv_n
\Struct{N}$} او~\LR{$\Struct{M} \models_L !E$} وي، نو
\LR{$\Struct{N} \models_L !E$} هم وي۔ بيا~\LR{$!E$} د لومړۍ درجې له يوې
تړلے فارمول سره هم‌ارزه ده؛ يعنې د لومړۍ درجې داسې~\LR{$!D$} شته چې
\LR{$\Mod(L){!E} = \Mod(L){!D}$} وي۔
\end{lem}
\olpseganchor{OLP-0207-B006}{end}

\olpseganchor{OLP-0207-B007}{start}
\begin{proof}
\LR{$n$} داسې ونيسئ چې هر دوه~\LR{$n$}-هم‌ارز جوړښتونه~\LR{$\Struct{M}$} او
\LR{$\Struct{N}$}،~\LR{$!E$} ته د ورکړي ارزښت په اړه سره سلا وي۔
\olref[bas][pis]{prop:qr-finite} راياد کړئ: په يوه متناهي
ژبه کښې، له منطقي هم‌ارزۍ پرته، د لومړۍ درجې يوازې متناهي
ډېرې داسې تړلي فارمولونه شته چې د کميت ټاکونکو رتبه ئې تر~\LR{$n$} زياته
نۀ وي۔ اوس د هر ټاکلي جوړښت~\LR{$\Struct{M}$} دپاره
\LR{$!D_{\Struct{M}}$} د هغو د لومړۍ درجې تړلي فارمولونه~\LR{$!E$} د هر منطقي
هم‌ارزۍ له ټولګي څخه د يوه استازي تړون ونيسئ چې په~\LR{$\Struct{M}$} کښې
رښتونې وي او~\LR{$\QuantRank{!E} \le n$} ولري؛ دا تړون متناهي دے، او
\LR{$\Struct{N} \models !D_{\Struct{M}}$} هله او يوازې هله چې
\LR{$\Struct{N} \equiv_n \Struct{M}$} وي۔ بيا
\LR{$!D = \textstyle\bigvee
\Setabs{!D_{\Struct{M}}}{\Struct{M} \models_L !E}$} کېږدئ؛ له هر
منطقي هم‌ارزۍ ټولګي څخه د يوه استازي په ساتلو دا بېلتون هم متناهي
دے۔
\textbf{د ژباړې د نسخې سمون (OLMOD-028):} سرچينې د ټولو رښتونو
ټيټې مرتبې جملو تړون «متناهي» بللے و، که څه هم په هر منطقي هم‌ارزۍ
ټولګي کښې بې‌شمېره نحوي بڼې راتللے شي؛ دلته په څرګنده له هر ټولګي
څخه يو استازی اخيستل کېږي، او په وروستي بېلتون کښې هم تکراري هم‌ارزې
برخې لرې کېږي۔
\olpseganchor{OLP-0207-B007}{end}

\olpseganchor{OLP-0207-B008}{start}
پايله~\LR{$\Mod(L){!E} = \Mod(L){!D}$} تر لاسه کېږي۔ په حقيقت کښې، که
\LR{$\Struct{N} \models_L !D$} وي، نو د کوم
\LR{$\Struct{M} \models_L !E$} دپاره~\LR{$\Struct{N} \models
!D_{\Struct{M}}$}؛ نو د لمې د فرض له مخې
\LR{$\Struct{N} \models_L !E$} هم۔ برعکس، که
\LR{$\Struct{N} \models_L !E$} وي، نو~\LR{$!D_\Struct{N}$} په~\LR{$!D$} کښې يوه
بېلېدونکې برخه ده، او دا چې~\LR{$\Struct{N} \models !D_\Struct{N}$}، نو
\LR{$\Struct{N} \models_L !D$} هم۔
\end{proof}
\olpseganchor{OLP-0207-B008}{end}

\olpseganchor{OLP-0207-B009}{start}
\begin{thm}[د ليندستروم قضيه]
  \ollabel{thm:lindstrom} فرض کړئ~\LR{$\tuple{L, \models_L}$} يو نورمال
  انتزاعي منطق دے چې د متناهي شاهد او لوېنهايم--سکولم خاصيتونه لري۔
  بيا~\LR{$\tuple{L, \models_L} \le \tuple{F, \models}$}؛ نو
  \LR{$\tuple{L, \models_L}$} د لومړۍ درجې له منطق سره هم‌ارز دے۔
  \textbf{د ژباړې د نسخې سمون (OLMOD-029):} سرچينې د قضيې په وينا
  کښې د نورمالوالي شرط پرېښے و، حال دا چې ثبوت د غځونې،
  آيزومورفيزم، نسبي کولو او نورو نورمالو خاصيتونو کار اخلي، او معکوسه
  څرګندونکې اړيکه هم له همدې شرط څخه راځي۔
\end{thm}
\olpseganchor{OLP-0207-B009}{end}

\olpseganchor{OLP-0207-B010}{start}
\begin{proof}
د~\olref{lem:lindstrom} له مخې، دومره ښودل کافي دي چې د هر
\LR{$!E \in L(\Lang{L})$} دپاره، چې~\LR{$\Lang{L}$} متناهي وي، داسې
\LR{$n \in \Nat$} شته چې د هر دوو جوړښتونه~\LR{$\Struct{M}$} او
\LR{$\Struct{N}$} دپاره: که~\LR{$\Struct{M} \equiv_n \Struct{N}$} وي، نو
\LR{$\Struct{M}$} او~\LR{$\Struct{N}$} پر~\LR{$!E$} سره سلا وي۔ بيا~\LR{$!E$} د لومړۍ
درجې له يوې تړلے فارمول سره هم‌ارزه ده، او له دې څخه
\LR{$\tuple{L, \models_L} \le \tuple{F, \models}$} پايله کېږي۔ ځکه چې په
يوه متناهي او سوچه اړيکيزه ژبه کښې کار کوو، د
\olref[bas][pis]{thm:b-n-f} له مخې دا وينا چې
\LR{$\Struct{M} \equiv_n \Struct{N}$} ده، په اړوندې الجبري وينا
\LR{$I_n(\emptyset,\emptyset)$} بدلولے شو۔
\olpseganchor{OLP-0207-B010}{end}

\olpseganchor{OLP-0207-B011}{start}
يو~\LR{$!E$} راکړل شي، او د تناقض دپاره فرض کړئ چې د هر~\LR{$n$} دپاره داسې
جوړښتونه~\LR{$\Struct{M}_n$} او~\LR{$\Struct{N}_n$} شته چې
\LR{$I_n(\emptyset, \emptyset)$}، خو، د بېلګې په توګه،
\LR{$\Struct{M}_n \models_L !E$} او~\LR{$\Struct{N}_n \not\models_L !E$} وي۔
د آيزومورفيزم د خاصيت له مخې فرض کولے شو چې ټول~\LR{$\Struct{M}_n$} د
ژبه ثابتونه په هماغو څيزونو تفسيروي؛ سربېره پر دې، ځکه په ژبه کښې
يوازې متناهي ډېرې اټومي تړلي فارمولونه شته، فرض کولے شو چې هماغه
اټومي تړلي فارمولونه پوره کوي؛ که داسې نۀ وي، د~\LR{$\Struct{M}$}ګانو يوه
فرعي لړۍ اخيستلے شو۔ \LR{$\Struct{M}$} د ټولو~\LR{$\Struct{M}_n$}ګانو اتحاد
ونيسئ؛ يعنې هغه يوازينی کوچنی جوړښت چې هر~\LR{$\Struct{M}_n$} د
فرعي جوړښت په توګه لري۔ لکه د~\olref[lsp]{thm:abstract-p-isom} په
ثبوت کښې،~\LR{$\Struct{M}^*$} د~\LR{$\Struct{M}$} هغه غځونه ونيسئ چې د تعريف ساحه
ئې~\LR{$\Domain{M} \cup \Domain{M}^{<\omega}$} وي، او په پراخه شوې
ژبه کښې د نښلون محمولونه~\LR{$P$} او~\LR{$Q$} ولري۔
\olpseganchor{OLP-0207-B011}{end}

\olpseganchor{OLP-0207-B012}{start}
په ورته ډول~\LR{$\Struct{N}_n$}،~\LR{$\Struct{N}$} او~\LR{$\Struct{N}^*$} تعريف
کړئ۔ اوس~\LR{$\Struct{K}$} هغه جوړښت ونيسئ چې د تعريف ساحه ئې د
\LR{$\Struct{M}^*$} او~\LR{$\Struct{N}^*$} د د تعريف ساحې تر څنګ طبيعي
عددونه~\LR{$\Nat$} د خپل طبيعي ترتيب~\LR{$\le$} سره رانغاړي۔ ژبه ئې داسې
اضافي محمولونه لري چې د اړوندو د تعريف ساحې~\LR{$\Domain{M}$}،~\LR{$\Domain{N}$}،
\LR{$\Domain{M}^{<\omega}$} او~\LR{$\Domain{N}^{<\omega}$} ښيي، او داسې
محمولونه هم لري چې د~\LR{$\Struct{M}_n$} او~\LR{$\Struct{N}_n$} د تعريف
ساحې په لاندې معنا کوډوي:
\begin{align*}
  \Domain{M_n} & = \Setabs{a \in \Domain{M}}{R(a, n)}; &
  \Domain{N_n} & = \Setabs{a \in \Domain{N}}{S(a,n)}; \\
  \Domain{M}^{<\omega}_n & = \Setabs{a \in \Domain{M}^{<\omega}}{R(a,n)}; &
  \Domain{N}^{<\omega}_n & = \Setabs{a \in \Domain{N}^{<\omega}}{S(a,n)}.
\end{align*}
جوړښت~\LR{$\Struct{K}$} يوه درې‌ځاييزه اړيکه~\LR{$J$} هم لري، داسې چې
\LR{$J(n, \mathbf{a}, \mathbf{b})$} هله او يوازې هله صدق کوي چې
\LR{$I_n(\mathbf{a}, \mathbf{b})$} وي۔
\textbf{د ژباړې د نسخې سمون (OLMOD-030):} سرچينې دلته ګډ کوډوونکی
جوړښت د مخکېني اتحاد په هماغه نښه ښودلے، او وروسته ئې د متناهي شاهد
مدل هم د پخوانۍ غځونې په نښه ښودلے و؛ دلته ګډ جوړښت او د هغه نوی مدل
په بېلو نښو ښودل کېږي۔
\olpseganchor{OLP-0207-B012}{end}

\olpseganchor{OLP-0207-B013}{start}
اوس د~\LR{$\Lang{L}$} په هغه ژبه کښې چې په~\LR{$R$}،~\LR{$S$}،~\LR{$J$} او
نورو نښو پراخه شوې ده، داسې تړلے فارمول~\LR{$!D$} شته چې وايي:~\LR{$\le$}
يو منفصل خطي ترتيب دے چې لومړے خو وروستی غړے نۀ لري؛ د ترتيب
د هر~\LR{$n$} دپاره~\LR{$\Struct{M}_n \models_L !E$} او
\LR{$\Struct{N}_n \not\models_L !E$}؛ او
\LR{$J(n, \mathbf{a}, \mathbf{b})$} هله او يوازې هله چې
\LR{$I_n(\mathbf{a}, \mathbf{b})$} وي۔ د انتزاعي جملې اړوندې برخې د نسبي کولو او
د کميت ټاکونکي له خاصيتونو څخه تر لاسه کېږي۔
\olpseganchor{OLP-0207-B013}{end}

\olpseganchor{OLP-0207-B014}{start}
ژبه په نويو ثابتونو~\LR{$c$}،~\LR{$c_0$}،~\LR{$c_1$}،~\dots پراخه کړئ، او له
\LR{$!D$} سره د لومړۍ درجې هغه جملې يوځاے کړئ چې~\LR{$c_0$} د ترتيب لومړے
غړے ټاکي، هر~\LR{$c_{k+1}$} د~\LR{$c_k$} سملاسي جانشين ټاکي، او د هر معياري
\LR{$k \in \Nat$} دپاره~\LR{$c_k < c$} وايي۔ د دې تيورۍ هر متناهي فرعي سټ
په~\LR{$\Struct{K}$} کښې داسې پوره کېدے شي چې~\LR{$c$} د هغو متناهي ډېرو
يادو شوو جانشينانو څخه پورته يو طبيعي عدد وټاکي۔ د متناهي شاهد د
خاصيت له مخې ټوله تيوري يو مدل~\LR{$\Struct{K}^*$} لري۔ نو
\LR{$n^* = \Assign{c}{K^*}$} د هر معياري~\LR{$k$} دپاره د~\LR{$c_k$} له تفسير څخه
پورته دے، او د کوډ شوي منفصل ترتيب يو نامعياري غړے دے۔ په
ځانګړي ډول،~\LR{$\Struct{K}^*$} داسې فرعي جوړښتونه
\LR{$\Struct{M_{n^*}}$} او~\LR{$\Struct{N_{n^*}}$} لري چې
\LR{$\Struct{M_{n^*}} \models_L !E$} او
\LR{$\Struct{N_{n^*}} \not\models_L !E$}۔ اوس د~\LR{$\Domain{M_{n^*}}$} او
\LR{$\Domain{N_{n^*}}$} د~\LR{$k$}-توکيزو لړيو د جوړو يو سټ~\LR{$\mathcal{I}$}
داسې تعريفولے شو:~\LR{$\tuple{\mathbf{a}, \mathbf{b}} \in \mathcal{I}$}
هله او يوازې هله چې~\LR{$J(n^*-k, \mathbf{a}, \mathbf{b})$} وي؛ دلته~\LR{$k$}
د~\LR{$\mathbf{a}$} او~\LR{$\mathbf{b}$} اوږدوالی دے۔ ځکه~\LR{$n^*$} نامعياري
دے، د هر معياري~\LR{$k$} دپاره~\LR{$n^* - k >0$}؛ نو سټ~\LR{$\mathcal{I}$} د دې
خبرې شاهد دے چې~\LR{$\Struct{M_{n^*}} \simeq_p
\Struct{N_{n^*}}$}۔ خو د~\olref[lsp]{thm:abstract-p-isom} له مخې
\LR{$\Struct{M_{n^*}}$} له~\LR{$\Struct{N_{n^*}}$} سره~\LR{$L$}-هم‌ارز دے، او دا
تناقض دے۔
\textbf{د ژباړې د نسخې سمون (OLMOD-031):} سرچينې ادعا کړې وه چې د
متناهي شاهد خاصيت د يوازېنۍ جملې داسې مدل ورکوي چې نامعياري پړاو
ولري؛ هغه جمله پخپله داسې پړاو نۀ لازموي۔ دلته ژبه په يو پورني ثابت
او د معياري جانشينانو په نومونو غځول شوې، هر متناهي فرعي سټ په اصلي
کوډوونکي جوړښت کښې پوره کېږي، او متناهي شاهد ټولې تيورۍ ته له هر
معياري پړاو څخه پورته غړے ورکوي۔ د انتزاعي جملې د پوره کولو نښې هم
له اړوند منطق سره څرګندې شوې دي۔
\end{proof}
\olpseganchor{OLP-0207-B014}{end}

\olpseganchor{OLP-0208-B004}{start}
\olpart{cmp}{محاسبه کېدنه}
\olpseganchor{OLP-0208-B004}{end}

\olpseganchor{OLP-0208-B005}{start}
\begin{editorial}
  دا برخه د محاسبه کېدنې د تيورۍ په اړه د جرېمي اېوېګېډ پر يادښتونو
  ولاړه ده۔ اوس‌مهال يوازې د بازګشتي تابعو څپرکی تمرينونه لري، او
  ټوله برخه لا د هڅونې، بېلګو، جزيياتو او تمرينونو په زياتولو سره
  پراخېدو ته اړتيا لري۔
\end{editorial}
\olpseganchor{OLP-0208-B005}{end}

\olpseganchor{OLP-0209-B004}{start}
\olchapter{cmp}{rec}{بازګشتي تابعې}
\olpseganchor{OLP-0209-B004}{end}

\olpseganchor{OLP-0209-B005}{start}
\begin{editorial}
  دا د بازګشتي تابعو په اړه د جرېمي اېوېګېډ يادښتونه دي چې ريچرډ
  زاک بيا کتلي او پراخ کړي دي۔ دا څپرکی ځينې تمرينونه لري، او په
  خپلواکه توګه هم شاملېدے شي څو د نحو د حسابي کولو د بحث بنسټ برابر
  کړي۔
\end{editorial}
\olpseganchor{OLP-0209-B005}{end}

\olpseganchor{OLP-0210-B004}{start}
\olfileid{cmp}{rec}{int}
\olsection{پېژندګلو}
\olpseganchor{OLP-0210-B004}{end}

\olpseganchor{OLP-0210-B005}{start}
د محاسبه کېدنې د يوې رياضي تيورۍ د جوړولو دپاره، تر هر څه وړاندې د
محاسبه کېدنې يو \emph{مدل} جوړول پکار دي۔ اوس موږ محاسبه کېدنه د هغه
ډول کار په توګه ګڼو چې کمپيوټرونه ئې کوي، او کمپيوټرونه له نښو سره
کار کوي۔ خو د محاسبه کېدنې د تيوريو د پرمختګ په پيل کښې د محاسبې
نمونه‌يي بېلګه \emph{عددي} محاسبه وه۔ رياضيپوهانو تل د عددونو د
تيورۍ له تابعو سره مينه لرلې، يعنې له داسې تابعو \LR{$f\colon \Nat^n \to
\Nat$} سره چې محاسبه کېدے شي۔ نو دا د حيرانتيا خبره نۀ ده چې د
محاسبه کېدنې د تيورۍ په پيل کښې همدا تابعې څېړل کېدې۔ د محاسبه
کېدونکو عددي تابعو تر ټولو پېژندل شوې بېلګې، لکه جمع، ضرب او توان
(د طبيعي عددونو)، يو په زړه پورې خاصيت لري: دا په
\emph{بازګشتي} ډول تعريفېدے شي۔ ځکه نو دا ډېره طبيعي ده چې د
بازګشتي تعريفونو پر بنسټ د \emph{محاسبه کېدونکې تابع} يو عمومي تعريف
جوړ کړو۔ د عددونو د تيورۍ د تابعو د بازګشتي تعريف له ډېرو ممکنو لارو
څخه دلته يوه په ځانګړي ډول ساده نمونه مرکزي کېږي: تش په نامه
\emph{بنسټيز بازګښت}۔
\olpseganchor{OLP-0210-B005}{end}

\olpseganchor{OLP-0210-B006}{start}
له محاسبه کېدونکو تابعو سره سره ښايي محاسبه کېدونکي سټونه او اړيکې
هم راته د پام وړ وي۔ يو سټ هغه وخت محاسبه کېدونکی دے چې دا محاسبه
کولے شو چې يو ورکړل شوے عدد د سټ غړے دے که نۀ؛ او يوه
اړيکه هله او يوازې هله محاسبه کېدونکې ده چې دا محاسبه کولے شو چې
لړۍ \LR{$\tuple{n_1, \dots, n_k}$} د اړيکې غړے ده که نۀ۔ د يو
سټ يا اړيکې د \emph{ځانګړونکې تابع} په پام کښې نيولو سره، د محاسبه
کېدونکو سټونو او اړيکو بحث د محاسبه کېدونکو تابعو تر بحث لاندې
راوستل کېدے شي۔ په دې توګه بنسټيزې بازګشتي اړيکې هم تعريفولے شو؛ د
بېلګې په توګه، دا اړيکه چې «\LR{$n$}،~\LR{$m$} بې له پاتې شوني وېشي» يوه
بنسټيزه بازګشتي اړيکه ده۔
\olpseganchor{OLP-0210-B006}{end}

\olpseganchor{OLP-0210-B007}{start}
خو بنسټيزې بازګشتي تابعې---يعنې هغه چې له بنسټيزو تابعو څخه د
ترکيب او بنسټيز بازګښت په کارولو تعريفېدے شي---يوازينۍ محاسبه
کېدونکې د عددونو د تيورۍ تابعې نۀ دي۔ د بنسټيز بازګښت ډېرې عمومونې
څېړل شوې دي، خو د تابعو د محاسبې تر ټولو پياوړې او په پراخه کچه منل
شوې زياته لاره بې‌پوله پلټنه ده۔ دا د \emph{جزوي بازګشتي تابعو}
تعريف ته، او له دې سره تړلي د \emph{عمومي بازګشتي تابعو} تعريف ته
لار هواروي۔ عمومي بازګشتي تابعې محاسبه کېدونکې او هرځاے تعريف شوې
دي، او دا تعريف په کره توګه هغه جزوي بازګشتي تابعې ځانګړې کوي چې
هرځاے تعريف شوې وي۔ بازګشتي تابعې د محاسبې د هر بل مدل (ټيورينګ
ماشينونو، لمبډا حساب، او داسې نورو) تقليد کولے شي؛ ځکه نو د محاسبې
له ډېرو منل شوو مدلونو څخه يو مدل وړاندې کوي۔
\textbf{د ژباړې د نسخې سمون (OLCMP-001):} سرچينې بنسټيزې بازګشتي
تابعې يوازې د بنسټيز بازګښت په کارولو تعريفېدونکې بللې وې؛ خو رسمي
تعريف بنسټيزې تابعې او ترکيب هم غواړي۔ دلته لنډه ځانګړنه له رسمي
تعريف سره سمه شوې ده۔
\olpseganchor{OLP-0210-B007}{end}

\olpseganchor{OLP-0211-B004}{start}
\olfileid{cmp}{rec}{pre}
\olsection{بنسټيز بازګښت}
\olpseganchor{OLP-0211-B004}{end}

\olpseganchor{OLP-0211-B005}{start}
د طبيعي عددونو يو خاصيت دا دے چې هر طبيعي عدد پر~\LR{$0$} د راتلونکي
عمل~\LR{$+1$} په متناهي څو ځله کارولو سره تر لاسه کېدے شي---هر طبيعي عدد
يا~\LR{$0$} دے، يا د \dots{} د~\LR{$0$} راتلونکی دے۔ د دې حقيقت په کارولو سره
د يوې تابع~\LR{$h\colon\Nat \to \Nat$} د ټاکلو يوه لاره داسې ده:
(الف)~وټاکئ چې د آرګومېنټ~\LR{$0$} دپاره د~\LR{$h$} ارزښت څۀ دے، او
(ب)~دا هم وټاکئ چې د~\LR{$h(x)$} ارزښت له ورکړې وروسته د~\LR{$h(x+1)$} ارزښت
څنګه محاسبه شي۔ ځکه (الف) په نېغه راته وايي چې~\LR{$h(0)$} څۀ دے، نو
\LR{$h$} د~\LR{$0$} دپاره تعريف شوه۔ اوس د~\LR{$x=0$} دپاره د (ب) لارښوونه په
کارولو سره، له~\LR{$h(0)$} څخه \LR{$h(1) = h(0+1)$} محاسبه کولے شو۔ هماغه
لارښوونه د~\LR{$x=1$} دپاره په کارولو سره، له~\LR{$h(1)$} څخه
\LR{$h(2) = h(1+1)$} محاسبه کوو، او همداسې مخکښې ځو۔ د هر طبيعي
عدد~\LR{$x$} دپاره به بالاخره هغه پړاو ته ورسېږو چې~\LR{$h(x+1)$} له~\LR{$h(x)$}
څخه تعريفوو؛ نو~\LR{$h(x)$} د ټولو~\LR{$x \in \Nat$} دپاره تعريف شوې ده۔
\textbf{د ژباړې د نسخې سمون (OLCMP-002):} سرچينې د دې جملې د
بازګښت لوری سرچپه ليکلی و او اوسنی ارزښت ئې له راتلونکي ارزښت څخه
تعريفاوه؛ دلته د مخکښې ورکړې لارښوونې له مخې راتلونکی ارزښت له اوسني
ارزښت څخه تعريف شوے دے۔
\olpseganchor{OLP-0211-B005}{end}

\olpseganchor{OLP-0211-B006}{start}
د بېلګې په توګه، فرض کړئ~\LR{$h\colon \Nat \to \Nat$} په لاندې دوو
معادلو ټاکو:
\begin{align*}
h(0) & =  1\\
h(x+1) & =  2 \cdot h(x)
\end{align*}
که د ضرب طريقه له مخکښې را معلومه وي، نو دا معادلې د پورته (الف)
او~(ب) دپاره اړين معلومات راکوي۔ د دويمې معادلې پرله‌پسې په کارولو
سره لاس ته راوړو چې
\begin{align*}
  h(1) & = 2\cdot h(0) = 2,\\
  h(2) & = 2\cdot h(1) = 2\cdot 2,\\
  h(3) & = 2 \cdot h(2) = 2\cdot 2 \cdot 2,\\
  & \vdots
\end{align*}
وينو چې هغه تابع~\LR{$h$} چې ټاکلې مو ده، \LR{$h(x) = 2^x$} ده۔
\olpseganchor{OLP-0211-B006}{end}

\olpseganchor{OLP-0211-B007}{start}
د طبيعي عددونو ځانګړے خاصيت ضمانت کوي چې يوازې يوه تابع~\LR{$h$} شته
چې دا دواړه معيارونه پوره کوي۔ د دې غوندې معادلو جوړې ته د
تابع~\LR{$h$} \emph{په بنسټيز بازګښت تعريف} وايي۔ دې ته ځکه
«بازګشتي» وايي چې~\LR{$h$} په بازګشتي ډول تعريفوو، يعنې تعريف، او په
ځانګړي ډول دويمه معادله، پخپله~\LR{$h$}
په ښي اړخ کښې لري۔ دې ته ځکه «بنسټيز» وايي چې د~\LR{$h(x+1)$} په تعريف
کښې يوازې د~\LR{$h(x)$} ارزښت، يعنې سملاسي مخکينی ارزښت، کاروو۔ دا
پر~\LR{$\Nat$} د يوې تابع د بازګشتي تعريف تر ټولو ساده لاره ده۔
\olpseganchor{OLP-0211-B007}{end}

\olpseganchor{OLP-0211-B008}{start}
د جمع او ضرب غوندې لا بنسټيزې تابعې هم په بنسټيز بازګښت تعريفولے
شو۔ خو په دې حالتونو کښې اړوندې تابعې~\LR{$2$}-ځاييزې دي۔ د آرګومېنټونو
يو ځاے ثابت نيسو او بل د بازګښت دپاره کاروو۔ د بېلګې په توګه، د
\LR{$\Add(x, y)$} د تعريف دپاره~\LR{$x$} ثابت نيولے شو، ارزښت لومړے د~\LR{$y=0$}
دپاره تعريفولے شو، او بيا د~\LR{$y+1$} دپاره ئې د~\LR{$y$} له مخې تعريفولے
شو۔ ځکه چې~\LR{$x$} ثابت دے، د تعريفي معادلو په کيڼ او ښي اړخ دواړو
کښې به راځي۔
\begin{align*}
\Add(x,0) & =  x\\
\Add(x,y+1) & =  \Add(x,y)+1
\end{align*}
دا معادلې د ټولو~\LR{$x$} \emph{او}~\LR{$y$} دپاره د~\LR{$\Add$} ارزښت ټاکي۔ د
بېلګې په توګه، د~\LR{$\Add(2,3)$} د موندلو دپاره تعريفي معادلې د~\LR{$x = 2$}
دپاره کاروو: لومړۍ معادله کاروو چې~\LR{$\Add(2,0) = 2$} ومومو، بيا
دويمه معادله پرله‌پسې کاروو چې~\LR{$\Add(2,1) = 2 + 1 = 3$}،
\LR{$\Add(2, 2) = 3 + 1 = 4$}، او~\LR{$\Add(2, 3) = 4 + 1 = 5$} ومومو۔
\olpseganchor{OLP-0211-B008}{end}

\olpseganchor{OLP-0211-B009}{start}
د~\LR{$\Add$} په تعريف کښې مو د دويمې معادلې په ښي اړخ کښې~\LR{$+$} وکاراوه،
خو يوازې د~\LR{$1$} د زياتولو دپاره۔ په بله وينا، د راتلونکي
تابع~\LR{$\Succ(z) = z+1$} مو وکاروله او پر مخکيني ارزښت~\LR{$\Add(x,y)$} مو
عملي کړه، څو~\LR{$\Add(x,y+1)$} تعريف کړو۔ نو بازګشتي تعريف داسې ګڼلے شو
چې د يوې تابع له مخې ورکړل شوے دے او موږ ئې پر مخکيني ارزښت عملي
کوو۔ خو که تابع يوازې پر مخکيني ارزښت نۀ، بلکې پر~\LR{$x$} او~\LR{$y$} هم
تابع وګرځوو، زيان نۀ لري---او کله ناکله اړينه وي۔ دا تعريف وګورئ:
\begin{align*}
  \Mult(x,0) & =  0 \\
  \Mult(x,y+1) & =  \Add(\Mult(x,y),x)
\end{align*}
دا د تابع~\LR{$\Mult$} يو بنسټيز بازګشتي تعريف دے چې تابع~\LR{$\Add$} هم پر
مخکيني ارزښت~\LR{$\Mult(x,y)$} او هم پر لومړي آرګومېنټ~\LR{$x$} عملي کوي۔
دا هم تابع~\LR{$\Mult(x,y)$} د ټولو آرګومېنټونو~\LR{$x$} او~\LR{$y$} دپاره
تعريفوي۔ د بېلګې په توګه، \LR{$\Mult(2,3)$} د~\LR{$\Mult(2,0)$}،
\LR{$\Mult(2,1)$}، \LR{$\Mult(2,2)$} او~\LR{$\Mult(2,3)$} په پرله‌پسې محاسبې
سره ټاکل کېږي:
\begin{align*}
  \Mult(2,0) & = 0\\
  \Mult(2,1) & = \Mult(2,0+1) =
  \Add(\Mult(2,0), 2) = \Add(0, 2) = 2\\
  \Mult(2,2) & = \Mult(2,1+1) =
  \Add(\Mult(2,1), 2) = \Add(2, 2) = 4\\
  \Mult(2,3) & = \Mult(2,2+1) =
  \Add(\Mult(2,2), 2) = \Add(4, 2) = 6
\end{align*}
\olpseganchor{OLP-0211-B009}{end}

\olpseganchor{OLP-0211-B010}{start}
نو عمومي نمونه داسې ده: د يوې تابع~\LR{$h(x_0, \dots, x_{k-1}, y)$}
د بنسټيز بازګشتي تعريف دپاره دوه معادلې ورکوو۔ لومړۍ د~\LR{$h(x_0,
\dots, x_{k-1}, 0)$} ارزښت~\LR{$h$} ته له اشارې پرته تعريفوي۔ دويمه د
\LR{$h(x_0, \dots, x_{k-1}, y+1)$} ارزښت د~\LR{$h(x_0, \dots,
x_{k-1}, y)$}، نورو آرګومېنټونو~\LR{$x_0$}، \dots،~\LR{$x_{k-1}$}، او~\LR{$y$}
له مخې تعريفوي۔ په دې دويمه معادله کښې يوازې د~\LR{$h$} سملاسي مخکينی
ارزښت کارېدے شي۔ که د دې دوو معادلو د ښي اړخونو ورکړي عملونه هم
تابعې~\LR{$f$} او~\LR{$g$} وګڼو، نو په بنسټيز بازګښت د نوې تابع~\LR{$h$} د تعريف
عمومي نمونه داسې ده:
\begin{align*}
  h(x_0, \dots, x_{k-1}, 0) & = f(x_0, \dots, x_{k-1})\\
  h(x_0, \dots, x_{k-1}, y+1) & =
  g(x_0, \dots, x_{k-1}, y, h(x_0, \dots, x_{k-1}, y))
\end{align*}
د~\LR{$\Add$} په حالت کښې، \LR{$k=1$} او~\LR{$f(x_0) = x_0$} (عينيت تابع) دي،
او~\LR{$g(x_0, y, z) = z + 1$} (هغه~\LR{$3$}-ځاييزه تابع چې د خپل درېيم
آرګومېنټ راتلونکی ورکوي) ده:
\begin{align*}
  \Add(x_0, 0) & = f(x_0) = x_0\\
  \Add(x_0, y+1) & = g(x_0, y, \Add(x_0, y)) =
  \Succ(\Add(x_0, y))
\end{align*}
د~\LR{$\Mult$} په حالت کښې، \LR{$f(x_0) = 0$} (هغه ثابت تابع چې تل~\LR{$0$}
ورکوي) او~\LR{$g(x_0, y, z) = \Add(z,x_0)$} (هغه~\LR{$3$}-ځاييزه تابع چې د
خپل وروستي او لومړي آرګومېنټ مجموعه ورکوي) دي:
\begin{align*}
  \Mult(x_0,0) & =  f(x_0) = 0 \\
  \Mult(x_0,y+1) & =  g(x_0, y, \Mult(x_0,y)) =
  \Add(\Mult(x_0,y), x_0)
\end{align*}
\olpseganchor{OLP-0211-B010}{end}

\olpseganchor{OLP-0212-B004}{start}
\olfileid{cmp}{rec}{com}
\olsection{ترکيب}
\olpseganchor{OLP-0212-B004}{end}

\olpseganchor{OLP-0212-B005}{start}
که~\LR{$f$} او~\LR{$g$} د طبيعي عددونو دوه يوځاييزې تابعې وي، نو ترکيب ئې
کولے شو: \LR{$h(x) = f(g(x))$}۔ په دې حالت کښې نوې تابع~\LR{$h(x)$} له
تابعو~\LR{$f$} او~\LR{$g$} څخه په \emph{ترکيب} تعريف شوې ده۔ غواړو دا د يو
څخه د زياتو آرګومېنټونو تابعو ته عمومي کړو۔
\olpseganchor{OLP-0212-B005}{end}

\olpseganchor{OLP-0212-B006}{start}
د دې يوه لاره داسې ده: فرض کړئ~\LR{$f$} يوه~\LR{$k$}-ځاييزه تابع ده، او
\LR{$g_0$}، \dots،~\LR{$g_{k-1}$} هغه~\LR{$k$} تابعې دي چې ټولې~\LR{$n$}-ځاييزې دي۔
بيا يوه نوې~\LR{$n$}-ځاييزه تابع~\LR{$h$} داسې تعريفولے شو:
\[
h(x_0, \dots, x_{n-1}) =
f(g_0(x_0, \dots, x_{n-1}), \dots, g_{k-1}(x_0, \dots, x_{n-1}))
\]
که~\LR{$f$} او ټولې~\LR{$g_i$} محاسبه کېدونکې وي، نو~\LR{$h$} هم محاسبه کېدونکې
ده: د~\LR{$h(x_0, \dots, x_{n-1})$} د محاسبې دپاره لومړے د هر
\LR{$i = 0$}، \dots،~\LR{$k-1$} دپاره ارزښتونه~\LR{$y_i = g_i(x_0, \dots,
x_{n-1})$} محاسبه کړئ۔ بيا دا ارزښتونه~\LR{$f$} ته ورکړئ، څو
\LR{$h(x_0, \dots, x_{n-1}) = f(y_0, \dots, y_{k-1})$} محاسبه کړئ۔
\textbf{د ژباړې د نسخې سمون (OLCMP-003):} سرچينې په وروستۍ محاسبې
کښې د تابع د آرګومېنټونو لړۍ د تابعو د شمېر له مخې پای ته رسولې وه؛
دلته لړۍ د تابع د ځاييزې کچې له مخې بشپړه شوې ده۔
\olpseganchor{OLP-0212-B006}{end}

\olpseganchor{OLP-0212-B007}{start}
ښايي دا د دې خبرې له حده زيات محدود بيان ښکاره شي چې له ځينو شته
تابعو څخه د نوې تابع د محاسبې پر مهال څۀ پېښېږي۔ يوه خبره دا ده چې
کله ناکله د يوې تابع ټول آرګومېنټونه نۀ کاروو؛ لکه هغه وخت چې د
\LR{$\Add$} په بنسټيز بازګشتي تعريف کښې مو د کارولو دپاره
\LR{$g(x, y, z) = \Succ(z)$} تعريف کړه۔ فرض کړئ د لاندې تابعو د کارولو
اجازه لرو:
\[
\Proj{n}{i}(x_0, \dots, x_{n-1}) = x_i
\]
تابعو~\LR{$\Proj{n}{i}$} ته \emph{پروجکشن} تابعې وايي:
\LR{$\Proj{n}{i}$} يوه~\LR{$n$}-ځاييزه تابع ده۔ بيا~\LR{$g$} داسې تعريفېدے شي:
\[
g(x, y, z) = \Succ(\Proj{3}{2}(x, y, z)).
\]
دلته د~\LR{$f$} ونډه~\LR{$1$}-ځاييزه تابع~\LR{$\Succ$} ترسره کوي، نو~\LR{$k=1$} دے۔ او
يوه~\LR{$3$}-ځاييزه تابع~\LR{$\Proj{3}{2}$} لرو چې د~\LR{$g_0$} ونډه ترسره کوي۔
پايله يوه~\LR{$3$}-ځاييزه تابع ده چې د درېيم آرګومېنټ راتلونکی ورکوي۔
\textbf{د ژباړې د نسخې سمون (OLCMP-004):} سرچينې د پروجکشن تابعو
له تعريف سملاسي وروسته په نوم کښې د تابعو شمېر د ځاييزې کچې پر ځای
کارولی و؛ دلته نوم له هماغې تعريف شوې ځاييزې کچې سره سم شوے دے۔
\olpseganchor{OLP-0212-B007}{end}

\olpseganchor{OLP-0212-B008}{start}
پروجکشن تابعې دا اجازه هم راکوي چې د آرګومېنټونو په بيا ترتيبولو يا
سره يو کولو نوې تابعې تعريف کړو۔ د بېلګې په توګه، تابع
\LR{$h(x) = \Add(x, x)$} داسې تعريفېدے شي:
\[
h(x_0) = \Add(\Proj{1}{0}(x_0),\Proj{1}{0}(x_0)).
\]
دلته~\LR{$k=2$}، \LR{$n=1$}، د~\LR{$f(y_0,y_1)$} ونډه~\LR{$\Add$} ترسره کوي، او د
\LR{$g_0(x_0)$} او~\LR{$g_1(x_0)$} دواړو ونډې~\LR{$\Proj{1}{0}(x_0)$}، يعنې
يوځاييزه پروجکشن تابع (يا عينيت تابع)، ترسره کوي۔
\olpseganchor{OLP-0212-B008}{end}

\olpseganchor{OLP-0212-B009}{start}
که~\LR{$f(y_0, y_1)$} هغه تابع وي چې له مخکښې ئې لرو، نو تابع
\LR{$h(x_0, x_1) = f(x_1, x_0)$} داسې تعريفولے شو:
\[
h(x_0, x_1) = f(\Proj{2}{1}(x_0, x_1),\Proj{2}{0}(x_0, x_1)).
\]
دلته~\LR{$k=2$}، \LR{$n = 2$}، او د~\LR{$g_0$} او~\LR{$g_1$} ونډې په ترتيب سره
\LR{$\Proj{2}{1}$} او~\LR{$\Proj{2}{0}$} ترسره کوي۔
\olpseganchor{OLP-0212-B009}{end}

\olpseganchor{OLP-0212-B010}{start}
ښايي دا اندېښنه هم ولرئ چې د ټولو~\LR{$g_0$}، \dots،~\LR{$g_{k-1}$} ځاييزه
کچه بايد يو شان~\LR{$n$} وي۔ (په ياد ولرئ چې د يوې تابع
\emph{ځاييزه کچه} د آرګومېنټونو شمېر دے؛ د~\LR{$n$}-ځاييزې تابع ځاييزه
کچه~\LR{$n$} ده۔) خو د پروجکشن تابعو زياتول اړينه انعطاف‌پذيري برابروي۔
د بېلګې په توګه، فرض کړئ~\LR{$f$} او~\LR{$g$}، \LR{$3$}-ځاييزې تابعې دي او~\LR{$h$}
هغه~\LR{$2$}-ځاييزه تابع ده چې داسې تعريف شوې ده:
\[
h(x,y) = f(x,g(x,x,y),y).
\]
د~\LR{$h$} تعريف له پروجکشن تابعو سره داسې بيا ليکل کېدے شي:
\[
h(x,y) = f(\Proj{2}{0}(x,y), g(\Proj{2}{0}(x,y), \Proj{2}{0}(x,y),
\Proj{2}{1}(x,y)), \Proj{2}{1}(x,y)).
\]
بيا~\LR{$h$} له~\LR{$\Proj{2}{0}$}، \LR{$l$} او~\LR{$\Proj{2}{1}$} سره د~\LR{$f$} ترکيب
دے، چې
\[
l(x,y) = g(\Proj{2}{0}(x,y),\Proj{2}{0}(x,y),\Proj{2}{1}(x,y)),
\]
يعنې~\LR{$l$} له~\LR{$\Proj{2}{0}$}، \LR{$\Proj{2}{0}$} او~\LR{$\Proj{2}{1}$} سره د
\LR{$g$} ترکيب دے۔
\olpseganchor{OLP-0212-B010}{end}

\olpseganchor{OLP-0213-B004}{start}
\olfileid{cmp}{rec}{prf}
\olsection{بنسټيزې بازګشتي تابعې}
\olpseganchor{OLP-0213-B004}{end}

\olpseganchor{OLP-0213-B005}{start}
راځئ يو ځل بيا وليکو چې له شته تابعو څخه د بنسټيز بازګښت او ترکيب
په کارولو نوې تابعې څنګه تعريفولے شو۔
\olpseganchor{OLP-0213-B005}{end}

\olpseganchor{OLP-0213-B006}{start}
\begin{defn}
  \ollabel{defn:primitive-recursion} فرض کړئ~\LR{$f$} يوه~\LR{$k$}-ځاييزه
  تابع (\LR{$k\ge 1$}) او~\LR{$g$} يوه~\LR{$(k+2)$}-ځاييزه تابع ده۔ هغه تابع چې
  \emph{له~\LR{$f$} او~\LR{$g$} څخه په بنسټيز بازګښت تعريفېږي}،
  \LR{$(k+1)$}-ځاييزه تابع~\LR{$h$} ده چې لاندې معادلې ئې تعريفوي:
  \begin{align*}
    h(x_0,\dots,x_{k-1},0) & =  f(x_0,\dots,x_{k-1}) \\
    h(x_0,\dots,x_{k-1},y+1) & = g(x_0,\dots,x_{k-1}, y, h(x_0,\dots,x_{k-1}, y))
  \end{align*}
\end{defn}
\olpseganchor{OLP-0213-B006}{end}

\olpseganchor{OLP-0213-B007}{start}
\begin{defn}
  \ollabel{defn:composition} فرض کړئ~\LR{$f$} يوه~\LR{$k$}-ځاييزه تابع ده، او
  \LR{$g_0$}، \dots،~\LR{$g_{k-1}$} هغه~\LR{$k$} تابعې دي چې ټولې~\LR{$n$}-ځاييزې
  دي۔ هغه تابع چې \emph{له~\LR{$f$} او~\LR{$g_0$}، \dots،~\LR{$g_{k-1}$} څخه په
    ترکيب تعريفېږي}، \LR{$n$}-ځاييزه تابع~\LR{$h$} ده چې داسې تعريفېږي:
  \[
  h(x_0, \dots, x_{n-1}) =
  f(g_0(x_0, \dots, x_{n-1}), \dots, g_{k-1}(x_0, \dots, x_{n-1})).
  \]
\end{defn}
\olpseganchor{OLP-0213-B007}{end}

\olpseganchor{OLP-0213-B008}{start}
له~\LR{$\Succ$} او پروجکشن تابعو
\[
\Proj{n}{i}(x_0,\dots,x_{n-1}) = x_i,
\]
سره سره، د هر طبيعي عدد~\LR{$n$} او~\LR{$i < n$} دپاره تابع~\LR{$\Zero(x) = 0$}
هم په بنسټيزو بازګشتي تابعو کښې شاملوو۔
\olpseganchor{OLP-0213-B008}{end}

\olpseganchor{OLP-0213-B009}{start}
\begin{defn}
  د بنسټيزو بازګشتي تابعو سټ له~\LR{$\Nat^n$} څخه~\LR{$\Nat$} ته د تابعو هغه
  سټ دے چې په لاندې مادو استقرايي ډول تعريفېږي:
  \begin{enumerate}
  \item \LR{$\Zero$} بنسټيزه بازګشتي ده۔
  \item \LR{$\Succ$} بنسټيزه بازګشتي ده۔
  \item هره پروجکشن تابع~\LR{$\Proj{n}{i}$} بنسټيزه بازګشتي ده۔
  \item که~\LR{$f$} يوه~\LR{$k$}-ځاييزه بنسټيزه بازګشتي تابع وي او~\LR{$g_0$}،
    \dots،~\LR{$g_{k-1}$}، \LR{$n$}-ځاييزې بنسټيزې بازګشتي تابعې وي، نو له
    \LR{$g_0$}، \dots،~\LR{$g_{k-1}$} سره د~\LR{$f$} ترکيب بنسټيزه بازګشتي دے۔
  \item که~\LR{$f$} يوه~\LR{$k$}-ځاييزه بنسټيزه بازګشتي تابع او~\LR{$g$} يوه
    \LR{$k+2$}-ځاييزه بنسټيزه بازګشتي تابع وي، نو له~\LR{$f$} او~\LR{$g$} څخه په
    بنسټيز بازګښت تعريف شوې تابع بنسټيزه بازګشتي ده۔
\end{enumerate}
\end{defn}
\olpseganchor{OLP-0213-B009}{end}

\olpseganchor{OLP-0213-B010}{start}
\begin{explain}
په لنډه وينا، د بنسټيزو بازګشتي تابعو سټ هغه تر ټولو وړوکے سټ دے
چې~\LR{$\Zero$}، \LR{$\Succ$} او پروجکشن تابعې~\LR{$\Proj{n}{j}$} لري، او د ترکيب
او بنسټيز بازګښت له مخې بند دے۔
\olpseganchor{OLP-0213-B010}{end}

\olpseganchor{OLP-0213-B011}{start}
د بنسټيزو بازګشتي تابعو سټ د «پړاوونو» له مخې هم بيانولے شو۔
\LR{$S_0$} دې د پيلوونکو تابعو، يعنې~\LR{$\Zero$}، \LR{$\Succ$} او پروجکشنونو، سټ
وښيي۔ دا د~\LR{$0$} پړاو بنسټيزې بازګشتي تابعې دي۔ کله چې يو
پړاو~\LR{$S_i$} تعريف شو، \LR{$S_{i+1}$} دې د پخپله~\LR{$S_i$} د تابعو او هغو
ټولو تابعو سټ وي چې د همدې پړاو پر تابعو د ترکيب يا بنسټيز بازګښت
يوه بېلګه په کارولو لاس ته راځي۔ بيا
\[
S = \bigcup_{i \in \Nat} S_i
\]
د ټولو بنسټيزو بازګشتي تابعو سټ دے۔
\textbf{د ژباړې د نسخې سمون (OLCMP-005):} سرچينې راتلونکی پړاو
يوازې د مخکيني پړاو پر تابعو د يو عمل له پايلو جوړاوه او پخپله ئې
مخکينۍ تابعې نۀ ساتلې؛ دلته پړاوونه په څرګنده تزايدي کړل شول، څو
اتحاد د تعريف شوو تابعو له مخې بند وي۔
\end{explain}
\olpseganchor{OLP-0213-B011}{end}

\olpseganchor{OLP-0213-B012}{start}
راځئ تصديق کړو چې~\LR{$\Add$} بنسټيزه بازګشتي تابع ده۔
\olpseganchor{OLP-0213-B012}{end}

\olpseganchor{OLP-0213-B013}{start}
\begin{prop}
  د جمع تابع~\LR{$\Add(x,y) = x+y$} بنسټيزه بازګشتي ده۔
\end{prop}
\olpseganchor{OLP-0213-B013}{end}

\olpseganchor{OLP-0213-B014}{start}
\begin{proof}
د~\LR{$\Add$} يو بنسټيز بازګشتي تعريف مو لا له دوو تابعو~\LR{$f$} او~\LR{$g$} څخه
لرو چې د~\olref{defn:primitive-recursion} له شکل سره سمون خوري:
\begin{align*}
  \Add(x_0, 0) & = f(x_0) = x_0\\
  \Add(x_0, y+1) & = g(x_0, y, \Add(x_0, y)) =
  \Succ(\Add(x_0, y))
\end{align*}
نو~\LR{$\Add$} په دې شرط بنسټيزه بازګشتي ده چې~\LR{$f$} او~\LR{$g$} هم داسې وي۔
\LR{$f(x_0) = x_0 = \Proj{1}{0}(x_0)$} ده، او پروجکشن تابعې بنسټيزې
بازګشتي شمېرل کېږي، نو~\LR{$f$} بنسټيزه بازګشتي ده۔ تابع~\LR{$g$} هغه
درې‌ځاييزه تابع~\LR{$g(x_0, y, z)$} ده چې داسې تعريفېږي:
\[
g(x_0, y, z) = \Succ(z).
\]
دا لا نۀ ښيي چې~\LR{$g$} بنسټيزه بازګشتي ده، ځکه~\LR{$g$} او~\LR{$\Succ$} په کره
توګه يوه تابع نۀ دي: \LR{$\Succ$} يوځاييزه ده، خو~\LR{$g$} بايد درې‌ځاييزه
وي۔ مګر~\LR{$g$} په ترکيب په «رسمي» ډول داسې تعريفولے شو:
\[
g(x_0, y, z) = \Succ(\Proj{3}{2}(x_0, y, z))
\]
څرنګه چې~\LR{$\Succ$} او~\LR{$\Proj{3}{2}$} بنسټيزې بازګشتي تابعې شمېرل
کېږي، نو~\LR{$g$} هم داسې ده، ځکه چې له بنسټيزو بازګشتي تابعو څخه په
ترکيب تعريفېدے شي۔
\end{proof}
\olpseganchor{OLP-0213-B014}{end}

\olpseganchor{OLP-0213-B015}{start}
\begin{prop}
  \ollabel{prop:mult-pr}
  د ضرب تابع~\LR{$\Mult(x,y) = x \cdot y$} بنسټيزه بازګشتي ده۔
\end{prop}
\olpseganchor{OLP-0213-B015}{end}

\olpseganchor{OLP-0213-B016}{start}
\begin{proof}
  تمرين۔
\end{proof}
\olpseganchor{OLP-0213-B016}{end}

\olpseganchor{OLP-0213-B017}{start}
\begin{prob}
  \olref[cmp][rec][prf]{prop:mult-pr} په دې ښودلو ثابته کړئ چې د
  \LR{$\Mult$} بنسټيز بازګشتي تعريف د
  \olref[cmp][rec][prf]{defn:primitive-recursion} غوښتل شوي شکل ته
  راوستل کېدے شي، او اړوندې تابعې~\LR{$f$} او~\LR{$g$} بنسټيزې بازګشتي دي۔
\end{prob}
\olpseganchor{OLP-0213-B017}{end}

\olpseganchor{OLP-0213-B018}{start}
\begin{ex}
دا مو د بنسټيز بازګشتي تعريف تر ټولو لومړۍ بېلګه ده:
\begin{align*}
h(0) & =  1 \\
h(y+1) & =  2 \cdot h(y).
\end{align*}
دا تابع د~\olref{defn:primitive-recursion} غوښتل شوي شکل ته نۀ شي
راتلے، ځکه~\LR{$k=0$} دے۔ په تعريف کښې ثابتونه~\LR{$1$} او~\LR{$2$} هم راځي۔ د
لومړۍ ستونزې د هوارولو دپاره يو بې‌کاره آرګومېنټ ورزياتوو او
تابع~\LR{$h'$} تعريفوو:
\begin{align*}
h'(x_0, 0) & =  f(x_0) = 1 \\
h'(x_0, y+1) & =  g(x_0, y, h'(x_0, y)) = 2 \cdot h'(x_0, y).
\end{align*}
تابع~\LR{$f(x_0) = 1$} له~\LR{$\Succ$} او~\LR{$\Zero$} څخه په ترکيب تعريفېدے شي:
\LR{$f(x_0) = \Succ(\Zero(x_0))$}۔ تابع~\LR{$g$} له~\LR{$g'(z) = 2 \cdot z$} او
پروجکشنونو څخه په ترکيب تعريفېدے شي:
\begin{align*}
g(x_0, y, z) & = g'(\Proj{3}{2}(x_0, y, z))
\intertext{\RL{او~\LR{$g'$} بيا پخپله په ترکيب داسې تعريفېدے شي:}}
g'(z) & = \Mult(g''(z), \Proj{1}{0}(z))
\intertext{\RL{او}}
g''(z) & = \Succ(f(z)),
\end{align*}
چې~\LR{$f$} هماغه پورته تابع ده: \LR{$f(z) = \Succ(\Zero(z))$}۔ اوس
چې~\LR{$h'$} لرو، بيا ترکيب کارولے شو او~\LR{$h(y) =
h'(\Proj{1}{0}(y),\Proj{1}{0}(y))$} ټاکلے شو۔ دا ښيي چې~\LR{$h$} له
بنسټيزو تابعو څخه د ترکيبونو او بنسټيزو بازګښتونو د يوې لړۍ په
کارولو تعريفېدے شي، نو~\LR{$h$} بنسټيزه بازګشتي ده۔
\end{ex}
\olpseganchor{OLP-0213-B018}{end}

\olpseganchor{OLP-0214-B004}{start}
\olfileid{cmp}{rec}{not}
\olsection{د بنسټيز بازګښت نښې}
\olpseganchor{OLP-0214-B004}{end}

\olpseganchor{OLP-0214-B005}{start}
د بنسټيزو بازګشتي تابعو د کره استقرايي بيان يوه ګټه دا ده چې دا
تابعې په منظم ډول بيانولے شو۔ د بېلګې په توګه، هرې داسې تابع ته
لاندې يوه «نښه» ټاکلے شو۔ د صفر، راتلونکي او پروجکشنونو دپاره په
ترتيب سره نښې~\LR{$\Zero$}، \LR{$\Succ$} او~\LR{$\Proj{n}{i}$} وکاروئ۔ اوس فرض
کړئ~\LR{$h$} له يوې~\LR{$k$}-ځاييزې تابع~\LR{$f$} او~\LR{$n$}-ځاييزو تابعو~\LR{$g_0$}،
\dots،~\LR{$g_{k-1}$} څخه په ترکيب تعريف شوې ده، او وروستيو تابعو ته
مو نښې~\LR{$F$}، \LR{$G_0$}، \dots،~\LR{$G_{k-1}$} ټاکلې دي۔ بيا د يوې نوې
نښې~\LR{$\fn{Comp}_{k,n}$} په کارولو سره، تابع~\LR{$h$} په
\LR{$\fn{Comp}_{k,n}[F,G_0,\dots,G_{k-1}]$} ښودلے شو۔
\olpseganchor{OLP-0214-B005}{end}

\olpseganchor{OLP-0214-B006}{start}
په بنسټيز بازګښت د تعريف شوو تابعو دپاره ورته نښې کارولے شو۔ فرض
کړئ~\LR{$(k+1)$}-ځاييزه تابع~\LR{$h$} له~\LR{$k$}-ځاييزې تابع~\LR{$f$} او
\LR{$(k+2)$}-ځاييزې تابع~\LR{$g$} څخه په بنسټيز بازګښت تعريف شوې ده، او~\LR{$f$}
او~\LR{$g$} ته ټاکل شوې نښې په ترتيب سره~\LR{$F$} او~\LR{$G$} دي۔ بيا~\LR{$h$} ته
ټاکل شوې نښه~\LR{$\fn{Rec}_k[F,G]$} ده۔
\olpseganchor{OLP-0214-B006}{end}

\olpseganchor{OLP-0214-B007}{start}
په ياد ولرئ چې د جمع تابع په بنسټيز بازګښت داسې تعريف شوې ده:
\begin{align*}
  \Add(x_0, 0) & = \Proj{1}{0}(x_0) = x_0\\
  \Add(x_0, y+1) & = \Succ(\Proj{3}{2}(x_0, y, \Add(x_0, y))) = \Add(x_0, y) +1
\end{align*}
دلته د~\LR{$f$} ونډه~\LR{$\Proj{1}{0}$} ترسره کوي، او د~\LR{$g$} ونډه
\LR{$\Succ(\Proj{3}{2}(x_0, y, z))$} ترسره کوي۔ وروستۍ تابع ته
\LR{$\fn{Comp}_{1,3}[\Succ,\Proj{3}{2}]$} نښه ټاکل شوې، ځکه دا له
يوې~\LR{$1$}-ځاييزې تابع~\LR{$\Succ$} او~\LR{$3$}-ځاييزې تابع~\LR{$\Proj{3}{2}$} څخه په
ترکيب د يوې تابع د تعريف پايله ده۔ په دې ترتيب، د جمع تابع په لاندې
نښه ښودلے شو:
\[
\fn{Rec}_1[\Proj{1}{0},\fn{Comp}_{1,3}[\Succ,\Proj{3}{2}]].
\]
د دې نښو لرل کله ناکله ګټور وي؛ د بېلګې په توګه، د بنسټيزو
بازګشتي تابعو د شمېرنې پر مهال۔
\olpseganchor{OLP-0214-B007}{end}

\olpseganchor{OLP-0214-B008}{start}
\begin{prob}
  د~\LR{$\Mult$} بشپړه بنسټيزه بازګشتي نښه ورکړئ۔
\end{prob}
\olpseganchor{OLP-0214-B008}{end}

\olpseganchor{OLP-0215-B004}{start}
\olfileid{cmp}{rec}{cmp}
\olsection{بنسټيزې بازګشتي تابعې محاسبه کېدونکې دي}
\olpseganchor{OLP-0215-B004}{end}

\olpseganchor{OLP-0215-B005}{start}
فرض کړئ يوه تابع~\LR{$h$} په بنسټيز بازګښت تعريف شوې ده:
\begin{eqnarray*}
h(\vec x, 0)     & = & f(\vec x) \\
h(\vec x, y + 1) & = & g(\vec x, y, h(\vec x, y))
\end{eqnarray*}
او فرض کړئ تابعې~\LR{$f$} او~\LR{$g$} محاسبه کېدونکې دي۔ (موږ~\LR{$\vec x$} د
\LR{$x_0$}، \dots،~\LR{$x_{k-1}$} د لنډون دپاره کاروو۔) بيا
\LR{$h(\vec x, 0)$} په څرګند ډول محاسبه کېدے شي، ځکه دا يوازې
\LR{$f(\vec x)$} ده چې محاسبه کېدونکې مو فرض کړې ده۔ بيا
\LR{$h(\vec x, 1)$} هم محاسبه کېدے شي، ځکه~\LR{$1 = 0 + 1$} دے او په دې ډول
\LR{$h(\vec x, 1)$} يوازې دا ده:
\begin{align*}
 h(\vec x, 1) & = g(\vec x, 0, h(\vec x, 0)) =  g(\vec x, 0, f(\vec x)).
\intertext{\RL{په همدې ډول مخکښې تللے شو او دا محاسبه کولے شو:}}
h(\vec x, 2) & = g(\vec x, 1, h(\vec x, 1)) = g(\vec x, 1, g(\vec x, 0, f(\vec x)))\\
h(\vec x, 3) & = g(\vec x, 2, h(\vec x, 2)) = g(\vec x, 2, g(\vec x, 1, g(\vec x, 0, f(\vec x))))\\
h(\vec x, 4) & = g(\vec x, 3, h(\vec x, 3)) = g(\vec x, 3, g(\vec x, 2, g(\vec x, 1, g(\vec x, 0, f(\vec x)))))\\
& \vdots
\end{align*}
نو په عمومي ډول د~\LR{$h(\vec x, y)$} د محاسبې دپاره په پرله‌پسې توګه
\LR{$h(\vec x, 0)$}، \LR{$h(\vec x, 1)$}، \dots محاسبه کوو، تر هغه چې
\LR{$h(\vec x, y)$} ته ورسېږو۔
\olpseganchor{OLP-0215-B005}{end}

\olpseganchor{OLP-0215-B006}{start}
نو که تابعې~\LR{$f$} او~\LR{$g$} محاسبه کېدونکې وي، يو بنسټيز بازګشتي تعريف
يوه نوې محاسبه کېدونکې تابع ورکوي۔ همدارنګه، که تابعې~\LR{$f$} او~\LR{$g_i$}
محاسبه کېدونکې وي، د تابعو ترکيب هم محاسبه کېدونکې تابع ورکوي۔
\olpseganchor{OLP-0215-B006}{end}

\olpseganchor{OLP-0215-B007}{start}
څرنګه چې بنسټيزې تابعې~\LR{$\Zero$}، \LR{$\Succ$} او~\LR{$\Proj{n}{i}$} محاسبه
کېدونکې دي، او ترکيب او بنسټيز بازګښت له محاسبه کېدونکو تابعو څخه
محاسبه کېدونکې تابعې ورکوي، نو هره بنسټيزه بازګشتي تابع محاسبه
کېدونکې ده۔
\olpseganchor{OLP-0215-B007}{end}

\olpseganchor{OLP-0216-B004}{start}
\olfileid{cmp}{rec}{exa}
\olsection{د بنسټيزو بازګشتي تابعو بېلګې}
\olpseganchor{OLP-0216-B004}{end}

\olpseganchor{OLP-0216-B005}{start}
موږ لا د بنسټيزو بازګشتي تابعو ځينې بېلګې لرو: د جمع او ضرب تابعې~\LR{$\Add$} او~\LR{$\Mult$}۔ د عينيت تابع~\LR{$\fn{id}(x) = x$} بنسټيزه بازګشتي ده، ځکه دا يوازې~\LR{$\Proj{1}{0}$} ده۔ ثابتې تابعې \LR{$\fn{const}_n(x) = n$} هم بنسټيزې بازګشتي دي، ځکه هغوئ له~\LR{$\Zero$} او~\LR{$\Succ$} څخه په پرله‌پسې ترکيب تعريفېدے شي۔ دا هغه وخت ګټور دي چې په بنسټيزو بازګشتي تعريفونو کښې ثابتونه کارول غواړو۔ مثلاً، که تابع~\LR{$f(x) = 2 \cdot x$} تعريفول غواړو، نو له~\LR{$\fn{const}_n(x)$} او ضرب څخه ئې په ترکيب داسې ترلاسه کولے شو:
\LR{$f(x) = \Mult(\fn{const}_2(x), \Proj{1}{0}(x))$}۔ له دې وروسته به دا چل کاروو۔
\olpseganchor{OLP-0216-B005}{end}

\olpseganchor{OLP-0216-B006}{start}
\begin{prop}
  د توان تابع~\LR{$\fn{exp}(x, y) = x^y$} بنسټيزه بازګشتي ده۔
\end{prop}
\olpseganchor{OLP-0216-B006}{end}

\olpseganchor{OLP-0216-B007}{start}
\begin{proof}
  \LR{$\fn{exp}$} په بنسټيز بازګښت داسې تعريفولے شو:
  \begin{align*}
    \fn{exp}(x, 0) & = 1\\
    \fn{exp}(x, y+1) & = \Mult(x, \fn{exp}(x,y)).
    \intertext{\RL{په کره توګه دا لا له بنسټيزو بازګشتي تابعو څخه يو بازګشتي تعريف نۀ دے۔ خو رسمي بڼه ئې داسې ده:}}
    \fn{exp}(x, 0) & = f(x)\\
    \fn{exp}(x, y+1) & = g(x, y, \fn{exp}(x,y)).
    \intertext{\RL{چې پکښې}}
    f(x) & = \Succ(\Zero(x)) = 1\\
    g(x, y, z) & = \Mult(\Proj{3}{0}(x, y, z), \Proj{3}{2}(x, y, z)) = x \cdot z
  \end{align*}
  او په دې ډول~\LR{$f$} او~\LR{$g$} له بنسټيزو بازګشتي تابعو څخه په ترکيب تعريف شوې دي۔
\end{proof}
\olpseganchor{OLP-0216-B007}{end}

\olpseganchor{OLP-0216-B008}{start}
\begin{prop}
  مخکينۍ تابع~\LR{$\fn{pred}(y)$} چې داسې تعريف شوې ده،
  \[
  \fn{pred}(y) = \begin{cases}
    0 & \text{\RL{که \LR{$y=0$} وي}}\\
    y-1 & \text{\RL{که نۀ}}
  \end{cases}
  \]
  بنسټيزه بازګشتي ده۔
\end{prop}
\olpseganchor{OLP-0216-B008}{end}

\olpseganchor{OLP-0216-B009}{start}
\begin{proof}
 پام وکړئ چې
 \begin{align*}
   \fn{pred}(0) & = 0 \text{\RL{ او}}\\
   \fn{pred}(y+1) & = y.
 \end{align*}
 دا نژدې يو بنسټيز بازګشتي تعريف دے۔ په کره توګه د بنسټيز بازګښت د تعريف له نمونې سره ځکه نۀ سمېږي چې هغه نمونه لږ تر لږه يو زيات آرګومېنټ~\LR{$x$} غواړي۔ دا په دې لحاظ هم نابلده ده چې د \LR{$\fn{pred}(y+1)$} په تعريف کښې په رښتيا~\LR{$\fn{pred}(y)$} نۀ کاروي۔ خو لومړے~\LR{$\fn{pred}'(x, y)$} داسې تعريفولے شو:
 \begin{align*}
   \fn{pred}'(x, 0) & = \Zero(x) = 0,\\
   \fn{pred}'(x, y+1) & = \Proj{3}{1}(x, y, \fn{pred'}(x, y)) = y.
 \end{align*}
او بيا~\LR{$\fn{pred}$} ترې په ترکيب تعريفولے شو، مثلاً داسې:
\LR{$\fn{pred}(x) = \fn{pred}'(\Zero(x), \Proj{1}{0}(x))$}۔
\end{proof}
\olpseganchor{OLP-0216-B009}{end}

\olpseganchor{OLP-0216-B010}{start}
\begin{prop}
  د فاکتوريل تابع~\LR{$\fn{fac}(x) = \fact{x} = 1 \cdot 2 \cdot 3 \cdot \dots \cdot x$} بنسټيزه بازګشتي ده۔
\end{prop}
\olpseganchor{OLP-0216-B010}{end}

\olpseganchor{OLP-0216-B011}{start}
\begin{proof}
  څرګند بنسټيز بازګشتي تعريف دا دے:
  \begin{align*}
    \fn{fac}(0) &= 1\\
    \fn{fac}(y+1) & = \fn{fac}(y) \cdot (y+1).
    \intertext{\RL{په رسمي ډول لومړے دوه ځاييزه تابع~\LR{$h$} داسې تعريفوو:}}
    h(x, 0) & = \fn{const}_1(x)\\
    h(x, y+1) & = g(x, y, h(x, y))
    \intertext{\RL{چې پکښې \LR{$g(x, y, z) = \Mult(\Proj{3}{2}(x, y, z), \Succ(\Proj{3}{1}(x, y, z)))$} ده؛ بيا ږدو چې}}
    \fn{fac}(y) & = h(\Proj{1}{0}(y), \Proj{1}{0}(y)) = h(y,y).
  \end{align*}
  له دې وروسته به لږ کم رسمي واوسو او هر ځل به د ترکيب او بنسټيز بازګښت بشپړ رسمي تعريفونه نۀ ورکوو۔
\end{proof}
\olpseganchor{OLP-0216-B011}{end}

\olpseganchor{OLP-0216-B012}{start}
\begin{prop}
  پرې شوې تفريق~\LR{$x \tsub y$}، چې داسې تعريف شوې ده،
  \[
  x \tsub y = \begin{cases}
    0 & \text{\RL{که \LR{$x < y$} وي}}\\
    x-y & \text{\RL{که نۀ}}
  \end{cases}
  \]
  بنسټيزه بازګشتي ده۔
\end{prop}
\olpseganchor{OLP-0216-B012}{end}

\olpseganchor{OLP-0216-B013}{start}
\begin{proof}
  لرو چې:
  \begin{align*}
    x \tsub 0 & = x\\
    x \tsub (y+1) & = \fn{pred}(x \tsub y)
  \end{align*}
\end{proof}
\olpseganchor{OLP-0216-B013}{end}

\olpseganchor{OLP-0216-B014}{start}
\begin{prop}
 د~\LR{$x$} او~\LR{$y$} تر منځ واټن~\LR{$\left|x-y\right|$} بنسټيز بازګشتي دے۔
\end{prop}
\olpseganchor{OLP-0216-B014}{end}

\olpseganchor{OLP-0216-B015}{start}
\begin{proof}
  لرو چې~\LR{$\left| x-y \right| = (x \tsub y) + (y \tsub x)$}؛ نو واټن له~\LR{$+$} او~\LR{$\tsub$} څخه په ترکيب تعريفېدے شي، او دواړه بنسټيزې بازګشتي دي۔
\end{proof}
\olpseganchor{OLP-0216-B015}{end}

\olpseganchor{OLP-0216-B016}{start}
\begin{prop}
  د~\LR{$x$} او~\LR{$y$} تر ټولو لوے ارزښت~\LR{$\fn{max}(x,y)$} بنسټيز بازګشتي دے۔
\end{prop}
\olpseganchor{OLP-0216-B016}{end}

\olpseganchor{OLP-0216-B017}{start}
\begin{proof}
  \LR{$\fn{max}(x,y)$} له~\LR{$+$} او~\LR{$\tsub$} څخه په ترکيب داسې تعريفولے شو:
  \[
  \fn{max}(x,y) \defis x + (y \tsub x).
  \]
  که~\LR{$x$} تر ټولو لوے وي، يعنې~\LR{$x \ge y$}، نو~\LR{$y \tsub x = 0$}؛ ځکه \LR{$x + (y \tsub x) = x + 0 = x$}۔ که~\LR{$y$} تر ټولو لوے وي، نو \LR{$y \tsub x = y - x$}، او ځکه~\LR{$x + (y \tsub x) = x + (y - x) = y$}۔
\end{proof}
\olpseganchor{OLP-0216-B017}{end}

\olpseganchor{OLP-0216-B018}{start}
\begin{prop}
  \ollabel{prop:min-pr}
  د~\LR{$x$} او~\LR{$y$} تر ټولو وړوکے ارزښت~\LR{$\fn{min}(x,y)$} بنسټيز بازګشتي دے۔
\end{prop}
\olpseganchor{OLP-0216-B018}{end}

\olpseganchor{OLP-0216-B019}{start}
\begin{proof}
  تمرين۔
\end{proof}
\olpseganchor{OLP-0216-B019}{end}

\olpseganchor{OLP-0216-B020}{start}
\begin{prob}
  \olref[cmp][rec][exa]{prop:min-pr} ثابت کړئ۔
\end{prob}
\olpseganchor{OLP-0216-B020}{end}

\olpseganchor{OLP-0216-B021}{start}
\begin{prob}
وښيئ چې
\[
f(x, y) =
2^{(2^{\iddots^{2^{x}}})}\raisebox{1ex}{\bigg\rbrace}
\raisebox{1ex}{\text {\LR{$y$} ځله \LR{$2$}}}
\]
بنسټيزه بازګشتي ده۔
\end{prob}
\olpseganchor{OLP-0216-B021}{end}

\olpseganchor{OLP-0216-B022}{start}
\begin{prob}
وښيئ چې صحيح وېش~\LR{$d(x, y) = \lfloor x/y \rfloor$} (يعنې هغه وېش چې د اعشاري ټکي وروسته هر څه پکښې غورځوئ) بنسټيز بازګشتي دے۔ کله چې~\LR{$y = 0$} وي، شرط ږدو چې~\LR{$d(x, y) = 0$} وي۔ د بنسټيز بازګښت او ترکيب په کارولو د~\LR{$d$} يو څرګند تعريف ورکړئ۔
\end{prob}
\olpseganchor{OLP-0216-B022}{end}

\olpseganchor{OLP-0216-B023}{start}
\begin{prop}
د بنسټيزو بازګشتي تابعو سټ د لاندې دوو عملونو له مخې بند دے:
\begin{enumerate}
\item متناهي جمعې: که~\LR{$f(\vec x, z)$} بنسټيزه بازګشتي وي، نو لاندې تابع هم داسې ده:
\[
g(\vec x, y) \defis \sum_{z = 0}^y f(\vec x, z).
\]
\item متناهي حاصل‌ضربونه: که~\LR{$f(\vec x, z)$} بنسټيزه بازګشتي وي، نو لاندې تابع هم داسې ده:
\[
h(\vec x, y) \defis \prod_{z = 0}^y f(\vec x, z).
\]
\end{enumerate}
\end{prop}
\olpseganchor{OLP-0216-B023}{end}

\olpseganchor{OLP-0216-B024}{start}
\begin{proof}
د بېلګې په توګه، متناهي جمعې په لاندې معادلو بازګشتي تعريفېږي:
\begin{align*}
  g(\vec x, 0) & = f(\vec x, 0)\\
  g(\vec x, y+1) & = g(\vec x, y) + f(\vec x, y+1).
\end{align*}
\end{proof}
\olpseganchor{OLP-0216-B024}{end}

\olpseganchor{OLP-0217-B004}{start}
\olfileid{cmp}{rec}{prr}
\olsection{بنسټيزې بازګشتي اړيکې}
\olpseganchor{OLP-0217-B004}{end}

\olpseganchor{OLP-0217-B005}{start}
\begin{defn}
اړيکه~\LR{$R(\vec x)$} هله بنسټيزه بازګشتي بلل کېږي چې ځانګړونکې تابع ئې
\[
\Char{R}(\vec x) = \left\{
  \begin{array}{ll}
  1 & \mbox{که \LR{$R(\vec x)$} وي} \\
  0 & \mbox{که نۀ}
  \end{array}
\right.
\]
بنسټيزه بازګشتي وي۔
\end{defn}
\olpseganchor{OLP-0217-B005}{end}

\olpseganchor{OLP-0217-B006}{start}
په بله وينا، کله چې د بنسټيزې بازګشتي اړيکې~\LR{$R(\vec x)$} خبره کوو، مطلب مو د~\LR{$\Char{R}(\vec x) = 1$} په بڼه اړيکه وي، چې~\LR{$\Char{R}$} يوه بنسټيزه بازګشتي تابع ده او پر هر ننوت 1 يا 0 راګرځوي۔ د بېلګې په توګه، اړيکه~\LR{$\fn{IsZero}(x)$} هله او يوازې هله صدق کوي چې \LR{$x = 0$} وي، او له هغې سره سمون لرونکې تابع~\LR{$\Char{\fn{IsZero}}$} په بنسټيز بازګښت داسې تعريفېږي:
\begin{align*}
\Char{\fn{IsZero}}(0) & = 1,\\
\Char{\fn{IsZero}}(x+1) & = 0.
\end{align*}
\olpseganchor{OLP-0217-B006}{end}

\olpseganchor{OLP-0217-B007}{start}
بايد څرګنده وي چې اړيکې له نورو بنسټيزو بازګشتي تابعو سره ترکيبولے شو۔ نو لاندې اړيکې هم بنسټيزې بازګشتي دي:
\begin{enumerate}
\item د برابرۍ اړيکه~\LR{$x = y$}، چې په~\LR{$\fn{IsZero}(\left|x-y\right|)$} تعريفېږي؛
\item د کوچني-يا-برابر اړيکه~\LR{$x \leq y$}، چې په~\LR{$\fn{IsZero}(x \tsub y)$} تعريفېږي۔
\end{enumerate}
\textbf{د ژباړې د نسخې سمون (OLCMP-006):} سرچينې دويمه اړيکه «کوچني» بللې وه، خو ورپسې نښه او ځانګړونکې تابع د کوچني-يا-برابر رسمي شرط ټاکي۔ دلته نوم له رسمي شرط سره برابر شوے دے۔
\olpseganchor{OLP-0217-B007}{end}

\olpseganchor{OLP-0217-B008}{start}
\begin{prop}
  د بنسټيزو بازګشتي اړيکو سټ د بولي عملونو له مخې بند دے؛ يعنې که \LR{$P(\vec x)$} او~\LR{$Q(\vec x)$} بنسټيزې بازګشتي وي، نو لاندې اړيکې هم داسې دي:
  \begin{enumerate}
  \item \LR{$\lnot P(\vec x)$}
  \item \LR{$P(\vec x) \land Q(\vec x)$}
  \item \LR{$P(\vec x) \lor Q(\vec x)$}
  \item \LR{$P(\vec x) \lif Q(\vec x)$}
  \end{enumerate}
\end{prop}
\olpseganchor{OLP-0217-B008}{end}

\olpseganchor{OLP-0217-B009}{start}
\begin{proof}
  فرض کړئ~\LR{$P(\vec x)$} او~\LR{$Q(\vec x)$} بنسټيزې بازګشتي دي، يعنې ځانګړونکې تابعې~\LR{$\Char{P}$} او~\LR{$\Char{Q}$} ئې بنسټيزې بازګشتي دي۔ بايد وښيو چې د~\LR{$\lnot P(\vec x)$} او نورو ځانګړونکې تابعې هم بنسټيزې بازګشتي دي۔
  \[
  \Char{\lnot P}(\vec x) = \begin{cases}
    0 & \text{\RL{که \LR{$\Char{P}(\vec x) = 1$} وي}}\\
    1 & \text{\RL{که نۀ}}
  \end{cases}
  \]
  \LR{$\Char{\lnot P}(\vec x)$} د~\LR{$1 \tsub \Char{P}(\vec x)$} په توګه تعريفولے شو۔
  \[
  \Char{P \land Q}(\vec x) = \begin{cases}
    1 & \text{\RL{که \LR{$\Char{P}(\vec x) = \Char{Q}(\vec x) = 1$} وي}}\\
    0 & \text{\RL{که نۀ}}
  \end{cases}
  \]
  \LR{$\Char{P \land Q}(\vec x)$} د~\LR{$\Char{P}(\vec x) \cdot \Char{Q}(\vec x)$}، يا د~\LR{$\fn{min}(\Char{P}(\vec x), \Char{Q}(\vec x))$} په توګه تعريفولے شو۔ په همدې ډول،
  \begin{align*}
    \Char{P \lor Q}(\vec x) & = \fn{max}(\Char{P}(\vec x), \Char{Q}(\vec x)) \text{\RL{ او}}\\
    \Char{P \lif Q}(\vec x) & = \fn{max}(1 \tsub \Char{P}(\vec x), \Char{Q}(\vec x)).
  \end{align*}
\end{proof}
\olpseganchor{OLP-0217-B009}{end}

\olpseganchor{OLP-0217-B010}{start}
\begin{prop}
  د بنسټيزو بازګشتي اړيکو سټ د محدودې کميت ټاکنې له مخې بند دے؛ يعنې که~\LR{$R(\vec x, z)$} بنسټيزه بازګشتي اړيکه وي، نو لاندې اړيکې هم داسې دي:
  \begin{align*}
    & \bforall{z < y}{R(\vec x, z)} \text{\RL{ او}}\\
    & \bexists{z < y}{R(\vec x, z)}.
  \end{align*}
  \LR{$\bforall{z < y}{R(\vec x, z)}$} د~\LR{$\vec x$} او~\LR{$y$} په باب هله او يوازې هله صدق کوي چې~\LR{$R(\vec x, z)$} له~\LR{$y$} څخه د هر کوچني~\LR{$z$} دپاره صدق وکړي؛ د~\LR{$\bexists{z < y}{R(\vec x, z)}$} دپاره هم ورته تشريح ده۔
\end{prop}
\olpseganchor{OLP-0217-B010}{end}

\olpseganchor{OLP-0217-B011}{start}
\begin{proof}
  د دود له مخې~\LR{$\bforall{z < 0}{R(\vec x, z)}$} رښتونے نيسو (په ساده دليل چې له~\LR{$0$} څخه کوچنے~\LR{$z$} نشته)، او \LR{$\bexists{z < 0}{R(\vec x, z)}$} دروغ نيسو۔ محدود کلي کميت ټاکونکی کټ مټ د متناهي حاصل‌ضرب يا تکراري کمينې په شان کار کوي؛ يعنې که \LR{$P(\vec x, y) \defiff \bforall{z < y}{R(\vec x, z)}$} وي، نو \LR{$\Char{P}(\vec x, y)$} داسې تعريفېدے شي:
  \begin{align*}
    \Char{P}(\vec x, 0) & = 1\\
    \Char{P}(\vec x, y+1) & = \fn{min}(\Char{P}(\vec x, y), \Char{R}(\vec x, y)).
  \end{align*}
  محدوده وجودي کميت ټاکنه هم په ورته ډول د~\LR{$\fn{max}$} په کارولو تعريفېدے شي۔ يا ئې له محدودې کلي کميت ټاکنې څخه د دې برابرۍ په کارولو تعريفولے شو:
  \LR{$\bexists{z < y}{R(\vec x, z)} \liff \lnot \bforall{z < y}{\lnot R(\vec x, z)}$}۔ پام وکړئ چې مثلاً د~\LR{$\bexists{x \leq y}{\dots x\dots}$} په بڼه محدود کميت ټاکونکی له~\LR{$\bexists{x < y+1}{\dots x \dots}$} سره برابر دے۔
\end{proof}
\olpseganchor{OLP-0217-B011}{end}

\olpseganchor{OLP-0217-B012}{start}
\begin{prob}
  وښيئ چې درې ځاييزه اړيکه~\LR{$x \equiv y \mod n$} (مودولو~\LR{$n$} تطابق) بنسټيزه بازګشتي ده۔
\end{prob}
\olpseganchor{OLP-0217-B012}{end}

\olpseganchor{OLP-0217-B013}{start}
يوه بله ګټوره بنسټيزه بازګشتي تابع شرطي تابع~\LR{$\fn{cond}(x,y,z)$} ده، چې داسې تعريفېږي:
\begin{align*}
  \fn{cond}(x,y,z) & = \begin{cases}
  y & \text{\RL{که \LR{$x = 0$} وي}} \\
  z & \text{\RL{که نۀ}}.
\end{cases}
\intertext{\RL{دا په بازګشتي ډول داسې تعريفېږي:}}
\fn{cond}(0,y,z) & = y,\\
\fn{cond}(x+1,y,z) & = z.
\end{align*}
له دې څخه د بنسټيزو بازګشتي اړيکو پر بنسټ په حالتونو د بنسټيزو بازګشتي تابعو تعريفونه توجيه کولے شو:
\olpseganchor{OLP-0217-B013}{end}

\olpseganchor{OLP-0217-B014}{start}
\begin{prop}
که~\LR{$g_0(\vec x)$}، \dots،~\LR{$g_m(\vec x)$} بنسټيزې بازګشتي تابعې وي، او~\LR{$R_0(\vec x)$}، \dots،~\LR{$R_{m-1}(\vec x)$} بنسټيزې بازګشتي اړيکې وي، نو لاندې تعريف شوې تابع~\LR{$f$} هم بنسټيزه بازګشتي ده:
\[
f(\vec x) = \begin{cases}
    g_0(\vec x) & \text{\RL{که \LR{$R_0(\vec{x})$} وي}} \\
    g_1(\vec x) & \text{\RL{که \LR{$R_1(\vec{x})$} وي او \LR{$R_0(\vec{x})$} نۀ وي}} \\
    \vdots & \\
    g_{m-1}(\vec x) & \text{\RL{که \LR{$R_{m-1}(\vec{x})$} وي او مخکيني هېڅ يو نۀ وي}} \\
    g_m(\vec x) & \mbox{که نۀ}
\end{cases}
\]
\end{prop}
\olpseganchor{OLP-0217-B014}{end}

\olpseganchor{OLP-0217-B015}{start}
\begin{proof}
  کله چې~\LR{$m = 1$} وي، دا يوازې لاندې تابع ده:
  \[
  f(\vec x) = \fn{cond}(\Char{\lnot R_0}(\vec x),g_0(\vec x),g_1(\vec x)).
  \]
  کله چې~\LR{$m$} له~\LR{$1$} څخه لوے وي، يوازې د همدې بڼې تعريفونه سره ترکيبولے شو۔
\end{proof}
\olpseganchor{OLP-0217-B015}{end}

\olpseganchor{OLP-0218-B004}{start}
\olfileid{cmp}{rec}{bmi}
\olsection{محدوده کمينه‌موندنه}
\olpseganchor{OLP-0218-B004}{end}

\olpseganchor{OLP-0218-B005}{start}
\begin{explain}
ډېر وخت دا ګټوره وي چې يوه تابع د هغه تر ټولو کوچني عدد په توګه تعريف کړو چې کوم خاصيت يا اړيکه~\LR{$P$} پوره کوي۔ که~\LR{$P$} پرېکړه کېدونکې وي، دا تابع يوازې په دې ډول محاسبه کولے شو چې ټول ممکن عددونه~\LR{$0$}، \LR{$1$}، \LR{$2$}، \dots يو په بل پسې وازمېيو، تر هغه چې تر ټولو کوچنے عدد ومومو چې~\LR{$P$} پوره کوي۔ دا ډول بې‌پولې پلټنه مو د بنسټيزو بازګشتي تابعو له چوکاټه باسي۔ خو که يوازې په هغه تر ټولو کوچني عدد کښې دلچسپي ولرو چې
\emph{له کوم خپلواکه ورکړل شوي حد څخه کوچنے وي}، نو بنسټيز بازګشتي پاتې کېږو۔ په بله او لږ عمومي وينا، فرض کړئ يوه بنسټيزه بازګشتي اړيکه~\LR{$R(x,z)$} لرو۔ هغه تابع په پام کښې ونيسئ چې \LR{$x$} او~\LR{$y$} هغه تر ټولو کوچني~\LR{$z < y$} ته لېږي چې~\LR{$R(x, z)$} پوره کوي۔ دا تابع هم د~\LR{$R(x, 0)$}، \LR{$R(x, 1)$}، \dots، \LR{$R(x, y-1)$} په ازمېيلو محاسبه کېدے شي۔ خو دا ولې بنسټيزه بازګشتي ده؟
\end{explain}
\olpseganchor{OLP-0218-B005}{end}

\olpseganchor{OLP-0218-B006}{start}
\begin{prop}
که~\LR{$R(\vec x, z)$} بنسټيزه بازګشتي وي، نو تابع~\LR{$m_R(\vec{x}, y)$} هم داسې ده؛ دا تابع، که داسې عدد وي، له~\LR{$y$} څخه تر ټولو کوچنے~\LR{$z$} راګرځوي چې~\LR{$R(\vec x, z)$} پوره کړي، او که داسې عدد نۀ وي نو~\LR{$y$} راګرځوي۔ تابع~\LR{$m_R$} به داسې ليکو:
\[
\bmin{z < y}{R(\vec{x}, z)},
\]
\end{prop}
\olpseganchor{OLP-0218-B006}{end}

\olpseganchor{OLP-0218-B007}{start}
\begin{proof}
داسې هېڅ~\LR{$z < 0$} نشته چې~\LR{$R(\vec{x}, z)$} پوره کړي، ځکه په اصل کښې هېڅ~\LR{$z < 0$} نشته۔ نو~\LR{$m_R(\vec x, 0) = 0$}۔
\olpseganchor{OLP-0218-B007}{end}

\olpseganchor{OLP-0218-B008}{start}
کله چې حد د~\LR{$y + 1$} په بڼه وي، درې حالتونه لرو:
\begin{enumerate}
\item داسې~\LR{$z < y$} شته چې~\LR{$R(\vec{x}, z)$} پوره کوي؛ په دې حالت کښې \LR{$m_R(\vec{x}, y+1) = m_R(\vec{x}, y)$}۔
\item داسې~\LR{$z<y$} نشته، خو~\LR{$R(\vec{x}, y)$} پوره کېږي؛ بيا \LR{$m_R(\vec{x}, y+1) = y$}۔
\item داسې~\LR{$z < y+1$} نشته چې~\LR{$R(\vec{x}, z)$} پوره کړي؛ بيا \LR{$m_R(\vec{x}, y+1) = y+1$}۔
\end{enumerate}
\textbf{د ژباړې د نسخې سمون (OLCMP-007):} سرچينې په درېيم حالت کښې د تابع د پاراميټرونو غلط وکټور ليکلی و، حال دا چې په ټول تعريف کښې هماغه لومړنی وکټور ثابت پاتې کېږي۔ دلته آرګومېنټ بېرته هماغه ثابت وکټور ته سم شوے دے۔
نو د~\LR{$m_R(\vec x, 0)$} پېلنی حالت په بنسټيز بازګښت د لاندې بشپړ تعريف په ترڅ کښې ورکولے شو:
\begin{align*}
m_R(\vec{x}, 0) & = 0\\
m_R(\vec{x}, y+1) & =
\begin{cases}
m_R(\vec{x}, y) & \text{\RL{که \LR{$m_R(\vec x, y) \neq y$} وي}}\\
y & \text{\RL{که \LR{$m_R(\vec x, y) = y$} او \LR{$R(\vec{x}, y)$} وي}}\\
y+1 & \text{\RL{که نۀ۔}}
\end{cases}
\end{align*}
پام وکړئ چې داسې~\LR{$z<y$} شته چې~\LR{$R(\vec x, z)$} پوره کړي، هله او يوازې هله چې~\LR{$m_R(\vec x, y) \neq y$} وي۔
\end{proof}
\olpseganchor{OLP-0218-B008}{end}

\olpseganchor{OLP-0218-B009}{start}
\begin{prob}
فرض کړئ~\LR{$R(\vec x, z)$} بنسټيزه بازګشتي ده۔ تابع~\LR{$m'_R(\vec{x}, y)$} چې، که داسې عدد وي، له~\LR{$y$} څخه تر ټولو کوچنے~\LR{$z$} راګرځوي چې \LR{$R(\vec{x}, z)$} پوره کړي، او که داسې عدد نۀ وي نو~\LR{$0$} راګرځوي، له~\LR{$\Char{R}$} څخه په بنسټيز بازګښت تعريف کړئ۔
\end{prob}
\olpseganchor{OLP-0218-B009}{end}

\olpseganchor{OLP-0219-B004}{start}
\olfileid{cmp}{rec}{pri}
\olsection{اوليه عددونه}
\olpseganchor{OLP-0219-B004}{end}

\olpseganchor{OLP-0219-B005}{start}
محدوده کميت ټاکنه او محدوده کمينه‌موندنه د دې دپاره ډېر وسايل راکوي چې وښيو طبيعي تابعې او اړيکې بنسټيزې بازګشتي دي۔ د بېلګې په توګه اړيکه «\LR{$x$}، \LR{$y$} وېشي» په پام کښې ونيسئ، چې په~\LR{$x \mid y$} ليکل کېږي۔ اړيکه~\LR{$x \mid y$} هله صدق کوي چې~\LR{$y$} پر~\LR{$x$} بې له پاتې شوني ووېشل شي، يعنې~\LR{$y$} د~\LR{$x$} صحيح مضرب وي۔ که دا اړيکه صدق ونۀ کړي، يعنې کله چې د~\LR{$y$} پر~\LR{$x$} وېش د پاتې شوني په ژبه بيانېدے شي او پاتې شونی~\LR{$> 0$} وي، نو~\LR{$x \nmid y$} ليکو۔ په بله وينا، \LR{$x \mid y$} هله او يوازې هله چې د کوم~\LR{$z$} دپاره~\LR{$x \cdot z = y$} وي۔ ښکاره ده چې هر داسې~\LR{$z$}، که موجود وي، بايد~\LR{$\leq y$} وي۔ نو~\LR{$x \mid y$} هله او يوازې هله چې د کوم~\LR{$z \le y$} دپاره~\LR{$x \cdot z = y$} وي۔ له~\LR{$=$} او ضرب څخه په محدودې وجودي کميت ټاکنې اړيکه~\LR{$x \mid y$} داسې تعريفولے شو:
\[
x \mid y \defiff \bexists{z \leq y}{(x \cdot z) = y}.
\]
نو مو وښودله چې~\LR{$x \mid y$} بنسټيزه بازګشتي ده۔
\textbf{د ژباړې د نسخې سمون (OLCMP-008):} سرچينې د وېش د اړيکې د ناکامېدو په تشريح کښې مقسوم او مقسوم‌عليه سرچپه کړي وو، او د صفر وېشونکي دپاره ئې هم د پاتې شوني ژبه کارولې وه۔ دلته لوری سم شوے او عمومي حالت ورپسې وجودي شرط ټاکي۔
\olpseganchor{OLP-0219-B005}{end}

\olpseganchor{OLP-0219-B006}{start}
طبيعي عدد~\LR{$x$} هله
\emph{اوليه} دے چې نۀ~\LR{$0$} وي، نۀ~\LR{$1$}، او يوازې پر~\LR{$1$} او پر خپل ځان وېشل کېږي۔ په بله وينا، اوليه عددونه داسې دي چې کله هم~\LR{$y \mid x$} وي، نو يا~\LR{$y = 1$} وي يا~\LR{$y=x$}۔ د دې د ازمېيلو دپاره چې~\LR{$x$} اوليه دے که نۀ، يوازې د ټولو~\LR{$y \le x$} دپاره \LR{$y \mid x$} ازمېيل پکار دي، ځکه که~\LR{$y > x$} وي، نو په خپله \LR{$y \nmid x$} وي۔ نو اړيکه~\LR{$\fn{Prime}(x)$}، چې هله او يوازې هله صدق کوي چې~\LR{$x$} اوليه وي، داسې تعريفېدے شي:
\[
\fn{Prime}(x) \defiff x \geq 2 \land \bforall{y \leq x}{(y \mid x \lif y = 1 \lor y = x)}
\]
او ځکه بنسټيزه بازګشتي ده۔
\olpseganchor{OLP-0219-B006}{end}

\olpseganchor{OLP-0219-B007}{start}
اوليه عددونه~\LR{$2$}، \LR{$3$}، \LR{$5$}، \LR{$7$}، \LR{$11$}، او داسې نور دي۔ تابع~\LR{$p(x)$} په پام کښې ونيسئ چې په دې لړۍ کښې~\LR{$x$}م اوليه عدد راګرځوي؛ يعنې \LR{$p(0) = 2$}، \LR{$p(1) = 3$}، \LR{$p(2) = 5$}، او داسې نور۔ (د اسانتيا دپاره به ډېر وخت~\LR{$p(x)$} د~\LR{$p_x$} په توګه ليکو: \LR{$p_0=2$}، \LR{$p_1=3$}، او داسې نور۔)
\olpseganchor{OLP-0219-B007}{end}

\olpseganchor{OLP-0219-B008}{start}
که يوه تابع~\LR{$\fn{nextPrime}(x)$} مو لرله چې له~\LR{$x$} څخه لومړنے لوے اوليه عدد راګرځوي، نو~\LR{$p$} په بنسټيز بازګښت په اسانۍ تعريفېدے شي:
\begin{align*}
  p(0) & = 2\\
  p(x+1) & = \fn{nextPrime}(p(x))
\end{align*}
څرنګه چې~\LR{$\fn{nextPrime}(x)$} هغه تر ټولو کوچنے~\LR{$y$} دے چې~\LR{$y > x$} وي او~\LR{$y$} اوليه وي، دا په بې‌پولې پلټنه په اسانۍ محاسبه کېدے شي۔ خو د اقليدس د يوې پايلې له امله ئې په محدودې کمينه‌موندنې هم تعريفولے شو: تل داسې يو اوليه عدد شته چې له~\LR{$x$} څخه لوے او له \LR{$\fact{x}+1$} څخه کوچنے يا برابر وي۔
\[
  \fn{nextPrime}(x) =
  \bmin{y \leq \fact{x}+1}{(y > x \land \fn{Prime}(y))}.
\]
دا ښيي چې~\LR{$\fn{nextPrime}(x)$}، او ځکه~\LR{$p(x)$}، يوازې محاسبه کېدونکې نۀ بلکې بنسټيزه بازګشتي ده۔
\textbf{د ژباړې د نسخې سمون (OLCMP-009):} سرچينې د راتلونکي اوليه عدد د تابع آرګومېنټ په دوو ځايونو کښې د تابع د نوم دننه ايښی و، حال دا چې نورې کارونې نوم او آرګومېنټ جلا ليکي۔ دلته دواړه ځایونه له هماغې رسمي کارونې سره يو شان شول۔
\olpseganchor{OLP-0219-B008}{end}

\olpseganchor{OLP-0219-B009}{start}
(که لېواله يئ، دا د اقليدس د قضيې يو لنډ ثبوت دے۔ که ورکړل شوے عدد له دوو کوچنے وي، 2 نېغ په نېغه مطلوب اوليه عدد دے۔ په پاتې حالت کښې فرض کړئ~\LR{$p_n$} تر ټولو لوے اوليه عدد دے چې~\LR{$\le x$} وي۔ د ټولو هغو اوليه عددونو حاصل‌ضرب \LR{$p = p_0\cdot p_1 \cdot \dots \cdot p_n$} په پام کښې ونيسئ چې \LR{$\le x$} دي۔ يا~\LR{$p+1$} اوليه دے، يا د~\LR{$x$} او~\LR{$p+1$} تر منځ بل اوليه عدد شته۔ ولې؟ فرض کړئ~\LR{$p+1$} اوليه نۀ دے۔ بيا کوم اوليه عدد \LR{$q \mid p+1$} شته چې~\LR{$q < p+1$} وي۔ اوليه عددونه~\LR{$\le x$} يو هم~\LR{$p+1$} نۀ وېشي۔ (د~\LR{$p$} د تعريف له مخې هر اوليه \LR{$p_i \le x$}، \LR{$p$} وېشي، يعنې پاتې شونی~\LR{$0$} ورکوي۔ نو هر \LR{$p_i \le x$} د~\LR{$p+1$} په وېش کښې پاتې شونی~\LR{$1$} ورکوي، او ځکه \LR{$p_i \nmid p+1$}۔) نو~\LR{$q$} يو اوليه عدد دے چې~\LR{$>x$} او~\LR{$<p+1$} دے۔ او \LR{$p \le \fact{x}$}، نو داسې اوليه عدد شته چې~\LR{$>x$} او \LR{$\le \fact{x}+1$} وي۔)
\textbf{د ژباړې د نسخې سمون (OLCMP-010):} د سرچينې لنډ ثبوت سمدستي تر ورکړل شوي عدد پورې تر ټولو لوے اوليه عدد ټاکه، خو دا څيز د لومړنيو دوو طبيعي عددونو دپاره نشته۔ دلته واړه حالتونه جلا وارزول شول او د حاصل‌ضرب استدلال پاتې حالتونو ته محدود شو؛ د قضيې حد او د راتلونکي اوليه عدد د تابع تعريف نۀ دي بدل شوي۔
\olpseganchor{OLP-0219-B009}{end}

\olpseganchor{OLP-0219-B010}{start}
\begin{prob}
صحيح وېش~\LR{$d(x, y)$} د محدودې کمينه‌موندنې په کارولو تعريف کړئ۔
\end{prob}
\olpseganchor{OLP-0219-B010}{end}

\olpseganchor{OLP-0220-B004}{start}
\olfileid{cmp}{rec}{seq}
\olsection{لړۍ}
\olpseganchor{OLP-0220-B004}{end}

\olpseganchor{OLP-0220-B005}{start}
د بنسټيزو بازګشتي تابعو سټ د پام وړ ټينګ دے۔ خو کله چې د \emph{لړيو} د سمبالولو يوه مناسبه طريقه جوړه کړو، لا ډېر څه کولے شو۔ د طبيعي عددونو متناهي لړۍ به له طبيعي عددونو سره په لاندې ډول وپېژنو: لړۍ~\LR{$\langle a_0, a_1, a_2, \dots, a_k \rangle$} له دې عدد سره سمون لري:
\[
p_0^{a_0+1} \cdot p_1^{a_1+1} \cdot p_2^{a_2+1} \cdot \dots \cdot p_k^{a_k+1}.
\]
توانونو ته يو ځکه زياتوو چې مثلاً لړۍ~\LR{$\langle 2, 7, 3\rangle$} او \LR{$\langle 2, 7, 3, 0, 0 \rangle$} بېل عددي کوډونه ولري۔ \LR{$0$} او~\LR{$1$} دواړه د تشې لړۍ کوډونه نيولے شو؛ د کره‌والي دپاره دې~\LR{$\emptyseq$} د~\LR{$0$} ښودنه وکړي۔
\olpseganchor{OLP-0220-B005}{end}

\olpseganchor{OLP-0220-B006}{start}
د لړيو دا کوډول د حساب د تش په نامه بنسټيزې قضيې له امله کار کوي: هر طبيعي عدد~\LR{$n \ge 2$} په پوره يوه طريقه په لاندې بڼه ليکل کېدے شي:
\[
n = p_0^{a_0} \cdot p_1^{a_1} \cdot \dots \cdot p_k^{a_k}
\]
چې~\LR{$a_k \ge 1$} وي۔ دا ضمانت کوي چې نقشه~\LR{$\tuple{}(a_0, \dots, a_k) = \tuple{a_0, \dots, a_k}$} يو په يو ده: بېلې لړۍ بېلو عددونو ته لېږل کېږي، او له هر عدد سره تر ټولو زياته يوه لړۍ سمون لري۔
\olpseganchor{OLP-0220-B006}{end}

\olpseganchor{OLP-0220-B007}{start}
اوس به وښيو چې د يوې لړۍ د اوږدوالي موندل، د هغې د~\LR{$i$}م غړي موندل، د لړۍ په پاے کښې د يوه غړي زياتول، او د دوو لړيو نښلول ټولې بنسټيزې بازګشتي چارې دي۔
\olpseganchor{OLP-0220-B007}{end}

\olpseganchor{OLP-0220-B008}{start}
\begin{prop}
  تابع~\LR{$\len{s}$}، چې د لړۍ~\LR{$s$} اوږدوالی راګرځوي، بنسټيزه بازګشتي ده۔
\end{prop}
\olpseganchor{OLP-0220-B008}{end}

\olpseganchor{OLP-0220-B009}{start}
\begin{proof}
  اړيکه~\LR{$R(i, s)$} داسې تعريف کړئ:
  \[
  R(i, s) \text{\RL{ هله او يوازې هله چې }} p_i \mid s \land p_{i+1} \nmid s.
  \]
  \LR{$R$} په څرګند ډول بنسټيزه بازګشتي ده۔ کله چې~\LR{$s$} د يوې ناتشې لړۍ کوډ وي، يعنې
  \[
  s = p_0^{a_0+1} \cdot \dots \cdot p_{k}^{a_{k}+1},
  \]
  \LR{$R(i,s)$} هله صدق کوي چې~\LR{$p_i$} تر ټولو لوے اوليه عدد وي چې \LR{$p_i \mid s$}، يعنې~\LR{$i = k$} وي۔ نو د~\LR{$s$} اوږدوالی هله او يوازې هله~\LR{$i+1$} دے چې~\LR{$p_i$} تر ټولو لوے اوليه عدد وي چې~\LR{$s$} وېشي؛ ځکه دا تعريف ورکولے شو:
  \[
  \len{s} =
  \begin{cases}
    0 & \text{\RL{که \LR{$s = 0$} يا \LR{$s = 1$} وي}} \\
    1 + \bmin{i < s}{R(i, s)} & \text{\RL{که نۀ}}
  \end{cases}
  \]
  دلته محدوده کمينه‌موندنه کارولے شو، ځکه کله چې~\LR{$s$} د يوې لړۍ کوډ وي، يوازې يو~\LR{$i$} اړيکه~\LR{$R(i,s)$} پوره کوي، او که~\LR{$i$} موجود وي نو پخپله له~\LR{$s$} څخه کوچنے دے۔
\end{proof}
\olpseganchor{OLP-0220-B009}{end}

\olpseganchor{OLP-0220-B010}{start}
\begin{prop}
  تابع~\LR{$\fn{append}(s,a)$}، چې د لړۍ~\LR{$s$} په پاے کښې د~\LR{$a$} د ورزياتولو پايله راګرځوي، بنسټيزه بازګشتي ده۔
\end{prop}
\olpseganchor{OLP-0220-B010}{end}

\olpseganchor{OLP-0220-B011}{start}
\begin{proof}
  \LR{$\fn{append}$} داسې تعريفېدے شي:
  \[
  \fn{append}(s,a) =
  \begin{cases}
    2^{a+1} & \text{\RL{که \LR{$s = 0$} يا \LR{$s = 1$} وي}} \\
    s \cdot p_{\len{s}}^{a+1} & \text{\RL{که نۀ۔}}
  \end{cases}
  \]
\end{proof}
\olpseganchor{OLP-0220-B011}{end}

\olpseganchor{OLP-0220-B012}{start}
\begin{prop}
  تابع~\LR{$\fn{element}(s,i)$}، چې د~\LR{$s$} \LR{$i$}م غړے راګرځوي (لومړني غړي ته~\LR{$0$}م وايي)، او که~\LR{$i$} د~\LR{$s$} له اوږدوالي څخه لوے يا ورسره برابر وي نو~\LR{$0$} راګرځوي، بنسټيزه بازګشتي ده۔
\end{prop}
\olpseganchor{OLP-0220-B012}{end}

\olpseganchor{OLP-0220-B013}{start}
\begin{proof}
پام وکړئ چې~\LR{$a$} هله او يوازې هله د~\LR{$s$} \LR{$i$}م غړے دے چې \LR{$p_i^{a+1}$} د~\LR{$p_i$} تر ټولو لوے توان وي چې~\LR{$s$} وېشي؛ يعنې \LR{$p_i^{a+1} \mid s$} وي خو~\LR{$p_i^{a+2} \nmid s$} وي۔ نو:
  \[
  \fn{element}(s,i) =
  \begin{cases}
    0 & \mbox{که \LR{$i \geq \len{s}$} وي} \\
    \bmin{a < s}{(p_i^{a+2} \nmid s)} & \text{\RL{که نۀ۔}}
  \end{cases}
  \]
\end{proof}
\olpseganchor{OLP-0220-B013}{end}

\olpseganchor{OLP-0220-B014}{start}
د پورته تعريف شوو تابعو د رسمي نومونو پر ځاے لا لنډې نښې کاروو۔ د~\LR{$\fn{element}(s,i)$} پر ځاے به~\LR{$(s)_i$} کاروو، او \LR{$\tuple{s_0, \dots, s_k}$} به د لاندې عبارت لنډون وي:
\[
\fn{append}(\fn{append}(\dots \fn{append}(\emptyseq,s_0) \dots),s_k).
\]
پام وکړئ که~\LR{$s$} د~\LR{$k$} اوږدوالی ولري، نو د~\LR{$s$} غړي \LR{$(s)_0$}، \dots،~\LR{$(s)_{k-1}$} دي۔
\olpseganchor{OLP-0220-B014}{end}

\olpseganchor{OLP-0220-B015}{start}
\begin{prop}
تابع~\LR{$\fn{concat}(s,t)$}، چې دوه لړۍ سره نښلوي، بنسټيزه بازګشتي ده۔
\end{prop}
\olpseganchor{OLP-0220-B015}{end}

\olpseganchor{OLP-0220-B016}{start}
\begin{proof}
  داسې تابع~\LR{$\fn{concat}$} غواړو چې دا خاصيت ولري:
  \[
    \fn{concat}(\tuple{a_0, \dots, a_k}, \tuple{b_0, \dots, b_l}) = \tuple{a_0, \dots, a_k, b_0, \dots, b_l}.
  \]
  مرستندويه تابع~\LR{$\fn{hconcat}(s,t,n)$} کاروو چې د~\LR{$t$} لومړۍ~\LR{$n$} نښې له~\LR{$s$} سره نښلوي۔ دا تابع په بنسټيز بازګښت داسې تعريفېدے شي:
  \begin{align*}
    \fn{hconcat}(s,t,0) & = s\\
    \fn{hconcat}(s,t,n+1) & = \fn{append}(\fn{hconcat}(s,t,n),(t)_n)
    \intertext{\RL{بيا~\LR{$\fn{concat}$} داسې تعريفولے شو:}}
    \fn{concat}(s,t) & = \fn{hconcat}(s,t,\len{t}).
  \end{align*}
\end{proof}
\olpseganchor{OLP-0220-B016}{end}

\olpseganchor{OLP-0220-B017}{start}
د~\LR{$\fn{concat}(s,t)$} پر ځاے به~\LR{$s \concat t$} ليکو۔
\olpseganchor{OLP-0220-B017}{end}

\olpseganchor{OLP-0220-B018}{start}
دا به راته ګټوره وي چې د يوې لړۍ عددي کوډ د هغې د اوږدوالي او تر ټولو لوے غړي له مخې محدود کړے شو۔ فرض کړئ~\LR{$s$} د~\LR{$k$} اوږدوالي يوه لړۍ ده او هر غړے ئې له کوم عدد~\LR{$x$} څخه کوچنے يا ورسره برابر دے۔ که لړۍ ناتشه وي، نو~\LR{$s$} تر ټولو زيات~\LR{$k$} اوليه عوامل لري، هر يو ئې تر ټولو زيات~\LR{$p_{k-1}$} دے، او هر يو د~\LR{$s$} په اوليه تجزيه کښې تر ټولو زيات د~\LR{$x+1$} توان ته پورته شوے دے۔ نو داسې تعريف کوو:
\[
\fn{sequenceBound}(x,k) =
\fn{cond}(k,1,p_{\fn{pred}(k)}^{k \cdot (x+1)})
\]
بيا د پورته بيان شوې لړۍ~\LR{$s$} عددي کوډ تر ټولو زيات \LR{$\fn{sequenceBound}(x,k)$} دے۔
\textbf{د ژباړې د نسخې سمون (OLCMP-011):} سرچينې د لړۍ په حد کښې د تشې لړۍ دپاره د منفي شاخص اوليه عدد کاراوه؛ ورپسې پلټنې هم نوماند له حد څخه سخت کوچنے غوښته او له حد سره برابر قانوني کوډ ئې پرېښود۔ دلته شرطي تابع د تشې لړۍ حد يو ټاکي او پلټنه تر حد پورې ټول‌شموله شوې ده۔
\olpseganchor{OLP-0220-B018}{end}

\olpseganchor{OLP-0220-B019}{start}
پر لړيو د داسې حد په لرلو نوې تابعې په محدودې پلټنې تعريفولے شو۔ مثلاً~\LR{$\fn{concat}$} په محدودې پلټنې تعريفولے شو۔ يوازې د هغه څيز (د نښلول شوې لړۍ د عدد) يو بنسټيز بازګشتي \emph{ځانګړتيايي شرط} ليکل او د پلټنې حد ټاکل پکار دي۔ لاندې تعريف کار کوي:
\begin{align*}
  \fn{concat}(s,t) = {} & \bmin{v \leq \fn{sequenceBound}(s+t,\len{s} + \len{t})}{} \\
  & \quad(\len{v} = \len{s} + \len{t} \land {}\\
  & \qquad \bforall{i < \len{s}}{((v)_i = (s)_i)} \land {} \\
  & \qquad \bforall{j < \len{t}}{((v)_{\len{s}+j} = (t)_j)})
\end{align*}
\textbf{د ژباړې د نسخې سمون (OLCMP-012):} سرچينې د دويمې لړۍ د غړو شرط د لومړۍ لړۍ پر غړو د کلي کميت ټاکونکي دننه ايښی و؛ نو د تشې لومړۍ لړۍ په حالت کښې هغه شرط تش‌صدقه کېده او د دويمې لړۍ هېڅ غړے ئې نۀ ټاکه۔ دلته د دواړو لړيو شرطونه د اوږدوالي له شرط سره په خپلواکو عطفونو تړل شوي دي۔
\olpseganchor{OLP-0220-B019}{end}

\olpseganchor{OLP-0220-B020}{start}
\begin{prob}
وښيئ چې داسې بنسټيزه بازګشتي تابع~\LR{$\fn{sconcat}(s)$} شته چې دا خاصيت ولري:
\[
\fn{sconcat}(\tuple{s_0, \dots, s_k}) = s_0 \concat \dots \concat s_k.
\]
\end{prob}
\olpseganchor{OLP-0220-B020}{end}

\olpseganchor{OLP-0220-B021}{start}
\begin{prob}
وښيئ چې داسې بنسټيزه بازګشتي تابع~\LR{$\fn{tail}(s)$} شته چې دا خاصيت ولري:
\begin{align*}
  \fn{tail}(\emptyseq) & = 0 \text{\RL{ او}}\\
  \fn{tail}(\tuple{s_0, \dots, s_{k}}) & = \tuple{s_1, \dots, s_{k}}.
\end{align*}
\end{prob}
\olpseganchor{OLP-0220-B021}{end}

\olpseganchor{OLP-0220-B022}{start}
\begin{prop}
  \ollabel{prop:subseq}
  تابع~\LR{$\fn{subseq}(s, i, n)$} چې د~\LR{$s$} له~\LR{$i$}م غړي څخه پيلېدونکې د \LR{$n$} اوږدوالي فرعي لړۍ راګرځوي، بنسټيزه بازګشتي ده۔
\end{prop}
\olpseganchor{OLP-0220-B022}{end}

\olpseganchor{OLP-0220-B023}{start}
\begin{proof}
  تمرين۔
\end{proof}
\olpseganchor{OLP-0220-B023}{end}

\olpseganchor{OLP-0220-B024}{start}
\begin{prob}
  \olref[cmp][rec][seq]{prop:subseq} ثابت کړئ۔
\end{prob}
\olpseganchor{OLP-0220-B024}{end}

\olpseganchor{OLP-0221-B004}{start}
\olfileid{cmp}{rec}{tre}
\olsection{ونې}
\olpseganchor{OLP-0221-B004}{end}

\olpseganchor{OLP-0221-B005}{start}
ځينې وخت د ونو ښودل د طبيعي عددونو په توګه ګټور وي، کټ مټ لکه لړۍ چې په عددونو ښيو او د هغو خاصيتونه او عملونه د کوډونو پر بنسټيزو بازګشتي اړيکو او تابعو ښيو۔ د دې دپاره لړۍ او د هغو کوډونه کاروو۔ يوه ونه يا يوازې يوه غوټه (چې ښايي لېبل ولري) وي، يا يوه غوټه (چې ښايي لېبل ولري) وي چې له يو شمېر فرعي ونو سره تړلې وي۔ دې غوټې ته د ونې \emph{ريښه}، او هغو فرعي ونو ته چې ورسره تړلې وي د هغې \emph{سمدستي فرعي ونې} وايي۔
\olpseganchor{OLP-0221-B005}{end}

\olpseganchor{OLP-0221-B006}{start}
ونې په بازګشتي ډول د~\LR{$\tuple{k, d_1, \dots, d_k}$} په لړۍ کوډوو، چې \LR{$k$} د سمدستي فرعي ونو شمېر او~\LR{$d_1$}، \dots،~\LR{$d_k$} د سمدستي فرعي ونو کوډونه دي۔ که غوټې لېبلونه ولري، له سمدستي فرعي ونو وروسته ئې شاملولے شو۔ نو هغه ونه چې يوازې يوه د~\LR{$l$} لېبل لرونکې غوټه ولري، په~\LR{$\tuple{0,l}$} کوډېږي؛ او هغه ونه چې يوه د~\LR{$l_1$} لېبل لرونکې ريښه ئې له دوو يوازېنيو غوټو سره تړلې وي، چې لېبلونه ئې~\LR{$l_2$} او~\LR{$l_3$} دي، په~\LR{$\tuple{2, \tuple{0, l_2}, \tuple{0, l_3}, l_1}$} کوډېږي۔
\olpseganchor{OLP-0221-B006}{end}

\olpseganchor{OLP-0221-B007}{start}
\begin{prop}
  \ollabel{prop:subtreeseq}
  تابع~\LR{$\fn{SubtreeSeq}(t)$}، چې د داسې لړۍ کوډ راګرځوي چې غړي ئې د \LR{$t$} کوډ لرونکې ونې د ټولو فرعي ونو کوډونه دي، بنسټيزه بازګشتي ده۔
\end{prop}
\olpseganchor{OLP-0221-B007}{end}

\olpseganchor{OLP-0221-B008}{start}
\begin{proof}
  لومړے پام وکړئ چې~\LR{$\fn{ISubtrees}(t) = \fn{subseq}(t, 1, (t)_0)$} بنسټيزه بازګشتي ده او د ونې~\LR{$t$} د سمدستي فرعي ونو کوډونه راګرځوي۔ اوس مرستندويه تابع~\LR{$\fn{hSubtreeSeq}(t,n)$} داسې تعريفوو چې د ټولو هغو فرعي ونو لړۍ محاسبه کړي چې ريښې ئې د~\LR{$t$} له ريښې څخه تر ټولو زيات~\LR{$n$} څنډې لرې وي۔ د~\LR{$t$} د هغو فرعي ونو لړۍ چې له ريښې څخه تر ټولو زيات~\LR{$0$} څنډې لرې دي---يعنې پخپله د~\LR{$t$} په ريښه پيلېږي---يوازې له~\LR{$t$} څخه جوړه لړۍ ده۔ د~\LR{$n+1$} تر کچې د ټولو فرعي ونو د لړۍ د ترلاسه کولو دپاره، د~\LR{$n$} تر کچې شته فرعي ونې د هغو ټولو سمدستي فرعي ونو له لړۍ سره نښلوو چې د~\LR{$n$} تر کچې له فرعي ونو راځي۔
  \textbf{د ژباړې د نسخې سمون (OLCMP-013):} سرچينې مرستندويه تابع له ريښې څخه په کټ مټ ټاکلي واټن د فرعي ونو لړۍ بلله، خو لاندې بازګشتي ګام مخکينۍ لړۍ هم ساتي او تر ټاکلي واټن پورې ټولې کچې راټولوي۔ دلته تشريح د هماغې معادلې له مجموعي چلند سره برابره شوې ده۔
  د ټولو د لست د ترلاسه کولو دپاره پام وکړئ: که~\LR{$f(x)$} يوه بنسټيزه بازګشتي تابع وي چې د لړيو کوډونه راګرځوي، نو د غړو د مثبت شمېر دپاره تابع~\LR{$g_f(s,k)=f((s)_0) \concat \dots \concat f((s)_{k-1})$} هم بنسټيزه بازګشتي ده؛ د صفر غړو نښلون تشه لړۍ نيسو:
    \begin{align*}
      g(s, 0) & = \emptyseq\\
      g(s, k+1) & = g(s, k) \concat f((s)_k)
    \end{align*}
    مثلاً، که~\LR{$s$} د ونو يوه لړۍ وي، نو \LR{$h(s) = g_{\fn{ISubtrees}}(s, \len{s})$} د~\LR{$s$} د غړو د ټولو سمدستي فرعي ونو لړۍ راکوي۔
    \textbf{د ژباړې د نسخې سمون (OLCMP-014):} سرچينې مرستندويه نښلون د صفرم غړي په تابع پيل کاوه او بيا ئې تر اوږدوالي پورې ټول‌شموله شاخص غځاوه؛ په دې سره ئې د لړۍ له وروستي قانوني شاخص وروسته يو زيات غړے لوسته۔ دلته پاراميټر د لومړيو غړو شمېر دے: صفر تش نښلون ورکوي او هر راتلونکی ګام سملاسي ورپسې قانوني غړے زياتوي۔
    له دې څخه~\LR{$\fn{hSubtreeSeq}$} داسې تعريفولے شو:
    \begin{align*}
      \fn{hSubtreeSeq}(t, 0) & = \tuple{t} \\
      \fn{hSubtreeSeq}(t, n+1) & = \fn{hSubtreeSeq}(t, n) \concat h(\fn{hSubtreeSeq}(t, n)).
    \end{align*}
    د~\LR{$t$} په کوډ شوې ونه کښې د فرعي ونو تر ټولو لوړه کچه، يعنې د ريښې او يوې پاڼه غوټې تر منځ تر ټولو لوے واټن، د کوډ~\LR{$t$} له مخې محدوده ده۔ نو د~\LR{$t$} په کوډ شوې ونه کښې د ټولو فرعي ونو د کوډونو لړۍ~\LR{$\fn{hSubtreeSeq}(t, t)$} ورکوي۔
\end{proof}
\olpseganchor{OLP-0221-B008}{end}

\olpseganchor{OLP-0221-B009}{start}
\begin{prob}
  د~\olref[cmp][rec][tre]{prop:subtreeseq} په ثبوت کښې د \LR{$\fn{hSubtreeSeq}$} تعريف په عمومي ډول تکرارونه راولي۔ يو داسې بل تعريف ورکړئ چې ضمانت وکړي د يوې فرعي ونې کوډ په پايله لست کښې يوازې يو ځل راځي۔
\end{prob}
\olpseganchor{OLP-0221-B009}{end}

\olpseganchor{OLP-0222-B004}{start}
\olfileid{cmp}{rec}{ore}
\olsection{نور بازګښتونه}
\olpseganchor{OLP-0222-B004}{end}

\olpseganchor{OLP-0222-B005}{start}
د جوړو او لړيو په کارولو د بنسټيز بازګښت لا نابلده (او ګټورې) بڼې توجيه کولے شو۔ مثلاً ډېر وخت دوه تابعې يوځاے تعريفول ګټور وي، لکه په لاندې تعريف کښې:
\begin{align*}
h_0(\vec x, 0) & = f_0(\vec x) \\
h_1(\vec x, 0) & = f_1(\vec x) \\
h_0(\vec x, y+1) & = g_0(\vec x, y, h_0(\vec x, y), h_1(\vec x, y)) \\
h_1(\vec x, y+1) & = g_1(\vec x, y, h_0(\vec x, y), h_1(\vec x, y))
\end{align*}
دا د \emph{هممهاله بازګښت} يوه بېلګه ده۔ د تابعو د تعريف بله ګټوره لار دا ده چې د~\LR{$h(\vec x, y+1)$} ارزښت د ټولو ارزښتونو \LR{$h(\vec x, 0)$}، \dots،~\LR{$h(\vec x, y)$} له مخې ورکړو، لکه په لاندې تعريف کښې:
\begin{align*}
h(\vec x, 0) & = f(\vec x) \\
h(\vec x, y+1) & = g(\vec x, y, \tuple{h(\vec x, 0), \dots, h(\vec x, y)}).
\end{align*}
لاندې چوکاټ دا مفکوره لا لنډه بيانوي:
\[
h(\vec x, y) = g(\vec x, y, \tuple{h(\vec x, 0), \dots, h(\vec x, y-1)})
\]
په دې پوهېدو سره چې کله~\LR{$y$}، \LR{$0$} وي، د~\LR{$g$} وروستے آرګومېنټ يوازې تشه لړۍ وي۔ په دواړو بڼو کښې مفکوره دا ده چې د «راتلونکي ګام» د محاسبې پر مهال تابع~\LR{$h$} تر دې ځايه د محاسبه شوو ارزښتونو ټوله لړۍ کارولے شي۔ دې ته \emph{د ارزښتونو د بهير بازګښت} وايي۔ د يوې ځانګړې بېلګې دپاره لاندې ډول تعريف توجيه کولے شي:
\begin{align*}
h(\vec x, y) & = \begin{cases}
  g(\vec x, y, h(\vec x, k(\vec x, y))) & \text{\RL{که \LR{$k(\vec x, y) < y$} وي}} \\
  f(\vec x) & \text{\RL{که نۀ}}
\end{cases}
\end{align*}
په بله وينا، په~\LR{$y$} د~\LR{$h$} ارزښت د~\LR{$h$} د هر مخکيني ارزښت له مخې محاسبه کېدے شي، او هغه مخکينی ارزښت~\LR{$k$} ورکوي۔
\olpseganchor{OLP-0222-B005}{end}

\olpseganchor{OLP-0222-B006}{start}
\begin{prob}
  د پاتې شوني تابع~\LR{$r(x,y)$} د ارزښتونو د بهير په بازګښت تعريف کړئ۔ (که~\LR{$x$} او~\LR{$y$} طبيعي عددونه او~\LR{$y > 0$} وي، نو~\LR{$r(x,y)$} له~\LR{$y$} څخه هغه کوچنے عدد دے چې د کوم~\LR{$z$} دپاره \LR{$x = z\times y + r(x,y)$} وي۔ د کره‌والي دپاره وايو که~\LR{$y=0$} وي، نو~\LR{$r(x,0) = 0$}۔)
\end{prob}
\olpseganchor{OLP-0222-B006}{end}

\olpseganchor{OLP-0222-B007}{start}
بايد په دې فکر وکړئ چې دا تابعې په عادي بنسټيز بازګښت څنګه ترلاسه کېږي۔ د بنسټيز بازګښت يوه وروستۍ بڼه په دې لحاظ لا انعطاف منونکې ده چې د لارې په اوږدو کښې د \emph{پاراميټرونو} (اړخيزو ارزښتونو) بدلولو اجازه ورکوي:
\begin{align*}
h(\vec x, 0) & = f(\vec x) \\
h(\vec x, y+1) & = g(\vec x, y, h(k(\vec x), y))
\end{align*}
دا هم په عادي بنسټيز بازګښت تقليدېدے شي۔ (دا کول ستونزمن دي۔ د يوې لارښوونې دپاره، محاسبه په لاس له پاے څخه بېرته پرانيزئ۔)
\olpseganchor{OLP-0222-B007}{end}

\olpseganchor{OLP-0223-B004}{start}
\olfileid{cmp}{rec}{npr}
\olsection{نابنسټيزې بازګشتي تابعې}
\olpseganchor{OLP-0223-B004}{end}

\olpseganchor{OLP-0223-B005}{start}
بنسټيزې بازګشتي تابعې ټولې هغه تابعې نۀ رانغاړي چې په شهودي ډول محاسبه کېدونکې دي۔ په شهودي ډول بايد څرګنده وي چې د ټولو يوځاييزو بنسټيزو بازګشتي تابعو يو لست~\LR{$f_0$}، \LR{$f_1$}، \LR{$f_2$}،~\dots جوړولے شو، داسې چې د~\LR{$y$} پر ننوت د~\LR{$f_x$} ارزښت په اغېزمنه توګه محاسبه کړے شو؛ په بله وينا، لاندې تعريف شوې تابع~\LR{$g(x,y)$} محاسبه کېدونکې ده:
\[
g(x,y) = f_x(y)
\]
خو بيا لاندې تابع هم محاسبه کېدونکې ده:
\begin{eqnarray*}
h(x) & = & g(x,x) + 1 \\
& = & f_x(x) +1.
\end{eqnarray*}
د هرې بنسټيزې بازګشتي تابع~\LR{$f_i$} دپاره د~\LR{$h$} او~\LR{$f_i$} ارزښتونه په~\LR{$i$} کښې سره بېل دي۔ نو~\LR{$h$} محاسبه کېدونکې خو بنسټيزه بازګشتي نۀ ده؛ د~\LR{$g$} په باب هم همدا ويلے شو۔ دا د کانتور د قطري‌کولو د استدلال يوه «اغېزمنه» بڼه ده۔
\olpseganchor{OLP-0223-B005}{end}

\olpseganchor{OLP-0223-B006}{start}
د هغو محاسبه کېدونکو تابعو لا څرګندې بېلګې هم ورکولے شو چې بنسټيزې بازګشتي نۀ دي۔ مثلاً نښه~\LR{$g^n(x)$} دې~\LR{$g(g(\dots g(x)))$} وښيي، چې ټولې~\LR{$n$} دانې~\LR{$g$}ګانې لري؛ او د تابعو لړۍ~\LR{$g_0,g_1,\dots$} داسې تعريف کړئ:
\begin{eqnarray*}
g_0(x) & = & x+1 \\
g_{n + 1}(x) & = & g_n^x(x)
\end{eqnarray*}
تاسو ارزولے شئ چې هره تابع~\LR{$g_n$} بنسټيزه بازګشتي ده۔ هره ورپسې تابع له مخکينۍ څخه ډېره چټکه لوېږي؛ \LR{$g_1(x)$} له~\LR{$2x$} سره برابره ده، \LR{$g_2(x)$} له~\LR{$2^x \cdot x$} سره برابره ده، او~\LR{$g_3(x)$} نژدې د~\LR{$x$} دانو~\LR{$2$}ګانو د تواني برج په شان لوېږي۔ د اکرمن--پېټر تابع په اصل کښې همدا تابع~\LR{$G(x) = g_x(x)$} ده، او ښودل کېدے شي چې دا له هرې بنسټيزې بازګشتي تابع څخه چټکه لوېږي۔
\olpseganchor{OLP-0223-B006}{end}

\olpseganchor{OLP-0223-B007}{start}
راځئ د بنسټيزو بازګشتي تابعو د شمېرنې مسئلې ته بېرته راشو۔ په ياد ولرئ چې هرې بنسټيزې بازګشتي تابع ته مو سمبوليکه نښه ټاکلې ده؛ نو د نښو شمېرل بس دي۔ هرې نښې~\LR{$F$} ته يو طبيعي عدد~\LR{$\#(F)$} په بازګشتي ډول داسې ټاکلے شو:
\begin{eqnarray*}
\#(\Zero) & = & \langle 0 \rangle \\
\#(\Succ) & = & \langle 1 \rangle \\
\#(\Proj{n}{i}) & = & \langle 2, n, i \rangle \\
\#(\fn{Comp}_{k,l}[H,G_0,\dots,G_{k-1}]) & = & \langle 3,k,l,\#(H),\#(G_0),\dots,\#(G_{k-1}) \rangle \\
\#(\fn{Rec}_l[G,H]) & = & \langle 4, l, \#(G), \#(H) \rangle
\end{eqnarray*}
\textbf{د ژباړې د نسخې سمون (OLCMP-015):} سرچينې په دې کوډولو کښې د پخوانۍ برخې ټاکلې صفر او راتلونکې نښې په ناڅاپي ډول په دوو نورو نښو بدلې کړې وې۔ دلته هماغه رسمي نښې بېرته کارول شوې دي؛ د هغو عددي ټاګونه هماغسې پاتې دي۔
دلته دا حقيقت کاروو چې د عددونو هره لړۍ د لړيو د مخکيني کوډ له لارې د يوه طبيعي عدد په توګه ګڼل کېدے شي۔ پايله دا ده چې هرې نښې ته يو طبيعي عددي کوډ ټاکل شوے دے۔ البته ځينې لړۍ (او ځکه ځينې عددونه) له هېڅ نښې سره سمون نۀ لري؛ خو که~\LR{$i$} د يوې يوځاييزې بنسټيزې بازګشتي تابع د نښې کوډ وي، \LR{$f_i$} هماغه تابع نيسو، او که~\LR{$i$} داسې نښه نۀ کوډوي نو د ثابتې صفر تابع ئې نيسو۔ وروستۍ پايله دا ده چې د يوځاييزو بنسټيزو بازګشتي تابعو د شمېرلو يوه څرګنده طريقه لرو۔
\textbf{د ژباړې د نسخې سمون (OLCMP-016):} سرچينې د کوډ ټاکنې پايله په سرچپه لوري بيان کړې وه، حال دا چې پورته بازګشتي معادلې هرې \emph{نښې} ته د لړۍ يو عددي کوډ ټاکي۔ دلته د نقشې لوری او په ورپسې شمېرنه کښې د يوځاييزې نښې شرط څرګند کړل شول۔
\olpseganchor{OLP-0223-B007}{end}

\olpseganchor{OLP-0223-B008}{start}
(په حقيقت کښې ځينې تابعې، لکه ثابته~\LR{$0$} تابع، په لست کښې تر يو ځل زياتې راځي۔ دا يوازې زموږ د کوډولو مصنوعي پايله نۀ ده، بلکې د دې حقيقت پايله هم ده چې ثابته صفر تابع له يوې څخه زياتې نښې لري۔ وروسته به ووينو چې دا تکرارونه په محاسبه کېدونکې توګه نۀ شي مخنيوے کېدے؛ مثلاً هېڅ داسې محاسبه کېدونکې تابع نشته چې پرېکړه وکړي يوه ورکړل شوې نښه ثابته صفر تابع ښيي که نۀ۔)
\olpseganchor{OLP-0223-B008}{end}

\olpseganchor{OLP-0223-B009}{start}
اوس تابع~\LR{$g(x,y)$} د~\LR{$f_x(y)$} په توګه اخلو، چې~\LR{$f_x$} هماغې شمېرنې ته اشاره کوي چې اوس مو بيان کړه۔ څنګه پوهېږو چې~\LR{$g(x,y)$} محاسبه کېدونکې ده؟ په شهودي ډول خبره څرګنده ده: د~\LR{$g(x,y)$} د محاسبې دپاره لومړے~\LR{$x$} «پرانيزئ» او وګورئ چې ايا دا د يوې يوځاييزې تابع نښه ده که نۀ۔ که وي، پر ننوت~\LR{$y$} د هغې تابع ارزښت محاسبه کړئ۔
\olpseganchor{OLP-0223-B009}{end}

\olpseganchor{OLP-0223-B010}{start}

\begin{digress}
کېدے شي لا مو منلې وي چې (په لږ کار سره!) داسې يو پروګرام، مثلاً په جاوا يا~C++ کښې، ليکلے شو چې دا کار وکړي؛ بيا د چرچ--ټيورينګ اصل ته مراجعه کولے شو، چې وايي هر هغه څه چې په شهودي ډول محاسبه کېدونکې وي، په ټيورينګ ماشين محاسبه کېدے شي۔
\olpseganchor{OLP-0223-B010}{end}

\olpseganchor{OLP-0223-B011}{start}
البته د~\LR{$g(x,y)$} د محاسبه کېدو د ښودلو لا نېغه لار دا ده چې هغه ټيورينګ ماشين په څرګند ډول بيان کړو چې دا تابع محاسبه کوي۔ په ځانګړي ډول، دا به د چرچ--ټيورينګ اصل او شهود ته د مراجعې مخه ونيسي۔ ډېر ژر به دومره وسايل جوړ کړي وي چې د~\LR{$g(x,y)$} محاسبه کېدل د محاسبې د داسې مدل په کارولو وښيو چې پر ټيورينګ ماشين \emph{تقليدېدے} شي: يعنې بازګشتي تابعې۔
\end{digress}

\olpseganchor{OLP-0223-B011}{end}

\olpseganchor{OLP-0224-B004}{start}
\olfileid{cmp}{rec}{par}
\olsection{جزوي بازګشتي تابعې}
\olpseganchor{OLP-0224-B004}{end}


\olpseganchor{OLP-0224-B005}{start}
د بازګشتي تابعو د تعريف د انګېزې دپاره پام وکړئ چې زموږ هغه ثبوت، چې
محاسبه کېدونکې خو نابنسټيزې بازګشتي تابعې شته، په حقيقت کښې تر دې ډېر
څه ثابتوي۔ استدلال ساده و: موږ يوازې دا حقيقت وکاراوه چې تابعې
\LR{$f_0,f_1,\dots$} په داسې ډول په لړ کښې راوستل کېدے شي چې د~\LR{$x$} او~\LR{$y$}
د تابع په توګه \LR{$f_x(y)$} محاسبه کېدونکې وي۔ نو دا استدلال د تابعو پر
\emph{هر هغه ټولګي عملي کېږي چې په دې ډول په لړ کښې راوستل کېدے شي}۔
له دې سره له يوې ستونزې سره مخ کېږو: غواړو محاسبه کېدونکې تابعې په
څرګند ډول بيان کړو؛ خو د محاسبه کېدونکو تابعو د يوې ټولګې هېڅ څرګند
بيان بشپړ کېدے نۀ شي!
\olpseganchor{OLP-0224-B005}{end}

\olpseganchor{OLP-0224-B006}{start}
د وتلو لار دا ده چې \emph{جزوي} تابعو ته هم ځای ورکړو۔ وبه وينو چې
جزوي محاسبه کېدونکې تابعې \emph{په رښتيا} په لړ کښې راوستل کېدے شي۔
په حقيقت کښې نږدې له وړاندې پوهېږو چې داسې ده، ځکه ټيورينګ
  ماشينونه په منظم ډول په لړ کښې راوستل کېدے شي۔ وروسته به خپل قطري
استدلال ته راوګرځو او وبه څېړو چې ولې د جزوي تابعو په ګډون هغه استدلال
مخ ته نۀ ځي۔
\olpseganchor{OLP-0224-B006}{end}

\olpseganchor{OLP-0224-B007}{start}
اوس پوښتنه دا ده: بنسټيزو بازګشتي تابعو ته څه ورزيات کړو چې ټولې جزوي
بازګشتي تابعې تر لاسه کړو؟ دوه کارونه پکار دي:
\begin{enumerate}
\item د بنسټيزو بازګشتي تابعو خپل تعريف داسې بدل کړو چې جزوي تابعو ته
  هم ځای ورکړي۔
\item تعريف ته يو څه \emph{ورزيات} کړو، څو ځينې نوې جزوي تابعې هم پکې
  شاملې شي۔
\end{enumerate}
\olpseganchor{OLP-0224-B007}{end}

\olpseganchor{OLP-0224-B008}{start}
لومړے کار اسان دے۔ د پخوا په شان له صفر، جانشين او پروجکشنونو څخه پيل
کوو، او ټولګے د ترکيب او بنسټيز بازګښت لاندې بندوو۔ يوازينے توپير دا
دے چې د ترکيب او بنسټيز بازګښت تعريفونه بايد د دې امکان دپاره بدل کړو
چې په تعريف کښې ځينې ترمونه تعريف شوي نۀ وي۔ که \LR{$f$} او~\LR{$g$} جزوي تابعې
وي، \LR{$f(x) \fdefined$} به په دې مانا ليکو چې \LR{$f$} پر~\LR{$x$} تعريف شوې ده،
يعنې \LR{$x$} د~\LR{$f$} په تعريف ساحه کښې دے؛ او \LR{$f(x) \fundefined$} به د دې
مخالفې مانا، يعنې دا چې \LR{$f$} پر~\LR{$x$} تعريف شوې نۀ ده، څرګندوي۔
\LR{$f(x) \simeq g(x)$} به په دې مانا کاروو چې يا \LR{$f(x)$} او~\LR{$g(x)$} دواړه
تعريف شوي نۀ دي، يا دواړه تعريف شوي او سره برابر دي۔ دا نښې به د لا
پېچلو ترمونو دپاره هم کاروو۔ دا دود به منو چې که \LR{$h$} او \LR{$g_0$}،
\dots،~\LR{$g_k$} ټولې جزوي تابعې وي، نو
\[
h(g_0(\vec x),\dots,g_k(\vec x))
\]
هله او يوازې هله تعريف شوے دے چې هر~\LR{$g_i$} پر~\LR{$\vec x$} تعريف شوے وي،
او \LR{$h$} پر \LR{$g_0(\vec x)$}، \dots،~\LR{$g_k(\vec x)$} تعريف شوې وي۔ په دې
پوهه د جزوي تابعو د ترکيب او بنسټيز بازګښت تعريفونه هماغه پورته
تعريفونه دي، خو ``\LR{$=$}'' بايد پر ``\LR{$\simeq$}'' بدل کړو۔
\olpseganchor{OLP-0224-B008}{end}

\olpseganchor{OLP-0224-B009}{start}
بنسټيزو بازګشتي تابعو ته چې څه ورزياتوو څو جزوي تابعې تر لاسه کړو،
هغه د \emph{بې‌حده پلټنې عامل} دے۔ که \LR{$f(x,\vec z)$} پر طبيعي عددونو
هره جزوي تابع وي، \LR{$\mu x \; f(x,\vec z)$} داسې تعريف کړئ:
\begin{quote}
  تر ټولو وړوکے \LR{$x$} چې \LR{$f(0,\vec z), f(1,\vec z), \dots, f(x,\vec
  z)$} ټول تعريف شوي وي، او \LR{$f(x,\vec z) = 0$} وي، که داسې~\LR{$x$} موجود وي؛
\end{quote}
او که داسې نۀ وي، \LR{$\mu x \; f(x,\vec z)$} تعريف شوے نۀ دے۔ دا
\LR{$\mu x \; f(x,\vec z)$} په يکتا ډول تعريفوي۔
\olpseganchor{OLP-0224-B009}{end}

\olpseganchor{OLP-0224-B010}{start}
\begin{explain}
پام وکړئ چې زموږ تعريف ټيورينګ ماشينونو، يا الګوريتمونو
يا د محاسبې کوم ځانګړي مدل ته هېڅ اشاره نۀ کوي۔ خو د ترکيب او بنسټيز
بازګښت په شان، د بې‌حده پلټنې تر شا هم عملياتي، محاسبوي شهود شته۔ د
يوې جزوي تابع د محاسبه کېدو په بحث کښې هغه آرګومېنټونه، چې تابع پرې
تعريف شوې نۀ وي، له هغو ننوتونو سره سمون خوري چې محاسبه پرې نۀ درېږي۔
د~\LR{$\mu x \; f(x,\vec z)$} د محاسبې کړنلار داسې ده:
\LR{$f(0,\vec z), f(1,\vec z), f(2,\vec z)$} او ورپسې قيمتونه محاسبه کړئ، تر څو د
صفر قيمت راوګرځي۔ خو که له منځنيو محاسبو څخه کومه يوه ونۀ درېږي، د
\LR{$\mu x \; f(x,\vec z)$} محاسبه هم نۀ درېږي۔
\end{explain}
\olpseganchor{OLP-0224-B010}{end}

\olpseganchor{OLP-0224-B011}{start}
که \LR{$R(x,\vec z)$} هره اړيکه وي، \LR{$\mu x \; R(x,\vec z)$} داسې تعريفېږي:
\LR{$\mu x \; (1 \tsub \Char{R}(x,\vec z))$}۔ په بله وينا، \LR{$\mu x \;
R(x,\vec z)$} د~\LR{$x$} تر ټولو وړوکے قيمت راګرځوي چې \LR{$R(x,\vec z)$} پوره
کوي۔ نو که \LR{$f(x,\vec z)$} هرځاے تعريف شوې تابع وي، \LR{$\mu x \;
f(x,\vec z)$} له \LR{$\mu x \; (f(x,\vec z) = 0)$} سره يو شان دے۔ خو پام
وکړئ چې زموږ لومړنی تعريف لا عمومي دے، ځکه دا امکان هم پرېږدي چې
\LR{$f(x,\vec z)$} په هر ځای کښې تعريف شوې نۀ وي (خو د دې پر خلاف د يوې
اړيکې ځانګړونکې تابع تل هرځاے تعريف شوې وي)۔
\olpseganchor{OLP-0224-B011}{end}

\olpseganchor{OLP-0224-B012}{start}
\begin{defn}
د \emph{جزوي بازګشتي تابعو} سټ له طبيعي عددونو څخه طبيعي عددونو ته د
بېلابېلو ځاييزو کچو د جزوي تابعو تر ټولو کوچنے سټ دے چې صفر، جانشين او
پروجکشنونه لري، او د ترکيب، بنسټيز بازګښت او بې‌حده پلټنې لاندې بند دے۔
\end{defn}
\olpseganchor{OLP-0224-B012}{end}

\olpseganchor{OLP-0224-B013}{start}
البته ځينې جزوي بازګشتي تابعې به په اتفاقي ډول هرځاے تعريف شوې وي،
يعنې د هر آرګومېنټ دپاره تعريف شوې وي۔
\olpseganchor{OLP-0224-B013}{end}

\olpseganchor{OLP-0224-B014}{start}
\begin{defn}
\ollabel{defn:recursive-fn}
د \emph{بازګشتي تابعو} سټ د هغو جزوي بازګشتي تابعو سټ دے چې هرځاے
تعريف شوې وي۔
\end{defn}
\olpseganchor{OLP-0224-B014}{end}

\olpseganchor{OLP-0224-B015}{start}
پر دې د ټينګار دپاره چې يوه بازګشتي تابع په هر ځای کښې تعريف شوې ده،
کله کله ورته ``هرځاے تعريف شوې بازګشتي'' تابع هم ويل کېږي۔
\olpseganchor{OLP-0224-B015}{end}

\olpseganchor{OLP-0225-B004}{start}
\olfileid{cmp}{rec}{nft}
\olsection{د نورمال شکل قضيه}
\olpseganchor{OLP-0225-B004}{end}

\olpseganchor{OLP-0225-B005}{start}
\begin{thm}[د کليني د نورمال شکل قضيه]
\ollabel{thm:kleene-nf}
يوه بنسټيزه بازګشتي اړيکه \LR{$T(e, x, s)$} او يوه بنسټيزه بازګشتي تابع
\LR{$U(s)$} شته چې دا خاصيت لري: که \LR{$f$} هره جزوي بازګشتي تابع وي، نو د
کوم~\LR{$e$} دپاره
\[
f(x) \simeq U(\umin{s}{T(e, x, s)})
\]
د هر~\LR{$x$} دپاره پوره کېږي۔
\end{thm}
\olpseganchor{OLP-0225-B005}{end}

\olpseganchor{OLP-0225-B006}{start}
\begin{explain}
د نورمال شکل د قضيې ثبوت اوږد دے، خو بنسټيزه مفکوره ئې ساده ده۔ هره
جزوي بازګشتي تابع يو \emph{شاخص}~\LR{$e$} لري؛ په شهودي ډول دا هغه عدد دے
چې د تابع پروګرام يا تعريف کوډوي۔ که \LR{$f(x) \fdefined$} وي، محاسبه په
منظم ډول ثبتېدے او په کوم عدد~\LR{$s$} کوډېدے شي، او دا حقيقت، چې \LR{$s$} پر
ننوت~\LR{$x$} د~\LR{$f$} محاسبه کوډوي، يوازې د~\LR{$x$} او تعريف~\LR{$e$} په کارولو سره
په بنسټيز بازګشتي ډول کتل کېدے شي۔ نو اړيکه~\LR{$T$}، يعنې ``د شاخص~\LR{$e$}
تابع د ننوت~\LR{$x$} دپاره يوه محاسبه لري، او \LR{$s$} دا محاسبه کوډوي''،
بنسټيزه بازګشتي ده۔ د محاسبې د بشپړ ريکارډ~\LR{$s$} په ورکړې سره د~\LR{$s$}
``پايله'' د~\LR{$f(x)$} قيمت دے، او هغه هم له~\LR{$s$} څخه په بنسټيز بازګشتي
ډول تر لاسه کېدے شي۔
\olpseganchor{OLP-0225-B006}{end}

\olpseganchor{OLP-0225-B007}{start}
د نورمال شکل قضيه ښيي چې د هرې جزوي بازګشتي تابع د تعريف دپاره يوازې
يوه بې‌حده پلټنه پکار ده۔ په بنسټيز ډول د ټولو عددونو له منځه پلټنه
کوو تر څو داسې عدد ومومو چې د شاخص~\LR{$e$} د تابع محاسبه پر ننوت~\LR{$x$}
کوډوي۔ عددونه~\LR{$e$} د جزوي بازګشتي تابعو د ``نومونو'' په توګه کارولے
شو، او د قضيې په معادله تعريف شوې تابع~\LR{$f$} ته \LR{$\cfind{e}$} ليکلے شو۔
پام وکړئ چې يوه جزوي بازګشتي تابع تر يوه ډېر شاخص لرلے شي---په حقيقت
کښې هره جزوي بازګشتي تابع بې‌شمېره ډېر شاخصونه لري۔
\end{explain}
\olpseganchor{OLP-0225-B007}{end}

\olpseganchor{OLP-0226-B004}{start}
\olfileid{cmp}{rec}{hlt}
\olsection{د درېدنې مسئله}
\olpseganchor{OLP-0226-B004}{end}

\olpseganchor{OLP-0226-B005}{start}
په عمومي ډول د \emph{درېدنې مسئله} دا ده چې د يوې محاسبه کېدونکې
تابع د مشخصې~\LR{$e$} (لکه پروګرام) او يو عدد~\LR{$n$} په ورکړې سره پرېکړه وشي
چې د دې تابع محاسبه پر ننوت~\LR{$n$} درېږي، يعنې کومه پايله ورکوي، که نه۔
اېلن ټيورينګ په نوميالي ډول ثابته کړه چې دا مسئله پخپله د يوې محاسبه
کېدونکې تابع له لارې نۀ شي حل کېدے؛ يعنې تابع
\[
h(e, n) =
\begin{cases}
1 & \text{\RL{که محاسبه \LR{$e$} پر ننوت \LR{$n$} ودرېږي}}\\
0 & \text{\RL{که داسې نۀ وي،}}
\end{cases}
\]
محاسبه کېدونکې نۀ ده۔
\olpseganchor{OLP-0226-B005}{end}

\olpseganchor{OLP-0226-B006}{start}
د جزوي بازګشتي تابعو په چوکاټ کښې د يو پروګرام د مشخصې رول د کليني
د نورمال شکل په قضيه کښې ورکړل شوے شاخص~\LR{$e$} لوبولے شي۔ که \LR{$f$} جزوي
بازګشتي تابع وي، هر هغه~\LR{$e$} چې د نورمال شکل د قضيې معادله پوره کوي،
د~\LR{$f$} يو شاخص دے۔ د عدد~\LR{$e$} په ورکړې سره د نورمال شکل قضيه وايي چې
\[
\cfind{e}(x) \simeq U(\mu s \; T(e, x, s))
\]
جزوي بازګشتي ده، او د هرې جزوي بازګشتي \LR{$f\colon \Nat \to \Nat$}
دپاره يو \LR{$e \in \Nat$} شته چې د ټولو~\LR{$x \in \Nat$} دپاره
\LR{$\cfind{e}(x) \simeq f(x)$} وي۔ په حقيقت کښې د هرې داسې~\LR{$f$} دپاره
يوازې يو نه، بلکې بې‌شمېره ډېر داسې~\LR{$e$} شته۔ \emph{د درېدنې تابع}~\LR{$h$}
داسې تعريفېږي:
\[
h(e, x) =
\begin{cases}
1 & \text{\RL{که \LR{$\cfind{e}(x) \fdefined$} وي}}\\
0 & \text{\RL{که داسې نۀ وي۔}}
\end{cases}
\]
نو \LR{$h(e,x)=0$} هله او يوازې هله چې \LR{$\cfind{e}(x) \fundefined$} وي؛ د
نورمال شکل په ټاکلې شمېرنه کښې هر طبيعي عدد~\LR{$e$} د کومې جزوي بازګشتي
تابع شاخص دے، که هغه تابع په هر ننوت هم تعريف شوې نۀ وي۔
\textbf{د ژباړې د نسخې سمون (OLCMP-017):} سرچينې د نورمال شکل د
شمېرنې تر ټاکلو وروسته هم د داسې عدد بديل ساتلے و چې ګنې د هېڅ جزوي
بازګشتي تابع شاخص نۀ وي۔ خو پورته معادله هر طبيعي شاخص ته يوه جزوي
بازګشتي تابع ټاکي۔ دلته هغه ناممکن بديل له تشريح او قطري ثبوت څخه
لرې کړے شوے دے؛ د درېدنې د تابع تعريف او د قضيې دعوه هماغسې دي۔
\olpseganchor{OLP-0226-B006}{end}

\olpseganchor{OLP-0226-B007}{start}
\begin{thm}
\ollabel{thm:halting-problem}
د درېدنې تابع~\LR{$h$} جزوي بازګشتي نۀ ده۔
\end{thm}
\olpseganchor{OLP-0226-B007}{end}

\olpseganchor{OLP-0226-B008}{start}
\begin{proof}
که \LR{$h$} جزوي بازګشتي وای، دا تابع مو تعريفولے شوه:
\[
d(y) =
\begin{cases}
1 & \text{\RL{که \LR{$h(y, y) = 0$} وي}}\\
\umin{x}{x \neq x} & \text{\RL{که داسې نۀ وي۔}}
\end{cases}
\]
څرنګه چې هېڅ عدد~\LR{$x$} شرط \LR{$x \neq x$} نۀ پوره کوي،
\LR{$\umin{x}{x \neq x}$} نشته؛ نو \LR{$d(y) \fundefined$} هله او يوازې هله
چې \LR{$h(y,y) \neq 0$} وي۔ له دې تعريف څخه دا پايلې راځي:
\begin{enumerate}
\item د شاخص~\LR{$y$} دپاره، \LR{$d(y) \fdefined$} هله او يوازې هله چې
  \LR{$\cfind{y}(y) \fundefined$} وي۔
\item \LR{$d(y) \fundefined$} هله او يوازې هله چې \LR{$\cfind{y}(y)
  \fdefined$} وي۔
\end{enumerate}
که \LR{$h$} جزوي بازګشتي وای، \LR{$d$} به هم جزوي بازګشتي وه۔ نو د کليني د
نورمال شکل د قضيې له مخې دا يو شاخص~\LR{$e_d$} لري۔ د~\LR{$h(e_d,e_d)$} قيمت
وګورئ۔ دوه ممکن حالتونه شته، \LR{$0$} او~\LR{$1$}۔
\begin{enumerate}
\item که \LR{$h(e_d, e_d) = 1$} وي، نو \LR{$\cfind{e_d}(e_d) \fdefined$} دے۔
  خو \LR{$\cfind{e_d} \simeq d$}، او \LR{$d(e_d)$} هله تعريف شوے دے چې
  \LR{$h(e_d,e_d)=0$} وي۔ نو \LR{$h(e_d,e_d) \neq 1$}۔
\item که \LR{$h(e_d, e_d) = 0$} وي، نو د شاخص~\LR{$e_d$} تابع
  \LR{$\cfind{e_d}(e_d) \fundefined$} ده۔ خو بيا هم
  \LR{$\cfind{e_d} \simeq d$}، حال دا چې \LR{$d(e_d)$} تعريف شوے دے؛ دغه
  مساوات به بيا \LR{$\cfind{e_d}(e_d) \fdefined$} غوښتل۔
\end{enumerate}
د شاخص~\LR{$e_d$} په دواړو ممکنو حالتونو کښې تناقض راځي۔ که \LR{$h$} جزوي بازګشتي وای، \LR{$d$} به هم
جزوي بازګشتي وه او د هغې شاخص~\LR{$e_d$} به منل شوی وای؛ خو هېڅ ممکن قيمت
د دغه شاخص له تعريف سره نۀ جوړېږي۔ نو بايد دې پايلې ته ورسېږو چې \LR{$h$}
جزوي بازګشتي کېدے نۀ شي۔
\end{proof}
\olpseganchor{OLP-0226-B008}{end}

\olpseganchor{OLP-0227-B004}{start}
\olfileid{cmp}{rec}{gen}
\olsection{عمومي بازګشتي تابعې}
\olpseganchor{OLP-0227-B004}{end}

\olpseganchor{OLP-0227-B005}{start}
د هرځاے تعريف شويو تابعو د يو سټ د تر لاسه کولو بله لار هم شته۔
هرځاے تعريف شوې تابع \LR{$f(x,\vec z)$} هغه وخت \emph{منظمه} وبولئ چې د
طبيعي عددونو د هرې لړۍ~\LR{$\vec z$} دپاره يو~\LR{$x$} موجود وي چې
\LR{$f(x,\vec z)=0$} وي۔ په بله وينا، منظمې تابعې کټ مټ هغه تابعې دي چې
بې‌حده پلټنه پرې عملي کېدے شي او په پايله کښې يوه هرځاے تعريف شوې
تابع تر لاسه کېږي۔ بې‌حده پلټنه په محتاط ډول يوازې منظمو تابعو ته
محدودولے شو:
\olpseganchor{OLP-0227-B005}{end}

\olpseganchor{OLP-0227-B006}{start}
\begin{defn}
\ollabel{defn:general-recursive}
د \emph{عمومي بازګشتي تابعو} سټ له طبيعي عددونو څخه طبيعي عددونو ته د
بېلابېلو ځاييزو کچو د تابعو تر ټولو کوچنے سټ دے چې صفر، جانشين او
پروجکشنونه لري، او د ترکيب، بنسټيز بازګښت او پر \emph{منظمو} تابعو د
عملي شوې بې‌حده پلټنې لاندې بند دے۔
\end{defn}
\olpseganchor{OLP-0227-B006}{end}

\olpseganchor{OLP-0227-B007}{start}
په څرګند ډول هره عمومي بازګشتي تابع هرځاے تعريف شوې ده۔ د
\olref{defn:general-recursive} او \olref[par]{defn:recursive-fn} تر
منځ توپير دا دے چې په وروستي تعريف کښې د جوړونې په منځ کښې د جزوي
بازګشتي تابعو کارول هم روا دي؛ يوازينے شرط دا دے چې په پاے کښې تر لاسه
شوې تابع هرځاے تعريف شوې وي۔ نو د ``عمومي'' ټکی، چې يو تاريخي پاتې
شونی دے، سم نوم نۀ دے؛ په ظاهره \olref{defn:general-recursive} تر
\olref[par]{defn:recursive-fn} \emph{لږ} عمومي دے۔ خو له نېکه مرغه دا
توپير ظاهري دے؛ که څه هم تعريفونه بېل دي، د عمومي بازګشتي تابعو سټ او
د بازګشتي تابعو سټ يو او هماغه سټ دے۔
\olpseganchor{OLP-0227-B007}{end}

\olpseganchor{OLP-0228-B004}{start}
\olchapter{cmp}{thy}{د محاسبه کېدنې تيوري}
\olpseganchor{OLP-0228-B004}{end}

\olpseganchor{OLP-0228-B005}{start}
\begin{editorial}
  د دې څپرکي مواد بايد وکتل او پراخ کړل شي۔ په ځانګړي ډول تر اوسه
  هېڅ تمرين نۀ لري۔
\end{editorial}
\olpseganchor{OLP-0228-B005}{end}

\olpseganchor{OLP-0229-B004}{start}
\olfileid{cmp}{thy}{int}
\olsection{پېژندنه}
\olpseganchor{OLP-0229-B004}{end}

\olpseganchor{OLP-0229-B005}{start}
د منطق هغه څانګه چې \emph{د محاسبه کېدنې تيوري} بلل کېږي، د تابعو او
سټونو د محاسبه کېدنې يا نسبي محاسبه کېدنې اړوندې مسئلې څېړي۔ دا د
کليني د اغېز نښه ده چې دې مضمون ته پخوا \emph{د بازګښت تيوري} ويل
کېدل، او نن دواړه نومونه په عام ډول کارېږي۔
\olpseganchor{OLP-0229-B005}{end}

\olpseganchor{OLP-0229-B006}{start}
تابع~\LR{$f\colon \Nat \pto \Nat$} هغه وخت \emph{جزوي محاسبه کېدونکې}
وبولئ چې د محاسبې په کوم مدل کښې محاسبه کېدے شي۔ که \LR{$f$} هرځاے تعريف
شوې وي، په ساده ډول به وايو چې \LR{$f$} \emph{محاسبه کېدونکې} ده۔ هغه
اړيکه~\LR{$R$} هم محاسبه کېدونکې بلل کېږي چې ځانګړونکې تابع~\LR{$\Char{R}$} ئې
محاسبه کېدونکې وي۔ که \LR{$f$} او~\LR{$g$} جزوي تابعې وي، \LR{$f(x) \fdefined$} به
په دې مانا ليکو چې \LR{$f$} پر~\LR{$x$} تعريف شوې ده، يعنې \LR{$x$} د~\LR{$f$} په تعريف
ساحه کښې دے؛ او \LR{$f(x) \fundefined$} به د دې مخالفه مانا، يعنې دا چې
\LR{$f$} پر~\LR{$x$} تعريف شوې نۀ ده، څرګندوي۔ \LR{$f(x) \simeq g(x)$} به په دې
مانا کاروو چې يا \LR{$f(x)$} او~\LR{$g(x)$} دواړه تعريف شوي نۀ دي، يا دواړه
تعريف شوي او سره برابر دي۔
\olpseganchor{OLP-0229-B006}{end}

\olpseganchor{OLP-0229-B007}{start}
د محاسبه کېدنې تيوري د محاسبې کوم ځانګړي مدل ته له مراجعې پرته هم
څېړل کېدے شي۔ د دې دپاره ښودل کېږي چې يوه نړيواله جزوي محاسبه
کېدونکې تابع~\LR{$\fn{Un}(k,x)$} شته۔ په دې سره جزوي محاسبه کېدونکې تابعې
په لړ کښې راوستل کېدے شي۔ موږ به \LR{$\cfind{k}$} د~\LR{$k$}-مې يوځاييزې جزوي
محاسبه کېدونکې تابع دپاره وکاروو، چې د
\LR{$\cfind{k}(x) \simeq \fn{Un}(k,x)$} له لارې تعريفېږي۔ (کليني د دې
موخې دپاره \LR{$\{ k \}$} کارولی و، خو په وروستيو کښې دا نښه هومره نۀ ده
کارول شوې۔) لږ لا په عمومي ډول، د هرې ځاييزې کچې جزوي محاسبه کېدونکې
تابعې په يو شان ډول په لړ کښې راوستل کېدے شي، او \LR{$\cfind{k}[n]$} به د
\LR{$k$}-مې \LR{$n$}-ځاييزې جزوي محاسبه کېدونکې تابع دپاره کاروو۔
\textbf{د ژباړې د نسخې سمون (OLCMP-018):} سرچينې په همدې پراګراف کښې
لومړے نړيواله شمېرنه د جزوي محاسبه کېدونکو تابعو دپاره ټاکلې وه، خو
په وروستۍ فقره کښې ئې بې له کومې پېژندنې نوم جزوي بازګشتي تابعو ته
واړاوه۔ دلته نوم د پراګراف له ټاکلي، مدل-خپلواک ټولګي سره يو شان ساتل
شوے دے؛ نښه او ځاييزه کچه نۀ دي بدل شوي۔
\olpseganchor{OLP-0229-B007}{end}

\olpseganchor{OLP-0229-B008}{start}
که \LR{$f(\vec x,y)$} هرځاے تعريف شوې يا جزوي تابع وي، نو
\LR{$\umin{y}{f(\vec x,y)}$} د~\LR{$\vec x$} هغه تابع ده چې تر ټولو وړوکے~\LR{$y$}
راګرځوي چې \LR{$f(\vec x,y)=0$} وي، په دې شرط چې \LR{$f(\vec x,0)$}، \dots،
\LR{$f(\vec x,y-1)$} ټول تعريف شوي وي؛ که داسې~\LR{$y$} نۀ وي،
\LR{$\umin{y}{f(\vec x,y)}$} تعريف شوے نۀ دے۔ که \LR{$R(\vec x,y)$} يوه اړيکه
وي، \LR{$\umin{y}{R(\vec x,y)}$} تر ټولو وړوکے~\LR{$y$} تعريفېږي چې
\LR{$R(\vec x,y)$} رښتيا وي؛ په بله وينا، تر ټولو وړوکے~\LR{$y$} چې
\LR{$1 \tsub \Char{R}(\vec x,y)=0$} وي۔
\olpseganchor{OLP-0229-B008}{end}

\olpseganchor{OLP-0229-B009}{start}
د دې د ښودلو دپاره چې يوه تابع محاسبه کېدونکې ده، دوه لارې شته:
\begin{enumerate}
\item په کلک ډول: يو ټيورينګ ماشين يا جزوي بازګشتي تابع په څرګند ډول
  بيان کړئ، او وښيئ چې همدا موخه شوې تابع محاسبه کوي؛
\item په غير رسمي ډول: هغه الګوريتم بيان کړئ چې تابع محاسبه کوي، او د
  چرچ اصل ته مراجعه وکړئ۔
\end{enumerate}
د دواړو تر منځ تېره پوله نشته؛ د يو الګوريتم مفصل بيان بايد دومره
معلومات ورکړي چې په اصل کښې د سم ټيورينګ ماشين يا د جزوي بازګشتي
تعريفونو د لړۍ طرحه نسبتاً څرګنده وي۔ بشپړ کلک تعريفونه ښايي معلوماتي
نۀ وي، او هڅه به کوو چې د دغو دوو لارو تر منځ مناسب منځنی دريځ ومومو؛
په لنډ ډول، هڅه به کوو داسې شهودي خو کلک ثبوتونه ومومو چې ترې څرګنده
وي کره تعريفونه تر لاسه کېدے شي۔
\olpseganchor{OLP-0229-B009}{end}

\olpseganchor{OLP-0230-B004}{start}
\olfileid{cmp}{thy}{cod}
\olsection{د محاسبو کوډول}
\olpseganchor{OLP-0230-B004}{end}

\olpseganchor{OLP-0230-B005}{start}
د محاسبې په هر مدل کښې لاندې کارونه کېدے شي:
\begin{enumerate}
\item د محاسبه کېدونکو تابعو \emph{تعريفونه} په منظم ډول بيانول۔ د
  بېلګې په توګه، د ټيورينګ ماشين مشخصې، بازګشتي تعريفونه يا د پروګرام
  ليکلو په يوه ژبه کښې پروګرامونه د دې تعريفونو ورکوونکي ګڼلے شئ۔
\item د يو ورکړل شوي ننوت دپاره، د کوم تعريف له مخې د يوې تابع د
  محاسبې بشپړ ريکارډ بيانول۔ د بېلګې په توګه، د ټيورينګ ماشين محاسبه
  د محاسبې په هر ګام کښې د بڼو د لړۍ له لارې بيانېدے شي (د ماشين
  حالت او د ټېپ منځپانګه)۔
\item دا کتل چې د يوې محاسبې يو اټکلي ريکارډ په رښتيا د دې ريکارډ دے
  که نه چې د ورکړل شوي تعريف محاسبه کېدونکې تابع به پر ورکړل شوي ننوت
  څنګه محاسبه شي (هغه ننوت چې تابع پرې تعريف شوې وي، يعنې محاسبه
  درېږي)۔
\item د يوې محاسبې د بشپړ ريکارډ له داسې بيان څخه د ورکړل شوي ننوت
  دپاره د تابع قيمت راايستل۔ د بېلګې په توګه، د يوې درېدونکې ټيورينګ
  ماشين محاسبې په وروستي ګام کښې د ټېپ منځپانګه قيمت دے۔
\end{enumerate}
\olpseganchor{OLP-0230-B005}{end}

\olpseganchor{OLP-0230-B006}{start}
د کوډولو په کارولو سره د محاسبه کېدونکې تابع هر بيان ته داسې عددي
\emph{شاخص} ټاکل کېدے شي چې لارښوونې په محاسبه کېدونکي ډول له شاخصه
بېرته تر لاسه شي۔ په ورته ډول د يوې محاسبې بشپړ ريکارډ هم په يوه عدد
کوډېدے شي۔ په پايله کښې پيدا شوې حسابي اړيکه، ``\LR{$s$} د شاخص~\LR{$e$} د تابع
د محاسبې ريکارډ پر ننوت~\LR{$x$} کوډوي''، او تابع، ``د کوډ~\LR{$s$} لرونکې د
محاسبې لړۍ وت''، بيا محاسبه کېدونکې دي؛ په حقيقت کښې دواړه بنسټيز
بازګشتي دي۔
\olpseganchor{OLP-0230-B006}{end}

\olpseganchor{OLP-0230-B007}{start}
دا بنسټيز حقيقت ډېر ځواکمن دے، او پرېږدي چې د محاسبه کېدنې په اړه ډېرې
حيرانوونکې او مهمې پايلې د محاسبې له غوره شوي مدل څخه په خپلواک ډول
ثابتې کړو۔
\olpseganchor{OLP-0230-B007}{end}

\olpseganchor{OLP-0231-B004}{start}
\olfileid{cmp}{thy}{nfm}
\olsection{د نورمال شکل قضيه}
\olpseganchor{OLP-0231-B004}{end}

\olpseganchor{OLP-0231-B005}{start}
فرض کړئ چې د محاسبه کېدونکو تابعو تعريفونه بيانولے شو، او کتلے شو
چې پر کوم ننوت د هغې تابع د محاسبې د بشپړ ريکارډ يو اټکلي بيان سم دے
که نه۔ نو له دې سره سمه ښکاري چې د محاسبې له مدل څخه په خپلواک ډول د
هرې محاسبه کېدونکې تابع~\LR{$f$} قيمت پر هر ننوت~\LR{$x$} داسې ټاکلے شو:
\begin{enumerate}
  \item د محاسبې د ريکارډونو د ټولو ممکنه بيانونو له منځه پلټنه وکړئ۔
  \item وازمويئ چې ورکړل شوے ريکارډ د~\LR{$f(x)$} د محاسبې ريکارډ دے که نه۔
  \item که سم ريکارډ مو موندلے وي، د~\LR{$f(x)$} قيمت ترې راوباسئ۔
\end{enumerate}
دا چې خبره په رښتيا همداسې ده، د کليني د نورمال شکل د قضيې منځپانګه
ده۔
\olpseganchor{OLP-0231-B005}{end}

\olpseganchor{OLP-0231-B006}{start}
\begin{thm}[د کليني د نورمال شکل قضيه]
\ollabel{thm:normal-form}
يوه بنسټيزه بازګشتي اړيکه~\LR{$T(e,x,s)$} او يوه بنسټيزه بازګشتي تابع
~\LR{$U(s)$} شته چې دا خاصيت لري: که \LR{$f$} هره جزوي محاسبه کېدونکې تابع وي،
نو د کوم~\LR{$e$} دپاره
\[
f(x) \simeq U(\umin{s}{T(e, x, s)})
\]
د هر~\LR{$x$} دپاره پوره کېږي۔
\end{thm}
\olpseganchor{OLP-0231-B006}{end}

\olpseganchor{OLP-0231-B007}{start}
\begin{proof}[د ثبوت خاکه]
د محاسبې د هر مدل دپاره د محاسبه کېدونکې تابع~\LR{$f$} يو بيان په کلک ډول
تعريفېدے او داسې بيان په يوه طبيعي عدد~\LR{$e$} کوډېدے شي۔ د ``محاسبې
لړۍ'' يو مفهوم هم په کلک ډول تعريفېدے شي چې د شاخص~\LR{$e$} تابع پر
ننوت~\LR{$x$} د محاسبه کولو بهير ثبتوي۔ دا ډول د محاسبې لړۍ هم په ورته ډول
په يوه عدد~\LR{$s$} کوډېدے شي۔ دا هر څه په داسې ډول کېدے شي چې
\begin{enumerate}
  \item اړيکه \LR{$T(e,x,s)$}، چې هله او يوازې هله پوره کېږي چې عدد~\LR{$s$} د
  شاخص~\LR{$e$} د تابع د محاسبې لړۍ پر ننوت~\LR{$x$} کوډوي، او
  \item تابع \LR{$U(s)$}، چې په~\LR{$s$} کوډ شوې د محاسبې لړۍ د هغې محاسبې
  وروستۍ پايلې ته وړي،
\end{enumerate}
دواړه محاسبه کېدونکي وي۔ په حقيقت کښې اړيکه~\LR{$T$} او تابع~\LR{$U$} دواړه
بنسټيز بازګشتي دي۔
\end{proof}
\olpseganchor{OLP-0231-B007}{end}

\olpseganchor{OLP-0231-B008}{start}
\begin{explain}
د نورمال شکل د قضيې د کلک ثبوت دپاره بايد د محاسبې يو مدل وټاکو، د
محاسبه کېدونکو تابعو د بيانونو او د محاسبې د لړيو کوډول په تفصيل عملي
کړو، او تصديق کړو چې اړيکه~\LR{$T$} او تابع~\LR{$U$} بنسټيز بازګشتي دي۔ د ډېرو
کارونو دپاره دومره بسنه کوي چې \LR{$T$} او~\LR{$U$} محاسبه کېدونکي او \LR{$U$}
هرځاے تعريف شوې وي۔
\olpseganchor{OLP-0231-B008}{end}

\olpseganchor{OLP-0231-B009}{start}
ښايي د نورمال شکل د قضيې ثبوت په لنډ شعار کښې داسې ښه ياد پاتې شي:
\LR{$\umin{s}{T(e,x,s)}$} پر ننوت~\LR{$x$} د شاخص~\LR{$e$} د تابع د محاسبې يوه لړۍ
لټوي، او که داسې لړۍ وموندل شي، \LR{$U$} ئې وت راګرځوي۔
\end{explain}
\olpseganchor{OLP-0231-B009}{end}

\olpseganchor{OLP-0231-B010}{start}
که د محاسبې مدل جزوي بازګشتي تابعې وي (چې کليني په اصل کښې همدا
کارولې)، قضيه ښيي چې د هرې تابع په تعريف کښې د بې‌حده پلټنې يوازې يوه
کارونه، يعنې يوازې يو \LR{$\umin{y}{f(\vec x,y)}$} عامل، پکار دے۔ په دې
مانا قضيه ښيي چې هره جزوي بازګشتي تابع يو نورمال شکل لري۔
\olpseganchor{OLP-0231-B010}{end}

\olpseganchor{OLP-0231-B011}{start}
\LR{$T$} او~\LR{$U$} د شمېرنې \LR{$\cfind{0}$}، \LR{$\cfind{1}$}، \LR{$\cfind{2}$}، \dots
د تعريف دپاره کارېدے شي۔ له دې وروسته به فرض کوو چې د~\LR{$T$} او~\LR{$U$} يو
مناسب انتخاب مو ثابت کړے دے، او معادله
\[
\cfind{e}(x) \simeq U(\umin{s}{T(e,x,s)})
\]
به د~\LR{$\cfind{e}$} \emph{تعريف} ګڼو۔
\olpseganchor{OLP-0231-B011}{end}

\olpseganchor{OLP-0231-B012}{start}
دا يو بل ګټور حقيقت دے:
\olpseganchor{OLP-0231-B012}{end}

\olpseganchor{OLP-0231-B013}{start}
\begin{thm}
هره جزوي محاسبه کېدونکې تابع بې‌شمېره ډېر شاخصونه لري۔
\end{thm}
\olpseganchor{OLP-0231-B013}{end}

\olpseganchor{OLP-0231-B014}{start}
دا خبره بيا هم په شهودي ډول څرګنده ده۔ د يوې محاسبه کېدونکې تابع له
هر ورکړل شوي بيان څخه يو بل بيان جوړولے شو چې هماغه تابع (د ننوت او
وت هماغه جوړه) محاسبه کوي، خو مثلاً لومړے داسې څه کوي چې پر محاسبه
هېڅ اغېز نۀ لري (لکه دا کتل چې \LR{$0=0$} دے، يا تر~\LR{$5$} پورې شمېرل، او
داسې نور)۔ د بدل شوي بيان شاخص به تل له اصلي شاخصه بېل وي۔ دواړه د
هماغې تابع شاخصونه دي، خو تابع په لږ بېل ډول محاسبه کوي۔
\olpseganchor{OLP-0231-B014}{end}

\olpseganchor{OLP-0232-B004}{start}
\olfileid{cmp}{thy}{smn}
\olsection{د \LR{$s$}-\LR{$m$}-\LR{$n$} قضيه}
\olpseganchor{OLP-0232-B004}{end}

\olpseganchor{OLP-0232-B005}{start}
\begin{explain}
راتلونکې قضيه د داسې دليل له مخې د ``\LR{$s$}-\LR{$m$}-\LR{$n$} قضيې'' په نوم
پېژندل کېږي چې لږ وروسته به څرګند شي۔ ستونزمنه برخه دا پوهېدل دي چې
قضيه په کره ډول څه وايي؛ کله چې ئې بيان درک کړئ، قضيه به بيخي څرګنده
وبرېښي۔
\end{explain}
\olpseganchor{OLP-0232-B005}{end}

\olpseganchor{OLP-0232-B006}{start}
\begin{thm}
\ollabel{thm:s-m-n}
  د طبيعي عددونو د هرې جوړې \LR{$n$} او~\LR{$m$} دپاره يوه بنسټيزه بازګشتي
  تابع~\LR{$s^m_n$} شته، داسې چې د هرې لړۍ
  \LR{$e$}، \LR{$a_0$}، \dots، \LR{$a_{m-1}$}، \LR{$y_0$}، \dots، \LR{$y_{n-1}$} دپاره لرو:
  \[
  \cfind{s^m_n(e, a_0, \dots, a_{m-1})}[n](y_0, \dots, y_{n-1}) \simeq
  \cfind{e}[m+n](a_0, \dots, a_{m-1}, y_0, \dots, y_{n-1}).
\]
\end{thm}
\olpseganchor{OLP-0232-B006}{end}

\olpseganchor{OLP-0232-B007}{start}
\begin{explain}
ګټوره ده چې \LR{$s^m_n$} پر \emph{پروګرامونو} د عمل کوونکې تابع په توګه
وګڼو۔ يعنې \LR{$s^m_n$} د يوې \LR{$(m+n)$}-ځاييزې تابع پروګرام~\LR{$e$} او ثابت
ننوتونه \LR{$a_0$}، \dots، \LR{$a_{m-1}$} اخلي؛ بيا د پاتې دليلونو د
\LR{$n$}-ځاييزې تابع پروګرام \LR{$s^m_n(e, a_0, \dots, a_{m-1})$} راګرځوي۔
که~\LR{$e$} د يوه ټيورينګ ماشين تشريح وګڼئ، نو
\LR{$s^m_n(e, a_0, \dots, a_{m-1})$} هغه ټيورينګ ماشين دے چې د
\LR{$y_0$}، \dots،~\LR{$y_{n-1}$} په ننوت کښې \LR{$a_0$}، \dots،~\LR{$a_{m-1}$} د
ننوت د لړۍ په سر کښې ور زياتوي، او~\LR{$e$} چلوي۔ نو هره~\LR{$s^m_n$} يوازې
هغه بنسټيزه بازګشتي تابع ده چې د اړوند ټيورينګ ماشين کوډ پيدا کوي۔
\textbf{د ژباړې د نسخې سمون (OLCMP-019):} سرچينې په همدې تشريح کښې
لومړے يو پروګرامي شاخص ټاکلے و، خو د راګرځېدونکي پروګرام په نښه کښې
ئې ناڅاپه بل متغير ليکلے و، او ورپسې د ټيورينګ ماشين تشريح هم د هماغه
بل متغير په نوم پيلوله۔ دلته دواړه ځايونه د قضيې له ټاکلي شاخص سره
يو شان ساتل شوي؛ د قضيې معادله نۀ ده بدله شوې۔
\end{explain}
\olpseganchor{OLP-0232-B007}{end}

\olpseganchor{OLP-0233-B004}{start}
\olfileid{cmp}{thy}{uni}
\olsection{نړيواله جزوي محاسبه کېدونکې تابع}
\olpseganchor{OLP-0233-B004}{end}

\olpseganchor{OLP-0233-B005}{start}
\begin{thm}
\ollabel{thm:univ-comp}
يوه نړيواله جزوي محاسبه کېدونکې تابع~\LR{$\fn{Un}(e,x)$} شته۔ په بله
وينا، داسې تابع \LR{$\fn{Un}(e,x)$} شته چې:
\begin{enumerate}
\item \LR{$\fn{Un}(e,x)$} جزوي محاسبه کېدونکې ده۔
\item که \LR{$f(x)$} هره جزوي محاسبه کېدونکې تابع وي، نو يو طبيعي عدد
\LR{$e$} شته، داسې چې د هر~\LR{$x$} دپاره \LR{$f(x) \simeq \fn{Un}(e,x)$} وي۔
\end{enumerate}
\end{thm}
\olpseganchor{OLP-0233-B005}{end}

\olpseganchor{OLP-0233-B006}{start}
\begin{proof}
\LR{$\fn{Un}(e,x) \simeq U(\umin{s}{T(e,x,s)})$} وټاکئ، چې~\LR{$U$} او~\LR{$T$}
د کليني د نورمال شکل په قضيه کښې د
(\olref[nfm]{thm:normal-form}) په شان دي۔
\end{proof}
\olpseganchor{OLP-0233-B006}{end}

\olpseganchor{OLP-0233-B007}{start}
\begin{explain}
دا يوازې د دې خبرې کره بڼه ده چې موږ د جزوي محاسبه کېدونکو تابعو
اغېزمنه شمېرنه لرو۔ نظر دا دے چې که د هغې تابع دپاره~\LR{$f_e$} وليکو چې
\LR{$f_e(x) = \fn{Un}(e,x)$} ئې تعريفوي، نو لړۍ \LR{$f_0$}، \LR{$f_1$}، \LR{$f_2$}،
\dots ټولې جزوي محاسبه کېدونکې تابعې رانغاړي، او \LR{$f_e(x)$} د~\LR{$e$}
او~\LR{$x$} په نسبت ``په يو شان'' محاسبه کېدے شي۔ د ساده‌والي دپاره داسې
دوه‌ځاييزه تابع کاروو چې د يوځاييزو تابعو دپاره نړيواله ده؛ خو د
عددونو د لړيو په کوډولو سره دا خبره په اسانۍ ډېرو دليلونو ته عمومي
کولے شو۔ د بېلګې په توګه، پام وکړئ چې که \LR{$f(x,y,z)$} يوه \LR{$3$}-ځاييزه
جزوي بازګشتي تابع وي، نو تابع
\LR{$g(x) \simeq f((x)_0, (x)_1, (x)_2)$} يوه يوځاييزه جزوي بازګشتي
تابع ده۔
\textbf{د ژباړې د نسخې سمون (OLCMP-020):} سرچينې په وروستۍ جمله کښې
له ممکنه جزوي درې‌ځاييزې تابع څخه جوړه شوې يوځاييزه تابع بې له
``جزوي'' قيده بازګشتي بللې وه، حال دا چې جوړه شوې تابع هماغه وخت
ناتعريفه پاتې کېږي چې اړوند اصلي تابع ناتعريفه وي۔ دلته د تابع جزوي
والی څرګند ساتل شوے دے؛ فورمول هماغه دے۔
\end{explain}
\olpseganchor{OLP-0233-B007}{end}

\olpseganchor{OLP-0234-B004}{start}
\olfileid{cmp}{thy}{nou}
\olsection{هېڅ نړيواله محاسبه کېدونکې تابع نشته}
\olpseganchor{OLP-0234-B004}{end}

\olpseganchor{OLP-0234-B005}{start}
که څه هم داسې جزوي محاسبه کېدونکې تابع شته چې د جزوي محاسبه کېدونکو
تابعو دپاره نړيواله ده، داسې هرځاے تعريف شوې محاسبه کېدونکې تابع
نشته چې د ټولو هرځاے تعريف شويو محاسبه کېدونکو تابعو دپاره نړيواله
وي۔ \textbf{د ژباړې د نسخې سمون (OLCMP-021):} د سرچينې په لومړۍ
جمله کښې لومړۍ تابع د جزوي محاسبه کېدونکو تابعو دپاره ``هرځاے تعريف
شوې'' بلل شوې وه، که څه هم ورپسې قضيه او تشريح ئې په سمه توګه
``نړيواله'' بولي او څرګندوي چې دا تابع پخپله جزوي ده۔ دلته هماغه
موخه شوې ځانګړنه، يعنې نړيوالتوب، ليکل شوې ده۔
\olpseganchor{OLP-0234-B005}{end}

\olpseganchor{OLP-0234-B006}{start}
\begin{thm}
\ollabel{thm:no-univ}
هېڅ نړيواله محاسبه کېدونکې تابع نشته۔ په بله وينا، هره تابع
\LR{$\fn{Un}'(k, x)$} چې دا ځانګړنه ولري چې که \LR{$f(x)$} هرځاے تعريف شوې
محاسبه کېدونکې تابع وي، نو يو طبيعي عدد~\LR{$k$} شته، داسې چې د هر~\LR{$x$}
دپاره \LR{$f(x) = \fn{Un}'(k,x)$} وي، محاسبه کېدونکې نۀ ده۔
\end{thm}
\olpseganchor{OLP-0234-B006}{end}

\olpseganchor{OLP-0234-B007}{start}
\begin{proof}
ثبوت يو ساده قطري استدلال دے: که \LR{$\fn{Un}'(k,x)$} هرځاے تعريف شوې او
محاسبه کېدونکې وای، نو
\[
d(x) = \fn{Un}'(x, x) + 1
\]
به هم هرځاے تعريف شوې او محاسبه کېدونکې وه۔ خو د تعريف له مخې
\LR{$d(k)$} له~\LR{$\fn{Un}'(k,k)$} سره برابره نۀ ده۔ نو د هر~\LR{$k$} دپاره د
\LR{$d(x)$} او~\LR{$\fn{Un}'(k, x)$} قيمتونه لږ تر لږه پر يوه~\LR{$x$}، يعنې
\LR{$x = k$}، سره بېل دي۔
\end{proof}
\olpseganchor{OLP-0234-B007}{end}

\olpseganchor{OLP-0234-B008}{start}
\begin{explain}
پورته \olref[uni]{thm:univ-comp} ښيي چې له دې قطري استدلاله وتلے شو،
خو يوازې په دې بيه چې نړيواله تابع جزوي ومنو۔ يعنې~\LR{$\fn{Un}$} د هرځاے
تعريف شويو محاسبه کېدونکو تابعو دپاره نړيواله ده، خو پخپله هرځاے
تعريف شوې نۀ ده۔ قطري استدلال په جزوي حالت کښې کار نۀ کوي۔
\end{explain}
\olpseganchor{OLP-0234-B008}{end}

\olpseganchor{OLP-0234-B009}{start}
\begin{prob}
  د دې د پوهېدو دپاره چې د \olref{thm:no-univ} د ثبوت قطري استدلال
  ولې په جزوي حالت کښې کار نۀ کوي، تابع
  \LR{$f(x) \simeq \fn{Un}(x,x)+1$} وګورئ۔ ايا دا جزوي محاسبه کېدونکې
  ده؟ که وي، يو شاخص~\LR{$e$} لري، يعنې
  \LR{$f(x) \simeq \fn{Un}(e,x)$}۔ د~\LR{$f(e)$} په اړه څه ويلے شئ؟
\end{prob}
\olpseganchor{OLP-0234-B009}{end}

\olpseganchor{OLP-0235-B004}{start}
\olfileid{cmp}{thy}{hlt}
\olsection{د درېدنې مسئله}
\olpseganchor{OLP-0235-B004}{end}

\olpseganchor{OLP-0235-B005}{start}
د جوړښت له مخې نړيواله جزوي محاسبه کېدونکې تابع~\LR{$\fn{Un}(e,x)$} هله
او يوازې هله تعريف شوې ده چې په~\LR{$e$} کوډ شوې د تابع محاسبه د ننوت~\LR{$x$}
دپاره يو قيمت پيدا کړي۔ طبيعي پوښتنه دا ده چې ايا موږ پرېکړه کولے شو
چې داسې ده که نه۔ په حقيقت کښې دا پرېکړه نۀ شو کولے۔ د محاسبې د
ټيورينګ ماشين په مدل کښې دا په دې مانا ده چې ايا ورکړل شوے ټيورينګ
ماشين پر ورکړل شوي ننوت درېږي که نه، په محاسبوي ډول د پرېکړې وړ نۀ
ده۔ له همدې کبله لاندې قضيه د ``درېدنې د مسئلې د ناپرېکړتيا'' په نوم
پېژندل کېږي۔ لاندې به دوه ثبوتونه ورکړو۔ لومړے زموږ د مخکنۍ بحث لړۍ
پسې غځوي، او دويم ډېر نېغ دے۔
\olpseganchor{OLP-0235-B005}{end}

\olpseganchor{OLP-0235-B006}{start}
\begin{thm}
\ollabel{thm:halting-problem}
فرض کړئ
\[
h(e, x)  =
\begin{cases}
1 & \text{\RL{که \LR{$\fn{Un}(e, x)$} تعريف شوې وي}} \\
0 & \text{\RL{که داسې نۀ وي۔}}
\end{cases}
\]
نو \LR{$h$} محاسبه کېدونکې نۀ ده۔
\end{thm}
\olpseganchor{OLP-0235-B006}{end}

\olpseganchor{OLP-0235-B007}{start}
\begin{proof}
فرض کړئ \LR{$h$} محاسبه کېدونکې ده۔ ښيو چې په دې سره يوه نړيواله محاسبه
کېدونکې تابع تعريفولے شو۔ داسې تعريف کړئ:
\[
\fn{Un'}(e,x) =
\begin{cases}
\fn{Un}(e,x) & \text{\RL{که \LR{$h(e,x) = 1$} وي}} \\
0 & \text{\RL{که داسې نۀ وي۔}}
\end{cases}
\]
خو اوس \LR{$\fn{Un'}(e, x)$} هرځاے تعريف شوې تابع ده، او که~\LR{$h$} محاسبه
کېدونکې وي نو دا هم محاسبه کېدونکې ده۔ د بېلګې په توګه، د بنسټيز
بازګښت په کارولو سره~\LR{$g$} داسې تعريفولے شو:
\begin{align*}
g(0, e, x) & \simeq 0 \\
g(y+1, e, x) & \simeq \fn{Un}(e,x);
\end{align*}
بيا
\[
\fn{Un'}(e,x) \simeq g(h(e,x),e,x).
\]
څرنګه چې \LR{$\fn{Un'}(e,x)$} په هر هغه ځاے کښې له~\LR{$\fn{Un}(e,x)$} سره
برابره ده چې وروستۍ تابع تعريف شوې وي، نو~\LR{$\fn{Un'}$} د هغو جزوي
محاسبه کېدونکو تابعو دپاره نړيواله ده چې هرځاے تعريف شوې وي۔ خو دا له
\olref[nou]{thm:no-univ} سره تناقض دے۔
\end{proof}
\olpseganchor{OLP-0235-B007}{end}

\olpseganchor{OLP-0235-B008}{start}
\begin{proof}
فرض کړئ \LR{$h(e,x)$} محاسبه کېدونکې ده۔ تابع~\LR{$g$} داسې تعريف کړئ:
\[
g(x) =
\begin{cases}
  0                & \text{\RL{که \LR{$h(x,x) = 0$} وي}} \\
  \fundefined & \text{\RL{که داسې نۀ وي۔}}
\end{cases}
\]
تابع~\LR{$g$} جزوي محاسبه کېدونکې ده۔ د بېلګې په توګه، دا د
\LR{$\umin{y}{h(x,x) = 0}$} په توګه تعريفولے شو۔ نو د کوم~\LR{$e$} دپاره د هر
~\LR{$x$} په نسبت \LR{$g(x) \simeq \fn{Un}(e,x)$} ده۔ ايا~\LR{$g$} پر~\LR{$e$} تعريف
شوې ده؟ که تعريف شوې وي، نو د~\LR{$g$} د تعريف له مخې \LR{$h(e,e) = 0$} دے
(کله چې~\LR{$g$} تعريف شوې وي، يوازې د~\LR{$0$} قيمت اخلي)۔ د~\LR{$h$} د تعريف له
مخې دا په دې مانا ده چې \LR{$\fn{Un}(e, e)$} تعريف شوې نۀ ده۔ زموږ د دې
فرض له مخې چې د هر~\LR{$x$} دپاره \LR{$g(x) \simeq \fn{Un}(e, x)$} ده،
\LR{$g(e)$} هم تعريف شوې نۀ ده، چې تناقض دے۔ بل پلو، که \LR{$g(e)$} تعريف شوې
نۀ وي، نو \LR{$h(e,e) \neq 0$}، او ځکه \LR{$h(e,e) = 1$} دے۔ له دې څخه راځي
چې \LR{$\fn{Un}(e, e)$} تعريف شوې ده۔ خو څرنګه چې
\LR{$g(x) \simeq \fn{Un}(e, x)$} ده، نو \LR{$g(e)$} به هم تعريف شوې وه۔ بيا
تناقض راغے۔
\textbf{د ژباړې د نسخې سمون (OLCMP-022):} سرچينې په لومړي حالت کښې
د قوس دننه دا ځانګړنه د درېدنې تابع ته منسوبه کړې وه، حال دا چې هغه
د قضيې له تعريفه هرځاے تعريف شوې ده۔ د استدلال اړوند خبره دا ده چې
جوړه شوې قطري تابع، کله تعريف شوې وي، يوازې صفر راګرځوي۔ دلته همدا
فاعل سم شوے دے؛ د ثبوت فورمولونه نۀ دي بدل شوي۔
\end{proof}
\olpseganchor{OLP-0235-B008}{end}

\olpseganchor{OLP-0235-B009}{start}

\begin{explain}
دا استدلال د ټيورينګ ماشينونو په ژبه هم بيانولے شو۔ فرض کړئ داسې
ټيورينګ ماشين~\LR{$H$} شته چې د يو ټيورينګ ماشين~\LR{$E$} تشريح او يو
ننوت~\LR{$x$} اخلي، او پرېکړه کوي چې~\LR{$E$} پر~\LR{$x$} درېږي که نه۔ بيا به يو
بل ټيورينګ ماشين~\LR{$G$} جوړولے شو چې يو ننوت~\LR{$x$} اخلي، \LR{$H$} چلوي څو
پرېکړه وکړي چې د شاخص~\LR{$x$} ماشين~\LR{$M_x$} پر ننوت~\LR{$x$} درېږي که نه، او
بيا د هغې مخالف کار کوي۔ په بله وينا، که~\LR{$H$} ووايي چې~\LR{$M_x$} پر~\LR{$x$}
درېږي، \LR{$G$} په ناپايه کړۍ کښې ننوځي؛ او که~\LR{$H$} ووايي چې~\LR{$M_x$} پر~\LR{$x$}
نۀ درېږي، نو~\LR{$G$} بس درېږي۔ ايا~\LR{$G$} د خپل شاخص په ننوت درېږي؟ پورته
استدلال ښيي چې دا هله او يوازې هله درېږي چې نۀ درېږي—يو تناقض۔ نو
زموږ دا فرض چې داسې ټيورينګ ماشين~\LR{$H$} شته، بايد ناسم وي۔
\end{explain}

\olpseganchor{OLP-0235-B009}{end}

\olpseganchor{OLP-0236-B004}{start}
\olfileid{cmp}{thy}{rus}
\olsection{د رسل له تناقض سره پرتله}
\olpseganchor{OLP-0236-B004}{end}

\olpseganchor{OLP-0236-B005}{start}
ګټوره ده چې د دې برخې استدلالونه د رسل له تناقض سره پرتله او بېل
کړو:
\begin{enumerate}
\item د رسل تناقض: فرض کړئ \LR{$S = \Setabs{x}{x \notin x}$}۔ بيا \LR{$S
  \in S$} هله او يوازې هله چې \LR{$S \notin S$} وي، چې تناقض دے۔
\olpseganchor{OLP-0236-B005}{end}

\olpseganchor{OLP-0236-B006}{start}
  \emph{پايله:} داسې سټ~\LR{$S$} نشته۔ د ``ټولو سټونو د سټ'' شتون منل د
  سټ تيورۍ له نورو بديهي اصولو سره ناسازګار دي۔
  \textbf{د ژباړې د نسخې سمون (OLCMP-023):} د سرچينې د دوه‌اړخيز
  شرط په دويم اړخ کښې ناڅاپه يو بل لوے توری راغلے و، که څه هم يوازې
  د سټ ټاکلے نوم تعريف شوے او د رسل شرط هم د هماغه سټ خپل غړيتوب
  ازمويي۔ دلته بل توری د ټاکلي سټ په نوم سم شوے دے۔
\olpseganchor{OLP-0236-B006}{end}

\olpseganchor{OLP-0236-B007}{start}
\item د رسل د تناقض يوه بدله بڼه: \LR{$F$} له ټولو تابعو له سټ څخه
  \LR{$\{ 0, 1 \}$} ته هغه ``تابع'' وګڼئ چې داسې تعريفېږي:
  \[
  F(f) =
  \begin{cases}
    1 & \text{\RL{که \LR{$f$} د~\LR{$f$} د تعريف ساحې غړے وي، او \LR{$f(f) = 0$} وي}} \\
    0 & \text{\RL{که داسې نۀ وي}}
  \end{cases}
  \]
  ورته استدلال ښيي چې \LR{$F(F) = 0$} هله او يوازې هله چې \LR{$F(F) = 1$}
  وي، چې تناقض دے۔
\olpseganchor{OLP-0236-B007}{end}

\olpseganchor{OLP-0236-B008}{start}
  \emph{پايله:} \LR{$F$} تابع نۀ ده۔ د ``ټولو تابعو سټ'' دومره لوے دے
  چې د يوې تابع تعريف ساحه نۀ شي کېدے۔
\olpseganchor{OLP-0236-B008}{end}

\olpseganchor{OLP-0236-B009}{start}
\item قطري استدلال: \LR{$f_0$}، \LR{$f_1$}، \dots د جزوي محاسبه کېدونکو تابعو
  شمېرنه وي، او \LR{$G \colon \Nat \to \{ 0, 1 \}$} داسې تعريف کړئ:
  \[
  G(x) =
  \begin{cases}
    1 & \text{\RL{که \LR{$f_x(x)\downarrow = 0$} وي}} \\
    0 & \text{\RL{که داسې نۀ وي}}
  \end{cases}
  \]
  که \LR{$G$} محاسبه کېدونکې وي، نو د کوم~\LR{$k$} دپاره تابع~\LR{$f_k$} ده۔ خو
  بيا \LR{$G(k) = 1$} هله او يوازې هله چې \LR{$G(k) = 0$} وي، چې تناقض دے۔
\olpseganchor{OLP-0236-B009}{end}

\olpseganchor{OLP-0236-B010}{start}
  \emph{پايله:} \LR{$G$} محاسبه کېدونکې نۀ ده۔ پام وکړئ چې د سټ تيورۍ د
  بديهي اصولو له مخې~\LR{$G$} بيا هم تابع ده؛ دلته تناقض نشته، يوازې يوه
  څرګندونه ده۔
\end{enumerate}
\olpseganchor{OLP-0236-B010}{end}

\olpseganchor{OLP-0236-B011}{start}
د جزوي تابعو، محاسبه کېدونکو تابعو، جزوي محاسبه کېدونکو تابعو او
داسې نورو يادونه ګډوډوونکې کېدے شي۔ له~\LR{$\Nat$} څخه~\LR{$\Nat$} ته د ټولو
جزوي تابعو سټ د څيزونو يوه لويه ټولګه ده۔ ځينې ئې هرځاے تعريف شوې،
ځينې ئې محاسبه کېدونکې، ځينې دواړه هرځاے تعريف شوې او محاسبه
کېدونکې، او ځينې يوه هم نۀ دي۔ په ياد ولرئ چې کله بې له قيده
``تابع'' وايو، موخه مو هرځاے تعريف شوې تابع وي۔ نو دا ډولونه لرو:
\begin{enumerate}
\item محاسبه کېدونکې تابعې
\item هغه جزوي محاسبه کېدونکې تابعې چې هرځاے تعريف شوې نۀ دي
\item هغه تابعې چې محاسبه کېدونکې نۀ دي
\item هغه جزوي تابعې چې نۀ هرځاے تعريف شوې او نۀ محاسبه کېدونکې دي
\end{enumerate}
د دې د بېلولو دپاره ښايي داسې لوے مربع رسمول ګټور وي چې له~\LR{$\Nat$}
څخه~\LR{$\Nat$} ته ټولې جزوي تابعې وښيي؛ بيا پکښې دوه يو پر بل ورغلي
سيمې په نښه کړئ، چې يوه د هرځاے تعريف شويو تابعو او بله د جزوي
محاسبه کېدونکو تابعو وي۔ ښه تمرين دے چې وګورئ د نقشې په هره جوړه شوې
سيمه کښې يو څيز بيانولے شئ که نه۔
\olpseganchor{OLP-0236-B011}{end}

\olpseganchor{OLP-0237-B004}{start}
\olfileid{cmp}{thy}{cps}
\olsection{محاسبه کېدونکي سټونه}
\olpseganchor{OLP-0237-B004}{end}

\olpseganchor{OLP-0237-B005}{start}
د محاسبه کېدنې مفهوم له محاسبه کېدونکو تابعو څخه محاسبه کېدونکو
سټونو ته غځولے شو:
\olpseganchor{OLP-0237-B005}{end}

\olpseganchor{OLP-0237-B006}{start}
\begin{defn}
  فرض کړئ \LR{$S$} د طبيعي عددونو يو سټ دے۔ \LR{$S$} هله او يوازې هله
  \emph{محاسبه کېدونکے} دے چې ځانګړونکې تابع~\LR{$\Char{S}$} ئې محاسبه
  کېدونکې وي۔ په بله وينا، \LR{$S$} هله او يوازې هله محاسبه کېدونکے دے
  چې تابع
\[
\Char{S}(x) =
\begin{cases}
1 & \text{\RL{که \LR{$x \in S$} وي}} \\
0 & \text{\RL{که داسې نۀ وي}}
\end{cases}
\]
محاسبه کېدونکې وي۔ همدارنګه، اړيکه
\LR{$R(x_0, \dots, x_{k-1})$} هله او يوازې هله محاسبه کېدونکې ده چې
ځانګړونکې تابع ئې محاسبه کېدونکې وي۔
\olpseganchor{OLP-0237-B006}{end}

\olpseganchor{OLP-0237-B007}{start}
محاسبه کېدونکي سټونه او اړيکې \emph{د پرېکړې وړ} هم بلل کېږي۔
\end{defn}
\olpseganchor{OLP-0237-B007}{end}

\olpseganchor{OLP-0237-B008}{start}
\begin{explain}
پام وکړئ چې اوس د محاسبه کېدنې څو مفهومونه لرو: د جزوي تابعو، تابعو
او سټونو دپاره۔ دا سره ګډ نۀ کړئ!{} هغه ټيورينګ ماشين چې
  جزوي تابع محاسبه کوي، د تابع په هغو ننوتونو کښې ئې وت راګرځوي چې
  تابع پرې تعريف شوې وي؛ هغه ټيورينګ ماشين چې يو سټ محاسبه کوي، تر
  دې پرېکړې وروسته چې ننوت په سټ کښې دے که نه، يا~\LR{$1$} يا~\LR{$0$}
  راګرځوي۔
\end{explain}
\olpseganchor{OLP-0237-B008}{end}

\olpseganchor{OLP-0238-B004}{start}
\olfileid{cmp}{thy}{ces}
\olsection{په محاسبوي ډول د شمېر وړ سټونه}
\olpseganchor{OLP-0238-B004}{end}

\olpseganchor{OLP-0238-B005}{start}
\begin{defn}
يو سټ هغه وخت \emph{په محاسبوي ډول د شمېر وړ} دے چې تش وي يا د يوې
محاسبه کېدونکې تابع د قيمتونو سټ وي۔
\end{defn}
\olpseganchor{OLP-0238-B005}{end}

\olpseganchor{OLP-0238-B006}{start}
\begin{history}
په محاسبوي ډول د شمېر وړ سټونه د دې پر ځاے \emph{په بازګشتي ډول د
  شمېر وړ} سټونه هم بلل کېږي۔ دا اصلي اصطلاح ده، او نن دواړه نومونه،
او ورسره لنډونونه ``c.e.'' او ``r.e.''، په عام ډول کارېږي۔
\end{history}
\olpseganchor{OLP-0238-B006}{end}

\olpseganchor{OLP-0238-B007}{start}
\begin{explain}
بايد په دې فکر وکړئ چې تعريف څه مانا لري او اصطلاح ولې مناسبه ده۔ نظر
دا دے چې که~\LR{$S$} د محاسبه کېدونکې تابع~\LR{$f$} د قيمتونو سټ وي، نو
\[
S = \{ f(0), f(1), f(2), \dots \},
\]
او ځکه~\LR{$f$} د~\LR{$S$} د غړو د ``شمېرلو'' په توګه ګڼل کېدے شي۔ پام وکړئ
چې د تعريف له مخې~\LR{$f$} لازمه نۀ ده زياتېدونکې تابع وي، يعنې شمېرنه
لازمه نۀ ده په زياتېدونکي ترتيب وي۔ په حقيقت کښې~\LR{$f$} لازمه نۀ ده
يو په يو هم وي، يعنې د~\LR{$S$} په شمېرنه \LR{$f(0)$}، \LR{$f(1)$}، \LR{$f(2)$}،
\dots{} کښې تکرار منل شوے دے۔ د بېلګې په توګه، ثابته تابع
\LR{$f(x) = 0$} سټ~\LR{$\{ 0 \}$} شمېري۔
\end{explain}
\olpseganchor{OLP-0238-B007}{end}

\olpseganchor{OLP-0238-B008}{start}
هر محاسبه کېدونکے سټ په محاسبوي ډول د شمېر وړ دے۔ د دې د ليدلو
دپاره فرض کړئ~\LR{$S$} محاسبه کېدونکے دے۔ که~\LR{$S$} تش وي، نو د تعريف له
مخې په محاسبوي ډول د شمېر وړ دے۔ که نه، \LR{$a$} د~\LR{$S$} هر يو غړے وي۔
تابع~\LR{$f$} داسې تعريف کړئ:
\[
f(x) =
\begin{cases}
x & \text{\RL{که \LR{$\Char{S}(x) = 1$} وي}} \\
a & \text{\RL{که داسې نۀ وي۔}}
\end{cases}
\]
بيا~\LR{$f$} محاسبه کېدونکې تابع ده، او~\LR{$S$} د~\LR{$f$} د قيمتونو سټ دے۔
\olpseganchor{OLP-0238-B008}{end}

\olpseganchor{OLP-0239-B004}{start}
\olfileid{cmp}{thy}{eqc}
\olpseganchor{OLP-0239-B004}{end}

\olpseganchor{OLP-0239-B005}{start}
\olsection[د c.e. سټونو تعريفونه]{په محاسبوي ډول د شمېر وړ سټونو
  معادل تعريفونه}
\textbf{د ژباړې د نسخې سمون (OLCMP-024):} د سرچينې په اوږده سرليک
کښې د تعريفونو د انګرېزي نوم املا دوه توري زيات لرل، او لاندې تشريح کښې د
``هغه ننوتونه چې محاسبه پرې درېږي'' پر ځاے دوه سربلونه سره نښتي وو۔
په پښتو متن کښې دواړه څرګندې سمونيزې تېروتنې سمې شوې دي؛ هېڅ فني
دعوه يا نښه نۀ ده بدله شوې۔
\olpseganchor{OLP-0239-B005}{end}


\olpseganchor{OLP-0239-B006}{start}
لاندې قضيه د دې خبرې څو مهم معادل بيانونه راکوي چې يو سټ په محاسبوي
ډول د شمېر وړ وي۔
\olpseganchor{OLP-0239-B006}{end}

\olpseganchor{OLP-0239-B007}{start}
\begin{thm}
\ollabel{thm:ce-equiv}
فرض کړئ \LR{$S$} د طبيعي عددونو يو سټ دے۔ بيا لاندې خبرې سره معادلې دي:
\begin{enumerate}
\item\ollabel{case:ce} \LR{$S$} په محاسبوي ډول د شمېر وړ دے۔
\item\ollabel{case:ran-pc} \LR{$S$} د يوې \emph{جزوي} محاسبه کېدونکې تابع
د قيمتونو سټ دے۔
\item\ollabel{case:ran-prim} \LR{$S$} تش دے يا د يوې بنسټيزې بازګشتي تابع
د قيمتونو سټ دے۔
\item\ollabel{case:ce-domain} \LR{$S$} د يوې جزوي محاسبه کېدونکې تابع
\emph{د تعريف ساحه} ده۔
\end{enumerate}
\end{thm}
\olpseganchor{OLP-0239-B007}{end}

\olpseganchor{OLP-0239-B008}{start}
\begin{explain}
لومړۍ درې فقرې وايي چې هر ناتش په محاسبوي ډول د شمېر وړ سټ په معادل
ډول د يوې محاسبه کېدونکې تابع، جزوي محاسبه کېدونکې تابع، يا بنسټيزې
بازګشتي تابع په وسيله شمېرل کېدے شي۔ څلورمه فقره وايي چې که~\LR{$S$} په
محاسبوي ډول د شمېر وړ وي، نو د کوم شاخص~\LR{$e$} دپاره
\[
S = \Setabs{x}{\cfind{e}(x) \fdefined}.
\]
په بله وينا، \LR{$S$} د هغو ننوتونو سټ دے چې د~\LR{$\cfind{e}$} محاسبه پرې
درېږي۔ له همدې کبله په محاسبوي ډول د شمېر وړ سټونه کله ناکله
\emph{نيمه د پرېکړې وړ} بلل کېږي: که يو عدد په سټ کښې وي، بالاخره
``هو'' تر لاسه کوئ؛ خو که پکښې نۀ وي، هېڅکله ``نه'' نۀ تر لاسه
کوئ!{}
\end{explain}
\olpseganchor{OLP-0239-B008}{end}

\olpseganchor{OLP-0239-B009}{start}
\begin{proof}
څرنګه چې هره بنسټيزه بازګشتي تابع محاسبه کېدونکې ده او هره محاسبه
کېدونکې تابع جزوي محاسبه کېدونکې ده، \olref{case:ran-prim} پر
\olref{case:ce} دلالت کوي او \olref{case:ce} پر
\olref{case:ran-pc} دلالت کوي۔ (پام وکړئ چې که~\LR{$S$} تش وي، نو~\LR{$S$} د
هغې جزوي محاسبه کېدونکې تابع د قيمتونو سټ دے چې هېڅ ځاے تعريف شوې
نۀ ده۔) که وښيو چې \olref{case:ran-pc} پر
\olref{case:ran-prim} دلالت کوي، لومړۍ درې فقرې مو سره معادلې ښودلې
دي۔
\olpseganchor{OLP-0239-B009}{end}

\olpseganchor{OLP-0239-B010}{start}
نو فرض کړئ \LR{$S$} د جزوي محاسبه کېدونکې تابع~\LR{$\cfind{e}$} د قيمتونو سټ
دے۔ که~\LR{$S$} تش وي، کار بشپړ دے۔ که نه، \LR{$a$} د~\LR{$S$} هر يو غړے وي۔ د
کليني د نورمال شکل د قضيې له مخې ليکلے شو:
\[
\cfind{e}(x) = U(\umin{s}{T(e, x, s)}).
\]
په ځانګړي ډول، \LR{$\cfind{e}(x) \fdefined$} او \LR{$= y$} هله او يوازې هله
چې داسې~\LR{$s$} وي چې \LR{$T(e, x, s)$} او \LR{$U(s) = y$} وي۔ تابع~\LR{$f(z)$} داسې
تعريف کړئ:
\[
f(z) = \begin{cases}
  U((z)_1) & \text{\RL{که \LR{$T(e, (z)_0, (z)_1)$} وي}} \\
  a        & \text{\RL{که داسې نۀ وي۔}}
\end{cases}
\]
بيا~\LR{$f$} بنسټيزه بازګشتي ده، ځکه~\LR{$T$} او~\LR{$U$} داسې دي۔
د ټيورينګ ماشينونو په ژبه، که~\LR{$z$} داسې جوړه
  \LR{$\tuple{(z)_0, (z)_1}$} کوډوي چې \LR{$(z)_1$} د ماشين~\LR{$M_e$} پر
  ننوت~\LR{$(z)_0$} يوه درېدلې محاسبه وي، نو~\LR{$f$} د محاسبې وت راګرځوي؛
  که نه، \LR{$a$} راګرځوي۔
\olpseganchor{OLP-0239-B010}{end}

\olpseganchor{OLP-0239-B011}{start}
بايد وښيو چې~\LR{$S$} د~\LR{$f$} د قيمتونو سټ دے؛ يعنې د هر طبيعي عدد~\LR{$y$}
دپاره، \LR{$y \in S$} هله او يوازې هله چې د~\LR{$f$} په قيمتونو کښې وي۔ په
مخکښ لوري کښې فرض کړئ \LR{$y \in S$}۔ بيا~\LR{$y$} د~\LR{$\cfind{e}$} په قيمتونو
کښې دے، نو د ځينو~\LR{$x$} او~\LR{$s$} دپاره \LR{$T(e,x,s)$} پوره کېږي او
\LR{$U(s) = y$} ده۔ خو بيا \LR{$y = f(\tuple{x,s})$}۔ بل پلو، فرض کړئ~\LR{$y$} د
~\LR{$f$} په قيمتونو کښې دے۔ بيا يا \LR{$y = a$}، يا د کوم~\LR{$z$} دپاره
\LR{$T(e,(z)_0,(z)_1)$} او \LR{$U((z)_1) = y$} دي۔ څرنګه چې په وروستي حالت
کښې \LR{$\cfind{e}((z)_0) \fdefined = y$} ده، په دواړو حالونو کښې~\LR{$y$}
په~\LR{$S$} کښې دے۔
\textbf{د ژباړې د نسخې سمون (OLCMP-025):} سرچينې په وروستۍ نتيجه
کښې يو ناتړلے متغير کارولے و، که څه هم د شاهد کوډ لومړۍ برخه د
محاسبې ننوت ټاکي۔ دلته د نورمال شکل له مخکنۍ جملې سره سم د شاهد له
کوډ څخه اخيستل شوے ننوت ليکل شوے دے۔
\olpseganchor{OLP-0239-B011}{end}

\olpseganchor{OLP-0239-B012}{start}
(نښه \LR{$\cfind{e}(x) \fdefined = y$} دا مانا لري چې ``\LR{$\cfind{e}(x)$}
تعريف شوې او له~\LR{$y$} سره برابره ده۔'' موږ په همدې اسانتيا
\LR{$\cfind{e}(x) = y$} هم کارولے شو؛ خو زيات غشے کله ناکله دا رايادوي
چې له يوې جزوي تابع سره کار لرو۔)
\olpseganchor{OLP-0239-B012}{end}

\olpseganchor{OLP-0239-B013}{start}
د \olref{thm:ce-equiv} د ثبوت د بشپړولو دپاره دومره بسنه کوي چې
وښيو \olref{case:ce} او~\olref{case:ce-domain} سره معادل دي۔ لومړے
ښيو چې \olref{case:ce} پر~\olref{case:ce-domain} دلالت کوي۔ که ورکړل
شوے سټ تش وي، د هغې جزوي محاسبه کېدونکې تابع د تعريف ساحه ده چې هېڅ
ځاے تعريف شوې نۀ ده۔ اوس د ناتش حالت دپاره فرض کړئ~\LR{$S$} د يوې محاسبه
کېدونکې تابع~\LR{$f$} د قيمتونو سټ دے، يعنې
\[
S = \Setabs{y}{\text{\RL{د کوم \LR{$x$} دپاره، }} f(x) = y}.
\]
وټاکئ:
\[
g(y) = \umin{x}{(f(x) = y)}.
\]
بيا~\LR{$g$} جزوي محاسبه کېدونکې تابع ده، او \LR{$g(y)$} هله او يوازې هله
تعريف شوې ده چې د کوم~\LR{$x$} دپاره \LR{$f(x) = y$} وي۔ په بله وينا، د~\LR{$g$}
د تعريف ساحه د~\LR{$f$} د قيمتونو سټ دے۔
د ټيورينګ ماشينونو په ژبه: که ټيورينګ ماشين~\LR{$F$} د~\LR{$S$}
غړي شمېري، نو ټيورينګ ماشين~\LR{$G$} داسې وټاکئ چې په~\LR{$S$} کښې د ورکړل
شوي غړي د موندلو دپاره د~\LR{$F$} په وتونو کښې پلټنه وکړي؛ که ئې ومومي
ودرېږي، او که ئې نۀ مومي د تل دپاره پلټنې ته دوام ورکړي۔
\textbf{د ژباړې د نسخې سمون (OLCMP-026):} سرچينې په دې دلالت کښې
سمدستي فرض کړې وه چې ورکړل شوے سټ د يوې هرځاے تعريف شوې محاسبه
کېدونکې تابع د قيمتونو سټ دے؛ دا د تش سټ دپاره ناشوني دي، که څه هم
تعريف تش سټ په څرګند ډول c.e. مني۔ دلته تش حالت لومړے د هېڅ ځاے
تعريف شوې جزوي تابع په وسيله حل شوے، او د تابع د قيمتونو استدلال
ناتش حالت ته کارول شوے دے۔
\olpseganchor{OLP-0239-B013}{end}

\olpseganchor{OLP-0239-B014}{start}
په پای کښې، د دې د ښودلو دپاره چې \olref{case:ce-domain} پر
\olref{case:ce} دلالت کوي، فرض کړئ~\LR{$S$} د جزوي محاسبه کېدونکې
تابع~\LR{$\cfind{e}$} د تعريف ساحه ده، يعنې
\[
S = \Setabs{x}{\cfind{e}(x) \fdefined}.
\]
که~\LR{$S$} تش وي، کار بشپړ دے؛ که نه، \LR{$a$} د~\LR{$S$} هر يو غړے وي۔ تابع~\LR{$f$}
داسې تعريف کړئ:
\[
f(z) = \begin{cases}
(z)_0 & \text{\RL{که \LR{$T(e,(z)_0,(z)_1)$} وي}} \\
a & \text{\RL{که داسې نۀ وي۔}}
\end{cases}
\]
بيا، لکه پورته، عدد~\LR{$x$} هله او يوازې هله د~\LR{$f$} په قيمتونو کښې دے
چې \LR{$\cfind{e}(x) \fdefined$} وي، يعنې هله او يوازې هله چې \LR{$x \in S$}
وي۔ د ټيورينګ ماشينونو په ژبه: که ماشين~\LR{$M_e$} د~\LR{$S$}
نيمه پرېکړه کوي، د~\LR{$S$} غړي د ټيورينګ ماشينونو د ټولو ممکنه محاسبو
په کتلو او د درېدلو محاسبو د اړوندو ننوتونو په راګرځولو سره وشمېرئ۔
\end{proof}
\olpseganchor{OLP-0239-B014}{end}

\olpseganchor{OLP-0239-B015}{start}
د \olref{thm:ce-equiv} فقره~\olref{case:ce-domain} موږ ته په
محاسبوي ډول د شمېر وړ سټونو د شمېرلو يوه اسانه لار راکوي: د هر~\LR{$e$}
دپاره~\LR{$W_e$} د~\LR{$\cfind{e}$} د تعريف ساحه وټاکئ، يعنې
\[
W_e = \Setabs{x}{\cfind{e}(x) \fdefined}.
\]
بيا که~\LR{$A$} هر په محاسبوي ډول د شمېر وړ سټ وي، د کوم~\LR{$e$} دپاره
\LR{$A = W_e$} وي۔
\olpseganchor{OLP-0239-B015}{end}

\olpseganchor{OLP-0239-B016}{start}
لاندې قضيه په محاسبوي ډول د شمېر وړ سټونو يوه بله ځانګړنه راکوي۔
\olpseganchor{OLP-0239-B016}{end}

\olpseganchor{OLP-0239-B017}{start}
\begin{thm}
\ollabel{thm:exists-char}
سټ~\LR{$S$} هله او يوازې هله په محاسبوي ډول د شمېر وړ دے چې داسې محاسبه
کېدونکې اړيکه~\LR{$R(x,y)$} وي چې
\[
S = \Setabs{ x }{ \lexists[y][R(x,y)] }.
\]
\end{thm}
\olpseganchor{OLP-0239-B017}{end}

\olpseganchor{OLP-0239-B018}{start}
\begin{proof}
په مخکښ لوري کښې فرض کړئ~\LR{$S$} په محاسبوي ډول د شمېر وړ دے۔ نو د کوم
~\LR{$e$} دپاره \LR{$S = W_e$}۔ د~\LR{$e$} د همدې قيمت دپاره~\LR{$S$} داسې ليکلے شو:
\[
S = \Setabs{ x }{ \lexists[y][T(e, x, y)] }.
\]
په بېرته لوري کښې فرض کړئ \LR{$S = \Setabs{ x }{
  \lexists[y][R(x, y)] }$}۔ تابع~\LR{$f$} داسې تعريف کړئ:
\[
f(x) \simeq \umin{y}{R(x, y)}.
\]
بيا~\LR{$f$} جزوي محاسبه کېدونکې ده، او~\LR{$S$} د~\LR{$f$} د تعريف ساحه ده۔
\end{proof}
\olpseganchor{OLP-0239-B018}{end}

\olpseganchor{OLP-0240-B004}{start}
\olfileid{cmp}{thy}{ncp}
\olsection{محاسبه نۀ کېدونکي سټونه شته}
\olpseganchor{OLP-0240-B004}{end}


\olpseganchor{OLP-0240-B005}{start}
پورته مو وليدل چې هر محاسبه کېدونکے سټ په محاسبوي ډول د شمېر وړ
دے۔ ايا د دې معکوس هم رښتيا دے؟ لاندې قضيه ښيي چې په عمومي ډول داسې
نۀ ده۔
\olpseganchor{OLP-0240-B005}{end}

\olpseganchor{OLP-0240-B006}{start}
\begin{thm}\ollabel{thm:K-0}
فرض کړئ \LR{$K_0$} سټ \LR{$\Setabs{\tuple{e, x}}{\cfind{e}(x) \fdefined}$}
دے۔ بيا \LR{$K_0$} په محاسبوي ډول د شمېر وړ دے، خو محاسبه کېدونکے نۀ دے۔
\end{thm}
\olpseganchor{OLP-0240-B006}{end}

\olpseganchor{OLP-0240-B007}{start}
\begin{proof}
د دې د ليدلو دپاره چې \LR{$K_0$} په محاسبوي ډول د شمېر وړ دے، پام وکړئ
چې دا د هغې تابع~\LR{$f$} د تعريف ساحه ده چې داسې تعريفېږي:
\[
f(z) = \umin{y}{(\len{z} = 2 \land T((z)_0, (z)_1, y))}.
\]
ځکه که \LR{$\cfind{e}(x)$} تعريف شوې وي، \LR{$f(\tuple{e, x})$} د درېدلې
محاسبې يوه لړۍ مومي؛ که \LR{$\cfind{e}(x)$} تعريف شوې نۀ وي، نو
\LR{$f(\tuple{e, x})$} هم تعريف شوې نۀ ده؛ او که~\LR{$z$} ان يوه جوړه هم نۀ
کوډوي، نو \LR{$f(z)$} بيا هم تعريف شوې نۀ ده۔
\olpseganchor{OLP-0240-B007}{end}

\olpseganchor{OLP-0240-B008}{start}
دا حقيقت چې \LR{$K_0$} محاسبه کېدونکے نۀ دے، هماغه د درېدنې د مسئلې
ناپرېکړتيا ده، \olref[hlt]{thm:halting-problem}۔
\end{proof}
\olpseganchor{OLP-0240-B008}{end}

\olpseganchor{OLP-0240-B009}{start}
سټ~\LR{$K_0$} د هغو جوړو~\LR{$\tuple{e,x}$} سټ دے چې
\LR{$\cfind{e}(x) \fdefined$} وي، يعنې \LR{$\tuple{e,x} \in K_0$} هله او
يوازې هله چې~\LR{$\cfind{e}$} پر ننوت~\LR{$x$} تعريف شوې وي (ودرېږي)، نو دې
ته ``د درېدنې سټ'' هم ويل کېږي۔ سټ
\LR{$K = \Setabs{e}{\cfind{e}(e) \fdefined}$} د ``پر خپل ځان د درېدنې
سټ'' دے۔ دا ډېر ځله د يوې معياري ناپرېکړه کېدونکې بېلګې په توګه
کارول کېږي۔
\olpseganchor{OLP-0240-B009}{end}

\olpseganchor{OLP-0240-B010}{start}
\begin{thm}\ollabel{thm:K}
پر خپل ځان د درېدنې سټ
\LR{$K = \Setabs{e}{\cfind{e}(e) \fdefined}$}، c.e. دے خو د پرېکړې
وړ نۀ دے۔
\end{thm}
\olpseganchor{OLP-0240-B010}{end}

\olpseganchor{OLP-0240-B011}{start}
\begin{proof}
  فرض کړئ \LR{$K$} د پرېکړې وړ دے، يعنې ځانګړونکې تابع~\LR{$\Char{K}$} ئې
  محاسبه کېدونکې ده۔ وټاکئ:
  \[d(e) = \begin{cases}
  1 & \text{\RL{که \LR{$\Char{K}(e) = 0$} وي}}\\
  \fundefined & \text{\RL{که داسې نۀ وي۔}}
  \end{cases}
  \]
  فرض کړئ \LR{$k$} د~\LR{$d$} شاخص دے، يعنې \LR{$d \simeq \cfind{k}$}۔ بيا
  \LR{$d(k) \simeq \cfind{k}(k)$} ده۔ دا له هغه حقيقت سره تناقض لري چې
  \LR{$d(k) \fdefined$} هله او يوازې هله چې
  \LR{$\cfind{k}(k) \fundefined$} وي، او دا د~\LR{$d$} له تعريفه راځي۔
\olpseganchor{OLP-0240-B011}{end}

\olpseganchor{OLP-0240-B012}{start}
  \LR{$K$} د تابع \LR{$f(x) = \umin{y}{T(x,x,y)}$} د تعريف ساحه ده او ځکه
  c.e. دے۔
\end{proof}
\olpseganchor{OLP-0240-B012}{end}

\olpseganchor{OLP-0241-B004}{start}
\olfileid{cmp}{thy}{clo}
\olpseganchor{OLP-0241-B004}{end}

\olpseganchor{OLP-0241-B005}{start}
\olsection[د c.e. سټونو اتحاد او اشتراک]{په محاسبوي ډول د شمېر وړ
  سټونه د اتحاد او اشتراک لاندې تړلي دي}
\olpseganchor{OLP-0241-B005}{end}

\olpseganchor{OLP-0241-B006}{start}
لاندې قضيه د محاسبوي ډول د شمېر وړ سټونو د ټولګي ځينې تړلتيايي
خاصيتونه راکوي۔
\olpseganchor{OLP-0241-B006}{end}

\olpseganchor{OLP-0241-B007}{start}
\begin{thm}
فرض کړئ \LR{$A$} او~\LR{$B$} په محاسبوي ډول د شمېر وړ دي۔ بيا \LR{$A \cap B$} او
\LR{$A \cup B$} هم په محاسبوي ډول د شمېر وړ دي۔
\end{thm}
\olpseganchor{OLP-0241-B007}{end}

\olpseganchor{OLP-0241-B008}{start}
\begin{proof}
\olref[eqc]{thm:ce-equiv} اجازه راکوي چې په محاسبوي ډول د شمېر وړ
سټونو بېلابېلې ځانګړنې وکاروو۔ د روښانتيا دپاره به څو بېل ثبوتونه
وړاندې کړو۔
\olpseganchor{OLP-0241-B008}{end}

\olpseganchor{OLP-0241-B009}{start}
د لومړي ثبوت دپاره فرض کړئ~\LR{$A$} د محاسبه کېدونکې تابع~\LR{$f$} په وسيله
شمېرل کېږي، او~\LR{$B$} د محاسبه کېدونکې تابع~\LR{$g$} په وسيله شمېرل کېږي۔
وټاکئ:
\begin{align*}
h(x) & = \umin{y}{(f(y) = x \lor g(y) = x)} \text{\RL{ او}}\\
j(x) & = \umin{y}{(f((y)_0) = x \land g((y)_1) = x)}.
\end{align*}
بيا \LR{$A \cup B$} د~\LR{$h$} د تعريف ساحه، او \LR{$A \cap B$} د~\LR{$j$} د تعريف
ساحه ده۔
\olpseganchor{OLP-0241-B009}{end}

\olpseganchor{OLP-0241-B010}{start}
\begin{explain}
په محاسبوي ژبه کښې خبره داسې ده: که داسې کړنلارې راکړل شوې وي چې
\LR{$A$} او~\LR{$B$} شمېري، نو دا چې يو غړے~\LR{$x$} په \LR{$A \cup B$} کښې دے که نه،
په دواړو شمېرنو کښې د~\LR{$x$} په لټولو نيمه پرېکړه کولے شو؛ او دا چې
\LR{$x$} په \LR{$A \cap B$} کښې دے که نه، په دواړو شمېرنو کښې په يو وخت د~\LR{$x$}
په لټولو نيمه پرېکړه کولے شو۔
\end{explain}
\olpseganchor{OLP-0241-B010}{end}

\olpseganchor{OLP-0241-B011}{start}
د دويم ثبوت دپاره بيا فرض کړئ چې~\LR{$A$} د~\LR{$f$} او~\LR{$B$} د~\LR{$g$} په وسيله
شمېرل کېږي۔ وټاکئ:
\[
k(x) = \begin{cases}
f(x/2) & \text{\RL{که \LR{$x$} جفت وي}} \\
g((x-1)/2) & \text{\RL{که \LR{$x$} طاق وي۔}}
\end{cases}
\]
بيا~\LR{$k$}، \LR{$A \cup B$} شمېري؛ نظر دا دے چې~\LR{$k$} په نوبت د~\LR{$f$} او~\LR{$g$}
شمېرنې کاروي۔ د~\LR{$A \cap B$} شمېرل لږ ستونزمن دي۔ که \LR{$A \cap B$} تش
وي، په څرګند ډول په محاسبوي ډول د شمېر وړ دے۔ که نه، \LR{$c$} د
\LR{$A \cap B$} هر يو غړے وي، او~\LR{$l$} داسې تعريف کړئ:
\[
l(x) = \begin{cases}
f((x)_0) & \text{\RL{که \LR{$f((x)_0) = g((x)_1)$} وي}} \\
c & \text{\RL{که داسې نۀ وي۔}}
\end{cases}
\]
په محاسبوي ژبه کښې~\LR{$l$} د~\LR{$f$} او~\LR{$g$} د شمېرنو د غړو پر جوړو تېرېږي
او هره برابره جوړه چې ومومي، وت ئې راګرځوي؛ که جوړه برابره نۀ وي،
ثابت غړے~\LR{$c$} راګرځوي۔
\textbf{د ژباړې د نسخې سمون (OLCMP-027):} سرچينې دې وروستي بديل ته
``د ثابت قيمت په راګرځولو تم کېدل'' ويلي و، خو تابع په هر ننوت يو
قيمت راګرځوي او محاسبه ئې بې پايلې نۀ پاتې کېږي۔ دلته دا د سم ثابت
شاتګ په توګه بيان شوے؛ په همدې برخه کښې دوه څرګندې انګرېزي ليکدوديزې
تېروتنې هم په طبيعي پښتو کښې نۀ دي نقل شوې۔ فورمولونه نۀ دي بدل شوي۔
\olpseganchor{OLP-0241-B011}{end}

\olpseganchor{OLP-0241-B012}{start}
د وروستي ثبوت دپاره فرض کړئ~\LR{$A$} د جزوي تابع~\LR{$m(x)$} \emph{د تعريف
ساحه} او~\LR{$B$} د جزوي تابع~\LR{$n(x)$} د تعريف ساحه ده۔ بيا \LR{$A \cap B$} د
جزوي تابع \LR{$m(x) + n(x)$} د تعريف ساحه ده۔
\olpseganchor{OLP-0241-B012}{end}

\olpseganchor{OLP-0241-B013}{start}
\begin{explain}
په محاسبوي ژبه کښې، که~\LR{$A$} د هغو قيمتونو سټ وي چې~\LR{$m$} پرې درېږي او
\LR{$B$} د هغو قيمتونو سټ وي چې~\LR{$n$} پرې درېږي، نو \LR{$A \cap B$} د هغو
قيمتونو سټ دے چې دواړه کړنلارې پرې درېږي۔
\end{explain}
\olpseganchor{OLP-0241-B013}{end}

\olpseganchor{OLP-0241-B014}{start}
د درېدنې د قيمتونو د سټ په توګه د~\LR{$A \cup B$} څرګندول لږ ستونزمن دي،
ځکه~\LR{$m$} او~\LR{$n$} بايد په موازي ډول تقليد شي۔ \LR{$d$} د~\LR{$m$} يو شاخص او~\LR{$e$}
د~\LR{$n$} يو شاخص وي؛ په بله وينا، \LR{$m = \cfind{d}$} او
\LR{$n = \cfind{e}$}۔ بيا \LR{$A \cup B$} د دې تابع د تعريف ساحه ده:
\[
p(x) = \umin{y}{(T(d,x,y) \lor T(e,x,y))}.
\]
\begin{explain}
په محاسبوي ژبه کښې، پر ننوت~\LR{$x$} تابع~\LR{$p$} يا د~\LR{$m$} يا د~\LR{$n$} درېدلې
محاسبه لټوي، او که يوه ئې ومومي درېږي۔
\end{explain}
\end{proof}
\olpseganchor{OLP-0241-B014}{end}

\olpseganchor{OLP-0242-B004}{start}
\olfileid{cmp}{thy}{cmp}
\olsection{په محاسبوي ډول د شمېر وړ سټونه د متمم لاندې تړلي نۀ دي}
\olpseganchor{OLP-0242-B004}{end}

\olpseganchor{OLP-0242-B005}{start}
فرض کړئ \LR{$A$} په محاسبوي ډول د شمېر وړ دے۔ ايا د~\LR{$A$} متمم،
\LR{$\Complement{A} = \Nat \setminus A$}، همېشه په محاسبوي ډول د شمېر وړ
هم دے؟ لاندې قضيه او نتيجه ښيي چې ځواب ``نه'' دے۔
\olpseganchor{OLP-0242-B005}{end}

\olpseganchor{OLP-0242-B006}{start}
\begin{thm}
\ollabel{thm:ce-comp}
فرض کړئ \LR{$A$} د طبيعي عددونو هر يو سټ دے۔ بيا \LR{$A$} هله او يوازې هله
محاسبه کېدونکے دے چې \LR{$A$} او~\LR{$\Complement{A}$} دواړه په محاسبوي ډول
د شمېر وړ وي۔
\end{thm}
\olpseganchor{OLP-0242-B006}{end}

\olpseganchor{OLP-0242-B007}{start}
\begin{proof}
مخکښ لورے اسان دے: که~\LR{$A$} محاسبه کېدونکے وي، نو~\LR{$\Complement{A}$}
هم محاسبه کېدونکے دے (\LR{$\Char{A} = 1 \tsub
\Char{\Complement{A}}$})، او ځکه دواړه په محاسبوي ډول د شمېر وړ دي۔
\olpseganchor{OLP-0242-B007}{end}

\olpseganchor{OLP-0242-B008}{start}
په بل لوري کښې فرض کړئ~\LR{$A$} او~\LR{$\Complement{A}$} دواړه په محاسبوي ډول
د شمېر وړ دي۔ \LR{$A$} د~\LR{$\cfind{d}$} د تعريف ساحه، او
\LR{$\Complement{A}$} د~\LR{$\cfind{e}$} د تعريف ساحه وي۔ تابع~\LR{$h$} داسې تعريف
کړئ:
\[
h(x) = \umin{s}{(T(d,x,s) \lor T(e,x,s))}.
\]
په بله وينا، پر ننوت~\LR{$x$} تابع~\LR{$h$} يا د~\LR{$\cfind{d}$} يا
د~\LR{$\cfind{e}$} يوه درېدلې محاسبه لټوي۔ اوس که \LR{$x \in A$} وي، په
لومړي حالت کښې به بريالۍ شي، او که \LR{$x \in \Complement{A}$} وي، په
دويم حالت کښې به بريالۍ شي۔ نو~\LR{$h$} هرځاے تعريف شوې محاسبه کېدونکې
تابع ده۔ خو اوس د هر~\LR{$x$} دپاره \LR{$x \in A$} هله او يوازې هله چې
\LR{$T(d, x, h(x))$} وي، يعنې که~\LR{$\cfind{d}$} پر همدې ننوت تعريف شوې وي۔
څرنګه چې \LR{$T(d, x, h(x))$} محاسبه کېدونکې اړيکه ده، نو~\LR{$A$}
محاسبه کېدونکے دے۔
\textbf{د ژباړې د نسخې سمون (OLCMP-028):} سرچينې په وروستۍ
دوه‌اړخيزه نتيجه کښې د متمم شاخص کارولے و، حال دا چې سمدستي مخکښې
ئې اصلي سټ د لومړي شاخص د تعريف ساحه ټاکلې ده۔ دلته هماغه لومړے شاخص
کارول شوے، څو ځانګړونکې ازموينه د اصلي سټ غړيتوب وښيي۔
\end{proof}
\olpseganchor{OLP-0242-B008}{end}

\olpseganchor{OLP-0242-B009}{start}
\begin{explain}
په غير رسمي محاسبوي ژبه کښې خبره لا اسانه پوهېدل کېږي: د~\LR{$A$} د
پرېکړې دپاره پر ننوت~\LR{$x$} د~\LR{$\cfind{d}$} او~\LR{$\cfind{e}$} درېدلې
محاسبې ولټوئ۔ يوه به خامخا ودرېږي؛ که~\LR{$\cfind{d}$} ودرېږي، نو~\LR{$x$}
په~\LR{$A$} کښې دے، او که نه، \LR{$x$} په~\LR{$\Complement{A}$} کښې دے۔
\textbf{د ژباړې د نسخې سمون (OLCMP-029):} سرچينې په دې غير رسمي
لنډيز کښې د لومړي ثبوت له شاخصونو سره نۀ سمون خوړونکے نوی، ناتعريف
شوے توری راوړے و، او بيا ئې د متمم شاخص درېدنه د اصلي سټ له غړيتوب
سره تړله۔ دلته د ثبوت هماغه دوه ټاکلي شاخصونه او سمه غړيتوبي څانګه
ساتل شوې ده۔
\end{explain}
\olpseganchor{OLP-0242-B009}{end}

\olpseganchor{OLP-0242-B010}{start}
\begin{cor}
\ollabel{cor:comp-k}
\LR{$\Complement{K_0}$} په محاسبوي ډول د شمېر وړ نۀ دے۔
\end{cor}
\olpseganchor{OLP-0242-B010}{end}

\olpseganchor{OLP-0242-B011}{start}
\begin{proof}
موږ پوهېږو چې~\LR{$K_0$} په محاسبوي ډول د شمېر وړ دے، خو محاسبه کېدونکے
نۀ دے۔ که~\LR{$\Complement{K_0}$} په محاسبوي ډول د شمېر وړ وای، نو د
\olref{thm:ce-comp} له مخې به~\LR{$K_0$} محاسبه کېدونکے و، چې له
\olref[ncp]{thm:K-0} سره تناقض دے۔
\end{proof}
\olpseganchor{OLP-0242-B011}{end}

\olpseganchor{OLP-0243-B004}{start}
\olfileid{cmp}{thy}{red}
\olsection{راکمېدنه}
\olpseganchor{OLP-0243-B004}{end}

\olpseganchor{OLP-0243-B005}{start}
\begin{explain}
اوس پوهېږو چې لږ تر لږه يو سټ، \LR{$K_0$}، شته چې په محاسبوي ډول د شمېر
وړ دے خو محاسبه کېدونکے نۀ دے۔ بايد څرګنده وي چې نور هم شته۔ د
راکمېدنې طريقه يوه پياوړې لار راکوي چې د نورو سټونو دا خاصيتونه، بې
له دې چې هر ځل بنسټيزو اصولو ته ستانۀ شو، ثابت کړو۔
\olpseganchor{OLP-0243-B005}{end}

\olpseganchor{OLP-0243-B006}{start}
په عمومي ډول، له سټ~\LR{$A$} څخه سټ~\LR{$B$} ته يوه ``راکمونه'' داسې طريقه ده
چې د دې پوښتنې ځوابونه چې غړي په~\LR{$B$} کښې دي که نه، د دې
پوښتنې په ځوابونو بدلوي چې غړي په~\LR{$A$} کښې دي که نه۔ موږ به
د ``ډېر-پر-يو راکمېدنې'' په نوم مفهوم ته پام وکړو، خو د بېلابېلو
خاصيتونو لرونکي د راکمېدنې ډېر نور مفهومونه هم شته۔ د راکمېدنې
مفهومونه د محاسبوي پېچلتيا په څېړنه کښې هم مرکزي دي، هلته د اغېزمنتيا
مسئلې هم بايد په پام کښې ونيول شي۔ د بېلګې په توګه، يو سټ هغه وخت
``اېن‌پي-بشپړ'' بلل کېږي چې په اېن‌پي کښې وي او د اېن‌پي هره مسئله
ورته راکمول کېدے شي؛ دلته د لاندې مفهوم ته ورته راکمونه کارېږي، خو
دا زيات شرط لري چې راکمونه په څوحدي وخت کښې محاسبه شي۔
\olpseganchor{OLP-0243-B006}{end}

\olpseganchor{OLP-0243-B007}{start}
موږ د راکمونې مفهوم لا له مخکښې په ضمني ډول کارولے دے۔ سټ~\LR{$K$} داسې
تعريف کړئ:
\[
K = \Setabs{x}{\cfind{x}(x) \fdefined},
\]
يعنې \LR{$K = \Setabs{x}{x \in W_x}$}۔ زموږ هغه ثبوت چې د درېدنې مسئله
نۀ حل کېږي (\olref[hlt]{thm:halting-problem})، تر ټولو نېغ دا ښيي چې
\LR{$K$} محاسبه کېدونکے نۀ دے۔ راياد کړئ چې~\LR{$K_0$} دا سټ دے:
\[
K_0 = \Setabs{\tuple{e, x}}{\cfind{e}(x) \fdefined },
\]
يعنې \LR{$K_0 = \Setabs{\tuple{e,x}}{x \in W_e}$}۔ د~\LR{$K$} د نۀ محاسبه
کېدنې هر ثبوت د~\LR{$K_0$} نۀ محاسبه کېدنې ته غځول اسان دي: که~\LR{$K_0$}
محاسبه کېدونکے وای، نو دا مو پرېکړه کولے شوه چې غړے~\LR{$x$} په
\LR{$K$} کښې دے که نه، يوازې په دې پوښتنې سره چې جوړه~\LR{$\tuple{x, x}$} په
\LR{$K_0$} کښې ده که نه۔ هغه تابع~\LR{$f$} چې~\LR{$x$} په~\LR{$\tuple{x, x}$} بدلوي،
له~\LR{$K$} څخه~\LR{$K_0$} ته د \emph{راکمونې} يوه بېلګه ده۔
\textbf{د ژباړې د نسخې سمون (OLCMP-030):} سرچينې د دې بحث په پيل
کښې د ``راکمونې مفهوم'' نوم دوه ځله پرله‌پسې ليکلے او د درېدنې مسئله
ئې د ``نۀ حل کېدونکې'' پر ځاے د يوې ليکدوديزې تېروتنې له امله بله
کلمه بللې وه۔ په پښتو کښې دواړه څرګندې تحريرې تېروتنې نۀ دي نقل شوې۔
\textbf{د ژباړې د نسخې سمون (OLCMP-031):} د درېدنې د جوړو دويم
سټ-جوړونکي بيان په سرچينه کښې د جوړې دوه مختصات اړولي وو، خو شرط ئې
بيا هم دويم عدد د لومړي عدد د شاخص په توګه کاراوه۔ دلته جوړه د لومړي،
کره تعريف په ترتيب ساتل شوې ده؛ ورپسې قطري راکمونه هم همدغه ترتيب
کاروي۔
\end{explain}
\olpseganchor{OLP-0243-B007}{end}

\olpseganchor{OLP-0243-B008}{start}
\begin{defn}
فرض کړئ \LR{$A$} او~\LR{$B$} د طبيعي عددونو سټونه دي۔ محاسبه کېدونکې تابع
~\LR{$f\colon \Nat \to \Nat$} له~\LR{$A$} څخه~\LR{$B$} ته \emph{ډېر-پر-يو راکمونه}
ده، هله او يوازې هله چې د هر طبيعي عدد~\LR{$x$} دپاره
\[
x \in A \quad \text{\RL{هله او يوازې هله چې}} \quad f(x) \in B.
\]
که داسې راکمونه~\LR{$f$} شته، وايو~\LR{$A$}، \LR{$B$} ته \emph{ډېر-پر-يو راکمېدونکے}
دے، چې \LR{$A \leq_m B$} ليکل کېږي۔ که~\LR{$A$}، \LR{$B$} ته ډېر-پر-يو راکمېدونکے
وي او برعکس هم، نو~\LR{$A$} او~\LR{$B$} \emph{ډېر-پر-يو معادل} بلل کېږي، چې
\LR{$A \equiv_m B$} ليکل کېږي۔
\end{defn}
\olpseganchor{OLP-0243-B008}{end}

\olpseganchor{OLP-0243-B009}{start}
که د پورته تعريف تابع~\LR{$f$} اتفاقاً يو په يو وي، نو~\LR{$A$}، \LR{$B$} ته
\emph{يو-پر-يو راکمېدونکے} بلل کېږي۔ لاندې ډېرې راکمونې دا پياوړے
شرط پوره کوي، خو موږ به دا حقيقت نۀ کاروو۔
\olpseganchor{OLP-0243-B009}{end}

\olpseganchor{OLP-0243-B010}{start}
\begin{digress}
دا رښتيا ده، خو هېڅ څرګنده نۀ ده، چې يو-پر-يو راکمېدنه په رښتيا د
ډېر-پر-يو راکمېدنې په پرتله پياوړے شرط دے۔ په بله وينا، داسې نامتناهي
سټونه~\LR{$A$} او~\LR{$B$} شته چې~\LR{$A$}، \LR{$B$} ته ډېر-پر-يو راکمېدونکے وي خو
~\LR{$B$} ته يو-پر-يو راکمېدونکے نۀ وي۔
\end{digress}
\olpseganchor{OLP-0243-B010}{end}

\olpseganchor{OLP-0244-B004}{start}
\olfileid{cmp}{thy}{ppr}
\olsection{د راکمېدنې خاصيتونه}
\olpseganchor{OLP-0244-B004}{end}

\olpseganchor{OLP-0244-B005}{start}
که~\LR{$A$}، \LR{$B$} ته راکمېږي، \LR{$A \leq_m B$} ليکو۔ دا نښه وړانديز کوي چې
که \LR{$A \leq_m B$} وي، نو~\LR{$A$} له~\LR{$B$} څخه ``زيات ستونزمن نۀ دے''، او
\LR{$B$} د~\LR{$A$} ``هومره يا ترې زيات ستونزمن دے''؛ او~\LR{$\le_m$} د عددونو د
معمولي~\LR{$\le$} په شان سټونه (يا د پرېکړې مسئلې) د پېچلتيا له مخې
مرتبوي۔ لاندې دوه قضيې د دې شهود ملاتړ کوي۔ لومړۍ وايي چې~\LR{$\le_m$}
انتقالي ده۔
\textbf{د ژباړې د نسخې سمون (OLCMP-032):} د سرچينې د دې وروستۍ
جملې په دوو کلمو کښې څرګندې ليکدوديزې تېروتنې وې؛ په طبيعي پښتو کښې
سمه موخه شوې جمله ژباړل شوې ده۔
\olpseganchor{OLP-0244-B005}{end}

\olpseganchor{OLP-0244-B006}{start}
\begin{prop}
\ollabel{prop:trans-red}
که \LR{$A \leq_m B$} او \LR{$B \leq_m C$} وي، نو \LR{$A \leq_m C$}۔
\end{prop}
\olpseganchor{OLP-0244-B006}{end}

\olpseganchor{OLP-0244-B007}{start}
\begin{proof}
له~\LR{$A$} څخه~\LR{$B$} ته د راکمونې او له~\LR{$B$} څخه~\LR{$C$} ته د راکمونې ترکيب
له~\LR{$A$} څخه~\LR{$C$} ته راکمونه ورکوي۔
\end{proof}
\olpseganchor{OLP-0244-B007}{end}

\olpseganchor{OLP-0244-B008}{start}
\begin{prob}
\olref{prop:trans-red} په دې ښودلو ثابت کړئ چې که~\LR{$f$} او~\LR{$g$} په
ترتيب سره له~\LR{$A$} څخه~\LR{$B$} او له~\LR{$B$} څخه~\LR{$C$} ته ډېر-پر-يو راکمونې
وي، نو~\LR{$\comp{f}{g}$} له~\LR{$A$} څخه~\LR{$C$} ته ډېر-پر-يو راکمونه ده۔
\end{prob}
\olpseganchor{OLP-0244-B008}{end}

\olpseganchor{OLP-0244-B009}{start}
\begin{prop}
\ollabel{prop:reduce}
فرض کړئ \LR{$A$} او~\LR{$B$} هر سټونه دي، او \LR{$A \leq_m B$}۔
\begin{enumerate}
\item که~\LR{$B$} په محاسبوي ډول د شمېر وړ وي، نو~\LR{$A$} هم داسې دے۔
\item که~\LR{$B$} محاسبه کېدونکے وي، نو~\LR{$A$} هم داسې دے۔
\end{enumerate}
\end{prop}
\olpseganchor{OLP-0244-B009}{end}

\olpseganchor{OLP-0244-B010}{start}
\begin{proof}
\LR{$f$} له~\LR{$A$} څخه~\LR{$B$} ته ډېر-پر-يو راکمونه وي۔ د لومړۍ دعوې دپاره بس
دا وازمويئ چې که~\LR{$B$} د جزوي تابع~\LR{$g$} د تعريف ساحه وي، نو~\LR{$A$} د
~\LR{$\comp{f}{g}$} د تعريف ساحه ده:
\begin{align*}
x \in A & \text{\RL{ هله او يوازې هله چې }} f(x) \in B \\
& \text{\RL{ هله او يوازې هله چې }}  g(f(x)) \fdefined.
\end{align*}
\olpseganchor{OLP-0244-B010}{end}

\olpseganchor{OLP-0244-B011}{start}
د دويمې دعوې دپاره راياد کړئ چې که~\LR{$B$} محاسبه کېدونکے وي، نو~\LR{$B$}
او~\LR{$\Complement{B}$} په محاسبوي ډول د شمېر وړ دي
(\olref[cmp]{thm:ce-comp})۔ دا ازمويل ستونزمن نۀ دي چې~\LR{$f$} له
\LR{$\Complement{A}$} څخه~\LR{$\Complement{B}$} ته هم ډېر-پر-يو راکمونه ده؛
نو د دې ثبوت د لومړۍ برخې له مخې~\LR{$A$} او~\LR{$\Complement{A}$} دواړه
په محاسبوي ډول د شمېر وړ دي۔ ځکه~\LR{$A$} هم د
\olref[cmp]{thm:ce-comp} له مخې محاسبه کېدونکے دے۔ (په بديل ډول،
کتلے شئ چې \LR{$\Char{A} = \comp{f}{\Char{B}}$}؛ نو که~\LR{$\Char{B}$} محاسبه
کېدونکې وي، \LR{$\Char{A}$} هم محاسبه کېدونکې ده۔)
\end{proof}
\olpseganchor{OLP-0244-B011}{end}

\olpseganchor{OLP-0244-B012}{start}
\begin{prob}
فرض کړئ~\LR{$f$} له~\LR{$A$} څخه~\LR{$B$} ته ډېر-پر-يو راکمونه ده۔ وښيئ چې~\LR{$f$} له
\LR{$\Complement{A}$} څخه~\LR{$\Complement{B}$} ته هم ډېر-پر-يو راکمونه ده۔
\end{prob}
\olpseganchor{OLP-0244-B012}{end}

\olpseganchor{OLP-0244-B013}{start}
\begin{prob}
وښيئ چې که \LR{$f\colon \Nat \to \Nat$} ډېر-پر-يو راکمونه وي، نو
\LR{$\Char{A} = \comp{f}{\Char{B}}$}۔
\textbf{د ژباړې د نسخې سمون (OLCMP-033):} سرچينې دلته د راکمونې
تابع يوازې د لومړي سټ له غړو څخه دويم سټ ته ټايپ کړې وه، حال دا چې
تعريف ئې د ټولو طبيعي عددونو دپاره تابع غواړي او د سټ غړيتوب د
دوه‌اړخيز شرط په وسيله ساتي۔ دلته د تعريف سمه طبيعي-عددونو تعريف
ساحه او هم‌ساحه کارول شوې ده۔
\end{prob}
\olpseganchor{OLP-0244-B013}{end}

\olpseganchor{OLP-0244-B014}{start}
\begin{digress}
د \emph{ټيورينګ راکمېدنې} په نوم د راکمېدنې يو ډېر عمومي مفهوم په
نورو سياقونو کښې ګټور دے، په ځانګړي ډول د ناپرېکړتيا د پايلو په
ثبوت کښې۔ پام وکړئ چې د \olref[cmp]{cor:comp-k} له مخې د~\LR{$K_0$} متمم
~\LR{$K_0$} ته نۀ راکمېږي، ځکه په محاسبوي ډول د شمېر وړ نۀ دے۔ خو په
شهودي ډول، که د~\LR{$K_0$} په اړه د پوښتنو ځوابونه مو پېژندل، د هغې د
متمم په اړه د پوښتنو ځوابونه به مو هم پېژندل۔ سټ~\LR{$A$} هغه وخت~\LR{$B$} ته
د ټيورينګ له مخې راکمېدونکے بلل کېږي چې په~\LR{$A$} کښې د پوښتنو ځوابونه
د داسې محاسبه کېدونکې کړنلارې په وسيله وټاکل شي چې د~\LR{$B$} په اړه
پوښتنې کولے شي۔ دا له ډېر-پر-يو راکمېدنې زيات ازاد دے، ځکه هلته
(1)~يوازې د~\LR{$B$} په اړه يوه پوښتنه کولے شئ، او (2)~د ``هو'' ځواب بايد
د~\LR{$A$} د پوښتنې په ``هو'' ځواب، او همداسې ``نه'' په ``نه''، بدل شي۔
بيا هم که~\LR{$A$}، \LR{$B$} ته د ټيورينګ له مخې راکمېدونکے وي او~\LR{$B$} محاسبه
کېدونکے وي، نو~\LR{$A$} هم محاسبه کېدونکے دے (که څه هم، لکه وليدل مو،
ورته خبره د محاسبوي شمېر وړتيا دپاره نۀ پوره کېږي)۔
\olpseganchor{OLP-0244-B014}{end}

\olpseganchor{OLP-0244-B015}{start}
بايد د راکمېدنې د هغو بېلابېلو مفهومونو په اړه فکر وکړئ چې بحث مو
پرې وکړ، او توپيرونه ئې درک کړئ۔ خو په دې څپرکي کښې به يوازې له
ډېر-پر-يو راکمېدنې سره کار کوو۔ په ضمني ډول، په تېر پراګراف کښې د
راکمېدنې دواړه ډولونه په محاسبوي پېچلتيا کښې ورته بڼې لري، خو دا
زيات شرط ورسره وي چې ټيورينګ ماشينونه په څوحدي وخت کښې وچلېږي: د
ډېر-پر-يو راکمېدنې پېچلتيايي بڼه \emph{کارپ راکمېدنه}، او د ټيورينګ
راکمېدنې پېچلتيايي بڼه \emph{کوک راکمېدنه} بلل کېږي۔
\end{digress}
\olpseganchor{OLP-0244-B015}{end}

\olpseganchor{OLP-0245-B004}{start}
\olfileid{cmp}{thy}{cce}
\olsection{بشپړ په محاسبوي ډول د شمېر وړ سټونه}
\olpseganchor{OLP-0245-B004}{end}

\olpseganchor{OLP-0245-B005}{start}
\begin{defn}
سټ~\LR{$A$} (د ډېر-پر-يو راکمېدنې لاندې) يو
\emph{بشپړ په محاسبوي ډول د شمېر وړ سټ} دے، که
\begin{enumerate}
\item \LR{$A$} په محاسبوي ډول د شمېر وړ وي، او
\item د هر بل په محاسبوي ډول د شمېر وړ سټ~\LR{$B$} دپاره \LR{$B \leq_m A$} وي۔
\end{enumerate}
\end{defn}
\olpseganchor{OLP-0245-B005}{end}

\olpseganchor{OLP-0245-B006}{start}
په بله وينا، بشپړ په محاسبوي ډول د شمېر وړ سټونه د ټولو ممکنه په
محاسبوي ډول د شمېر وړ سټونو ``تر ټولو ستونزمن'' دي۔ دوي اجازه راکوي
چې د \emph{هر} په محاسبوي ډول د شمېر وړ سټ په اړه پوښتنو ته ځواب
ووايو۔
\olpseganchor{OLP-0245-B006}{end}

\olpseganchor{OLP-0245-B007}{start}
\begin{thm}
\LR{$K$}، \LR{$K_0$} او~\LR{$K_1$} ټول بشپړ په محاسبوي ډول د شمېر وړ سټونه دي۔
\end{thm}
\olpseganchor{OLP-0245-B007}{end}

\olpseganchor{OLP-0245-B008}{start}
\begin{proof}
د دې د ليدلو دپاره چې~\LR{$K_0$} بشپړ دے، \LR{$B$} هر په محاسبوي ډول د شمېر
وړ سټ وي۔ نو د کوم شاخص~\LR{$e$} دپاره
\[
B = W_e = \Setabs{x}{\cfind{e}(x) \fdefined}.
\]
\LR{$f$} هغه تابع وي چې \LR{$f(x) = \tuple{e, x}$}۔ بيا د هر طبيعي عدد~\LR{$x$}
دپاره \LR{$x \in B$} هله او يوازې هله چې \LR{$f(x) \in K_0$} وي۔ په بله
وينا، \LR{$f$}، \LR{$B$}، \LR{$K_0$} ته راکموي۔
\olpseganchor{OLP-0245-B008}{end}

\olpseganchor{OLP-0245-B009}{start}
د دې د ليدلو دپاره چې~\LR{$K_1$} بشپړ دے، پام وکړئ چې د
\olref[k1]{prop:k1} په ثبوت کښې مو~\LR{$K_0$} ورته راکم کړے و۔ نو د
\olref[ppr]{prop:trans-red} له مخې هر په محاسبوي ډول د شمېر وړ سټ
~\LR{$K_1$} ته هم راکمول کېدے شي۔
\olpseganchor{OLP-0245-B009}{end}

\olpseganchor{OLP-0245-B010}{start}
\LR{$K_0$} هم په ډېره ورته طريقه~\LR{$K$} ته راکمول کېدے شي۔
\textbf{د ژباړې د نسخې سمون (OLCMP-034):} د سرچينې وروستۍ جمله د
راکمونې لورے سرچپه ليکلے و: له قطري سټ څخه د جوړو سټ ته راکمونه د
لومړي سټ بشپړتيا نۀ ثابتوي۔ دلته د جوړو بشپړ سټ قطري سټ ته راکمېږي،
څو د انتقال له مخې هر c.e. سټ هم قطري سټ ته راکم شي۔
\end{proof}
\olpseganchor{OLP-0245-B010}{end}

\olpseganchor{OLP-0245-B011}{start}
\begin{prob}
له~\LR{$K$} څخه~\LR{$K_0$} ته يوه راکمونه ورکړئ۔
\end{prob}
\olpseganchor{OLP-0245-B011}{end}

\olpseganchor{OLP-0245-B012}{start}
\begin{digress}
نو څرګنده شوه چې تر اوسه مو د په محاسبوي ډول د شمېر وړ سټونو هره
کتلې بېلګه يا محاسبه کېدونکې ده، يا بشپړه۔ دا بايد عجيبه وبرېښي!{}
ايا داسې په محاسبوي ډول د شمېر وړ سټونه شته چې نۀ محاسبه کېدونکي او
نۀ بشپړ وي؟ ځواب هو دے، خو دا تر هغه وخته نۀ وه ثابته شوې چې د
۱۹۵۰مو کلونو په منځ کښې فريډبرګ او موچنيک په خپلواکه توګه ثابته کړه۔
\end{digress}
\olpseganchor{OLP-0245-B012}{end}

\olpseganchor{OLP-0246-B004}{start}
\olfileid{cmp}{thy}{k1}
\olsection{د راکمېدنې يوه بېلګه}
\olpseganchor{OLP-0246-B004}{end}

\olpseganchor{OLP-0246-B005}{start}
راځئ د \olref[ppr]{prop:reduce} يوه کارونه وګورو۔
\olpseganchor{OLP-0246-B005}{end}

\olpseganchor{OLP-0246-B006}{start}
\begin{prop}
\ollabel{prop:k1}
فرض کړئ
\[
K_1 = \Setabs{e}{\cfind{e}(0) \fdefined}.
\]
بيا \LR{$K_1$} په محاسبوي ډول د شمېر وړ دے، خو محاسبه کېدونکے نۀ دے۔
\end{prop}
\olpseganchor{OLP-0246-B006}{end}

\olpseganchor{OLP-0246-B007}{start}
\begin{proof}
څرنګه چې \LR{$K_1 = \Setabs{e}{\lexists[s][T(e,0,s)]}$}، نو د
\olref[eqc]{thm:exists-char} له مخې~\LR{$K_1$} په محاسبوي ډول د شمېر وړ
دے۔
\olpseganchor{OLP-0246-B007}{end}

\olpseganchor{OLP-0246-B008}{start}
د دې د ښودلو دپاره چې~\LR{$K_1$} محاسبه کېدونکے نۀ دے، راځئ وښيو چې
\LR{$K_0$} ورته راکمېدونکے دے۔
\olpseganchor{OLP-0246-B008}{end}

\olpseganchor{OLP-0246-B009}{start}
\begin{explain}
دا لږ ستونزمن دے، ځکه د~\LR{$K_1$} په کارولو سره يوازې د هغو محاسبو په
اړه پوښتنې کولے شو چې له ځانګړي ننوت، \LR{$0$}، څخه پيلېږي۔ فرض کړئ يو
هوښيار ملګرے لرئ چې د دې ډول پوښتنو ځوابونه درکولے شي (داسې ملګري د
``اوراکلونو'' په نوم پېژندل کېږي)۔ بيا فرض کړئ څوک درته راشي او
وپوښتي چې~\LR{$\tuple{e, x}$} په~\LR{$K_0$} کښې دے که نه، يعنې ماشين~\LR{$e$} پر
ننوت~\LR{$x$} درېږي که نه۔ يو کار دا کولے شئ چې بل ماشين~\LR{$e_x$} جوړ کړئ؛
دا ماشين د \emph{هر} ننوت دپاره هغه ننوت له پامه غورځوي او پر ځاے ئې
~\LR{$e$} پر ننوت~\LR{$x$} چلوي۔ نو څرګنده ده چې د ماشين~\LR{$e$} پر ننوت~\LR{$x$}
درېدل له دې پوښتنې سره معادل دي چې ماشين~\LR{$e_x$} پر ننوت~\LR{$0$} (يا هر
بل ننوت) درېږي که نه۔ بيا له خپل ملګري وپوښتئ چې دا نوی ماشين~\LR{$e_x$}
پر ننوت~\LR{$0$} درېږي که نه؛ بدلې شوې پوښتنې ته د ملګري ځواب د اصلي
پوښتنې ځواب درکوي۔ دا له~\LR{$K_0$} څخه~\LR{$K_1$} ته موخه شوې راکمونه راکوي۔
\end{explain}
\olpseganchor{OLP-0246-B009}{end}

\olpseganchor{OLP-0246-B010}{start}
د نړيوالې جزوي محاسبه کېدونکې تابع په کارولو سره، \LR{$f$} هغه
درې‌ځاييزه تابع وي چې داسې تعريفېږي:
\[
f(x,y,z) \simeq \cfind{x}(y).
\]
پام وکړئ چې~\LR{$f$} خپل درېيم ننوت بيخي له پامه غورځوي۔ داسې شاخص~\LR{$e$}
وټاکئ چې \LR{$f = \cfind{e}[3]$} وي؛ نو لرو:
\[
\cfind{e}[3](x,y,z) \simeq \cfind{x}(y).
\]
د \LR{$s$}-\LR{$m$}-\LR{$n$} قضيې له مخې داسې تابع~\LR{$s(e,x,y)$} شته چې د هر~\LR{$z$}
دپاره
\begin{align*}
\cfind{s(e,x,y)}(z) & \simeq \cfind{e}[3](x,y,z) \\
& \simeq \cfind{x}(y).
\end{align*}
\olpseganchor{OLP-0246-B010}{end}

\olpseganchor{OLP-0246-B011}{start}
\begin{explain}
د پورته غير رسمي استدلال په ژبه، \LR{$s(e,x,y)$} د هغه ماشين شاخص دے چې
د هر ننوت~\LR{$z$} دپاره هغه ننوت له پامه غورځوي او~\LR{$\cfind{x}(y)$}
محاسبه کوي۔
\end{explain}
\olpseganchor{OLP-0246-B011}{end}

\olpseganchor{OLP-0246-B012}{start}
په ځانګړي ډول، لرو:
\[
\cfind{s(e,x,y)}(0) \fdefined \quad \text{\RL{هله او يوازې هله چې}} \quad
\cfind{x}(y) \fdefined.
\]
په بله وينا، \LR{$\tuple{x, y} \in K_0$} هله او يوازې هله چې
\LR{$s(e,x,y) \in K_1$} وي۔ نو هغه تابع~\LR{$g$} چې داسې تعريفېږي:
\[
g(w) = s(e,(w)_0,(w)_1)
\]
له~\LR{$K_0$} څخه~\LR{$K_1$} ته راکمونه ده۔
\end{proof}
\olpseganchor{OLP-0246-B012}{end}

\olpseganchor{OLP-0247-B004}{start}
\olfileid{cmp}{thy}{tot}
\olsection{د هرځاے تعريف کېدو مسئله ناپرېکړه کېدونکې ده}
\olpseganchor{OLP-0247-B004}{end}

\olpseganchor{OLP-0247-B005}{start}
راځئ د~\LR{$s$}-\LR{$m$}-\LR{$n$} قضيې د کارولو يوه بله بېلګه وګورو، څو وښيو چې
يو څيز محاسبه کېدونکے نۀ دے۔ \LR{$\fn{Tot}$} د هرځاے تعريف شويو محاسبه
کېدونکو تابعو د شاخصونو سټ وي، يعنې
\[
\fn{Tot} = \Setabs{x}{\text{\RL{د هر \LR{$y$} دپاره، \LR{$\cfind{x}(y)\fdefined$}}}}.
\]
\olpseganchor{OLP-0247-B005}{end}

\olpseganchor{OLP-0247-B006}{start}
\begin{prop}
\ollabel{prop:total}
\LR{$\fn{Tot}$} محاسبه کېدونکے نۀ دے۔
\end{prop}
\olpseganchor{OLP-0247-B006}{end}

\olpseganchor{OLP-0247-B007}{start}
\begin{proof}
د دې د ليدلو دپاره چې~\LR{$\fn{Tot}$} محاسبه کېدونکے نۀ دے، دا ښودل بس
دي چې~\LR{$K$} ورته راکمېدونکے دے۔ \LR{$h(x,y)$} داسې تعريف کړئ:
\[
h(x,y) \simeq
\begin{cases}
0 & \text{\RL{که \LR{$x \in K$} وي}} \\
\fundefined & \text{\RL{که داسې نۀ وي۔}}
\end{cases}
\]
پام وکړئ چې~\LR{$h(x,y)$} بيخي پر~\LR{$y$} تکيه نۀ لري۔ دا ليدل بايد ستونزمن
نۀ وي چې~\LR{$h$} جزوي محاسبه کېدونکې ده: پر ننوتونو~\LR{$x, y$} دا~\LR{$h$} په
دې ډول محاسبه کوو چې لومړے تابع~\LR{$\cfind{x}$} پر ننوت~\LR{$x$} تقليدوو؛
که دا محاسبه ودرېږي، \LR{$h(x,y)$}، \LR{$0$} راګرځوي او درېږي۔ نو \LR{$h(x,y)$} يوازې
\LR{$\Zero(\umin{s}{T(x,x,s)})$} ده، چېرته چې~\LR{$\Zero$} ثابت صفر تابع ده۔
\textbf{د ژباړې د نسخې سمون (OLCMP-035):} سرچينې د دې محاسبې د
بيان په پيل کښې يو ناسمه زياته کلمه لرله؛ طبيعي پښتو جمله د هغې له
نقلولو پرته ژباړل شوې ده۔
\olpseganchor{OLP-0247-B007}{end}

\olpseganchor{OLP-0247-B008}{start}
د~\LR{$s$}-\LR{$m$}-\LR{$n$} قضيې په کارولو داسې بنسټيزه بازګشتي تابع~\LR{$k(x)$} شته
چې د هر~\LR{$x$} او~\LR{$y$} دپاره
\[
\cfind{k(x)}(y) =
\begin{cases}
0 & \text{\RL{که \LR{$x \in K$} وي}} \\
\fundefined & \text{\RL{که داسې نۀ وي۔}}
\end{cases}
\]
نو~\LR{$\cfind{k(x)}$} هله هرځاے تعريف شوې ده چې~\LR{$x \in K$} وي، او که نه
نو تعريف شوې نۀ ده۔ په دې توګه~\LR{$k$} له~\LR{$K$} څخه~\LR{$\fn{Tot}$} ته راکمونه
ده۔
\end{proof}
\olpseganchor{OLP-0247-B008}{end}

\olpseganchor{OLP-0247-B009}{start}
\begin{digress}
په حقيقت کښې~\LR{$\fn{Tot}$} آن په محاسبوي ډول د شمېر وړ هم نۀ دے؛
پېچلتيا ئې د ``حسابي پوړيز نظام'' په لا لوړه برخه کښې پرته ده۔ خو
دلته به د دې پياوړتيا په اړه اندېښنه نۀ کوو۔
\end{digress}
\olpseganchor{OLP-0247-B009}{end}

\olpseganchor{OLP-0248-B004}{start}
\olfileid{cmp}{thy}{rce}
\olsection{د رايس قضيه}
\olpseganchor{OLP-0248-B004}{end}

\olpseganchor{OLP-0248-B005}{start}
که لږ فکر وکړئ، وبه وينئ چې د~\LR{$\fn{Tot}$} ځانګړي جزئيات د
\olref[tot]{prop:total} په ثبوت کښې هېڅ رول نۀ لري۔ موږ~\LR{$h(x,y)$}
داسې جوړ کړے و چې کټ مټ هغه وخت د ثابتې تابع~\LR{$j(y) = 0$} په شان عمل
وکړي چې~\LR{$x$} په~\LR{$K$} کښې وي؛ خو په همدغو حالاتو کښې مو د هرې بلې
جزوي محاسبه کېدونکې تابع په شان هم عمل پرې کولے شو۔ دا مشاهده موږ ته
اجازه راکوي يوه لا عمومي قضيه بيان کړو چې، په لنډ ډول، وايي د محاسبه
کېدونکو تابعو هېڅ غيرابتذالي خاصيت د پرېکړې وړ نۀ دے۔
\olpseganchor{OLP-0248-B005}{end}

\olpseganchor{OLP-0248-B006}{start}
په ياد ولرئ چې~\LR{$\cfind{0}$}، \LR{$\cfind{1}$}، \LR{$\cfind{2}$}،~\dots زموږ د
جزوي محاسبه کېدونکو تابعو معياري شمېرنه ده۔
\olpseganchor{OLP-0248-B006}{end}

\olpseganchor{OLP-0248-B007}{start}
\begin{thm}[د رايس قضيه]
  \LR{$C$} د جزوي محاسبه کېدونکو تابعو هر يو سټ وي، او \LR{$A =
  \Setabs{n}{\cfind{n} \in C}$} وي۔ که~\LR{$A$} محاسبه کېدونکے وي، نو يا
  خو~\LR{$C$} تش دے، يا~\LR{$C$} د ټولو جزوي محاسبه کېدونکو تابعو سټ دے۔
\end{thm}
\olpseganchor{OLP-0248-B007}{end}

\olpseganchor{OLP-0248-B008}{start}
\emph{شاخصي سټ} داسې سټ~\LR{$A$} دے چې که~\LR{$n$} او~\LR{$m$} هغه شاخصونه وي چې
يوه تابع ``محاسبه کوي''، نو يا~\LR{$n$} او~\LR{$m$} دواړه په~\LR{$A$} کښې دي، يا
يو هم نۀ دے۔ دا ليدل ستونزمن نۀ دي چې په قضيه کښې سټ~\LR{$A$} دا خاصيت
لري۔ برعکس، که~\LR{$A$} شاخصي سټ وي او~\LR{$C$} د هغو تابعو سټ وي چې دا
شاخصونه ئې محاسبه کوي، نو \LR{$A = \Setabs{n}{\cfind{n} \in C}$}۔
\olpseganchor{OLP-0248-B008}{end}

\olpseganchor{OLP-0248-B009}{start}
\begin{explain}
په دې اصطلاح کښې د رايس قضيه له دې وينا سره معادله ده چې هېڅ
غيرابتذالي شاخصي سټ د پرېکړې وړ نۀ دے۔ د قضيې د مطلب د پوهېدو دپاره
ګټوره ده چې د \emph{پروګرامونو} (مثلاً ستاسو په خوښه پروګرامي ژبه
کښې) او د هغو د محاسبه کوونکو تابعو تر منځ توپير روښانه کړو۔ د
پروګرامونو (شاخصونو) په اړه، چې نحوي څيزونه دي، يقيناً داسې پوښتنې
شته چې محاسبه کېدونکې دي: ايا دا پروګرام تر~150 زياتې نښې لري؟ ايا
تر~22 زياتې کرښې لري؟ ايا د ``وايل'' جمله لري؟ ايا رشته ``هېلو
ورلډ'' کله هم د ``پرنټ'' جملې د دليل په توګه راځي؟ د رايس قضيه وايي
چې د پروګرام د \emph{چلند} په اړه هېڅ غيرابتذالي پوښتنه محاسبه
کېدونکې نۀ ده۔ په دې کښې دا پوښتنې هم راځي: ايا پروګرام پر ننوت~\LR{$0$}
درېږي؟ ايا کله هم درېږي؟ ايا کله هم جفت عدد راګرځوي؟
\end{explain}
\olpseganchor{OLP-0248-B009}{end}

\olpseganchor{OLP-0248-B010}{start}
\begin{proof}[د رايس د قضيې ثبوت]
فرض کړئ~\LR{$C$} نۀ تش دے او نۀ د ټولو جزوي محاسبه کېدونکو تابعو سټ، او
~\LR{$A$} په~\LR{$C$} کښې د تابعو د شاخصونو سټ وي۔ وبه ښيو چې که~\LR{$A$} محاسبه
کېدونکے وای، د درېدنې مسئله مو حل کولے شوه؛ نو~\LR{$A$} محاسبه کېدونکے
نۀ دے۔
\olpseganchor{OLP-0248-B010}{end}

\olpseganchor{OLP-0248-B011}{start}
بې له دې چې عموميت له لاسه ورکړو، فرض کولے شو چې هغه تابع~\LR{$f$} چې
هېچرته تعريف شوې نۀ ده، په~\LR{$C$} کښې نۀ ده (که نه، په لاندې استدلال
کښې~\LR{$C$} او د هغې متمم سره بدل کړئ)۔ \LR{$g$} په~\LR{$C$} کښې هره تابع وي۔ نظر
دا دے چې که مو~\LR{$A$} پرېکړه کولے شو، د هغو شاخصونو چې~\LR{$f$} محاسبه کوي
او هغو چې~\LR{$g$} محاسبه کوي، توپير مو کولے شو؛ بيا مو همدغه وړتيا د
درېدنې د مسئلې د حل دپاره کارولے شوه۔
\olpseganchor{OLP-0248-B011}{end}

\olpseganchor{OLP-0248-B012}{start}
طريقه ئې دا ده۔ د نړيوالو جزوي محاسبه کېدونکو تابعو په کارولو دا
تابع تعريفولے شو:
\[
h(x,y) \simeq
\begin{cases}
\text{\RL{تعريف شوې نۀ ده}} & \text{\RL{که \LR{$\cfind{x}(x) \fundefined$} وي}} \\
g(y) & \text{\RL{که داسې نۀ وي۔}}
\end{cases}
\]
د~\LR{$h$} د محاسبې دپاره لومړے د~\LR{$\cfind{x}(x)$} د محاسبې هڅه کوو؛ که
دا محاسبه ودرېږي، بيا~\LR{$g(y)$} محاسبه کوو؛ او که \emph{هغه} محاسبه
ودرېږي، وت ئې راګرځوو۔ په لا صوري ډول ليکلے شو:
\[
h(x,y) \simeq \Proj{2}{0}(g(y),\fn{Un}(x,x)).
\]
چېرته چې~\LR{$\Proj{2}{0}(z_0, z_1) = z_0$} هغه~\LR{$2$}-ځاييزه پروجکشن تابع
ده چې~\LR{$0$}-م دليل راګرځوي، او محاسبه کېدونکې ده۔
\olpseganchor{OLP-0248-B012}{end}

\olpseganchor{OLP-0248-B013}{start}
بيا~\LR{$h$} د جزوي محاسبه کېدونکو تابعو ترکيب دے، او ښے اړخ کټ مټ هغه
وخت تعريف شوے او له~\LR{$g(y)$} سره برابر دے چې~\LR{$\fn{Un}(x,x)$} او~\LR{$g(y)$}
دواړه تعريف شوي وي۔
\olpseganchor{OLP-0248-B013}{end}

\olpseganchor{OLP-0248-B014}{start}
پام وکړئ چې د هر ثابت~\LR{$x$} دپاره، که~\LR{$\cfind{x}(x)$} تعريف شوې نۀ
وي، نو~\LR{$h(x,y)$} د هر~\LR{$y$} دپاره تعريف شوې نۀ ده؛ او که~\LR{$\cfind{x}(x)$}
تعريف شوې وي، نو \LR{$h(x,y) \simeq g(y)$}۔ نو د~\LR{$x$} د هر ثابت قيمت
دپاره~\LR{$h(x,y)$} يا کټ مټ د~\LR{$f$}، يا کټ مټ د~\LR{$g$} په شان عمل کوي، او
دا پرېکړه چې~\LR{$\cfind{x}(x)$} تعريف شوې ده که نه، د دې پرېکړې سره
برابره ده چې له دې دوو څخه کوم حالت دے۔ خو دا بيا د دې پرېکړې سره
برابره ده چې \LR{$h_x(y) \simeq h(x,y)$} په~\LR{$C$} کښې ده که نه، او که~\LR{$A$}
محاسبه کېدونکے وای، همدا کار مو کولے شو۔
\olpseganchor{OLP-0248-B014}{end}

\olpseganchor{OLP-0248-B015}{start}
په لا صوري ډول، څرنګه چې~\LR{$h$} جزوي محاسبه کېدونکې ده، د کوم شاخص~\LR{$e$}
دپاره له تابع~\LR{$\cfind{e}$} سره برابره ده۔ د~\LR{$s$}-\LR{$m$}-\LR{$n$} قضيې له مخې
داسې بنسټيزه بازګشتي تابع~\LR{$s$} شته چې د هر~\LR{$x$} دپاره
\LR{$\cfind{s(e,x)}(y) \simeq h_x(y)$} وي۔
\textbf{د ژباړې د نسخې سمون (OLCMP-043):} سرچينې دلته د ممکنه
ناتعريفو جزوي تابعو د ارزښتونو تر منځ د عادي برابرۍ نښه کارولې وه؛
د پاراميټري کولو د قضيې له خپل بيان سره سم دلته د جزوي برابرۍ نښه
ساتل شوې ده۔
اوس د هر~\LR{$x$} دپاره لرو چې که
\LR{$\cfind{x}(x) \fdefined$} وي، نو~\LR{$\cfind{s(e,x)}$} هماغه تابع ده چې
~\LR{$g$} ده، او ځکه~\LR{$s(e,x)$} په~\LR{$A$} کښې دے۔ له بلې خوا، که
\LR{$\cfind{x}(x) \uparrow$} وي، نو~\LR{$\cfind{s(e,x)}$} هماغه تابع ده چې
~\LR{$f$} ده، او ځکه~\LR{$s(e,x)$} په~\LR{$A$} کښې نۀ دے۔ په بله وينا، د هر~\LR{$x$}
دپاره \LR{$x \in K$} هله او يوازې هله چې \LR{$s(e,x) \in A$} وي۔ که~\LR{$A$}
محاسبه کېدونکے وای، ~\LR{$K$} هم محاسبه کېدونکے و، چې تناقض دے۔ نو~\LR{$A$}
محاسبه کېدونکے نۀ دے۔
\end{proof}
\olpseganchor{OLP-0248-B015}{end}

\olpseganchor{OLP-0248-B016}{start}
د رايس قضيه ډېره پياوړې ده۔ لاندې سمدستي نتيجه ئې د کارونې څو
بېلګې ښيي۔
\begin{cor}
لاندې سټونه د پرېکړې وړ نۀ دي۔
\begin{enumerate}
\item \LR{$\Setabs{x}{\text{\RL{\LR{$17$} د \LR{$\cfind{x}$} په قيمتونو کښې دے}}}$}
\item \LR{$\Setabs{x}{\text{\RL{\LR{$\cfind{x}$} ثابته ده}}}$}
\item \LR{$\Setabs{x}{\text{\RL{\LR{$\cfind{x}$} هرځاے تعريف شوې ده}}}$}
\item \LR{$\Setabs{x}{\text{\RL{هر کله چې \LR{$y < y'$}، \LR{$\cfind{x}(y) \fdefined$}، او
    که \LR{$\cfind{x}(y') \fdefined$}، نو \LR{$\cfind{x}(y) < \cfind{x}(y')$}}}}$}
\end{enumerate}
\end{cor}
\olpseganchor{OLP-0248-B016}{end}

\olpseganchor{OLP-0248-B017}{start}
\begin{proof}
  دا ټول غيرابتذالي شاخصي سټونه دي۔
\end{proof}
\olpseganchor{OLP-0248-B017}{end}

\olpseganchor{OLP-0249-B004}{start}
\olfileid{cmp}{thy}{fix}
\olsection{د ثابت ټکي قضيه}
\olpseganchor{OLP-0249-B004}{end}

\olpseganchor{OLP-0249-B005}{start}
راځئ يو ځل بيا د درېدنې مسئله وګورو۔ د لنډمهالې نښې په توګه
~\LR{$\gn{\cfind{x}(y)}$} د~\LR{$\tuple{x, y}$} دپاره وليکو؛ دا د
\LR{$\cfind{x}(y)$} د قيمت د يو ``نوم'' استازيتوب وګڼئ۔ په دې نښه د
درېدنې د مسئلې د ناپرېکړتيا يو ثبوت بيا بيانولے شو۔
\olpseganchor{OLP-0249-B005}{end}

\olpseganchor{OLP-0249-B006}{start}
پوښتنه: ايا داسې محاسبه کېدونکې تابع~\LR{$h$} شته چې لاندې خاصيت ولري؟ د
هر~\LR{$x$} او~\LR{$y$} دپاره
\[
h(\gn{\cfind{x}(y)}) =
\begin{cases}
1 & \text{\RL{که \LR{$\cfind{x}(y) \fdefined$} وي}} \\
0 & \text{\RL{که داسې نۀ وي۔}}
\end{cases}
\]
\olpseganchor{OLP-0249-B006}{end}

\olpseganchor{OLP-0249-B007}{start}
ځواب: نه؛ که نه، دا جزوي تابع
\[
g(x) \simeq
\begin{cases}
0 & \text{\RL{که \LR{$h(\gn{\cfind{x}(x)}) = 0$} وي}} \\
\text{\RL{تعريف شوې نۀ ده}} & \text{\RL{که داسې نۀ وي}}
\end{cases}
\]
به محاسبه کېدونکې وه، او ځکه به ئې کوم شاخص~\LR{$e$} درلود۔ خو بيا لرو:
\[
\cfind{e}(e) \simeq
\begin{cases}
0 & \text{\RL{که \LR{$h(\gn{\cfind{e}(e)}) = 0$} وي}} \\
\text{\RL{تعريف شوې نۀ ده}} & \text{\RL{که داسې نۀ وي،}}
\end{cases}
\]
چې په دې صورت کښې~\LR{$\cfind{e}(e)$} هله او يوازې هله تعريف شوې ده چې
تعريف شوې نۀ وي، او دا تناقض دے۔
\olpseganchor{OLP-0249-B007}{end}

\olpseganchor{OLP-0249-B008}{start}
اوس هغه معادله وګورئ چې~\LR{$\cfind{e}$} لري۔ دلته په يوه معنا د
ځان-مراجعې يوه بېلګه شته: موږ د~\LR{$\cfind{e}(e)$} قيمت داسې برابر کړے
چې په يوه ځانګړې طريقه پر~\LR{$\gn{\cfind{e}(e)}$} تکيه وکړي۔ د ثابت ټکي
قضيه وايي چې دا کار په عمومي ډول \emph{کولے شو}---يوازې د تناقضونو د
ثابتولو دپاره نۀ۔
\olpseganchor{OLP-0249-B008}{end}

\olpseganchor{OLP-0249-B009}{start}
\olref{lem:fixed-equiv} د ثابت ټکي د قضيې د بيانولو دوه معادل شکلونه
راکوي۔ د منطقي نظره، د دې بيانونو معادلتوب له دې راځي چې دواړه رښتيا
دي؛ خو زموږ اصلي موخه دا ده چې هر يو په نېغه توګه له بل څخه راځي، نو
د همدې قضيې د بديلو بيانونو په توګه اخيستل کېدے شي۔
\olpseganchor{OLP-0249-B009}{end}

\olpseganchor{OLP-0249-B010}{start}
\begin{lem}
\ollabel{lem:fixed-equiv}
لاندې بيانونه معادل دي:
\begin{enumerate}
\item د هرې جزوي محاسبه کېدونکې تابع~\LR{$g(x,y)$} دپاره داسې شاخص~\LR{$e$}
  شته چې د هر~\LR{$y$} دپاره
\[
\cfind{e}(y) \simeq g(e,y).
\]
\item د هرې محاسبه کېدونکې تابع~\LR{$f(x)$} دپاره داسې شاخص~\LR{$e$} شته چې
  د هر~\LR{$y$} دپاره
\[
\cfind{e}(y) \simeq \cfind{f(e)}(y).
\]
\end{enumerate}
\end{lem}
\olpseganchor{OLP-0249-B010}{end}

\olpseganchor{OLP-0249-B011}{start}
\begin{proof}
\LR{$(1) \Rightarrow (2)$}: ورکړې~\LR{$f$} ته~\LR{$g$} داسې تعريف کړئ چې
\LR{$g(x,y) \simeq \fn{Un}(f(x),y)$} وي۔ (1) وکاروئ څو داسې شاخص~\LR{$e$}
تر لاسه کړئ چې د هر~\LR{$y$} دپاره
\begin{align*}
\cfind{e}(y) & \simeq \fn{Un}(f(e),y) \\
& \simeq \cfind{f(e)}(y).
\end{align*}
\olpseganchor{OLP-0249-B011}{end}

\olpseganchor{OLP-0249-B012}{start}
\LR{$(2) \Rightarrow (1)$}: ورکړې~\LR{$g$} ته د~\LR{$s$}-\LR{$m$}-\LR{$n$} قضيه وکاروئ څو
داسې~\LR{$f$} تر لاسه کړئ چې د هر~\LR{$x$} او~\LR{$y$} دپاره
\LR{$\cfind{f(x)}(y) \simeq g(x,y)$} وي۔ (2) وکاروئ څو داسې شاخص~\LR{$e$} تر
لاسه کړئ چې
\begin{align*}
\cfind{e}(y) & \simeq \cfind{f(e)}(y) \\
& \simeq g(e,y).
\end{align*}
په دې سره ثبوت بشپړېږي۔
\textbf{د ژباړې د نسخې سمون (OLCMP-038):} سرچينې په دې دوو
استنتاجي لړيو کښې د جزوي تابعو د قيمتونو تر منځ د عادي برابرۍ نښه
کارولې وه، سره له دې چې اړخونه تعريف شوي نۀ هم کېدے شي۔ دلته د همدې
څپرکي ټاکلې جزوي برابرۍ نښه کارول شوې ده، څو دواړه تعريف شوي او
ناتعريف حالتونه په کره ډول پوښل شي۔
\end{proof}
\olpseganchor{OLP-0249-B012}{end}

\olpseganchor{OLP-0249-B013}{start}
\begin{explain}
مخکښې له دې چې وښيو بيان~(1) رښتيا دے (او ځکه~(2) هم)، فکر وکړئ چې
دا څومره عجيبه خبره ده۔ \LR{$e$} يو کمپيوټري پروګرام وګڼئ؛ بيان~(1) وايي
چې د هرې ورکړې جزوي محاسبه کېدونکې~\LR{$g(x,y)$} دپاره داسې کمپيوټري
پروګرام~\LR{$e$} موندلے شئ چې \LR{$g_e(y) \simeq g(e,y)$} محاسبه کوي۔ په بله
وينا، داسې کمپيوټري پروګرام موندلے شئ چې يوه پر خپل پروګرام مراجعه
کوونکې تابع محاسبه کوي۔
\end{explain}
\olpseganchor{OLP-0249-B013}{end}

\olpseganchor{OLP-0249-B014}{start}
\begin{thm}
په~\olref{lem:fixed-equiv} کښې دواړه بيانونه رښتيا دي۔ په ځانګړي
ډول، د هرې جزوي محاسبه کېدونکې تابع~\LR{$g(x,y)$} دپاره داسې شاخص~\LR{$e$}
شته چې د هر~\LR{$y$} دپاره
\[
\cfind{e}(y) \simeq g(e,y).
\]
\end{thm}
\olpseganchor{OLP-0249-B014}{end}

\olpseganchor{OLP-0249-B015}{start}
\begin{proof}
اړين توکي لا د پورته درېدنې مسئلې په بحث کښې نغښتي دي۔
\LR{$\fn{diag}(x)$} داسې محاسبه کېدونکې تابع وي چې د هر~\LR{$x$} دپاره د تابع
\LR{$f_x(y) \simeq \cfind{x}(x,y)$} يو شاخص راګرځوي، يعنې
\[
\cfind{\fn{diag}(x)}(y) \simeq \cfind{x}(x,y).
\]
\LR{$\fn{diag}$} داسې تابع وګڼئ چې د يوې دوه‌ځاييزې تابع پروګرام د يوې
يوځاييزې تابع په پروګرام بدلوي، او دا د اصلي پروګرام په لومړي دليل
کښې د هغه خپل شاخص په ثابتولو تر لاسه کېږي۔ تابع~\LR{$\fn{diag}$} په صوري
ډول داسې تعريفېدے شي: لومړے~\LR{$s$} داسې تعريف کړئ:
\[
s(x,y) \simeq \fn{Un}^2(x,x,y),
\]
چېرته چې~\LR{$\fn{Un}^2$} يوه درې‌ځاييزه تابع ده چې د جزوي محاسبه
کېدونکو دوه‌ځاييزو تابعو دپاره نړيواله ده۔ بيا د~\LR{$s$}-\LR{$m$}-\LR{$n$} قضيې
له مخې داسې بنسټيزه بازګشتي تابع~\LR{$\fn{diag}$} موندلے شو چې دا پوره
کړي:
\[
\cfind{\fn{diag}(x)}(y) \simeq s(x,y).
\]
\olpseganchor{OLP-0249-B015}{end}

\olpseganchor{OLP-0249-B016}{start}
اوس تابع~\LR{$l$} داسې تعريف کړئ:
\[
l(x,y) \simeq g(\fn{diag}(x),y).
\]
او~\LR{$\gn{l}$} د~\LR{$l$} يو شاخص وي۔ په پای کښې
\LR{$e = \fn{diag}(\gn{l})$} وټاکئ۔ بيا د هر~\LR{$y$} دپاره لرو:
\begin{align*}
\cfind{e}(y) & \simeq \cfind{\fn{diag}(\gn{l})}(y) \\
& \simeq \cfind{\gn{l}}(\gn{l}, y) \\
& \simeq l(\gn{l}, y) \\
& \simeq g(\fn{diag}(\gn{l}),y) \\
& \simeq g(e, y),
\end{align*}
لکه چې غوښتل شوي وو۔
\end{proof}
\olpseganchor{OLP-0249-B016}{end}

\olpseganchor{OLP-0249-B017}{start}
\begin{explain}
دلته څه روان دي؟ فرض کړئ تاسو ته د داسې کمپيوټري پروګرام د ليکلو
دنده درکړل شوې چې خپل متن چاپ کړي۔ بيا فرض کړئ چې په داسې پروګرامي
ژبه کار کوئ چې د رشتو يو بډای او عجيب کتابتون لري۔ په ځانګړي ډول،
فرض کړئ ستاسو پروګرامي ژبه داسې تابع~\LR{$\fn{diag}$} لري: تابع
\LR{$\fn{diag}$} يوه ننوت رشته~\LR{$s$} اخلي، په~\LR{$s$} کښې د `x' نښې هره بېلګه
مومي، او هغه د اصلي
رشتې په اقتباس شوې بڼه بدلوي۔ د بېلګې په توګه، که دا رشته
\begin{quote}
\begin{verbatim}
hello x world
\end{verbatim}
\end{quote}
د ننوت په توګه ورکړل شي، تابع دا
\begin{quote}
\begin{verbatim}
hello 'hello x world' world
\end{verbatim}
\end{quote}
د وت په توګه راګرځوي۔ په دې صورت کښې موخه شوے پروګرام ليکل اسان
دي؛ ازمويلے شئ چې
\begin{quote}
\begin{verbatim}
print(diag('print(diag(x))'))
\end{verbatim}
\end{quote}
کار سر ته رسوي۔ د~C++ او جاوا په شان ډېرو عامو پروګرامي ژبو کښې هم
همدغه نظر (په يو څه پېچلې پليينې سره) کار کوي۔
\olpseganchor{OLP-0249-B017}{end}

\olpseganchor{OLP-0249-B018}{start}
موږ د ثابت ټکي د قضيې له ثبوت څخه يوازې څو ګامه لرې يو۔ فرض کړئ د
چاپ تابعې يوه بڼه~\LR{$\fn{print}(x,y)$} يوه رشته~\LR{$x$} او بل عددي دليل~\LR{$y$}
اخلي، او رشته~\LR{$x$}، \LR{$y$} ځله چاپ کوي۔ بيا دا ``پروګرام''
\begin{quote}
\begin{verbatim}
getinput(y);
print(diag('getinput(y); print(diag(x), y)'), y)
\end{verbatim}
\end{quote}
پر ننوت~\LR{$y$} خپل ځان~\LR{$y$} ځله چاپ کوي۔ که د
\LR{$\fn{getinput}$}---\LR{$\fn{print}$}---\LR{$\fn{diag}$} چوکاټ په هرې تابع
\LR{$g(x,y)$} بدل کړو، دا تر لاسه کېږي:
\begin{quote}
\begin{verbatim}
g(diag('g(diag(x), y)'), y)
\end{verbatim}
\end{quote}
دا داسې پروګرام دے چې پر ننوت~\LR{$y$}، تابع~\LR{$g$} پر خپل پروګرام او~\LR{$y$}
چلوي۔ که ``اقتباس کول'' د ``يو شاخص کارولو'' په توګه وګڼو، پورته
ثبوت تر لاسه کوو۔
\olpseganchor{OLP-0249-B018}{end}

\olpseganchor{OLP-0249-B019}{start}
اوس دپاره که غواړئ ثبوت صوري چل، يا توره جادوګري، وګڼئ پروا نۀ کوي۔
خو بايد د پورته ورکړل شوي استدلال جزئيات بيا جوړ کړے شئ۔ کله چې د
نابشپړتيا قضيې (او اړونده ``د ثابت ټکي قضيه'') ثابتوو، د دې د کار
کولو د علت د پوهېدو په نورو لارو به بحث وکړو۔
\end{explain}
\olpseganchor{OLP-0249-B019}{end}

\olpseganchor{OLP-0249-B020}{start}

\begin{digress}
همدغه نظر د ``ثابت ټکي'' يو ترکيبګر هم راکولے شي۔ فرض کړئ يو لامبډا
ترم~\LR{$g$} لرئ، او داسې بل ترم~\LR{$k$} غواړئ چې~\LR{$k$} له~\LR{$gk$} سره
\LR{$\beta$}-معادل وي۔ زموږ د نښيزو دودونو په کارولو ترمونه داسې تعريف
کړئ:
\[
\fn{diag}(x) = xx
\]
او
\[
l(x) = g(\fn{diag}(x))
\]
يعنې~\LR{$l$} دا ترم دے: \LR{$\lambd[x][g(xx)]$}۔ \LR{$k$} ترم~\LR{$ll$} وي۔ بيا لرو:
\begin{align*}
k & = (\lambd[x][g(xx)])(\lambd[x][g(xx)]) \\
& \red  g((\lambd[x][g(xx)])(\lambd[x][g(xx)])) \\
& = gk.
\end{align*}
که
\[
Y = \lambd[g][((\lambd[x][g(xx)])(\lambd[x][g(xx)]))]
\]
واخلو، نو~\LR{$Yg$} او~\LR{$g(Yg)$} يو ګډ ترم ته راکمېږي؛ ځکه
\LR{$Yg \equiv_\beta g(Yg)$}۔ دې ته ``د کري ترکيبګر'' ويل کېږي۔ که پر
ځاے ئې
\[
Y = (\lambd[xg][g(xxg)])(\lambd[xg][g(xxg)])
\]
واخلو، نو په حقيقت کښې~\LR{$Yg$}، \LR{$g(Yg)$} ته راکمېږي، چې لا پياوړے بيان
دے۔ د~\LR{$Y$} دا وروستۍ بڼه ``د ټيورينګ ترکيبګر'' بلل کېږي۔
\end{digress}

\olpseganchor{OLP-0249-B020}{end}

\olpseganchor{OLP-0250-B004}{start}
\olfileid{cmp}{thy}{apf}
\olsection{د ثابت ټکي د قضيې کارول}
\olpseganchor{OLP-0250-B004}{end}

\olpseganchor{OLP-0250-B005}{start}
د ثابت ټکي قضيه په اصل کښې اجازه راکوي چې جزوي محاسبه کېدونکې تابعې
د هغو د شاخصونو په وسيله تعريف کړو۔ د بېلګې په توګه، داسې شاخص~\LR{$e$}
موندلے شو چې د هر~\LR{$y$} دپاره
\[
\cfind{e}(y) = e + y.
\]
د بلې بېلګې په توګه، د ثابت ټکي د قضيې ثبوت کارولے شو څو په جاوا يا
C++ کښې داسې پروګرام جوړ کړو چې خپل متن چاپ کړي۔
\olpseganchor{OLP-0250-B005}{end}

\olpseganchor{OLP-0250-B006}{start}
راياد کړئ چې که د هر~\LR{$e$} دپاره~\LR{$W_e$} د~\LR{$\cfind{e}$} د تعريف ساحه
وټاکو، نو لړۍ~\LR{$W_0$}، \LR{$W_1$}، \LR{$W_2$}،~\dots په محاسبوي ډول د شمېر وړ
سټونه شمېري۔ ځينې دا سټونه محاسبه کېدونکي دي۔ پوښتلے شو چې ايا داسې
الګوريتم شته چې قيمت~\LR{$x$} د ننوت په توګه واخلي، او که~\LR{$W_x$} محاسبه
کېدونکے وي، د هغه د ځانګړونکې تابع يو شاخص را وګرځوي؟ ځواب ``نه''
دے، داسې الګوريتم نشته:
\olpseganchor{OLP-0250-B006}{end}

\olpseganchor{OLP-0250-B007}{start}
\begin{thm}
داسې جزوي محاسبه کېدونکې تابع~\LR{$f$} نشته چې دا خاصيت ولري: هر کله چې
~\LR{$W_e$} محاسبه کېدونکے وي، نو~\LR{$f(e)$} تعريف شوے وي او~\LR{$\cfind{f(e)}$}
ئې ځانګړونکې تابع وي۔
\end{thm}
\olpseganchor{OLP-0250-B007}{end}

\olpseganchor{OLP-0250-B008}{start}
\begin{proof}
\LR{$f$} هره جزوي محاسبه کېدونکې تابع وي؛ داسې~\LR{$e$} به جوړ کړو چې~\LR{$W_e$}
محاسبه کېدونکے وي، خو~\LR{$\cfind{f(e)}$} ئې ځانګړونکې تابع نۀ وي۔ د ثابت
ټکي د قضيې په کارولو داسې شاخص~\LR{$e$} موندلے شو چې
\[
\cfind{e}(y) \simeq
\begin{cases}
0 & \text{\RL{که \LR{$y=0$} او \LR{$\cfind{f(e)}(0) \fdefined = 0$} وي}} \\
\text{\RL{تعريف شوې نۀ ده}} & \text{\RL{که داسې نۀ وي۔}}
\end{cases}
\]
يعنې~\LR{$e$} د ثابت ټکي د قضيې په دې تابع له پلي کولو تر لاسه کېږي:
\[
g(x,y) \simeq
\begin{cases}
0 & \text{\RL{که \LR{$y=0$} او \LR{$\cfind{f(x)}(0) \fdefined = 0$} وي}} \\
\text{\RL{تعريف شوې نۀ ده}} & \text{\RL{که داسې نۀ وي۔}}
\end{cases}
\]
په غير رسمي ډول داسې ليدلے شو چې~\LR{$g$} جزوي محاسبه کېدونکې ده: پر
ننوتونو~\LR{$x$} او~\LR{$y$} الګوريتم لومړے ګوري چې~\LR{$y$} له~\LR{$0$} سره برابر دے
که نه۔ که برابر وي، الګوريتم~\LR{$f(x)$} محاسبه کوي، او بيا نړيوال ماشين
کاروي څو~\LR{$\cfind{f(x)}(0)$} محاسبه کړي۔ که دا وروستۍ محاسبه ودرېږي او
~\LR{$0$} را وګرځوي، الګوريتم~\LR{$0$} راګرځوي؛ که نه، الګوريتم نۀ درېږي۔
\textbf{د ژباړې د نسخې سمون (OLCMP-036):} سرچينې د ثبوت په پيل کښې
خپل اختياري سيال بشپړه محاسبه کېدونکې تابع بللې وه، حال دا چې د
قضيې دعوه د ټولو جزوي محاسبه کېدونکو تابعو په اړه ده۔ دلته د قضيې
هماغه جزوي ټولګے ساتل شوے؛ ورپسې جوړښت د ناتعريف حالت په ګډون هم
کار کوي۔
\olpseganchor{OLP-0250-B008}{end}

\olpseganchor{OLP-0250-B009}{start}
خو اوس پام وکړئ چې که~\LR{$\cfind{f(e)}(0)$} تعريف شوې او له~\LR{$0$} سره
برابره وي، نو~\LR{$\cfind{e}(y)$} کټ مټ هغه وخت تعريف شوې ده چې~\LR{$y$} له
~\LR{$0$} سره برابر وي، نو \LR{$W_e = \{ 0 \}$}۔ که~\LR{$\cfind{f(e)}(0)$} تعريف
شوې نۀ وي، يا تعريف شوې خو له~\LR{$0$} سره برابره نۀ وي، نو
\LR{$W_e = \emptyset$}۔ په دواړو حالاتو کښې~\LR{$\cfind{f(e)}$} د~\LR{$W_e$}
ځانګړونکې تابع نۀ ده، ځکه پر ننوت~\LR{$0$} ناسم ځواب ورکوي۔
\end{proof}
\olpseganchor{OLP-0250-B009}{end}

\olpseganchor{OLP-0251-B004}{start}
\olfileid{cmp}{thy}{slf}
\olsection{د ځان-مراجعې په کارولو د تابعو تعريف}
\olpseganchor{OLP-0251-B004}{end}

\olpseganchor{OLP-0251-B005}{start}
په عمومي ډول دا ګټوره ده چې تابعې د هغو په خپلو اصطلاحاتو کښې تعريف
کړے شو۔ د بېلګې په توګه، ورکړو محاسبه کېدونکو تابعو~\LR{$k$}، \LR{$l$} او~\LR{$m$}
ته د ثابت ټکي لمه وايي چې داسې جزوي محاسبه کېدونکې تابع~\LR{$f$} شته چې
د هر~\LR{$y$} دپاره دا معادله پوره کوي:
\[
f(y) \simeq
\begin{cases}
k(y) & \text{\RL{که \LR{$l(y) = 0$} وي}} \\
f(m(y)) & \text{\RL{که داسې نۀ وي۔}}
\end{cases}
\]
په لا ځانګړي ډول، ~\LR{$f$} بيا داسې تر لاسه کېږي:
\[
g(x,y) \simeq
\begin{cases}
k(y) & \text{\RL{که \LR{$l(y) = 0$} وي}} \\
\cfind{x}(m(y)) & \text{\RL{که داسې نۀ وي}}
\end{cases}
\]
او بيا د ثابت ټکي لمه کارول کېږي څو داسې شاخص~\LR{$e$} وموندل شي چې
\LR{$\cfind{e}(y) \simeq g(e,y)$} وي۔
\textbf{د ژباړې د نسخې سمون (OLCMP-044):} سرچينې د ثابت ټکي په دې
پليينه کښې د دوو جزوي تابعو د ممکنه ناتعريفو ارزښتونو تر منځ عادي
برابري کارولې وه؛ دلته د قضيې ټاکلې جزوي برابري ساتل شوې ده۔
\olpseganchor{OLP-0251-B005}{end}

\olpseganchor{OLP-0251-B006}{start}
د يوې محسوسې بېلګې دپاره د ``تر ټولو لوی ګډ مقسوم عليه'' تابع
\LR{$\fn{gcd}(u,v)$} داسې تعريفېدے شي:
\[
\fn{gcd}(u,v) \simeq
\begin{cases}
v & \text{\RL{که \LR{$u = 0$} وي}} \\
\fn{gcd}(\fn{mod}(v, u), u) & \text{\RL{که داسې نۀ وي}}
\end{cases}
\]
چېرته چې~\LR{$\fn{mod}(v, u)$} پر~\LR{$u$} د~\LR{$v$} د وېش پاتې شونې څرګندوي۔ د
ثابت ټکي لمې ته مراجعه ښيي چې~\LR{$\fn{gcd}$} جزوي محاسبه کېدونکې ده۔
(په حقيقت کښې دا په پورته بڼه راوستل کېدے شي، چې~\LR{$y$} د جوړې
\LR{$\tuple{u, v}$} کوډ وګڼل شي۔) بيا پر~\LR{$u$} استقرا ښيي چې په حقيقت کښې
~\LR{$\fn{gcd}$} هرځاے تعريف شوې ده۔
\olpseganchor{OLP-0251-B006}{end}

\olpseganchor{OLP-0251-B007}{start}
البته داسې ځان-مراجع تعريفونه جوړولے شو چې له پورته بېلګو ډېر پېچلي
وي۔ زياتره پروګرامي ژبې په يوه يا بله طريقه د تابعو تعريف په خپلو
اصطلاحاتو کښې مني۔ پام وکړئ چې دا د هغې وړتيا په پرتله لږ حيرانوونکے
دے چې يوه تابع د هغه الګوريتم د \emph{شاخص} په وسيله تعريف کړو چې
هماغه تابع محاسبه کوي؛ او د ثابت ټکي قضيه په بشپړ عموميت کښې همدا
اجازه راکوي۔
\textbf{د ژباړې د نسخې سمون (OLCMP-037):} سرچينې د دې وروستۍ جملې
په پای کښې د يو الګوريتم له خوا محاسبه کېدونکې يوه تابع په جمع نوم
ياده کړې وه۔ دلته شمېر د هماغې ټاکلې تابع سره يو شان ساتل شوے دے۔
\olpseganchor{OLP-0251-B007}{end}

\olpseganchor{OLP-0252-B004}{start}
\olpart{tur}{ټيورينګ ماشينونه}
\olpseganchor{OLP-0252-B004}{end}

\olpseganchor{OLP-0253-B004}{start}
\olchapter{tur}{mac}{د ټيورينګ ماشين محاسبې}
\olpseganchor{OLP-0253-B004}{end}

\olpseganchor{OLP-0254-B004}{start}
\olfileid{tur}{mac}{int}
\olsection{پېژندنه}
\olpseganchor{OLP-0254-B004}{end}

\olpseganchor{OLP-0254-B005}{start}
دا څه معنا لري چې، مثلاً، له~\LR{$\Nat$} څخه~\LR{$\Nat$} ته يوه تابع
\emph{محاسبه کېدونکې} وي؟ د لومړنيو او تر ټولو نامتو ځوابونو څخه يو
دا دے چې تابع هغه وخت محاسبه کېدونکې ده چې د ټيورينګ ماشين په وسيله
محاسبه کېدے شي۔ دا مفهوم الن ټيورينګ په~1936 کښې وړاندې کړ۔ د ټيورينګ
ماشينونه د \emph{محاسبې د يو مدل} بېلګه ده---دوي د ``محاسبوي کړنلارې''
نظر د تعريفولو يوه رياضياتي کره لار ده۔ د دې کره معنا تر بحث لاندې
ده، خو په پراخه کچه منل شوې چې ټيورينګ ماشينونه د محاسبوي کړنلارو د
ټاکلو يوه لار ده۔ که څه هم د ``ټيورينګ ماشين'' اصطلاح د خوځنده پرزو
لرونکي فزيکي ماشين انځور ذهن ته راولي، په کره معنا ټيورينګ ماشين يو
سوچه رياضياتي جوړښت دے، او په همدې توګه د محاسبوي کړنلارې نظر مثالي
کوي۔ مثلاً، موږ د ټيورينګ ماشين د وخت يا حافظې پر اړتياوو هېڅ بنديز
نۀ ږدو: ټيورينګ ماشينونه يو څيز ان هغه وخت محاسبه کولے شي چې محاسبه
له هغو ټولو ذرو څخه زياته زېرمه يا زيات ګامونه وغواړي چې په کاينات
کښې شته۔
\olpseganchor{OLP-0254-B005}{end}

\olpseganchor{OLP-0254-B006}{start}
\begin{explain}
ښايي غوره وي چې ټيورينګ ماشين د يو ځانګړي ډول خيالي ميکانيزم د
پروګرام په توګه وګڼو۔ دا ميکانيزم له يوې \emph{پټې} او يو
\emph{لوست-ليک سر} څخه جوړ دے۔ زموږ د ټيورينګ ماشينونو په بڼه کښې
پټه په يوه لوري (ښي لوري ته) نامتناهي ده، او په \emph{مربعونو} وېشل
شوې، چې هر يو ئې د يوې متناهي \emph{الفبا} يوه نښه لرلے شي۔ دا ډول
الفباوې هر شمېر بېلابېلې نښې لرلے شي، خو موږ به تر ډېره په درېيو بسنه
کوو: \LR{$\TMendtape$}، \LR{$\TMblank$} او~\LR{$\TMstroke$}۔ کله چې ميکانيزم پيل
شي، پټه تشه وي (يعنې هر مربع د~\LR{$\TMblank$} نښه لري)، پرته له تر ټولو
کيڼ مربع څخه چې~\LR{$\TMendtape$} لري، او يو متناهي شمېر مربعونو څخه چې
\emph{ننوت} لري۔ په هر وخت کښې ميکانيزم له متناهي شمېر
\emph{حالتونو} څخه په يوه کښې وي۔ په پيل کښې سر د ننوت تر ټولو کيڼ
مربع، يعنې د پټې د پای نښې ښي لوري ته مربع، لولي او په يو ټاکلي
\emph{لومړني حالت} کښې وي۔ د ميکانيزم د چلېدو په هر ګام کښې د اوس
لوستل کېدونکي مربع منځپانګه، د ميکانيزم حالت او د ټيورينګ ماشين
پروګرام په ګډه ټاکي چې ورپسې څه کېږي۔ د ټيورينګ ماشين پروګرام د يوې
جزوي تابع په توګه ورکول کېږي چې حالت~\LR{$q$} او نښه~\LR{$\sigma$} د ننوت په
توګه اخلي، او درې‌يز~\LR{$\tuple{q', \sigma', D}$} راګرځوي۔ هر کله چې
ميکانيزم په حالت~\LR{$q$} کښې وي او نښه~\LR{$\sigma$} لولي، د اوسني مربع نښه
په~\LR{$\sigma'$} بدلوي، سر د دې له مخې کيڼ، ښي لوري ته ځي يا پر ځاے پاتې
کېږي چې~\LR{$D$} په ترتيب سره~\LR{$\TMleft$}، \LR{$\TMright$} يا~\LR{$\TMstay$} وي، او
ميکانيزم حالت~\LR{$q'$} ته ځي۔
\textbf{د ژباړې د نسخې سمون (OLCMP-039):} سرچينې د پټې تر ټولو کيڼ
مربع د پای نښې د ځاے په توګه ټاکلے و، خو بيا ئې سر هم په پيل کښې پر
هماغه مربع ايښے و۔ د همدې څپرکي صوري لومړنی ترتيب سر د پای نښې ښي
لوري ته، د ننوت پر لومړي مربع ږدي؛ دلته هماغه ټاکلے دود په څرګنده
ساتل شوے دے۔
\olpseganchor{OLP-0254-B006}{end}

\olpseganchor{OLP-0254-B007}{start}
د بېلګې په توګه، په~\olref{fig:tm} کښې حالت وګورئ۔
\begin{figure}
  \olasset{../upstream/assets/diagrams/turing-machine.tikz}
  \caption{يو ټيورينګ ماشين چې خپل پروګرام عملي کوي۔}
  \ollabel{fig:tm}
\end{figure}
د ټيورينګ ماشين د پټې ليدونکې برخه په تر ټولو کيڼ مربع کښې د پټې د
پای نښه~\LR{$\TMendtape$} لري، ورپسې درې~\LR{$1$} ګانې، يو~\LR{$0$}، او بيا څلور
نورې~\LR{$1$} ګانې دي۔ سر له کيڼ لوري درېيم مربع لولي، چې~\LR{$1$} لري، او په
حالت~\LR{$q_1$} کښې دے---وايو ``ماشين په حالت~\LR{$q_1$} کښې~\LR{$1$} لولي۔'' که
د ټيورينګ ماشين پروګرام د ننوت~\LR{$\tuple{q_1, 1}$} دپاره درې‌يز
\LR{$\tuple{q_2, 0, \TMstay}$} را وګرځوي، ماشين به اوس په درېيم مربع کښې
~\LR{$1$} په~\LR{$0$} بدل کړي، لوست-ليک سر به پر خپل ځاے پرېږدي، او حالت~\LR{$q_2$}
ته به واوړي۔ که بيا پروګرام د ننوت~\LR{$\tuple{q_2, 0}$} دپاره
\LR{$\tuple{q_3, 0, \TMright}$} را وګرځوي، ماشين به~\LR{$0$} په بل~\LR{$0$} وليکي
(يعنې په عمل کښې به د لوست-ليک سر لاندې د پټې منځپانګه هماغسې
پرېږدي)، يو مربع ښي لوري ته به لاړ شي، او حالت~\LR{$q_3$} ته به ننوځي۔ او
همداسې نور۔
\olpseganchor{OLP-0254-B007}{end}

\olpseganchor{OLP-0254-B008}{start}
وايو ماشين هغه وخت \emph{درېږي} چې له داسې حالت~\LR{$q_n$} او نښې
~\LR{$\sigma$} سره مخ شي چې د~\LR{$\tuple{q_n, \sigma}$} دپاره هېڅ لارښوونه
نۀ وي، يعنې د ننوت~\LR{$\tuple{q_n,\sigma}$} دپاره د حالت بدلون تابع
تعريف شوې نۀ وي۔ په بله وينا، ماشين د عملي کولو دپاره هېڅ لارښوونه
نۀ لري، او په هماغه ځاے کار بندوي۔ کله ناکله درېدل د ځانګړي درېدني
حالت~\LR{$h$} په وسيله څرګندېږي۔ دا به وروسته په زياتو جزئياتو وښودل شي۔
\end{explain}
\olpseganchor{OLP-0254-B008}{end}

\olpseganchor{OLP-0254-B009}{start}
\begin{digress}
د ټيورينګ د مقالې، ``د محاسبه کېدونکو عددونو په اړه،'' ښکلا دا ده چې
هغه نۀ يوازې يو صوري تعريف وړاندې کوي، بلکې دا استدلال هم کوي چې
تعريف د محاسبه کېدنې شهودي مفهوم رانغاړي۔ له تعريفه بايد څرګنده وي
چې هره تابع چې د ټيورينګ ماشين په وسيله محاسبه کېږي، په شهودي معنا
محاسبه کېدونکې ده۔ ټيورينګ درې ډوله استدلالونه وړاندې کوي چې معکوسه
خبره هم رښتيا ده، يعنې هره تابع چې موږ ئې په طبيعي ډول محاسبه کېدونکې
ګڼو، د داسې ماشين په وسيله محاسبه کېدونکې ده۔ د ټيورينګ په خپلو
ټکو، دا دي:
\begin{enumerate}
\item شهود ته نېغه مراجعه۔
\item د دوو تعريفونو د معادلتوب ثبوت (په هغه صورت کښې چې نوی تعريف
  زيات شهودي جاذبيت ولري)۔
\item د محاسبه کېدونکو عددونو د لويو ټولګيو بېلګې ورکول۔
\end{enumerate}
زموږ موخه دا ده چې د محاسبه کېدنې مفهوم ``په اصل کښې'' تعريف کړو،
يعنې د وخت او ځاے عملي محدوديتونه په پام کښې وانۀ خلو۔ البته، کله چې
د محاسبه کېدنې تر ټولو پراخ تعريف ولرو، بيا د محدودو سرچينو لرونکې
محاسبه څېړلے شو؛ دا د ``محاسبوي پېچلتيا'' په نوم مضمون زړه جوړوي۔
\end{digress}
\olpseganchor{OLP-0254-B009}{end}

\olpseganchor{OLP-0254-B010}{start}
\begin{history}
الن ټيورينګ په~1936 کښې ټيورينګ ماشينونه اختراع کړل۔ که څه هم په هغه
وخت کښې د هغه پاملرنه د لومړۍ درجې منطق پرېکړه کېدنې ته وه، دا مقاله
د کمپيوټر د جوړښت د بنسټونو يوه پرېکنده مقاله بلل شوې ده۔ په مقاله
کښې ټيورينګ پر محاسبه کېدونکو حقيقي عددونو، يعنې هغو حقيقي عددونو چې
اعشاري غځونې ئې محاسبه کېدونکې دي، تمرکز کوي؛ خو يادونه کوي چې د هغه
مفهومونه پر طبيعي عددونو محاسبه کېدونکو تابعو او نورو ته اړول ستونزمن
نۀ دي۔ پام وکړئ چې دا په~1941 کښې د لومړي کارکوونکي عمومي-موخې
کمپيوټر له جوړېدو پوره پنځه کاله مخکښې وو (الماني کونراد زوزه دا د
خپلو والدينو د کور په ناسته کوټه کښې جوړ کړ)، په بلېچلي پارک کښې د
ټيورينګ او د هغه د ملګرو د کوډ ماتوونکي کولوسس له جوړېدو اووه کاله
مخکښې (1943)، د امريکايي اېنياک څخه نهه کاله مخکښې (1945)، په
مانچسټر کښې د لومړي برتانوي عمومي-موخې کمپيوټر---د مانچسټر د کوچنۍ
کچې ازمېښتي ماشين---له جوړېدو دولس کاله مخکښې (1948)، او د
امريکايانو له خوا د بيناک له لومړۍ ازموينې ديارلس کاله مخکښې (1949)۔
د مانچسټر کوچنۍ کچې ازمېښتی ماشين د لومړي زېرمه-پروګرام کمپيوټر وياړ
لري---پخواني ماشينونه د هرې نوې دندې دپاره په لاس بيا سيمبندي کېدل۔
\end{history}
\olpseganchor{OLP-0254-B010}{end}

\olpseganchor{OLP-0255-B004}{start}
\olfileid{tur}{mac}{rep}
\olsection{د ټيورينګ ماشينونو څرګندونه}
\olpseganchor{OLP-0255-B004}{end}

\olpseganchor{OLP-0255-B005}{start}
\begin{explain}
ټيورينګ ماشينونه په ليديز ډول د \emph{حالت ډياګرامونو} په وسيله
ښودل کېدے شي۔ دا ډياګرامونه د حالت له هغو خانو جوړ دي چې غشي ئې سره
نښلوي۔ لکه چې تمه کېږي، د حالت هره خانه د ماشين د يو حالت استازيتوب
کوي۔ هر غشی د هغې لارښوونې استازيتوب کوي چې له هماغه حالت څخه عملي
کېدے شي، او د لارښوونې جزئيات د اړوند غشي دپاسه يا لاندې ليکل شوي
وي۔ لاندې ماشين وګورئ، چې يوازې دوه دننني حالتونه، \LR{$q_0$} او~\LR{$q_1$}،
او يوه لارښوونه لري:
\[
\begin{tikzpicture}[->,>=stealth',shorten >=1pt,auto,node distance=2.8cm,
                    semithick]
  \tikzstyle{every state}=[fill=none,draw=black,text=black]
\olpseganchor{OLP-0255-B005}{end}

\olpseganchor{OLP-0255-B006}{start}
  \node[initial,state]   (A)              {\LR{$q_0$}};
  \node[state]   (B) [right of=A] {\LR{$q_1$}};
\olpseganchor{OLP-0255-B006}{end}

\olpseganchor{OLP-0255-B007}{start}
  \path (A) edge  node {\TMtrans{\TMblank}{\TMstroke}{\TMright}} (B);
\end{tikzpicture}
\]
په ياد ولرئ چې ټيورينګ ماشين يو لوست-ليک سر او يوه پټه لري چې ننوت
پرې ليکل شوے وي۔ لارښوونه داسې لوستل کېږي: \emph{که په حالت~\LR{$q_0$}
کښې~\LR{$\TMblank$} لولئ، يو~\LR{$\TMstroke$} وليکئ، ښي لوري ته لاړ شئ، او
حالت~\LR{$q_1$} ته واوړئ}۔ دا له دې سره معادل ده چې د حالت بدلون تابع
\LR{$\tuple{q_0, \TMblank}$} په~\LR{$\tuple{q_1, \TMstroke,
\TMright}$} واړوي۔
\end{explain}
\olpseganchor{OLP-0255-B007}{end}

\olpseganchor{OLP-0255-B008}{start}
\begin{ex}
\emph{جفت ماشين}: لاندې ټيورينګ ماشين هله او يوازې هله درېږي چې پر
پټه د~\LR{$\TMstroke$} ګانو شمېر جفت وي (په دې فرض چې د پټې ټول
\LR{$\TMstroke$} ګانې له لومړي~\LR{$\TMblank$} څخه مخکښې راځي)۔
\[
\begin{tikzpicture}[->,>=stealth',shorten >=1pt,auto,node distance=2.8cm,
                    semithick]
  \tikzstyle{every state}=[fill=none,draw=black,text=black]
\olpseganchor{OLP-0255-B008}{end}

\olpseganchor{OLP-0255-B009}{start}
  \node[initial,state]         (A)              {\LR{$q_0$}};
  \node[state]         (B) [right of=A] {\LR{$q_1$}};
\olpseganchor{OLP-0255-B009}{end}

\olpseganchor{OLP-0255-B010}{start}
  \path (A) edge [bend left] node {\TMtrans{\TMstroke}{\TMstroke}{\TMright}} (B)
        (B) edge [loop above] node {\TMtrans{\TMblank}{\TMblank}{\TMright}} (B)
            edge [bend left] node {\TMtrans{\TMstroke}{\TMstroke}{\TMright}} (A);
\end{tikzpicture}
\]
\olpseganchor{OLP-0255-B010}{end}

\olpseganchor{OLP-0255-B011}{start}
د حالت ډياګرام له لاندې حالت بدلون تابع سره سمون لري:
\begin{align*}
\delta(q_0, \TMstroke) & = \tuple{q_1, \TMstroke, \TMright},\\
\delta(q_1, \TMstroke) & = \tuple{q_0, \TMstroke, \TMright},\\
\delta(q_1, \TMblank)  & = \tuple{q_1, \TMblank, \TMright}
\end{align*}
\end{ex}
\olpseganchor{OLP-0255-B011}{end}

\olpseganchor{OLP-0255-B012}{start}
\begin{explain}
پورته ماشين يوازې هغه وخت درېږي چې ننوت د کرښو جفت شمېر وي۔ که نه،
ماشين (په نظري ډول) تر نامعلومې مودې کار ته دوام ورکوي۔ د هر ماشين
او ننوت دپاره د ماشين \emph{تشکيلونه} پرله‌پسې څارلے شو، څو وت ئې
وټاکو۔ د تشکيلونو صوري تعريف به وروسته ورکړو۔ اوس د تشکيلونو په اړه
په شهودي ډول د ډياګرامونو د يوې لړۍ په توګه فکر کولے شو، چې د کار په
بهير کښې په هره شېبه د ماشين حالت ښيي۔ تشکيلونه د پټې منځپانګه، د
ماشين حالت او د لوست-ليک سر ځاے ښيي۔
\olpseganchor{OLP-0255-B012}{end}

\olpseganchor{OLP-0255-B013}{start}
راځئ د جفت ماشين تشکيلونه هغه وخت پرله‌پسې وڅارو چې د څلورو
\LR{$\TMstroke$} ګانو پر ننوت پيل شي۔ په دې صورت کښې تمه لرو چې ماشين
ودرېږي۔ بيا به ماشين د درېيو~\LR{$\TMstroke$} ګانو پر ننوت وچلوو، چې
هلته به ماشين تل روان وي۔
\olpseganchor{OLP-0255-B013}{end}

\olpseganchor{OLP-0255-B014}{start}
ماشين په حالت~\LR{$q_0$} کښې، د تر ټولو کيڼ~\LR{$\TMstroke$} په لوستلو پيل
کوي۔ د ماشين لومړنی حالت داسې ښودلے شو:
\[
\TMendtape \TMstroke_0 \TMstroke \TMstroke \TMstroke \TMblank \ldots
\]
پورته تشکيل روښانه دے۔ لکه چې ليدل کېږي، ماشين په صفرم حالت کښې د
تر ټولو کيڼ~\LR{$\TMstroke$} په لوستلو پيل کوي۔ دا پر لومړي~\LR{$\TMstroke$}
د حالت د نوم په فرعي نښه ښودل کېږي۔
\textbf{د ژباړې د نسخې سمون (OLCMP-040):} سرچينې په دې ځاے کښې د
ماشين حالت په نثر کښې لومړی حالت بللے و، خو ډياګرام او فرعي نښه دواړه
صفرم حالت ټاکي؛ دلته هماغه سم حالت راوړل شوے دے۔
په دې ځاے کښې د عملي کېدو وړ لارښوونه
\LR{$\delta(q_0, \TMstroke) = \tuple{q_1, \TMstroke, \TMright}$} ده، نو
ماشين پر پټه ښي لوري ته ځي او حالت~\LR{$q_1$} ته اوړي۔
\[
\TMendtape \TMstroke \TMstroke_1 \TMstroke \TMstroke \TMblank \ldots
\]
څرنګه چې ماشين اوس په حالت~\LR{$q_1$} کښې يو~\LR{$\TMstroke$} لولي، بايد
لارښوونه~\LR{$\delta(q_1, \TMstroke) = \tuple{q_0,
  \TMstroke, \TMright}$} ``تعقيب'' کړو۔ په نتيجه کښې دا تشکيل لاس ته
راځي:
\[
\TMendtape \TMstroke \TMstroke \TMstroke_0 \TMstroke \TMblank \ldots
\]
د ماشين د دوام سره قواعد يو ځل بيا په هماغه ترتيب عملي کېږي، او
لاندې دوه تشکيلونه لاس ته راځي:
\[
\TMendtape \TMstroke \TMstroke \TMstroke \TMstroke_1 \TMblank \ldots
\]
\[
\TMendtape \TMstroke \TMstroke \TMstroke \TMstroke \TMblank_0 \ldots
\]
ماشين اوس په حالت~\LR{$q_0$} کښې يو~\LR{$\TMblank$} لولي۔ د حالت بدلون د
ډياګرام له مخې په اسانه وينو چې د عملي کولو هېڅ لارښوونه نشته، نو
ماشين درېدلے دے۔ معنا دا چې ننوت منل شوے دے۔
\olpseganchor{OLP-0255-B014}{end}

\olpseganchor{OLP-0255-B015}{start}
اوس فرض کړئ ماشين د درېيو~\LR{$\TMstroke$} ګانو پر ننوت پيلوو۔ لومړني
څو تشکيلونه ورته دي، ځکه هماغه لارښوونې عملي کېږي او يوازې د پټې په
ننوت کښې لږ توپير شته:
\[
\TMendtape \TMstroke_0 \TMstroke \TMstroke \TMblank \ldots
\]
\[
\TMendtape \TMstroke \TMstroke_1 \TMstroke \TMblank \ldots
\]
\[
\TMendtape \TMstroke \TMstroke \TMstroke_0 \TMblank \ldots
\]
\[
\TMendtape \TMstroke \TMstroke \TMstroke \TMblank_1 \ldots
\]
ماشين اوس له ټولو~\LR{$\TMstroke$} ګانو څخه تېر شوے، او په حالت~\LR{$q_1$}
کښې يو~\LR{$\TMblank$} لولي۔ لکه چې په ډياګرام کښې ښودل شوې، د
\LR{$\delta(q_1, \TMblank) =\tuple{q_1,
\TMblank, \TMright}$} په بڼه يوه لارښوونه شته۔ څرنګه چې پټه ښي لوري
ته تر نامعلومې مودې پورې له~\LR{$\TMblank$} څخه ډکه ده، ماشين به دا
لارښوونه \emph{تل} عملي کوي، په حالت~\LR{$q_1$} کښې به پاتې کېږي او لا
پسې به ښي لوري ته ځي۔ ماشين به هېڅکله ونۀ درېږي، او ننوت نۀ مني۔
\end{explain}
\olpseganchor{OLP-0255-B015}{end}

\olpseganchor{OLP-0255-B016}{start}
\begin{explain}
دا خبره مهمه ده چې ټول ماشينونه نۀ درېږي۔ که درېدل په دې معنا وي چې
د ماشين د عملي کولو لارښوونې خلاصې شي، نو داسې ماشين جوړولے شو چې
هېڅکله نۀ درېږي: بس دا باوري کړو چې په هر حالت کښې د هرې نښې دپاره
يو وتونکے غشی شته۔ په حالت~\LR{$q_0$} کښې د~\LR{$\TMblank$} د لوستلو يوه
لارښوونه په ور زياتولو جفت ماشين داسې بدلولے شو چې تل روان وي۔
\end{explain}
\olpseganchor{OLP-0255-B016}{end}

\olpseganchor{OLP-0255-B017}{start}
\begin{ex}
\[
\begin{tikzpicture}[->,>=stealth',shorten >=1pt,auto,node distance=2.8cm,
                    semithick]
  \tikzstyle{every state}=[fill=none,draw=black,text=black]
\olpseganchor{OLP-0255-B017}{end}

\olpseganchor{OLP-0255-B018}{start}
  \node[initial,state] (A)              {\LR{$q_0$}};
  \node[state]         (B) [right of=A] {\LR{$q_1$}};
\olpseganchor{OLP-0255-B018}{end}

\olpseganchor{OLP-0255-B019}{start}
  \path
  (A) edge [bend left] node {\TMtrans{\TMstroke}{\TMstroke}{\TMright}} (B)
      edge [loop above] node {\TMtrans{\TMblank}{\TMblank}{\TMright}} (A)
  (B) edge [loop above] node {\TMtrans{\TMblank}{\TMblank}{\TMright}} (B)
      edge [bend left] node {\TMtrans{\TMstroke}{\TMstroke}{\TMright}} (A);
\end{tikzpicture}
\]
\end{ex}
\olpseganchor{OLP-0255-B019}{end}

\olpseganchor{OLP-0255-B020}{start}
\begin{explain}
د ماشين جدولونه د ټيورينګ ماشينونو د څرګندولو بله لار ده۔ د ماشين
جدولونه د پټې الفبا پر~\LR{$x$}-محور، او د ماشين د حالتونو سټ پر~\LR{$y$}-محور
ښيي۔ د جدول دننه، د هر حالت او نښې په ګډ ځاے کښې د لارښوونې پاتې
برخه---نوې نښه، نوی حالت او د خوځښت لوری---ليکل کېږي۔ د ماشين جدولونه
دا اسانه کوي چې ماشين په کوم حالت او د کومې نښې دپاره درېږي۔ په جدول
کښې هر تش ځاے د ماشين د درېدو يو ممکن ټکے دے۔ د حالت د ډياګرامونو او
د لارښوونو د سټونو پر خلاف، چې د ماشين د درېدو ځايونه پکښې تل سملاسي
څرګند نۀ وي، د ماشين په جدول کښې د تشو ځايونو په موندلو هر د درېدو
ټکے ژر پېژندل کېږي۔
\textbf{د ژباړې د نسخې سمون (OLCMP-041):} سرچينې د جدول د حجرې
توکي د نوي حالت او نوې نښې په سرچپه ترتيب ياد کړي وو، حال دا چې
ښودل شوې حجرې لومړے نوې نښه او بيا نوی حالت ليکي؛ د هر تش ځاے په
اړه ورپسې جمله کښې ورکه تړونکې کلمه هم ور زياته شوه۔
\end{explain}
\olpseganchor{OLP-0255-B020}{end}

\olpseganchor{OLP-0255-B021}{start}
\begin{ex}
د جفت ماشين جدول دا دے:
\[
\centering
\begin{tabular}{lllll}
\cline{2-4}
\multicolumn{1}{l|}{}      & \multicolumn{1}{c|}{\LR{$\TMblank$}}
& \multicolumn{1}{c|}{\LR{$\TMstroke$}}     & \multicolumn{1}{c|}{\LR{$\TMendtape$}}   &  \\ \cline{1-4}
\multicolumn{1}{|l|}{\LR{$q_0$}} & \multicolumn{1}{l|}{}
& \multicolumn{1}{l|}{\TMtrans{\TMstroke}{q_1}{\TMright}} &
\multicolumn{1}{l|}{\phantom{\TMtrans{\TMstroke}{q_1}{\TMright}}} &  \\ \cline{1-4}
\multicolumn{1}{|l|}{\LR{$q_1$}} & \multicolumn{1}{l|}{\TMtrans{\TMblank}{q_1}{\TMright}}
& \multicolumn{1}{l|}{\TMtrans{\TMstroke}{q_0}{\TMright}} &
\multicolumn{1}{l|}{\phantom{\TMtrans{\TMstroke}{q_1}{\TMright}}} &  \\ \cline{1-4}
\end{tabular}
\]
لکه چې وينو، ماشين هغه وخت درېږي چې په حالت~\LR{$q_0$} کښې يو~\LR{$\TMblank$} لولي۔
\end{ex}
\olpseganchor{OLP-0255-B021}{end}

\olpseganchor{OLP-0255-B022}{start}
\begin{explain}
تر اوسه مو يوازې هغه ماشينونه کتلي چې ننوت لولي او مني ئې۔ خو
ټيورينګ ماشينونه د لوستلو او ليکلو دواړو وړتيا لري۔ د داسې ماشين يوه
بېلګه (که څه هم ډېرې، ډېرې بېلګې شته) \emph{دوه‌چنده کوونکے} دے۔
دوه‌چنده کوونکے چې پر پټه د~\LR{$n$} شمېر~\LR{$\TMstroke$} ګانو له يوې ډلې سره
پيل شي، د~\LR{$2n$} شمېر~\LR{$\TMstroke$} ګانو يوه ډله د وت په توګه ورکوي۔
\end{explain}
\olpseganchor{OLP-0255-B022}{end}

\olpseganchor{OLP-0255-B023}{start}
\begin{ex}
  \ollabel{ex:doubler} مخکښې له دې چې دوه‌چنده کوونکے ماشين جوړ کړو،
  د مسئلې د حل دپاره يوه \emph{تګلاره} جوړول مهم دي۔ څرنګه چې ماشين
  (لکه چې موږ تعريف کړے) دا نۀ شي ياد ساتلے چې څو~\LR{$\TMstroke$} ګانې
  ئې لوستې دي، بايد د پټې د ټولو~\LR{$\TMstroke$} ګانو د شمېر ساتلو يوه
  لار پيدا کړو۔ يوه داسې لار دا ده چې وت له ننوت څخه په يو~\LR{$\TMblank$}
  جلا کړو۔ بيا ماشين له ننوت څخه لومړے~\LR{$\TMstroke$} پاکولے، د ننوت پر
  پاتې برخې تېرېدلے، يو~\LR{$\TMblank$} پرېښودلے، او دوه نوي~\LR{$\TMstroke$}
  ګانې ليکلے شي۔ بيا به ماشين بېرته لاړ شي، په ننوت کښې به دويم
  \LR{$\TMstroke$} پيدا کړي، او هغه به هم دوه‌چنده کړي۔ د ننوت د هر يو
  \LR{$\TMstroke$} په بدل کښې به د وت دوه~\LR{$\TMstroke$} ګانې وليکي۔ د
  ماشين د تګ په ترڅ کښې د ننوت په پاکولو باوري کولے شو چې هېڅ
  \LR{$\TMstroke$} پاتې نۀ شي او نۀ دوه ځله دوه‌چنده شي۔ کله چې ټول ننوت
  پاک شي، پر پټه به~\LR{$2n$} شمېر~\LR{$\TMstroke$} ګانې پاتې وي۔ د لاس ته
  راغلي ټيورينګ ماشين د حالت ډياګرام په~\olref{fig:doubler} کښې
  ښودل شوے دے۔
  \begin{figure}
\[
\begin{tikzpicture}[->,>=stealth',shorten >=1pt,auto,node distance=2.8cm,
                    semithick]
  \tikzstyle{every state}=[fill=none,draw=black,text=black]
\olpseganchor{OLP-0255-B023}{end}

\olpseganchor{OLP-0255-B024}{start}
  \node[initial,state] (1)              {\LR{$q_0$}};
  \node[state]         (2) [right of=1] {\LR{$q_1$}};
  \node[state]         (3) [right of=2] {\LR{$q_2$}};
  \node[state]         (4) [below of=3] {\LR{$q_3$}};
  \node[state]         (5) [left of=4]  {\LR{$q_4$}};
  \node[state]         (6) [left of=5]  {\LR{$q_5$}};
\olpseganchor{OLP-0255-B024}{end}

\olpseganchor{OLP-0255-B025}{start}
  \path (1) edge node {\TMtrans{\TMstroke}{\TMblank}{\TMright}} (2)
    (2) edge [loop above] node {\TMtrans{\TMstroke}{\TMstroke}{\TMright}} (2)
      edge node {\TMtrans{\TMblank}{\TMblank}{\TMright}} (3)
    (3) edge [loop above] node {\TMtrans{\TMstroke}{\TMstroke}{\TMright}} (3)
        edge node {\TMtrans{\TMblank}{\TMstroke}{\TMright}} (4)
    (4) edge [loop below] node {\TMtrans{\TMblank}{\TMstroke}{\TMleft}} (4)
        edge node {\TMtrans{\TMstroke}{\TMstroke}{\TMleft}} (5)
    (5) edge [loop below]  node {\TMtrans{\TMstroke}{\TMstroke}{\TMleft}} (5)
        edge              node {\TMtrans{\TMblank}{\TMblank}{\TMleft}} (6)
    (6) edge [loop below] node {\TMtrans{\TMstroke}{\TMstroke}{\TMleft}} (6)
        edge              node {\TMtrans{\TMblank}{\TMblank}{\TMright}} (1);
\end{tikzpicture}
\]
\caption{يو دوه‌چنده کوونکے ماشين}
\ollabel{fig:doubler}
\end{figure}
\end{ex}
\olpseganchor{OLP-0255-B025}{end}

\olpseganchor{OLP-0255-B026}{start}
\begin{prob}
يو خپل‌سری ننوت وټاکئ او په~\olref[tur][mac][rep]{ex:doubler} کښې
د دوه‌چنده کوونکي ماشين تشکيلونه پرله‌پسې وڅارئ۔
\end{prob}
\olpseganchor{OLP-0255-B026}{end}

\olpseganchor{OLP-0255-B027}{start}
\begin{prob}
د~\LR{$\{\TMendtape,\TMblank, A, B\}$} الفبا لرونکے ټيورينګ ماشين جوړ
کړئ چې د~\LR{$A$} ګانو او~\LR{$B$} ګانو هره هغه رشته مني، يعنې پرې درېږي، چې
د~\LR{$A$} ګانو شمېر ئې د~\LR{$B$} ګانو له شمېر سره برابر وي \emph{او} ټول
\LR{$A$} ګانې له ټولو~\LR{$B$} ګانو مخکښې راشي؛ او هره هغه رشته ردوي، يعنې
پرې نۀ درېږي، چې د~\LR{$A$} ګانو شمېر د~\LR{$B$} ګانو له شمېر سره برابر نۀ وي
يا ټولې~\LR{$A$} ګانې له ټولو~\LR{$B$} ګانو مخکښې رانۀ شي۔ (مثلاً، ماشين بايد
\LR{$AABB$} او~\LR{$AAABBB$} ومني، خو~\LR{$AAB$} او~\LR{$AABBAABB$} دواړه رد کړي۔)
\end{prob}
\olpseganchor{OLP-0255-B027}{end}

\olpseganchor{OLP-0255-B028}{start}
\begin{prob}
د~\LR{$\{\TMendtape,\TMblank, A, B\}$} الفبا لرونکے ټيورينګ ماشين جوړ
کړئ چې د~\LR{$A$} ګانو او~\LR{$B$} ګانو هره رشته~\LR{$\alpha$} د ننوت په توګه
واخلي او نقل ئې کړي، څو د~\LR{$\alpha\alpha$} په بڼه وت پيدا کړي۔ (مثلاً،
د~\LR{$ABBA$} ننوت بايد د~\LR{$ABBAABBA$} وت ورکړي۔)
\end{prob}
\olpseganchor{OLP-0255-B028}{end}

\olpseganchor{OLP-0255-B029}{start}
\begin{prob}
\emph{الفبايي؟}: د~\LR{$\{\TMendtape,\TMblank, A, B\}$} الفبا لرونکے
ټيورينګ ماشين جوړ کړئ چې د~\LR{$A$} ګانو او~\LR{$B$} ګانو يوه متناهي لړۍ د
ننوت په توګه ورکړل شي، نو وګوري چې ټولې~\LR{$A$} ګانې د ټولو~\LR{$B$} ګانو
کيڼ لوري ته راځي که نه۔ ماشين بايد د ننوت رشته پر پټه پرېږدي، او که
رشته ``الفبايي'' وي ودرېږي، او که نۀ وي تل په کړۍ کښې وچلېږي۔
\end{prob}
\olpseganchor{OLP-0255-B029}{end}

\olpseganchor{OLP-0255-B030}{start}
\begin{prob}
\emph{الفبايي کوونکے}: د~\LR{$\{\TMendtape,\TMblank, A, B\}$} الفبا
لرونکے ټيورينګ ماشين جوړ کړئ چې د~\LR{$A$} ګانو او~\LR{$B$} ګانو يوه متناهي
لړۍ د ننوت په توګه واخلي او داسې ئې بيا ترتيب کړي چې ټولې~\LR{$A$} ګانې
د ټولو~\LR{$B$} ګانو کيڼ لوري ته راشي۔ (مثلاً، لړۍ~\LR{$BABAA$} بايد په
\LR{$AAABB$}، او لړۍ~\LR{$ABBABB$} بايد په~\LR{$AABBBB$} بدله شي۔)
\textbf{د ژباړې د نسخې سمون (OLCMP-042):} سرچينې د ننوت د اخيستلو
او بيا ترتيبولو تر منځ د عطف کلمه پرېښې وه؛ دلته بشپړه جمله راوړل
شوې ده۔
\end{prob}
\olpseganchor{OLP-0255-B030}{end}

\olpseganchor{OLP-0256-B004}{start}
\olfileid{tur}{mac}{tur}
\olsection{ټيورينګ ماشينونه}
\olpseganchor{OLP-0256-B004}{end}

\olpseganchor{OLP-0256-B005}{start}
\begin{explain}
د ټيورينګ ماشين د جوړښت صوري تعريف مجرد ښکاري، خو په اصل کښې ساده
دے: دا يوازې هغه ټول معلومات په يوه رياضياتي جوړښت کښې رانغاړي چې د
ټيورينګ ماشين د کار د ټاکلو دپاره پکار دي۔ په دې کښې دا شامل دي: (1)
ماشين په کومو حالتونو کښې اوسېدلے شي، (2) پر پټه کومې نښې روا دي، (3)
ماشين بايد په کوم حالت کښې پيل وکړي، او (4) د ماشين د لارښوونو سټ څه
دے۔
\end{explain}
\olpseganchor{OLP-0256-B005}{end}

\olpseganchor{OLP-0256-B006}{start}
\begin{defn}[ټيورينګ ماشين]
يو \emph{ټيورينګ ماشين} \LR{$M$} هغه څلور‌يز \LR{$\langle Q, \Sigma, q_0,
\delta\rangle$} دے چې له لاندې توکو جوړ دے:
\begin{enumerate}
\item د \emph{حالتونو} يو متناهي سټ~\LR{$Q$}،
\item يوه متناهي \emph{الفبا} \LR{$\Sigma$} چې \LR{$\TMendtape$} او
  \LR{$\TMblank$} پکښې شامل دي،
\item يو \emph{لومړنی حالت}~\LR{$q_0 \in Q$}،
\item د \emph{لارښوونو} يو متناهي سټ~\LR{$\delta\colon Q \times \Sigma
  \pto Q \times \Sigma \times \{\TMleft, \TMright, \TMstay\}$}۔
\end{enumerate}
جزوي تابع~\LR{$\delta$} ته د~\LR{$M$} \emph{د حالت بدلون تابع} هم ويل کېږي۔
\end{defn}
\olpseganchor{OLP-0256-B006}{end}

\olpseganchor{OLP-0256-B007}{start}
\begin{explain}
موږ فرض کوو چې پټه يوازې په يوه لوري کښې نامتناهي ده۔ ځکه نو ګټوره
ده چې ځانګړې نښه~\LR{$\TMendtape$} د پټې د کيڼ پای د نښې په توګه وټاکو۔
په دې سره د ټيورينګ ماشين پروګرامونه په اسانۍ پوهېدے شي چې کله له پټې
څخه د ``لوېدو په خطر'' کښې دي۔ موږ دا فرض کولے شو چې دا نښه هېڅکله
نۀ بدلېږي، يعنې که \LR{$\delta(q,\TMendtape)$} تعريف شوې وي، نو
\LR{$\delta(q,\TMendtape) = \tuple{q', \TMendtape, x}$} وي۔ ځينې درسي
کتابونه دا کار کوي؛ موږ ئې نۀ کوو۔ د خپل ټيورينګ ماشين د جوړولو پر
وخت يوازې دومره پام وکړئ چې~\LR{$\TMendtape$} هېڅکله بدله نۀ کړي۔ سربېره
پر دې، داسې حالتونه هم شته چې د دې نښې د بدلولو اجازه څه ګټور
انعطاف برابروي۔
\end{explain}
\olpseganchor{OLP-0256-B007}{end}

\olpseganchor{OLP-0256-B008}{start}
\begin{ex}
\emph{جفت ماشين:} جفت ماشين په صوري ډول هغه څلور‌يز
\LR{$\tuple{Q, \Sigma, q_0, \delta}$} دے چې
\begin{align*}
Q & = \{ q_0, q_1 \} \\
\Sigma & = \{ \TMendtape, \TMblank, \TMstroke \}, \\
\delta(q_0, \TMstroke) & = \tuple{q_1, \TMstroke, \TMright},\\
\delta(q_1, \TMstroke) & = \tuple{q_0, \TMstroke, \TMright},\\
\delta(q_1, \TMblank)  & = \tuple{q_1, \TMblank, \TMright}.
\end{align*}
\end{ex}
\olpseganchor{OLP-0256-B008}{end}

\olpseganchor{OLP-0257-B004}{start}
\olfileid{cmp}{tur}{con}
\olsection{تشکيلونه او محاسبې}
\olpseganchor{OLP-0257-B004}{end}

\olpseganchor{OLP-0257-B005}{start}
\begin{explain}
هغه مخکښېنی بحث را ياد کړئ چې د جفت ماشين تشکيلونه مو پکښې يو په بل
پسې څارل۔ د پټې، لوست-ليک سر او ټيورينګ ماشين پروګرام خيالي ميکانيزم
په حقيقت کښې يوازې د ټيورينګ ماشين د محاسبې د ليدلو يوه شهودي لار ده۔
په صوري ډول، پر يوه ټاکلي ننوت د ټيورينګ ماشين محاسبه د
\emph{تشکيلونو} د يوې لړۍ په توګه تعريفولے شو---او تشکيل په خپل وار د
نښو يوه لړۍ (چې د محاسبې په يوه ټاکلي پړاو کښې د پټې له منځپانګې سره
سمون لري)، د لوست-ليک سر د ځاے ښودونکے يو عدد، او يو حالت دے۔ د دې
توکو په وسيله ټاکلے شو چې ټيورينګ ماشين~\LR{$M$} پر يوه ټاکلي ننوت څه
محاسبه کوي۔
\end{explain}
\olpseganchor{OLP-0257-B005}{end}

\olpseganchor{OLP-0257-B006}{start}
\begin{defn}[تشکيل]
د ټيورينګ ماشين \LR{$M = \tuple{Q, \Sigma, q_0,
\delta}$} يو \emph{تشکيل} هغه درې‌يز \LR{$\tuple{C, m, q}$} دے چې
\begin{enumerate}
\item \LR{$C \in \Sigma^*$} د~\LR{$\Sigma$} د نښو يوه متناهي لړۍ ده،
\item \LR{$m \in \Nat$} يو داسې عدد دے چې \LR{$< \len{C}$}، او
\item \LR{$q \in Q$}۔
\end{enumerate}
په شهودي ډول، لړۍ~\LR{$C$} د پټې منځپانګه ده (له تر ټولو کيڼ مربع څخه تر
وروستي ناتش يا مخکښې ليدل شوي مربع پورې د ټولو مربعونو نښې)، \LR{$m$} د
هغه مربع شمېره ده چې لوست-ليک سر ئې لولي (تر ټولو کيڼ مربع په~\LR{$0$}
شمېرلو سره)، او~\LR{$q$} د ماشين اوسنی حالت دے۔
\end{defn}
\olpseganchor{OLP-0257-B006}{end}

\olpseganchor{OLP-0257-B007}{start}
\begin{explain}
د ټيورينګ ماشين ممکن ننوت د نښو يوه لړۍ ده، او عموماً داسې لړۍ وي چې
يو عدد په يوه بڼه کوډوي۔ د ټيورينګ ماشين لومړنی تشکيل هغه تشکيل دے
چې ماشين پکښې پر همدغه ننوت کار پيلوي: پټه د پټې د پای نښه لري او
سملاسي ورپسې ښي لوري ته پر مربعونو ننوت ليکل شوے وي، لوست-ليک سر د
ننوت تر ټولو کيڼ مربع (يعنې د کيڼ پای د نښې ښي لوري ته مربع) لولي،
او ميکانيزم په ټاکلي لومړني حالت~\LR{$q_0$} کښې وي۔
\end{explain}
\olpseganchor{OLP-0257-B007}{end}

\olpseganchor{OLP-0257-B008}{start}
\begin{defn}[لومړنی تشکيل]
د ناتش ننوت \LR{$I \in \Sigma^*$} دپاره د~\LR{$M$} \emph{لومړنی تشکيل} دا دے:
\[
\tuple{\TMendtape \frown I, 1, q_0}.
\]
که ننوت تشه لړۍ وي، لومړنی تشکيل دا دے:
\[
\tuple{\TMendtape \frown \TMblank, 1, q_0}.
\]
\textbf{د ژباړې د نسخې سمون (OLCMP-059):} سرچينه د هرې ممکنې
ننوت-لړۍ دپاره لومړنی تشکيل يوازې د پای نښې او ننوت په يوځاے کولو
جوړوي؛ د تشې ننوت په حالت کښې دې سره د سر شمېره د لړۍ له اوږدوالي
سره برابره شي او د تشکيل شرط ماتېږي۔ دلته د لوستل کېدونکې تشې مربع
په ور زياتولو تشکيل بشپړ شوے دے۔
\end{defn}
\olpseganchor{OLP-0257-B008}{end}

\olpseganchor{OLP-0257-B009}{start}
\begin{explain}
نښه~\LR{$\frown$} د \emph{يوځاے کولو} دپاره ده---د ننوت رشته د کيڼ پای د
نښې سملاسي ښي لوري ته پيلېږي۔
\textbf{د ژباړې د نسخې سمون (OLCMP-045):} د سرچينې عبارت دلته نحوي
نيمګړے دے او د ننوت او د کيڼ پای نښې تر منځ سم لوريز تړاو نۀ
څرګندوي؛ هم پورته نثر او هم صوري لومړنی تشکيل ننوت د نښې ښي لوري ته
ږدي، نو دلته همدغه ټاکلے لوری راوړل شوے دے۔
\end{explain}
\olpseganchor{OLP-0257-B009}{end}

\olpseganchor{OLP-0257-B010}{start}
\begin{defn}
وايو چې تشکيل \LR{$\tuple{C, m, q}$} \emph{په يوه ګام کښې تشکيل
\LR{$\tuple{C', m', q'}$} ورکوي} (د~\LR{$M$} له مخې)، هله او يوازې هله چې
\begin{enumerate}
\item د~\LR{$C$} \LR{$m$}-مه نښه~\LR{$\sigma$} وي،
\item د~\LR{$M$} د لارښوونو سټ وټاکي چې \LR{$\delta(q, \sigma) =
  \tuple{q', \sigma', D}$}،
\item د~\LR{$C'$} \LR{$m$}-مه نښه~\LR{$\sigma'$} وي، او
\item
\begin{enumerate}
\item که \LR{$m>0$} وي، \LR{$D = L$} او \LR{$m' = m - 1$} وي، او که نۀ وي \LR{$m'=0$}، يا
\item \LR{$D = R$} او \LR{$m' = m + 1$} وي، يا
\item \LR{$D = N$} او \LR{$m' = m$} وي،
\end{enumerate}
\item که \LR{$m' = \len{C}$} وي، نو \LR{$\len{C'} = \len{C} + 1$} او د~\LR{$C'$}
  \LR{$m'$}-مه نښه~\LR{$\TMblank$} ده۔ که نۀ وي، \LR{$\len{C'}=\len{C}$}۔
\item د هر داسې~\LR{$i$} دپاره چې \LR{$i < \len{C}$} او \LR{$i \neq m$} وي، \LR{$C'(i) = C(i)$}،
\end{enumerate}
\end{defn}
\olpseganchor{OLP-0257-B010}{end}

\olpseganchor{OLP-0257-B011}{start}
\begin{defn}\ollabel{defn:run-output}
پر ننوت~\LR{$I$} د~\LR{$M$} يوه \emph{چلونه} د~\LR{$M$} د تشکيلونو داسې لړۍ~\LR{$C_i$}
ده چې~\LR{$C_0$} د ننوت~\LR{$I$} دپاره د~\LR{$M$} لومړنی تشکيل وي، او هر~\LR{$C_i$} په
يوه ګام کښې~\LR{$C_{i+1}$} ورکوي۔
\olpseganchor{OLP-0257-B011}{end}

\olpseganchor{OLP-0257-B012}{start}
وايو چې~\LR{$M$} \emph{پر ننوت~\LR{$I$} له~\LR{$k$} ګامونو وروسته درېږي} که \LR{$C_k =
\tuple{C, m, q}$} وي، د~\LR{$C$} \LR{$m$}-مه نښه~\LR{$\sigma$} وي، او \LR{$\delta(q,
\sigma)$} تعريف شوې نۀ وي۔ په دې صورت کښې د ننوت~\LR{$I$} دپاره د~\LR{$M$}
\emph{وت}~\LR{$O$} هله ټاکل کېږي چې~\LR{$O$} د نښو داسې رشته وي چې
په~\LR{$\TMblank$} نۀ ختمېږي او د کوم~\LR{$j \in \Nat$} دپاره \LR{$C =
\TMendtape \concat O \concat \TMblank^j$} وي۔ که داسې رشته او طبيعي
عدد شتون ونۀ لري، دې درېدو ته وت نۀ ټاکل کېږي۔ (\LR{$\TMblank^j$} د~\LR{$j$}
شمېر~\LR{$\TMblank$} نښو يوه لړۍ ده۔)
\textbf{د ژباړې د نسخې سمون (OLCMP-060):} سرچينه هرې درېدلې چلونې
ته وت ټاکي، که څه هم مخکښې د پټې د پای نښې بدلول روا شوي دي؛ داسې
بدلون ښايي د وت د اړينې بڼې شتون ناممکن کړي۔ دلته وت يوازې هغه وخت
ټاکل کېږي چې دغه بڼه په رښتيا شتون ولري۔
\end{defn}
\olpseganchor{OLP-0257-B012}{end}

\olpseganchor{OLP-0257-B013}{start}
\begin{explain}
د دې تعريف له مخې د~\LR{$M$} هر ټاکل شوی وت~\LR{$O$} تل له~\LR{$\TMblank$} پرته په
بلې نښې ختمېږي؛ يا که په وخت~\LR{$k$} کښې ټوله پټه له~\LR{$\TMblank$} څخه ډکه
وي (له تر ټولو کيڼې~\LR{$\TMendtape$} پرته)، نو~\LR{$O$} تشه رشته ده۔ هغه
درېدلې چلونه چې د پای نښې له اړينې بڼې سره نۀ وي، هېڅ وت نۀ لري۔
\end{explain}
\olpseganchor{OLP-0257-B013}{end}

\olpseganchor{OLP-0258-B004}{start}
\olfileid{tur}{mac}{una}
\olsection{د عددونو يوګونې څرګندونه}
\olpseganchor{OLP-0258-B004}{end}

\olpseganchor{OLP-0258-B005}{start}
\begin{explain}
ټيورينګ ماشينونه پر خپله پټه د ليکل شوو نښو پر لړيو کار کوي۔ د
ټيورينګ ماشين د الفبا له مخې، د نښو دا لړۍ بېلابېل ننوتونه او وتونه
څرګندولے شي۔ البته، هغه ټيورينګ ماشينونه ځانګړې پاملرنه غواړي چې
\emph{حسابي} تابعې، يعنې د طبيعي عددونو تابعې، محاسبه کوي۔ د مثبتو
صحيح عددونو د څرګندولو يوه ساده لار دا ده چې د يوې نښې~\LR{$\TMstroke$}
د لړيو په توګه کوډ شي۔ که \LR{$n \in \Nat$} وي، نو \LR{$\TMstroke^n$} دې هغه
وخت تشه لړۍ وي چې \LR{$n = 0$} وي؛ او که نۀ وي، هغه لړۍ دې وي چې کره
\LR{$n$} شمېر~\LR{$\TMstroke$} نښې لري۔
\end{explain}
\olpseganchor{OLP-0258-B005}{end}

\olpseganchor{OLP-0258-B006}{start}
\begin{defn}[محاسبه]
ټيورينګ ماشين~\LR{$M$} هله او يوازې هله تابع \LR{$f\colon \Nat^k \to \Nat$}
\emph{محاسبه کوي} چې~\LR{$M$} پر دې ننوت ودرېږي:
\[
\TMstroke^{n_1} \TMblank \TMstroke^{n_2} \TMblank \dots \TMblank \TMstroke^{n_k}
\]
او وت ئې \LR{$\TMstroke^{f(n_1, \dots, n_k)}$} وي۔
\end{defn}
\olpseganchor{OLP-0258-B006}{end}

\olpseganchor{OLP-0258-B007}{start}
\begin{prob}
  دا تعريف کړئ چې ټيورينګ ماشين~\LR{$M$} کله تابع
  \LR{$f\colon \Nat^k \to \Nat^m$} محاسبه کوي۔
\end{prob}
\olpseganchor{OLP-0258-B007}{end}

\olpseganchor{OLP-0258-B008}{start}
\begin{ex}\ollabel{ex:adder}
\emph{جمع:}
راځئ داسې ماشين جوړ کړو چې تابع \LR{$f(n,m) = n + m$} محاسبه کړي۔ د دې
دپاره داسې ماشين پکار دے چې پر پټه د~\LR{$n$} او~\LR{$m$} اوږدوالي لرونکو د
\LR{$\TMstroke$} دوو غونډونو سره پيل وکړي، او د~\LR{$n+m$} شمېر~\LR{$\TMstroke$}
ګانو په يوه غونډ ودرېږي۔ د~\LR{$\TMstroke$} د ننوت دواړه غونډونه په يوه
\LR{$\TMblank$} سره بېلېږي؛ نو يوه لار دا ده چې په هغه مربع کښې يوه کرښه
وليکو چې~\LR{$\TMblank$} لري، او وروستۍ~\LR{$\TMstroke$} پاکه کړو۔
\begin{figure}\[
\begin{tikzpicture}[->,>=stealth',shorten >=1pt,auto,node distance=2.8cm,
                    semithick]
  \tikzstyle{every state}=[fill=none,draw=black,text=black]
\olpseganchor{OLP-0258-B008}{end}

\olpseganchor{OLP-0258-B009}{start}
  \node[initial,state]         (A)              {\LR{$q_0$}};
  \node[state]         (B) [right of=A] {\LR{$q_1$}};
  \node[state]         (C) [right of=B] {\LR{$q_2$}};
\olpseganchor{OLP-0258-B009}{end}

\olpseganchor{OLP-0258-B010}{start}
  \path (A) edge node {\TMtrans{\TMblank}{\TMstroke}{\TMstay}} (B)
           edge [loop above] node {\TMtrans{\TMstroke}{\TMstroke}{\TMright}} (B)
        (B) edge [loop above] node {\TMtrans{\TMstroke}{\TMstroke}{\TMright}} (B)
            edge node {\TMtrans{\TMblank}{\TMblank}{\TMleft}} (C)
        (C) edge [loop above] node {\TMtrans{\TMstroke}{\TMblank}{\TMstay}} (C);
\end{tikzpicture}
\]\caption{هغه ماشين چې \LR{$f(x,y) = x+y$} محاسبه کوي}
\ollabel{fig:adder}
\end{figure}
\end{ex}
\olpseganchor{OLP-0258-B010}{end}

\olpseganchor{OLP-0258-B011}{start}
\begin{prob}
په~\olref[tur][mac][una]{ex:adder} کښې د ماشين تشکيلونه د
ننوت~\LR{$\tuple{3,2}$} دپاره يو په بل پسې وڅارئ۔ که ماشين \LR{$0+0$} محاسبه
کړي، څه پېښېږي؟
\end{prob}
\olpseganchor{OLP-0258-B011}{end}

\olpseganchor{OLP-0258-B012}{start}
\begin{explain}
په~\olref[rep]{ex:doubler} کښې مو د داسې ټيورينګ ماشين بېلګه ورکړه
چې د~\LR{$\TMstroke$} نښو يوه لړۍ د ننوت په توګه اخلي او پر پټه د هغې د
دوه‌چنده شمېر~\LR{$\TMstroke$} نښو په لړۍ درېږي---يعنې دوه‌چنده کوونکے
ماشين۔ خو څرنګه چې په وت کښې د~\LR{$\TMstroke$} نښو د دوه‌چنده غونډ کيڼ
لوري ته~\LR{$\TMblank$} نښې شته، نو دا په حقيقت کښې تابع \LR{$f(x) = 2x$} نۀ
محاسبه کوي، که څه هم ښايي تاسو داسې ګڼلې وي۔ موږ به د دې د سمولو
دوه لارې وښيو۔
\end{explain}
\olpseganchor{OLP-0258-B012}{end}

\olpseganchor{OLP-0258-B013}{start}
\begin{ex}
په~\olref{fig:doubler-disc} کښې ماشين تابع \LR{$f(x) = 2x$} محاسبه کوي۔
دا ماشين د~\olref[rep]{ex:doubler} د ماشين په شان د ننوت د هرې
\LR{$\TMstroke$} په بدل کښې ننوت نۀ پاکوي او په لرې ښي لوري کښې دوه
\LR{$\TMstroke$} نښې نۀ ليکي، بلکې د ننوت د هرې~\LR{$\TMstroke$} دپاره ښي
لوري ته يوه~\LR{$\TMstroke$} ور زياتوي۔ دا بايد په ياد وساتي چې ننوت چېرته
ختمېږي؛ نو د ننوت او ور زياتو شوو کرښو تر منځ يوه~\LR{$\TMblank$} پرېږدي
او په پای کښې ئې په~\LR{$\TMstroke$} ډکوي۔ موږ بايد دا هم ``په ياد'' ولرو
چې په ننوت کښې چېرته يو؛ نو د ننوت په غونډ کښې يوه~\LR{$\TMstroke$} په
لنډمهاله توګه په~\LR{$\TMblank$} بدلوو۔
  \begin{figure}
\[
\begin{tikzpicture}[->,>=stealth',shorten >=1pt,auto,node distance=2.8cm,
                    semithick]
  \tikzstyle{every state}=[fill=none,draw=black,text=black]
\olpseganchor{OLP-0258-B013}{end}

\olpseganchor{OLP-0258-B014}{start}
  \node[initial,state] (0)              {\LR{$q_0$}};
  \node[state]         (1) [above of=0] {\LR{$q_1$}};
  \node[state]         (2) [above of=1] {\LR{$q_2$}};
  \node[state]         (3) [right of=2] {\LR{$q_3$}};
  \node[state]         (4) [below of=3] {\LR{$q_4$}};
  \node[state]         (5) [below of=4] {\LR{$q_5$}};
  \node[state]         (6) [above right of=4] {\LR{$q_6$}};
  \node[state]         (7) [below of=6] {\LR{$q_7$}};
  \node[state]         (8) [below of=7] {\LR{$q_8$}};
\olpseganchor{OLP-0258-B014}{end}

\olpseganchor{OLP-0258-B015}{start}
  \path (0) edge node {\TMtrans{\TMstroke}{\TMblank}{\TMright}} (1)
%    (0) edge node {\TMtrans{\TMblank}{\TMblank}{\TMstay}} (h)
    (1) edge [loop left] node {\TMtrans{\TMstroke}{\TMstroke}{\TMright}} (1)
      edge node {\TMtrans{\TMblank}{\TMblank}{\TMright}} (2)
    (2) edge [loop above] node {\TMtrans{\TMstroke}{\TMstroke}{\TMright}} (2)
        edge node {\TMtrans{\TMblank}{\TMstroke}{\TMleft}} (3)
    (3) edge [loop above] node {\TMtrans{\TMstroke}{\TMstroke}{\TMleft}} (3)
        edge node[left] {\TMtrans{\TMblank}{\TMblank}{\TMleft}} (4)
    (4) edge        node[left] {\TMtrans{\TMstroke}{\TMstroke}{\TMleft}} (5)
    (5) edge [loop below] node {\TMtrans{\TMstroke}{\TMstroke}{\TMleft}} (5)
        edge              node {\TMtrans{\TMblank}{\TMstroke}{\TMright}} (0)
    (4) edge      node[sloped] {\TMtrans{\TMblank}{\TMstroke}{\TMright}} (6)
    (6) edge              node {\TMtrans{\TMblank}{\TMstroke}{\TMright}} (7)
    (7) edge [loop right] node {\TMtrans{\TMstroke}{\TMstroke}{\TMright}} (7)
        edge              node {\TMtrans{\TMblank}{\TMblank}{\TMleft}} (8)
    (8) edge [loop right] node {\TMtrans{\TMstroke}{\TMblank}{\TMstay}} (8);
    \end{tikzpicture}
\]
\caption{هغه ماشين چې \LR{$f(x) = 2x$} محاسبه کوي}
\ollabel{fig:doubler-disc}
\end{figure}
\end{ex}
\olpseganchor{OLP-0258-B015}{end}

\olpseganchor{OLP-0258-B016}{start}
\begin{ex}\ollabel{ex:mover}
د تابع \LR{$f(x) = 2x$} د محاسبې دويمه شونتيا دا ده چې اصلي دوه‌چنده
کوونکے ماشين وساتو، خو په پای کښې داسې حالتونه او لارښوونې ور زياتې
کړو چې د کرښو دوه‌چنده غونډ د پټې تر ټولو کيڼ لوري ته يوسي۔ په
\olref{fig:mover} کښې ماشين همدا وروستۍ برخه ترسره کوي: که دا پر داسې
پټه پيل شي چې لومړے د~\LR{$\TMblank$} نښو يو غونډ او بيا د~\LR{$\TMstroke$}
نښو يو غونډ لري (او سر د~\LR{$\TMblank$} نښو د غونډ په هر ځاے کښې وي)،
نو~\LR{$\TMstroke$} نښې يوه يوه پاکوي او د پټې په پيل کښې ئې ليکي۔ د دې
د معلومولو دپاره چې کار کله بشپړ شو، ماشين لومړے د~\LR{$\TMstroke$} نښو
د غونډ پای په~\LR{$\TMendtape$} نښه نښانوي، او دا نښه په پای کښې پاکېږي۔
موږ د حالتونو شمېرل له~\LR{$q_6$} څخه پيل کړي، څو دوي دوه‌چنده کوونکي
ماشين ته ور زياتېدے شي۔ تاسو يوازې يوې اضافي لارښوونې ته اړتيا لرئ:
\LR{$\delta(q_0, \TMblank) = \tuple{q_6,\TMblank,\TMstay}$}؛ يعنې له~\LR{$q_0$}
څخه~\LR{$q_6$} ته يو داسې غشی چې نښه ئې~\LR{$\TMtrans{\TMblank}{\TMblank}{\TMstay}$}
وي۔ (يوه نرۍ ستونزه شته: لاس ته راغلے ماشين د ننوت~\LR{$x=0$} دپاره کار
نۀ کوي۔ موږ دا د تمرين په توګه پرېږدو۔)
\begin{figure}
  \[
  \begin{tikzpicture}[->,>=stealth',shorten >=1pt,auto,node distance=2.8cm,
                      semithick]
    \tikzstyle{every state}=[fill=none,draw=black,text=black]
    \node[initial,state] (6)              {\LR{$q_6$}};
    \node[state]         (7) [right of=6] {\LR{$q_7$}};
    \node[state]         (8) [right of=7] {\LR{$q_8$}};
    \node[state]         (9) [below of=8] {\LR{$q_9$}};
    \node[state]         (10) [left of=9] {\LR{$q_{10}$}};
    \node[state]         (11) [left of=10]  {\LR{$q_{11}$}};
    \node[state]         (12) [below of=11]  {\LR{$q_{12}$}};
    \node[state]         (13) [right of=12] {\LR{$q_{13}$}};
    \node[state]         (14) [right of=13] {\LR{$q_{14}$}};
    \path
    (6)  edge node {\TMtrans{\TMblank}{\TMblank}{\TMright}} (7)
    (7)  edge [loop above] node {\TMtrans{\TMblank}{\TMblank}{\TMright}} (7)
         edge node {\TMtrans{\TMstroke}{\TMstroke}{\TMright}} (8)
    (8)  edge [loop above] node {\TMtrans{\TMstroke}{\TMstroke}{\TMright}} (8)
         edge node {\TMtrans{\TMblank}{\TMendtape}{\TMleft}} (9)
    (9)  edge [loop right] node {\TMtrans{\TMstroke}{\TMstroke}{\TMleft}} (9)
         edge node {\TMtrans{\TMblank}{\TMblank}{\TMright}} (10)
    (10) edge node[above] {\TMtrans{\TMstroke}{\TMblank}{\TMleft}} (11)
    (11) edge [loop above] node {\TMtrans{\TMblank}{\TMblank}{\TMleft}} (11)
         edge node[left] {\begin{tabular}{@{}l@{}}
          \TMtrans{\TMendtape}{\TMendtape}{\TMright}\\
          \TMtrans{\TMstroke}{\TMstroke}{\TMright}
         \end{tabular}} (12)
    (12) edge node[] {\TMtrans{\TMblank}{\TMstroke}{\TMright}} (13)
    (13) edge [loop below] node {\TMtrans{\TMblank}{\TMblank}{\TMright}} (13)
         edge node[sloped] {\TMtrans{\TMstroke}{\TMblank}{\TMleft}} (11)
    (13) edge node {\TMtrans{\TMendtape}{\TMblank}{\TMstay}} (14);
  \end{tikzpicture}
  \]
  \caption{د~\LR{$\TMstroke$} نښو د يو غونډ کيڼ لوري ته وړل}
  \ollabel{fig:mover}
\end{figure}
\end{ex}
\olpseganchor{OLP-0258-B016}{end}

\olpseganchor{OLP-0258-B017}{start}
\begin{prob}
په~\olref[tur][mac][una]{ex:mover} کښې مو داسې ماشين بيان کړ چې د
\olref[tur][mac][rep]{fig:doubler} د دوه‌چنده کوونکي ماشين او د
\olref[tur][mac][una]{fig:mover} د لېږدوونکي ماشين له يوځاے کولو څخه
جوړ دے۔ که دا يوځاے شوے ماشين پر ننوت~\LR{$x=0$}، يعنې پر تشه پټه پيل
کړئ، څه پېښېږي؟ ماشين به څنګه سم کړئ چې په دې صورت کښې د~\LR{$2x=0$} له
وت سره ودرېږي؟ (تاسو بايد دا د يوه حالت او يوه حالت بدلون په ور
زياتولو وکړے شئ۔)
\textbf{د ژباړې د نسخې سمون (OLCMP-046):} سرچينې دلته د هغه بل
دوه‌چنده کوونکي شکل حواله ورکړې وه چې لا پخپله بشپړ وت جوړوي او
حالتونه ئې له لېږدوونکي سره ټکر کوي؛ د مخکښېني بيان او د حالتونو له
شمېرنې سره سم د اصلي شپږ-حالته دوه‌چنده کوونکي حواله راوړل شوه۔
\end{prob}
\olpseganchor{OLP-0258-B017}{end}

\olpseganchor{OLP-0258-B018}{start}
\begin{prob}
\emph{تفريق:} داسې ټيورينګ ماشين طرح کړئ چې د~\LR{$n$} او~\LR{$m$} اوږدوالي
لرونکې د کرښو دوه ناتشې رشتې د ننوت په توګه ورکړل شي، چې \LR{$n > m$}،
نو تابع \LR{$f(n,m) = n - m$} محاسبه کړي۔
\end{prob}
\olpseganchor{OLP-0258-B018}{end}

\olpseganchor{OLP-0258-B019}{start}
\begin{prob}
\emph{مساوات:} داسې ټيورينګ ماشين طرح کړئ چې لاندې تابع محاسبه کړي:
\[
\fn{equality}(n,m) =
\begin{cases}
  \text{\RL{1}} & \text{\RL{که~\LR{$n = m$} وي}} \\
  \text{\RL{0}} & \text{\RL{که~\LR{$n \neq m$} وي}}
\end{cases}
\]
چې پکښې~\LR{$n$} او~\LR{$m \in \PosInt$} دي۔
\end{prob}
\olpseganchor{OLP-0258-B019}{end}

\olpseganchor{OLP-0258-B020}{start}
\begin{prob}
د تابع \LR{$\min(x,y)$} د محاسبې دپاره يو ټيورينګ ماشين طرح کړئ، چې
\LR{$x$} او~\LR{$y$} داسې مثبت صحيح عددونه وي چې پر پټه د~\LR{$\TMstroke$} نښو په
رشتو څرګند شوي او په يوه~\LR{$\TMblank$} بېل شوي وي۔ تاسو د ماشين په
الفبا کښې اضافي نښې کارولے شئ۔
\olpseganchor{OLP-0258-B020}{end}

\olpseganchor{OLP-0258-B021}{start}
تابع~\LR{$\min$} له خپلو ارزښتمنو څخه تر ټولو کوچنی ارزښت ټاکي؛ نو
\LR{$\min(3,5)=3$}، \LR{$\min(20,16)=16$}، او \LR{$\min(4,4)=4$}، او داسې نور۔
\end{prob}
\olpseganchor{OLP-0258-B021}{end}

\olpseganchor{OLP-0258-B022}{start}
\begin{defn}
  ټيورينګ ماشين~\LR{$M$} هله او يوازې هله جزوي تابع \LR{$f\colon \Nat^k
  \pto \Nat$} محاسبه کوي چې
  \begin{enumerate}
    \item که \LR{$f(n_1, \dots, n_k) = m$} وي، نو~\LR{$M$} پر ننوت
    \LR{$\TMstroke^{n_1}\concat\TMblank\concat
    \dots \concat\TMblank\concat\TMstroke^{n_k}$} د وت~\LR{$\TMstroke^{m}$} سره درېږي۔
    \item که \LR{$f(n_1, \dots, n_k)$} تعريف شوې نۀ وي، نو~\LR{$M$} بېخي نۀ
    درېږي، يا داسې درېږي چې معتبر وت ورته نۀ ټاکل کېږي، يا له داسې
    وت سره درېږي چې د~\LR{$\TMstroke$} نښو يوازينی غونډ نۀ وي۔
  \end{enumerate}
\textbf{د ژباړې د نسخې سمون (OLCMP-061):} د سرچينې دويم شرط دا
فرضوي چې هره درېدلې چلونه يو وت لري۔ د مخکښېني مشروط وت له تعريف
سره د سمون دپاره، هغه درېدنه هم دلته شامله شوه چې معتبر وت ورته نۀ
ټاکل کېږي۔
\end{defn}
\olpseganchor{OLP-0258-B022}{end}

\olpseganchor{OLP-0259-B004}{start}
\olfileid{tur}{mac}{hal}
\olsection{درېدني حالتونه}
\olpseganchor{OLP-0259-B004}{end}

\olpseganchor{OLP-0259-B005}{start}
\begin{explain}
که څه هم موږ خپل ماشينونه داسې تعريف کړي چې يوازې د عملي کېدونکې
لارښوونې د نشتوالي پر وخت درېږي، د ټيورينګ ماشينونو په عامو
څرګندونو کښې يو ځانګړی \emph{درېدنی حالت}~\LR{$h$} وي، داسې چې
\LR{$h \in Q$}۔
\olpseganchor{OLP-0259-B005}{end}

\olpseganchor{OLP-0259-B006}{start}
د درېدني حالت تر شا نظر ساده دے: کله چې ماشين خپل کار بشپړ کړي (د
ننوت منلو ته چمتو وي، يا د وت ليکل ئې بشپړ کړي وي)، داسې حالت~\LR{$h$}
ته ننوځي چې هلته درېږي۔ ځينې ماشينونه دوه درېدني حالتونه لري: يو
ننوت مني او بل ننوت ردوي۔
\end{explain}
\olpseganchor{OLP-0259-B006}{end}

\olpseganchor{OLP-0259-B007}{start}
\begin{ex}\emph{درېدني حالتونه}۔
د دې مفهوم د روښانولو دپاره راځئ د جفت ماشين له يو بدلون څخه پيل
وکړو۔ د دې پر ځاے چې ماشين د جفت ننوت په صورت کښې په حالت~\LR{$q_0$}
کښې ودرېږي، يوه لارښوونه ور زياتولے شو چې ماشين درېدني حالت ته
واستوي۔
\[
\begin{tikzpicture}[->,>=stealth',shorten >=1pt,auto,node distance=2.8cm,
                    semithick]
  \tikzstyle{every state}=[fill=none,draw=black,text=black]
\olpseganchor{OLP-0259-B007}{end}

\olpseganchor{OLP-0259-B008}{start}
  \node[initial, state]         (A)                     {\LR{$q_0$}};
  \node[state]         (B) [right of=A] {\LR{$q_1$}};
  \node[state]         (C) [below of=A] {\LR{$h$}};
\olpseganchor{OLP-0259-B008}{end}

\olpseganchor{OLP-0259-B009}{start}
  \path (A) edge [bend left] node {\TMtrans{\TMstroke}{\TMstroke}{R}} (B)
  	    edge node {\TMtrans{\TMblank}{\TMblank}{N}} (C)
        (B) edge [loop above] node {\TMtrans{\TMblank}{\TMblank}{R}} (B)
            edge [bend left] node {\TMtrans{\TMstroke}{\TMstroke}{R}} (A);
\end{tikzpicture}
\]
\olpseganchor{OLP-0259-B009}{end}

\olpseganchor{OLP-0259-B010}{start}
راځئ بېلګه نوره هم پراخه کړو۔ کله چې ماشين ومومي چې ننوت طاق دے،
هېڅکله نۀ درېږي۔ موږ د کړۍ وهونکې لارښوونې پر ځاے داسې لارښوونه
ايښودلے شو چې رد حالت~\LR{$r$} ته ځي، او په دې توګه ماشين ته يو
\emph{رد} حالت ور زياتولے شو۔
\[
\begin{tikzpicture}[->,>=stealth',shorten >=1pt,auto,node distance=2.8cm,
                    semithick]
  \tikzstyle{every state}=[fill=none,draw=black,text=black]
\olpseganchor{OLP-0259-B010}{end}

\olpseganchor{OLP-0259-B011}{start}
  \node[initial,state]         (A)                     {\LR{$q_0$}};
  \node[state]         (B) [right of=A] {\LR{$q_1$}};
  \node[state]         (C) [below of=A] {\LR{$h$}};
  \node[state]         (D) [below of=B] {\LR{$r$}};
\olpseganchor{OLP-0259-B011}{end}

\olpseganchor{OLP-0259-B012}{start}
  \path (A) edge [bend left] node {\TMtrans{\TMstroke}{\TMstroke}{R}} (B)
  	    edge node {\TMtrans{\TMblank}{\TMblank}{N}} (C)
        (B) edge node {\TMtrans{\TMblank}{\TMblank}{N}} (D)
            edge [bend left] node {\TMtrans{\TMstroke}{\TMstroke}{R}} (A);
\end{tikzpicture}
\]
\end{ex}
\olpseganchor{OLP-0259-B012}{end}

\olpseganchor{OLP-0259-B013}{start}
\begin{explain}
د ځانګړي درېدني حالت زياتول د دې په څېر حالتونو کښې ګټور وي، ځکه په
څرګنده ښيي چې ماشين ټاکلي ننوتونه کله مني يا ردوي۔ خو پام کول مهم دي
چې د ځانګړي درېدني حالت په زياتولو هېڅ محاسبوي ځواک نۀ زياتېږي۔ همداراز
د درېدو لږ صوري مفهوم خپلې ګټې لري۔ د درېدو هغه تعريف چې تر اوسه مو
په دې څپرکي کښې کارولے، د \emph{درېدو مسئلې} ثبوت شهودي او څرګندول ئې
اسانه کوي۔ له همدې امله موږ خپل اصلي تعريف ته دوام ورکوو۔
\end{explain}
\olpseganchor{OLP-0259-B013}{end}

\olpseganchor{OLP-0260-B004}{start}
\olfileid{tur}{mac}{dis}
\olsection{منضبط ماشينونه}
\olpseganchor{OLP-0260-B004}{end}

\olpseganchor{OLP-0260-B005}{start}
\begin{explain}
په برخه~\olref[hal]{sec} کښې مو هغه ټيورينګ ماشينونه وڅېړل چې يو
ټاکلی درېدنی حالت~\LR{$h$} لري---داسې ماشينونه، که بېخي درېږي، هرومرو په
حالت~\LR{$h$} کښې درېږي۔ په دې معنا د يوه درېدني حالت لرونکي ماشينونه د
هغو ټيورينګ ماشينونو په پرتله ډېر ``منضبط'' دي چې موږ ئې په عمومي
ډول منو۔ نور بنديزونه هم د ټيورينګ ماشينونو پر چلند لګولے شو۔ مثلاً،
موږ ټيورينګ ماشينونه له دې هم نۀ دي منع کړي چې کله د پټې د پای نښه
پر مربع~\LR{$0$} پاکه کړي، يا له مربع~\LR{$0$} څخه کيڼ لوري ته د تلو هڅه وکړي۔
(زموږ تعريف وايي چې په دې صورت کښې سر بس پر مربع~\LR{$0$} پاتې کېږي؛ په
نورو تعريفونو کښې ماشين درېږي۔) همداراز کله ناکله دا فرض هم ګټور وي
چې ټيورينګ ماشين، که بېخي درېږي، پر مربع~\LR{$1$} ودرېږي۔
\end{explain}
\olpseganchor{OLP-0260-B005}{end}

\olpseganchor{OLP-0260-B006}{start}
\begin{defn}\ollabel{defn:disciplined}
ټيورينګ ماشين~\LR{$M$} هله او يوازې هله \emph{منضبط} دے چې
\begin{enumerate}
    \item يو ټاکلی يوازينی درېدنی حالت~\LR{$h$} ولري،
    \item که بېخي درېږي، د مربع~\LR{$1$} د لوستلو پر وخت ودرېږي،
    \item پر مربع~\LR{$0$} د~\LR{$\TMendtape$} نښه هېڅکله پاکه نۀ کړي، او
    \item له مربع~\LR{$0$} څخه کيڼ لوري ته د تلو هڅه هېڅکله ونۀ کړي۔
\end{enumerate}
\end{defn}
\olpseganchor{OLP-0260-B006}{end}

\olpseganchor{OLP-0260-B007}{start}
\begin{explain}
موږ مخکښې بحث کړے چې هر ټيورينګ ماشين په داسې ماشين بدلېدے شي چې
هماغه چلند او يو ټاکلی درېدنی حالت ولري۔ دا کار يوازې د نوي حالت~\LR{$h$}
په زياتولو، او د هر داسې جوړې~\LR{$\tuple{q,\sigma}$} دپاره د لارښوونې
\LR{$\delta(q, \sigma) = \tuple{h, \sigma, N}$} په زياتولو کېږي چې اصلي
\LR{$\delta$} پرې تعريف شوې نۀ وي۔ که څه هم ثبوت ئې ستړے کوونکے دے، خو
دا رښتيا ده چې هر ټيورينګ ماشين~\LR{$M$} په داسې منضبط ټيورينګ ماشين~\LR{$M'$}
بدلېدے شي چې پر هماغو ننوتونو درېږي او هماغه وت پيدا کوي۔ مثلاً، که
ټيورينګ ماشين ودرېږي او پر مربع~\LR{$1$} نۀ وي، داسې لارښوونې ور زياتولے
شو چې سر تر هغه کيڼ لوري ته بوځي څو د پټې د پای نښه پيدا کړي، بيا يو
مربع ښي لوري ته لاړ شي، او بيا ودرېږي۔ د نورو شرطونو د پوره کولو پر
لار به تاسو پخپله فکر وکړئ۔
\end{explain}
\olpseganchor{OLP-0260-B007}{end}

\olpseganchor{OLP-0260-B008}{start}
\begin{ex}
    په~\olref{fig:adder-disc} کښې د جمع ماشين چې په
    \olref[una]{ex:adder} کښې و، په منضبط ماشين بدلوو۔
    \begin{figure}\[
\begin{tikzpicture}[->,>=stealth',shorten >=1pt,auto,node distance=2.8cm,
                    semithick]
  \tikzstyle{every state}=[fill=none,draw=black,text=black]
\olpseganchor{OLP-0260-B008}{end}

\olpseganchor{OLP-0260-B009}{start}
  \node[initial,state]         (A)              {\LR{$q_0$}};
  \node[state]         (B) [right of=A] {\LR{$q_1$}};
  \node[state]         (C) [below right of=B] {\LR{$q_2$}};
  \node[state]         (D) [below left of=C] {\LR{$q_3$}};
  \node[state]         (H) [left of=D] {\LR{$h$}};
\olpseganchor{OLP-0260-B009}{end}

\olpseganchor{OLP-0260-B010}{start}
  \path (A) edge node {\TMtrans{\TMblank}{\TMstroke}{\TMstay}} (B)
           edge [loop below] node {\TMtrans{\TMstroke}{\TMstroke}{\TMright}} (B)
        (B) edge [loop below] node {\TMtrans{\TMstroke}{\TMstroke}{\TMright}} (B)
            edge node {\TMtrans{\TMblank}{\TMblank}{\TMleft}} (C)
        (C) edge node {\TMtrans{\TMstroke}{\TMblank}{\TMleft}} (D)
        (D) edge [loop above] node {\TMtrans{\TMstroke}{\TMstroke}{\TMleft}} (D)
        edge node {\TMtrans{\TMendtape}{\TMendtape}{\TMright}} (H);
\end{tikzpicture}
\]\caption{يو منضبط جمع کوونکے ماشين}
\ollabel{fig:adder-disc}
\end{figure}
\end{ex}
\olpseganchor{OLP-0260-B010}{end}

\olpseganchor{OLP-0260-B011}{start}
\begin{prop}\ollabel{prop:disciplined} د هر ټيورينګ ماشين~\LR{$M$} دپاره
داسې منضبط ټيورينګ ماشين~\LR{$M'$} شته چې که~\LR{$M$} د وت~\LR{$O$} سره درېږي،
هغه هم د وت~\LR{$O$} سره درېږي؛ او که~\LR{$M$} نۀ درېږي، هغه هم نۀ درېږي۔ په
ځانګړې توګه، هره تابع \LR{$f\colon\Nat^n \to \Nat$} چې د ټيورينګ ماشين
په وسيله محاسبه کېدونکې وي، د منضبط ټيورينګ ماشين په وسيله هم محاسبه
کېدونکې ده۔
\end{prop}
\olpseganchor{OLP-0260-B011}{end}

\olpseganchor{OLP-0260-B012}{start}
\begin{prob}\label{tur:mac:dis:prob:disc-succ}
داسې منضبط ماشين ورکړئ چې \LR{$f(x) = x+1$} محاسبه کړي۔
\end{prob}
\olpseganchor{OLP-0260-B012}{end}


\olpseganchor{OLP-0260-B013}{start}
\begin{prob}\label{tur:mac:dis:prob:copier}
داسې منضبط ماشين پيدا کړئ چې پر ننوت \LR{$\TMstroke^n$} پيل شي او وت
\LR{$\TMstroke^n \concat \TMblank \concat \TMstroke^n$} پيدا کړي۔
\end{prob}
\olpseganchor{OLP-0260-B013}{end}

\olpseganchor{OLP-0261-B004}{start}
\olfileid{tur}{mac}{cmb}
\olsection{د ټيورينګ ماشينونو يوځاے کول}
\olpseganchor{OLP-0261-B004}{end}

\olpseganchor{OLP-0261-B005}{start}
\begin{explain}
د ټيورينګ ماشينونو هغه بېلګې چې تر اوسه مو ليدلې، نسبتاً ساده وې۔
خو په حقيقت کښې هره مسئله چې په هره اوسنۍ پروګرام‌ليکنې ژبه حل کېدے
شي، د ټيورينګ ماشينونو په وسيله هم حل کېدے شي۔ د لا پېچلو ټيورينګ
ماشينونو د جوړولو دپاره بايد ځان ډاډه کړو چې ماشينونه يوځاے کولے شو؛
په دې توګه به يوه کړنلاره په ساده برخو ووېشو او د لا پېچلو مسئلو د
حل ماشينونه به جوړ کړو۔ که د پېچلې مسئلې په رغوونکو برخو د وېشلو
طبيعي لار پيدا کړو، مسئله په څو پړاوونو کښې حل کولے شو: څو ساده
ټيورينګ ماشينونه جوړوو او بيا ئې په يوه داسې ماشين کښې يوځاے کوو چې
ټوله مسئله حل کړي۔ دا خبره په بله برخه کښې د درېدو مسئلې د څېړلو پر
وخت ځانګړې مهمه ده۔
\olpseganchor{OLP-0261-B005}{end}

\olpseganchor{OLP-0261-B006}{start}
ټيورينګ ماشينونه \LR{$M = \tuple{Q, \Sigma, q_0, \delta}$} او
~\LR{$M' = \tuple{Q', \Sigma', q_0', \delta'}$} څنګه يوځاے کړو؟ اوس د
~\LR{$M$} له درېدو وروسته د پټې تشکيل د ماشين~\LR{$M'$} د يوې چلونې د ننوت
تشکيل په توګه کاروو۔ د داسې يوه ټيورينګ ماشين \LR{$M \frown M'$} د لاس
ته راوړلو دپاره لاندې کارونه وکړئ:
\begin{enumerate}
    \item د~\LR{$M'$} ټول حالتونه~\LR{$Q'$} بيا وشمېرئ (يا بيا ونوموئ)، څو
    \LR{$M$} او~\LR{$M'$} هېڅ ګډ حالت ونۀ لري (\LR{$Q \cap Q' = \emptyset$})۔
    \item د~\LR{$M \frown M'$} حالتونه \LR{$Q \cup Q'$} دي۔
    \item د پټې الفبا \LR{$\Sigma \cup \Sigma'$} ده۔
    \item لومړنی حالت~\LR{$q_0$} دے۔
    \item د حالت بدلون تابع هغه تابع~\LR{$\delta''$} ده چې داسې ورکول کېږي:
    \[\delta''(q,\sigma) =
    \begin{cases}
      \delta(q,\sigma) & \text{\RL{که \LR{$q \in Q$} وي او دا ارزښت تعريف شوے وي}}\\
      \delta'(q,\sigma) & \text{\RL{که \LR{$q \in Q'$} وي}}\\
      \tuple{q_0', \sigma, \TMstay} & \text{\RL{که \LR{$q \in Q$} وي او
      \LR{$\delta(q,\sigma)$} تعريف شوې نۀ وي}}
    \end{cases}\]
\end{enumerate}
لاس ته راغلے ماشين چې په حالت \LR{$q \in Q$} کښې وي، د~\LR{$M$} لارښوونې
کاروي، او چې په حالت~\LR{$q \in Q'$} کښې وي، د~\LR{$M'$} لارښوونې کاروي۔ کله
چې دا په داسې حالت \LR{$q \in Q$} کښې وي او داسې نښه~\LR{$\sigma$} لولي چې
~\LR{$M$} ورته د حالت بدلون نۀ لري (يعنې~\LR{$M$} به درېدلے و)، د~\LR{$M'$} لومړني
حالت ته ننوځي (او د پټې منځپانګه او د سر ځاے هماغسې پرېږدي)۔
\textbf{د ژباړې د نسخې سمون (OLCMP-047):} د سرچينې د ټوټه‌ييز تعريف
لومړي شرط يوازې دا ويل چې حالت د لومړي ماشين دے، او په دې توګه ئې له
درېيم شرط سره ټکر کاوه۔ دلته څرګنده شوه چې لومړی شرط هغه وخت دے چې
د لومړي ماشين ارزښت تعريف شوے وي؛ درېيم شرط ئې ناتعريف شوی حالت
سمبالوي۔
\olpseganchor{OLP-0261-B006}{end}

\olpseganchor{OLP-0261-B007}{start}
پام وکړئ، تر څو چې ماشين~\LR{$M$} منضبط نۀ وي، موږ نۀ پوهېږو چې د~\LR{$M$}
د درېدو پر وخت د پټې سر چېرته وي؛ نو د~\LR{$M$} په درېدني تشکيل کښې سر
اړين نۀ دے چې مربع~\LR{$1$} ولولي۔ د ماشينونو د يوځاے کولو پر وخت بايد
دا خبره په پام کښې وساتو۔
\end{explain}
\olpseganchor{OLP-0261-B007}{end}

\olpseganchor{OLP-0261-B008}{start}
\begin{ex}
\emph{د ماشينونو يوځاے کول:} موږ به داسې ماشين طرح کړو چې د~\LR{$n$} او
~\LR{$m$} اوږدوالي لرونکو د~\LR{$\TMstroke$} نښو له دوو غونډونو جوړ ننوت باندې
پيل شي او پر پټه د~\LR{$2(m+n)$} شمېر~\LR{$\TMstroke$} نښو په يوه غونډ
ودرېږي۔ د دې ماشين د جوړولو دپاره دوه اشنا ماشينونه يوځاے کولے شو:
د جمع ماشين او دوه‌چنده کوونکے۔ لومړے د جمع ماشين د حالت ډياګرام
کاږو۔
\[
\begin{tikzpicture}[->,>=stealth',shorten >=1pt,auto,node distance=2.8cm,
                    semithick]
  \tikzstyle{every state}=[fill=none,draw=black,text=black]
\olpseganchor{OLP-0261-B008}{end}

\olpseganchor{OLP-0261-B009}{start}
  \node[initial,state] (A)              {\LR{$q_0$}};
  \node[state]         (B) [right of=A] {\LR{$q_1$}};
  \node[state]         (C) [right of=B] {\LR{$q_2$}};
\olpseganchor{OLP-0261-B009}{end}

\olpseganchor{OLP-0261-B010}{start}
  \path (A) edge node {\TMtrans{\TMblank}{\TMstroke}{\TMstay}} (B)
            edge [loop above] node {\TMtrans{\TMstroke}{\TMstroke}{\TMright}} (B)
        (B) edge [loop above] node {\TMtrans{\TMstroke}{\TMstroke}{\TMright}} (B)
            edge node {\TMtrans{\TMblank}{\TMblank}{\TMleft}} (C)
        (C) edge [loop above] node {\TMtrans{\TMstroke}{\TMblank}{\TMstay}} (C);
\end{tikzpicture}
\]
موږ نۀ غواړو چې ماشين په حالت~\LR{$q_2$} کښې ودرېږي، بلکې د وت د
دوه‌چنده کولو دپاره بايد کار ته دوام ورکړي۔ را ياد کړئ چې دوه‌چنده
کوونکے ماشين د ننوت لومړۍ کرښه پاکوي او په بېل وت کښې دوه کرښې ليکي۔
راځئ يوه لارښوونه ور زياته کړو، څو ډاډمن شو چې د پټې سر د جمع ماشين
د وت لومړۍ کرښه لولي۔
\[
\begin{tikzpicture}[->,>=stealth',shorten >=1pt,auto,node distance=2.8cm,
                    semithick]
  \tikzstyle{every state}=[fill=none,draw=black,text=black]
\olpseganchor{OLP-0261-B010}{end}

\olpseganchor{OLP-0261-B011}{start}
  \node[initial,state] (A)              {\LR{$q_0$}};
  \node[state]         (B) [right of=A] {\LR{$q_1$}};
  \node[state]         (C) [right of=B] {\LR{$q_2$}};
  \node[state]         (D) [below left of=C] {\LR{$q_3$}};
  \node[state]         (E) [below left of=D] {\LR{$q_4$}};
\olpseganchor{OLP-0261-B011}{end}

\olpseganchor{OLP-0261-B012}{start}
  \path (A) edge node {\TMtrans{\TMblank}{\TMstroke}{\TMstay}} (B)
            edge [loop above] node {\TMtrans{\TMstroke}{\TMstroke}{\TMright}} (B)
        (B) edge [loop above] node {\TMtrans{\TMstroke}{\TMstroke}{\TMright}} (B)
            edge node {\TMtrans{\TMblank}{\TMblank}{\TMleft}} (C)
        (C) edge node {\TMtrans{\TMstroke}{\TMblank}{\TMleft}} (D)
        (D) edge [loop left] node {\TMtrans{\TMstroke}{\TMstroke}{\TMleft}} (D)
            edge node[left, xshift=-2mm] {\TMtrans{\TMendtape}{\TMendtape}{\TMright}} (E);
\end{tikzpicture}
\]
اوس د ننوت دوه‌چنده کول اسانه دي---يوازې دوه‌چنده کوونکے ماشين له
حالت~\LR{$q_4$} سره نښلوو۔ د دې دپاره د دوه‌چنده کوونکي حالتونه داسې بيا
نومول پکار دي چې د~\LR{$q_0$} پر ځاے له~\LR{$q_4$} څخه پيل شي؛ په دې توګه دوه
لومړني حالتونه نۀ جوړېږي۔ وروستنی ډياګرام بايد د
\olref{fig:combined} په شان وي۔
\begin{figure}
\[
\begin{tikzpicture}[->,>=stealth',shorten >=1pt,auto,node distance=2.8cm,
                    semithick]
  \tikzstyle{every state}=[fill=none,draw=black,text=black]
  \node[initial,state] (A)              {\LR{$q_0$}};
  \node[state]         (B) [right of=A] {\LR{$q_1$}};
  \node[state]         (C) [right of=B] {\LR{$q_2$}};
  \node[state]         (D) [below left of=C] {\LR{$q_3$}};
  \node[state]         (E) [below left of=D] {\LR{$q_4$}};
  \node[state]         (2) [right of=E] {\LR{$q_5$}};
  \node[state]         (3) [right of=2] {\LR{$q_6$}};
  \node[state]         (4) [below of=3] {\LR{$q_7$}};
  \node[state]         (5) [left of=4]  {\LR{$q_8$}};
  \node[state]         (6) [left of=5]  {\LR{$q_9$}};
\olpseganchor{OLP-0261-B012}{end}

\olpseganchor{OLP-0261-B013}{start}
  \path (A) edge node {\TMtrans{\TMblank}{\TMstroke}{\TMstay}} (B)
            edge [loop above] node {\TMtrans{\TMstroke}{\TMstroke}{\TMright}} (B)
        (B) edge [loop above] node {\TMtrans{\TMstroke}{\TMstroke}{\TMright}} (B)
            edge node {\TMtrans{\TMblank}{\TMblank}{\TMleft}} (C)
        (C) edge node {\TMtrans{\TMstroke}{\TMblank}{\TMleft}} (D)
        (D) edge [loop left] node {\TMtrans{\TMstroke}{\TMstroke}{\TMleft}} (D)
            edge node[left, xshift=-2mm] {\TMtrans{\TMendtape}{\TMendtape}{\TMright}} (E)
    (E) edge node {\TMtrans{\TMstroke}{\TMblank}{\TMright}} (2)
    (2) edge [loop above] node {\TMtrans{\TMstroke}{\TMstroke}{\TMright}} (2)
      edge node {\TMtrans{\TMblank}{\TMblank}{\TMright}} (3)
    (3) edge [loop above] node {\TMtrans{\TMstroke}{\TMstroke}{\TMright}} (3)
        edge node {\TMtrans{\TMblank}{\TMstroke}{\TMright}} (4)
    (4) edge [loop below] node {\TMtrans{\TMblank}{\TMstroke}{\TMleft}} (4)
        edge node {\TMtrans{\TMstroke}{\TMstroke}{\TMleft}} (5)
    (5) edge [loop below]  node {\TMtrans{\TMstroke}{\TMstroke}{\TMleft}} (5)
        edge              node {\TMtrans{\TMblank}{\TMblank}{\TMleft}} (6)
    (6) edge [loop below] node {\TMtrans{\TMstroke}{\TMstroke}{\TMleft}} (6)
        edge              node {\TMtrans{\TMblank}{\TMblank}{\TMright}} (E);
\end{tikzpicture}
\]
\caption{د جمع کوونکي او دوه‌چنده کوونکي ماشينونو يوځاے کول}
\ollabel{fig:combined}
\end{figure}
\end{ex}
\olpseganchor{OLP-0261-B013}{end}

\olpseganchor{OLP-0261-B014}{start}
\begin{prop}
    که~\LR{$M$} او~\LR{$M'$} منضبط وي او په ترتيب سره تابعې \LR{$f\colon
    \Nat^k \to \Nat$} او \LR{$f'\colon \Nat \to \Nat$} محاسبه کوي، نو
    \LR{$M \frown M'$} منضبط دے او~\LR{$\comp{f}{f'}$} محاسبه کوي۔
\end{prop}
\olpseganchor{OLP-0261-B014}{end}

\olpseganchor{OLP-0261-B015}{start}
\begin{proof}
    څرنګه چې~\LR{$M$} منضبط دے، کله چې د وت
    ~\LR{$f(n_1,\dots,n_k) = m$} سره درېږي، سر مربع~\LR{$1$} لولي۔ که اوس د
    ~\LR{$M'$} لومړني حالت ته ننوځو، ماشين به د وت~\LR{$f'(m)$} سره ودرېږي او
    بيا به هم مربع~\LR{$1$} لولي۔ د~\olref[dis]{defn:disciplined} نور
    شرطونه هم پوره کېږي۔
\end{proof}
\olpseganchor{OLP-0261-B015}{end}

\olpseganchor{OLP-0261-B016}{start}
\begin{prob}
    د~\cref{tur:mac:dis:prob:disc-succ} ماشين~\LR{$M$} واخلئ،
    \LR{$M \frown M$} جوړ کړئ، او په دې توګه داسې منضبط ټيورينګ ماشين
    ورکړئ چې \LR{$f(x) = x+2$} محاسبه کوي۔
\end{prob}
\olpseganchor{OLP-0261-B016}{end}

\olpseganchor{OLP-0262-B004}{start}
\olfileid{tur}{mac}{var}
\olsection{د ټيورينګ ماشينونو بڼې}
\olpseganchor{OLP-0262-B004}{end}

\olpseganchor{OLP-0262-B005}{start}
په حقيقت کښې د ټيورينګ ماشينونو د تعريفولو ډېرې شونې لارې شته، او
زموږ تعريف يوازې يوه لار ده۔ زموږ تعريف له ځينو اړخونو څخه تر نورو
ازاد دے۔ موږ هره متناهي الفبا منو؛ يو لا محدود تعريف ښايي يوازې د
پټې دوې نښې، \LR{$\TMstroke$} او~\LR{$\TMblank$}، ومني۔ موږ ماشين ته اجازه
ورکوو چې يوه نښه پر پټه وليکي او په هماغه وخت کښې وخوځېږي؛ نور
تعريفونه يا ليکل او يا خوځېدل مني۔ موږ د سر له خوځولو پرته د نښې د
ليکلو شونتيا منو؛ نور تعريفونه د~\LR{$\TMstay$} ``لارښوونه'' نۀ لري۔ زموږ
تعريف له نورو اړخونو څخه زيات محدود دے۔ موږ فرض کړې چې پټه يوازې په
يوه لوري کښې نامتناهي ده؛ نور تعريفونه پټه کيڼ او ښي دواړو لوريو ته
نامتناهي پرېږدي۔ ان هر شمېر بېلې پټې، يا د مربعونو نامتناهي جال هم
منل کېدے شي۔ موږ د ټيورينګ ماشين د لارښوونو سټ د حالت بدلون تابع په
وسيله څرګندوو؛ نور تعريفونه د حالت بدلون اړيکه کاروي چې ماشين ته په
هر ټاکلي حالت کښې تر يوې زياتې ممکنې لارښوونې ورکوي۔
\olpseganchor{OLP-0262-B005}{end}

\olpseganchor{OLP-0262-B006}{start}
د تعريف دا وروستنۍ نرمونه ځانګړې په زړه پورې ده۔ زموږ په تعريف کښې
چې ماشين په حالت~\LR{$q$} کښې نښه~\LR{$\sigma$} لولي، \LR{$\delta(q, \sigma)$}
ټاکي چې نوې نښه، حالت او د پټې د سر ځاے څه دے۔ خو که د لارښوونو سټ
د اوسنيو حالت-نښه جوړو \LR{$\tuple{q,
  \sigma}$} او نويو حالت-نښه-لوري درې‌يزونو \LR{$\tuple{q', \sigma',
  D}$} تر منځ يوه اړيکه وګڼو، نو د ټيورينګ ماشين عمل ښايي په يوازيني
ډول ټاکلے نۀ وي---د لارښوونو اړيکه هم \LR{$\tuple{q,
  \sigma, q', \sigma', D}$} او هم \LR{$\tuple{q, \sigma, q'', \sigma'', D'}$}
لرلے شي۔ په دې صورت کښې موږ يو \emph{ناتعييني} ټيورينګ ماشين لرو۔
دا ماشينونه د محاسبوي پېچلتيا په تيورۍ کښې مهم رول لري۔
\olpseganchor{OLP-0262-B006}{end}

\olpseganchor{OLP-0262-B007}{start}
د ټيورينګ ماشين د درېدو دپاره هم بېلابېل دودونه شته: موږ وايو هغه
وخت درېږي چې د حالت بدلون تابع تعريف شوې نۀ وي؛ نور تعريفونه غواړي
چې ماشين په يوه ځانګړي ټاکلي درېدني حالت کښې وي۔ په~\olref[hal]{sec}
کښې مو څرګنده کړې چې د ټاکلي درېدني حالت غوښتل داسې بنديز نۀ دے چې
د ټيورينګ ماشينونو محاسبه کېدونکي څيزونه بدل کړي۔ څرنګه چې زموږ د
ټيورينګ ماشينونو پټې يوازې په يوه لوري کښې نامتناهي دي، داسې حالتونه
شته چې ټيورينګ ماشين يوه لارښوونه په سمه توګه نۀ شي عملي کولے: کله
چې تر ټولو کيڼ مربع لولي او بايد کيڼ لوري ته لاړ شي۔ زموږ د تعريف له
مخې ماشين د ``لوېدو'' پر ځاے هماغلته پاتې کېږي، خو موږ داسې تعريفولے
شو چې په دې حالت کښې ودرېږي۔ دا تعريف هم معادل دے: د هغه ټيورينګ
ماشين چلند چې له مربع~\LR{$0$} څخه کيڼ لوري ته د تلو پر هڅه درېږي، د هر
حالت بدلون \LR{$\delta(q,\TMendtape) =
\tuple{q',\sigma,\TMleft}$} په ړنګولو تقليدولے شو---بيا ماشين پر
~\LR{$\TMendtape$} د کيڼ لوري ته د تلو د هڅې پر ځاے درېږي۔\footnote{دا
\emph{بالکل} کار نۀ کوي، ځکه هېڅ څيز موږ د مربع~\LR{$0$} پرته پر نورو
مربعونو د~\LR{$\TMendtape$} له ليکلو او لوستلو څخه نۀ منع کوي
(\olref[una]{ex:mover} وګورئ)۔ د دې د مخنيوي دپاره د دې موخې لپاره
يوه دويمه نښه~\LR{$\TMendtape'$} ور زياتولے شو۔}
\olpseganchor{OLP-0262-B007}{end}

\olpseganchor{OLP-0262-B008}{start}
د عددونو د څرګندولو بېلابېلې لارې هم شته (او په دې توګه د ټيورينګ
ماشين محاسبه شوې ننوت-وت تابع هم بدلېږي): موږ يوګونې څرګندونه کاروو،
خو دوه‌ګونې څرګندونه هم کارېدے شي۔ د دې دپاره پر~\LR{$\TMblank$} او
~\LR{$\TMendtape$} سربېره دوو نښو ته اړتيا ده۔
\olpseganchor{OLP-0262-B008}{end}

\olpseganchor{OLP-0262-B009}{start}
اوس يوه په زړه پورې خبره دا ده: له دې بڼو څخه يوه هم دا نۀ بدلوي چې
کومې تابعې د ټيورينګ په معنا محاسبه کېدونکې دي۔ \emph{که يوه تابع د
يو تعريف له مخې د ټيورينګ په معنا محاسبه کېدونکې وي، نو د ټولو
تعريفونو له مخې د ټيورينګ په معنا محاسبه کېدونکې ده۔}
\olpseganchor{OLP-0262-B009}{end}

\olpseganchor{OLP-0262-B010}{start}
موږ د دې د پخلي جزئياتو ته نۀ ځو۔ يوازې يوه بېلګه وګورئ: کيڼ او ښي
دواړو لوريو ته د نامتناهي پټې، يا د څو پټو، په منلو هېڅ اضافي
محاسبوي ځواک نۀ لاس ته راوړو۔ په لنډه توګه دليل دا دے چې د يوې
يو-لوري نامتناهي پټې لرونکے ټيورينګ ماشين د څو پټو يا دوه-لوري
نامتناهي پټو تقليد کولے شي۔ مثلاً، د اضافي حالتونو او لارښوونو په
کارولو د څو پټو يا دوه-لوري نامتناهي پټې لرونکي ماشين پروګرام په
داسې پروګرام ``ژباړلے'' شو چې يوه يو-لوري نامتناهي پټه لري۔ ژباړل
شوے ماشين جفت مربعونه د پټې~\LR{$1$} مربعونو دپاره (يا د دوه-لوري
نامتناهي پټې د ``مثبتو'' مربعونو دپاره)، او طاق مربعونه د پټې~\LR{$2$}
مربعونو دپاره (يا د ``منفي'' مربعونو دپاره) کارولے شي۔
\olpseganchor{OLP-0262-B010}{end}

\olpseganchor{OLP-0263-B004}{start}
\olfileid{tur}{mac}{ctt}
\olsection{د چرچ--ټيورينګ اصل}
\olpseganchor{OLP-0263-B004}{end}

\olpseganchor{OLP-0263-B005}{start}
ټيورينګ ماشينونه بايد د اغېزمنې کړنلارې د مفهوم کره بديل وي۔ ټيورينګ
فکر کاوه چې هر څوک چې هم د اغېزمنې کړنلارې او هم د ټيورينګ ماشين په
مفهوم پوه شي، دا شهود به ولري چې هر هغه کار چې د اغېزمنې کړنلارې په
وسيله کېدے شي، د ټيورينګ ماشين په وسيله هم کېدے شي۔ دا ادعا له دې
خبرې ملاتړ اخلي چې د اغېزمنې کړنلارې د مفهوم ټول نور وړانديز شوي کره
بديلونه په امتدادي ډول د ټيورينګ ماشين له مفهوم سره معادل راووځي---
يعنې د تابعو کټ مټ هماغه سټ محاسبه کولے شي۔ دې ادعا ته
\emph{د چرچ--ټيورينګ اصل} ويل کېږي۔
\olpseganchor{OLP-0263-B005}{end}

\olpseganchor{OLP-0263-B006}{start}
\begin{defn}[د چرچ--ټيورينګ اصل]
\emph{د چرچ--ټيورينګ اصل} وايي چې هر هغه څيز چې د اغېزمنې کړنلارې
په وسيله محاسبه کېدے شي، د ټيورينګ په معنا محاسبه کېدونکے دے۔
\end{defn}
\olpseganchor{OLP-0263-B006}{end}

\olpseganchor{OLP-0263-B007}{start}
د چرچ--ټيورينګ اصل په دوو لارو کارول کېږي۔ لومړے ډول کارول ئې د
سستۍ يوه پلمه ده۔ فرض کړئ د يوه څيز د محاسبې د اغېزمنې کړنلارې
توضيح لرو، مثلاً په ``کاذب کوډ'' کښې۔ بيا د چرچ--ټيورينګ اصل کارولے
شو، څو دا ادعا توجيه کړو چې هماغه تابع يو ټيورينګ ماشين محاسبه کوي،
ان که موږ هغه ماشين په عمل کښې نۀ وي جوړ کړے۔
\olpseganchor{OLP-0263-B007}{end}

\olpseganchor{OLP-0263-B008}{start}
د چرچ--ټيورينګ اصل بل استعمال له فلسفي پلوه لا په زړه پورې دے۔
ثابتېدے شي چې داسې تابعې شته چې ټيورينګ ماشينونه ئې نۀ شي محاسبه
کولے۔ له دې څخه د چرچ--ټيورينګ اصل په کارولو پايله اخيستلے شو چې
هغه تابع په هېڅ ډول کړنلارې سره په اغېزمنه توګه نۀ شي محاسبه کېدے۔
ځکه که داسې کړنلاره موجوده وے، د چرچ--ټيورينګ اصل له مخې به ورته يو
ټيورينګ ماشين هم موجود وے۔ نو که ثابت کړے شو چې هېڅ ټيورينګ ماشين
هغه نۀ شي محاسبه کولے، اغېزمنه کړنلاره هم ورته نشته۔ په ځانګړې توګه،
د چرچ--ټيورينګ اصل د دې ادعا دپاره کارول کېږي چې د درېدو په نوم
مسئله نۀ يوازې د ټيورينګ ماشينونو په وسيله نۀ شي حل کېدے، بلکې په
هېڅ ډول په اغېزمنه توګه نۀ شي حل کېدے۔
\olpseganchor{OLP-0263-B008}{end}

\olpseganchor{OLP-0264-B004}{start}
\olchapter{tur}{und}{ناپرېکړتيا}
\olpseganchor{OLP-0264-B004}{end}

\olpseganchor{OLP-0265-B004}{start}
\olfileid{tur}{und}{int}
\olsection{پېژندنه}
\olpseganchor{OLP-0265-B004}{end}

\olpseganchor{OLP-0265-B005}{start}
ښايي دا څرګنده وبرېښي چې هره تابع، ان هره حسابي تابع، محاسبه کېدونکې
نۀ شي کېدے۔ تابعې دومره ډېرې دي او چلند ئې دومره پېچلے دے۔ مثلاً هغه
تابعې چې د راډيو‌اکټيفو ذرو له تجزيې، يا له بل ګډوډ يا تصادفي چلند
څخه تعريفېږي۔ فرض کړئ له يوه ټاکلي وخت څخه يوه-ثانيه‌يي واټنونه شمېرو،
او تابع~\LR{$f(n)$} د کاينات د هغو ذرو د شمېر په توګه تعريفوو چې له هغې
لومړنۍ شېبې وروسته په~\LR{$n$}-م يوه-ثانيه‌يي واټن کښې تجزيه کېږي۔ دا د
داسې تابعې غوندې ښکاري چې هېڅکله ئې د محاسبې هيله نۀ شو کولے۔
\olpseganchor{OLP-0265-B005}{end}

\olpseganchor{OLP-0265-B006}{start}
خو دا چې د داسې تابعو د محاسبې لار تصور نۀ شو کولے يوه خبره ده، او
دا چې په رښتيا ثابته کړو چې دوي نۀ محاسبه کېږي، بله خبره ده۔ په اصل
کښې ان هغه تابعې چې بېخي ناهيلې کوونکې پېچلې ښکاري، په مجردې معنا
محاسبه کېدونکې کېدے شي۔ مثلاً فرض کړئ کاينات د وخت له پلوه متناهي دے
---لکه ځينې کيهان‌پوهنيزې تيورۍ چې وړاندوينه کوي، يوه ورځ به په ډېر
لرې راتلونکي کښې کاينات په يوه ټکي کښې راټول شي۔ بيا له هغې لومړنۍ
شېبې څخه يوازې د متناهي (خو بېخي لوے) شمېر ثانيو دپاره~\LR{$f(n)$} تعريف
شوې وي۔ او هره تابع چې يوازې د متناهي شمېر ننوتونو دپاره تعريف شوې
وي محاسبه کېدونکې ده: وتونه په يوه لوے جدول کښې لست کولے، يا ئې د
ټيورينګ ماشين د حالت بدلون په يوه ډېر لوے ډياګرام کښې کوډولے شو۔
\olpseganchor{OLP-0265-B006}{end}

\olpseganchor{OLP-0265-B007}{start}
موږ ډېر ځله د هغو تابعو ځانګړو حالتونو ته پام کوو چې ارزښتونه ئې د
هو/نه پوښتنو ځوابونه ورکوي۔ مثلاً پوښتنه ``ايا~\LR{$n$} لومړنی عدد دے؟''
له دې تابع سره تړلې ده:
\[
\fn{isprime}(n) = \begin{cases}
  1 & \text{\RL{که \LR{$n$} لومړنی وي}}\\
  0 & \text{\RL{که نۀ وي۔}}
  \end{cases}
\]
وايو يوه هو/نه پوښتنه په \emph{اغېزمنه توګه پرېکړه کېدے} شي، که ور
سره تړلې د~\LR{$1/0$} ارزښت لرونکې تابع په اغېزمنه توګه محاسبه کېدونکې وي۔
\olpseganchor{OLP-0265-B007}{end}

\olpseganchor{OLP-0265-B008}{start}
د دې د رياضياتي ثبوت دپاره چې داسې تابعې شته چې په اغېزمنه توګه نۀ
شي محاسبه کېدے، يا داسې مسئلې شته چې په اغېزمنه توګه نۀ شي پرېکړه
کېدے، اړينه ده چې د محاسبې يو ټاکلی مدل وټاکو، او وښيو چې داسې تابعې
شته چې دا ئې نۀ شي محاسبه کولے يا داسې مسئلې شته چې دا ئې نۀ شي
پرېکړه کولے۔ مثلاً ښودلے شو چې هره تابع د ټيورينګ ماشينونو په وسيله
نۀ شي محاسبه کېدے او هره مسئله د ټيورينګ ماشينونو په وسيله نۀ شي
پرېکړه کېدے۔ بيا د چرچ--ټيورينګ اصل کارولے شو، څو پايله واخلو چې نۀ
يوازې ټيورينګ ماشينونه د هرې تابع د محاسبې دپاره پوره ځواکمن نۀ دي،
بلکې هېڅ اغېزمنه کړنلاره دا کار نۀ شي کولے۔
\textbf{د ژباړې د نسخې سمون (OLCMP-048):} په سرچينه کښې د ``په
اغېزمنه توګه پرېکړه کېدے'' په غونډله کښې مرستندوی فعل پاتې شوے و؛
دلته بشپړ ګرامري جوړښت راوړل شوے دے او معنا نۀ ده بدله شوې۔
\olpseganchor{OLP-0265-B008}{end}

\olpseganchor{OLP-0265-B009}{start}
د داسې منفي پايلو د ثبوت کلي دا حقيقت دے چې موږ پخپله ټيورينګ
ماشينونو ته عددونه ټاکلے شو۔ تر ټولو اسانه لار ئې دا ده چې دوي وشمېرو:
ښايي د ټيورينګ ماشينونو او د هغو د پروګرامونو د ليکلو يوه ټاکلې لار
وټاکو، او بيا ئې په منظم ډول لست کړو۔ کله چې ووينو دا کار کېدے شي،
نو د ټيورينګ-نامحاسبه کېدونکو تابعو شتون د شمېر کچې له ساده دلايلو
څخه راځي: له~\LR{$\Nat$} څخه~\LR{$\Nat$} ته د تابعو سټ (په اصل کښې ان يوازې
له~\LR{$\Nat$} څخه~\LR{$\{0, 1\}$} ته) د شمېر نۀ وړ دے، خو څرنګه چې ټول
ټيورينګ ماشينونه شمېرلے شو، د ټيورينګ-محاسبه کېدونکو تابعو سټ يوازې
شمېرېدونکے لامتناهي دے۔
\olpseganchor{OLP-0265-B009}{end}

\olpseganchor{OLP-0265-B010}{start}
موږ داسې \emph{ټاکلې} تابعې او مسئلې هم تعريفولے شو چې په ترتيب سره
ئې نۀ محاسبه کېدل او ناپرېکړتيا ثابتولے شو۔ يوه داسې مسئله د
\emph{درېدو مسئلې} په نوم ده۔ د ټيورينګ ماشين لارښوونې په لست کولو
هغه په متناهي ډول بيانېدے شي۔ د ټيورينګ ماشين دا ډول بيان، يعنې د
ټيورينګ ماشين پروګرام، البته بل ټيورينګ ماشين ته د ننوت په توګه
کارېدے شي۔ نو هغه ټيورينګ ماشينونه څېړلے شو چې د نورو ټيورينګ
ماشينونو په اړه پوښتنې پرېکړه کوي۔ يوه په زړه پورې پوښتنه دا ده:
``ايا ورکړل شوے ټيورينګ ماشين به پر ننوت~\LR{$n$} له پيل وروسته بالاخره
ودرېږي؟'' ښه به وے که داسې ټيورينګ ماشين موجود وے چې دا پوښتنه ئې
پرېکړه کولے: دا د کيفيت-کنټرول يو ټيورينګ ماشين وګڼئ چې ډاډ ورکوي
ټيورينګ ماشينونه په نامتناهي کړيو او داسې نورو کښې نۀ بندېږي۔ په زړه
پورې حقيقت، چې ټيورينګ ثابت کړ، دا دے چې داسې ټيورينګ ماشين نشي
موجودېدے۔ هېڅ يو ټيورينګ ماشين نشته چې د يوه ټيورينګ ماشين~\LR{$M$} له
بيان او يوه عدد~\LR{$n$} څخه جوړ ننوت باندې پيل شي او تل له~\LR{$1$} يا~\LR{$0$}
وت سره د دې له مخې ودرېږي چې~\LR{$M$} به پر ننوت~\LR{$n$} درېدلے و که نه۔
\textbf{د ژباړې د نسخې سمون (OLCMP-049):} سرچينې په وروستۍ پوښتنه
کښې د اېم له نوم وروسته د ``ماشين'' کلمه په تکرار راوړې وه؛ دلته
زائده تکرار لرې شوے او منطقي شرط هماغسې ساتل شوے دے۔
\olpseganchor{OLP-0265-B010}{end}

\olpseganchor{OLP-0265-B011}{start}
کله چې د ټاکلو ناپرېکړه کېدونکو مسئلو بېلګې ولرو، د نورو مسئلو د
ناپرېکړتيا د ښودلو دپاره ئې هم کارولے شو۔ مثلاً يوه مشهوره
ناپرېکړه کېدونکې مسئله دا پوښتنه ده: ``ايا د لومړۍ درجې
فارمول~\LR{$!A$} معتبر دے؟'' داسې ټيورينګ ماشين نشته چې د لومړۍ درجې
فارمول~\LR{$!A$} د ننوت په توګه ورکړل شي او هرومرو له~\LR{$1$} يا~\LR{$0$} وت
سره د دې له مخې ودرېږي چې~\LR{$!A$} معتبر دے که نه۔ په تاريخي ډول د دې
مسئلې د اغېزمن حل د کړنلارې د موندلو پوښتنې ته يوازې ``د'' پرېکړې
مسئله ويل کېده؛ نو وايو چې د پرېکړې مسئله ناحل کېدونکې ده۔ ټيورينګ
او چرچ دا پايله نږدې په يوه وخت کښې په خپلواکه توګه ثابته کړه، ځکه
نو دې ته د چرچ--ټيورينګ قضيه هم ويل کېږي۔
\olpseganchor{OLP-0265-B011}{end}

\olpseganchor{OLP-0266-B004}{start}
\olfileid{tur}{und}{enu}
\olsection{د ټيورينګ ماشينونو شمېرل}
\olpseganchor{OLP-0266-B004}{end}

\olpseganchor{OLP-0266-B005}{start}
\begin{explain}
موږ ښودلے شو چې د ټولو ټيورينګ ماشينونو سټ د شمېر وړ دے۔ دا
له دې حقيقت څخه راځي چې هر ټيورينګ ماشين په متناهي ډول بيانېدے شي۔
د حالتونو سټ او د پټې لغت‌زېرمه متناهي سټونه دي۔ د حالت بدلون تابع
له~\LR{$Q \times \Sigma$} څخه~\LR{$Q \times \Sigma \times \{\TMleft,
\TMright, \TMstay\}$} ته يوه قسمي تابع ده، نو دا هم د هغو متناهي
شمېر آرګوماني جوړو د ارزښتونو په لست کولو ټاکل کېدے شي چې تابع پرې
تعريف شوې ده۔
\olpseganchor{OLP-0266-B005}{end}

\olpseganchor{OLP-0266-B006}{start}
تر دې ځايه دا خبره سمه ده، خو يو نازک توپير شته۔ د ټيورينګ ماشينونو
تعريف پر دې هېڅ بنديز نۀ لګاوه چې د حالتونو د سټ او د پټې د الفبا
غړي څه شيان کېدے شي۔ نو مثلاً د هر حقيقي عدد دپاره له تخنيکي
پلوه يو ټيورينګ ماشين شته چې هماغه عدد د حالت په توګه کاروي۔ خو د
ټيورينګ ماشين \emph{چلند} له دې څخه خپلواک دے چې کوم څيزونه د حالتونو
او لغت‌زېرمه په توګه کارول کېږي۔ په~\olref{fig:variants} کښې دا دوه
ټيورينګ ماشينونه وګورئ۔
\begin{figure}
\begin{center}
\begin{tikzpicture}[->,>=stealth',shorten >=1pt,auto,node distance=2.8cm,
                    semithick]
  \tikzstyle{every state}=[fill=none,draw=black,text=black]
\olpseganchor{OLP-0266-B006}{end}

\olpseganchor{OLP-0266-B007}{start}
  \node[initial,state]         (A)              {\LR{$q_0$}};
  \node[state]         (B) [right of=A] {\LR{$q_1$}};
\olpseganchor{OLP-0266-B007}{end}

\olpseganchor{OLP-0266-B008}{start}
  \path (A) edge [bend left] node {\TMtrans{\TMstroke}{\TMstroke}{\TMright}} (B)
        (B) edge [loop above] node {\TMtrans{\TMblank}{\TMblank}{\TMright}} (B)
            edge [bend left] node {\TMtrans{\TMstroke}{\TMstroke}{\TMright}} (A);
\end{tikzpicture}\\
\begin{tikzpicture}[->,>=stealth',shorten >=1pt,auto,node distance=2.8cm,
  semithick]
\tikzstyle{every state}=[fill=none,draw=black,text=black]
\olpseganchor{OLP-0266-B008}{end}

\olpseganchor{OLP-0266-B009}{start}
\node[initial,state]         (A)              {\LR{$s$}};
\node[state]         (B) [right of=A] {\LR{$h$}};
\olpseganchor{OLP-0266-B009}{end}

\olpseganchor{OLP-0266-B010}{start}
\path (A) edge [bend left] node {\TMtrans{A}{A}{\TMright}} (B)
(B) edge [loop above] node {\TMtrans{\TMblank}{\TMblank}{\TMright}} (B)
edge [bend left] node {\TMtrans{A}{A}{\TMright}} (A);
\end{tikzpicture}
\end{center}
\caption{د \emph{جفت} ماشين بڼې}
\ollabel{fig:variants}
\end{figure}
دا دوه ډياګرامونه له دوو ماشينونو سره سمون لري: يو~\LR{$M$} چې د پټې
الفبا ئې~\LR{$\Sigma = \{\TMendtape,\TMblank,\TMstroke\}$} او د حالتونو
سټ ئې~\LR{$\{q_0,q_1\}$} دے، او بل~\LR{$M'$} چې الفبا ئې~\LR{$\Sigma' =
\{\TMendtape,\TMblank,A\}$} او حالتونه ئې~\LR{$\{s,h\}$} دي۔ خو نورې
لارښوونې ئې يو شان دي: \LR{$M$} به د~\LR{$n$} دانو~\LR{$\TMstroke$} پر لړۍ هله او
يوازې هله ودرېږي چې~\LR{$n$} جفت وي، او~\LR{$M'$} به د~\LR{$n$} دانو~\LR{$A$} پر لړۍ
هله او يوازې هله ودرېږي چې~\LR{$n$} جفت وي۔ موږ يوازې~\LR{$\TMstroke$} په~\LR{$A$}،
\LR{$q_0$} په~\LR{$s$}، او~\LR{$q_1$} په~\LR{$h$} بدل نومولي دي۔ دا بېلګه عموميت مومي:
موږ ټيورينګ ماشينونه تر هغه يو شان ګڼلے شو چې يو د نښو او حالتونو د
داسې نوم‌بدلون له لارې له بل څخه تر لاسه کېږي۔ په اصل کښې، موږ د
ټيورينګ ماشين نښې او حالتونه په ساده توګه مثبت صحيح اعداد ګڼلے شو:
د~\LR{$\sigma_0$} پر ځاے~\LR{$1$}، د~\LR{$\sigma_1$} پر ځاے~\LR{$2$}، او داسې نور
وګڼئ؛ \LR{$\TMendtape$}~\LR{$1$} دے، \LR{$\TMblank$}~\LR{$2$} دے، او داسې نور۔ په دې
ډول د \emph{جفت} ماشين په~\olref{fig:standard-even} کښې انځور شوے
ماشين ګرځي۔
\begin{figure}
\[\begin{tikzpicture}[->,>=stealth',shorten >=1pt,auto,node distance=2.8cm,
  semithick]
\tikzstyle{every state}=[fill=none,draw=black,text=black]
\olpseganchor{OLP-0266-B010}{end}

\olpseganchor{OLP-0266-B011}{start}
\node[initial,state]         (A)              {\LR{$1$}};
\node[state]         (B) [right of=A] {\LR{$2$}};
\olpseganchor{OLP-0266-B011}{end}

\olpseganchor{OLP-0266-B012}{start}
\path (A) edge [bend left] node {\TMtrans{3}{3}{\TMright}} (B)
(B) edge [loop above] node {\TMtrans{2}{2}{\TMright}} (B)
edge [bend left] node {\TMtrans{3}{3}{\TMright}} (A);
\end{tikzpicture}
\]
\caption{يو معياري \emph{جفت} ماشين}
\ollabel{fig:standard-even}
\end{figure}
هغه ټيورينګ ماشين چې حالتونه او نښې ئې مثبت صحيح اعداد وي ښايي
\emph{معياري} ماشين وبولو، او له دې وروسته يوازې معياري ماشينونه
وڅېړو۔\footnote{د ``معياري ماشين'' اصطلاح معياري نۀ ده۔}
\olpseganchor{OLP-0266-B012}{end}

\olpseganchor{OLP-0266-B013}{start}
موږ غوښتل وښيو چې د ټيورينګ ماشينونو سټ د شمېر وړ دے، او پورتنيو
خبرو ته په پام کافي ده چې وښيو د معياري ټيورينګ ماشينونو سټ
د شمېر وړ دے۔ فرض کړئ يو معياري ټيورينګ ماشين~\LR{$M = \tuple{Q,
\Sigma, q_0, \delta}$} راکړل شوے دے۔ هغه د مثبتو صحيح اعدادو په يوه
متناهي مزي څنګه بيانولے شو؟ لومړے به د حالتونو شمېر، پخپله حالتونه،
د نښو شمېر، پخپله نښې، او پيلنی حالت لست کړو۔ (په ياد ولرئ، دا ټول
مثبت صحيح اعداد دي، ځکه~\LR{$M$} معياري ماشين دے۔) د~\LR{$\delta$} په اړه څه؟
د ممکنو آرګومانونو سټ، يعنې د~\LR{$\tuple{q,\sigma}$} جوړې، متناهي دے،
ځکه~\LR{$Q$} او~\LR{$\Sigma$} متناهي دي۔ نو په~\LR{$\delta$} کښې معلومات يوازې د
ټولو هغو \LR{$5$}-تاييونو~\LR{$\tuple{q, \sigma, q', \sigma', d}$} متناهي
لست دے چې~\LR{$\delta(q,\sigma) = \tuple{q', \sigma', D}$} وي، او~\LR{$d$}
هغه عدد دے چې لوری~\LR{$D$} کوډوي (مثلاً د~\LR{$\TMleft$} دپاره~\LR{$1$}،
د~\LR{$\TMright$} دپاره~\LR{$2$}، او د~\LR{$\TMstay$} دپاره~\LR{$3$})۔
\olpseganchor{OLP-0266-B013}{end}

\olpseganchor{OLP-0266-B014}{start}
په دې ډول هر معياري ټيورينګ ماشين د مثبتو صحيح اعدادو په يوه متناهي
لست، يعنې د~\LR{$s_M \in (\PosInt)^*$} په لړۍ، بيانېدے شي۔ مثلاً معياري
\emph{جفت} ماشين په دې لړۍ کوډ کېږي:
\[
2, \underbrace{1, 2}_Q, 3, \overbrace{1, 2, 3}^\Sigma, 1, \underbrace{1, 3, 2, 3, 2}_{\delta(1,3) = \tuple{2,3,R}},
\overbrace{2, 2, 2, 2, 2}^{\delta(2,2) = \tuple{2,2,R}},
\underbrace{2, 3, 1, 3, 2}_{\delta(2,3) = \tuple{1,3,R}}.
\]
\end{explain}
\olpseganchor{OLP-0266-B014}{end}

\olpseganchor{OLP-0266-B015}{start}
\begin{thm}
له~\LR{$\Nat$} څخه~\LR{$\Nat$} ته داسې تابعې شته چې د ټيورينګ په وسيله
محاسبه کېدونکې نۀ دي۔
\end{thm}
\olpseganchor{OLP-0266-B015}{end}

\olpseganchor{OLP-0266-B016}{start}
\begin{proof}
موږ پوهېږو چې د مثبتو صحيح اعدادو د متناهي لړيو سټ~\LR{$(\PosInt)^*$}
د شمېر وړ دے (\cref{sfr:siz:zigzag:prob:posint-star})۔ له دې څخه
راځي چې د معياري ټيورينګ ماشينونو د بيانونو سټ، د~\LR{$(\PosInt)^*$} د
يوه فرعي سټ په توګه، پخپله هم شمېرل کېدونکے دے۔ له~\LR{$\Nat$} څخه~\LR{$\Nat$}
ته هره د ټيورينګ په وسيله محاسبه کېدونکې تابع د کوم (په اصل کښې د
ډېرو) ټيورينګ ماشين له خوا محاسبه کېږي۔ د هغې د حالتونو او نښو نومونه
په مثبتو صحيح اعدادو بدلولو سره (په ځانګړي ډول~\LR{$\TMendtape$} په~\LR{$1$}،
\LR{$\TMblank$} په~\LR{$2$}، او~\LR{$\TMstroke$} په~\LR{$3$}) وينو چې هره د ټيورينګ په
وسيله محاسبه کېدونکې تابع د يوه معياري ټيورينګ ماشين له خوا محاسبه
کېږي۔ معنا دا چې له~\LR{$\Nat$} څخه~\LR{$\Nat$} ته د ټولو د ټيورينګ په وسيله
محاسبه کېدونکو تابعو سټ هم شمېرل کېدونکے دے۔
\olpseganchor{OLP-0266-B016}{end}

\olpseganchor{OLP-0266-B017}{start}
له بلې خوا، له~\LR{$\Nat$} څخه~\LR{$\Nat$} ته د ټولو تابعو سټ د شمېر وړ
نۀ دے (\cref{sfr:siz:red:prob:nat-nat})۔ که ټولې تابعې د کوم ټيورينګ
ماشين په وسيله محاسبه کېدے، نو د هغو ټيورينګ ماشينونو د ټولو بيانونو
په لست کولو مو د ټولو تابعو سټ شمېرلے شو چې دا تابعې محاسبه کوي۔ نو
ځينې تابعې شته چې د ټيورينګ په وسيله محاسبه کېدونکې نۀ دي۔
\end{proof}
\olpseganchor{OLP-0266-B017}{end}

\olpseganchor{OLP-0266-B018}{start}
\begin{prob}
  ايا د ټيورينګ ماشينونو د بيانولو داسې لار درته ذهن ته راځي چې د
  حالتونو او د الفبا د نښو څرګند لست کولو ته اړتيا ونۀ لري؟ تاسو د
  ``معياري'' ماشين خپل مفهوم تعريفولے شئ، خو د دې په اړه څه ووايئ چې
  ولې هر ټيورينګ ماشين ستاسو په نوې معنا د يوه ``معياري'' ماشين له
  خوا تمثيلېدے شي۔
  \textbf{د ژباړې د نسخې سمون (OLCMP-050):} سرچينې دلته ويل چې هر
  ټيورينګ ماشين د معياري ماشين له خوا ``محاسبه'' کېدے شي؛ ماشين تابع
  نۀ ده، نو دلته موخه د هغه تمثيلول يا مشابه‌چلول دي۔
\end{prob}
\olpseganchor{OLP-0266-B018}{end}

\olpseganchor{OLP-0267-B004}{start}
\olfileid{tur}{und}{uni}
\olsection{نړيوال ټيورينګ ماشينونه}
\olpseganchor{OLP-0267-B004}{end}

\olpseganchor{OLP-0267-B005}{start}
په~\olref[enu]{sec} کښې مو وڅېړل چې هر ټيورينګ ماشين د صحيح اعدادو
په يوه متناهي لړۍ بيانېدے شي۔ دا لړۍ د ټيورينګ ماشين حالتونه، الفبا،
پيلنی حالت، او لارښوونې کوډوي۔ موږ دا هم په ګوته کړه چې د دې ټولو
بيانونو سټ د شمېر وړ دے۔ څرنګه چې د داسې بيانونو سټ
شمېرېدونکے لامتناهي دے، معنا ئې دا ده چې له~\LR{$\Nat$} څخه دې بيانونو ته يوه
پر هدف سټ بشپړه تابع شته۔ داسې پر هدف سټ بشپړه تابع مثلاً د کانتور د
زېګ‌زېګ ميتود په کارولو تر لاسه کېدے شي۔ دا موږ ته د ټيورينګ
ماشينونو د ټولو (بيانونو) د شمېرلو لار راکوي۔ که يوه داسې شمېره
ثابته کړو، نو اوس د~\LR{$1$}-م، \LR{$2$}-م، \dots، او~\LR{$e$}-م ټيورينګ ماشين په
اړه خبرې معنا لري۔ دې عددونو ته \emph{شاخصونه} ويل کېږي۔
\olpseganchor{OLP-0267-B005}{end}

\olpseganchor{OLP-0267-B006}{start}
\begin{defn}
که~\LR{$M$} (زموږ په ثابته شوې شمېره کښې) \LR{$e$}-م ټيورينګ ماشين وي، وايو
چې~\LR{$e$} د~\LR{$M$} يو \emph{شاخص} دے۔ د~\LR{$e$}-م ټيورينګ ماشين دپاره~\LR{$M_e$}
ليکو۔
\end{defn}
\olpseganchor{OLP-0267-B006}{end}

\olpseganchor{OLP-0267-B007}{start}
يو ماشين ښايي له يوه څخه ډېر شاخصونه ولري؛ مثلاً د~\LR{$M$} دوه بيانونه
ښايي د هغو د لارښوونو د لست کولو په ترتيب کښې توپير ولري، او دا
بېلابېل بيانونه به بېلابېل شاخصونه ولري۔
\olpseganchor{OLP-0267-B007}{end}

\olpseganchor{OLP-0267-B008}{start}
مهمه دا ده چې د ټيورينګ ماشين د بيانونو شمېره داسې ورکول کېدے شي
چې موږ د~\LR{$M$} بيان د هغه له شاخص څخه، او د يوه ماشين~\LR{$M$} يو شاخص د
هغه له بيان څخه په اغېزمنه توګه محاسبه کړے شو۔ بيا د چرچ--ټيورينګ د
اصل له مخې داسې ټيورينګ ماشين موندل کېدے شي چې د شاخص~\LR{$e$} لرونکي
ټيورينګ ماشين بيان بېرته تر لاسه کړي او اړوند بيان پر خپله پټه د وت
په توګه وليکي۔ بيان به د~\LR{$\TMstroke$} د بلاکونو يوه لړۍ وي (چې د~\LR{$M_e$}
په بيانوونکې لړۍ کښې مثبت صحيح اعداد تمثيلوي)۔
\olpseganchor{OLP-0267-B008}{end}

\olpseganchor{OLP-0267-B009}{start}
دې ته په پام اوس دا پوښتنه طبيعي ده: د ټيورينګ ماشين د شاخصونو کومې
تابعې پخپله د ټيورينګ ماشينونو په وسيله محاسبه کېدونکې دي؟ د ټيورينګ
ماشين د شاخصونو کوم خاصيتونه د ټيورينګ ماشينونو په وسيله پرېکړه کېدے
شي؟ يوه بېلګه: هغه تابع چې شاخص~\LR{$e$} د هغه شاخص~\LR{$e$} لرونکي ټيورينګ ماشين
د حالتونو شمېر ته وړي، د يوه ټيورينګ ماشين په وسيله محاسبه کېدونکې
ده۔ داسې ټيورينګ ماشين به دا کار کوي: پر هغې پټې له پيل وروسته چې د
\LR{$e$} دانو~\LR{$\TMstroke$} يو بلاک لري، لومړے به~\LR{$e$} په خپل بيان ناکوډ
کړي۔ بيان اوس پر پټه د~\LR{$\TMstroke$} د بلاکونو په يوه لړۍ تمثيلېږي۔
څرنګه چې په دې لړۍ کښې لومړے غړے د حالتونو شمېر دے، اوس
يوازې دا پاتې دي چې د~\LR{$\TMstroke$} له لومړي بلاک پرته نور هر څه پاک
کړي او بيا ودرېږي۔
\olpseganchor{OLP-0267-B009}{end}

\olpseganchor{OLP-0267-B010}{start}
يوه د پام وړ پايله دا ده:
\olpseganchor{OLP-0267-B010}{end}

\olpseganchor{OLP-0267-B011}{start}
\begin{thm}\ollabel{thm:universal-tm} داسې يو \emph{نړيوال ټيورينګ
  ماشين}~\LR{$U$} شته چې کله پر ننوت~\LR{$\tuple{e,n}$} پيل شي،
  \begin{enumerate}
    \item هله او يوازې هله درېږي چې~\LR{$M_e$} پر ننوت~\LR{$n$} ودرېږي، او
    \item که~\LR{$M_e$} له وت~\LR{$m$} سره ودرېږي، نو~\LR{$U$} هم همداسې کوي۔
  \end{enumerate}
  نو~\LR{$U$} هغه تابع~\LR{$f\colon \Nat \times \Nat \pto \Nat$} محاسبه کوي
  چې~\LR{$f(e,n) = m$} وي که~\LR{$M_e$} پر ننوت~\LR{$n$} له پيل وروسته له وت~\LR{$m$}
  سره ودرېږي، او که نۀ وي نو ناتعريفه ده۔
\end{thm}
\olpseganchor{OLP-0267-B011}{end}

\olpseganchor{OLP-0267-B012}{start}
\begin{proof}
  په عمل کښې د~\LR{$U$} جوړول بنسټيز ډول ناشوني دي، ځکه دا يو بېخي پېچلے
  ماشين دے۔ خو موږ د هغه د کار عمومي طرحه بيانولے، او بيا د
  چرچ--ټيورينګ اصل رابللے شو۔ کله چې~\LR{$U$} پيل شي، پټه ئې د~\LR{$e$} دانو
  \LR{$\TMstroke$} يو بلاک لري چې ورپسې د~\LR{$n$} دانو~\LR{$\TMstroke$} يو بلاک
  دے۔ لومړے شاخص~\LR{$e$} د ننوت~\LR{$n$} ښي لور ته ``ناکوډوي''۔ له دې څخه د
  اعدادو يو لست جوړېږي (يعنې د~\LR{$\TMstroke$} هغه بلاکونه چې په
  \LR{$\TMblank$} سره بېلېږي) چې د~\LR{$M_e$} ماشين لارښوونې بيانوي۔ بيا~\LR{$U$}
  د~\LR{$M_e$} د پيلني حالت عدد او عدد~\LR{$1$} د~\LR{$M_e$} د بيان ښي لور ته پر
  پټه ليکي۔ (دا هم په يوګوني ډول، د~\LR{$\TMstroke$} د بلاکونو په توګه،
  تمثيلېږي۔) ورپسې ننوت (د~\LR{$n$} دانو~\LR{$\TMstroke$} بلاک) ښي لور ته
  کاپي کوي---خو هر~\LR{$\TMstroke$} د درې دانو~\LR{$\TMstroke$} په يوه بلاک
  بدلوي (په ياد ولرئ، د~\LR{$\TMstroke$} نښې عدد~\LR{$3$} دے، \LR{$1$} د~\LR{$\TMendtape$}
  عدد او~\LR{$2$} د~\LR{$\TMblank$} عدد دے)۔ د دې بلاکونو د لړۍ په کيڼ سر کښې
  (چې د~\LR{$U$} پر پټه د~\LR{$\TMblank$} نښې سره بېلېږي) يو~\LR{$\TMstroke$}، يعنې
  د~\LR{$\TMendtape$} کوډ، ليکي۔
\olpseganchor{OLP-0267-B012}{end}

\olpseganchor{OLP-0267-B013}{start}
  اوس د~\LR{$U$} پر پټه دا شيان دي: شاخص~\LR{$e$}، عدد~\LR{$n$}، د پيلني حالت کوډ
  عدد (``اوسنی حالت'')، د سر د پيلني موقعيت عدد~\LR{$1$} (``د سر اوسنی
  موقعيت'')، او د ``پټې'' پيلنۍ منځپانګه (د~\LR{$\TMstroke$} د بلاکونو
  يوه لړۍ چې د~\LR{$M_e$} د نښو کوډ عددونه---``نښې''---تمثيلوي او په
  \LR{$\TMblank$} سره بېلېږي)۔
\olpseganchor{OLP-0267-B013}{end}

\olpseganchor{OLP-0267-B014}{start}
  اوس هغه څه په دې ډول مشابه چلوي چې~\LR{$M_e$} به پر ننوت~\LR{$n$} له پيل
  وروسته کړي وے:
  \begin{enumerate}
    \item د ``سر د اوسني موقعيت'' عدد~\LR{$k$} پيدا کول (په پيل کښې دا~\LR{$1$}
    دے)،
    \item د ``پټې'' \LR{$k$}-م بلاک ته تلل، څو وګوري چې هلته کومه ``نښه''
    ده،
    \item\ollabel{find-inst}%
    له اوسني ``حالت'' او ``نښې'' سره سمه لارښوونه پيدا کول،
    \item د ``پټې'' \LR{$k$}-م بلاک ته بېرته تلل او هلته ``نښه'' د هغې
    نښې په کوډ عدد بدلول چې~\LR{$M_e$} به ليکلې وے،
    \item سر هغه ځاے ته وړل چې اوسنی ``حالت'' پکې ثبتېږي، او هلته
    عدد د نوي حالت په عدد بدلول،
    \item هغه ځاے ته تلل چې د ``پټې موقعيت'' پکې ثبتېږي، او يو
    \LR{$\TMstroke$} پاکول يا يو~\LR{$\TMstroke$} زياتول (په ترتيب سره که
    لارښوونه د کيڼ يا ښي لور د حرکت خبره کوي)۔
    \item تکرارول۔\footnote{موږ دلته پر ځينو نازکو ستونزو سرسري
    تېرېږو۔ مثلاً کله چې~\LR{$U$} هغه شمېرند زياتوي چې د ``سر اوسنی
    موقعيت'' پرې ساتي، ښايي څه زيات ځاے ته اړ شي---په هغه حالت کښې
    به ټوله ``پټه'' ښي لور ته وخوځوي۔}
  \end{enumerate}
  که~\LR{$M_e$} پر ننوت~\LR{$n$} له پيل وروسته هېڅکله ونۀ درېږي، نو~\LR{$U$} هم
  هېڅکله نۀ درېږي، ځکه وت ئې ناتعريفه دے۔
\olpseganchor{OLP-0267-B014}{end}

\olpseganchor{OLP-0267-B015}{start}
  که په ګام~\olref{find-inst} کښې څرګنده شي چې د~\LR{$M_e$} بيان د اوسني
  ``حالت''/``نښې'' جوړې دپاره هېڅ لارښوونه نۀ لري، نو~\LR{$M_e$} به
  درېدلے وے۔ که دا پېښه شي، \LR{$U$} د خپلې پټې هغه برخه پاکوي چې د
  ``پټې'' کيڼ لور ته ده۔ د درې دانو~\LR{$\TMstroke$} د هر بلاک دپاره
  (چې د~\LR{$M_e$} پر پټه يو~\LR{$\TMstroke$} تمثيلوي) د خپلې پټې پر کيڼ سر
  يو~\LR{$\TMstroke$} ليکي، او ``پټه'' پرله‌پسې پاکوي۔ کله چې دا کار بشپړ
  شي، د~\LR{$U$} پټه د~\LR{$m$} اوږدوالي د~\LR{$\TMstroke$} يوازې يو بلاک لري۔
\olpseganchor{OLP-0267-B015}{end}

\olpseganchor{OLP-0267-B016}{start}
  که~\LR{$U$} پر ``پټه'' د درې دانو~\LR{$\TMstroke$} له بلاک پرته له بل څه
  سره مخ شي، سملاسي درېږي۔ څرنګه چې په دې حالت کښې د~\LR{$U$} پټه د
  \LR{$\TMstroke$} يوازې يو بلاک نۀ لري، وت ئې طبيعي عدد نۀ دے؛ يعنې په
  دې حالت کښې~\LR{$f(e,n)$} ناتعريفه ده۔
\end{proof}
\olpseganchor{OLP-0267-B016}{end}

\olpseganchor{OLP-0268-B004}{start}
\olfileid{tur}{und}{hal}
\olsection{د درېدنې مسئله}
\olpseganchor{OLP-0268-B004}{end}

\olpseganchor{OLP-0268-B005}{start}
\begin{explain}
فرض کړئ موږ د ټيورينګ ماشين د بيانونو کومه شمېره ثابته کړې ده۔ نو
هر ټيورينګ ماشين يو \emph{شاخص} اخلي: د ټيورينګ ماشين د بيانونو په
شمېره~\LR{$M_1$}، \LR{$M_2$}، \LR{$M_3$}، \dots{} کښې د هغه ځاے۔
\olpseganchor{OLP-0268-B005}{end}

\olpseganchor{OLP-0268-B006}{start}
موږ پوهېږو چې د ټيورينګ په وسيله نۀ محاسبه کېدونکې تابعې بايد موجودې
وي: د ټيورينګ ماشين د بيانونو سټ---او له دې سره د ټيورينګ ماشينونو
سټ---د شمېر وړ دے، خو له~\LR{$\Nat$} څخه~\LR{$\Nat$} ته د ټولو تابعو سټ
داسې نۀ دے۔ خو موږ د نامحاسبه کېدونکو تابعو ټاکلې بېلګې هم موندلے
شو۔ يوه داسې تابع د درېدنې تابع ده۔
\end{explain}
\olpseganchor{OLP-0268-B006}{end}

\olpseganchor{OLP-0268-B007}{start}
\begin{defn}[د درېدنې تابع]
 د \emph{درېدنې تابع}~\LR{$h$} داسې تعريفېږي:
\[
h(e,n) =
\begin{cases}
  \text{\RL{0}} & \text{\RL{که ماشين~\LR{$M_e$} د ننوت \LR{$n$} دپاره ونۀ درېږي}} \\
  \text{\RL{1}} & \text{\RL{که ماشين~\LR{$M_e$} د ننوت \LR{$n$} دپاره ودرېږي}}
\end{cases}
\]
\end{defn}
\olpseganchor{OLP-0268-B007}{end}

\olpseganchor{OLP-0268-B008}{start}
\begin{defn}[د درېدنې مسئله]
د \emph{درېدنې مسئله} د دې ټاکلو مسئله ده چې (د هر~\LR{$e$} او~\LR{$n$} دپاره)
ايا ټيورينګ ماشين~\LR{$M_e$} د~\LR{$n$} دانو کرښو پر ننوت درېږي که نه۔
\end{defn}
\olpseganchor{OLP-0268-B008}{end}

\olpseganchor{OLP-0268-B009}{start}
\begin{explain}
موږ د دې په ښودلو ثابتوو چې~\LR{$h$} د ټيورينګ په وسيله محاسبه کېدونکې
نۀ ده چې ور پورې تړلې تابع~\LR{$s$} د ټيورينګ په وسيله محاسبه کېدونکې
نۀ ده۔ دا ثبوت پر دې حقيقت تکيه کوي چې هر هغه څه چې يو ټيورينګ ماشين
ئې محاسبه کولے شي د يوه منضبط ټيورينګ ماشين په وسيله هم محاسبه کېدے
شي (\olref[mac][dis]{sec})، او پر دې حقيقت چې دوه ټيورينګ ماشينونه
سره نښلول کېدے شي، څو يو واحد ماشين جوړ کړي (\olref[mac][cmb]{sec})۔
\end{explain}
\olpseganchor{OLP-0268-B009}{end}

\olpseganchor{OLP-0268-B010}{start}
\begin{defn} تابع~\LR{$s$} داسې تعريفېږي:
\[
s(e) =
\begin{cases}
  \text{\RL{0}} & \text{\RL{که ماشين~\LR{$M_e$} د ننوت \LR{$e$} دپاره ونۀ درېږي}} \\
  \text{\RL{1}} & \text{\RL{که ماشين~\LR{$M_e$} د ننوت \LR{$e$} دپاره ودرېږي}}
\end{cases}
\]
\end{defn}
\olpseganchor{OLP-0268-B010}{end}

\olpseganchor{OLP-0268-B011}{start}
\begin{lem}
تابع~\LR{$s$} د ټيورينګ په وسيله محاسبه کېدونکې نۀ ده۔
\end{lem}
\olpseganchor{OLP-0268-B011}{end}

\olpseganchor{OLP-0268-B012}{start}
\begin{proof}
د تناقض دپاره فرضوو چې تابع~\LR{$s$} د ټيورينګ په وسيله محاسبه کېدونکې
ده۔ بيا به داسې ټيورينګ ماشين~\LR{$S$} موجود وے چې~\LR{$s$} محاسبه کوي۔ د
عموميت له زيان پرته فرضولے شو چې کله~\LR{$S$} درېږي، د لومړۍ مربع د
لوستلو پر مهال درېږي (يعنې منضبط دے)۔ دا ماشين له بل ماشين~\LR{$J$} سره
``نښلېدے'' شي چې که پر ننوت~\LR{$0$} پيل شي درېږي (يعنې که په پيلني
حالت کښې د پټې د پاي-نښې ښي لور ته مربع لوستلو پر مهال~\LR{$\TMblank$}
ولولي)، او که نۀ وي ښي لور ته روانېږي او هېڅکله نۀ درېږي۔
\LR{$S \concat J$}، يعنې هغه ماشين چې د~\LR{$S$} له~\LR{$J$} سره په نښلولو جوړېږي،
يو ټيورينګ ماشين دے، نو د کوم~\LR{$e$} دپاره~\LR{$M_e$} دے (يعنې په شمېره کښې
يو ځاے راځي)۔ \LR{$M_e$} د~\LR{$e$} دانو~\LR{$\TMstroke$} پر ننوت پيل کړئ۔ دوه
امکانونه دي: يا~\LR{$M_e$} درېږي يا نۀ درېږي۔
\begin{enumerate}
\item فرض کړئ~\LR{$M_e$} د~\LR{$e$} دانو~\LR{$\TMstroke$} پر ننوت درېږي۔ بيا~\LR{$s(e)
  = 1$}۔ نو کله چې~\LR{$S$} پر~\LR{$e$} پيل شي، د وت په توګه پر پټه له يوه
  \LR{$\TMstroke$} سره درېږي۔ بيا~\LR{$J$} پر داسې پټه پيلېږي چې يو~\LR{$\TMstroke$}
  لري۔ په هغه حالت کښې~\LR{$J$} نۀ درېږي۔ خو~\LR{$M_e$} هماغه ماشين~\LR{$S
  \concat J$} دے، نو بايد کټ مټ هغه څه وکړي چې لومړے~\LR{$S$} او ورپسې~\LR{$J$}
  ئې کوي (يعنې په دې حالت کښې ښي لور ته روان شي او هېڅکله ونۀ درېږي)۔
  نو~\LR{$M_e$} د~\LR{$e$} دانو~\LR{$\TMstroke$} پر ننوت نۀ شي درېدې۔
\olpseganchor{OLP-0268-B012}{end}

\olpseganchor{OLP-0268-B013}{start}
\item اوس فرض کړئ~\LR{$M_e$} د~\LR{$e$} دانو~\LR{$\TMstroke$} پر ننوت نۀ درېږي۔
  بيا~\LR{$s(e) = 0$}، او کله چې~\LR{$S$} پر ننوت~\LR{$e$} پيل شي، له تشې پټې سره
  درېږي۔ کله چې~\LR{$J$} پر تشه پټه پيل شي، سملاسي درېږي۔ بيا هم~\LR{$M_e$}
  هغه څه کوي چې لومړے~\LR{$S$} او ورپسې~\LR{$J$} ئې کوي، نو~\LR{$M_e$} بايد د~\LR{$e$}
  دانو~\LR{$\TMstroke$} پر ننوت ودرېږي۔
\end{enumerate}
په هر حالت کښې له خپل فرض سره تناقض ته رسېږو۔ دا ښيي چې داسې ټيورينګ
ماشين~\LR{$S$} نشي موجودېدے: \LR{$s$} د ټيورينګ په وسيله محاسبه کېدونکې نۀ
ده۔
\end{proof}
\olpseganchor{OLP-0268-B013}{end}

\olpseganchor{OLP-0268-B014}{start}
\begin{thm}[د درېدنې د مسئلې ناحلي]
\ollabel{thm:halting-problem} د درېدنې مسئله ناحل ده، يعنې تابع~\LR{$h$}
د ټيورينګ په وسيله محاسبه کېدونکې نۀ ده۔
\end{thm}
\olpseganchor{OLP-0268-B014}{end}

\olpseganchor{OLP-0268-B015}{start}
\begin{proof}
فرض کړئ~\LR{$h$} د ټيورينګ په وسيله محاسبه کېدونکې وے، مثلاً د ټيورينګ
ماشين~\LR{$H$} په وسيله۔ موږ له~\LR{$H$} څخه داسې ټيورينګ ماشين جوړولے شو چې
\LR{$s$} محاسبه کړي: لومړے د ننوت يوه کاپي جوړه کړئ (چې په~\LR{$\TMblank$}
نښه بېلېږي)۔ بيا بېرته پيل ته لاړ شئ او~\LR{$H$} وچلوئ۔ موږ په څرګند ډول
داسې ماشين جوړولے شو چې لومړے کار کوي
(\cref{tur:mac:dis:prob:copier} وګورئ)، او که~\LR{$H$} موجود وے، له داسې
کاپي‌کوونکي ماشين سره مو هغه ``نښلولے'' او داسې نوے ماشين مو ترې
جوړولے شو چې ټاکي ايا~\LR{$M_e$} پر ننوت~\LR{$e$} درېږي؛ يعنې~\LR{$s$} محاسبه کوي۔
خو موږ لا ښودلې ده چې داسې هېڅ ماشين نشي موجودېدے۔ نو~\LR{$h$} هم د
ټيورينګ په وسيله محاسبه کېدونکې نۀ ده۔
\end{proof}
\olpseganchor{OLP-0268-B015}{end}

\olpseganchor{OLP-0268-B016}{start}
\begin{prob}
د درې-درېدنې مسئله د داسې پرېکړيزې کړنلارې د ورکولو مسئله ده چې
وټاکي ايا په خپلسر غوره شوے ټيورينګ ماشين پر داسې پټه چې له دې پرته
تشه وي، د
درې دانو~\LR{$\TMstroke$} پر ننوت درېږي که نه۔ ثابته کړئ چې د درې-درېدنې
مسئله ناحل ده۔
\end{prob}
\olpseganchor{OLP-0268-B016}{end}

\olpseganchor{OLP-0268-B017}{start}
\begin{prob}
وښيئ چې که د ټيورينګ ماشين او ننوت د هغو جوړو~\LR{$M_e$} او~\LR{$n$} دپاره
چې~\LR{$e \neq n$} وي د درېدنې مسئله حل کېدونکې وي، نو د هغو حالتونو
دپاره هم حل کېدونکې ده چې~\LR{$e = n$} وي۔
\end{prob}
\olpseganchor{OLP-0268-B017}{end}

\olpseganchor{OLP-0268-B018}{start}
\begin{prob}
موږ ثابته کړه چې که ننوت داسې عدد~\LR{$e$} وي چې د ټولو ټيورينګ ماشينونو
د يوې شمېرې له لارې ټيورينګ ماشين~\LR{$M_e$} پېژني، نو د درېدنې مسئله
ناحل ده۔ که پر ځاے ئې د~\olref[tur][und][enu]{sec} د ټيورينګ
ماشينونو بيانونه نېغ په نېغه د ننوت په توګه ومنو، نو څه؟ ايا داسې
ټيورينګ ماشين موجودېدے شي چې د درېدنې مسئله پرېکړه کړي خو د شاخصونو
پر ځاے د ټيورينګ ماشينونو بيانونه د ننوت په توګه واخلي؟ څرګنده کړئ
چې ولې يا ولې نه۔
\end{prob}
\olpseganchor{OLP-0268-B018}{end}

\olpseganchor{OLP-0268-B019}{start}
\begin{prob} وښيئ چې \emph{قسمي} تابع~\LR{$s'$} چې داسې تعريفېږي:
  \[
  s'(e) =
  \begin{cases}
    \text{\RL{1}} & \text{\RL{که ماشين~\LR{$M_e$} د ننوت \LR{$e$} دپاره ودرېږي}}\\
    \text{\RL{ناتعريفه}} & \text{\RL{که ماشين~\LR{$M_e$} د ننوت \LR{$e$} دپاره ونۀ درېږي}}
  \end{cases}
  \]
  د ټيورينګ په وسيله محاسبه کېدونکې \emph{ده}۔
\end{prob}
\olpseganchor{OLP-0268-B019}{end}

\olpseganchor{OLP-0269-B004}{start}
\olfileid{tur}{und}{dec}
\olsection{د پرېکړې مسئله}
\olpseganchor{OLP-0269-B004}{end}

\olpseganchor{OLP-0269-B005}{start}
وايو د لومړۍ درجې منطق هله او يوازې هله \emph{پرېکړه کېدونکے} دے
چې د دې ټاکلو اغېزمن ميتود موجود وي چې يو ورکړل شوے تړلے فارمول
معتبر دے که نه۔ لکه چې څرګندېږي، داسې ميتود نشته: د لومړۍ درجې د
غونډلو د اعتبار د پرېکړې مسئله ناحل ده۔
\olpseganchor{OLP-0269-B005}{end}

\olpseganchor{OLP-0269-B006}{start}
د دې مهمې منفي پايلې د ټينګولو دپاره ثابتوو چې د پرېکړې مسئله د يوه
ټيورينګ ماشين په وسيله نۀ شي حل کېدے۔ يعنې ښيو چې داسې ټيورينګ ماشين
نشته چې هر کله پر داسې پټه پيل شي چې د لومړۍ درجې يو تړلے فارمول
لري، بالاخره ودرېږي او د دې له مخې~\LR{$1$} يا~\LR{$0$} وباسي چې تړلے فارمول
معتبر دے که نه۔ د چرچ--ټيورينګ د اصل له مخې هره محاسبه کېدونکې تابع
د ټيورينګ په وسيله محاسبه کېدونکې ده۔ نو که دا ``د اعتبار تابع''
بيخي په اغېزمنه توګه محاسبه کېدونکې وے، د ټيورينګ په وسيله به محاسبه
کېدونکې وے۔ که د ټيورينګ په وسيله محاسبه کېدونکې نۀ وي، نو په
اغېزمنه توګه هم نۀ شي محاسبه کېدے۔
\olpseganchor{OLP-0269-B006}{end}

\olpseganchor{OLP-0269-B007}{start}
د پرېکړې مسئلې د ناحلي د ثبوت تګلاره مو دا ده چې د درېدنې مسئله هغې
ته راکمه کړو۔ معنا ئې دا ده: موږ ثابته کړې چې هغه تابع~\LR{$h(e,w)$} چې
که په~\LR{$e$} بيان شوے ټيورينګ ماشين پر ننوت~\LR{$w$} ودرېږي، له وت~\LR{$1$} سره
درېږي، او که نۀ وي~\LR{$0$} وباسي، د ټيورينګ په وسيله محاسبه کېدونکې نۀ
ده۔ موږ به وښيو چې که داسې ټيورينګ ماشين وے چې د لومړۍ درجې د غونډلو
اعتبار پرېکړه کوي، نو داسې ټيورينګ ماشين هم شته چې~\LR{$h$} محاسبه کوي۔
څرنګه چې~\LR{$h$} د ټيورينګ ماشين په وسيله نۀ شي محاسبه کېدے، داسې
ټيورينګ ماشين هم نشته چې اعتبار پرېکړه کړي۔
\olpseganchor{OLP-0269-B007}{end}

\olpseganchor{OLP-0269-B008}{start}
په دې تګلاره کښې لومړے ګام دا ښودل دي چې د هر ننوت~\LR{$w$} او يوه
ټيورينګ ماشين~\LR{$M$} دپاره موږ په اغېزمنه توګه هغه تړلے فارمول
\LR{$!T(M, w)$} بيانولے شو چې د~\LR{$M$} د لارښوونو سټ او ننوت~\LR{$w$} تمثيلوي،
او هغه تړلے فارمول~\LR{$!E(M, w)$} چې ``\LR{$M$} بالاخره درېږي'' څرګندوي،
داسې چې:
\begin{quote}
  \LR{$\Entails !T(M, w) \lif !E(M,w)$} هله او يوازې هله چې~\LR{$M$} د
  ننوت~\LR{$w$} دپاره ودرېږي۔
\end{quote}
زموږ د ثبوت لويه برخه به د دې غونډلو~\LR{$!T(M, w)$} او~\LR{$!E(M, w)$} په
بيانولو او د دې خبرې په تاييدولو کښې وي چې~\LR{$!T(M, w) \lif !E(M, w)$}
هله او يوازې هله معتبر دے چې~\LR{$M$} پر ننوت~\LR{$w$} ودرېږي۔
\olpseganchor{OLP-0269-B008}{end}

\olpseganchor{OLP-0270-B004}{start}
\olfileid{tur}{und}{rep}
\olsection{د ټيورينګ ماشينونو تمثيلول}
\olpseganchor{OLP-0270-B004}{end}

\olpseganchor{OLP-0270-B005}{start}
\begin{explain}
د لومړۍ درجې منطق په يوه تړلے فارمول د ټيورينګ ماشينونو او د هغو
د چلند د تمثيلولو دپاره بايد يوه مناسبه ژبه تعريف کړو۔ ژبه دوه برخې
لري: د ماشين د تشکيلونو د بيانولو دپاره پريديکاتونه، او د اجرا د
ګامونو (``شېبو'') او پر پټه د موقعيتونو د شمېرلو دپاره تعبيرونه۔
\olpseganchor{OLP-0270-B005}{end}

\olpseganchor{OLP-0270-B006}{start}
موږ دوه ډوله پريديکاتونه ورپېژنو چې دواړه دوه-ځايه دي: د هر
حالت~\LR{$q$} دپاره يو پريديکات~\LR{$\Obj Q_q$}، او د پټې د هرې
نښې~\LR{$\sigma$} دپاره يو پريديکات~\LR{$\Obj S_\sigma$}۔ لومړي موږ ته
اجازه راکوي چې د~\LR{$M$} حالت او د هغه د پټې د سر موقعيت بيان کړو، او
دويم موږ ته اجازه راکوي چې د پټې منځپانګه بيان کړو۔
\olpseganchor{OLP-0270-B006}{end}

\olpseganchor{OLP-0270-B007}{start}
د پټې د سر د موقعيتونو او د اجرا شوو ګامونو د شمېر د څرګندولو دپاره
د اعدادو د څرګندولو يوې لارې ته اړ يو۔ دا د يوه ثابت سمبول~\LR{$\Obj
0$}، او د يوې \LR{$1$}-ځايه تابع~\LR{$\prime$}، يعنې د تالي تابع، په وسيله کېږي۔
د دود له مخې دا د خپل آرګومان \emph{وروسته} ليکل کېږي (او قوسونه
ترې پرېږدو)۔
\olpseganchor{OLP-0270-B007}{end}

\olpseganchor{OLP-0270-B008}{start}
د هر عدد~\LR{$n$} دپاره يو معياري ترم~\LR{$\num{n}$} شته، يعنې د~\LR{$n$}
\emph{عددنښه}، چې په~\LR{$\Lang L_M$} کښې ئې تمثيلوي۔ \LR{$\num{0}$} هم
\LR{$\Obj 0$} دے، \LR{$\num{1}$} هم~\LR{$\Obj 0'$} دے، \LR{$\num{2}$} هم~\LR{$\Obj 0''$}
دے، او داسې نور۔ په لا صوري ډول:
\begin{align*}
\num{0} & = \Obj 0 \\
\num{n+1} &= \num{n}'
\end{align*}
ترم~\LR{$\num{0}$}، يعنې~\LR{$\Obj 0$}، هم پر پټه تر ټولو کيڼ موقعيت او هم د
اجرا له لومړي ګام څخه مخکې وخت (پيلنی تشکيل) نوموي۔ ترم~\LR{$\num{1}$}،
يعنې~\LR{$\Obj 0'$}، د تر ټولو کيڼې مربع ښي لور ته مربع او د اجرا له لومړي
ګام څخه وروسته وخت نوموي، او داسې نور۔
\olpseganchor{OLP-0270-B008}{end}

\olpseganchor{OLP-0270-B009}{start}
موږ يو پريديکات~\LR{$<$} هم ورپېژنو، څو هم د پټې د موقعيتونو ترتيب
(چې هلته معنا ئې ``کيڼ لور ته'' ده) او هم د اجرا د ګامونو ترتيب
(چې هلته معنا ئې ``مخکې'' ده) څرګند کړو۔
\olpseganchor{OLP-0270-B009}{end}

\olpseganchor{OLP-0270-B010}{start}
کله چې ژبه مو تياره شوه، د~\LR{$!T(M, w)$} ``اکسيومونه'' لست کوو؛ يعنې
هغه تړلي فارمولونه چې په ګډه د~\LR{$M$} چلند بيانوي کله چې پر ننوت~\LR{$w$}
وچلول شي۔ داسې تړلي فارمولونه به وي چې پر~\LR{$\Obj 0$}، \LR{$\prime$}، او~\LR{$<$}
شرطونه ږدي، داسې تړلي فارمولونه چې د ننوت تشکيل بيانوي، او داسې
تړلي فارمولونه چې د~\LR{$M$} تشکيل د يوې ځانګړې لارښوونې له اجرا وروسته
بيانوي۔
\end{explain}
\olpseganchor{OLP-0270-B010}{end}

\olpseganchor{OLP-0270-B011}{start}
\begin{defn}
  \ollabel{defn:tm-descr}
يو ټيورينګ ماشين~\LR{$M = \tuple{Q, \Sigma, q_0, \delta}$} چې راکړل شي،
ژبه~\LR{$\Lang L_M$} له دې شيانو جوړه ده:
\begin{enumerate}
\item د هر حالت~\LR{$q \in Q$} دپاره يو دوه-ځايه پريديکات
  \LR{$\Obj Q_q(x, y)$}۔ په شهودي ډول~\LR{$\Obj Q_q(\num{m}, \num{n})$} دا
  څرګندوي: ``له~\LR{$n$} ګامونو وروسته، \LR{$M$} په حالت~\LR{$q$} کښې دے او~\LR{$m$}-مه
  مربع لولي۔''
\item د هرې نښې~\LR{$\sigma\in \Sigma$} دپاره يو دوه-ځايه پريديکات
  \LR{$\Obj S_\sigma(x, y)$}۔ په شهودي ډول~\LR{$\Obj S_\sigma(\num{m},
  \num{n})$} دا څرګندوي: ``له~\LR{$n$} ګامونو وروسته، \LR{$m$}-مه مربع نښه
  \LR{$\sigma$} لري۔''
\item يو ثابت سمبول~\LR{$\Obj 0$}
\item يوه يو-ځايه تابع~\LR{$\prime$}
\item يو دوه-ځايه پريديکات~\LR{$<$}
\end{enumerate}
\end{defn}
\olpseganchor{OLP-0270-B011}{end}

\olpseganchor{OLP-0270-B012}{start}
هغه تړلي فارمولونه چې پر ننوت~\LR{$w = \sigma_{i_1}\dots\sigma_{i_k}$}
د ټيورينګ ماشين~\LR{$M$} عمل بيانوي دا دي:
\begin{enumerate}
\item هغه اکسيومونه چې اعداد او~\LR{$<$} بيانوي:
\begin{enumerate}
\item يو تړلے فارمول چې وايي هر عدد له خپل تالي څخه کوچنی دے:
\[
\lforall[x][x < x']
\]
\item يو تړلے فارمول چې د~\LR{$<$} تعدي تضمينوي:
\[
\lforall[x][\lforall[y][\lforall[z][
      ((x < y \land y < z) \lif x < z)]]]
\]
\end{enumerate}
\item هغه اکسيومونه چې د ننوت تشکيل بيانوي:
\begin{enumerate}
\item له~\LR{$0$} ګامونو وروسته---مخکې له دې چې ماشين پيل شي---\LR{$M$} په
  پيلني حالت~\LR{$q_0$} کښې دے او مربع~\LR{$1$} لولي:
\[
\Obj Q_{q_0}(\num{1}, \num{0})
\]
\item لومړۍ \LR{$k+1$} مربعې نښې~\LR{$\TMendtape$}، \LR{$\sigma_{i_1}$}، \dots،
  \LR{$\sigma_{i_k}$} لري:
\[
\Obj S_\TMendtape(\num{0}, \num{0}) \land
\Obj S_{\sigma_{i_1}}(\num{1}, \num{0}) \land
\dots \land
\Obj S_{\sigma_{i_k}}(\num{k}, \num{0})
\]
\item له دې پرته پټه تشه ده:
\[
\lforall[x][(\num{k} < x \lif \Obj S_\TMblank(x, \num{0}))]
\]
\end{enumerate}
\item هغه اکسيومونه چې له يوه تشکيل څخه بل ته د حالت بدلون بيانوي:
\olpseganchor{OLP-0270-B012}{end}

\olpseganchor{OLP-0270-B013}{start}
د لاندې برخې دپاره، \LR{$!A(x, y)$} د ټولو هغو تړلي فارمولونه عطف وي چې
دا بڼه لري:
\[
\lforall[z][
  (((z < x \lor x < z) \land \Obj S_\sigma(z, y))
  \lif \Obj S_\sigma(z, y'))]
\]
چې~\LR{$\sigma \in \Sigma$} وي۔ موږ~\LR{$!A(\num{m},\num{n})$} د دې څرګندولو
دپاره کاروو: ``له مربع~\LR{$m$} پرته، له~\LR{$n+1$} ګامونو وروسته پټه هماغسې
ده لکه له~\LR{$n$} ګامونو وروسته۔''
\begin{enumerate}
\item \ollabel{rep-right} د هرې لارښوونې~\LR{$\delta(q_i, \sigma) =
  \tuple{q_j, \sigma', \TMright}$} دپاره دا تړلے فارمول:
\begin{align*}
& \lforall[x][\lforall[y][(
   (\Obj Q_{q_i}(x, y) \land \Obj S_{\sigma}(x, y)) \lif {}]] \\
&\qquad   (\Obj Q_{q_j}(x', y') \land \Obj S_{\sigma'}(x, y') \land
!A(x, y)))
\end{align*}
دا وايي چې که ماشين له~\LR{$y$} ګامونو وروسته په حالت~\LR{$q_i$} کښې وي او
هغه مربع~\LR{$x$} لولي چې نښه~\LR{$\sigma$} لري، نو له~\LR{$y+1$} ګامونو وروسته
مربع~\LR{$x+1$} لولي، په حالت~\LR{$q_j$} کښې وي، مربع~\LR{$x$} اوس~\LR{$\sigma'$} لري،
او له~\LR{$x$} پرته هره مربع هماغه نښه لري چې له~\LR{$y$} ګامونو وروسته ئې
درلوده۔
\olpseganchor{OLP-0270-B013}{end}

\olpseganchor{OLP-0270-B014}{start}
\item \ollabel{rep-left} د هرې لارښوونې~\LR{$\delta(q_i, \sigma) =
  \tuple{q_j, \sigma', \TMleft}$} دپاره دا تړلے فارمول:
\begin{align*}
& \lforall[x][\lforall[y][
    ((\Obj Q_{q_i}(x', y) \land \Obj S_{\sigma}(x', y)) \lif {}]]\\
& \qquad   (\Obj Q_{q_j}(x, y') \land \Obj S_{\sigma'}(x', y') \land
!A(x', y))) \land {}\\
& \lforall[y][((\Obj Q_{q_i}(\num{0}, y) \land \Obj S_{\sigma}(\num{0},
    y)) \lif {}]\\
& \qquad (\Obj Q_{q_j}(\num{0}, y') \land \Obj S_{\sigma'}(\num{0},
  y') \land !A(\num{0}, y)))
\end{align*}
يوه شېبه فکر وکړئ چې دا څنګه کار کوي: اوس له ``که مربع~\LR{$x$} لولي
\dots'' څخه نۀ، بلکې له ``که مربع~\LR{$x+1$} لولي~\dots'' څخه پيل
کوو۔ کيڼ لور ته حرکت معنا دا چې په بل ګام کښې ماشين مربع~\LR{$x$} لولي۔
خو هغه مربع~\LR{$x+1$} ده چې څه پرې ليکل کېږي۔ دا کار ځکه داسې کوو چې
منفي کول يا د مخکيني تابع نۀ لرو۔
\textbf{د ژباړې د نسخې سمون (OLCMP-051):} د سرچينې په لومړي
کيڼ-حرکت اکسيوم کښې د نابدلې پټې شرط ناسمه مربع مستثنا کوله؛ دلته
هغه مربع مستثنا شوې چې ماشين پرې ليکي، څو له سملاسي ورپسې تشريح سره
او د حالت بدلون له معنا سره سمون ولري۔
\olpseganchor{OLP-0270-B014}{end}

\olpseganchor{OLP-0270-B015}{start}
په ياد ولرئ چې د~\LR{$x+1$} په بڼه اعداد~\LR{$1$}، \LR{$2$}، \dots، دي؛ يعنې دا
هغه حالت نۀ رانغاړي چې ماشين مربع~\LR{$0$} لولي او بايد کيڼ لور ته وخوځېږي
(چې البته نۀ شي کولے---يوازې پر ځاے پاتې کېږي)۔ هغه ځانګړے حالت د
دويم عطف له خوا رانغاړل کېږي: وايي چې که ماشين له~\LR{$y$} ګامونو وروسته
په حالت~\LR{$q_i$} کښې مربع~\LR{$0$} لولي او مربع~\LR{$0$} نښه~\LR{$\sigma$} لري، نو
له~\LR{$y+1$} ګامونو وروسته لا هم مربع~\LR{$0$} لولي، اوس په حالت~\LR{$q_j$} کښې
دے، پر مربع~\LR{$0$} نښه~\LR{$\sigma'$} ده، او له مربع~\LR{$0$} پرته مربعې هماغه
نښې لري چې له~\LR{$y$} ګامونو وروسته ئې لرلې۔
\textbf{د ژباړې د نسخې سمون (OLCMP-052):} په سرچينه کښې د ``له
ګامونو وروسته'' پر ځاے يوه چاپي تېروتنه راغلې وه؛ دلته سمه شوې ده۔
\item \ollabel{rep-stay} د هرې لارښوونې~\LR{$\delta(q_i, \sigma) =
  \tuple{q_j, \sigma', \TMstay}$} دپاره دا تړلے فارمول:
\begin{align*}
& \lforall[x][\lforall[y][(
   (\Obj Q_{q_i}(x, y) \land \Obj S_{\sigma}(x, y)) \lif {}]] \\
&\qquad   (\Obj Q_{q_j}(x, y') \land \Obj S_{\sigma'}(x, y') \land
!A(x, y)))
\end{align*}
\end{enumerate}
\end{enumerate}
\LR{$!T(M, w)$} د ټيورينګ ماشين~\LR{$M$} او ننوت~\LR{$w$} دپاره د پورتنيو ټولو
تړلي فارمولونه عطف وي۔
\olpseganchor{OLP-0270-B015}{end}

\olpseganchor{OLP-0270-B016}{start}
د دې څرګندولو دپاره چې~\LR{$M$} بالاخره درېږي، بايد داسې تړلے فارمول
ومومو چې وايي: ``له يو شمېر ګامونو وروسته به د حالت بدلون تابع
ناتعريفه وي۔'' \LR{$X$} د ټولو هغو جوړو~\LR{$\tuple{q, \sigma}$} سټ وي چې
\LR{$\delta(q, \sigma)$} ناتعريفه وي۔ بيا~\LR{$!E(M, w)$} دا تړلے فارمول وي:
\[
\lexists[x][\lexists[y][(\bigvee_{\tuple{q, \sigma} \in
      X}(\Obj Q_q(x, y) \land \Obj S_\sigma(x, y)))]]
\]
\olpseganchor{OLP-0270-B016}{end}

\olpseganchor{OLP-0270-B017}{start}
که داسې ټيورينګ ماشين وکاروو چې ټاکل شوے درېدنی حالت~\LR{$h$} ولري، خبره
لا اسانه ده: بيا دا تړلے فارمول~\LR{$!E(M, w)$}
\[
\lexists[x][\lexists[y][\Obj Q_h(x, y)]]
\]
دا څرګندوي چې ماشين بالاخره درېږي۔
\olpseganchor{OLP-0270-B017}{end}

\olpseganchor{OLP-0270-B018}{start}
\begin{prop}
\ollabel{prop:mlessk}
که~\LR{$m < k$} وي، نو~\LR{$!T(M, w) \Entails \num{m} < \num{k}$}۔
\end{prop}
\olpseganchor{OLP-0270-B018}{end}

\olpseganchor{OLP-0270-B019}{start}
\begin{proof}
تمرین۔
\end{proof}
\olpseganchor{OLP-0270-B019}{end}

\olpseganchor{OLP-0270-B020}{start}
\begin{prob}
\olref[tur][und][rep]{prop:mlessk} ثابت کړئ۔
(اشاره: پر~\LR{$k-m$} استقرا وکاروئ)۔
\end{prob}
\olpseganchor{OLP-0270-B020}{end}

\olpseganchor{OLP-0271-B004}{start}
\olfileid{tur}{und}{ver}
\olsection{د تمثيل تاييدول}
\olpseganchor{OLP-0271-B004}{end}

\olpseganchor{OLP-0271-B005}{start}
\begin{explain}
د دې د تاييدولو دپاره چې زموږ تمثيل کار کوي، بايد دوه خبرې ثابتې
کړو۔ لومړے بايد وښيو چې که~\LR{$M$} پر ننوت~\LR{$w$} ودرېږي، نو~\LR{$!T(M, w)
\lif !E(M, w)$} معتبر دے۔ بيا بايد معکوسه خبره وښيو، يعنې که~\LR{$!T(M,
w) \lif !E(M, w)$} معتبر وي، نو~\LR{$M$} په رښتيا هم پر ننوت~\LR{$w$} د چلولو
پر مهال بالاخره درېږي۔
\olpseganchor{OLP-0271-B005}{end}

\olpseganchor{OLP-0271-B006}{start}
د دې دواړو د ثبوت تګلارې بېخي سره توپير لري۔ د لومړۍ پايلې دپاره
بايد وښيو چې د لومړۍ درجې منطق يو تړلے فارمول (يعنې~\LR{$!T(M, w)
\lif !E(M, w)$}) معتبر دے۔ د دې تر ټولو اسانه لار د يوه
اشتقاق ورکول دي۔ خو زموږ ثبوت بايد د ټولو~\LR{$M$} او~\LR{$w$} دپاره
کار وکړي، نو په اصل کښې داسې يو واحد تړلے فارمول نشته چې موږ ئې
دپاره اشتقاق ورکړو، بلکې نامتناهي ډېر دي۔ نو تر ټولو ښه کار
چې کولے ئې شو دا دے چې په استقرا ثابته کړو: \LR{$M$} او~\LR{$w$} هر ډول وي، او
پر ننوت~\LR{$w$} د~\LR{$M$} درېدل هر څومره ګامونه اخلي، د~\LR{$!T(M, w) \lif
!E(M, w)$} يو اشتقاق به موجود وي۔
\olpseganchor{OLP-0271-B006}{end}

\olpseganchor{OLP-0271-B007}{start}
په طبيعي ډول به زموږ استقرا د هغو ګامونو پر شمېر مخته ځي چې~\LR{$M$} ئې
درېدني تشکيل ته تر رسېدو مخکې اخلي۔ په خپل استقرايي ثبوت کښې به ټينګه
کړو چې پر ننوت~\LR{$w$} د~\LR{$M$} د چلولو د هر ګام~\LR{$n$} دپاره~\LR{$!T(M, w)
\Entails !C(M, w, n)$}، چې~\LR{$!C(M, w, n)$} له~\LR{$n$} ګامونو وروسته پر~\LR{$w$}
د چلېدلي~\LR{$M$} تشکيل په سمه توګه بيانوي۔ اوس که~\LR{$M$} پر ننوت~\LR{$w$}، مثلاً
له~\LR{$n$} ګامونو وروسته، ودرېږي، \LR{$!C(M, w, n)$} به يو درېدنی تشکيل بيان
کړي۔ دا به هم وښيو چې هر کله~\LR{$!C(M, w, n)$} يو درېدنی تشکيل بيانوي،
\LR{$!C(M, w, n) \Entails !E(M, w)$}۔ نو که~\LR{$M$} پر ننوت~\LR{$w$} ودرېږي، د
کوم~\LR{$n$} دپاره به~\LR{$M$} له~\LR{$n$} ګامونو وروسته په درېدني تشکيل کښې وي۔
له همدې امله~\LR{$!T(M, w) \Entails !C(M, w, n)$}، چې~\LR{$!C(M, w, n)$} يو
درېدنی تشکيل بيانوي؛ او څرنګه چې په هغه حالت کښې~\LR{$!C(M, w, n)
\Entails !E(M, w)$}، نو~\LR{$!T(M, w) \Entails !E(M, w)$} تر لاسه کوو؛
يعنې~\LR{$\Entails !T(M, w) \lif !E(M, w)$}۔
\textbf{د ژباړې د نسخې سمون (OLCMP-053):} سرچينې د دې استدلال په
وروستۍ نتيجه کښې د تمثيلوونکې غونډلې د نښې لومړے سمبول پرې ايښے و؛
دلته بېرته راوړل شوے دے۔
\olpseganchor{OLP-0271-B007}{end}

\olpseganchor{OLP-0271-B008}{start}
د معکوسې خبرې تګلاره بېخي توپير لري۔ دلته فرضوو چې~\LR{$\Entails !T(M,
w) \lif !E(M, w)$}، او بايد ثابته کړو چې~\LR{$M$} پر ننوت~\LR{$w$} درېږي۔ له
فرض څخه مو تر لاسه کېږي چې~\LR{$!T(M, w) \Entails !E(M, w)$}، يعنې
\LR{$!E(M, w)$} په هر هغه جوړښت کښې رښتيا دے چې~\LR{$!T(M, w)$} پکې
رښتيا دے۔ نو داسې يو جوړښت~\LR{$\Struct{M}$} بيان کوو چې~\LR{$!T(M,
w)$} پکې رښتيا دے: ډومېن به ئې~\LR{$\Nat$} وي، او د ټولو~\LR{$\Obj Q_q$} او
\LR{$\Obj S_\sigma$} تعبير به پر ننوت~\LR{$w$} د~\LR{$M$} د چلولو پر مهال د تشکيلونو
له خوا ورکول کېږي۔ نو مثلاً~\LR{$\Sat{M}{\Obj Q_q(\num{m}, \num{n})}$}
هله او يوازې هله چې~\LR{$M$} پر ننوت~\LR{$w$} د~\LR{$n$} ګامونو دپاره وچلول شي،
په حالت~\LR{$q$} کښې وي او مربع~\LR{$m$} لولي۔ اوس څرنګه چې د فرض له مخې
\LR{$!T(M, w) \Entails !E(M, w)$}، او د جوړښت له مخې~\LR{$\Sat{M}{!T(M, w)}$}،
نو~\LR{$\Sat{M}{!E(M, w)}$}۔ خو~\LR{$\Sat{M}{!E(M, w)}$} هله او يوازې هله چې
کوم~\LR{$n \in \Domain{M} = \Nat$} موجود وي داسې چې پر ننوت~\LR{$w$} چلېدلے
\LR{$M$} له~\LR{$n$} ګامونو وروسته په درېدني تشکيل کښې وي۔
\textbf{د ژباړې د نسخې سمون (OLCMP-054):} سرچينې د جوړښت د صدق په
تشريح کښې په تېروتنه د ماشين پر ځاے د تمثيلوونکې غونډلې نوم راوړے و؛
دلته فاعل د بحث له ماشين سره سم شوے دے۔
\end{explain}
\olpseganchor{OLP-0271-B008}{end}

\olpseganchor{OLP-0271-B009}{start}
\begin{defn}
\LR{$!C(M, w, n)$} دا تړلے فارمول وي:
\[
\Obj Q_q(\num{m}, \num{n}) \land \Obj S_{\sigma_0}(\num{0}, \num{n})
\land \dots \land \Obj S_{\sigma_k}(\num{k}, \num{n}) \land
\lforall[x][(\num{k} < x \lif \Obj S_\TMblank(x, \num{n}))]
\]
چې~\LR{$q$} په وخت~\LR{$n$} کښې د~\LR{$M$} حالت دے، \LR{$M$} په وخت~\LR{$n$} کښې مربع~\LR{$m$}
لولي، مربع~\LR{$i$} په وخت~\LR{$n$} کښې نښه~\LR{$\sigma_i$} لري چې~\LR{$0 \le i \le
k$}، او~\LR{$k$} يا په وخت~\LR{$0$} کښې د پټې تر ټولو ښي ناتشه مربع ده، يا هغه
تر ټولو ښي مربع ده چې د پټې سر له~\LR{$n$} ګامونو وروسته ليدلې ده؛ په دې
دواړو کښې هره يوه چې لويه وي۔
\end{defn}
\olpseganchor{OLP-0271-B009}{end}

\olpseganchor{OLP-0271-B010}{start}
\begin{lem}\ollabel{lem:halt-config-implies-halt}
که پر ننوت~\LR{$w$} چلېدلے~\LR{$M$} له~\LR{$n$} ګامونو وروسته په يوه درېدني تشکيل
کښې وي، نو~\LR{$!C(M, w, n) \Entails !E(M, w)$}۔
\end{lem}
\olpseganchor{OLP-0271-B010}{end}

\olpseganchor{OLP-0271-B011}{start}
\begin{proof}
فرض کړئ~\LR{$M$} د ننوت~\LR{$w$} دپاره له~\LR{$n$} ګامونو وروسته درېږي۔ داسې يو
حالت~\LR{$q$}، مربع~\LR{$m$}، او نښه~\LR{$\sigma$} شته چې:
\begin{enumerate}
\item له~\LR{$n$} ګامونو وروسته~\LR{$M$} په حالت~\LR{$q$} کښې دے، هغه مربع~\LR{$m$}
  لولي چې~\LR{$\sigma$} پرې ښکاري۔
\item د حالت بدلون تابع~\LR{$\delta(q, \sigma)$} ناتعريفه ده۔
\end{enumerate}
\LR{$!C(M, w, n)$} د دې تشکيل بيان دے او دا فقرې~\LR{$\Obj Q_{q}(\num{m},
\num{n})$} او~\LR{$\Obj S_{\sigma}(\num{m}, \num{n})$} به پکې شاملې وي۔
دا فقرې په ګډه~\LR{$!E(M, w)$} راايجابوي:
\[
\lexists[x][\lexists[y][(\bigvee_{\tuple{q, \sigma} \in
      X}(\Obj Q_q(x, y) \land \Obj S_\sigma(x, y)))]]
\]
ځکه~\LR{$\Obj Q_{q}(\num{m}, \num{n}) \land \Obj S_{\sigma}(\num{m},
\num{n}) \Entails \bigvee_{\tuple{q, \sigma} \in X} (\Obj Q_q(\num{m},
\num{n}) \land \Obj S_{\sigma}(\num{m}, \num{n}))$}، ځکه
\LR{$\tuple{q,\sigma} \in X$}۔
\textbf{د ژباړې د نسخې سمون (OLCMP-055):} سرچينې په وروستي
استدلال کښې داسې لومړني حالت او نښه وکارول چې په فرض کښې نۀ وو ټاکل
شوي، او د نښې له محمول څخه ئې د ژبې نښه هم پرې ايښې وه؛ دلته هماغه
ټاکل شوي حالت او نښه او بشپړ محمول کارول شوي دي۔
\end{proof}
\olpseganchor{OLP-0271-B011}{end}

\olpseganchor{OLP-0271-B012}{start}
\begin{explain}
نو که~\LR{$M$} د ننوت~\LR{$w$} دپاره ودرېږي، داسې~\LR{$n$} شته چې~\LR{$!C(M, w, n)
\Entails !E(M,w)$}۔ اوس به وښيو چې د هر وخت~\LR{$n$} دپاره~\LR{$!T(M, w)
\Entails !C(M, w, n)$}۔
\end{explain}
\olpseganchor{OLP-0271-B012}{end}

\olpseganchor{OLP-0271-B013}{start}
\begin{lem}
\ollabel{lem:config}
د هر~\LR{$n$} دپاره، که~\LR{$M$} له~\LR{$n$} ګامونو وروسته لا نۀ وي درېدلے، نو
\LR{$!T(M, w) \Entails !C(M, w, n)$}۔
\end{lem}
\olpseganchor{OLP-0271-B013}{end}

\olpseganchor{OLP-0271-B014}{start}
\begin{proof}
د استقرا بنسټ: که~\LR{$n = 0$} وي، د~\LR{$!C(M, w, 0)$} عاطفې د~\LR{$!T(M, w)$}
هم عاطفې دي، نو د هغه له خوا ايجابېږي۔
\olpseganchor{OLP-0271-B014}{end}

\olpseganchor{OLP-0271-B015}{start}
استقرايي فرض: که~\LR{$M$} له~\LR{$n$}-م ګام څخه مخکې نۀ وي درېدلے، نو
\LR{$!T(M,w) \Entails !C(M, w, n)$}۔ بايد وښيو چې (تر هغه چې~\LR{$!C(M, w,
n)$} يو درېدنی تشکيل نۀ بيانوي) \LR{$!T(M, w) \Entails !C(M, w, n+1)$}۔
\olpseganchor{OLP-0271-B015}{end}

\olpseganchor{OLP-0271-B016}{start}
فرض کړئ~\LR{$n \ge 0$} او پر~\LR{$w$} پيل شوي~\LR{$M$} له~\LR{$n$} ګامونو وروسته په
حالت~\LR{$q$} کښې مربع~\LR{$m$} لولي۔ څرنګه چې~\LR{$M$} له~\LR{$n$} ګامونو وروسته نۀ
درېږي، د~\LR{$M$} په پروګرام کښې بايد له لاندې درېو بڼو څخه يوه لارښوونه
وي:
\textbf{د ژباړې د نسخې سمون (OLCMP-056):} سرچينې استقرايي ګام يوازې
د مثبتو شاخصونو دپاره پيل کړے و او له صفرم تشکيل څخه لومړي تشکيل ته
بدلون ئې پرې ايښے و؛ دلته شرط صفر هم رانغاړي۔
\olpseganchor{OLP-0271-B016}{end}

\olpseganchor{OLP-0271-B017}{start}
\begin{enumerate}
\item \ollabel{right} \LR{$\delta(q, \sigma) = \tuple{q', \sigma', \TMright}$}
\olpseganchor{OLP-0271-B017}{end}

\olpseganchor{OLP-0271-B018}{start}
\item \ollabel{left} \LR{$\delta(q, \sigma) = \tuple{q', \sigma', \TMleft}$}
\olpseganchor{OLP-0271-B018}{end}

\olpseganchor{OLP-0271-B019}{start}
\item \ollabel{stay} \LR{$\delta(q, \sigma) = \tuple{q', \sigma', \TMstay}$}
\end{enumerate}
\olpseganchor{OLP-0271-B019}{end}

\olpseganchor{OLP-0271-B020}{start}
موږ به دا درې واړه حالتونه په وار سره وڅېړو۔
\olpseganchor{OLP-0271-B020}{end}

\olpseganchor{OLP-0271-B021}{start}
\begin{enumerate}
\item فرض کړئ د~\olref{right} په بڼه يوه لارښوونه شته۔ د
  \olref[rep]{defn:tm-descr}\olref[rep]{rep-right} له مخې معنا ئې دا
  ده چې
\begin{align*}
& \lforall[x][\lforall[y][((\Obj Q_{q}(x,
  y) \land \Obj S_\sigma(x, y)) \lif {}]]\\
& \qquad (\Obj Q_{q'}(x',
  y') \land \Obj S_{\sigma'}(x, y') \land !A(x, y)))
\intertext{\RL{د~\LR{$!T(M,w)$} يوه عاطفه ده۔ دا لاندې تړلے فارمول راايجابوي
  (کلي منځته راوړنه، \LR{$\num{m}$} د~\LR{$x$} دپاره او~\LR{$\num{n}$} د~\LR{$y$}
  دپاره):}}
& (\Obj Q_{q}(\num{m}, \num{n}) \land \Obj S_{\sigma}(\num{m},
\num{n})) \lif {}\\
& \qquad (\Obj Q_{q'}(\num{m}', \num{n}') \land
\Obj S_{\sigma'}(\num{m}, \num{n}') \land !A(\num{m}, \num{n})).
\intertext{\RL{د استقرايي فرض له مخې~\LR{$!T(M, w) \Entails !C(M, w, n)$}؛
  يعنې:}}
& \Obj Q_q(\num{m}, \num{n}) \land \Obj S_{\sigma_0}(\num{0}, \num{n})
\land \dots \land \Obj S_{\sigma_k}(\num{k}, \num{n}) \land \\
& \qquad\lforall[x][(\num{k} < x \lif \Obj S_\TMblank(x, \num{n}))]\\
\intertext{\RL{څرنګه چې له~\LR{$n$} ګامونو وروسته د پټې مربع~\LR{$m$} نښه~\LR{$\sigma$}
  لري، اړونده عاطفه~\LR{$\Obj S_\sigma(\num{m}, \num{n})$} ده، نو دا
  راايجابوي:}}
& \Obj Q_{q}(\num{m}, \num{n}) \land \Obj S_{\sigma}(\num{m},
   \num{n})
\intertext{\RL{اوس تر لاسه کوو:}}
& \Obj Q_{q'}(\num{m}', \num{n}') \land \Obj S_{\sigma'}(\num{m},
  \num{n}') \land {}\\
& \qquad\Obj S_{\sigma_0}(\num{0}, \num{n}') \land \dots \land
  \Obj S_{\sigma_k}(\num{k}, \num{n}') \land {}\\
& \qquad \lforall[x][(\num{k} < x \lif \Obj S_\TMblank(x, \num{n}'))]
\end{align*}
په دې ډول: لومړۍ کرښه نېغ په نېغه د مخکيني شرطي د تالي څخه د مودوس
پونېنس په وسيله راځي۔ د منځنۍ کرښې هره عاطفه---چې
\LR{$S_{\sigma_m}(\num{m},\num{n}')$} ترې وتلې ده---په~\LR{$!C(M, w, n)$} کښې
له اړوندې عاطفې او~\LR{$!A(\num{m}, \num{n})$} څخه راځي۔
\olpseganchor{OLP-0271-B021}{end}

\olpseganchor{OLP-0271-B022}{start}
که~\LR{$m < k$} وي، \LR{$!T(M,w) \Entails \num{m} < \num {k}$}
(\olref[rep]{prop:mlessk}) او د~\LR{$<$} د تعدي له مخې
\LR{$\lforall[x][(\num{k} < x \lif \num{m} < x)]$} لرو۔ که~\LR{$m = k$} وي،
نو~\LR{$\lforall[x][(\num{k} < x \lif \num{m} < x)]$} يوازې د منطق له
مخې راځي۔ بيا وروستۍ کرښه په~\LR{$!C(M, w, n)$} کښې له اړوندې عاطفې،
\LR{$\lforall[x][(\num{k} < x \lif \num{m} < x)]$}، او~\LR{$!A(\num{m},
\num{n})$} څخه راځي۔ که~\LR{$m<k$} وي، دا لا همدا~\LR{$!C(M, w, n+1)$} دے۔
\olpseganchor{OLP-0271-B022}{end}

\olpseganchor{OLP-0271-B023}{start}
اوس فرض کړئ~\LR{$m=k$}۔ په هغه حالت کښې د پټې سر له~\LR{$n+1$} ګامونو وروسته
مربع~\LR{$k+1$} هم ليدلې ده، چې اوس تر ټولو ښي ليدل شوې مربع ده۔ نو
\LR{$!C(M, w, n+1)$} يوه نوې عاطفه~\LR{$\Obj S_\TMblank(\num{k}',\num{n}')$}
لري، او وروستۍ عاطفه ئې~\LR{$\lforall[x][(\num{k}' < x \lif \Obj
S_\TMblank(x, \num{n}'))]$} ده۔ بايد تاييد کړو چې دا دوه تړلي فارمولونه
هم ايجابېږي۔
\olpseganchor{OLP-0271-B023}{end}

\olpseganchor{OLP-0271-B024}{start}
موږ لا~\LR{$\lforall[x][(\num{k} < x \lif \Obj S_\TMblank(x,
  \num{n}'))]$} لرو۔ په ځانګړي ډول، له دې څخه~\LR{$\num{k} < \num{k}'
\lif \Obj S_\TMblank(\num{k}', \num{n}')$} تر لاسه کوو۔ له اکسيوم
\LR{$\lforall[x][x < x']$} څخه~\LR{$\num{k} < \num{k}'$} تر لاسه کوو۔ د مودوس
پونېنس له مخې~\LR{$\Obj S_\TMblank(\num{k}',\num{n}')$} راځي۔
\olpseganchor{OLP-0271-B024}{end}

\olpseganchor{OLP-0271-B025}{start}
همدارنګه، څرنګه چې~\LR{$!T(M,w) \Entails \num{k} < \num{k}'$}، د~\LR{$<$}
د تعدي اکسيوم موږ ته~\LR{$\lforall[x][(\num{k}' < x \lif \Obj
  S_\TMblank(x, \num{n}'))]$} راکوي۔ (د دې تاييد د تمرین په توګه
پرېږدو۔)
\textbf{د ژباړې د نسخې سمون (OLCMP-057):} سرچينې په دې دوو ځايونو
کښې نحوي ثابتېدنه کارولې وه، حال دا چې راجع شوې قضيه او د ټولې برخې
ادعا معنايي ايجاب بيانوي؛ دواړه اړيکې له همدې مخې يو شان شوې دي۔
\olpseganchor{OLP-0271-B025}{end}

\olpseganchor{OLP-0271-B026}{start}
\item فرض کړئ د~\olref{left} په بڼه يوه لارښوونه شته۔ بيا د
  \olref[rep]{defn:tm-descr}\olref[rep]{rep-left} له مخې،
\begin{align*}
& \lforall[x][\lforall[y][((\Obj Q_{q}(x', y) \land \Obj
    S_{\sigma}(x', y)) \lif {}]]\\
& \qquad (\Obj Q_{q'}(x, y') \land \Obj
  S_{\sigma'}(x', y') \land !A(x', y))) \land {}\\
& \lforall[y][((\Obj Q_{q}(\num{0}, y) \land \Obj S_{\sigma}(\num{0},
    y)) \lif {}]\\
& \qquad (\Obj Q_{q'}(\num{0}, y') \land \Obj S_{\sigma'}(\num{0}, y')
  \land !A(\num{0}, y)))
\intertext{\RL{د~\LR{$!T(M,w)$} يوه عاطفه ده۔ که~\LR{$m>0$} وي، نو~\LR{$l = m - 1$}
  پرېږدئ (يعنې~\LR{$m = l+1$})۔ د پورتني تړلے فارمول لومړۍ عاطفه دا
  راايجابوي:}}
& (\Obj Q_{q}(\num{l}', \num{n}) \land \Obj S_{\sigma}(\num{l}', \num{n}))
\lif {} \\
& \qquad (\Obj Q_{q'}(\num{l}, \num{n}') \land \Obj S_{\sigma'}(\num{l}',
\num{n}') \land !A(\num{l}', \num{n}))
\intertext{\RL{که نۀ وي، \LR{$l = m = 0$} پرېږدئ او لاندې هغه تړلے فارمول
  وګورئ چې دويمه عاطفه ئې راايجابوي:}}
& ((\Obj Q_{q}(\num{0}, \num{n}) \land \Obj S_{\sigma}(\num{0}, \num{n})) \lif {}\\
& \qquad   (\Obj Q_{q'}(\num{0}, \num{n}') \land \Obj S_{\sigma'}(\num{0}, \num{n}') \land
!A(\num{0}, \num{n})))
\intertext{\RL{هره يوه غونډله دا راايجابوي:}}
&  \Obj Q_{q'}(\num{l}, \num{n}') \land \Obj S_{\sigma'}(\num{m},
  \num{n}') \land {}\\
&\qquad  \Obj S_{\sigma_0}(\num{0}, \num{n}')\land \dots \land
  \Obj S_{\sigma_k}(\num{k}, \num{n}')  \land {} \\
& \qquad  \lforall[x][(\num{k} < x
  \lif \Obj S_\TMblank(x, \num{n}'))]
\end{align*}
لکه مخکې۔ (په ياد ولرئ چې په لومړي حالت کښې~\LR{$\num{l}' \ident
\num{l+1} \ident \num{m}$}، او په دويم حالت کښې~\LR{$\num{l} \ident
\num{0}$}۔) خو دا کټ مټ~\LR{$!C(M, w, n+1)$} دے۔
\textbf{د ژباړې د نسخې سمون (OLCMP-058):} سرچينې د کيڼ-حرکت په
لومړۍ څانګه کښې د ليکل شوې مربع پر ځاے ګاونډۍ مربع له بدلون څخه
مستثنا کړې وه، او د صفرمې مربع په څانګه کښې ئې د روانې لارښوونې پر
ځاے بې‌اړيکې شاخصي حالتونه کارولي وو۔ اکسيوم او دواړه منځته راوړنې
دلته د هماغې روانې لارښوونې له حالتونو او ليکل شوې مربع سره سمې شوې
دي۔
\olpseganchor{OLP-0271-B026}{end}

\olpseganchor{OLP-0271-B027}{start}
\item د~\olref{stay} حالت د تمرین په توګه پرېښوول شوے دے۔
\end{enumerate}
موږ وښودل چې د هر~\LR{$n$} دپاره~\LR{$!T(M, w) \Entails !C(M, w, n)$}۔
\end{proof}
\olpseganchor{OLP-0271-B027}{end}

\olpseganchor{OLP-0271-B028}{start}
\begin{prob}
د~\olref[tur][und][ver]{lem:config} د ثبوت حالت
\olref[tur][und][ver]{stay} بشپړ کړئ۔
\end{prob}
\olpseganchor{OLP-0271-B028}{end}

\olpseganchor{OLP-0271-B029}{start}
\begin{prob}
له~\LR{$\Obj S_{\sigma_i}(\num{i}, \num{n})$} او~\LR{$!A(m, n)$} څخه د
\LR{$\Obj S_{\sigma_i}(\num{i}, \num{n}')$} يو اشتقاق ورکړئ
(په دې فرض چې~\LR{$i \neq m$}؛ يعنې يا~\LR{$i < m$} يا~\LR{$m < i$})۔
\end{prob}
\olpseganchor{OLP-0271-B029}{end}

\olpseganchor{OLP-0271-B030}{start}
\begin{prob}
له~\LR{$\lforall[x][(\num{k} < x \lif \Obj S_\TMblank(x, \num{n}'))]$}،
\LR{$\lforall[x][x < x']$}، او
\LR{$\lforall[x][\lforall[y][\lforall[z][ ((x < y \land y < z) \lif x <
      z)]]]$} څخه د~\LR{$\lforall[x][(\num{k}' < x \lif \Obj
  S_\TMblank(x, \num{n}'))]$} يو اشتقاق ورکړئ۔)
\end{prob}
\olpseganchor{OLP-0271-B030}{end}


\olpseganchor{OLP-0271-B031}{start}
\begin{lem}
\ollabel{lem:valid-if-halt}
که~\LR{$M$} پر ننوت~\LR{$w$} ودرېږي، نو~\LR{$!T(M, w) \lif
!E(M, w)$} معتبر دے۔
\end{lem}
\olpseganchor{OLP-0271-B031}{end}

\olpseganchor{OLP-0271-B032}{start}
\begin{proof}
د~\olref{lem:config} له مخې پوهېږو چې د هر وخت~\LR{$n$} دپاره په
وخت~\LR{$n$} کښې د~\LR{$M$} د تشکيل بيان~\LR{$!C(M, w, n)$} د~\LR{$!T(M, w)$} له خوا
ايجابېږي۔ فرض کړئ~\LR{$M$} له~\LR{$k$} ګامونو وروسته درېږي۔ په هغه وخت کښې به
د کوم~\LR{$m \in \Nat$} دپاره مربع~\LR{$m$} لولي۔ بيا~\LR{$!C(M, w, k)$} د~\LR{$M$}
يو درېدنی تشکيل بيانوي؛ يعنې~\LR{$\Obj Q_q(\num{m}, \num{k})$} او~\LR{$\Obj
S_\sigma(\num{m}, \num{k})$} دواړه د عاطفو په توګه لري، چې
\LR{$\delta(q,\sigma)$} ناتعريفه ده۔ نو د
\olref{lem:halt-config-implies-halt} له مخې~\LR{$!C(M, w, k) \Entails
!E(M, w)$}۔ خو څرنګه چې~\LR{$!T(M, w) \Entails !C(M, w, k)$}، نو~\LR{$!T(M,
w) \Entails !E(M, w)$}، او ځکه~\LR{$!T(M, w) \lif !E(M, w)$} معتبر دے۔
\end{proof}
\olpseganchor{OLP-0271-B032}{end}


\olpseganchor{OLP-0271-B033}{start}
\begin{explain}
د خپلې ادعا د تاييد بشپړولو دپاره بايد معکوس لوری هم ټينګ کړو: که
\LR{$!T(M, w) \lif !E(M, w)$} معتبر وي، نو~\LR{$M$} په رښتيا هم پر ننوت~\LR{$w$}
له پيل وروسته درېږي۔
\end{explain}
\olpseganchor{OLP-0271-B033}{end}

\olpseganchor{OLP-0271-B034}{start}
\begin{lem}
\ollabel{lem:halt-if-valid}
که~\LR{$\Entails !T(M, w) \lif !E(M, w)$}، نو~\LR{$M$} پر ننوت~\LR{$w$} درېږي۔
\end{lem}
\olpseganchor{OLP-0271-B034}{end}

\olpseganchor{OLP-0271-B035}{start}
\begin{proof}
د~\LR{$\Lang L_M$} هغه جوړښت~\LR{$\Struct M$} وګورئ چې ډومېن ئې
\LR{$\Nat$} دے، \LR{$\Obj 0$} په~\LR{$0$}، \LR{$\prime$} په تالي تابع، او~\LR{$<$} د کوچنيوالي
په اړيکه تعبيروي، او محمولونه~\LR{$\Obj Q_q$} او~\LR{$\Obj S_\sigma$} په دې
ډول تعبيروي:
\begin{align*}
  \Assign{\Obj Q_q}{M} & =
\Setabs{\tuple{m, n}}{\begin{array}{ll}\text{\RL{پر~\LR{$w$} له پيل وروسته، له~\LR{$n$} ګامونو وروسته،}}\\ \text{\RL{\LR{$M$} په حالت \LR{$q$}
  کښې دے او مربع~\LR{$m$} لولي}}
\end{array}} \\
\Assign{\Obj S_\sigma}{M} & = \Setabs{\tuple{m, n}}{\begin{array}{ll}
\text{\RL{پر~\LR{$w$} له پيل وروسته، له \LR{$n$} ګامونو وروسته،}}\\ \text{\RL{مربع~\LR{$m$} د~\LR{$M$}
  نښه~\LR{$\sigma$} لري}}
\end{array}}
\end{align*}
په بله وينا، موږ جوړښت~\LR{$\Struct{M}$} داسې جوړوو چې هغه څه ګام
په ګام بيان کړي چې پر ننوت~\LR{$w$} پيل شوے~\LR{$M$} په رښتيا کوي۔ په څرګند
ډول~\LR{$\Sat{M}{!T(M, w)}$}۔ که~\LR{$\Entails !T(M, w) \lif !E(M, w)$}، نو
\LR{$\Sat{M}{!E(M, w)}$} هم؛ يعنې:
\[
\Sat{M}{\lexists[x][\lexists[y][(\bigvee_{\tuple{q, \sigma} \in
      X}(\Obj Q_q(x, y) \land \Obj S_\sigma(x, y)))]]}.
\]
څرنګه چې~\LR{$\Domain{M} = \Nat$}، بايد داسې~\LR{$m$}، \LR{$n \in \Nat$} موجود
وي چې د کوم~\LR{$q$} او~\LR{$\sigma$} دپاره~\LR{$\Sat{M}{\Obj Q_q(\num{m},
\num{n}) \land \Obj S_\sigma(\num{m}, \num{n})}$}، داسې چې
\LR{$\delta(q, \sigma)$} ناتعريفه وي۔ د~\LR{$\Struct M$} د تعريف له مخې، معنا
ئې دا ده چې پر ننوت~\LR{$w$} پيل شوے~\LR{$M$} له~\LR{$n$} ګامونو وروسته په
حالت~\LR{$q$} کښې دے او نښه~\LR{$\sigma$} لولي، او د حالت بدلون تابع ناتعريفه
ده؛ يعنې~\LR{$M$} درېدلے دے۔
\end{proof}
\olpseganchor{OLP-0271-B035}{end}

\olpseganchor{OLP-0272-B004}{start}
\olfileid{tur}{und}{uns}
\olsection{د پرېکړې مسئله ناحل ده}
\olpseganchor{OLP-0272-B004}{end}

\olpseganchor{OLP-0272-B005}{start}
\begin{thm}
\ollabel{thm:decision-prob}
د پرېکړې مسئله ناحل ده: داسې ټيورينګ ماشين~\LR{$D$} نشته چې پر هغې پټې
له پيلېدو وروسته چې د لومړۍ درجې منطق يو تړلے فارمول~\LR{$!B$} د ننوت
په توګه لري، \LR{$D$}~بالاخره ودرېږي او هله او يوازې هله~\LR{$1$} وباسي چې~\LR{$!B$}
معتبر وي، او که نۀ وي~\LR{$0$} وباسي۔
\end{thm}
\olpseganchor{OLP-0272-B005}{end}

\olpseganchor{OLP-0272-B006}{start}
\begin{proof}
فرض کړئ د پرېکړې مسئله حل کېدونکې وے، يعنې فرض کړئ چې يو ټيورينګ
ماشين~\LR{$D$} شته۔ بيا به مو د درېدنې مسئله په لاندې ډول حل کړې وے۔ يو
ټيورينګ ماشين~\LR{$E$} جوړوو چې د ټيورينګ ماشين~\LR{$M_e$} عدد~\LR{$e$} او
ننوت~\LR{$w$} ور کړل شي، اړوند تړلے فارمول~\LR{$!T(M_e, w) \lif !E(M_e,
w)$} محاسبه کړي او په داسې حال کښې ودرېږي چې پر پټه تر ټولو کيڼه
مربع لولي۔ بيا به ماشين~\LR{$E \concat D$} ته چې ننوتونه~\LR{$e$} او~\LR{$w$} ور
کړل شي، لومړے~\LR{$!T(M_e, w) \lif !E(M_e, w)$} محاسبه کړي او ورپسې به
د پرېکړې مسئلې ماشين~\LR{$D$} پر همدې ننوت وچلوي۔ \LR{$D$} هله او يوازې هله
له وت~\LR{$1$} سره درېږي چې~\LR{$!T(M_e, w) \lif !E(M_e, w)$} معتبر وي، او
که نۀ وي~\LR{$0$} وباسي۔ د~\olref[ver]{lem:halt-if-valid} او
\olref[ver]{lem:valid-if-halt} له مخې، \LR{$!T(M_e, w) \lif !E(M_e, w)$}
هله او يوازې هله معتبر دے چې~\LR{$M_e$} پر ننوت~\LR{$w$} ودرېږي۔ نو~\LR{$E
\concat D$} چې ننوتونه~\LR{$e$} او~\LR{$w$} ور کړل شي، هله او يوازې هله له
وت~\LR{$1$} سره درېږي چې~\LR{$M_e$} پر ننوت~\LR{$w$} ودرېږي، او که نۀ وي له
وت~\LR{$0$} سره درېږي۔ په بله وينا، \LR{$E \concat D$} به د درېدنې مسئله حل
کړې وے۔ خو د~\olref[hal]{thm:halting-problem} له مخې پوهېږو چې داسې
ټيورينګ ماشين شتون نۀ شي لرلے۔
\end{proof}
\olpseganchor{OLP-0272-B006}{end}

\olpseganchor{OLP-0272-B007}{start}
\begin{cor}\ollabel{cor:undecidable-sat}%
دا پرېکړه کېدونکې نۀ ده چې آيا د لومړۍ درجې منطق کوم خپل‌سری
تړلے فارمول د صدق وړ دے۔
\end{cor}
\olpseganchor{OLP-0272-B007}{end}

\olpseganchor{OLP-0272-B008}{start}
\begin{proof}
  فرض کړئ چې د صدق وړتيا يو ټيورينګ ماشين~\LR{$S$} پرېکړه کولے شوه۔ بيا
  به مو د پرېکړې مسئله په لاندې ډول حل کړې وے: يو
  تړلے فارمول~\LR{$!B$} د ننوت په توګه واخلئ او~\LR{$!B$} يوه مربع ښي لوري
  ته يوسئ۔ بېرته مربع~\LR{$1$} ته راشئ او نښه~\LR{$\lnot$} وليکئ۔
\olpseganchor{OLP-0272-B008}{end}

\olpseganchor{OLP-0272-B009}{start}
  اوس ټيورينګ ماشين~\LR{$S$} وچلوئ۔ دا بالاخره درېږي او پر پټه يا~\LR{$1$}
  وباسي (که~\LR{$\lnot !B$} د صدق وړ وي) يا~\LR{$0$} (که~\LR{$\lnot !B$} د صدق
  وړ نۀ وي)۔ که پر مربع~\LR{$1$} يوه~\LR{$\TMstroke$} وي، پاکه ئې کړئ؛ که
  مربع~\LR{$1$} تشه وي، يوه~\LR{$\TMstroke$} پرې وليکئ، او بيا ودرېږئ۔
\olpseganchor{OLP-0272-B009}{end}

\olpseganchor{OLP-0272-B010}{start}
  دا ټيورينګ ماشين تل درېږي، او وت ئې هله او يوازې هله~\LR{$1$} وي چې
  \LR{$\lnot !B$} د صدق وړ نۀ وي، او که نۀ وي~\LR{$0$} وي۔ څرنګه چې~\LR{$!B$} هله
  او يوازې هله معتبر دے چې~\LR{$\lnot !B$} د صدق وړ نۀ وي، ماشين هله او
  يوازې هله~\LR{$1$} وباسي چې~\LR{$!B$} معتبر وي، او که نۀ وي~\LR{$0$} وباسي؛
  يعنې د پرېکړې مسئله به حل کړي۔
  \textbf{د ژباړې د نسخې سمون (OLCMP-062):} سرچينې د ننوت د غونډلې
  په لومړۍ يادونه کښې د ټاکلې څيز-ژبې نښه پرېښې وه؛ دلته له ورپسې
  ټولو يادونو سره سمه بشپړه نښه کارول شوې ده۔
\end{proof}
\olpseganchor{OLP-0272-B010}{end}

\olpseganchor{OLP-0272-B011}{start}
\begin{explain}
نو داسې ټيورينګ ماشين نشته چې دې پوښتنې ته تل سم ``هو'' يا ``نه''
ځواب ورکړي: ``آيا~\LR{$!B$} د لومړۍ درجې منطق يو معتبر تړلے فارمول
دے؟'' خو داسې ټيورينګ ماشين \emph{شته} چې تل سم ``هو'' ځواب
ورکوي---مګر که ځواب ``نه'' وي، هېڅ نۀ درېږي۔ دا د لومړۍ درجې منطق
د سموالي او بشپړتيا له قضيې، او له دې حقيقت څخه راځي چې
اشتقاقونه په اغېزمنه توګه شمېرل کېدے شي۔
\end{explain}
\olpseganchor{OLP-0272-B011}{end}

\olpseganchor{OLP-0272-B012}{start}
\begin{thm}
  \ollabel{thm:valid-ce}%
  د لومړۍ درجې د تړلي فارمولونه اعتبار نيمه د پرېکړې وړ دے: داسې
  ټيورينګ ماشين~\LR{$E$} شته چې پر هغې پټې له پيلېدو وروسته چې د لومړۍ
  درجې منطق يو تړلے فارمول~\LR{$!B$} د ننوت په توګه لري، \LR{$E$}~هله او يوازې
  هله بالاخره درېږي او~\LR{$1$} وباسي چې~\LR{$!B$} معتبر وي، او که نۀ وي نۀ
  درېږي۔
\end{thm}
\olpseganchor{OLP-0272-B012}{end}

\olpseganchor{OLP-0272-B013}{start}
\begin{proof}
  د لومړۍ درجې منطق ټول ممکن اشتقاقونه د يوه اغېزمن الګوريتم
  په وسيله يو په بل پسې توليدېدے شي۔ ماشين~\LR{$E$} همدا کار کوي، او کله
  چې داسې اشتقاق ومومي چې~\LR{$\Proves !B$} ښيي، له وت~\LR{$1$} سره
  درېږي۔ د سموالي د قضيې له مخې، که~\LR{$E$} له وت~\LR{$1$} سره درېږي، لامل
  ئې دا دے چې~\LR{$\Entails !B$}۔ د بشپړتيا د قضيې له مخې، که~\LR{$\Entails
  !B$} وي، داسې اشتقاق شته چې~\LR{$\Proves !B$} ښيي۔ څرنګه
  چې~\LR{$E$} ټول ممکن اشتقاقونه په منظم ډول توليدوي، بالاخره به
  داسې يو ومومي چې~\LR{$\Proves !B$} ښيي؛ نو بالاخره به له وت~\LR{$1$} سره
  ودرېږي۔
\end{proof}
\olpseganchor{OLP-0272-B013}{end}

\olpseganchor{OLP-0273-B004}{start}
\olfileid{tur}{und}{tra}
\olsection{د تراختن‌بروت قضيه}
\olpseganchor{OLP-0273-B004}{end}

\olpseganchor{OLP-0273-B005}{start}
\begin{explain}
  په~\olref[rep]{sec} کښې مو د يوه ټيورينګ ماشين~\LR{$M$} او ننوت
  رشته~\LR{$w$} دپاره تړلي فارمولونه~\LR{$!T(M,w)$} او~\LR{$!E(M,w)$} تعريف کړل۔
  بيا مو په~\olref[ver]{lem:valid-if-halt} او
  \olref[ver]{lem:halt-if-valid} کښې وښودل چې~\LR{$!T(M,w) \lif
  !E(M,w)$} هله او يوازې هله معتبر دے چې~\LR{$M$} پر ننوت~\LR{$w$} له پيلېدو
  وروسته بالاخره ودرېږي۔ څرنګه چې د درېدنې مسئله پرېکړه کېدونکې نۀ
  ده، له دې څخه پايله اخلو چې د لومړۍ درجې منطق د تړلي فارمولونه
  اعتبار او د صدق وړتيا پرېکړه کېدونکي نۀ دي
  (\cref{tur:und:uns:thm:decision-prob,tur:und:uns:cor:undecidable-sat})۔
\olpseganchor{OLP-0273-B005}{end}

\olpseganchor{OLP-0273-B006}{start}
  خو د جملو اعتبار او د صدق وړتيا د خپلو‌سرو
  جوړښتونه دپاره تعريف شوي، که هغه متناهي وي او که نامتناهي۔
  ښايي تاسو ګومان وکړئ چې دا پرېکړه اسانه ده چې آيا يو
  تړلے فارمول په يوه متناهي جوړښت کښې د صدق وړ دے (يا په
  ټولو متناهي جوړښتونه کښې معتبر دے)۔ د پرېکړې مسئلې د ناحلي
  ثبوت داسې برابرولے شو چې وښيي خبره داسې نۀ ده۔
\olpseganchor{OLP-0273-B006}{end}

\olpseganchor{OLP-0273-B007}{start}
  لومړے، که د~\olref[ver]{lem:halt-if-valid} ثبوت ته بېرته ورشئ،
  وبه وينئ چې هلته مو د~\LR{$!T(M,w)$} داسې مدل~\LR{$\Struct M$} جوړ کړ چې
  بالکل دا بيانوي چې ماشين~\LR{$M$} پر ننوت~\LR{$w$} له پيلېدو وروسته څه
  کوي۔ د هغه مدل ډومېن~\LR{$\Nat$} و، يعنې نامتناهي و۔ خو که~\LR{$M$} په
  رښتيا پر ننوت~\LR{$w$} ودرېږي، په همدې ډول يو متناهي مدل~\LR{$\Struct M'$}
  جوړولے شو۔ فرض کړئ~\LR{$M$} پر ننوت~\LR{$w$} له پيلېدو وروسته له~\LR{$k$}
  ګامونو وروسته درېږي۔ د ډومېن~\LR{$\Domain{M'}$} دپاره سټ~\LR{$\{0,
  \dots, n\}$} واخلئ، چې~\LR{$n$} د~\LR{$k$} او د~\LR{$w$} د اوږدوالي تر منځ لوے
  عدد دے، او داسې وټاکئ:
  \[
    \Assign{\prime}{M'}(x) =
    \begin{cases}
      x + 1 &\text{\RL{که \LR{$x < n$} وي}}\\
      n &\text{\RL{که نۀ وي،}}
    \end{cases}
  \]
  او~\LR{$\tuple{x,y} \in \Assign{<}{M'}$} هله او يوازې هله چې~\LR{$x < y$}
  يا~\LR{$x = y = n$} وي۔ له دې پرته~\LR{$\Struct{M'}$} بالکل د~\LR{$\Struct{M}$}
  په شان تعريفېږي۔ د~\LR{$\Struct{M'}$} د تعريف له مخې، لکه د
  \olref[ver]{lem:halt-if-valid} په ثبوت کښې،
  \LR{$\Sat{M'}{!T(M,w)}$}۔ او څرنګه چې مو فرض کړې ده چې~\LR{$M$} پر
  ننوت~\LR{$w$} درېږي، \LR{$\Sat{M'}{!E(M,w)}$}۔ نو~\LR{$\Struct{M'}$} د~\LR{$!T(M,w)
  \land !E(M,w)$} يو متناهي مدل دے (پام وکړئ چې~\LR{$\lif$} مو
  په~\LR{$\land$} بدل کړے دے)۔
\olpseganchor{OLP-0273-B007}{end}

\olpseganchor{OLP-0273-B008}{start}
  موږ د ثبوت نيمې لارې ته رسېدلي يو: ښودلې مو ده چې که~\LR{$M$} پر
  ننوت~\LR{$w$} ودرېږي، نو~\LR{$!T(M,w) \land !E(M,w)$} يو متناهي مدل لري۔
  له بده مرغه د دې نقيض عکس سم نۀ دے؛ يعنې داسې ټيورينګ ماشينونه
  شته چې پر کوم ننوت~\LR{$w$} نۀ درېږي، خو~\LR{$!T(M,w) \land !E(M,w)$} بيا
  هم يو متناهي مدل لري۔ د بېلګې په توګه هغه ماشين~\LR{$M$} په پام کښې
  ونيسئ چې يوازې يو حالت~\LR{$q_0$} او دا لارښوونه لري:
  \LR{$\delta(q_0,\TMblank) = \tuple{q_0,\TMblank,\TMstay}$}۔ دا ماشين
  پر تش ننوت~\LR{$w = \emptyseq$} له پيلېدو وروسته هېڅکله نۀ درېږي: په
  يوه نامتناهي کړۍ کښې دے، خو پټه نۀ بدلوي او سر نۀ خوځوي۔ ټول
  تشکيلونه يو شان دي (هماغه حالت، هماغه د سر موقعيت، هماغه د پټې
  منځپانګه)۔ داسې متناهي جوړښت~\LR{$\Struct{M''}$} تعريفولے شو
  چې~\LR{$!T(M,\emptyseq) \land !E(M,\emptyseq)$} پوره کړي (تمرين)۔ خو
  \LR{$!T(M,w)$} په مناسبه توګه بدلولے شو، څو داسې جوړښتونه ترې
  بهر کړو۔
\olpseganchor{OLP-0273-B008}{end}

\olpseganchor{OLP-0273-B009}{start}
  \begin{prob}
    فرض کړئ~\LR{$M$} داسې ټيورينګ ماشين دے چې يوازې حالت~\LR{$q_0$} او يوازې
    لارښوونه~\LR{$\delta(q_0,\TMblank) =
    \tuple{q_0,\TMblank,\TMstay}$} لري۔ داسې وټاکئ:
    \LR{$\Domain{M''} = \{0, 1, 2\}$}، \LR{$\Assign{\prime}{M''}(0) =
    \Assign{\prime}{M''}(1) = 1$} او~\LR{$\Assign{\prime}{M''}(2) = 2$}،
    او~\LR{$\Assign{<}{M''} = \{\tuple{0,1}, \tuple{1,1},
    \tuple{2,2}\}$}۔ \LR{$\Assign{\Obj Q_{q_0}}{M''}$}، \LR{$\Assign{\Obj
    S_{\TMblank}}{M''}$}، او~\LR{$\Assign{\Obj S_{\TMendtape}}{M''}$}
    داسې تعريف کړئ چې~\LR{$!T(M,\emptyseq)$} او~\LR{$!E(M,\emptyseq)$} رښتوني
    شي، او څرګنده کړئ چې ولې رښتوني دي۔ اشاره: پام وکړئ چې
    \LR{$\delta(q_0, \TMendtape)$} ناتعريفه ده۔ ډاډ ترلاسه کړئ چې
    \begin{align*}
    & \Obj Q_{q_0}(\num{1}, \num{n}) \land \Obj S_{\TMendtape}(\num{0}, \num{n})
      \land
      \lforall[x][(\num{0} < x \lif \Obj S_\TMblank(x, \num{n}))] \text{\RL{\quad د هر \LR{$n \in \Nat$} دپاره}}\\
    & \lexists[y][(\Obj Q_{q_0}(\num{0}, y) \land \Obj S_{\TMendtape}(\num{0}, y))]
    \end{align*}
    دواړه په~\LR{$\Struct{M''}$} کښې رښتوني وي۔
    \textbf{د ژباړې د نسخې سمون (OLCMP-063):} سرچينې د يوازېني
    حالت پر ځاے په لارښوونه کښې يو ناتعريفه حالت راوړے و؛ دلته هماغه
    ټاکل شوے يوازېنی حالت کارول شوے دے۔
    \textbf{د ژباړې د نسخې سمون (OLCMP-064):} سرچينې د يو د تالي
    ارزښت په تېروتنه د پخواني جوړښت په نوم پورې تړلے و؛ دلته د تمرين
    روان جوړښت په يو شان ډول کارول شوے دے۔
  \end{prob}
\end{explain}
\olpseganchor{OLP-0273-B009}{end}

\olpseganchor{OLP-0273-B010}{start}
هغه تړلي فارمولونه په پام کښې ونيسئ چې د ټيورينګ ماشين~\LR{$M$} عمل پر
ننوت~\LR{$w = \sigma_{i_1}\dots\sigma_{i_k}$} بيانوي:
\begin{enumerate}
\item هغه اکسيومونه چې اعداد او~\LR{$<$} بيانوي (بالکل لکه په
\olref[rep]{sec} کښې د~\LR{$!T(M,w)$} په تعريف کښې)۔
\item هغه اکسيومونه چې د ننوت تشکيل بيانوي: بالکل لکه د~\LR{$!T(M,w)$}
په تعريف کښې۔
\item هغه اکسيومونه چې له يوه تشکيل څخه بل ته د حالت بدلون بيانوي:
\olpseganchor{OLP-0273-B010}{end}

\olpseganchor{OLP-0273-B011}{start}
د لاندې برخې دپاره، \LR{$!A(x, y)$} دې د پخوا په شان وي، او داسې دې وي:
\[
  !B(y) \ident \lforall[x][(x < y \lif \eq/[x][y])].
\]
\begin{enumerate}
\item \ollabel{rep-right} د هرې لارښوونې~\LR{$\delta(q_i, \sigma) =
  \tuple{q_j, \sigma', \TMright}$} دپاره دا تړلے فارمول:
\begin{align*}
& \lforall[x][\lforall[y][(
   (\Obj Q_{q_i}(x, y) \land \Obj S_{\sigma}(x, y)) \lif {}]] \\
&\qquad   (\Obj Q_{q_j}(x', y') \land \Obj S_{\sigma'}(x, y') \land
!A(x, y) \land !B(y')))
\end{align*}
\item \ollabel{rep-left} د هرې لارښوونې~\LR{$\delta(q_i, \sigma) =
  \tuple{q_j, \sigma', \TMleft}$} دپاره دا تړلے فارمول:
\begin{align*}
& \lforall[x][\lforall[y][
    ((\Obj Q_{q_i}(x', y) \land \Obj S_{\sigma}(x', y)) \lif {}]]\\
& \qquad   (\Obj Q_{q_j}(x, y') \land \Obj S_{\sigma'}(x', y') \land
!A(x', y) \land !B(y'))) \land {}\\
& \lforall[y][((\Obj Q_{q_i}(\num{0}, y) \land \Obj S_{\sigma}(\num{0},
    y)) \lif {}]\\
& \qquad (\Obj Q_{q_j}(\num{0}, y') \land \Obj S_{\sigma'}(\num{0},
  y') \land !A(\num{0}, y) \land !B(y')))
\end{align*}
\textbf{د ژباړې د نسخې سمون (OLCMP-065):} د سرچينې د کيڼ حرکت
عادي څانګې بيا د نابدلې پټې په شرط کښې د ليکل کېدونکې مربع پر ځاے
منزل مستثنا کاوه؛ دلته د ليکل کېدونکې مربع سم آرګومان کارول شوے دے۔
\textbf{د ژباړې د نسخې سمون (OLCMP-066):} د سرچينې همدې عادي څانګې
هغه د بېلوالي شرط پرېښے و چې د دې درې واړو حالت-بدلونونو موخه ده؛
دلته ور زيات شوے، څو د بل تشکيل عدد له ټولو مخکښېنيو بېل وي۔
\item \ollabel{rep-stay} د هرې لارښوونې~\LR{$\delta(q_i, \sigma) =
  \tuple{q_j, \sigma', \TMstay}$} دپاره دا تړلے فارمول:
\begin{align*}
& \lforall[x][\lforall[y][(
   (\Obj Q_{q_i}(x, y) \land \Obj S_{\sigma}(x, y)) \lif {}]] \\
&\qquad (\Obj Q_{q_j}(x, y') \land \Obj S_{\sigma'}(x, y') \land
!A(x, y) \land !B(y')))
\end{align*}
\end{enumerate}
لکه چې وينئ، هغه تړلي فارمولونه چې د~\LR{$M$} حالت-بدلونونه بيانوي د
\LR{$!T(M,w)$} له اړوندو تړلے فارمول سره يو شان دي، يوازې په پای کښې
\LR{$!B(y')$} ور زياتوو۔ \LR{$!B(y')$} دا تضمينوي چې د ``بل'' تشکيل عدد~\LR{$y'$}
له ټولو مخکښېنيو عددونو~\LR{$\Obj 0$}، \LR{$\Obj 0'$}، \dots، څخه بېل وي۔
% \item Sentences that express consistency conditions:
% \begin{enumerate}
%   \item\ollabel{rep-con-symb}%
%   تړلے فارمول that says that no square can have more than one
%   symbol on it at any given time:
%   \[\lforall[x][\lforall[y][(\Obj S_{\sigma}(x, y) \lif \lnot \Obj
%   S_{\sigma'}(x, y))]],\]
%   for every pair $\sigma \neq \sigma'$ of tape symbols of~$T$.
%   \item\ollabel{rep-con-state}%
%   تړلے فارمول that says that $M$ cannot be in more than one state
%   at any given time:
%   \[\lforall[x][\lforall[y][(\Obj Q_{q}(x, y) \lif
%   \lforall[z][\lnot \Obj
%   Q_{q'}(z, y))]],\]
%   for every pair $q \neq q'$ of states of~$T$.
%   \item\ollabel{rep-con-tape}%
%   تړلے فارمول that says that $M$ cannot be on more than one tape
%   square at any given time:
%   \[\lforall[x][\lforall[y][(\Obj Q_{q}(x, y) \lif
%   \lforall[z][\eq/[z][y]\lif \lnot \Obj
%   Q_{q'}(z, y))]],\]
%   for every pair $q$, $q'$ of states of~$T$.
% \end{enumerate}
\end{enumerate}
\LR{$!T'(M, w)$} د ټيورينګ ماشين~\LR{$M$} او ننوت~\LR{$w$} دپاره د پورتنيو ټولو
تړلي فارمولونه عطف وي۔
\olpseganchor{OLP-0273-B011}{end}

\olpseganchor{OLP-0273-B012}{start}
\begin{lem}\ollabel{lem:halts-sat}
  که~\LR{$M$} پر ننوت~\LR{$w$} له پيلېدو وروسته ودرېږي، نو~\LR{$!T'(M,w) \land
  !E(M,w)$} يو متناهي مدل لري۔
\end{lem}
\olpseganchor{OLP-0273-B012}{end}

\olpseganchor{OLP-0273-B013}{start}
\begin{proof}
  \LR{$\Struct{M'}$} د~\olref[ver]{lem:halt-if-valid} په ثبوت کښې د
  جوړښت په شان واخلئ، خو دا توپيرونه ولري:
  \begin{align*}
    \Domain{M'} & = \{0, \dots, n\},\\
    \Assign{\prime}{M'}(x) & =
    \begin{cases}
      x + 1 &\text{\RL{که \LR{$x < n$} وي}}\\
      n &\text{\RL{که نۀ وي،}}
    \end{cases} \\
    \tuple{x,y} \in \Assign{<}{M'} &\text{\RL{هله او يوازې هله چې \LR{$x < y$} يا \LR{$x = y = n$} وي،}}
  \end{align*}
  چې~\LR{$n = \max(k,\len{w})$} او~\LR{$k$} تر ټولو کوچنی داسې عدد دے چې~\LR{$M$}
  پر ننوت~\LR{$w$} له پيلېدو وروسته له~\LR{$k$} ګامونو وروسته درېدلے وي۔ د دې
  خبرې تاييد د تمرين په توګه پرېږدو چې~\LR{$\Sat{M'}{!T'(M,w) \land
  !E(M,w)}$}۔
  \textbf{د ژباړې د نسخې سمون (OLCMP-067):} سرچينې په وروستي
  پوره-کېدو حکم کښې د درېدنې ټاکلې څيز-ژبې نښه پرېښې وه؛ دلته بشپړه
  جمله راوړل شوې ده۔
\end{proof}
\olpseganchor{OLP-0273-B013}{end}

\olpseganchor{OLP-0273-B014}{start}
\begin{prob}
  د~\olref[tur][und][tra]{lem:halts-sat} ثبوت په دې ښودلو بشپړ کړئ
  چې~\LR{$\Sat{M'}{!T'(M,w) \land !E(M,w)}$}۔
  \textbf{د ژباړې د نسخې سمون (OLCMP-068):} سرچينې په دې تمرين کښې
  د بدلې شوې تمثيلي جملې پر ځاے پخوانۍ جمله راوړې وه او د درېدنې د
  جملې ټاکلې نښه ئې پرېښې وه؛ دلته دواړه د لمې له بيان سره سم شوي دي۔
\end{prob}
\olpseganchor{OLP-0273-B014}{end}

\olpseganchor{OLP-0273-B015}{start}
\begin{lem}\ollabel{lem:sat-halts}
  که~\LR{$!T'(M,w) \land !E(M,w)$} يو متناهي مدل ولري، نو~\LR{$M$} پر
  ننوت~\LR{$w$} له پيلېدو وروسته درېږي۔
\end{lem}
\olpseganchor{OLP-0273-B015}{end}

\olpseganchor{OLP-0273-B016}{start}
\begin{proof}
  نقيض عکس ښيو۔ فرض کړئ~\LR{$M$} پر ننوت~\LR{$w$} له پيلېدو وروسته نۀ
  درېږي۔ که~\LR{$!T'(M,w) \land !E(M,w)$} هېڅ مدل ونۀ لري، خبره بشپړه
  ده۔ نو فرض کړئ~\LR{$\Struct{M}$} د~\LR{$!T'(M,w) \land !E(M, w)$} يو مدل
  دے۔ بايد وښيو چې دا متناهي کېدے نۀ شي۔
  \textbf{د ژباړې د نسخې سمون (OLCMP-069):} سرچينې دلته د روانې
  بدلې شوې جملې پر ځاے پخوانۍ تمثيلي جمله د مدل فرض ګرځولې وه؛ دلته
  فرض د لمې او د پاتې ثبوت له بدلې شوې جملې سره يو شان شوے دے۔
\olpseganchor{OLP-0273-B016}{end}

\olpseganchor{OLP-0273-B017}{start}
  بالکل د~\olref[ver]{lem:config} په شان ثابتولے شو چې که~\LR{$M$} پر
  ننوت~\LR{$w$} له پيلېدو وروسته تر~\LR{$n$} ګامونو پورې نۀ وي درېدلے، نو
  \LR{$!T'(M,w) \Entails !C(M, w, n) \land !B(\num{n})$}۔ څرنګه چې~\LR{$M$}
  پر ننوت~\LR{$w$} له پيلېدو وروسته نۀ درېږي، د هر~\LR{$n \in \Nat$} دپاره
  \LR{$!T'(M,w) \Entails !C(M, w, n) \land !B(\num{n})$}۔ پام وکړئ چې
  د~\olref[rep]{prop:mlessk} له مخې، د ټولو~\LR{$k < n$} دپاره~\LR{$!T'(M,w)
  \Entails \num{k} < \num{n}$}۔ همدارنګه~\LR{$!B(\num{n}) \Entails
  \num{k} < \num{n} \lif \eq/[\num{k}][\num{n}]$}۔ نو د ټولو~\LR{$k <
  n$} دپاره~\LR{$\Sat{M}{\eq/[\num{k}][\num{n}]}$}؛ يعنې نامتناهي ډېر
  ترمونه~\LR{$\num{k}$} بايد په~\LR{$\Struct{M}$} کښې ټول بېلابېل ارزښتونه
  ولري۔ خو دا غوښتنه کوي چې~\LR{$\Domain{M}$} نامتناهي وي؛ نو
  \LR{$\Struct{M}$} د~\LR{$!T'(M,w) \land !E(M, w)$} متناهي مدل نۀ شي
  کېدے۔
\end{proof}
\olpseganchor{OLP-0273-B017}{end}

\olpseganchor{OLP-0273-B018}{start}
\begin{prob}
  د~\olref[tur][und][tra]{lem:sat-halts} ثبوت په دې ښودلو بشپړ کړئ
  چې که~\LR{$M$} پر ننوت~\LR{$w$} له پيلېدو وروسته تر~\LR{$n$} ګامونو پورې نۀ وي
  درېدلے، نو~\LR{$!T'(M,w) \Entails !B(\num{n})$}۔
\end{prob}
\olpseganchor{OLP-0273-B018}{end}

\olpseganchor{OLP-0273-B019}{start}
\begin{thm}[د تراختن‌بروت قضيه]
  \ollabel{thm:trakhtenbrodt}
  دا پرېکړه کېدونکې نۀ ده چې آيا د لومړۍ درجې منطق کوم خپل‌سری
  تړلے فارمول يو متناهي مدل لري (يعنې په يوه متناهي جوړښت کښې د
  صدق وړ دے)۔
\end{thm}
\olpseganchor{OLP-0273-B019}{end}

\olpseganchor{OLP-0273-B020}{start}
\begin{proof}
  فرض کړئ داسې ټيورينګ ماشين~\LR{$F$} وے چې په يوه متناهي جوړښت کښې د
  صدق وړتيا پرېکړه کوي۔ بيا به مو د هر ټيورينګ ماشين~\LR{$M$} او
  ننوت~\LR{$w$} دپاره جمله~\LR{$!T'(M,w) \land !E(M,w)$} محاسبه کړې وے، او
  د~\LR{$F$} په وسيله به مو پرېکړه کړې وے چې آيا دا يو متناهي مدل لري۔
  د~\cref{tur:und:tra:lem:halts-sat,tur:und:tra:lem:sat-halts} له
  مخې، دا هله او يوازې هله متناهي مدل لري چې~\LR{$M$} پر ننوت~\LR{$w$} له
  پيلېدو وروسته ودرېږي۔ نو د~\LR{$F$} په وسيله به مو د درېدنې مسئله حل
  کړې وے، حال دا چې پوهېږو دا ناحل ده۔
\end{proof}
\olpseganchor{OLP-0273-B020}{end}

\olpseganchor{OLP-0273-B021}{start}
\begin{cor}\ollabel{cor:fproof-incomp}%
  داسې هېڅ اشتقاق نظام نشي موجودېدے چې د \emph{متناهي}
  اعتبار دپاره سم او بشپړ وي؛ يعنې داسې اشتقاق نظام چې
  \LR{$\Proves !B$} هله او يوازې هله ولري چې د هرې متناهي
  جوړښت~\LR{$\Struct{M}$} دپاره~\LR{$\Sat{M}{!B}$} وي۔
\end{cor}
\olpseganchor{OLP-0273-B021}{end}

\olpseganchor{OLP-0273-B022}{start}
\begin{proof}
  تمرين۔
\end{proof}
\olpseganchor{OLP-0273-B022}{end}

\olpseganchor{OLP-0273-B023}{start}
\begin{prob}
  \olref[tur][und][tra]{cor:fproof-incomp} ثابت کړئ۔ پام وکړئ چې~\LR{$!B$}
  په هره متناهي جوړښت کښې هله او يوازې هله پوره کېږي چې
  \LR{$\lnot !B$} په يوه متناهي جوړښت کښې د صدق وړ نۀ وي۔ څرګنده کړئ چې
  ولې په يوه متناهي جوړښت کښې د صدق وړتيا د
  \olref[tur][und][uns]{thm:valid-ce} په معنا نيمه د پرېکړې وړ ده۔
  له دې څخه دا استدلال وکړئ چې که د متناهي اعتبار دپاره کوم
  اشتقاق نظام وے، نو په يوه متناهي جوړښت کښې د صدق وړتيا
  به پرېکړه کېدونکې وه۔
\end{prob}
\olpseganchor{OLP-0273-B023}{end}

\olpseganchor{OLP-0274-B004}{start}
\olpart{inc}{ناتکميلي}
\olpseganchor{OLP-0274-B004}{end}

\olpseganchor{OLP-0274-B005}{start}
\begin{editorial}
  د دې برخې مواد د ناتکميلۍ قضيې رانغاړي۔ دا د لومړۍ درجې منطق د
  برخو (په ځانګړي ډول د ثبوت نظام) او د محاسبه‌پوهې په برخه کښې د
  بازګشتي تابعو پر موادو تکيه کوي۔ دا د جېرېمي اېوېګاډ پر يادښتونو
  ولاړ دي چې ريچرډ زاک سمونونه پکې کړي دي۔
\end{editorial}
\olpseganchor{OLP-0274-B005}{end}

\olpseganchor{OLP-0275-B004}{start}
\olchapter{inc}{int}{ناتکميلۍ ته پېژندنه}
\olpseganchor{OLP-0275-B004}{end}

\olpseganchor{OLP-0276-B004}{start}
\olfileid{inc}{int}{bgr}
\olpseganchor{OLP-0276-B004}{end}

\olpseganchor{OLP-0276-B005}{start}
\olsection{تاريخي شاليد}
\olpseganchor{OLP-0276-B005}{end}

\olpseganchor{OLP-0276-B006}{start}
په دې برخه کښې به پر هغو تاريخي پرمختګونو لنډ بحث وکړو چې د
ناتکميلۍ د قضيو د چاپېريال په پوهېدو کښې مرسته کوي۔ په ځانګړي ډول به
د رياضياتي منطق د تاريخ يوه ډېره لنډه ټوليزه کتنه وړاندې کړو، او بيا
به د رياضياتو د بنسټونو د تاريخ په باب څو خبرې وکړو۔
\olpseganchor{OLP-0276-B006}{end}

\olpseganchor{OLP-0276-B007}{start}
\begin{digress}
  د «رياضياتي منطق» غونډله دوه‌معناييزه ده۔ د «رياضياتي» ويے د موضوع
  د بيان په توګه تعبيرېدے شي، لکه «د رياضياتو منطق»، چې د رياضياتي
  استدلال اصول ښيي؛ يا د لارو د بيان په توګه، لکه «د منطق رياضيات»،
  چې د استدلال د اصولو رياضياتي مطالعه ښيي۔ لاندې روايت رياضياتي منطق
  په دواړو معناوو، او ډېر ځله په يوه وخت کښې، رانغاړي۔
\end{digress}
\olpseganchor{OLP-0276-B007}{end}

\olpseganchor{OLP-0276-B008}{start}
د منطق مطالعه په اصل کښې له ارسطو څخه پيل شوه، چې نږدې په
384--322 \textsc{ق‌م} کښې ژوندے و۔ د هغه \emph{مقولات}،
\emph{لومړۍ تحليلات} او \emph{دويمې تحليلات} د علمي استدلال د اصولو
منظمې مطالعې رانغاړي، چې د قياس بشپړه او منظمه مطالعه هم پکې شامله ده۔
\olpseganchor{OLP-0276-B008}{end}

\olpseganchor{OLP-0276-B009}{start}
د ارسطو منطق د منځنيو پېړيو په اوږدو کښې پر مدرسي فلسفې واکمن و؛ ان
تر اتلسمې پېړۍ پورې کانټ ټينګار کاوه چې د ارسطو منطق بشپړ دے او
سمون ته اړتيا نۀ لري۔ خو د قياس تيوري ډېره محدوده ده او د رياضياتي
استدلال له تر ټولو سطحي اړخونو پرته نور څه نۀ شي مدل کولے۔ يوه پېړۍ
تر دې وړاندې، د نيوټن هم‌مهالي لايبنيڅ د منطقي استدلال دپاره د يوه
بشپړ «حساب» تصور وکړ او د داسې حساب د جوړولو پر لور ئې ځينې لومړني
ګامونه واخيستل؛ په اصل کښې ئې د قضييز منطق يوه بڼه بيان کړه۔
\olpseganchor{OLP-0276-B009}{end}

\olpseganchor{OLP-0276-B010}{start}
نولسمه پېړۍ د منطق دپاره يو بنسټيز بدلون و۔ په 1854 کښې جورج بول
\emph{د فکر قانونونه} وليکل، چې د قضييز منطق بشپړه الجبري مطالعه
پکې وه او له اوسنيو وړاندېونو څخه ډېره لرې نۀ ده۔ په 1879 کښې ګوټلوب
فرېګې خپله \emph{بېګريفشريفټ} (مفهومي ليکدود) خپره کړه، چې قضييز
منطق د کميت‌ټاکونکو او اړيکو په ورزياتولو غځوي او په دې توګه د لومړۍ
درجې منطق هم رانغاړي۔ په حقيقت کښې د فرېګې منطقي نظامونو د لوړې درجې
منطق او نور څه هم رانغاړل۔ فرېګې په خپلو \emph{د حساب بنسټيزو
قانونونو} کښې هڅه وکړه وښيي چې ټول حساب د هغه په بېګريفشريفټ کښې له
سوچه منطقي فرضيو څخه استنتاجېدے شي۔ له بده مرغه دا فرضيې ناسازګارې
وختې، لکه رسل چې په 1902 کښې وښودل۔ خو که ناسازګار بديهي اصل څنګ ته
کړو، فرېګې نږدې په يوازې ځان اوسنی منطق اختراع کړ، چې يوه حيرانوونکې
لاسته راوړنه ده۔ له بول وروسته الجبري فکر لرونکو پوهانو، لکه پيرس او
شروډر، کميتي منطق په خپلواکه توګه هم جوړ کړ۔
\olpseganchor{OLP-0276-B010}{end}

\olpseganchor{OLP-0276-B011}{start}
اوس راځئ د رياضياتو د بنسټونو پرمختګونو ته مخه کړو۔ البته څرنګه چې
منطق په رياضياتو کښې مهم رول لري، له پاس بيان شويو پرمختګونو سره ئې
ډېر متقابل اغېز شته۔ د بېلګې په توګه، فرېګې خپل منطق په دې څرګنده
موخه جوړ کړ چې وښيي ټول رياضيات يوازې پر هغه منطقي چوکاټ ولاړېدے شي؛
په ځانګړي ډول غوښتل ئې وښيي چې رياضيات، د کانټ د نظر پر خلاف، له
تجربې وړاندې \emph{تحليلي} حقايق دي، نه له تجربې وړاندې
\emph{ترکيبي} حقايق۔
\olpseganchor{OLP-0276-B011}{end}

\olpseganchor{OLP-0276-B012}{start}
ډېر خلک د رياضياتو اصلي زېږون له يونانيانو سره تړي۔ د اقليدس
\emph{اصول}، چې نږدې 300 \textsc{ق‌م} کښې ليکل شوي، لا له وړاندې د
يوناني رياضياتو يوه پخه بېلګه ده او پر کره‌والي او دقت ټينګار کوي۔
د اقليدس د \emph{اصولو} تعريفونه او ثبوتونه د نن ورځې د لېسې د هندسې
په درسي کتابونو کښې نږدې هماغسې پاتې دي (تر هغه ځايه چې لا په لېسو
کښې هندسه تدريسېږي)۔ د رياضياتي استدلال دا مدل نه يوازې په رياضياتو
بلکې د فلسفې په څانګو کښې هم د کره استدلال بېلګه ګڼل شوې ده۔
(سپينوزا ان اخلاقي او ديني استدلالونه د اقليدسي سبک له مخې وړاندې
کړل، چې ليدل ئې عجيبه دي!)
\olpseganchor{OLP-0276-B012}{end}

\olpseganchor{OLP-0276-B013}{start}
نيوټن او لايبنيڅ په اوولسمه پېړۍ کښې کالکولس اختراع کړ۔ (د لومړيتوب
يوه سخته شخړه پېړۍ پېړۍ روانه وه، خو اوس زياتره پوهان مني چې دواړه
پرمختګونه تر ډېره خپلواک وو۔) په کالکولس کښې، د بېلګې په توګه، د
نامتناهي وړو مقدارونو د نامتناهي جمعو په باب استدلال راځي؛ دغو
ځانګړتياوو د بيشپ برکلي نيوکې راوپارولې، چې استدلال ئې کاوه پر خدای
باور د هغه د وخت تر رياضياتو لږ معقول نۀ دے۔ په اتلسمه پېړۍ کښې د
کالکولس لارې په پراخه توګه وکارول شوې؛ مثلاً ليونارډ اويلر د
نامتناهي جمعو محاسبې په حيرانوونکو پايلو وکارولې۔
\olpseganchor{OLP-0276-B013}{end}

\olpseganchor{OLP-0276-B014}{start}
په نولسمه پېړۍ کښې رياضيپوهانو هڅه وکړه د کالکولس د لا ټينګ بنسټ په
جوړولو د برکلي نيوکو ته ځواب ووايي۔ د کوشي، وايرشتراس، بولتسانو او
نورو هڅو موږ د حد، پيوستون، تفاضل او تکامل اوسنيو تعريفونو ته
ورسولو، چې د «اېپسيلونونو او ډېلټاګانو» په اصطلاحاتو کښې دي؛ يعنې
نامتناهي وړو مقدارونو ته هېڅ مراجعه نۀ کوي۔ د پېړۍ په وروستيو کښې
رياضيپوهانو هڅه وکړه لا وړاندې لاړ شي او د حقيقي عددونو په ګډون د
کالکولس ټول اړخونه د طبيعي عددونو په اصطلاحاتو کښې بيان کړي۔
(کرونېکر: «خدای صحيح عددونه پيدا کړل، نور هر څه د انسان کار دے۔»)
په 1872 کښې ډېدېکينډ «پيوستون او غيرناطق عددونه» وليکل، چې وئې ښودل
څنګه حقيقي عددونه د ناطقو عددونو د سټونو په توګه «جوړېدے» شي (او
ناطق عددونه، لکه چې پوهېږئ، د طبيعي عددونو د جوړو په توګه ليدل کېدے
شي)؛ په 1888 کښې ئې «Was sind und was sollen die Zahlen» (نږدې معنا
ئې «طبيعي عددونه څه دي، او څه بايد وي؟») وليکل، چې موخه ئې په سوچه
«منطقي» اصطلاحاتو کښې د طبيعي عددونو بيان و۔ په 1887 کښې کرونېکر
«\"Uber den Zahlbegriff» («د عدد د مفهوم په باب») وليکل او د صحيح
عددونو په اصطلاحاتو کښې ئې د ټولو رياضياتي څيزونو د تمثيل خبره وکړه؛
په 1889 کښې جوزېپې پيانو د طبيعي عددونو صوري، نښيز بديهي اصول ورکړل۔
\olpseganchor{OLP-0276-B014}{end}

\olpseganchor{OLP-0276-B015}{start}
د نولسمې پېړۍ پاى د نامتناهي په چلند کښې يو نوی زړورتوب هم راوړ۔
تر دې وړاندې له نامتناهي څيزونو او جوړښتونو (لکه د طبيعي عددونو له
سټ) سره ډېر په احتياط چلند کېده؛ «نامتناهي ډېر» د «څومره چې وغواړئ
هومره» په معنا، او «په حد کښې ورنږدې کېږي» د «څومره چې وغواړئ هومره
نږدې کېږي» په معنا ګڼل کېده۔ خو ګيورګ کانتور وښودل چې نامتناهي په
خپله ظاهري معنا اخيستل کېدے شي۔ د کانتور، ډېدېکينډ او نورو کار د
رياضياتو هغه عمومي سټ‌تيوريکي پوهاوی رامنځته کړ چې اوس په پراخه توګه
منل شوے دے۔
\olpseganchor{OLP-0276-B015}{end}

\olpseganchor{OLP-0276-B016}{start}
دا مو په منطق او بنسټونو کښې د شلمې پېړۍ پرمختګونو ته رسوي۔ په 1902
کښې رسل د فرېګې په منطقي نظام کښې پاراډوکس وموند۔ په 1904 کښې
څېر مېلو د کانتور د ښه‌ترتيب اصل د تش په نوم «د انتخاب بديهي اصل» په
کارولو ثابت کړ؛ د دې بديهي اصل د رواوالي په باب ډېر بحث وشو۔ د 1910
او 1913 تر منځ د رسل او وايټ‌هېډ د \emph{پرېنسيپيا مېتېماتيکا} درې
ټوکونه خپاره شول، چې د رياضياتو پر منطقي بنسټونو د درولو د فرېګې
پروګرام ئې وغځاوه۔ له بده مرغه رسل او وايټ‌هېډ اړ شول دوه داسې اصول
ومني چې د سوچه منطقي اصولو په توګه ئې توجيه ستونزمنه ښکارېده: د
نامتناهي والي بديهي اصل او د «راکمېدنې» بديهي اصل۔ په 1900مو کلونو
کښې پوانکاره په رياضياتو کښې د «غيرپريديکاتي تعريفونو» پر کارولو
نيوکه وکړه، او په 1910مو کلونو کښې بروور د ټولو رياضياتو د بيا بنسټيز
کولو دپاره د داسې «شهودي» بنسټ وړانديز پيل کړ چې د خارج وسط اصل
(\LR{$!A \lor \lnot !A$}) پکې نۀ کارېده۔
\olpseganchor{OLP-0276-B016}{end}

\olpseganchor{OLP-0276-B017}{start}
په رښتيا هم عجيب وختونه وو!{} منطق ته د ټولو رياضياتو د راکمولو
پروګرام اوس «منطقپالنه» بلل کېږي او د پاسنيو ستونزو له امله عموماً
ناکام ګڼل کېږي۔ د شهودي ذهني جوړونو په اصطلاحاتو کښې د رياضياتو د
جوړولو پروګرام «شهوديت» بلل کېږي او پر ورځنيو رياضياتو ډېر سخت
محدوديتونه ايښوونکی ګڼل کېږي۔ د پېړۍ د بدلېدو شاوخوا ډېوېډ هېلبرټ،
چې د هر وخت له تر ټولو اغېزمنو رياضيپوهانو څخه و، د کانتور او
ډېدېکينډ د نويو مجردو لارو ټينګ ملاتړی و: «هېڅوک به موږ له هغه جنت
څخه ونه باسي چې کانتور راته جوړ کړے دے۔» په همدې وخت کښې د دې نويو
لارو (چې عجيبه دا ده اوس «کلاسیکې» بلل کېږي) بنسټيزو نيوکو ته هم
متوجه و۔ هغه داسې لار وړانديز کړه چې دواړه موخې تر لاسه کړي:
\begin{enumerate}
\item کلاسیکې لارې په صوري بديهي اصولو او قاعدو تمثيل کړئ؛ رياضياتي
  پوښتنې په يوه بديهي نظام کښې د فارمولونه په توګه تمثيل کړئ۔
\item په خوندي، «متناهي‌پالو» لارو ثابته کړئ چې دا صوري استنتاجي
  نظامونه سازګار دي۔
\end{enumerate}
\olpseganchor{OLP-0276-B017}{end}

\olpseganchor{OLP-0276-B018}{start}
د هېلبرټ کار د لومړۍ موخې د پوره کولو پر لور ډېر وړاندې لاړ۔ په 1899
کښې ئې دا کار په خپل نامتو کتاب \emph{د هندسې بنسټونه} کښې د هندسې
دپاره کړے و۔ په ورپسې کلونو کښې هغه او يو شمېر زده‌کوونکو او
همکارانو ئې د رياضياتو پر نورو څانګو کار وکړ، څو هغه څه وکړي چې
هېلبرټ د هندسې دپاره کړي وو۔ هېلبرټ پخپله د حساب او تحليل دپاره
بديهي نظامونه ورکړل۔ څېر مېلو د سټ تيوري بديهي کول ورکړل، چې
فرېنکل، سکولېم، فون نويمان او نورو وغځول۔ د 1920مو کلونو تر نيمايي
دوه داسې لارې وې چې د «ټولو» رياضياتو د بديهي کولو نوم ئې غوښت:
د رسل او وايټ‌هېډ \emph{پرېنسيپيا مېتېماتيکا} او هغه څه چې وروسته
د څېر مېلو--فرېنکل سټ تيوري په نوم وپېژندل شول۔
\olpseganchor{OLP-0276-B018}{end}

\olpseganchor{OLP-0276-B019}{start}
په 1921 کښې هېلبرټ يوه څېړنيزه پروژه پيل کړه څو د دې نظامونو د
سازګارۍ د ثبوت موخه پوره کړي۔ د هغه څو زده‌کوونکو، په ځانګړي ډول
برنايس، اکرمان، او وروسته ګېنڅن، په دې پروژه کښې ورسره مرسته وکړه۔
د دې موخې بنسټيز نظر دا و چې په رياضياتو کښې د ناسازګارۍ د
اشتقاق د امکان پوښتنه د نښو د ممکنه لړيو په باب يو
ترکيبي مسئله وګرځوي؛ يعنې د هغو جملو ممکنه لړۍ چې د حساب، تحليل يا
سټ تيورۍ د يوه بديهي نظام له بديهي اصولو څخه د، مثلاً،
\LR{$!A \land \lnot !A$} د سم اشتقاق معيار پوره کوي۔ د نښو د
داسې لړۍ د ناممکن والي ثبوت، څرنګه چې پخپله رياضياتي ثبوت دے، په دې
بديهي نظامونو کښې صوري کېدے شي۔ په بله وينا، داسې يوه جمله~\LR{$\OCon$}
به وي چې وايي، مثلاً، حساب سازګار دے۔ سربېره پر دې، دا جمله بايد په
اړوندو نظامونو کښې د اثبات وړ وي، په ځانګړي ډول که ثبوت ئې يوازې ډېرې
محدودې، «متناهي‌پالې» لارې غواړي۔
\olpseganchor{OLP-0276-B019}{end}

\olpseganchor{OLP-0276-B020}{start}
دويمه موخه، چې جوړ شوي بديهي نظامونه به هره رياضياتي پوښتنه حل کړي،
په دوو لارو کره کېدے شي۔ په يوه لار داسې بيانېدے شي: د رياضياتو د
بديهي نظام د ژبې د هرې جملې~\LR{$!A$} دپاره يا~\LR{$!A$} يا~\LR{$\lnot !A$} له
بديهي اصولو څخه د اثبات وړ وي۔ که دا رښتيا وے، داسې جمله به نۀ وه چې
د بديهي اصولو پر بنسټ نه ثابته او نه ردېدے شي، او داسې پوښتنه به نۀ
وه چې بديهي اصول ئې حل نۀ کړي۔ دې خاصيت لرونکی بديهي نظام
\emph{بشپړ} بلل کېږي۔ البته د هرې ټاکلې جملې دپاره لا هم دا معلومول
ستونزمن کار کېدے شي چې له دوو بديلونو څخه کوم يو سم دے۔ خو په اصل
کښې بايد د معلومولو يوه لار وي۔ په حقيقت کښې د هېلبرټ تر څېړنې لاندې
د بديهي اصولو او اشتقاق نظامونو دپاره بشپړتيا دا لازموله چې
داسې يوه لار شته، که څه هم هېلبرټ په دې نۀ پوهېده۔ د پوښتنې دويم
تعبير دا لا پياوړې غوښتنه ده: داسې ميخانيکي، محاسبوي لار موجوده وي
چې د يوې ورکړل شوې جملې~\LR{$!A$} دپاره معلومه کړي له بديهي اصولو څخه
استنتاجېدونکې ده که نۀ۔
\olpseganchor{OLP-0276-B020}{end}

\olpseganchor{OLP-0276-B021}{start}
په 1931 کښې ګوډل د «ناتکميلۍ دوه قضيې» ثابتې کړې، چې وئې ښودل دا
پروګرام بريالی کېدے نۀ شي۔ د رياضياتو دپاره هېڅ داسې بديهي نظام نشته
چې بشپړ وي؛ په ځانګړي ډول هغه جمله چې د بديهي اصولو سازګاري بيانوي
نه ثابتېدے او نه ردېدے شي۔
\olpseganchor{OLP-0276-B021}{end}

\olpseganchor{OLP-0276-B022}{start}
دې د هېلبرټ اصلي پروګرام ته وژونکی ګوزار ورکړ۔ خو لکه په رياضياتو
کښې چې ډېر ځله کېږي، د څېړنې په زړه پورې نوې لارې ئې هم پرانيستې۔
که د ټولو رياضياتو يو واحد، هر اړخيز صوري نظام نشته، نو معنا لري چې
لا محدود نظامونه جوړ او وڅېړل شي چې څه پکې ثابتېدے شي۔ همدارنګه معنا
لري چې د دې نظامونو د سازګارۍ د ثابتولو دپاره لږ محدودې ثبوتي لارې
جوړې، او د هغو د سازګارۍ د ثبوت د سختوالي د اندازه کولو لارې وموندل
شي۔ څرنګه چې ګوډل وښودل (نږدې) هر صوري نظام داسې پوښتنې لري چې نۀ
ئې شي حل کولے، نو معنا لري چې د هغو «په زړه پورې» پوښتنو لټون وشي
چې يو ورکړل شوے صوري نظام ئې نۀ شي حل کولے، او معلومه شي چې د حل
دپاره ئې صوري نظام بايد څومره پياوړے وي۔ منطقپوهان تر نن ورځې دا
پوښتنې د رياضياتو په يوه نوې څانګه، د ثبوتونو په تيورۍ، کښې څېړي۔
\olpseganchor{OLP-0276-B022}{end}

\olpseganchor{OLP-0277-B004}{start}
\olfileid{inc}{int}{def}
\olpseganchor{OLP-0277-B004}{end}

\olpseganchor{OLP-0277-B005}{start}
\olsection{تعريفونه}
\olpseganchor{OLP-0277-B005}{end}

\olpseganchor{OLP-0277-B006}{start}
د رياضياتو د صوري کولو، او د داسې صوري کولو د سازګارۍ او بشپړتيا د
ښودلو د هېلبرټ پروژې د عملي کولو دپاره لومړنی کار دا و چې يوه ژبه،
منطقي چوکاټ او د بديهي اصولو نظام وټاکل شي۔ د خپلې موخې دپاره به فرض
کړو چې رياضيات د لومړۍ درجې په يوه ژبه کښې صوري کېدے شي؛ يعنې د
ثابت سمبولونه، تابعې او پريديکاتونه داسې يو سټ شته چې د
لومړۍ درجې منطق له تړونکو او کميت‌ټاکونکو سره يوځاے موږ ته د رياضياتو
د ادعاوو د بيان وس راکوي۔ زياتره خلک مني چې داسې ژبه شته: د سټ تيورۍ
ژبه، چې~\LR{$\in$} پکې يوازينۍ غيرمنطقي نښه ده۔ دا چې دومره ساده ژبه دومره
بيانوونکې ده، البته په لومړي نظر ډېره نامعقوله ادعا ښکاري، او د دې د
ثابتولو دپاره ډېر کار وشو چې نږدې ټول رياضيات په همدې ډېرې محدودې
اصطلاح‌زېرمه کښې بيانېدے شي۔ د ساده‌والي دپاره به اوس خپل بحث حساب
ته محدود کړو، يعنې د رياضياتو هغې برخې ته چې يوازې له طبيعي
عددونو~\LR{$\Nat$} سره کار لري۔ د حساب د حقايقو د بيان طبيعي ژبه~\LR{$\Lang
L_A$} ده۔ \LR{$\Lang L_A$} يو دوه‌ځاييز پريديکات~\LR{$<$}، يو
ثابت سمبول~\LR{$\Obj 0$}، يوه يوځاييزه تابع~\LR{$\prime$}، او دوه
دوه‌ځاييزې تابعې~\LR{$+$} او~\LR{$\times$} لري۔
\olpseganchor{OLP-0277-B006}{end}

\olpseganchor{OLP-0277-B007}{start}
\begin{defn}
د تړلي فارمولونه سټ~\LR{$\Gamma$} يوه \emph{تيوري} ده که د لازمې پايلې له
مخې تړلې وي، يعنې که
\LR{$\Gamma = \Setabs{!A}{\Gamma \Entails !A}$} وي۔
\end{defn}
\olpseganchor{OLP-0277-B007}{end}

\olpseganchor{OLP-0277-B008}{start}
د تيوريو د مشخصولو دوه اسانه لارې شته۔ يوه دا ده چې تيوري په يوه
جوړښت کښې د رښتينو تړلي فارمولونه سټ وټاکو۔ د بېلګې په توګه،
د~\LR{$\Lang L_A$} هغه جوړښت وګورئ چې د تعريف ساحه ئې~\LR{$\Nat$} وي او
ټولې غيرمنطقي نښې پکې هغسې تفسير شوې وي لکه تمه چې ترې کېږي۔
\olpseganchor{OLP-0277-B008}{end}

\olpseganchor{OLP-0277-B009}{start}
\begin{defn}\ollabel{def:standard-model}
د حساب \emph{معياري مدل} هغه جوړښت~\LR{$\Struct N$} دے چې داسې
تعريفېږي:
\begin{enumerate}
\item \LR{$\Domain N = \Nat$}
\item \LR{$\Assign{\Obj 0}{N} = 0$}
\item \LR{$\Assign{\Obj \prime}{N}(n) = n + 1$} د ټولو \LR{$n \in \Nat$} دپاره
\item \LR{$\Assign{\Obj +}{N}(n, m) = n + m$} د ټولو \LR{$n, m \in \Nat$} دپاره
\item \LR{$\Assign{\Obj \times}{N}(n, m) = n\cdot m$} د ټولو \LR{$n, m \in \Nat$} دپاره
\item \LR{$\Assign{\Obj <}{N} = \Setabs{\tuple{n, m}}{n \in \Nat, m \in
  \Nat, n < m}$}
\end{enumerate}
\end{defn}
\olpseganchor{OLP-0277-B009}{end}

\olpseganchor{OLP-0277-B010}{start}
د~`\LR{$\times$}' او~`\LR{$\cdot$}' تر منځ توپير ته پام وکړئ: \LR{$\times$} د حساب
د ژبې يوه نښه ده۔ البته دا مو داسې ټاکلې چې ضرب راياد کړي، خو
\LR{$\times$} د ضرب عمل نۀ دے، بلکې يوه دوه‌ځاييزه تابع‌نښه ده (په رسمي
ډول~\LR{$\Obj f^2_1$})۔ د دې پر خلاف، \LR{$\cdot$} پخپله د ضرب عادي تابع
\emph{ده}۔ کله چې د~\LR{$n \cdot m$} په شان څه وينئ، موخه مو د~\LR{$n$} او~\LR{$m$}
عددونو حاصل‌ضرب وي؛ کله چې د~\LR{$x \times y$} په شان څه وينئ، د حساب د
ژبې د يوې اصطلاح خبره کوو۔ په معياري مدل کښې تابع‌نښه~\LR{$\times$} پر
طبيعي عددونو د تابع~\LR{$\cdot$} په توګه تفسيرېږي۔ د جمع دپاره~\LR{$+$} هم د
حساب د ژبې د تابع‌نښې او هم پر طبيعي عددونو د جمع د تابع په توګه
کاروو۔ دلته بايد له چاپېرياله معلومه کړئ چې کومه معنا موخه ده۔
\olpseganchor{OLP-0277-B010}{end}

\olpseganchor{OLP-0277-B011}{start}
\begin{defn}
د \emph{رښتيني حساب} تيوري د هغو تړلي فارمولونه سټ دے چې د حساب په
معياري مدل کښې پوره کېږي، يعنې
\[
\Th{TA} = \Setabs{!A}{\Sat{N}{!A}}.
\]
\end{defn}
\olpseganchor{OLP-0277-B011}{end}

\olpseganchor{OLP-0277-B012}{start}
\LR{$\Th{TA}$} يوه تيوري ده، ځکه هر کله چې~\LR{$\Th{TA} \Entails !A$} وي،
\LR{$!A$} په هر هغه جوړښت کښې پوره کېږي چې~\LR{$\Th{TA}$} پوره کوي۔
څرنګه چې~\LR{$\Sat{N}{\Th{TA}}$}، نو~\LR{$\Sat{N}{!A}$} لرو، او له دې امله
\LR{$!A \in \Th{TA}$}۔
\olpseganchor{OLP-0277-B012}{end}

\olpseganchor{OLP-0277-B013}{start}
د تيورۍ~\LR{$\Gamma$} د مشخصولو بله لار دا ده چې هغه د هغو تړلي فارمولونه سټ وټاکو
چې د جملو له يوه سټ~\LR{$\Gamma_0$} څخه لازمېږي۔ په دې حالت کښې~\LR{$\Gamma$}
د لازمې پايلې له مخې د~\LR{$\Gamma_0$} «تړون» دے۔ د تيورۍ دا ډول مشخصول
يوازې هغه وخت په زړه پورې دي چې~\LR{$\Gamma_0$} په څرګند ډول مشخص شوے وي؛
مثلاً د~\LR{$\Gamma_0$} غړي لست شوي وي۔ لږ تر لږه~\LR{$\Gamma_0$}
بايد د پرېکړې وړ وي؛ يعنې داسې محاسبه کېدونکې ازموينه بايد وي چې يوه
تړلے فارمول د~\LR{$\Gamma_0$} غړې ده که نۀ۔ د~\LR{$\Gamma_0$} تړلي فارمولونه
د~\LR{$\Gamma$} دپاره \emph{بديهي اصول} بولو، او وايو چې~\LR{$\Gamma$} د
\LR{$\Gamma_0$} په وسيله \emph{بديهي شوې} ده۔
\olpseganchor{OLP-0277-B013}{end}

\olpseganchor{OLP-0277-B014}{start}
\begin{defn}
تيوري~\LR{$\Gamma$} هله او يوازې هله د~\LR{$\Gamma_0$} په وسيله بديهي شوې ده
چې
\[
\Gamma = \Setabs{!A}{\Gamma_0 \Entails !A}
\]
\end{defn}
\olpseganchor{OLP-0277-B014}{end}

\olpseganchor{OLP-0277-B015}{start}
\begin{defn}
هغه تيوري~\LR{$\Th{Q}$} چې په لاندې جملو بديهي شوې ده د «رابنسن
\LR{$\Th{Q}$}» په نوم پېژندل کېږي او د حساب يوه ډېره ساده تيوري ده۔
\begin{align*}
& \lforall[x][\lforall[y][(\eq[x'][y'] \lif \eq[x][y])]] \tag{\LR{$!Q_1$}}\\
& \lforall[x][\eq/[\Obj 0][x']] \tag{\LR{$!Q_2$}}\\
& \lforall[x][(\eq[x][\Obj 0] \lor \lexists[y][\eq[x][y']])] \tag{\LR{$!Q_3$}}\\
& \lforall[x][\eq[(x + \Obj 0)][x]] \tag{\LR{$!Q_4$}}\\
& \lforall[x][\lforall[y][\eq[(x + y')][(x + y)']]] \tag{\LR{$!Q_5$}}\\
& \lforall[x][\eq[(x \times \Obj 0)][\Obj 0]] \tag{\LR{$!Q_6$}}\\
& \lforall[x][\lforall[y][\eq[(x \times y')][((x \times y) + x)]]] \tag{\LR{$!Q_7$}}\\
& \lforall[x][\lforall[y][(x < y \liff \lexists[z][\eq[(z' + x)][y]])]] \tag{\LR{$!Q_8$}}
\end{align*}
د تړلي فارمولونه سټ~\LR{$\{!Q_1, \dots, !Q_8\}$} د~\LR{$\Th{Q}$} بديهي اصول
دي، نو~\LR{$\Th{Q}$} له هغو څخه لازمېدونکي ټول تړلي فارمولونه رانغاړي:
\[
\Th{Q} = \Setabs{!A}{\{!Q_1, \dots, !Q_8\} \Entails !A}.
\]
\end{defn}
\olpseganchor{OLP-0277-B015}{end}

\olpseganchor{OLP-0277-B016}{start}
\begin{defn}
فرض کړئ~\LR{$!A(x)$} په~\LR{$\Lang L_A$} کښې داسې فارمول ده چې آزاد
متغيرونه~\LR{$x$} او~\LR{$y_1$}، \dots،~\LR{$y_n$} لري۔ بيا د لاندې بڼې هره
تړلے فارمول
\[
\lforall[y_1][\dots\lforall[y_n][((!A(\Obj 0) \land \lforall[x][(!A(x)
\lif !A(x'))]) \lif \lforall[x][!A(x)])]]
\]
د \emph{استقرا شېما} يوه بېلګه ده۔
\olpseganchor{OLP-0277-B016}{end}

\olpseganchor{OLP-0277-B017}{start}
\emph{د پيانو حساب}~\LR{$\Th{PA}$} هغه تيوري ده چې د~\LR{$\Th{Q}$} په بديهي
اصولو او د استقرا شېما په ټولو بېلګو بديهي شوې ده۔
\end{defn}
\olpseganchor{OLP-0277-B017}{end}

\olpseganchor{OLP-0277-B018}{start}
\begin{explain}
د استقرا شېما هره بېلګه په~\LR{$\Struct{N}$} کښې رښتونې ده۔ دا خبره تر
ټولو اسانه هغه وخت ليدل کېږي چې فارمول~\LR{$!A$} يوازې يو آزاد
متغير~\LR{$x$} ولري۔ بيا~\LR{$!A(x)$} په~\LR{$\Struct{N}$} کښې د~\LR{$\Nat$} يو
فرعي سټ~\LR{$X_{!A}$} تعريفوي۔ \LR{$X_{!A}$} د ټولو هغو~\LR{$n \in \Nat$} سټ دے
چې~\LR{$s(x) = n$} وي او~\LR{$\Sat{N}{!A(x)}[s]$}۔ د استقرا شېما اړونده بېلګه
دا ده:
\[
((!A(\Obj 0) \land \lforall[x][(!A(x) \lif !A(x'))]) \lif
  \lforall[x][!A(x)]).
\]
که مقدمه ئې په~\LR{$\Struct{N}$} کښې رښتونې وي، نو~\LR{$0 \in X_{!A}$}، او
هر کله چې~\LR{$n \in X_{!A}$} وي، \LR{$n+1$} هم پکې وي۔ له~\LR{$0 \in X_{!A}$}
څخه~\LR{$1 \in X_{!A}$} ترلاسه کوو۔ له~\LR{$1 \in X_{!A}$} څخه~\LR{$2 \in X_{!A}$}
ترلاسه کوو، او همداسې وړاندې۔ نو د هر~\LR{$n \in \Nat$} دپاره
\LR{$n \in X_{!A}$}۔ خو دا په دې معنا ده چې~\LR{$\lforall[x][!A(x)]$} په
\LR{$\Struct{N}$} کښې پوره کېږي۔
\end{explain}
\olpseganchor{OLP-0277-B018}{end}

\olpseganchor{OLP-0277-B019}{start}
\LR{$\Th{Q}$} او~\LR{$\Th{PA}$} دواړه بديهي شوې تيورۍ دي۔ لويه پوښتنه دا ده
چې څومره پياوړې دي؟ مثلاً آيا~\LR{$\Th{PA}$} د~\LR{$\Nat$} په باب ټول هغه
حقايق ثابتولے شي چې په~\LR{$\Lang L_A$} کښې بيانېدے شي؟ په کره توګه، آيا
د~\LR{$\Th{PA}$} بديهي اصول ټولې هغه پوښتنې حل کوي چې په~\LR{$\Lang L_A$} کښې
جوړېدے شي؟
\olpseganchor{OLP-0277-B019}{end}

\olpseganchor{OLP-0277-B020}{start}
د دې پوښتنې بله بڼه دا ده: آيا~\LR{$\Th{PA} = \Th{TA}$}؟ \LR{$\Th{TA}$} په
څرګند ډول د~\LR{$\Nat$} په باب ټول حقايق ثابتوي (يعنې ټول رانغاړي)، او
ټولې هغه پوښتنې حل کوي چې په~\LR{$\Lang L_A$} کښې جوړېدے شي؛ ځکه که~\LR{$!A$}
په~\LR{$\Lang L_A$} کښې تړلے فارمول وي، نو يا~\LR{$\Sat{N}{!A}$} يا
\LR{$\Sat{N}{\lnot !A}$}، او له دې امله يا~\LR{$\Th{TA} \Entails !A$} يا
\LR{$\Th{TA} \Entails \lnot !A$}۔ داسې تيوري \emph{بشپړ} بولو۔
\olpseganchor{OLP-0277-B020}{end}

\olpseganchor{OLP-0277-B021}{start}
\begin{defn}
تيوري~\LR{$\Gamma$} هله او يوازې هله \emph{بشپړ} ده چې د هغې د
ژبې د هرې تړلے فارمول~\LR{$!A$} دپاره يا~\LR{$\Gamma \Entails !A$} يا
\LR{$\Gamma \Entails \lnot !A$} وي۔
\end{defn}
\olpseganchor{OLP-0277-B021}{end}

\olpseganchor{OLP-0277-B022}{start}
\begin{explain}
د بشپړتيا د قضيې له مخې~\LR{$\Gamma \Entails !A$} هله او يوازې هله چې
\LR{$\Gamma \Proves !A$}، نو~\LR{$\Gamma$} هله او يوازې هله بشپړه ده چې د
هغې د ژبې د هرې تړلے فارمول~\LR{$!A$} دپاره يا~\LR{$\Gamma \Proves !A$} يا
\LR{$\Gamma \Proves \lnot !A$} وي۔
\end{explain}
\olpseganchor{OLP-0277-B022}{end}

\olpseganchor{OLP-0277-B023}{start}
بله پوښتنه چې راپورته کېږي دا ده: آيا داسې محاسبوي کړنلاره شته چې
وآزمويو يوه تړلے فارمول په~\LR{$\Th{TA}$}، په~\LR{$\Th{PA}$}، يا ان يوازې
په~\LR{$\Th{Q}$} کښې ده که نۀ؟ دا پوښتنه د دې په تعريفولو لا کره کولے شو
چې يو سټ (مثلاً د تړلي فارمولونه سټ) کله د پرېکړې وړ وي۔
\olpseganchor{OLP-0277-B023}{end}

\olpseganchor{OLP-0277-B024}{start}
\begin{defn}
سټ~\LR{$X$} هله او يوازې هله \emph{د پرېکړې وړ} دے چې داسې محاسبوي
کړنلاره وي چې پر ننوت~\LR{$x$}، که~\LR{$x \in X$} وي~\LR{$1$} او که نۀ وي~\LR{$0$}
بېرته ورکوي۔
\end{defn}
\olpseganchor{OLP-0277-B024}{end}

\olpseganchor{OLP-0277-B025}{start}
نو زموږ پوښتنه داسې شوه: آيا~\LR{$\Th{TA}$} (\LR{$\Th{PA}$}، \LR{$\Th{Q}$})
د پرېکړې وړ ده؟
\olpseganchor{OLP-0277-B025}{end}

\olpseganchor{OLP-0277-B026}{start}
د دې ټولو پوښتنو ځواب به «نه» وي۔ له دې تيوريو څخه يوه هم د پرېکړې
وړ نۀ ده۔ خو دا پديده يوازې په همدې ځانګړو تيوريو پورې محدوده نۀ ده۔
په حقيقت کښې \emph{هره} هغه تيوري چې ځينې شرطونه پوره کوي له همدغو
پايلو سره مخ کېږي۔ له دغو شرطونو څخه يو، چې~\LR{$\Th{Q}$} او~\LR{$\Th{PA}$}
ئې پوره کوي، دا دے چې د بديهي اصولو د د پرېکړې وړ سټ په وسيله
بديهي شوې وي۔
\olpseganchor{OLP-0277-B026}{end}

\olpseganchor{OLP-0277-B027}{start}
\begin{defn}
يوه تيوري \emph{په محاسبوي ډول بديهي کېدونکې} ده که د بديهي اصولو د
د پرېکړې وړ سټ په وسيله بديهي شوې وي۔
\end{defn}
\olpseganchor{OLP-0277-B027}{end}

\olpseganchor{OLP-0277-B028}{start}
\begin{ex}
هره تيوري چې د تړلي فارمولونه د يوه متناهي سټ په وسيله بديهي شوې وي
په محاسبوي ډول بديهي کېدونکې ده، ځکه هر متناهي سټ د پرېکړې وړ دے۔ نو مثلاً
\LR{$\Th{Q}$}، په محاسبوي ډول بديهي کېدونکې ده۔
\olpseganchor{OLP-0277-B028}{end}

\olpseganchor{OLP-0277-B029}{start}
د~\LR{$\Th{PA}$} په شان په شېمايي ډول بديهي شوې تيورۍ هم
په محاسبوي ډول بديهي کېدونکې دي۔ د دې د ازمويلو دپاره چې~\LR{$!B$} د~\LR{$\Th{PA}$} په
بديهي اصولو کښې شامله ده که نۀ، يعنې د تابع~\LR{$\Char{X}$} د محاسبې
دپاره، چې که~\LR{$!B$} د~\LR{$\Th{PA}$} يو بديهي اصل وي
\LR{$\Char{X}(!B) = 1$} او که نۀ وي~\LR{$=0$}، داسې کولے شو: لومړے وګورئ
\LR{$!B$} د~\LR{$\Th{Q}$} له بديهي اصولو څخه ده که نۀ۔ که وي، ځواب «هو» دے
او~\LR{$\Char{X}(!B) = 1$}۔ که نۀ وي، وازمويئ چې د استقرا شېما بېلګه ده
که نۀ۔ دا په منظمه توګه کېدے شي؛ په دې حالت کښې ښايي تر ټولو اسانه
دا وي چې داسې ئې ووينو: د استقرا شېما هره بېلګه په يو شمېر کلي
کميت‌ټاکونکو پيلېږي، او بيا داسې فرعي-فارمول راځي چې يو شرطي
دے۔ د هغه شرطي پايله~\LR{$\lforall[x][!A(x, y_1, \dots, y_n)]$} ده، چې
\LR{$x$} او~\LR{$y_1$}، \dots،~\LR{$y_n$} د~\LR{$!A$} ټول آزاد متغيرونه دي او د~\LR{$!B$}
په پيل کښې کميت‌ټاکونکي متغيرونه~\LR{$y_1$}، \dots،~\LR{$y_n$} تړي۔ کله چې دا
\LR{$!A$} راوباسو او وګورو چې آزاد متغيرونه ئې د لومړنيو کلي
کميت‌ټاکونکو او~\LR{$\lforall[x]$} له تړلو متغيرونو سره برابر دي، بيا
ګورو چې د شرطي مقدمه له لاندې سره برابره ده که نۀ:
\[
!A(\Obj 0, y_1, \dots, y_n) \land \lforall[x][(!A(x, y_1, \dots, y_n)
\lif !A(x', y_1, \dots, y_n))]
\]
بيا هم، که برابره وي، \LR{$!B$} د استقرا شېما بېلګه ده؛ او که برابره نۀ
وي، \LR{$!B$} د استقرا شېما بېلګه نۀ ده۔
\end{ex}
\olpseganchor{OLP-0277-B029}{end}

\olpseganchor{OLP-0277-B030}{start}
د دې پوښتنې، او د دې لا عامې پوښتنې د ځوابولو دپاره چې کومې تيورۍ
بشپړې يا د پرېکړې وړ دي، لاندې تعريف هم ګټور دے۔ په ياد ولرئ چې سټ~\LR{$X$}
هله او يوازې هله د شمېر وړ دے چې تش وي يا داسې پر هدف سټ بشپړه
تابع~\LR{$f \colon \Nat \to X$} موجوده وي۔ داسې تابع د~\LR{$X$} شمېره بلل کېږي۔
\olpseganchor{OLP-0277-B030}{end}

\olpseganchor{OLP-0277-B031}{start}
\begin{defn}
سټ~\LR{$X$} هله او يوازې هله \emph{په محاسبوي ډول د شمېر وړ} (په لنډ ډول
c.e.) بلل کېږي چې تش وي يا محاسبه کېدونکې شمېره ولري۔
\end{defn}
\olpseganchor{OLP-0277-B031}{end}

\olpseganchor{OLP-0277-B032}{start}
د د محاسبوي بديهي کېدو وړتيا تر څنګ، پر هغو تيوريو بل شرط چې د ناتکميلۍ
قضيې پرې پلي کېږي دا دے چې د محاسبه کېدونکو تابعو او
د پرېکړې وړ اړيکو په باب بنسټيز حقايق ثابتولو ته کافي پياوړې وي۔
له «بنسټيزو حقايقو» څخه موخه هغه تړلي فارمولونه دي چې د هرې تابعې د
ارګومېنټونو دپاره ئې ارزښتونه بيانوي۔ او له «کافي پياوړې» څخه موخه
دا ده چې اړوندې تيورۍ دا جملې په خپلو قضيو کښې شمېري۔ د بېلګې په
توګه يوه نمونه‌يي محاسبه کېدونکې تابع، يعنې جمع، وګورئ۔ د~\LR{$2$} او~\LR{$3$}
ارګومېنټونو دپاره د~\LR{$+$} ارزښت~\LR{$5$} دے، يعنې~\LR{$2+3 = 5$}۔ د حساب په
ژبه کښې يوه جمله چې وايي د~\LR{$2$} او~\LR{$3$} ارګومېنټونو دپاره د~\LR{$+$} ارزښت
\LR{$5$} دے، دا ده: \LR{$(\num{2} + \num{3}) = \num{5}$}۔ او مثلاً~\LR{$\Th{Q}$}
دا جمله ثابتوي۔ په عام ډول غواړو د هرې محاسبه کېدونکې تابع
\LR{$f(x_1, x_2)$} دپاره په~\LR{$\Lang{L_A}$} کښې داسې فارمول
\LR{$!A_f(x_1, x_2, y)$} وي چې هر کله~\LR{$f(n_1, n_2) = m$}،
\LR{$\Th{Q} \Proves !A_f(\num{n_1}, \num{n_2}, \num{m})$} وي۔ په دې ډول
\LR{$\Th{Q}$} ثابتوي چې د~\LR{$n_1$}، \LR{$n_2$} ارګومېنټونو دپاره د~\LR{$f$} ارزښت~\LR{$m$}
دے۔ په حقيقت کښې لږ څه زيات غواړو، يعنې دا هم ثابته کړي چې د~\LR{$n_1$}،
\LR{$n_2$} ارګومېنټونو دپاره بل هېڅ عدد د~\LR{$f$} ارزښت نۀ دے۔ د
د پرېکړې وړ اړيکو دپاره هم همداسې ده۔ لاندې دوه تعريفونه دا خبره
کره کوي۔
\olpseganchor{OLP-0277-B032}{end}

\olpseganchor{OLP-0277-B033}{start}
\begin{defn}
فارمول~\LR{$!A(x_1, \dots, x_k, y)$} په~\LR{$\Gamma$} کښې تابع
\LR{$f\colon \Nat^k \to \Nat$} هله او يوازې هله \emph{تمثيل}
کوي چې هر کله~\LR{$f(n_1, \dots, n_k) = m$} وي، نو
\begin{enumerate}
\item \LR{$\Gamma \Proves !A(\num{n_1}, \dots, \num{n_k}, \num{m})$}، او
\item \LR{$\Gamma \Proves \lforall[y](!A(\num{n_1}, \dots, \num{n_k},
y) \lif y = \num{m})$}۔
\end{enumerate}
\end{defn}
\olpseganchor{OLP-0277-B033}{end}

\olpseganchor{OLP-0277-B034}{start}
\begin{defn}
فارمول~\LR{$!A(x_1, \dots, x_k)$} اړيکه~\LR{$R \subseteq \Nat^k$} هله او
يوازې هله \emph{تمثيل} کوي چې
\begin{enumerate}
\item هر کله~\LR{$R(n_1, \dots, n_k)$} وي،
\LR{$\Gamma \Proves !A(\num{n_1}, \dots, \num{n_k})$}، او
\item هر کله~\LR{$R(n_1, \dots, n_k)$} نۀ وي،
\LR{$\Gamma \Proves \lnot !A(\num{n_1}, \dots, \num{n_k})$}۔
\end{enumerate}
\end{defn}
\olpseganchor{OLP-0277-B034}{end}

\olpseganchor{OLP-0277-B035}{start}
يوه تيوري د ناتکميلۍ د قضيو د پلي کېدو دپاره «کافي پياوړې» ده که
ټولې محاسبه کېدونکې تابعې او ټولې د پرېکړې وړ اړيکې تمثيل
کوي۔ \LR{$\Th{Q}$} او د هغې غځونې دا شرط پوره کوي، خو د دې ثابتول به لږ
وخت واخلي؛ دا د هغو خبرو په باب يو غيرابتذالي حقيقت دے چې~\LR{$\Th{Q}$}
ئې ثابتولے شي، او ښودل ئې ځکه ستونزمن دي چې~\LR{$\Th{Q}$} يوازې څو بديهي
اصول لري چې بايد دا ټول حقايق ترې ثابت کړو۔ خو~\LR{$\Th{Q}$} ډېره کمزورې
تيوري ده۔ نو که څه هم دا ثابتول ستونزمن دي چې~\LR{$\Th{Q}$} ټولې محاسبه
کېدونکې تابعې تمثيلوي، زياتره په زړه پورې تيورۍ تر~\LR{$\Th{Q}$} پياوړې
دي؛ يعنې تر~\LR{$\Th{Q}$} زيات څه ثابتوي۔ او هر څه چې~\LR{$\Th{Q}$} ثابتوي،
هره پياوړې تيوري ئې هم ثابتوي؛ څرنګه چې~\LR{$\Th{Q}$} ټولې محاسبه کېدونکې
تابعې تمثيلوي، هره پياوړې تيوري ئې هم تمثيلوي۔ په دې معنا ډېرې په
زړه پورې تيورۍ د ناتکميلۍ د قضيو دا شرط پوره کوي۔ نو زموږ سخت کار به
ګټه وکړي، ځکه ښيي چې د ناتکميلۍ قضيې پر ډېرو تيوريو پلي کېږي۔ په
يقيني ډول هره هغه تيوري چې موخه ئې د «ټولو رياضياتو» صوري کول وي،
بايد هر څه ثابت کړي چې~\LR{$\Th{Q}$} ئې ثابتوي، ځکه لږ تر لږه بايد د
لومړنيو محاسبو پايلې رانغاړلے شي۔ نو هره هغه تيوري چې د «ټولو
رياضياتو» د تيورۍ نومانده وي، د ناتکميلۍ قضيې به پرې پلي کېږي۔
\olpseganchor{OLP-0277-B035}{end}

\olpseganchor{OLP-0278-B004}{start}
\olfileid{inc}{int}{ovr}
\olpseganchor{OLP-0278-B004}{end}

\olpseganchor{OLP-0278-B005}{start}
\olsection{د ناتکميلۍ د پايلو ټوليزه کتنه}
\olpseganchor{OLP-0278-B005}{end}

\olpseganchor{OLP-0278-B006}{start}
هېلبرټ تمه لرله چې رياضيات په داسې په محاسبوي ډول بديهي کېدونکې تيورۍ کښې
صوري کېدے شي چې بشپړ او د پرېکړې وړ ثابتېدے شي۔ سربېره پر
دې، موخه ئې دا وه چې د دې تيورۍ سازګاري په ډېرو کمزورو،
«متناهي‌پالو» لارو ثابته کړي، څو د شهوديت د ننګونو پر وړاندې له
کلاسیکو رياضياتو دفاع وکړي۔ د ګوډل د ناتکميلۍ قضيو وښودل چې دا موخې
نۀ شي تر لاسه کېدے۔
\olpseganchor{OLP-0278-B006}{end}

\olpseganchor{OLP-0278-B007}{start}
د ګوډل د ناتکميلۍ لومړۍ قضيې وښودل چې د رسل او وايټ‌هېډ د
\emph{پرېنسيپيا مېتېماتيکا} يوه بڼه بشپړ نۀ ده۔ خو ثبوت په
اصل کښې ډېر عام و او پر ډېرو بېلابېلو تيوريو پلي کېږي۔ معنا ئې دا ده
چې ستونزه يوازې دا نۀ وه چې \emph{پرېنسيپيا مېتېماتيکا} رياضيات په
بشپړ ډول رانۀ غښتل، بلکې \emph{هېڅ} منل کېدونکې تيوري ئې نۀ شي
رانغاړلے۔ څه وخت ولګېد چې د تيوريو هغه ځانګړتياوې بېلې شي چې د
ناتکميلۍ د قضيو د پلي کېدو دپاره بس دي، او د ګوډل ثبوت داسې عام شي
چې يوازې په همدې ځانګړتياوو تکيه وکړي۔ خو اوس کولے شو د حساب د
ژبې~\LR{$\Lang L_A$} د تيوريو دپاره د ناتکميلۍ د لومړۍ قضيې يوه ډېره عامه
بڼه بيان کړو۔
\olpseganchor{OLP-0278-B007}{end}

\olpseganchor{OLP-0278-B008}{start}
\begin{thm}
که~\LR{$\Gamma$} په~\LR{$\Lang L_A$} کښې داسې سازګاره او په محاسبوي ډول بديهي کېدونکې
تيوري وي چې ټولې محاسبه کېدونکې تابعې او ټولې د پرېکړې وړ اړيکې
تمثيل کوي، نو~\LR{$\Gamma$} بشپړ نۀ ده۔
\end{thm}
\olpseganchor{OLP-0278-B008}{end}

\olpseganchor{OLP-0278-B009}{start}
دا ويل چې~\LR{$\Gamma$} بشپړ نۀ ده، په دې معنا دي چې لږ تر لږه
د يوې تړلے فارمول~\LR{$!A$} دپاره~\LR{$\Gamma \Proves/ !A$} او
\LR{$\Gamma \Proves/ \lnot !A$} وي۔ داسې تړلے فارمول (له~\LR{$\Gamma$}
څخه) \emph{خپلواکه} بلل کېږي۔
\textbf{د ژباړې د نسخې سمون (OLCMP-070):} سرچينې د ګاما له نوم
سره د تړون قوس په ناسمه توګه د رياضياتي چاپېريال دننه ليکلے و؛ دلته
قوس د جملې په خپل ځای کښې راوړل شوے او د تيورۍ نوم نۀ دے بدل شوے۔
په حقيقت کښې په نسبتاً چټکۍ ثابتولے شو
چې خپلواکې جملې بايد موجودې وي۔ خو د قضيې د ګوډلي ثبوت ځواک په دې کښې
دے چې د داسې خپلواکې تړلے فارمول يوه \emph{ټاکلې بېلګه} وړاندې کوي۔
دا په زړه پورې جوړونه يوه تړلے فارمول~\LR{$!G_\Gamma$} پيدا کوي، چې د
\LR{$\Gamma$} \emph{ګوډل جمله} بلل کېږي او ځکه ناثابتېدونکې ده چې په
\LR{$\Gamma$} کښې~\LR{$!G_\Gamma$} له دې ادعا سره برابره ده چې~\LR{$!G_\Gamma$} په
\LR{$\Gamma$} کښې ناثابتېدونکې ده۔ جوړونه دا کار \emph{تعميري} کوي؛ يعنې
که د~\LR{$\Gamma$} بديهي کول او د اشتقاق نظام بيان راکړل شي، ثبوت
په رښتيا د~\LR{$!G_\Gamma$} د ليکلو يوه لار راکوي۔
\olpseganchor{OLP-0278-B009}{end}

\olpseganchor{OLP-0278-B010}{start}
د ګوډل په ثبوت کښې جوړونه غواړي چې د~\LR{$\Lang L_A$} د اصطلاحاتو او
فارمولونه ځانګړتياوې او عملونه پخپله په~\LR{$\Lang L_A$} کښې د بيانولو
لار پيدا کړو۔ په دې کښې داسې ځانګړتياوې شاملې دي لکه «\LR{$!A$}
تړلے فارمول ده»، «\LR{$\delta$} د~\LR{$!A$} اشتقاق ده»، او داسې
عملونه لکه~\LR{$\Subst{!A}{t}{x}$}۔ دا لار بايد (الف) دا ځانګړتياوې او
اړيکې د نښو او د هغو د لړيو د «کوډولو» په وسيله د طبيعي عددونو په
توګه بيان کړي (ځکه اصطلاحات، فارمولونه، اشتقاقونه او نور د
نښو لړۍ دي، او~\LR{$\Lang L_A$} د طبيعي عددونو په باب خبرې کوي)۔ بايد
(ب) دا کار داسې وکړي چې~\LR{$\Gamma$} اړوند حقايق ثابت کړي؛ نو بايد وښيو
چې دا ځانګړتياوې د طبيعي عددونو په د پرېکړې وړ خاصيتونو کوډېږي او
عملونه د طبيعي عددونو له محاسبه کېدونکو تابعو سره سمون لري۔ دې ته
«د نحو حسابي کول» ويل کېږي۔
\olpseganchor{OLP-0278-B010}{end}

\olpseganchor{OLP-0278-B011}{start}
خو د نحو د حسابي کولو تر څېړلو وړاندې به دا شرط وګورو چې~\LR{$\Gamma$}
«کافي پياوړې» وي؛ يعنې ټولې محاسبه کېدونکې تابعې او ټولې
د پرېکړې وړ اړيکې تمثيل کړي۔ دا غواړي چې «محاسبه کېدونکي»
ته يو کره تعريف ورکړو۔ دا په بېلابېلو لارو کېدے شي؛ مثلاً د ټيورينګ
ماشينونو د مدل په وسيله، يا د هغو تابعو په توګه چې په کومې عمومي
موخې پروګرام‌ليک ژبه کښې پروګرامونه ئې محاسبه کوي۔ خو څرنګه چې زموږ
موخه د~\LR{$\Lang L_A$} د ژبې په يوه تيورۍ کښې د دې تابعو او اړيکو
تمثيل کول دي، غوره ده چې يوازې د عددي تابعو د محاسبه کېدنې يو ساده تعريف
وټاکو۔ دا د \emph{بازګشتي تابعې} مفهوم دے۔ نو لومړے به د بازګشتي
تابعو بحث وکړو۔ بيا به وښيو چې~\LR{$\Th{Q}$} لا له وړاندې ټولې بازګشتي
تابعې او اړيکې تمثيل کوي۔ دا به موږ ته اجازه راکړي چې د
ناتکميلۍ قضيه پر~\LR{$\Th{Q}$} او~\LR{$\Th{PA}$} غوندې ټاکلو تيوريو پلې کړو،
ځکه ثابته کړې به مو وي چې دا د «کافي پياوړو» تيوريو بېلګې دي۔
\textbf{د ژباړې د نسخې سمون (OLCMP-071):} سرچينې د دې فقرې په
موخه‌بيانوونکې جمله کښې د اصطلاحي نښې په پای کښې زائده انګرېزي~s
لرله؛ دلته هماغه ټاکل شوې اصطلاح له زائدې نښې پرته راوړل شوې ده۔
\olpseganchor{OLP-0278-B011}{end}

\olpseganchor{OLP-0278-B012}{start}
د نحو د حسابي کولو وروستۍ پايله داسې فارمول~\LR{$\OProv[\Gamma](x)$}
ده چې د فارمولونه د عددونو په توګه د کوډولو په وسيله د~\LR{$\Gamma$}
له بديهي اصولو څخه د اثبات وړتيا بيانوي۔ په ځانګړي ډول، که~\LR{$!A$} په
عدد~\LR{$n$} کوډ شوې وي او~\LR{$\Gamma \Proves !A$}، نو
\LR{$\Gamma \Proves \Prov[\Gamma](\num{n})$}۔ د~\LR{$\Gamma$} دپاره دا
«اثبات‌وړتيا پريديکات» موږ ته دا وس هم راکوي چې په يوه معنا د~\LR{$\Gamma$}
سازګاري د~\LR{$\Lang L_A$} د يوې تړلے فارمول په توګه بيان کړو: د
\LR{$\Gamma$} دپاره د «سازګارۍ جمله» په نوم هغه تړلے فارمول
\LR{$\lnot \Prov[\Gamma](\num{n})$} ونيسئ، چې~\LR{$n$} پکې د يوه تناقض،
مثلاً~\LR{$\lfalse$}، کوډ وي۔ د ناتکميلۍ
دويمه قضيه وايي چې سازګارې په محاسبوي ډول بديهي کېدونکې تيورۍ خپلې د سازګارۍ
جملې هم نۀ ثابتوي۔ د دې قضيې د پلي کېدو شرطونه تر دې لږ سخت دي چې
تيوري دې يوازې ټولې محاسبه کېدونکې تابعې او د پرېکړې وړ اړيکې
تمثيل کړي، خو وبه ښيو چې~\LR{$\Th{PA}$} دا شرطونه پوره کوي۔
\olpseganchor{OLP-0278-B012}{end}

\olpseganchor{OLP-0279-B004}{start}
\olfileid{inc}{int}{dec}
\olpseganchor{OLP-0279-B004}{end}

\olpseganchor{OLP-0279-B005}{start}
\olsection{ناپرېکړتيا او ناتکميلي}
\olpseganchor{OLP-0279-B005}{end}

\olpseganchor{OLP-0279-B006}{start}
د ناتکميلۍ د قضيو د ګوډلي ثبوت د نحو حسابي کول غواړي۔ خو ان له دې
پرته، يوازې په دې فرض چې يوه تيوري ټولې د پرېکړې وړ اړيکې
تمثيل کوي، ځينې ښې پايلې تر لاسه کولے شو۔ ثبوت يو قطري
استدلال دے چې د درېدنې د مسئلې د ناپرېکړتيا ثبوت ته ورته دے۔
\olpseganchor{OLP-0279-B006}{end}

\olpseganchor{OLP-0279-B007}{start}
\begin{thm}
که~\LR{$\Gamma$} يوه داسې سازګاره تيوري وي چې هره د پرېکړې وړ اړيکه
تمثيل کوي، نو~\LR{$\Gamma$} د پرېکړې وړ نۀ ده۔
\end{thm}
\olpseganchor{OLP-0279-B007}{end}

\olpseganchor{OLP-0279-B008}{start}
\begin{proof}
فرض کړئ~\LR{$\Gamma$} د پرېکړې وړ وي۔ ښيو چې که~\LR{$\Gamma$} هره
د پرېکړې وړ اړيکه تمثيل کړي، نو بايد ناسازګاره وي۔
\olpseganchor{OLP-0279-B008}{end}

\olpseganchor{OLP-0279-B009}{start}
د پرېکړې وړ خاصيتونه (يوځاييزې اړيکې) په هغو فارمولونه تمثيلېږي
چې يو آزاد متغير لري۔ فرض کړئ~\LR{$!A_0(x)$}، \LR{$!A_1(x)$}، \dots، د ټولو
داسې فارمولونه يوه محاسبه کېدونکې شمېره وي۔ اوس لاندې سټ
\LR{$D \subseteq \Nat$} وګورئ:
\[
D = \Setabs{n}{\Gamma \Proves \lnot !A_n(\num{n})}
\]
سټ~\LR{$D$} د پرېکړې وړ دے، ځکه دا ازمويلے شو چې~\LR{$n \in D$} دے که نۀ:
لومړے~\LR{$!A_n(x)$}، او له هغې څخه~\LR{$\lnot !A_n(\num{n})$} محاسبه کوو۔ په
څرګند ډول په~\LR{$!A_n(x)$} کښې د~\LR{$x$} د هر آزاد پېښېدنې پر ځاے د
اصطلاح~\LR{$\num{n}$} ايښوول
او د~\LR{$!A_n(\num{n})$} مخې ته~\LR{$\lnot$} ليکل يو ميخانيکي کار دے۔ د فرض
له مخې~\LR{$\Gamma$} د پرېکړې وړ ده، نو دا ازمويلے شو چې
\LR{$\lnot !A_n(\num{n}) \in \Gamma$} ده که نۀ۔ که وي، \LR{$n \in D$}؛ او که
نۀ وي، \LR{$n \notin D$}۔ نو~\LR{$D$} هم د پرېکړې وړ دے۔
\textbf{د ژباړې د نسخې سمون (OLCMP-072):} سرچينې د قطري جوړونې په
دوو پرله‌پسې ګامونو کښې د فورمول شاخص غورځولے و؛ دلته په دواړو
ځايونو کښې شاخص بېرته راوړل شوے، څو ازموينه د قطري سټ له تعريف سره
هماغه فورمول وکاروي۔
\olpseganchor{OLP-0279-B009}{end}

\olpseganchor{OLP-0279-B010}{start}
څرنګه چې~\LR{$\Gamma$} ټول د پرېکړې وړ خاصيتونه تمثيل کوي، نو~\LR{$D$} هم
تمثيل کوي۔ او هغه فارمولونه چې~\LR{$D$} په~\LR{$\Gamma$} کښې
تمثيل کوي ټول د~\LR{$!A_0(x)$}، \LR{$!A_1(x)$}، \dots، په لړۍ کښې
دي۔ نو داسې عدد~\LR{$d$} ونيسئ چې~\LR{$!A_d(x)$} په~\LR{$\Gamma$} کښې~\LR{$D$}
تمثيل کوي۔ که~\LR{$d \notin D$} وي، نو څرنګه چې~\LR{$!A_d(x)$}
تمثيلوي~\LR{$D$}، \LR{$\Gamma \Proves \lnot !A_d(\num{d})$}۔ خو دا په دې معنا ده
چې~\LR{$d$} د~\LR{$D$} ټاکونکی شرط پوره کوي، او له دې امله~\LR{$d \in D$}۔ دا له
\LR{$d \notin D$} سره تناقض لري۔ نو د غيرمستقيم ثبوت له مخې~\LR{$d \in D$}۔
\textbf{د ژباړې د نسخې سمون (OLCMP-073):} سرچينې په دې واحد کښې د
تمثيل د اصطلاحي نښې په پای کښې زائده انګرېزي~s څلور ځله لرله؛ دلته
ټولې پېښې له زائدې نښې پرته د ټاکل شوې اصطلاح په توګه راوړل شوې دي۔
\olpseganchor{OLP-0279-B010}{end}

\olpseganchor{OLP-0279-B011}{start}
څرنګه چې~\LR{$d \in D$}، د~\LR{$D$} د تعريف له مخې
\LR{$\Gamma \Proves \lnot !A_d(\num{d})$}۔ له بلې خوا، څرنګه چې~\LR{$!A_d(x)$}
په~\LR{$\Gamma$} کښې~\LR{$D$} تمثيل کوي،
\LR{$\Gamma \Proves !A_d(\num{d})$}۔ نو~\LR{$\Gamma$} ناسازګاره ده۔
\end{proof}
\olpseganchor{OLP-0279-B011}{end}

\olpseganchor{OLP-0279-B012}{start}
\begin{explain}
پاسنۍ قضيه ښيي چې هېڅ سازګاره تيوري چې ټولې د پرېکړې وړ اړيکې
تمثيل کوي، د پرېکړې وړ کېدے نۀ شي۔ موږ به وښيو چې~\LR{$\Th{Q}$}
ټولې د پرېکړې وړ اړيکې تمثيل کوي؛ معنا ئې دا ده چې ټولې
هغه تيورۍ چې~\LR{$\Th{Q}$} رانغاړي، لکه~\LR{$\Th{PA}$} او~\LR{$\Th{TA}$}، هم دا
کار کوي او له دې امله هغوى هم د پرېکړې وړ نۀ دي۔ (څرنګه چې دا ټولې
تيورۍ په معياري مدل کښې رښتونې دي، ټولې سازګارې دي۔)
\olpseganchor{OLP-0279-B012}{end}

\olpseganchor{OLP-0279-B013}{start}
دا پايله د ناتکميلۍ د لومړۍ قضيې د يوې کمزورې بڼې د ترلاسه کولو
دپاره هم کارولے شو۔ هره هغه تيوري چې په محاسبوي ډول بديهي کېدونکې او
بشپړ وي، د پرېکړې وړ ده۔ نو سازګارې تيورۍ چې
په محاسبوي ډول بديهي کېدونکې وي او ټول د پرېکړې وړ خاصيتونه تمثيل
کوي، بشپړ کېدے نۀ شي۔
\end{explain}
\olpseganchor{OLP-0279-B013}{end}

\olpseganchor{OLP-0279-B014}{start}
\begin{thm}
که~\LR{$\Gamma$} په محاسبوي ډول بديهي کېدونکې او بشپړ وي، نو د پرېکړې وړ
ده۔
\end{thm}
\olpseganchor{OLP-0279-B014}{end}

\olpseganchor{OLP-0279-B015}{start}
\begin{proof}
هره ناسازګاره تيوري د پرېکړې وړ ده، ځکه ناسازګارې تيورۍ ټول
تړلي فارمولونه رانغاړي؛ نو د «آيا~\LR{$!A \in \Gamma$}؟» پوښتنې ځواب تل
«هو» وي، يعنې پرېکړه ئې کېدے شي۔
\olpseganchor{OLP-0279-B015}{end}

\olpseganchor{OLP-0279-B016}{start}
نو فرض کړئ~\LR{$\Gamma$} سازګاره، او ورسره په محاسبوي ډول بديهي کېدونکې او
بشپړ وي۔ څرنګه چې~\LR{$\Gamma$} په محاسبوي ډول بديهي کېدونکې ده، نو
په محاسبوي ډول د شمېر وړ ده۔ ځکه له~\LR{$\Gamma$} د بديهي اصولو څخه ټول
سم اشتقاقونه په يوه محاسبه کېدونکې تابع شمېرلے شو۔ له يوه سم
اشتقاق څخه هغه تړلے فارمول محاسبه کولے شو چې دا ئې
اشتقاق؛ نو په ګډه داسې محاسبه کېدونکې تابع شته چې د~\LR{$\Gamma$} ټولې
قضيې شمېري۔ څرنګه چې~\LR{$\Gamma$} سازګاره او بشپړ ده،
تړلے فارمول هله او يوازې هله د~\LR{$\Gamma$} قضيه ده چې~\LR{$\lnot !A$}
قضيه نۀ وي۔ نو داسې پرېکړه کولے شو چې~\LR{$!A \in \Gamma$} ده که نۀ۔ د
\LR{$\Gamma$} ټولې قضيې وشمېرئ۔ کله چې~\LR{$!A$} په لست کښې راشي، پوهېږو چې
\LR{$\Gamma \Proves !A$}۔ کله چې~\LR{$\lnot !A$} په لست کښې راشي، پوهېږو چې
\LR{$\Gamma \Proves/ !A$}۔ څرنګه چې~\LR{$\Gamma$} بشپړ ده، له دې دوو
حالتونو يو بالاخره رامنځته کېږي؛ نو کړنلاره بالاخره ځواب ورکوي۔
\end{proof}
\olpseganchor{OLP-0279-B016}{end}

\olpseganchor{OLP-0279-B017}{start}
\begin{cor}
\ollabel{cor:incompleteness}
که~\LR{$\Gamma$} سازګاره او په محاسبوي ډول بديهي کېدونکې وي او هر د پرېکړې وړ
خاصيت تمثيل کړي، نو بشپړ نۀ ده۔
\end{cor}
\olpseganchor{OLP-0279-B017}{end}

\olpseganchor{OLP-0279-B018}{start}
\begin{proof}
که~\LR{$\Gamma$} بشپړ وے، د تېرې قضيې له مخې به د پرېکړې وړ وه
(ځکه په محاسبوي ډول بديهي کېدونکې او سازګاره ده)۔ خو څرنګه چې~\LR{$\Gamma$} هر د پرېکړې وړ
خاصيت تمثيل کوي، د لومړۍ قضيې له مخې د پرېکړې وړ نۀ ده۔
\end{proof}
\olpseganchor{OLP-0279-B018}{end}

\olpseganchor{OLP-0279-B019}{start}
\begin{prob}
وښيئ چې~\LR{$\Th{TA} = \Setabs{!A}{\Sat{N}{!A}}$}، په محاسبوي ډول بديهي کېدونکې نۀ
ده۔ دا فرضولے شئ چې~\LR{$\Th{TA}$} ټول د پرېکړې وړ خاصيتونه تمثيلوي۔
\end{prob}
\olpseganchor{OLP-0279-B019}{end}

\olpseganchor{OLP-0279-B020}{start}
کله چې ثابته کړو، مثلاً، \LR{$\Th{Q}$} ټول د پرېکړې وړ خاصيتونه
تمثيل کوي، ملازمه راته وايي چې~\LR{$\Th{Q}$} بايد ناتکميله وي۔
خو ثبوت ئې د خپلواکې تړلے فارمول بېلګه نۀ راکوي؛ يوازې ښيي چې داسې
تړلے فارمول بايد موجوده وي۔ د بېلګې د ترلاسه کولو دپاره بايد نحو
حسابي کړو او د ګوډل د اصلي ثبوت نظر تعقيب کړو۔ او البته لا هم بايد
لومړۍ ادعا ثابته کړو، يعنې دا چې~\LR{$\Th{Q}$} په رښتيا ټول د پرېکړې وړ
خاصيتونه تمثيل کوي۔
\olpseganchor{OLP-0279-B020}{end}

\olpseganchor{OLP-0279-B021}{start}
بايد پام وشي چې هره \emph{په زړه پورې} تيوري ناتکميله يا ناپرېکړه
کېدونکې نۀ ده۔ ډېرې داسې تيورۍ شته چې د په زړه پورې رياضياتي حقايقو
بيانولو ته کافي پياوړې دي، خو د ګوډل د پايلې شرطونه نۀ پوره کوي۔ د
بېلګې په توګه،
\LR{$\Th{Pres} = \Setabs{!A \in \Lang{L_{A^+}}}{\Sat{N}{!A}}$}، يعنې د
ضرب~\LR{$\times$} پرته د حساب د ژبې د هغو تړلي فارمولونه سټ چې په معياري
مدل کښې رښتونې دي، هم بشپړه او هم د پرېکړې وړ ده۔ دې تيورۍ ته د
پرېسبورګر حساب ويل کېږي، او د طبيعي عددونو په باب ټول هغه حقايق
ثابتوي چې يوازې په~\LR{$\Obj{0}$}، \LR{$\prime$} او~\LR{$+$} جوړېدے شي۔
\olpseganchor{OLP-0279-B021}{end}

\olpseganchor{OLP-0280-B004}{start}
\olchapter{inc}{art}{د نحو حسابي کول}
\olpseganchor{OLP-0280-B004}{end}

\olpseganchor{OLP-0280-B005}{start}
\begin{editorial}
  يادونه وکړئ چې د نښه‌لرونکو تابلوګانو دپاره حسابي کول لا نۀ دي
  چمتو شوي۔
\end{editorial}
\olpseganchor{OLP-0280-B005}{end}











\olpseganchor{OLP-0280-B011}{start}
%
  

\olpseganchor{OLP-0280-B011}{end}

\olpseganchor{OLP-0280-B012}{start}
%
  

\olpseganchor{OLP-0280-B012}{end}

\olpseganchor{OLP-0281-B004}{start}
\olfileid{inc}{art}{int}
\olsection{پېژندګلو}
\olpseganchor{OLP-0281-B004}{end}

\olpseganchor{OLP-0281-B005}{start}
د محاسبه‌پوهې او منطق د نښلولو دپاره داسې لار پکار ده چې د منطق د
څيزونو (سمبولونو، ترمونو، فارمولونه او اشتقاقونه)، پر هغو د
عملياتو، او د هغو د خاصيتونو او اړيکو په اړه داسې خبرې پرې وکړو چې
محاسبوي چلند ته برابرې وي۔ دا کار نېغ په نېغه کولے شو: پر سمبولونو،
د سمبولونو پر لړيو او له هغو څخه پر جوړو شويو نورو څيزونو محاسبه
کېدونکې تابعې او اړيکې څېړو۔ ځکه چې د منطقي نحو ټول څيزونه متناهي دي
او د سمبولونو له د شمېر وړ سټ څخه جوړ شوي، نو د محاسبې په ځينو
مدلونو کښې دا شوني دي۔ خو د محاسبې نور مدلونه---لکه بازګشتي تابعې---
يوازې عددونو او د هغو اړيکو او تابعو ته محدود دي۔ برسېره پر دې، په
پاې کښې غواړو چې نحو د ځينو تيوريو دننه هم وڅېړو، په ځانګړي ډول په
هغو تيوريو کښې چې د حساب په ژبه کښې بيان شوې وي۔ په دې حالتونو کښې
د نحو \emph{حسابي کول} اړين دي؛ يعنې نحوي څيزونه، پر هغو عمليات او د
هغو اړيکې په ترتيب د عددونو، حسابي تابعو او حسابي اړيکو په توګه
تمثيل کړو۔ دا مفکوره چې لايبنېڅ ته رسېږي، نحوي څيزونو ته د عددونو
ټاکل دي۔
\olpseganchor{OLP-0281-B005}{end}

\olpseganchor{OLP-0281-B006}{start}
سمبولونو ته د هغو د «کوډونو» په توګه د عددونو ټاکل نسبتاً ساده دي۔
ځينې سمبولونه لږه ستونزه پيدا کوي، ځکه د بېلګې په توګه بې‌شمېره
متغيرونه، او ان د هرې ځاييزې کچې~\LR{$n$} بې‌شمېره تابعې شته۔ خو
البته سمبولونو ته په منظم ډول داسې عددونه ټاکل کېدے شي چې، د بېلګې
په توګه، \LR{$\Obj v_2$} او \LR{$\Obj v_3$} بېل کوډونه ولري۔ د سمبولونو لړۍ
(لکه ترمونه او فارمولونه) تر دې لويه ستونزه ده۔ خو که د عددونو له
لړيو سره سوچه حسابي چلند کولے شو (د بېلګې په توګه د لړيو د اولي
عددونو-د-قوتونو کوډول)، نو د جلا سمبولونو کوډول د سمبولونو د لړيو تر
کوډولو غځولے شو، او بيا ئې د فارمولونه تر لړيو يا نورو ترتيبونو،
لکه اشتقاقونه، نور هم غځولے شو۔ دا غځېدلے کوډول «ګوډل شمېرنه»
بلل کېږي۔ هر ترم، فارمول او اشتقاق ته يوه ګوډل شمېره
ټاکل کېږي۔
\olpseganchor{OLP-0281-B006}{end}

\olpseganchor{OLP-0281-B007}{start}
کله چې د سمبولونو لړۍ د هغو د کوډونو د لړۍ په توګه کوډوو، او د لړيو
د کوډولو داسې نظام غوره کوو چې په محاسبه کېدونکو تابعو پرې چلند کېدے
شي، نو پر ګوډل شمېرو هم په محاسبه کېدونکو تابعو چلند کولے شو۔ په عمل
کښې به ټولې اړوندې تابعې بنسټيزې بازګشتي وي۔ د بېلګې په توګه، د يوې
لړۍ اوږدوالے محاسبه کول او د لړۍ له کوډ څخه د هغې \LR{$i$}-م غړے محاسبه
کول دواړه بنسټيز بازګشتي دي۔ که د لړۍ کوډوونکے عدد، د بېلګې په توګه،
د فارمول~\LR{$!A$} ګوډل شمېره وي، سمدستي وينو چې د فارمول
اوږدوالے او په فارمول کښې د \LR{$i$}-م سمبول (کوډ) هم د~\LR{$!A$} له
ګوډل شمېرې څخه محاسبه کېدے شي۔ د دې ثابتول لږ ستونزمن دي چې، د بېلګې
په توګه، د سم جوړ شوي ترم يا د سم اشتقاق د ګوډل شمېرې په
توګه د يو عدد خاصيت بنسټيز بازګشتي دے۔ بيا هم دا شوني دي، ځکه د پام
وړ لړۍ (ترمونه، فارمولونه او اشتقاقونه) په استقرايي ډول
تعريف شوې دي۔
\olpseganchor{OLP-0281-B007}{end}

\olpseganchor{OLP-0281-B008}{start}
د بېلګې په توګه د تعويض عمليه وګورئ۔ که \LR{$!A$} يو فارمول، \LR{$x$} يو
متغير او \LR{$t$} يو ترم وي، نو \LR{$\Subst{!A}{t}{x}$} په~\LR{$!A$} کښې د~\LR{$x$} د
هرې ازادې پېښې په~\LR{$t$} بدلولو پايله ده۔ اوس فرض کړئ چې \LR{$!A$}، \LR{$x$} او
\LR{$t$} ته مو په ترتيب \LR{$k$}، \LR{$l$} او \LR{$m$} ګوډل شمېرې ټاکلې دي۔ هماغه
طرحه \LR{$\Subst{!A}{t}{x}$} ته، مثلاً، ګوډل شمېره~\LR{$n$} ټاکي۔ دا تطبيق---
چې \LR{$k$}، \LR{$l$} او \LR{$m$} پر~\LR{$n$} تطبيقوي---د تعويض د عمليې حسابي سيال دے۔
کله چې د تعويض عمليه \LR{$!A$}، \LR{$x$} او \LR{$t$} پر
\LR{$\Subst{!A}{t}{x}$} تطبيقوي، د حسابي شوي تعويض تابع د ګوډل شمېرو
\LR{$k$}، \LR{$l$} او \LR{$m$} پر ګوډل شمېرې~\LR{$n$} تطبيقوي۔ وبه وينو چې دا تابع
بنسټيزه بازګشتي ده۔
\olpseganchor{OLP-0281-B008}{end}

\olpseganchor{OLP-0281-B009}{start}
د نحو حسابي کول يوازې انتزاعي ارزښت نۀ لري، که څه هم په لومړي ځل دا
يوه نااشنا ژوره مفکوره وه چې د حساب د ژبې په شان ژبې، چې د ژبو «په
اړه د خبرو» کومه اله ورسره نۀ وي، بيا هم د عباراتو پېچلي خاصيتونه
صوري کولے شي۔ له دې وروسته يوازې يو وړوکے ګام پاتې کېږي چې وپوښتو په
دې ژبه کښې يوه تيوري، لکه د پېانو حساب، د خپلې ژبې په اړه څه
\emph{ثابتولے} شي (د بېلګې په توګه دا چې تړلي فارمولونه د ثبوت وړ يا
رښتيا دي که نۀ)۔ دا موږ د ګوډل (د ناثبوت‌وړتيا په اړه) او تارسکي (د
رښتيا د ناتعريفېدنې په اړه) نومياليو محدودوونکو قضيو ته رسوي۔ خو د
نحو د حسابي کولو چل د محاسبه‌پوهې د ځينو مهمو پايلو د ثبوت دپاره هم
مهم دے، د بېلګې په توګه د تيوريو د محاسبوي ځواک يا د محاسبې د بېلو
مدلونو ترمنځ د اړيکې په اړه۔ د نحو حسابي کول د نورو څيزونو او
خاصيتونو د حسابي کولو د يو مدل په توګه کار کوي۔ د بېلګې په توګه، په
ورته ډول د حالتونو او محاسبو (مثلاً د ټيورينګ ماشينونو د محاسبو)
حسابي کول هم شوني دي۔ په دې سره د محاسبې په يو مدل (لکه ټيورينګ
ماشينونو) کښې محاسبې په بل مدل (لکه بازګشتي تابعو) کښې شبيه کولے شو۔
\olpseganchor{OLP-0281-B009}{end}

\olpseganchor{OLP-0282-B004}{start}
\olfileid{inc}{art}{cod}
\olsection{د سمبولونو کوډول}
\olpseganchor{OLP-0282-B004}{end}

\olpseganchor{OLP-0282-B005}{start}
د لومړۍ درجې منطق بنسټيزه ژبه~\LR{$\Lang L$} دا سمبولونه کاروي:
\[
\lfalse \quad \lnot \quad \lor \quad \land \quad \lif \quad \lforall
\quad \lexists \quad \eq \quad ( \quad ) \quad ,
\]
له دې سره د متغيرونو او ثابت سمبولونه د شمېر وړ سټونه، او د هرې
ځاييزې کچې د تابعې او پريديکاتونه، د شمېر وړ سټونه هم
کاروي۔ موږ
دې سمبولونو هر يوه ته داسې \emph{کوډونه} ټاکلے شو چې هر سمبول ته د
هغه د کوډ په توګه يو ځانګړے عدد ورسېږي، او دوو بېلو سمبولونو ته يو
شان عدد ور نۀ کړل شي۔ پوهېږو چې دا شوني دي، ځکه د ټولو سمبولونو سټ
د شمېر وړ دے، نو د هغه او د طبيعي عددونو د سټ ترمنځ
دوه اړخيزه يو پر يو تابع شته۔ خو دا هم باوري کول غواړو چې سمبول (او د هغه په
اړه ځينې معلومات، لکه د تابع ځاييزه کچه) له کوډ څخه په
محاسبه کېدونکې لار بېرته ترلاسه کړے شو۔ البته د دې کار ډېرې ممکنې
لارې شته۔ دلته يوه داسې لار ده چې بنسټيزې بازګشتي تابعې کاروي۔
(ياد کړئ چې \LR{$\tuple{n_0, \dots, n_k}$} هغه عدد دے چې د عددونو لړۍ
\LR{$n_0$}، \dots، \LR{$n_k$} کوډوي۔)
\olpseganchor{OLP-0282-B005}{end}

\olpseganchor{OLP-0282-B006}{start}
\begin{defn}
که \LR{$s$} د~\LR{$\Lang L$} يو سمبول وي، د هغه \emph{سمبول-کوډ}~\LR{$\scode s$}
داسې تعريف کړئ:
\begin{enumerate}
\item که \LR{$s$} په منطقي سمبولونو کښې وي، \LR{$\scode s$} د لاندې جدول له
  مخې ورکول کېږي:
\[
\begin{array}{cccccccccc}
  \lfalse & \lnot & \lor & \land & \lif & \lforall \\
\tuple{0, 0} & \tuple{0, 1} & \tuple{0, 2} & \tuple{0, 3} &
\tuple{0, 4} & \tuple{0, 5} \\
\lexists & \eq & ( & ) & ,\\
\tuple{0, 6} & \tuple{0, 7} &
\tuple{0, 8} & \tuple{0, 9} & \tuple{0, 10}
\end{array}
\]
\item که \LR{$s$}، \LR{$i$}-م متغير~\LR{$\Obj v_i$} وي، نو
  \LR{$\scode s = \tuple{1, i}$}۔
\item که \LR{$s$}، \LR{$i$}-م ثابت سمبول~\LR{$\Obj c_i$} وي، نو
  \LR{$\scode s = \tuple{2, i}$}۔
\item که \LR{$s$}، \LR{$i$}-م \LR{$n$}-ځايه تابع~\LR{$\Obj f_i^n$} وي، نو
  \LR{$\scode s = \tuple{3, n, i}$}۔
\item که \LR{$s$}، \LR{$i$}-م \LR{$n$}-ځايه پريديکات~\LR{$\Obj P_i^n$} وي، نو
  \LR{$\scode s = \tuple{4, n, i}$}۔
\end{enumerate}
\end{defn}
\olpseganchor{OLP-0282-B006}{end}

\olpseganchor{OLP-0282-B007}{start}
\begin{prop}
لاندې اړيکې بنسټيزې بازګشتي دي:
\begin{enumerate}
\item \LR{$\fn{Fn}(x, n)$} هله صدق کوي چې \LR{$x$} د کوم~\LR{$i$} دپاره د
  \LR{$\Obj f^n_i$} کوډ وي؛ يعنې \LR{$x$} د يوې \LR{$n$}-ځايه تابع کوډ وي۔
\item \LR{$\fn{Pred}(x, n)$} هله صدق کوي چې \LR{$x$} د کوم~\LR{$i$} دپاره د
  \LR{$\Obj P^n_i$} کوډ وي، يا \LR{$x$} د~\LR{$\eq$} کوډ او \LR{$n = 2$} وي؛ يعنې \LR{$x$}
  د يوه \LR{$n$}-ځايه پريديکات کوډ وي۔
\end{enumerate}
\end{prop}
\olpseganchor{OLP-0282-B007}{end}

\olpseganchor{OLP-0282-B008}{start}
\begin{defn}
که \LR{$s_0, \dots, s_{n-1}$} د سمبولونو يوه لړۍ وي، د هغې
\emph{ګوډل شمېره} \LR{$\tuple{\scode{s_0}, \dots, \scode{s_{n-1}}}$} ده۔
\end{defn}
\olpseganchor{OLP-0282-B008}{end}

\olpseganchor{OLP-0282-B009}{start}
\begin{explain}
پام وکړئ چې \emph{کوډونه} او \emph{ګوډل شمېرې} بېل څيزونه دي۔ د
بېلګې په توګه متغير~\LR{$\Obj v_5$} کوډ~\LR{$\scode{\Obj v_5} = \tuple{1, 5}
= 2^2\cdot 3^6$} لري۔ خو که متغير~\LR{$\Obj v_5$} د يوه ترم په توګه
وګڼو، دا د سمبولونو يوه لړۍ (د اوږدوالي~\LR{$1$}) هم ده۔ د
\emph{ترم}~\LR{$\Obj v_5$} \emph{ګوډل شمېره}~\LR{$\Gn{\Obj v_5}$} دا ده:
\LR{$\tuple{\scode{\Obj v_5}} = 2^{\scode{\Obj v_5} + 1} =
2^{2^2\cdot 3^6 + 1}$}۔
\end{explain}
\olpseganchor{OLP-0282-B009}{end}

\olpseganchor{OLP-0282-B010}{start}
\begin{ex}
ياد کړئ چې که \LR{$k_0$}، \dots، \LR{$k_{n-1}$} د عددونو يوه لړۍ وي، نو د
اولي عددونو-د-قوتونو په کوډول کښې د لړۍ کوډ
\LR{$\tuple{k_0, \dots, k_{n-1}}$} دا دے:
\[
2^{k_0+1}\cdot3^{k_1+1}\cdot \dots \cdot p_{n-1}^{k_{n-1}+1},
\]
چې \LR{$p_i$}، \LR{$i$}-م اولي عدد دے (له \LR{$p_0 = 2$} څخه پيلېږي)۔ نو د بېلګې
په توګه فارمول~\LR{$\eq[\Obj v_0][\Obj 0]$}، يا په څرګنده توګه
\LR{${\eq}(\Obj v_0,\Obj c_0)$}، دا ګوډل شمېره لري:
\[
\tuple{\scode{\eq},\scode{(},\scode{\Obj v_0},\scode{,}, \scode{\Obj
    c_0},\scode{)}}.
\]
دلته \LR{$\scode{\eq}$} دا دے: \LR{$\tuple{0,7} = 2^{0+1}\cdot
3^{7+1}$}؛ \LR{$\scode{\Obj v_0}$} دا دے: \LR{$\tuple{1,0} =
2^{1+1}\cdot3^{0+1}$}؛ او همداسې نور۔ نو \LR{$\Gn{=(\Obj v_0,\Obj c_0)}$}
دا دے:
\begin{multline*}
2^{\scode{=} + 1}\cdot 3^{\scode{(}+1}\cdot 5^{\scode{\Obj v_0}+1}
\cdot 7^{\scode{,} + 1} \cdot 11^{\scode{\Obj c_0}+1} \cdot
13^{\scode{)}+1} = \\
2^{2^1\cdot 3^8 + 1}\cdot 3^{2^1\cdot 3^9+1}\cdot 5^{2^2\cdot 3^1+1}
\cdot 7^{2^1\cdot 3^{11} + 1} \cdot 11^{2^3\cdot3^1+1} \cdot
13^{2^1\cdot3^{10}+1} = \\
2^{13\,123}\cdot 3^{39\,367}\cdot 5^{13}\cdot 7^{354\,295}\cdot11^{25}\cdot13^{118\,099}.
\end{multline*}
\end{ex}
\olpseganchor{OLP-0282-B010}{end}

\olpseganchor{OLP-0283-B004}{start}
\olfileid{inc}{art}{trm}
\olsection{د ترمونو کوډول}
\olpseganchor{OLP-0283-B004}{end}

\olpseganchor{OLP-0283-B005}{start}
\begin{explain}
يو ترم يوازې د سمبولونو يو ځانګړے ډول لړۍ ده: دا د ترمونو د جوړښتي
قاعدو له مخې له ثابتونو او متغيرونو څخه په استقرايي ډول جوړېږي۔ ځکه
چې د سمبولونو لړۍ د عددونو په توګه کوډېدلے شي---د سمبولونو د کوډولو
له يوې طرحې او د عددونو د لړيو د کوډولو له يوې لارې سره---نو ترمونو
ته د ګوډل شمېرو ټاکل ستونزمن نۀ دي۔ اصلي ستونزه دا ښودل دي چې يو عدد
چې د سم جوړ شوي ترم ګوډل شمېره وي، دا خاصيت ئې محاسبه کېدونکے، بلکې
بنسټيز بازګشتي دے۔
\end{explain}
\olpseganchor{OLP-0283-B005}{end}

\olpseganchor{OLP-0283-B006}{start}
متغيرونه او ثابت سمبولونه تر ټولو ساده ترمونه دي، او دا ازمويل
چې \LR{$x$} د داسې يوه ترم ګوډل شمېره ده، اسانه دي: \LR{$\fn{Var}(x)$} هله
صدق کوي چې \LR{$x$} د کوم~\LR{$i$} دپاره \LR{$\Gn{\Obj v_i}$} وي۔ په بله وينا،
\LR{$x$} د اوږدوالي~\LR{$1$} يوه لړۍ ده او يوازينے غړے ئې \LR{$(x)_0$} د کوم
متغير~\LR{$\Obj v_i$} کوډ دے؛ يعنې \LR{$x$} د کوم~\LR{$i$} دپاره
\LR{$\tuple{\tuple{1, i}}$} دے۔ په ورته ډول \LR{$\fn{Const}(x)$} هله صدق کوي
چې \LR{$x$} د کوم~\LR{$i$} دپاره \LR{$\Gn{\Obj c_i}$} وي۔ دا دواړه اړيکې بنسټيزې
بازګشتي دي، ځکه که داسې \LR{$i$} موجود وي، بايد \LR{$<x$} وي:
\begin{align*}
  \fn{Var}(x) & \defiff \bexists{i<x}{x = \tuple{\tuple{1, i}}}\\
  \fn{Const}(x) & \defiff \bexists{i<x}{x = \tuple{\tuple{2, i}}}
\end{align*}
\olpseganchor{OLP-0283-B006}{end}

\olpseganchor{OLP-0283-B007}{start}
\begin{prop}
\ollabel{prop:term-primrec}
هغه اړيکې \LR{$\fn{Term}(x)$} او \LR{$\fn{ClTerm}(x)$} چې په ترتيب هله صدق کوي
چې \LR{$x$} د يوه ترم يا تړلي ترم ګوډل شمېره وي، بنسټيزې بازګشتي دي۔
\end{prop}
\olpseganchor{OLP-0283-B007}{end}

\olpseganchor{OLP-0283-B008}{start}
\begin{proof}
د سمبولونو يوه لړۍ~\LR{$s$} هله ترم دے چې د ترمونو داسې لړۍ~\LR{$s_0$}،
\dots، \LR{$s_{k-1}=s$} موجوده وي چې ثبتوي ترم~\LR{$s$} د ترمونو د جوړښتي
قاعدو له مخې له ثابت سمبولونه او متغيرونه څخه څنګه جوړ شوے دے۔
د دې د بيانولو دپاره چې داسې يوه اټکلي جوړښتي لړۍ جوړښتي قاعدې
تعقيبوي، بايد د هر \LR{$i<k$} دپاره له لاندې څخه يو حالت سم وي:
\begin{enumerate}
\item \LR{$s_i$} يو متغير~\LR{$\Obj v_j$} دے؛ يا
\item \LR{$s_i$} يو ثابت سمبول~\LR{$\Obj c_j$} دے؛ يا
\item \LR{$s_i$} له هغو \LR{$n$} ترمونو \LR{$t_1$}، \dots، \LR{$t_n$} څخه، چې له
  ځاے~\LR{$i$} وړاندې راغلي وي، د يوه \LR{$n$}-ځايه
  تابع~\LR{$\Obj f^n_j$} په کارولو جوړ شوے دے۔
\end{enumerate}
د دې د ښودلو دپاره چې د ګوډل شمېرو اړونده اړيکه بنسټيزه بازګشتي
ده، دا شرط بايد په بنسټيز بازګشتي ډول بيان کړو؛ يعنې بنسټيزې
بازګشتي تابعې، اړيکې او محدود کمیت ايښودنه وکاروو۔
\olpseganchor{OLP-0283-B008}{end}

\olpseganchor{OLP-0283-B009}{start}
فرض کړئ \LR{$y$} هغه عدد دے چې لړۍ~\LR{$s_0$}، \dots، \LR{$s_{k-1}$} کوډوي؛ يعنې
\LR{$y = \tuple{\Gn{s_0}, \dots, \Gn{s_{k-1}}}$}۔ دا د ګوډل شمېرې~\LR{$x$}
لرونکي ترم دپاره هله جوړښتي لړۍ کوډوي چې د هر \LR{$i<k$} دپاره:
\begin{enumerate}
\item \LR{$\fn{Var}((y)_i)$}؛ يا
\item \LR{$\fn{Const}((y)_i)$}؛ يا
\item داسې \LR{$n$} او عدد~\LR{$z=\tuple{z_1, \dots, z_n}$} شته چې هر \LR{$z_l$}
  د کوم \LR{$i'<i$} دپاره له \LR{$(y)_{i'}$} سره برابر وي، او
\[
(y)_i = \Gn{\Obj f^n_j(} \concat \fn{flatten}(z) \concat \Gn{)},
\]
\end{enumerate}
او سربېره پر دې \LR{$(y)_{k-1}=x$} وي۔ (تابع \LR{$\fn{flatten}(z)$} لړۍ
\LR{$\tuple{\Gn{t_1}, \dots, \Gn{t_n}}$} په \LR{$\Gn{t_1, \dots, t_n}$}
بدلوي او بنسټيزه بازګشتي ده۔)
\olpseganchor{OLP-0283-B009}{end}

\olpseganchor{OLP-0283-B010}{start}
په (3) کښې شاخصونه \LR{$j$} او \LR{$n$}، د ترمونو \LR{$t_l$} ګوډل شمېرې~\LR{$z_l$}،
او د لړۍ~\LR{$\tuple{z_1, \dots, z_n}$} کوډ~\LR{$z$} ټول له~\LR{$y$} څخه واړه دي۔
پورته \LR{$k$} په \LR{$\len{y}$} بدلولے شو۔ نو «\LR{$y$} د ګوډل شمېرې~\LR{$x$} لرونکي
ترم د جوړښتي لړۍ کوډ دے» په داسې ډول بيانولے شو چې وښيي دا اړيکه
بنسټيزه بازګشتي ده۔
\olpseganchor{OLP-0283-B010}{end}

\olpseganchor{OLP-0283-B011}{start}
اوس يوازې دا باور ترلاسه کول پکار دي چې پر~\LR{$y$} يو بنسټيز بازګشتي حد
شته۔ خو که \LR{$x$} د يوه ترم ګوډل شمېره وي، نو هغه بايد داسې جوړښتي لړۍ
ولري چې تر ټولو زيات \LR{$\len{x}$} ترمونه ولري (ځکه د~\LR{$s$} په جوړښتي
لړۍ کښې هر ترم بايد د~\LR{$s$} په کوم ځاے کښې پيل شي، او دوه فرعي ترمونه
په يوه ځاے نۀ شي پيلېدے)۔ البته د~\LR{$s$} د هر فرعي ترم ګوډل شمېره
\LR{$\le x$} ده۔ نو تل داسې جوړښتي لړۍ شته چې کوډ ئې
\LR{$\le p_{k-1}^{k(x+1)}$} وي، چې \LR{$k=\len{x}$} دے۔
\olpseganchor{OLP-0283-B011}{end}

\olpseganchor{OLP-0283-B012}{start}
د \LR{$\fn{ClTerm}$} دپاره يوازې د متغيرونه فقره پرېږدئ۔
\end{proof}
\olpseganchor{OLP-0283-B012}{end}

\olpseganchor{OLP-0283-B013}{start}
\begin{prob}
وښيئ چې تابع \LR{$\fn{flatten}(z)$}، چې لړۍ
\LR{$\tuple{\Gn{t_1}, \dots, \Gn{t_n}}$} په \LR{$\Gn{t_1, \dots, t_n}$}
بدلوي، بنسټيزه بازګشتي ده۔
\end{prob}
\olpseganchor{OLP-0283-B013}{end}

\olpseganchor{OLP-0283-B014}{start}
\begin{prop}\ollabel{prop:num-primrec}
  تابع \LR{$\fn{num}(n) = \Gn{\num{n}}$} بنسټيزه بازګشتي ده۔
\end{prop}
\olpseganchor{OLP-0283-B014}{end}

\olpseganchor{OLP-0283-B015}{start}
\begin{proof}
  \LR{$\fn{num}(n)$} په بنسټيز بازګښت داسې تعريفوو:
  \begin{align*}
    \fn{num}(0) & = \Gn{\Obj 0}\\
    \fn{num}(n+1) & = \Gn{\prime(} \concat \fn{num}(n) \concat \Gn{)}.
  \end{align*}
\end{proof}
\olpseganchor{OLP-0283-B015}{end}

\olpseganchor{OLP-0284-B004}{start}
\olfileid{inc}{art}{frm}
\olsection{د فارمولونه کوډول}
\olpseganchor{OLP-0284-B004}{end}

\olpseganchor{OLP-0284-B005}{start}
کله چې اړيکه \LR{$\fn{Term}(x)$} مو په بنسټيز بازګشتي ډول تعريف کړه، د
هغې په کارولو د فارمولونه اړونده اړيکه \LR{$\fn{Frm}(x)$} هم په
بنسټيز بازګشتي ډول تعريفولے شو۔
\olpseganchor{OLP-0284-B005}{end}

\olpseganchor{OLP-0284-B006}{start}
\begin{prop}
هغه اړيکه \LR{$\fn{Atom}(x)$} چې هله صدق کوي چې \LR{$x$} د يوه اټومي
فارمول ګوډل شمېره وي، بنسټيزه بازګشتي ده۔
\end{prop}
\olpseganchor{OLP-0284-B006}{end}

\olpseganchor{OLP-0284-B007}{start}
\begin{proof}
عدد \LR{$x$} هله د يوه اټومي فارمول ګوډل شمېره ده چې له لاندې څخه
يو حالت صدق وکړي:
\begin{enumerate}
\item داسې \LR{$n$}، \LR{$j<x$} او \LR{$z<x$} شته چې د هر \LR{$i<n$} دپاره
  \LR{$\fn{Term}((z)_i)$} وي او \LR{$x = $}
\[
\Gn{\Obj P^n_j(} \concat \fn{flatten}(z) \concat \Gn{)}.
\]
\item داسې \LR{$z_1,z_2<x$} شته چې \LR{$\fn{Term}(z_1)$}،
  \LR{$\fn{Term}(z_2)$} وي، او \LR{$x = $}
\[
\Gn{{\eq}(} \concat z_1 \concat \Gn{,} \concat z_2 \concat{\Gn{)}}.
\]
  \item \LR{$x = \Gn{\lfalse}$}۔
  \item \LR{$x = \Gn{\ltrue}$}۔
\end{enumerate}
\end{proof}
\olpseganchor{OLP-0284-B007}{end}

\olpseganchor{OLP-0284-B008}{start}
\begin{prop}
\ollabel{prop:frm-primrec}
هغه اړيکه \LR{$\fn{Frm}(x)$} چې هله صدق کوي چې \LR{$x$} د فارمول ګوډل
شمېره وي، بنسټيزه بازګشتي ده۔
\end{prop}
\olpseganchor{OLP-0284-B008}{end}

\olpseganchor{OLP-0284-B009}{start}
\begin{proof}
د سمبولونو لړۍ~\LR{$s$} هله فارمول ده چې د فارمول داسې
جوړښتي لړۍ~\LR{$s_0$}، \dots، \LR{$s_{k-1}=s$} موجوده وي چې ثبتوي~\LR{$s$} د
جوړښتي قاعدو له مخې له اټومي فارمولونه څخه څنګه جوړ شوے دے۔ د هر
\LR{$s_i$} کوډ (او په حقيقت کښې د لړۍ \LR{$\tuple{s_0, \dots, s_{k-1}}$} کوډ
هم) د~\LR{$s$} له کوډ~\LR{$x$} څخه وړوکے دے۔
\end{proof}
\olpseganchor{OLP-0284-B009}{end}

\olpseganchor{OLP-0284-B010}{start}
\begin{prob}
د \olref[inc][art][frm]{prop:frm-primrec} تفصيلي ثبوت د
\olref[inc][art][trm]{prop:term-primrec} د لومړي ثبوت په څېر ورکړئ۔
\end{prob}
\olpseganchor{OLP-0284-B010}{end}

\olpseganchor{OLP-0284-B011}{start}
\begin{prop}
\ollabel{prop:freeocc-primrec}
اړيکه \LR{$\fn{FreeOcc}(x,z,i)$}، چې هله صدق کوي چې د ګوډل شمېرې~\LR{$x$}
لرونکي فارمول \LR{$i$}-م سمبول د ګوډل شمېرې~\LR{$z$} لرونکي متغير يوه ازاده
پېښه وي، بنسټيزه بازګشتي ده۔
\end{prop}
\olpseganchor{OLP-0284-B011}{end}

\olpseganchor{OLP-0284-B012}{start}
\begin{proof}
تمرین۔
\end{proof}
\olpseganchor{OLP-0284-B012}{end}

\olpseganchor{OLP-0284-B013}{start}
\begin{prob}
\olref[inc][art][frm]{prop:freeocc-primrec} ثابت کړئ۔ دا حقيقت
کارولے شئ چې د فارمول هر هغه فرعي تار چې پخپله فارمول
وي، د هغه فرعي-فارمول دے۔
\end{prob}
\olpseganchor{OLP-0284-B013}{end}

\olpseganchor{OLP-0284-B014}{start}
\begin{prop}
هغه خاصيت \LR{$\fn{Sent}(x)$} چې هله صدق کوي چې \LR{$x$} د تړلے فارمول
ګوډل شمېره وي، بنسټيز بازګشتي دے۔
\end{prop}
\olpseganchor{OLP-0284-B014}{end}

\olpseganchor{OLP-0284-B015}{start}
\begin{proof}
تړلے فارمول داسې فارمول دے چې د متغيرونه ازادې پېښې
نۀ لري۔ نو \LR{$\fn{Sent}(x)$} هله صدق کوي چې
\begin{multline*}
\bforall{i<\len{x}}{\bforall{z<x}{}}\\
(\bexists{j<z}{z=\Gn{\Obj v_j}} \lif \lnot\fn{FreeOcc}(x,z,i)).
\end{multline*}
\end{proof}
\olpseganchor{OLP-0284-B015}{end}

\olpseganchor{OLP-0285-B004}{start}
\olfileid{inc}{art}{sub}
\olsection{تعويض}
\olpseganchor{OLP-0285-B004}{end}

\olpseganchor{OLP-0285-B005}{start}
ياد کړئ چې تعويض هغه عمليه ده چې په فارمول~\LR{$!A$} کښې د
متغير~\LR{$u$} ټولې ازادې پېښې په ترم~\LR{$t$} بدلوي، او
\LR{$\Subst{!A}{t}{u}$} ليکل کېږي۔ کله چې دا عمليه د متغيرونه،
فارمولونه او ترمونو پر ګوډل شمېرو ترسره شي، بنسټيزه بازګشتي ده۔
\olpseganchor{OLP-0285-B005}{end}

\olpseganchor{OLP-0285-B006}{start}
\begin{prop}\ollabel{prop:subst-primrec}
داسې بنسټيزه بازګشتي تابع \LR{$\fn{Subst}(x,y,z)$} شته چې دا خاصيت لري:
\[
\fn{Subst}(\Gn{!A}, \Gn{t}, \Gn{u}) = \Gn{\Subst{!A}{t}{u}}.
\]
\end{prop}
\olpseganchor{OLP-0285-B006}{end}

\olpseganchor{OLP-0285-B007}{start}
\begin{proof}
بيا تابع \LR{$\fn{hSubst}$} په بنسټيز بازګښت داسې تعريفولے شو:
\begin{multline*}
\begin{aligned}
\fn{hSubst}(x, y, z, 0) & = \emptyseq \\
\fn{hSubst}(x, y, z, i+1) & =
\end{aligned}\\
\begin{cases}
\fn{hSubst}(x, y, z, i) \concat y & \text{\RL{که \LR{$\fn{FreeOcc}(x, z, i)$} وي}} \\
\fn{append}(\fn{hSubst}(x, y, z, i), (x)_{i}) & \text{\RL{که نۀ وي۔}}
\end{cases}
\end{multline*}
اوس \LR{$\fn{Subst}(x,y,z)$} د \LR{$\fn{hSubst}(x,y,z,\len{x})$} په توګه
تعريفېدے شي۔
\end{proof}
\olpseganchor{OLP-0285-B007}{end}

\olpseganchor{OLP-0285-B008}{start}
\begin{prop}
\ollabel{prop:free-for}
اړيکه \LR{$\fn{FreeFor}(x,y,z)$}، چې هله صدق کوي چې د ګوډل شمېرې~\LR{$y$}
لرونکے ترم د ګوډل شمېرې~\LR{$z$} لرونکي متغير دپاره د ګوډل شمېرې~\LR{$x$}
لرونکي فارمول کښې ازاد وي، بنسټيزه بازګشتي ده۔
\end{prop}
\olpseganchor{OLP-0285-B008}{end}

\olpseganchor{OLP-0285-B009}{start}
\begin{proof} تمرین۔ \end{proof}
\olpseganchor{OLP-0285-B009}{end}

\olpseganchor{OLP-0285-B010}{start}
\begin{prob}
\olref[inc][art][sub]{prop:free-for} ثابت کړئ۔
\end{prob}
\olpseganchor{OLP-0285-B010}{end}

\olpseganchor{OLP-0286-B004}{start}
\olfileid{inc}{art}{plk}
\olsection{په \LR{$\Log{LK}$} کښې اشتقاقونه}
\olpseganchor{OLP-0286-B004}{end}

\olpseganchor{OLP-0286-B005}{start}
\begin{explain}
د اشتقاقونه د حسابي کولو دپاره بايد اشتقاقونه د عددونو
په توګه تمثيل کړو۔ ځکه چې اشتقاقونه د سېکوېنټونو داسې ونې دي
چې هر استنتاج يوه نښه هم لري، نو بازګشتي تمثيل تر ټولو څرګنده لار
ده: يو اشتقاق د يوې څوګونې په توګه تمثيلوو چې اجزا ئې
پاې-سېکوېنټ، نښه، او د وروستي استنتاج مقدماتو ته د رسېدونکو
فرعي-اشتقاقونه تمثيلونه دي۔
\end{explain}
\olpseganchor{OLP-0286-B005}{end}

\olpseganchor{OLP-0286-B006}{start}
\begin{defn}
که \LR{$\Gamma$} د تړلي فارمولونه يوه متناهي لړۍ وي،
\LR{$\Gamma=\tuple{!A_1, \dots, !A_n}$}، نو
\LR{$\Gn{\Gamma}=\tuple{\Gn{!A_1}, \dots, \Gn{!A_n}}$}۔
\olpseganchor{OLP-0286-B006}{end}

\olpseganchor{OLP-0286-B007}{start}
که \LR{$\Gamma \Sequent \Delta$} يو سېکوېنټ وي، نو د
\LR{$\Gamma \Sequent \Delta$} ګوډل شمېره دا ده:
\[
\Gn{\Gamma \Sequent \Delta} = \tuple{\Gn{\Gamma}, \Gn{\Delta}}
\]
\olpseganchor{OLP-0286-B007}{end}

\olpseganchor{OLP-0286-B008}{start}
که \LR{$\pi$} په \LR{$\Log{LK}$} کښې اشتقاق وي، \LR{$\Gn{\pi}$} داسې
تعريفېږي:
\begin{enumerate}
\item که \LR{$\pi$} يوازې له ابتدايي سېکوېنټ \LR{$\Gamma \Sequent \Delta$}
  څخه جوړ وي، نو \LR{$\Gn{\pi}$} دا دے:
  \[
    \tuple{0, \Gn{\Gamma \Sequent \Delta}}.
  \]
\item که \LR{$\pi$} په داسې استنتاج پاې ته رسېږي چې يو يا دوه مقدمات
  لري، پايله ئې \LR{$\Gamma \Sequent \Delta$} وي، او \LR{$\pi_1$} او \LR{$\pi_2$}
  هغه سمدستي فرعي ثبوتونه وي چې د وروستي استنتاج پر مقدماتو پاې ته
  رسېږي، نو \LR{$\Gn{\pi}$} په ترتيب له دې دوو څخه يو دے:
  \begin{align*}
    & \tuple{1, \Gn{\pi_1}, \Gn{\Gamma \Sequent \Delta}, k} \text{\RL{ يا}}\\
    & \tuple{2, \Gn{\pi_1}, \Gn{\pi_2}, \Gn{\Gamma \Sequent \Delta}, k},
  \end{align*}
  چې \LR{$k$} د لاندې جدول له مخې د وروستي استنتاج د کارول شوې قاعدې
  له مخې ټاکل کېږي:
\olpseganchor{OLP-0286-B008}{end}

\olpseganchor{OLP-0286-B009}{start}
  \begin{tabular}{lccccccc}
    \text{\RL{قاعده:}} & \LeftR{\Weakening} & \RightR{\Weakening} &
    \LeftR{\Contraction} & \RightR{\Contraction} &
    \LeftR{\Exchange} & \RightR{\Exchange} \\
    \LR{$k$}: & 1 & 2 & 3 & 4 & 5 & 6 \\[2ex]
    \text{\RL{قاعده:}} & \LeftR{\lnot} & \RightR{\lnot} &
    \LeftR{\land} & \RightR{\land} &
    \LeftR{\lor} & \RightR{\lor} \\
    \LR{$k$}: & 7 & 8 & 9 & 10 & 11 & 12 \\[2ex]
    \text{\RL{قاعده:}} & \LeftR{\lif} & \RightR{\lif} &
    \LeftR{\lforall} & \RightR{\lforall} &
    \LeftR{\lexists} & \RightR{\lexists} \\
    \LR{$k$}: & 13 & 14 & 15 & 16 & 17 & 18 \\[2ex]
    \text{\RL{قاعده:}} & \Cut & = \\
    \LR{$k$}: & 19 & 20
  \end{tabular}
\end{enumerate}
\end{defn}
\olpseganchor{OLP-0286-B009}{end}

\olpseganchor{OLP-0286-B010}{start}
\begin{ex}
  دا ډېر ساده اشتقاق وګورئ:
  \begin{prooftree}
    \Axiom$!A \fCenter !A$
    \RightLabel{\LeftR{\land}}
    \UnaryInf$!A \land !B \fCenter !A$
    \RightLabel{\RightR{\lif}}
    \UnaryInf$\fCenter (!A \land !B) \lif !A$
  \end{prooftree}
  د ابتدايي سېکوېنټ ګوډل شمېره به
  \LR{$p_0=\tuple{0, \Gn{!A \Sequent !A}}$} وي۔ د هغه اشتقاق
  ګوډل شمېره چې د~\LR{$\LeftR{\land}$} پر پايله ختمېږي،
  \LR{$p_1=\tuple{1,p_0,\Gn{!A \land !B \Sequent !A},9}$} ده
  (\LR{$1$} ځکه چې \LR{$\LeftR{\land}$} يو مقدمه لري؛ بيا د پايلې
  \LR{$!A \land !B \Sequent !A$} ګوډل شمېره؛ او \LR{$9$} هغه عدد دے چې
  \LR{$\LeftR{\land}$} کوډوي)۔ بيا د ټول اشتقاق ګوډل شمېره
  \LR{$\tuple{1,p_1,\Gn{\Sequent (!A \land !B) \lif !A},14}$} ده؛ يعنې
  \[
  \tuple{1, \tuple{1, \tuple{0, \Gn{!A \Sequent !A}}, \Gn{!A \land !B \Sequent !A}, 9},
    \Gn{\Sequent (!A \land !B) \lif !A}, 14}.
  \]
  \textbf{د ژباړې د نسخې سمون (OLCMP-074):} سرچينې د~\LR{$p_0$} په
  دننني ګوډل عبارت کښې د~\LR{$!A$} وروسته يو بې‌جوړې تړونکے ګرد قوس
  ايښے و؛ دلته لرې شو څو کوډ هماغه ابتدايي سېکوېنټ ونيسي۔
\end{ex}
\olpseganchor{OLP-0286-B010}{end}

\olpseganchor{OLP-0286-B011}{start}
\begin{explain}
چې د اشتقاقونه پر يو تمثيل هوکړې ته ورسېدو، دا هم بايد وښيو چې
داسې اشتقاقونه په بنسټيز بازګشتي ډول اداره کولے شو، او د هغو
اړين خاصيتونه او اړيکې هم په همدې ډول بيانولے شو۔ ځينې عمليات ساده
دي: د بېلګې په توګه، د اشتقاق له ګوډل شمېرې~\LR{$p$} څخه
\LR{$\fn{EndSequent}(p)=(p)_{(p)_0+1}$} د هغه د پاې-سېکوېنټ ګوډل شمېره
راکوي، او \LR{$\fn{LastRule}(p)=(p)_{(p)_0+2}$} د هغه د وروستۍ قاعدې
کوډ راکوي۔ خاصيت \LR{$\fn{Sequent}(s)$} چې داسې تعريفېږي:
\[
  \len{s} = 2 \land \bforall{i<\len{(s)_0} + \len{(s)_1}}{\fn{Sent}(((s)_0 \concat (s)_1)_i)}
\]
هله د~\LR{$s$} په اړه صدق کوي چې \LR{$s$} د تړلي فارمولونه لرونکي يوه
سېکوېنټ ګوډل شمېره وي۔ ځينې نور عمليات ډېر ستونزمن دي۔ لږ تر لږه د
هغو د ترسره کولو لار به راونغاړو۔ موخه دا ده چې وښيو اړيکه «\LR{$\pi$}
له~\LR{$\Gamma$} څخه د~\LR{$!A$} يو اشتقاق دے» د~\LR{$\pi$} او~\LR{$!A$} د
ګوډل شمېرو بنسټيزه بازګشتي اړيکه ده۔
\textbf{د ژباړې د نسخې سمون (OLCMP-075):} سرچينې دې تابع ته په
لومړي تعريف کښې \LR{$\fn{EndSeq}$} ويلي، خو په همدې پاتې فصل کښې هر ځاے
\LR{$\fn{EndSequent}$} کاروي؛ دلته يو ثابت نوم کارول شوے دے۔
\end{explain}
\olpseganchor{OLP-0286-B011}{end}

\olpseganchor{OLP-0286-B012}{start}
\begin{prop}\ollabel{prop:followsby}
  خاصيت \LR{$\fn{Correct}(p)$}، چې هله صدق کوي چې د ګوډل شمېرې~\LR{$p$}
  لرونکي اشتقاق~\LR{$\pi$} وروستے استنتاج سم وي، بنسټيز بازګشتي
  دے۔
\end{prop}
\olpseganchor{OLP-0286-B012}{end}

\olpseganchor{OLP-0286-B013}{start}
\begin{proof}
  \LR{$\Gamma \Sequent \Delta$} هله ابتدايي سېکوېنټ دے چې يا داسې
  تړلے فارمول~\LR{$!A$} وي چې \LR{$\Gamma \Sequent \Delta$} همدا
  \LR{$!A \Sequent !A$} وي، يا داسې ترم~\LR{$t$} وي چې
  \LR{$\Gamma \Sequent \Delta$} همدا \LR{$\emptyset \Sequent \eq[t][t]$} وي۔
  د ګوډل شمېرو په ژبه، \LR{$\fn{InitSeq}(s)$} هله صدق کوي چې
  \begin{align*}
  \bexists{x < s}{} (\fn{Sent}(x) & \land
  s = \tuple{\tuple{x},\tuple{x}})
  \lor {}\\
  \bexists{t<s}{} (\fn{Term}(t) & \land
  s = \tuple{0, \tuple{\Gn{{\eq}(} \concat t \concat \Gn{,} \concat t \concat \Gn{)}}}).
  \end{align*}
\olpseganchor{OLP-0286-B013}{end}

\olpseganchor{OLP-0286-B014}{start}
  د استنتاج د هرې قاعدې~\LR{$R$} دپاره دا هم بايد وښيو چې اړيکه
  \LR{$\fn{FollowsBy}_R(p)$} بنسټيزه بازګشتي ده؛
  \LR{$\fn{FollowsBy}_R(p)$} هله صدق کوي چې \LR{$p$}
  د اشتقاق~\LR{$\pi$} ګوډل شمېره وي، او د~\LR{$\pi$} پاې-سېکوېنټ د
  \LR{$R$} په سم پلي کولو د~\LR{$\pi$} د سمدستي فرعي-اشتقاقونه له
  پاې-سېکوېنټونو څخه راشي۔
\olpseganchor{OLP-0286-B014}{end}

\olpseganchor{OLP-0286-B015}{start}
  يو ساده حالت د \RightR{\land} قاعده ده۔ که \LR{$\pi$} په يوه سم
  \LR{$\RightR{\land}$} استنتاج ختم شي، داسې ښکاري:
  \begin{prooftree}
    \AxiomC{}
    \RightLabel{\LR{$\pi_1$}}
    \Deduce$\Gamma \fCenter \Delta, !A$
\olpseganchor{OLP-0286-B015}{end}

\olpseganchor{OLP-0286-B016}{start}
    \AxiomC{}
    \RightLabel{\LR{$\pi_2$}}
    \Deduce$\Gamma \fCenter \Delta, !B$
\olpseganchor{OLP-0286-B016}{end}

\olpseganchor{OLP-0286-B017}{start}
    \RightLabel{\RightR\land}
    \BinaryInf$\Gamma \fCenter \Delta, !A \land !B$
  \end{prooftree}
  نو په اشتقاق~\LR{$\pi$} کښې وروستے استنتاج هله د
  \LR{$\RightR{\land}$} سم پلي کول دي چې د تړلي فارمولونه داسې لړۍ
  \LR{$\Gamma$} او \LR{$\Delta$}، او داسې دوه تړلي فارمولونه \LR{$!A$} او~\LR{$!B$} شته
  چې د~\LR{$\pi_1$} پاې-سېکوېنټ \LR{$\Gamma \Sequent \Delta,!A$}، د~\LR{$\pi_2$}
  پاې-سېکوېنټ \LR{$\Gamma \Sequent \Delta,!B$}، او د~\LR{$\pi$} پاې-سېکوېنټ
  \LR{$\Gamma \Sequent \Delta,!A\land !B$} وي۔ دا يوازې د ګوډل شمېرو
  ژبې ته ژباړل پکار دي۔ که \LR{$s=\Gn{\Gamma \Sequent \Delta}$} وي، نو
  \LR{$(s)_0=\Gn{\Gamma}$} او \LR{$(s)_1=\Gn{\Delta}$} دي۔ نو
  \LR{$\fn{FollowsBy}_{\RightR{\land}}(p)$} هله صدق کوي چې
\begin{align*}
& \bexists{g < p}{\bexists{d < p}{\bexists{a < p}{\bexists{b < p}{\quad}}}} \\
& \qquad \fn{EndSequent}(p) =
  \tuple{g, d \concat
    \tuple{\Gn{(} \concat a \concat \Gn{\land} \concat b \concat \Gn{)}}} \land {} \\
& \qquad \fn{EndSequent}((p)_1) = \tuple{g, d \concat \tuple{a}} \land {}\\
& \qquad \fn{EndSequent}((p)_2) = \tuple{g, d \concat \tuple{b}} \land {}\\
& \qquad (p)_0 = 2 \land \fn{LastRule}(p) = 10.
\end{align*}
جلا کرښې په ترتيب دا بيانوي: «داسې لړۍ~(\LR{$\Gamma$}) شته چې ګوډل
شمېره ئې~\LR{$g$} ده، داسې لړۍ~(\LR{$\Delta$}) شته چې ګوډل شمېره ئې~\LR{$d$} ده،
داسې فارمول~(\LR{$!A$}) شته چې ګوډل شمېره ئې~\LR{$a$} ده، او داسې
فارمول~(\LR{$!B$}) شته چې ګوډل شمېره ئې~\LR{$b$} ده»؛ داسې چې «د~\LR{$\pi$}
پاې-سېکوېنټ \LR{$\Gamma \Sequent \Delta,!A\land !B$} دے»، «د~\LR{$\pi_1$}
پاې-سېکوېنټ \LR{$\Gamma \Sequent \Delta,!A$} دے»، «د~\LR{$\pi_2$}
پاې-سېکوېنټ \LR{$\Gamma \Sequent \Delta,!B$} دے»، او «\LR{$\pi$} دوه
سمدستي فرعي ثبوتونه لري او وروستۍ استنتاجي قاعده
\LR{$\RightR\land$} (په شمېره~\LR{$10$}) ده»۔
\olpseganchor{OLP-0286-B017}{end}

\olpseganchor{OLP-0286-B018}{start}
په~\LR{$\pi$} کښې وروستے استنتاج هله د \LR{$\RightR\lexists$} سم پلي کول
دي چې داسې لړۍ \LR{$\Gamma$} او \LR{$\Delta$}، داسې فارمول~\LR{$!A$}، متغير
\LR{$x$} او ترم~\LR{$t$} شته چې د~\LR{$\pi$} پاې-سېکوېنټ
\LR{$\Gamma \Sequent \Delta,\lexists[x][!A]$} او د~\LR{$\pi_1$}
پاې-سېکوېنټ \LR{$\Gamma \Sequent \Delta,\Subst{!A}{t}{x}$} وي۔ نو د
ګوډل شمېرو په ژبه \LR{$\fn{FollowsBy}_{\RightR\lexists}(p)$} هله صدق کوي
چې
\begin{align*}
  & \bexists{g<p}{\bexists{d<p}{\bexists{a<p}{\bexists{x<p}{\bexists{t<p}{\quad}}}}}\\
  & \qquad \fn{EndSequent}(p) =
  \tuple{
    g,
    d \concat \tuple{\Gn{\lexists} \concat x \concat a}
  } \land {}\\
  & \qquad \fn{EndSequent}((p)_1) =
  \tuple{
    g,
    d \concat \tuple{\fn{Subst}(a, t, x)}
  } \land {}\\
  & \qquad (p)_0 = 1 \land \fn{LastRule}(p) = 18.
\end{align*}
بيا \LR{$\fn{Correct}(p)$} داسې تعريفوو:
\begin{multline*}
  \fn{Sequent}(\fn{EndSequent}(p)) \land {}\\
  [(\fn{LastRule}(p) = 1 \land
  \fn{FollowsBy}_{\LeftR\Weakening}(p)) \lor \dots \lor {}\\
  (\fn{LastRule}(p) = 20 \land \fn{FollowsBy}_{\eq}(p)) \lor {}\\
  (p)_0 = 0 \land \fn{InitialSeq}(\fn{EndSequent}(p))]
\end{multline*}
لومړۍ کرښه باوري کوي چې د~\LR{$p$} پاې-سېکوېنټ په رښتيا د
تړلي فارمولونه يو سېکوېنټ دے۔ وروستۍ کرښه هغه حالت رانغاړي چې \LR{$p$}
يوازې يو ابتدايي سېکوېنټ وي۔
\textbf{د ژباړې د نسخې سمون (OLCMP-076):} سرچينې د دې تعريف په
تشريح کښې ناڅاپه د~\LR{$d$} پاې-سېکوېنټ ياد کړے و، خو ټول تعريف او قضيه
د~\LR{$p$} په اړه دي؛ دلته هماغه يو شان ميتا-متغير ساتل شوے دے۔
\end{proof}
\olpseganchor{OLP-0286-B018}{end}

\olpseganchor{OLP-0286-B019}{start}
\begin{prob}
  لاندې خاصيتونه د \olref[inc][art][plk]{prop:followsby} په څېر
  تعريف کړئ:
  \begin{enumerate}
  \item \LR{$\fn{FollowsBy}_{\Cut}(p)$}؛
  \item \LR{$\fn{FollowsBy}_{\LeftR\lif}(p)$}؛
  \item \LR{$\fn{FollowsBy}_{\eq}(p)$}؛
  \item \LR{$\fn{FollowsBy}_{\RightR{\lforall}}(p)$}۔
  \end{enumerate}
  د وروستي خاصيت دپاره دا هم بايد وښيئ چې په بنسټيز بازګشتي ډول
  ازمويلے شئ چې د ګوډل شمېرې~\LR{$p$} لرونکي اشتقاق وروستے
  استنتاج د ځانګړي متغير شرط پوره کوي؛ يعنې د~\LR{$\RightR{\lforall}$}
  ځانګړے متغير~\LR{$a$} په پاې-سېکوېنټ کښې نۀ راځي۔
\end{prob}
\olpseganchor{OLP-0286-B019}{end}

\olpseganchor{OLP-0286-B020}{start}
\begin{prop}
  \ollabel{prop:deriv}
  اړيکه \LR{$\fn{Deriv}(p)$}، چې هله صدق کوي چې \LR{$p$} د يوه سم
  اشتقاق~\LR{$\pi$} ګوډل شمېره وي، بنسټيزه بازګشتي ده۔
\end{prop}
\olpseganchor{OLP-0286-B020}{end}

\olpseganchor{OLP-0286-B021}{start}
\begin{proof}
  اشتقاق~\LR{$\pi$} هله سم دے چې هر استنتاج ئې د يوې قاعدې سم
  پلي کول وي؛ يعنې د هغه هر فرعي-اشتقاقونه په سم استنتاج پاې ته
  ورسېږي۔ نو \LR{$\fn{Deriv}(p)$} هله صدق کوي چې
  \[
  \bforall{i<\len{\fn{SubtreeSeq}(p)}}{\fn{Correct}((\fn{SubtreeSeq}(p))_i)}.
  \]
  \textbf{د ژباړې د نسخې سمون (OLCMP-077):} سرچينې په دې کرښه کښې
  اړيکه \LR{$\fn{Deriv}(d)$} بللې وه، سره له دې چې قضيه او د ښي لوري ټول
  استعمالونه~\LR{$p$} لري؛ د \LR{$\fn{Correct}$} وروستے تړونکے قوس هم ورک و۔
  دواړه دلته د هماغه ټاکلي اشتقاق د سم شرط دپاره برابر شول۔
\end{proof}
\olpseganchor{OLP-0286-B021}{end}

\olpseganchor{OLP-0286-B022}{start}
\begin{prop}
فرض کړئ \LR{$\Gamma$} د تړلي فارمولونه يو بنسټيز بازګشتي سټ دے۔ بيا
اړيکه \LR{$\Prf[\Gamma](x,y)$}، چې وايي «\LR{$x$} د
\LR{$\Gamma_0 \Sequent !A$} د اشتقاق~\LR{$\pi$} کوډ دے، چې د کوم
متناهي \LR{$\Gamma_0\subseteq\Gamma$} دپاره وي، او \LR{$y$} د~\LR{$!A$} ګوډل
شمېره ده»، بنسټيزه بازګشتي ده۔
\end{prop}
\olpseganchor{OLP-0286-B022}{end}

\olpseganchor{OLP-0286-B023}{start}
\begin{proof}
فرض کړئ «\LR{$y\in\Gamma$}» د بنسټيز بازګشتي
پريديکات~\LR{$R_\Gamma(y)$} په واسطه ورکول کېږي۔ بايد وښيو چې
\LR{$\Prf[\Gamma](x,y)$} بنسټيزه بازګشتي ده؛ دا اړيکه هله صدق کوي چې
\LR{$y$} د يوې جملې~\LR{$!A$} ګوډل شمېره وي او \LR{$x$} د داسې
\LR{$\Log{LK}$}-اشتقاق کوډ وي چې پاې-سېکوېنټ ئې
\LR{$\Gamma_0\Sequent !A$} وي۔
\olpseganchor{OLP-0286-B023}{end}

\olpseganchor{OLP-0286-B024}{start}
د تېرې قضيې له مخې خاصيت \LR{$\fn{Deriv}(x)$}، چې هله صدق کوي چې \LR{$x$} په
\LR{$\Log{LK}$} کښې د يوه سم اشتقاق~\LR{$\pi$} کوډ وي، بنسټيز بازګشتي
دے۔ که \LR{$x$} داسې کوډ وي، \LR{$\fn{EndSequent}(x)$} د~\LR{$\pi$} د
پاې-سېکوېنټ کوډ دے، نو \LR{$(\fn{EndSequent}(x))_0$} د پاې-سېکوېنټ د
کيڼ لوري کوډ او \LR{$(\fn{EndSequent}(x))_1$} د ښي لوري کوډ دے۔ نو «د
~\LR{$\pi$} د پاې-سېکوېنټ ښے لور~\LR{$!A$} دے» داسې بيانولے شو:
\LR{$\len{(\fn{EndSequent}(x))_1}=1\land
((\fn{EndSequent}(x))_1)_0=y$}۔
\textbf{د ژباړې د نسخې سمون (OLCMP-078):} سرچينې د ښي لوري د
يوازينۍ جملې کوډ له~\LR{$x$} سره برابر کړے و، خو \LR{$x$} د ټول
اشتقاق کوډ دے او د قضيې له مخې د پاې جملې کوډ~\LR{$y$} دے؛ دلته
هماغه \LR{$y$} کارول شوے دے۔
د~\LR{$\pi$} د پاې-سېکوېنټ کيڼ لور البته پخپله متناهي دے؛ يوازې دا بيانول
پکار دي چې هره جمله ئې په~\LR{$\Gamma$} کښې ده۔ نو
\LR{$\Prf[\Gamma](x,y)$} داسې تعريفولے شو:
\begin{align*}
\Prf[\Gamma](x, y) \defiff {}&
\fn{Deriv}(x) \land {} \\
& \bforall{i <
  \len{(\fn{EndSequent}(x))_0}}{R_\Gamma(((\fn{EndSequent}(x))_0)_i)} \land {}\\
& \len{(\fn{EndSequent}(x))_1} = 1 \land ((\fn{EndSequent}(x))_1)_0 = y.
\end{align*}
\end{proof}
\olpseganchor{OLP-0286-B024}{end}

\olpseganchor{OLP-0287-B004}{start}
\olfileid{inc}{art}{pnd}
\olsection{په طبيعي استنتاج کښې اشتقاقونه}
\olpseganchor{OLP-0287-B004}{end}

\olpseganchor{OLP-0287-B005}{start}
\begin{explain}
د اشتقاقونه د حسابي کولو دپاره بايد اشتقاقونه د عددونو
په توګه تمثيل کړو۔ ځکه چې اشتقاقونه د فارمولونه داسې ونې
دي چې هر استنتاج يوه يا دوه نښې لري، نو بازګشتي تمثيل تر ټولو څرګنده
لار ده: يو اشتقاق د يوې څوګونې په توګه تمثيلوو چې اجزا ئې د
وروستي استنتاج مقدماتو ته د رسېدونکو سمدستي فرعي-اشتقاقونه
شمېر، د دغو فرعي-اشتقاقونه تمثيلونه، پاې-فارمول، د وروستي
استنتاج د ختمولو نښه، او يو داسې عدد دي چې د وروستي استنتاج ډول ښيي۔
\end{explain}
\olpseganchor{OLP-0287-B005}{end}

\olpseganchor{OLP-0287-B006}{start}
\begin{defn}
  که \LR{$\delta$} په طبيعي استنتاج کښې اشتقاق وي، نو
  \LR{$\Gn{\delta}$} په استقرايي ډول داسې تعريفېږي:
  \begin{enumerate}
  \item
    که \LR{$\delta$} يوازې له فرض~\LR{$!A$} څخه جوړ وي، نو \LR{$\Gn{\delta}$}
    دا دے: \LR{$\tuple{0,\Gn{!A},n}$}۔ عدد~\LR{$n$} که دا يو
    نۀ ايستل شوې فرض وي، \LR{$0$} دے؛ او که نۀ، عددي نښه ده۔
  \item
    که \LR{$\delta$} په داسې استنتاج پاې ته رسېږي چې صفر، يو، دوه يا درې
    مقدمات لري، نو \LR{$\Gn{\delta}$} په ترتيب له دې څخه يو دے:
    \begin{align*}
      & \tuple{0, \Gn{!A}, n, k}, \\
      & \tuple{1, \Gn{\delta_1}, \Gn{!A}, n, k}, \\
      & \tuple{2, \Gn{\delta_1}, \Gn{\delta_2}, \Gn{!A}, n, k}, \text{\RL{ يا}}\\
      & \tuple{3, \Gn{\delta_1}, \Gn{\delta_2}, \Gn{\delta_3}, \Gn{!A}, n,
        k},
    \end{align*}
    دلته \LR{$\delta_1$}، \LR{$\delta_2$} او \LR{$\delta_3$} هغه
    فرعي-اشتقاقونه دي چې په~\LR{$\delta$} کښې د وروستي استنتاج پر
    مقدمه يا مقدماتو پاې ته رسېږي؛ \LR{$!A$} په~\LR{$\delta$} کښې د وروستي
    استنتاج پايله ده؛
    \LR{$n$} د وروستي استنتاج د ختمولو نښه ده (که استنتاج هېڅ فرض نۀ
    ختموي، \LR{$0$} ده)؛ او \LR{$k$} د لاندې جدول له مخې د وروستي استنتاج د
    کارول شوې قاعدې له مخې ټاکل کېږي۔
\olpseganchor{OLP-0287-B006}{end}

\olpseganchor{OLP-0287-B007}{start}
    \begin{tabular}{lccccccc}
      \text{\RL{قاعده:}} & \Intro{\land} & \Elim{\land} & \Intro{\lor} & \Elim{\lor} \\
      \LR{$k$}: & 1 & 2 & 3 & 4  \\[1.5ex]
      \text{\RL{قاعده:}} &  \Intro{\lif} & \Elim{\lif} & \Intro{\lnot} & \Elim{\lnot} \\
      \LR{$k$}: & 5 & 6 & 7 & 8 \\[1.5ex]
      \text{\RL{قاعده:}}  &  \FalseInt & \FalseCl & \Intro{\lforall} & \Elim{\lforall} \\
      \LR{$k$}: & 9 & 10 & 11 & 12 \\[1.5ex]
      \text{\RL{قاعده:}} & \Intro{\lexists} & \Elim{\lexists} & \Intro{\eq} & \Elim{\eq} \\
      \LR{$k$}: & 13 & 14 & 15 & 16
    \end{tabular}
  \end{enumerate}
\end{defn}
\olpseganchor{OLP-0287-B007}{end}

\olpseganchor{OLP-0287-B008}{start}
\begin{ex}
  دا ډېر ساده اشتقاق وګورئ:
  \begin{prooftree}
    \AxiomC{\LR{$\Discharge{!A \land !B}{1}$}}
    \RightLabel{\Elim{\land}}
    \UnaryInfC{\LR{$!A$}}
    \DischargeRule{\Intro{\lif}}{1}
    \UnaryInfC{\LR{$(!A \land !B) \lif !A$}}
  \end{prooftree}
  د فرض ګوډل شمېره به
  \LR{$d_0=\tuple{0,\Gn{!A \land !B},1}$} وي۔ د هغه اشتقاق ګوډل
  شمېره چې د~\LR{$\Elim{\land}$} پر پايله ختمېږي،
  \LR{$d_1=\tuple{1,d_0,\Gn{!A},0,2}$} ده (\LR{$1$} ځکه چې
  \LR{$\Elim{\land}$} يوه مقدمه لري؛ بيا د پايلې~\LR{$!A$} ګوډل شمېره؛ \LR{$0$}
  ځکه چې هېڅ فرض نۀ ختمېږي؛ او \LR{$2$} هغه عدد دے چې \LR{$\Elim{\land}$}
  کوډوي)۔ بيا د ټول اشتقاق ګوډل شمېره
  \LR{$\tuple{1,d_1,\Gn{((!A \land !B) \lif !A)},1,5}$} ده؛ يعنې
  \[
  \tuple{1, \tuple{1, \tuple{0, \Gn{(!A \land !B)}, 1}, \Gn{!A}, 0, 2},
    \Gn{((!A \land !B) \lif !A)}, 1, 5}.
  \]
\end{ex}
\olpseganchor{OLP-0287-B008}{end}

\olpseganchor{OLP-0287-B009}{start}
\begin{explain}
چې د اشتقاقونه پر يو تمثيل هوکړې ته ورسېدو، دا هم بايد وښيو
چې د داسې اشتقاقونه ګوډل شمېرې په بنسټيز بازګشتي ډول اداره
کولے شو، او د هغو اړين خاصيتونه او اړيکې هم بيانولے شو۔ ځينې عمليات
ساده دي: د بېلګې په توګه، د اشتقاق له ګوډل شمېرې~\LR{$d$} څخه
\LR{$\fn{EndFmla}(d)=(d)_{(d)_0+1}$} د هغه د پاې-فارمول ګوډل
شمېره راکوي، \LR{$\fn{DischargeLabel}(d)=(d)_{(d)_0+2}$} د ختمولو نښه
راکوي، او \LR{$\fn{LastRule}(d)=(d)_{(d)_0+3}$} هغه عدد راکوي چې د وروستي
استنتاج ډول ښيي۔ ځينې نور عمليات ډېر ستونزمن دي۔ لږ تر لږه د هغو د
ترسره کولو لار به راونغاړو۔ موخه دا ده چې وښيو اړيکه «\LR{$\delta$}
له~\LR{$\Gamma$} څخه د~\LR{$!A$} يو اشتقاق دے» د~\LR{$\delta$} او~\LR{$!A$}
د ګوډل شمېرو بنسټيزه بازګشتي اړيکه ده۔
\end{explain}
\olpseganchor{OLP-0287-B009}{end}

\olpseganchor{OLP-0287-B010}{start}
\begin{prop}
  لاندې اړيکې بنسټيزې بازګشتي دي:
  \begin{enumerate}
  \item \LR{$!A$} په~\LR{$\delta$} کښې د نښې~\LR{$n$} لرونکي فرض په توګه راځي۔
  \item په~\LR{$\delta$} کښې د نښې~\LR{$n$} لرونکي ټول فرضونه د~\LR{$!A$} په بڼه
    دي (يعنې په~\LR{$\delta$} کښې د نښې~\LR{$n$} په کارولو فرض~\LR{$!A$}
    ايستل کولے شو)۔
  \end{enumerate}
\end{prop}
\olpseganchor{OLP-0287-B010}{end}

\olpseganchor{OLP-0287-B011}{start}
\begin{proof}
  بايد وښيو چې د فارمولونه د ګوډل شمېرو او د اشتقاقونه د
  ګوډل شمېرو ترمنځ اړوندې اړيکې بنسټيزې بازګشتي دي۔
  \begin{enumerate}
  \item غواړو وښيو چې \LR{$\fn{Assum}(x,d,n)$} بنسټيزه بازګشتي ده؛ دا
    هله صدق کوي چې \LR{$x$} د ګوډل شمېرې~\LR{$d$} لرونکي اشتقاق د
    نښې~\LR{$n$} لرونکي فرض ګوډل شمېره وي۔ دا هله سم دي چې د ګوډل شمېرې
    \LR{$\tuple{0,x,n}$} لرونکے اشتقاق د~\LR{$d$} يو
    فرعي-اشتقاق وي۔ پام وکړئ، زموږ د اشتقاقونه د کوډولو
    لار د هغو ونو د کوډولو ځانګړے حالت دے چې په
    \olref[cmp][rec][tre]{sec} کښې معرفي شوې؛ نو بنسټيزه بازګشتي
    تابع \LR{$\fn{SubtreeSeq}(d)$} د~\LR{$d$} د ټولو فرعي-اشتقاقونه د
    ګوډل شمېرو لړۍ (چې اوږدوالے ئې تر ټولو زيات~\LR{$d$} دے) راکوي۔ نو
    داسې تعريفولے شو:
    \[
    \fn{Assum}(x, d, n) \defiff \bexists{i<d}{(\fn{SubtreeSeq}(d))_i =
      \tuple{0, x, n}}.
    \]
  \item غواړو وښيو چې \LR{$\fn{Discharge}(x,d,n)$} بنسټيزه بازګشتي ده؛
    دا هله صدق کوي چې د ګوډل شمېرې~\LR{$d$} لرونکي اشتقاق ټول
    نښه~\LR{$n$} لرونکي فرضونه د ګوډل شمېرې~\LR{$x$} لرونکي فارمول وي۔
    خو دا اړيکه هله صدق کوي چې
    \LR{$\bforall{y<d}{(\fn{Assum}(y,d,n)\lif y=x)}$}۔
  \end{enumerate}
\end{proof}
\olpseganchor{OLP-0287-B011}{end}

\olpseganchor{OLP-0287-B012}{start}
\begin{prop}
  \ollabel{prop:followsby}
  خاصيت \LR{$\fn{Correct}(d)$}، چې هله صدق کوي چې د ګوډل شمېرې~\LR{$d$}
  لرونکي اشتقاق~\LR{$\delta$} وروستے استنتاج سم وي، بنسټيز
  بازګشتي دے۔
\end{prop}
\olpseganchor{OLP-0287-B012}{end}

\olpseganchor{OLP-0287-B013}{start}
\begin{proof}
  دلته د استنتاج د هرې قاعدې~\LR{$R$} دپاره بايد وښيو چې اړيکه
  \LR{$\fn{FollowsBy}_R(d)$} بنسټيزه بازګشتي ده؛
  \LR{$\fn{FollowsBy}_R(d)$} هله صدق کوي چې \LR{$d$} د
  اشتقاق~\LR{$\delta$} ګوډل شمېره وي او د~\LR{$\delta$}
  پاې-فارمول د~\LR{$R$} په سم پلي کولو د~\LR{$\delta$} د سمدستي
  فرعي-اشتقاقونه له پاې-فارمولونو څخه راشي۔
\olpseganchor{OLP-0287-B013}{end}

\olpseganchor{OLP-0287-B014}{start}
  يو ساده حالت د \Intro{\land} قاعده ده۔ که \LR{$\delta$} په يوه سم
  \LR{$\Intro{\land}$} استنتاج ختم شي، داسې ښکاري:
  \begin{prooftree}
    \AxiomC{}
    \RightLabel{\LR{$\delta_1$}}
    \DeduceC{\LR{$!A$}}
\olpseganchor{OLP-0287-B014}{end}

\olpseganchor{OLP-0287-B015}{start}
    \AxiomC{}
    \RightLabel{\LR{$\delta_2$}}
    \DeduceC{\LR{$!B$}}
\olpseganchor{OLP-0287-B015}{end}

\olpseganchor{OLP-0287-B016}{start}
    \RightLabel{\Intro\land}
    \BinaryInfC{\LR{$!A \land !B$}}
  \end{prooftree}
  بيا د~\LR{$\delta$} ګوډل شمېره~\LR{$d$} دا ده:
  \LR{$\tuple{2,d_1,d_2,\Gn{(!A \land !B)},0,k}$}، چې
  \LR{$\fn{EndFmla}(d_1)=\Gn{!A}$}، \LR{$\fn{EndFmla}(d_2)=\Gn{!B}$}،
  \LR{$n=0$} او \LR{$k=1$} دي۔ نو \LR{$\fn{FollowsBy}_{\Intro\land}(d)$} داسې
  تعريفولے شو:
  \begin{multline*}
    (d)_0 = 2 \land \fn{DischargeLabel}(d) = 0 \land \fn{LastRule}(d) = 1 \land {}\\
    \fn{EndFmla}(d) = {}\\
    \Gn{(} \concat \fn{EndFmla}((d)_1) \concat \Gn{\land}
    \concat \fn{EndFmla}((d)_2) \concat \Gn{)}.
  \end{multline*}
\olpseganchor{OLP-0287-B016}{end}

\olpseganchor{OLP-0287-B017}{start}
  بله ساده بېلګه د \LR{$\Intro\eq$} قاعده ده۔ دا هېڅ مقدمه نۀ لري، نو د
  فرضونو په شان \LR{$(d)_0=0$} دے۔ دا د ختمولو نښه هم نۀ لري؛ يعنې \LR{$n=0$}
  دے۔ خو \LR{$!A$} بايد د کوم تړلي ترم~\LR{$t$} دپاره د~\LR{$\eq[t][t]$} په بڼه
  وي۔ يو بنسټيز بازګشتي تعريف دا دے:
  \begin{multline*}
    (d)_0 = 0 \land \fn{DischargeLabel}(d) = 0 \land {}\\
    \bexists{t<d}{(\fn{ClTerm}(t) \land
      \fn{EndFmla}(d) = {}\\
      \Gn{{\eq}(} \concat t \concat \Gn{,} \concat t \concat \Gn{)})}.
  \end{multline*}
\olpseganchor{OLP-0287-B017}{end}

\olpseganchor{OLP-0287-B018}{start}
  د يوې لا پېچلې بېلګې دپاره \LR{$\fn{FollowsBy}_{\Intro{\lif}}(d)$} هله
  صدق کوي چې د~\LR{$\delta$} پاې-فارمول د~\LR{$(!A\lif !B)$} په بڼه وي،
  د~\LR{$\delta_1$} پاې-فارمول هم~\LR{$!B$} وي، او په~\LR{$\delta$} کښې هر
  نښه~\LR{$n$} لرونکے فرض د~\LR{$!A$} په بڼه وي۔ دا په بنسټيز بازګشتي ډول
  داسې بيانولے شو:
  \begin{multline*}
    (d)_0 = 1 \land {}\\
    \bexists{a<d}{(\fn{Discharge}(a, (d)_1, \fn{DischargeLabel}(d)) \land {}}\\
      \fn{EndFmla}(d) = (\Gn{(} \concat a \concat \Gn{\lif}
      \concat \fn{EndFmla}((d)_1) \concat \Gn{)}))
  \end{multline*}
  (دلته \LR{$a$} د~\LR{$!A$} ګوډل شمېره وګڼئ)۔
\olpseganchor{OLP-0287-B018}{end}

\olpseganchor{OLP-0287-B019}{start}
  د بلې بېلګې دپاره \Intro{\lexists} وګورئ۔ دلته په~\LR{$\delta$} کښې
  وروستے استنتاج هله سم دے چې داسې فارمول~\LR{$!A$}، تړلے ترم~\LR{$t$}
  او متغير~\LR{$x$} شته چې \LR{$\Subst{!A}{t}{x}$} د
  اشتقاق~\LR{$\delta_1$} پاې-فارمول وي او
  \LR{$\lexists[x][!A]$} د وروستي استنتاج پايله وي۔ نو
  \LR{$\fn{FollowsBy}_{\Intro{\lexists}}(d)$} هله صدق کوي چې
  \begin{multline*}
    (d)_0 = 1 \land \fn{DischargeLabel}(d) = 0 \land {} \\
    \bexists{a < d}{\bexists{x<d}{\bexists{t<d}{
          (\fn{ClTerm}(t) \land \fn{Var}(x)  \land {}}}}\\
    \fn{Subst}(a,t,x) = \fn{EndFmla}((d)_1) \land
    \fn{EndFmla}(d) = (\Gn{\lexists} \concat x \concat a)).
  \end{multline*}
\olpseganchor{OLP-0287-B019}{end}

\olpseganchor{OLP-0287-B020}{start}
  بيا \LR{$\fn{Correct}(d)$} داسې تعريفوو:
  \begin{multline*}
    \fn{Sent}(\fn{EndFmla}(d)) \land {}\\
    [(\fn{LastRule}(d) = 1 \land
    \fn{FollowsBy}_{\Intro\land}(d)) \lor \dots \lor {}\\
    (\fn{LastRule}(d) = 16 \land \fn{FollowsBy}_{\Elim\eq}(d)) \lor {}\\
    \bexists{n<d}{\bexists{x<d}{(d = \tuple{0, x, n})}}].
  \end{multline*}
  لومړۍ کرښه باوري کوي چې د~\LR{$d$} پاې-فارمول يوه جمله ده۔ وروستۍ
  کرښه هغه حالت رانغاړي چې \LR{$d$} يوازې يو فرض وي۔
  \textbf{د ژباړې د نسخې سمون (OLCMP-079):} سرچينې د قاعدو او فرض
  د فصل د ګروپ قوسونه پرېښي وو، نو د جملې شرط يوازې پر لومړي
  بديل پلي کېده۔ دلته ټول بديلونه د هماغه \LR{$\fn{Sent}$} شرط لاندې
  ګروپ شوي دي، لکه چې ورپسې تشريح ئې غواړي۔
\end{proof}
\olpseganchor{OLP-0287-B020}{end}

\olpseganchor{OLP-0287-B021}{start}
\begin{prob}
  لاندې خاصيتونه د \olref[inc][art][pnd]{prop:followsby} په څېر
  تعريف کړئ:
  \begin{enumerate}
  \item \LR{$\fn{FollowsBy}_{\Elim{\lif}}(d)$}؛
  \item \LR{$\fn{FollowsBy}_{\Elim{\eq}}(d)$}؛
  \item \LR{$\fn{FollowsBy}_{\Elim{\lor}}(d)$}؛
  \item \LR{$\fn{FollowsBy}_{\Intro{\lforall}}(d)$}۔
  \end{enumerate}
  د وروستي خاصيت دپاره دا هم بايد وښيئ چې په بنسټيز بازګشتي ډول
  ازمويلے شئ چې د ګوډل شمېرې~\LR{$d$} لرونکي اشتقاق وروستے
  استنتاج د ځانګړي متغير شرط پوره کوي؛ يعنې د
  \LR{$\Intro{\lforall}$} استنتاج ځانګړے متغير~\LR{$a$} نۀ د~\LR{$d$} په
  پاې-فارمول کښې راځي او نۀ د~\LR{$d$} په يوه پرانيستي فرض کښې۔ د
  دې دپاره د \olref[inc][art][pnd]{prop:openassum} بنسټيز بازګشتي
  پريديکات~\LR{$\fn{OpenAssum}$} کارولے شئ۔
\end{prob}
\olpseganchor{OLP-0287-B021}{end}

\olpseganchor{OLP-0287-B022}{start}
\begin{prop}
  \ollabel{prop:deriv}
  اړيکه \LR{$\fn{Deriv}(d)$}، چې هله صدق کوي چې \LR{$d$} د يوه سم
  اشتقاق~\LR{$\delta$} ګوډل شمېره وي، بنسټيزه بازګشتي ده۔
\end{prop}
\olpseganchor{OLP-0287-B022}{end}

\olpseganchor{OLP-0287-B023}{start}
\begin{proof}
  اشتقاق~\LR{$\delta$} هله سم دے چې هر استنتاج ئې د يوې قاعدې
  سم پلي کول وي؛ يعنې د هغه هر فرعي-اشتقاقونه په سم استنتاج پاې
  ته ورسېږي۔ نو \LR{$\fn{Deriv}(d)$} هله صدق کوي چې
  \[
  \bforall{i<\len{\fn{SubtreeSeq}(d)}}{\fn{Correct}((\fn{SubtreeSeq}(d))_i)}
  \]
\end{proof}
\olpseganchor{OLP-0287-B023}{end}

\olpseganchor{OLP-0287-B024}{start}
\begin{prop}
  \ollabel{prop:openassum} اړيکه \LR{$\fn{OpenAssum}(z,d)$}، چې هله صدق
  کوي چې \LR{$z$} د ګوډل شمېرې~\LR{$d$} لرونکي اشتقاق~\LR{$\delta$} د يوه
  نۀ ايستل شوې فرض~\LR{$!A$} ګوډل شمېره وي، بنسټيزه بازګشتي ده۔
\end{prop}
\olpseganchor{OLP-0287-B024}{end}

\olpseganchor{OLP-0287-B025}{start}
\begin{proof}
  د يوه فرض پېښه هله ايستل ده چې د نښې~\LR{$n$} سره د~\LR{$\delta$}
  په داسې فرعي-اشتقاق کښې راشي چې د ختمولو په نښه~\LR{$n$} لرونکې
  قاعده ختمېږي۔ نو \LR{$!A$} هله د~\LR{$\delta$} يو نۀ ايستل شوې فرض دے
  چې لږ تر لږه يوه پېښه ئې په~\LR{$\delta$} کښې ايستل نۀ وي۔
  بايد پام وکړو: \LR{$\delta$} کېدے شي د~\LR{$!A$} دواړه ايستل او
  نۀ ايستل شوې پېښې ولري۔
\olpseganchor{OLP-0287-B025}{end}

\olpseganchor{OLP-0287-B026}{start}
  داسې لړۍ \LR{$\delta_0$}، \dots، \LR{$\delta_k$} وګورئ چې
  \LR{$\delta_0=\delta$}، \LR{$\delta_k$} فرض \LR{$\Discharge{!A}{n}$} (د کوم~\LR{$n$}
  دپاره) وي، او \LR{$\delta_{i+1}$} د~\LR{$\delta_i$} سمدستي
  فرعي-اشتقاق وي۔ که داسې لړۍ موجوده وي چې په هغې کښې هېڅ
  \LR{$\delta_i$} د ختمولو نښه~\LR{$n$} لرونکي استنتاج باندې نۀ ختمېږي، نو
  \LR{$!A$} د~\LR{$\delta$} يو نۀ ايستل شوې فرض دے۔
\olpseganchor{OLP-0287-B026}{end}

\olpseganchor{OLP-0287-B027}{start}
  بنسټيزه بازګشتي تابع \LR{$\fn{SubtreeSeq}(d)$} د~\LR{$\delta$} د ټولو
  فرعي-اشتقاقونه د ګوډل شمېرو لړۍ راکوي۔ د~\LR{$\delta$} د
  فرعي-اشتقاقونه د ګوډل شمېرو هره لړۍ د هغې يوه فرعي لړۍ ده۔
  د فرعي لړۍ کېدل يوه بنسټيزه بازګشتي اړيکه ده:
  \LR{$\fn{Subseq}(s,s')$} هله صدق کوي چې
  \LR{$\bforall{i<\len{s}}{\lexists[j<\len{s'}][(s)_i=(s')_j]}$}۔ د
  سمدستي فرعي-اشتقاق کېدل هم داسې دي:
  \LR{$\fn{Subderiv}(d,d')$} هله صدق کوي چې
  \LR{$\bexists{j<(d')_0}{d=(d')_{j+1}}$}۔
  \textbf{د ژباړې د نسخې سمون (OLCMP-080):} د سرچينې ازموينې
  \LR{$(d')_j$} کاراوه، خو په همدې تعريف شوي کوډ کښې \LR{$(d')_0$} د سمدستي
  فرعي-اشتقاقونو شمېر دے او د هغو کوډونه له شاخص~\LR{$1$} څخه
  پيلېږي۔ دلته \LR{$j+1$} کارول شوے، څو لومړے عددي سرليک ونه نيسي او
  وروستے فرعي-اشتقاق پرې نۀ ږدي۔
  نو \LR{$\fn{OpenAssum}(z,d)$} داسې تعريفولے شو:
  \begin{multline*}
    \bexists{s<\fn{SubtreeSeq}(d)}{(\fn{Subseq}(s, \fn{SubtreeSeq}(d))
      \land (s)_0 = d \land {}} \\
    \bexists{n<d}{((s)_{\len{s} \tsub 1} = \tuple{0, z, n} \land {}}\\
      \bforall{i<(\len{s} \tsub 1)}{(\fn{Subderiv}((s)_{i+1}, (s)_i) \land {}}\\
      \fn{DischargeLabel}((s)_i) \neq n))).
  \end{multline*}
\end{proof}
\olpseganchor{OLP-0287-B027}{end}

\olpseganchor{OLP-0287-B028}{start}
\begin{prop}
  \ollabel{prop:prf-prim-rec}
  فرض کړئ \LR{$\Gamma$} د تړلي فارمولونه يو بنسټيز بازګشتي سټ دے۔ بيا
  اړيکه \LR{$\Prf[\Gamma](x,y)$}، چې وايي «\LR{$x$} له~\LR{$\Gamma$} کښې
  نۀ ايستل شوې فرضونو څخه د~\LR{$!A$} د اشتقاق~\LR{$\delta$} کوډ
  دے او \LR{$y$} د~\LR{$!A$} ګوډل شمېره ده»، بنسټيزه بازګشتي ده۔
\end{prop}
\olpseganchor{OLP-0287-B028}{end}

\olpseganchor{OLP-0287-B029}{start}
\begin{proof}
  فرض کړئ «\LR{$y\in\Gamma$}» د بنسټيز بازګشتي
  پريديکات~\LR{$R_\Gamma(y)$} په واسطه ورکول کېږي۔ بايد وښيو چې
  \LR{$\Prf[\Gamma](x,y)$} بنسټيزه بازګشتي ده؛ دا اړيکه هله صدق کوي چې
  \LR{$y$} د يوې جملې~\LR{$!A$} ګوډل شمېره وي، او \LR{$x$} د طبيعي استنتاج د داسې
  اشتقاق کوډ وي چې پاې-فارمول ئې~\LR{$!A$} او ټول
  نۀ ايستل شوي فرضونه ئې په~\LR{$\Gamma$} کښې وي۔
\olpseganchor{OLP-0287-B029}{end}

\olpseganchor{OLP-0287-B030}{start}
  د \olref{prop:deriv} له مخې خاصيت \LR{$\fn{Deriv}(x)$}، چې هله صدق کوي
  چې \LR{$x$} په طبيعي استنتاج کښې د يوه سم اشتقاق~\LR{$\delta$} ګوډل
  شمېره وي، بنسټيز بازګشتي دے۔ نو \LR{$\Prf[\Gamma](x,y)$} داسې تعريفولے
  شو:
  \begin{align*}
    \Prf[\Gamma](x, y) \defiff {}
    & \fn{Deriv}(x) \land \fn{EndFmla}(x) = y \land {} \\
    & \bforall{z < x}{(\fn{OpenAssum}(z, x) \lif R_\Gamma(z))}.
  \end{align*}
\end{proof}
\olpseganchor{OLP-0287-B030}{end}

\olpseganchor{OLP-0288-B004}{start}
\olfileid{inc}{art}{pax}
\olsection{بديهي اشتقاقونه}
\olpseganchor{OLP-0288-B004}{end}

\olpseganchor{OLP-0288-B005}{start}
\begin{explain}
د بديهي اشتقاقونه د حسابي کولو دپاره بايد اشتقاقونه د
عددونو په توګه تمثيل کړو۔ ځکه چې اشتقاقونه يوازې د
فارمولونه لړۍ دي، نو څرګنده لار دا ده چې هر اشتقاق د هغه
د فارمولونه د کوډونو د لړۍ په کوډ سره کوډ کړو۔
\end{explain}
\olpseganchor{OLP-0288-B005}{end}

\olpseganchor{OLP-0288-B006}{start}
\begin{defn}
که \LR{$\delta$} داسې بديهي اشتقاق وي چې له فارمولونه
\LR{$!A_1$}، \dots،~\LR{$!A_n$} څخه جوړ وي، نو \LR{$\Gn{\delta}$} دا دے:
\[
\tuple{\Gn{!A_1}, \dots, \Gn{!A_n}}.
\]
\end{defn}
\olpseganchor{OLP-0288-B006}{end}

\olpseganchor{OLP-0288-B007}{start}
\begin{ex}
  دا ډېر ساده اشتقاق وګورئ:
  \begin{derivation}
  1. & \LR{$!B \lif (!B \lor !A)$} \\
  2. & \LR{$(!B \lif (!B \lor !A)) \lif (!A  \lif (!B \lif (!B \lor !A)))$}\\
  3. & \LR{$!A  \lif (!B \lif (!B \lor !A))$}
  \end{derivation}
  د دې اشتقاق ګوډل شمېره به دا وي:
  \begin{align*}
  \openTuple\,
    & \Gn{!B \lif (!B \lor !A)}, \\
    &\Gn{(!B \lif (!B \lor !A)) \lif (!A  \lif (!B \lif (!B \lor !A)))},\\
    & \Gn{!A  \lif (!B \lif (!B \lor !A))} \,\closeTuple.
  \end{align*}
\end{ex}
\olpseganchor{OLP-0288-B007}{end}

\olpseganchor{OLP-0288-B008}{start}
\begin{explain}
چې د اشتقاقونه پر يو تمثيل هوکړې ته ورسېدو، دا هم بايد وښيو
چې داسې اشتقاقونه په بنسټيز بازګشتي ډول اداره کولے شو، او د
هغو اړين خاصيتونه او اړيکې هم په همدې ډول بيانولے شو۔ ځينې عمليات
ساده دي: د بېلګې په توګه، د اشتقاق له ګوډل شمېرې~\LR{$d$} څخه
\LR{$(d)_{\len{d}-1}$} د هغه د پاې-فارمول ګوډل شمېره راکوي۔ ځينې
نور عمليات ډېر ستونزمن دي۔ لږ تر لږه د هغو د ترسره کولو لار به
راونغاړو۔ موخه دا ده چې وښيو اړيکه «\LR{$\delta$} له~\LR{$\Gamma$} څخه د~\LR{$!A$}
يو اشتقاق دے» د~\LR{$\delta$} او~\LR{$!A$} په ګوډل شمېرو کښې
بنسټيزه بازګشتي ده۔
\end{explain}
\olpseganchor{OLP-0288-B008}{end}

\olpseganchor{OLP-0288-B009}{start}
\begin{prop}
\ollabel{prop:followsby}
لاندې اړيکې بنسټيزې بازګشتي دي:
\begin{enumerate}
\item \LR{$!A$} يو بديهي اصل دے۔
\item په~\LR{$\delta$} کښې \LR{$i$}-مه کرښه په مودوس پوننس توجيه شوې ده۔
\item په~\LR{$\delta$} کښې \LR{$i$}-مه کرښه په \QR{} توجيه شوې ده۔
\item \LR{$\delta$} يو سم اشتقاق دے۔
\end{enumerate}
\end{prop}
\olpseganchor{OLP-0288-B009}{end}

\olpseganchor{OLP-0288-B010}{start}
\begin{proof}
بايد وښيو چې د فارمولونه د ګوډل شمېرو او د اشتقاقونه د
ګوډل شمېرو ترمنځ اړوندې اړيکې بنسټيزې بازګشتي دي۔
\begin{enumerate}
\item موږ ته د بديهي شېماوو يو ټاکلے لست راکړل شوے، او \LR{$!A$} هله يو
  بديهي اصل دے چې د دې شېماوو څخه د يوې په ورکړې بڼه وي۔ ځکه چې د
  شېماوو لست متناهي دے، بس ده چې د هرې بديهي شېما دپاره په بنسټيز
  بازګشتي ډول وازميوو چې \LR{$!A$} د هغې په بڼه دے که نۀ۔ د بېلګې په
  توګه دا بديهي شېما وګورئ:
  \[
  !B \lif (!C \lif !B).
  \]
  \LR{$!A$} هله د دې بديهي شېما يوه بېلګه ده چې داسې فارمولونه \LR{$!B$}
  او \LR{$!C$} شته چې \LR{$!A$} د «\LR{$($}»، بيا \LR{$!B$}، بيا «\LR{$\lif$}»، بيا «\LR{$($}»،
  بيا \LR{$!C$}، بيا «\LR{$\lif$}»، بيا \LR{$!B$}، او بيا «\LR{$))$}» په نښلولو ترلاسه
  شي۔ د~\LR{$!A$} د ګوډل شمېرې~\LR{$n$} اړوند خاصيت ازمويلے شو، ځکه د لړيو
  نښلول بنسټيز بازګشتي دي، او د~\LR{$!B$} او~\LR{$!C$} ګوډل شمېرې بايد د~\LR{$!A$}
  له ګوډل شمېرې څخه وړې وي؛ ځکه کله چې اړيکه صدق کوي، \LR{$!B$} او \LR{$!C$}
  دواړه د~\LR{$!A$} فرعي-فارمولونه دي۔ نو داسې تعريفولے شو:
  \begin{multline*}
  \fn{IsAx}_{!B \lif (!C \lif !B)}(n) \defiff \bexists{b< n}{
    \bexists{c < n}{(\fn{Sent}(b) \land \fn{Sent}(c) \land {}}}\\ n =
  \Gn{(} \concat b \concat \Gn{\lif} \concat \Gn{(} \concat c \concat
  \Gn{\lif} \concat b \concat \Gn{))}).
  \end{multline*}
  که د هرې بديهي شېما دپاره داسې تعريف ولرو، د هغو د فصل خاصيت
  \LR{$\fn{IsAx}(n)$} تعريفوي: «\LR{$n$} د يوه بديهي اصل ګوډل شمېره ده»۔
\item په~\LR{$\delta$} کښې \LR{$i$}-مه کرښه هله په مودوس پوننس توجيه شوې ده
  چې داسې کرښې \LR{$j$} او \LR{$k<i$} شته چې د~\LR{$j$} پر کرښه تړلے فارمول کوم
  فارمول~\LR{$!A$}، د~\LR{$k$} پر کرښه جمله \LR{$!A\lif !B$}، او د~\LR{$i$} پر کرښه
  جمله~\LR{$!B$} وي۔
  \begin{multline*}
    \fn{MP}(d, i) \defiff \bexists{j < i}{\bexists{k < i}{}}\\
    (d)_k = \Gn{(} \concat (d)_j
      \concat \Gn{\lif} \concat (d)_i \concat \Gn{)}
  \end{multline*}
  ځکه محدود کمیت ايښودنه، نښلول او~\LR{$=$} بنسټيز بازګشتي دي، نو دا
  يوه بنسټيزه بازګشتي اړيکه تعريفوي۔
\item په~\LR{$\delta$} کښې يوه کرښه هله په \QR{} توجيه شوې ده چې د
  \LR{$!B\lif\lforall[x][!A(x)]$} په بڼه وي، يوه مخکښې کرښه د کوم
  ثابت سمبول~\LR{$c$} دپاره \LR{$!B\lif !A(c)$} وي، او \LR{$c$} په~\LR{$!B$} کښې نۀ
  راځي۔ دا هله کېږي چې:
  \begin{enumerate}
  \item داسې تړلے فارمول~\LR{$!B$} وي؛ او
  \item داسې فارمول~\LR{$!A(x)$} وي چې يوازې يو ازاد متغير~\LR{$x$}
    ولري؛ داسې چې
  \item کرښه~\LR{$i$}، \LR{$!B\lif\lforall[x][!A(x)]$} ولري؛
  \item کومه کرښه \LR{$j<i$} د کوم ثابت~\LR{$c$} دپاره
    \LR{$!B\lif\Subst{!A}{c}{x}$} ولري؛ او
  \item \LR{$c$} په~\LR{$!B$} کښې نۀ راځي۔
  \end{enumerate}
  دا ټول په بنسټيز بازګشتي ډول ازمويلے شو، ځکه د~\LR{$!B$}، \LR{$!A(x)$} او
  \LR{$x$} ګوډل شمېرې د~\LR{$i$} پر کرښه د فارمول له ګوډل شمېرې څخه وړې دي،
  او د~\LR{$c$} ګوډل شمېره د~\LR{$j$} پر کرښه د فارمول له ګوډل شمېرې څخه
  وړه ده:
  \begin{multline*}
    \fn{QR}_1(d, i) \defiff \bexists{j<i}{}\bexists{b < (d)_i}{\bexists{x <
        (d)_i}{\bexists{a < (d)_i}{\bexists{c < (d)_j}{(}}}}
    \\ \fn{Var}(x) \land \fn{Const}(c) \land {} \\
    (d)_i = \Gn{(} \concat
    b \concat \Gn{\lif} \concat \Gn{\lforall} \concat x \concat a
    \concat \Gn{)} \land {}\\
    (d)_j = \Gn{(} \concat b \concat
    \Gn{\lif} \concat \fn{Subst}(a,c,x) \concat \Gn{)} \land {}\\ \fn{Sent}(b)
    \land \fn{Sent}(\fn{Subst}(a,c,x)) \land {} \bforall{k <
      \len{b}}{(b)_k \neq (c)_0})
  \end{multline*}
  دلته فرض کوو چې \LR{$c$} او \LR{$x$} په ترتيب د ثابت او متغير ګوډل شمېرې
  دي چې د ترمونو په توګه ورته کتل شوي (يعنې د هغو سمبول-کوډونه نۀ
  دي)۔ دا چې \LR{$x$} د~\LR{$!A(x)$} يوازينے ازاد متغير دے، په دې ازميوو چې
  \LR{$\Subst{!A(x)}{c}{x}$} يوه تړلے فارمول ده؛ او دا چې \LR{$c$} په~\LR{$!B$}
  کښې نۀ راځي، په دې شرط باوروو چې د~\LR{$!B$} هر سمبول له~\LR{$c$} څخه بېل
  وي۔
  \textbf{د ژباړې د نسخې سمون (OLCMP-081):} سرچينې د مخکښې کرښې
  شاخص~\LR{$j$} کاراوه خو \LR{$j<i$} ئې کمیت ايښودلے نۀ و، او ورپسې تشريح کښې
  ئې په تېروتنه د~\LR{$a$} پر ځاے د ثابت~\LR{$c$} حد ياد کړے و۔ دلته ورک
  محدود کمیت ورزيات او تشريح له ښودل شوي شرط سره برابره شوې ده۔
\olpseganchor{OLP-0288-B010}{end}

\olpseganchor{OLP-0288-B011}{start}
  د \QR{} بله بڼه د تمرین دپاره پرېږدو۔
\item \LR{$d$} هله د يوه سم اشتقاق ګوډل شمېره ده چې هره کرښه ئې
  بديهي اصل وي، يا په مودوس پوننس يا~\QR{} توجيه شوې وي۔ نو:
  \[
  \fn{Deriv}(d) \defiff \bforall{i < \len{d}}{(\fn{IsAx}((d)_i) \lor
    \fn{MP}(d,i) \lor \fn{QR}(d, i))}
  \]
\end{enumerate}
\end{proof}
\olpseganchor{OLP-0288-B011}{end}

\olpseganchor{OLP-0288-B012}{start}
\begin{prob}
لاندې اړيکې د \olref[inc][art][pax]{prop:followsby} په څېر تعريف
کړئ:
\begin{enumerate}
\item \LR{$\fn{IsAx}_{!A \lif (!B \lif (!A \land !B))}(n)$}؛
\item \LR{$\fn{IsAx}_{\lforall[x][!A(x)] \lif !A(t)}(n)$}؛
\item \LR{$\fn{QR}_{2}(d, i)$} (د \QR{} د بلې بڼې دپاره)۔
\end{enumerate}
\end{prob}
\olpseganchor{OLP-0288-B012}{end}

\olpseganchor{OLP-0288-B013}{start}
\begin{prop}
فرض کړئ \LR{$\Gamma$} د تړلي فارمولونه يو بنسټيز بازګشتي سټ دے۔ بيا
اړيکه \LR{$\Prf[\Gamma](x,y)$}، چې وايي «\LR{$x$} له~\LR{$\Gamma$} څخه د~\LR{$!A$} د
اشتقاق~\LR{$\delta$} کوډ دے او \LR{$y$} د~\LR{$!A$} ګوډل شمېره ده»،
بنسټيزه بازګشتي ده۔
\end{prop}
\olpseganchor{OLP-0288-B013}{end}

\olpseganchor{OLP-0288-B014}{start}
\begin{proof}
فرض کړئ «\LR{$y\in\Gamma$}» د بنسټيز بازګشتي
پريديکات~\LR{$R_\Gamma(y)$} په واسطه ورکول کېږي۔ بايد وښيو چې اړيکه
\LR{$\Prf[\Gamma](x,y)$} بنسټيزه بازګشتي ده؛
\LR{$\Prf[\Gamma](x,y)$} هله صدق کوي چې
\LR{$y$} د تړلے فارمول~\LR{$!A$} ګوډل شمېره وي او \LR{$x$} له~\LR{$\Gamma$} څخه
د~\LR{$!A$} د اشتقاق کوډ وي۔
\olpseganchor{OLP-0288-B014}{end}

\olpseganchor{OLP-0288-B015}{start}
د تېرې قضيې له مخې خاصيت \LR{$\fn{Deriv}(x)$}، چې هله صدق کوي چې \LR{$x$} د
يوه سم اشتقاق~\LR{$\delta$} کوډ وي، بنسټيز بازګشتي دے۔ خو هغه
تعريف سټ~\LR{$\Gamma$} د اشتقاق د کرښو د توجيه د يوې اضافي لارې
په توګه په پام کښې نۀ و نيولے۔ زموږ بنسټيز بازګشتي ازموينې چې يوه
کرښه په \QR{} توجيه شوې که نۀ، دا شرط هم پرېښے و چې ثابت~\LR{$c$} ته په
\LR{$\Gamma$} کښې د راتلو اجازه نشته۔ تعريف داسې سمولے شو چې \LR{$\Gamma$}
نېغ په نېغه په پام کښې ونيسي، خو د \LR{$\fn{Deriv}$} او د استنتاج د قضيې
کارول اسانه دي۔ \LR{$\Gamma\Proves !A$} هله صدق کوي چې د تړلي فارمولونه
داسې کوم متناهي لست \LR{$!B_1$}، \dots، \LR{$!B_n\in\Gamma$} وي چې
\LR{$\{!B_1,\dots,!B_n\}\Proves !A$}۔ او د استنتاج د قضيې له مخې دا هله
سم دي چې
\LR{$\Proves(!B_1\lif(!B_2\lif\cdots(!B_n\lif !A)\cdots))$} وي۔ دا چې
د ګوډل شمېرې~\LR{$z$} لرونکې تړلے فارمول په دې بڼه ده که نۀ، په
بنسټيز بازګشتي ډول ازمويلے شو۔ نو \LR{$x$} له~\LR{$\Gamma$} څخه د ګوډل
شمېرې~\LR{$y$} لرونکې تړلے فارمول د اشتقاق د ګوډل شمېرې په توګه
نۀ اخلو، بلکې \LR{$x$} له~\LR{$\emptyset$} څخه د پورته بڼې د يوه ځالي شرطي د
اشتقاق ګوډل شمېره اخلو۔
\olpseganchor{OLP-0288-B015}{end}

\olpseganchor{OLP-0288-B016}{start}
لومړے، که د تړلي فارمولونه يوه لړۍ ولرو، د ټولو دغو جملو په مقدماتو
او د يوې ورکړې تړلے فارمول په تالي سره شرطي په بنسټيز بازګشتي ډول
جوړولے شو:
\begin{align*}
  \fn{hCond}(s, y, 0) & = y \\
  \fn{hCond}(s, y, n+1) & = \Gn{(} \concat (s)_{n} \concat \Gn{\lif}
  \concat \fn{hCond}(s, y, n) \concat \Gn{)}\\
  \fn{Cond}(s, y) & = \fn{hCond}(s, y, \len{s})\\
  \intertext{\RL{نو \LR{$\Prf[\Gamma](x,y)$} داسې تعريفولے شو:}}
  \Prf[\Gamma](x, y) & \defiff \bexists{s < \fn{sequenceBound}(x,x)}{(} \\
  &\qquad (x)_{\len{x}-1} = \fn{Cond}(s,y) \land {} \\
  &\qquad\bforall{i<\len{s}}{(s)_i \in \Gamma} \land {} \\
  &\qquad\fn{Deriv}(x)).
\end{align*}
\textbf{د ژباړې د نسخې سمون (OLCMP-082):} سرچينې د
\LR{$\fn{hCond}$} په بازګشتي معادله کښې پر ښي لوري \LR{$\fn{Cond}$} کاراوه،
سره له دې چې \LR{$\fn{Cond}$} يوازې په ورپسې کرښه د بشپړې مرستندويې تابع
له لارې تعريفېږي۔ دلته بازګشتي غږ بېرته \LR{$\fn{hCond}$} ته شوے دے۔
د~\LR{$s$} حد له دې څخه راځي چې هر \LR{$(s)_i$} د اشتقاق د وروستۍ
کرښې د يوه فرعي-فارمول ګوډل شمېره ده؛ يعنې له
\LR{$(x)_{\len{x}-1}$} څخه وړه ده۔ د مقدماتو~\LR{$!B \in \Gamma$} شمېر،
يعنې د~\LR{$s$} اوږدوالے، د~\LR{$x$} د وروستۍ کرښې له اوږدوالي څخه کم دے۔
\end{proof}
\olpseganchor{OLP-0288-B016}{end}

\olpseganchor{OLP-0289-B004}{start}
\olchapter{inc}{req}{په~\LR{$\Th{Q}$} کښې تمثيلېدنه}
\olpseganchor{OLP-0289-B004}{end}

\olpseganchor{OLP-0290-B004}{start}
\olfileid{inc}{req}{int}
\olsection{پېژندګلو}
\olpseganchor{OLP-0290-B004}{end}

\olpseganchor{OLP-0290-B005}{start}
د ناتکميلۍ قضيې پر هغو تيوريو پلي کېږي چې د محاسبه کېدونکو تابعو
په اړه بنسټيز حقايق پکې بيان او ثابتېدے شي۔ موږ به يوه ډېره لږه
داسې تيوري بيان کړو چې «\LR{$\Th{Q}$}» (يا کله کله د رافايل رابنسن په
نوم «د رابنسن~\LR{$Q$}») بلل کېږي۔ وبه وايو چې يوه تابعه په~\LR{$\Th{Q}$}
کښې
\emph{تمثيلېدونکې} وي څه معنا لري، او بيا به لاندې خبره ثابته کړو:
\begin{quote}
  يوه تابعه په~\LR{$\Th{Q}$} کښې هله او يوازې هله تمثيلېدونکې ده چې
  محاسبه کېدونکې وي۔
\end{quote}
دا له يوې خوا د محاسبې يو بل مدل راکوي۔ خو له بلې خوا به ئې د دې د
ښودلو دپاره هم وکاروو چې سټ~\LR{$\Setabs{!A}{\Th{Q} \Proves !A}$} د
پرېکړې وړ نۀ دے؛ د درېدنې مسئله به ورته راکمه کړو۔ تر هغه وخته چې
بحث بشپړوو، تر دې به ډېرې پياوړې پايلې مو ثابتې کړې وي۔
\olpseganchor{OLP-0290-B005}{end}

\olpseganchor{OLP-0290-B006}{start}
د~\LR{$\Th{Q}$} ژبه د حساب ژبه ده؛ \LR{$\Th{Q}$} له لاندې بديهي اصولو څخه
جوړه ده (چې د عينيت لرونکي لومړۍ درجې منطق له نورو بديهي
اصولو او قاعدو سره يوځای کارول کېږي):
\begin{align*}
& \lforall[x][\lforall[y][(\eq[x'][y'] \lif \eq[x][y])]] \tag{\LR{$!Q_1$}}\\
& \lforall[x][\eq/[\Obj 0][x']] \tag{\LR{$!Q_2$}}\\
& \lforall[x][(\eq[x][\Obj 0] \lor \lexists[y][\eq[x][y']])] \tag{\LR{$!Q_3$}}\\
& \lforall[x][\eq[(x + \Obj 0)][x]] \tag{\LR{$!Q_4$}}\\
& \lforall[x][\lforall[y][\eq[(x + y')][(x + y)']]] \tag{\LR{$!Q_5$}}\\
& \lforall[x][\eq[(x \times \Obj 0)][\Obj 0]] \tag{\LR{$!Q_6$}}\\
& \lforall[x][\lforall[y][\eq[(x \times y')][((x \times y) + x)]]] \tag{\LR{$!Q_7$}}\\
& \lforall[x][\lforall[y][(x < y \liff \lexists[z][\eq[(z' + x)][y]])]] \tag{\LR{$!Q_8$}}
\end{align*}
د هر طبيعي عدد~\LR{$n$} دپاره عددنښه~\LR{$\num{n}$} داسې تعريف کړئ چې ترم
\LR{$\Obj{0}^{\prime\prime\ldots\prime}$} وي او د پريم نښو ټول شمېر ئې~\LR{$n$} وي۔ نو
\LR{$\num{0}$} يوازې ثابت سمبول~\LR{$\Obj{0}$} دے، \LR{$\num{1}$}،
\LR{$\Obj{0}'$} دے، \LR{$\num{2}$}، \LR{$\Obj{0}''$} دے، او همداسې وړاندې۔
\olpseganchor{OLP-0290-B006}{end}

\olpseganchor{OLP-0290-B007}{start}
د حساب د يوې تيورۍ په توګه~\LR{$\Th{Q}$} \emph{بېخي} کمزورې ده؛ د
بېلګې په توګه ان د~\LR{$\lforall[x][\eq/[x][x']]$} يا
\LR{$\lforall[x][\lforall[y][(x + y) = (y + x)]]$} په شان ډېر ساده
حقايق هم پکې نۀ شئ ثابتولے۔ خو وبه وينو چې د~\LR{$\Th{Q}$} د پام وړتيا
تر ډېره \emph{له همدې} کمزورۍ څخه راځي۔ په حقيقت کښې دومره پياوړې
ده چې د ناتکميلۍ قضيه پرې پلې شي، خو له دې زياته نۀ ده۔ د~\LR{$\Th{Q}$}
د پام وړتيا بل لامل دا دے چې د بديهي اصولو سټ ئې \emph{متناهي} دے۔
\olpseganchor{OLP-0290-B007}{end}

\olpseganchor{OLP-0290-B008}{start}
له~\LR{$\Th{Q}$} څخه يوه پياوړې تيوري (چې د \emph{پيانو حساب}
\LR{$\Th{PA}$} بلل کېږي) په~\LR{$\Th{Q}$} کښې د استقرا شېما په زياتولو
ترلاسه کېږي:
\[
(!A(\Obj 0) \land \lforall[x][(!A(x) \lif !A(x'))]) \lif \lforall[x][!A(x)]
\]
دلته~\LR{$!A(x)$} هر فارمول دے۔ که~\LR{$!A(x)$} له~\LR{$x$} پرته نور ازاد
متغيرونه ولري، په سر کښې کلي کمیت ايښوونکي ورزياتوو چې ټول وتړي
(څو د استقرا د شېما اړونده بېلګه تړلے فارمول وي)۔ د بېلګې په
توګه، که~\LR{$!A(x, y)$} کښې متغير~\LR{$y$} هم ازاد وي، اړونده بېلګه
دا ده:
\[
\lforall[y][((!A(\Obj 0) \land \lforall[x][(!A(x) \lif !A(x'))]) \lif
  \lforall[x][!A(x)])]
\]
د استقرا د شېما د بېلګو په کارولو سره د~\LR{$\Th{PA}$} له بديهي اصولو
څخه د~\LR{$\Th{Q}$} په پرتله ډېر څه ثابتېدے شي۔ په حقيقت کښې د طبيعي
عددونو په اړه د داسې «طبيعي» بيانونو موندل ډېر کار غواړي چې د پيانو
په حساب کښې نۀ ثابتېږي!{}
\olpseganchor{OLP-0290-B008}{end}

\olpseganchor{OLP-0290-B009}{start}
\begin{defn}
\ollabel{defn:representable-fn}
  له طبيعي عددونو څخه طبيعي عددونو ته تابعه
  \LR{$f(x_0,\ldots,x_k)$} هله \emph{په~\LR{$\Th{Q}$} کښې تمثيلېدونکې}
  بلل کېږي چې داسې فارمول~\LR{$!A_f(x_0,\dots,x_k,y)$} وي چې هر
  کله~\LR{$f(n_0,\dots,n_k) = m$} وي، \LR{$\Th{Q}$} لاندې دواړه ثابت کړي:
\begin{enumerate}
\item\ollabel{defn:rep:a} \LR{$!A_f(\num{n_0}, \dots, \num{n_k}, \num{m})$}
\item\ollabel{defn:rep:b} \LR{$\lforall[y][(!A_f(\num{n_0}, \dots,
\num{n_k}, y) \lif \num{m} = y)]$}۔
\end{enumerate}
\end{defn}
\olpseganchor{OLP-0290-B009}{end}

\olpseganchor{OLP-0290-B010}{start}
د تعريف د بيان نورې برابرې لارې هم شته؛ د بېلګې په توګه په برابر
ډول دا غوښتلے شو چې~\LR{$\Th{Q}$}،
\LR{$\lforall[y][(!A_f(\num{n_0}, \dots,
\num{n_k}, y) \liff \eq[y][\num{m}])]$} ثابت کړي۔
\olpseganchor{OLP-0290-B010}{end}

\olpseganchor{OLP-0290-B011}{start}
\begin{thm}
\ollabel{thm:representable-iff-comp}
يوه تابعه په~\LR{$\Th{Q}$} کښې هله او يوازې هله تمثيلېدونکې ده چې
محاسبه کېدونکې وي۔
\end{thm}
\olpseganchor{OLP-0290-B011}{end}

\olpseganchor{OLP-0290-B012}{start}
د قضيې ثبوت دوه لوري لري۔ له کيڼ څخه ښي لور ته ثبوت، کله چې د نحو
حسابي کول جوړ شي، نسبتاً ساده دے۔ بل لور زيات کار غواړي۔ بنسټيزه
مفکوره دا ده: «عمومي بازګشتي» د «محاسبه کېدونکي» د کره کولو يوه لار
ټاکو، او ښيو چې هره عمومي بازګشتي تابعه په~\LR{$\Th{Q}$} کښې
تمثيلېدونکې ده۔ په ياد ولرئ، تابعه هله عمومي بازګشتي ده چې له
\LR{$\Zero$}، تالي تابع~\LR{$\Succ$}، او پروجکشن تابعو~\LR{$\Proj{n}{i}$} څخه د
ترکيب، بنسټيز بازګښت، او منظم اقل موندلو په کارولو سره تعريفېدے شي۔
نو د دې د ښودلو يوه لار چې هره عمومي بازګشتي تابعه په~\LR{$\Th{Q}$} کښې
تمثيلېدونکې ده، دا ده چې وښيو بنسټيزې تابعې تمثيلېدونکې دي، او هر
کله چې ځينې تابعې تمثيلېدونکې وي، د هغو له ترکيب، بنسټيز بازګښت، او
منظم اقل موندلو څخه تعريف شوې تابعې هم تمثيلېدونکې وي۔ په بله وينا،
ښودلے شو چې بنسټيزې تابعې تمثيلېدونکې دي، او تمثيلېدونکې تابعې د
ترکيب، بنسټيز بازګښت، او منظم اقل موندلو «لاندې تړلې» دي۔ دا خبره
باور راکوي چې هره عمومي بازګشتي تابعه تمثيلېدونکې ده۔
\olpseganchor{OLP-0290-B012}{end}

\olpseganchor{OLP-0290-B013}{start}
خو دا ګام چې وښيو تمثيلېدونکې تابعې د بنسټيز بازګښت لاندې تړلې دي،
ستونزمن راوځي۔ له دې ګام څخه د ډډې دپاره لومړے ښيو چې په حقيقت کښې
له بنسټيز بازګښت پرته کار کولے شو۔ يعنې ښيو چې هره عمومي بازګشتي
تابعه يوازې د بنسټيزو تابعو، ترکيب، او منظم اقل موندلو په کارولو
تعريفېدے شي۔ د دې دپاره ښيو چې بنسټيز بازګښت په رښتيا د يوه ځانګړي
منظم اقل موندلو په وسيله ترسره کېدے شي۔ خو د دې د سم کار دپاره بايد
يو شمېر نورې بنسټيزې تابعې هم ورزياتې کړو: جمع، ضرب، او د عينيت د
اړيکې ځانګړونکې تابعه~\LR{$\Char{=}$}۔ بيا قضيه په دې ښودلو ثابتولے شو
چې دا \emph{ټولې} بنسټيزې تابعې په~\LR{$\Th{Q}$} کښې تمثيلېدونکې دي، او
تمثيلېدونکې تابعې د ترکيب او منظم اقل موندلو لاندې تړلې دي۔
\olpseganchor{OLP-0290-B013}{end}

\olpseganchor{OLP-0291-B004}{start}
\olfileid{inc}{req}{rpc}
\olsection{په~\LR{$\Th{Q}$} کښې تمثيلېدونکې تابعې محاسبه کېدونکې دي}
\olpseganchor{OLP-0291-B004}{end}

\olpseganchor{OLP-0291-B005}{start}
موږ به ثابته کړو چې هره هغه تابعه چې په~\LR{$\Th{Q}$} کښې تمثيلېدونکې
وي، محاسبه کېدونکې ده۔ لومړے بايد په~\LR{$\Th{Q}$} کښې د تمثيلېدونکو
تابعو په اړه يوه لمه ثابته کړو۔
\olpseganchor{OLP-0291-B005}{end}

\olpseganchor{OLP-0291-B006}{start}
\begin{lem}\ollabel{lem:rep-q}
  که~\LR{$f(x_0,
  \dots, x_k)$} په~\LR{$\Th{Q}$} کښې تمثيلېدونکې وي، داسې
  فارمول~\LR{$!A_f(x_0, \dots, x_k, y)$} شته چې
  \[
  \Th{Q} \Proves !A_f(\num{n_0}, \dots, \num{n_k}, \num{m})
  \quad\text{\RL{هله او يوازې هله چې}}\quad m = f(n_0, \dots, n_k).
  \]
\end{lem}
\textbf{د ژباړې د نسخې سمون (OLCMP-083):} سرچينې په مقدمې کښې د
تمثيلوونکي فارمول فرعي نښه پرېښې وه، خو په نمايش او ټول ورپسې ثبوت
کښې ئې د تابعې له نوم سره فرعي نښه کارولې؛ دلته نښه يو شان شوې ده۔
\olpseganchor{OLP-0291-B006}{end}

\olpseganchor{OLP-0291-B007}{start}
\begin{proof}
  «که» برخه
\olref[int]{defn:representable-fn}\olref[int]{defn:rep:a} ده۔ «يوازې
که» برخه داسې ليدل کېږي: فرض کړئ
\LR{$\Th{Q} \Proves !A_f(\num{n_0},
\dots, \num{n_k}, \num{m})$}، خو~\LR{$m \neq f(n_0, \dots, n_k)$}۔
\LR{$l = f(n_0, \dots, n_k)$} ونيسئ۔ د
\olref[int]{defn:representable-fn}\olref[int]{defn:rep:a} له مخې،
\LR{$\Th{Q} \Proves !A_f(\num{n_0}, \dots, \num{n_k}, \num{l})$}۔ د
\olref[int]{defn:representable-fn}\olref[int]{defn:rep:b} له مخې،
\LR{$\lforall[y][(!A_f(\num{n_0}, \dots, \num{n_k}, y) \lif \num{l} =
y)]$}۔ د منطق او دې فرض په کارولو سره چې~\LR{$\Th{Q} \Proves
!A_f(\num{n_0}, \dots, \num{n_k}, \num{m})$}، ترلاسه کوو چې
\LR{$\Th{Q} \Proves \eq[\num{l}][\num{m}]$}۔ له بلې خوا، د
\olref[bre]{lem:q-proves-neq} له مخې، \LR{$\Th{Q} \Proves
\eq/[\num{l}][\num{m}]$}۔ نو~\LR{$\Th{Q}$} ناسازګاره ده۔ خو دا ناشونې
ده، ځکه~\LR{$\Th{Q}$} په معياري مدل کښې پوره کېږي (وګورئ
\olref[int][def]{def:standard-model})، \LR{$\Sat{N}{\Th{Q}}$}، او د صحت
د قضيې له مخې هره د پوره کېدو وړ تيوري تل سازګاره وي
(\tagrefs{prfAX/{fol:axd:sou:cor:consistency-soundness},
prfSC/{fol:seq:sou:cor:consistency-soundness},
prfND/{fol:ntd:sou:cor:consistency-soundness},
prfTab/{fol:tab:sou:cor:consistency-soundness}})۔
\end{proof}
\olpseganchor{OLP-0291-B007}{end}

\olpseganchor{OLP-0291-B008}{start}
\begin{lem}
هره هغه تابعه چې په~\LR{$\Th{Q}$} کښې تمثيلېدونکې وي، محاسبه کېدونکې
ده۔
\end{lem}
\olpseganchor{OLP-0291-B008}{end}

\olpseganchor{OLP-0291-B009}{start}
\begin{proof}
لومړے د دې خبرې شهودي مفکوره وړاندې کړو۔ د~\LR{$f$} د محاسبې دپاره
لاندې کار کوو۔ د حساب د ژبې ټول ممکن اشتقاقونه~\LR{$\delta$} په
لړ کښې راوړو۔ دا کار په ميخانيکي ډول کېدے شي۔ د هر يوه دپاره ګورو
چې آيا د~\LR{$!A_f(\num{n_0}, \dots, \num{n_k}, \num{m})$} په بڼه د
فارمول کوم اشتقاق دے که نۀ (هماغه فارمول چې د
\olref{lem:rep-q} له مخې په~\LR{$\Th{Q}$} کښې~\LR{$f$} تمثيلوي)۔ که وي، د
\olref{lem:rep-q} له مخې~\LR{$m = f(n_0, \dots, n_k)$}، نو د~\LR{$f$} ارزښت
مو وموند۔ پلټنه پاې ته رسېږي، ځکه
\LR{$\Th{Q} \Proves !A_f(\num{n_0}, \dots, \num{n_k},
\num{f(n_0, \dots, n_k)})$}؛ نو بالاخره د سم ډول کوم~\LR{$\delta$} مومو۔
\olpseganchor{OLP-0291-B009}{end}

\olpseganchor{OLP-0291-B010}{start}
دا لا بشپړ کره نۀ دي، ځکه زموږ کړنلاره يوازې پر عددونو د کار پر ځاے
پر اشتقاقونه او فارمولونه کار کوي، او لا مو په کره توګه نۀ
دي ويلي چې ولې «د ټولو ممکنو اشتقاقونه لړل» په ميخانيکي ډول
شوني دي۔ خو لکه چې ليدلي مو دي، ترمونه، فارمولونه او
اشتقاقونه په ګوډل شمېرو کوډېدے شي۔ د محاسبې يو کره مدل، يعنې
عمومي بازګشتي تابعې، مو هم معرفي کړے دے۔ او ليدلي مو دي چې اړيکه
\LR{$\Prf[\Th{Q}](d,y)$}، چې هله صدق کوي چې~\LR{$d$} د~\LR{$\Th{Q}$} له بديهي
اصولو څخه د ګوډل شمېرې~\LR{$y$} لرونکي فارمول د کوم
اشتقاق ګوډل شمېره وي، (بنسټيزه) بازګشتي ده۔ نورې بنسټيزې
بازګشتي تابعې چې پکار مو دي~\LR{$\fn{num}$}
(\olref[art][trm]{prop:num-primrec}) او~\LR{$\fn{Subst}$}
(\olref[art][sub]{prop:subst-primrec}) دي۔ له دغو څخه~\LR{$f$} د اقل
موندلو په وسيله تعريفېدے شي؛ نو~\LR{$f$} بازګشتي ده۔
\olpseganchor{OLP-0291-B010}{end}

\olpseganchor{OLP-0291-B011}{start}
لومړے دا تعريف کړئ:
\begin{multline*}
  A(n_0, \dots, n_k, m) = \\
  \fn{Subst}(\fn{Subst}(\dots\fn{Subst}(\Gn{!A_f}, \fn{num}(n_0), \Gn{x_0}),\\ \dots),
  \fn{num}(n_k),  \Gn{x_k}), \fn{num}(m), \Gn{y})
\end{multline*}
دا پېچلې ښکاري، خو يوازې همدا تابعه ده:
\LR{$A(n_0, \dots, n_k,
m) = \Gn{!A_f(\num{n_0}, \dots, \num{n_k}, \num{m})}$}۔
\olpseganchor{OLP-0291-B011}{end}

\olpseganchor{OLP-0291-B012}{start}
اوس اړيکه~\LR{$R(n_0, \dots, n_k, s)$} وګورئ چې هله صدق کوي چې
\LR{$(s)_0$} د~\LR{$\Th{Q}$} څخه د
\LR{$!A_f(\num{n_0}, \dots, \num{n_k}, \num{(s)_1})$} د کوم
اشتقاق ګوډل شمېره وي:
\[
R(n_0, \dots, n_k, s) \quad\text{\RL{هله او يوازې هله چې}}\quad \Prf[\Th{Q}]((s)_0, A(n_0,
\dots, n_k, (s)_1))
\]
که داسې~\LR{$s$} ومومو چې~\LR{$R(n_0, \dots, n_k, s)$} پرې صدق کوي، د عددونو
داسې جوړه---\LR{$(s)_0$} او~\LR{$(s)_1$}---مو موندلې چې~\LR{$(s)_0$} د
\LR{$A_f(\num{n_0}, \dots, \num{n_k}, (s)_1)$} د کوم اشتقاق
ګوډل شمېره ده۔ نو د~\LR{$s$} لټول په غيرصوري ثبوت کښې د~\LR{$d$} او~\LR{$m$} د
جوړې د لټولو په شان دي۔ يوه محاسبه کېدونکې تابعه چې داسې~\LR{$s$}
«لټوي»، په منظم اقل موندلو تعريفېدے شي۔ په پام کښې ونيسئ چې~\LR{$R$}
منظمه ده: د هر~\LR{$n_0$}، \dots،~\LR{$n_k$} دپاره د
\LR{$\Th{Q} \Proves !A_f(\num{n_0}, \dots, \num{n_k},
\num{f(n_0, \dots, n_k)})$} کوم اشتقاق~\LR{$\delta$} شته؛ نو
\LR{$R(n_0, \dots, n_k, s)$} د
\LR{$s = \tuple{\Gn{\delta}, f(n_0, \dots, n_k)}$} دپاره صدق کوي۔ نو~\LR{$f$}
داسې ليکلے شو:
\[
f(n_0,\dots,n_{k}) = (\umin{s}{R(n_0, \dots, n_k, s)})_1.
\]
\end{proof}
\olpseganchor{OLP-0291-B012}{end}

\olpseganchor{OLP-0292-B004}{start}
\olfileid{inc}{req}{bet}
\olsection{د بېټا تابعې لمه}
\olpseganchor{OLP-0292-B004}{end}


\olpseganchor{OLP-0292-B005}{start}
د دې د ښودلو دپاره چې که جمع، ضرب، او~\LR{$\Char{=}$} راسره وي نو
بنسټيز بازګښت ترسره کولے شو، بايد داسې تابعې جوړې کړو چې لړۍ سمبال
کړي۔ (که توان هم راسره وے، کار به مو اسان و۔) کله چې بنسټيز بازګښت
راسره و، د «\LR{$n$}م اول عدد» په شان څيزونه مو تعريفولے شول او نسبتاً
ساده کوډونه مو غوره کولے شول۔ خو دلته بنسټيز بازګښت راسره نشته---په
حقيقت کښې غواړو وښيو چې بنسټيز بازګښت د اقل موندلو په کارولو ترسره
کولے شو---نو بايد څه زياته ځيرکتيا وکاروو۔
\olpseganchor{OLP-0292-B005}{end}

\olpseganchor{OLP-0292-B006}{start}
\begin{lem}
\ollabel{lem:beta}
داسې تابعه~\LR{$\beta(d,i)$} شته چې د هرې لړۍ \LR{$a_0$}، \dots،~\LR{$a_n$}
دپاره داسې عدد~\LR{$d$} شته چې د هر~\LR{$i \le n$} دپاره
\LR{$\beta(d,i) = a_i$} وي۔ سربېره پر دې،~\LR{$\beta$} له بنسټيزو تابعو څخه
يوازې د ترکيب او منظم اقل موندلو په کارولو تعريفېدے شي۔
\end{lem}
\olpseganchor{OLP-0292-B006}{end}

\olpseganchor{OLP-0292-B007}{start}
\LR{$d$} داسې وګڼئ چې لړۍ~\LR{$\tuple{a_0, \dots, a_n}$} کوډوي، او
\LR{$\beta(d,i)$} ئې~\LR{$i$}م غړے بېرته ورکوي۔ (پام وکړئ چې دا «کوډول» هغه
د اول عددونو د توانونو کوډول \emph{نۀ} کاروي چې لا له وړاندې ورسره
اشنا يو!) لمه ډېره لږه ادعا کوي؛ دا نۀ وايي چې لړۍ يو ځاے کولے شو،
غړي ورزياتولے شو، يا ان دا چې د ترکيب او منظم اقل موندلو له لارې
تعريفېدونکو تابعو په کارولو له~\LR{$a_0$}، \dots،~\LR{$a_n$} څخه~\LR{$d$}
\emph{محاسبه} کولے شو۔ يوازې دومره وايي چې د «کوډ پرانيستلو» داسې
تابعه شته چې هره لړۍ «کوډ شوې» وي۔
\olpseganchor{OLP-0292-B007}{end}

\olpseganchor{OLP-0292-B008}{start}
د~\LR{$\beta$} نښه د ګوډل ده۔ بيا وايو، د لمې د ثبوت ګرانه برخه دا ده
چې د ظاهراً محدودو وسيلو، يعنې يوازې د ترکيب او اقل موندلو، په
کارولو مناسبه~\LR{$\beta$} تعريف کړو---خو اجازه لرو چې جمع، ضرب،
او~\LR{$\Char{=}$} وکاروو۔ د دې لمې د ثبوت بېلابېلې لارې شته، خو يوه له
تر ټولو پاکو لارو څخه لا هم د ګوډل اصلي لاره ده، چې د عددونو د تيورۍ
هغه حقيقت ئې کاراوه چې د سونزي قضيه بلل کېږي (په دوديز ډول، «د چين
د پاتې شونو قضيه»)۔
\olpseganchor{OLP-0292-B008}{end}

\olpseganchor{OLP-0292-B009}{start}
\begin{defn}
دوه طبيعي عددونه~\LR{$a$} او~\LR{$b$} هله \emph{خپلمنځي اول} دي چې ستر
ګډ مقسوم‌عليه ئې~\LR{$1$} وي؛ په بله وينا، له دې پرته بل ګډ مقسوم‌عليه
نۀ لري۔
\end{defn}
\olpseganchor{OLP-0292-B009}{end}

\olpseganchor{OLP-0292-B010}{start}
\begin{defn}
طبيعي عددونه~\LR{$a$} او~\LR{$b$} \emph{د~\LR{$c$} په مودولو سره مطابق} دي،
\LR{$a \equiv b \mod c$}، هله چې~\LR{$c \mid (a-b)$}؛ يعنې کله چې پر~\LR{$c$}
ووېشل شي،~\LR{$a$} او~\LR{$b$} يو شان پاتې شونې لري۔
\end{defn}
\olpseganchor{OLP-0292-B010}{end}

\olpseganchor{OLP-0292-B011}{start}
دا د سونزي قضيه ده:
\begin{thm}
فرض کړئ \LR{$x_0$}، \dots،~\LR{$x_n$} (جوړه په جوړه) خپلمنځي اول دي۔
\LR{$y_0$}، \dots،~\LR{$y_n$} هر عددونه وي۔ بيا داسې عدد~\LR{$z$} شته چې
\begin{align*}
z & \equiv y_0 \mod x_0 \\
z & \equiv y_1 \mod x_1 \\
& \vdots  \\
z & \equiv y_n \mod x_n.
\end{align*}
\end{thm}
\olpseganchor{OLP-0292-B011}{end}

\olpseganchor{OLP-0292-B012}{start}
د سونزي قضيه به داسې کاروو: که \LR{$x_0$}، \dots،~\LR{$x_n$} په ترتيب سره
له~\LR{$y_0$}، \dots،~\LR{$y_n$} څخه لوی وي،~\LR{$z$} د لړۍ
\LR{$\tuple{y_0, \dots,y_n}$} کوډ ګڼلے شو۔ د~\LR{$y_i$} د بېرته موندلو
دپاره يوازې~\LR{$z$} پر~\LR{$x_i$} وېشو او پاتې شونې ئې اخلو۔ د دې کوډ کارولو
دپاره بايد د~\LR{$x_0$}، \dots،~\LR{$x_n$} مناسب ارزښتونه ومومو۔
\olpseganchor{OLP-0292-B012}{end}

\olpseganchor{OLP-0292-B013}{start}
په دې برخه کښې به څو کتنې راسره مرسته وکړي۔ د ورکړل شوو
\LR{$y_0$}، \dots،~\LR{$y_n$} دپاره دا ونيسئ:
\begin{align*}
j &= \max(n, y_0 + 1, \dots, y_n + 1), \\
m &= \lcm(1,\dots,j),
\end{align*}
او دا ونيسئ:
\begin{align*}
x_0 & = 1 + m \\
x_1 & = 1 + 2 \cdot m \\
x_2 & = 1 + 3 \cdot m \\
& \vdots  \\
x_n & = 1 + (n+1) \cdot m
\end{align*}
بيا دوه خبرې سمې دي:
\begin{enumerate}
\item\ollabel{rel-prime} \LR{$x_0,\dots,x_n$} خپلمنځي اول دي۔
\item\ollabel{less} د هر~\LR{$i$} دپاره، \LR{$y_i < x_i$}۔
\end{enumerate}
د دې ليدلو دپاره چې~\olref{rel-prime} سمه ده، پام وکړئ چې که~\LR{$p$}
اول عدد وي او \LR{$p \mid x_i$} او~\LR{$p \mid x_k$} وي، نو
\LR{$p \mid 1 + (i+1) m$} او~\LR{$p \mid 1 + (k+1) m$} دي۔ خو بيا~\LR{$p$} د
دواړو توپير وېشي:
\[
(1 + (i+1)m) - (1+ (k+1)m) = (i-k) m.
\]
څرنګه چې~\LR{$p$}، \LR{$1 + (i+1)m$} وېشي، په عين وخت کښې~\LR{$m$} نشي وېشلے
(که ئې وېشلے، لومړۍ وېشنې به~\LR{$1$} پاتې شونې پرېښې وه)۔ نو ځکه چې
\LR{$p$}، \LR{$(i-k)m$} وېشي،~\LR{$p$} بايد~\LR{$i-k$} ووېشي۔ خو~\LR{$\left|i-k\right|$} تر
ډېره~\LR{$n$} دے، او موږ~\LR{$j \geq n$} غوره کړے دے، نو له دې څخه بيا
\LR{$p \mid m$} لازمېږي، چې تناقض دے۔ نو هېڅ اول عدد~\LR{$x_i$} او~\LR{$x_k$}
دواړه نۀ وېشي۔ فقره~\olref{less} اسانه ده:
\LR{$y_i < j \leq m < x_i$} لرو۔
\olpseganchor{OLP-0292-B013}{end}

\olpseganchor{OLP-0292-B014}{start}
اوس د~\LR{$\beta$} تابعې لمه ثابته کړو۔ په ياد ولرئ چې~\LR{$0$}، تالي،
جمع، ضرب،~\LR{$\Char{=}$}، پروجکشنونه، او هره هغه تابعه کارولے شو چې
له دوى څخه د ترکيب او پر منظمو تابعو د اقل موندلو په کارولو تعريف
شوې وي۔ که د کومې اړيکې ځانګړونکې تابعه همداسې تعريفېدونکې وي، هغه
اړيکه هم کارولے شو۔ د پخوا په شان ښودلے شو چې دا اړيکې د بولي
ترکيبونو او محدود کمیت ايښودنې په وړاندې تړلې دي؛ د بېلګې په توګه:
\begin{align*}
\fn{not}(x) & \defis \Char{=}(x,0)\\
\bmin{x \leq z}{R(x,y)} & \defis \umin{x}{(R(x,y) \lor x = z)}\\
\bexists{x \leq z}{R(x,y)} & \defiff R(\bmin{x \leq z}{R(x,y)}, y)
\end{align*}
بيا ښودلے شو چې دا ټول هم له بنسټيز بازګښت پرته تعريفېدونکي دي:
\begin{enumerate}
\item جوړوونکې تابعه، \LR{$J(x,y) = \frac{1}{2}[(x+y)(x+y+1)] + x$}؛
% maybe explain more what is going on here, a bit confusing.
\item د پروجکشن تابعې
\begin{align*}
K(z) & = \bmin{x \leq z}{\bexists{y \leq z}{z = J(x,y)}},\\
L(z) & = \bmin{y \leq z}{\bexists{x \leq z}{z = J(x,y)}};
\end{align*}
\item د کوچني والي اړيکه~\LR{$x < y$}؛
\item د وېش اړيکه~\LR{$x \mid y$}؛
% \item $x \tsub y$
% \item $\fn{Prime}(x)$
% \item Assuming $p$ is prime, the relation ``$x$ is a power of $p$'':
% \[
% \bforall{y \leq x}{(y \mid x \lif y = 1 \lor y = x)}.
% \]
\item تابعه~\LR{$\fn{rem}(x,y)$}، چې د~\LR{$y$} پر~\LR{$x$} د وېش پاتې شونې
  بېرته ورکوي۔
\end{enumerate}
اوس تعريف کړئ:
\begin{align*}
\beta^*(d_0,d_1,i) & = \fn{rem}(1+(i+1) d_1,d_0) \text{\RL{ او}}\\
\beta(d,i) & = \beta^*(K(d),L(d),i).
\end{align*}
دا هماغه تابعه ده چې غواړو ئې۔ پورته غوندې ورکړل شوو
\LR{$a_0,\dots,a_n$} دپاره دا ونيسئ:
\[
j = \max(n,a_0+1,\dots,a_n+1),
\]
او~\LR{$d_1 = \lcm(1,\dots,j)$} ونيسئ۔ له پورته
\olref{rel-prime} څخه پوهېږو چې \LR{$1+d_1$}، \LR{$1+2 d_1$}، \dots،
\LR{$1+(n+1) d_1$} خپلمنځي اول دي، او له~\olref{less} څخه چې دا ټول
له~\LR{$a_0,\dots,a_n$} څخه لوی دي۔ د سونزي د قضيې له مخې داسې ارزښت
\LR{$d_0$} شته چې د هر~\LR{$i$} دپاره
\[
d_0 \equiv a_i \mod (1+(i+1)d_1)
\]
وي، او نو (ځکه~\LR{$d_1$} له~\LR{$a_i$} څخه لوی دے)،
\[
a_i = \fn{rem}(1+(i+1)d_1,d_0).
\]
\LR{$d = J(d_0,d_1)$} ونيسئ۔ بيا د هر~\LR{$i \le n$} دپاره لرو:
\begin{align*}
\beta(d,i) & =  \beta^*(d_0,d_1,i) \\
& =  \fn{rem}(1+(i+1) d_1,d_0) \\
& =  a_i
\end{align*}
او دا هماغه څه دي چې پکار مو وو۔ په دې سره د~\LR{$\beta$} تابعې د لمې
ثبوت بشپړېږي۔
\olpseganchor{OLP-0292-B014}{end}

\olpseganchor{OLP-0292-B015}{start}
\begin{prob}
  وښيئ چې اړيکې~\LR{$x < y$}، \LR{$x \mid y$}، او تابعه
  \LR{$\fn{rem}(x,y)$} له بنسټيز بازګښت پرته تعريفېدے شي۔ \LR{$0$}، تالي،
  جمع، ضرب،~\LR{$\Char{=}$}، پروجکشنونه، او محدوده کمينه‌موندنه او کمیت
  ايښودنه کارولے شئ۔
\end{prob}
\olpseganchor{OLP-0292-B015}{end}

\olpseganchor{OLP-0293-B004}{start}
\olfileid{inc}{req}{pri}
\olsection{د بنسټيز بازګښت شبيه کول}
\olpseganchor{OLP-0293-B004}{end}

\olpseganchor{OLP-0293-B005}{start}
اوس ښودلے شو چې د بنسټيز بازګښت په وسيله تعريف د بېټا تابعې په
کارولو سره د منظم اقل موندلو له لارې «شبيه» کېدے شي۔ فرض کړئ
\LR{$f(\vec x)$} او~\LR{$g(\vec x, y, z)$} لرو۔ بيا هغه تابعه~\LR{$h(\vec x,y)$}
چې له~\LR{$f$} او~\LR{$g$} څخه په بنسټيز بازګښت تعريف شوې وي، دا ده:
\begin{align*}
h(\vec x, 0) & =  f(\vec x) \\
h(\vec x, y+1) & =  g(\vec x, y, h(\vec x, y)).
\end{align*}
\textbf{د ژباړې د نسخې سمون (OLCMP-084):} سرچينې له دې نمايش څخه
مخکښې د تابعې د ارګومېنټونو نااړوند لړ ليکلے و؛ دلته د دواړو
بازګشتي معادلو له ارګومېنټونو سره سم لړ کارول شوے دے۔
بايد وښيو چې~\LR{$h$} يوازې د ترکيب او منظم اقل موندلو په کارولو، د
بنسټيزو تابعو او هغو تابعو په مرسته چې له بنسټيزو تابعو څخه د ترکيب
او منظم اقل موندلو له لارې تعريف شوې وي (لکه~\LR{$\beta$})، له~\LR{$f$} او
\LR{$g$} څخه تعريفېدے شي۔
\olpseganchor{OLP-0293-B005}{end}

\olpseganchor{OLP-0293-B006}{start}
\begin{lem}
\ollabel{lem:prim-rec}
که~\LR{$h$} له~\LR{$f$} او~\LR{$g$} څخه په بنسټيز بازګښت تعريفېدے شي، نو له~\LR{$f$}،
\LR{$g$}، تابعو~\LR{$\Zero$}، \LR{$\Succ$}، \LR{$\Proj{n}{i}$}، \LR{$\Add$}، \LR{$\Mult$}،
\LR{$\Char{=}$} څخه د ترکيب او منظم اقل موندلو په کارولو تعريفېدے شي۔
\end{lem}
\olpseganchor{OLP-0293-B006}{end}

\olpseganchor{OLP-0293-B007}{start}
\begin{proof}
لومړے يوه مرستندويه تابعه~\LR{$\hat h(\vec x, y)$} تعريف کړئ چې تر ټولو
وړوکے عدد~\LR{$d$} بېرته ورکوي چې~\LR{$d$} داسې لړۍ کوډوي چې لاندې شرطونه پوره
کوي:
\begin{enumerate}
\item \LR{$(d)_0 = f(\vec x)$}، او
\item د هر~\LR{$i < y$} دپاره، \LR{$(d)_{i+1} = g(\vec x, i, (d)_i)$}،
\end{enumerate}
چېرته چې اوس~\LR{$(d)_i$} د~\LR{$\beta(d,i)$} لنډون دے۔ په بله وينا،
\LR{$\hat h$} دا لړۍ بېرته ورکوي:
\LR{$\tuple{h(\vec x, 0), h(\vec x, 1), \dots, h(\vec x, y)}$}۔
\LR{$\hat h$} داسې ليکلے شو:
\[
\hat h(\vec x, y) = \umin{d}{(\beta(d,0) = f(\vec x) \land \bforall{i <
  y}{\beta(d,i+1) = g(\vec x, i,\beta(d,i)})}.
\]
پام وکړئ: دلته هېڅ بنسټيز بازګښت پکار نۀ دے، يوازې اقل موندل دي۔
هغه تابعه چې اقل ئې مومو د بېټا تابعې د لمې
\olref[bet]{lem:beta} له مخې منظمه ده۔
\olpseganchor{OLP-0293-B007}{end}

\olpseganchor{OLP-0293-B008}{start}
خو اوس لرو:
\[
h(\vec x, y) = \beta(\hat h(\vec x, y), y),
\]
نو~\LR{$h$} له بنسټيزو تابعو څخه يوازې د ترکيب او منظم اقل موندلو په
کارولو تعريفېدے شي۔
\end{proof}
\olpseganchor{OLP-0293-B008}{end}

\olpseganchor{OLP-0294-B004}{start}
\olfileid{inc}{req}{bre}
\olsection{بنسټيزې تابعې په~\LR{$\Th{Q}$} کښې تمثيلېدونکې دي}
\olpseganchor{OLP-0294-B004}{end}

\olpseganchor{OLP-0294-B005}{start}
لومړے بايد وښيو چې ټولې بنسټيزې تابعې په~\LR{$\Th{Q}$} کښې
تمثيلېدونکې دي۔ په پاې کښې بايد وښيو چې هرې~\LR{$k$}-ځايه بنسټيزې
تابعې~\LR{$f(x_0,\dots,x_{k-1})$} ته داسې فارمول
\LR{$!A_f(x_0,\dots,x_{k-1},y)$} څنګه وټاکو چې هغه تمثيل کړي۔
\olpseganchor{OLP-0294-B005}{end}

\olpseganchor{OLP-0294-B006}{start}
صفر، تالي، جمع، ضرب، د مساواتو ځانګړونکې تابعه، او پروجکشنونه
تمثيلولے شو۔ په هره قضيه کښې مناسب تمثيلوونکی فارمول په بشپړه توګه
ساده دے؛ د بېلګې په توګه، صفر د~\LR{$y = \Obj 0$} فارمول په وسيله
تمثيلېږي، تالي د~\LR{$x_0' = y$} فارمول په وسيله تمثيلېږي، او جمع
د~\LR{$(x_0 + x_1) = y$} فارمول په وسيله تمثيلېږي۔ کار په دې کښې دے چې وښيو
\LR{$\Th{Q}$} اړوند تړلي فارمولونه ثابتولے شي؛ د بېلګې په توګه، دا ويل
چې جمع د پورته فارمول په وسيله تمثيل شوې ده، دا غواړي چې وښيو
د طبيعي عددونو د هرې جوړې~\LR{$m$} او~\LR{$n$} دپاره~\LR{$\Th{Q}$} دا ثابتوي:
\begin{align*}
& \eq[\num n + \num m][\num {n+m}] \text{\RL{ او}}\\
& \lforall[y][(\eq[(\num n + \num m)][y] \lif \eq[y][\num{n+m}])].
\end{align*}
\olpseganchor{OLP-0294-B006}{end}

\olpseganchor{OLP-0294-B007}{start}
\begin{prop}
\ollabel{prop:rep-zero}
صفر تابعه~\LR{$\Zero(x) = 0$} په~\LR{$\Th{Q}$} کښې د
\LR{$!A_{\Zero}(x,y) \ident \eq[y][\Obj 0]$} په وسيله تمثيلېږي۔
\end{prop}
\olpseganchor{OLP-0294-B007}{end}

\olpseganchor{OLP-0294-B008}{start}
\begin{prop}
\ollabel{prop:rep-succ}
تالي تابعه~\LR{$\Succ(x) = x+1$} په~\LR{$\Th{Q}$} کښې د
\LR{$!A_{\Succ}(x,y) \ident \eq[y][x']$} په وسيله تمثيلېږي۔
\end{prop}
\olpseganchor{OLP-0294-B008}{end}

\olpseganchor{OLP-0294-B009}{start}
\begin{prop}
\ollabel{prop:rep-proj}
پروجکشن تابعه~\LR{$\Proj{n}{i}(x_0, \dots, x_{n-1}) = x_i$} په
\LR{$\Th{Q}$} کښې د دې په وسيله تمثيلېږي:
\[!A_{\Proj{n}{i}}(x_0, \dots, x_{n-1}, y) \ident \eq[y][x_i].\]
\end{prop}
\olpseganchor{OLP-0294-B009}{end}

\olpseganchor{OLP-0294-B010}{start}
\begin{prob}
ثابته کړئ چې~\LR{$\eq[y][\Obj 0]$}، \LR{$\eq[y][x']$}، او~\LR{$\eq[y][x_i]$} په
ترتيب سره~\LR{$\Zero$}، \LR{$\Succ$}، او~\LR{$\Proj{n}{i}$} تمثيلوي۔
\end{prob}
\olpseganchor{OLP-0294-B010}{end}

\olpseganchor{OLP-0294-B011}{start}
\begin{prop}
\ollabel{prop:rep-id}
د~\LR{$=$} ځانګړونکې تابعه،
\[
\Char{=}(x_0, x_1) =
\begin{cases}
  1 & \text{\RL{که }} x_0 =x_1\\
  0 & \text{\RL{بل ډول}}
\end{cases}
\]
په~\LR{$\Th{Q}$} کښې د دې په وسيله تمثيلېږي:
\[
  !A_{\Char{=}}(x_0, x_1, y) \ident (\eq[x_0][x_1] \land \eq[y][\num{1}]) \lor (\eq/[x_0][x_1] \land
\eq[y][\num{0}]).
\]
\end{prop}
\olpseganchor{OLP-0294-B011}{end}

\olpseganchor{OLP-0294-B012}{start}
ثبوت لاندې لمې ته اړتيا لري۔
\olpseganchor{OLP-0294-B012}{end}

\olpseganchor{OLP-0294-B013}{start}
\begin{lem}
\ollabel{lem:q-proves-neq} که طبيعي عددونه~\LR{$n$} او~\LR{$m$} داسې وي چې
\LR{$n \neq m$}، نو \LR{$\Th{Q} \Proves \eq/[\num n][\num m]$}۔
\end{lem}
\olpseganchor{OLP-0294-B013}{end}

\olpseganchor{OLP-0294-B014}{start}
\begin{proof}
پر~\LR{$n$} استقرا وکاروئ څو وښيئ چې د هر~\LR{$m$} دپاره، که~\LR{$n \neq m$}
وي، نو \LR{$Q \Proves \eq/[\num n][\num m]$}۔
\olpseganchor{OLP-0294-B014}{end}

\olpseganchor{OLP-0294-B015}{start}
په بنسټيزه قضيه کښې~\LR{$n = 0$} دے۔ که~\LR{$m$} له~\LR{$0$} سره مساوي نۀ وي، نو
د کوم طبيعي عدد~\LR{$k$} دپاره~\LR{$m = k + 1$}۔ يو داسې بديهي اصل لرو چې
\LR{$\lforall[x][\eq/[0][x']]$} وايي۔ د کمیت ايښوونکي د يوه بديهي اصل له
مخې، د~\LR{$x$} پر ځاے د~\LR{$\num k$} په ايښودلو، دا پايله اخلو:
\LR{$\eq/[0][\num k']$}۔ خو~\LR{$\num k'$} هماغه~\LR{$\num m$} دے۔
\olpseganchor{OLP-0294-B015}{end}

\olpseganchor{OLP-0294-B016}{start}
د استقرا په پړاو کښې فرضولے شو چې ادعا د~\LR{$n$} دپاره سمه ده، او
\LR{$n+1$} ګورو۔ \LR{$m$} هر طبيعي عدد ونيسئ۔ دوه امکانونه شته: يا~\LR{$m = 0$}
دے يا د کوم~\LR{$k$} دپاره~\LR{$m = k+1$}۔ لومړۍ قضيه لکه پورته حل کېږي۔ په
دويمه قضيه کښې فرض کړئ~\LR{$n+1 \neq k+1$}۔ بيا~\LR{$n \neq k$}۔ د~\LR{$n$}
دپاره د استقرايي فرضيې له مخې
\LR{$\Th{Q} \Proves \eq/[\num n][\num k]$}۔ داسې بديهي اصل لرو چې
\LR{$\lforall[x][\lforall[y][\eq[x'][y'] \lif \eq[x][y]]]$} وايي۔ د
کمیت ايښوونکي د يوه بديهي اصل په کارولو، لرو:
\LR{$\eq[\num n'][\num k'] \lif \eq[\num n][\num k]$}۔ د قضيو د منطق په
کارولو، په~\LR{$\Th{Q}$} کښې دا پايله اخلو:
\LR{$\eq/[\num n][\num k] \lif \eq/[\num n'][\num k']$}۔ د وضع مقدم په
کارولو، \LR{$\eq/[\num n'][\num k']$} اخلو، چې هماغه غوښتنه ده، ځکه
\LR{$\num k'$}، \LR{$\num m$} دے۔
\end{proof}
\olpseganchor{OLP-0294-B016}{end}

\olpseganchor{OLP-0294-B017}{start}
\begin{explain}
پام وکړئ چې لمه ډېره خبره نۀ کوي: په اصل کښې وايي چې~\LR{$\Th{Q}$}
ثابتولے شي چې بېلابېلې عددنښې بېلابېل څيزونه ښيي۔ د بېلګې په
توګه،~\LR{$\Th{Q}$}، \LR{$0'' \neq 0'''$} ثابتوي۔ خو دا ښودل چې خبره په عام
ډول سمه ده څه پاملرنه غواړي۔ دا هم په پام کښې ونيسئ چې که څه هم
استقرا کاروو، خو دا د~\LR{$\Th{Q}$} \emph{بهر} استقرا ده۔
\end{explain}
\olpseganchor{OLP-0294-B017}{end}

\olpseganchor{OLP-0294-B018}{start}
\begin{proof}[د \olref{prop:rep-id} ثبوت]
که~\LR{$n = m$} وي، نو~\LR{$\num{n}$} او~\LR{$\num{m}$} يو شان ترمونه دي، او
\LR{$\Char{=}(n, m) = 1$}۔ خو \LR{$\Th{Q} \Proves (\eq[\num{n}][\num{m}] \land
\eq[\num{1}][\num{1}])$}، نو~\LR{$!A_=(\num{n}, \num{m},
\num{1})$} ثابتوي۔ که~\LR{$n \neq m$} وي، نو~\LR{$\Char=(n, m) = 0$}۔ د
\olref{lem:q-proves-neq} له مخې،
\LR{$\Th{Q} \Proves \eq/[\num{n}][\num{m}]$}، او نو دا هم ثابتوي:
\LR{$(\eq/[\num{n}][\num{m}] \land \Obj 0 = \Obj 0)$}۔ له دې امله
\LR{$\Th{Q} \Proves !A_=(\num{n}, \num{m}, \num{0})$}۔
\olpseganchor{OLP-0294-B018}{end}

\olpseganchor{OLP-0294-B019}{start}
د دويمې برخې دپاره هم دوه قضيې لرو۔ که~\LR{$n = m$} وي، بايد وښيو:
\LR{$\Th{Q} \Proves \lforall[y][(!A_=(\num{n}, \num{m}, y)
  \lif \eq[y][\num{1}])]$}۔ په غيرصوري ډول استدلال کوو، فرض کړئ
\LR{$!A_=(\num{n}, \num{m}, y)$}، يعنې
\[
(\eq[\num{n}][\num{n}] \land \eq[y][\num{1}]) \lor
(\eq/[\num{n}][\num{n}] \land \eq[y][\num{0}])
\]
کيڼ منفصل د منطق له مخې~\LR{$\eq[y][\num{1}]$} لازموي؛ ښے منفصل له
\LR{$\eq[\num{n}][\num{n}]$} سره تناقض لري، چې په منطق ثابتېدونکے دے۔
\olpseganchor{OLP-0294-B019}{end}

\olpseganchor{OLP-0294-B020}{start}
بل پلو، فرض کړئ~\LR{$n \neq m$}۔ بيا~\LR{$!A_=(\num{n},
\num{m}, y)$} دا دے:
\[
(\eq[\num{n}][\num{m}] \land \eq[y][\num{1}]) \lor
(\eq/[\num{n}][\num{m}] \land \eq[y][\num{0}])
\]
دلته کيڼ منفصل له~\LR{$\eq/[\num{n}][\num{m}]$} سره تناقض لري، چې د
\olref{lem:q-proves-neq} له مخې په~\LR{$\Th{Q}$} کښې ثابتېدونکے دے؛ ښے
منفصل~\LR{$\eq[y][\num{0}]$} لازموي۔
\end{proof}
\olpseganchor{OLP-0294-B020}{end}

\olpseganchor{OLP-0294-B021}{start}
\begin{prop}
\ollabel{prop:rep-add}
د جمع تابعه~\LR{$\Add(x_0, x_1) = x_0+x_1$} په~\LR{$\Th{Q}$} کښې د دې په
وسيله تمثيلېږي:
\[
  !A_{\Add}(x_0, x_1, y) \ident \eq[y][(x_0 + x_1)].
\]
\end{prop}
\olpseganchor{OLP-0294-B021}{end}

\olpseganchor{OLP-0294-B022}{start}
\begin{lem}
\ollabel{lem:q-proves-add}
\LR{$\Th{Q} \Proves \eq[(\num{n} + \num{m})][\num{n+m}]$}
\end{lem}
\olpseganchor{OLP-0294-B022}{end}

\olpseganchor{OLP-0294-B023}{start}
\begin{proof}
دا پر~\LR{$m$} په استقرا ثابتوو۔ که~\LR{$m = 0$} وي، ادعا دا ده:
\LR{$\Th{Q} \Proves \eq[(\num{n} + \Obj 0)][\num{n}]$}۔ دا له بديهي
اصل~\LR{$!Q_4$} څخه راځي۔ اوس د~\LR{$m$} دپاره ادعا فرض کړئ؛ راځئ د~\LR{$m+1$}
دپاره ادعا ثابته کړو، يعنې وښيو:
\LR{$\Th{Q} \Proves \eq[(\num{n} +
  \num{m+1})][\num{n+m+1}]$}۔ پام وکړئ چې~\LR{$\num{m+1}$} هماغه
\LR{$\num{m}'$} دے، او~\LR{$\num{n+m+1}$} هماغه~\LR{$\num{n+m}'$} دے۔ د بديهي
اصل~\LR{$!Q_5$} له مخې، \LR{$\Th{Q} \Proves
\eq[(\num{n} + \num{m}')][(\num{n}+\num{m})']$}۔ د استقرايي فرضيې
له مخې، \LR{$\Th{Q} \Proves
\eq[(\num{n} + \num{m})][\num{n+m}]$}۔ نو
\LR{$\Th{Q} \Proves \eq[(\num{n} + \num{m}')][\num{n+m}']$}۔
\end{proof}
\olpseganchor{OLP-0294-B023}{end}

\olpseganchor{OLP-0294-B024}{start}
\begin{proof}[د \olref{prop:rep-add} ثبوت]
هغه فارمول~\LR{$!A_\Add(x_0, x_1, y)$} چې~\LR{$\Add$} تمثيلوي،
\LR{$\eq[y][(x_0 + x_1)]$} ده۔ لومړے ښيو چې که~\LR{$\Add(n, m) = k$} وي،
نو \LR{$\Th{Q} \Proves !A_\Add(\num{n}, \num{m}, \num{k})$}؛ يعنې
\LR{$\Th{Q} \Proves \eq[\num{k}][(\num{n} + \num{m})]$}۔ خو ځکه چې
\LR{$k = n + m$}، نو~\LR{$\num{k}$} هماغه~\LR{$\num{n+m}$} دے، او په
\olref{lem:q-proves-add} کښې مو ښودلې ده چې
\LR{$\Th{Q} \Proves \eq[(\num{n} + \num{m})][\num{n+m}]$}۔
\olpseganchor{OLP-0294-B024}{end}

\olpseganchor{OLP-0294-B025}{start}
دا هم بايد وښيو چې که~\LR{$\Add(n, m) = k$} وي، نو
\[
\Th{Q} \Proves \lforall[y][(!A_\Add(\num{n}, \num{m}, y) \lif
  \eq[y][\num{k}])].
\]
فرض کړئ~\LR{$\eq[(\num{n} + \num{m})][y]$} لرو۔ ځکه چې
\[
\Th{Q} \Proves \eq[(\num{n}+\num{m})][\num{n+m}],
\]
کيڼ اړخ په~\LR{$\num{n+m}$} بدلولے شو او د خپل سري~\LR{$y$} دپاره
\LR{$\eq[\num{n+m}][y]$} تر لاسه کوو۔
\end{proof}
\olpseganchor{OLP-0294-B025}{end}

\olpseganchor{OLP-0294-B026}{start}
\begin{prop}
\ollabel{prop:rep-mult}
د ضرب تابعه~\LR{$\Mult(x_0, x_1) = x_0 \cdot x_1$} په~\LR{$\Th{Q}$} کښې د
دې په وسيله تمثيلېږي:
\[
  !A_{\Mult}(x_0, x_1, y) \ident y = (x_0 \times x_1).
\]
\end{prop}
\olpseganchor{OLP-0294-B026}{end}

\olpseganchor{OLP-0294-B027}{start}
\begin{proof} تمرین۔ \end{proof}
\olpseganchor{OLP-0294-B027}{end}

\olpseganchor{OLP-0294-B028}{start}
\begin{lem}
\ollabel{lem:q-proves-mult}
\LR{$\Th{Q} \Proves \eq[(\num{n} \times \num{m})][\num{n \cdot m}]$}
\end{lem}
\olpseganchor{OLP-0294-B028}{end}

\olpseganchor{OLP-0294-B029}{start}
\begin{proof} تمرین۔ \end{proof}
\olpseganchor{OLP-0294-B029}{end}

\olpseganchor{OLP-0294-B030}{start}
\begin{prob}
\olref[inc][req][bre]{lem:q-proves-mult} ثابته کړئ۔
\end{prob}
\olpseganchor{OLP-0294-B030}{end}

\olpseganchor{OLP-0294-B031}{start}
\begin{prob}
له \olref[inc][req][bre]{lem:q-proves-mult} څخه د
\olref[inc][req][bre]{prop:rep-mult} په ثبوت کښې کار واخلئ۔
\end{prob}
\olpseganchor{OLP-0294-B031}{end}

\olpseganchor{OLP-0294-B032}{start}
\begin{explain}
  په ياد ولرئ چې د حساب د ژبې د تابعې نښې دپاره~\LR{$\times$}، او پر
  عددونو د عادي ضرب د عمل دپاره~\LR{$\cdot$} کاروو۔ نو~\LR{$\cdot$} د عددونو
  د تعبيرونو ترمنځ راتلے شي (لکه په~\LR{$m \cdot n$} کښې)، خو~\LR{$\times$}
  يوازې د حساب د ژبې د ترمونو ترمنځ راځي (لکه په
  \LR{$(\num{m} \times \num{n})$} کښې)۔ تر دې هم مغشوشوونکې دا ده چې
  \LR{$+$} هم د تابع او هم د جمع د عمل دپاره کارېږي۔ کله چې د
  ترمونو ترمنځ راځي---لکه په~\LR{$(\num{n} + \num{m})$} کښې---دا د حساب
  د ژبې~\LR{$2$}-ځايه تابع ده، او کله چې د عددونو ترمنځ راځي---
  لکه په~\LR{$n+m$} کښې---دا د جمع عمل دے۔ په دې کښې د~\LR{$\num{n+m}$}
  قضيه هم شامله ده: دا له عدد~\LR{$n+m$} سره سمون لرونکې معياري عددنښه
  ده۔
\end{explain}
\olpseganchor{OLP-0294-B032}{end}

\olpseganchor{OLP-0295-B004}{start}
\olfileid{inc}{req}{cmp}
\olsection{ترکيب په~\LR{$\Th{Q}$} کښې تمثيلېدونکے دے}
\olpseganchor{OLP-0295-B004}{end}

\olpseganchor{OLP-0295-B005}{start}
فرض کړئ~\LR{$h$} داسې تعريف شوې ده:
\[
h(x_0,\dots,x_{l-1}) = f(g_0(x_0,\dots,x_{l-1}), \dots,
g_{k-1}(x_0,\dots,x_{l-1})).
\]
چېرته چې موږ لا وړاندې داسې فارمولونه \LR{$!A_f, !A_{g_0}, \dots,
!A_{g_{k-1}}$} موندلي دي چې په ترتيب سره تابعې~\LR{$f$}، او~\LR{$g_0$}،
\dots،~\LR{$g_{k-1}$} تمثيلوي۔ بايد داسې فارمول~\LR{$!A_h$} ومومو چې
\LR{$h$} تمثيل کړي۔
\olpseganchor{OLP-0295-B005}{end}

\olpseganchor{OLP-0295-B006}{start}
له يوې ساده قضيې پيل وکړو، چې ټولې تابعې~\LR{$1$}-ځايه وي؛ يعنې
\LR{$h(x) = f(g(x))$} وګورئ۔ که~\LR{$!A_f(y, z)$}، \LR{$f$} تمثيلوي، او
\LR{$!A_g(x, y)$}، \LR{$g$} تمثيلوي، داسې فارمول~\LR{$!A_h(x, z)$} پکار
ده چې~\LR{$h$} تمثيل کړي۔ پام وکړئ چې~\LR{$h(x) = z$} هله او يوازې هله چې
داسې~\LR{$y$} شته چې \LR{$z = f(y)$} او~\LR{$y = g(x)$} دواړه وي۔ (که
\LR{$h(x) = z$} وي، نو~\LR{$g(x)$} داسې~\LR{$y$} دے؛ که داسې~\LR{$y$} موجود وي، نو
څرنګه چې \LR{$y = g(x)$} او~\LR{$z = f(y)$} دي، \LR{$z = f(g(x))$}۔) له دې څخه
ښکاري چې \LR{$\lexists[y][(!A_g(x, y) \land !A_f(y,
  z))]$} د~\LR{$!A_h(x, z)$} دپاره ښه نوماند دے۔ يوازې بايد دا تصديق کړو
چې~\LR{$\Th{Q}$} اړوند فارمولونه ثابتوي۔
\olpseganchor{OLP-0295-B006}{end}

\olpseganchor{OLP-0295-B007}{start}
\begin{prop}
\ollabel{prop:rep1}
که~\LR{$h(n) = m$} وي، نو \LR{$\Th{Q} \Proves !A_h(\num{n}, \num{m})$}۔
\end{prop}
\olpseganchor{OLP-0295-B007}{end}

\olpseganchor{OLP-0295-B008}{start}
\begin{proof}
فرض کړئ~\LR{$h(n) = m$}، يعنې~\LR{$f(g(n)) = m$}۔ \LR{$k = g(n)$} ونيسئ۔ بيا
\begin{align*}
  \Th{Q} & \Proves !A_g(\num{n}, \num{k})
  \intertext{\RL{ځکه~\LR{$!A_g$}، \LR{$g$} تمثيلوي، او}}
  \Th{Q} & \Proves !A_f(\num{k}, \num{m})
  \intertext{\RL{ځکه~\LR{$!A_f$}، \LR{$f$} تمثيلوي۔ نو}}
  \Th{Q} & \Proves !A_g(\num{n}, \num{k}) \land !A_f(\num{k}, \num{m})
  \intertext{\RL{او په پايله کښې دا هم}}
  \Th{Q} & \Proves \lexists[y][(!A_g(\num{n}, y) \land !A_f(y, \num{m}))],
\end{align*}
يعنې \LR{$\Th{Q} \Proves !A_h(\num{n}, \num{m})$}۔
\end{proof}
\olpseganchor{OLP-0295-B008}{end}

\olpseganchor{OLP-0295-B009}{start}
\begin{prop}
\ollabel{prop:rep2}
که~\LR{$h(n) = m$} وي، نو \LR{$\Th{Q} \Proves \lforall[z][(!A_h(\num{n}, z) \lif
  z = \num{m})]$}۔
\end{prop}
\olpseganchor{OLP-0295-B009}{end}

\olpseganchor{OLP-0295-B010}{start}
\begin{proof}
فرض کړئ~\LR{$h(n) = m$}، يعنې~\LR{$f(g(n)) = m$}۔ \LR{$k = g(n)$} ونيسئ۔ بيا
\begin{align*}
  \Th{Q} & \Proves \lforall[y][(!A_g(\num{n}, y) \lif \eq[y][\num{k}])]
  \intertext{\RL{ځکه~\LR{$!A_g$}، \LR{$g$} تمثيلوي، او}}
  \Th{Q} & \Proves \lforall[z][(!A_f(\num{k}, z) \lif \eq[z][\num{m}])]
  \intertext{\RL{ځکه~\LR{$!A_f$}، \LR{$f$} تمثيلوي۔ يوازې د لږ منطق په کارولو
    ښودلے شو چې دا هم}}
  \Th{Q} & \Proves \lforall[z][(\lexists[y][(!A_g(\num{n}, y) \land
      !A_f(y, z))] \lif \eq[z][\num{m}])].
\end{align*}
يعنې \LR{$\Th{Q} \Proves \lforall[y][(!A_h(\num n, y) \lif \eq[y][\num m])]$}۔
\end{proof}
\olpseganchor{OLP-0295-B010}{end}

\olpseganchor{OLP-0295-B011}{start}
همدا مفکوره په هغې زياتې پېچلې قضيې کښې هم کار کوي چې~\LR{$f$} او
\LR{$g_i$} له~\LR{$1$} څخه زيات ځایونه لري۔
\olpseganchor{OLP-0295-B011}{end}

\olpseganchor{OLP-0295-B012}{start}
\begin{prop}
\ollabel{prop:rep-composition}
که~\LR{$!A_f(y_0, \dots, y_{k-1}, z)$} په~\LR{$\Th{Q}$} کښې
\LR{$f(y_0, \dots, y_{k-1})$} تمثيلوي، او
\LR{$!A_{g_i}(x_0, \dots, x_{l-1}, y)$} په~\LR{$\Th{Q}$} کښې
\LR{$g_i(x_0, \dots, x_{l-1})$} تمثيلوي، نو
\begin{multline*}
  \lexists[y_0\dots][\lexists[y_{k-1}][(!A_{g_0}(x_0,\dots,x_{l-1},y_0) \land
      \dots \land {}]]\\
  !A_{g_{k-1}}(x_0,\dots,x_{l-1},y_{k-1}) \land !A_f(y_0,\dots,y_{k-1},z))
\end{multline*}
دا تمثيلوي:
\[
h(x_0, \dots, x_{l-1}) = f(g_0(x_0, \dots, x_{l-1}), \dots, g_{k-1}(x_0,
\dots, x_{l-1})).
\]
\end{prop}
\olpseganchor{OLP-0295-B012}{end}

\olpseganchor{OLP-0295-B013}{start}
\begin{proof}
تمرین۔
\end{proof}
\olpseganchor{OLP-0295-B013}{end}

\olpseganchor{OLP-0295-B014}{start}
\begin{prob}
د \olref[inc][req][cmp]{prop:rep1} او
\olref[inc][req][cmp]{prop:rep2} ثبوتونه د لارښود په توګه وکاروئ او
د \olref[inc][req][cmp]{prop:rep-composition} ثبوت په تفصيل ترسره
کړئ۔
\end{prob}
\textbf{د ژباړې د نسخې سمون (OLCMP-085):} سرچينې په تمرين کښې
دويمه قضيه دوه ځله راجع کړې وه؛ دلته د دليل دواړه مخکښ پړاوونه په
خپل سم ترتيب راجع شوي دي۔
\olpseganchor{OLP-0295-B014}{end}

\olpseganchor{OLP-0296-B004}{start}
\olfileid{inc}{req}{min}
\olsection{منظم اقل موندل په~\LR{$\Th{Q}$} کښې تمثيلېدونکے دے}
\olpseganchor{OLP-0296-B004}{end}

\olpseganchor{OLP-0296-B005}{start}
بې‌حده پلټنه وګورو۔ فرض کړئ~\LR{$g(x, z)$} منظمه او په~\LR{$\Th{Q}$} کښې
تمثيلېدونکې ده، د بېلګې په توګه د~فارمول~\LR{$!A_g(x, z, y)$} په
وسيله۔ فرض کړئ~\LR{$f$} د~\LR{$f(z) = \umin{x}{[g(x, z) = 0]}$} په ډول تعريف
شوې ده۔ غواړو داسې فارمول~\LR{$!A_f(z, y)$} ومومو چې~\LR{$f$} تمثيل
کړي۔ د~\LR{$f(z)$} ارزښت هغه عدد~\LR{$x$} دے چې (الف)~\LR{$g(x, z) = 0$} پوره
کوي او (ب) تر ټولو وړوکے داسې عدد وي، يعنې د هر~\LR{$w < x$} دپاره
\LR{$g(w, z) \neq 0$} وي۔ نو دا طبيعي انتخاب دے:
\[
!A_f(z,y) \ident !A_g(y, z, \Obj 0) \land \lforall[w][(w < y \lif \lnot
  !A_g(w, z, \Obj 0))].
\]
البته، په عامه قضيه کښې بايد~\LR{$z$} په~\LR{$z_0$}، \dots،~\LR{$z_k$} بدل کړو۔
\olpseganchor{OLP-0296-B005}{end}

\olpseganchor{OLP-0296-B006}{start}
ثبوت به يو ځل بيا د هغو څيزونو په اړه څو لمې وکاروي چې~\LR{$\Th{Q}$}
ئې د ثابتولو دپاره پوره ځواک لري۔
\olpseganchor{OLP-0296-B006}{end}

\olpseganchor{OLP-0296-B007}{start}
\begin{lem}
\ollabel{lem:succ} د هر~ثابت سمبول~\LR{$a$} او هر طبيعي عدد~\LR{$n$} دپاره،
\[
\Th{Q} \Proves \eq[(a' + \num n)][(a + \num n)'].
\]
\end{lem}
\olpseganchor{OLP-0296-B007}{end}

\olpseganchor{OLP-0296-B008}{start}
\begin{proof}
ثبوت، د معمول په شان، پر~\LR{$n$} په استقرا دے۔ په بنسټيزه قضيه کښې
\LR{$n = 0$}؛ بايد وښيو چې~\LR{$\Th{Q}$}،
\LR{$\eq[(a' + \Obj 0)][(a + \Obj 0)']$} ثابتوي۔ خو لرو:
\begin{align}
  \Th{Q} & \Proves \eq[(a' + \Obj 0)][a'] \quad \text{\RL{د بديهي اصل~\LR{$Q_4$} له مخې}}
  \ollabel{step1}\\
  \Th{Q} & \Proves  \eq[(a + \Obj 0)][a] \quad \text{\RL{د بديهي اصل~\LR{$Q_4$} له مخې}}
  \ollabel{step2} \\
  \Th{Q} & \Proves \eq[(a + \Obj 0)'][a'] \quad \text{\RL{د \olref{step2} له مخې}}
  \ollabel{step3} \\
  \Th{Q} & \Proves \eq[(a' + \Obj 0)][(a + \Obj 0)'] \quad
  \text{\RL{د \olref{step1} او \olref{step3} له مخې}}\notag
\end{align}
د استقرا په پړاو کښې فرضولے شو چې ښودلې مو ده:
\LR{$\Th{Q} \Proves \eq[(a' + \num n)][(a + \num n)']$}۔ ځکه چې
\LR{$\num{n+1}$}، \LR{$\num{n}'$} دے، بايد وښيو چې~\LR{$\Th{Q}$}،
\LR{$\eq[(a' + \num{n}')][(a + \num{n}')']$} ثابتوي۔ لرو:
\begin{align}
  \Th{Q} & \Proves \eq[(a' + \num n')][(a' + \num n)'] \quad
  \text{\RL{د بديهي اصل~\LR{$!Q_5$} له مخې}} \ollabel{step5}\\
  \Th{Q} & \Proves \eq[(a' + \num n)'][(a + \num n')'] \quad
  \text{\RL{د استقرايي فرضيې او بديهي اصل~\LR{$!Q_5$} له مخې}} \ollabel{step6}\\
  \Th{Q} & \Proves \eq[(a' + \num n')][(a + \num n')'] \quad
  \text{\RL{د \olref{step5} او \olref{step6} له مخې۔}} \notag
\end{align}
\textbf{د ژباړې د نسخې سمون (OLCMP-086):} سرچينې په دويمه کرښه
کښې هماغه وروستۍ ادعا د استقرايي فرضيې په نوم تکرار کړې او په وروستۍ
کرښه کښې ئې يوازې منځنۍ مساوات اخيستې وه۔ دلته منځنۍ مساوات لومړے
له استقرايي فرضيې او د جمع له بديهي اصل څخه، او وروستۍ ادعا بيا د
دواړو کرښو له تعدي څخه اخيستل شوې ده۔
\end{proof}
\olpseganchor{OLP-0296-B008}{end}

\olpseganchor{OLP-0296-B009}{start}
بيا يادول په کار دي چې دا له دې خبرې کمزورې ده چې~\LR{$\Th{Q}$}،
\LR{$\lforall[x][\lforall[y][(x' + y) = (x + y)']]$} ثابتوي۔ که څه هم
دا تړلے فارمول په~\LR{$\Struct{N}$} کښې رښتينې ده،~\LR{$\Th{Q}$} ئې نۀ
ثابتوي۔
\olpseganchor{OLP-0296-B009}{end}

\olpseganchor{OLP-0296-B010}{start}
\begin{lem}
\ollabel{lem:less-zero}
\LR{$\Th{Q} \Proves \lforall[x][\lnot x < \Obj 0]$}.
\end{lem}
\olpseganchor{OLP-0296-B010}{end}

\olpseganchor{OLP-0296-B011}{start}
\begin{proof}
ثبوت په غيرصوري ډول ورکوو (يعنې يوازې دومره نښې ورکوو چې صوري
اشتقاق څنګه جوړېږي)۔
\olpseganchor{OLP-0296-B011}{end}

\olpseganchor{OLP-0296-B012}{start}
د خپل سري~\LR{$a$} دپاره بايد~\LR{$\lnot a < \Obj 0$} ثابت کړو۔ د~\LR{$<$} د
تعريف له مخې، بايد په~\LR{$\Th{Q}$} کښې
\LR{$\lnot \lexists[y][\eq[(y'+a)][\Obj 0]]$} ثابت کړو۔
\LR{$\lexists[y][\eq[(y'+a)][\Obj 0]]$} فرضوو او تناقض ثابتوو۔ فرض کړئ
\LR{$\eq[(b'+a)][\Obj 0]$}۔ د~\LR{$!Q_3$} په کارولو لرو:
\LR{$\eq[a][\Obj 0] \lor \lexists[y][\eq[a][y']]$}۔ قضيې بېلوو۔
\olpseganchor{OLP-0296-B012}{end}

\olpseganchor{OLP-0296-B013}{start}
لومړۍ قضيه:~\LR{$\eq[a][\Obj 0]$} صدق کوي۔ له
\LR{$\eq[(b'+a)][\Obj 0]$} څخه~\LR{$\eq[(b' + \Obj 0)][\Obj 0]$} لرو۔ د
\LR{$\Th{Q}$} د بديهي اصل~\LR{$!Q_4$} له مخې،
\LR{$\eq[(b' + \Obj 0)][b']$} لرو، او له دې امله
\LR{$\eq[b'][\Obj 0]$}۔ خو د بديهي اصل~\LR{$!Q_2$} له مخې
\LR{$\eq/[b'][\Obj 0]$} هم لرو، چې تناقض دے۔
\olpseganchor{OLP-0296-B013}{end}

\olpseganchor{OLP-0296-B014}{start}
دويمه قضيه: د کوم~\LR{$c$} دپاره~\LR{$\eq[a][c']$}۔ خو بيا
\LR{$\eq[(b' + c')][\Obj 0]$} لرو۔ د بديهي اصل~\LR{$!Q_5$} له مخې،
\LR{$\eq[(b' + c)'][\Obj 0]$} لرو، چې يو ځل بيا له بديهي اصل~\LR{$Q_2$} سره
تناقض لري۔
\end{proof}
\olpseganchor{OLP-0296-B014}{end}

\olpseganchor{OLP-0296-B015}{start}
\begin{lem}
\ollabel{lem:less-nsucc}
د هر طبيعي عدد~\LR{$n$} دپاره،
  \[
  \Th{Q} \Proves
  \lforall[x][(x < \num {n+1} \lif (\eq[x][\Obj 0] \lor \dots \lor
    \eq[x][\num n]))].
  \]
\end{lem}
\olpseganchor{OLP-0296-B015}{end}

\olpseganchor{OLP-0296-B016}{start}
\begin{proof}
پر~\LR{$n$} استقرا کاروو۔ لومړے بنسټيزه قضيه وګورو، چې~\LR{$n = 0$} وي۔ په
دې قضيه کښې بايد د خپل سري~\LR{$a$} دپاره
\LR{$a < \num 1 \lif \eq[a][\Obj 0]$} وښيو۔ فرض کړئ~\LR{$a < \num 1$}۔ بيا
د~\LR{$<$} د تعريفوونکي بديهي اصل له مخې،
\LR{$\lexists[y][\eq[(y'+a)][\Obj 0']]$} لرو (ځکه~\LR{$\num 1 \ident
\Obj 0'$})۔
\olpseganchor{OLP-0296-B016}{end}

\olpseganchor{OLP-0296-B017}{start}
فرض کړئ~\LR{$b$} دا ځانګړنه لري، يعنې
\LR{$\eq[(b'+a)][\Obj 0']$} لرو۔ بايد~\LR{$\eq[a][\Obj 0]$} وښيو۔ د بديهي
اصل~\LR{$!Q_3$} له مخې، يا~\LR{$\eq[a][\Obj 0]$} لرو يا داسې~\LR{$c$} شته چې
\LR{$\eq[a][c']$}۔ په لومړۍ قضيه کښې د ښودلو څه نشته۔ نو
\LR{$\eq[a][c']$} فرض کړئ۔ بيا~\LR{$\eq[(b' + c')][\Obj 0']$} لرو۔ د
\LR{$\Th{Q}$} د بديهي اصل~\LR{$!Q_5$} له مخې،
\LR{$\eq[(b'+c)'][\Obj 0']$} لرو۔ د بديهي اصل~\LR{$!Q_1$} له مخې،
\LR{$\eq[(b' + c)][\Obj 0]$} لرو۔ خو د بديهي اصل~\LR{$!Q_8$} له مخې د دې معنا
دا ده چې~\LR{$c < \Obj 0$}، چې له~\olref{lem:less-zero} سره تناقض لري۔
\olpseganchor{OLP-0296-B017}{end}

\olpseganchor{OLP-0296-B018}{start}
اوس د استقرا پړاو ته راځو۔ د~\LR{$n$} قضيه فرضوو او د~\LR{$n+1$} قضيه
ثابتوو۔ نو فرض کړئ~\LR{$a < \num {n+2}$}۔ يو ځل بيا د~\LR{$!Q_3$} په کارولو
دوه قضيې بېلولے شو:~\LR{$\eq[a][\Obj 0]$}، او د کوم~\LR{$b$} دپاره
\LR{$\eq[a][b']$}۔ په لومړۍ قضيه کښې،
\LR{$\eq[a][\Obj 0] \lor \dots \lor \eq[a][\num{n+1}]$} په ښکاره لازمېږي۔
په دويمه قضيه کښې~\LR{$b' < \num {n+2}$} لرو، يعنې
\LR{$b' < \num{n+1}'$}۔ د بديهي اصل~\LR{$!Q_8$} له مخې، د کوم~\LR{$c$} دپاره
\LR{$\eq[(c'+b')][\num{n+1}']$}۔ د بديهي اصل~\LR{$!Q_5$} له مخې،
\LR{$\eq[(c'+b)'][\num{n+1}']$}۔ د بديهي اصل~\LR{$!Q_1$} له مخې،
\LR{$\eq[(c'+b)][\num{n+1}]$}، او نو د بديهي اصل~\LR{$!Q_8$} له مخې
\LR{$b < \num{n+1}$}۔ د استقرايي فرضيې له مخې،
\LR{$\eq[b][\Obj 0] \lor \dots \lor \eq[b][\num{n}]$}۔ له دې څخه د منطق
په وسيله~\LR{$\eq[b'][\Obj 0'] \lor \dots \lor \eq[b'][\num{n}']$} اخلو،
او ځکه چې~\LR{$\eq[a][b']$}، نو
\LR{$\eq[a][\num{1}] \lor \dots \lor \eq[a][\num{n+1}]$}۔
\end{proof}
\olpseganchor{OLP-0296-B018}{end}

\olpseganchor{OLP-0296-B019}{start}
\begin{lem}
  \ollabel{lem:trichotomy} د هر طبيعي عدد~\LR{$m$} دپاره،
  \[
  \Th{Q} \Proves
  \lforall[y][((y < \num{m} \lor \num{m} < y) \lor \eq[y][\num{m}])].
  \]
\end{lem}
\olpseganchor{OLP-0296-B019}{end}

\olpseganchor{OLP-0296-B020}{start}
\begin{proof}
پر~\LR{$m$} په استقرا۔ لومړے د~\LR{$m=0$} قضيه وګورئ۔ د~\LR{$!Q_3$} له مخې
\LR{$\Th{Q} \Proves
\lforall[y][(\eq[y][\Obj 0] \lor \lexists[z][\eq[y][z']])]$}۔
\LR{$a$} خپل سرے ونيسئ۔ بيا يا~\LR{$\eq[a][\Obj 0]$} وي يا د کوم~\LR{$b$} دپاره
\LR{$\eq[a][b']$}۔ په لومړۍ قضيه کښې
\LR{$(a < \Obj 0 \lor \Obj 0 < a) \lor \eq[a][\Obj 0]$} هم لرو۔ خو که
\LR{$\eq[a][b']$} وي، نو د~\LR{$\eq$} د منطق له مخې
\LR{$\eq[(b' + \Obj 0)][(a + \Obj 0)]$}۔ د~\LR{$!Q_4$} له مخې
\LR{$\eq[(a + \Obj 0)][a]$}، نو~\LR{$\eq[(b' + \Obj 0)][a]$} لرو، او له دې
امله~\LR{$\lexists[z][\eq[(z' + \Obj 0)][a]]$}۔ په~\LR{$!Q_8$} کښې د~\LR{$<$} د
تعريف له مخې،~\LR{$\Obj 0 < a$}۔ که~\LR{$\Obj 0 < a$} وي، نو
\LR{$(\Obj 0 < a \lor a < \Obj 0) \lor \eq[a][\Obj 0]$} هم لرو۔
\olpseganchor{OLP-0296-B020}{end}

\olpseganchor{OLP-0296-B021}{start}
اوس فرض کړئ چې لرو:
\begin{align*}
  \Th{Q} & \Proves \lforall[y][((y < \num{m} \lor \num{m} < y) \lor
    \eq[y][\num{m}])]
  \intertext{\RL{او غواړو وښيو:}}
  \Th{Q} & \Proves \lforall[y][((y < \num{m+1} \lor \num{m+1} < y) \lor
    \eq[y][\num{m+1}])]
\end{align*}
\LR{$a$} خپل سرے ونيسئ۔ د~\LR{$!Q_3$} له مخې، يا~\LR{$\eq[a][\Obj 0]$} وي يا د
کوم~\LR{$b$} دپاره~\LR{$\eq[a][b']$}۔ په لومړۍ قضيه کښې، د~\LR{$!Q_4$} له مخې
\LR{$\eq[\num{m}' + a][\num{m+1}]$} لرو، او نو د~\LR{$!Q_8$} له مخې
\LR{$a < \num{m+1}$}۔
\olpseganchor{OLP-0296-B021}{end}

\olpseganchor{OLP-0296-B022}{start}
اوس دويمه قضيه،~\LR{$\eq[a][b']$}، وګورئ۔ د استقرايي فرضيې له مخې،
\LR{$(b < \num{m} \lor \num{m} < b) \lor
    \eq[b][\num{m}]$}۔
\olpseganchor{OLP-0296-B022}{end}

\olpseganchor{OLP-0296-B023}{start}
لومړے منفصل~\LR{$b < \num{m}$} د~\LR{$!Q_8$} له مخې له
\LR{$\lexists[z][\eq[(z' + b)][\num{m}]]$} سره معادل دے۔ فرض کړئ~\LR{$c$}
دا ځانګړنه لري۔ که~\LR{$\eq[(c' + b)][\num{m}]$} وي، نو
\LR{$\eq[(c' + b)'][\num{m}']$} هم وي۔ د~\LR{$!Q_5$} له مخې،
\LR{$\eq[(c' + b)'][(c' + b')]$}۔ له دې امله
\LR{$\eq[(c' + b')][\num{m}']$}۔ پر~\LR{$c'$} د وجودي تعميم په کولو او په دې
پام سره چې~\LR{$\num{m}' \ident \num{m+1}$}،
\LR{$\lexists[u][\eq[(u' + b')][\num{m+1}]]$} اخلو۔ له دې امله، که
\LR{$b < \num{m}$} وي، نو~\LR{$b' < \num{m+1}$}، او نو
\LR{$a < \num{m+1}$}۔
\olpseganchor{OLP-0296-B023}{end}

\olpseganchor{OLP-0296-B024}{start}
اوس فرض کړئ~\LR{$\num{m} < b$}، يعنې
\LR{$\lexists[z][\eq[(z' + \num{m})][b]]$}۔ فرض کړئ~\LR{$c$} داسې~\LR{$z$} دے،
يعنې~\LR{$\eq[(c' + \num{m})][b]$}۔ د منطق له مخې،
\LR{$\eq[(c' + \num{m})'][b']$}۔ د~\LR{$!Q_5$} له مخې،
\LR{$\eq[(c' + \num{m}')][b']$}۔ ځکه چې~\LR{$\eq[a][b']$} او
\LR{$\num{m}' \ident \num{m+1}$}، نو~\LR{$\eq[(c' + \num{m+1})][a]$}۔ د
\LR{$!Q_8$} له مخې،~\LR{$\num{m+1} < a$}۔
\olpseganchor{OLP-0296-B024}{end}

\olpseganchor{OLP-0296-B025}{start}
په پاې کښې،~\LR{$\eq[b][\num{m}]$} فرض کړئ۔ بيا د منطق له مخې
\LR{$\eq[b'][\num{m}']$}، او نو~\LR{$\eq[a][\num{m+1}]$}۔
\olpseganchor{OLP-0296-B025}{end}

\olpseganchor{OLP-0296-B026}{start}
نو د~\LR{$m$} او~\LR{$b$} د قضيې له هر منفصل څخه د~\LR{$m+1$} او~\LR{$a$} اړوند منفصل
تر لاسه کولے شو۔
\end{proof}
\olpseganchor{OLP-0296-B026}{end}

\olpseganchor{OLP-0296-B027}{start}
\begin{prop}
\ollabel{prop:rep-minimization}
که~\LR{$!A_g(x, z, y)$} په~\LR{$\Th{Q}$} کښې~\LR{$g(x, z)$} تمثيلوي، نو
\[
!A_f(z,y) \ident !A_g(y, z, \Obj 0) \land \lforall[w][(w < y \lif \lnot
  !A_g(w, z, \Obj 0))]
\]
\LR{$f(z) = \umin{x}{[g(x, z) = 0]}$} تمثيلوي۔
\end{prop}
\olpseganchor{OLP-0296-B027}{end}

\olpseganchor{OLP-0296-B028}{start}
\begin{proof}
لومړے ښيو چې که~\LR{$f(n) = m$} وي، نو
\LR{$\Th{Q} \Proves !A_f(\num n, \num m)$}؛ يعنې
\begin{align}
  \Th{Q} & \Proves !A_g(\num{m}, \num{n}, \Obj 0) \land \lforall[w][(w <
    \num{m} \lif \lnot !A_g(w, \num{n}, \Obj 0))]. \notag
  \intertext{\RL{ځکه چې~\LR{$!A_g(x, z, y)$}، \LR{$g(x, z)$} تمثيلوي او
    \LR{$g(m, n) = 0$} که~\LR{$f(n) = m$} وي، لرو:}}
\Th{Q} & \Proves !A_g(\num{m}, \num{n}, \Obj 0). \notag
\intertext{\RL{که~\LR{$f(n) = m$} وي، نو د هر~\LR{$k < m$} دپاره~\LR{$g(k, n) \neq 0$}۔ نو}}
\Th{Q} & \Proves \lnot !A_g(\num{k}, \num{n}, \Obj 0). \notag
\intertext{\RL{تر لاسه کوو:}}
\Th{Q} & \Proves \lforall[w][(w < \num{m} \lif \lnot
  !A_g(w, \num{n}, \Obj 0))]. \ollabel{rep-less}
\end{align}
که~\LR{$m = 0$} وي، دا د~\olref{lem:less-zero} له مخې، او که نۀ وي د
\olref{lem:less-nsucc} له مخې راځي۔
\olpseganchor{OLP-0296-B028}{end}

\olpseganchor{OLP-0296-B029}{start}
اوس وښيو چې که~\LR{$f(n) = m$} وي، نو
\LR{$\Th{Q} \Proves
\lforall[y][(!A_f(\num{n}, y) \lif \eq[y][\num{m}])]$}۔ يو ځل بيا
استدلال په غيرصوري ډول لنډوو او صوري کول ئې لوستونکي ته پرېږدو۔
\olpseganchor{OLP-0296-B029}{end}

\olpseganchor{OLP-0296-B030}{start}
فرض کړئ~\LR{$!A_f(\num{n}, b)$}۔ له دې څخه (الف)
\LR{$!A_g(b, \num{n}, \Obj 0)$} او (ب)
\LR{$\lforall[w][(w < b \lif \lnot !A_g(w, \num{n}, \Obj
  0))]$} اخلو۔ د~\olref{lem:trichotomy} له مخې،
\LR{$(b < \num{m} \lor \num{m} < b) \lor \eq[b][\num{m}]$}۔ وبه ښيو چې
\LR{$b < \num{m}$} او~\LR{$\num{m} < b$} دواړه تناقض ته رسېږي۔
\olpseganchor{OLP-0296-B030}{end}

\olpseganchor{OLP-0296-B031}{start}
که~\LR{$\num{m} < b$} وي، نو له~(ب) څخه
\LR{$\lnot !A_g(\num{m}, \num{n}, \Obj 0)$}۔ خو~\LR{$m = f(n)$}، نو
\LR{$g(m, n) = 0$}؛ او ځکه چې~\LR{$!A_g$}، \LR{$g$} تمثيلوي،
\LR{$\Th{Q} \Proves !A_g(\num{m}, \num{n}, \Obj 0)$}۔ نو تناقض لرو۔
\olpseganchor{OLP-0296-B031}{end}

\olpseganchor{OLP-0296-B032}{start}
اوس~\LR{$b < \num{m}$} فرض کړئ۔ بيا ځکه چې د~\olref{rep-less} له مخې
\LR{$\Th{Q} \Proves \lforall[w][(w <
  \num{m} \lif \lnot !A_g(w, \num{n}, \Obj 0))]$}، نو
\LR{$\lnot !A_g(b, \num{n}, \Obj 0)$} اخلو۔ دا يو ځل بيا له~(الف) سره
تناقض لري۔
\end{proof}
\olpseganchor{OLP-0296-B032}{end}

\olpseganchor{OLP-0297-B004}{start}
\olfileid{inc}{req}{crq}
\olsection{محاسبه کېدونکې تابعې په~\LR{$\Th{Q}$} کښې تمثيلېدونکې دي}
\olpseganchor{OLP-0297-B004}{end}

\olpseganchor{OLP-0297-B005}{start}
\begin{thm}
هره محاسبه کېدونکې تابعه په~\LR{$\Th{Q}$} کښې تمثيلېدونکې ده۔
\end{thm}
\olpseganchor{OLP-0297-B005}{end}

\olpseganchor{OLP-0297-B006}{start}
\begin{proof}
د کره والي دپاره، او د چرچ--ټيورينګ د دعوې په کارولو سره، ووايو چې
يوه تابعه هله او يوازې هله محاسبه کېدونکې ده چې عمومي بازګشتي وي۔
عمومي بازګشتي تابعې هغه دي چې له صفر تابع~\LR{$\Zero$}، تالي
تابع~\LR{$\Succ$}، او پروجکشن تابع~\LR{$\Proj{n}{i}$} څخه د ترکيب، بنسټيز
بازګښت، او منظم اقل موندلو په کارولو تعريفېدے شي۔ د
\olref[pri]{lem:prim-rec} له مخې، هره تابعه~\LR{$h$} چې له~\LR{$f$} او~\LR{$g$}
څخه تعريفېدے شي، له~\LR{$f$}، \LR{$g$}، او~\LR{$\Zero$}، \LR{$\Succ$}،
\LR{$\Proj{n}{i}$}، \LR{$\Add$}، \LR{$\Mult$}، \LR{$\Char{=}$} څخه هم د ترکيب او منظم
اقل موندلو په کارولو تعريفېدے شي۔ له دې امله، يوه تابعه هله او يوازې
هله عمومي بازګشتي ده چې له~\LR{$\Zero$}، \LR{$\Succ$}، \LR{$\Proj{n}{i}$}،
\LR{$\Add$}، \LR{$\Mult$}، \LR{$\Char{=}$} څخه د ترکيب او منظم اقل موندلو په
کارولو تعريفېدے شي۔
\olpseganchor{OLP-0297-B006}{end}

\olpseganchor{OLP-0297-B007}{start}
سربېره پر دې، ښودلې مو ده چې د پام وړ بنسټيزې تابعې په~\LR{$\Th{Q}$}
کښې تمثيلېدونکې دي
(\cref{inc:req:bre:prop:rep-zero,inc:req:bre:prop:rep-succ,inc:req:bre:prop:rep-proj,inc:req:bre:prop:rep-id,inc:req:bre:prop:rep-add,inc:req:bre:prop:rep-mult})،
او هره هغه تابعه چې له تمثيلېدونکو تابعو څخه په ترکيب يا منظم اقل
موندلو تعريف شوې وي
(\olref[cmp]{prop:rep-composition}،
\olref[min]{prop:rep-minimization}) هم تمثيلېدونکې ده۔ نو هره عمومي
بازګشتي تابعه په~\LR{$\Th{Q}$} کښې تمثيلېدونکې ده۔
\end{proof}
\olpseganchor{OLP-0297-B007}{end}

\olpseganchor{OLP-0297-B008}{start}
\begin{explain}
موږ ښودلې ده چې د محاسبه کېدونکو تابعو سټ د هغو تابعو د سټ په توګه
ځانګړے کېدے شي چې په~\LR{$\Th{Q}$} کښې تمثيلېدونکې دي۔ په حقيقت کښې
ثبوت تر دې عام دے۔ د تمثيلېدنې له تعريف څخه په اسانۍ ليدل کېږي چې
هره هغه تيوري چې~\LR{$\Th{Q}$} غځوي (يا~\LR{$\Th{Q}$} پکې تفسيرېدے شي)،
محاسبه کېدونکې تابعې تمثيلولے شي۔ خو برعکس، په هر هغه
اشتقاق نظام کښې چې د اشتقاق مفهوم ئې محاسبه کېدونکے
وي، هره تمثيلېدونکې تابعه محاسبه کېدونکې ده۔ نو د بېلګې په توګه، د
محاسبه کېدونکو تابعو سټ د هغو تابعو د سټ په توګه ځانګړے کېدے شي چې
د پيانو په حساب، يا ان د زرمیلو--فرېنکل په سټ تيورۍ کښې تمثيلېدونکې
دي۔ لکه ګوډل چې يادونه کړې، دا يو څه هېښوونکې ده۔ وبه وينو چې د
ثبوت‌وړتيا په برخه کښې پوښتنې د ټاکلې تيورۍ پر وړاندې ډېرې حساسې
دي؛ په ټوليز ډول هر څومره چې بديهي اصول پياوړي وي، هغومره ډېر څه
ثابتولے شئ۔ خو د بديهي تيوريو په يوه پراخه ډله کښې تمثيلېدونکې
تابعې کټ مټ محاسبه کېدونکې تابعې دي؛ تر هغه چې تيورۍ
بديهي کېدونکې وي، پياوړې تيورۍ زياتې تابعې نۀ تمثيلوي۔
\end{explain}
\olpseganchor{OLP-0297-B008}{end}

\olpseganchor{OLP-0298-B004}{start}
\olfileid{inc}{req}{rel}
\olpseganchor{OLP-0298-B004}{end}

\olpseganchor{OLP-0298-B005}{start}
\olsection{د اړيکو تمثيلول}
\olpseganchor{OLP-0298-B005}{end}

\olpseganchor{OLP-0298-B006}{start}
دا بيان کړو چې يوه \emph{اړيکه} تمثيلېدونکې وي څه معنا لري۔
\olpseganchor{OLP-0298-B006}{end}

\olpseganchor{OLP-0298-B007}{start}
\begin{defn}
\ollabel{defn:representing-relations} پر طبيعي عددونو اړيکه
\LR{$R(x_0,\dots,x_k)$} هله \emph{په~\LR{$\Th{Q}$} کښې تمثيلېدونکې} ده چې
داسې فارمول~\LR{$!A_R(x_0,\dots,x_k)$} وي چې هر کله
\LR{$R(n_0,\dots,n_k)$} رښتيا وي، \LR{$\Th{Q}$}،
\LR{$!A_R(\num{n_0},\dots,\num{n_k})$} ثابت کړي، او هر کله
\LR{$R(n_0,\dots,n_k)$} ناسم وي، \LR{$\Th{Q}$}،
\LR{$\lnot !A_R(\num{n_0}, \dots, \num{n_k})$} ثابت کړي۔
\end{defn}
\olpseganchor{OLP-0298-B007}{end}

\olpseganchor{OLP-0298-B008}{start}
\begin{thm}
\ollabel{thm:representing-rels} يوه اړيکه په~\LR{$\Th{Q}$} کښې هله او
يوازې هله تمثيلېدونکې ده چې محاسبه کېدونکې وي۔
\end{thm}
\olpseganchor{OLP-0298-B008}{end}

\olpseganchor{OLP-0298-B009}{start}
\begin{proof}
د مخکښ لوري دپاره فرض کړئ اړيکه~\LR{$R(x_0,\dots,x_k)$} په
فارمول~\LR{$!A_R(x_0,\dots,x_k)$} تمثيلېږي۔ د~\LR{$R$} د محاسبې الګوريتم دا
دے: پر ننوتونو~\LR{$n_0$}، \dots،~\LR{$n_k$} په هممهاله ډول د
\LR{$!A_R(\num{n_0}, \dots, \num{n_k})$} د ثبوت او د
\LR{$\lnot !A_R(\num{n_0}, \dots, \num{n_k})$} د ثبوت پلټنه وکړئ۔ زموږ
د فرض له مخې پلټنه خامخا يو له دغو دواړو مومي؛ که لومړے وي، «هو»
راپور کړئ، او که نۀ، «نه» راپور کړئ۔
\olpseganchor{OLP-0298-B009}{end}

\olpseganchor{OLP-0298-B010}{start}
په بل لوري کښې فرض کړئ~\LR{$R(x_0, \dots, x_k)$} محاسبه کېدونکې ده۔ د
تعريف له مخې دا په دې معنا ده چې تابعه~\LR{$\Char{R}(x_0, \dots, x_k)$}
محاسبه کېدونکې ده۔ د~\olref[int]{thm:representable-iff-comp} له
مخې، \LR{$\Char{R}$} په کوم فارمول، مثلاً
\LR{$!A_{\Char{R}}(x_0, \dots, x_k, y)$}، تمثيلېږي۔
\LR{$!A_R(x_0, \dots, x_k)$} هغه فارمول ونيسئ چې
\LR{$!A_{\Char{R}}(x_0, \dots, x_k, \num{1})$} وي۔ بيا د هر~\LR{$n_0$}،
\dots،~\LR{$n_k$} دپاره، که~\LR{$R(n_0, \dots, n_k)$} رښتيا وي، نو
\LR{$\Char{R}(n_0, \dots, n_k) = 1$}؛ په دې حالت کښې~\LR{$\Th{Q}$}،
\LR{$!A_{\Char{R}}(\num{n_0}, \dots, \num{n_k}, \num{1})$} ثابتوي، او
نو~\LR{$\Th{Q}$}، \LR{$!A_R(\num{n_0}, \dots, \num{n_k})$} ثابتوي۔ له بلې
خوا، که~\LR{$R(n_0, \dots, n_k)$} ناسم وي، نو
\LR{$\Char{R}(n_0, \dots, n_k) = 0$}۔ دا په دې معنا ده چې~\LR{$\Th{Q}$}
لاندې ثابتوي:
\[
\lforall[y][(!A_{\Char{R}}(\num{n_0}, \dots, \num{n_k}, y) \lif y =
  \num{0})].
\]
څرنګه چې~\LR{$\Th{Q}$}، \LR{$\eq/[\num{0}][\num{1}]$} ثابتوي، \LR{$\Th{Q}$}،
\LR{$\lnot !A_{\Char{R}}(\num{n_0}, \dots, \num{n_k}, \num{1})$}
ثابتوي، او نو~\LR{$\lnot !A_R(\num{n_0}, \dots, \num{n_k})$} هم
ثابتوي۔
\end{proof}
\olpseganchor{OLP-0298-B010}{end}

\olpseganchor{OLP-0298-B011}{start}
\begin{prob}
وښيئ چې که~\LR{$R$} په~\LR{$\Th{Q}$} کښې تمثيلېدونکې وي، نو~\LR{$\Char{R}$} هم
همداسې ده۔
\end{prob}
\olpseganchor{OLP-0298-B011}{end}

\olpseganchor{OLP-0299-B004}{start}
\olfileid{inc}{req}{und}
\olsection{ناپرېکړتيا}
\olpseganchor{OLP-0299-B004}{end}

\olpseganchor{OLP-0299-B005}{start}
تيوري~\LR{$\Th{T}$} هغه وخت \emph{ناپرېکړه‌وړ} بولو چې داسې محاسبوي
کړنلاره نۀ وي چې له متناهي شمېر ګامونو وروسته د~\LR{$\Th{T}$} د ژبې د
هرې جملې~\LR{$!A$} دپاره په باوري ډول دې پوښتنې ته سم ځواب ورکړي:
«آيا~\LR{$\Th{T}$}، \LR{$!A$} ثابتوي؟» نو~\LR{$\Th{Q}$} به هله او يوازې هله د
پرېکړې وړ وه چې داسې محاسبوي کړنلاره وای چې د حساب د ژبې د ورکړل شوې
جملې~\LR{$!A$} دپاره پرېکړه وکړي چې~\LR{$\Th{Q} \Proves !A$} ده که نۀ۔ دا
پوښتنه داسې لا کره کولے شو: آيا اړيکه~\LR{$\Prov[\Th{Q}](y)$}، چې پر~\LR{$y$}
هله صدق کوي چې~\LR{$y$} په~\LR{$\Th{Q}$} کښې د ثابتېدونکې جملې ګوډل شمېره
وي، بازګشتي ده؟ ځواب «نه» دے۔
\olpseganchor{OLP-0299-B005}{end}

\olpseganchor{OLP-0299-B006}{start}
\begin{thm}
\LR{$\Th{Q}$} ناپرېکړه‌وړ ده؛ يعنې اړيکه
\[
\Prov[\Th{Q}](y) \defiff \fn{Sent}(y) \land
\lexists[x][\Prf[\Th{Q}](x, y)]
\]
بازګشتي نۀ ده۔
\end{thm}
\olpseganchor{OLP-0299-B006}{end}

\olpseganchor{OLP-0299-B007}{start}
\begin{proof}
فرض کړئ بازګشتي وای۔ بيا د درېدنې مسئله داسې حلولے شو: د ورکړو
\LR{$e$} او~\LR{$n$} دپاره پوهېږو چې~\LR{$\cfind{e}(n) \fdefined$} هله او يوازې
هله چې داسې~\LR{$s$} وي چې~\LR{$T(e, n, s)$}، چې~\LR{$T$} د
\olref[cmp][rec][nft]{thm:kleene-nf} څخه د کلېني پريديکات دے۔ ځکه
\LR{$T$} بنسټيزه بازګشتي ده، په~\LR{$\Th{Q}$} کښې د کوم فارمول~\LR{$!B_T$} په
وسيله تمثيلېږي؛ يعنې~\LR{$\Th{Q} \Proves
!B_T(\num{e}, \num{n}, \num{s})$} هله او يوازې هله چې~\LR{$T(e,n,s)$}۔
که~\LR{$\Th{Q} \Proves !B_T(\num{e}, \num{n}, \num{s})$}، نو
\LR{$ \Th{Q} \Proves \lexists[y][!B_T(\num{e}, \num{n}, y)]$} هم۔ که
داسې~\LR{$s$} نۀ وي، نو د هر~\LR{$s$} دپاره
\LR{$\Th{Q} \Proves \lnot !B_T(\num{e}, \num{n}, \num{s})$}۔ خو~\LR{$\Th{Q}$}،
\LR{$\omega$}-سازګاره ده؛ يعنې که د هر~\LR{$n \in \Nat$} دپاره
\LR{$\Th{Q} \Proves \lnot !A(\num{n})$}، نو
\LR{$\Th{Q} \Proves/ \lexists[y][!A(y)]$}۔ دا ځکه پوهېږو چې د~\LR{$\Th{Q}$}
بديهي اصول په معياري مدل~\LR{$\Struct{N}$} کښې رښتوني دي۔ نو
\LR{$\Th{Q} \Proves/ \lexists[y][!B_T(\num{e}, \num{n}, y)]$}۔ په بله
وينا، \LR{$\Th{Q} \Proves \lexists[y][!B_T(\num{e}, \num{n}, y)]$}
هله او يوازې هله چې داسې~\LR{$s$} وي چې~\LR{$T(e,n,s)$}؛ يعنې هله او يوازې
هله چې~\LR{$\cfind{e}(n) \fdefined$}۔ له~\LR{$e$} او~\LR{$n$} څخه
\LR{$\Gn{\lexists[y][!B_T(\num{e}, \num{n}, y)]}$} محاسبه کولے شو؛
\LR{$g(e,n)$} هغه بنسټيزه بازګشتي تابعه ونيسئ چې دا کار کوي۔ نو
\[
h(e, n) =
\begin{cases}
1 & \text{\RL{که~\LR{$\Prov[\Th{Q}](g(e, n))$} وي}}\\
0 & \text{\RL{که نۀ وي}}.
\end{cases}
\]
دا به ښودلې وای چې که~\LR{$\Prov[\Th{Q}]$} بازګشتي وي، \LR{$h$} هم بازګشتي
ده۔ خو د~\olref[cmp][rec][hlt]{thm:halting-problem} له مخې~\LR{$h$}
بازګشتي نۀ ده، نو~\LR{$\Prov[\Th{Q}]$} هم بازګشتي نۀ شي کېدے۔
\end{proof}
\olpseganchor{OLP-0299-B007}{end}

\olpseganchor{OLP-0299-B008}{start}
\begin{cor}
د لومړۍ درجې منطق ناپرېکړه‌وړ دے۔
\end{cor}
\olpseganchor{OLP-0299-B008}{end}

\olpseganchor{OLP-0299-B009}{start}
\begin{proof}
که د لومړۍ درجې منطق د پرېکړې وړ وای، په~\LR{$\Th{Q}$} کښې ثبوت‌وړتيا
هم د پرېکړې وړ وه، ځکه~\LR{$\Th{Q} \Proves !A$} هله او يوازې هله چې
\LR{$\Proves !T \lif !A$}، چېرته چې~\LR{$!T$} د~\LR{$\Th{Q}$} د بديهي اصولو
عطف دے۔
\end{proof}
\olpseganchor{OLP-0299-B009}{end}

\olpseganchor{OLP-0300-B004}{start}
\olfileid{inc}{inp}{s1c}
\olsection{د~\texorpdfstring{\LR{$\Sigma_1$}}{سيګما-۱} بشپړتيا}
\olpseganchor{OLP-0300-B004}{end}

\olpseganchor{OLP-0300-B005}{start}
د~\LR{$\Th{Q}$} او د هغې د سازګارو، بديهي کېدونکو غځونو د ناتکميلۍ
سره سره، ليدلي مو دي چې~\LR{$\Th{Q}$} د عددنښو په اړه ډېر بنسټيز
حقيقتونه ثابتوي۔ په حقيقت کښې دا خبره په پام وړ ډول غځېدے شي۔ د دې
د پوهېدو دپاره چې په~\LR{$\Th{Q}$} کښې کوم څيزونه ثابتېدے شي، د
\LR{$\Delta_0$}، \LR{$\Sigma_1$}، او~\LR{$\Pi_1$} فارمولونه مفهومونه معرفي
کوو۔ په نا کره ډول،~\LR{$\Sigma_1$} فارمول د
\LR{$\lexists[x][!B(x)]$} په بڼه وي، چې~\LR{$!B$} يوازې د قضيو د نښلوونکو او
محدودو کمیت ايښوونکو په کارولو جوړ شوے وي۔ وبه ښيو چې که~\LR{$!A$} داسې
\LR{$\Sigma_1$} تړلے فارمول وي چې په~\LR{$\Struct{N}$} کښې رښتينې وي، نو
\LR{$\Th{Q} \Proves !A$} (\olref{thm:sigma1-completeness})۔
\olpseganchor{OLP-0300-B005}{end}

\olpseganchor{OLP-0300-B006}{start}
\begin{defn}
\ollabel{defn:bd-quant}
يوه \emph{محدوده وجودي فارمول} د
\LR{$\lexists[x][(x < t \land !A(x))]$} په بڼه وي، چې~\LR{$t$} هر ترم دے، او
موږ ئې په دوديز ډول~\LR{$\bexists{x < t}{!A(x)}$} ليکو۔
%
يوه \emph{محدوده کلي فارمول} د
\LR{$\lforall[x][(x < t \lif !A(x))]$} په بڼه وي، چې~\LR{$t$} هر ترم دے، او
موږ ئې په دوديز ډول~\LR{$\bforall{x < t}{!A(x)}$} ليکو۔
\end{defn}
\olpseganchor{OLP-0300-B006}{end}

\olpseganchor{OLP-0300-B007}{start}
\begin{defn}
\ollabel{defn:delta0-sigma1-pi1-frm}
يوه فارمول~\LR{$!B$} هله~\LR{$\Delta_0$} ده چې له اټومي
فارمولونه څخه يوازې د قضيو د نښلوونکو او محدودې کمیت ايښودنې په
کارولو جوړه شوې وي۔
%
يوه فارمول~\LR{$!A$} هله~\LR{$\Sigma_1$} ده چې
\LR{$!A \ident \lexists[x][!B(x)]$} وي او~\LR{$!B$}، \LR{$\Delta_0$} وي۔
%
يوه فارمول~\LR{$!A$} هله~\LR{$\Pi_1$} ده چې
\LR{$!A \ident \lforall[x][!B(x)]$} وي او~\LR{$!B$}، \LR{$\Delta_0$} وي۔
\end{defn}
\olpseganchor{OLP-0300-B007}{end}

\olpseganchor{OLP-0300-B008}{start}
\begin{lem}
\ollabel{lem:q-proves-clterm-id} فرض کړئ~\LR{$t$} تړلے ترم دے چې
\LR{$\Value{t}{N} = n$}۔ بيا \LR{$\Th{Q} \Proves \eq[t][\num n]$}۔
\end{lem}
\olpseganchor{OLP-0300-B008}{end}

\olpseganchor{OLP-0300-B009}{start}
\begin{proof}
دا د~\LR{$t$} پر پېچلتيا په استقرا ثابتوو۔ په بنسټيزه قضيه کښې
\LR{$\Value{\Obj 0}{N} = 0$}، او
\LR{$\Th{Q} \Proves \eq[\Obj 0][\num 0]$}، ځکه
\LR{$\num 0 \ident \Obj 0$}۔
%
د استقرايي قضيې دپاره، فرض کړئ~\LR{$t_1$} او~\LR{$t_2$} داسې ترمونه دي چې
\LR{$\Value{t_1}{N} = n_1$}، \LR{$\Value{t_2}{N} = n_2$}،
\LR{$\Th{Q} \Proves \eq[t_1][\num n_1]$}، او
\LR{$\Th{Q} \Proves \eq[t_2][\num n_2]$}۔
\olpseganchor{OLP-0300-B009}{end}

\olpseganchor{OLP-0300-B010}{start}
بيا~\LR{$\Value{(t_1')}{N} = n_1 + 1$}، او د عينيت د لومړي-رتبه قواعدو
له مخې، چې پر استقرايي فرضيه او
فارمول~\LR{$\eq[\num{n_1}'][\num{n_1}']$} پلي شوي، لرو چې
\LR{$\Th{Q} \Proves \eq[t_1'][{\num n_1}']$}؛ نو د عددنښو د تعريف له
مخې \LR{$\Th{Q} \Proves \eq[t_1'][\num{n_1 + 1}]$}۔
\olpseganchor{OLP-0300-B010}{end}

\olpseganchor{OLP-0300-B011}{start}
د جمعو دپاره لرو:
\[
      \Value{(t_1 + t_2)}{N}
    = \Value{t_1}{N} + \Value{t_2}{N}
    = n_1 + n_2.
\]
د استقرايي فرضيې او د عينيت د قواعدو له مخې،
\LR{$\Th{Q} \Proves \eq[t_1 + t_2][\num{n_1} + t_2]$}، او بيا د عينيت
د قواعدو په دويم پلي کولو
\LR{$\Th{Q} \Proves \eq[t_1 + t_2][\num{n_1} + \num{n_2}]$}۔ د
\olref[inc][req][bre]{lem:q-proves-add} له مخې،
\LR{$\Th{Q} \Proves \eq[\num{n_1} + \num{n_2}][\num{n_1 + n_2}]$}،
نو \LR{$\Th{Q} \Proves \eq[t_1 + t_2][\num{n_1 + n_2}]$}۔
\olpseganchor{OLP-0300-B011}{end}

\olpseganchor{OLP-0300-B012}{start}
همداسې استدلال د~\LR{$\times$} دپاره هم کار کوي، چې
\olref[inc][req][bre]{lem:q-proves-mult} پکې کاروو۔
%
ځکه چې دا د حساب ټول تړلي ترمونه رانغاړي، نو د هر داسې تړلي
ترم~\LR{$t$} دپاره چې~\LR{$\Value{t}{N} = n$} وي،
\LR{$\Th{Q} \Proves \eq[t][\num n]$}۔
\end{proof}
\olpseganchor{OLP-0300-B012}{end}

\olpseganchor{OLP-0300-B013}{start}
\begin{prob}
په \olref{lem:q-proves-clterm-id} کښې د~\LR{$\times$} اړوند برخه په
تفصيل ثابته کړئ۔
\end{prob}
\olpseganchor{OLP-0300-B013}{end}

\olpseganchor{OLP-0300-B014}{start}
\begin{lem}
\ollabel{lem:atomic-completeness}
فرض کړئ~\LR{$t_1$} او~\LR{$t_2$} تړلي ترمونه دي۔ بيا
\begin{enumerate}
\item که~\LR{$\Value{t_1}{N} = \Value{t_2}{N}$} وي،
    نو \LR{$\Th{Q} \Proves \eq[t_1][t_2]$}۔
\item که~\LR{$\Value{t_1}{N} \neq \Value{t_2}{N}$} وي،
    نو \LR{$\Th{Q} \Proves \eq/[t_1][t_2]$}۔
\item که~\LR{$\Value{t_1}{N} < \Value{t_2}{N}$} وي،
    نو \LR{$\Th{Q} \Proves t_1 < t_2$}۔
\item که~\LR{$\Value{t_2}{N} \leq \Value{t_1}{N}$} وي،
    نو \LR{$\Th{Q} \Proves \lnot(t_1 < t_2)$}۔
\end{enumerate}
\end{lem}
\olpseganchor{OLP-0300-B014}{end}

\olpseganchor{OLP-0300-B015}{start}
\begin{proof}
د ورکړل شوو ترمونو~\LR{$t_1$} او~\LR{$t_2$} دپاره
\LR{$n = \Value{t_1}{N}$} او~\LR{$m = \Value{t_2}{N}$} ټاکو۔
\olpseganchor{OLP-0300-B015}{end}

\olpseganchor{OLP-0300-B016}{start}
فرض کړئ~\LR{$!A \ident t_1 = t_2$}۔ د
\olref{lem:q-proves-clterm-id} له مخې،
\LR{$\Th{Q} \Proves \eq[t_1][\num n]$} او
\LR{$\Th{Q} \Proves \eq[t_2][\num m]$}۔ که~\LR{$n = m$} وي، نو
\LR{$\Th{Q} \Proves \eq[\num n][\num m]$}، او له دې امله د عينيت د
تعدي له مخې~\LR{$\Th{Q} \Proves \eq[t_1][t_2]$}۔ که~\LR{$n \neq m$} وي، نو
\LR{$\Th{Q} \Proves \eq/[\num n][\num m]$}، او يو ځل بيا د عينيت د
تعدي له مخې~\LR{$\Th{Q} \Proves \eq/[t_1][t_2]$}۔
\textbf{د ژباړې د نسخې سمون (OLCMP-087):} سرچينې دويم تړلي ترم
هم د لومړي ترم له ارزښت سره برابر کړے و؛ دلته هر ترم له خپل ټاکل
شوي ارزښت سره برابر شوے دے۔
\olpseganchor{OLP-0300-B016}{end}

\olpseganchor{OLP-0300-B017}{start}
اوس~\LR{$!A \ident t_1 < t_2$} ونيسئ۔ په دواړو قضيو کښې پر بديهي
اصل~\LR{$!Q_8$} تکيه کوو، چې وايي د ټولو~\LR{$x,y$} دپاره
\LR{$x < y \liff \lexists[z][\eq[z' + x][y]]$}۔
\olpseganchor{OLP-0300-B017}{end}

\olpseganchor{OLP-0300-B018}{start}
فرض کړئ~\LR{$\Sat{N}{t_1 < t_2}$}۔ بيا داسې~\LR{$k \in \Nat$} شته چې
\LR{$n + k + 1 = m$}۔ د~\olref{lem:q-proves-clterm-id} له مخې،
\LR{$\Th{Q} \Proves \eq[t_1][\num n]$} او
\LR{$\Th{Q} \Proves \eq[t_2][\num m]$}، او د دې لمې د لومړۍ برخې له
مخې، \LR{$\Th{Q} \Proves \eq[{\num k}' + \num n][\num m]$}۔ د عينيت د
تعدي له مخې، \LR{$\Th{Q} \Proves \eq[{\num k}' + t_1][t_2]$}، نو
\LR{$\Th{Q} \Proves \lexists[z][\eq[z' + t_1][t_2]]$}۔ د~\LR{$!Q_8$} د
ښي-څخه-کيڼ لوري له مخې،~\LR{$\Th{Q} \Proves t_1 < t_2$}۔
\textbf{د ژباړې د نسخې سمون (OLCMP-088):} سرچينې په لومړۍ ثابته
مساوات کښې د جمع غړي د هغه وجودي شاهد پر خلاف اړولي وو چې په بله
کرښه او د کوچني والي په بديهي اصل کښې کارېږي؛ دلته ترتيب د هماغه
شاهد له بڼې سره برابر شوے دے۔
\olpseganchor{OLP-0300-B018}{end}

\olpseganchor{OLP-0300-B019}{start}
د دې پر ځاے فرض کړئ~\LR{$\Sat/{N}{t_1 < t_2}$}، يعنې~\LR{$m \leq n$}۔
%
په~\LR{$\Th{Q}$} کښې کار کوو او~\LR{$t_1 < t_2$} فرضوو۔ د~\LR{$!Q_8$} د
کيڼ-څخه-ښي لوري له مخې، داسې~\LR{$z$} شته چې
\LR{$\eq[z' + t_1][t_2]$}۔ ځکه چې
\LR{$\Th{Q} \Proves \eq[t_1][\num n]$} او
\LR{$\Th{Q} \Proves \eq[t_2][\num m]$}، نو
\LR{$\eq[z' + \num n][\num m]$}۔
%
پر~\LR{$m$} د~\LR{$!Q_5$} په کارولو د يوې بهرنۍ استقرا له مخې،
\LR{$\eq[z' + \num{n - m}][\Obj 0]$}۔
که~\LR{$m = n$} وي، نو~\LR{$\eq[z'][\Obj 0]$}، چې د~\LR{$!Q_2$} له لارې تناقض
ورکوي۔ که~\LR{$m < n$} وي، نو د~\LR{$!Q_5$} په بيا کارولو
\LR{$\eq[(z' + \num{n - m - 1})'][\Obj 0]$}، چې د~\LR{$!Q_2$} له لارې
تناقض ورکوي۔ نو~\LR{$\Th{Q} \Proves \lnot(t_1 < t_2)$}۔
\textbf{د ژباړې د نسخې سمون (OLCMP-089):} سرچينې په لومړۍ قضيه
کښې د تر لاسه شوې مساوات پر ځاے نابرابري ليکلې وه، او د صفر له تالي
سره د نابرابرۍ دپاره ئې په دواړو قضيو کښې ناسم بديهي اصل راجع کړے
و؛ دلته نښه او دواړه راجعې له اړوند بديهي اصل سره سمې شوې دي۔
\end{proof}
\olpseganchor{OLP-0300-B019}{end}

\olpseganchor{OLP-0300-B020}{start}
\begin{lem}
\ollabel{lem:bounded-quant-equiv}
فرض کړئ~\LR{$!A$} کومه فارمول ده،~\LR{$t$} تړلے ترم دے، او
\LR{$k=\Value{t}{N}$}۔ بيا
\begin{enumerate}
\item \LR{$\Th{Q} \Proves \bforall{x<t}{!A(x)}$} هله او يوازې هله چې
\LR{$\Th{Q} \Proves !A(\num 0) \land \dots \land !A(\num{k-1})$}۔
\item \LR{$\Th{Q} \Proves \bexists{x<t}{!A(x)}$} هله او يوازې هله چې
\LR{$\Th{Q} \Proves !A(\num 0) \lor \dots \lor !A(\num{k-1})$}۔
\end{enumerate}
\end{lem}
\olpseganchor{OLP-0300-B020}{end}

\olpseganchor{OLP-0300-B021}{start}
\begin{proof}
    د محدود کلي کمیت ايښوونکي قضيه ثابتوو۔
    که~\LR{$\Value{t}{N} = 0$} وي، نو د معادلۍ کيڼ اړخ په~\LR{$\Th{Q}$}
    کښې ثابتېدونکے دے، ځکه د
    \olref[inc][req][min]{lem:less-zero} له مخې هېڅ~\LR{$x<\num 0$}
    نشته۔ همداسې، تش عطف په ساده ډول~\LR{$\ltrue$} ګڼلے شو، چې په
    \LR{$\Th{Q}$} کښې هم ثابتېدونکے دے۔
    \textbf{د ژباړې د نسخې سمون (OLCMP-090):} سرچينې دلته تش منفصل
    ياد کړے و، خو د کلي کمیت ايښوونکي ښے اړخ عطف دے؛ تش عطف رښتيا
    وي، لکه چې ورپسې نښه هم ښيي۔
    %
    نو فرضوو چې د کوم طبيعي عدد~\LR{$k$} دپاره
    \LR{$\Value{t}{N} = k+1$}۔ د
    \olref{lem:q-proves-clterm-id} له مخې فرضولے شو چې د
    \LR{$\bforall{x<\num{k+1}}{!A(x)}$} په بڼه پر فارمول کار کوو۔
\olpseganchor{OLP-0300-B021}{end}

\olpseganchor{OLP-0300-B022}{start}
    فرض کړئ~\LR{$\Th{Q} \Proves \bforall{x<\num{k+1}}{!A(x)}$}، او
    \LR{$n \leq k$} ونيسئ۔ ځکه چې د~\olref{lem:atomic-completeness} له
    مخې~\LR{$\Th{Q} \Proves \num n < \num{k+1}$}، نو د منطق له مخې
    \LR{$\Th{Q} \Proves !A(\num n)$}۔ د هر~\LR{$n \leq k$} دپاره د دې حقيقت
    په~\LR{$k+1$} ځله پلي کولو، لکه څنګه چې غوښتل شوي وو، اخلو:
    \LR{$\Th{Q} \Proves !A(\num 0) \land \dots \land !A(\num k)$}۔
\olpseganchor{OLP-0300-B022}{end}

\olpseganchor{OLP-0300-B023}{start}
    د بل لوري دپاره فرض کړئ چې
    \LR{$\Th{Q} \Proves !A(\num 0) \land \dots \land !A(\num k)$}۔ په
    \LR{$\Th{Q}$} کښې کار کوو؛ فرض کړئ~\LR{$x < \num{k+1}$}۔ د
    \olref[inc][req][min]{lem:less-nsucc} له مخې لرو:
    \LR{$x = \num 0 \lor \dots \lor x = \num k$}، نو د منطق له مخې
    \LR{$!A(x)$}؛ او له دې څخه کلي ادعا
    \LR{$\bforall{x<\num{k+1}}{!A(x)}$} راځي۔
\olpseganchor{OLP-0300-B023}{end}

\olpseganchor{OLP-0300-B024}{start}
    د محدودې وجودي کمیت ايښودنې لرونکو فارمولونه د معادلۍ ثبوت
    ورته دے۔
\end{proof}
\olpseganchor{OLP-0300-B024}{end}

\olpseganchor{OLP-0300-B025}{start}
\begin{prob}
په \olref{lem:bounded-quant-equiv} کښې د وجودي قضيې تفصيلي ثبوت
ورکړئ۔
\end{prob}
\olpseganchor{OLP-0300-B025}{end}

\olpseganchor{OLP-0300-B026}{start}
\begin{lem}
\ollabel{lem:delta0-completeness}
که~\LR{$!A$} داسې~\LR{$\Delta_0$} تړلے فارمول وي چې په~\LR{$\Struct{N}$} کښې
رښتينې وي، نو~\LR{$\Th{Q} \Proves !A$}۔
\end{lem}
\olpseganchor{OLP-0300-B026}{end}

\olpseganchor{OLP-0300-B027}{start}
\begin{proof}
دا د~فارمول پر پېچلتيا په استقرا ثابتوو۔
%
بنسټيزه قضيه~\olref{lem:atomic-completeness} ورکوي، نو استقرايي
پړاو ته ځو۔ د اسانۍ دپاره د نفي قضيه د هغې~فارمول د جوړښت له
مخې پر فرعي قضيو وېشو چې نفي ورباندې پلي شوې ده۔
\olpseganchor{OLP-0300-B027}{end}

\olpseganchor{OLP-0300-B028}{start}
\begin{enumerate}
\item فرض کړئ~\LR{$(!A \land !B)$} په~\LR{$\Struct{N}$} کښې رښتينې ده، نو
\LR{$!A$} او~\LR{$!B$} په~\LR{$\Struct{N}$} کښې رښتينې دي۔ د استقرايي فرضيې له
مخې، \LR{$\Th{Q} \Proves !A$} او~\LR{$\Th{Q} \Proves !B$}، نو د منطق له مخې
\LR{$\Th{Q} \Proves (!A \land !B)$}۔
%
\item فرض کړئ~\LR{$\lnot (!A \land !B)$} په~\LR{$\Struct{N}$} کښې رښتينې ده،
نو يا~\LR{$\lnot !A$} يا~\LR{$\lnot !B$} په~\LR{$\Struct{N}$} کښې رښتينې ده۔ د
عموميت له زيان پرته، لومړۍ فرض کړئ۔ د استقرايي فرضيې له مخې
\LR{$\Th{Q} \Proves \lnot !A$}، او له دې امله د منطق له مخې
\LR{$\Th{Q} \Proves \lnot (!A \land !B)$}۔
%
\item فرض کړئ~\LR{$(!A \lor !B)$} په~\LR{$\Struct{N}$} کښې رښتينې ده، نو يا
\LR{$!A$} په~\LR{$\Struct{N}$} کښې يا~\LR{$!B$} په~\LR{$\Struct{N}$} کښې رښتينې ده۔ د عموميت له زيان
پرته، لومړۍ فرض کړئ۔ د استقرايي فرضيې له مخې
\LR{$\Th{Q} \Proves !A$}، او له دې امله د منطق له مخې
\LR{$\Th{Q} \Proves (!A \lor !B)$}۔
%
\item فرض کړئ~\LR{$\lnot(!A \lor !B)$} په~\LR{$\Struct{N}$} کښې رښتينې ده،
نو~\LR{$\lnot !A$} او~\LR{$\lnot !B$} په~\LR{$\Struct{N}$} کښې رښتينې دي۔ د
استقرايي فرضيې له مخې، \LR{$\Th{Q} \Proves \lnot !A$} او
\LR{$\Th{Q} \Proves \lnot !B$}۔ په پايله کښې، د منطق له مخې
\LR{$\Th{Q} \Proves \lnot(!A \lor !B)$}۔
%
\item فرض کړئ~\LR{$\bforall{x<t}{!A(x)}$} په~\LR{$\Struct{N}$} کښې رښتينې
ده، چې~\LR{$t$} تړلے ترم او~\LR{$k=\Value{t}{N}$} دے۔ د استقرايي فرضيې او
منطق له مخې، که~\LR{$!A(\num n)$} د هر~\LR{$n < \Value{t}{N}$} دپاره په
\LR{$\Struct{N}$} کښې رښتينې وي، نو
\LR{$\Th{Q} \Proves !A(\num 0) \land \dots \land !A(\num{k-1})$}۔ د
\olref{lem:bounded-quant-equiv} له مخې،
\LR{$\Th{Q} \Proves \bforall{x<t}{!A(x)}$}۔
%
\item د محدود وجودي کمیت ايښوونکي قضيه، چې د
\LR{$\bexists{x < t}{!A(x)}$} په بڼه تړلے فارمول لرو، د محدود کلي
کمیت ايښوونکي قضيې ته ورته ده۔
%
\item فرض کړئ~\LR{$\lnot \bforall{x<t}{!A(x)}$} په~\LR{$\Struct{N}$} کښې
رښتينې ده، چې~\LR{$t$} تړلے ترم دے۔ دا تړلے فارمول له
تړلے فارمول~\LR{$\bexists{x<t}{\lnot !A(x)}$} سره معادله ده، او دا معادله په
\LR{$\Th{Q}$} کښې استنتاجېدے شي؛ نو د محدودو وجودي کمیت ايښوونکو
استدلال پلي کولے شو۔
%
\item همداسې، فرض کړئ~\LR{$\lnot \bexists{x<t}{!A(x)}$} په~\LR{$\Struct{N}$}
کښې رښتينې ده، چې~\LR{$t$} تړلے ترم دے۔ دا تړلے فارمول په~\LR{$\Th{Q}$} کښې
له~\LR{$\bforall{x<t}{\lnot!A(x)}$} سره معادله ده، نو د محدودو کلي کمیت
ايښوونکو استدلال پلي کولے شو۔
\textbf{د ژباړې د نسخې سمون (OLCMP-097):} سرچينې د محدود وجودي
کمیت ايښوونکي ماکرو ته فورمولي ماتريس د اړين دويم آرګومېنټ په بڼه
نۀ و ورکړے؛ دلته هماغه ماتريس په څرګندو قوسونو کښې ورکړل شوے دے۔
%
\item په پاې کښې، فرض کړئ~\LR{$\lnot !A$} په~\LR{$\Struct{N}$} کښې رښتينې
ده۔ يوازې هغه قضيې پاتې دي چې~\LR{$!A$} اټومي وي او چې د کومې
\LR{$\Delta_0$} تړلے فارمول~\LR{$!B$} دپاره
\LR{$\lnot !A \ident \lnot\lnot !B$} وي۔ که~\LR{$!A$} اټومي وي، نو د
\olref{lem:atomic-completeness} له مخې
\LR{$\Th{Q} \Proves \lnot !A$}۔ که
\LR{$\lnot !A \ident \lnot\lnot !B$} وي، نو د منطق له مخې دا په
\LR{$\Th{Q}$} کښې په ثابتېدونکي ډول له~\LR{$!B$} سره معادله ده، چې په
\LR{$\Struct{N}$} کښې رښتينې ده، ځکه~\LR{$\lnot !A$} په~\LR{$\Struct{N}$} کښې
رښتينې ده۔ نو د استقرايي فرضيې له مخې
\LR{$\Th{Q} \Proves \lnot !A$}۔
\end{enumerate}
\end{proof}
\olpseganchor{OLP-0300-B028}{end}

\olpseganchor{OLP-0300-B029}{start}
\begin{prob}
په \olref{lem:delta0-completeness} کښې د وجودي قضيې تفصيلي ثبوت
ورکړئ۔
\end{prob}
\olpseganchor{OLP-0300-B029}{end}

\olpseganchor{OLP-0300-B030}{start}
\begin{thm}
\ollabel{thm:sigma1-completeness}
که~\LR{$!A$} داسې~\LR{$\Sigma_1$} تړلے فارمول وي چې په~\LR{$\Struct{N}$} کښې
رښتينې وي، نو~\LR{$\Th{Q} \Proves !A$}۔
\end{thm}
\olpseganchor{OLP-0300-B030}{end}

\olpseganchor{OLP-0300-B031}{start}
\begin{proof}
که~\LR{$\lexists[x][!A(x)]$} داسې~\LR{$\Sigma_1$} تړلے فارمول وي چې په
\LR{$\Struct{N}$} کښې رښتينې وي، نو داسې طبيعي عدد~\LR{$n$} او د متحولونو
د ارزښت ټاکنه~\LR{$s$} شته چې~\LR{$s(x) = n$} او~\LR{$\Sat{N}{!A(x)}[s]$}۔ د
پوره کېدو د اړيکې په اړه له معياري حقيقتونو څخه
\LR{$\Sat{N}{!A(\num n)}$} لازمېږي۔
\textbf{د ژباړې د نسخې سمون (OLCMP-098):} سرچينې دلته د اوږده
وجودي کمیت ايښوونکي ماکرو يو-قوسي غږ کارولے و؛ د سټايل د اعلان شوې
منځګړتيا له مخې دلته متحول او ماتريس په اختياري قوسونو کښې جلا ورکړل
شوي دي۔
خو~\LR{$!A(\num n)$} يوه
\LR{$\Delta_0$} فارمول ده، نو د~\olref{lem:delta0-completeness} له
مخې~\LR{$\Th{Q} \Proves !A(\num n)$}؛ او له دې امله د منطق له مخې
\LR{$\Th{Q} \Proves \lexists[x][!A(x)]$} هم لرو۔
\end{proof}
\olpseganchor{OLP-0300-B031}{end}

\olpseganchor{OLP-0301-B004}{start}
\begin{editorial}
  دا څپرکے د محاسبه کېدنې د تيورۍ د څپرکي پر موادو تکيه لري، خو که
  هغه څپرکے نۀ وي لوستل شوے، دا پرېښودل کېدے شي۔ اوس مهال دا د
  جېرېمي اېويګېډ د يادښتونو يوه بنسټيزه اړونه ده، بياکتنه ئې نۀ ده
  شوې، او تمرينونه پکې نشته۔
\end{editorial}
\olpseganchor{OLP-0301-B004}{end}

\olpseganchor{OLP-0301-B005}{start}
\olchapter{inc}{tcp}{تيورۍ او محاسبه کېدنه}
\olpseganchor{OLP-0301-B005}{end}

\olpseganchor{OLP-0302-B004}{start}
\olfileid{inc}{tcp}{int}
\olpseganchor{OLP-0302-B004}{end}

\olpseganchor{OLP-0302-B005}{start}
\olsection{پېژندګلو}
\olpseganchor{OLP-0302-B005}{end}

\olpseganchor{OLP-0302-B006}{start}
\begin{editorial}
دا برخه بايد له سره وليکل شي۔
\end{editorial}
\olpseganchor{OLP-0302-B006}{end}

\olpseganchor{OLP-0302-B007}{start}
موږ لاندې شيان لرو:
\begin{enumerate}
\item دا تعريف چې يوه تابعه په~\LR{$\Th{Q}$} کښې تمثيلېدونکې وي څه معنا
  لري (\olref[req][int]{defn:representable-fn})؛
\item دا تعريف چې يوه اړيکه په~\LR{$\Th{Q}$} کښې تمثيلېدونکې وي څه معنا
  لري (\olref[req][rel]{defn:representing-relations})؛
\item هغه قضيه چې وايي د~\LR{$\Th{Q}$} تمثيلېدونکې تابعې کټ مټ هماغه
  محاسبه کېدونکې تابعې دي
  (\olref[req][int]{thm:representable-iff-comp})؛
\item او هغه قضيه چې وايي د~\LR{$\Th{Q}$} تمثيلېدونکې اړيکې کټ مټ هماغه
  محاسبه کېدونکې اړيکې دي
  \olref[req][rel]{thm:representing-rels}۔
\end{enumerate}
\olpseganchor{OLP-0302-B007}{end}

\olpseganchor{OLP-0302-B008}{start}
يوه \emph{تيوري} د جملو داسې سټ دے چې د اشتقاق له مخې تړلے وي؛
يعنې دا خاصيت لري چې هر کله~\LR{$T$}، \LR{$!A$} ثابتوي، نو~\LR{$!A$} په~\LR{$T$} کښې
وي۔ غوره به دا وي چې تيوري د جملو د يوې ټولګې په توګه وګڼو چې د
هغو جملو ټولې لازمې پايلې هم ورسره وي۔ له دې وروسته به~\LR{$\Th{Q}$} د
هغې \emph{تيورۍ} دپاره کاروو چې په~\olref[req][int]{sec} کښې له
اتو بديهي اصولو څخه د اشتقاق وړ جملو له سټ څخه جوړه ده۔ په ياد ولرئ
چې د~\LR{$\Th{Q}$} فارمولونه د عددونو په توګه کوډولے شو؛ که~\LR{$!A$} داسې
فارمول وي، پرېږدئ~\LR{$\Gn{!A}$} هغه عدد وښيي چې~\LR{$!A$} کوډوي۔ د همدې
کوډونې له مخې اوس پوښتلے شو چې د فارمولونو بېلابېل سټونه محاسبه
کېدونکي دي که نۀ۔
\olpseganchor{OLP-0302-B008}{end}

\olpseganchor{OLP-0303-B004}{start}
\olfileid{inc}{tcp}{qce}
\olpseganchor{OLP-0303-B004}{end}

\olpseganchor{OLP-0303-B005}{start}
\olsection{\LR{$\Th{Q}$}، c.e.-بشپړه ده}
\olpseganchor{OLP-0303-B005}{end}


\olpseganchor{OLP-0303-B006}{start}
\begin{thm}
\LR{$\Th{Q}$}، c.e. ده خو د پرېکړې وړ نۀ ده۔ په حقيقت کښې دا يو
بشپړ c.e. سټ دے۔
\end{thm}
\olpseganchor{OLP-0303-B006}{end}

\olpseganchor{OLP-0303-B007}{start}
\begin{proof}
دا ليدل ستونزمن نۀ دي چې~\LR{$\Th{Q}$}، c.e. ده، ځکه دا د هغو
جملو~\LR{$y$} د (کوډونو) سټ دے چې داسې~\LR{$x$} وي چې د~\LR{$y$} په~\LR{$\Th{Q}$} کښې
ثبوت وي:
\[
Q = \Setabs{y}{\lexists[x][\Prf[\Th{Q}](x,y)]}.
\]
خو پوهېږو چې~\LR{$\Prf[\Th{Q}](x,y)$} محاسبه کېدونکې ده (په حقيقت کښې
بنسټيزه بازګشتي ده)، او هر هغه سټ چې په پورتنۍ بڼه وليکل شي، په
محاسبوي ډول د شمېر وړ وي۔
\olpseganchor{OLP-0303-B007}{end}

\olpseganchor{OLP-0303-B008}{start}
دا ويل چې دا يو بشپړ په محاسبوي ډول د شمېر وړ سټ دے، له دې خبرې
سره برابر دي چې~\LR{$K \leq_m Q$}، چېرته چې
\LR{$K = \Setabs{x}{!A_x(x) \downarrow}$}۔ نو راځئ وښيو چې~\LR{$K$}،
\LR{$\Th{Q}$} ته راکمېدونکے دے۔ ځکه چې د کلېني پريديکات~\LR{$T(e,x,s)$}
بنسټيزه بازګشتي دے، په~\LR{$\Th{Q}$} کښې، فرض کړئ، د~\LR{$!A_T$} په وسيله
تمثيلېږي۔ بيا د هر~\LR{$x$} دپاره لرو:
\begin{align*}
x \in K & \lif \lexists[s][T(x,x,s)] \\
& \lif \lexists[s][(\Th{Q} \Proves !A_T(\num x,\num x, \num s))]
  \\
& \lif \Th{Q} \Proves \lexists[s][!A_T(\num x, \num x, s)].
\end{align*}
برعکس، که~\LR{$\Th{Q} \Proves \lexists[s][!A_T(\num x, \num x, s)]$}،
نو په حقيقت کښې د کوم طبيعي عدد~\LR{$n$} دپاره بايد فارمول
\LR{$!A_T(\num x, \num x, \num n)$} رښتيا وي۔ اوس که~\LR{$T(x,x,n)$} ناسم
وای، نو ځکه چې~\LR{$!A_T$}، \LR{$T$} تمثيلوي، \LR{$\Th{Q}$} به
\LR{$\lnot !A_T(\num x, \num x, \num n)$} ثابت کړے وای۔ خو بيا به
\LR{$\Th{Q}$} يو ناسم فارمول ثابت کړے وای، چې تناقض دے۔ نو~\LR{$T(x,x,n)$}
بايد رښتيا وي، او له دې څخه~\LR{$!A_x(x) \downarrow$} لازمېږي۔
\olpseganchor{OLP-0303-B008}{end}

\olpseganchor{OLP-0303-B009}{start}
په لنډ ډول، د هر~\LR{$x$} دپاره، \LR{$x$} هله او يوازې هله په~\LR{$K$} کښې دے
چې~\LR{$\Th{Q}$}، \LR{$\lexists[s][!A_T(\num x,\num x,s)]$} ثابت کړي۔ نو هغه
تابعه~\LR{$f$} چې~\LR{$x$} د جملې~\LR{$\lexists[s][!A_T(\num x,
  \num x, s)]$} (يو کوډ) ته وړي، له~\LR{$K$} څخه~\LR{$\Th{Q}$} ته راکمونه ده۔
\textbf{د ژباړې د نسخې سمون (OLCMP-091):} سرچينې په دې وروستۍ
پايله کښې د تمثيلوونکي فارمول پر ځاے د ماوراءژبې اړيکه ليکلې وه؛
دلته هماغه فارمول کارول شوے چې په مخکښې استدلال کښې معرفي او ثابت
شوے دے۔
\end{proof}
\olpseganchor{OLP-0303-B009}{end}

\olpseganchor{OLP-0304-B004}{start}
\olfileid{inc}{tcp}{oqn}
\olpseganchor{OLP-0304-B004}{end}

\olpseganchor{OLP-0304-B005}{start}
\olsection{د~\LR{$\Th{Q}$}، \LR{$\omega$}-سازګارې غځونې ناپرېکړه‌وړې دي}
\olpseganchor{OLP-0304-B005}{end}

\olpseganchor{OLP-0304-B006}{start}
\begin{explain}
دا ثبوت چې~\LR{$\Th{Q}$} په محاسبوي ډول د شمېر وړ بشپړه ده، پر دې حقيقت
تکيه درلوده چې په~\LR{$\Th{Q}$} کښې هره ثابتېدونکې جمله د طبيعي عددونو
په اړه «رښتيا» ده۔ راتلونکے تعريف او قضيه دا پايله پياوړې کوي، او
کټ مټ د «رښتياوالي» هغه اړخونه ټاکي چې په پورتني ثبوت کښې ورته
اړتيا وه۔ خو پر دې قضيه ډېر مه درېږئ، ځکه ژر به ئې لا نوره هم
پياوړې کړو۔ موږ دا زياتره د تاريخي موخې دپاره راوړو: د ګوډل اصلي
مقالې د~\LR{$\omega$}-سازګارۍ مفهوم کارولے و، خو لږ وروسته ئې پايله د
\LR{$\omega$}-سازګارۍ پر ځاے د عادي سازګارۍ په کارولو پياوړې شوه۔
\end{explain}
\olpseganchor{OLP-0304-B006}{end}

\olpseganchor{OLP-0304-B007}{start}
\begin{defn}
\ollabel{thm:oconsis-q}
يوه تيوري~\LR{$\Th{T}$} هله~\LR{$\omega$}-سازګاره ده چې لاندې شرط پوره کړي:
که~\LR{$\lexists[x][!A(x)]$} هره جمله وي او~\LR{$\Th{T}$}،
\LR{$\lnot !A(\num 0)$}، \LR{$\lnot !A(\num 1)$}، \LR{$\lnot !A(\num 2)$}،
\dots ثابتوي، نو~\LR{$\Th{T}$}، \LR{$\lexists[x][!A(x)]$} نۀ ثابتوي۔
\end{defn}
\olpseganchor{OLP-0304-B007}{end}

\olpseganchor{OLP-0304-B008}{start}
\begin{thm}
  پرېږدئ~\LR{$\Th{T}$} هره هغه~\LR{$\omega$}-سازګاره تيوري وي چې~\LR{$\Th{Q}$}
  پکې شامله وي۔ بيا~\LR{$\Th{T}$} د پرېکړې وړ نۀ ده۔
\end{thm}
\olpseganchor{OLP-0304-B008}{end}

\olpseganchor{OLP-0304-B009}{start}
\begin{proof}
که~\LR{$\Th{Q}$} په~\LR{$\Th{T}$} کښې شامله وي، نو~\LR{$\Th{T}$} محاسبه کېدونکې
تابعې او اړيکې تمثيلوي۔ يوازې پخوانی ثبوت لږ بدلول پکار دي۔ لکه
پورته، که~\LR{$x \in K$} وي، نو~\LR{$\Th{T}$}،
\LR{$\lexists[s][!A_T(\num x,\num x, s)]$} ثابتوي۔ برعکس، فرض کړئ
\LR{$\Th{T}$}، \LR{$\lexists[s][!A_T(\num x, \num x, s)]$} ثابتوي۔ بيا بايد
\LR{$x$} په~\LR{$K$} کښې وي: که نۀ وي، د ماشين~\LR{$x$} داسې محاسبه نشته چې پر
ننوت~\LR{$x$} ودرېږي؛ ځکه~\LR{$!A_T$} د کلېني اړيکه~\LR{$T$} تمثيلوي، \LR{$\Th{T}$}،
\LR{$\lnot !A_T(\num x, \num x, \num 0)$}،
\LR{$\lnot !A_T(\num x, \num x, \num 1)$}، \dots ثابتوي، او په دې سره
\LR{$\Th{T}$}، \LR{$\omega$}-ناسازګاره کېږي۔
\end{proof}
\olpseganchor{OLP-0304-B009}{end}

\olpseganchor{OLP-0305-B004}{start}
\olfileid{inc}{tcp}{cqn}
\olpseganchor{OLP-0305-B004}{end}

\olpseganchor{OLP-0305-B005}{start}
\olsection{د~\LR{$\Th{Q}$} سازګارې غځونې ناپرېکړه‌وړې دي}
\olpseganchor{OLP-0305-B005}{end}

\olpseganchor{OLP-0305-B006}{start}
\begin{explain}
په ياد ولرئ چې يوه تيوري هله~\emph{سازګاره} ده چې د هېڅ
فارمول~\LR{$!A$} دپاره هم~\LR{$!A$} او هم~\LR{$\lnot !A$} دواړه ثابت نۀ کړي۔ ځکه
له تناقض څخه هر څه لازمېږي، يوه ناسازګاره تيوري ابتذالي وي: هره
جمله پکې ثابتېدونکې وي۔ څرګنده ده چې که يوه تيوري~\LR{$\omega$}-سازګاره
وي، نو سازګاره هم ده۔ خو يوازې سازګاره اوسېدل کمزورے شرط دے (يعنې
داسې تيورۍ شته چې سازګارې دي خو~\LR{$\omega$}-سازګارې نۀ دي)۔ په
\olref[oqn]{thm:oconsis-q} کښې فرضيه ساده سازګارۍ ته کمزورې کولے
شو، او په دې توګه يوه پياوړې قضيه ترلاسه کوو۔
\end{explain}
\olpseganchor{OLP-0305-B006}{end}

\olpseganchor{OLP-0305-B007}{start}
\begin{lem}
هېڅ «نړيواله محاسبه کېدونکې اړيکه» نشته۔ يعنې داسې دوه‌ځاييزه
محاسبه کېدونکې اړيکه~\LR{$R(x,y)$} نشته چې لاندې خاصيت ولري: هر کله
\LR{$S(y)$} يوه يوځاييزه محاسبه کېدونکې اړيکه وي، داسې~\LR{$k$} وي چې د هر
\LR{$y$} دپاره~\LR{$S(y)$} هله او يوازې هله رښتيا وي چې~\LR{$R(k,y)$} رښتيا وي۔
\end{lem}
\olpseganchor{OLP-0305-B007}{end}

\olpseganchor{OLP-0305-B008}{start}
\begin{proof}
فرض کړئ~\LR{$R(x,y)$} يوه نړيواله محاسبه کېدونکې اړيکه وي۔ پرېږدئ
\LR{$S(y)$}، اړيکه~\LR{$\lnot R(y,y)$} وي۔ ځکه~\LR{$S(y)$} محاسبه کېدونکې ده، د
کوم~\LR{$k$} دپاره~\LR{$S(y)$} له~\LR{$R(k,y)$} سره هم‌ارزه ده۔ خو بيا~\LR{$S(k)$} هم
له~\LR{$R(k,k)$} او هم له~\LR{$\lnot R(k,k)$} سره هم‌ارزه ده، چې تناقض دے۔
\end{proof}
\olpseganchor{OLP-0305-B008}{end}


\olpseganchor{OLP-0305-B009}{start}
\begin{thm}
پرېږدئ~\LR{$\Th{T}$} هره هغه سازګاره تيوري وي چې~\LR{$\Th{Q}$} پکې شامله
وي۔ بيا~\LR{$\Th{T}$} د پرېکړې وړ نۀ ده۔
\end{thm}
\olpseganchor{OLP-0305-B009}{end}


\olpseganchor{OLP-0305-B010}{start}
\begin{proof}
فرض کړئ~\LR{$\Th{T}$} د~\LR{$\Th{Q}$} يوه سازګاره او د پرېکړې وړ غځونه وي۔
موږ به د~\LR{$\Th{T}$} په کارولو د يوې نړيوالې محاسبه کېدونکې اړيکې په
تعريفولو سره تناقض ترلاسه کړو۔
\olpseganchor{OLP-0305-B010}{end}

\olpseganchor{OLP-0305-B011}{start}
پرېږدئ~\LR{$R(x,y)$} هله او يوازې هله صدق وکړي چې
\begin{quote}
\LR{$x$} د~\LR{$!D(u)$} فارمول کوډوي، او~\LR{$\Th{T}$}، \LR{$!D(\num y)$} ثابتوي۔
\end{quote}
ځکه چې فرض کوو~\LR{$\Th{T}$} د پرېکړې وړ ده، \LR{$R$} محاسبه کېدونکې ده۔
راځئ وښيو چې~\LR{$R$} نړيواله ده۔ که~\LR{$S(y)$} هره محاسبه کېدونکې اړيکه
وي، نو په~\LR{$\Th{Q}$} (او له دې کبله په~\LR{$\Th{T}$}) کښې د کوم
فارمول~\LR{$!D_S(u)$} په وسيله تمثيلېږي۔ بيا د هر~\LR{$n$} دپاره لرو:
\begin{eqnarray*}
S(n) & \lif & T \vdash !D_S(\num n) \\
& \lif & R(\Gn{!D_S(u)}, n)
\end{eqnarray*}
او
\begin{eqnarray*}
\lnot S(n) & \lif & T \vdash \lnot !D_S(\num n) \\
& \lif & T \not\vdash !D_S(\num n) \quad \text{\RL{(ځکه~\LR{$\Th{T}$} سازګاره ده)}} \\
& \lif & \lnot R(\Gn{!D_S(u)}, n).
\end{eqnarray*}
\textbf{د ژباړې د نسخې سمون (OLCMP-092):} سرچينې په دواړو
ماوراءژبنیو اړيکو کښې طبيعي عدد د څيز-ژبې د عددنښې په بڼه ليکلے و؛
دلته اړيکه پر طبيعي عدد تطبيق شوې، او عددنښه يوازې د تمثيلوونکي
فارمول دننه پاتې ده۔
يعنې د هر~\LR{$y$} دپاره~\LR{$S(y)$} هله او يوازې هله رښتيا وي چې
\LR{$R(\Gn{!D_S(u)}, y)$} رښتيا وي۔ نو~\LR{$R$} نړيواله ده، او هغه تناقض مو
ترلاسه کړ چې غوښت۔
\end{proof}
\olpseganchor{OLP-0305-B011}{end}

\olpseganchor{OLP-0305-B012}{start}
پرېږدئ «رښتينی حساب» تيوري~\LR{$\Setabs{!A}{ \Sat{\Nat}{!A}}$} وي؛
يعنې د حساب د ژبې د هغو جملو سټ چې په معياري تفسير کښې رښتيا دي۔
\olpseganchor{OLP-0305-B012}{end}

\olpseganchor{OLP-0305-B013}{start}
\begin{cor}
رښتينی حساب د پرېکړې وړ نۀ دے۔
\end{cor}
\olpseganchor{OLP-0305-B013}{end}

\olpseganchor{OLP-0305-B014}{start}
%In Section~\olref{undefinability:truth:section} we will state a stronger
%result, due to Tarski.
\olpseganchor{OLP-0305-B014}{end}

\olpseganchor{OLP-0306-B004}{start}
\olfileid{inc}{tcp}{cax}
\olpseganchor{OLP-0306-B004}{end}

\olpseganchor{OLP-0306-B005}{start}
\olsection{په محاسبوي ډول بديهي کېدونکې تيورۍ}
\olpseganchor{OLP-0306-B005}{end}

\olpseganchor{OLP-0306-B006}{start}
يوه تيوري~\LR{$\Th{T}$} هله~\emph{په محاسبوي ډول بديهي کېدونکې} بلل کېږي چې د بديهي
اصولو يو محاسبه کېدونکے سټ~\LR{$A$} ولري۔ (دا ويل چې~\LR{$A$} د~\LR{$\Th{T}$}
د بديهي اصولو سټ دے، معنا ئې دا ده چې
\LR{$T = \Setabs{!B}{A \Proves !B}$}۔)
\textbf{د ژباړې د نسخې سمون (OLCMP-093):} سرچينې په دې سټ-جوړوونکې
بڼه کښې يوه نښه هم د بديهي اصولو د سټ او هم د اړوندې جملې دپاره
کارولې وه؛ دلته جمله په بله نښه بېله شوې ده، څو تعريف هماغه تيوري
وټاکي چې له ټاکلي بديهي سټ څخه د اشتقاق وړ ده۔
د طبيعي عددونو هره «معقوله» بديهي کول به دا خاصيت لري۔ په ځانګړې
توګه، هره هغه تيوري چې د بديهي اصولو متناهي سټ ولري،
په محاسبوي ډول بديهي کېدونکې وي۔
\olpseganchor{OLP-0306-B006}{end}

\olpseganchor{OLP-0306-B007}{start}
\begin{lem}
فرض کړئ~\LR{$\Th{T}$}، په محاسبوي ډول بديهي کېدونکې ده۔ بيا~\LR{$\Th{T}$}،
په محاسبوي ډول د شمېر وړ ده۔
\end{lem}
\olpseganchor{OLP-0306-B007}{end}

\olpseganchor{OLP-0306-B008}{start}
\begin{proof}
فرض کړئ~\LR{$A$} د~\LR{$\Th{T}$} د بديهي اصولو يو محاسبه کېدونکے سټ دے۔ د
دې د ټاکلو دپاره چې~\LR{$!A \in T$} ده که نۀ، يوازې د بديهي اصولو څخه
د~\LR{$!A$} د اشتقاق پلټنه وکړئ۔
\olpseganchor{OLP-0306-B008}{end}

\olpseganchor{OLP-0306-B009}{start}
لږ په بله وينا، \LR{$!A$} هله او يوازې هله په~\LR{$\Th{T}$} کښې ده چې
په~\LR{$A$} کښې د بديهي اصولو يو متناهي لست~\LR{$!B_1$}، \dots، \LR{$!B_k$} او
په لومړۍ درجې منطق کښې د~\LR{$(!B_1 \land \dots \land !B_k) \lif !A$}
يو اشتقاق وي۔ خو موږ لا پوهېږو چې هر هغه سټ چې تعريف ئې د
«داسې \dots شته چې \dots» په بڼه وي، c.e. وي، په دې شرط چې
دويم «\dots» محاسبه کېدونکے وي۔
\end{proof}
\olpseganchor{OLP-0306-B009}{end}

\olpseganchor{OLP-0307-B004}{start}
\olfileid{inc}{tcp}{cdc}
\olpseganchor{OLP-0307-B004}{end}

\olpseganchor{OLP-0307-B005}{start}
\olsection{په محاسبوي ډول بديهي کېدونکې بشپړې تيورۍ د پرېکړې وړ دي}
\olpseganchor{OLP-0307-B005}{end}

\olpseganchor{OLP-0307-B006}{start}
يوه تيوري هله~\emph{بشپړه} بلل کېږي چې د هرې جملې~\LR{$!A$} دپاره يا
\LR{$!A$} او يا~\LR{$\lnot !A$} ثابتېدونکې وي۔
\olpseganchor{OLP-0307-B006}{end}

\olpseganchor{OLP-0307-B007}{start}
\begin{lem}
فرض کړئ يوه تيوري~\LR{$\Th{T}$} بشپړه او په محاسبوي ډول بديهي کېدونکې ده۔ بيا
\LR{$\Th{T}$} د پرېکړې وړ ده۔
\end{lem}
\olpseganchor{OLP-0307-B007}{end}

\olpseganchor{OLP-0307-B008}{start}
\begin{proof}
فرض کړئ~\LR{$\Th{T}$} بشپړه ده او~\LR{$A$} د بديهي اصولو يو محاسبه کېدونکے
سټ دے۔ که~\LR{$\Th{T}$} ناسازګاره وي، نو څرګنده ده چې محاسبه کېدونکې
ده۔ (الګوريتم: «تل هو ووايه۔») نو فرض کولے شو چې~\LR{$\Th{T}$} سازګاره
هم ده۔
\olpseganchor{OLP-0307-B008}{end}

\olpseganchor{OLP-0307-B009}{start}
د دې پرېکړې دپاره چې يوه جمله~\LR{$!A$} په~\LR{$\Th{T}$} کښې ده که نۀ،
په يو وخت له~\LR{$\Th{T}$} څخه هم د~\LR{$!A$} د اشتقاق او هم د
\LR{$\lnot !A$} د اشتقاق پلټنه وکړئ۔ ځکه~\LR{$\Th{T}$} بشپړه ده،
خامخا به يو يا بل ومومئ؛ او ځکه~\LR{$\Th{T}$} سازګاره ده، که د
\LR{$\lnot !A$} يو اشتقاق ومومئ، نو د~\LR{$!A$} هېڅ اشتقاق
نشته۔
\olpseganchor{OLP-0307-B009}{end}

\olpseganchor{OLP-0307-B010}{start}
په بله وينا، موږ لا پوهېږو چې~\LR{$\Th{T}$}، c.e. ده؛ نو د هغې
قضيې له مخې چې مخکښې مو ثابته کړې، بس دا ښودل پکار دي چې د~\LR{$\Th{T}$}
متمم هم~c.e. دے۔ خو فارمول~\LR{$!A$} هله او يوازې هله په
\LR{$\Th{\bar T}$} کښې دے چې~\LR{$\lnot !A$} په~\LR{$\Th{T}$} کښې وي؛ نو
\LR{$\Th{\bar T} \leq_m \Th{T}$}۔
\end{proof}
\olpseganchor{OLP-0307-B010}{end}

\olpseganchor{OLP-0308-B004}{start}
\olfileid{inc}{tcp}{inc}
\olpseganchor{OLP-0308-B004}{end}

\olpseganchor{OLP-0308-B005}{start}
\olsection{\LR{$\Th{Q}$} هېڅ بشپړه، سازګاره او په محاسبوي ډول بديهي کېدونکې
  غځونه نۀ لري}
\olpseganchor{OLP-0308-B005}{end}

\olpseganchor{OLP-0308-B006}{start}
\begin{thm}
\ollabel{thm:first-incompleteness}
د~\LR{$\Th{Q}$} هېڅ بشپړه، سازګاره او په محاسبوي ډول بديهي کېدونکې غځونه نشته۔
\end{thm}
\olpseganchor{OLP-0308-B006}{end}

\olpseganchor{OLP-0308-B007}{start}
\begin{proof}
موږ لا پوهېږو چې د~\LR{$\Th{Q}$} هېڅ سازګاره او د پرېکړې وړ غځونه
نشته۔ خو که~\LR{$\Th{T}$} بشپړه او بديهي‌کړې وي، نو د پرېکړې وړ
ده۔
\end{proof}
\olpseganchor{OLP-0308-B007}{end}

\olpseganchor{OLP-0308-B008}{start}
\begin{explain}
دا قضيه د ګوډل د ناتکميلۍ د لومړۍ قضيې له اصلي ۱۹۳۱ز بيان څخه ډېره
لرې نۀ ده۔ له نوې اصطلاح پرته، مهم توپيرونه دا دي: ګوډل د
«سازګار» پر ځاے «\LR{$\omega$}-سازګار» لري؛ او هغه نۀ شو کولے په بشپړ
عموميت «په محاسبوي ډول بديهي کېدونکې» ووايي، ځکه هغه مهال د محاسبه کېدنې صوري
مفهوم لا جوړ شوے نۀ و۔ (د محاسبه کېدنې صوري مدلونه په راتلونکې
لسيزه کښې رامنځته شول، چې ګوډل هم پکې برخه لرله، او تر ډېره د دې
دپاره چې د هغو تيوريو ډولونه ځانګړي کړي چې د ګوډل ښکارندې پرې پلې
کېږي۔)
\olpseganchor{OLP-0308-B008}{end}

\olpseganchor{OLP-0308-B009}{start}
قضيه وايي چې هر څه يوځای نۀ شئ لرلے؛ يعنې بشپړتيا، سازګاري او
د محاسبوي بديهي کېدو وړتيا۔ خو که له دې هر يو پرېږدئ، نور دوه لرلے شئ:
\LR{$\Th{Q}$} سازګاره او د محاسبوي بديهي اصولو لرونکې ده، خو بشپړه نۀ
ده؛ ناسازګاره تيوري بشپړه ده او محاسبوي بديهي اصول لري (مثلاً د
\LR{$\{ \eq/[0][0] \}$} په وسيله)، خو سازګاره نۀ ده؛ او د حساب د رښتينو
جملو سټ بشپړ او سازګار دے، خو د محاسبوي بديهي اصولو لرونکے نۀ دے۔
\end{explain}
\olpseganchor{OLP-0308-B009}{end}

\olpseganchor{OLP-0309-B004}{start}
\olfileid{inc}{tcp}{ins}
\olpseganchor{OLP-0309-B004}{end}

\olpseganchor{OLP-0309-B005}{start}
\olsection{په~\LR{$\Th{Q}$} کښې ثابتېدونکې او ردېدونکې جملې په محاسبوي
  ډول نه بېلېدونکې دي}
\olpseganchor{OLP-0309-B005}{end}


\olpseganchor{OLP-0309-B006}{start}
پرېږدئ~\LR{$\Th{\bar Q}$} د هغو جملو سټ وي چې \emph{نفيانې} ئې په
\LR{$\Th{Q}$} کښې ثابتېدونکې دي؛ يعنې
\LR{$\Th{\bar Q} = \Setabs{!A}{\Th{Q} \Proves \lnot !A}$}۔ په ياد ولرئ
چې بېل سټونه~\LR{$A$} او~\LR{$B$} هله په محاسبوي ډول نه بېلېدونکي بلل کېږي
چې داسې محاسبه کېدونکے سټ~\LR{$C$} نۀ وي چې~\LR{$A \subseteq C$} او
\LR{$B \subseteq \Complement{C}$}۔
\olpseganchor{OLP-0309-B006}{end}

\olpseganchor{OLP-0309-B007}{start}
\begin{lem}
\LR{$\Th{Q}$} او~\LR{$\Th{\bar Q}$} په محاسبوي ډول نه بېلېدونکي دي۔
\end{lem}
\olpseganchor{OLP-0309-B007}{end}

\olpseganchor{OLP-0309-B008}{start}
\begin{proof}
فرض کړئ~\LR{$C$} داسې محاسبه کېدونکے سټ دے چې~\LR{$\Th{Q} \subseteq C$} او
\LR{$\Th{\bar Q} \subseteq \Complement{C}$}۔ پرېږدئ~\LR{$R(x,y)$} لاندې
اړيکه وي:
\begin{quote}
\LR{$x$} د~\LR{$!D(u)$} فارمول کوډوي او~\LR{$!D(\num y)$} په~\LR{$C$} کښې دے۔
\end{quote}
موږ به وښيو چې~\LR{$R(x,y)$} يوه نړيواله محاسبه کېدونکې اړيکه ده، چې
تناقض ترې ترلاسه کېږي۔
\olpseganchor{OLP-0309-B008}{end}

\olpseganchor{OLP-0309-B009}{start}
فرض کړئ~\LR{$S(y)$} محاسبه کېدونکې ده، او په~\LR{$\Th{Q}$} کښې د~\LR{$!D_S(u)$}
په وسيله تمثيلېږي۔ بيا
\begin{eqnarray*}
S(n) & \lif & \Th{Q} \Proves !D_S(\num n) \\
& \lif & !D_S(\num n) \in C
\end{eqnarray*}
او
\begin{eqnarray*}
\lnot S(n) & \lif & \Th{Q} \Proves \lnot !D_S(\num n) \\
& \lif & !D_S(\num n) \in \Th{\bar Q} \\
& \lif & !D_S(\num n) \not\in C
\end{eqnarray*}
\textbf{د ژباړې د نسخې سمون (OLCMP-094):} سرچينې په دواړو
ماوراءژبنیو اړيکو کښې طبيعي عدد د څيز-ژبې د عددنښې په بڼه ليکلے و؛
دلته عددنښه يوازې د تمثيلوونکي فارمول دننه کارول شوې ده۔
نو~\LR{$S(y)$} له~\LR{$R(\Gn{!D_S(u)},y)$} سره هم‌ارزه ده۔
\textbf{د ژباړې د نسخې سمون (OLCMP-095):} سرچينې په وروستۍ کرښه
کښې د فارمول کوډ په داسې بڼه ليکلے و چې ازاد متغير ئې په نامعتبره
عددنښه بدلوله او د څپرکي د ګوډل-کوډ نښه ئې نۀ کاروله؛ دلته د هماغه
يو-ازاد-متغير فارمول سم کوډ راغلے دے۔
\end{proof}
\olpseganchor{OLP-0309-B009}{end}

\olpseganchor{OLP-0310-B004}{start}
\olfileid{inc}{tcp}{con}
\olpseganchor{OLP-0310-B004}{end}

\olpseganchor{OLP-0310-B005}{start}
\olsection{له~\LR{$\Th{Q}$} سره سازګارې تيورۍ ناپرېکړه‌وړې دي}
\olpseganchor{OLP-0310-B005}{end}

\olpseganchor{OLP-0310-B006}{start}
لاندې قضيه وايي چې يوازې~\LR{$\Th{Q}$} ناپرېکړه‌وړې نۀ ده، بلکې په
حقيقت کښې هره هغه تيوري ناپرېکړه‌وړې ده چې له~\LR{$\Th{Q}$} سره اختلاف
ونۀ لري۔
\olpseganchor{OLP-0310-B006}{end}

\olpseganchor{OLP-0310-B007}{start}
\begin{thm}
پرېږدئ~\LR{$\Th{T}$} د حساب په ژبه کښې هره هغه تيوري وي چې له~\LR{$\Th{Q}$}
سره سازګاره وي (يعنې~\LR{$\Th{T} \cup \Th{Q}$} سازګاره وي)۔ بيا
\LR{$\Th{T}$} ناپرېکړه‌وړې ده۔
\end{thm}
\olpseganchor{OLP-0310-B007}{end}

\olpseganchor{OLP-0310-B008}{start}
\begin{proof}
په ياد ولرئ چې~\LR{$\Th{Q}$} د بديهي اصولو متناهي سټ~\LR{$!Q_1$}، \dots،
\LR{$!Q_8$} لري۔ ان دا ټول په يوه بديهي اصل بدلولے شو:
\LR{$!E = !Q_1 \land \dots \land !Q_8$}۔
\olpseganchor{OLP-0310-B008}{end}

\olpseganchor{OLP-0310-B009}{start}
فرض کړئ~\LR{$\Th{T}$} يوه د پرېکړې وړ تيوري ده چې له~\LR{$\Th{Q}$} سره
سازګاره ده۔ پرېږدئ
\[
C = \Setabs{!A}{\Th{T} \Proves !E \lif !A}.
\]
موږ ښيو چې~\LR{$C$} به د~\LR{$\Th{Q}$} او~\LR{$\Th{\bar Q}$} يوه محاسبه کېدونکې
بېلوونکې وي، چې تناقض دے۔ لومړے، که~\LR{$!A$} په~\LR{$\Th{Q}$} کښې وي، نو
\LR{$!A$} د~\LR{$\Th{Q}$} له بديهي اصولو څخه ثابتېدونکې ده؛ د استنتاج د
قضيې له مخې په لومړۍ درجې منطق کښې د~\LR{$!E \lif !A$} يو
اشتقاق شته۔ نو~\LR{$!A$} په~\LR{$C$} کښې ده۔
\olpseganchor{OLP-0310-B009}{end}

\olpseganchor{OLP-0310-B010}{start}
له بلې خوا، که~\LR{$!A$} په~\LR{$\Th{\bar Q}$} کښې وي، نو په لومړۍ درجې
منطق کښې د~\LR{$!E \lif \lnot !A$} يو ثبوت شته۔ که~\LR{$\Th{T}$}،
\LR{$!E \lif !A$} هم ثابت کړي، نو~\LR{$\Th{T}$}، \LR{$\lnot !E$} ثابتوي، او په
دې حالت کښې~\LR{$\Th{T} \cup \Th{Q}$} ناسازګاره ده۔ خو موږ فرض کوو
چې~\LR{$\Th{T} \cup \Th{Q}$} سازګاره ده، نو~\LR{$\Th{T}$}، \LR{$!E \lif !A$}
نۀ ثابتوي، او له دې امله~\LR{$!A$} په~\LR{$C$} کښې نۀ ده۔
\olpseganchor{OLP-0310-B010}{end}

\olpseganchor{OLP-0310-B011}{start}
موږ وښودله چې که~\LR{$!A$} په~\LR{$\Th{Q}$} کښې وي، نو په~\LR{$C$} کښې ده، او
که~\LR{$!A$} په~\LR{$\Th{\bar Q}$} کښې وي، نو په~\LR{$\Complement{C}$} کښې ده۔
نو~\LR{$C$} يوه محاسبه کېدونکې بېلوونکې ده، او همدا هغه تناقض دے چې
غوښت مو۔
\end{proof}
\olpseganchor{OLP-0310-B011}{end}

\olpseganchor{OLP-0310-B012}{start}
دا قضيه ډېره پياوړې ده۔ د بېلګې په توګه، له دې څخه لاندې پايله
لازمېږي:
\olpseganchor{OLP-0310-B012}{end}

\olpseganchor{OLP-0310-B013}{start}
\begin{cor}
  د حساب د ژبې دپاره لومړۍ درجې منطق (يعنې سټ
  \LR{$\Setabs{!A}{\text{\RL{\LR{$!A$} په لومړۍ درجې منطق کښې ثابتېدونکې ده}}}$})
  ناپرېکړه‌وړ دے۔
\end{cor}
\olpseganchor{OLP-0310-B013}{end}

\olpseganchor{OLP-0310-B014}{start}
\begin{proof}
لومړۍ درجې منطق د~\LR{$\emptyset$} د پايلو سټ دے، چې له~\LR{$\Th{Q}$} سره
سازګار دے۔
\end{proof}
\olpseganchor{OLP-0310-B014}{end}

\olpseganchor{OLP-0311-B004}{start}
\olfileid{inc}{tcp}{itp}
\olpseganchor{OLP-0311-B004}{end}

\olpseganchor{OLP-0311-B005}{start}
\olsection{هغه تيورۍ ناپرېکړه‌وړې دي چې~\LR{$\Th{Q}$} پکې تفسيرېدونکې وي}
\olpseganchor{OLP-0311-B005}{end}

\olpseganchor{OLP-0311-B006}{start}
دا پايلې لا نورې هم پياوړې کولے شو۔ په غيرصوري ډول، په بله
ژبه~\LR{$\Lang{L_2}$} کښې د يوې ژبې~\LR{$\Lang{L_1}$} تفسيرول دا غواړي چې د
\LR{$\Lang{L_1}$} کائنات، اړيکيزې نښې او تابعي نښې په~\LR{$\Lang{L_2}$} کښې
د فارمولونه په وسيله تعريف شي۔ که څه هم دلته ورته وخت نۀ ورکوو،
دا تعريف په کره بڼه جوړېدے شي۔
\olpseganchor{OLP-0311-B006}{end}

\olpseganchor{OLP-0311-B007}{start}
\begin{thm}
  فرض کړئ~\LR{$\Th{T}$} په داسې ژبه کښې يوه تيوري ده چې د حساب ژبه پکې
  داسې تفسيرېدے شي چې~\LR{$\Th{T}$} د~\LR{$\Th{Q}$} له تفسير سره سازګاره وي۔
  بيا~\LR{$\Th{T}$} ناپرېکړه‌وړې ده۔ که~\LR{$\Th{T}$} د~\LR{$\Th{Q}$} د بديهي
  اصولو تفسير ثابتوي، نو د~\LR{$\Th{T}$} هېڅ سازګاره غځونه د پرېکړې وړ
  نۀ ده۔
\end{thm}
\olpseganchor{OLP-0311-B007}{end}

\olpseganchor{OLP-0311-B008}{start}
ثبوت د تېرې قضيې د ثبوت يوازې يو وړوکے بدلون دے؛ د يوې خلاف بېلګې
په کارولو د~\LR{$\Th{Q}$} او~\LR{$\Th{\bar Q}$} بېلوونکې ترلاسه کېدے شي۔
\LR{$\Th{ZFC}$}، يعنې د انتخاب له بديهي اصل سره د زر مېلو--فرېنکل سټ
تيوري، داسې بديهي بنسټ ګڼلے شو چې د عادي رياضياتو ډېره برخه پکې
ترسره کېدے شي۔ په~\LR{$\Th{ZFC}$} کښې طبيعي عددونه تعريفولے شو، او د
دې تفسير له مخې د~\LR{$\Th{Q}$} بديهي اصول رښتيا دي۔ نو لرو:
\olpseganchor{OLP-0311-B008}{end}

\olpseganchor{OLP-0311-B009}{start}
\begin{cor}
د~\LR{$\Th{ZFC}$} هېڅ د پرېکړې وړ غځونه نشته۔
\end{cor}
\olpseganchor{OLP-0311-B009}{end}

\olpseganchor{OLP-0311-B010}{start}
\begin{cor}
د~\LR{$\Th{ZFC}$} هېڅ بشپړه، سازګاره او په محاسبوي ډول
په محاسبوي ډول بديهي کېدونکې غځونه نشته۔
\end{cor}
\olpseganchor{OLP-0311-B010}{end}

\olpseganchor{OLP-0311-B011}{start}
د~\LR{$\Th{ZFC}$} ژبه يوازې يوه دوه‌ځاييزه اړيکه~\LR{$\in$} لري۔ (په حقيقت
کښې ان عينيت ته هم اړتيا نشته۔) نو لرو:
\olpseganchor{OLP-0311-B011}{end}

\olpseganchor{OLP-0311-B012}{start}
\begin{cor}
د هرې هغې ژبې دپاره لومړۍ درجې منطق ناپرېکړه‌وړ دے چې يوه
دوه‌ځاييزه اړيکيزه نښه ولري۔
\end{cor}
\olpseganchor{OLP-0311-B012}{end}

\olpseganchor{OLP-0311-B013}{start}
دا پايله هرې هغې ژبې ته هم غځېږي چې دوه يوځاييزې تابعي نښې ولري،
ځکه د هغو په وسيله يوه دوه‌ځاييزه اړيکيزه نښه شبيه کېدے شي۔ همدا
اوس يادې شوې پايلې کره پولې دي: په حقيقت کښې د هرې هغې ژبې دپاره
لومړۍ درجې منطق د پرېکړې وړ دے چې يوازې \emph{يوځاييزې} اړيکيزې
نښې او تر ډېره يوه \emph{يوځاييزه} تابعي نښه ولري۔
\olpseganchor{OLP-0311-B013}{end}

\olpseganchor{OLP-0311-B014}{start}
يو بل په زړه پورې حقيقت۔ پوهېږو چې په معياري مدل کښې د ژبې
\LR{$\Obj 0$}، \LR{$'$}، \LR{$+$}، \LR{$\times$}، \LR{$<$} د رښتينو جملو سټ ناپرېکړه‌وړ
دے۔ په حقيقت کښې~\LR{$<$} د نورو نښو په وسيله تعريفولے شو، او بيا~\LR{$+$}
د~\LR{$\times$} او~\LR{$'$} په وسيله تعريفولے شو۔ نو د ژبې~\LR{$\Obj 0$}، \LR{$'$}،
\LR{$\times$} د رښتينو جملو سټ ناپرېکړه‌وړ دے۔ له بلې خوا، پرېسبورګر
ښودلې ده چې د ژبې~\LR{$\Obj 0$}، \LR{$'$}، \LR{$+$} د هغو جملو سټ چې د حساب په
معياري تفسير کښې رښتيا دي، د پرېکړې وړ دے۔
\textbf{د ژباړې د نسخې سمون (OLCMP-096):} سرچينې وروستے صدق «د
حساب په ژبه کښې» بللے و؛ ژبه يوازې نحو ټاکي، نو دلته هغه معياري
تفسير په څرګنده ياد شوے چې د پرېسبورګر د پرېکړې د پايلې موخه ده۔
خو دا کړنلاره له محاسبوي پلوه عملي نۀ ده۔
\olpseganchor{OLP-0311-B014}{end}

\olpseganchor{OLP-0312-B004}{start}
\olchapter{inc}{inp}{ناتکميلي او اثبات‌وړتيا}
\olpseganchor{OLP-0312-B004}{end}

\olpseganchor{OLP-0313-B004}{start}
\olfileid{inc}{inp}{int}
\olpseganchor{OLP-0313-B004}{end}

\olpseganchor{OLP-0313-B005}{start}
\olsection{پېژندنه}
\olpseganchor{OLP-0313-B005}{end}

\olpseganchor{OLP-0313-B006}{start}
هېلبرټ په دې اند و چې د يو رياضيکي جوړښت، لکه د طبيعي عددونو، دپاره
د بديهي اصولو نظام تر هغه وخته بسيا نۀ دے چې د هماغه جوړښت ټولې
رښتينې جملې ترې استنتاج نۀ شي۔ د صوري استنتاج له نظامونو سره د هغه
وروستۍ لېوالتيا چې ورسره يوځاے کړو، نو داسې ښکاري چې هغه غوښتل د
طبيعي عددونو په باب زموږ صوري استدلالي نظامونه يوازې سازګار نۀ، بلکې
\emph{بشپړ} هم وي؛ يعنې د نظام د ژبې هره جمله يا پخپله د اشتقاق وړ
وي يا ئې نفي همداسې وي۔ د ګوډل د ناتکميلۍ لومړۍ قضيه ښيي چې
داسې د بديهي اصولو نظام نشته: د حساب دپاره هېڅ بشپړ، سازګار او
په محاسبوي ډول بديهي کېدونکې صوري نظام نشته۔ په حقيقت کښې، هېڅ «په کافي اندازه
پياوړې»، سازګاره او په محاسبوي ډول بديهي کېدونکې رياضيکي تيوري بشپړه نۀ ده۔
\olpseganchor{OLP-0313-B006}{end}

\olpseganchor{OLP-0313-B007}{start}
د هېلبرټ يو لا مهم هدف، چې د نوې («کلاسيکې») رياضي د توجيه دپاره
د هغه د پروګرام منځنے ټکی و، د کلاسيک استدلال استازيتوب کوونکو صوري
نظامونو د سازګارۍ متناهي‌پاله ثبوتونه موندل وو۔ نو د هېلبرټ د
پروګرام له پلوه د ګوډل د ناتکميلۍ دويمه قضيه تر لومړۍ ډېر دروند ګوزار
و۔ دويمه قضيه هم د لومړۍ په شان په لږ ناڅرګندو ټکو بيانېدای شي:
نږدې معنا ئې دا ده چې د حساب هېڅ په کافي اندازه پياوړې تيوري خپله
سازګاري نۀ شي ثابتولے۔ د «په کافي اندازه پياوړې» مطلب به له~\LR{$\Th{Q}$}
څخه لږ څه زيات وغواړي۔
\olpseganchor{OLP-0313-B007}{end}

\olpseganchor{OLP-0313-B008}{start}
د ناتکميلۍ د قضيې د ګوډل د اصلي ثبوت مفکوره د اپيمنيدس په تناقض
کښې موندل کېږي۔ اپيمنيدس، چې د کريت اوسېدونکے و، وويل چې د کريت
ټول اوسېدونکي دروغجن دي؛ د دې تناقض لا نېغه بڼه «دا جمله ناسمه ده»
ويل دي۔ په اصل کښې ګوډل د رښتياوالي پر ځاے د اشتقاق وړتيا وکاروله
او تړلے فارمول ئې صوري کړه چې په غير مستقيم ډول وايي: زه پخپله
د اشتقاق وړ نۀ يم۔ که دغه تړلے فارمول د اشتقاق وړ واے، تيوري به
ناسازګاره وه۔ ګوډل دا هم وښودله چې د هغه د بديهي اصولو له نظامه د
دغې تړلے فارمول نفي هم د اشتقاق وړ نۀ ده۔ (د دې دويمې برخې دپاره
ګوډل اړ و چې فرض کړي تيوري~\LR{$\Th{T}$} «\LR{$\omega$}-سازګاره» ده۔
\LR{$\omega$}-سازګاري له سازګارۍ سره تړاو لري، خو تر هغې پياوړے خاصيت
دے۔\footnote{يعنې هره \LR{$\omega$}-سازګاره تيوري سازګاره ده، خو
برعکس لازمه نۀ ده۔} له ګوډل څو کاله وروسته
روسر وښودله چې يوازې د~\LR{$\Th{T}$} عادي سازګاري فرض کول بس دي۔)
\olpseganchor{OLP-0313-B008}{end}

\olpseganchor{OLP-0313-B009}{start}
لومړۍ ستونزه دا پوهېدل دي چې داسې تړلے فارمول څنګه جوړېږي چې
ځان ته مراجعه وکړي۔ د~\LR{$\Th{Q}$} په ژبه کښې د هر فارمول~\LR{$!A$}
دپاره دې~\LR{$\gn{!A}$} هغه عدديه نښه وي چې له~\LR{$\Gn{!A}$} سره سمون خوري۔
د دې معنا ته ځير شئ: \LR{$!A$}~د~\LR{$\Th{Q}$} په ژبه کښې فارمول دے؛
\LR{$\Gn{!A}$}~يو طبيعي عدد دے؛ او \LR{$\gn{!A}$}~د~\LR{$\Th{Q}$} په ژبه کښې يو
\emph{ترم} دے۔ نو د~\LR{$\Th{Q}$} د ژبې هر فارمول~\LR{$!A$} يو
\emph{نوم}، يعنې~\LR{$\gn{!A}$}، لري چې د~\LR{$\Th{Q}$} په ژبه کښې ترم
دے۔ دا موږ ته
هغه مفهومي چوکاټ راکوي چې په هغه کښې د~\LR{$\Th{Q}$} د ژبې
فارمولونه د نورو فارمولونه په باب څه «ويلے» شي۔ لاندې لمه
د ثابت ټکي د لمې په نوم يادېږي۔
\olpseganchor{OLP-0313-B009}{end}

\olpseganchor{OLP-0313-B010}{start}
\begin{lem}
فرض کړئ \LR{$\Th{T}$} هره هغه تيوري ده چې~\LR{$\Th{Q}$} غځوي، او \LR{$!B(x)$}
هر هغه فارمول دے چې يوازې متغير~\LR{$x$} پکې آزاد وي۔ نو داسې
تړلے فارمول~\LR{$!A$} شته چې
\LR{$\Th{T} \Proves!A \liff !B(\gn{!A})$}۔
\end{lem}
\olpseganchor{OLP-0313-B010}{end}

\olpseganchor{OLP-0313-B011}{start}
لمه وايي چې د هر خاصيت~\LR{$!B(x)$} دپاره داسې تړلے فارمول~\LR{$!A$}
شته چې وايي «\LR{$!B(x)$} زما په باب رښتيا دے»، او~\LR{$\Th{T}$} هم دغه
برابري «پېژني»۔
\olpseganchor{OLP-0313-B011}{end}

\olpseganchor{OLP-0313-B012}{start}
داسې تړلے فارمول څنګه جوړولے شو؟ د اپيمنيدس د تناقض لاندې بڼه
وګورئ چې کواین وړاندې کړې ده:
\begin{quote}
«که يې نقل‌قول ترې وړاندې کېښودل شي، دروغ وايي» که يې نقل‌قول
ترې وړاندې کېښودل شي، دروغ وايي۔
\end{quote}
دا جمله په نېغه ځان ته مراجعه نۀ کوي۔ يوازې د نقل‌قول د نښو ترمنځ
د نحوي شيانو په باب خبره کوي، نو له دې جملو سره په يوه کچه ده:
\begin{enumerate}
\item «رابرټ» يو ښۀ نوم دے۔
\item «زه وځغاستم۔» يوه لنډه جمله ده۔
\item «درې کلمې لري» درې کلمې لري۔
\end{enumerate}
خو که «که يې نقل‌قول ترې وړاندې کېښودل شي، دروغ وايي» عبارت واخلو
او د همدغه عبارت نقل شوې بڼه ترې مخکې کېږدو، څه پېښېږي؟ هماغه
لومړۍ جمله ترې جوړېږي!{} په لنډه توګه، جمله وايي چې پخپله ناسمه ده۔
\olpseganchor{OLP-0313-B012}{end}

\olpseganchor{OLP-0314-B004}{start}
\olfileid{inc}{inp}{fix}
\olpseganchor{OLP-0314-B004}{end}

\olpseganchor{OLP-0314-B005}{start}
\olsection{د ثابت ټکي لمه}
\olpseganchor{OLP-0314-B005}{end}

\olpseganchor{OLP-0314-B006}{start}
\begin{explain}
د ثابت ټکي لمه وايي چې د هر فارمول~\LR{$!B(x)$} دپاره داسې
تړلے فارمول~\LR{$!A$} شته چې
\LR{$\Th{T} \Proves !A \liff !B(\gn{!A})$}، په دې شرط چې~\LR{$\Th{T}$}
تيوري~\LR{$\Th{Q}$} غځوي۔ د دروغجنې جملې په بېلګه کښې غواړو چې
\LR{$!A$} په~\LR{$\Th{T}$} کښې په ثابتېدونکي ډول له «\LR{$\gn{!A}$} ناسمه ده»
سره معادله وي؛ يعنې له هغې جملې سره چې وايي~\LR{$\Gn{!A}$} د يوې
ناسمې تړلے فارمول ګوډل شمېره ده۔ د ثبوت د مفکورې د پوهېدو دپاره
دا ګټوره ده چې هغه د~\LR{$!A$} په باب د کواین له دې غير صوري څرګندونې
سره پرتله کړو:
«که يې خپله نقل شوې بڼه ترې مخکې کېښودل شي، دروغ وايي» که يې
خپله نقل شوې بڼه ترې مخکې کېښودل شي، دروغ وايي۔ له يوه عبارت څخه
د هغه خپله نقل شوې بڼه مخکې ايښودل او بيا ترې جمله جوړول د هغه
عبارت \emph{قطرول} بللے شو، او پايله ئې د هغه قطري بڼه ده۔ نو د
«که يې خپله نقل شوې بڼه ترې مخکې کېښودل شي، دروغ وايي» قطري بڼه
همدا «که يې خپله نقل شوې بڼه ترې مخکې کېښودل شي، دروغ وايي» که يې
خپله نقل شوې بڼه ترې مخکې کېښودل شي، دروغ وايي، ده۔ اوس پام وکړئ
چې د کواین دروغجنه جمله د «دروغ وايي» قطري بڼه نۀ ده، بلکې د
«که يې خپله نقل شوې بڼه ترې مخکې کېښودل شي، دروغ وايي» قطري بڼه
ده۔ نو هغه خاصيت چې د دروغجنې جملې د جوړولو دپاره قطرول کېږي،
پخپله هم قطرول رانغاړي!{}
\olpseganchor{OLP-0314-B006}{end}

\olpseganchor{OLP-0314-B007}{start}
د حساب په ژبه کښې د يو آزاد متغير لرونکي فارمول نقل شوې بڼه
د هغه د ګوډل شمېرې په حسابولو او بيا د دغې شمېرې د معياري عدديې
په آزاد متغير کښې په ځاے کولو جوړوو۔ د~\LR{$!E(x)$} قطري بڼه
\LR{$!E(\num{n})$} ده، چېرته چې \LR{$n = \Gn{!E(x)}$} دے۔ (له دې وروسته
\LR{$\num{\Gn{!E(x)}}$} په~\LR{$\gn{!E(x)}$} لنډوو۔) نو که \LR{$!B(x)$} د
«ناسمه ده» خاصيت وي، «که خپله نقل شوې بڼه ترې مخکې راشي، دروغ وايي»
به دا وي: «کله چې د خپلې قطري بڼې پر ګوډل شمېرې پلے شي، دروغ وايي»۔
که د هغې تابع~\LR{$\fn{diag}(n)$} دپاره چې د ګوډل شمېرې~\LR{$n$} لرونکي
فارمول د قطري بڼې ګوډل شمېره حسابوي، يوه نښه~\LR{$\Obj{diag}$}
ولرو، \LR{$!E(x)$} به د~\LR{$!B(\Obj{diag}(x))$} په بڼه وليکو۔ د کواین
د دروغجنې جملې بڼه به بيا د همدې عبارت قطري بڼه وي، يعنې
\LR{$!E(\gn{!E(x)})$} يا
\LR{$!B(\Obj{diag}(\gn{!B(\Obj{diag}(x))}))$}۔ البته اوس \LR{$!B(x)$}
هر بل خاصيت هم کېدے شي، او همدا جوړښت به بيا هم کار وکړي۔ د
ناتکميلۍ د قضيې دپاره \LR{$!B(x)$} داسې اخلو: «\LR{$x$} په~\LR{$\Th{T}$} کښې
د اشتقاق وړ نۀ دے»۔ نو \LR{$!E(x)$} به ووايي: «کله چې د خپلې قطري
بڼې پر ګوډل شمېرې پلے شي، داسې تړلے فارمول جوړوي چې په~\LR{$\Th{T}$}
کښې د اشتقاق وړ نۀ ده»۔
\olpseganchor{OLP-0314-B007}{end}

\olpseganchor{OLP-0314-B008}{start}
د دې خبرې په~\LR{$\Th{T}$} کښې د صوري کولو دپاره بايد د~\LR{$\fn{diag}$}
د صوري کولو لار ومومو۔ تابع~\LR{$\fn{diag}(n)$} محاسبه کېدونکې ده، بلکې
ابتدايي بازګشتي ده: که \LR{$n$} د يوه فورمول~\LR{$!E(x)$} ګوډل شمېره وي،
\LR{$\fn{diag}(n)$} د~\LR{$!E(\gn{!E(x)})$} ګوډل شمېره ورکوي۔ (ياد ولرئ
چې~\LR{$\gn{!E(x)}$} د~\LR{$!E(x)$} د ګوډل شمېرې معياري عدديه، يعنې
\LR{$\num{\Gn{!E(x)}}$}، ده۔) که \LR{$\Obj{diag}$} په~\LR{$\Th{T}$} کښې د
تابع~\LR{$\fn{diag}$} تمثيلوونکې تابع‌نښه واے، موږ به \LR{$!A$} د
\LR{$!B(\Obj{diag}(\gn{!B(\Obj{diag}(x))}))$} فورمول اخيستے واے۔
پام وکړئ چې
\begin{align*}
\fn{diag}(\Gn{!B(\Obj{diag}(x))}) & =
\Gn{!B(\Obj{diag}(\gn{!B(\Obj{diag}(x))}))} \\
& = \Gn{!A}.
\end{align*}
که فرض کړو چې~\LR{$\Th{T}$}
\[
\Obj{diag}(\gn{!B(\Obj{diag}(x))}) = \gn{!A},
\]
اشتقاق کولے شي، نو دا هم اشتقاق کولے شي چې
\LR{$!B(\Obj{diag}(\gn{!B(\Obj{diag}(x))}))
\liff !B(\gn{!A})$}۔ خو د دې معادلۍ کيڼ اړخ د تعريف له مخې~\LR{$!A$}
دے۔
\olpseganchor{OLP-0314-B008}{end}

\olpseganchor{OLP-0314-B009}{start}
البته \LR{$\Obj{diag}$} په عمومي توګه د~\LR{$\Th{T}$} تابع‌نښه نۀ ده، او
يقيناً د~\LR{$\Th{Q}$} تابع‌نښه نۀ ده۔ خو ځکه چې~\LR{$\fn{diag}$} محاسبه
کېدونکې ده، په~\LR{$\Th{Q}$} کښې د کوم فورمول
\LR{$!D_{\fn{diag}}(x,y)$} په وسيله \emph{تمثيلېږي}۔ نو د
\LR{$!B(\Obj{diag}(x))$} پر ځاے
\LR{$\lexists[y][(!D_{\fn{diag}}(x,y) \land !B(y))]$} ليکلے شو۔
له دې پرته، پورته لنډ بيان شوے ثبوت بشپړېږي، او په حقيقت کښې
لا په~\LR{$\Th{Q}$} کښې بشپړېږي۔
\end{explain}
\olpseganchor{OLP-0314-B009}{end}

\olpseganchor{OLP-0314-B010}{start}
\begin{lem}
\ollabel{lem:fixed-point} فرض کړئ \LR{$!B(x)$} هر هغه فورمول دے چې
يو آزاد متغير~\LR{$x$} لري۔ نو داسې جمله~\LR{$!A$} شته چې
\LR{$\Th{Q} \Proves !A \liff !B(\gn{!A})$}۔
\end{lem}
\olpseganchor{OLP-0314-B010}{end}

\olpseganchor{OLP-0314-B011}{start}
\begin{proof}
د ورکړي~\LR{$!B(x)$} دپاره دې \LR{$!E(x)$} فورمول
\LR{$\lexists[y][(!D_{\fn{diag}}(x,y) \land !B(y))]$} وي، او
\LR{$!A$} دې د هغه قطري بڼه، يعنې \LR{$!E(\gn{!E(x)})$} فورمول وي۔
\olpseganchor{OLP-0314-B011}{end}

\olpseganchor{OLP-0314-B012}{start}
ځکه چې \LR{$!D_{\fn{diag}}$} د~\LR{$\fn{diag}$} تمثيل کوي، او
\LR{$\fn{diag}(\Gn{!E(x)}) = \Gn{!A}$}، نو~\LR{$\Th{Q}$} دا
اشتقاق کولے شي:
\begin{align}
  & !D_{\fn{diag}}(\gn{!E(x)}, \gn{!A}) \ollabel{repdiag1} \\
  & \lforall[y][(!D_{\fn{diag}}(\gn{!E(x)},y) \lif
  \eq[y][\gn{!A}])]. \ollabel{repdiag2}
\end{align}
اوس ښيو چې \LR{$\Th{Q} \Proves !A \liff !B(\gn{!A})$}۔ موږ يوازې د
منطق او په~\LR{$\Th{Q}$} کښې د د اشتقاق وړ حقايقو په کارولو، په غير
صوري ډول استدلال کوو۔
\olpseganchor{OLP-0314-B012}{end}

\olpseganchor{OLP-0314-B013}{start}
لومړے فرض کړئ~\LR{$!A$}، يعنې \LR{$!E(\gn{!E(x)})$}۔ د~\LR{$!E(x)$} تعريف
ته په ستنېدو وينو چې \LR{$!E(\gn{!E(x)})$} همدا دے:
\[
\lexists[y][(!D_{\fn{diag}}(\gn{!E(x)},y) \land !B(y))].
\]
داسې يو~\LR{$y$} وګورئ۔ ځکه چې
\LR{$!D_{\fn{diag}}(\gn{!E(x)},y)$} دے، د~\olref{repdiag2} له مخې
\LR{$y = \gn{!A}$}۔ نو له~\LR{$!B(y)$} څخه~\LR{$!B(\gn{!A})$} اخلو۔
\olpseganchor{OLP-0314-B013}{end}

\olpseganchor{OLP-0314-B014}{start}
اوس فرض کړئ \LR{$!B(\gn{!A})$}۔ د~\olref{repdiag1} له مخې لرو:
\begin{align*}
& !D_{\fn{diag}}(\gn{!E(x)}, \gn{!A}) \land !B(\gn{!A}).
\intertext{\RL{نو ورڅخه راځي چې}}
& \lexists[y][(!D_{\fn{diag}}(\gn{!E(x)},y) \land !B(y))].
\end{align*}
خو دا هماغه \LR{$!E(\gn{!E(x)})$}، يعنې~\LR{$!A$}، دے۔
\end{proof}
\olpseganchor{OLP-0314-B014}{end}

\olpseganchor{OLP-0314-B015}{start}
\begin{digress}
دا ثبوت د محاسبه کېدنې په تيورۍ کښې د ثابت ټکي د لمې له ثبوت سره
پرتله کړئ۔ توپير دا دے چې دلته غواړو يو \emph{بيان} د خپل ځان په
وسيله تعريف کړو، خو هلته مو يوه \emph{تابعه} د خپل ځان په وسيله
تعريفوله؛ له دې توپيره پرته مفکوره په رښتيا هماغه ده۔
\end{digress}
\olpseganchor{OLP-0314-B015}{end}

\olpseganchor{OLP-0314-B016}{start}
\begin{prob}
  فارمول~\LR{$!A(x)$} د \emph{رښتياوالي تعريف} دے، که د هرې
  تړلي فارمولونه~\LR{$!B$} دپاره \LR{$\Th{Q} \Proves !B \liff !A(\gn{!B})$} وي۔
  د ثابت ټکي د لمې په کارولو وښيئ چې هېڅ فارمول د رښتياوالي
  تعريف نۀ دے۔
\end{prob}
\olpseganchor{OLP-0314-B016}{end}

\olpseganchor{OLP-0315-B004}{start}
\olfileid{inc}{inp}{1in}
\olpseganchor{OLP-0315-B004}{end}

\olpseganchor{OLP-0315-B005}{start}
\olsection{د ناتکميلۍ لومړۍ قضيه}
\olpseganchor{OLP-0315-B005}{end}

\olpseganchor{OLP-0315-B006}{start}
اوس د ګوډل د ناتکميلۍ د لومړۍ قضيې اصلي ثبوت بيانولے شو۔ فرض کړئ
\LR{$\Th{T}$} په هغې ژبه کښې هره په محاسبوي ډول بديهي‌کړې تيوري ده چې
د حساب ژبه غځوي، او~\LR{$\Th{T}$} د~\LR{$\Th{Q}$} بديهي اصول رانغاړي۔ دا په
ځانګړي ډول دا معنا لري چې~\LR{$\Th{T}$} محاسبه کېدونکې تابعې او اړيکې
تمثيلوي۔
\olpseganchor{OLP-0315-B006}{end}

\olpseganchor{OLP-0315-B007}{start}
موږ استدلال کړے دے چې که فورمولونه او ثبوتونه په مناسب ډول د عددونو
په توګه کوډ شي، نو اړيکه~\LR{$\Prf[\Th{T}](x,y)$} محاسبه کېدونکې ده۔
\LR{$\Prf[\Th{T}](x,y)$} هله او يوازې هله برقرارېږي چې~\LR{$x$} په~\LR{$\Th{T}$}
کښې د هغه فارمول د اشتقاق ګوډل شمېره وي چې ګوډل شمېره
ئې~\LR{$y$} ده۔ په حقيقت کښې ګوډل د خپلې ټاکلې تيورۍ دپاره وښودله چې
دا اړيکه ابتدايي بازګشتي ده؛ د دې دپاره ئې په خپله مقاله کښې د
۴۵ تابعو او اړيکو لړ وکاراوه۔ پينځه څلوېښتمه اړيکه، \LR{$x B y$}، د
ګوډل د~\LR{$\Th{T}$} د هماغې ټاکنې دپاره همدا~\LR{$\Prf[\Th{T}](x,y)$} ده۔
ياد ولرئ چې ګوډل په خپله مقاله کښې «بازګشتي» وايي، خو موږ به اوس
ورته «ابتدايي بازګشتي» وايو۔
\olpseganchor{OLP-0315-B007}{end}

\olpseganchor{OLP-0315-B008}{start}
څرنګه چې~\LR{$\Prf[\Th{T}](x,y)$} محاسبه کېدونکې ده، په~\LR{$\Th{T}$} کښې
تمثيلېدے شي۔ د هغې تمثيلوونکي فورمول دپاره~\LR{$\OPrf[\Th{T}](x,y)$}
کاروو۔ \LR{$\OProv[\Th{T}](y)$} دې فورمول
\LR{$\lexists[x][\OPrf[\Th{T}](x,y)]$} وي۔ دا د ګوډل په لړ کښې
شپږڅلوېښتمه اړيکه، \LR{$\fn{Bew}(y)$}، بيانوي۔ لکه څنګه چې ګوډل يادونه
کوي، دا يوازينۍ اړيکه ده چې «بازګشتي نۀ شو بللے»۔ ښايي د هغه مطلب
دا و: له تعريفه ئې محاسبه کېدنه نۀ څرګندېږي؛ او وروستيو پرمختګونو
په حقيقت کښې وښودله چې محاسبه کېدونکې هم نۀ ده۔
\olpseganchor{OLP-0315-B008}{end}

\olpseganchor{OLP-0315-B009}{start}
فرض کړئ \LR{$\Th{T}$} داسې په محاسبوي ډول بديهي کېدونکې تيوري ده چې~\LR{$\Th{Q}$}
رانغاړي۔ نو~\LR{$\Prf[\Th{T}](x, y)$} د پرېکړې وړ ده، او ځکه په~\LR{$\Th{Q}$}
کښې د يوه فارمول~\LR{$\OPrf[\Th{T}](x, y)$} په وسيله تمثيلېږي۔
\LR{$\OProv[\Th{T}](y)$} دې هغه فورمول وي چې پورته مو بيان کړ۔ د ثابت
ټکي د لمې له مخې داسې فورمول~\LR{$!G_\Th{T}$} شته چې~\LR{$\Th{Q}$}، او ځکه
\LR{$\Th{T}$} هم، دا اشتقاق کوي:
\begin{equation}
\ollabel{eqn:qpf}
!G_\Th{T} \liff \lnot \OProv[\Th{T}](\gn{!G_\Th{T}}).
\end{equation}
پام وکړئ چې~\LR{$!G_\Th{T}$} په اصل کښې وايي: «\LR{$!G_\Th{T}$} په~\LR{$\Th{T}$}
کښې د اشتقاق وړ نۀ دے»۔
\olpseganchor{OLP-0315-B009}{end}

\olpseganchor{OLP-0315-B010}{start}
\begin{lem}\ollabel{lem:cons-G-unprov}
که~\LR{$\Th{T}$} سازګاره او په محاسبوي ډول بديهي کېدونکې تيوري وي چې~\LR{$\Th{Q}$}
غځوي، نو \LR{$\Th{T} \Proves/ !G_\Th{T}$}۔
\end{lem}
\olpseganchor{OLP-0315-B010}{end}

\olpseganchor{OLP-0315-B011}{start}
\begin{proof}
فرض کړئ \LR{$\Th{T}$}، \LR{$!G_\Th{T}$} اشتقاق کوي۔ نو د هغه
اشتقاق \emph{شته} دے، او په دې توګه د کوم عدد~\LR{$m$} دپاره
اړيکه~\LR{$\Prf[\Th{T}](m,
\Gn{!G_\Th{T}})$} برقرارېږي۔ خو بيا~\LR{$\Th{Q}$} جمله
\LR{$\OPrf[\Th{T}](\num m, \gn{!G_\Th{T}})$} اشتقاق کوي۔ نو~\LR{$\Th{Q}$}
\LR{$\lexists[x][\OPrf[\Th{T}](x,\gn{!G_\Th{T}})]$} هم اشتقاق کوي،
چې د تعريف له مخې~\LR{$\OProv[\Th{T}](\gn{!G_\Th{T}})$} دے۔ د
\olref{eqn:qpf} له مخې~\LR{$\Th{Q}$}، \LR{$\lnot !G_\Th{T}$} اشتقاق کوي؛
او ځکه چې~\LR{$\Th{T}$}، \LR{$\Th{Q}$} غځوي، \LR{$\Th{T}$} ئې هم ثابتوي۔
موږ وښودله چې که~\LR{$\Th{T}$}، \LR{$!G_\Th{T}$} اشتقاق کړي، نو د هغه
نفي~\LR{$\lnot !G_\Th{T}$} هم اشتقاق کوي؛ په دې صورت کښې به تيوري
ناسازګاره وي۔
\end{proof}
\olpseganchor{OLP-0315-B011}{end}

\olpseganchor{OLP-0315-B012}{start}
\begin{defn}
\ollabel{thm:oconsis-q}
تيوري~\LR{$\Th{T}$} \emph{\LR{$\omega$}-سازګاره} ده که لاندې شرط پوره کړي:
که~\LR{$\lexists[x][!A(x)]$} هره جمله وي او~\LR{$\Th{T}$} د
\LR{$\lnot !A(\num 0)$}، \LR{$\lnot !A(\num 1)$}، \LR{$\lnot !A(\num 2)$}،
\dots هره يوه اشتقاق کړي، نو~\LR{$\Th{T}$}،
\LR{$\lexists[x][!A(x)]$} نۀ ثابتوي۔
\end{defn}
\olpseganchor{OLP-0315-B012}{end}

\olpseganchor{OLP-0315-B013}{start}
پام وکړئ چې هره \LR{$\omega$}-سازګاره تيوري سازګاره هم ده۔ دا نېغ له
دې خبرې راځي چې که~\LR{$\Th{T}$} ناسازګاره وي، نو د هر~\LR{$!A$} دپاره
\LR{$\Th{T} \Proves !A$}۔ په ځانګړي ډول که~\LR{$\Th{T}$} ناسازګاره وي، د
هر~\LR{$n$} دپاره \LR{$\lnot !A(\num n)$} اشتقاق کوي او ورسره
\LR{$\lexists[x][!A(x)]$} هم اشتقاق کوي۔ نو ناسازګاره~\LR{$\Th{T}$}
\LR{$\omega$}-ناسازګاره ده۔ د متقابل عکس له مخې، که~\LR{$\Th{T}$}
\LR{$\omega$}-سازګاره وي، نو بايد سازګاره وي۔
\olpseganchor{OLP-0315-B013}{end}

\olpseganchor{OLP-0315-B014}{start}
\begin{lem}\ollabel{lem:omega-cons-G-unref}
که~\LR{$\Th{T}$}، \LR{$\omega$}-سازګاره او په محاسبوي ډول بديهي کېدونکې تيوري وي چې
\LR{$\Th{Q}$} غځوي، نو \LR{$\Th{T} \Proves/ \lnot !G_\Th{T}$}۔
\end{lem}
\olpseganchor{OLP-0315-B014}{end}

\olpseganchor{OLP-0315-B015}{start}
\begin{proof}
ښيو چې که~\LR{$\Th{T}$}، \LR{$\lnot !G_\Th{T}$} اشتقاق کړي، نو
\LR{$\omega$}-ناسازګاره ده۔ فرض کړئ~\LR{$\Th{T}$}، \LR{$\lnot !G_\Th{T}$}
اشتقاق کوي۔ که~\LR{$\Th{T}$} ناسازګاره وي، نو \LR{$\omega$}-ناسازګاره
ده، او ثبوت بشپړ دے۔ که نۀ، نو~\LR{$\Th{T}$} سازګاره ده او د
\olref{lem:cons-G-unprov} له مخې~\LR{$!G_\Th{T}$} نۀ اشتقاق کوي۔
څرنګه چې په~\LR{$\Th{T}$} کښې د~\LR{$!G_\Th{T}$} هېڅ اشتقاق نشته،
نو~\LR{$\Th{Q}$} دا اشتقاق کوي:
\[
\lnot \OPrf[\Th{T}](\num 0, \gn{!G_\Th{T}}), \lnot \OPrf[\Th{T}](\num 1,
\gn{!G_\Th{T}}), \lnot \OPrf[\Th{T}](\num 2, \gn{!G_\Th{T}}), \dots
\]
او~\LR{$\Th{T}$} هم همداسې کوي۔ له بلې خوا د~\olref{eqn:qpf} له مخې
\LR{$\lnot !G_\Th{T}$} له
\LR{$\lexists[x][\OPrf[\Th{T}](x,\gn{!G_\Th{T}})]$} سره معادل دے۔
نو~\LR{$\Th{T}$}، \LR{$\omega$}-ناسازګاره ده۔
\end{proof}
\olpseganchor{OLP-0315-B015}{end}

\olpseganchor{OLP-0315-B016}{start}
\begin{prob}
  هره \LR{$\omega$}-سازګاره تيوري سازګاره ده۔ وښيئ چې برعکس نۀ ده؛
  يعنې داسې تيورۍ شته چې سازګارې خو \LR{$\omega$}-ناسازګارې دي۔ دا په
  ښودلو ثابت کړئ چې \LR{$\Th{Q} \cup
  \{\lnot !G_\Th{Q}\}$} سازګاره خو \LR{$\omega$}-ناسازګاره ده۔
\end{prob}
\olpseganchor{OLP-0315-B016}{end}

\olpseganchor{OLP-0315-B017}{start}
\begin{thm}
\ollabel{thm:first-incompleteness} فرض کړئ \LR{$\Th{T}$} هره
\LR{$\omega$}-سازګاره، په محاسبوي ډول بديهي کېدونکې تيوري ده چې~\LR{$\Th{Q}$} غځوي۔ نو
\LR{$\Th{T}$} بشپړه نۀ ده۔
\end{thm}
\olpseganchor{OLP-0315-B017}{end}

\olpseganchor{OLP-0315-B018}{start}
\begin{proof}
  که~\LR{$\Th{T}$}، \LR{$\omega$}-سازګاره وي، نو سازګاره ده؛ ځکه د
  \olref{lem:cons-G-unprov} له مخې \LR{$\Th{T}
  \Proves/ !G_\Th{T}$}۔ د~\olref{lem:omega-cons-G-unref} له مخې
  \LR{$\Th{T} \Proves/ \lnot !G_\Th{T}$}۔ دا معنا لري چې~\LR{$\Th{T}$}
  ناتکميله ده، ځکه نۀ~\LR{$!G_\Th{T}$} او نۀ~\LR{$\lnot !G_\Th{T}$}
  اشتقاق کوي۔
\end{proof}
\olpseganchor{OLP-0315-B018}{end}

\olpseganchor{OLP-0316-B004}{start}
\olfileid{inc}{inp}{ros}
\olpseganchor{OLP-0316-B004}{end}

\olpseganchor{OLP-0316-B005}{start}
\olsection{د روسر قضيه}
\olpseganchor{OLP-0316-B005}{end}

\olpseganchor{OLP-0316-B006}{start}
ايا د ګوډل ثبوت داسې بدلولے شو چې د «\LR{$\omega$}-سازګارې» پر ځاے
يوازې «سازګاره» تيوري فرض کړو او پياوړې پايله واخلو؟ ځواب «هو» دے؛
د روسر موندلې چاره دا کار کوي۔ د روسر چاره دا ده چې د
\LR{$\OProv[\Th{T}](y)$} پر ځاے د د اشتقاق وړتيا يو «بدل شوے»
پريديکات~\LR{$\ORProv_T(y)$} وکاروو۔
\olpseganchor{OLP-0316-B006}{end}

\olpseganchor{OLP-0316-B007}{start}
\begin{thm}
\ollabel{thm:rosser}
فرض کړئ \LR{$\Th{T}$} هره سازګاره او په محاسبوي ډول بديهي کېدونکې تيوري ده چې
\LR{$\Th{Q}$} غځوي۔ نو~\LR{$\Th{T}$} بشپړه نۀ ده۔
\end{thm}
\olpseganchor{OLP-0316-B007}{end}

\olpseganchor{OLP-0316-B008}{start}
\begin{proof}
ياد ولرئ چې \LR{$\OProv[\Th{T}](y)$} د
\LR{$\lexists[x][\OPrf[\Th{T}](x,
  y)]$} په توګه تعريف شوے دے۔ دلته~\LR{$\OPrf[\Th{T}](x, y)$} هغه
د پرېکړې وړ اړيکه تمثيلوي چې هله او يوازې هله برقرارېږي چې~\LR{$x$}
د ګوډل شمېرې~\LR{$y$} لرونکې تړلے فارمول د اشتقاق ګوډل
شمېره وي۔ هغه اړيکه هم د پرېکړې وړ ده چې د~\LR{$x$} او~\LR{$y$} تر منځ
هله برقرارېږي چې~\LR{$x$} د ګوډل شمېرې~\LR{$y$} لرونکې جملې د
\emph{ردونې} ګوډل شمېره وي۔ \LR{$\fn{not}(x)$} دې هغه ابتدايي
بازګشتي تابعه وي چې داسې کار کوي: که~\LR{$x$} د يوه فورمول~\LR{$!A$} کوډ
وي، نو~\LR{$\fn{not}(x)$} د~\LR{$\lnot
!A$} کوډ وي۔ په دې توګه \LR{$\Refut[\Th{T}](x, y)$} هله او يوازې هله
برقرارېږي چې \LR{$\Prf[\Th{T}](x, \fn{not}(y))$} برقرار وي۔
\LR{$\ORefut[\Th{T}](x, y)$} دې دا اړيکه تمثيل کړي۔ نو که
\LR{$\Th{T} \Proves \lnot !A$} وي او~\LR{$\delta$} ئې اړوند اشتقاق
وي، \LR{$\Th{Q} \Proves
\ORefut[\Th{T}](\gn{\delta}, \gn{!A})$}۔ موږ
\LR{$\ORProv[\Th{T}](y)$} داسې تعريفوو:
\[
\lexists[x][(\OPrf[\Th{T}](x,y) \land \lforall[z][(z < x \lif \lnot
  \ORefut[\Th{T}](z,y))])].
\]
په تقريبي ډول~\LR{$\ORProv[\Th{T}](y)$} وايي: «په~\LR{$\Th{T}$} کښې د~\LR{$y$}
ثبوت شته، او د~\LR{$y$} تر دې لنډه ردونه نشته»۔ که~\LR{$\Th{T}$} سازګاره
فرض کړو، نو~\LR{$\ORProv[\Th{T}](y)$} د هماغو عددونو دپاره رښتيا دے
چې~\LR{$\OProv[\Th{T}](y)$} ئې دپاره رښتيا دے؛ خو په~\LR{$\Th{T}$} کښې د
\emph{اثبات‌وړتيا} له پلوه (او اوس پوهېږو چې د رښتياوالي او
اثبات‌وړتيا تر منځ توپير شته!) دواړه بېل خاصيتونه لري۔ که
\LR{$\Th{T}$} \emph{نا}سازګاره وي، نو دواړه د هماغو عددونو دپاره
\emph{نۀ} برقرارېږي!{} (\LR{$\ORProv[\Th{T}](y)$} ته ډېر ځله داسې
لوستل کېږي: «\LR{$y$} د روسر له مخې اثبات‌وړ دے»۔ خو لکه چې مو همدا
اوس وليدل، د روسر اثبات‌وړتيا د اثبات‌وړتيا کوم ځانګړے ډول نۀ
دے: په ناسازګارو تيوريو کښې داسې تړلي فارمولونه شته چې اثبات‌وړ
دي، خو د روسر له مخې اثبات‌وړ نۀ دي۔ دا نوم لوستونکے ګډوډولے شي؛
د ګډوډۍ د مخنيوي دپاره ورته يو بېل مصنوعي نوم، لکه «\LR{$y$} شمووړ دے»، هم
ويلے شئ۔)
\olpseganchor{OLP-0316-B008}{end}

\olpseganchor{OLP-0316-B009}{start}
د ثابت ټکي د لمې له مخې داسې فورمول~\LR{$!R_\Th{T}$} شته چې
\begin{equation}
  \Th{Q} \Proves !R_\Th{T} \liff \lnot \ORProv[\Th{T}](\gn{!R_\Th{T}}).
  \ollabel{RT}
\end{equation}
د~\olref[1in]{thm:first-incompleteness} د ثبوت برعکس، دلته دعوه
دا ده چې که~\LR{$\Th{T}$} سازګاره وي، نو~\LR{$\Th{T}$}، \LR{$!R_\Th{T}$}
نۀ اشتقاق کوي، او~\LR{$\Th{T}$}، \LR{$\lnot !R_\Th{T}$} هم نۀ اشتقاق کوي۔
(په نورو ټکو، د \LR{$\omega$}-سازګارۍ فرض ته اړتيا نشته۔)
\olpseganchor{OLP-0316-B009}{end}

\olpseganchor{OLP-0316-B010}{start}
لومړے ښيو چې \LR{$\Th{T} \Proves/ !R_{\Th{T}}$}۔ که ئې ثابت کړے
واے، نو له~\LR{$T$} څخه به د~\LR{$!R_{\Th{T}}$} يو اشتقاق و؛ د
هغې ګوډل شمېره دې~\LR{$n$} وي۔ بيا
\LR{$\Th{Q} \Proves \OPrf[\Th{T}](\num{n}, \gn{!R_{\Th{T}}})$}، ځکه چې
\LR{$\OPrf[\Th{T}]$} په~\LR{$\Th{Q}$} کښې \LR{$\Prf[\Th{T}]$} تمثيلوي۔ همدا
راز، د هر~\LR{$k < n$} دپاره~\LR{$k$} د~\LR{$\lnot !R_{\Th{T}}$} د يوه
اشتقاق ګوډل شمېره نۀ ده، ځکه~\LR{$\Th{T}$} سازګاره ده۔ نو د
هر~\LR{$k < n$} دپاره
\LR{$\Th{Q} \Proves \lnot
\ORefut[\Th{T}](\num{k}, \gn{!R_{\Th{T}}})$}۔ د
\olref[req][min]{lem:less-nsucc} له مخې
\LR{$\Th{Q} \Proves \lforall[z][(z < \num{n} \lif \lnot \ORefut[\Th{T}](z,
  \gn{!R_{\Th{T}}}))]$}۔ نو
\[
\Th{Q} \Proves \lexists[x][(\OPrf[\Th{T}](x,\gn{!R_{\Th{T}}}) \land \lforall[z][(z
    < x \lif \lnot \ORefut[\Th{T}](z,\gn{!R_{\Th{T}}}))])],
\]
خو دا هماغه~\LR{$\ORProv[\Th{T}](\gn{!R_{\Th{T}}})$} دے۔ د~\olref{RT}
له مخې \LR{$\Th{Q}
\Proves \lnot !R_{\Th{T}}$}۔ ځکه چې~\LR{$\Th{T}$}، \LR{$\Th{Q}$} غځوي،
نو \LR{$\Th{T}
\Proves \lnot !R_{\Th{T}}$} هم دے۔ موږ فرض کړې وه چې
\LR{$\Th{T} \Proves !R_{\Th{T}}$}؛ نو~\LR{$\Th{T}$} به ناسازګاره شي، چې د
قضيې له فرض سره ټکر لري۔
\olpseganchor{OLP-0316-B010}{end}

\olpseganchor{OLP-0316-B011}{start}
اوس ښيو چې \LR{$\Th{T} \Proves/ \lnot !R_{\Th{T}}$}۔ بيا فرض کړئ چې
دا نفي ثابتوي او~\LR{$n$} د~\LR{$\lnot !R_{\Th{T}}$} د يوه
اشتقاق ګوډل شمېره ده۔ نو
\LR{$\Refut[\Th{T}](n, \Gn{!R_{\Th{T}}})$} برقرارېږي۔ ځکه چې
\LR{$\ORefut[\Th{T}]$} په~\LR{$\Th{Q}$} کښې \LR{$\Refut[\Th{T}]$} تمثيلوي،
نو \LR{$\Th{Q} \Proves
\ORefut[\Th{T}](\num{n}, \gn{!R_{\Th{T}}})$}۔ بيا به وښيو چې
\LR{$\Th{T}$} ناسازګاره کېږي، ځکه~\LR{$!R_{\Th{T}}$} به هم اشتقاق کوي۔
څرنګه چې
\begin{align*}
\Th{Q} & \Proves !R_{\Th{T}} \liff \lnot \ORProv[\Th{T}](\gn{!R_{\Th{T}}}),
\intertext{\RL{او ځکه چې \LR{$\Th{T}$}، \LR{$\Th{Q}$} غځوي، بس دا وښيو چې}}
\Th{Q} & \Proves \lnot \ORProv[\Th{T}](\gn{!R_{\Th{T}}}).
\end{align*}
تړلے فارمول~\LR{$\lnot
\ORProv[\Th{T}](\gn{!R_{\Th{T}}})$}، يعنې
\begin{align*}
  \lnot & \lexists[x][(\OPrf[\Th{T}](x,\gn{!R_{\Th{T}}}) \land
    \lforall[z][(z < x \lif \lnot \ORefut[\Th{T}](z,\gn{!R_{\Th{T}}}))])],
  \intertext{\RL{په منطقي ډول له دې سره معادله ده:}}
  & \lforall[x][(\OPrf[\Th{T}](x,\gn{!R_{\Th{T}}}) \lif
    \lexists[z][(z < x \land \ORefut[\Th{T}](z,\gn{!R_{\Th{T}}}))])].
\end{align*}
موږ د منطق او په~\LR{$\Th{Q}$} کښې د اشتقاق کېدونکو حقايقو په
کارولو غير صوري استدلال کوو۔ فرض کړئ~\LR{$x$} خپل‌سرے دے او
\LR{$\OPrf[\Th{T}](x,
\gn{!R_{\Th{T}}})$} برقرارېږي۔ موږ لا پوهېږو چې
\LR{$\Th{T} \Proves/ !R_{\Th{T}}$}، نو د هر~\LR{$k$} دپاره
\LR{$\Th{Q} \Proves \lnot \OPrf[\Th{T}](\num{k}, \gn{!R_{\Th{T}}})$}۔
له دې د هر~\LR{$k$} دپاره \LR{$\eq/[x][\num{k}]$} راځي۔ په ځانګړي ډول (الف)
\LR{$\eq/[x][\num{n}]$} لرو۔ همدا راز
\LR{$\lnot(\eq[x][\num{0}]
\lor \eq[x][\num{1}] \lor \dots \lor \eq[x][\num{n-1}])$} لرو،
نو د~\olref[req][min]{lem:less-nsucc} له مخې (ب)
\LR{$\lnot(x < \num{n})$}۔ د~\olref[req][min]{lem:trichotomy} له مخې
\LR{$\num{n} < x$}۔ ځکه چې
\LR{$\Th{Q} \Proves
\ORefut[\Th{T}](\num{n}, \gn{!R_{\Th{T}}})$}، لرو چې
\LR{$\num{n} < x \land
\ORefut[\Th{T}](\num{n}, \gn{!R_{\Th{T}}})$}؛ له دې څخه
\LR{$\lexists[z][(z < x
\land \ORefut[\Th{T}](z,\gn{!R_{\Th{T}}}))]$} راځي۔ ځکه چې~\LR{$x$}
خپل‌سرے و، اړينه پايله دا ده:
\[
\lforall[x][(\OPrf[\Th{T}](x,\gn{!R_{\Th{T}}}) \lif
  \lexists[z][(z < x \land \ORefut[\Th{T}](z,\gn{!R_{\Th{T}}}))])].
\]
\end{proof}
\olpseganchor{OLP-0316-B011}{end}

\olpseganchor{OLP-0316-B012}{start}
\begin{prob}
د طبيعي عددونو دوه سټونه~\LR{$A$} او~\LR{$B$} هغه وخت
\emph{په محاسبوي ډول نه بېلېدونکي} بلل کېږي چې هېڅ د پرېکړې وړ
سټ~\LR{$X$} داسې نۀ وي چې \LR{$A \subseteq X$} او
\LR{$B \subseteq \Complement{X}$} وي (\LR{$\Complement{X}$} د~\LR{$X$}
متمم، يعنې~\LR{$\Nat \setminus X$}، دے)۔ فرض کړئ~\LR{$\Th{T}$} د~\LR{$\Th{Q}$}
سازګاره او په محاسبوي ډول بديهي کېدونکې غځونه ده۔ \LR{$A$} دې په~\LR{$\Th{T}$} کښې د
اثبات‌وړو تړلي فارمولونه د ګوډل شمېرو سټ وي، او~\LR{$B$} دې په~\LR{$\Th{T}$}
کښې د ردېدونکو جملو د ګوډل شمېرو سټ وي۔ ثابت کړئ چې~\LR{$A$} او~\LR{$B$}
په محاسبوي ډول نه بېلېدونکي دي۔
\end{prob}
\olpseganchor{OLP-0316-B012}{end}

\olpseganchor{OLP-0317-B004}{start}
\olfileid{inc}{inp}{gop}
\olpseganchor{OLP-0317-B004}{end}

\olpseganchor{OLP-0317-B005}{start}
\olsection{د ګوډل له اصلي مقالې سره پرتله}
\olpseganchor{OLP-0317-B005}{end}

\olpseganchor{OLP-0317-B006}{start}
د ګوډل د ۱۹۳۱ز کال مقالې ته لږ وخت ورکول ګټور دي۔ د هغې پېژندنه
د هغو مفکورو لنډ بيان کوي چې موږ همدا اوس وڅېړلې۔ که مقاله يوازې
په تېره هم وګورئ، بيا هم د هر پړاو کار څرګندېږي: لومړے ګوډل صوري
نظام~\LR{$P$} بيانوي (نحو، بديهي اصول او د ثبوت قاعدې)؛ بيا ابتدايي
بازګشتي تابعې او اړيکې تعريفوي؛ ورپسې ښيي چې~\LR{$x B y$} ابتدايي
بازګشتي ده او استدلال کوي چې ابتدايي بازګشتي تابعې او اړيکې
په~\LR{$\Th{P}$} کښې تمثيلېږي۔ بيا د ناتکميلۍ قضيه لکه پورته ثابتوي۔
په درېيمه برخه کښې ښيي چې ناثابتېدونکې دعوه د حساب د ژبې يوه جمله
اخيستے شو۔ د~\LR{$\beta$}-لمې اصل همدلته دے؛ موږ هم دا لمه د لړيو
د سمبالولو دپاره وکاروله، کله چې مو ښودله چې بازګشتي تابعې
په~\LR{$\Th{Q}$} کښې تمثيلېږي۔ ګوډل دومره وړاندې نۀ ځي چې د بسيا
بديهي اصولو يو لږترين سټ بېل کړي، خو اوس پوهېږو چې~\LR{$\Th{Q}$} د
دې کار دپاره بسيا ده۔ په پاې کښې، په څلورمه برخه کښې د ناتکميلۍ
د دويمې قضيې د ثبوت لنډه نقشه ورکوي۔
\olpseganchor{OLP-0317-B006}{end}

\olpseganchor{OLP-0318-B004}{start}
\olfileid{inc}{inp}{prc}
\olpseganchor{OLP-0318-B004}{end}

\olpseganchor{OLP-0318-B005}{start}
\olsection{د~\LR{$\Th{PA}$} د د اشتقاق وړتيا شرطونه}
\olpseganchor{OLP-0318-B005}{end}

\olpseganchor{OLP-0318-B006}{start}
د پېانو حساب، يا~\LR{$\Th{PA}$}، هغه تيوري ده چې~\LR{$\Th{Q}$} د ټولو
فارمولونه دپاره د استقرا د بديهي اصولو په زياتولو غځوي۔ په نورو
ټکو، \LR{$\Th{Q}$} ته د دې بڼې بديهي اصول ورزياتېږي:
\[
(!A(0) \land \lforall[x][(!A(x) \lif !A(x'))]) \lif \lforall[x][!A(x)]
\]
د هر فارمول~\LR{$!A$} دپاره۔ پام وکړئ چې دا په حقيقت کښې يوه
\emph{شېما} ده، يعنې بې‌شمېره بديهي اصول دي (او دا هم څرګندېږي
چې~\LR{$\Th{PA}$} په متناهي ډول {\em نۀ} شي بديهي کېدے)۔ خو ځکه چې
په مؤثره طريقه معلومولے شو چې د نښو يوه لړۍ د استقرا د کوم بديهي
اصل بېلګه ده که نۀ، د~\LR{$\Th{PA}$} د بديهي اصولو سټ محاسبه کېدونکے
دے۔ \LR{$\Th{PA}$} له~\LR{$\Th{Q}$} څخه ډېره پياوړې تيوري ده۔ مثلاً د
عادي استقرا په کارولو په آسانۍ ثابتولے شو چې جمع او ضرب تبادلي
دي۔ په حقيقت کښې د عددونو د تيورۍ او ترکيبپوهنې زياتره
متناهي‌پاله استدلالونه په~\LR{$\Th{PA}$} کښې ترسره کېدے شي۔
\olpseganchor{OLP-0318-B006}{end}

\olpseganchor{OLP-0318-B007}{start}
څرنګه چې~\LR{$\Th{PA}$} په محاسبوي ډول بديهي‌کړې ده، د
د اشتقاق وړتيا پريديکات~\LR{$\Prf[\Th{PA}](x,y)$} محاسبه کېدونکے
دے او ځکه په~\LR{$\Th{Q}$} کښې، او په دې توګه په~\LR{$\Th{PA}$} کښې هم،
تمثيلېږي۔ د پخوا په شان د دې اړيکې د تمثيلوونکي فورمول دپاره
\LR{$\OPrf[\Th{PA}](x,y)$} کاروو۔ \LR{$\OProv[\Th{PA}](y)$} دې فورمول
\LR{$\lexists[x][\OPrf[\Th{PA}](x,y)]$} وي؛ دا په مفهومي ډول وايي:
«\LR{$y$} د~\LR{$\Th{PA}$} له بديهي اصولو څخه د اشتقاق وړ دے»۔
\textbf{د ژباړې د نسخې سمون (OLPRV-001):} په اصلي سرچينه کښې د
همدې تعريف دننه~\LR{$\Prf[\Th{PA}](x,y)$} ليکل شوے، خو دا د عددونو
تر منځ بهرنۍ اړيکه ده؛ دلته د هغې تمثيلوونکے د ژبې فورمول
\LR{$\OPrf[\Th{PA}](x,y)$} پکار دے، لکه چې مخکينۍ جمله ئې په ښکاره
نوموي۔ له~\LR{$\Th{Q}$} د بديهي اصولو څخه لږ څه زيات ته ځکه اړ يو
چې خپله تيوري بايد دومره پياوړې وي چې د دې د اشتقاق وړتيا
پريديکات په باب څو بنسټيز حقايق اشتقاق کړي۔ اړين حقايق دا دي:
\begin{enumerate}
\item[P1.] که \LR{$\Th{PA} \Proves !A$} وي، نو \LR{$\Th{PA} \Proves
  \OProv[\Th{PA}](\gn{!A})$}۔
\item[P2.] د ټولو فارمولونه \LR{$!A$} او \LR{$!B$} دپاره:
  \[
  \Th{PA} \Proves \OProv[\Th{PA}](\gn{!A \lif !B}) \lif
  (\OProv[\Th{PA}](\gn{!A}) \lif \OProv[\Th{PA}](\gn{!B})).
  \]
\item[P3.] د هر فارمول~\LR{$!A$} دپاره:
  \[
  \Th{PA} \Proves \OProv[\Th{PA}](\gn{!A})
  \lif \OProv[\Th{PA}](\gn{\OProv[\Th{PA}](\gn{!A})}).
  \]
\end{enumerate}
دا درې خاصيتونه يوازې په دې توګه کتلے شو چې فارمول
\LR{$\OProv[\Th{PA}](y)$} په دقت بيان کړو او د~\LR{$\Th{PA}$} بديهي اصول
د اړوندو صوري اشتقاقونه د بيانولو دپاره وکاروو۔ شرطونه (۱)
او~(۲) آسانه دي؛ اصلي کار شرط~(۳) غواړي۔ (فکر وکړئ چې دا به
څه ډول کار وي \dots) د تفصيلاتو بشپړول ستړي کوونکي او دلته
بې‌ګټې وي، نو اوس له تاسو غواړو چې دا ومنئ چې~\LR{$\Th{PA}$} همدا
درې پورته خاصيتونه لري۔ د~\LR{$\OProv[\Th{PA}](y)$} يوه مناسبه ټاکنه
به دا هم پوره کوي:
\begin{enumerate}
\item[P4.] که \LR{$\Th{PA} \Proves \OProv[\Th{PA}](\gn{!A})$} وي، نو
  \LR{$\Th{PA} \Proves !A$}۔
\end{enumerate}
خو موږ به دې حقيقت ته اړ نۀ شو۔
\olpseganchor{OLP-0318-B007}{end}

\olpseganchor{OLP-0318-B008}{start}
\begin{digress}
په ضمن کښې ګوډل هم دلته زموږ په شان د تفصيلاتو په ليکلو کښې
تنبلي وکړه۔ د خپلې ۱۹۳۱ز مقالې په پاې کښې ئې د ناتکميلۍ د دويمې
قضيې د ثبوت لنډه نقشه ورکړه او ژمنه ئې وکړه چې تفصيلات به په بله
مقاله کښې راوړي۔ بيا ورته وخت ورنۀ غے؛ ځکه هر څوک چې په استدلال
پوهېدل، باور ئې درلود چې تفصيلات بشپړېدے شي (او هغه اړ نۀ و چې
پخپله ئې وليکي۔)
\end{digress}
\olpseganchor{OLP-0318-B008}{end}

\olpseganchor{OLP-0319-B004}{start}
\olfileid{inc}{inp}{2in}
\olpseganchor{OLP-0319-B004}{end}

\olpseganchor{OLP-0319-B005}{start}
\olsection{د ناتکميلۍ دويمه قضيه}
\olpseganchor{OLP-0319-B005}{end}

\olpseganchor{OLP-0319-B006}{start}
دا دعوه څنګه بيان کړو چې~\LR{$\Th{PA}$} خپله سازګاري نۀ شي ثابتولے؟
دا ويل چې~\LR{$\Th{PA}$} ناسازګاره ده، د دې ويلو برابر دي چې
\LR{$\Th{PA} \Proves \eq[0][1]$}۔ نو د سازګارۍ جمله~\LR{$\OCon[\Th{PA}]$}
د تړلے فارمول~\LR{$\lnot
\OProv[\Th{PA}](\gn{\eq[0][1]})$} په توګه اخيستے شو، او بيا
لاندې قضيه اړينه پايله راکوي:
\olpseganchor{OLP-0319-B006}{end}

\olpseganchor{OLP-0319-B007}{start}
\begin{thm}
\ollabel{thm:second-incompleteness}
که~\LR{$\Th{PA}$} سازګاره وي، نو~\LR{$\Th{PA}$}،
\LR{$\OCon[\Th{PA}]$} نۀ اشتقاق کوي۔
\end{thm}
\olpseganchor{OLP-0319-B007}{end}

\olpseganchor{OLP-0319-B008}{start}
دا خبره مهمه ده چې قضيه د~\LR{$\OCon[\Th{PA}]$} په ټاکلي تمثيل پورې
تړلې ده، يعنې د~\LR{$\OProv[\Th{PA}](y)$} په ټاکلي تمثيل پورې۔ موږ
يوازې دومره کاروو چې د~\LR{$\OProv[\Th{PA}](y)$} تمثيل د
د اشتقاق وړتيا درې شرطونه پوره کوي؛ نو قضيه هرې هغې تيورۍ ته
غځېږي چې د د اشتقاق وړتيا پريديکات ئې دغه خاصيتونه ولري۔
\olpseganchor{OLP-0319-B008}{end}

\olpseganchor{OLP-0319-B009}{start}
د ګوډل د استدلال لنډه نقشه لوستل ګټور دي، ځکه قضيه پکې لکه د ښې
ټوکې اغېزمنه وروستۍ جمله راځي۔ استدلال داسې دے: \LR{$!G_\Th{PA}$}
دې د ګوډل هغه جمله وي چې د~\olref[1in]{thm:first-incompleteness}
په ثبوت کښې مو جوړه کړه۔ موږ وښودله: «که~\LR{$\Th{PA}$} سازګاره وي،
نو~\LR{$\Th{PA}$}، \LR{$!G_\Th{PA}$} نۀ اشتقاق کوي»۔ که دا خبره
پخپله \emph{په}~\LR{$\Th{PA}$} کښې صوري کړو، لاندې ثبوت لرو:
\[
\OCon[\Th{PA}] \lif \lnot \OProv[\Th{PA}](\gn{!G_\Th{PA}}).
\]
اوس فرض کړئ~\LR{$\Th{PA}$}، \LR{$\OCon[\Th{PA}]$} اشتقاق کوي۔ بيا
\LR{$\lnot
\OProv[\Th{PA}](\gn{!G_\Th{PA}})$} هم اشتقاق کوي۔
\textbf{د ژباړې د نسخې سمون (OLPRV-002):} اصلي سرچينه په دې
جمله کښې~\LR{$\Prov[\Th{PA}]$} ليکي، خو د صوري جملې دپاره
\LR{$\OProv[\Th{PA}]$} پکار دے؛ همدا نښه پورته په ثبوتېدونکې
پايلي کښې راغلې ده۔ ځکه چې~\LR{$!G_\Th{PA}$} د ګوډل جمله ده، دا له
\LR{$!G_\Th{PA}$} سره معادله ده۔ نو~\LR{$\Th{PA}$}، \LR{$!G_\Th{PA}$}
اشتقاق کوي۔
\olpseganchor{OLP-0319-B009}{end}

\olpseganchor{OLP-0319-B010}{start}
خو موږ پوهېږو چې که~\LR{$\Th{PA}$} سازګاره وي، \LR{$!G_\Th{PA}$} نۀ
اشتقاق کوي!{} نو که~\LR{$\Th{PA}$} سازګاره وي،
\LR{$\OCon[\Th{PA}]$} هم نۀ اشتقاق کولے شي۔
\olpseganchor{OLP-0319-B010}{end}

\olpseganchor{OLP-0319-B011}{start}
د استدلال د لا کره کولو دپاره~\LR{$!G_\Th{PA}$} د~\LR{$\Th{PA}$} د ګوډل
جمله اخلو او د د اشتقاق وړتيا شرطونه (P1)--(P3) کاروو، څو
وښيو چې~\LR{$\Th{PA}$}، \LR{$\OCon[\Th{PA}] \lif
!G_\Th{PA}$} اشتقاق کوي۔ له دې به څرګنده شي چې~\LR{$\Th{PA}$}،
\LR{$\OCon[\Th{PA}]$} نۀ اشتقاق کوي۔ دا په~\LR{$\Th{PA}$} کښې د
ثبوت لنډه نقشه ده۔ (د آسانتيا دپاره د~\LR{$\Th{PA}$} لاندې ليکونه
پرېږدو۔)
\begin{align}
& !G \liff \lnot \OProv(\gn{!G}) \ollabel{G2-1}\\
& \qquad\text{\RL{\LR{$!G$} د ګوډل جمله ده}}\notag \\
& !G \lif \lnot \OProv(\gn{!G}) \ollabel{G2-2}\\
  & \qquad\text{\RL{له \olref{G2-1} څخه}} \notag\\
& !G \lif
  (\OProv(\gn{!G}) \lif \lfalse) \ollabel{G2-3}\\
  & \qquad\text{\RL{له \olref{G2-2} څخه د منطق له مخې}}\notag\\
& \OProv(\gn{
    !G \lif
    (\OProv(\gn{!G}) \lif \lfalse)
  }) \ollabel{G2-4}\\
  & \qquad\text{\RL{له \olref{G2-3} څخه د P1 شرط له مخې}} \notag\\
& \OProv(\gn{!G}) \lif
  \OProv(\gn{
    (\OProv(\gn{!G}) \lif \lfalse)
    }) \ollabel{G2-5}\\
  & \qquad\text{\RL{له \olref{G2-4} څخه د P2 شرط له مخې}} \notag\\
& \OProv(\gn{!G}) \lif (\OProv(\gn{\OProv(\gn{!G})}) \lif \OProv(\gn{\lfalse})) \ollabel{G2-6}\\
  & \qquad\text{\RL{له \olref{G2-5} څخه د P2 شرط او منطق له مخې}} \notag\\
& \OProv(\gn{!G}) \lif
  \OProv(\gn{\OProv(\gn{!G})}) \ollabel{G2-7}\\
   & \qquad\text{\RL{د P3 له مخې}} \notag\\
& \OProv(\gn{!G}) \lif \OProv(\gn{\lfalse}) \ollabel{G2-8}\\
  & \qquad \text{\RL{له \olref{G2-6} او \olref{G2-7} څخه د منطق له مخې}}\notag\\
& \OCon \lif \lnot \OProv(\gn{!G}) \ollabel{G2-9}\\
  & \qquad\text{\RL{د \olref{G2-8} متقابل عکس او \LR{$\OCon \ident \lnot \OProv(\gn{\lfalse})$}}}\notag \\
& \OCon \lif !G \notag\\
  & \qquad\text{\RL{له \olref{G2-1} او \olref{G2-9} څخه د منطق له مخې}}\notag
\end{align}
په پورته ثبوت کښې د منطق کارول يوازې د قضيې د منطق بنسټيز
حقايق دي: مثلاً~\olref{G2-3} له
\LR{$\Proves \lnot!A \liff (!A\lif
\lfalse)$} څخه کار اخلي، او~\olref{G2-8} له
\LR{$!A \lif (!B \lif !C), !A \lif !B
\Proves !A \lif !C$} څخه۔ په~\olref{G2-5} او~\olref{G2-6}
کښې د P2 شرط کارول د هغه پر بېلګو
\LR{$\OProv(\gn{!A \lif !B}) \lif
(\OProv(\gn{!A}) \lif \OProv(\gn{!B}))$} ولاړ دي۔ په لومړۍ
بېلګه کښې \LR{$!A \ident
!G$} او \LR{$!B \ident \OProv(\gn{!G}) \lif \lfalse$} دي؛ په
دويمه کښې \LR{$!A
\ident \OProv(\gn{!G})$} او \LR{$!B \ident \lfalse$} دي۔
\textbf{د ژباړې د نسخې سمون (OLPRV-003):} د دويمې بېلګې د
\LR{$!A$} په اصلي نښه کښې~\LR{$\gn{G}$} ليکل شوي وو؛ دلته د لومړۍ
بېلګې او د ثبوت د فورمولونو په شان د جملې نښه~\LR{$!G$} ده، نو
\LR{$\gn{!G}$} ساتل پکار دي۔
\olpseganchor{OLP-0319-B011}{end}

\olpseganchor{OLP-0319-B012}{start}
د ناتکميلۍ د دويمې قضيې لا انتزاعي بڼه دا ده:
\olpseganchor{OLP-0319-B012}{end}

\olpseganchor{OLP-0319-B013}{start}
\begin{thm}
\ollabel{thm:second-incompleteness-gen} فرض کړئ~\LR{$\Th{T}$} هره
سازګاره، بديهي‌کړې تيوري ده چې~\LR{$\Th{Q}$} غځوي، او
\LR{$\OProv[\Th{T}](y)$} هر هغه فورمول دے چې د~\LR{$\Th{T}$} دپاره د
د اشتقاق وړتيا شرطونه P1--P3 پوره کوي۔ نو~\LR{$\Th{T}$}،
\LR{$\OCon[T]$} نۀ اشتقاق کوي۔
\end{thm}
\olpseganchor{OLP-0319-B013}{end}

\olpseganchor{OLP-0319-B014}{start}
\begin{prob}
وښيئ چې~\LR{$\Th{PA}$}، \LR{$!G_{\Th{PA}} \lif \OCon[\Th{PA}]$}
اشتقاق کوي۔
\end{prob}
\olpseganchor{OLP-0319-B014}{end}

\olpseganchor{OLP-0319-B015}{start}
\begin{digress}
د دې کيسې درس دا دے چې د رياضي هېڅ «معقوله» سازګاره تيوري
خپله د سازګارۍ جمله نۀ شي اشتقاق کولے۔ فرض کړئ~\LR{$\Th{T}$}
د رياضي داسې تيوري ده چې~\LR{$\Th{Q}$} او د هېلبرټ «متناهي‌پاله»
استدلال (هر څه چې مطلب ترې وي) رانغاړي۔ نو ټوله~\LR{$\Th{T}$} د
\LR{$\Th{T}$} د سازګارۍ جمله نۀ شي اشتقاق کولے؛ ځکه ئې متناهي‌پاله
برخه هم په لا پياوړي دليل سره د~\LR{$\Th{T}$} د سازګارۍ جمله نۀ شي
اشتقاق کولے۔ په دې معنا د «ټولې رياضي» دپاره د سازګارۍ
متناهي‌پاله ثبوت نشته۔
\olpseganchor{OLP-0319-B015}{end}

\olpseganchor{OLP-0319-B016}{start}
د «متناهي‌پاله» لفظ په تفسير کښې لږ ګنجايش شته۔ ګوډل په خپلې
۱۹۳۱ز مقالې کښې دا شونتيا مني چې هغه څه چې موږ ئې «متناهي‌پاله»
بولو، ښايي د رياضي له هغو ډولونو بهر وي چې هېلبرټ غوښتل صوري
کړي۔ خو ګوډل په دې کښې د نرمۍ څخه کار اخيست؛ نن دا ګرانه ده چې
داسې څه ومومو چې په معقول ډول متناهي‌پاله وبلل شي، خو مثلاً
په~\LR{$\Th{ZFC}$}، يعنې د ټاکنې له بديهي اصل سره د زرملو--فرانکل
د سټونو په تيورۍ کښې، صوري نۀ شي۔
\end{digress}
\olpseganchor{OLP-0319-B016}{end}

\olpseganchor{OLP-0320-B004}{start}
\olfileid{inc}{inp}{lob}
\olpseganchor{OLP-0320-B004}{end}

\olpseganchor{OLP-0320-B005}{start}
\olsection{د لوب قضيه}
\olpseganchor{OLP-0320-B005}{end}

\olpseganchor{OLP-0320-B006}{start}
د يوې تيورۍ~\LR{$\Th{T}$} د ګوډل جمله د~\LR{$\lnot
\OProv[\Th{T}](y)$} يو ثابت ټکی دے؛ يعنې داسې تړلے فارمول~\LR{$!G$}
چې
\[
\Th{T} \Proves \lnot \OProv[\Th{T}](\gn{!G}) \liff !G.
\]
دا جمله د اشتقاق وړ نۀ ده۔ ځکه که~\LR{$\Th{T} \Proves !G$} وي، نو
(الف) د د اشتقاق وړتيا د شرط~(۱) له مخې
\LR{$\Th{T} \Proves \OProv[\Th{T}](\gn{!G})$}؛ او (ب)
\LR{$\Th{T}
\Proves !G$} له
\LR{$\Th{T} \Proves \lnot \OProv[\Th{T}](\gn{!G})
\liff !G$} سره يوځاے دا راکوي چې
\LR{$\Th{T} \Proves \lnot \OProv[\Th{T}](\gn{!G})$}۔ په دې صورت کښې
به~\LR{$\Th{T}$} ناسازګاره وي۔ اوس طبيعي پوښتنه دا ده چې د
\LR{$\OProv[\Th{T}](y)$} د ثابت ټکي حالت څه دے؛ يعنې د داسې
تړلے فارمول~\LR{$!H$} حالت چې
\[
\Th{T} \Proves \OProv[\Th{T}](\gn{!H}) \liff !H.
\]
که دا جمله د اشتقاق وړ واے، د شرط~(۱) له مخې به
\LR{$\Th{T} \Proves \OProv[\Th{T}](\gn{!H})$} و؛ خو همدا پايله د
پورته معادلۍ د مودس پوننس په کارولو هم ترلاسه کېږي۔ نو لږ تر لږه
د ګوډل د جملې د پخواني استدلال له لارې د~\LR{$\Th{T}$} ناسازګاري نۀ
راځي۔ دا، البته، نۀ ښيي چې~\LR{$\Th{T}$} په رښتيا~\LR{$!H$}
اشتقاق \emph{کوي}۔
\olpseganchor{OLP-0320-B006}{end}

\olpseganchor{OLP-0320-B007}{start}
که پوښتنه لږ عمومي کړو، پرمختګ کولے شو۔ د ثابت ټکي د معادلۍ
له کيڼ څخه ښي ته لوری،
\LR{$\OProv[\Th{T}](\gn{!H}) \lif !H$}، د يوې عمومي شېما بېلګه ده
چې د \emph{انعکاس اصل} بلل کېږي:
\LR{$\OProv[\Th{T}](\gn{!A}) \lif !A$}۔ دا نوم ځکه ورکول کېږي چې
په يوه معنا ښيي~\LR{$\Th{T}$} د هغو څه په باب «انعکاس» کولے شي چې
اشتقاق کوي؛ په بنسټيز ډول د هر~\LR{$!A$} دپاره وايي: «که~\LR{$\Th{T}$}،
\LR{$!A$} اشتقاق کولے شي، نو~\LR{$!A$} رښتيا دے»۔ دا خبره، البته،
يوازې د سمې تيورۍ دپاره رښتيا ده؛ نو داسې ښکاري چې تيورۍ به
په عمومي توګه د دې شېما ټولې بېلګې نۀ اشتقاق کوي۔ نو يوه
کافي پياوړې تيوري چې د د اشتقاق وړتيا شرطونه پوره کوي، کومې
بېلګې اشتقاق کولے شي؟ هغه ټولې بېلګې چې~\LR{$!A$} پخپله
د اشتقاق وړ وي، يقيناً۔ لاندې پايله ښيي چې بس همدومره دي۔
\olpseganchor{OLP-0320-B007}{end}

\olpseganchor{OLP-0320-B008}{start}
\begin{thm}\ollabel{thm:lob}
فرض کړئ~\LR{$\Th{T}$} داسې په محاسبوي ډول بديهي کېدونکې تيوري ده چې~\LR{$\Th{Q}$}
غځوي، او~\LR{$\OProv[\Th{T}](y)$} هغه فورمول دے چې د
\olref[2in]{sec} شرطونه P1--P3 پوره کوي۔ که~\LR{$\Th{T}$}،
\LR{$\OProv[\Th{T}](\gn{!A}) \lif !A$} اشتقاق کړي، نو په حقيقت
کښې~\LR{$\Th{T}$}، \LR{$!A$} هم اشتقاق کوي۔
\end{thm}
\olpseganchor{OLP-0320-B008}{end}

\olpseganchor{OLP-0320-B009}{start}
په بل ډول: که \LR{$\Th{T} \Proves/ !A$} وي، نو
\LR{$\Th{T} \Proves/
\OProv[\Th{T}](\gn{!A}) \lif !A$}۔ دا پايله د لوب د قضيې په نوم
يادېږي۔
\olpseganchor{OLP-0320-B009}{end}

\olpseganchor{OLP-0320-B010}{start}
\begin{explain}
د لوب د قضيې د ثبوت لارښوده مفکوره د دې يو ځيرک «ثبوت» دے چې
سانتا کلاز شته۔ (که دغه پايله مو نۀ خوښېږي، هره بله خوښه پايله
پر ځاے کولے شئ۔) استدلال داسې دے:
\begin{enumerate}
\item \LR{$X$} دې دا جمله وي: «که~\LR{$X$} رښتيا وي، نو سانتا کلاز شته»۔
\item فرض کړئ~\LR{$X$} رښتيا ده۔
\item نو د هغې ويل شوې خبره برقرار ده؛ يعنې که~\LR{$X$} رښتيا وي،
  سانتا کلاز شته۔
\item ځکه چې فرض مو کړے~\LR{$X$} رښتيا ده، له (۲) او~(۳) څخه د
  مودس پوننس په کارولو پايله اخلو چې سانتا کلاز شته۔
\item موږ (۴)، «سانتا کلاز شته»، له (۲)، «\LR{$X$} رښتيا ده»، څخه
  وايستله۔ د شرطي ثبوت له مخې مو وښودله: «که~\LR{$X$} رښتيا وي، نو
  سانتا کلاز شته»۔
\item خو دا هماغه جمله~\LR{$X$} ده۔ نو ومو ښودله چې~\LR{$X$} رښتيا ده۔
\item بيا د پورته (۲)--(۴) استدلال له مخې سانتا کلاز شته۔
\end{enumerate}
که په دې مفکوره کښې «رښتيا ده» په «د اشتقاق وړ ده» او «سانتا
کلاز شته» په~\LR{$!A$} واړوو او هغه صوري کړو، د لوب د قضيې ثبوت
ترې جوړېږي۔ چاره دا ده چې د ثابت ټکي لمه په
فارمول~\LR{$\OProv[\Th{T}](y) \lif !A$} پلې کړو۔ د دې فورمول
ثابت ټکی په پورته لنډ استدلال کښې له~\LR{$X$} جملې سره سمون لري۔
\end{explain}
\olpseganchor{OLP-0320-B010}{end}

\olpseganchor{OLP-0320-B011}{start}
\begin{proof}[د \olref{thm:lob} ثبوت]
فرض کړئ~\LR{$!A$} داسې تړلے فارمول ده چې~\LR{$\Th{T}$}،
\LR{$\OProv[\Th{T}](\gn{!A}) \lif !A$} اشتقاق کوي۔ \LR{$!B(y)$}
دې فارمول~\LR{$\OProv[\Th{T}](y)
\lif !A$} وي، او د ثابت ټکي د لمې په کارولو داسې
تړلے فارمول~\LR{$!D$} ومومئ چې~\LR{$\Th{T}$}،
\LR{$!D \liff !B(\gn{!D})$} اشتقاق کوي۔ بيا لاندې هره يوه
په~\LR{$\Th{T}$} کښې د اشتقاق وړ ده:
\begin{align}
  & !D \liff (\OProv[\Th{T}](\gn{!D}) \lif !A) \ollabel{L-1}\\
  & \qquad \text{\RL{\LR{$!D$} د~\LR{$!B(y)$} ثابت ټکی دے}}\notag \\
  & !D \lif (\OProv[\Th{T}](\gn{!D}) \lif !A) \ollabel{L-2}\\
  & \qquad\text{\RL{له \olref{L-1} څخه}}\notag\\
  & \OProv[\Th{T}](\gn{!D \lif (\OProv[\Th{T}](\gn{!D}) \lif !A)}) \ollabel{L-3}\\
  & \qquad \text{\RL{له \olref{L-2} څخه د P1 شرط له مخې}}\notag \\
  & \OProv[\Th{T}](\gn{!D}) \lif \OProv[\Th{T}](\gn{\OProv[\Th{T}](\gn{!D}) \lif !A})
  \ollabel{L-4}\\
  &\qquad \text{\RL{له \olref{L-3} څخه د P2 شرط له مخې}}\notag \\
  & \OProv[\Th{T}](\gn{!D}) \lif (\OProv[\Th{T}](\gn{\OProv[\Th{T}](\gn{!D})}) \lif \OProv[\Th{T}](\gn{!A})) \ollabel{L-5}\\
  &\qquad \text{\RL{له \olref{L-4} څخه د P2 شرط له مخې بيا}} \notag\\
& \OProv[\Th{T}](\gn{!D}) \lif \OProv[\Th{T}](\gn{\OProv[\Th{T}](\gn{!D})}) \ollabel{L-6}\\
  & \qquad\text{\RL{د د اشتقاق وړتيا د P3 شرط له مخې}} \notag\\
  & \OProv[\Th{T}](\gn{!D}) \lif \OProv[\Th{T}](\gn{!A}) \ollabel{L-7} \\
  &\qquad\text{\RL{له \olref{L-5} او \olref{L-6} څخه}}\notag\\
  & \OProv[\Th{T}](\gn{!A}) \lif !A \ollabel{L-8}\\
  &\qquad\text{\RL{د قضيې د فرض له مخې}} \notag\\
  & \OProv[\Th{T}](\gn{!D}) \lif !A \ollabel{L-9}\\
  &\qquad\text{\RL{له \olref{L-7} او \olref{L-8} څخه}}\notag\\
  & (\OProv[\Th{T}](\gn{!D}) \lif !A) \lif !D \ollabel{L-10}\\
  & \qquad \text{\RL{له \olref{L-1} څخه}}\notag \\
  & !D \ollabel{L-11}\\
  & \qquad\text{\RL{له \olref{L-9} او \olref{L-10} څخه}}\notag \\
  & \OProv[\Th{T}](\gn{!D}) \ollabel{L-12}\\
  & \qquad\text{\RL{له \olref{L-11} څخه د P1 شرط له مخې}}\notag \\
  & !A \qquad\qquad\text{\RL{له \olref{L-8} او \olref{L-12} څخه}}\notag
\end{align}
\end{proof}
\olpseganchor{OLP-0320-B011}{end}

\olpseganchor{OLP-0320-B012}{start}
د لوب قضيه چې ولرو، د ناتکميلۍ د دويمې قضيې لنډ ثبوت ترې راځي
(د هغو تيوريو دپاره چې د P1--P3 شرطونه پوره کوونکے
د اشتقاق وړتيا پريديکات لري): که
\LR{$\Th{T} \Proves
\OProv[\Th{T}](\gn{\lfalse}) \lif \lfalse$} وي، نو
\LR{$\Th{T} \Proves \lfalse$}۔ که~\LR{$\Th{T}$} سازګاره وي،
\LR{$\Th{T} \Proves/ \lfalse$}۔ نو
\LR{$\Th{T}
\Proves/ \OProv[\Th{T}](\gn{\lfalse}) \lif \lfalse$}، يعنې
\LR{$\Th{T} \Proves/
\OCon[\Th{T}]$}۔ له دې څخه دا هم ښودلے شو چې~\LR{$!H$}، د
\LR{$\OProv[\Th{T}](x)$} ثابت ټکی، د اشتقاق وړ دے۔ ځکه چې
\begin{align*}
  \Th{T} & \Proves \OProv[\Th{T}](\gn{!H}) \liff !H\\
  \intertext{\RL{په ځانګړي ډول}}
    \Th{T} & \Proves \OProv[\Th{T}](\gn{!H}) \lif !H
\end{align*}
او ځکه د لوب د قضيې له مخې \LR{$\Th{T} \Proves !H$}۔
\olpseganchor{OLP-0320-B012}{end}

\olpseganchor{OLP-0320-B013}{start}
% Going in the other direction, for homework I
% may ask you to work through a short proof of L\"ob's theorem, using
% the second incompleteness theorem instead of the fixed-point lemma.
\olpseganchor{OLP-0320-B013}{end}

\olpseganchor{OLP-0320-B014}{start}
\begin{prob}
فرض کړئ~\LR{$\Th{T}$} په محاسبوي ډول بديهي‌کړې تيوري ده، او
\LR{$\OProv[\Th{T}]$} د~\LR{$\Th{T}$} دپاره د اشتقاق وړتيا پريديکات
دے۔ لاندې څلور بيانونه وګورئ:
\begin{enumerate}
\item که \LR{$T \Proves !A$} وي، نو \LR{$T \Proves \OProv[\Th{T}](\gn{!A})$}۔
\item \LR{$T \Proves !A \lif \OProv[\Th{T}](\gn{!A})$}۔
\item که \LR{$T \Proves \OProv[\Th{T}](\gn{!A})$} وي، نو \LR{$T \Proves !A$}۔
\item \LR{$T \Proves \OProv[\Th{T}](\gn{!A}) \lif !A$}۔
\end{enumerate}
له کومو شرطونو سره له دې بيانونو هر يو رښتيا دے؟
\end{prob}
\olpseganchor{OLP-0320-B014}{end}

\olpseganchor{OLP-0321-B004}{start}
\olfileid{inc}{inp}{tar}
\olpseganchor{OLP-0321-B004}{end}

\olpseganchor{OLP-0321-B005}{start}
\olsection{د رښتياوالي نه تعريفېدنه}
\olpseganchor{OLP-0321-B005}{end}

\olpseganchor{OLP-0321-B006}{start}
د \emph{تعريفېدنې} مفهوم د حساب د ژبې پر صوري معناپوهنې ولاړ دے۔
موږ د حساب په ژبه کښې د فورمولونو او جملو يو سټ بيان کړے دے۔
«مراد تفسير» دا دے چې دغه جملې د طبيعي عددونو په باب دعوې وبولو؛
داسې دعوه رښتيا يا ناسمه کېدے شي۔ \LR{$\Struct{N}$} دې هغه
جوړښت وي چې تعريف‌سيمه ئې~\LR{$\Nat$} او د حساب د ژبې
نښو تفسير ئې معياري دے۔ بيا~\LR{$\Sat{N}{!A}$} معنا لري: «\LR{$!A$}
په معياري تفسير کښې رښتيا ده»۔
\olpseganchor{OLP-0321-B006}{end}

\olpseganchor{OLP-0321-B007}{start}
\begin{defn}
د طبيعي عددونو اړيکه~\LR{$R(x_1,\dots,x_k)$} په~\LR{$\Struct{N}$} کښې
هله او يوازې هله \emph{تعريفېدونکې} ده چې د حساب د ژبې داسې
فورمول~\LR{$!A(x_1,\dots,x_k)$} وي چې د ټولو~\LR{$n_1,\dots,n_k$} دپاره
\LR{$R(n_1,\dots,n_k)$} هله او يوازې هله برقرار وي چې
\LR{$\Sat{N}{!A(\num n_1,\dots,\num
  n_k)}$}۔
\end{defn}
\olpseganchor{OLP-0321-B007}{end}

\olpseganchor{OLP-0321-B008}{start}
په بل ډول، يوه اړيکه په~\LR{$\Struct{N}$} کښې هله او يوازې هله
تعريفېدونکې ده چې په تيورۍ~\LR{$\Th{TA}$} کښې تمثيلېدونکې وي؛ دلته
\LR{$\Th{TA} =
\Setabs{!A}{\Sat{N}{!A}}$} د حساب د رښتينو جملو سټ دے۔ (که دا
خبره سملاسي روښانه نۀ وي، تعريفونو ته بېرته ورشئ او ځان پرې
قانع کړئ۔)
\olpseganchor{OLP-0321-B008}{end}

\olpseganchor{OLP-0321-B009}{start}
\begin{lem}
هره محاسبه کېدونکې اړيکه په~\LR{$\Struct{N}$} کښې تعريفېدونکې ده۔
\end{lem}
\olpseganchor{OLP-0321-B009}{end}

\olpseganchor{OLP-0321-B010}{start}
\begin{proof}
په آسانۍ کتل کېدے شي چې په~\LR{$\Th{Q}$} کښې د يوې اړيکې
تمثيلوونکے فورمول په~\LR{$\Struct{N}$} کښې هماغه اړيکه تعريفوي۔
\end{proof}
\olpseganchor{OLP-0321-B010}{end}

\olpseganchor{OLP-0321-B011}{start}
اوس پوښتنه دا ده چې آيا برعکس هم رښتيا دے؟ يعنې آيا په~\LR{$\Struct{N}$}
کښې هره تعريفېدونکې اړيکه محاسبه کېدونکې ده؟ ځواب نۀ دے۔ مثلاً:
\olpseganchor{OLP-0321-B011}{end}

\olpseganchor{OLP-0321-B012}{start}
\begin{lem}
د درېدنې اړيکه په~\LR{$\Struct{N}$} کښې تعريفېدونکې ده۔
\end{lem}
\olpseganchor{OLP-0321-B012}{end}

\olpseganchor{OLP-0321-B013}{start}
\begin{proof}
ياد ولرئ چې د کلين د عادي بڼې قضيه وايي: هره جزوي محاسبه
کېدونکې تابعه~\LR{$f$} داسې شاخص~\LR{$e$} لري چې د ټولو~\LR{$x \in \Nat$}
دپاره \LR{$f(x) =
U(\umin{s}{T(e,x,s)})$} وي؛ دلته~\LR{$U$} او~\LR{$T$} ابتدايي بازګشتي
او ځکه کلي دي۔ نو~\LR{$f(x)$} هله او يوازې هله تعريف شوې وي (يعنې
محاسبه ودرېږي) چې داسې~\LR{$s$} وي چې~\LR{$T(e,x,s)$} برقرار وي۔
\olpseganchor{OLP-0321-B013}{end}

\olpseganchor{OLP-0321-B014}{start}
اوس~\LR{$H$} د درېدنې اړيکه واخلو، يعنې
\[
H = \Setabs{\tuple{e,x}}{\lexists[s][T(e, x, s)]}.
\]
\LR{$!D_T$} دې په~\LR{$\Struct{N}$} کښې~\LR{$T$} تعريف کړي۔ نو
\[
H = \Setabs{\tuple{e,x}}{\Sat{N}{\lexists[s][!D_T(\num e, \num x, s)]}},
\]
او په دې توګه~\LR{$\lexists[s][!D_T(z, x, s)]$} په~\LR{$\Struct{N}$}
کښې~\LR{$H$} تعريفوي۔
\end{proof}
\olpseganchor{OLP-0321-B014}{end}

\olpseganchor{OLP-0321-B015}{start}
\begin{prob}
وښيئ چې \LR{$Q(n) \defiff n \in \Setabs{\Gn{!A}}{\Th{Q} \Proves !A}$}
  په حساب کښې تعريفېدونکې ده۔
\end{prob}
\olpseganchor{OLP-0321-B015}{end}

\olpseganchor{OLP-0321-B016}{start}
خو پخپله~\LR{$\Th{TA}$} څنګه ده؟ آيا په حساب کښې تعريفېدونکې ده؟
يعنې آيا سټ~\LR{$\Setabs{\Gn{!A}}{\Sat{N}{!A}}$} په حساب کښې
تعريفېدونکے دے؟ د تارسکي قضيه منفي ځواب ورکوي۔
\olpseganchor{OLP-0321-B016}{end}

\olpseganchor{OLP-0321-B017}{start}
\begin{thm}
\ollabel{thm:tarski}
د حساب د رښتينو تړلي فارمولونه سټ په حساب کښې نۀ شي تعريفېدے۔
\end{thm}
\olpseganchor{OLP-0321-B017}{end}

\olpseganchor{OLP-0321-B018}{start}
\begin{proof}
فرض کړئ~\LR{$!D(x)$} دا سټ تعريفوي؛ يعنې~\LR{$\Sat{N}{!A}$} هله او
يوازې هله چې~\LR{$\Sat{N}{!D(\gn{!A})}$}۔ د ثابت ټکي د لمې له مخې
داسې فورمول~\LR{$!A$} شته چې
\LR{$\Th{Q} \Proves !A \liff \lnot !D(\gn{!A})$}، او ځکه
\LR{$\Sat{N}{!A \liff \lnot !D(\gn{!A})}$}۔ خو بيا~\LR{$\Sat{N}{!A}$}
هله او يوازې هله چې~\LR{$\Sat{N}{\lnot !D(\gn{!A})}$}؛ دا له دې
فرض سره ټکر لري چې~\LR{$!D(y)$} د حساب د رښتينو جملو سټ تعريفوي۔
\end{proof}
\olpseganchor{OLP-0321-B018}{end}

\olpseganchor{OLP-0321-B019}{start}
تارسکي دا څېړنه د رښتياوالي پر يوه لا عمومي فلسفي مفهوم پلې کړه۔
د هرې ژبې~\LR{$L$} دپاره ئې استدلال وکړ چې د~\LR{$L$} د رښتياوالي مناسب
مفهوم بايد د هرې جملې~\LR{$X$} دپاره دا شرط پوره کړي:
\begin{quote}
«\LR{$X$}» هله او يوازې هله رښتيا ده چې \LR{$X$}۔
\end{quote}
د انګرېزۍ ژبې دپاره د تارسکي ډېره نقل شوې بېلګه دا جمله ده:
\begin{quote}
«واوره سپينه ده» هله او يوازې هله رښتيا ده چې واوره سپينه ده۔
\end{quote}
خو د هرې هغې ژبې دپاره چې د قطري تابع د تمثيل توان ولري، او د
هر ژبني پريديکات~\LR{$T(x)$} دپاره، داسې جمله~\LR{$X$} جوړولے شو چې وايي:
«\LR{$X$} هله او يوازې هله چې \LR{$T(\text{\RL{«\LR{$X$}»}})$} نۀ وي»۔ که موږ نۀ
غواړو چې د رښتياوالي پريديکات يوه جمله په يو وخت کښې هم رښتيا او
هم ناسمه وبولي، نو د تارسکي پايله دا وه چې د يوې ژبې د ټولو جملو
دپاره د رښتياوالي پريديکات د ژبې له پولو څخه له وتلو پرته نۀ شو
ټاکلے۔ په نورو ټکو، د يوې ژبې د رښتياوالي پريديکات په هماغه ژبه
کښې نۀ شي تعريفېدے۔
\olpseganchor{OLP-0321-B019}{end}

\clearpage
\begin{LTR}\latinfont
\begin{thebibliography}{99}
\bibitem[Benacerraf(1965)]{Benacerraf1965} Paul Benacerraf. 1965. ``What Numbers Could Not Be.'' \emph{The Philosophical Review} 74(1): 47--73.
\bibitem[Frege(1884)]{Frege1884} Gottlob Frege. 1884. \emph{Die Grundlagen der Arithmetik}. Breslau: Wilhelm Koebner.
\bibitem[Cantor(1892)]{Cantor1892} Georg Cantor. 1892. ``Uber eine elementare Frage der Mannigfaltigkeitslehre.'' \emph{Jahresbericht der Deutschen Mathematiker-Vereinigung} 1: 75--78.
\bibitem[Potter(2004)]{Potter2004} Michael Potter. 2004. \emph{Set Theory and Its Philosophy}. Oxford University Press.
\bibitem[Conway(2006)]{Conway2006} John Conway. 2006. ``The Power of Mathematics.'' In \emph{Power}, Cambridge University Press.
\bibitem[O'Connor and Robertson(2005)]{OConnorRobertson:RN} John J. O'Connor and Edmund F. Robertson. 2005. ``The Real Numbers: Stevin to Hilbert.''
\bibitem[Katz and Katz(2012)]{KatzKatz2012} Karin Usadi Katz and Mikhail G. Katz. 2012. ``Stevin Numbers and Reality.'' \emph{Foundations of Science} 17(2): 109--123.
\bibitem[Hilbert(2013)]{EwaldSieg2013} David Hilbert. 2013. \emph{Lectures on the Foundations of Arithmetic and Logic 1917--1933}. Edited by William Bragg Ewald and Wilfried Sieg. Springer.
\bibitem[Dedekind(1888)]{Dedekind1888} Richard Dedekind. 1888. \emph{Was sind und was sollen die Zahlen?} Braunschweig: Vieweg.
\bibitem[Magnus et al.(2021)]{Magnus2021} P. D. Magnus et al. 2021. \emph{Forall x: Calgary. An Introduction to Formal Logic}. Open Logic Project.
\bibitem[Smullyan(1968)]{Smullyan1968} Raymond M. Smullyan. 1968. \emph{First-Order Logic}. New York: Springer.
\bibitem[Zuckerman(1973)]{Zuckerman1973} Martin M. Zuckerman. 1973. ``Formation Sequences for Propositional Formulas.'' \emph{Notre Dame Journal of Formal Logic} 14(1): 134--138.
\end{thebibliography}
\end{LTR}

\end{document}
